\documentclass[11pt,a4paper]{article}
\pdfextension mapfile{+cm.map}
\pdfextension mapfile{+cmextra.map}
\pdfextension mapfile{+symbols.map}
\pdfextension mapfile{+euler.map}
\pdfextension mapfile{+latxfont.map}
\pdfextension mapfile{+rsfs.map}
\usepackage[margin=25mm,headheight=15pt]{geometry}
\usepackage{fontspec}
\setmainfont{Latin Modern Roman}
\usepackage{amsmath,amssymb,amsthm,amscd,mathtools,mathrsfs}
\usepackage{microtype,booktabs,longtable,array,xcolor,graphicx,enumitem,needspace,fancyhdr}
\usepackage[round,authoryear]{natbib}
\usepackage[unicode,colorlinks=true,linkcolor=blue!40!black,citecolor=blue!40!black,urlcolor=blue!40!black]{hyperref}
\usepackage{xurl}
\usepackage{listings}
\usepackage{tabularx}
\lstset{basicstyle=\ttfamily\small,breaklines=true,columns=fullflexible,keepspaces=true}
\usepackage{tikz}
\usepackage{tikz-cd}
\usetikzlibrary{arrows.meta,calc,patterns,positioning}
\usepackage{aliascnt}
\usepackage[nameinlink,noabbrev]{cleveref}
\allowdisplaybreaks[2]
\setcounter{MaxMatrixCols}{32}
\setlength{\emergencystretch}{3em}
\setlength{\parskip}{2pt}
\newtheorem{theorem}{Theorem}[section]
\newaliascnt{proposition}{theorem}
\newtheorem{proposition}[proposition]{Proposition}
\aliascntresetthe{proposition}
\crefname{proposition}{proposition}{propositions}
\newaliascnt{lemma}{theorem}
\newtheorem{lemma}[lemma]{Lemma}
\aliascntresetthe{lemma}
\crefname{lemma}{lemma}{lemmas}
\newaliascnt{corollary}{theorem}
\newtheorem{corollary}[corollary]{Corollary}
\aliascntresetthe{corollary}
\crefname{corollary}{corollary}{corollarys}
\newaliascnt{claim}{theorem}
\newtheorem{claim}[claim]{Claim}
\aliascntresetthe{claim}
\crefname{claim}{claim}{claims}
\newaliascnt{statement}{theorem}
\newtheorem{statement}[statement]{Statement}
\aliascntresetthe{statement}
\crefname{statement}{statement}{statements}
\theoremstyle{definition}
\newaliascnt{definition}{theorem}
\newtheorem{definition}[definition]{Definition}
\aliascntresetthe{definition}
\crefname{definition}{definition}{definitions}
\newaliascnt{remark}{theorem}
\newtheorem{remark}[remark]{Remark}
\aliascntresetthe{remark}
\crefname{remark}{remark}{remarks}
\newaliascnt{example}{theorem}
\newtheorem{example}[example]{Example}
\aliascntresetthe{example}
\crefname{example}{example}{examples}
\newaliascnt{problem}{theorem}
\newtheorem{problem}[problem]{Open problem}
\aliascntresetthe{problem}
\crefname{problem}{open problem}{open problems}
\newaliascnt{question}{theorem}
\newtheorem{question}[question]{Question}
\aliascntresetthe{question}
\crefname{question}{question}{questions}
\newaliascnt{auditrecord}{theorem}
\newtheorem{auditrecord}[auditrecord]{Verification record}
\aliascntresetthe{auditrecord}
\crefname{auditrecord}{verification record}{verification records}
\newaliascnt{test}{theorem}
\newtheorem{test}[test]{Verification}
\aliascntresetthe{test}
\crefname{test}{verification}{verifications}
\newaliascnt{conjecture}{theorem}
\newtheorem{conjecture}[conjecture]{Conjecture}
\aliascntresetthe{conjecture}
\crefname{conjecture}{conjecture}{conjectures}
\newcommand{\C}{\mathbb C}\newcommand{\R}{\mathbb R}\newcommand{\Z}{\mathbb Z}\newcommand{\Q}{\mathbb Q}\newcommand{\N}{\mathbb N}\newcommand{\F}{\mathbb F}
\newcommand{\CC}{\mathbb C}\newcommand{\RR}{\mathbb R}\newcommand{\ZZ}{\mathbb Z}\newcommand{\QQ}{\mathbb Q}\newcommand{\NN}{\mathbb N}
\newcommand{\PP}{\mathbb P}\newcommand{\Pp}{\mathbb P}\newcommand{\Pone}{\mathbb P^1}\newcommand{\HH}{\mathbb H}
\newcommand{\RP}{\mathbb{RP}}\newcommand{\CP}{\mathbb{CP}}
\newcommand{\OO}{\mathcal O}\newcommand{\cO}{\mathcal O}\newcommand{\cM}{\mathcal M}\newcommand{\cF}{\mathcal F}
\newcommand{\h}{\mathfrak h}\newcommand{\hz}{\mathfrak h_z}\newcommand{\Acon}{\mathfrak A}\newcommand{\wt}[1]{\widehat{#1}}
\newcommand{\red}{\mathrm{red}}\newcommand{\tor}{\mathrm{tor}}\newcommand{\dP}{\mathrm{dP}}
\newcommand{\gh}{\widehat\gamma}\newcommand{\uh}{\widehat u}\newcommand{\wh}{\widehat w}\newcommand{\dhat}{\widehat\delta}\renewcommand{\dh}{\widehat\delta}
\newcommand{\dd}{\mathrm d}\newcommand{\ii}{\mathrm i}\newcommand{\eps}{\varepsilon}\newcommand{\partt}{\mathrm{part}}\newcommand{\Bcirc}{B^\circ}\newcommand{\ddbar}{\partial\bar\partial}
\DeclareMathOperator{\Pic}{Pic}\DeclareMathOperator{\Alb}{Alb}\DeclareMathOperator{\Aut}{Aut}\DeclareMathOperator{\NS}{NS}
\DeclareMathOperator{\Hom}{Hom}\DeclareMathOperator{\End}{End}\DeclareMathOperator{\Ext}{Ext}\DeclareMathOperator{\Tor}{Tor}
\DeclareMathOperator{\Spec}{Spec}\DeclareMathOperator{\Proj}{Proj}\DeclareMathOperator{\Sym}{Sym}\DeclareMathOperator{\Diff}{Diff}
\DeclareMathOperator{\GL}{GL}\DeclareMathOperator{\SL}{SL}\DeclareMathOperator{\PSL}{PSL}\DeclareMathOperator{\Sp}{Sp}
\DeclareMathOperator{\im}{im}\DeclareMathOperator{\coker}{coker}\DeclareMathOperator{\rank}{rank}\DeclareMathOperator{\diag}{diag}
\DeclareMathOperator{\ord}{ord}\DeclareMathOperator{\vol}{vol}\DeclareMathOperator{\pr}{pr}\DeclareMathOperator{\id}{id}
\DeclareMathOperator{\divisor}{div}\DeclareMathOperator{\Tr}{Tr}\DeclareMathOperator{\Nm}{Nm}\DeclareMathOperator{\Sha}{Sha}
\DeclareMathOperator{\Imn}{Im}\DeclareMathOperator{\Ren}{Re}\DeclareMathOperator{\trdeg}{tr.deg}
\newcommand{\Uses}[1]{\par\smallskip\noindent\textit{Uses:} #1.\par\smallskip}
\newcommand{\Verified}{}\newcommand{\AuthorSelfCheck}{}\newcommand{\NeedsCheck}{}\newcommand{\Open}{\textbf{Open.}\ }

\newcommand{\StatementRecord}[4]{\par\smallskip\begingroup\upshape\raggedright\noindent\textbf{Label:} \nolinkurl{#1}.\par\noindent\textbf{Status:} #2.\par\noindent\textbf{Check record:} #3.\par\noindent\textbf{Uses:} #4.\par\endgroup\smallskip}
\newcommand{\RungZero}[1]{\par\noindent\textbf{Result and mechanism.} #1\par\smallskip}
\newcommand{\RungOne}{\par\Needspace{6\baselineskip}\noindent\textbf{Proof at the first rung.}\par\nobreak\smallskip}
\newcommand{\CurrentChecker}{GPT-6 Astra (OpenAI, ChatGPT Pro mode, effort Pro), 18 November 2026}
\newaliascnt{cited}{theorem}
\newtheorem{cited}[cited]{Cited result}
\aliascntresetthe{cited}
\crefname{cited}{cited result}{cited results}
\newaliascnt{refuted}{theorem}
\newtheorem{refuted}[refuted]{Correction}
\aliascntresetthe{refuted}
\crefname{refuted}{correction}{corrections}
\newaliascnt{computed}{theorem}
\newtheorem{computed}[computed]{Computation}
\aliascntresetthe{computed}
\crefname{computed}{computation}{computations}
\newaliascnt{openq}{theorem}
\newtheorem{openq}[openq]{Open question}
\aliascntresetthe{openq}
\crefname{openq}{open question}{open questions}

\providecommand{\rk}{\operatorname{rk}}
\providecommand{\Ann}{\operatorname{Ann}}
\providecommand{\Log}{\operatorname{Log}}
\providecommand{\vir}{\mathrm{vir}}

\begin{document}
\title{Maps, manifolds and cusp carriers}\author{AI-produced mathematical text}\date{19 November 2026}\maketitle
\section*{The results and their mechanisms}
We use the marked family of Paper~I, with its repaired circle quotient
$e:X\to S^3$ and algebraic reduction $f:X\to\PP^1$.
The integral lattice, its monodromies and the positive normal-circle
orientation are fixed in the transport dictionary below.  In particular,
an affine logarithmic parameter is an actual vector in that lattice;
its integral part remains in the boundary gluing map.  The maps in this
paper are specified with these markings.  Current proof checks and the
historical checks are distinguished in the record attached to each full
statement.

\subsection*{1. The repaired circle quotient and the freedom of its extension}
\RungZero{The quotient of the regular family extends across the cusp through
the explicit target solid torus constructed in Paper~I, with the
target and collar conditions fixing the extension used here.
Agreement only on the exterior permits all twelve elements of
$\pi_6(S^3)$, because the relative pinch of an interior six-ball
adds any chosen element (\cref{mapsI:src:A.2.1,maps:src:A.2.2}).}
\Needspace{5\baselineskip}
\RungOne
\begin{enumerate}[label=\arabic*.,leftmargin=*]
\item Use the primitive circle quotient and its explicit period
comparison from Paper~I.  The comparison preserves the signed primitive
covector, so the quotient and its boundary meridian are fixed maps.
\item The cusp target is the solid torus supplied by the local quotient
construction.  Its higher homotopy group in degree six is zero, which
gives the relative uniqueness within this prescribed target.
\item Choose a closed six-ball in the interior of the cusp
neighbourhood.  Collapse its boundary to obtain the relative pinch
$X\to X\vee S^6$, equal to the identity off the ball.
\item Compose the first summand with $e$ and the second with a
representative of $\alpha\in\pi_6(S^3)$.  The resulting map agrees
with $e$ on the exterior and represents $[e]+\alpha$.
This proves the exact scope of the extension assertion.
\end{enumerate}

\subsection*{2. The class of the map to the three-sphere}
\RungZero{The marked quotient has
$[e]=6\nu'=\eta_3\eta_4\eta_5$, the element of order two in
$\pi_6(S^3)=\Z/12$.  A regular inverse image is a three-torus;
the monodromy calculation identifies its normal framing with its Lie
framing, whose real $e$-invariant is $1/2$
(\cref{maps:src:B.1.1,maps:src:B.1.2}).}
\RungOne
\begin{enumerate}[label=\arabic*.,leftmargin=*]
\item Express the framing difference on the regular four-torus.
Its degree-one obstruction is the common invariant covector $\gamma$
modulo two.  Restriction to $\ker\gamma$ therefore lifts this
comparison to $\operatorname{Spin}$.
\item Twice the remaining degree-three obstruction is the restriction
of the common invariant class dual to the invariant vector $d$.
Since $d\in\ker\gamma$, this restriction is zero.  The integral
third cohomology of the three-torus is torsion-free, so the obstruction
itself is zero.
\item Pontryagin--Thom now represents $e$ by the Lie-framed
three-torus.  The periodic flat Dirac operator has paired nonzero
spectrum and a two-dimensional kernel.  The normalization in
\cite[Theorem~4.14, p.~423]{mapsAPS2} gives real $e$-invariant $1/2$.
\item Toda's stabilization calculation identifies the injected
$\pi_6(S^3)$ with the subgroup generated by $2\nu$ in the stable
three-stem.  The class with this invariant is $6\nu'$, also the
displayed composite of the three Hopf maps
\cite[Theorem~4.1 and Corollaries~4.2--4.3]{mapsToda52}.
\end{enumerate}

\subsection*{3. Null-homotopy of the map to the base}
\RungZero{The holomorphic algebraic reduction $f:X\to\PP^1$ is
null-homotopic.  Its explicit factorization through the circle quotient
contains the normalized power map of degree $3\cdot4=12$ on $S^3$,
which kills $\pi_6(S^3)$ (\cref{maps:src:B.2.1}).}
\RungOne
\begin{enumerate}[label=\arabic*.,leftmargin=*]
\item In the two multiplicative coordinates, the normalized power map
is $P_{3,4}(z_1,z_2)=(z_1^3,z_2^4)/\|(z_1^3,z_2^4)\|$.
A regular value has twelve positively oriented inverse images, so its
degree is twelve.
\item The marked quotient construction gives the equality of maps
$f=H\circ P_{3,4}\circ e$, where $H:S^3\to S^2$ is the Hopf
map and $\PP^1$ is identified with $S^2$ by the stated orientation.
\item The degree-twelve self-map acts by multiplication by twelve on
$\pi_6(S^3)=\Z/12$.  Thus $P_{3,4}\circ e$ is null-homotopic;
postcomposition with $H$ gives the required null-homotopy of $f$.
The body also gives the normalized-section homotopy directly.
\end{enumerate}

\subsection*{4. The cusp map retaining the angular coordinate}
\RungZero{The cusp collapse retaining the angular coordinate represents
$\eta_4\eta_5\ne0$ in $\pi_6(S^4)$.  At a triple point its
regular inverse image is a two-torus, and the explicit normal-frame
comparison makes this the Lie-framed torus of Arf invariant one
(\cref{maps:src:B.3.1,maps:src:B.3.2,maps:src:B.3.3}).}
\RungOne
\begin{enumerate}[label=\arabic*.,leftmargin=*]
\item Define the map in the triple-point chart using the retained
angle and the two radial coordinates.  Keep the bounded period
correction when differentiating the normal frame.
\item The resulting frame differs from the product frame by
$\operatorname{diag}(e^{-i(a+b)},e^{ia},e^{ib})$ in
$\operatorname{SU}(3)$.  The displayed two $\operatorname{SU}(2)$
factor homotopies contract this comparison.
\item The inverse-image torus therefore has its Lie framing.  Its
spin quadratic form has Arf invariant one, so its Pontryagin--Thom
class is the nonzero $\eta_4\eta_5$.
\item Two allowed real-period lifts agree on the exterior collar;
their straight-line interpolation gives a relative homotopy.
The core-collapse comparison factors through the specified
$S^3\vee S^2$ cofiber, with the map using its $S^3$ summand.
These are the explicit comparisons among the cusp maps.
\end{enumerate}

\subsection*{5. A finite cell model compatible with the fibration}
\RungZero{The regular torus mapping cylinders, the finite affine caps
and the marked honeycomb cusp fit into a finite CW model with cell
counts $(5,20,45,56,42,17,3)$.  The integral attaching maps give the
homotopy type of the six-sphere
(\cref{maps:src:B.5.1,maps:src:B.5.2}).}
\RungOne
\begin{enumerate}[label=\arabic*.,leftmargin=*]
\item Construct the regular mapping-cylinder bundle using the actual
integral monodromies.  In every finite affine cap retain the drift
from the logarithmic translation; it is part of the attaching map.
\item In the cusp write the seam map as
$[a,b,\theta]\mapsto[a,b-\theta(B_0a-v(a))]$.
The frozen-period comparison is an explicit logarithmic-coordinate
homeomorphism and agrees with the prescribed boundary markings.
\item The positive quotient has interior $T^2\times(0,1)$ and the
two marked boundary tori.  Its incompressible boundary and the
induced fundamental-group isomorphism meet Waldhausen's
boundary-fixed product criterion
\cite[Theorem~6.1, p.~77]{mapsWaldhausen68}.
\item Triangulate these marked pieces and their common faces.  The
resulting mapping-cylinder cells have the displayed counts, and
their alternating sum is two.
\item Van Kampen gives the relations $x^3=x^4=1$.  The primitive
finite pullback lattices and the two cusp sweeps give unimodular
integral attachments, so $H_1=H_2=0$.  Duality, the Euler number,
Hurewicz and Whitehead finish the sphere identification.
\end{enumerate}

\subsection*{6. The local Kummer cover and its global obstruction}
\RungZero{The prescribed degree-$N^2$ relative Kummer cover of the cusp,
for $N>1$, has no extension to a smooth finite cover of the whole
threefold, because its marked monodromy continuation has primitive
ray indices that violate the full-fan smoothness equations.
The explicit semidirect cover and its normal monomial completion
give the finite flat normal completion
(\cref{maps:src:C.1.1,maps:src:C.1.2,maps:src:C.1.3,maps:src:C.1.4,maps:src:C.1.5}).}
\RungOne
\begin{enumerate}[label=\arabic*.,leftmargin=*]
\item On the $N$th root of the base use the actual period
$C(v^N)/N$.  The monomial maps and their lattice inclusions
separate the absolute degree-$N^3$ map from its relative
degree-$N^2$ map.
\item Across adjacent unimodular triangles, smoothness imposes the
displayed divisibility relations on primitive ray indices.
The reverse relations force equality, so a smooth full-fan cover
must use the scalar toric sublattice $n\Z^3$.
\item Continue the prescribed local Kummer lattice by the marked
order-four monodromy.  Its toric intersection is
$\Z w\oplus N\Z d$.  Its primitive ray indices violate that
scalar condition for $N>1$.
\item The explicit word
$A_2^2A_1A_2^2A_1^{-1}$ also produces $N^2$ distinct based
lattices; their residues are exhibited in the body.  The finite
semidirect quotient supplies a global cover with the required
local lattice.
\item Its normal monomial completion has a free basis indexed by
fractional exponents in a half-open fundamental box.  This proves
finite flatness of that normal completion while recording the
smoothness obstruction above.
\end{enumerate}

\subsection*{7. The cyclic universal covers}
\RungZero{For the original marked family with fundamental group
$\Z/d$, where $d$ is prime to six, the universal cover has
$H_2=\Z^{d-1}\oplus\Z/d$, $H_3=\Z/d$, and Euler number
$2d$.  Its cusp divisor products and finite-coefficient cubic
evaluations determine the complete Zhubr invariant, hence one
oriented diffeomorphism type $M_d$ for each $d$
(\cref{maps:src:D.1.1,maps:src:D.1.4,maps:src:D.1.6,maps:src:D.1.9,maps:src:D.1.11}).}
\RungOne
\begin{enumerate}[label=\arabic*.,leftmargin=*]
\item The connected cusp cover has $d$ components in its periodic
normalization.  Its Laurent relation matrix reduces integrally
to $\operatorname{diag}(1,t-1,t-1,0)$.
\item The nontrivial deck characters make the regular torus Koszul
complex acyclic.  The primitive finite-cap degree-two pullbacks
and the cusp sweep with coefficient $d$ determine the integral
groups, including their torsion.
\item The $d$ cusp divisors give classes $x_i$ with
$\sum x_i=0$ and
$x_ix_j=(\mathbf1_{j=i+1}+\mathbf1_{j=i-1}-2\mathbf1_{j=i})\kappa$,
where $\kappa$ generates the degree-four torsion.
\item For the finite-coefficient class $u$ the body computes
$\langle ux_ix_j,[M_d]\rangle$ by the same cyclic matrix,
$\langle u^2x_i,[M_d]\rangle=-12$, and
$\langle u^3,[M_d]\rangle=0$ modulo $d$.  The support-function
restriction has square $-12$, which fixes the second value.
\item Transfer kills the integral pullback in degree four, so
$c_2=p_1=0$.  Multiplication by four is injective on
$\Z^{d-1}\oplus\Z/d$, giving zero refined Pontryagin class.
The full $H_6(K(H_2,2))$ calculation shows that these evaluations
specify the remaining invariant; Zhubr's classification applies
to the already constructed covers.
\item Euler multiplicativity gives $\chi=2d$.  A connected sum of
$g$ copies of $S^3\times S^3$ has Euler number $2-2g$; for
$d>1$ the equality would require $g=1-d<0$.  This supplies the
numerical obstruction to that proposed connected-sum type.
\end{enumerate}

\subsection*{8. The other elliptic surfaces and their six-manifolds}
\RungZero{The primitive specialization lattices and actual normal-circle-torus
maps classify the marked natural-line sphere outputs on all
twenty-one rank-one Mordell--Weil profiles: $59$ patterns on $50$
initial rows, with twelve profiles of unbounded sphere multiplicity,
four of logarithmic order one and five with no sphere output.
Their integral cokernels and full characteristic evaluations also
identify the minimal simply connected nonspherical outputs and
weight-$d$ Engel manifolds, and the free OS49 Gysin lattice retains
the extra torsion outside the coprime sector
(\cref{maps:src:D.2.28,maps:src:D.2.44,maps:src:D.2.45,maps:src:D.2.46,maps:src:D.2.47,maps:src:D.2.48,night:P3F.34}).}
\RungOne
\begin{enumerate}[label=\arabic*.,leftmargin=*]
\item Fix the initial section pair, the component on which the
linearization vanishes, the marked filling and every affine
logarithmic vector.  The profile is the elliptic-fibre configuration;
a pattern is the entire marked initial row, with its stated
arithmetic conditions.
\item Compute the primitive specialization lattices and the actual
normal-circle-torus maps.  Regular-fibre restriction is carried
through the boundary and Gysin maps before the finite-coefficient
classes are used.  Their integral cokernels retain all torsion.
\item The all-order calculations for the original, III-star,
Liu--Rollenske and Wang families give their explicit determinant
or normal-torus matrices.  The additive and multiplicative local
calculations close the remaining profiles; the complete table in
\cref{maps:src:D.2.28} records every pattern.
\item In the simply connected residual cases retain the spin
characteristic class, Zhubr's refined class and the full
fundamental-class image.  The conic surgery and the marked
$Z_{c,\epsilon}$ surgeries give explicit representatives of the
types in \cref{maps:src:D.2.44,maps:src:D.2.45}.
\item For a period sublattice of the cusp, intersect the component
restrictions on each degree-six del Pezzo normalization.  Keep
loops, repeated edges and repeated vertices in this calculation;
the complete trilinear expression is
\cref{maps:src:D.2.47}.
\item In the simply connected weight-$d$ sector the Seifert
determinant is $P=\pm1$.  The full invariant includes the
Pontryagin square modulo $d\gcd(d,2)$ and the third power
modulo $d\gcd(d,3)$, with evaluations $-12$ and $0$.
Together with the cyclic divisor matrix and refined class zero,
these give the single type $M_d$ in \cref{maps:src:D.2.48}.
\item For a free OS49 cap $(m,n,d)=(3k,n,d)$ with $n\ne0$,
put $g=\gcd(3k,n,3d)$.  The exact primitive Gysin lattice has
Smith invariants $g,|n|$, giving
$H_2=H_3=\Z/|n|\oplus\Z/g$ with $g\in\{1,3\}$.
For example $(9,3,1)$ gives $(\Z/3)^2$.
The zero-$n$ calculation and the precise freeness conditions are
retained in \cref{night:P3F.33,night:P3F.34}.
\end{enumerate}

\subsection*{9. The conifold transition}
\RungZero{Moving the zero of the line-bundle linearization gives the local
equation $xy=u(zv-s)$ and a proper analytic family with one ordinary
double point. The actual replacement of $S^2\times D^4$ by
$D^3\times S^3$ and its primitive boundary maps identify its two
smooth sides as the original six-sphere and $S^3\times S^3$,
the latter with algebraic dimension one
(\cref{maps:src:D.3.1,maps:src:D.3.2,maps:src:D.3.3}).}
\RungOne
\begin{enumerate}[label=\arabic*.,leftmargin=*]
\item Start with the marked elliptic surface and the line
$M=\OO(P-O)$.  Move the linearization as $\phi_s=\phi_0-s\phi_1$,
using its exact annular identification on the common exterior.
\item Near the collision the equation is
$xy=u(zv-s)$.  At $s=0$ this is the ordinary double point;
its small resolution has exceptional normal bundle
$\OO(-1)\oplus\OO(-1)$.  The root caps remain in the same
proper family because the collars were chosen uniformly.
\item On real manifolds the transition replaces $S^2\times D^4$
by $D^3\times S^3$.  General position makes the complement
simply connected; Alexander duality and the two primitive boundary
maps give $H_3=\Z^2$ and $H_2=H_4=0$.
\item The spin Wall invariant has the standard rank-two skew form
and zero normal invariant.  The classification, together with
$\Theta_6=0$, identifies the smoothing with $S^3\times S^3$.
\item The root canonical formula and the varying periods give
algebraic dimension one.  The formal cohomology of the exceptional
curve transports the stated structure-sheaf cohomology across the
small resolution.  The comparison with the published CDP quotient
uses its actual algebraic-reduction fibres in
\cref{maps:src:D.3.4}.
\end{enumerate}

\subsection*{10. A fourfold obtained by a transverse fibre product}
\RungZero{The explicit crepant small resolution of the transverse elliptic
fibre product gives a compact complex fourfold with
$K=\Pi^*L^{-5}$ and algebraic dimension two; its Todd and
Pontryagin calculations give signature two and middle form
$\langle1\rangle\oplus\langle1\rangle$.
The computed values $b_1=0$ and $h^{0,1}=1$, preserved by a
K\"ahler modification, obstruct membership in Fujiki's class
$\mathcal C$ through the K\"ahler equality $b_1=2h^{0,1}$
(\cref{maps:src:E.1.2,maps:src:E.2.1}).}
\RungOne
\begin{enumerate}[label=\arabic*.,leftmargin=*]
\item Pull back the Weierstrass family to the root curve of the
threefold fibration.  The integral matrices $A,B$ give the actual
period comparison $AZ_3B^{-1}=\operatorname{diag}(Z_2,s+c_E)$.
\item The unimodular Kuhn fan resolves the local fibre product
without adding rays.  Its period action is proper and its marked
section extends face by face.  The resolution is therefore small
and crepant.
\item The canonical formula on the root curve gives
$K=\Pi^*L^{-5}$.  The split relative elliptic complex and Leray
give $(h^{0,q})=(1,1,0,0,0)$; the explicit meromorphic product
comparison gives algebraic dimension two.
\item A K\"ahler modification in class $\mathcal C$ preserves
$b_1$ and $h^{0,1}$.  Here $b_1=0$ and $h^{0,1}=1$, contradicting
the K\"ahler equality $b_1=2h^{0,1}$.  This proves the stated
exclusion through these invariants.
\item The universal Todd and Pontryagin formulas give
$c_2^2=2$, $c_4=6$ and signature two.  The middle lattice is
positive definite unimodular of rank two, hence has the displayed
form.  The body records the precise scope of each cited formula.
\end{enumerate}

\paragraph{Supporting constructions and remaining scope.}
The later sections keep the complete transverse-pullback,
covering-space and cusp-carrier arguments, including their explicit
quotient maps, and the marked realization of the two collision
profiles.  The remaining classification question is recorded as
\cref{maps:src:D.2.4}; it is kept separate from the proved sectors
and from the current proof-review records.

\section{Common conventions and the transport dictionary}\label{sec:III:conventions}
The base is $B=\mathbb P^1_t$, with $(p_1,p_2,p_0)=(0,1,\infty)$,
and the regular-fibre lattice is
$\Lambda=\mathbb Z\langle\widehat\gamma,\widehat u,\widehat w,\widehat\delta\rangle$.
The covector $\gamma$ is dual to $\widehat\gamma$.
The family homology matrices are $A_1,A_2,M_0$, with
$A_1A_2M_0=I$; their cohomology transport matrices are
$T_j=A_j^{-t}$. An action of $A_j$ on a cohomology class means its
induced pullback, not multiplication of a homology column by that class.
Base meridians are counterclockwise. The analytic root generators are
clockwise, $s_j(g_jz)=e^{-2\pi i/m_j}s_j(z)$, and
$g_1g_2g_0=1$. The retained multiplicities are $(3,4)$, with
$f^*p_j=m_jS_j$ and reduced cusp $W=f^{-1}(p_0)$.
We use
\[
\Pi_c=\begin{pmatrix}6\mu&\tau&1&0\\\beta^{\rm part}+c&\mu&0&1\end{pmatrix},
\qquad w=e^{-\pi i c},\qquad
D_{\rm part}=\operatorname{Im}\beta^{\rm part}
-\frac{6(\operatorname{Im}\mu)^2}{\operatorname{Im}\tau}.
\]
The normalization is $(\beta^{\rm part}+\tau)(p_0)=0$ and
$M=\max D_{\rm part}=-3\log(6)/\pi$.
For a fixed member $D=D_{\rm part}+\operatorname{Im}c$;
admissibility is $\operatorname{Im}c<-M$.
The positive $w$-loop lifts to $c(r)=c_0-2r$.

Engel's indices $(0,1,\infty)$ are our $(1,2,0)$.
His line bundle is denoted $M_S=\mathcal O_S(P-O)$.
The common marking matrix is
\[
P_E=\begin{pmatrix}0&-2&-1&-2\\2&-1&0&-1\\0&2&1&3\\-1&0&0&0\end{pmatrix},
\qquad \psi P_E=-\gamma,\qquad a_{j-1}=-P_Ev_j/m_j.
\]
The standard twists are $v_1=(1,2,-4,0)^t$ and
$v_2=(-1,-3,3,0)^t$. For general twists
$p=12l_0-4l_1-3l_2$ and $\pi_1(X_p)=\mathbb Z/|p|$.
The vector shorthand $(g,u,w,d)$ in fibre calculations denotes the ordered
homology basis, never the scalar parameter $w$ in a period expression.

\paragraph{Rank-six matrices.}
In the fourfold construction we use $M_i$ for actual homology transport
and $T_i=M_i^{-t}$ for left cohomology transport. The explicitly constructed
matrices named $A,B$ in the rational normal-form calculation are geometric
cohomology pullbacks: their homology inputs are $(A^t,B^t)$ and their left
cohomology transport is $(A^{-1},B^{-1})$. Every peripheral product is kept
in its stated order; reversing generators also conjugates or inverts the
cusp as specified in the proof. The fixed Kuhn cap and its full marked
boundary are part of a gluing datum. Affine multiplicity, linear order,
translation numerator and ordinary clutching are separate inputs.

All homology and cohomology have integral coefficients unless a coefficient
group is printed. We use the complex orientation, outward stabilization,
and the displayed ordered normal coordinates. Local applications of letters
for an integer matrix, a cap or a lattice are explicitly defined at their
use; the complete symbol and source-label dictionary accompanies the papers.


\section{Sphere maps, framings and homotopy classes}\label{sec:III:maps}
\begin{proposition}[Counterexample: as an unrestricted uniqueness assertion]\label{maps:src:A.2.2}
\StatementRecord{maps:src:A.2.2}{proved}{Historical events: Claude Fable 5.2 (Anthropic), effort high, 16 November 2026: carried as background without fault found. GPT-6 Astra (OpenAI), ChatGPT Pro mode, effort Pro, S6\_09\_PAPERS, 17 November 2026: targeted proof check. Also checked in this writing. Pre-continuation writing events: GPT-6 Astra (OpenAI), ChatGPT Pro mode, effort Pro, S6\_10\_PAPERS, 18 November 2026: targeted proof check. Continuation metadata events: GPT-6 Astra (OpenAI), ChatGPT Pro mode, effort Pro, S6\_10\_PAPERS continuation, 19 November 2026: used and rechecked in this run; status-vocabulary reconciliation; metadata reconciliation only; no new mathematical check. \emph{Pre-continuation recorded proof scope.} GPT-6 Astra (OpenAI, ChatGPT Pro mode), effort Pro, 2026-11-17; Claude Fable 5.2 (Anthropic), effort high, 16 November 2026; also checked in this writing by GPT-6 Astra (OpenAI), ChatGPT Pro mode, effort Pro, 18 November 2026}{the hypotheses, notation and definitions stated here; the explicit constructions and calculations in the proof}


Agreement with the literal quotient outside a neighbourhood
of $W$, without a condition on the image of that neighbourhood,
does not determine a homotopy class $S^6\to S^3$.  All twelve
classes occur among maps satisfying that bare agreement
condition.

\end{proposition}


\begin{proof}
\begin{enumerate}[label=\arabic*.,leftmargin=*]
\item Take a closed embedded six-ball in the interior of the cusp
neighbourhood.  Homotope a given extension on a smaller ball
to be constant there, keeping its boundary and the complement
fixed; a target coordinate ball around the image of its centre
provides the contraction.  Pinch the smaller ball to
$S^6\vee D^6$ relative to its boundary and insert any based
map $\alpha:S^6\to S^3$ on the new sphere.  The resulting
global map agrees with the original map off that ball, and
represents $[e]+\alpha$ under $[S^6,S^3]=\pi_6(S^3)$.
Distinct $\alpha$ give distinct classes.  Smooth relative
approximation preserves this conclusion.  The insertion
cannot be performed within the cusp solid torus, since that
solid torus has $\pi_6=0$; this is exactly the distinction
in \cref{mapsI:src:A.2.1}.
\end{enumerate}
\end{proof}

The corrected content of \cite[Remark 8.3]{mapsEngel} is
therefore a quotient over the regular part and the two finite
fillings, followed by a replacement within its circle homotopy
class over a cusp collar.  The literal character does not
extend.  The replacement gives a canonical homotopy class
when the solid-torus image condition is retained.


\phantomsection\label{qa:III:layout:12}
\subsection{The framed three-torus and the class of \texorpdfstring{$e$}{e}}

We use $\pi_3^{\mathrm S}=\mathbb Z/24\langle\nu\rangle$, with
$e_{\mathbb R}(\nu)=1/24$, and
$\pi_6(S^3)=\mathbb Z/12\langle\nu'\rangle$, where the
stabilization of $\nu'$ is $2\nu$.  We write
$\eta_n:S^{n+1}\to S^n$ for the $(n-2)$-fold suspension of
the complex Hopf map.  The order-two classes below are
unaffected by reversal of either source or target orientation.
An orientation-preserving diffeomorphism to the standard sphere
transports its bounding stable framing to $TX\oplus\mathbb R$.
Its homotopy class does not depend on that diffeomorphism:
two oriented stable trivializations differ by a map
$S^6\to\mathrm{SO}$, and $\pi_6(\mathrm{SO})=0$
by \cite[Corollary to Theorem II and the table on p.~315]{mapsBott59}.
We use outward stabilization and the convention that the normal
frame precedes the tangential frame; a constant permutation is
corrected once to preserve orientation.

\begin{lemma}\label{maps:src:B.1.1}
\StatementRecord{maps:src:B.1.1}{proved}{Historical events: Claude Fable 5.2 (Anthropic), effort high, 16 November 2026: carried as background without fault found. GPT-6 Astra (OpenAI), ChatGPT Pro mode, effort Pro, S6\_09\_PAPERS, 17 November 2026: targeted proof check. Also checked in this writing. Pre-continuation writing events: GPT-6 Astra (OpenAI), ChatGPT Pro mode, effort Pro, S6\_10\_PAPERS, 18 November 2026: targeted proof check. Continuation metadata events: GPT-6 Astra (OpenAI), ChatGPT Pro mode, effort Pro, S6\_10\_PAPERS continuation, 19 November 2026: used and rechecked in this run; status-vocabulary reconciliation; metadata reconciliation only; no new mathematical check. \emph{Pre-continuation recorded proof scope.} GPT-6 Astra (OpenAI, ChatGPT Pro mode), effort Pro, 2026-11-17; Claude Fable 5.2 (Anthropic), effort high, 16 November 2026; also checked in this writing by GPT-6 Astra (OpenAI), ChatGPT Pro mode, effort Pro, 18 November 2026}{the hypotheses, notation and definitions stated here; the explicit constructions and calculations in the proof}


Let $F$ be a regular fibre of the original fibration.  Compare
the ambient stable framing of the oriented sphere $X$ with
the framing on $TX|_F$ consisting of four translation-invariant
vertical vector fields and an oriented lifted base frame.
The restriction of their stable difference to any coset of
$T_L\subset F$ is null-homotopic.

\end{lemma}


\begin{proof}
\begin{enumerate}[label=\arabic*.,leftmargin=*]
\item After a constant orientation adjustment the comparison is
a map $D:F\to\mathrm{SO}$, where stabilization is understood.
Its homotopy class is invariant under both monodromies.
Indeed, transport $F$ around a based loop in $B^\circ$ using
the flat real marking.  The translated vertical frame varies
continuously on $F\times[0,1]$; at the endpoint it changes by
the constant real matrix of the affine monodromy.  An oriented
base frame can be chosen periodic along the loop.  The space
of lifts of that frame to $TX$ is affine, so changing the lift
does not affect the comparison class.  The ambient stable
framing is already defined along this entire cylinder.  Thus
the comparison at its endpoints is $D$ and $D\circ A_j$
times a constant positive matrix.  Every such constant
matrix is path connected to the identity.  This proves
$[D]=A_j^*[D]$.  An affine translational part is isotopic
to the identity and has no effect.

\item The map on fundamental groups defines
$\sigma(D)\in H^1(F;\mathbb Z/2)$.  The exhaustive calculation
in \texttt{maps\_exact.py} gives
\[
 H^1(F;\mathbb Z/2)^{\langle A_1,A_2\rangle}
      =\{0,\gamma\bmod2\}.
\]
Both classes vanish on $L$, so $D|_{T_L}$ lifts to
$\mathrm{Spin}$.  Up to dimension four the latter has
$\pi_1=\pi_2=\pi_4=0$, $\pi_3=\mathbb Z$.
Obstruction theory consequently identifies maps from a
three- or four-dimensional torus to $\mathrm{Spin}$ with
$H^3$ of that torus.  The two possible lifts differ by
multiplication by the central element $-1$, which is
homotopic to the identity, so this identification has no
unrecorded lifting ambiguity.

\item The doubled stable difference $2[D]$ has zero $\sigma$,
and therefore corresponds to a monodromy-invariant class
in $H^3(F;\mathbb Z)$.  Integral Poincare duality, natural
for these determinant-one monodromies, identifies the
invariants in $H^3(F;\mathbb Z)$ with the fixed vectors
$\mathbb Zd\subset H_1(F;\mathbb Z)$.  Such a class
evaluates to zero on $T_L$, since $d\in L$:
its evaluation is the determinant with columns $d,u,w,d$.
Hence $2[D|_{T_L}]=0$ in $[T_L,\mathrm{Spin}]=\mathbb Z$.
This group has no torsion; thus $[D|_{T_L}]=0$.
The same argument applies to a coset since translation is
isotopic to the identity.
\end{enumerate}
\end{proof}

\begin{theorem}\label{maps:src:B.1.2}
\StatementRecord{maps:src:B.1.2}{proved}{Historical events: Claude Fable 5.2 (Anthropic), effort high, 16 November 2026: carried as background without fault found. Claude Fable 5.2 (Anthropic), effort high, 17 November 2026: carried as background without fault found. GPT-6 Astra (OpenAI), ChatGPT Pro mode, effort Pro, S6\_09\_PAPERS, 17 November 2026: targeted proof check. Also checked in this writing. Pre-continuation writing events: GPT-6 Astra (OpenAI), ChatGPT Pro mode, effort Pro, S6\_10\_PAPERS, 18 November 2026: targeted proof check. Continuation metadata events: GPT-6 Astra (OpenAI), ChatGPT Pro mode, effort Pro, S6\_10\_PAPERS continuation, 19 November 2026: used and rechecked in this run; status-vocabulary reconciliation; metadata reconciliation only; no new mathematical check. \emph{Pre-continuation recorded proof scope.} GPT-6 Astra (OpenAI, ChatGPT Pro mode), effort Pro, 2026-11-17; Claude Fable 5.2, effort high, 2026-11-17; Claude Fable 5.2 (Anthropic), effort high, 16 November 2026; also checked in this writing by GPT-6 Astra (OpenAI), ChatGPT Pro mode, effort Pro, 18 November 2026}{the hypotheses, notation and definitions stated here; \cref{maps:src:B.1.1}; \cite[Theorem (4.14), p.~423]{mapsAPS2}; \cite[Theorem 4.1 and Corollaries 4.2--4.3]{mapsToda52}}


For the repaired quotient of \cref{mapsI:src:A.2.1} with $|p|=1$,
\[
 [e]=6\nu'=\eta_3\eta_4\eta_5\ne0
       \quad\text{in }\pi_6(S^3)=\mathbb Z/12.
 \tag{B.1.1}\label{maps:eq:B.1.1}
\]
Its stable class is $12\nu$ and its real Adams invariant
is $1/2\in\mathbb Q/\mathbb Z$.

\end{theorem}


\begin{proof}
\begin{enumerate}[label=\arabic*.,leftmargin=*]
\item Choose a point of $Q$ over a regular base value outside
the modified cusp disc and outside the two finite root
discs.  Its entire inverse image is one coset of $T_L$.
Order a target frame as the positive real base frame followed
by the positive $\gamma$-direction.  Lifting this frame
under $de$ gives a normal three-framing.  Complete it by
the constant Lie frame of $L$ with the unique constant
orientation adjustment making the full frame positive.
This is precisely the restriction of the product frame
in Lemma \ref{maps:src:B.1.1}, up to a constant positive change of basis.
That lemma proves that the stable tangential framing of
the inverse image is the Lie framing of $T^3$.
Pontryagin--Thom therefore gives
\[
 \operatorname{st}[e]=[T^3,\mathrm{Lie}]
 = [S^1,\mathrm{Lie}]^3=\eta^3=12\nu.
 \tag{B.1.2}\label{maps:eq:B.1.2}
\]
The product equality follows by taking the product of the
three framed circles and their tubular collapse maps.
The Lie-framed circle is the nonbounding framed circle;
its class is the stabilized Hopf element $\eta$.

\item Here is an independent evaluation of the framed bordism
class, which does not use the ring relation $\eta^3=12\nu$.
Put the flat product metric on $\mathbb R^3/\mathbb Z^3$,
with the spin structure determined by its Lie framing.
Its spinors are periodic $\mathbb C^2$-valued functions.
On the Fourier mode $e^{2\pi i k\cdot x}$, $k\in\mathbb Z^3$,
the Dirac operator is the Hermitian matrix
$2\pi\sum k_j\sigma_j$, with eigenvalues
$\pm2\pi|k|$ for $k\ne0$ and a two-dimensional kernel
for $k=0$.  Thus its eta function is identically zero in
its convergent half-plane and hence at zero, while $h=2$.
The Levi-Civita connection and the connection induced
by the framing coincide, so their transgression term
vanishes.  The real formula of
\cite[Theorem (4.14), p.~423]{mapsAPS2}, with its factor
$\epsilon(1)=1/2$, gives
\[
 e_{\mathbb R}(T^3,\mathrm{Lie})
       =\frac12\frac{0+2}{2}=\frac12\pmod{\mathbb Z}.
 \tag{B.1.3}\label{maps:eq:B.1.3}
\]
This detects $12\nu$ in $\mathbb Z/24$.  The complex
$e$-invariant would be zero here; omitting the real
factor would therefore lose the desired detection.

\item Finally stabilization is injective on $\pi_6(S^3)$.
For the exact unstable input we use
\cite[Theorem 4.1 and Corollaries 4.2--4.3]{mapsToda52}:
the suspended generator is the torsion term in
$[\iota_4,\iota_4]=2\nu_4-E\nu'$, and stabilization
sends it to $2\nu$.  In the group of order twelve this
is injective.  Therefore the class with stable image
$12\nu$ is $6\nu'$.  The composite
$\eta_3\eta_4\eta_5$ has the same stable image, so the
same injection identifies it with $6\nu'$.
\end{enumerate}
\end{proof}

\subsection{The holomorphic map to the base}

\begin{theorem}\label{maps:src:B.2.1}
\StatementRecord{maps:src:B.2.1}{proved}{Historical events: Claude Fable 5.2 (Anthropic), effort high, 16 November 2026: carried as background without fault found. Claude Fable 5.2 (Anthropic), effort high, 17 November 2026: carried as background without fault found. GPT-6 Astra (OpenAI), ChatGPT Pro mode, effort Pro, S6\_09\_PAPERS, 17 November 2026: targeted proof check. Also checked in this writing. Pre-continuation writing events: GPT-6 Astra (OpenAI), ChatGPT Pro mode, effort Pro, S6\_10\_PAPERS, 18 November 2026: targeted proof check. Continuation metadata events: GPT-6 Astra (OpenAI), ChatGPT Pro mode, effort Pro, S6\_10\_PAPERS continuation, 19 November 2026: used and rechecked in this run; status-vocabulary reconciliation; metadata reconciliation only; no new mathematical check. \emph{Pre-continuation recorded proof scope.} GPT-6 Astra (OpenAI, ChatGPT Pro mode), effort Pro, 2026-11-17; Claude Fable 5.2, effort high, 2026-11-17; Claude Fable 5.2 (Anthropic), effort high, 16 November 2026; also checked in this writing by GPT-6 Astra (OpenAI), ChatGPT Pro mode, effort Pro, 18 November 2026}{the hypotheses, notation and definitions stated here; the explicit constructions and calculations in the proof}


The map $f:X\to\mathbb P^1$ of the original sphere
construction is null-homotopic.  This conclusion does not
require the choice of $e$.

\end{theorem}


\begin{proof}
\begin{enumerate}[label=\arabic*.,leftmargin=*]
\item The quotient projection $h:Q\to B$ is, after the
orientation-preserving identifications fixed in \cref{mapsI:src:A.1.3},
the weighted Hopf map
$[z_1^3:z_2^4]$, followed by a Mobius transformation
moving its two exceptional values to $0,1$.
It is the ordinary Hopf map after the self-map
\[
 P_{3,4}(z_1,z_2)=
 \frac{(z_1^3,z_2^4)}{\sqrt{|z_1|^6+|z_2|^8}}
 \quad\text{of }S^3.
 \tag{B.2.1}\label{maps:eq:B.2.1}
\]
At a value with both coordinates nonzero, its phase
map has $3\cdot4=12$ preimages, all with positive local
degree; its radial coordinate is strictly increasing.
Consequently its degree is twelve.  A degree-$k$ map
of $S^3$ is homotopic to the $k$th power map of the
topological group $S^3$.  On $\pi_6(S^3)$ the latter
is multiplication by $k$, since pointwise multiplication
and the homotopy-group addition agree.  Thus
$h_* = (\eta_2)_*\circ12=0$ on $\pi_6(S^3)$.
The construction of $e$ is over $B$, so $f=h\circ e$
as actual maps.  This proves $[f]=0$.

\item For a second proof, let $S_1,S_2$ be the reduced multiple
fibres.  The line bundles $\mathcal O_X(S_j)$ are
topologically trivial because $H^2(X;\mathbb Z)=0$.
The coordinate sections of $f^*\mathcal O(1)$ have
divisors $3S_1$ and $4S_2$ after a Mobius change of
base coordinate.  Choose a smooth frame of
$f^*\mathcal O(1)$.  Since $H^1(X;\mathbb Z)=0$, every
map $X\to\mathbb C^*$ admits a logarithm.  We can
therefore choose frames of $\mathcal O_X(S_1)$ and
$\mathcal O_X(S_2)$ whose third and fourth powers,
under the given divisor isomorphisms, equal the chosen
frame.  The canonical divisor sections then become
smooth functions $a,b$ with no common zero and
$f=[a^3:b^4]$.  Put
$G=(a,b)/\sqrt{|a|^2+|b|^2}:X\to S^3$.
Writing $r=\sqrt{|a|^2+|b|^2}>0$, the explicit homotopy
$[a^3/r^{3s}:b^4/r^{4s}]$, $0\leq s\leq1$, gives
$f\simeq\eta_2P_{3,4}G$.  It is defined everywhere because
$a$ and $b$ have no common zero.  The latter map is zero
by the same degree calculation and the exponent twelve
of $\pi_6(S^3)$.  This proof uses neither the cusp
replacement nor the computed class of $e$.
\end{enumerate}
\end{proof}

\subsection{The four workbench maps}

The exact maps, including their suspension conventions, are
needed here.  Put $\overline T=(\overline\Lambda\otimes\mathbb R)
/\overline\Lambda$ and let $\overline\pi:N_0\setminus W
\to\overline T$ be the long-period quotient.  On the punctured
cusp write $q=re^{i\phi}$.  The map $q_2$ takes $x$ to
$[r,e^{i\phi},\overline\pi(x)]$ in the unreduced suspension
of $S^1\times\overline T$, sends $W$ to its lower vertex,
and sends the complement of the cusp neighbourhood to its
upper vertex.  Collapsing the two-skeleton of
$S^1\times\overline T=T^3$ and then taking the reduced
suspension gives $H:X\to S^4$.  Collapsing the one-skeleton
of $\overline T=T^2$ and forgetting $\phi$ gives
$\widehat q_3:X\to S^3$.  Retaining only $\gamma$ and
forgetting $\phi$ gives $\widehat q_\gamma:X\to S^2$.
These formulas specify the marked maps, including the retained coordinates
and quotient identifications; arbitrary sphere quotients of the same targets
are not substituted for them.

\begin{theorem}\label{maps:src:B.3.1}
\StatementRecord{maps:src:B.3.1}{proved}{Historical events: Claude Fable 5.2 (Anthropic), effort high, 16 November 2026: carried as background without fault found. GPT-6 Astra (OpenAI), ChatGPT Pro mode, effort Pro, S6\_09\_PAPERS, 17 November 2026: targeted proof check. Also checked in this writing. Pre-continuation writing events: GPT-6 Astra (OpenAI), ChatGPT Pro mode, effort Pro, S6\_10\_PAPERS, 18 November 2026: targeted proof check. Continuation metadata events: GPT-6 Astra (OpenAI), ChatGPT Pro mode, effort Pro, S6\_10\_PAPERS continuation, 19 November 2026: used and rechecked in this run; status-vocabulary reconciliation; metadata reconciliation only; no new mathematical check. \emph{Pre-continuation recorded proof scope.} GPT-6 Astra (OpenAI, ChatGPT Pro mode), effort Pro, 2026-11-17; Claude Fable 5.2 (Anthropic), effort high, 16 November 2026; also checked in this writing by GPT-6 Astra (OpenAI), ChatGPT Pro mode, effort Pro, 18 November 2026}{the hypotheses, notation and definitions stated here; \cref{maps:eq:B.3.2}; \cref{maps:eq:B.3.3}}


For the exact retained-angle map just defined,
\[
 [H]=\eta_4\eta_5\ne0\in\pi_6(S^4)=\mathbb Z/2.
 \tag{B.3.1}\label{maps:eq:B.3.1}
\]
Its regular inverse image is the native vanishing
two-torus with its Lie stable framing and Arf invariant one.

\end{theorem}


\begin{proof}
\begin{enumerate}[label=\arabic*.,leftmargin=*]
\item Take $u_0=(-1/4,1/4)$, $q_0=r_0e^{i\pi}$, and
$p_0=B_0u_0=(1/4,1/4)$.  For $r_0$ sufficiently small
the full inverse image of the corresponding open-cell
value lies in the original chart of \cref{mapsI:src:A.1.2}.  To see this
without losing components, a lift of any point of that
inverse image has long coordinate $u_0+n$ for a unique
$n\in\mathbb Z^2$; applying the actual period deck
transformation $-n$ changes it to $u_0$.  Distinct lifts
with this coordinate cannot be deck equivalent.
Put $h_j(q)=\exp(-2\pi i(C(q)u_0)_j)$ and make the
holomorphic coordinate change
\[
 (w_0,w_1,w_2)
   =(z_0/(h_1h_2),h_1z_1,h_2z_2).
\]
It preserves the product $q$.  The full inverse image is
\[
 w(a,b)=
 (r_0^{1/2}e^{i(\pi-a-b)},
  r_0^{1/4}e^{ia},r_0^{1/4}e^{ib}),
 \quad a,b\in\mathbb R/2\pi\mathbb Z.
 \tag{B.3.2}\label{maps:eq:B.3.2}
\]
All coordinates tend uniformly to zero as $r_0\to0$.
The period quotient is injective on a neighbourhood of
the chart origin, so this torus is contained in a
contractible coordinate ball of $X$.
The map $H$ is smooth on an open neighbourhood of this
torus, although its suspension presentation is only
continuous elsewhere.  Choose a smaller open-cell target
ball whose full inverse image lies in that neighbourhood;
compactness permits this because the displayed torus is
the full inverse image of its centre.  Relative smooth
approximation, uniformly small off that neighbourhood,
then gives a smooth representative with exactly the same
inverse image and differential there.  Pontryagin--Thom
therefore applies to the framing computed below.

\item We compute its framing before discarding the bounded
period terms.  Write
\[
 \ell=\log|q|,\quad
 R(q)=-2\pi\operatorname{Im}C(q),\quad
 A(q)=\ell B_0+R(q),\quad
 v=(\log|w_1|,\log|w_2|)^t.
\]
The exact coordinate identity is
$v=\ell B_0u+R(q)(u-u_0)$, and on \eqref{maps:eq:B.3.2} it gives
\[
 du=A(q_0)^{-1}(dv-p_0\,d\ell).
 \tag{B.3.3}\label{maps:eq:B.3.3}
\]
Thus no derivative of a bounded correction has been
ignored.  Let $L_j$ and $Q_j$ denote the real tangent
vectors $w_j\partial_{w_j}$ and $iw_j\partial_{w_j}$
in the real identification $T\mathbb C^3=\mathbb C^3$.
With $A_0=A(q_0)$, dual lifts of the target coordinate
order $(\ell,\phi,u_1,u_2)$ are
\[
 N_\ell=\tfrac12L_0+\tfrac14L_1+\tfrac14L_2,
 \qquad N_\phi=Q_0,
\]
\[
 N_i=-((A_0)_{1i}+(A_0)_{2i})L_0
            +(A_0)_{1i}L_1+(A_0)_{2i}L_2,
 \quad i=1,2.
 \tag{B.3.4}\label{maps:eq:B.3.4}
\]
The tangent vectors are $T_a=-Q_0+Q_1$ and
$T_b=-Q_0+Q_2$.  Formula \eqref{maps:eq:B.3.3} checks that the
four normal vectors map to the four standard coordinate
vectors and both tangent vectors map to zero.
The full determinant is
$-\det(A_0)|q_0|^2<0$, so the positive tangent order is
$(-T_a,T_b)$.  The inequality holds for small $r_0$
because $A_0$ is joined to $\ell_0B_0$ through
$\ell_0B_0+sR(q_0)$, all of which are invertible.

\item Every displayed vector is a constant linear combination
of $L_j,Q_j$.  Its dependence on $(a,b)$ is therefore
exactly
\[
 D(a,b)=\operatorname{diag}(e^{-i(a+b)},e^{ia},e^{ib})
            \in\mathrm{SU}(3).
\]
This map is null-homotopic.  An explicit periodic
contraction uses
\[
 Q_s(\theta)=
 \begin{pmatrix}
 \sin(\pi s/2)e^{-i\theta}&-\cos(\pi s/2)\\
 \cos(\pi s/2)&\sin(\pi s/2)e^{i\theta}
 \end{pmatrix},
 \qquad K_s(\theta)=Q_s(\theta)Q_s(0)^*.
\]
Each matrix lies in $\mathrm{SU}(2)$,
$K_0=I$, and $K_1(\theta)=\operatorname{diag}
(e^{-i\theta},e^{i\theta})$.  The product of its
embeddings in coordinate planes $(0,1)$ and $(0,2)$,
evaluated at $a$ and $b$, contracts $D$.
Consequently the induced stable tangential frame of
\eqref{maps:eq:B.3.2} is the constant Lie frame.  The ambient framing
comparison on the containing ball is null-homotopic,
so the same conclusion holds in the framed sphere.

\item For the induced spin structure the two coordinate
circles are nonbounding.  Its quadratic form on
$H_1(T^2;\mathbb F_2)$ is
$q(xa+yb)=x+xy+y$, and its Arf invariant is
$q(a)q(b)=1$.  Hence its framed class is
$\eta^2\in\pi_2^{\mathrm S}$.  Freudenthal suspension
is an isomorphism for $\pi_6(S^4)$, so the unstable
class is exactly \eqref{maps:eq:B.3.1}.
\end{enumerate}
\end{proof}

\begin{theorem}\label{maps:src:B.3.2}
\StatementRecord{maps:src:B.3.2}{proved}{Historical events: Claude Fable 5.2 (Anthropic), effort high, 16 November 2026: carried as background without fault found. GPT-6 Astra (OpenAI), ChatGPT Pro mode, effort Pro, S6\_09\_PAPERS, 17 November 2026: targeted proof check. Also checked in this writing. Pre-continuation writing events: GPT-6 Astra (OpenAI), ChatGPT Pro mode, effort Pro, S6\_10\_PAPERS, 18 November 2026: targeted proof check. Continuation metadata events: GPT-6 Astra (OpenAI), ChatGPT Pro mode, effort Pro, S6\_10\_PAPERS continuation, 19 November 2026: used and rechecked in this run; status-vocabulary reconciliation; metadata reconciliation only; no new mathematical check. \emph{Pre-continuation recorded proof scope.} GPT-6 Astra (OpenAI, ChatGPT Pro mode), effort Pro, 2026-11-17; Claude Fable 5.2 (Anthropic), effort high, 16 November 2026; also checked in this writing by GPT-6 Astra (OpenAI), ChatGPT Pro mode, effort Pro, 18 November 2026}{the hypotheses, notation and definitions stated here; the explicit constructions and calculations in the proof}


The angle-forgetting workbench maps satisfy
$[\widehat q_3]=0$ and $[\widehat q_\gamma]=0$.

\end{theorem}


\begin{proof}
\begin{enumerate}[label=\arabic*.,leftmargin=*]
\item Both components of $\overline\pi$ extend as integral
$H^1$ classes over $N_0$, by (A.1.5).  Choose a map
$\pi':N_0\to T^2$ representing them and equal to
$\overline\pi$ on an outer collar; the construction
and relative uniqueness are as in \cref{mapsI:src:A.2.1}.  On the
punctured cusp the quotient $\overline\pi-\pi'$ has
a lift to $\mathbb R^2$ that is zero on the collar.
Indeed the punctured cusp retracts onto its boundary
mapping torus, and the difference is zero in relative
$H^1$.  Multiplying that lift by a homotopy parameter
gives a homotopy from $\overline\pi$ to $\pi'$.
After passage to either suspension target this
homotopy is continuous at $W$, uniformly in the
homotopy parameter: the whole compact torus of
angles is collapsed as $r\to0$.  It is also constant
on the outer boundary and extends by the basepoint
outside $N_0$.

\item After this homotopy the three-sphere map factors as
\[
 X\longrightarrow
 \frac{[0,\epsilon]\times T^2}{\{\epsilon\}\times T^2}
 \longrightarrow\Sigma(T^2/T^{2,(1)}),
 \qquad x\longmapsto[r(x),\pi'(x)]
\]
on $N_0$, with the complement mapping to the vertex.
The middle space is a cone and is contractible.
For the two-sphere map use the first component
$\theta':N_0\to S^1$ and the cone on $S^1$.
These give actual null-homotopies of the two maps.
This argument cannot discard the retained base
angle in $H$, whose intermediate space is instead
$(D^2\times T^2)/(\partial D^2\times T^2)$.
\end{enumerate}
\end{proof}

\begin{theorem}\label{maps:src:B.3.3}
\StatementRecord{maps:src:B.3.3}{proved}{Historical events: Claude Fable 5.2 (Anthropic), effort high, 16 November 2026: carried as background without fault found. GPT-6 Astra (OpenAI), ChatGPT Pro mode, effort Pro, S6\_09\_PAPERS, 17 November 2026: targeted proof check. Also checked in this writing. Pre-continuation writing events: GPT-6 Astra (OpenAI), ChatGPT Pro mode, effort Pro, S6\_10\_PAPERS, 18 November 2026: targeted proof check. Continuation metadata events: GPT-6 Astra (OpenAI), ChatGPT Pro mode, effort Pro, S6\_10\_PAPERS continuation, 19 November 2026: used and rechecked in this run; status-vocabulary reconciliation; metadata reconciliation only; no new mathematical check. \emph{Pre-continuation recorded proof scope.} GPT-6 Astra (OpenAI, ChatGPT Pro mode), effort Pro, 2026-11-17; Claude Fable 5.2 (Anthropic), effort high, 16 November 2026; also checked in this writing by GPT-6 Astra (OpenAI), ChatGPT Pro mode, effort Pro, 18 November 2026}{the hypotheses, notation and definitions stated here; \cref{maps:src:B.3.2}; \cref{maps:src:B.1.2}}


Let $C=M\cup\operatorname{Cone}(\partial M)$ and let
$q_1:X\to C$ be the workbench map retaining the
Seifert quotient away from the cusp and collapsing
$W$ to the cone point.  Under the cofiber splitting
$C\simeq S^3\vee S^2$ determined by the cusp core,
\[
 [q_1]=(j_3)_*(6\nu'),
 \tag{B.3.5}\label{maps:eq:B.3.5}
\]
with no $S^2$ or mixed Whitehead component.

\end{theorem}


\begin{proof}
\begin{enumerate}[label=\arabic*.,leftmargin=*]
\item Write $Q=M\cup(D^2\times S^1)$ and let $K$ be the
core $\{0\}\times S^1$ of its cusp solid torus.
Collapsing this core turns the solid torus into
the cone on its boundary, by its radial coordinate.
Thus $C=Q/K$ as a quotient space.  The same
homotopy used in \ref{maps:src:B.3.2}, now retaining the base angle,
shows $q_1\simeq\kappa e$, where $\kappa:Q\to Q/K$
is the quotient map.  The inclusion $K=S^1\to S^3$
is null-homotopic, and its cofiber is
$S^3\vee\Sigma S^1$.  We choose this equivalence
to restrict on $Q$ to $j_3$.  A null-homotopy of
the circle exists since $\pi_1(S^3)=0$; its
ambiguity is $\pi_2(S^3)=0$.  Hence the indicated
splitting is fixed in the only respect used here,
and \ref{maps:src:B.1.2} proves \eqref{maps:eq:B.3.5}.
\end{enumerate}
\end{proof}

\subsection{The other named sphere fibrations}

\begin{theorem}\label{maps:src:B.4.1}
\StatementRecord{maps:src:B.4.1}{proved}{Historical events: Claude Fable 5.2 (Anthropic), effort high, 16 November 2026: carried as background without fault found. GPT-6 Astra (OpenAI), ChatGPT Pro mode, effort Pro, S6\_09\_PAPERS, 17 November 2026: targeted proof check. Also checked in this writing. Pre-continuation writing events: GPT-6 Astra (OpenAI), ChatGPT Pro mode, effort Pro, S6\_10\_PAPERS, 18 November 2026: targeted proof check. Continuation metadata events: GPT-6 Astra (OpenAI), ChatGPT Pro mode, effort Pro, S6\_10\_PAPERS continuation, 19 November 2026: used and rechecked in this run; status-vocabulary reconciliation; metadata reconciliation only; no new mathematical check. \emph{Pre-continuation recorded proof scope.} GPT-6 Astra (OpenAI, ChatGPT Pro mode), effort Pro, 2026-11-17; Claude Fable 5.2 (Anthropic), effort high, 16 November 2026; also checked in this writing by GPT-6 Astra (OpenAI), ChatGPT Pro mode, effort Pro, 18 November 2026}{the hypotheses, notation and definitions stated here; \cite[Lemma 2.1 and equations (2.11)--(2.12)]{mapsViaLog}; \cref{maps:src:B.1.1}; \cref{maps:src:B.1.2}; \cite[Proposition 4.2]{mapsVia2}; \cite[Proposition 5.4]{mapsVia2}; \cite[Lemma 3.10 and Proposition 3.11]{mapsLR}; \cite[Proposition 4.8]{mapsLR}; \cite[Propositions 3.12--3.13 and Corollary 3.16]{mapsLR}; \cref{maps:eq:B.4.1}; \cref{maps:eq:B.4.3}; \cref{maps:src:B.3.1}; \cref{maps:src:B.3.2}}


For the sphere members of Viaclovsky's two-parameter family,
the Liu--Rollenske family, and every $X_{3k,4l}$ with
$\gcd(3k,4l)=1$, the primitive invariant-circle quotient
admits a canonical replacement over the degeneration
neighbourhoods as in \cref{mapsI:src:A.2.1}.  Its target is $S^3$ and
its class is $6\nu'$.  In each case the map to
$\mathbb P^1$ is null-homotopic.  The retained-angle
map at the primitive rank-two cusp is $\eta_4\eta_5$,
and the two corresponding angle-forgetting maps are
null-homotopic.  Thus the quotient class $[e]$ does
not distinguish these fibrations.

\end{theorem}


\begin{proof}
\begin{enumerate}[label=\arabic*.,leftmargin=*]
\item We supply the integral lattices, the local extension
comparisons, and the target calculation; a list of
Betti numbers would not imply this statement.

\item For $X_{3k,4l}$ the regular homological representation
is still $(A_1,A_2,M_0)$ of \cref{mapsI:src:A.1.1}.  We use \cref{push:src:E.1.1} and \cref{push:src:E.1.2} with
its original finite-root convention.  In Viaclovsky's
marking $(\delta,a,b,d)$ the comparison is
\[
 P_V=\begin{pmatrix}-2&1&0&1\\-6&0&-1&0\\0&-1&0&0\\-1&0&0&0\end{pmatrix}.
\]
Substitution in the matrices of
\cite[Lemma 2.1 and equations (2.11)--(2.12)]{mapsViaLog}
gives $V_iP_V=P_VA_{i+1}$, $\det P_V=-1$,
$P_Vd=\delta$, $P_Vv_1=-u_0$, and
$P_V(-v_2)=\delta-u_1$.  This checks the step of
\cref{push:src:E.1.1} and \cref{push:src:E.1.2} that we use, without redoing its analytic
classification.  With $m=3k,n=4l$, the actual
finite translations are the vectors in \cref{push:src:E.1.1} and \cref{push:src:E.1.2} and
the circle transgression is
$b=nN_0+mN_1=\pm1$ for a sphere member.
The two target fillings have multiplicities $m,n$;
their genus-one meridian determinant is $b$.
Thus their target is $S^3$.  Its map to the base
is a weighted Hopf map of degree product $mn=12kl$.
The common lattices and their mod-two counterparts
are those in \ref{maps:src:B.1.1}, and all the arguments of
\cref{mapsI:src:A.2.1} and \ref{maps:src:B.1.2} apply without change.

\item For the two-parameter family use the homology basis
$(\widehat\psi,\widehat u,\widehat w,\widehat\delta)$
of \cref{push:src:E.5.1} and \cref{push:src:E.5.2} and \cref{push:src:E.5.3}.  Its actual four matrices are
\[
 H_* =\begin{pmatrix}1&0&0&0\\-2&0&-1&0\\0&1&0&0\\1&1&0&1\end{pmatrix},
 \qquad
 H_1 =\begin{pmatrix}1&0&0&0\\0&1&0&0\\-2&-1&1&0\\-2&-1&0&1\end{pmatrix},
\]
\[
 H_2 =\begin{pmatrix}1&0&0&0\\0&1&1&0\\0&0&1&0\\0&0&0&1\end{pmatrix},
 \qquad
 H_3 =\begin{pmatrix}1&0&0&0\\0&1&0&0\\0&-1&1&0\\1&0&0&1\end{pmatrix}.
\]
They are inverse transposes of the matrices in
\cite[Proposition 4.2]{mapsVia2}.  The direct
calculation gives
\[
 \Lambda^\Gamma=\mathbb Z\widehat\delta,
 \quad (\Lambda^*)^\Gamma=\mathbb Z\psi,
 \quad (\Lambda^*\otimes\mathbb F_2)^\Gamma
          =\mathbb F_2\psi,
 \quad L=\ker\psi.
 \tag{B.4.1}\label{maps:eq:B.4.1}
\]
The finite root action has order four and translation
$\ell/4$, where $\ell=(1,-1,-1,0)^t$; hence
$\psi(\ell)=1$.  The three primitive vanishing
lattices are
\[
 K_1=\mathbb Z(\widehat w+\widehat\delta),\qquad
 K_2=\mathbb Z\widehat u,\qquad
 K_3=\mathbb Z\langle\widehat w,\widehat\delta\rangle.
 \tag{B.4.2}\label{maps:eq:B.4.2}
\]
They all lie in $L$.  To check that the needed
class really extends, the two rank-one local
models in the source are the quotient of an
infinite chain of ruled elliptic surfaces, with
the adjacent sections identified by a translation.
Their nearby fundamental group surjects to the
central one with precisely $K_i$ killed; the
base meridian bounds the extending section.
This follows directly by filling the one
vanishing circle with its nodal thimble, and
retaining the other three circle generators.
It is also the degree-one part of the source's
local-specialization proposition, labelled
\texttt{mum}.  At the rank-two cusp the same
assertion is (A.1.5).  Thus $\psi$ extends
over every reduced filling, and its boundary
restriction is injective on $H^1$.  \cref{mapsI:src:A.2.1}
constructs and compares the replacements.

\item There is a single multiple target circle of
multiplicity four.  The other three circle
fillings are ordinary solid tori, with no
meridian shift because the chosen section
extends there.  The finite logarithmic
shift has primitive circle component $1/4$.
Equivalently its integral transgression is
$d_2(4\psi)=1$, as computed directly in
\cite[Proposition 5.4]{mapsVia2}.
Hence the genus-one target has meridian
determinant one and is $S^3$.

\item For Liu--Rollenske use their common homology
basis $(e_1,e_2,e_3,e_4)$.  The two finite
monodromies are
\[
 R_0=\begin{pmatrix}-1&-1&-1&0\\2&1&1&0\\0&0&1&0\\0&0&0&1\end{pmatrix},
 \qquad
 R_1=\begin{pmatrix}1&1&0&1\\-2&0&-1&-1\\0&-1&1&-1\\0&0&0&1\end{pmatrix}.
\]
Their source uses $R_\infty=(R_1R_0)^{-1}$.
They have respective orders four and six,
but the fibre over $1$ is reduced: finite
monodromy does not mean it is a smooth
multiple fibre.  The common invariants are
\[
 \Lambda^\Gamma=\mathbb Z(-e_1+2e_3),
 \quad(\Lambda^*)^\Gamma=\mathbb Z e_4^*,
 \quad(\Lambda^*\otimes\mathbb F_2)^\Gamma
            =\mathbb F_2 e_4^*.
 \tag{B.4.3}\label{maps:eq:B.4.3}
\]
Here $\delta=-e_1+2e_3$ and $\xi=e_4^*$,
so again the fixed vector belongs to
$L=\ker\xi$.
The local product identification of
\cite[Lemma 3.10 and Proposition 3.11]{mapsLR},
$Y_{\Delta_1}\cong S_{\Delta_1}\times T^2$,
preserves the two circle classes $f_3,f_4$,
with $\xi=f_4^*$.  Consequently $\xi$
extends over this reduced type-II fibre
and is surjected onto by the nearby
fundamental group.  At infinity the
primitive vanishing lattice is
$\mathbb Z\langle h_1,h_3\rangle$ and
$\xi=h_4^*$, so (A.1.5) gives its
extension there.  The only logarithmic
shift is $a=e_4/4$ over $0$.
Its primitive transgression is
$d_2(4\xi)=1$ by \cite[Proposition 4.8]{mapsLR}.
Thus the same one-exceptional-circle
target is $S^3$.  These are applications
of the exact local product and integral
monodromy statements in
\cite[Propositions 3.12--3.13 and Corollary 3.16]{mapsLR},
with the comparison of classes written
explicitly; rational invariant cycles
alone would not suffice.

\item For both one-multiple-fibre families,
Lemma \ref{maps:src:B.1.1} applies verbatim with
$(\gamma,d)$ replaced by $(\psi,\widehat\delta)$
or $(\xi,\delta)$.  Equations \eqref{maps:eq:B.4.1}
and \eqref{maps:eq:B.4.3} are all of its hypotheses.
The chosen regular inverse image of
the quotient map is a single $T_L$,
because every replacement is over its
own degeneration disc.  Therefore
Theorem \ref{maps:src:B.1.2} gives $[e]=6\nu'$ in
each family.  The target projection
is the weighted Hopf map of weights
$(4,1)$, whose power-map factor has
degree four; it sends $6\nu'$ to zero.
For $X_{3k,4l}$ the degree is $12kl$,
which also sends it to zero.

\item Finally each named family has a
primitive hexagonal rank-two cusp.
The proof of \ref{maps:src:B.3.1} is local in its
normal-crossing chart, independent
of the other fibres and of the
bounded period matrix.  It gives
$H=\eta_4\eta_5$ with the same target
coordinate convention.  The two
long-period classes extend over
that cusp, so \ref{maps:src:B.3.2} proves the
two angle-forgetting assertions.
All kernel and mod-two calculations
above are executed in
\texttt{code/maps\_exact.py}.
\end{enumerate}
\end{proof}

\subsection{A cell presentation from the three fillings}

We give a topological cell construction, retaining the actual
boundary maps before making cellular approximations.  The
smaller cellular complex used for the integer calculation is
homotopy equivalent to this construction; its cell counts are
not asserted to be the counts of a smooth handle decomposition.
Throughout this subsection $p=-1$, so the finite twists are
$v_1=g+2u-4w$ and $v_2=-g-3u+3w$ and the cusp gluing has
zero translation drift.

For a matrix $A$ on $F=\mathbb R^4/\Lambda$, our mapping-torus
convention is
\[
 M_A=(F\times[0,1])/((x,1)\sim(A^{-1}x,0)).
 \tag{B.5.1}\label{maps:eq:B.5.1}
\]
The positive interval loop $\rho$ satisfies
$\rho t_\lambda\rho^{-1}=t_{A\lambda}$.  This fixes the
inverse in \eqref{maps:eq:B.5.1} and agrees with the clockwise meridians
used in A.1.

\begin{theorem}\label{maps:src:B.5.1}
\StatementRecord{maps:src:B.5.1}{proved}{Historical events: Claude Fable 5.2 (Anthropic), effort high, 16 November 2026: carried as background without fault found. Claude Opus 5.5 (Anthropic), effort max, sweep 1, scope s6, 18 November 2026: computation or formal certificate. GPT-6 Astra (OpenAI), ChatGPT Pro mode, effort Pro, S6\_09\_PAPERS, 17 November 2026: targeted proof check. Also checked in this writing. Pre-continuation writing events: GPT-6 Astra (OpenAI), ChatGPT Pro mode, effort Pro, S6\_10\_PAPERS, 18 November 2026: targeted proof check. Continuation metadata events: GPT-6 Astra (OpenAI), ChatGPT Pro mode, effort Pro, S6\_10\_PAPERS continuation, 19 November 2026: used and rechecked in this run; status-vocabulary reconciliation; metadata reconciliation only; no new mathematical check. \emph{Pre-continuation recorded proof scope.} GPT-6 Astra (OpenAI, ChatGPT Pro mode), effort Pro, 2026-11-17; Claude Fable 5.2 (Anthropic), effort high, 16 November 2026; also checked in this writing by GPT-6 Astra (OpenAI), ChatGPT Pro mode, effort Pro, 18 November 2026}{the hypotheses, notation and definitions stated here; \cref{maps:eq:B.5.3}; \cite{mapsAC26}; \cref{maps:eq:B.5.1}; \cref{maps:eq:B.5.6}; \cite[Theorem 6.1, p.~77]{mapsWaldhausen68}; \cref{maps:eq:B.5.4}; \cref{maps:eq:B.5.8}; \cite[Propositions 4G.1--4G.2]{mapsHatcher02}}


The underlying topological manifold of $X$ is obtained from
the flat $T^4$-bundle over a ribbon bouquet of two circles,
with monodromies $A_1,A_2$, and the three spaces $S_1,S_2,W$
by the three double mapping cylinders specified below.
Their maps admit finite polyhedral cell structures.  They
therefore give a CW structure adapted to the pair of pants
and the three special fibres, with all attaching maps
specified by integral affine formulae and primitive circle
quotients.  Replacing the maps by cellular maps on the smaller
torus cell structures gives a finite CW homotopy model with
cell counts
\[
 (c_0,c_1,c_2,c_3,c_4,c_5,c_6)
       =(5,20,45,56,42,17,3).
 \tag{B.5.2}\label{maps:eq:B.5.2}
\]

\end{theorem}


\begin{proof}
\begin{enumerate}[label=\arabic*.,leftmargin=*]
\item Give a bouquet with oriented loops $\rho_1,\rho_2$ the ribbon
order in which the two ends of each loop are consecutive.
Its oriented thickening is a pair of pants.  Collapse the
widths of its two bands to the edges and the remaining
vertex disc to the vertex.  The induced boundary maps are
piecewise linear paths with words
$\rho_1,\rho_2,(\rho_1\rho_2)^{-1}$; intervals between
successive letters are mapped to the vertex.  The thickening
is the mapping cylinder of this boundary map.  This assertion
can be checked on each band, where the cylinder consists of
the two half-bands, and on the vertex disc, where the cones
on the intervening boundary intervals are its sectors.
Lift this construction to the flat $F$-bundle.  Write $U_*$
for the resulting bundle over the bouquet.  Its boundary
spaces are $M_{A_1},M_{A_2},M_{M_0}$.  Their maps
$b_j:M_{A_j}\to U_*$ and $b_0:M_{M_0}\to U_*$ are the
lifted ribbon boundary maps.  They are the identity on a
marked fibre and send their interval loops respectively to
$\rho_1,\rho_2,\rho_2^{-1}\rho_1^{-1}$.
The equality $M_0=(A_1A_2)^{-1}$ verifies the last map
including its fibre monodromy.

\item At a finite fibre put
\[
 h_j(x)=A_jx+v_j/m_j,\qquad S_j=F/\langle h_j\rangle,
 \qquad
 r_j([x,\theta])=[x+\theta v_j/m_j].
 \tag{B.5.3}\label{maps:eq:B.5.3}
\]
The last formula is a well-defined map $M_{A_j}\to S_j$:
at $\theta=1$ its value agrees with $[A_j^{-1}x]$, because
$h_j(A_j^{-1}x)=x+v_j/m_j$.  The finite neighbourhood is
the associated disc bundle
$(F\times D^2)/\langle(x,s)\mapsto(h_jx,e^{-2\pi i/m_j}s)\rangle$.
Its boundary point corresponding to $[x,\theta]$ is
\[
 [x+\theta v_j/m_j,\ e^{-2\pi i\theta/m_j}].
\]
Radial collapse is exactly \eqref{maps:eq:B.5.3}; its mapping cylinder is
the disc bundle.  In particular this is a statement about
the underlying smooth neighbourhood and does not replace
the scheme-theoretic fibre $m_jS_j$ by a reduced fibre in
the holomorphic map.  The logarithmic gluing translation is
$\log(s)\Pi v_j/(2\pi i)$.  Its clockwise drift is precisely
$v_j/m_j$.  Subtracting that displayed drift leaves a
periodic translation with a lift to $\mathbb R^4$; linear
interpolation of this lift straightens it on a collar.
Thus the actual gluing is identified with \eqref{maps:eq:B.5.3}, without
discarding a loop of translations of nonzero lattice drift.
This also corrects the first paragraph on page 42 of
\cite{mapsAC26}: multiplicity does not prevent this
neighbourhood from being a smooth normal disc bundle.
We make no assertion of a holomorphic tubular projection.
The meridian relation proved in that paragraph remains
valid; its meridian is the normal circle winding $m_j$
times around the base, rather than a once-around base
loop.

\item Here is the complete cusp boundary map.  Put
$a=(a_1,a_2)\in\mathbb R^2/\mathbb Z^2$ for the $(g,u)$
coordinates and $b\in T_f=\mathbb R^2/\mathbb Z^2$ for
the $(w,d)$ coordinates.  Let $\mathcal H$ be the hexagonal
honeycomb dual to the $A_2$ triangulation.  For $B_0a$ in
the interior of the hexagon centred at $v\in\mathbb Z^2$,
write $v(a)=v$.  The central fibre has the explicit quotient
description
\[
 W_{\mathcal H}=(T_B\times T_f)/\sim,
 \qquad T_B=\mathbb R^2/\mathbb Z^2,
 \tag{B.5.4}\label{maps:eq:B.5.4}
\]
where over a honeycomb edge separating centres $v,v'$ one
collapses the primitive circle on $v'-v$, and over each
honeycomb vertex one collapses all of $T_f$.  No circle is
collapsed over a hexagon interior.  The three oriented edge
directions can be chosen as
\[
 n_1=w,\qquad n_2=d-w,\qquad n_3=-d.
 \tag{B.5.5}\label{maps:eq:B.5.5}
\]
The quotient is homeomorphic to the actual $W$: the positive
part of each toric component is its hexagon, the phase
stabilizers are exactly these primitive circles, and the
period action identifies hexagons by $B_0\mathbb Z^2$.
The constant phase cocycle is removed by the linear phase
gauge coming from $\operatorname{Re}C(0)a$.

\item With the interval convention \eqref{maps:eq:B.5.1}, the cusp attachment is
\[
 r_0([a,b,\theta])
   =[a,b-\theta(B_0a-v(a))]\ \in W_{\mathcal H}.
 \tag{B.5.6}\label{maps:eq:B.5.6}
\]
This formula includes its seam data.  On crossing an edge
the two expressions differ by $\theta(v-v')$, which lies
in the circle collapsed on that edge.  At a vertex every
phase is collapsed.  Replacing $a$ by $a+\lambda$ changes
$v(a)$ by $B_0\lambda$, so the expression is periodic.
At $\theta=1$, the integer vector $v(a)$ disappears modulo
$\mathbb Z^2$, leaving $[a,b-B_0a]$; this is the endpoint
identification by $M_0^{-1}$.  Hence \eqref{maps:eq:B.5.6} is continuous
on the entire mapping torus, not only on its open strata.

\item We justify that the cusp neighbourhood is its actual
mapping cylinder.  Straighten $C(q)$ to $C_0=C(0)$ by
the map on logarithmic coordinates
\[
 \xi\longmapsto
 \xi-(C(q)-C_0)(\operatorname{Im}Z_q)^{-1}
                       \operatorname{Im}\xi.
 \tag{B.5.7}\label{maps:eq:B.5.7}
\]
Writing $\xi=Z_qa+b$ proves that this sends $Z_qa+b$ to
$Z_q^0a+b$, fixes the real lattice, and conjugates the
period actions.  In a bounded-position toric chart the
correction multiplying $x$ tends to one as $q\to0$;
the same holds for its inverse.  Thus it extends by the
identity on $W$ to a homeomorphism of the closed cusp tube.

\item Let $A$ be the nonnegative real part of the frozen toric
cover over $|q|\leq\eta$, and $A_0$ its central part.
The local orthant descriptions make $A$ a three-manifold
with boundary; the corners may be rounded.  Its interior
has coordinates $(y,\rho)\in\mathbb R^2\times(0,\eta)$,
$y=\log|x|/\log\rho$.  The positive period action is
$y\mapsto y+B_\rho^+\lambda$, where
$B_\rho^+=B_0-2\pi\operatorname{Im}C_0/\log\rho$.
For small $\eta$ these matrices are invertible.  Thus the
interior of the positive quotient $Q^+$ is explicitly
$T^2\times(0,\eta)$, using $a=(B_\rho^+)^{-1}y$.
The bounded fundamental parallelogram meets only finitely
many of the toric bounded-position charts, including at
$\rho=0$, so $Q^+$ is compact.  Its two boundary components
are tori.  The inner one is the honeycomb quotient, and
the outer one is the positive regular-fibre torus.

\item Both boundary inclusions induce the identity on the period
deck group $\mathbb Z^2$.  Indeed $A_0$ is the honeycomb
plane and hence contractible; $A$ is contractible because
its interior is contractible and pushing a boundary collar
inward gives a homotopy equivalence with that interior.
The quotient is irreducible.  An embedded locally flat
two-sphere lies in its interior $T^2\times(0,\eta)$.
Lift it to $\mathbb R^2\times\mathbb R$.  Its bounded
ball is disjoint from all its nonzero lattice translates:
disjoint boundary spheres would otherwise give nested
balls, impossible for translates of equal finite volume.
The ball therefore descends.  This also proves
irreducibility of the product target $T^2\times I$.
The boundary tori are incompressible because their
fundamental groups inject.

\item There is a map of pairs $Q^+\to T^2\times I$ inducing
the identity on $\pi_1$ and restricting to the prescribed
homeomorphisms on the two boundaries.  For its first
coordinate, extend the classifying map to $T^2=K(\mathbb
Z^2,1)$ and use the homotopy extension property on the
boundary; use a collar height for the second coordinate.
Apply \cite[Theorem 6.1, p.~77]{mapsWaldhausen68} to this
map.  The source is sufficiently large, since a boundary
torus pushed into its interior is incompressible.  The
product target has the same property, and all the other
hypotheses were just verified.  The homotopy is fixed on
the boundary.  Its alternative of mapping the connected
source into the boundary is impossible, since the two
source boundaries map to distinct target components.
It therefore ends at a covering map.  The induced
fundamental-group isomorphism makes that covering
one-sheeted.  We obtain a product homeomorphism relative
to both prescribed boundary markings.

\item Lift this product to $A$ and spread it by compact phases.
The full period action on a phase $(\phi,\vartheta)$ is
$(\phi,\vartheta)\mapsto
(u_\lambda\phi\vartheta^{B_0\lambda},\vartheta)$;
it does not depend on the positive radial point.  The
lifted product is therefore equivariant for the full
action.  Phase stabilizers occur only on $A_0$ and are
exactly those of \eqref{maps:eq:B.5.4}.  Consequently the resulting
quotient is the mapping cylinder of the boundary phase
map.  On the component with centre $v$, the $q$-phase
acts on $T_f$ by $-v$ times that phase, since $(v,1)$ is
the stabilizer ray.  Substituting the clockwise interval
phase and the marked period $-\theta B_0a$ gives precisely
\eqref{maps:eq:B.5.6}.

\item Gluing the three mapping cylinders of $r_j$ to the
mapping cylinder of the ribbon boundary maps gives
\[
 U_*\ \cup\ \coprod_{j=0}^2(M_j\times[0,1])
       \ \cup\ (W\sqcup S_1\sqcup S_2),
 \tag{B.5.8}\label{maps:eq:B.5.8}
\]
with $(x,0)\sim b_j(x)$ and $(x,1)\sim r_j(x)$.
The preceding constructions identify \eqref{maps:eq:B.5.8}
homeomorphically with $X$.  The comparison need not
preserve the radial value of $f$ inside the cusp, but
it preserves the special strata and all three outer
collars; this is the sense of adaptation used here.

\item For a finite cell construction use rational hexagons,
unit cubes for tori, and the affine quotient \eqref{maps:eq:B.5.3}.
For \eqref{maps:eq:B.5.6} use toric phase coordinates
$b'=b-\theta B_0a$.  In these coordinates its expression
on a hexagon is $[a,b'+\theta v]$, and its face
identifications are affine:
$a\mapsto a+\lambda$ changes $b'$ by
$-\theta B_0\lambda$.  There are finitely many faces
in the fundamental polyhedra.  Cut their images at the
integer phase hyperplanes and subdivide; this makes
every face map a finite piecewise affine map.  Primitive
edge-circle quotients are the integral linear maps
$b\mapsto\det(n_k,b)\bmod\mathbb Z$.
Triangulate each occurrence of a torus separately and
interpolate through a common subdivision along each
mapping cylinder.  This does not require a finite
triangulation invariant under the infinite monodromy
group.  Products of simplices with intervals are
triangulated by their ordered vertex prisms.  These
rules supply finite characteristic maps for \eqref{maps:eq:B.5.8},
and pulling the cells back gives the asserted CW
structure of $X$.

\item For the smaller homotopy model put the usual product
cell structure on each torus and apply cellular
approximation to the displayed maps.  This changes
the double mapping cylinder by homotopy equivalence;
equivalently use \cite[Propositions 4G.1--4G.2]{mapsHatcher02}.
The two finite surfaces are mapping tori of $T_L$,
because $m_j\gamma$ is a primitive circle coordinate
on $S_j$; their linear monodromies are
$B_j=A_j|_L$ (an inverse convention gives the same
counts).  Their cell counts are $(1,4,6,4,1)$.
The counts for $U_*$ are $(1,6,14,16,9,2)$, and
each $M_j$ has counts $(1,5,10,10,5,1)$.
In \eqref{maps:eq:B.5.4} the quotient honeycomb has two vertices,
three edges, and one hexagon.  The surviving phase
spaces are respectively a point, a circle, and a
two-torus, giving counts $(2,3,4,2,1)$ for $W$.
Adding these cells and one prism cell for each
boundary cell gives \eqref{maps:eq:B.5.2}.
\end{enumerate}
\end{proof}

\begin{theorem}\label{maps:src:B.5.2}
\StatementRecord{maps:src:B.5.2}{proved}{Historical events: Claude Fable 5.2 (Anthropic), effort high, 16 November 2026: carried as background without fault found. GPT-6 Astra (OpenAI), ChatGPT Pro mode, effort Pro, S6\_09\_PAPERS, 17 November 2026: targeted proof check. Also checked in this writing. Pre-continuation writing events: GPT-6 Astra (OpenAI), ChatGPT Pro mode, effort Pro, S6\_10\_PAPERS, 18 November 2026: targeted proof check. Continuation metadata events: GPT-6 Astra (OpenAI), ChatGPT Pro mode, effort Pro, S6\_10\_PAPERS continuation, 19 November 2026: used and rechecked in this run; status-vocabulary reconciliation; metadata reconciliation only; no new mathematical check. \emph{Pre-continuation recorded proof scope.} GPT-6 Astra (OpenAI, ChatGPT Pro mode), effort Pro, 2026-11-17; Claude Fable 5.2 (Anthropic), effort high, 16 November 2026; also checked in this writing by GPT-6 Astra (OpenAI), ChatGPT Pro mode, effort Pro, 18 November 2026}{the hypotheses, notation and definitions stated here; \cref{maps:src:B.5.1}; \cref{maps:eq:B.5.3}; \cref{maps:eq:B.5.6}; \cref{maps:eq:B.5.11}; \cref{maps:eq:B.5.4}; \cref{maps:eq:B.5.15}; \cref{maps:eq:B.5.12}; \cref{maps:eq:B.5.18}; \cref{maps:eq:B.5.2}; \cref{maps:eq:B.5.9}; \cite[Theorems 4.32 and 4.5]{mapsHatcher02}}


The attaching data of Theorem \ref{maps:src:B.5.1} give
$\pi_1(X)=1$ and integral homology
$H_k(X;\mathbb Z)=\mathbb Z$ for $k=0,6$ and zero
otherwise.  They give an explicit recognition of $X$
as a homotopy sphere, and exhibit the normally framed
three-torus used in B.1.

\end{theorem}


\begin{proof}
\begin{enumerate}[label=\arabic*.,leftmargin=*]
\item The graph piece has presentation
\[
 \pi_1(U_*)=
 \langle\Lambda,\rho_1,\rho_2\mid
 [\Lambda,\Lambda]=1,\quad
 \rho_jt_\lambda\rho_j^{-1}=t_{A_j\lambda}\rangle.
\]
Formula \eqref{maps:eq:B.5.3} imposes $\rho_j^{m_j}=t_{v_j}$.
Formula \eqref{maps:eq:B.5.6} kills $w,d$ and the cusp meridian
$a_0=\rho_2^{-1}\rho_1^{-1}$.  Since $A_1w=u-w$,
the normal closure also kills $u$.  Writing $c$ for
the image of $g$, the remaining presentation is
\[
 \langle c,x,y\mid c\text{ central},\quad
 xy=1,\quad x^3=c,\quad y^4=c^{-1}\rangle.
 \tag{B.5.9}\label{maps:eq:B.5.9}
\]
Substituting $y=x^{-1}$ gives $x^3=x^4$, so $x=c=y=1$.
This calculates the nonabelian fundamental group,
not only its abelianization.

\item We record the integral homology maps in primitive
bases.  Write $\gamma,u^*,w^*,d^*$ for the cohomology
basis dual to $g,u,w,d$, and put
\[
 \eta_1=2u^*+w^*+3d^*,\quad
 \eta_2=u^*+w^*+2d^*,\quad
 q_2=u^*w^*+6\gamma d^*.
 \tag{B.5.10}\label{maps:eq:B.5.10}
\]
Products here are exterior products.  First attach
the two finite pieces over a disc containing $0,1$;
call the result $X^\circ$.  Its homotopy type is the
double mapping cylinder $S_1\leftarrow F\to S_2$.
The degree-one
map $(\pi_{1*},-\pi_{2*})$ is
\[
 \alpha_1=
 \begin{pmatrix}
 3&0&0&0\\0&2&1&3\\4&0&0&0\\0&-1&-1&-2
 \end{pmatrix}.
 \tag{B.5.11}\label{maps:eq:B.5.11}
\]
To verify its primitive target bases, abelianize
$\pi_1(S_j)=\langle\Lambda,\rho_j\mid
\rho_j\lambda\rho_j^{-1}=A_j\lambda,
\rho_j^{m_j}=t_{v_j}\rangle$.
For $S_1$ the images are
$g=3\rho_1,u=2w,d=3w$; for $S_2$ they are
$g=-4\rho_2,u=w,d=2w$.  Thus
$(\rho_1,w)$ and $(\rho_2,w)$ are free integral bases.
This proves \eqref{maps:eq:B.5.11} without a rational invariant-lattice
identification.  Its Smith form is $(1,1,1,0)$,
its kernel is generated by
$\nu_0=-u-w+d$, and its primitive left kernel is
$(4,0,-3,0)$.  Consequently
$H_1(X^\circ)=\mathbb Z\zeta$, normalized by
\[
 g\mapsto12\zeta,\quad
 \rho_1\mapsto4\zeta,\quad
 \rho_2\mapsto-3\zeta,\quad
 a_0\mapsto-\zeta.
 \tag{B.5.12}\label{maps:eq:B.5.12}
\]

\item For completeness we derive the finite degree-two
pullback lattices.  The mapping-torus matrices on
$L$ are
\[
 B_1=\begin{pmatrix}0&1&0\\-1&-1&0\\1&0&1\end{pmatrix},
 \qquad
 B_2=\begin{pmatrix}0&-1&0\\1&0&0\\0&1&1\end{pmatrix}.
\]
The exterior-power mapping-cone complexes show that
$H_*(S_j)$ is free, with ranks $(1,2,2,2,1)$:
the nonzero Smith entries of
$\bigwedge^r B_j-I$ are all one.  In particular the
intersection lattice in degree two is unimodular.
Direct exterior multiplication gives the primitive
invariant lattices
$I_1=\langle\gamma\eta_1,q_2\rangle$ and
$I_2=\langle\gamma\eta_2,q_2\rangle$.
Their intersection determinants have absolute values
$9$ and $4$.  The pullback under a covering of degree
$m_j$ multiplies the intersection pairing by $m_j$.
Its image therefore has determinant $m_j^2$ and index
one in $I_1$, respectively index two in $I_2$.
The embedded torus generated by $v_2,d$ is invariant
under $h_2$ and covers its quotient with degree four.
The class $a\gamma\eta_2+bq_2$ evaluates on it as
$-2a-6b$.  A pullback class must therefore satisfy
$a+b\equiv0\pmod2$.  This is already an index-two
condition and hence determines the image exactly.
Thus
\[
 \pi_1^*H^2(S_1)=\langle\gamma\eta_1,q_2\rangle,
 \qquad
 \pi_2^*H^2(S_2)=\langle2\gamma\eta_2,
                                  \gamma\eta_2+q_2\rangle.
 \tag{B.5.13}\label{maps:eq:B.5.13}
\]

\item Use bases with pullbacks
$(\gamma\eta_1,q_2-2\gamma\eta_1)$ and
$(2\gamma\eta_2,q_2-3\gamma\eta_2)$, and their
Kronecker-dual homology bases.  Evaluating on
$(gu,gw,gd,uw,ud,wd)$ gives
\[
 \alpha_2=
 \begin{pmatrix}
 2&1&3&0&0&0\\-4&-2&0&1&0&0\\
 -2&-2&-4&0&0&0\\3&3&0&-1&0&0
 \end{pmatrix}.
 \tag{B.5.14}\label{maps:eq:B.5.14}
\]
Its Smith form is again $(1,1,1,0)$ and its primitive
left kernel is $\Phi=(4,2,3,2)$.  Thus its cokernel is
$\mathbb Z e$, where the cohomology class $\mathcal A$
given by $\Phi$ satisfies
$\langle\mathcal A,e\rangle=1$ and restricts on $F$
to $2q_2$.  The exact sequence of this double mapping
cylinder gives
\[
 0\longrightarrow\mathbb Z e\longrightarrow
 H_2(X^\circ)\longrightarrow\mathbb Z\nu_0
 \longrightarrow0.
 \tag{B.5.15}\label{maps:eq:B.5.15}
\]

\item The central quotient \eqref{maps:eq:B.5.4} has cellular boundary
$\partial_1(e_k)=Q-P$ and all higher cellular
boundaries zero.  Indeed the hexagon attaching word
can be chosen as
$e_1e_2^{-1}e_3e_1^{-1}e_2e_3^{-1}$; every edge occurs
with both signs.  Crossing with a surviving phase
circle preserves this cancellation, and the torus
cell boundary is the commutator.  Hence
\[
 H_*(W)=(\mathbb Z,\mathbb Z^2,\mathbb Z^4,
                       \mathbb Z^2,\mathbb Z).
\]
At $\theta=0$ the cusp map is the quotient
$c(a,b)=[a,b]$.  On $H_1$ it is projection to
$\langle g,u\rangle$.  With degree-two target basis
$(T_B,\overline C_1,\overline C_2,\overline C_3)$,
its matrix is
\[
 c_2=\begin{pmatrix}
 1&0&0&0&0&0\\0&0&-1&0&0&0\\
 0&-1&-1&1&1&0\\0&0&0&1&0&0
 \end{pmatrix}.
 \tag{B.5.16}\label{maps:eq:B.5.16}
\]
For a mixed torus $\beta v$ this follows from the
degree formula
$\sum_k\det(n_k,B_0\beta)\det(n_k,v)[\overline C_k]$.
The first factor is its oriented honeycomb edge
coefficient, and the second is the degree of the
primitive circle quotient.  The map sends $gu$ to
$T_B$ and $wd$ to zero.  In degree three it sends
$guw,gud$ to a basis and kills $gwd,uwd$; in degree
four it has degree one on the open hexagon times
$T_f$.  All these maps are integrally surjective.
Their kernels are exactly
$\operatorname{im}(\bigwedge^r M_0-I)$: inclusion
follows by transport around the boundary, and the
displayed kernels and monodromy images have the
same primitive bases, as checked in
\texttt{code/maps\_cw.py}.

\item Let $K_i$ be the kernel of $H_i(M_0)\to H_i(W)$.
The mapping-torus sequence and the preceding
surjections give
\[
 K_1=\mathbb Z a_0,\qquad
 K_2=\mathbb Z\xi_w\oplus\mathbb Z\xi_d,
 \qquad \xi_v=[S^1_v\times S^1_{a_0}].
 \tag{B.5.17}\label{maps:eq:B.5.17}
\]
These are literal tori because $M_0$ fixes $w,d$.
Under the mapping-torus projection they go to $w,d$,
so their span is primitive.  In \eqref{maps:eq:B.5.6}, over a
hexagon centre the $a_0$ coordinate is collapsed;
thus their images in $H_2(W)$ vanish.  The ranks
from the mapping-torus sequence show that this is
the whole kernel.

\item Orient the connecting map in \eqref{maps:eq:B.5.15} so that
\[
 \partial j_*\xi_v=(I-A_2)v
                   =(I-A_1^{-1})v\quad(v\in\langle w,d\rangle).
\]
One verifies this on the two rectangles obtained
by cutting the swept torus along the finite band:
their unmatched boundary circles are $v$ and
$-A_2v$.  Since $(I-A_2)w=-\nu_0$ and
$(I-A_2)d=0$, the class $z=j_*\xi_w$ together with
$e$ is a basis of $H_2(X^\circ)$.
To determine $j_*\xi_d$, use the circle translation
in direction $d$ on $X^\circ$.  For this action the
integral sweep operator on cohomology restricts on
$F$ to contraction by $d$.  The class
$\mathcal A|_F=2q_2$ satisfies
\[
 \iota_d(2q_2)=-12\gamma.
\]
Restriction $H^1(X^\circ)\to H^1(F)$ is injective
by \eqref{maps:eq:B.5.12}, with image $\mathbb Z\,12\gamma$.
Thus $\iota_d\mathcal A=-\Theta$, where
$\Theta(\zeta)=1$.  Evaluating the swept torus and
using $\Theta(a_0)=-1=p$ gives
\[
 j_*\xi_d=-p e=e,
 \qquad
 [j_*|_{K_2}]_{(\xi_d,\xi_w)}^{(e,z)}
       =\begin{pmatrix}1&0\\0&1\end{pmatrix}.
 \tag{B.5.18}\label{maps:eq:B.5.18}
\]
This computes the attaching coefficient that cannot
be recovered from ranks alone.

\item In degree one the full cusp attachment map in bases
$(g,u,a_0)$ and $(\bar g,\bar u,\zeta)$ is
\[
 D_1=\begin{pmatrix}1&0&0\\0&1&0\\-12&0&1\end{pmatrix}.
\]
In degree two, split the surjection onto $H_2(W)$;
the remaining block is exactly \eqref{maps:eq:B.5.18}.  Thus the
full map is integrally equivalent to $I_4\oplus I_2$.
Both are unimodular.  Mayer--Vietoris proves
$H_1(X)=H_2(X)=0$, with no torsion.  The cell count
\eqref{maps:eq:B.5.2} gives $\chi(X)=2$.  Since $X$ is a closed
oriented six-manifold, duality and universal
coefficients give $H_4=H_5=0$ and make $H_3$ free:
$\operatorname{Ext}(H_3,\mathbb Z)$ injects into
$H^4(X)=H_2(X)=0$.  The Euler characteristic then
gives $\operatorname{rank}H_3=0$.  This proves all
the asserted integral groups.  The program checks
every matrix, Smith form, primitive coefficient
and cell count just used, using exact arithmetic.

\item Starting from \eqref{maps:eq:B.5.9}, successive applications of
the Hurewicz theorem make $\pi_i(X)=0$ for
$1\leq i\leq5$.  The same theorem produces a map
$S^6\to X$ carrying the fundamental class to the
chosen generator.  The mapping-cylinder pair of
this map is simply connected and has zero relative
homology.  If it had a first nonzero relative
homotopy group, the relative Hurewicz theorem
would make its relative homology nonzero.  All
relative homotopy groups therefore vanish, and
the CW Whitehead theorem makes the map a homotopy
equivalence.  These
are \cite[Theorems 4.32 and 4.5]{mapsHatcher02}
applied to the explicitly computed complex.

\item Finally keep a small regular disc in the pair of
pants uncollapsed while forming its ribbon spine.
The same construction relative to that disc retains
the block $F\times D^2$ with its flat vertical
frame and base frame.  The three-torus
$\{\gamma=\theta_0\}\subset F$ has normal frame
given by the two base directions and the
$\gamma$ direction.  This is exactly the regular
fibre and normal framing of B.1.  The CW
presentation therefore retains the geometric
framed fibre; no smooth framing of hypothetical
cancelled handles has been inferred from the
homology calculation.
\end{enumerate}
\end{proof}


\section{The local Kummer cover and its global obstruction}

We distinguish a finite holomorphic map with smooth source from a finite
map with normal, possibly singular, source. A branched cover of compact
complex manifolds has the former meaning in this section. Write
$V_{\rm tor}=\Z\langle\widehat w,\widehat\delta\rangle$.

\begin{theorem}[checked local Kummer theorem]\label{maps:src:C.1.1}
\StatementRecord{maps:src:C.1.1}{proved}{Historical events: Claude Fable 5.2 (Anthropic), effort high, 16 November 2026: carried as background without fault found. GPT-6 Astra (OpenAI), ChatGPT Pro mode, effort Pro, S6\_09\_PAPERS, 17 November 2026: targeted proof check. Also checked in this writing. Pre-continuation writing events: GPT-6 Astra (OpenAI), ChatGPT Pro mode, effort Pro, S6\_10\_PAPERS, 18 November 2026: targeted proof check. Continuation metadata events: GPT-6 Astra (OpenAI), ChatGPT Pro mode, effort Pro, S6\_10\_PAPERS continuation, 19 November 2026: used and rechecked in this run; status-vocabulary reconciliation; metadata reconciliation only; no new mathematical check. \emph{Pre-continuation recorded proof scope.} GPT-6 Astra (OpenAI, ChatGPT Pro mode), effort Pro, 2026-11-17; Claude Fable 5.2 (Anthropic), effort high, 16 November 2026; also checked in this writing by GPT-6 Astra (OpenAI), ChatGPT Pro mode, effort Pro, 18 November 2026}{the hypotheses, notation and definitions stated here; the explicit constructions and calculations in the proof}


Let $N\geq1$. On the cusp toric cover with period correction $C(q)$,
put $q=v^N$ and $C_N(v)=C(v^N)/N$. The map
\[
 P_N(\xi_1,\xi_2,v)=(\xi_1^N,\xi_2^N,v^N)
\]
descends to a finite flat holomorphic map
$\kappa_N:N_0^{\langle N\rangle}\to N_0$ of degree $N^3$ between
smooth threefolds. The source central fibre is reduced and irreducible,
with normalization $\mathrm{dP}_6$. Its ramification divisor is
$(N-1)W_N$, and $\kappa_N^*W=NW_N$.

The base change $N_0[N]=N_0\times_{\Delta_q}\Delta_v$ is already normal.
The induced finite map $Q_N:N_0^{\langle N\rangle}\to N_0[N]$ has
degree $N^2$. For $N>1$ it is quasi-etale, and its nonflat and non-etale
loci are precisely the inverse images of the double and triple strata.
It is not a normalization morphism. On a regular fibre its image lattice is
\[
 L_N=\Z\langle\widehat\gamma,\widehat u,
                   N\widehat w,N\widehat\delta\rangle.
 \tag{C.1.1}\label{maps:eq:C.1.1}
\]

\end{theorem}


\begin{proof}
\begin{enumerate}[label=\arabic*.,leftmargin=*]
\item On the target torus the period action is
$\Psi_\lambda(x,q)=(e^{2\pi iC(q)\lambda}q^{B_0\lambda}x,q)$,
where $B_0=\left(\begin{smallmatrix}0&1\\-1&0\end{smallmatrix}\right)$.
The source uses the same fan and the action
\[
 \Psi^N_\lambda(\xi,v)=
 \big(e^{2\pi iC(v^N)\lambda/N}v^{B_0\lambda}\xi,v\big).
\]
The integral shear preserves the $A_2$ triangulation and the
holomorphic unit factor acts in every toric chart. Additivity in $\lambda$
proves the action law. Raising the two torus coordinates to the $N$th
power proves $P_N\Psi^N_\lambda=\Psi_\lambda P_N$ exactly.
There is no root-selection ambiguity: the logarithm being divided by $N$
is the specified holomorphic function $2\pi iC(v^N)\lambda$.

\item In a maximal smooth toric chart the map is
$(w_0,w_1,w_2)\mapsto(w_0^N,w_1^N,w_2^N)$. Its convergent power-series
ring is free over the target ring on the $N^3$ monomials whose three
exponents lie in $[0,N-1]$. Thus $P_N$ is finite flat. These charts cover
the fan, and their inverse images are the corresponding charts because
scalar multiplication preserves every cone. Freeness and proper
discontinuity of the source lattice action now follow directly from the
target action: a source fixed point would map to a target fixed point;
if two source compact sets meet under a translate, their target compact
images meet under the same translate. Only finitely many such translates
exist. This transfers the established cusp action property without a new
estimate on $C_N$.

\item Over an evenly covered target open set choose one lift to the toric
cover. Every source lattice orbit over it has a unique representative
whose image is in that lift. Hence the descended morphism is locally the
same finite flat map, of degree $N^3$. It is proper. Since $f_0$ is proper,
the source map to $\Delta_v$ is proper as well: the inverse image of a
compact $K\subset\Delta_v$ is closed in
$\kappa_N^{-1}(f_0^{-1}(\{v^N:v\in K\}))$, a compact set.
The free action on a smooth toric space gives a smooth source. The fan and
its period action on vertices have not changed, so its central fibre has
one component and the same hexagonal normalization. In this chart
$q=z_0z_1z_2$, $v=w_0w_1w_2$, and the Jacobian is
$N^3(w_0w_1w_2)^{N-1}$, proving the asserted divisors.

\item At a triple point the relative factor has the exact ring map
\[
 \frac{\C\{z_0,z_1,z_2,v\}}{(z_0z_1z_2-v^N)}
 \longrightarrow\C\{w_0,w_1,w_2\},\qquad
 z_i\longmapsto w_i^N,\quad v\longmapsto w_0w_1w_2.
 \tag{C.1.2}\label{maps:eq:C.1.2}
\]
The source ring on the left is the invariant ring for
$G=\{(\epsilon_0,\epsilon_1,\epsilon_2)\in\mu_N^3:
\epsilon_0\epsilon_1\epsilon_2=1\}$ on the right. Indeed an invariant
monomial has three exponents congruent modulo $N$; factor out their
common residue as a power of $w_0w_1w_2$. The same expansion for convergent
series proves the analytic assertion. A finite-group invariant ring of a
normal domain is normal: an element of its fraction field integral over
the invariants is integral over the original ring, belongs to that ring
by normality, and is invariant. This proves normality here and on the
face charts, and the degree is $|G|=N^2$.

\item The stabilizer is trivial when at most one $w_i$ vanishes, has order $N$
when exactly two vanish, and is $G$ at the origin. The origin fibre has
length
\[
 N^3-(N-1)^3=3N^2-3N+1>N^2\quad(N>1).
\]
At a double point the two-variable factor has generic rank $N$ and
closed fibre length $2N-1>N$. These counts follow by excluding monomials
divisible by the product of the variables from the box of exponents
$0,\ldots,N-1$. They prove failure of flatness exactly where claimed;
elsewhere the free finite action gives etaleness and flatness.

\item Finally, in additive coordinates multiplication by $N$ is the map
\[
 \C^2/(N\Z^2+Z(Nr)\Z^2)\longrightarrow
 \C^2/(\Z^2+Z(Nr)\Z^2),\qquad v=e^{2\pi ir}.
\]
It proves \eqref{maps:eq:C.1.1}, degree $N^2$, and the fibre deck group
$(\Z/N)^2$. With $S_N=\operatorname{diag}(1,1,N,N)$ and
$M_0=I+R$, $R^2=0$, direct multiplication gives
$S_N^{-1}M_0^NS_N=M_0$. This also checks the marking in the local
workbench theorem, rather than identifying covers from their degrees
alone.
\end{enumerate}
\end{proof}

\begin{lemma}[smoothness of a cover of the entire fan]\label{maps:src:C.1.2}
\StatementRecord{maps:src:C.1.2}{proved}{Historical events: Claude Fable 5.2 (Anthropic), effort high, 16 November 2026: carried as background without fault found. GPT-6 Astra (OpenAI), ChatGPT Pro mode, effort Pro, S6\_09\_PAPERS, 17 November 2026: targeted proof check. Also checked in this writing. Pre-continuation writing events: GPT-6 Astra (OpenAI), ChatGPT Pro mode, effort Pro, S6\_10\_PAPERS, 18 November 2026: targeted proof check. Continuation metadata events: GPT-6 Astra (OpenAI), ChatGPT Pro mode, effort Pro, S6\_10\_PAPERS continuation, 19 November 2026: used and rechecked in this run; status-vocabulary reconciliation; metadata reconciliation only; no new mathematical check. \emph{Pre-continuation recorded proof scope.} GPT-6 Astra (OpenAI, ChatGPT Pro mode), effort Pro, 2026-11-17; Claude Fable 5.2 (Anthropic), effort high, 16 November 2026; also checked in this writing by GPT-6 Astra (OpenAI), ChatGPT Pro mode, effort Pro, 18 November 2026}{the hypotheses, notation and definitions stated here; the explicit constructions and calculations in the proof}


Let $\Sigma$ be the complete periodic $A_2$ triangulation at height one
used for the cusp, in $N'=\Z^2\oplus\Z$. Suppose a finite connected
holomorphic cover of its toric neighbourhood is etale over
$(\C^*)^2\times\Delta^*$ and has smooth total space. The cocharacter
sublattice defining its restriction to this torus is $nN'$ for some
$n\geq1$. In particular its intersection with the two-dimensional
vanishing lattice is $n\Z^2$.

\end{lemma}


\begin{proof}
\begin{enumerate}[label=\arabic*.,leftmargin=*]
\item A connected finite cover of $(\C^*)^2\times\Delta^*$ is classified by
a finite-index sublattice $K\subset N'$, since this space has fundamental
group $\Z^3$. The unique normal extension across the toric boundary is
the toric map with the same real cones and lattice $K$. This can be seen
on each affine chart by taking the integral closure of its monomial ring
in the finite torus field extension; the resulting monomials have
exponents in the dual of $K$. The convergent analytic version has the
same local finite quotient charts. Uniqueness of normal extension glues
these charts and applies to the given finite cover.

\item For a target ray $r_v=(v,1)$ let $a_vr_v$ be its primitive generator in
$K$. A smooth source chart over a triangle with vertices $v_0,v_1,v_2$
requires $a_0r_0,a_1r_1,a_2r_2$ to be a basis of $K$. This is the usual
smoothness criterion, also obtained directly by identifying that chart's
semigroup with the free monoid on its three dual basis vectors.
In the adjacent triangle across the edge $(r_1,r_2)$ the other vertex is
$r_3=r_1+r_2-r_0$. Both target triples are integral bases. Equality of
their indices gives
\[
 [N':K]=a_0a_1a_2=a_3a_1a_2,
 \qquad a_3=a_0.
\]
Expressing $a_3r_3$ in the first source basis gives coefficients
$-1,a_0/a_1,a_0/a_2$. These are integers. Applying the same argument
across the other two edges of the first triangle gives all the reverse
divisibilities. Thus $a_0=a_1=a_2=n$, and their three vectors generate
$K=nN'$. This already proves the result; connectivity of the triangles
also shows directly that every ray has ramification index $n$.

\item The assertion is about the entire cusp fan. A cover can be smooth in
one triangle and singular in its neighbour, so testing one triple point
cannot replace this argument.
\end{enumerate}
\end{proof}

\begin{theorem}[no smooth global Kummer extension]\label{maps:src:C.1.3}
\StatementRecord{maps:src:C.1.3}{proved}{Historical events: Claude Fable 5.2 (Anthropic), effort high, 16 November 2026: carried as background without fault found. GPT-6 Astra (OpenAI), ChatGPT Pro mode, effort Pro, S6\_09\_PAPERS, 17 November 2026: targeted proof check. Also checked in this writing. Pre-continuation writing events: GPT-6 Astra (OpenAI), ChatGPT Pro mode, effort Pro, S6\_10\_PAPERS, 18 November 2026: targeted proof check. Continuation metadata events: GPT-6 Astra (OpenAI), ChatGPT Pro mode, effort Pro, S6\_10\_PAPERS continuation, 19 November 2026: used and rechecked in this run; status-vocabulary reconciliation; metadata reconciliation only; no new mathematical check. \emph{Pre-continuation recorded proof scope.} GPT-6 Astra (OpenAI, ChatGPT Pro mode), effort Pro, 2026-11-17; Claude Fable 5.2 (Anthropic), effort high, 16 November 2026; also checked in this writing by GPT-6 Astra (OpenAI), ChatGPT Pro mode, effort Pro, 18 November 2026}{the hypotheses, notation and definitions stated here; \cref{maps:src:C.1.1}; \cref{maps:eq:C.1.1}; \cref{maps:src:C.1.2}}


For $N>1$, the local map $\kappa_N$ in Theorem \ref{maps:src:C.1.1} cannot be the
restriction, even to one component of the inverse image of a cusp
neighbourhood, of a finite holomorphic map $Y\to X$ from a connected
smooth compact complex manifold. Additional cusp sheets and additional
vertical branch fibres do not remove the obstruction. For $N=1$ the
displayed local map is the identity and has the identity global extension.

\end{theorem}


\begin{proof}
\begin{enumerate}[label=\arabic*.,leftmargin=*]
\item Suppose such a finite map $F:Y\to X$ exists. Form the Stein factorization
of the proper map $fF$, with connected fibres $g:Y\to C'$ and finite map
$h:C'\to B$. Near the specified cusp point, the fibres of $g$ are tori
and their maps to the corresponding fibres of $f$ are etale isogenies,
with image lattice $L_N$ of \eqref{maps:eq:C.1.1}. No component of the ramification
divisor of $F$ can dominate $C'$: a proper dominating component would
meet every nearby regular fibre in this open set where $F$ is etale.
Thus all ramification is vertical. Delete its finitely many base values,
the branch values of $h$, and $0,1,\infty$. Over the remaining connected
base the map is a cover of torus bundles. On every sheet the image
of a fibre fundamental group is a finite-index sublattice of $\Lambda$,
and continuation along a loop with monodromy $A$ carries it to $AL_N$.
Translations do not change this subgroup, since the fibre group is abelian.

\item Follow a based meridian of the second finite fibre, choosing it to avoid
the additional deleted points, and then return to the same small
punctured cusp disc. The resulting sheet has image lattice $A_2L_N$.
It lies in some component of $F^{-1}(N_0)$; it need not be the original
component. Pull that component back to the unramified toric covering
$Y_\Sigma\to N_0$. This is still a finite cover with smooth total space,
etale off the central fibre. Each connected component therefore satisfies
Lemma \ref{maps:src:C.1.2}. Consequently its fibre image lattice must meet
$V_{\rm tor}$ in a scalar sublattice $nV_{\rm tor}$.

\item The actual second monodromy gives
\[
 A_2\widehat u=\widehat w,\qquad
 A_2\widehat w=-\widehat u+\widehat\delta,\qquad
 A_2\widehat\delta=\widehat\delta.
\]
The first coordinate of $A_2\widehat\gamma$ is one. Taking the
intersection of $A_2L_N$ with $V_{\rm tor}$ therefore gives exactly
\[
 A_2L_N\cap V_{\rm tor}
   =\Z\widehat w\oplus N\Z\widehat\delta.
 \tag{C.1.3}\label{maps:eq:C.1.3}
\]
It cannot equal $nV_{\rm tor}$ when $N>1$: the primitive first generator
forces $n=1$ but the second generator then excludes
$\widehat\delta$. This contradiction proves the theorem. Reversing the
based meridian replaces $A_2$ by its inverse and gives the same
non-scalar intersection, so the obstruction does not depend on that
orientation convention.
\end{enumerate}
\end{proof}

\begin{proposition}[a separate degree obstruction]\label{maps:src:C.1.4}
\StatementRecord{maps:src:C.1.4}{proved}{Historical events: Claude Fable 5.2 (Anthropic), effort high, 16 November 2026: carried as background without fault found. GPT-6 Astra (OpenAI), ChatGPT Pro mode, effort Pro, S6\_09\_PAPERS, 17 November 2026: targeted proof check. Also checked in this writing. Pre-continuation writing events: GPT-6 Astra (OpenAI), ChatGPT Pro mode, effort Pro, S6\_10\_PAPERS, 18 November 2026: targeted proof check. Continuation metadata events: GPT-6 Astra (OpenAI), ChatGPT Pro mode, effort Pro, S6\_10\_PAPERS continuation, 19 November 2026: used and rechecked in this run; status-vocabulary reconciliation; metadata reconciliation only; no new mathematical check. \emph{Pre-continuation recorded proof scope.} GPT-6 Astra (OpenAI, ChatGPT Pro mode), effort Pro, 2026-11-17; Claude Fable 5.2 (Anthropic), effort high, 16 November 2026; also checked in this writing by GPT-6 Astra (OpenAI), ChatGPT Pro mode, effort Pro, 18 November 2026}{the hypotheses, notation and definitions stated here; the explicit constructions and calculations in the proof}


Even allowing a normal singular source, there is no finite global
extension whose whole inverse image of $N_0$ is the single prescribed
map $\kappa_N$, for $N>1$. The monodromy orbit of $L_N$ has at least
$N^2$ elements, whereas such an extension would have a Stein base cover
of degree $N$.

\end{proposition}


\begin{proof}
\begin{enumerate}[label=\arabic*.,leftmargin=*]
\item Put $D=A_2^2A_1A_2^2A_1^{-1}$. Exact multiplication gives
\[
 D=\begin{pmatrix}1&0&0&0\\0&1&0&0\\6&0&1&0\\-6&1&0&1\end{pmatrix}.
\]
The $N^2$ lattices $D^aM_0^bL_N$, $0\leq a,b<N$, are distinct.
Indeed their intersections with $\ker\gamma$, reduced modulo $N\Lambda$,
are generated respectively by
$\widehat u+b\widehat w+a\widehat\delta$; the coefficient of
$\widehat u$ is one, so equality of these cyclic submodules forces
equality of both residues $a,b$. A base covering supporting one fibre
lattice must have at least as many sheets as its monodromy orbit.
The unique prescribed cusp has local connected-fibre parameter $v$
with $q=v^N$ and is the entire preimage, so the base degree would be
$N$. This contradicts $N^2>N$.

\item The exact orbit enumeration for $2\leq N\leq8$ is supplied in
\texttt{code/kummer\_orbits.py}; it gives respectively
$12,36,96,150,432,392,768$ elements. These finite computations check the
formulas in small cases. The preceding argument proves the required
lower bound for every $N$ and does not infer it from a finite search.
\end{enumerate}
\end{proof}

\begin{theorem}[normal singular globalizations]\label{maps:src:C.1.5}
\StatementRecord{maps:src:C.1.5}{proved}{Historical events: Claude Fable 5.2 (Anthropic), effort high, 16 November 2026: carried as background without fault found. GPT-6 Astra (OpenAI), ChatGPT Pro mode, effort Pro, S6\_09\_PAPERS, 17 November 2026: targeted proof check. Also checked in this writing. Pre-continuation writing events: GPT-6 Astra (OpenAI), ChatGPT Pro mode, effort Pro, S6\_10\_PAPERS, 18 November 2026: targeted proof check. Continuation metadata events: GPT-6 Astra (OpenAI), ChatGPT Pro mode, effort Pro, S6\_10\_PAPERS continuation, 19 November 2026: used and rechecked in this run; status-vocabulary reconciliation; metadata reconciliation only; no new mathematical check. \emph{Pre-continuation recorded proof scope.} GPT-6 Astra (OpenAI, ChatGPT Pro mode), effort Pro, 2026-11-17; Claude Fable 5.2 (Anthropic), effort high, 16 November 2026; also checked in this writing by GPT-6 Astra (OpenAI), ChatGPT Pro mode, effort Pro, 18 November 2026}{the hypotheses, notation and definitions stated here; \cref{maps:src:C.1.1}; \cref{maps:src:C.1.3}}


If a normal compact complex source is allowed, the prescribed local
Kummer cover does extend as one component above the cusp, for every $N$.
There is a connected finite flat such cover, branched only over
$S_1\cup S_2\cup W$, of degree
\[
 N^2|G_N|,\qquad
 G_N=\langle(A_1,-1),(A_2,0)\rangle
 \subset\operatorname{GL}_4(\Z/N)\times\Z/N.
 \tag{C.1.4}\label{maps:eq:C.1.4}
\]
Here matrices and the second coordinates are reduced modulo $N$.
The component in question has degree $N^3$ over $N_0$ and is precisely
$\kappa_N$. For $N>1$ other cusp components include singularities;
this is a cover by a normal space, not by a complex manifold.

\end{theorem}


\begin{proof}
\begin{enumerate}[label=\arabic*.,leftmargin=*]
\item Remove the three special fibres and write the result as $X^*\to B^*$.
The base retracts to a graph. We choose a continuous section that agrees
with the zero section of the marked torus family on a cusp collar.
Such a choice exists: the only obstruction is in
$H^2(B^*,\text{cusp collar};\Lambda)$, which is zero since the based
cusp loop is a free factor in the fundamental group of the punctured
sphere; equivalently, extend the given section along a relative graph
retraction. The resulting group presentation is
\[
 \pi_1(X^*)=\Lambda\rtimes F(a_1,a_2),\qquad
 a_0=(a_1a_2)^{-1},\quad a_j\lambda a_j^{-1}=A_j\lambda.
\]
The restriction to the cusp collar is the specified splitting with
monodromy $M_0$. A topological finite covering of a complex manifold
has a unique lifted holomorphic atlas, so it is enough first to
construct its subgroup.

\item Map the free group to $G_N$ by the two generators in \eqref{maps:eq:C.1.4} and
map the fibre lattice to $(\Z/N)^4$. This gives an epimorphism
\[
 \pi_1(X^*)\longrightarrow
 K_N=(\Z/N)^4\rtimes G_N,
\]
where only the matrix coordinate acts on the first factor. Let
$U_N\subset(\Z/N)^4$ be the subgroup spanned by $g,u$, of order
$N^2$, and take the connected cover corresponding to the inverse
image of $U_N$. Its degree is $|K_N|/|U_N|=N^2|G_N|$.
An element $\lambda a_0^k$ of the cusp group lies in this subgroup
exactly when $k\equiv0\pmod N$ and
$\lambda\in L_N$. Indeed $a_0$ has second coordinate $1$, and
$M_0^N\equiv I\pmod N$ because $(M_0-I)^2=0$.
Consequently its cusp subgroup is exactly
$L_N\rtimes\langle a_0^N\rangle$. The power map of \ref{maps:src:C.1.1} has
this subgroup, using the zero-section lift of the base loop.
The selected connected collar cover is therefore holomorphically
isomorphic to the punctured part of $\kappa_N$.

\item We construct the extension, including its normality and flatness.
At any point of the omitted divisor choose a polydisc with divisor
$z_1\cdots z_r=0$, $1\leq r\leq3$. Each connected component of
the restricted finite cover is classified by a finite-index subgroup
$K\subset\Z^r$. Write $M'=\operatorname{Hom}(K,\Z)$, containing
$\Z^r$ by restriction. Its completion is the normal toric analytic
space with ring generated by convergent monomials with exponents
\[
 M'\cap\R_{\geq0}^r
\]
over the remaining polydisc coordinates. This semigroup is saturated,
so the ring is normal. For each coset in $M'/\Z^r$ choose its unique
representative whose coordinates belong to $[0,1)$. The corresponding
monomials form a free basis over the original polydisc ring: subtracting
the representative leaves nonnegative integer coordinates, and distinct
cosets cannot cancel. Thus the extension is finite flat. It is also
the integral closure in the finite covering's local sheaf of meromorphic
functions; we do not use the field of global meromorphic functions on
the compact nonalgebraic space.
Integral closure is unique, so these local extensions glue on all
overlaps, including the self-identifications of $W$.

\item The resulting space is normal and connected; finiteness over compact
$X$ makes it compact. The same uniqueness identifies the selected
cusp completion with the already normal smooth source in \ref{maps:src:C.1.1}.
Theorem \ref{maps:src:C.1.3} forces at least one other cusp component to be singular
when $N>1$. This proves the distinction between the two source
categories and the asserted degree. Exact enumeration gives
$|G_2|=48$, $|G_3|=216$, $|G_4|=3072$; the accompanying program
checks these finite examples, while the construction proves the
statement for arbitrary $N$.
\end{enumerate}
\end{proof}

The local cover constructions have the two different degrees just computed. Non-invariance of its fibre sublattice alone
does not prohibit a further base cover. The entire-fan argument excludes
smooth global sources; the finite quotient in \ref{maps:src:C.1.5} gives the different
answer for normal singular sources with additional cusp sheets.


\section{The other six-manifolds}
\subsection{The cyclic covers: a necessary correction}
\begin{theorem}[Counterexample: the proposed connected sum for cyclic covers]\label{maps:src:D.1.1}
\StatementRecord{maps:src:D.1.1}{proved}{Historical events: Claude Fable 5.2 (Anthropic), effort high, 16 November 2026: carried as background without fault found. GPT-6 Astra (OpenAI), ChatGPT Pro mode, effort Pro, S6\_09\_PAPERS, 17 November 2026: targeted proof check. Also checked in this writing. Pre-continuation writing events: GPT-6 Astra (OpenAI), ChatGPT Pro mode, effort Pro, S6\_10\_PAPERS, 18 November 2026: targeted proof check. Continuation metadata events: GPT-6 Astra (OpenAI), ChatGPT Pro mode, effort Pro, S6\_10\_PAPERS continuation, 19 November 2026: used and rechecked in this run; status-vocabulary reconciliation; metadata reconciliation only; no new mathematical check. \emph{Pre-continuation recorded proof scope.} GPT-6 Astra (OpenAI, ChatGPT Pro mode), effort Pro, 2026-11-17 UTC; Claude Fable 5.2 (Anthropic), effort high, 16 November 2026; also checked in this writing by GPT-6 Astra (OpenAI), ChatGPT Pro mode, effort Pro, 18 November 2026}{the hypotheses, notation and definitions stated here; the explicit constructions and calculations in the proof}


Put $d=|12\ell_0-4\ell_1-3\ell_2|$ for an admissible member $X_p$.
The universal cover $\widetilde X_p$ is a compact complex threefold of
Euler characteristic $2d$. Its Betti numbers necessarily satisfy
\[
 b_3(\widetilde X_p)=2+2b_2(\widetilde X_p)-2d,
 \qquad b_2(\widetilde X_p)\geq d-1.
\]
Consequently, for $d>1$ its second integral homology cannot vanish,
and it cannot be diffeomorphic to $\#^{d-1}(S^3\times S^3)$.
The proposed connected-sum classification under vanishing second homology has an impossible premise for $d>1$;
it is not a classification of the covers.
\end{theorem}


\begin{proof}
\begin{enumerate}[label=\arabic*.,leftmargin=*]
\item The covering has $d$ sheets, is unramified, and pulls back the complex
atlas, so is compact, complex, and oriented. Multiplicativity of Euler
characteristic for a finite covering follows by lifting a finite CW
structure: every cell has exactly $d$ lifts. Thus
$\chi(\widetilde X_p)=d\chi(X_p)=2d$.
Simple connectivity gives $b_1=b_5=0$, and oriented Poincare duality gives
$b_4=b_2$, $b_0=b_6=1$. Substitution gives the displayed identity. The
inequality follows from $b_3\geq0$. Finally a connected sum of $r$
copies of $S^3\times S^3$ has Euler characteristic $2-2r$; for $r=d-1$
this equals $4-2d$, which equals $2d$ only when $d=1$.
For a concrete admissible example take $(\ell_0,\ell_1,\ell_2)=(0,2,-1)$.
Then $d=5$, the cover has Euler characteristic $10$, whereas
$\#^4(S^3\times S^3)$ has Euler characteristic $-6$.
\end{enumerate}
\end{proof}

\begin{theorem}[the cusp of the cyclic cover]\label{maps:src:D.1.2}
\StatementRecord{maps:src:D.1.2}{proved}{Historical events: Claude Fable 5.2 (Anthropic), effort high, 16 November 2026: carried as background without fault found. GPT-6 Astra (OpenAI), ChatGPT Pro mode, effort Pro, S6\_09\_PAPERS, 17 November 2026: targeted proof check. Also checked in this writing. Pre-continuation writing events: GPT-6 Astra (OpenAI), ChatGPT Pro mode, effort Pro, S6\_10\_PAPERS, 18 November 2026: targeted proof check. Continuation metadata events: GPT-6 Astra (OpenAI), ChatGPT Pro mode, effort Pro, S6\_10\_PAPERS continuation, 19 November 2026: used and rechecked in this run; status-vocabulary reconciliation; metadata reconciliation only; no new mathematical check. \emph{Pre-continuation recorded proof scope.} GPT-6 Astra (OpenAI, ChatGPT Pro mode), effort Pro, 2026-11-17 UTC; Claude Fable 5.2 (Anthropic), effort high, 16 November 2026; also checked in this writing by GPT-6 Astra (OpenAI), ChatGPT Pro mode, effort Pro, 18 November 2026}{the hypotheses, notation and definitions stated here; \cite[Theorem~7.17]{AC26}}


The inverse image $\widetilde W_p$ of the cusp is connected. Its
normalization has $d$ components, each isomorphic to $dP_6$. Its
integral homology is torsion free, with ranks
\[
 (b_0,b_1,b_2,b_3,b_4)(\widetilde W_p)=(1,2,d+3,2,d).
\]
Here the covering of the dual torus is the sublattice
$\mathbb Z\widehat u\oplus d\mathbb Z\widehat\gamma$ of
$\overline\Lambda=\mathbb Z\widehat\gamma\oplus\mathbb Z\widehat u$.
\end{theorem}


\begin{proof}
\begin{enumerate}[label=\arabic*.,leftmargin=*]
\item The presentation of the fundamental group in
\cite[Theorem~7.17]{AC26} sends the regular-fibre lattice onto
$\mathbb Z/d$ by its primitive $\gamma$ coordinate: $12$ is a unit
modulo $d$, and $3x=\ell_1c$, $4y=\ell_2c$ express the other two
meridians as multiples of $c$. On the cusp $\pi_1(W)=\overline\Lambda$,
and the same map is surjective. Hence the inverse image is connected,
with the asserted lattice of deck periods. Since
$B_0(\widehat u)=e_1$ and $B_0(\widehat\gamma)=-e_2$, its toric
components are indexed by $\mathbb Z^2/(\mathbb Ze_1+d\mathbb Ze_2)$.

\item We compute the normalization resolution, retaining integral coefficients.
Let $R=\mathbb Z[t]/(t^d-1)$, where $t$ moves one component in the
$e_2$ direction. The degree-zero row of the resolution is the cellular
cochain complex of the dual torus with this $d$-sheeted marking; its
cohomology is $(\mathbb Z,\mathbb Z^2,\mathbb Z)$ in degrees $0,1,2$.
The degree-four row is $R$ in horizontal degree zero. In degree two,
write the six boundary curves on $dP_6$ in cyclic order as
\[
 E_1,\quad H-E_1-E_2,\quad E_2,\quad H-E_2-E_3,
 \quad E_3,\quad H-E_1-E_3.
\]
The restriction differential from the $d$ normalizations to the $3d$
double curves, in the basis $(H,E_1,E_2,E_3)$, is
\[
 B(t)=
 \begin{pmatrix}
 -1&-1&-1&-1\\
 1&1&1&t\\
 -t&-t&-1&-t
 \end{pmatrix}:R^4\longrightarrow R^3.
\]
Replacing the second row by its sum with the first and the third by
its difference with $t$ times the first gives rows with sole entries
$t-1$ in columns four and three. The first-row unit eliminates the
remaining first-row entries. Thus this matrix is equivalent over $R$
to a matrix with diagonal entries $1,t-1,t-1$ and one zero column.
Multiplication by $t-1$ on $R$ has kernel $\mathbb Z(1+t+\cdots+t^{d-1})$
and cokernel $\mathbb Z$. Therefore $\ker B$ is free abelian of rank
$d+2$, and $\operatorname{coker}B\simeq\mathbb Z^2$.
There can be no higher differential: all rows with odd vertical degree
are zero and the resolution has horizontal degrees only $0,1,2$.
The only extension in degree two is an extension of free abelian
groups, hence splits as an abelian group. This gives cohomology ranks
$(1,2,d+3,2,d)$ with no torsion. The universal coefficient theorem,
applied from the top degree downwards, gives the same assertion for
homology. The normalization resolution used here is exact on stalks:
at a point with $r$ local branches it is the augmented cochain complex
of the $(r-1)$-simplex, and hence computes ordinary cohomology even at
the self-identifications and the triple points.
\end{enumerate}
\end{proof}

\begin{theorem}[the rational homology of the cover]\label{maps:src:D.1.3}
\StatementRecord{maps:src:D.1.3}{proved}{Historical events: Claude Fable 5.2 (Anthropic), effort high, 16 November 2026: carried as background without fault found. GPT-6 Astra (OpenAI), ChatGPT Pro mode, effort Pro, S6\_09\_PAPERS, 17 November 2026: targeted proof check. Also checked in this writing. Pre-continuation writing events: GPT-6 Astra (OpenAI), ChatGPT Pro mode, effort Pro, S6\_10\_PAPERS, 18 November 2026: targeted proof check. Continuation metadata events: GPT-6 Astra (OpenAI), ChatGPT Pro mode, effort Pro, S6\_10\_PAPERS continuation, 19 November 2026: used and rechecked in this run; status-vocabulary reconciliation; metadata reconciliation only; no new mathematical check. \emph{Pre-continuation recorded proof scope.} GPT-6 Astra (OpenAI, ChatGPT Pro mode), effort Pro, 2026-11-17 UTC; Claude Fable 5.2 (Anthropic), effort high, 16 November 2026; also checked in this writing by GPT-6 Astra (OpenAI), ChatGPT Pro mode, effort Pro, 18 November 2026}{the hypotheses, notation and definitions stated here; the explicit constructions and calculations in the proof}


For every admissible $p$, the rational homology of $\widetilde X_p$ has
ranks
\[
 (1,0,d-1,0,d-1,0,1).
\]
The deck group on either $H^2(\widetilde X_p;\mathbb C)$ or
$H^4(\widetilde X_p;\mathbb C)$ is the direct sum of its nontrivial
characters, each with multiplicity one. This statement does not
assert the absence of integral torsion.
\end{theorem}


\begin{proof}
\begin{enumerate}[label=\arabic*.,leftmargin=*]
\item Let $\lambda\ne1$ be a $d$th root of unity and let
$\mathcal L_\lambda$ be the associated rank-one complex local system
on $X_p$. Its restriction to a regular fibre is nontrivial on
$\widehat\gamma$. The cellular cochain complex of a torus with this
local system is a Koszul complex containing the invertible scalar
$\lambda-1$ and is therefore contractible. On a reduced multiple
fibre the finite torus cover has the same property. Transfer, or
averaging over the finite deck group in characteristic zero, shows
that its cohomology also vanishes. The regular region and all three
boundary overlaps consequently have zero cohomology with these
coefficients, as follows either by their bundle spectral sequences
or by contraction of the fibre cochain complex. Mayer--Vietoris
therefore identifies $H^*(X_p;\mathcal L_\lambda)$ with
$H^*(W;\mathcal L_\lambda|_W)$.
In the preceding normalization calculation substitute $t=\lambda$.
The degree-zero row is acyclic, the matrix $B(\lambda)$ has rank three
and one-dimensional kernel, and the degree-four row is one-dimensional.
Thus this local system contributes one copy of $\mathbb C$ in degrees
two and four and nothing else. The trivial character contributes
$H^*(X_p;\mathbb C)$, which is concentrated in degrees zero and six by
\textup{[AC26, Theorem 7.22]}. The decomposition of the regular
representation of $\mathbb Z/d$ into its characters computes the
cohomology of the covering and proves the claim.
\end{enumerate}
\end{proof}

\begin{theorem}[the integral homology of the universal cover]\label{maps:src:D.1.4}
\StatementRecord{maps:src:D.1.4}{proved}{Historical events: Claude Fable 5.2 (Anthropic), effort high, 16 November 2026: carried as background without fault found. Claude Opus 5.5 (Anthropic), effort max, sweep 1, scope s6, 18 November 2026: computation or formal certificate. GPT-6 Astra (OpenAI), ChatGPT Pro mode, effort Pro, S6\_09\_PAPERS, 17 November 2026: targeted proof check. Also checked in this writing. Pre-continuation writing events: GPT-6 Astra (OpenAI), ChatGPT Pro mode, effort Pro, S6\_10\_PAPERS, 18 November 2026: targeted proof check. Continuation metadata events: GPT-6 Astra (OpenAI), ChatGPT Pro mode, effort Pro, S6\_10\_PAPERS continuation, 19 November 2026: used and rechecked in this run; status-vocabulary reconciliation; metadata reconciliation only; no new mathematical check. \emph{Pre-continuation recorded proof scope.} GPT-6 Astra (OpenAI, ChatGPT Pro mode), effort Pro, 2026-11-17 UTC; Claude Fable 5.2 (Anthropic), effort high, 16 November 2026; also checked in this writing by GPT-6 Astra (OpenAI), ChatGPT Pro mode, effort Pro, 18 November 2026}{the hypotheses, notation and definitions stated here; \cref{maps:src:D.1.3}}


For every admissible triple, with $d=|p|$, one has
\[
 H_i(\widetilde X_p;\mathbb Z)=
 \begin{cases}
 \mathbb Z,&i=0,6,\\
 \mathbb Z^{d-1}\oplus\mathbb Z/d,&i=2,\\
 \mathbb Z/d,&i=3,\\
 \mathbb Z^{d-1},&i=4,\\
 0,&i=1,5.
 \end{cases}
\]
Thus for $d>1$ the cover is also not a connected sum of copies of
$S^2\times S^4$. In particular a torsion-free classification theorem
cannot be applied to these covers.
\end{theorem}


\begin{proof}
\begin{enumerate}[label=\arabic*.,leftmargin=*]
\item We give the integral attachment calculation. Write
$\gamma'=\gamma/d$ on the covering lattice
$\Lambda_d=\mathbb Z(d\widehat\gamma)\oplus\mathbb Z\widehat u
\oplus\mathbb Z\widehat w\oplus\mathbb Z\widehat\delta$.
Choose integers $h_j$ with
$m_jh_j\equiv\ell_j\pmod d$. Replace each meridian lift by its
product with the translation $-h_j\widehat\gamma$. Its $m_j$th power
has translation vector
$v_j-h_j\mathcal N_j\widehat\gamma$, where
\[
 \mathcal N_1\widehat\gamma=3\varepsilon,
 \qquad \mathcal N_2\widehat\gamma=4\varepsilon'.
\]
The non-invariant part of an affine translation is removed by
conjugation by a real torus translation: averaging projects to the
invariant part, and its complementary part is in the image of $A_j-I$
over $\mathbb R$. This produces precisely the same smooth local
quotient with invariant translation vector. Its new integers are
\[
 a_j=(\ell_j-m_jh_j)/d,\qquad
 a_0=(\ell_0-h_1-h_2)/d.
\]
The numerator defining $a_0$ is divisible by $d$ by the original
meridian relation. They satisfy
\[
 p'=12a_0-4a_1-3a_2=p/d\in\{1,-1\}.
\]
In the basis of $\Lambda_d$, the two matrices are
\[
 A_1^{(d)}=\begin{pmatrix}1&0&0&0\\6d&0&1&0\\-6d&-1&-1&0\\-2d&1&0&1\end{pmatrix},
 \quad
 A_2^{(d)}=\begin{pmatrix}1&0&0&0\\0&0&-1&0\\-6d&1&0&0\\3d&0&1&1\end{pmatrix}.
\]
Their coinvariant lattices are the same saturated rank-two quotients
as for $d=1$. The invariant forms are
$\psi_1=2u+w+3\delta$, $\psi_2=u+w+2\delta$,
and their common invariant alternating form is
$q_d=u\wedge w+6d\gamma'\wedge\delta$.

\item Let $\widetilde X^\circ$ be the complement of the cusp tube. The
finite reduced fibres, denoted $S'_1,S'_2$, have torsion-free first
homology: their abelianized affine presentation is the free rank-two
coinvariant lattice with one meridian adjoined, modulo
$m_ja_j^{\rm mer}=[v'_j]$, a primitive relation since
$\gcd(a_j,m_j)=1$. Poincare duality, Euler characteristic zero and
the universal coefficient theorem give $H_2(S'_j)=\mathbb Z^2$.
The intersection form is even and unimodular: the canonical class is
torsion for a bielliptic surface and $H^2(S'_j)$ is torsion free,
so its mod-two reduction, the second Stiefel--Whitney class, is zero.

\item The pullback images in $H^2(F_d)$ have respective integral bases
\[
 \begin{aligned}
 L_1&=\langle\gamma'\psi_1,\ q_d-2d\gamma'\psi_1\rangle,\\
 L_2&=\langle2\gamma'\psi_2,\ q_d-3d\gamma'\psi_2\rangle.
 \end{aligned}
\]
Here is a verification of the sublattice assertion, including its
index. The invariant lattices are
$\langle\gamma'\psi_j,q_d\rangle$, with intersection determinants
$-9$ and $-4$. Pullback multiplies intersection pairings by $3$ and
$4$, so its indices are respectively one and two. At the second
fibre, $q_d$ cannot descend since its square divided by four is
$3d$, odd. The class $\gamma'\psi_2$ cannot descend either: on the
torus swept by $\delta$ and the affine meridian its integral would
be $-a_2/2$, which is not integral. Of the three index-two sublattices
of the invariant lattice, only the parity sublattice
$\{a\gamma'\psi_2+bq_d:a\equiv b\pmod2\}$ excludes both classes.
The displayed basis generates that sublattice because $d$ is odd.
Its pairing divided by four, and the first pairing divided by three,
are the hyperbolic plane.

\item In homology bases dual to these cohomology bases, the map from the
common regular fibre to the two finite pieces is
\[
 \alpha_2=
 \begin{pmatrix}
 2&1&3&0&0&0\\
 -4d&-2d&0&1&0&0\\
 -2&-2&-4&0&0&0\\
 3d&3d&0&-1&0&0
 \end{pmatrix}.
\]
Its rank is three. Among its maximal minors are $2$, $3d$, $4d$;
their greatest common divisor is one because $d$ is odd. Therefore
its cokernel is $\mathbb Z$, whose primitive dual is
\[
 \Phi=(4d,2,3d,2).
\]
The degree-one map has kernel
$\mathbb Z(-\widehat u-\widehat w+\widehat\delta)$, exactly as at
$d=1$. Mayer--Vietoris for the two finite pieces therefore gives
\[
 0\longrightarrow\mathbb Z\longrightarrow H_2(\widetilde X^\circ)
 \longrightarrow\mathbb Z\longrightarrow0.
\]
Furthermore $H_1(\widetilde X^\circ)=\mathbb Z$, with primitive dual
$\Theta$ restricting to $12\gamma'$ on $F_d$; the cusp meridian
has value $p'=\pm1$ under $\Theta$.

\item Write $N$ for the cusp tube and $M=\partial N$. The exterior-square
matrix of its monodromy has saturated image generated by
$\widehat w\wedge\widehat\delta$ and
$(d\widehat\gamma)\wedge\widehat w+d\widehat u\wedge\widehat\delta$.
Thus the Wang sequence makes $H_2(M)$ free of rank six. Sweeping the
two angular circles $\widehat w,\widehat\delta$ along the toric
meridian gives two classes $\xi_w,\xi_\delta$ which map under the
Wang projection to a basis of the invariant lattice. Both bound
integrally in $N$: sweep the corresponding circle along the extended
zero-section disc in the toric chart. Their span is primitive in
$H_2(M)$. On the other hand
$H_3(N,M)\simeq H^3(N)\simeq H^3(\widetilde W_p)=\mathbb Z^2$;
exactness bounds the kernel of $H_2(M)\to H_2(N)$ by rank two.
Consequently that kernel is exactly
$\mathbb Z\xi_w\oplus\mathbb Z\xi_\delta$.

\item Cutting $\xi_w$ into the two finite pieces gives connecting boundary
$\pm(-\widehat u-\widehat w+\widehat\delta)$, a primitive generator,
whereas $\xi_\delta$ has zero connecting boundary. The latter belongs
to $\operatorname{coker}\alpha_2$. If $\mathcal A$ is the class dual
to $\Phi$, then $\mathcal A|_{F_d}=2q_d$. Integration along the
free $\delta$-circle on $\widetilde X^\circ$ gives
\[
 (\iota_\delta\mathcal A)|_{F_d}=\pm12d\gamma',
 \qquad \iota_\delta\mathcal A=\pm d\Theta.
\]
The second equality follows from injectivity of restriction in degree
one. Evaluation on the meridian therefore gives
$\Phi(j_*\xi_\delta)=\pm d p'=\pm d$. Hence
\[
 H_2(\widetilde X^\circ)/j_*\ker(H_2(M)\to H_2(N))
 \simeq\mathbb Z/d.
\]
This evaluates the actual swept torus, rather than inferring its
coefficient from rational homology.

\item Finally, the exact sequence of the pair
$(\widetilde X_p,\widetilde X^\circ)$ and excision give
\[
 0\longrightarrow\mathbb Z/d\longrightarrow H_2(\widetilde X_p)
 \longrightarrow\ker\{H_2(N,M)\to H_1(\widetilde X^\circ)\}
 \longrightarrow0.
\]
Poincare--Lefschetz duality identifies $H_2(N,M)$ with
$H^4(\widetilde W_p)=\mathbb Z^d$. The extended section disc has
boundary the meridian, whose value is $p'=\pm1$, so the final map is
onto $\mathbb Z$ and its kernel is free of rank $d-1$. The sequence
splits as an abelian-group sequence. The rational calculation of
\ref{maps:src:D.1.3} gives $b_3=0$; duality and universal coefficients now yield
$H_3=\operatorname{Ext}(H_2,\mathbb Z)=\mathbb Z/d$,
$H_4=H^2=\mathbb Z^{d-1}$ and $H_5=0$.
\end{enumerate}
\end{proof}

\begin{proposition}[characteristic and rational ring constraints]\label{maps:src:D.1.5}
\StatementRecord{maps:src:D.1.5}{proved}{Historical events: Claude Fable 5.2 (Anthropic), effort high, 16 November 2026: carried as background without fault found. GPT-6 Astra (OpenAI), ChatGPT Pro mode, effort Pro, S6\_09\_PAPERS, 17 November 2026: targeted proof check. Also checked in this writing. Pre-continuation writing events: GPT-6 Astra (OpenAI), ChatGPT Pro mode, effort Pro, S6\_10\_PAPERS, 18 November 2026: targeted proof check. Continuation metadata events: GPT-6 Astra (OpenAI), ChatGPT Pro mode, effort Pro, S6\_10\_PAPERS continuation, 19 November 2026: used and rechecked in this run; status-vocabulary reconciliation; metadata reconciliation only; no new mathematical check. \emph{Pre-continuation recorded proof scope.} GPT-6 Astra (OpenAI, ChatGPT Pro mode), effort Pro, 2026-11-17 UTC; Claude Fable 5.2 (Anthropic), effort high, 16 November 2026; also checked in this writing by GPT-6 Astra (OpenAI), ChatGPT Pro mode, effort Pro, 18 November 2026}{the hypotheses, notation and definitions stated here; \cref{maps:src:D.1.3}}


For the complex orientation, $w_2(\widetilde X_p)=0$,
$c_1(\widetilde X_p)=0$, and $p_1(\widetilde X_p)$ is torsion of
order dividing $d$. The rational cup product
$H^2(\widetilde X_p;\mathbb Q)\otimes H^2(\widetilde X_p;\mathbb Q)
\to H^4(\widetilde X_p;\mathbb Q)$ is zero. The integral Pontryagin
class is not determined by these assertions.
\end{proposition}


\begin{proof}
\begin{enumerate}[label=\arabic*.,leftmargin=*]
\item The mod-two cohomology $H^2(X_p;\mathbb Z/2)$ vanishes, because all
the integral torsion in degrees one and two has odd order. Thus
$w_2(X_p)=0$ and the same holds on its cover. The class $c_1(X_p)$
is torsion and its pullback lies in the torsion-free group
$H^2(\widetilde X_p;\mathbb Z)$; hence the pullback is zero.
Similarly $H^4(X_p;\mathbb Z)=\mathbb Z/d$, so the pullback $p_1$
is annihilated by $d$.
For the ring statement, the restriction of every nontrivial-character
class in degree two to the normalization of a component of the cusp
lies in the kernel of $B(t)$ for $t\ne1$. This kernel is generated
by $-H+E_1$, whose square is zero in $H^4(dP_6)$.
The product of any two such restricted classes is therefore zero.
The restriction from degree-four cohomology of the covering to
$\bigoplus H^4(dP_6;\mathbb Q)$ is injective: it is the characterwise
isomorphism of \ref{maps:src:D.1.3} on the nontrivial summands, and the source has
no invariant summand. This proves that the original product is zero.
\end{enumerate}
\end{proof}

\begin{theorem}[the integral cohomology ring]\label{maps:src:D.1.6}
\StatementRecord{maps:src:D.1.6}{proved}{Historical events: Claude Fable 5.2 (Anthropic), effort high, 16 November 2026: carried as background without fault found. GPT-6 Astra (OpenAI), ChatGPT Pro mode, effort Pro, S6\_09\_PAPERS, 17 November 2026: targeted proof check. Also checked in this writing. Pre-continuation writing events: GPT-6 Astra (OpenAI), ChatGPT Pro mode, effort Pro, S6\_10\_PAPERS, 18 November 2026: targeted proof check. Continuation metadata events: GPT-6 Astra (OpenAI), ChatGPT Pro mode, effort Pro, S6\_10\_PAPERS continuation, 19 November 2026: used and rechecked in this run; status-vocabulary reconciliation; metadata reconciliation only; no new mathematical check. \emph{Pre-continuation recorded proof scope.} GPT-6 Astra (OpenAI, ChatGPT Pro mode), effort Pro, 2026-11-17 UTC; Claude Fable 5.2 (Anthropic), effort high, 16 November 2026; also checked in this writing by GPT-6 Astra (OpenAI), ChatGPT Pro mode, effort Pro, 18 November 2026}{the hypotheses, notation and definitions stated here; \cref{maps:src:D.1.4}; \cref{maps:src:D.1.3}; \cref{maps:src:D.1.2}}


Suppose $d>1$, so in this construction $d\geq5$. Number the cusp
components $D_j$, $j\in\mathbb Z/d$, by the positive $e_2$ deck
translation, and give all complex subvarieties their complex orientations.
Write $x_j=\operatorname{PD}[D_j]\in H^2(\widetilde X_p;\mathbb Z)$.
These classes generate, with the sole relation $\sum_jx_j=0$.
There is a distinguished generator $\kappa$ of the torsion subgroup of
$H^4(\widetilde X_p;\mathbb Z)$, namely the dual of the complex-oriented
ruling class $H-E_1$ in any normalized cusp component. The products are
\[
 x_i x_j=\begin{cases}
 -2\kappa,&i=j,\\
 \kappa,&i-j=1\text{ or }-1\pmod d,\\
 0,&\text{otherwise}.
 \end{cases}
\]
Together with the orientation pairing between the free parts of $H^2$
and $H^4$, these are every product in positive degree. For $d=1$ the
ring is the ring of $S^6$.
\end{theorem}


\begin{proof}
\begin{enumerate}[label=\arabic*.,leftmargin=*]
\item The relative cohomology group $H^2(N,M)=H_4(N)=\mathbb Z^d$ is
generated by the Thom classes of the irreducible Cartier divisors $D_j$.
This remains true for their self-identifications, since their fundamental
cycles generate $H_4(N)$ and Poincare--Lefschetz duality applies to the
smooth tube. The boundary homomorphism from $H^1(\widetilde X^\circ)$
takes its primitive generator to $(\pm1,\ldots,\pm1)$, with a common
sign. Indeed the toric meridians at adjacent height-one rays differ by
an angular $\widehat w$ or $\widehat\delta$ circle, and both have
zero image under $\Theta$; \ref{maps:src:D.1.4} evaluates any one of them as $p'=\pm1$.
Consequently $\mathbb Z^d/\mathbb Z(1,\ldots,1)$ injects into
$H^2(\widetilde X_p)$. The cokernel embeds in the torsion-free group
$H^2(\widetilde X^\circ)=\mathbb Z^2$ and has rank zero by \ref{maps:src:D.1.3}.
It is therefore zero. This proves the assertion about the generators
integrally, including saturation.

\item For clarity we compute the relevant curve relations before their
images in the whole manifold. The first opposite pair on component
$j$ identifies $E_{1j}$ with $H_j-E_{2j}-E_{3j}$. The other two
opposite pairs identify $E_{2j}$ with $E_{2,j+1}$ and $E_{3j}$ with
$E_{3,j+1}$, after using the first relation. These are precisely the
transpose relations of $B(t)$ in \ref{maps:src:D.1.2}. Thus the ruling classes
$f_j=H_j-E_{1j}=E_{2j}+E_{3j}$ all have the same image $f$ in
$H_2(N)$. Components $j$ and $j+1$ intersect along the two external
boundary curves, with multiplicity one and total class $f$; no
nonadjacent components meet. Hence $x_jx_{j+1}=\operatorname{PD}(f)$.
Multiplying $\sum_ix_i=0$ by $x_j$ gives $x_j^2=-2\operatorname{PD}(f)$.

\item It remains to verify that $f$ is a torsion generator, rather than a
proper multiple. The torus $\widehat u\wedge\widehat\delta$ specializes
to one of the two external boundary-curve classes, up to sign, and
$\widehat u\wedge\widehat w$ specializes to their signed sum. One
can check this in a single strip of the covering tiling: the base
circle has direction $B_0\widehat u=e_1$; it crosses the two edges
with stabilizer directions $e_2-e_1$ and $-e_2$, and its degree on
each surviving angular circle is the determinant with that direction.
These determinants are units. This is the integral annulus-collapse
calculation of AC26 (7.9), applied to the unchanged $e_1$ period.
In the first finite piece, $\widehat u\wedge\widehat\delta$ has
coordinates $(0,0)$ in the dual bases of \ref{maps:src:D.1.4}, whereas
$\widehat u\wedge\widehat w$ has coordinates $(0,1)$.
The functional $\Phi=(4d,2,3d,2)$ therefore evaluates their images
as $0$ and $2$, respectively. It follows that one external curve
has image zero and the other has image $\pm2$ in
$\mathbb Z/d\subset H_2(\widetilde X_p)$. Their complex-oriented
sum $f$ is consequently a generator, since $d$ is odd. Its sign
defines the stated $\kappa$, so the intersection formulas themselves
have no remaining sign ambiguity.

\item All products involving a torsion class and landing in $H^6=\mathbb Z$
vanish. The groups in degrees one and five vanish, and a product of
two degree-three classes is torsion and hence zero. Finally Poincare
duality gives the perfect integral pairing of the free degree-two
and degree-four groups. This exhausts all degree possibilities.
\end{enumerate}
\end{proof}

\begin{theorem}[dependence of the smooth type and the count]\label{maps:src:D.1.7}
\StatementRecord{maps:src:D.1.7}{proved}{Historical events: Claude Fable 5.2 (Anthropic), effort high, 16 November 2026: carried as background without fault found. GPT-6 Astra (OpenAI), ChatGPT Pro mode, effort Pro, S6\_09\_PAPERS, 17 November 2026: targeted proof check. Also checked in this writing. Pre-continuation writing events: GPT-6 Astra (OpenAI), ChatGPT Pro mode, effort Pro, S6\_10\_PAPERS, 18 November 2026: targeted proof check. Continuation metadata events: GPT-6 Astra (OpenAI), ChatGPT Pro mode, effort Pro, S6\_10\_PAPERS continuation, 19 November 2026: used and rechecked in this run; status-vocabulary reconciliation; metadata reconciliation only; no new mathematical check. \emph{Pre-continuation recorded proof scope.} GPT-6 Astra (OpenAI, ChatGPT Pro mode), effort Pro, 2026-11-17 UTC; Claude Fable 5.2 (Anthropic), effort high, 16 November 2026; also checked in this writing by GPT-6 Astra (OpenAI), ChatGPT Pro mode, effort Pro, 18 November 2026}{the hypotheses, notation and definitions stated here; \cref{maps:src:D.1.9}; \cref{maps:src:D.1.12}}


Within the multiplicity-$(3,4)$ construction, the oriented
diffeomorphism type of $X_p$ depends exactly on $d=|p|$.
For each $d$ coprime to six it carries a punctured-disc family of
pairwise non-biholomorphic complex structures, parametrized by
$W=e^{-\pi i(c+k_2)}$, $0<|W|<216$. At each fixed nonzero admissible
linearization scale $u=e^{-2\pi ic}$ it carries exactly two of these
biholomorphism classes. Its universal covers likewise have one
oriented diffeomorphism type for each $d$. This statement identifies
dependence on the construction data. The complete Zhubr classifying
invariant and a canonical cyclic action are supplied in \ref{maps:src:D.1.9}--\ref{maps:src:D.1.12}.
\end{theorem}


\begin{proof}
\begin{enumerate}[label=\arabic*.,leftmargin=*]
\item We use the all-twist comparison in \cref{fam:src:B.2.1}, whose sufficient step is the integral norm equation and analytic torsor gluing. Specifically, equality of $|p|$ and of $W$ gives
$p'=ap$, $a=\pm1$, and $c'-c\equiv k_2-k'_2\pmod2$.
The finite norm equations have integral solutions; the remaining
cusp defect is corrected in its entire rank-two quotient lattice,
and $H^1(\mathcal O\oplus\mathcal O(-1))=0$ glues the translations.
The cusp map preserves the prescribed fan. These are the steps that
give actual biholomorphisms, including every discrete twist, rather
than equality of abstract numerical data.

\item For fixed discrete data, varying $c$ along any compact path in the
admissible half-plane gives a proper smooth family. Here the
properness argument of \cref{fam:src:A.3.1} uses a bounded correction
$C(t,c)$ in $Z=sB_0+C$: on a compact parameter path the effective
logarithmic lattice is uniformly close to $B_0$, so finitely many
closed toric polydiscs provide compact representatives. The finite
free quotients and the ordinary torus parallelepipeds are uniform
as well. All charts are products with the parameter before their
free identifications, so projection is a submersion. The same proof
does not use $p=\pm1$, and therefore applies to the present data.
Integrating a horizontal vector field over the compact path gives
an orientation-preserving diffeomorphism of its end fibres.

\item Given two data with the same $d$, first change one parameter by an
integer to align its value of $c+k_2$ with that of the other
discrete datum; use the biholomorphism of \cref{fam:src:B.2.1} at this aligned parameter,
and connect the resulting parameter to the desired one inside the
admissible half-plane. This proves independence of every datum
except $d$. Distinct $d$ give nonisomorphic fundamental groups.
Every resulting diffeomorphism lifts to an oriented diffeomorphism
of universal covers; different $d$ are distinguished upstairs by
Euler characteristic $2d$.

\item By \cref{fam:src:B.2.1} and \cref{fam:src:C.2.3}, the map to $W$ is an exact biholomorphic classification.
Since $u=W^2$, its fibre at fixed $u\ne0$ is precisely the two
points $W$ and $-W$. No two of these are equal. The asserted count
and the punctured-disc family follow.
\end{enumerate}
\end{proof}

\begin{theorem}[the integral Pontryagin class]\label{maps:src:D.1.8}
\StatementRecord{maps:src:D.1.8}{proved}{Historical events: Claude Fable 5.2 (Anthropic), effort high, 16 November 2026: carried as background without fault found. GPT-6 Astra (OpenAI), ChatGPT Pro mode, effort Pro, S6\_09\_PAPERS, 17 November 2026: targeted proof check. Also checked in this writing. Pre-continuation writing events: GPT-6 Astra (OpenAI), ChatGPT Pro mode, effort Pro, S6\_10\_PAPERS, 18 November 2026: targeted proof check. Continuation metadata events: GPT-6 Astra (OpenAI), ChatGPT Pro mode, effort Pro, S6\_10\_PAPERS continuation, 19 November 2026: used and rechecked in this run; status-vocabulary reconciliation; metadata reconciliation only; no new mathematical check. \emph{Pre-continuation recorded proof scope.} GPT-6 Astra (OpenAI, ChatGPT Pro mode), effort Pro, 2026-11-17 UTC; Claude Fable 5.2 (Anthropic), effort high, 16 November 2026; also checked in this writing by GPT-6 Astra (OpenAI), ChatGPT Pro mode, effort Pro, 18 November 2026}{the hypotheses, notation and definitions stated here; \cref{maps:src:D.1.6}}


The pullback
\[
 \pi^*:H^4(X_p;\mathbb Z)\longrightarrow
 H^4(\widetilde X_p;\mathbb Z)
\]
is zero. In particular
$c_2(\widetilde X_p)=p_1(\widetilde X_p)=0$ integrally.
\end{theorem}


\begin{proof}
\begin{enumerate}[label=\arabic*.,leftmargin=*]
\item We apply the external-curve calculation in \ref{maps:src:D.1.6} first downstairs.
The two mixed regular-fibre tori specialize to one external curve
and a signed sum of the two; their values under the primitive
functional $\Phi=(4,2,3,2)$ are $0$ and $2$.
The first external curve therefore has zero class in
$H_2(X_p)=\mathbb Z/d$, and the other, which we call $C$, has
class $\pm2$ and generates because $d$ is odd.
The covering has exactly $d$ lifts $C_j$ of $C$, each mapped
orientation-preservingly to $C$, and the transpose relations of
$B(t)$ identify all $[C_j]$ in $H_2(\widetilde W_p)$.
Their image is torsion of order $d$ in $H_2(\widetilde X_p)$
by the same mixed-torus calculation. Homology transfer consequently
takes the generator $[C]$ to
\[
 \operatorname{tr}[C]=\sum_j[C_j]=d[C_0]=0.
\]
For a finite orientation-preserving cover, Poincare duality
intertwines cohomological pullback with homological transfer;
this follows by lifting a transverse cycle or from the cap-product
chain identity with the lifted fundamental class. Thus $\pi^*$
vanishes in degree four. Naturality of Chern and Pontryagin
classes, applied to $T\widetilde X_p\simeq\pi^*TX_p$, proves
both assertions. For $d=1$ the target degree-four group already
vanishes.
\end{enumerate}
\end{proof}

\begin{theorem}[the finite-coefficient cubic]\label{maps:src:D.1.9}
\StatementRecord{maps:src:D.1.9}{proved}{Historical events: Claude Fable 5.2 (Anthropic), effort high, 16 November 2026: carried as background without fault found. GPT-6 Astra (OpenAI), ChatGPT Pro mode, effort Pro, S6\_09\_PAPERS, 17 November 2026: targeted proof check. Also checked in this writing. Pre-continuation writing events: GPT-6 Astra (OpenAI), ChatGPT Pro mode, effort Pro, S6\_10\_PAPERS, 18 November 2026: targeted proof check. Continuation metadata events: GPT-6 Astra (OpenAI), ChatGPT Pro mode, effort Pro, S6\_10\_PAPERS continuation, 19 November 2026: used and rechecked in this run; status-vocabulary reconciliation; metadata reconciliation only; no new mathematical check. \emph{Pre-continuation recorded proof scope.} GPT-6 Astra (OpenAI, ChatGPT Pro mode), effort Pro, 2026-11-17 UTC; Claude Fable 5.2 (Anthropic), effort high, 16 November 2026; also checked in this writing by GPT-6 Astra (OpenAI), ChatGPT Pro mode, effort Pro, 18 November 2026}{the hypotheses, notation and definitions stated here; \cref{maps:src:D.1.6}}


Let $d>1$ and let $C_\delta$ be the complex-oriented fixed curve of
the $\delta$ action downstairs. There is a class
$v\in H^2(X_p;\mathbb Z/d)$ with $v(C_\delta)=1$.
For $u=\pi^*v$ and the reductions of the classes $x_j$ in \ref{maps:src:D.1.6},
the full degree-two cubic, evaluated on the complex fundamental
class of $\widetilde X_p$, is
\[
 \begin{aligned}
 \langle x_i x_j x_k,[\widetilde X_p]\rangle&=0,
 &\langle u x_i x_j,[\widetilde X_p]\rangle&=C_{ij},\\
 \langle u^2x_i,[\widetilde X_p]\rangle&=-12,
 &\langle u^3,[\widetilde X_p]\rangle&=0
 \quad\text{in }\mathbb Z/d.
 \end{aligned}
\]
where $C_{ii}=-2$, $C_{i,i\pm1}=1$ and all other entries are zero.
The $d-1$ classes $x_0,\ldots,x_{d-2}$ together with $u$ form a
$(\mathbb Z/d)$-basis of $H^2(\widetilde X_p;\mathbb Z/d)$.
In particular $u$ is deck invariant. No nontrivial extension of the
deck representation is needed to compute this cubic.
\end{theorem}


\begin{proof}
\begin{enumerate}[label=\arabic*.,leftmargin=*]
\item The fixed curve generates $H_2(X_p)=\mathbb Z/d$ by \ref{maps:src:D.1.6}'s
downstairs calculation. The universal coefficient epimorphism
\[
 H^2(X_p;\mathbb Z/d)\longrightarrow
 \operatorname{Hom}(H_2(X_p),\mathbb Z/d)
\]
therefore supplies $v$. Each lift of $C_\delta$ maps to it with
degree one, so $\langle u\kappa,[\widetilde X_p]\rangle=1$.
The diagonal-direction external curve has zero homology class,
and the ruling is its sum with the fixed curve. Thus the assertion
about this pairing uses the complex ruling orientation, not a
choice of a torsion generator up to sign. Universal coefficients
upstairs identify degree-two cohomology with
$\operatorname{Hom}(\mathbb Z^{d-1}\oplus\mathbb Z/d,\mathbb Z/d)$.
Evaluation on its torsion summand shows that $u$ completes the
reductions of the free integral basis to a basis. We choose free
lifts in $H_2$ annihilated by $u$, subtracting multiples of its
torsion generator if necessary; this fixes the splitting used in
the classifying coordinates below. The first two
cubic formulas follow from \ref{maps:src:D.1.6}.

\item We compute the remaining component restriction with its orientation.
In the toric model with lattice coordinates $(r,s)$ on its height-one
vertices, take the integral support function
\[
 \psi(r,s)=-\frac{r(r-1)}2+3s(s-1),
\]
interpolated linearly over each unimodular cone. It defines a toric
line bundle. Its changes under the two period translations are
\[
 \psi(r+1,s)-\psi(r,s)=-r,
 \qquad
 \psi(r,s-1)-\psi(r,s)=-6s+6.
\]
These affine functions give commuting character multipliers
$x_1^{-1}$ and $x_2^{-6}t^6$, so the bundle descends over the
period lattice generated by $e_1$ and $-e_2$. The holomorphic-unit
corrections in the actual cusp action can be deformed to one in a
sufficiently small tube, as in the checked cusp isotopy of AC26,
Section~7.2. This transports the integral topological class and
does not change any of the following pairings.

\item For a wall with endpoints $v,v'$ and the two adjacent third vertices
$w,w'$, the relation of height-one rays is $w+w'=v+v'$.
The two toric trivializations on the curve differ by a power of its
coordinate of exponent
$\psi(v)+\psi(v')-\psi(w)-\psi(w')$. This is its degree for the
complex orientation. Using respectively
\[
 \begin{array}{c|c|c}
 (v,v')&(w,w')&\deg\\\hline
 ((0,0),(1,0))&((0,1),(1,-1))&-6\\
 ((1,0),(0,1))&((0,0),(1,1))&0\\
 ((0,0),(0,1))&((1,0),(-1,1))&1
 \end{array}
\]
gives the internal, diagonal and $\delta$-fixed curve degrees
$(-6,0,1)$. The regular-fibre first Chern class of this bundle is
$-(u\wedge w+6\gamma\wedge\delta)$: the winding of the first
multiplier on the ordered $(\widehat u,\widehat w)$ rectangle is
$-1$, and that of the second on
$(\widehat\gamma,\widehat\delta)$ is $-6$.
There are no other windings. This also fixes the relative signs
that AC26 (7.9), when used with separately oriented curve generators,
leaves implicit.

\item The complement class dual to $\Phi=(4,2,3,2)$ restricts to twice
$q=u\wedge w+6\gamma\wedge\delta$. The specialization from the
regular fibre to the downstairs cusp is onto in degree two
(AC26, Proposition~7.12, whose integral annulus proof is used here).
Explicitly, for any of the three complex curves $C$, choose
$y\in H_2(F;\mathbb Z)$ specializing to $[C]$. The descended
toric line bundle evaluates on $y$ as $-q(y)=\deg_L(C)$,
while its image in the exterior cyclic generator $g$ has
coefficient $2q(y)$. Thus, integrally modulo $d$,
\[
 [C]=-2\deg_L(C)\,g\quad\text{in }H_2(X_p;\mathbb Z).
\]
Since $2$ is a unit modulo $d$, the images of the three
complex curves have precisely the ratios $-6:0:1$ after
normalizing by the complex $\delta$-fixed curve. On the normalized $dP_6$ the
internal curve is $E_1$, the zero external curve is $E_2$, and the
fixed external curve is $E_3$. The opposite-pair relation is
$H=E_1+E_2+E_3$ after mapping to the cusp. Thus the restriction of
$u$ to every normalized component is
\[
 \eta=-5H+6E_1-E_3=-5(H-E_1)+(E_1-E_3)
 \quad\text{in }H^2(dP_6;\mathbb Z/d).
\]
Indeed its evaluations on $E_1,E_2,E_3$ are $-6,0,1$, and on $H$
the evaluation is $-5$. Hence $\eta^2=25-36-1=-12$.
The projection formula for the component fundamental cycle gives
$\langle u^2x_i,[\widetilde X_p]\rangle=-12$, also when that
component has a self-identification: its normalization pushes its
fundamental class with degree one.

\item Finally $u$ is an actual pullback from downstairs. Therefore
\[
 \langle u^3,[\widetilde X_p]\rangle
 =\langle v^3,\pi_*[\widetilde X_p]\rangle
 =d\langle v^3,[X_p]\rangle=0\quad\text{in }\mathbb Z/d.
\]
This evaluates the otherwise undetermined finite-coefficient cube.
The exact support-function, wall-degree and intersection computations
are executed in \texttt{code/manifolds\_mu.py}.
\end{enumerate}
\end{proof}

\begin{lemma}[the cubic detects the complete fundamental-class invariant]\label{maps:src:D.1.10}
\StatementRecord{maps:src:D.1.10}{proved}{Historical events: Claude Fable 5.2 (Anthropic), effort high, 16 November 2026: carried as background without fault found. GPT-6 Astra (OpenAI), ChatGPT Pro mode, effort Pro, S6\_09\_PAPERS, 17 November 2026: targeted proof check. Also checked in this writing. Pre-continuation writing events: GPT-6 Astra (OpenAI), ChatGPT Pro mode, effort Pro, S6\_10\_PAPERS, 18 November 2026: targeted proof check. Continuation metadata events: GPT-6 Astra (OpenAI), ChatGPT Pro mode, effort Pro, S6\_10\_PAPERS continuation, 19 November 2026: used and rechecked in this run; status-vocabulary reconciliation; metadata reconciliation only; no new mathematical check. \emph{Pre-continuation recorded proof scope.} GPT-6 Astra (OpenAI, ChatGPT Pro mode), effort Pro, 2026-11-17 UTC; Claude Fable 5.2 (Anthropic), effort high, 16 November 2026; also checked in this writing by GPT-6 Astra (OpenAI), ChatGPT Pro mode, effort Pro, 18 November 2026}{the hypotheses, notation and definitions stated here; \cref{maps:src:D.1.9}}


Let $A=\mathbb Z^r\oplus\mathbb Z/d$ with $\gcd(d,6)=1$.
The integral homology group of $K(A,2)$ in degree six is
\[
 H_6(K(A,2);\mathbb Z)\simeq
 \mathbb Z^{\binom{r+2}{3}}\oplus
 (\mathbb Z/d)^{\binom{r+1}{2}+r+1}.
\]
Its free coordinates are detected by integral degree-two cubics;
its torsion coordinates are detected by the finite-coefficient
monomials $u x_i x_j$ ($i\leq j$), $u^2x_i$ and $u^3$, where
$u$ evaluates to one on the cyclic summand of $A$. Consequently
\ref{maps:src:D.1.9} and the integral cubic determine the complete class
$\mu(\widetilde X_p)$ in Zhubr's theorem, not merely a quotient of it.
\end{lemma}


\begin{proof}
\begin{enumerate}[label=\arabic*.,leftmargin=*]
\item We give the low-degree calculation to exclude invisible torsion.
Put $R=\mathbb Z[1/6]$. In the path-loop Serre spectral sequence of
$K(\mathbb Z,3)$ the fibre is $K(\mathbb Z,2)=\mathbb CP^\infty$.
If $t$ has degree two, its transgression is the degree-three
generator $z$. Multiplicativity gives
$d_3(t^2)=2tz$ and $d_3(t^3)=3t^2z$.
Since two and three are units and the total path space is
contractible, induction in total degree forces the base cohomology
through degree seven to be the exterior algebra on $z$:
there is no surviving fibre class that could kill an additional
base class in degrees four through seven. Also $z^2=0$ because
$z$ is odd and two is a unit. This proves exactly the low-degree
base calculation needed below.

\item The coefficient extension
$0\to\mathbb Z\xrightarrow{d}\mathbb Z\to\mathbb Z/d\to0$
gives the fibration
\[
 K(\mathbb Z,2)\longrightarrow K(\mathbb Z/d,2)
 \longrightarrow K(\mathbb Z,3)
\]
with transgression $d_3(t)=dz$. Its cohomology through degree seven
is therefore computed by $R[t]\otimes\Lambda(z)$ with this
 differential. The groups in degrees three, five and seven are
$R/(d)$, $R/(2d)$ and $R/(3d)$, all $\mathbb Z/d$, and the positive even
groups through degree six vanish. There is no higher differential
in this range, by the base calculation. Universal coefficients,
starting in degree two, give
\[
 H_2=H_4=H_6=\mathbb Z/d,\qquad H_3=H_5=0.
\]
With coefficients $\mathbb Z/d$ the indicated transgressions vanish,
and $t,t^2,t^3$ survive; in particular the cube evaluates as a
generator of the dual of $H_6$. The same reasoning applies after
localization at every prime dividing $d$. At all other primes,
$K(\mathbb Z/d,2)$ is homologically trivial after localization,
as follows from its single nonzero homotopy group and the localized
Hurewicz theorem. This rules out additional two- or three-primary
classes omitted by passing to $R$.

\item Finally we use the product decomposition
\[
 K(A,2)=(\mathbb CP^\infty)^r\times K(\mathbb Z/d,2).
\]
The first factor has torsion-free homology concentrated in even
degrees, so the integral Kunneth formula has no Tor term here.
Its degree-six free part has rank $\binom{r+2}{3}$; tensoring its
degree-four, degree-two and degree-zero groups with the three
cyclic groups just computed gives respectively
$\binom{r+1}{2}$, $r$ and one cyclic summands. The stated products
evaluate on these distinct summands, and each is a perfect
detector in its coordinate. This proves the lemma.
\end{enumerate}
\end{proof}

The complete credited statement and re-derivation are in \cref{maps:src:D.1.11}.
\begin{theorem}[a canonical cyclic action and its quotient]\label{maps:src:D.1.12}
\StatementRecord{maps:src:D.1.12}{proved}{Historical events: Claude Fable 5.2 (Anthropic), effort high, 16 November 2026: carried as background without fault found. GPT-6 Astra (OpenAI), ChatGPT Pro mode, effort Pro, S6\_09\_PAPERS, 17 November 2026: targeted proof check. Also checked in this writing. Pre-continuation writing events: GPT-6 Astra (OpenAI), ChatGPT Pro mode, effort Pro, S6\_10\_PAPERS, 18 November 2026: targeted proof check. Continuation metadata events: GPT-6 Astra (OpenAI), ChatGPT Pro mode, effort Pro, S6\_10\_PAPERS continuation, 19 November 2026: used and rechecked in this run; status-vocabulary reconciliation; metadata reconciliation only; no new mathematical check. \emph{Pre-continuation recorded proof scope.} GPT-6 Astra (OpenAI, ChatGPT Pro mode), effort Pro, 2026-11-17 UTC; Claude Fable 5.2 (Anthropic), effort high, 16 November 2026; also checked in this writing by GPT-6 Astra (OpenAI), ChatGPT Pro mode, effort Pro, 18 November 2026}{the hypotheses, notation and definitions stated here; \cref{maps:src:D.1.7}; \cref{maps:src:D.1.4}; \cref{maps:src:D.1.11}}


There is an explicit free action $\sigma_d$ of $\mathbb Z/d$ on
$M_d$, defined by the affine and toric data below, such that every
admissible multiplicity-$(3,4)$ member with $|p|=d$ is
orientation-preservingly diffeomorphic to $M_d/\langle\sigma_d\rangle$.
All cyclic actions arising from these members are equivariantly
diffeomorphic to this action, allowing an automorphism of the
abstract deck group. Their quotients carry exactly the complex
structures and fixed-$u$ count in \ref{maps:src:D.1.7}.
\end{theorem}


\begin{proof}
\begin{enumerate}[label=\arabic*.,leftmargin=*]
\item Choose the specific original datum
$(\ell_0,\ell_1,\ell_2)=(0,-d,d)$ and scalar period $c=-i$.
It is admissible and has $p=d$. Its cover has lattice $\Lambda_d$
and matrices $A_1^{(d)},A_2^{(d)}$ of \ref{maps:src:D.1.4}. Use the respective
finite invariant translation vectors
\[
 v'_1=(-1,-2d,4d,0),\qquad v'_2=(1,3d,-3d,0),
\]
divided by three and four in the affine finite actions, and use
the cusp period lattice $\langle e_1,d e_2\rangle$. The invariant
meridian value is $p'=1$. This specifies a single smooth model,
independently of any input twist ambiguity; \ref{maps:src:D.1.11} identifies its
underlying manifold with $M_d$.

\item On a regular torus define $\sigma_d$ by translation through the
original $\widehat\gamma$, whose coordinate in $\Lambda_d$ is
$(1/d,0,0,0)$. The matrices $A_j^{(d)}$ change this vector by an
integral vector, so this translation is globally defined on the
regular family and commutes with the finite quotient actions.
It has order $d$. On a finite reduced fibre a fixed point of its
$a$th power would require, in the invariant $\gamma'$ coordinate,
$a/d\equiv\pm k/m_j\pmod1$ for some $k$. Since $\gcd(d,m_j)=1$,
this forces $d\mid a$. On the cusp it is the affine period
translation by $-e_2$, including the prescribed holomorphic unit,
and cycles the $d$ components. No nonzero power below $d$ fixes a
stratum: after subtracting a covering period it would translate a
finite triangle, edge or vertex to itself by a nonzero vector.
Thus it extends over the entire toric tube and the action is free.
The quotient restores the original lattice, twists $(-d,d)$ and
period $c=-i$, so is exactly that fixed representative.

\item \ref{maps:src:D.1.7} gives an oriented diffeomorphism from any other constructed
quotient with the same $d$ to this one. Lift it to universal covers.
Uniqueness of path lifting makes the lift intertwine deck
transformations by the induced automorphism of $\pi_1=\mathbb Z/d$.
This proves the claimed equivariant classification within the
entire specified construction by a direct proof, without assuming
a classification of arbitrary free actions on $M_d$. Different
$d$ are distinguished downstairs by fundamental group and upstairs
by Euler characteristic.

\item The fixed-$u$ count is two for the present multiplicities. In the
coordinates of Viaclovsky's higher-multiplicity sphere family,
the dictionary in \cref{push:src:E.1.1} and the classification in \cref{fam:src:B.3.1} give $2k\ell$ at fixed $u$ for the pair $(3k,4\ell)$. The tangent computation in \cref{push:src:E.2.1} gives $(1,1,1,0)$, and \cref{push:src:E.2.2} determines the Kodaira dimension by the sign of
$k\ell-3k-4\ell$ for those sphere members; at $(k,\ell)=(1,1)$
this specializes to the count two and negative Kodaira dimension.
These sphere results concern $d=1$ and do not assert tangent
cohomology for the torsion covers $M_d$ with $d>1$.
\end{enumerate}
\end{proof}

\subsection{Rational elliptic starting configurations and exact computations}
The complete credited statement and re-derivation are in \cref{maps:src:D.2.1}.
\begin{computed}[as a bounded exact computation]\label{maps:src:D.2.2}
\StatementRecord{maps:src:D.2.2}{proved}{Historical events: Claude Fable 5.2 (Anthropic), effort high, 16 November 2026: carried as background without fault found. Claude Fable 5.2 (Anthropic), effort high, 17 November 2026: carried as background without fault found. GPT-6 Astra (OpenAI), ChatGPT Pro mode, effort Pro, S6\_09\_PAPERS, 17 November 2026: targeted proof check. Also checked in this writing. Pre-continuation writing events: GPT-6 Astra (OpenAI), ChatGPT Pro mode, effort Pro, S6\_10\_PAPERS, 18 November 2026: computation or formal certificate. Continuation metadata events: GPT-6 Astra (OpenAI), ChatGPT Pro mode, effort Pro, S6\_10\_PAPERS continuation, 19 November 2026: used and rechecked in this run; status-vocabulary reconciliation; metadata reconciliation only; no new mathematical check. \emph{Pre-continuation recorded proof scope.} GPT-6 Astra (OpenAI, ChatGPT Pro mode), effort Pro, 2026-11-17 UTC; Claude Fable 5.2, effort high, 2026-11-17; GPT-6 Astra (OpenAI, ChatGPT Pro mode), S6\_06\_BLIND, effort Pro, 2026-11-17; Claude Fable 5.2 (Anthropic), effort high, 16 November 2026; exact archived inputs and replies checked in this writing by GPT-6 Astra (OpenAI), ChatGPT Pro mode, effort Pro, 18 November 2026}{the hypotheses, notation and definitions stated here; \cite{manExplorer26}; \cref{maps:src:D.2.1}; \cref{mapsD:repair:collision-realization}}


Fifty specified inputs were evaluated in the public exact integral-cochain engine
\cite{manExplorer26}. In every configuration of \ref{maps:src:D.2.1} we take
$P=G$, $Q=nG$, one simple zero at the first $I_1$ in the stored
marking (or at the smallest multiplicative fibre if no $I_1$ exists),
and independently take at each additive fibre either its original
filling or the invariant log vector $(0,1/m)$, where $m$ is the
minimal semistable-reduction order. Within this exact slice the
engine recognizes 24 simply connected inputs and 26 non-simply-connected
inputs. Exactly three have the integral cohomology of $S^6$:
OS43 with $III^*+3I_1$ and the order-four modification;
OS43 with $III^*+II+I_1$ and only the $III^*$ modification;
OS49 with $IV^*+III+I_1$ and both modifications. The executable is
\path{code/manifolds_explorer_21.py}.

Every exact input, fundamental-group result, and group in degrees
zero through six is retained in
\texttt{records/manifolds\_explorer\_21.csv}, its JSON companion,
and the individual raw responses. The homology groups are the
cohomology groups in reverse order, by oriented Poincare duality.
The integer operations use Smith forms, not floating-point ranks.
The public implementation checks endpoint groups, Poincare and
torsion duality, and agreement of abelianized fundamental group
with the universal coefficient theorem. Re-running the script
with its recorded inputs produces the table; cached responses
retain their URL and API version~2.

The proposition here concerns the executed computation, with these
bounds. Its interpretation as topology of each smooth model uses
the author's supplied marked local-model database. The public scope
document explicitly leaves detailed geometric validation certificates
in the author's research workspace. We therefore do not promote this
run into an independent proof of all its geometric attachment maps,
nor into the unbounded classification asked for in D.2.

For the archived computation, all fifty exact request payloads and
saved API-v2 replies were regenerated and compared, giving the stated
$24$ simply connected outputs, $26$ other outputs and three outputs
with the integral cohomology of a sphere; the eight specified clutch
comparisons were checked as well.
This is an audit of the retained requests and responses.  The source audit records a timeout of its fresh remote request. In the present writing the fifty saved requests and API-v2 replies, and the eight integer-clutch comparisons, were regenerated and matched exactly by \path{code/III/audit_explorer_archive.py}. This is the stated bounded archive audit.  The independent geometric input for the two
collision profiles is \cref{mapsD:repair:collision-realization}.


\end{computed}


\begin{proposition}[Euler checks and parameter sensitivity]\label{maps:src:D.2.3}
\StatementRecord{maps:src:D.2.3}{proved}{Historical events: Claude Fable 5.2 (Anthropic), effort high, 16 November 2026: carried as background without fault found. GPT-6 Astra (OpenAI), ChatGPT Pro mode, effort Pro, S6\_09\_PAPERS, 17 November 2026: targeted proof check. Also checked in this writing. Pre-continuation writing events: GPT-6 Astra (OpenAI), ChatGPT Pro mode, effort Pro, S6\_10\_PAPERS, 18 November 2026: targeted proof check. Continuation metadata events: GPT-6 Astra (OpenAI), ChatGPT Pro mode, effort Pro, S6\_10\_PAPERS continuation, 19 November 2026: used and rechecked in this run; status-vocabulary reconciliation; metadata reconciliation only; no new mathematical check. \emph{Pre-continuation recorded proof scope.} GPT-6 Astra (OpenAI, ChatGPT Pro mode), effort Pro, 2026-11-17 UTC; Claude Fable 5.2 (Anthropic), effort high, 16 November 2026; also checked in this writing by GPT-6 Astra (OpenAI), ChatGPT Pro mode, effort Pro, 18 November 2026}{the hypotheses, notation and definitions stated here; the explicit constructions and calculations in the proof}


For any of the semistable triangular rank-two Mumford fillings, its
reduced central fibre has Euler number twice the number of vertices
in the dual torus triangulation. A rank-one wheel over an elliptic
curve, a smooth torus, and a free finite quotient of a torus have
Euler number zero. Consequently the total Euler number is the sum
of these rank-two contributions whenever all other fillings are of
these types. The fibre configuration and multiplicities alone do
not specify the input topology: integer clutching must also be kept.
\end{proposition}


\begin{proof}
\begin{enumerate}[label=\arabic*.,leftmargin=*]
\item Inclusion--exclusion over the toric strata counts the zero-dimensional
strata, since every positive-dimensional complex torus has Euler
number zero. These strata correspond to triangles. A torus
triangulation satisfies $3F=2E$ and $V-E+F=0$, hence $F=2V$.
A rank-one wheel is stratified by bundles over elliptic curves, and
each stratum has Euler number zero. A finite torus quotient has Euler
number zero by multiplicativity of Euler number for covers.
The regular locus of the total fibration also has Euler number zero;
Euler additivity gives the claim, applying the localization argument to the stated cells.

\item For the last assertion, even the finite computation needs the
full rational vector, not its class modulo the lattice. A smooth
weight-zero log vector $\theta$ of denominator $m$ has boundary
relation $\mathrm{meridian}^{m}=\mathrm{fibre}^{m\theta}$.
Replacing $\theta$ by $\theta+v$ for an integral $v$ changes that
relation and its marked cohomological boundary map. The executed
semistable examples OS45,47,55,56 have $\pi_1=\mathbb Z$ for the
zero clutch and trivial fundamental group for the recorded primitive
integer clutch, while the fibre multiplicity is one in both cases.
These four exact comparisons are saved individually in
\texttt{records/manifolds\_explorer\_runs/}.
\end{enumerate}
\end{proof}

The exact remaining classification questions and all established subcases are in \cref{maps:src:D.2.4}.
\begin{theorem}[arbitrary multiplicities in the original configuration]\label{maps:src:D.2.5}
\StatementRecord{maps:src:D.2.5}{proved}{Historical events: Claude Fable 5.2 (Anthropic), effort high, 16 November 2026: carried as background without fault found. GPT-6 Astra (OpenAI), ChatGPT Pro mode, effort Pro, S6\_09\_PAPERS, 17 November 2026: targeted proof check. Also checked in this writing. Pre-continuation writing events: GPT-6 Astra (OpenAI), ChatGPT Pro mode, effort Pro, S6\_10\_PAPERS, 18 November 2026: targeted proof check. Continuation metadata events: GPT-6 Astra (OpenAI), ChatGPT Pro mode, effort Pro, S6\_10\_PAPERS continuation, 19 November 2026: used and rechecked in this run; status-vocabulary reconciliation; metadata reconciliation only; no new mathematical check. \emph{Pre-continuation recorded proof scope.} GPT-6 Astra (OpenAI, ChatGPT Pro mode), effort Pro, 2026-11-17 UTC; Claude Fable 5.2 (Anthropic), effort high, 16 November 2026; also checked in this writing by GPT-6 Astra (OpenAI), ChatGPT Pro mode, effort Pro, 18 November 2026}{the hypotheses, notation and definitions stated here; \cite[\S2.1, (2.11)--(2.12)]{manViaclovskyLogs26}; \cite[Propositions 4.1, 4.3, 4.5]{manViaclovskyLogs26}; \cref{maps:src:D.3.2}}


Consider the $IV^*,III,I_1$ configuration with its rank-two Mumford
filling, and the logarithmic data of
\cite[\S2.1, (2.11)--(2.12)]{manViaclovskyLogs26}:
\[
 m_0=3k,\qquad m_1=4\ell,\qquad
 v_i=r_i\delta+n_i u_i,\qquad
 \gcd(n_i,m_i)=1,\qquad k,\ell>0.
\]
Here $r_i,n_i$ are arbitrary integers, not residue classes. Put
$P=m_1n_0+m_0n_1$. There is no upper bound on these parameters.
If $P\ne0$, the resulting smooth compact threefold has
\[
 \pi_1(X)=\mathbb Z/|P|,\qquad
 H_j(X;\mathbb Z)=
 \begin{cases}
 \mathbb Z&j=0,6,\\
 \mathbb Z/|P|&j=1,2,3,4,\\
 0&j=5.
 \end{cases}
\]
If $P=0$, then $\pi_1(X)=\mathbb Z$ and all integral homology
is free, with Betti numbers $(1,1,2,2,2,1,1)$.
In particular, in this admissible family $X$ is simply connected
exactly when $|P|=1$, and then it is diffeomorphic to $S^6$.
Such sphere data exist exactly when $\gcd(3k,4\ell)=1$.
This proves a parametric table for every multiplicity and every
integer twist in this family; it does not impose the determinant-one
hypothesis when computing the other manifolds.
\end{theorem}


\begin{proof}
\begin{enumerate}[label=\arabic*.,leftmargin=*]
\item We give the integral reduction explicitly, including the step that
extends beyond the simply connected hypothesis in
\cite[Propositions 4.1, 4.3, 4.5]{manViaclovskyLogs26}.
In the integral basis $(\delta,a,b,d)$, the three homological
monodromies are
\[
 A_0=\begin{pmatrix}1&-1&0&2\\0&-1&-1&6\\0&1&0&0\\0&0&0&1\end{pmatrix},\quad
 A_1=\begin{pmatrix}1&0&1&-3\\0&0&1&0\\0&-1&0&6\\0&0&0&1\end{pmatrix},\quad
 A_\infty=\begin{pmatrix}1&0&0&1\\0&1&1&0\\0&0&1&0\\0&0&0&1\end{pmatrix}.
\]
They satisfy $A_0A_1A_\infty=I$, $A_0^3=A_1^4=I$,
$u_0=d+2a+2b$, $u_1=d+3a+3b$ and
$\Lambda^{A_i}=\mathbb Z\delta\oplus\mathbb Zu_i$ for $i=0,1$.
The good reduction followed by the degree-$k$ or degree-$\ell$
base change remains smooth. The affine generator has translation
$v_i/m_i$ and rotates the base by $e^{2\pi i/m_i}$.
The invariant covector $d^\vee$ changes by $n_i/m_i$;
$\gcd(n_i,m_i)=1$ therefore proves that every nonidentity power
is free even on the central torus. Quotienting gives a smooth
filling retracting onto a bielliptic surface $F_i$.
The logarithm translation used in the punctured-disc gluing has
normal-circle relation $x_i^{m_i}=v_i$. These are precisely the
charts and gluing of Proposition 2.3 of the cited paper; the
freeness argument uses no relation between the two multiplicities.

\item The Mumford filling kills $\delta,a,x_\infty$ in the fundamental
group. The conjugation relation $x_1ax_1^{-1}=b^{-1}$ then kills
$b$, and $x_1=x_0^{-1}$. Thus the full group is already abelian,
generated by $d$ and $x=x_0$, with relation matrix
\[
 R=\begin{pmatrix}-n_0&m_0\\-n_1&-m_1\end{pmatrix}.
\]
Its entries have greatest common divisor one because
$\gcd(n_0,m_0)=1$, and $\det R=P$. Its Smith form is therefore
$(1,|P|)$ when $P\ne0$, or $(1,0)$ when $P=0$.
This is a computation of the group, not just its abelianization.

\item Let $V$ be the inverse image of the thrice-punctured sphere, and
$M_i=\partial N_i$ the boundary of each filling. All $H_1(M_i)$
are free of rank three: the individual coinvariants give
$b=-a,\delta=-3a$ at $0$, and $b=-a,\delta=-2a$ at $1$;
the cusp has primitive rank-two vanishing lattice.
The common coinvariants give $H_1(V)=\mathbb Z^3$.
Independently of $P$, the finite central surfaces have
\[
 H_1(F_0)=\mathbb Z\langle a,d,g_0\rangle/
   \langle3r_0a-n_0d+m_0g_0\rangle,
\quad
 H_1(F_1)=\mathbb Z\langle a,d,g_1\rangle/
   \langle2r_1a-n_1d+m_1g_1\rangle.
\]
These are free of rank two. Choose
$m_i\alpha_i+n_i\beta_i=1$ and put
$z_i=\alpha_i d+\beta_i g_i$; $(a,z_i)$ is an integral basis.
Also $H_1(F_\infty)=\mathbb Z^2$.
Consequently the degree-one Mayer--Vietoris map is a map between
free rank-nine groups with cokernel $H_1(X)$.
If $P\ne0$ its kernel is zero, even though its cokernel may be
nonzero. If $P=0$ its kernel is free of rank one.
This replaces the simply connected assumption in the cited
Proposition 4.1 by the actual rank computation.

\item The map $H_2(M_i)\to H_2(N_i)$ is onto for every filling.
Indeed its cokernel maps to
$H_2(N_i,M_i)\simeq H^4(F_i)=\mathbb Z$ by intersection
with $F_i$; the multiple $m_iF_i$ is principal, so that
intersection is zero. By the pair sequence, its kernel consists
of the normal-circle tori of classes in $H^3(F_i)$.
The homology Serre sequence over the graph underlying the base gives
\[
 0\longrightarrow (\bigwedge^2\Lambda)_{A_0,A_1}
 \longrightarrow H_2(V)\longrightarrow
 \ker\bigl[(A_0-I)\ (A_1-I)\bigr]\longrightarrow0.
\]
The left group is $\mathbb Z[a\wedge b]$: the exterior relations
give $\delta\wedge a=\delta\wedge b=a\wedge d=b\wedge d=0$
and $\delta\wedge d=-6a\wedge b$.
The invariant form
$\xi=a^\vee\wedge b^\vee+6d^\vee\wedge\delta^\vee$
evaluates to one on $a\wedge b$, proving primitivity.
The right group is free of rank five.
The two cusp normal tori are the sweeps of $\delta,a$.
In the graph kernel they are
$(-\delta,-\delta)$ and $(b,-a)$; quotienting by them leaves
the integral basis $(\delta,0),(u_0,0),(0,u_1)$.
For verification, the equation in the $A_0$ coinvariants forces
$b^\vee(v_1)=3d^\vee(v_1)$; subtracting the appropriate cusp
vector then puts $v_1$ in $\mathbb Zu_1$ and $v_0$ in
$\mathbb Z\delta\oplus\mathbb Zu_0$. Uniqueness follows
by comparing the second components.
Thus, after the cusp relations, the target for the remaining
normal tori is a free rank-four group with basis $(s,t,t_0,t_1)$,
where $s=[a\wedge b]$ and the last three classes are these sweeps.

\item For a loop $\lambda$ in the finite covering torus, the normal
torus cuts the reference fibre in
$\mathcal N_i\lambda$, where
$\mathcal N_i=\sum_{j=0}^{m_i-1}A_i^j$, and its $s$ coordinate
is $-\xi(\lambda,v_i)$. The normal torus of $g_i$ has
fibre cut $v_i$ and $s$ coordinate zero; both follow by applying
the homomorphisms to $m_i g_i=v_i$ and cancelling in free groups.
The norms are
\[
 \mathcal N_0a=-k\delta,\quad \mathcal N_0d=m_0u_0,
 \qquad
 \mathcal N_1a=-2\ell\delta,\quad \mathcal N_1d=m_1u_1.
\]
Moreover $\xi(a,u_0)=2$, $\xi(a,u_1)=3$ and
$\xi(d,\delta)=6$. Since $(0,\delta)=-(\delta,0)$ modulo
the cusp relations, the normal-torus map in the bases
$(a,z_0,a,z_1)$ and $(s,t,t_0,t_1)$ is
\[
 B=\begin{pmatrix}
 -2n_0&-6r_0\alpha_0&-3n_1&-6r_1\alpha_1\\
 -k&r_0\beta_0&2\ell&-r_1\beta_1\\
 0&1&0&0\\
 0&0&0&1
 \end{pmatrix}.
\]
Removing the two unit pivots reduces its cokernel to that of
\[
 C=\begin{pmatrix}-2n_0&-3n_1\\-k&2\ell\end{pmatrix},
 \qquad \det C=-P.
\]
The entries of $C$ have greatest common divisor one: any common
divisor divides $\gcd(k,2n_0)$ and hence divides $2$, whereas
$n_1$ is odd because $\gcd(n_1,4\ell)=1$.
Thus $\operatorname{coker}B=\mathbb Z/|P|$ for $P\ne0$
and $\mathbb Z$ for $P=0$. Mayer--Vietoris gives $H_2(X)$
equal to this cokernel for $P\ne0$; for $P=0$ it is an extension
of $\mathbb Z$ by $\mathbb Z$, hence $\mathbb Z^2$.

\item Every fibre except the rank-two cusp has Euler number zero, so
$\chi(X)=2$. Integral Poincare duality and the universal
coefficient theorem give $H_5\simeq\operatorname{Hom}(H_1,\mathbb Z)$,
$H_4\simeq\operatorname{Hom}(H_2,\mathbb Z)
\oplus\operatorname{Ext}(H_1,\mathbb Z)$, and
$\operatorname{Tor}H_3\simeq\operatorname{Tor}H_2$.
The Euler characteristic determines the remaining rank, giving
all stated groups. For $|P|=1$, sphere recognition follows as in
\ref{maps:src:D.3.2} from Hurewicz, Whitehead and $\Theta_6=0$.
Finally the Bezout equation $m_1n_0+m_0n_1=\pm1$ is solvable
exactly when $\gcd(m_0,m_1)=1$, and it automatically implies both
local admissibility conditions.

\item The unrestricted exact implementation is the script
\texttt{manifolds\_unbounded\_logs.py} in \texttt{code/}. It symbolically checks
the determinant and norm formulas and actually computes the Smith
forms on the stated finite validation grid; its all-parameter
correctness is the integral argument above. For example,
$(m_0,m_1;n_0,n_1)=(9,4;-2,1)$ gives $P=1$ and $S^6$;
$(9,4;1,1)$ gives all four torsion groups $\mathbb Z/13$;
$(12,12;1,-1)$ gives the stated infinite-cyclic case.
\end{enumerate}
\end{proof}

The logarithmic quotient and its marked local filling are those of \cite[Proposition~2.3]{manViaclovskyLogs26}; the original integral homology case is \cite[Proposition~4.1]{manViaclovskyLogs26}. The unrestricted determinant and torsion calculation above retains the complete integer boundary maps.

\begin{theorem}[infinitely many orders for $III^*+3I_1$]\label{maps:src:D.2.6}
\StatementRecord{maps:src:D.2.6}{proved}{Historical events: Claude Fable 5.2 (Anthropic), effort high, 16 November 2026: carried as background without fault found. GPT-6 Astra (OpenAI), ChatGPT Pro mode, effort Pro, S6\_09\_PAPERS, 17 November 2026: targeted proof check. Also checked in this writing. Pre-continuation writing events: GPT-6 Astra (OpenAI), ChatGPT Pro mode, effort Pro, S6\_10\_PAPERS, 18 November 2026: targeted proof check. Continuation metadata events: GPT-6 Astra (OpenAI), ChatGPT Pro mode, effort Pro, S6\_10\_PAPERS continuation, 19 November 2026: used and rechecked in this run; status-vocabulary reconciliation; metadata reconciliation only; no new mathematical check. \emph{Pre-continuation recorded proof scope.} GPT-6 Astra (OpenAI, ChatGPT Pro mode), effort Pro, 2026-11-17 UTC; Claude Fable 5.2 (Anthropic), effort high, 16 November 2026; also checked in this writing by GPT-6 Astra (OpenAI), ChatGPT Pro mode, effort Pro, 18 November 2026}{the hypotheses, notation and definitions stated here; \cite[Proposition 3.1, Lemma 3.2, Proposition 4.2]{manViaclovskyTwo26}; \cref{maps:src:D.2.5}}


Use the corrected $III^*,I_1,I_1,I_1$ construction of \cref{push:src:E.5.2} and \cref{push:src:E.5.3}, with
two primitive rank-one Mumford fillings and one rank-two hexagonal
filling. In the ordinary homology basis $(e_0,e_1,e_2,e_3)$ dual
to $(\psi,u,w,\delta)$, put $\ell=e_0-e_1-e_2$.
At $III^*$ one may make a logarithmic transformation with
\[
 m=4k,\qquad v=n\ell+r e_3,\qquad
 k>0,\quad n,r\in\mathbb Z,\quad \gcd(n,4k)=1.
\]
For every such datum the resulting compact complex threefold has
\[
 \pi_1(X)=\mathbb Z/|n|,\qquad
 H_0=H_6=\mathbb Z,\quad H_5=0,\quad
 H_1=H_2=H_3=H_4=\mathbb Z/|n|.
\]
Consequently $n=1$ or $n=-1$ gives a complex structure on $S^6$
for every $k>0$, with no coprimality restriction on $k$.
Thus this configuration answers Viaclovsky's Question 4
affirmatively. This assertion concerns the displayed analytic
fillings and integer twists, not all modifications of that surface.
\end{theorem}


\begin{proof}
\begin{enumerate}[label=\arabic*.,leftmargin=*]
\item The geometric input is the construction in \cref{push:src:E.5.2} and \cref{push:src:E.5.3}, with the original geometric statements in
\cite[Proposition 3.1, Lemma 3.2, Proposition 4.2]{manViaclovskyTwo26}.
The order-four good-reduction family is smooth and proper and its
punctured comparison preserves the distinguished section.
Pull it back by a degree-$k$ map of the covering disc, so that the
original base coordinate is $s^{4k}$. The linear deck action still
has order four on the torus lattice. It fixes $\ell$ and $e_3$.
Translation by $v/(4k)$, combined with rotation of $s$, defines
an order-$4k$ action. On the central torus, the invariant map
$\psi:T^4\to\mathbb R/\mathbb Z$ changes by $n/(4k)$;
coprimality proves freeness. The punctured comparison is changed
by the holomorphic translation
$[\log(s)v(s)/(2\pi i)]$, exactly as in \ref{maps:src:D.2.5}.
The quotient is therefore smooth and glues to the unchanged
three Mumford pieces. Properness follows from properness of the
four fillings and compactness of the base. Its meridian relation
is $g^m=v$. This constructs the claimed analytic threefold for
every displayed parameter, rather than extrapolating a finite table.

\item Here is the integral topology. The cohomological matrices, in the
stated marking, are
\[
 \begin{aligned}
 T_0&=\begin{pmatrix}1&0&-2&-1\\0&0&-1&0\\0&1&0&-1\\0&0&0&1\end{pmatrix},
 &T_1&=\begin{pmatrix}1&0&2&2\\0&1&1&1\\0&0&1&0\\0&0&0&1\end{pmatrix},\\
 T_2&=\begin{pmatrix}1&0&0&0\\0&1&0&0\\0&-1&1&0\\0&0&0&1\end{pmatrix},
 &T_3&=\begin{pmatrix}1&0&0&-1\\0&1&1&0\\0&0&1&0\\0&0&0&1\end{pmatrix}.
 \end{aligned}
\]
Let $A_i=T_i^{-t}$. The primitive vanishing subgroups of the
three Mumford pieces are respectively
$\mathbb Z(e_2+e_3)$, $\mathbb Ze_1$, and
$\mathbb Ze_2\oplus\mathbb Ze_3$.
They kill $e_1,e_2,e_3$. Their meridians are killed by the
extending section, and the sphere relation then kills the finite
meridian as well. The finite relation $g^m=v$ leaves exactly
$ne_0=0$. All boundary fundamental-group maps are onto, so
no generators have been omitted. This proves $\pi_1(X)=\mathbb Z/|n|$.

\item We justify the rank-one normal tori used in the degree-two
calculation. In a primitive rank-one Mumford filling, the
normalization is a flat $\mathbb P^1$-bundle over an elliptic
curve, with its two sections identified by a translation.
Topologically trivialize the flat line bundle and isotope the
translation to the identity, keeping the extending section fixed.
The collar pair is then the product of the primitive nodal
elliptic filling with $T^2$; its local equation is $xy=q$.
It follows that $H^3$ of its central fibre is free of rank three.
The normal tori are the sweeps of all three integral invariant
loops: the vanishing circle, whose rotation extends over $xy=q$,
and the two loops of the remaining elliptic factor. They are
boundaries of relative three-chains obtained from the section
disc and those rotations or translations. The unit-factor
deformation leaves the marked boundary and section comparison
unchanged up to isotopy. For the rank-two filling, the same
argument gives the two invariant angular loops and $H^3=\mathbb Z^2$.
Thus these are integral normal-torus classes; retaining only the
vanishing loop in a rank-one piece would omit two relations.

\item For completeness we give primitive-minor certificates for the
global lattice calculation. Set
\[
 D=\bigl[(A_0-I)\ (A_1-I)\ (A_2-I)\bigr],
 \qquad
 j_3(x)=-(A_1A_2x,A_2x,x).
\]
For $i=0,1,2$ let $j_i(x)$ place $x$ in block $i$ of
$\Lambda^3$. The fixed lattices at the three Mumford points have
the respective integral bases
\[
 (-e_0+2e_1,e_2,e_3),\qquad (e_0,e_1,e_3),\qquad(e_2,e_3).
\]
Form a $12\times9$ matrix $K$ whose first eight columns are
their images under $j_1,j_2,j_3$, in that order, and whose ninth
column is $j_0(\ell)$. Direct multiplication gives $DK=0$ and
$\operatorname{rank}D=3$. The minor of $K$ using rows
\[
 (1,2,4,5,7,8,9,10,12)
\]
is $-1$, with rows numbered from one. Hence its columns form
an integral basis of $\ker D$. Moreover
$j_0(e_3)=-K_3-K_6-K_8$, so the finite $e_3$ sweep vanishes
after the Mumford relations.

\item In the exterior basis
$(e_0e_1,e_0e_2,e_0e_3,e_1e_2,e_1e_3,e_2e_3)$,
the matrix $[\bigwedge^2A_i-I]_{i=0}^3$ has primitive rank five:
rows $(1,2,3,5,6)$ and columns $(1,2,4,5,7)$ have minor $-1$.
Its cokernel is detected by
$\xi=u\wedge w+2\psi\wedge\delta$, whose value on $e_1e_2$
is one. Thus the torus-bundle Serre sequence gives
$H_2(V)=\mathbb Z^{10}$, generated by this fibre class and
the nine sweep classes defining $K$. Each sweep has zero
integral against the flat form $\xi$; these classes consequently
fix an integral splitting, not only a rational one.
The eight Mumford normal tori span a primitive subgroup.
Their domain has rank eight, and the preceding local construction
realizes all eight columns, so their image is exactly that subgroup.
The quotient is
\[
 Q=\mathbb Z\langle s,t\rangle,
 \qquad s=[e_1\wedge e_2],\quad t=[T_{\ell,0}].
\]

\item Abelianization of the finite fibre group gives
$e_2=-e_1$, $e_3=2e_1$ and
$mg=ne_0+2r e_1$. Since $\gcd(n,m)=1$, it is free of rank
two with basis $e_1,z$, where
$z=\alpha e_0+\beta g$ and $m\alpha+n\beta=1$.
The normal-circle formulas proved in \ref{maps:src:D.2.5} apply. Here
\[
 \mathcal N_me_1=2k e_3,\qquad
 \mathcal N_me_0=m\ell,\qquad
 \xi(e_1,v)=-n,\qquad \xi(e_0,v)=2r.
\]
The first norm dies in $Q$. The normal tori of $e_1,e_0,g$
therefore have coordinates $(n,0)$, $(-2r,m)$ and $(0,n)$.
In the basis $e_1,z$, the remaining attachment matrix is
\[
 \begin{pmatrix}n&-2r\alpha\\0&1\end{pmatrix}.
\]
It has Smith form $(1,|n|)$ for every $k$ and $r$.

\item To check that this cokernel is actually $H_2(X)$, the boundary
$H_1$ groups have ranks $3,4,4,3$ and are free; the target
of their Mayer--Vietoris map has rank
$4+2+3+3+2=14$ and is also free. Its cokernel is the finite
group $H_1(X)=\mathbb Z/|n|$, so its kernel is zero.
Every central fibre is irreducible and a multiple of it is
principal; the pair-intersection argument in \ref{maps:src:D.2.5} makes each
$H_2(M_i)\to H_2(N_i)$ onto. Thus $H_2(X)$ is precisely
the displayed normal-torus cokernel. Euler localization gives
$\chi(X)=2$, and integral duality gives all remaining groups
as in \ref{maps:src:D.2.5}. Sphere recognition for $|n|=1$ is the same argument.

\item The script \texttt{code/manifolds\_iii\_star.py} executes all
stated lattice certificates, the norm identities, and 312 actual
$10\times10$ Smith forms for $1\leq k\leq12$,
$-15\leq n\leq15$, $\gcd(n,4k)=1$, and $r=0,3$.
The proof is independent of that finite validation bound.
\end{enumerate}
\end{proof}

\begin{theorem}[all orders $3k$ for Wang's primitive $I_2$ filling]\label{maps:src:D.2.7}
\StatementRecord{maps:src:D.2.7}{proved}{Historical events: Claude Fable 5.2 (Anthropic), effort high, 16 November 2026: carried as background without fault found. GPT-6 Astra (OpenAI), ChatGPT Pro mode, effort Pro, S6\_09\_PAPERS, 17 November 2026: targeted proof check. Also checked in this writing. Pre-continuation writing events: GPT-6 Astra (OpenAI), ChatGPT Pro mode, effort Pro, S6\_10\_PAPERS, 18 November 2026: targeted proof check. Continuation metadata events: GPT-6 Astra (OpenAI), ChatGPT Pro mode, effort Pro, S6\_10\_PAPERS continuation, 19 November 2026: used and rechecked in this run; status-vocabulary reconciliation; metadata reconciliation only; no new mathematical check. \emph{Pre-continuation recorded proof scope.} GPT-6 Astra (OpenAI, ChatGPT Pro mode), effort Pro, 2026-11-17 UTC; Claude Fable 5.2 (Anthropic), effort high, 16 November 2026; also checked in this writing by GPT-6 Astra (OpenAI), ChatGPT Pro mode, effort Pro, 18 November 2026}{the hypotheses, notation and definitions stated here; \cite{manWang26}; \cref{maps:src:D.2.2}; \cref{maps:src:D.2.6}; \cref{maps:src:D.2.5}}


Use the rational elliptic surface
\[
 y^2=x^3-3(1+t)x+(2+3t+2t^2),
 \qquad P=(1,\sqrt2\,t),\qquad Q=6P.
\]
Its fibres are $IV^*,I_2,I_1,I_1$, and the degree-one
linearization vanishes at the second $I_1$ point. Use Wang's
primitive one-component rank-one Mumford model at $I_2$, the
rank-one model at the other $I_1$, and the hexagonal rank-two
model at the zero of the linearization. In the homology marking
dual to $(\psi,u,w,\delta)$ put
\[
 \ell=e_0-2e_1-4e_2,\qquad
 m=3k,\qquad v=n\ell+r e_3,
 \qquad k>0,\quad \gcd(n,3k)=1.
\]
The order-$m$ logarithmic quotient at $IV^*$ is smooth for every
integer $r$, and the resulting threefold satisfies
\[
 \begin{gathered}
 \pi_1(X)=\mathbb Z/|n|,\qquad H_0=H_6=\mathbb Z,
 \qquad H_5=0,\\
 H_1=H_4=\mathbb Z/|n|,\qquad
 H_2=H_3=\mathbb Z/(2|n|).
 \end{gathered}
\]
In particular every $k>0$, with $n=\pm1$, gives a simply
connected complex rational homology six-sphere with two-torsion.
The case $(k,n,r)=(1,1,0)$ independently reproduces Wang's
homology theorem. The other orders are included by the proof,
without a finite parameter bound.
\end{theorem}


\begin{proof}
\begin{enumerate}[label=\arabic*.,leftmargin=*]
\item The source is \cite{manWang26}: Propositions 2.1 and
2.7--2.10 verify the surface and local models, and
Sections 4--5 fix their markings. In particular the $I_2$ period
exponent is $(2,1)$, not $(2,0)$. The unimodular change
$(x,y)=(v,v^2/z)$ makes it $(1,0)$ and gives a single
elliptic ruled normalization. This is the geometric input
missing from the zero-weight ``original $I_2$'' API run in \ref{maps:src:D.2.2}.
It is a different compactification, so the two computations
must not be equated.

\item Pull the smooth cubic good-reduction family back by degree $k$.
The order-three linear monodromy fixes $\ell,e_3$.
Combine base rotation of order $3k$ with translation by
$(n\ell+r e_3)/(3k)$, and use the section-preserving punctured
comparison and logarithmic clutching from \ref{maps:src:D.2.6}. The map $\psi$
changes by $n/(3k)$, so the action is free. This supplies the
analytic construction, properness and the marked relation
$g^{3k}=v$ for all the asserted data.

\item For the integral computation, use the fibre order
$(IV^*,I_2,I_1,I_1)$ and cohomological matrices
\[
 \begin{aligned}
 T_0&=\begin{pmatrix}1&0&-6&-2\\0&-1&-1&1\\0&1&0&-1\\0&0&0&1\end{pmatrix},
 &T_1&=\begin{pmatrix}1&0&6&3\\0&1&2&1\\0&0&1&0\\0&0&0&1\end{pmatrix},\\
 T_2&=\begin{pmatrix}1&0&0&0\\0&2&1&0\\0&-1&0&0\\0&0&0&1\end{pmatrix},
 &T_3&=\begin{pmatrix}1&0&0&-1\\0&1&1&0\\0&0&1&0\\0&0&0&1\end{pmatrix}.
 \end{aligned}
\]
Set $A_i=T_i^{-t}$. Their product is the identity and
$A_0^3=I$. The first two primitive vanishing directions are
$2e_2+e_3$ and $e_1+e_2$; the last filling kills $e_2,e_3$.
These generate $\langle e_1,e_2,e_3\rangle$. As in \ref{maps:src:D.2.6},
the extending sections kill all three Mumford meridians and
the sphere relation kills the fourth. The remaining finite
relation is $ne_0=0$, which proves the actual fundamental-group
assertion. Before the finite filling, the same calculation
gives $\pi_1(X_D)=\mathbb Z\langle e_0\rangle$ for the
inverse image $X_D$ of a disc containing the three Mumford points.

\item Use $D$ and the peripheral maps $j_i$ from \ref{maps:src:D.2.6}. The three
Mumford invariant lattices have bases
\[
 (-e_0+3e_1,e_2,e_3),\qquad
 (e_0,e_1+e_2,e_3),\qquad(e_2,e_3).
\]
Let $K$ consist of their eight peripheral columns, followed by
$j_0(\ell)$. The submatrix on rows
\[
 (1,2,4,5,7,8,9,10,12)
\]
has determinant $-1$. Since $DK=0$ and $\operatorname{rank}D=3$, these columns
are a basis of the entire integral kernel. Also
$j_0(e_3)=-K_3-K_6-K_8$.
For the exterior matrix $[\bigwedge^2 A_i-I]$, rows
$(1,2,3,5,6)$ and columns $(1,2,4,6,7)$ have determinant $1$.
Its primitive cokernel is $\mathbb Z$, detected by
\[
 \xi=u\wedge w+6\psi\wedge\delta.
\]
The normal-circle argument of \ref{maps:src:D.2.6} applies to these actual
primitive Mumford models and accounts for all eight columns.
Consequently
\[
 H_2(X_D)=Q=\mathbb Z\langle s,t\rangle,
 \quad s=[e_1\wedge e_2],\quad t=[T_{\ell,0}].
\]
Here the degree-one Mayer--Vietoris map has free domain of rank
$11$, free target of rank $12$, and cokernel $\mathbb Z$;
its kernel is zero. The stated generators therefore also
describe $H_2(X_D)$, rather than only a quotient of it.
The flat form $\xi$ has values $(1,0)$ on these generators.

\item In the finite linear coinvariants,
$e_1=-2e_2$ and $e_3=-3e_2$. Hence its reduced fibre has
\[
 H_1(F)=
 \mathbb Z\langle e_0,e_2,g\rangle/
 \langle mg-ne_0+3r e_2\rangle.
\]
This is free, with basis $e_2,z$, where
$z=\alpha e_0+\beta g$ and $m\alpha+n\beta=1$.
The exact norm and contraction identities are
\[
 \mathcal N_m e_2=-k e_3,\qquad
 \mathcal N_m e_0=m\ell,\qquad
 \xi(e_2,v)=2n,\qquad \xi(e_0,v)=6r.
\]
The first norm dies in $Q$. Thus the normal tori of
$e_2,e_0,g$ have coordinates $(-2n,0)$, $(-6r,m)$,
$(0,n)$. In the free basis $e_2,z$ the remaining matrix is
\[
 B_{k,n,r}=\begin{pmatrix}-2n&-6r\alpha\\0&1\end{pmatrix}.
\]
Its Smith form is $(1,2|n|)$.

\item All four boundary $H_1$ groups are free of total rank $14$;
the target has the same rank and cokernel $\mathbb Z/|n|$.
The degree-one kernel is zero. Each central fibre is irreducible
and a multiple is principal, so the boundary-to-filling map
in degree two is onto, by \ref{maps:src:D.2.5}. Therefore
$H_2(X)=\operatorname{coker}B_{k,n,r}$.
Euler localization gives $\chi(X)=2$; rational duality gives
$b_3=0$. Integral duality and the universal coefficient theorem
then give exactly the groups in the statement.

\item The program \texttt{code/manifolds\_wang.py} executes both
primitive-minor certificates, the norm identities and 336 full
$10\times10$ Smith forms with $1\leq k\leq12$,
$-15\leq n\leq15$, $\gcd(n,3k)=1$, and $r=0,3$.
Its input matrices and outputs are retained. The symbolic
reduction proves the formula for arbitrary displayed parameters.
\end{enumerate}
\end{proof}

The order-three geometric model and its original surgery calculation are given in \cite[Theorem~7.1 and Appendix~B]{manWang26}. The present comparison retains the primitive $I_2$ period $(2,1)$, which becomes $(1,0)$ under $(x,y)=(v,v^2/z)$; it is not replaced by the zero-weight $I_2$ archival API choice.

\begin{theorem}[the oriented type of all simply connected Wang members]\label{maps:src:D.2.8}
\StatementRecord{maps:src:D.2.8}{proved}{Historical events: Claude Fable 5.2 (Anthropic), effort high, 16 November 2026: carried as background without fault found. GPT-6 Astra (OpenAI), ChatGPT Pro mode, effort Pro, S6\_09\_PAPERS, 17 November 2026: targeted proof check. Also checked in this writing. Pre-continuation writing events: GPT-6 Astra (OpenAI), ChatGPT Pro mode, effort Pro, S6\_10\_PAPERS, 18 November 2026: targeted proof check. Continuation metadata events: GPT-6 Astra (OpenAI), ChatGPT Pro mode, effort Pro, S6\_10\_PAPERS continuation, 19 November 2026: used and rechecked in this run; status-vocabulary reconciliation; metadata reconciliation only; no new mathematical check. \emph{Pre-continuation recorded proof scope.} GPT-6 Astra (OpenAI, ChatGPT Pro mode), effort Pro, 2026-11-17 UTC; Claude Fable 5.2 (Anthropic), effort high, 16 November 2026; also checked in this writing by GPT-6 Astra (OpenAI), ChatGPT Pro mode, effort Pro, 18 November 2026}{the hypotheses, notation and definitions stated here; \cref{maps:src:D.2.7}; \cite[Theorem 7.1 and Appendix B]{manWang26}; \cite[Theorems 3, 2.10, pp.795,803, and \S3.12, p.807]{manZhubr75}}


Every member of \ref{maps:src:D.2.7} with $n=\pm1$, arbitrary $k>0$ and
arbitrary $r\in\mathbb Z$, is orientation-preservingly
diffeomorphic to
\[
 M_{\mathrm{conic}}=
 \bigl(\mathbb {CP}^3\setminus\operatorname{int}
       (S^2\times D^4)\bigr)
 \mathop{\cup}_{S^2\times S^3}(D^3\times S^3),
\]
where the removed sphere is the conic
$[s:t]\mapsto[s^2:st:t^2:0]$ with its oriented normal framing.
For the nonzero $a\in H^2(X;\mathbb F_2)$,
\[
 \langle a^3,[X]_2\rangle=1.
\]
This also supplies an independent check of the smooth
classification asserted for order three in
\cite[Theorem 7.1 and Appendix B]{manWang26}.
\end{theorem}


\begin{proof}
\begin{enumerate}[label=\arabic*.,leftmargin=*]
\item We first state precisely the classification being used, including
its two-primary information. Zhubr's Theorem 3 identifies the
oriented simply connected spin six-manifolds with fixed $H_2=G$
and $b_3=0$ with the corresponding spin-bordism classes over
$K(G,2)$, modulo automorphisms of $G$.
For $G=\mathbb Z/2$ his Theorem 2.10 and the calculation in
\S3.12 give
\[
 \Omega^{\mathrm{Spin}}_6(K(\mathbb Z/2,2))=\mathbb Z/2;
\]
the generator is the image of
$(\mathbb {CP}^3,\mathbb {CP}^3\hookrightarrow
K(\mathbb Z,2))$ under reduction $\mathbb Z\to\mathbb Z/2$
\cite[Theorems 3, 2.10, pp.795,803, and \S3.12, p.807]{manZhubr75}.
The homomorphism evaluating the cube of the universal
degree-two mod-two class is nonzero on this generator, and
therefore is an isomorphism to $\mathbb F_2$.
The automorphism group of $\mathbb Z/2$ is trivial. Hence there
are precisely two such oriented smooth types and their
distinguishing invariant is $\langle a^3,[X]_2\rangle$.
In particular, the vanishing of the ordinary integral $p_1$
does not need to be mistaken for a calculation of every refined
two-primary invariant.

\item The conic has normal bundle $\mathcal O(4)\oplus\mathcal O(2)$.
As an oriented real rank-four bundle on $S^2$ it is trivial:
its class is detected by $w_2$, which is zero. Its framing is
unique up to homotopy since $\pi_2SO(4)=0$.
Surgery preserves simple connectivity and spin, and its integral
homology sequence, with the attaching sphere representing twice
the generator, gives $H_2=H_3=\mathbb Z/2$ and all other reduced
groups zero except $H_6=\mathbb Z$.
The mod-two hyperplane class extends across the surgery trace,
because it evaluates to zero on the conic. Bordism invariance
therefore gives cube one on $M_{\mathrm{conic}}$.
It remains to compute this cube on the entire family.

\item Keep the common piece $X_D$ from \ref{maps:src:D.2.7}. The scalar involution
$\iota(z,v)=(z,-v)$ has exactly two fixed curves there: the
double elliptic curve $E$ of the primitive $I_2$ model, and the
rational curve $C$ corresponding to a vertical edge in the
rank-two triangulation. At $I_2$ the monomial coordinates
$(v,v^2/z)$ have involution $(-x,y)$, so both normal directions
at $E$ have sign $-1$. At the other rank-one point, the induced
translation of the elliptic factor is a nonzero two-torsion
point, with no fixed points. In a rank-two triangle the weights
are two odd normal weights and one even tangential weight.
There is one vertical edge orbit and its two endpoints are the
two distinct triple points, giving the smooth curve $C$.

\item The involution extends over the finite logarithmic quotient,
and introduces no additional fixed points there even when $k$
is even. Indeed, an equality $\iota x=g^j x$ on its central
covering torus would imply $jn/(3k)=0$ in $\mathbb R/\mathbb Z$
after applying $\psi$. Thus $3k\mid j$, since $n=\pm1$.
The equality would then assert that translation by $e_3/2$
fixes a point on a torus, which is impossible. Away from the
central torus the base rotation gives the same conclusion.
This replaces the odd-order shortcut in Wang's order-three
proof and establishes the fixed-set statement for every $k,r$.

\item We give the integral line-bundle calculation on the unchanged
piece $X_D$. In multiplicative coordinates let $h_1,h_2$ be
the elliptic period and the representative of $P$, and let $R$
be the multiplier of the contraction $F$. On an auxiliary line
with coordinate $\zeta$, use
\[
 \begin{aligned}
 h(z,v,\zeta)&=(h_1z,h_2v,-z^{-1}\zeta),\\
 F(z,v,\zeta)&=(h_2^{-6}z,Rv,h_2^{-3}R^3v^6\zeta).
 \end{aligned}
\]
The two lifts commute and their alternating Chern form is
$\xi=uw+6\psi\delta$. Their semicharacter modulo two is
$u+w+uw$, with zero value on the radical spanned by $e_0,e_3$;
direct substitution shows that all four monodromies preserve it.
The symmetry is checked, not inferred from vertical curvature.
Writing $B=h_2^{-3}R^3$, the lift
\[
 J(z,v,\zeta)=(z^{-1},v^{-1},-z^{-1}\zeta)
\]
satisfies $J^2=1$, $JhJ=h^{-1}$ and $JFJ=F^{-1}$;
the last identity is exactly $B^2=h_2^{-6}R^6$.
Multiplying $J$ by $-1$ normalizes it on the rigidifying section.
Thus these are symmetric rigidified Appell--Humbert data.
The displayed semicharacter is unchanged on overlap markings,
so the normalized fibrewise identifications patch.

\item Compare with the prequantum model on the flat real universal
cover, whose connection potential is $\pi i\xi(x,dx)$ and
whose lattice multipliers have that same semicharacter.
It is preserved by the linear monodromies. Its fibrewise
holomorphic structure is the Appell--Humbert structure for $\xi$.
Both families are symmetric and rigidified and have the same
alternating form. Their ratio is therefore a family of flat
line bundles of order dividing two. Its square is canonically
trivial on each fibre after rigidification: the trivializing
holomorphic section is unique with value one at the origin.
Those unique sections patch over the base. Hence the comparison
line is globally two-torsion on the regular family, rather than
merely having trivial fibrewise Chern class. Since $H^2(V;\mathbb Z)$
is free by the preceding integral graph calculation, its Chern
class is zero. The actual holomorphic line consequently has
exactly the prequantum Chern class on $V$, whose values on
fibre and swept cycles are respectively one and zero.
Finally $H_2(V)\to H_2(X_D)$ is onto by \ref{maps:src:D.2.7} and
$H_1(X_D)=\mathbb Z$, so restriction
$H^2(X_D;\mathbb Z)\to H^2(V;\mathbb Z)$ is injective.
The extension $\mathcal L$ below therefore has values $(1,0)$
on $(s,t)$, including the mixed sweep value.

\item These factors extend over the three Mumford pieces. At $I_2$
the three valuations are
\[
 (\operatorname{ord}h_1,\operatorname{ord}h_2,
   \operatorname{ord}R)=(2,1,-3).
\]
Replacing $F$ by $Fh^3$ gives line multipliers, up to units,
$-x^{-2}y$ and $y^3$. The support function in the degenerating
coordinate is $-j(j-1)$; all its slopes are even. Hence the
stabilizer character of the resulting line $\mathcal L$ is
trivial along $E$. At the rank-two point the valuations are
$(1,0,1)$ and a support function is
\[
 \varphi(a,b)=-\frac{a(a-1)}2+3b^2.
\]
Its period differences are $-a$ and $6b+3$, as required by
the line multipliers $z^{-1}$ and $q^3v^6$.
Its vertical slopes are odd, giving the nontrivial character
along $C$; the adjacent horizontal slopes differ by one, so
$\deg(\mathcal L|_C)$ is odd. At the remaining rank-one
point the same factors extend over its primitive fan.
Units contribute no winding or parity, and translation to the
section vertex changes the vertical slopes by a multiple of six.

\item We now check the global mod-two extension without assuming that
the degree-$3k$ transfer is invertible. \ref{maps:src:D.2.7} proved
$H_2(X_D)=\mathbb Z(s,t)$ and identified
$c_1(\mathcal L)$ with $(1,0)$ on that basis.
For $n=\pm1$ both first entries of
\[
 B_{k,n,r}=\begin{pmatrix}-2n&-6r\alpha\\0&1\end{pmatrix}
\]
are even. Its cokernel is generated by $s$ of order two, and
its second relation sends $t$ to an even multiple of $s$.
The nonzero global class $a$ therefore restricts exactly to
$\rho_2c_1(\mathcal L)$ on $X_D$. There is no universal
coefficient ambiguity on $X_D$ since $H_1(X_D)=\mathbb Z$.
This is the integral attachment step that makes the argument
valid for all $k$ and every $r$.

\item Finally $H^2(X;\mathbb Z)=0$ implies $c_1(X)=0$, hence spin.
Write $G=\langle\iota\rangle$ and
$\tau\in H^1(BG;\mathbb F_2)$ for the generator.
The class $a$ admits an equivariant lift $\widetilde a$:
the degree-two Borel obstruction with target
$H^2(BG;H^1(X;\mathbb F_2))$ vanishes, and a fixed point
splits the map to $BG$, killing the remaining differential
into $H^3(BG;\mathbb F_2)$.
The circle containing $\iota$ acts trivially on
$\pi_1(X_D)=\mathbb Z$, and a fixed point splits the extension,
so $\pi_1((X_D)_G)=\mathbb Z\times\mathbb Z/2$.
The difference between $\widetilde a$ and the equivariant
class of $\mathcal L$ vanishes in ordinary cohomology and thus
comes from this group's degree-two cohomology, by the
universal-cover low-degree sequence. Every fixed point induces
the same order-two element $(0,1)$ in this group. A single
adjustment by a multiple of $\tau^2$ consequently gives
\[
 \widetilde a|_C=a|_C+\tau^2,
 \qquad \widetilde a|_E=a|_E+\tau\eta,
 \quad\eta\in H^1(E;\mathbb F_2).
\]
The normal bundle of each fixed curve has $w_2=0$, since
$c_1(X)=0$ and $c_1(TC),c_1(TE)$ are even.
Both normal complex directions have sign $-1$, so their
equivariant Euler class is $\tau^4$.
The equivariant Thom isomorphism after inverting $\tau$ gives
\[
 \langle a^3,[X]_2\rangle
 =\int_C\frac{(a|_C+\tau^2)^3}{\tau^4}
  +\int_E\frac{(a|_E+\tau\eta)^3}{\tau^4}
 =\int_C a|_C=1.
\]
For the second equality, squares of degree-one classes on the
oriented elliptic curve vanish, and all other unwanted terms
have the wrong ordinary degree on a curve. The stated
classification now identifies $X$ with $M_{\mathrm{conic}}$.
\end{enumerate}
\end{proof}

\begin{theorem}[all orders $4k$ in the Liu--Rollenske configuration]\label{maps:src:D.2.9}
\StatementRecord{maps:src:D.2.9}{proved}{Historical events: Claude Fable 5.2 (Anthropic), effort high, 16 November 2026: carried as background without fault found. GPT-6 Astra (OpenAI), ChatGPT Pro mode, effort Pro, S6\_09\_PAPERS, 17 November 2026: targeted proof check. Also checked in this writing. Pre-continuation writing events: GPT-6 Astra (OpenAI), ChatGPT Pro mode, effort Pro, S6\_10\_PAPERS, 18 November 2026: targeted proof check. Continuation metadata events: GPT-6 Astra (OpenAI), ChatGPT Pro mode, effort Pro, S6\_10\_PAPERS continuation, 19 November 2026: used and rechecked in this run; status-vocabulary reconciliation; metadata reconciliation only; no new mathematical check. \emph{Pre-continuation recorded proof scope.} GPT-6 Astra (OpenAI, ChatGPT Pro mode), effort Pro, 2026-11-17 UTC; Claude Fable 5.2 (Anthropic), effort high, 16 November 2026; also checked in this writing by GPT-6 Astra (OpenAI), ChatGPT Pro mode, effort Pro, 18 November 2026}{the hypotheses, notation and definitions stated here; \cref{maps:src:D.2.6}; \cite[Lemma 3.10, Proposition 3.11]{manLR26}; \cite[Proposition 3.13]{manLR26}; \cref{maps:src:D.2.5}}


Use the $III^*,II,I_1$ surface and the unchanged type-$II$
filling of Liu--Rollenske, with the degree-one linearization zero
at $I_1$ and the primitive hexagonal filling there.
In their homology basis $(e_1,e_2,e_3,e_4)$ put
\[
 \delta=-e_1+2e_3,\qquad \ell=e_4,\qquad
 m=4k,\qquad v=n\ell+r\delta,
 \quad k>0,\quad\gcd(n,4k)=1.
\]
The logarithmic transformation at $III^*$ exists for every
such integer datum, and its total space satisfies
\[
 \pi_1(X)=\mathbb Z/|n|,\qquad
 H_0=H_6=\mathbb Z,\quad H_5=0,\quad
 H_1=H_2=H_3=H_4=\mathbb Z/|n|.
\]
Thus $n=\pm1$ gives $S^6$ for every $k>0$. In particular,
the Liu--Rollenske configuration also answers Viaclovsky's
Question 4 affirmatively. The type-$II$ fibre is unchanged
in this statement.
\end{theorem}


\begin{proof}
\begin{enumerate}[label=\arabic*.,leftmargin=*]
\item The source's good-reduction construction and section-preserving
comparison at $III^*$ give the analytic model. The degree-$k$
pullback and free affine quotient argument of \ref{maps:src:D.2.6} applies
with invariant covector $e_4^*$, which takes value $n$ on $v$.
It proves existence, smoothness, properness and the relation
$g^m=v$ without restricting $k$.

\item The extra local input, absent from \ref{maps:src:D.2.6}, is the unchanged
type-$II$ filling. By \cite[Lemma 3.10, Proposition 3.11]{manLR26},
its marked neighbourhood is homeomorphic to the type-$II$
elliptic neighbourhood times $T^2$, with the section preserved.
The cuspidal rational curve is homeomorphic to $S^2$.
The section evaluates as one on both this central two-cycle
and the nearby elliptic curve, so the degree-two specialization
is primitive. The product description makes its two normal
tori precisely the meridian swept by each of the two transverse
$T^2$ loops: each bounds the section disc times that loop.
K\"unneth, or the relative sequence of this product pair,
shows that these generate the entire rank-two kernel of
the boundary map in degree two. This justifies the integral
attachment used below without treating the cusp as a normal
crossing curve.

\item Here are all matrices and primitive certificates. The source
uses $T_\infty T_1T_0=I$ in homology. We order its matrices
as $(A_0,A_1,A_2)=(T_0,T_\infty,T_1)$, so their product is $I$:
\[
 \begin{gathered}
 A_0=\begin{pmatrix}-1&-1&-1&0\\2&1&1&0\\0&0&1&0\\0&0&0&1\end{pmatrix},\qquad
 A_1=\begin{pmatrix}1&0&0&-1\\-2&0&-1&1\\2&1&2&1\\0&0&0&1\end{pmatrix},\\
 A_2=\begin{pmatrix}1&1&0&1\\-2&0&-1&-1\\0&-1&1&-1\\0&0&0&1\end{pmatrix}.
 \end{gathered}
\]
These are \cite[Proposition 3.13]{manLR26} in its actual
integral marking. The cusp's primitive vanishing lattice is
$\langle a,\delta\rangle$, where $a=e_1-e_2-e_3$.
Its monodromy-normal closure contains
$A_2a=-e_2$, and hence all of
$\langle e_1,e_2,e_3\rangle$. The two unchanged meridians
are killed by the extending section. The sphere relation kills
the third, and $g^m=v$ then leaves only $ne_4=0$.
Surjectivity of the boundary maps on fundamental groups gives
$\pi_1(X)=\mathbb Z/|n|$.

\item Set $D=[A_0-I\ \ A_1-I]$ and, in $\Lambda^2$, define
\[
 j_0(x)=(x,0),\quad j_1(x)=(0,x),\quad
 j_2(x)=-(A_1x,x).
\]
The cusp fixed basis is $(a,\delta)$ and the type-$II$ fixed
basis is $(\delta,e_4-e_2)$. Form
\[
 K=[j_1(a)\ j_1(\delta)\ j_2(\delta)\
       j_2(e_4-e_2)\ j_0(\ell)].
\]
Then $DK=0$, $\operatorname{rank}D=3$, and the $5\times5$
minor on rows $(1,2,4,5,6)$ is $1$.
Moreover $j_0(\delta)=-K_2-K_3$.
The exterior matrix $[\bigwedge^2A_i-I]_{i=0}^2$ has rank
five and the minor on rows $(1,2,3,4,5)$ and columns
$(2,4,5,6,11)$ is $1$. Its primitive cokernel is detected by
\[
 \omega=2e_1^*e_2^*-e_2^*e_3^*-e_3^*e_4^*.
\]
Thus $H_2(V)$ is free of rank six, split by a fibre class
$s=-e_3e_4$, of $\omega$-value one, and the five zero-value
sweeps $K$. The four normal tori of the two unchanged fillings
are exactly its first four sweep columns, by the preceding
product argument and the rank-two argument in \ref{maps:src:D.2.6}.
After quotienting by them the group is
$Q=\mathbb Z\langle s,t\rangle$, $t=j_0(\ell)$.

\item The finite linear coinvariants satisfy $e_1=e_2=0$, so
$\delta=2e_3$. The finite reduced fibre therefore has free
$H_1$ with basis $e_3,z$, where
\[
 mg=ne_4+2r e_3,\qquad z=\alpha e_4+\beta g,
 \qquad m\alpha+n\beta=1.
\]
Direct calculation gives
\[
 \mathcal N_me_3=2k\delta,\quad
 \mathcal N_me_4=m\ell,\quad
 \omega(e_3,v)=-n,\quad \omega(e_4,v)=2r.
\]
The $\delta$ sweep dies in $Q$, and the normal-circle formula
therefore reduces the last map to
\[
 \begin{pmatrix}n&-2r\alpha\\0&1\end{pmatrix}.
\]
All boundary $H_1$ groups are free of total rank nine, and their
target is free of rank nine with finite cokernel
$\mathbb Z/|n|$; the degree-one kernel is zero.
Every central fibre is irreducible and a multiple is principal,
so the degree-two boundary maps are onto as in \ref{maps:src:D.2.5}.
It follows that $H_2(X)=\mathbb Z/|n|$.
The unchanged type-$II$ piece and the finite quotient have
Euler number zero, whereas the hexagonal fibre has Euler
number two. Duality yields all the other displayed groups.
Sphere recognition when $|n|=1$ is as in \ref{maps:src:D.2.6}.

\item The independent script
\texttt{code/manifolds\_liu\_rollenske.py} checks the two
unit minors, the norm identities, and 312 full $6\times6$
Smith forms for the same $(k,n,r)$ bounds as \ref{maps:src:D.2.6}.
The proof establishes the formula without those bounds.
\end{enumerate}
\end{proof}

The product type-$II$ cap and the original marked torus monodromy are established in \cite[Lemma~3.10 and Propositions~3.11, 3.13]{manLR26}. The proof above gives the full indicated $3k$ family in that marking.

\begin{theorem}[unrestricted orders in both Yang families]\label{maps:src:D.2.10}
\StatementRecord{maps:src:D.2.10}{proved}{Historical events: Claude Fable 5.2 (Anthropic), effort high, 16 November 2026: carried as background without fault found. GPT-6 Astra (OpenAI), ChatGPT Pro mode, effort Pro, S6\_09\_PAPERS, 17 November 2026: targeted proof check. Also checked in this writing. Pre-continuation writing events: GPT-6 Astra (OpenAI), ChatGPT Pro mode, effort Pro, S6\_10\_PAPERS, 18 November 2026: targeted proof check. Continuation metadata events: GPT-6 Astra (OpenAI), ChatGPT Pro mode, effort Pro, S6\_10\_PAPERS continuation, 19 November 2026: used and rechecked in this run; status-vocabulary reconciliation; metadata reconciliation only; no new mathematical check. \emph{Pre-continuation recorded proof scope.} GPT-6 Astra (OpenAI, ChatGPT Pro mode), effort Pro, 2026-11-17 UTC; Claude Fable 5.2 (Anthropic), effort high, 16 November 2026; also checked in this writing by GPT-6 Astra (OpenAI), ChatGPT Pro mode, effort Pro, 18 November 2026}{the hypotheses, notation and definitions stated here; \cite[Lemma 4.1, Lemma 4.3]{manYang26}; \cite[Proposition 4.4]{manYang26}; \cite[Theorem 7.8]{manYang26}; \cite[Lemma 7.10, Appendix B.4]{manYang26}; \cite[Theorem~7.13]{manYang26}}


Fix a positive cusp index $R$ in Yang's Type B or Type C period
family, with its canonical cusp collar. Use the following constants:
\[
 \begin{array}{c|ccc|cc}
 \text{type}&c&d&a=d/c&m_1&m_2\\\hline
 B&2&4&2&4k&6l\\
 C&3&3&1&3k&6l
 \end{array}
 \qquad k,l\geq1.
\]
In the marking $\Lambda=\mathbb Z\langle\gamma,u,w,\delta\rangle$,
put $\ell_1=\gamma-R(u+w)$ and choose the invariant finite
translation vectors
\[
 v_1=n_1\ell_1+r_1\delta,\qquad
 v_2=n_2\gamma+r_2\delta,\qquad
 \gcd(n_j,m_j)=1,
 \qquad n_j,r_j\in\mathbb Z.
\]
These data give compact smooth complex threefolds for every displayed
integer choice. Define
\[
 P=m_2n_1+m_1n_2,\qquad h=\gcd(R,ak,P)>0.
\]
If $P\ne0$, set
$G=\mathbb Z/h\oplus\mathbb Z/(R|P|/h)$, omitting trivial factors.
Then
\[
 \begin{array}{c|ccccccc}
 i&0&1&2&3&4&5&6\\\hline
 H_i(X;\mathbb Z)&\mathbb Z&\mathbb Z/|P|&
 \mathbb Z^{R-1}\oplus G&G&
 \mathbb Z^{R-1}\oplus\mathbb Z/|P|&0&\mathbb Z
 \end{array}
\]
and $\pi_1(X)=\mathbb Z/|P|$.
If $P=0$, necessarily $m_1=m_2$ and $n_2=-n_1$; in this case
$h=\gcd(R,ak)$, $\pi_1(X)=\mathbb Z$, and
\[
 \begin{array}{c|ccccccc}
 i&0&1&2&3&4&5&6\\\hline
 H_i(X;\mathbb Z)&\mathbb Z&\mathbb Z&
 \mathbb Z^{R+1}\oplus\mathbb Z/h&
 \mathbb Z^2\oplus\mathbb Z/h&\mathbb Z^{R+1}&
 \mathbb Z&\mathbb Z.
 \end{array}
\]
In both cases $\chi(X)=2R$; the displayed homology is independent
of $r_1,r_2$. No member is simply connected. More generally, with
these two finite multiplicities, this last obstruction holds for
every free affine choice of the two logarithmic fillings, even if
the sufficient coprimality condition above is dropped.

For $k=l=1$, $(n_1,n_2)=(1,-1)$ and $r_1=r_2=0$, this independently
recovers Yang's original groups: Type B has fundamental group
$\mathbb Z/2$ and middle torsion $\mathbb Z/R\oplus\mathbb Z/2$;
Type C has fundamental group $\mathbb Z/3$ and middle torsion
$\mathbb Z/(3R)$.
\end{theorem}


\begin{proof}
\begin{enumerate}[label=\arabic*.,leftmargin=*]
\item \emph{Analytic construction and exact marking.}
The actual dual deck matrices of the primary source are
\[
 A_1=\begin{pmatrix}
 1&0&0&0\\0&0&1&0\\-cR&-1&2-c&0\\-R&0&-1&1
 \end{pmatrix},\qquad
 A_2=\begin{pmatrix}
 1&0&0&0\\0&1&-1&0\\0&1&0&0\\0&0&0&1
 \end{pmatrix}.
\]
Their orders are $d,6$, and their integral fixed lattices are
$\langle\ell_1,\delta\rangle$ and
$\langle\gamma,\delta\rangle$, respectively.
Pull the good-reduction root discs back by degrees $k,l$.
On the resulting root family the affine action is
\[
 (s,x)\longmapsto(\zeta_{m_j}^{-1}s,A_jx+v_j/m_j).
\]
The holomorphic period formula and the logarithmic collar maps of
\cite[Lemma 4.1, Lemma 4.3]{manYang26} apply unchanged: in complex
coordinates the translation is the holomorphic period section
$\Pi(\zeta_{m_j}^{-1}s)v_j/m_j$.
The invariant functional $\chi=\gamma^*$ changes by $n_j/m_j$.
Applying it to a purported fixed point of a nontrivial power
forces $m_j\mid qn_j$, impossible for $0<q<m_j$.
The actions are therefore free, their quotients are proper, and
the central divisors have multiplicities $m_j$.
The collar isomorphism adds or subtracts
$(\log s)/(2\pi i)\,\Pi(s)v_j$; changing the logarithm adds an
integral period. This proves existence with the stated marking
and leaves the toroidal cusp unchanged. Compactness follows from
the same three compact tubes and compact base core as in the
source's global gluing argument.

\item Write $g_j$ for the transferred clockwise meridian. The finite
boundary map imposes exactly $g_j^{m_j}=v_j$, by
\cite[Proposition 4.4]{manYang26}; this can also be checked from
the path $q\mapsto qv_j/m_j$ in the central affine quotient.
The cusp kills $w,\delta$ and $g_1g_2$.
Since $A_1w=u+(2-c)w-\delta$, their normal closure kills $u$.
The remaining generators $\gamma,g_1$ commute and satisfy
\[
 m_1g_1=n_1\gamma,\qquad -m_2g_1=n_2\gamma.
\]
The entries of this relation matrix have gcd one, and its
determinant has absolute value $|P|$. This proves the asserted
fundamental groups, including the rank-one group when $P=0$.
For arbitrary free affine fillings one can instead kill the
entire fibre subgroup: the resulting quotient is
\[
 \langle g_1,g_2\mid g_1g_2=1,
                  g_1^{m_1}=g_2^{m_2}=1\rangle
 =\mathbb Z/\gcd(m_1,m_2).
\]
It is nontrivial, because its order is divisible by two in Type B
and by three in Type C. This proves the stronger obstruction.

\item \emph{Integral finite-quotient lattices.}
Before the cusp is attached, let $X^\circ$ be the union of the two
finite pieces. Each retracts to its reduced quotient surface $S_j$,
and their intersection retracts to a regular torus $F$.
The first linear coinvariants satisfy $w=-u$, $\delta=cu$;
the second kill $u,w$. Thus
\[
 \begin{split}
 H_1(S_1)&=\langle\gamma,u,g_1\mid
             m_1g_1=n_1\gamma+cr_1u\rangle,\\
 H_1(S_2)&=\langle\gamma,\delta,g_2\mid
             m_2g_2=n_2\gamma+r_2\delta\rangle.
 \end{split}
\]
Each relation is primitive by $\gcd(n_j,m_j)=1$; both groups are
$\mathbb Z^2$. Euler characteristic zero, duality and the universal
coefficient theorem give $H^2(S_j)=\mathbb Z^2$ with a unimodular
intersection pairing.

\item Use the ordered regular-fibre cohomology basis
\[
 (\gamma^*u^*,\gamma^*w^*,\gamma^*\delta^*,
            u^*w^*,u^*\delta^*,w^*\delta^*).
\]
Solving the integral fixed equations gives the first invariant
lattice with basis
\[
 \omega_1=-R\gamma^*u^*+R\gamma^*w^*+u^*w^*,\qquad
 \omega_2=\gamma^*u^*-\gamma^*w^*+c\gamma^*\delta^*.
\]
Both squares vanish and $\omega_1\omega_2=c\,\mathrm{vol}_F$.
The second invariant lattice has basis
$u^*w^*,\gamma^*\delta^*$ and pairing $U$.
If $L_j=p_j^*H^2(S_j)$, transfer gives injectivity and the
discriminant formula gives the exact indices
\[
 [H^2(F)^{A_1}:L_1]=m_1/c=ak,
 \qquad [H^2(F)^{A_2}:L_2]=m_2.
\]
Indeed, the quotient pairing is unimodular and pullback multiplies
it by the covering degree $m_j$.

\item The square $[q\delta+t v_j/m_j]$ closes to a two-torus in $S_j$.
For a pulled-back integral class $\alpha$ its evaluation is
$\alpha(\delta,v_j)/m_j\in\mathbb Z$.
For $\alpha=A\omega_1+B\omega_2$ this is
$-cn_1B/m_1$; coprimality forces $B\in ak\mathbb Z$.
At the second fibre it forces the coefficient of
$\gamma^*\delta^*$ to lie in $m_2\mathbb Z$.
These candidate sublattices have exactly the required indices,
so they are the entire pullback lattices, without a hidden
saturation factor. Consequently their integral basis matrices are
\[
 P_1=\begin{pmatrix}
 -R&ak\\R&-ak\\0&m_1\\1&0\\0&0\\0&0
 \end{pmatrix},\qquad
 P_2=\begin{pmatrix}
 0&0\\0&0\\0&-m_2\\1&0\\0&0\\0&0
 \end{pmatrix}.
\]
Choose the homology bases dual to these integral cohomology
bases. In particular, when $|n_j|>1$ the finite $\delta$ sweep is
not declared to be a basis vector. Its actual coordinates are
$(0,-n_1)$ and $(0,n_2)$ in the respective surfaces. This is the
integral distinction needed for arbitrary twists.
The degree-two Mayer--Vietoris map is therefore
\[
 \alpha_2=\begin{pmatrix}
 -R&R&0&1&0&0\\ak&-ak&m_1&0&0&0\\
 0&0&0&-1&0&0\\0&0&m_2&0&0&0
 \end{pmatrix}.
\]
The degree-one kernel is exactly the primitive line
$\ker\alpha_1=\mathbb Z(u+w)$: the second inclusion kills only
the $u,w$ directions of $\Lambda$, and the first then forces their
coefficients to agree. The extra relation involving $g_j$ cannot
create another kernel vector from $\Lambda$, because its
$g_j$ coefficient is the nonzero integer $m_j$.

\item \emph{The two cusp relations are primitive.}
Let $N$ be the unchanged cusp tube, $M=\partial N$, and
$P_0=A_1A_2$ its clockwise forward transport.
Direct exterior-square calculation gives rank two and a unit
$2\times2$ minor of $\bigwedge^2P_0-I$: use columns $(1,4)$
and rows $(2,6)$ in Type B or $(3,6)$ in Type C.
In both cases the determinant is $-1$, for every $R$.
Thus its cokernel is $\mathbb Z^4$. The Wang sequence then makes
$H_2(M)$ free, and the two swept tori $\xi_w,\xi_\delta$ project
to the integral basis $w,\delta$ of $\ker(P_0-I)$.
They generate a primitive direct summand.

\item For completeness, the normalization calculation used here is
integral. There are $R$ normalized $dP_6$ components, $3R$ double
curves and $2R$ triple points. With $\epsilon=1$ in Type B and
$\epsilon=2$ in Type C, component labels are
$j=a-\epsilon b\pmod R$. If a component class has coordinates
$(A_j,b_{1j},b_{2j},b_{3j})$, its three restriction differences are
\[
 \begin{split}
 X_j&=b_{1j}-A_{j+1}+b_{2,j+1}+b_{3,j+1},\\
 Y_j&=A_j-b_{1j}-b_{2j}-b_{3,j-\epsilon},\\
 Z_j&=b_{2j}-A_{j-\epsilon-1}
                  +b_{1,j-\epsilon-1}+b_{3,j-\epsilon-1}.
 \end{split}
\]
Its image is exactly the kernel of the surjective map
$(X,Y,Z)\mapsto(\sum(X+Y),\sum(Y+Z))$ to $\mathbb Z^2$.
To verify equality, put $T=S^{-1}$ for cyclic shift and use $Y$
to eliminate $A$. Solve
$Z+T^{\epsilon+1}Y=(I-T)q$, possible precisely when the second
sum vanishes; set $b_2=q$, $b_3=-T^{-\epsilon}q$.
The identity
$(I-T^{\epsilon+1})-T(I-T^\epsilon)=I-T$ solves that equation.
The other one is an integer cyclic difference equation for $b_1$
with zero coordinate sum, and is therefore soluble.
This proves saturation for all $R$, including self-identifications
when $R=1$. The normalization spectral sequence, whose other
rows are the dual two-torus and $\mathbb Z^R$ in degree four,
now gives $H^3(N)=\mathbb Z^2$, $H^4(N)=\mathbb Z^R$ and
$\chi(N)=2R$; this checks \cite[Theorem 7.8]{manYang26} directly.

\item The angular torus action on $N$ sweeps the extended zero-section
disc into integral solid-torus chains bounding $\xi_w,\xi_\delta$.
Poincar\'e--Lefschetz duality gives
$H_3(N,M)=H^3(N)=\mathbb Z^2$, so the kernel of
$H_2(M)\to H_2(N)$ has rank at most two. It contains the primitive
rank-two summand just constructed; a finite enlargement would
embed in the free quotient of $H_2(M)$ by that summand.
Hence
\[
 K_2=\ker(H_2(M)\to H_2(N))
       =\mathbb Z\xi_w\oplus\mathbb Z\xi_\delta.
\]
This establishes the complete kernel over $\mathbb Z$.

\item The two finite forward transports are $P_j=A_j^{-1}$.
They send $w$ first to $-u$ and then back to $w$.
Cutting the clockwise cusp sweep into the two cylinders gives
connecting class $-(u+w)$ in the convention
$\alpha=(i_{1*},-i_{2*})$, exactly as the explicit cylinder
boundaries in \cite[Lemma 7.10, Appendix B.4]{manYang26}.
It is a primitive generator of $\ker\alpha_1$, irrespective
of the finite translation vectors. Sweeping $\delta$ over the
pair of pants gives
$j_*\xi_\delta=-s_{1,\delta}-s_{2,\delta}$.
The evaluations already computed show that $-j_*\xi_\delta$
has the column $b=(0,-n_1,0,n_2)^{\mathsf T}$ in the integral
dual bases. Therefore
\[
 H_2(X^\circ)/j_*K_2
    =\operatorname{coker}[\alpha_2\mid b].
\]
Eliminate the unit relation from the $uw$ column. The remaining
three-generator presentation is
\[
 C=\begin{pmatrix}
 R&0&0\\-ak&m_1&-n_1\\0&m_2&n_2
 \end{pmatrix},\qquad \det C=RP.
\]
Its entry gcd is one. Its nonzero $2\times2$ minors, up to signs,
are
\[
 Rm_1,\ Rn_1,\ Rm_2,\ Rn_2,\ akm_2,\ akn_2,\ P.
\]
Their gcd is exactly $h=\gcd(R,ak,P)$, by the two coprimality
conditions. Thus for $P\ne0$ its Smith factors are
$(1,h,R|P|/h)$; for $P=0$ its rank is two and its factors are
$(1,h,0)$. The latter cokernel is $\mathbb Z\oplus\mathbb Z/h$.

\item \emph{The relative term, including $P=0$.}
Excision and duality give $H_2(N,M)=H^4(N)=\mathbb Z^R$.
Before the cusp is attached, the joint monodromy images kill
$\ker\chi$ in the abelianization. Hence
\[
 H_1(X^\circ)=\left\langle\gamma,g_1,g_2\ \middle|\
 m_1g_1=n_1\gamma,\quad m_2g_2=n_2\gamma\right\rangle
       =\mathbb Z\oplus\mathbb Z/d_0,
 \quad d_0=\gcd(m_1,m_2).
\]
The equality follows integrally: the relation matrix has entry
gcd one and maximal-minor gcd $d_0$.
The full image of the relative boundary map is
$\langle a_0\rangle$, where $a_0=-(g_1+g_2)$.
Indeed, the cusp boundary kernel is generated by its meridian
and its vanishing lattice; the latter is already zero in
$H_1(X^\circ)$, and the extended section disc supplies $a_0$.

\item When $P\ne0$, the primitive homomorphism from $H_1(X^\circ)$
to its free quotient is
\[
 \Phi(\gamma,g_1,g_2)
 =\frac1{d_0}(m_1m_2,\ m_2n_1,\ m_1n_2).
\]
The same gcd-of-minors calculation proves primitivity.
It sends $a_0$ to $-P/d_0\ne0$, so $\langle a_0\rangle$
is infinite cyclic. The kernel of the relative boundary map
is free of rank $R-1$. This does not assume that $a_0$ is
primitive in $H_1(X^\circ)$.

\item When $P=0$, equality of the reduced fractions
$n_1/m_1=-n_2/m_2$ forces $m_1=m_2=m$ and $n_2=-n_1$.
Use $a_0$ and $g_1$ as meridian generators: the relations are
$ma_0=0$ and $mg_1=n_1\gamma$. The second is primitive,
so $a_0$ has exact order $m$. The relative boundary image is
therefore the finite group $\mathbb Z/m$, and its kernel is
free of rank $R$, of index $m$ in $\mathbb Z^R$.
This is the extra free contribution in the degenerate case.

\item The exact sequence for $(X,X^\circ)$ now reads
\[
 0\longrightarrow\operatorname{coker}C
 \longrightarrow H_2(X)
 \longrightarrow\ker\bigl(H_2(N,M)\to H_1(X^\circ)\bigr)
 \longrightarrow0.
\]
Its right group is free, so it splits. This gives the displayed
degree-two groups in both cases. All pieces other than the cusp
have Euler characteristic zero, so $\chi(X)=2R$.
For $P\ne0$, duality gives $b_1=b_5=0$, $b_2=b_4=R-1$,
and Euler characteristic gives $b_3=0$.
For $P=0$, it gives $b_1=b_5=1$, $b_2=b_4=R+1$ and then
$b_3=2$. In either case the universal coefficient theorem and
duality identify the degree-three torsion with the degree-two
torsion and the degree-four torsion with the degree-one torsion.
This proves all the integral homology assertions separately
in the two cases.

\item The script \texttt{code/manifolds\_yang.py} verifies the symbolic
monodromy and primitive-minor identities, 32 full cusp
normalization Smith forms, 162 finite-quotient and sweep
checks, and 2592 full augmented $4\times7$ Smith forms,
including 72 with $P=0$. Its executable also accepts arbitrary
admissible integer inputs. These finite checks accompany the
unbounded integral proof; they impose no hypothesis on it.
The original cases are proved in \cite[Theorem~7.13]{manYang26}.
\end{enumerate}
\end{proof}

\begin{theorem}[all orders $3k$ for $IV^*+II+I_2$]\label{maps:src:D.2.11}
\StatementRecord{maps:src:D.2.11}{proved}{Historical events: Claude Fable 5.2 (Anthropic), effort high, 16 November 2026: carried as background without fault found. Claude Fable 5.2 (Anthropic), effort high, 17 November 2026: carried as background without fault found. Claude Opus 5.5 (Anthropic), effort max, sweep 1, scope s6, 18 November 2026: used and rechecked in this run. GPT-6 Astra (OpenAI), ChatGPT Pro mode, effort Pro, S6\_09\_PAPERS, 17 November 2026: targeted proof check. Also checked in this writing. Pre-continuation writing events: GPT-6 Astra (OpenAI), ChatGPT Pro mode, effort Pro, S6\_10\_PAPERS, 18 November 2026: targeted proof check. Continuation metadata events: GPT-6 Astra (OpenAI), ChatGPT Pro mode, effort Pro, S6\_10\_PAPERS continuation, 19 November 2026: used and rechecked in this run; status-vocabulary reconciliation; metadata reconciliation only; no new mathematical check. \emph{Pre-continuation recorded proof scope.} GPT-6 Astra (OpenAI, ChatGPT Pro mode), effort Pro, 2026-11-17 UTC; Claude Fable 5.2, effort high, 2026-11-17; GPT-6 Astra (OpenAI, ChatGPT Pro mode), S6\_06\_BLIND, effort Pro, 2026-11-17; Claude Fable 5.2 (Anthropic), effort high, 16 November 2026; also checked in this writing by GPT-6 Astra (OpenAI), ChatGPT Pro mode, effort Pro, 18 November 2026}{the hypotheses, notation and definitions stated here; \cref{maps:src:D.2.1}; \cite[Proposition 3.1]{manYang26}; \cite[Proposition 3.4]{manYang26}; \cite[Proposition 3.6]{manYang26}; \cite[Lemma 3.10]{manLR26}; \cite[Proposition 3.11]{manLR26}; \cref{maps:src:D.2.10}; \cref{maps:src:D.2.6}}


The rational elliptic surface
\[
 y^2=x^3-3t^3(t-1)x+t^4(t-1)(2t-1),
 \qquad P=(t^2,t^2),
 \tag{D.C1}\label{maps:eq:D.C1}
\]
has fibres $IV^*,II,I_2$ at $0,1,\infty$ and Mordell--Weil group
$\mathbb Z$ generated by $P$. Put $M=\mathcal O(P-O)$ and
$Q=3P$. Let $C_0$ be the identity component of the $I_2$ fibre.
There is a linearization with divisor $C_0$, and its torus
family has Yang's index-one Type C marking. Keep the type-$II$
filling unchanged and use the index-one toroidal filling at
$I_2$. At $IV^*$ choose any logarithmic order and invariant vector
\[
 m=3k,\qquad v=n(\gamma-u-w)+r\delta,
 \qquad k>0,\quad n,r\in\mathbb Z,
 \quad\gcd(n,3k)=1.
\]
The resulting compact smooth complex threefold has
\[
 \pi_1(X)=\mathbb Z/|n|,\qquad
 H_i(X;\mathbb Z)=
 \begin{cases}
 \mathbb Z&i=0,6,\\
 \mathbb Z/|n|&1\leq i\leq4,\\
 0&i=5.
 \end{cases}
\]
Consequently $n=\pm1$ gives the standard $S^6$ for every positive
$k$. This is an affirmative unbounded configuration for
Viaclovsky's Question 4. The linearization is supported on one
$I_2$ component, so the local-narrowness restriction of the public
engine does not include this input.
\end{theorem}


\begin{proof}
\begin{enumerate}[label=\arabic*.,leftmargin=*]
\item \emph{The explicit elliptic surface and linearization.}
For the coefficients $f,g$ of \eqref{maps:eq:D.C1}, direct calculation gives
\[
 4f^3+27g^2=27t^8(t-1)^2,\qquad
 j=6912t(1-t).
\]
The minimal valuations give $IV^*$ at zero and $II$ at one.
In the chart at infinity the discriminant has order two and
the cubic has a node, giving $I_2$. These fibres have Euler
sum twelve; the degree-four and degree-six Weierstrass
coefficients give the rational elliptic surface. Substitution
checks the section $P$, which is disjoint from $O$ and meets
the nonidentity components at $IV^*$ and $I_2$.
Its height is $2-4/3-1/2=1/6$. The Shioda--Tate rank is one,
and OS entry 49 has height lattice $\langle1/6\rangle$ and
no torsion, as checked in \ref{maps:src:D.2.1}; hence $P$ is a generator.

\item The section $Q=3P$ meets the identity component of $IV^*$
and the nonidentity component of $I_2$. Its height is $3/2$,
so the height formula gives $Q\cdot O=0$.
The mixed height $\langle P,Q\rangle=1/2$, with the sole
local contribution $1/2$ at $I_2$, gives $P\cdot Q=0$.
Let $J=M\otimes t_Q^*M^{-1}$. Its multidegree vanishes at
$IV^*$ and $II$. At $I_2$, translation by $Q$ exchanges the
two components, so its multidegrees are $(-2,2)$ on
$(C_0,C_1)$, exactly those of $\mathcal O(C_0)$.
Their difference is thus a pullback from the base. Restriction
to $O$ gives
\[
 J\cdot O=Q\cdot O-P\cdot Q+1=1=C_0\cdot O,
\]
so the base degree of that difference is zero. Therefore
$J\simeq\mathcal O(C_0)$, supplying the asserted effective
linearization. Its norm under the component exchange has divisor
$C_0+C_1=F_\infty$. This explains why a half full-fibre weight
appears in period coordinates while the geometric divisor is
integral.

\item \emph{Identification with the actual Type C periods.}
The elliptic $j$-function is the one in
\cite[Proposition 3.1]{manYang26}. Over $t=s^3$ the good
model has central curve $Y^2=X^3+1$, and $P$ specializes to
$(0,1)$, a nonzero fixed point of order three.
Over $t-1=s^6$, the same section specializes to the origin
of the good elliptic model. At $I_2$ it meets the nonidentity
component. Choose the sign of the elliptic marking so that its
Abel--Jacobi coordinate at the first point is $(\tau+1)/3$.
Its normalized lift therefore obeys the two Type C affine
laws and has $\mu-\tau/2$ regular at infinity.
The uniqueness proof of \cite[Proposition 3.4]{manYang26},
using $H^0(\mathcal O(-1))=H^1(\mathcal O(-1))=0$, identifies
it with $\mu_C$.
The extension line $M$ has this same Picard coordinate,
whereas the translation has coordinate $3\mu_C$.
The norm divisor just calculated fixes the cusp central
period. The remaining scalar period is the affine constant
of \cite[Proposition 3.6]{manYang26}. Consequently the open
two-torus family is the index-one Type C family, for a suitable
admissible constant, and its primitive toroidal cusp has one
normalized $dP_6$ component.

\item An exact integral check of the marking is also available.
In the order $(II,IV^*,I_2)$, the homology matrices are
\[
 \begin{split}
 M_{II}&=\begin{pmatrix}1&1&0&0\\-1&0&0&0\\0&0&1&0\\0&0&0&1\end{pmatrix},\\
 M_{IV}&=\begin{pmatrix}0&-1&0&0\\1&-1&0&-3\\0&1&1&1\\0&0&0&1\end{pmatrix},\\
 M_{I_2}&=\begin{pmatrix}1&2&0&3\\0&1&0&0\\0&-1&1&-1\\0&0&0&1\end{pmatrix}.
 \end{split}
\]
The first two follow from the locally narrow finite-point
section cocycles $(0,0)$ and $(0,-1)$ for $P$, multiplied by
three for $Q$. Their product determines the third matrix;
no local-narrowness hypothesis at $I_2$ is used.
Put $U=M_{IV}$, $V=U^{-1}M_{II}U$ and
\[
 S=\begin{pmatrix}
 0&0&-1&0\\0&-1&1&1\\0&0&0&1\\-1&0&3&0
 \end{pmatrix}.
\]
Then $\det S=-1$, $U^{-t}S=ST_1$ and $V^{-t}S=ST_2$,
where $T_1,T_2$ are precisely Yang's displayed Type C
cohomological matrices. This is the finite Hurwitz reorder
followed by an integral change of basis, not merely a rational
comparison of eigenvalues.

\item \emph{The unchanged type-$II$ filling with its integral marking.}
Here the product proof of \cite[Lemma 3.10]{manLR26} applies
with $3P$ in place of its translation section; we check all the
needed ingredients. A closed elliptic neighborhood $E$ of the
type-$II$ fibre retracts to its cuspidal cubic, homeomorphic to
$S^2$. The restriction of $M$ has degree zero and is therefore
differentiably trivial. In a smooth trivialization the lifted
contracting generator has form
\[
 (x,z)\longmapsto(t_{-Q}(x),\lambda(x)z),\qquad |\lambda|<1.
\]
As $E$ is simply connected, $\lambda$ has a smooth logarithm.
The relative exponential of the additive N\'eron group is an
isomorphism on the central smooth locus. After shrinking the
disc it provides a holomorphic logarithm of $-Q$ and hence a
fibre-preserving isotopy of $t_{-Q}$ to the identity.
The translations extend over the cuspidal cubic and fix its
cusp. Combining this isotopy with the logarithm of $\lambda$
trivializes the circle-bundle mapping torus of a fundamental
annulus. Thus this filling is homeomorphic over its base disc
to $E\times T^2$, with angular circle $\delta$ and a primitive
second circle of $\gamma^*$-value one.

\item Its regular specialization sends the elliptic area $uw$ to
the fundamental class of the cuspidal cubic and $\gamma\delta$
to the area of this second torus; it kills the mixed classes.
Equivalently the exact degree-two pullback lattice is
$\langle u^*w^*,\gamma^*\delta^*\rangle$.
There is no index-six factor: this is the unchanged filling,
not a free order-six affine quotient.
Both transverse normal tori bound products of a section disc
with a circle. In particular its $\delta$ meridian sweep is
zero. These assertions also follow by applying the marked
specialization proof of \cite[Proposition 3.11]{manLR26}
to the product just constructed.

\item \emph{All orders and all displayed twists.}
At $IV^*$ the order-$3k$ free quotient and its marked collar
are those of \ref{maps:src:D.2.10} with $R=1$, $n_1=n$, $r_1=r$.
The two other compact fillings have just been specified,
so the same finite-core gluing proves compactness and
smoothness. The unchanged type-$II$ piece kills $u,w$ and
its meridian; the cusp kills $\delta$ and the remaining
meridian. Van Kampen leaves exactly $n\gamma=0$, proving
$\pi_1(X)=\mathbb Z/|n|$.

\item The first finite pullback matrix is unchanged. For the second
piece the primitive product lattice replaces its coefficient
$m_2$ by one, and its vanishing $\delta$ sweep replaces $n_2$
by zero. Thus the complete augmented integral map is
\[
 D=\begin{pmatrix}
 -1&1&0&1&0&0&0\\
 k&-k&3k&0&0&0&-n\\
 0&0&0&-1&0&0&0\\
 0&0&1&0&0&0&0
 \end{pmatrix}.
\]
The primitive $w$ cusp sweep still kills the degree-one
Mayer--Vietoris kernel. Unit elimination reduces $D$ to
one unit factor and
\[
 \begin{pmatrix}1&0&0\\-k&3k&-n\\0&1&0\end{pmatrix}.
\]
Its Smith factors are $(1,1,|n|)$ for arbitrary allowed
integers, independently of $r$.
Before attaching the cusp, the remaining degree-one group is
$\langle\gamma,g\mid3kg=n\gamma\rangle=\mathbb Z$.
Its cusp meridian is $-g$, a nonzero multiple $-n$ of a
primitive generator. Hence the relative map
$H_2(N,M)=\mathbb Z\to H_1(X^\circ)=\mathbb Z$ is injective,
so it contributes no extra free degree-two class.
We conclude $H_2(X)=\mathbb Z/|n|$.
The cusp has Euler number two and the other pieces have
Euler number zero. Duality and the universal coefficient
theorem give the remaining groups. When $|n|=1$, sphere
recognition applies exactly as in \ref{maps:src:D.2.6}.

\item The independent executable
\texttt{code/manifolds\_type\_c\_ii.py} checks the Weierstrass
identity, section, $j$-function, integral monodromy conjugacy,
symbolic determinant, and 168 full augmented Smith forms.
The preceding proof has no upper bound on $k,n,r$.
\end{enumerate}
\end{proof}

The local Type-C geometric input is supplied by \cite[Propositions~3.1, 3.4, 3.6, 4.4 and Appendix~B.4]{manYang26}; the product type-$II$ marking is supplied by \cite[Lemma~3.10 and Propositions~3.11, 3.13]{manLR26}. The unrestricted $3P$ calculations use the explicitly displayed boundary maps.

\begin{theorem}[seven further unbounded one-log families]\label{maps:src:D.2.12}
\StatementRecord{maps:src:D.2.12}{proved}{Historical events: Claude Fable 5.2 (Anthropic), effort high, 16 November 2026: carried as background without fault found. GPT-6 Astra (OpenAI), ChatGPT Pro mode, effort Pro, S6\_09\_PAPERS, 17 November 2026: targeted proof check. Also checked in this writing. Pre-continuation writing events: GPT-6 Astra (OpenAI), ChatGPT Pro mode, effort Pro, S6\_10\_PAPERS, 18 November 2026: targeted proof check. Continuation metadata events: GPT-6 Astra (OpenAI), ChatGPT Pro mode, effort Pro, S6\_10\_PAPERS continuation, 19 November 2026: used and rechecked in this run; status-vocabulary reconciliation; metadata reconciliation only; no new mathematical check. \emph{Pre-continuation recorded proof scope.} GPT-6 Astra (OpenAI, ChatGPT Pro mode), effort Pro, 2026-11-17 UTC; Claude Fable 5.2 (Anthropic), effort high, 16 November 2026; also checked in this writing by GPT-6 Astra (OpenAI), ChatGPT Pro mode, effort Pro, 18 November 2026}{the hypotheses, notation and definitions stated here; \cref{maps:src:D.2.1}; \cref{maps:src:D.2.6}; \cite{manExplorer26}; \cref{maps:eq:D.R1}; \cref{maps:eq:D.R3}; \cref{maps:eq:D.R4}}


Let $G$ be the marked Mordell--Weil generator of the indicated
rational elliptic surface. In the table, the superscript $\bullet$
marks the fibre supporting the degree-one linearization, and $d$
is the linear order at the sole additive fibre.
\[
 \begin{array}{c|l|rr|r}
 \mathrm{OS}&\text{ordered fibre configuration}&P/G&Q/G&d\\\hline
 45&I_8,II,I_1^\bullet,I_1&8&1&6\\
 47&I_7,I_2,I_1^\bullet,II&14&1&6\\
 47&I_7,I_1^\bullet,I_1,III&7&2&4\\
 55&I_5,I_4,II,I_1^\bullet&20&1&6\\
 56&I_5,I_2,IV,I_1^\bullet&10&3&3\\
 56&I_5,I_3,III,I_1^\bullet&15&2&4\\
 49&I_2,IV^*,I_1^\bullet,I_1&2&3&3
 \end{array}
\]
Use $M=\mathcal O(P-O)$ and translation by $Q$.
Keep the two weight-zero multiplicative fillings and the
primitive toroidal filling at $I_1^\bullet$.
In the integral marking specified below, put
\[
 m=dk,\qquad v=n\ell+r\delta,
 \qquad k>0,\quad\gcd(n,dk)=1,
 \qquad n,r\in\mathbb Z.
\]
Every row admits these logarithmic fillings for every displayed
integer choice, and the resulting compact smooth threefold has
\[
 \begin{gathered}
 \pi_1(X)=\mathbb Z/|n|,\qquad
 H_i(X;\mathbb Z)=\mathbb Z/|n|\quad(1\leq i\leq4),\\
 H_5(X)=0,\qquad H_0(X)=H_6(X)=\mathbb Z.
 \end{gathered}
\]
In particular $n=\pm1$ gives $S^6$ for every $k>0$ in all seven
rows. Within the one-log construction, any simply connected
choice of affine vector must have $n=\pm1$, so the coprimality
hypothesis discards none of its sphere cases.
\end{theorem}


\begin{proof}
\begin{enumerate}[label=\arabic*.,leftmargin=*]
\item \emph{Geometric data and the actual multiplicative fillings.}
The configurations and height lattices are the primary
Oguiso--Shioda/Fukae data verified in \ref{maps:src:D.2.1}. Each displayed
pair has height pairing one. At every fibre at least one of
$P,Q$ is locally narrow, and at each weight-zero $I_j$ the
section $P$ is narrow while $Q$ generates its component group.
In the marked component groups the latter residues are
\[
 \begin{array}{c|l}
 \text{row}&Q\text{ on its higher multiplicative fibres}\\\hline
 1&3\pmod8\\
 2&5\pmod7,\quad1\pmod2\\
 3&3\pmod7\\
 4&2\pmod5,\quad3\pmod4\\
 5&3\pmod5,\quad1\pmod2\\
 6&2\pmod5,\quad1\pmod3\\
 7&1\pmod2.
 \end{array}
\]
All residues are units. At the additive fibre $Q$ is narrow;
its component residue is zero. The exact $P$ residues as well
as these $Q$ residues are retained with the matrices in the
accompanying \texttt{manifolds\_rankone\_families.json} record.

\item The line $J=M\otimes t_Q^*M^{-1}$ has multidegree zero on
every fibre: either $t_Q$ preserves its components, or $M$
already has multidegree zero there. Translation acts trivially on the degree-zero Picard group of
the generic elliptic fibre, so $J$ is generically trivial. Its
vertical divisor and zero component degrees then make it a
pullback from the base; all original fibres have multiplicity one.
Its degree is $\langle P,Q\rangle=1$, as all local mixed
height corrections vanish. Choose its section with its simple
zero at the marked $I_1^\bullet$. The contracting-quotient
construction and its properness estimate are the same as in
Engel's construction.

\item At a weight-zero $I_j$, the restriction of $M$ to each normalized
component $\mathbb P^1$ has degree zero and is trivial.
Translation by $Q$ cyclically permutes all $j$ components,
because its component residue is a unit. The quotient therefore
has one normalization component and one double curve.
After $j$ steps the first component returns; its normalization
is the quotient of $\mathbb P^1\times\mathbb C^*$ by a
translation on the second factor and a constant multiplicative
automorphism on the first. It is a $\mathbb P^1$ bundle over an
elliptic curve defined by a degree-zero line, hence is
topologically trivial. Its zero and infinity sections are
identified by a translation and flat units. These are isotopic
to the identity in the marked circle directions.
Each node has the original semistable chart $xy=t$ up to a
holomorphic unit: the infinite contracting action has no point
stabilizer and does not ramify this chart.
Thus the filled neighborhood and its normal-circle chains are
the primitive one-node model $N_{\rm node}\times T^2$ used
in \ref{maps:src:D.2.6}. In particular all three fixed one-cycles supply its
three normal tori, not just the vanishing circle.
This argument identifies the geometric local model, rather
than inferring it from a rational monodromy matrix.

\item At the marked $I_1$, the simple zero gives the usual one-component
hexagonal toroidal filling. Its two angular normal tori are
primitive. At the additive fibre, take the good-reduction root
model and then its degree-$k$ base change. An invariant vector
with primitive invariant-coordinate value $n$ defines a free
order-$dk$ affine action because $\gcd(n,dk)=1$.
The marked logarithmic collar and compact gluing are exactly
those proved in \ref{maps:src:D.2.6}. These steps prove existence and properness
for every allowed $k,n,r$.

\item \emph{Exact integral certificates, with no parameter bounds.}
Use the geometric open-family monodromy lifts from the section
cocycles and degree-one linearization, as in \ref{maps:src:D.2.1} and the
public primary monodromy derivation of \cite{manExplorer26}.
Cyclically reorder the four matrices so that the additive
one is first, and denote them by $A_0,A_1,A_2,A_3$, with
$A_0A_1A_2A_3=I$.
The marking here is $(e_1,e_2,\delta,\gamma)$, with
$\chi=\gamma^*$. The following integral vectors and constants
are obtained by exact unimodular reduction:
\[
 \begin{array}{c|r|r|l}
 \text{row}&d&q&\ell\\\hline
 1&6&1&(-1,0,0,1)\\
 2&6&1&(0,-1,0,1)\\
 3&4&2&(-3,-2,0,1)\\
 4&6&1&(0,0,0,1)\\
 5&3&3&(-1,3,0,1)\\
 6&4&2&(-1,3,0,1)\\
 7&3&3&(-3,-1,0,1).
 \end{array}
\]
In every row $A_0\ell=\ell$, $A_i\delta=\delta$ and
$\chi(\ell)=1$.
All four matrices, integral fixed bases and the integer
left inverses used next are written in the retained JSON;
the executable derives and verifies them using only the
MIT monodromy source, without a local-model database or API.

\item Let $D=[A_0-I\ A_1-I\ A_2-I]$ and set
\[
 j_0(x)=(x,0,0),\quad j_1(x)=(0,x,0),\quad
 j_2(x)=(0,0,x),\quad
 j_3(x)=-(A_1A_2x,A_2x,x).
\]
For each $i=1,2,3$, choose the full integral fixed basis
$F_i=\ker(A_i-I)$, of rank two at the cusp and rank three
at a one-node fibre. Form the $12\times9$ matrix
\[
 K=[j_1(F_1)\ j_2(F_2)\ j_3(F_3)\ j_0(\ell)].
\]
The exact certificates give
\[
 DK=0,\quad\operatorname{rank}D=3,\quad LK=I_9,\quad
 (Lj_0\delta)_9=0,
 \tag{D.R1}\label{maps:eq:D.R1}
\]
where $L$ is the explicitly retained integer left inverse.
Thus the nine actual sweep tori form a primitive integral
basis of the graph kernel. The exterior coinvariant matrix
$[\bigwedge^2A_i-I]_{i=0}^3$ has five nonzero Smith entries,
all one. Its primitive dual row is the invariant form
\[
 \omega=P_0\,e_1^*e_2^*-Q_0\,\delta^*\gamma^*,
 \qquad P=P_0G,\quad Q=Q_0G.
 \tag{D.R2}\label{maps:eq:D.R2}
\]
The regular-family $H_2$ is therefore free of rank ten:
one fibre coinvariant of $\omega$-value one and the nine
sweeps. The constant form $\omega$ has zero value on the
latter actual fixed-loop sweeps, so this is an integral
splitting. Quotienting by the eight normal tori of the
unchanged fillings gives
$\mathbb Z\langle s,t\rangle$, with $t=j_0(\ell)$.
Equation \eqref{maps:eq:D.R1} also proves $j_0(\delta)=0$ in this quotient.

\item In the vertical lattice $\ker\chi$, the finite coinvariants
are free of rank one. Let $b$ be a primitive representative;
the exact fixed-basis computation gives
\[
 [\delta]=q[b],\qquad
 \mathcal N_d b=(d/q)\delta,\qquad
 \mathcal N_d\ell=d\ell,\qquad
 -\omega(b,\ell)=1,\quad
 -\omega(\ell,\delta)=-q.
 \tag{D.R3}\label{maps:eq:D.R3}
\]
Here $\mathcal N_d=\sum_{i=0}^{d-1}A_0^i$; all are
identities of integral vectors, not rational averages.
The finite reduced fibre has free $H_1$, with generators
$b,z$, where
\[
 mg=n\ell+rq b,\qquad
 z=\alpha\ell+\beta g,\qquad m\alpha+n\beta=1.
\]
Apply the normal-circle torus formula of \ref{maps:src:D.2.6} to these two
cycles. The first column is the fibre area
$-\omega(b,v)$ and the sweep $j_0(\mathcal N_m b)$;
the second is $-\alpha\omega(\ell,v)$ and
$j_0(\alpha\mathcal N_m\ell+\beta v)$.
After the eight unchanged normal relations, \eqref{maps:eq:D.R1}--\eqref{maps:eq:D.R3}
give exactly
\[
 \begin{pmatrix}n&-qr\alpha\\0&1\end{pmatrix}.
 \tag{D.R4}\label{maps:eq:D.R4}
\]
Its Smith form is $(1,|n|)$ for all the asserted integers.

\item For van Kampen, the images of the three unchanged vanishing
lattices already generate the full primitive $\ker\chi$;
their combined matrix has Smith factors $(1,1,1)$, verified
in the same exact computation. Their meridians and the base
relation kill the remaining finite meridian. The only remaining
relation is $n\gamma=0$. This proves the nonabelian
fundamental-group assertion. It also shows, without assuming
coprimality, that any simply connected affine choice must have
$n=\pm1$.

\item The degree-one Mayer--Vietoris map has free source and target
of rank fourteen and finite cokernel $\mathbb Z/|n|$; hence
its kernel is zero. Every local degree-two boundary map is
onto, by the product models, the finite free quotient, and
the toroidal specialization proved above. Therefore \eqref{maps:eq:D.R4}
is the full remaining degree-two presentation and
$H_2(X)=\mathbb Z/|n|$.
Only the hexagonal cusp contributes Euler characteristic,
giving $\chi(X)=2$. Duality supplies all the other groups,
and homotopy-sphere recognition with $\Theta_6=0$ proves
the sphere statement.

\item The exact program \texttt{code/manifolds\_rankone\_families.py}
checks all constant matrices, integral bases and \eqref{maps:eq:D.R1}--\eqref{maps:eq:D.R3},
and also 234 full $10\times10$ Smith forms at displayed
sample parameters. Formula \eqref{maps:eq:D.R4} proves the result for
unrestricted orders and twists; the samples are corroboration.
\end{enumerate}
\end{proof}

\begin{theorem}[all smooth-fibre logs in four configurations]\label{maps:src:D.2.13}
\StatementRecord{maps:src:D.2.13}{proved}{Historical events: Claude Fable 5.2 (Anthropic), effort high, 16 November 2026: carried as background without fault found. GPT-6 Astra (OpenAI), ChatGPT Pro mode, effort Pro, S6\_09\_PAPERS, 17 November 2026: targeted proof check. Also checked in this writing. Pre-continuation writing events: GPT-6 Astra (OpenAI), ChatGPT Pro mode, effort Pro, S6\_10\_PAPERS, 18 November 2026: targeted proof check. Continuation metadata events: GPT-6 Astra (OpenAI), ChatGPT Pro mode, effort Pro, S6\_10\_PAPERS continuation, 19 November 2026: used and rechecked in this run; status-vocabulary reconciliation; metadata reconciliation only; no new mathematical check. \emph{Pre-continuation recorded proof scope.} GPT-6 Astra (OpenAI, ChatGPT Pro mode), effort Pro, 2026-11-17 UTC; Claude Fable 5.2 (Anthropic), effort high, 16 November 2026; also checked in this writing by GPT-6 Astra (OpenAI), ChatGPT Pro mode, effort Pro, 18 November 2026}{the hypotheses, notation and definitions stated here; \cref{maps:src:D.2.14}; \cref{maps:src:D.2.12}; \cite[Proposition 2.3]{manViaclovskyLogs26}; \cref{maps:eq:D.S6}; \cref{maps:eq:D.S5}; \cref{maps:eq:D.S8}}


Consider the following four marked rational elliptic surfaces,
with $G$ a Mordell--Weil generator:
\[
\begin{array}{c|l|r}
\mathrm{OS}&\text{ordered fibres}&p\\\hline
45&I_8,I_1,I_1,I_1,I_1^\bullet&8\\
47&I_7,I_2,I_1,I_1,I_1^\bullet&14\\
55&I_5,I_4,I_1,I_1,I_1^\bullet&20\\
56&I_5,I_3,I_2,I_1,I_1^\bullet&30.
\end{array}
\]
Use $P=pG$, $Q=G$, and the degree-one linearization with zero
at $I_1^\bullet$. Keep the original four primitive one-node
fillings and the one-component hexagonal cusp. At an additional
smooth fibre, perform the logarithmic transformation of order
$m>0$ and integral translation vector
\[
 v=a e_1+b e_2+c\delta+n\gamma,\qquad
 \gcd(m,a,b,c,n)=1.
 \tag{D.S1}\label{maps:eq:D.S1}
\]
This is the exact freeness condition; no coprimality condition
on $m,n$ alone is imposed. Put
\[
 t=\gcd(m,n,a,b)\,\gcd(m,n,p,c).
 \tag{D.S2}\label{maps:eq:D.S2}
\]
The two factors in \eqref{maps:eq:D.S2} are coprime, and $t\mid m$.
For $n\ne0$, also $t\mid n$, and the result has
\[
\begin{aligned}
 \pi_1(X)&=H_1(X)=\mathbb Z/|n|,\qquad H_5(X)=0,\\
 H_2(X)&=H_3(X)=G_{m,n,t},\qquad H_4(X)=\mathbb Z/|n|,\\
 G_{m,n,t}&=\mathbb Z/t\oplus\mathbb Z/m
               \oplus\mathbb Z/(m|n|/t).
\end{aligned}
\tag{D.S3}\label{maps:eq:D.S3}
\]
For $n=0$, the answer is instead
\[
\begin{aligned}
 \pi_1(X)&=H_1(X)=H_5(X)=\mathbb Z,\\
 H_2(X)&=H_3(X)=\mathbb Z^2\oplus\mathbb Z/t\oplus\mathbb Z/m,
 \qquad H_4(X)=\mathbb Z^2.
\end{aligned}
\tag{D.S4}\label{maps:eq:D.S4}
\]
In both cases $H_0=H_6=\mathbb Z$ and $\chi(X)=2$.
The simply connected members are precisely $n=\pm1$; they
have $H_2=H_3=(\mathbb Z/m)^2$. Consequently the sphere
members in these four families are precisely $m=1$, $n=\pm1$.
For $m>1$ the theorem gives simply connected complex rational
homology spheres with two cyclic factors of order $m$ in their
middle torsion. Their oriented smooth type is identified in
\ref{maps:src:D.2.14} below.
\end{theorem}


\begin{proof}
\begin{enumerate}[label=\arabic*.,leftmargin=*]
\item \emph{Geometric construction and complete affine scope.}
In these four configurations the height of $G$ is $1/p$.
The section $P$ is narrow at every fibre, whereas $Q$ has
unit component residue at each higher multiplicative fibre:
respectively $3$ modulo $8$; $(5,1)$ modulo $(7,2)$;
$(2,1)$ modulo $(5,4)$; and $(1,1,1)$ modulo $(5,3,2)$.
Thus $M=\mathcal O(P-O)$ has degree zero on each component,
and the quotient by $Q$ cycles through all components at each
weight-zero fibre. The geometric argument of \ref{maps:src:D.2.12} identifies
each such filling with the primitive one-node model, including
all three normal tori. The mixed height is one, so
$M\otimes t_Q^*M^{-1}$ is the pullback of $\mathcal O(1)$.
Its chosen simple zero gives the stated hexagonal cusp.

\item Over a small disc at the additional regular value, the torus
family has a holomorphic origin and trivial integral local
system. Pull it back by $z=s^m$ and quotient by
\[
 (x,s)\longmapsto (x+v(s)/m,e^{2\pi i/m}s).
\]
Here $v(s)$ is the holomorphic period corresponding to the
integral vector $v$. The action is free exactly when \eqref{maps:eq:D.S1}
holds: away from $s=0$ its nonidentity powers move the base,
and on the central torus translation by $v/m$ has exact
order $m$ precisely under \eqref{maps:eq:D.S1}. The usual logarithmic
collar, given by the translation
$(\log s)v(s)/(2\pi i)$, glues this proper quotient to the
original family; compare \cite[Proposition 2.3]{manViaclovskyLogs26}.
Changing $v$ by a multiple of $m$ retains its specified
integral collar clutch, rather than discarding that clutch.
The reduced central fibre is the torus with lattice
\[
 \Lambda'=\Lambda+\mathbb Z(v/m).
 \tag{D.S5}\label{maps:eq:D.S5}
\]
Its normal complex line is flat of finite order. Since the
integral $H^2$ of a torus is torsion-free, this line is
topologically trivial. Thus the boundary-to-fibre map is
onto on $H_2$, and its kernel consists of its four normal
circle tori.

\item \emph{Primitive global graph calculation.}
The marking is $(e_1,e_2,\delta,\gamma)$ and the invariant
coordinate is $\chi=\gamma^*$. Insert $A_0=I$ for the extra
regular fibre before the five original monodromies
$A_1,\ldots,A_5$. Their product is $I$.
All original vanishing lattices together generate precisely
the primitive $\ker\chi$; their combined matrix has three
unit Smith entries. Therefore van Kampen reduces the group
to the generator $\gamma$. All singular meridians vanish,
and the base relation kills the remaining regular meridian
$\xi$. The new relation $m\xi-v=0$ is then $n\gamma=0$.
This proves the assertion about the full fundamental group.

\item For completeness the integral degree-two calculation is
recorded explicitly. Set
\[
 D=[A_0-I\ A_1-I\ \cdots\ A_4-I]:\mathbb Z^{20}\to\mathbb Z^4.
\]
It has rank three and primitive image. As in \ref{maps:src:D.2.12}, denote
the five direct boundary embeddings into the graph chains by
$j_i$, and set
\[
 j_5(x)=-\bigl(A_1A_2A_3A_4x,\ A_2A_3A_4x,
                      \ A_3A_4x,\ A_4x,\ x\bigr).
\]
The unchanged fillings supply fourteen normal tori: three
at each of the four one-node fibres and two at the cusp.
If $B$ is their $20\times14$ graph matrix, then
\[
 K=[B\ j_0(e_1)\ j_0(e_2)\ j_0(\gamma)]
 \tag{D.S6}\label{maps:eq:D.S6}
\]
is a primitive integral basis of $\ker D$.
The retained integer matrix $L$ satisfies $LK=I_{17}$,
and the last three coordinates of $Lj_0(x)$ are
$(x_1,x_2,x_4)$. In particular the residual $\delta$ sweep
vanishes. These identities are verified for all four
constant marked monodromy tuples in the accompanying exact
program, which saves the matrices and their integer inverses.

\item The exterior coinvariants are free of rank one, detected by
\[
 \omega=p e_1^*e_2^*-\delta^*\gamma^*.
\]
Their presentation has five unit Smith entries. Every column
of \eqref{maps:eq:D.S6} is an actual fixed-loop sweep, on which the constant
vertical form $\omega$ has value zero. This specifies an
integral splitting of regular-family $H_2$.
After eliminating the fourteen unchanged normal tori,
the remaining group is freely generated by a fibre class
of $\omega$-value one and the three sweeps in \eqref{maps:eq:D.S6}.

\item \emph{The exact normal map.}
The lattice \eqref{maps:eq:D.S5} is generated by $e_1,e_2,\delta,\gamma,g=v/m$,
with its single primitive relation
$mg=a e_1+b e_2+c\delta+n\gamma$.
For an integral vector $x$, its normal torus has fibre
coordinate $-\omega(x,v)$ and sweep $j_0(mx)$.
The $g$ normal torus has fibre coordinate zero and sweep
$j_0(v)$. One can read these formulas directly on the
affine boundary torus: its normal circle is $m\xi-v$,
so the first contribution is $-x\wedge v$, and the second
is $mx$ swept around $\xi$. This gives the complete
$4\times5$ matrix
\[
 C=\begin{pmatrix}
 -pb&pa&n&-c&0\\
 m&0&0&0&a\\
 0&m&0&0&b\\
 0&0&0&m&n
 \end{pmatrix}.
 \tag{D.S7}\label{maps:eq:D.S7}
\]
Its columns satisfy exactly the stated lattice relation,
so the five-generator presentation introduces no additional
normal torus.

\item \emph{Smith reduction for every integer choice.}
The greatest common divisor of all entries of $C$ is one,
by \eqref{maps:eq:D.S1}. Its maximal minors, up to signs, are
\[
 m^3n,\quad m^2n^2,\quad m^2nc,\quad m^2nb,\quad m^2na.
\]
For $n\ne0$ their gcd is $m^2|n|$.
The lower determinantal divisors can be checked over each
local ring $\mathbb Z_{(\ell)}$, as follows.
If $m$ is a unit, three diagonal pivots leave just the
relation $n$.
If $n$ is a unit and $m$ is not, its pivot in the top row
and its pivot in the last column leave two copies of $m$.
In the remaining case $\ell\mid m,n$, condition \eqref{maps:eq:D.S1}
forces $c$ or one of $a,b$ to be a unit.
If $c$ is a unit, pivot on $-c$. Row operations give one
copy of $m$ together with the two-row presentation
\[
 \begin{pmatrix}m&0&h\\0&mn&n\end{pmatrix},
 \qquad (h)=(a,b)\text{ in }\mathbb Z_{(\ell)}.
\]
For $n\ne0$ the two nonzero determinantal divisors of this block
are $\gcd(m,n,a,b)$ and $m|n|$, up to local units.
For $n=0$ it has rank one and entry gcd $\gcd(m,a,b)$.
If instead $a$ or $b$ is a unit, pivot on that entry in
the last column. The remaining three-row presentation is
equivalent, up to units, to
\[
 \begin{pmatrix}p&n&-c\\m&0&0\\0&0&m\end{pmatrix}.
\]
Its first two determinantal divisors are
$\gcd(m,n,p,c)$ and $m\gcd(m,n,p,c)$.
The two gcd factors in \eqref{maps:eq:D.S2} have disjoint prime divisors:
a common divisor would divide $m,a,b,c,n$, contrary to
\eqref{maps:eq:D.S1}. These local calculations therefore give the global
Smith entries
\[
 (1,t,m,m|n|/t)\quad(n\ne0),\qquad
 (1,t,m,0)\quad(n=0).
 \tag{D.S8}\label{maps:eq:D.S8}
\]
The same unit-pivot argument proves the $n=0$ formula
directly, without treating a finite group order as a limit.

\item The degree-one Mayer--Vietoris map has free source and
target of rank twenty-four. Its cokernel is the already
computed $H_1(X)$. Thus its kernel is zero when $n\ne0$,
and is a free copy of $\mathbb Z$ when $n=0$.
All local degree-two boundary maps are onto, including
the smooth logarithmic piece by \eqref{maps:eq:D.S5}.
Consequently the full remaining degree-two cokernel is
$\operatorname{coker} C$, with precisely this extra free
relative term in the $n=0$ case. A free quotient splits
the resulting abelian-group extension. This proves $H_2$.
Only the hexagonal cusp contributes Euler characteristic,
giving two. Integral duality and the universal coefficient
theorem now give \eqref{maps:eq:D.S3}--\eqref{maps:eq:D.S4}; in particular $b_3=2$
when $n=0$, and torsion appears in both $H_2$ and $H_3$.
The sphere conclusion follows from $\Theta_6=0$.

\item The program \texttt{code/manifolds\_semistable.py} verifies
all four constant primitive graph and exterior calculations
and 2924 full $18\times19$ Smith presentations, including
388 cases with $n=0$. The proof of \eqref{maps:eq:D.S8} has no bounds
on any parameter. The source comparison is recorded in the verification index.
\end{enumerate}
\end{proof}

\begin{theorem}[the twisted Smale suspension]\label{maps:src:D.2.14}
\StatementRecord{maps:src:D.2.14}{proved}{Historical events: Claude Fable 5.2 (Anthropic), effort high, 16 November 2026: carried as background without fault found. GPT-6 Astra (OpenAI), ChatGPT Pro mode, effort Pro, S6\_09\_PAPERS, 17 November 2026: targeted proof check. Also checked in this writing. Pre-continuation writing events: GPT-6 Astra (OpenAI), ChatGPT Pro mode, effort Pro, S6\_10\_PAPERS, 18 November 2026: targeted proof check. Continuation metadata events: GPT-6 Astra (OpenAI), ChatGPT Pro mode, effort Pro, S6\_10\_PAPERS continuation, 19 November 2026: used and rechecked in this run; status-vocabulary reconciliation; metadata reconciliation only; no new mathematical check. \emph{Pre-continuation recorded proof scope.} GPT-6 Astra (OpenAI, ChatGPT Pro mode), effort Pro, 2026-11-17 UTC; Claude Fable 5.2 (Anthropic), effort high, 16 November 2026; also checked in this writing by GPT-6 Astra (OpenAI), ChatGPT Pro mode, effort Pro, 18 November 2026}{the hypotheses, notation and definitions stated here; \cite[Definition 5.1]{manGR25}; \cref{maps:src:D.2.13}; \cref{maps:src:D.2.12}; \cref{maps:eq:D.S7}; \cref{maps:eq:D.S11}; \cite[Theorem 2.9]{manGR25}; \cite[Theorems 2.2--2.3]{manBarden65}; \cite[Section 4.3]{manReiser24}; \cite[Theorem A]{manGR25}; \cref{maps:eq:D.S13}; \cite[Lemma 5.2]{manGR25}; \cite[Theorem 1]{manZhubr75}; \cref{maps:eq:D.S12}; \cref{maps:eq:D.S14}; \cite[Example 5.3(1)]{manGR25}}


Let $N_m$ be the unique closed simply connected spin five-manifold
with $H_2(N_m;\mathbb Z)=(\mathbb Z/m)^2$, so that $N_1=S^5$.
Define its twisted suspension by
\[
 \widetilde\Sigma_0 N_m
 =\bigl((N_m\setminus\mathring D^5)\times S^1\bigr)
       \cup_{\widetilde\alpha}(S^4\times D^2),
 \qquad \widetilde\alpha(x,z)=(\alpha_zx,z),
 \tag{D.S9}\label{maps:eq:D.S9}
\]
where $\alpha:S^1\to SO(5)$ represents its nontrivial
fundamental-group element. This specifies the framing in
\cite[Definition 5.1]{manGR25}; it is not replaced by an
unspecified choice between the two suspension framings.
Every simply connected member of the four families in
\ref{maps:src:D.2.13} is oriented diffeomorphic to
\[
                         \widetilde\Sigma_0 N_m.
 \tag{D.S10}\label{maps:eq:D.S10}
\]
In particular its oriented smooth type depends only on $m$,
and these types are distinct for distinct positive $m$.
It is spin and has $p_1=0$. The proof retains the full
refined invariant required when $m$ is even; vanishing of
$p_1$ alone is not used as a classification criterion.
\end{theorem}


\begin{proof}
\begin{enumerate}[label=\arabic*.,leftmargin=*]
\item \emph{The angular action and its fixed curve.}
Multiplication of the line-bundle coordinate by
$e^{i\theta}$ defines the global circle action of direction
$\delta$. It commutes with the contracting quotient,
the toroidal gluings and the logarithmic translations.
On the noncompact line-bundle core the action is free:
a nontrivial contracting iterate cannot agree with a unit
scalar, by the strict norm contraction. On the four
weight-zero one-node fillings it is free for the same
reason on the normalized $\mathbb P^1\times\mathbb C^*$
annulus, and it remains free on its double curve.

\item The primitive hexagonal cusp is different. In the height-one
toric lattice, use the cones on the two triangles of a unit
square, with the diagonal joining $(1,0)$ to $(0,1)$.
The cocharacter $\delta=(0,1,0)$ has weights $(-1,0,1)$
in each of their smooth affine charts, up to permutation.
Every nonzero weight is a unit, so the action is semifree
there. Its fixed curves correspond to vertical edges, which
form a single orbit under the two period translations.
They therefore give one embedded $C\simeq\mathbb P^1$
in the quotient. The adjacent wall relation is
\[
 (1,0,1)+(-1,1,1)=(0,0,1)+(0,1,1).
\]
It gives the two normal degrees $-1,-1$. With the circle
weights retained, the normal bundle is
\[
 \nu_C=\mathcal O(-1)_{+1}\oplus\mathcal O(-1)_{-1}.
 \tag{D.S11}\label{maps:eq:D.S11}
\]
Holomorphic unit factors in the actual annular charts
commute with the angular action and deform equivariantly
to the toric factors, so these are the actual smooth
normal data. This also gives a direct chart verification
of the fixed edge described in Engel's final remarks.

\item On the added smooth logarithmic fibre, the lattice is
$\Lambda'=\Lambda+\mathbb Z(v/m)$. If a scalar angle
$s\delta$ lies in this lattice, applying $\chi=\gamma^*$
shows that $kn/m$ is an integer for its $v/m$ coefficient
$k$. Since here $n=\pm1$, it follows that $m\mid k$,
and then $s$ is an integer. Thus there are no new finite
stabilizers at this fibre. The global circle action is
semifree with fixed locus exactly $C$.

\item \emph{Equivariant surgery and the spin quotient.}
The specialization $H_2(T^4)\to H_2$ of the cusp is onto,
by the integral annulus calculation used in \ref{maps:src:D.2.12}.
Consequently $[C]$ maps into the cyclic fibre-coinvariant
summand generated by $s$ in \eqref{maps:eq:D.S7}. Its $\delta$ normal
column imposes $ns=0$; hence $[C]=0$ when $n=\pm1$.
Since $X$ is simply connected, Hurewicz makes $C$
null-homotopic.

\item Conjugating the second complex coordinate in \eqref{maps:eq:D.S11}
identifies its equivariant real bundle with
$\mathcal O(-1)\oplus\mathcal O(1)$ carrying equal
scalar weights. Its transition loop is
$\operatorname{diag}(z^{-1},z)\in SU(2)$ and is
null-homotopic. Thus an equivariant tubular neighborhood
is $S^2\times D^4$, with trivial action on the first
factor and the free Hopf action on each punctured
normal sphere. Conjugating one normal coordinate reverses
normal orientation, so simultaneously reverse the model
$S^2$ orientation in this tubular identification; the
ambient orientation is then preserved.
Replace it equivariantly by $D^3\times S^3$.
The resulting $Y$ has a free circle action.

\item Two homotopic embeddings of $S^2$ in a six-manifold
are isotopic in this dimension; for instance apply
\cite[Theorem 2.9]{manGR25}, since $6>2\cdot2+1$.
Therefore the null-homotopic $C$ is isotopic to a small
unknotted sphere in a ball. Its framing is unique because
$\pi_2(SO(4))=0$. The underlying surgery is consequently
\[
                         Y\cong X\#(S^3\times S^3).
 \tag{D.S12}\label{maps:eq:D.S12}
\]
In particular $Y$ is simply connected.

\item The complex manifold $X$ is spin: $H^2(X;\mathbb Z)=0$
forces $c_1(X)=0$. Its spin structure is unique.
The circle action lifts with period $2\pi$ to this spin
structure. To see the period, lift the universal covering
action of $\mathbb R$ and inspect it at a point of $C$:
the parity of the isotropy weights is $1+(-1)=0$.
Thus the full $2\pi$ rotation acts trivially on the spin
lift, rather than by its nontrivial deck involution.
The equivariant trivialization of \eqref{maps:eq:D.S11} and its standard
Hopf replacement extend this same lift over the surgery.
The free action on $Y$ therefore has an equivariant spin
structure. Quotienting its horizontal spin frames shows
that the oriented orbit manifold $B=Y/S^1$ is spin.

\item The homotopy sequence gives $\pi_1(B)=0$ and the split
abelian-group sequence
\[
 0\longrightarrow H_2(Y)\longrightarrow H_2(B)
       \xrightarrow{\langle e,-\rangle}\mathbb Z
       \longrightarrow0,
\]
where $e$ is the Euler class of $Y\to B$. Thus
$H_2(B)=\mathbb Z\oplus(\mathbb Z/m)^2$ and $e$ is
primitive. The spin case of the oriented Smale--Barden
classification, with its realization of homology
isomorphisms, gives
\[
 (B,e)\cong\bigl((S^2\times S^3)\#N_m,(x,0)\bigr),
 \tag{D.S13}\label{maps:eq:D.S13}
\]
where the sign of the generator $x$ on the first summand
is chosen by the oriented circle-bundle convention.
The precise primary statements are
\cite[Theorems 2.2--2.3]{manBarden65}; the spin uniqueness
for $N_m$ is also recalled in \cite[Section 4.3]{manReiser24}.

\item \emph{The framing and cancellation.}
Apply \cite[Theorem A]{manGR25} to \eqref{maps:eq:D.S13}, with
$B_1=S^2\times S^3$ and $B_2=N_m$.
The Euler class on $B_1$ is primitive, its circle-bundle
total space is $S^3\times S^3$, and $B_1$ is spin.
The spin branch of that theorem specifies the twisted
framing, and yields
\[
          Y\cong(S^3\times S^3)\#\widetilde\Sigma_0 N_m.
 \tag{D.S14}\label{maps:eq:D.S14}
\]
By \cite[Lemma 5.2]{manGR25}, the latter suspension is
simply connected and spin, has the same $H_2$ as $N_m$,
and has $b_3=0$. Zhubr's complete oriented classification
\cite[Theorem 1]{manZhubr75} permits cancellation of the
$S^3\times S^3$ summand in \eqref{maps:eq:D.S12}--\eqref{maps:eq:D.S14}: connected
sum with it leaves $H_2$, the refined class $p$ and the
class $\mu$ unchanged, and adds two to $b_3$.
The diffeomorphism of the stabilized manifolds identifies
these full invariants; deleting the common increase of
$b_3$ identifies the classification data of the two
unstabilized manifolds. This proves \eqref{maps:eq:D.S10}, including
the two-primary data. No general cancellation theorem
outside this classified class is being assumed.

\item Finally $TY\simeq\pi^*TB\oplus\mathbb R$, and
$H^4(B;\mathbb Z)=H_1(B;\mathbb Z)=0$, so $p_1(Y)=0$.
Restriction across \eqref{maps:eq:D.S12} gives $p_1(X)=0$ even though
its target $H^4(X;\mathbb Z)=(\mathbb Z/m)^2$ can be
nonzero. Distinct $m$ have distinct integral torsion
orders. For $m=1$, \eqref{maps:eq:D.S9} is $S^6$, also directly by
\cite[Example 5.3(1)]{manGR25}.
The argument computes the smooth type rather than
deducing it from these additive groups or from $p_1$.
\end{enumerate}
\end{proof}

\begin{theorem}[all-order $S^2\times S^4$ outputs]\label{maps:src:D.2.15}
\StatementRecord{maps:src:D.2.15}{proved}{Historical events: Claude Fable 5.2 (Anthropic), effort high, 16 November 2026: carried as background without fault found. GPT-6 Astra (OpenAI), ChatGPT Pro mode, effort Pro, S6\_09\_PAPERS, 17 November 2026: targeted proof check. Also checked in this writing. Pre-continuation writing events: GPT-6 Astra (OpenAI), ChatGPT Pro mode, effort Pro, S6\_10\_PAPERS, 18 November 2026: targeted proof check. Continuation metadata events: GPT-6 Astra (OpenAI), ChatGPT Pro mode, effort Pro, S6\_10\_PAPERS continuation, 19 November 2026: used and rechecked in this run; status-vocabulary reconciliation; metadata reconciliation only; no new mathematical check. \emph{Pre-continuation recorded proof scope.} GPT-6 Astra (OpenAI, ChatGPT Pro mode), effort Pro, 2026-11-17 UTC; Claude Fable 5.2 (Anthropic), effort high, 16 November 2026; also checked in this writing by GPT-6 Astra (OpenAI), ChatGPT Pro mode, effort Pro, 18 November 2026}{the hypotheses, notation and definitions stated here; \cref{maps:src:D.2.12}; \cref{maps:eq:D.I1}; \cref{maps:src:D.1.2}; \cref{maps:eq:D.I2}; \cref{maps:eq:D.I3}; \cref{maps:eq:D.I4}; \cref{maps:eq:D.S11}; \cref{maps:src:D.2.14}; \cite[Theorem A and Theorem B(2)]{manGR25}}


Take the OS56 type-$II$ profile with ordered elliptic fibres
\[
                         I_2,\ II,\ I_5,\ I_3.
\]
Let $G$ be its marked Mordell--Weil generator, take
$P=30G$, $Q=G$, and put the degree-one linearization zero
at $I_2$. Use the two-component toroidal filling there,
the original primitive one-node fillings at $I_5,I_3$,
and the logarithmic filling at $II$ with
\[
 m=6k,\qquad v=n\gamma+r\delta,\qquad
 k>0,\quad\gcd(n,6k)=1,\quad r\in\mathbb Z.
\]
The resulting compact smooth threefold has
\[
\begin{aligned}
 \pi_1(X)&=H_1(X)=\mathbb Z/|n|,\qquad H_5(X)=0,\\
 H_2(X)&=H_4(X)=\mathbb Z\oplus\mathbb Z/|n|,
 \qquad H_3(X)=\mathbb Z/|n|,\\
 H_0(X)&=H_6(X)=\mathbb Z,\qquad \chi(X)=4.
\end{aligned}
\]
Every simply connected member, equivalently $n=\pm1$,
is oriented diffeomorphic to $S^2\times S^4$.
In particular this fixed configuration supplies that
smooth type for every logarithmic multiplicity $6k$.
\end{theorem}


\begin{proof}
\begin{enumerate}[label=\arabic*.,leftmargin=*]
\item \emph{The actual local geometries.}
The generator has height $1/30$, and $P$ is narrow at
every fibre. The component residues of $Q$ are $1$ in
each of $\mathbb Z/2$, $\mathbb Z/5$, $\mathbb Z/3$,
and zero at $II$. The mixed height is one, so the same
vertical-divisor argument as in \ref{maps:src:D.2.12} gives
$M\otimes t_Q^*M^{-1}=\pi^*\mathcal O(1)$.
Its section can have its simple zero on the whole $I_2$
fibre. The two weight-zero multiplicative fibres give
the primitive one-node models proved in \ref{maps:src:D.2.12}.
The type-$II$ good-reduction model has linear order six,
with invariant basis $\delta,\gamma$; its higher-order
free affine quotients are supplied by the same holomorphic
construction as in \ref{maps:src:D.2.12}.

\item At the zero of the linearization, the Tate parameter has
order two. Because $P$ is narrow, its line has order zero
in the angular Tate direction, while $Q$ has component
residue one. The two valuation periods of the torus
uniformization are therefore
\[
                 (2,0),\qquad(1,1).
 \tag{D.I1}\label{maps:eq:D.I1}
\]
The second coordinate one is the simple zero of the
linearization. Holomorphic units do not change this
integral lattice. Use the standard unit-triangle Mumford
fan and quotient by the lattice $\Gamma$ in \eqref{maps:eq:D.I1}.
It has index two, so the central fibre has two normalized
$dP_6$ components, six double curves and four triple
points. Its total space is smooth: all the fan cones are
unimodular, and the contracting period action is free
and properly discontinuous on the chosen tube.
This is the specified two-component filling, not a
primitive one-component replacement of its monodromy.

\item \emph{The integral cusp defect.}
Let $U$ be a compact tubular neighborhood of this cusp,
and $M=\partial U$. The fan symmetry
\[
 F=\begin{pmatrix}1&1\\0&-1\end{pmatrix},\qquad
 F\begin{pmatrix}2&1\\0&1\end{pmatrix}
     =\begin{pmatrix}2&2\\0&-1\end{pmatrix}
\]
preserves the triangular fan and identifies this lattice
with $2\mathbb Z\oplus\mathbb Z$.
Consequently the integral normalization resolution in
\ref{maps:src:D.1.2}, at $d=2$, applies directly and gives
\[
 \pi_1(U)=\mathbb Z^2,\qquad
 (H_0,H_1,H_2,H_3,H_4)(U)
   =(\mathbb Z,\mathbb Z^2,\mathbb Z^5,\mathbb Z^2,\mathbb Z^2).
 \tag{D.I2}\label{maps:eq:D.I2}
\]
This local normalization calculation does not require the
global covering degree to be coprime to six.

\item The cusp monodromy $A_3$ has $\operatorname{rank}(A_3-I)=2$
and Smith entries $(1,2)$, while
$\bigwedge^2A_3-I$ has two unit Smith entries.
The Wang sequences therefore give
\[
 H_1(M)=\mathbb Z^3\oplus\mathbb Z/2,
 \qquad H_2(M)=\mathbb Z^6.
\]
The cusp kills the full saturated angular lattice
$\ker(A_3-I)$, and its canonical meridian bounds the
extended section disc. Hence the map on $H_1$ is onto
$H_1(U)=\mathbb Z^2$, with kernel
$\mathbb Z\oplus\mathbb Z/2$.
Poincare--Lefschetz duality gives
$H_2(U,M)=H^4(U)=\mathbb Z^2$ and $H_1(U,M)=0$.
The relative sequence then shows that
\[
 H_2(U)/\operatorname{im}H_2(M)\cong\mathbb Z.
 \tag{D.I3}\label{maps:eq:D.I3}
\]
In particular this quotient has no hidden two-torsion.
The two angular normal tori bound in the toric model.
They project to the full primitive fixed lattice in the
Wang sequence for $H_2(M)$, so they form a primitive
rank-two subgroup. By \eqref{maps:eq:D.I2}--\eqref{maps:eq:D.I3} the kernel of
$H_2(M)\to H_2(U)$ has exactly rank two; these are
therefore its whole kernel.
Moreover the remaining image is exactly the regular-fibre
specialization image. Its cokernel is the free group
in \eqref{maps:eq:D.I3}.

\item \emph{The global integral map.}
Cyclically reorder the original four monodromies as
$A_0,A_1,A_2,A_3$ corresponding to $II,I_5,I_3,I_2$.
The full matrices are retained in the accompanying certificate
\texttt{manifolds\_i2\_cusp.json}.
They satisfy $A_0\cdots A_3=I$, $A_0^6=I$ and
fix $\delta$ and the covector $\gamma^*$.
Here $A_0\gamma=\gamma$ and
\[
 \omega=30e_1^*e_2^*-\delta^*\gamma^*
\]
is the primitive exterior coinvariant detector.
The exterior presentation has five unit Smith entries,
so its group is free of rank one.

\item There is one integral issue in degree one that cannot be
discarded. The combined matrix $[A_i-I]_{i=0}^3$ has
Smith entries $(1,1,2)$, so the regular family's fibre
coinvariants are $\mathbb Z\oplus\mathbb Z/2$.
The torsion in the cusp boundary is represented by
$-\delta$, and its image in this global $\mathbb Z/2$
has coefficient $-1$. Thus the torsion summands in
the degree-one Mayer--Vietoris map are mapped
isomorphically. Replacing the cusp image by its actual
saturated angular vanishing lattice gives Smith entries
$(1,1,1)$; it kills the entire primitive
$\ker\gamma^*$ in van Kampen. The base relation kills
the remaining finite meridian, leaving just
$n\gamma=0$. This proves $\pi_1(X)=\mathbb Z/|n|$.

\item Use the boundary maps $j_i$ of \ref{maps:src:D.2.12} and put
\[
 K=[j_1\ker(A_1-I)\ j_2\ker(A_2-I)
       \ j_3\ker(A_3-I)\ j_0(\gamma)].
\]
The exact integer certificate gives a left inverse
$LK=I_9$, and $(Lj_0\delta)_9=0$.
As before, these actual fixed-loop sweeps have zero
$\omega$-area. After the six one-node normal tori and
the two cusp normal tori, the regular-family $H_2$ is
therefore $\mathbb Z\langle s,t\rangle$ with
$t=j_0\gamma$.
The finite vertical coinvariants are generated by
$\delta$, and
$\mathcal N_m\delta=m\delta$,
$\mathcal N_m\gamma=m\gamma$.
Writing $m\alpha+n\beta=1$, its two normal tori give
\[
                        \begin{pmatrix}n&-r\alpha\\0&1\end{pmatrix}.
 \tag{D.I4}\label{maps:eq:D.I4}
\]
Their cokernel is $\mathbb Z/|n|$ for all parameters.
All these equalities are integral vector identities in
the retained matrices.

\item The degree-two boundary maps of the other fillings are
onto. At the cusp their cokernel is the free group
\eqref{maps:eq:D.I3}. Thus the full degree-two Mayer--Vietoris cokernel
is $\mathbb Z\oplus\mathbb Z/|n|$; the extension by
the free cusp defect splits. The degree-one source and
target both have free rank fourteen and one copy of
$\mathbb Z/2$. Their torsion map is the isomorphism
just checked, and their total cokernel is the finite
group $H_1(X)$. Their kernel is therefore zero.
This proves the asserted $H_2$. Four toric triple points
give Euler characteristic four; integral duality then
gives all the remaining groups.

\item \emph{Smooth recognition at $n=\pm1$.}
In this case $H_2(X)=\mathbb Z$, and the cusp map
$H_2(U)\to H_2(X)$ is onto, because \eqref{maps:eq:D.I4} kills the
entire regular-family contribution. The local Mumford
threefold has trivial canonical line. Hence $c_1(X)$
evaluates as zero on every $H_2$ class, so $c_1(X)=0$
and $X$ is spin.

\item The angular $\delta$ action is semifree with two fixed
rational curves. The one-node pieces are free as before.
At the finite logarithmic fibre, its invariant
$\gamma^*$ coordinate changes by $n/m$ under a deck
generator; for $n=\pm1$ this prevents any additional
stabilizer of the $\delta$ action. At the cusp there
are two vertical-edge orbits, each with the normal
weights and degrees \eqref{maps:eq:D.S11}.

\item It remains essential to check that these two fixed curves
have zero homology class; specialization is not onto in
\eqref{maps:eq:D.I3}. Colour a lattice vertex $(a,b)$ by $a+b$ modulo
two to label its normalized component. For a vertical
edge from $(a,b)$ to $(a,b+1)$, the two adjacent third
vertices are $(a+1,b)$ and $(a-1,b+1)$.
Each colour occurs once among the endpoints and once
among the third vertices. The toric wall formula thus
gives degree $-1+1=0$ for the line of either global
component on either fixed curve. Their images in
$H_2(U,M)$ are exactly these component intersection
numbers, so both images are zero. By the relative
sequence both curves lie in the boundary image, which
is precisely the regular-fibre specialization image.
Equation \eqref{maps:eq:D.I4}, with $n=\pm1$, now kills both classes
in $H_2(X)$.

\item The equivariant surgery and spin-lift argument of \ref{maps:src:D.2.14}
can therefore be performed on both disjoint curves.
Whitney isotopy applies also to their disjoint union,
and their normal framings are unique. It gives
\[
            Y\cong X\#_2(S^3\times S^3)
\]
with a free circle action and simply connected spin
quotient $B$. The homotopy sequence gives
$H_2(B)=\mathbb Z^2$ with primitive Euler class.
The oriented Smale--Barden theorem identifies the pair
with $((S^2\times S^3)\#(S^2\times S^3),(x,0))$.
Applying \cite[Theorem A and Theorem B(2)]{manGR25}
gives
\[
              Y\cong(S^2\times S^4)\#_2(S^3\times S^3).
\]
Zhubr cancellation as in \ref{maps:src:D.2.14} proves
$X\cong S^2\times S^4$ with its oriented smooth type.

\item The exact program \texttt{code/manifolds\_i2\_cusp.py}
verifies the primitive graph and exterior maps, the
nonzero boundary-torsion image, the actual saturated
vanishing lattice, the two fixed-curve degree vectors,
and 630 full $10\times10$ Smith forms.
The norm identities and \eqref{maps:eq:D.I4} prove the result without
any bounds on $k,n,r$.
\end{enumerate}
\end{proof}

\begin{theorem}[original odd-star fillings and all regular logs]\label{maps:src:D.2.16}
\StatementRecord{maps:src:D.2.16}{proved}{Historical events: Claude Fable 5.2 (Anthropic), effort high, 16 November 2026: carried as background without fault found. GPT-6 Astra (OpenAI), ChatGPT Pro mode, effort Pro, S6\_09\_PAPERS, 17 November 2026: targeted proof check. Also checked in this writing. Pre-continuation writing events: GPT-6 Astra (OpenAI), ChatGPT Pro mode, effort Pro, S6\_10\_PAPERS, 18 November 2026: targeted proof check. Continuation metadata events: GPT-6 Astra (OpenAI), ChatGPT Pro mode, effort Pro, S6\_10\_PAPERS continuation, 19 November 2026: used and rechecked in this run; status-vocabulary reconciliation; metadata reconciliation only; no new mathematical check. \emph{Pre-continuation recorded proof scope.} GPT-6 Astra (OpenAI, ChatGPT Pro mode), effort Pro, 2026-11-17 UTC; Claude Fable 5.2 (Anthropic), effort high, 16 November 2026; also checked in this writing by GPT-6 Astra (OpenAI), ChatGPT Pro mode, effort Pro, 18 November 2026}{the hypotheses, notation and definitions stated here; \cref{maps:src:D.2.14}; \cref{maps:src:D.2.12}; \cite[Section 2.4, (2.2)]{manUU21}; \cite[Sections 7.6--7.7]{manSS10}; \cite[Proposition 3.5(1),(3) and its proof]{manUU21}; \cref{maps:eq:D.T5}; \cref{maps:eq:D.T4}; \cref{maps:eq:D.T6}; \cref{maps:eq:D.T7}; \cref{maps:eq:D.T9}; \cref{maps:src:D.2.13}; \cref{maps:eq:D.T10}; \cite[Theorem A]{manGR25}; \cite[Theorem B(2)]{manGR25}}


Take one of the following marked configurations:
\[
\begin{array}{c|l|r|r}
 \mathrm{OS}&\text{ordered fibres}&p&b\\\hline
 46&I_3^*,I_1,I_1^\bullet,I_1&4&2\\
 50&I_1^*,I_3,I_1^\bullet,I_1&12&1.
\end{array}
\]
Set $P=pG$, $Q=G$, and place the degree-one linearization
zero at $I_1^\bullet$. Retain the original $I_j^*$ filling,
the primitive weight-zero nodal fillings and the hexagonal cusp.
At an additional regular fibre perform any free log of order
$m>0$ with vector $v=a e_1+b' e_2+c\delta+n\gamma$,
where $\gcd(m,a,b',c,n)=1$. Define
\[
 t=\gcd(m,n,a,b')\,\gcd(m,n,p,c),\qquad
 G=\mathbb Z/t\oplus\mathbb Z/m\oplus\mathbb Z/(m|n|/t)
 \quad(n\ne0).
\]
For $n\ne0$ the complete answer is
\[
\begin{aligned}
 \pi_1(X)&=H_1(X)=\mathbb Z/|n|, &H_5(X)&=0,\\
 H_2(X)&=\mathbb Z^b\oplus G,
 &H_3(X)&=\mathbb Z^{2b}\oplus G,\\
 H_4(X)&=\mathbb Z^b\oplus\mathbb Z/|n|.
\end{aligned}
\tag{D.T1}\label{maps:eq:D.T1}
\]
For $n=0$, put $T=\mathbb Z/t\oplus\mathbb Z/m$; then
\[
\begin{aligned}
 \pi_1(X)&=H_1(X)=H_5(X)=\mathbb Z,\\
 H_2(X)&=\mathbb Z^{b+2}\oplus T,
 &H_3(X)&=\mathbb Z^{2b+2}\oplus T,\\
 H_4(X)&=\mathbb Z^{b+2}.
\end{aligned}
\tag{D.T2}\label{maps:eq:D.T2}
\]
In both cases $H_0=H_6=\mathbb Z$ and $\chi(X)=2$.
All simply connected members, exactly those with $n=\pm1$,
have the oriented smooth type
\[
 X\cong \#^b(S^2\times S^4)\ \#\ \#^b(S^3\times S^3)
                   \ \#\ \widetilde\Sigma_0 N_m,
 \tag{D.T3}\label{maps:eq:D.T3}
\]
where the precisely framed twisted suspension and the spin
Smale manifold $N_m$ are as in \ref{maps:src:D.2.14}. In particular
$\widetilde\Sigma_0N_1=S^6$, so $m=1$ gives the corresponding
connected sums of sphere products.
\end{theorem}


\begin{proof}
\begin{enumerate}[label=\arabic*.,leftmargin=*]
\item \emph{The original filling, including its component action.}
The generator heights are $1/4$ and $1/12$ respectively.
In each case $P$ is narrow everywhere, $Q$ generates the
order-four component group at $I_j^*$, and in the second
case $Q$ also generates the component group of $I_3$.
Consequently the line $M=\mathcal O(P-O)$ has zero
multidegree at every fibre. As in \ref{maps:src:D.2.12} the line
$M\otimes t_Q^*M^{-1}$ is vertical and has zero multidegree,
hence is pulled back from the base; its degree is the mixed
height $\langle P,Q\rangle=1$. The specified linearization
exists. The $I_3$ quotient is the actual primitive nodal
model, by its unit component action. The simple zero on
$I_1^\bullet$ gives the primitive hexagonal cusp.

\item Here is an integral topological model of the original star
filling. Let $E\to\Delta$ be the minimal elliptic-surface
neighbourhood with fibre $I_j^*$, $j=3$ or $1$.
It retracts onto the affine $D_{j+4}$ tree of $j+5$ spheres.
A section in the identity component has a holomorphic local
logarithm. Explicitly the local exponential sequence is
$0\to\mathcal G\to\omega^{-1}\to E^{\rm id}\to0$,
and its stalk $\mathcal G_0$ is zero at additive reduction;
see \cite[Section 2.4, (2.2)]{manUU21}.
Thus choose $R$ in the identity component with $4R=4Q$.
Then $Q_0=Q-R$ has exact order four and the same component
as $Q$. The path $Q_0+sR$, $0\le s\le1$, is defined using
that logarithm and induces an isotopy of fibre-preserving
translations on the entire minimal surface. These
translations extend over all components, as also recalled in
\cite[Sections 7.6--7.7]{manSS10}. The order-four invariant
section can alternatively be read directly from
\cite[Proposition 3.5(1),(3) and its proof]{manUU21}.

\item Likewise the narrow section $P$ can be contracted to zero
through its local logarithm. Topologically this also follows
from the zero component degrees of $M$: $H^2(E;\mathbb Z)$
is free on those degrees, so $M$ is trivial. A unit in its
contraction lift has a logarithm on $E$, since $H^1(E;\mathbb Z)=0$.
Straightening this unit and the two section logarithms gives
a smooth family of contracting quotients with marked boundary.
A uniformly sufficiently small contraction parameter makes
the whole family proper and free. Ehresmann's theorem with
boundary therefore identifies the original filling with
\[
 U\cong S^1_\delta\times T(t_{Q_0}:E\to E),
 \tag{D.T4}\label{maps:eq:D.T4}
\]
where $T$ denotes the mapping torus. Integral fibre markings
are transported through this family, so this is a statement
about the actual original filling and its boundary map.
The line deformation is rigidified along $O$, and the unit
logarithm is chosen in that rigidification. Thus every member
has the same origin disc. On its regular fibres these are
deformations of group extensions and their translating period;
integral period loops are transported continuously, so their
integer monodromy is constant. The origin disc and fixed
period loops therefore give the transported meridian and sweep
marking, without an additional angular winding.

\item Translation by $Q_0$ cycles the four multiplicity-one end
components. The incidence tree then forces it to reflect the
chain of $j+1$ double components. All component maps preserve
complex orientation. Its permutation on $H_2(E)$ therefore
has
\[
                  \rho=1+(j+1)/2=b+1
\]
orbits. Invariants and coinvariants of a permutation lattice
are free on its orbits. The mapping-torus sequence and the
product with $S^1$ give, integrally,
\[
 \pi_1(U)=\mathbb Z^2,\qquad
 (H_0,H_1,H_2,H_3,H_4)(U)
  =(\mathbb Z,\mathbb Z^2,\mathbb Z^{\rho+1},
                         \mathbb Z^{2\rho},\mathbb Z^\rho).
 \tag{D.T5}\label{maps:eq:D.T5}
\]

\item \emph{The full boundary kernel and the free defect.}
In the global $(e_1,e_2,\delta,\gamma)$ marking, the star
monodromies are, respectively,
\[
 A_*^{46}=\begin{pmatrix}
 -1&-3&0&1\\0&-1&0&1\\-4&-8&1&3\\0&0&0&1
 \end{pmatrix},\qquad
 A_*^{50}=\begin{pmatrix}
 -1&-1&0&-1\\0&-1&0&1\\-12&-24&1&3\\0&0&0&1
 \end{pmatrix}.
 \tag{D.T6}\label{maps:eq:D.T6}
\]
In particular $(A_*^2-I)$ has rank one: these are not
finite-order monodromies. Both $A_*-I$ and
$\bigwedge^2A_*-I$ have only unit nonzero Smith entries,
with ranks two and four. For $M_*=\partial U$ the Wang
sequences give $H_1(M_*)=\mathbb Z^3$ and
$H_2(M_*)=\mathbb Z^4$. The map on $H_1$ is onto the
$\mathbb Z^2$ in \eqref{maps:eq:D.T5}, with primitive kernel generated
by the meridian, which bounds the extended origin disc.
Thus Poincare--Lefschetz duality and the relative sequence
imply
\[
 H_2(U)/\operatorname{im}H_2(M_*)=\mathbb Z^{\rho-1}
                                           =\mathbb Z^b.
 \tag{D.T7}\label{maps:eq:D.T7}
\]
The assertion is integral, since the map from
$H_2(U,M_*)=H^4(U)=\mathbb Z^\rho$ onto the meridian
kernel $\mathbb Z$ is surjective.

\item The two circle actions in \eqref{maps:eq:D.T4} give the entire normal-torus
kernel. One is $\delta$; the other is translation in the
mapping-torus coordinate with period four, using $Q_0^4=1$.
They sweep the extended origin disc, so their boundary tori
bound. Their orbit vectors on a regular fibre form the full
fixed lattice
\[
 \ker(A_*-I)=\mathbb Z\delta+\mathbb Z\ell,\qquad
 \ell=\begin{cases}(-1,2,0,4),&\mathrm{OS46},\\
                   (-3,2,0,4),&\mathrm{OS50}.
       \end{cases}
 \tag{D.T8}\label{maps:eq:D.T8}
\]
For an explicit check of the product marking, the integral
angular coordinate in \eqref{maps:eq:D.T4} is
$\delta^*-2e_1^*-e_2^*$ or
$\delta^*-6e_1^*-9e_2^*$ respectively. It is invariant
under \eqref{maps:eq:D.T6}, is one on $\delta$, and is zero on $\ell$.
The minimal nonzero $\gamma$-coordinate of a fixed vector
is four, exactly the period of the second action.
Consequently the two normal tori project isomorphically to
the full fixed lattice in the Wang sequence. They form a
primitive rank-two subgroup of $H_2(M_*)$. By \eqref{maps:eq:D.T5}--\eqref{maps:eq:D.T7}
the whole kernel has rank two, so they are its entire kernel.
It follows at once that the remaining image equals the
regular-fibre specialization image; its cokernel is \eqref{maps:eq:D.T7}.
This argument retains the star fibre's integral component
information. Replacing it by a product with an untwisted
elliptic surface, or by an order-two good-reduction filling,
would give a different boundary problem.

\item \emph{Global integral calculation for every affine vector.}
Insert $A_0=I$ before the four original monodromies. Set
$D=[A_0-I\ A_1-I\ A_2-I\ A_3-I]$, and use the direct
boundary maps $j_0,\ldots,j_3$ and
\[
 j_4(x)=-(A_1A_2A_3x,A_2A_3x,A_3x,x).
\]
The ten unchanged normal columns consist of two at the
star fibre, three at each of the two primitive nodal fibres,
and two at the cusp. Let their $16\times10$ matrix be $B$.
The exact constant integral certificates give
\[
 K=[B\ j_0(e_1)\ j_0(e_2)\ j_0(\gamma)],\qquad
 LK=I_{13},\qquad
 (Lj_0x)_{11,12,13}=(x_1,x_2,x_4).
 \tag{D.T9}\label{maps:eq:D.T9}
\]
The matrix $D$ has rank three with unit Smith entries;
$K$ is an integral basis of its kernel. The exterior
coinvariants are free of rank one, detected primitively by
$\omega=p e_1^*e_2^*-\delta^*\gamma^*$.
All ten normal tori are the marked fixed-loop sweeps proved
above and in \ref{maps:src:D.2.12}. The constant vertical form has zero
value on them. Thus \eqref{maps:eq:D.T9}, together with a fibre class of
$\omega$-value one, specifies the integral splitting,
without an unspecified fibre-area correction.

\item The original degree-one vanishing lattices generate the
primitive $\ker\gamma^*$. All original meridians bound.
Van Kampen therefore leaves only $\gamma$, and the regular
log imposes exactly $n\gamma=0$ as in \ref{maps:src:D.2.13}.
After eliminating the ten unchanged normals, the entire
remaining torsion presentation is
\[
 C=\begin{pmatrix}
 -pb'&pa&n&-c&0\\
 m&0&0&0&a\\
 0&m&0&0&b'\\
 0&0&0&m&n
 \end{pmatrix}.
 \tag{D.T10}\label{maps:eq:D.T10}
\]
Indeed the smooth-log lattice is
$\Lambda'=\Lambda+\mathbb Z(v/m)$, with primitive relation
$mg=v$; its four normal tori have coordinates
$(-\omega(x,v),j_0(mx))$, and the fifth lattice generator
has coordinates $(0,j_0v)$. This is the same integral
calculation as \ref{maps:src:D.2.13}, with $p=4$ or $12$.
Its all-parameter local-ring Smith proof gives entries
$(1,t,m,m|n|/t)$ for $n\ne0$, and $(1,t,m,0)$ for $n=0$.
The $b$ additional generators from \eqref{maps:eq:D.T7} are free and
split off, since their quotient is free.

\item The degree-one Mayer--Vietoris source and target are free
of the same rank nineteen. For $n\ne0$ their finite cokernel
has order $|n|$, so their kernel is zero. For $n=0$ the
cokernel and kernel each have free rank one; the latter
supplies the additional split free generator in $H_2$.
Thus $H_2$ is exactly as stated in \eqref{maps:eq:D.T1}--\eqref{maps:eq:D.T2}.
The original star filling has Euler zero by \eqref{maps:eq:D.T4}, and
the only nonzero Euler contribution is two from the cusp.
Integral duality, including its torsion terms, now gives
all the other displayed groups. In particular when $n=0$
one has $b_1=b_5=1$, $b_2=b_4=b+2$ and $b_3=2b+2$.

\item \emph{The exact smooth types.}
Assume $n=\pm1$. The free part of $H_2(X)$ is represented
in the original star filling. Its canonical line is trivial:
the canonical two-form on the minimal elliptic neighbourhood
is nonvanishing and translation-invariant, and its wedge
with $dw/w$ descends through the contraction. Thus $c_1$
evaluates trivially on every free generator. Since $X$ is
simply connected, $H^2(X;\mathbb Z)$ is torsion-free, so
$c_1(X)=0$ and $X$ is spin.

\item The global $\delta$ action is free on the original star
fill: in its contracting line model a nontrivial contraction
cannot equal a rotation. It is free on every other filling
except for the single cusp rational curve with normal
weights $(1,-1)$. This curve is in the cusp's fibre
specialization image, hence represents a multiple of the
regular class $s$. The normal relation $ns=0$ in \eqref{maps:eq:D.T10}
makes it zero. The equivariant nullsphere surgery of \ref{maps:src:D.2.14}
therefore gives a projectably spin free-circle manifold
$Y=X\#(S^3\times S^3)$.
Its simply connected spin orbit $B_5$ has
\[
 H_2(B_5)=\mathbb Z^{b+1}\oplus(\mathbb Z/m)^2,
\]
with primitive Euler class. Smale--Barden classifies the
oriented pair as
\[
 B_5=\#^{b+1}(S^2\times S^3)\#N_m,
 \qquad e=(x,0,\ldots,0).
\]
Apply \cite[Theorem A]{manGR25} successively, keeping the
primitive Euler class on the first $S^2\times S^3$ summand.
Its total space is $S^3\times S^3$. Each additional
$S^2\times S^3$ contributes its twisted suspension, which
is $(S^2\times S^4)\#(S^3\times S^3)$ by
\cite[Theorem B(2)]{manGR25}; $N_m$ contributes
$\widetilde\Sigma_0N_m$. Cancel the common $S^3\times S^3$
using the full Zhubr invariants exactly as in \ref{maps:src:D.2.14}.
This proves \eqref{maps:eq:D.T3}, including even $m$ and its framing.

\item The program \texttt{manifolds\_star\_smooth.py} retains all
constant marked matrices and integer left inverses,
checks \eqref{maps:eq:D.T6}--\eqref{maps:eq:D.T10}, and verifies 1,462 full $14\times15$
Smith forms, including 194 cases with $n=0$. The unrestricted
conclusion follows from the geometry and symbolic Smith
proof, not a bound on those parameters.
\end{enumerate}
\end{proof}

\begin{corollary}[two further original additive profiles]\label{maps:src:D.2.17}
\StatementRecord{maps:src:D.2.17}{proved}{Historical events: Claude Fable 5.2 (Anthropic), effort high, 16 November 2026: carried as background without fault found. GPT-6 Astra (OpenAI), ChatGPT Pro mode, effort Pro, S6\_09\_PAPERS, 17 November 2026: targeted proof check. Also checked in this writing. Pre-continuation writing events: GPT-6 Astra (OpenAI), ChatGPT Pro mode, effort Pro, S6\_10\_PAPERS, 18 November 2026: targeted proof check. Continuation metadata events: GPT-6 Astra (OpenAI), ChatGPT Pro mode, effort Pro, S6\_10\_PAPERS continuation, 19 November 2026: used and rechecked in this run; status-vocabulary reconciliation; metadata reconciliation only; no new mathematical check. \emph{Pre-continuation recorded proof scope.} GPT-6 Astra (OpenAI, ChatGPT Pro mode), effort Pro, 2026-11-17 UTC; Claude Fable 5.2 (Anthropic), effort high, 16 November 2026; also checked in this writing by GPT-6 Astra (OpenAI), ChatGPT Pro mode, effort Pro, 18 November 2026}{the hypotheses, notation and definitions stated here; \cref{maps:eq:D.T1}; \cref{maps:eq:D.T3}; \cref{maps:src:D.2.16}; \cref{maps:eq:D.T10}}


The complete formulas and oriented diffeomorphism classification
\eqref{maps:eq:D.T1}--\eqref{maps:eq:D.T3}, for every free smooth-fibre log, also hold for
\[
\begin{array}{c|c|l|r|r}
 \mathrm{OS}&\mathrm{profile}&\text{ordered fibres}&p&b\\\hline
 46&II&I_3^*,II,I_1^\bullet&4&2\\
 50&IV&I_1^*,IV,I_1^\bullet&12&1.
\end{array}
\]
Use $P=pG$, $Q=G$, the simple linearization zero shown,
and the original fillings at both additive fibres.
\end{corollary}


\begin{proof}
\begin{enumerate}[label=\arabic*.,leftmargin=*]
\item The section heights and component residues are the same at
the star fibre as in \ref{maps:src:D.2.16}. At $II$, both sections are
narrow. At $IV$, $P$ is narrow and $Q$ has residue two,
a generator of its order-three component group.
The mixed height is again one. The same zero-multidegree
argument gives the required degree-one linearization, with
no zero at either additive fibre.

\item For $II$, contract both section logarithms to zero. The
marked original filling is $E_{II}\times T^2$, where
$E_{II}$ retracts onto its cuspidal rational curve,
topologically $S^2$. For $IV$, divide $3Q$ in the additive
identity component as in \ref{maps:src:D.2.16}. The resulting section
$Q_0$ has order three and cycles all three components of
$IV$. Its central curve is three spheres meeting at one
point, so $H_2(E_{IV})=\mathbb Z^3$ and $H_1(E_{IV})=0$.
The marked fill is
$S^1_\delta\times T(t_{Q_0}:E_{IV}\to E_{IV})$.
The local logarithm, translation extension and boundary
rigidification arguments apply verbatim. In either case
there is one component orbit, and hence
\[
 \pi_1(U)=\mathbb Z^2,\qquad
 (H_0,H_1,H_2,H_3,H_4)(U)
       =(\mathbb Z,\mathbb Z^2,\mathbb Z^2,\mathbb Z^2,\mathbb Z).
\]

\item The two additive monodromies in the respective global
markings are
\[
 A_{II}=\begin{pmatrix}
 1&1&0&0\\-1&0&0&0\\0&0&1&0\\0&0&0&1
 \end{pmatrix},\qquad
 A_{IV}=\begin{pmatrix}
 -1&1&0&-2\\-1&0&0&0\\-24&0&1&-16\\0&0&0&1
 \end{pmatrix}.
\]
Both $A-I$ and $\bigwedge^2A-I$ have only unit nonzero
Smith entries, of ranks two and four. Thus their boundaries
have $H_1=\mathbb Z^3$, $H_2=\mathbb Z^4$. The two circle
actions sweep the origin disc and give the complete
primitive normal-torus kernel, with fixed vectors
\[
 \delta,\gamma\quad\text{at }II;\qquad
 \delta,(-2,2,0,3)\quad\text{at }IV.
\]
The second vector has minimal positive $\gamma$-coordinate
three, exactly the period of the order-three translation.
The relative sequence, with $H_2(U,\partial U)=\mathbb Z$
mapping onto the meridian kernel, shows that the boundary
specialization is onto $H_2(U)$. Thus these additional
original fillings contribute no free defect or torsion.

\item There are now four punctures including the regular log.
The graph differential has domain $\mathbb Z^{12}$ and
rank three with primitive image. The six unchanged normal
columns, together with $j_0(e_1),j_0(e_2),j_0(\gamma)$,
form a primitive basis of its rank-nine kernel. Its integer
left inverse has the same residual projection
$(x_1,x_2,x_4)$. Exterior coinvariants are still detected
primitively by $p e_1^*e_2^*-\delta^*\gamma^*$.
Consequently the entire residual normal matrix is \eqref{maps:eq:D.T10},
with the same free star defect $\mathbb Z^b$.
The degree-one Mayer--Vietoris source and target now have
rank fourteen; their kernel is zero for $n\ne0$ and is
$\mathbb Z$ for $n=0$. The full homology therefore follows
by the identical unrestricted Smith proof and duality.

\item Finally the global $\delta$ action remains free on both
original additive fillings. Its unique fixed curve is
still the null-homologous cusp curve for $n=\pm1$.
The free $H_2$ remains represented by the star neighbourhood,
where the canonical line is trivial. Thus the projectably
spin surgery, Smale--Barden orbit classification, repeated
Galaz-Garc\'ia--Reiser Theorem A and Zhubr cancellation in
\ref{maps:src:D.2.16} apply without modification and give \eqref{maps:eq:D.T3}.

\item The independent program
\texttt{manifolds\_star\_additive.py} records the full
constant integral certificates and 1,462 full $10\times11$
Smith checks, including 194 cases with $n=0$. The proof is
unrestricted in the affine vector and multiplicity.
\end{enumerate}
\end{proof}

\begin{theorem}[the remaining profile with a three-component cusp]\label{maps:src:D.2.18}
\StatementRecord{maps:src:D.2.18}{proved}{Historical events: Claude Fable 5.2 (Anthropic), effort high, 16 November 2026: carried as background without fault found. GPT-6 Astra (OpenAI), ChatGPT Pro mode, effort Pro, S6\_09\_PAPERS, 17 November 2026: targeted proof check. Also checked in this writing. Pre-continuation writing events: GPT-6 Astra (OpenAI), ChatGPT Pro mode, effort Pro, S6\_10\_PAPERS, 18 November 2026: targeted proof check. Continuation metadata events: GPT-6 Astra (OpenAI), ChatGPT Pro mode, effort Pro, S6\_10\_PAPERS continuation, 19 November 2026: used and rechecked in this run; status-vocabulary reconciliation; metadata reconciliation only; no new mathematical check. \emph{Pre-continuation recorded proof scope.} GPT-6 Astra (OpenAI, ChatGPT Pro mode), effort Pro, 2026-11-17 UTC; Claude Fable 5.2 (Anthropic), effort high, 16 November 2026; also checked in this writing by GPT-6 Astra (OpenAI), ChatGPT Pro mode, effort Pro, 18 November 2026}{the hypotheses, notation and definitions stated here; \cref{maps:src:D.2.16}; \cref{maps:src:D.2.15}; \cref{maps:src:D.1.2}; \cref{maps:src:D.2.17}; \cref{maps:eq:D.V5}; \cref{maps:eq:D.T10}; \cref{maps:eq:D.V4}; \cref{maps:src:D.2.14}}


For OS50 profile $II$, with ordered fibres $II,I_3,I_1^*$,
take $P=12G$, $Q=G$ and the degree-one linearization vanishing
on the whole $I_3$ fibre. Keep the original $II$ and $I_1^*$
fillings. At an additional regular fibre perform the free
order-$m$ log with vector $(a,b',c,n)$, where
$\gcd(m,a,b',c,n)=1$. Define $t$ and $G$ as in \ref{maps:src:D.2.16} with
$p=12$. For $n\ne0$ one has
\[
\begin{aligned}
 \pi_1(X)&=H_1(X)=\mathbb Z/|n|, &H_5(X)&=0,\\
 H_2(X)&=\mathbb Z^3\oplus G,
 &H_3(X)&=\mathbb Z^2\oplus G,\\
 H_4(X)&=\mathbb Z^3\oplus\mathbb Z/|n|.
\end{aligned}
\tag{D.V1}\label{maps:eq:D.V1}
\]
For $n=0$, put $T=\mathbb Z/t\oplus\mathbb Z/m$; the answer is
\[
\begin{aligned}
 \pi_1(X)&=H_1(X)=H_5(X)=\mathbb Z,\\
 H_2(X)&=\mathbb Z^5\oplus T,
 &H_3(X)&=\mathbb Z^4\oplus T,\\
 H_4(X)&=\mathbb Z^5.
\end{aligned}
\tag{D.V2}\label{maps:eq:D.V2}
\]
In both cases $H_0=H_6=\mathbb Z$ and $\chi(X)=6$.
Every simply connected member, exactly $n=\pm1$, has
\[
 X\cong \#^3(S^2\times S^4)\#(S^3\times S^3)
                          \#\widetilde\Sigma_0N_m
 \tag{D.V3}\label{maps:eq:D.V3}
\]
as an oriented smooth manifold.
\end{theorem}


\begin{proof}
\begin{enumerate}[label=\arabic*.,leftmargin=*]
\item \emph{The actual three-component cusp.}
The generator height is $1/12$. The section $P$ is narrow
at every fibre, while $Q$ has component residue two at
$I_3$ and is a generator at $I_1^*$. The mixed height is
one, so the same multidegree and vertical-line argument
as in \ref{maps:src:D.2.16} gives the required linearization.
The actual valuation periods at its simple zero are
\[
 \Gamma=\langle(3,0),(2,1)\rangle,
 \qquad
 \begin{pmatrix}1&1\\0&-1\end{pmatrix}
 \begin{pmatrix}3&2\\0&1\end{pmatrix}
       =\begin{pmatrix}3&3\\0&-1\end{pmatrix}.
 \tag{D.V4}\label{maps:eq:D.V4}
\]
Thus the same integral symmetry of the triangular fan as
in \ref{maps:src:D.2.15} identifies this lattice with
$3\mathbb Z\oplus\mathbb Z$. The smooth Mumford filling
has three normalized $dP_6$ components. The integral
normalization calculation of \ref{maps:src:D.1.2} gives
\[
 (H_0,H_1,H_2,H_3,H_4)(U_c)
  =(\mathbb Z,\mathbb Z^2,\mathbb Z^6,\mathbb Z^2,\mathbb Z^3),
 \qquad \pi_1(U_c)=\mathbb Z^2.
\]
The original $I_1^*$ filling has the order-four translation
model of \ref{maps:src:D.2.16}, with two component orbits and free
specialization defect one. The original $II$ filling has
the product model of \ref{maps:src:D.2.17} and no defect.

\item The cusp monodromy in the global marking is
\[
 A_c=\begin{pmatrix}
 7&12&0&-8\\-3&-5&0&4\\48&96&1&-63\\0&0&0&1
 \end{pmatrix}.
\]
The nonzero Smith entries of $A_c-I$ are $(1,3)$,
whereas those of $\bigwedge^2A_c-I$ are $(1,1)$.
Thus its boundary $M_c$ has
$H_1(M_c)=\mathbb Z^3\oplus\mathbb Z/3$ and
$H_2(M_c)=\mathbb Z^6$.
Its $H_1$ map kills precisely the meridian and the
$\mathbb Z/3$ angular torsion, and is onto $\mathbb Z^2$.
The relative sequence with
$H_2(U_c,M_c)=H^4(U_c)=\mathbb Z^3$ therefore gives
\[
 H_2(U_c)/\operatorname{im}H_2(M_c)=\mathbb Z^2.
 \tag{D.V5}\label{maps:eq:D.V5}
\]
Indeed this quotient is the kernel of a surjection from
$\mathbb Z^3$ to $\mathbb Z\oplus\mathbb Z/3$, hence is
free of rank two. The two toric angular normal tori bound
and project onto the full primitive fixed lattice of
$A_c$. The kernel of the degree-two boundary map has
rank two by \eqref{maps:eq:D.V5}, so these are its whole kernel.
The image is precisely the regular-fibre specialization
image. The total free local defect is consequently
$\mathbb Z^2\oplus\mathbb Z=\mathbb Z^3$.

\item \emph{The integral global map and its three-torsion.}
Insert $A_0=I$ before the three original monodromies.
The global fibre-coinvariant presentation has nonzero
Smith entries $(1,1,3)$. The cusp's boundary torsion maps
to this global $\mathbb Z/3$ by a unit; the exact integral
vector and its image are retained in the certificate.
Thus these torsion summands cancel isomorphically in
Mayer--Vietoris. They must not be discarded beforehand.
Replacing the cusp vanishing image by its actual saturated
angular lattice gives three unit Smith entries and kills
precisely $\ker\gamma^*$. All original meridians bound,
and van Kampen leaves exactly $n\gamma=0$.

\item The six unchanged normal tori consist of two from each
original filling. Their graph columns together with
$j_0(e_1),j_0(e_2),j_0(\gamma)$ form a primitive integral
basis of the rank-nine graph kernel. Its left inverse
has residual projection $(x_1,x_2,x_4)$, including zero
for the $\delta$ sweep. The exterior coinvariants are
free of rank one with primitive detector
$12e_1^*e_2^*-\delta^*\gamma^*$.
The actual fixed-loop sweeps specify the same integral
zero-area splitting as in \ref{maps:src:D.2.16}. Consequently after the
six unchanged normals are eliminated, the residual
regular-log presentation is exactly \eqref{maps:eq:D.T10} with $p=12$.
Its all-parameter Smith calculation supplies $G$ for
$n\ne0$, and $\mathbb Z\oplus T$ for $n=0$.
The free local defect \eqref{maps:eq:D.V5} and its star counterpart
split off. After cancelling the isomorphic degree-one
$\mathbb Z/3$ summands, the remaining source and target
have equal free rank fourteen. Their kernel is zero for
$n\ne0$ and is free of rank one for $n=0$, giving the
additional split generator in \eqref{maps:eq:D.V2}.
Only the three-component cusp contributes Euler
characteristic, namely six. Duality now proves all of
\eqref{maps:eq:D.V1}--\eqref{maps:eq:D.V2}, with the indicated distinct free ranks
in degree three.

\item \emph{Fixed-curve classes and smooth recognition.}
For $n=\pm1$, the free $H_2$ is generated by classes in
the cusp and star neighbourhoods. Both canonical lines
are locally trivial, by the Mumford volume form and
the translation-invariant form used in \ref{maps:src:D.2.16}.
Hence $c_1(X)=0$ and $X$ is spin.

\item There are three $\delta$-fixed rational curves, represented
by vertical edges of the triangular fan modulo $\Gamma$.
Their classes in the free cusp defect are zero, as the
following integral computation shows. Colour a vertex
$(a,b)$ by $a+b$ modulo three; this labels its global
component by \eqref{maps:eq:D.V4}. A vertical edge has endpoints
$(a,b),(a,b+1)$ and the two adjacent third vertices
$(a+1,b),(a-1,b+1)$. The colours of the endpoints are
$(c,c+1)$; those of the third vertices are $(c+1,c)$.
The toric wall formula therefore gives degree zero
against each global component divisor. Poincare--Lefschetz
identifies these degrees with the map
$H_2(U_c)\to H_2(U_c,M_c)=\mathbb Z^3$.
Each fixed curve is consequently in the boundary image,
which equals the regular-fibre image. It represents a
multiple of the regular class $s$, killed by the relation
$ns=0$. All three fixed curves are therefore null-homologous
in $X$, and hence null-homotopic. This checks their actual
integral classes, rather than inferring nullity from their
normal bundles.

\item Each normal bundle has the equivariant model
$\mathcal O(-1)_{+1}\oplus\mathcal O(-1)_{-1}$.
The three-component embedded link is isotopic to three
unknotted spheres: apply the dimension-six embedding
isotopy theorem of \ref{maps:src:D.2.14} to the disjoint union.
The same orientation-preserving equivariant surgeries
produce
$Y=X\#^3(S^3\times S^3)$ with a projectably spin free
circle action. Its simply connected spin orbit has
$H_2(B_5)=\mathbb Z^4\oplus(\mathbb Z/m)^2$ and primitive
Euler class. Thus Smale--Barden gives the oriented pair
\[
 (B_5,e)=\bigl(\#^4(S^2\times S^3)\#N_m,
                                      (x,0,0,0,0)\bigr).
\]
Repeated Galaz-Garc\'ia--Reiser Theorem A and Theorem B(2)
identify its total space with
$\#^3(S^2\times S^4)\#^4(S^3\times S^3)
\#\widetilde\Sigma_0N_m$.
Full Zhubr cancellation of three copies of $S^3\times S^3$
proves \eqref{maps:eq:D.V3}.

\item The program \texttt{manifolds\_star\_i3\_cusp.py} checks
the integral period lattice, all three fixed-curve degree
vectors, the boundary torsion map, saturated vanishing
lattice, full graph/exterior certificates, and 731 full
$10\times11$ Smith forms, including 97 cases with $n=0$.
The geometry and symbolic reduction prove the result
for every admissible affine vector and multiplicity.
\end{enumerate}
\end{proof}

\begin{proposition}[finite initial data for the whole-fibre problem]\label{maps:src:D.2.19}
\StatementRecord{maps:src:D.2.19}{proved}{Historical events: Claude Fable 5.2 (Anthropic), effort high, 16 November 2026: carried as background without fault found. GPT-6 Astra (OpenAI), ChatGPT Pro mode, effort Pro, S6\_09\_PAPERS, 17 November 2026: targeted proof check. Also checked in this writing. Pre-continuation writing events: GPT-6 Astra (OpenAI), ChatGPT Pro mode, effort Pro, S6\_10\_PAPERS, 18 November 2026: targeted proof check. Continuation metadata events: GPT-6 Astra (OpenAI), ChatGPT Pro mode, effort Pro, S6\_10\_PAPERS continuation, 19 November 2026: used and rechecked in this run; status-vocabulary reconciliation; metadata reconciliation only; no new mathematical check. \emph{Pre-continuation recorded proof scope.} GPT-6 Astra (OpenAI, ChatGPT Pro mode), effort Pro, 2026-11-17 UTC; Claude Fable 5.2 (Anthropic), effort high, 16 November 2026; also checked in this writing by GPT-6 Astra (OpenAI), ChatGPT Pro mode, effort Pro, 18 November 2026}{the hypotheses, notation and definitions stated here; \cref{maps:src:D.2.1}; \cref{maps:eq:D.F2}; \cite[Section 11.8]{manSS10}; \cref{maps:src:D.2.11}}


Restrict to the uncorrected line
$M=\mathcal O(P-O)$ and require its linearization line
$J=M\otimes t_Q^*M^{-1}$ to be $f^*\mathcal O_{\mathbb P^1}(1)$.
This is the case of a simple zero on one whole fibre.
Write $P=aG$, $Q=bG$ and $h(G)=1/d$.
Up to simultaneous sign reversal, the possible ordered
pairs are exactly the positive pairs $ab=d$ for which,
at every reducible fibre, $P$ or $Q$ is narrow there.
For the marked profiles of \ref{maps:src:D.2.1} they are
\[
\begin{array}{c|r|l}
 \mathrm{OS}&d& (a,b)\\\hline
 43&2&(1,2),(2,1)\\
 45&8&(1,8),(8,1)\\
 46&4&(1,4),(4,1)\\
 47&14&(1,14),(2,7),(7,2),(14,1)\\
 49&6&(1,6),(2,3),(3,2),(6,1)\\
 50&12&(1,12),(3,4),(4,3),(12,1)\\
 55&20&(1,20),(4,5),(5,4),(20,1)\\
 56&30&(1,30),(2,15),(3,10),(5,6),\\
 &&(6,5),(10,3),(15,2),(30,1).
\end{array}
\tag{D.F1}\label{maps:eq:D.F1}
\]
The same row applies to every profile of its OS entry.
Across all 21 marked profiles there are 104 positive
ordered factor-pair candidates, exactly 88 of which
satisfy the stated requirement. Including each marked
multiplicative fibre and one regular-fibre position for
the simple zero gives 348 finite initial data. All their
marked integral monodromy tuples are recorded exactly.
These counts classify initial data in the stated scope,
not diffeomorphism types or all logarithmic fillings.
\end{proposition}


\begin{proof}
\begin{enumerate}[label=\arabic*.,leftmargin=*]
\item On a smooth elliptic fibre, translation acts trivially
on $\operatorname{Pic}^0$, so $J$ is generically trivial
and hence is represented by a vertical divisor.
At a reducible fibre, label its multiplicity-one
components by the Neron component group $\Phi$.
The multidegree of $M$ on those components is
$e_{\bar P}-e_0$. Translation by $Q$ acts regularly
on this component group, and therefore its pullback
preserves that vector if and only if
\[
                         \bar P=0\quad\hbox{or}\quad\bar Q=0.
 \tag{D.F2}\label{maps:eq:D.F2}
\]
Indeed, if $\bar P\ne0$, its unique positive and unique
negative entries must both be fixed; the regular
translation fixes a component only when $\bar Q=0$.
If $P$ is narrow, every component degree is zero.
If $Q$ is narrow, its translation fixes every component
(the connected identity group cannot permute a finite
set). Thus \eqref{maps:eq:D.F2} is also sufficient for all multidegrees
of $J$ to vanish, including the multiple components.

\item A vertical divisor of zero multidegree is an integral
combination of whole fibres: each fibre intersection
matrix has radical generated by its multiplicity vector,
and that vector has an entry one because the surface has
a section. Hence precisely under \eqref{maps:eq:D.F2}, $J=f^*L$ for
a line $L$ on the base. Intersecting with $O$ gives
\[
 \deg L=P\cdot O+1+Q\cdot O-P\cdot Q
        =\langle P,Q\rangle=ab/d.
\]
For the first equality use $O^2=-1$ and translation
invariance of intersection, so
$(P-Q)\cdot O=P\cdot Q$ and $(-Q)\cdot O=Q\cdot O$.
For the second, every local correction in the height
pairing vanishes under \eqref{maps:eq:D.F2}; see
\cite[Section 11.8]{manSS10}.
Thus $J=f^*\mathcal O(1)$ is equivalent to \eqref{maps:eq:D.F2} and
$ab=d$. Inversion on the elliptic surface takes
$(P,Q,M,J)$ to $(-P,-Q,[-1]^*M,[-1]^*J)$ and preserves
complex orientation. We may consequently take $a,b>0$,
without identifying the two ordered factors.

\item The finite calculation now needs no parameter bound.
For each of the 21 profiles, list the positive divisors
$a\mid d$, put $b=d/a$, and evaluate both section
component maps. Applying \eqref{maps:eq:D.F2} gives exactly \eqref{maps:eq:D.F1}.
The 16 excluded marked pairs fail the requirement for
this particular uncorrected line; this is not a
nonexistence theorem for other line bundles.
For each admitted pair, place the single weight one
at each multiplicative site in turn and at one
additional regular site. The exact program
\texttt{manifolds\_degree\_one.py} gives 348 tuples,
checks every matrix is integral with determinant one,
and verifies every ordered monodromy product is $I$.
It retains all component data and all matrices, as well
as the additive sites and their finite-part orders.

\item There are precise boundaries to this reduction.
Component-supported linearizations are not covered:
\ref{maps:src:D.2.11} already has $\langle G,3G\rangle=1/2$ and
$J=\mathcal O(C_0)$, rather than a whole-fibre zero.
Adding an integral vertical divisor to $M$ changes its
multidegrees and can alter \eqref{maps:eq:D.F2}.
Nor does the finite list bound the logarithmic orders
or affine clutching vectors, choose a compactification
at an additive zero, or classify all Mumford subdivisions.
These remain geometric and integral-attachment choices
in the full D.2 problem. This proposition provides an
exhaustive finite set of starting data only under its
explicit whole-fibre and uncorrected-line hypotheses.
\end{enumerate}
\end{proof}

\begin{theorem}[six original-filled families, including the untwisted Engel case]\label{maps:src:D.2.20}
\StatementRecord{maps:src:D.2.20}{proved}{Historical events: Claude Fable 5.2 (Anthropic), effort high, 16 November 2026: carried as background without fault found. GPT-6 Astra (OpenAI), ChatGPT Pro mode, effort Pro, S6\_09\_PAPERS, 17 November 2026: targeted proof check. Also checked in this writing. Pre-continuation writing events: GPT-6 Astra (OpenAI), ChatGPT Pro mode, effort Pro, S6\_10\_PAPERS, 18 November 2026: targeted proof check. Continuation metadata events: GPT-6 Astra (OpenAI), ChatGPT Pro mode, effort Pro, S6\_10\_PAPERS continuation, 19 November 2026: used and rechecked in this run; status-vocabulary reconciliation; metadata reconciliation only; no new mathematical check. \emph{Pre-continuation recorded proof scope.} GPT-6 Astra (OpenAI, ChatGPT Pro mode), effort Pro, 2026-11-17 UTC; Claude Fable 5.2 (Anthropic), effort high, 16 November 2026; also checked in this writing by GPT-6 Astra (OpenAI), ChatGPT Pro mode, effort Pro, 18 November 2026}{the hypotheses, notation and definitions stated here; \cref{maps:src:D.2.19}; \cref{maps:src:D.2.16}; \cref{maps:eq:D.O4}; \cref{maps:eq:D.O5}; \cref{maps:eq:D.O8}; \cref{maps:src:D.3.3}; \cref{maps:eq:D.O7}; \cref{maps:src:D.2.14}; \cite[Theorem A and Theorem B(2)]{manGR25}}


Use the following configurations and section choices:
\[
\begin{array}{c|c|l|r}
 \mathrm{OS}&\mathrm{profile}&\text{fibres}&d\\\hline
 43&\mathrm{default}&III^*,3I_1&2\\
 43&II&III^*,II,I_1&2\\
 46&\mathrm{default}&I_3^*,3I_1&4\\
 46&II&I_3^*,II,I_1&4\\
 49&III&IV^*,III,I_1&6\\
 50&IV&I_1^*,IV,I_1&12.
\end{array}
\]
Take $P=G$, $Q=dG$, retain every original additive filling,
and put the simple linearization zero at the first marked
$I_1$. The other $I_1$ fillings are unchanged. At an
additional regular fibre perform any free log
\[
 m>0,\qquad v=a e_1+b e_2+c\delta+n\gamma,
 \qquad\gcd(m,a,b,c,n)=1.
\]
Set
\[
        g=\gcd(m,a,b,n),\qquad
        T_{m,g}=\mathbb Z/g\oplus(\mathbb Z/m)^2.
\]
For $n\ne0$, the answer is
\[
\begin{aligned}
 \pi_1(X)&=H_1(X)=\mathbb Z/|n|, &H_5(X)&=0,\\
 H_2(X)&=\mathbb Z^7\oplus T_{m,g},
 &H_3(X)&=\mathbb Z^{14}\oplus T_{m,g},\\
 H_4(X)&=\mathbb Z^7\oplus\mathbb Z/|n|.
\end{aligned}
\tag{D.O1}\label{maps:eq:D.O1}
\]
For $n=0$, it is
\[
\begin{aligned}
 \pi_1(X)&=H_1(X)=H_5(X)=\mathbb Z,\\
 H_2(X)&=\mathbb Z^8\oplus T_{m,g},
 &H_3(X)&=\mathbb Z^{14}\oplus T_{m,g},\\
 H_4(X)&=\mathbb Z^8.
\end{aligned}
\tag{D.O2}\label{maps:eq:D.O2}
\]
In all cases $H_0=H_6=\mathbb Z$ and $\chi(X)=2$.
The simply connected members are exactly $n=\pm1$, and
all of them have oriented diffeomorphism type
\[
 X\cong\#^7(S^2\times S^4)\#^7(S^3\times S^3)
                                    \#\widetilde\Sigma_0N_m.
 \tag{D.O3}\label{maps:eq:D.O3}
\]
In particular none is a six-sphere. For the Engel row,
the trivial regular log $(m,v)=(1,0)$ recovers the
untwisted groups of \cref{push:src:A.4.2} $(1,1,8,14,8,1,1)$ integrally,
with no torsion.
\end{theorem}


\begin{proof}
\begin{enumerate}[label=\arabic*.,leftmargin=*]
\item \emph{The original additive bundles.}
The section $Q$ is narrow everywhere. At every additive
fibre except $II$, $P$ meets a nonidentity component.
Proposition \ref{maps:src:D.2.19} supplies the degree-one linearization.
The zero on $I_1$ gives the same primitive hexagonal
cusp as before. This proof retains the original additive
bundles, including their nonzero component degrees.

\item Let $E\to\Delta$ be one such additive elliptic-surface
neighbourhood, with $r$ components. It is simply connected
and retracts onto a connected union of $r$ rational spheres
with $H_2(E)=\mathbb Z^r$ and no other positive homology.
This applies also to the tangent pair of type $III$ and
the three spheres of type $IV$, which each meet at one
point. The component degree vector of $M=\mathcal O(P-O)$
is the primitive vector
\[
                   c=e_{\bar P}-e_0\in H^2(E;\mathbb Z).
\]
Because $Q$ is narrow, its local logarithm contracts it
to zero through translations. Rigidify the line along
$O$ as in \ref{maps:src:D.2.16}. The unit in the lift has a logarithm
since $H^1(E)=0$, and can be straightened. The actual
original filling is therefore, with its transported
boundary marking, the principal $T^2$ bundle over $E$
whose two Chern classes are $(c,0)$. Its two circle
directions are denoted $\delta,\gamma$.

\item The integral bundle spectral sequence has only base
rows zero and two. Its nonzero differentials are
\[
 H_2(E)\longrightarrow H_1(T^2),\quad x\longmapsto c(x)\delta,
 \qquad
 H_2(E)\otimes H_1(T^2)\longrightarrow H_2(T^2),\quad
 x\otimes y\longmapsto c(x)\delta\wedge y.
\]
Both images are primitive because $c$ is primitive.
There are no higher differentials, and all surviving
groups and extensions are free. Consequently
\[
 \pi_1(U)=\mathbb Z\gamma,\qquad
 (H_0,H_1,H_2,H_3,H_4)(U)
  =(\mathbb Z,\mathbb Z,\mathbb Z^{r-1},
                             \mathbb Z^{2r-1},\mathbb Z^r).
 \tag{D.O4}\label{maps:eq:D.O4}
\]
In particular $H_2(U)\to H_2(E)$ identifies its image
with $\ker c$.

\item \emph{The marked fibre kernel.}
Write the upper-left elliptic monodromy block as $B_i$
and its $\gamma$ column as $q_i$. Local narrowness of $Q$
means that
\[
                 q_i=(B_i-I)z_i,\qquad z_i\in\mathbb Z^2.
\]
The logarithm deformation identifies the elliptic
projection of a regular four-torus by
$e_1\mapsto e_1$, $e_2\mapsto e_2$, $\delta\mapsto0$,
$\gamma\mapsto z_i$. Thus its area covector, in the
usual ordered exterior basis, is
\[
           w_i=(1,0,(z_i)_2,0,-(z_i)_1,0).
 \tag{D.O5}\label{maps:eq:D.O5}
\]
The specialization of the elliptic fibre in $H_2(E)$
is its component-multiplicity vector. This vector is
primitive, lies in $\ker c$, and generates the image
of $H_2(F)\to H_2(U)$. The latter map is exactly
$w_i$ times that primitive vector, by \eqref{maps:eq:D.O4}.
Therefore its kernel is the full integral $\ker w_i$.

\item The boundary monodromies in these six cases have two
unit nonzero Smith entries in $A_i-I$ and four in
$\bigwedge^2A_i-I$. Hence $H_1(\partial U)=\mathbb Z^3$
and $H_2(\partial U)=\mathbb Z^4$.
The two principal circle actions sweep the origin
disc, giving normal tori with orbit vectors
$\delta$ and $(-z_i,0,1)$. They form the full fixed
lattice of $A_i$. Thus these two sweeps split the
Wang projection integrally, and the remaining
boundary image equals the regular-fibre image just
computed. It follows that the local free defect is
\[
                H_2(U)/\operatorname{im}H_2(\partial U)
                         =\mathbb Z^{r-2}.
 \tag{D.O6}\label{maps:eq:D.O6}
\]
The full boundary kernel consists of the two actual
normal tori and the fibre kernel \eqref{maps:eq:D.O5}. This is why
using only the fixed-loop normal tori would lose
relations in this case.

\item The exterior coinvariants of the global regular family
are free of rank one, with primitive detector
$\omega=e_1^*e_2^*-d\delta^*\gamma^*$.
If $s$ has $\omega$-value one, the image of $\ker w_i$
in this group is $h_i\mathbb Zs$, where
\[
                  h_i=\gcd(d,(z_i)_1,(z_i)_2).
\]
Indeed an integral basis of $\ker w_i$ has $\omega$
values $0,-(z_i)_2,0,(z_i)_1,-d$.
Here all local inputs are explicit:
\[
\begin{array}{c|c|r|r|r}
 \mathrm{OS/profile}&\text{non-narrow-}P\text{ fibre}
       &r&z_i&h_i\\\hline
 43/\mathrm{default},43/II&III^*&8&(3,1)&1\\
 46/\mathrm{default},46/II&I_3^*&8&(1,-2)&1\\
 49/III&IV^*&7&(6,2)&2\\
       &III&2&(9,3)&3\\
 50/IV&I_1^*&6&(9,-6)&3\\
      &IV&3&(8,-8)&4.
\end{array}
\tag{D.O7}\label{maps:eq:D.O7}
\]
In every row of the theorem the $h_i$ together have gcd
one. Thus the actual original fillings kill $s$.
The original $II$ and ordinary $I_1$ pieces have the
product models already proved and add no defect.
The total free local defect $\beta$ is six in the first
four configurations and five in the last two.

\item \emph{The unrestricted global matrices.}
Insert $A_0=I$ for the regular log. If the total number
of punctures is $N$, the graph differential has domain
$\mathbb Z^{4(N-1)}$, rank three and primitive image.
Write $B$ for all unchanged fixed-loop normal columns,
including the two at the cusp. The exact constant
certificates give an integral basis
\[
 K=[B\ j_0(e_1)\ j_0(e_2)\ j_0(\gamma)]
\]
of the graph kernel, with left inverse $L$ whose last
three coordinates on $j_0(x)$ are $(x_1,x_2,x_4)$.
The actual circle sweeps above fix their zero fibre-area
coordinates. After these normal columns and the relations
$h_is=0$ are eliminated, the regular-family contribution
is freely generated by those three sweeps.

\item The smooth log has lattice $\Lambda+\mathbb Z(v/m)$,
with primitive relation $m(v/m)=v$.
Its normal tori have graph components $j_0(mx)$;
the additional lattice generator has component $j_0v$.
The fibre components no longer matter, since $s=0$.
The resulting presentation is exactly
\[
  \begin{pmatrix}
    m&0&0&a\\0&m&0&b\\0&0&m&n
  \end{pmatrix}.
 \tag{D.O8}\label{maps:eq:D.O8}
\]
Its determinantal divisors are $g,mg,m^2g$,
so its cokernel is $T_{m,g}$ for all integers in the
statement. Notice that $c$ enters the geometric
freeness condition even though it drops out of \eqref{maps:eq:D.O8}.

\item The original fillings kill all base meridians and
precisely the primitive $\ker\gamma^*$ in the fibre
fundamental group. Their complement over the punctured
base therefore has fundamental group $\mathbb Z\gamma$.
The new meridian relation is $m\xi-v=0$, and the base
relation kills $\xi$, leaving $n\gamma=0$. This proves
the asserted full fundamental group.

\item The degree-one Mayer--Vietoris map must now be counted
with the groups \eqref{maps:eq:D.O4}, not with the rank-two groups of
the narrow-$P$ filling. Its source and target are free.
For the first and third rows of the theorem their ranks
are $(19,18)$; for the second and fourth, $(14,13)$;
for the last two, $(14,12)$.
Its cokernel is the abelian group already computed by
van Kampen. Thus for $n\ne0$ its kernel rank is
$\varepsilon=1$ in the first four rows and two in the
last two. For $n=0$ this rank increases by one.
These free kernels split off in the degree-two sequence.
Since $\beta+\varepsilon=7$, this proves exactly the
stated $H_2$ groups, including their torsion.
All local pieces except the primitive cusp have Euler
zero, while the cusp contributes two. Integral duality
then gives \eqref{maps:eq:D.O1}--\eqref{maps:eq:D.O2} in every degree. In particular
the Engel row at $(m,v)=(1,0)$ independently reproduces
\cref{push:src:A.4.2}, including its integral freeness.

\item \emph{Canonical class, nullsphere surgery and smooth type.}
Let $F_m$ denote the reduced new multiple fibre.
The canonical-form computation in \ref{maps:src:D.3.3} applies to the
original pieces: before the regular log the canonical
line is $f^*\mathcal O(-1)$, and the primitive toroidal
cusp adds no divisor. A regular order-$m$ log adds
$(m-1)F_m$. Since $f^*(\mathrm{point})=mF_m$, one obtains
\[
                         K_X=\mathcal O_X(-F_m).
 \tag{D.O9}\label{maps:eq:D.O9}
\]
Here this formula also controls the topological class,
including the free classes supplied by the degree-one
Mayer--Vietoris kernel. The complement of $F_m$ has
$\pi_1=\mathbb Z\gamma$. Its primitive normal meridian
is $m\xi-v$, whose image is $-n\gamma$.
The residue sequence for the divisor therefore shows
that $[F_m]\in H^2(X;\mathbb Z)$ has exact order $|n|$
when $n\ne0$. For $n=0$ it is primitive of infinite
order: the quotient by its image embeds in the torsion-free
$H^2(X\setminus F_m)$, whose torsion vanishes by
$H_1(X\setminus F_m)=\mathbb Z$. For $n=\pm1$ it vanishes. Thus $c_1(X)=0$ in the
simply connected cases, and $X$ is spin; local canonical
triviality alone would not have sufficed for this step.

\item The angular $\delta$ action is free on every original
bundle and on the smooth log when $n=\pm1$. Its sole
fixed curve is the cusp rational curve with normal
weights $(1,-1)$. This curve belongs to the regular-fibre
specialization image and hence is a multiple of $s$.
By \eqref{maps:eq:D.O7}, $s=0$ already in the exterior after the
original additive fillings. The fixed curve is therefore
null-homologous, and then null-homotopic in a simply
connected member.
Apply the oriented equivariant nullsphere surgery and
projectable spin lift of \ref{maps:src:D.2.14}. It gives
$Y=X\#(S^3\times S^3)$ as a free circle bundle over a
simply connected spin five-manifold with
\[
            H_2(B_5)=\mathbb Z^8\oplus(\mathbb Z/m)^2
\]
and primitive Euler class. Smale--Barden identifies the
oriented pair as
\[
 (\#^8(S^2\times S^3)\#N_m,(x,0,\ldots,0)).
\]
Repeated \cite[Theorem A and Theorem B(2)]{manGR25}
give
\[
 Y\cong\#^7(S^2\times S^4)\#^8(S^3\times S^3)
                                        \#\widetilde\Sigma_0N_m.
\]
Cancellation by the complete Zhubr invariants proves
\eqref{maps:eq:D.O3}, including the specified twisted suspension
framing at even $m$.

\item The exact program \texttt{manifolds\_original\_additive.py}
retains all monodromies, fixed-loop bases, local covectors,
primitive fibre relations and full integral gluing maps.
It verifies 4,386 full degree-two Smith forms and 4,386
full degree-one Smith forms, including 582 cases with
$n=0$ in each degree. It separately checks the base-degree
map on bridge classes in both rank regimes. Formula
\eqref{maps:eq:D.O8} and the proved integral local maps supply the
unrestricted theorem.
\end{enumerate}
\end{proof}

\begin{theorem}\label{maps:src:D.2.21}
\StatementRecord{maps:src:D.2.21}{proved}{Historical events: Claude Fable 5.2 (Anthropic), effort high, 16 November 2026: carried as background without fault found. GPT-6 Astra (OpenAI), ChatGPT Pro mode, effort Pro, S6\_09\_PAPERS, 17 November 2026: targeted proof check. Also checked in this writing. Pre-continuation writing events: GPT-6 Astra (OpenAI), ChatGPT Pro mode, effort Pro, S6\_10\_PAPERS, 18 November 2026: targeted proof check. Continuation metadata events: GPT-6 Astra (OpenAI), ChatGPT Pro mode, effort Pro, S6\_10\_PAPERS continuation, 19 November 2026: used and rechecked in this run; status-vocabulary reconciliation; metadata reconciliation only; no new mathematical check. \emph{Pre-continuation recorded proof scope.} GPT-6 Astra (OpenAI, ChatGPT Pro mode), effort Pro, 2026-11-17 UTC; Claude Fable 5.2 (Anthropic), effort high, 16 November 2026; also checked in this writing by GPT-6 Astra (OpenAI), ChatGPT Pro mode, effort Pro, 18 November 2026}{the hypotheses, notation and definitions stated here; \cref{maps:src:D.2.20}; \cref{maps:src:D.2.5}; \cref{maps:src:D.2.12}; \cref{maps:src:D.2.14}; \cref{maps:src:D.2.16}}


Use the uncorrected degree-one data $P=G,Q=dG$ and the simple zero at the
marked $I_1$ in the following five rows. Keep the indicated original
additive filling and replace the other additive fibre by its finite-good
logarithmic model of order $m=qk$, with invariant power vector
$n\ell+r\delta$, where $k>0$, $r\in\mathbb Z$ and $\gcd(n,m)=1$.
\[
\begin{array}{c|c|c|r|r|r|r}
\mathrm{OS/profile}&d&\mathrm{finite\ log}\ (q)&\mathrm{original}&b&h&c\\\hline
49/III&6&IV^*\ (3)&III&1&3&2\\
49/III&6&III\ (4)&IV^*&6&2&3\\
50/IV&12&IV\ (3)&I_1^*&5&3&4\\
43/II&2&II\ (6)&III^*&7&1&2\\
46/II&4&II\ (6)&I_3^*&7&1&4
\end{array}
\]
Then $n\ne0$, $\pi_1(X)=H_1(X)=\mathbb Z/|n|$, and
\[
 H_2(X)=\mathbb Z^b,\qquad H_3(X)=\mathbb Z^{2b},\qquad
 H_4(X)=\mathbb Z^b\oplus\mathbb Z/|n|,\qquad H_5(X)=0.
\]
The endpoint groups are $\mathbb Z$, and $\chi(X)=2$.
Every member with $n=\pm1$ is oriented diffeomorphic to
\[
                 \#^b(S^2\times S^4)\ \#\ \#^b(S^3\times S^3).
\]
In particular none of these five simply connected families is $S^6$.

\end{theorem}
\begin{proof}
\begin{enumerate}[label=\arabic*.,leftmargin=*]
\item The component data, original local models and geometric
finite-good logarithmic replacements are those already proved in
\ref{maps:src:D.2.20} and \ref{maps:src:D.2.5}--\ref{maps:src:D.2.12}. We retain their markings rather than replacing
an original additive fibre by its invariant torus cohomology. Order the
three monodromies cyclically so that the finite-good one is first, and
write them $A_0,A_1,A_2$. Their product is the identity. The two remaining
full fixed lattices provide four primitive normal-torus vectors $B$.
If $\ell$ is the finite invariant vector of $\gamma$-coefficient one,
then $K=[B,j_0\ell]$ is a primitive integral basis of the rank-five
kernel of the graph differential. The exact integer inverses, monodromies,
component projections and complete matrices are saved in
\texttt{code/manifolds\_mixed\_additive.py} and its JSON certificate.

\item The regular degree-two coinvariants are freely generated by $s$, detected
by $\omega=e_1^*e_2^*-d\delta^*\gamma^*$. The original additive cap has
$r_0=b+1$ components and the actual principal $T^2$-bundle model of
\ref{maps:src:D.2.20}. Its local free defect is $r_0-2=b-1$. Its full fibre kernel maps
to $h\mathbb Zs$, with $h$ in the table. The one-component toric cusp
has no free defect. The finite cap has no free defect either.

\item Let $\beta$ generate the finite vertical coinvariant lattice and write
$\delta=q_\delta\beta$. Then
\[
 N_m\beta=(m/q_\delta)\delta,\qquad N_m\ell=m\ell,
 \qquad -\omega(\beta,\ell)=c=d/q_\delta,\quad
 -\omega(\ell,\delta)=-d.
\]
Choose $m\alpha+n\zeta=1$. The actual finite normal-circle tori are
\[
 \binom{-\omega(\beta,n\ell+r\delta)}{Lj_0(N_m\beta)},\qquad
 \binom{-\alpha\omega(\ell,n\ell+r\delta)}
 {Lj_0(\alpha N_m\ell+\zeta(n\ell+r\delta))}.
\]
After eliminating the four primitive columns $B$, the entire remaining
integral cokernel has the presentation
\[
                 \begin{pmatrix}h&cn&-dr\alpha\\0&0&1\end{pmatrix}.
\]
In each row $\gcd(h,cn)=1$: the nontrivial possibilities use respectively
$3\nmid n$, $2\nmid n$, $3\nmid n$; the last two have $h=1$.
Thus this cokernel vanishes for every allowed integer input, without an
order bound. The companion code verifies every constant primitive lattice
and $1,068$ full $6\times7$ Smith matrices as an independent certificate.

\item For completeness, the first-degree Mayer--Vietoris source has rank nine
and target rank eight: each boundary has rank three; the exterior has
rank three, while the finite, original and cusp caps have ranks two,
one and two. The original cap kills the elliptic and $\delta$ loops,
leaving $\gamma$. All original and cusp meridians vanish. The finite
power relation is therefore $n\gamma=0$. This proves
$\pi_1(X)=\mathbb Z/|n|$ directly and makes the rational first-degree
Mayer--Vietoris map onto. Its kernel is consequently free of rank one.
Together with the local defect $b-1$ and the zero cokernel above, this
proves $H_2(X)=\mathbb Z^b$. Integral duality, the universal coefficient
theorem and $\chi=2$ give the remaining groups as displayed.

\item When $n=\pm1$, the reduced finite fibre divisor is zero in integral
cohomology: the complement has fundamental group $\mathbb Z\gamma$,
and its primitive normal meridian maps to $\pm n\gamma$, so the divisor
residue sequence gives divisor order $|n|$. The canonical divisor differs
from a base pullback by a multiple of this reduced fibre; the full base
fibre is $m$ times the same divisor. Thus $c_1(X)=0$, including on the free
Mayer--Vietoris bridge class. The $\delta$ action is free on the original
principal bundle and on the finite quotient, since $\gcd(n,m)=1$.
Its only fixed component is the standard cusp sphere with normal weights
$+1,-1$. As in \ref{maps:src:D.2.20} this sphere lies in the image of the regular fibre,
hence is a multiple of $s$ and is nullhomologous by the unimodular
presentation. The equivariant surgery, spin lift and Smale--Barden
argument of \ref{maps:src:D.2.14} now give a free circle action on
$X\#(S^3\times S^3)$ with orbit
$\#^{b+1}(S^2\times S^3)$ and primitive Euler class.
Repeated Galaz-Garc\'ia--Reiser Theorem A and Theorem B(2), followed by
full Zhubr cancellation exactly as in \ref{maps:src:D.2.16}, give the asserted oriented
diffeomorphism type.\hfill$\square$
\end{enumerate}
\end{proof}


\begin{theorem}[finite initial data for sphere outputs]\label{maps:src:D.2.22}
\StatementRecord{maps:src:D.2.22}{proved}{Historical events: Claude Fable 5.2 (Anthropic), effort high, 16 November 2026: carried as background without fault found. GPT-6 Astra (OpenAI), ChatGPT Pro mode, effort Pro, S6\_09\_PAPERS, 17 November 2026: targeted proof check. Also checked in this writing. Pre-continuation writing events: GPT-6 Astra (OpenAI), ChatGPT Pro mode, effort Pro, S6\_10\_PAPERS, 18 November 2026: targeted proof check. Continuation metadata events: GPT-6 Astra (OpenAI), ChatGPT Pro mode, effort Pro, S6\_10\_PAPERS continuation, 19 November 2026: used and rechecked in this run; status-vocabulary reconciliation; metadata reconciliation only; no new mathematical check. \emph{Pre-continuation recorded proof scope.} GPT-6 Astra (OpenAI, ChatGPT Pro mode), effort Pro, 2026-11-17 UTC; Claude Fable 5.2 (Anthropic), effort high, 16 November 2026; also checked in this writing by GPT-6 Astra (OpenAI), ChatGPT Pro mode, effort Pro, 18 November 2026}{the hypotheses, notation and definitions stated here; \cref{maps:src:D.2.2}; \cref{maps:src:D.2.11}; \cref{maps:src:D.1.2}; \cite[Theorem 6.3 and Corollary 6.13]{manSS10}; \cref{maps:eq:D.FI3}; \cref{maps:eq:D.FI4}; \cref{maps:eq:D.FI5}}


Consider the standard section-line construction over any of the twenty-one
rational elliptic profiles with Mordell--Weil group $\mathbb Z$ in \ref{maps:src:D.2.2}.
Write $G$ for its chosen generator, $\langle G,G\rangle=1/d$,
$P=aG$, $Q=bG$, and $M=\mathcal O(P-O)$. Use a nonzero holomorphic
linearization with effective vertical divisor $D$, and the following local
fillings: a zero-free original quotient or a free affine logarithmic
replacement at each additive fibre; an original quotient or a standard
smooth Mumford filling at each multiplicative fibre; and, if desired,
free logarithmic replacements and rank-one Mumford fillings at regular
fibres. The latter occur at the regular zeros of the linearization.
These are the filling types specified in D.2; a new additive degeneration
at a zero of the linearization is not one of these local types.

If the completed manifold is a homotopy six-sphere, then $D$ is one
reduced component of one multiplicative fibre $I_n$. Translation by $Q$
acts transitively on its components and
\[
             \langle P,Q\rangle=\frac1n,
             \qquad ab=\frac dn.                                      \tag{D.FI1}\label{maps:eq:D.FI1}
\]
After simultaneously changing the signs of $P,Q$, take $a,b>0$.
There are exactly $172$ marked triples (profile, ordered section pair,
zero position) satisfying the resulting integral component criterion.
They are listed completely in
\texttt{records/manifolds\_sphere\_natural.json}.
This is a finite list of possible initial data, not a claim that all
its filling choices are spheres.

The same necessary reduction holds in the larger construction with
$M=\mathcal O(P-O+V)$ for an arbitrary integral vertical divisor $V$.
Modulo simultaneous translation of the chosen zero component, there
are $226$ marked (profile, ordered section pair, zero-fibre) candidates;
all $226$ component equations have integral solutions $V$.
Their remaining solution freedom is exactly the $Q$-invariant vertical
divisor lattice. This optional enlargement is distinguished from the
natural section-line construction in the preceding paragraph.

\end{theorem}


\begin{proof}
\begin{enumerate}[label=\arabic*.,leftmargin=*]
\item We give the local Euler identity, the rational rank obstruction, and the
finite integral enumeration separately. No logarithmic multiplicity is
bounded in any step.

\item \emph{1. The local line and the determinant.}
Near a multiplicative fibre choose a Tate coordinate with $q=t^n$;
holomorphic unit factors do not affect any valuation below. If $P=t^p$
times a unit, choose $0\le p<n$. The product defining the Tate theta
function gives
\[
 \theta(qz;q)=-z^{-1}\theta(z;q),\qquad
 \sigma(z)=\frac{\theta(z/P;q)}{\theta(z;q)},\qquad
 \sigma(qz)=P\sigma(z).
\]
Thus $\sigma$ is the meromorphic section with divisor $P-O$ of the
line with factor of automorphy $P$. At a generic point of component
$i$, represented by $z=t^i$ times a unit, its valuation in the flat
frame is
\[
 f_i=\min(0,i-p)\quad(0\le i<n),\qquad f_{i+n}=f_i+p.
\]
Indeed, in this range only the factor $1-z/P$ of the numerator can
have negative valuation; the other theta factors are units. The
meromorphic section has no vertical divisor in $\mathcal O(P-O)$.
Consequently the flat frame, expressed in an extending local frame,
has valuation $F_i=-f_i$. Adding $V=\sum v_iC_i$ changes this to
$F_i=-f_i+v_i$, where $v_{i+n}=v_i$. Put $A=-p$, so
$F_{i+n}=F_i+A$.

\item If translation by $Q$ shifts the component index by $B$, let $w$ be
the valuation of its flat linearization multiplier. Comparing the
extending frames in $M$ and $t_Q^*M$ gives, with no rational
intersection-theory ambiguity,
\[
 z_i=\operatorname{ord}_{C_i}(D)=w+F_i-F_{i+B}.
\]
For a quasi-periodic sequence with increment $A$, telescoping yields
\[
 \sum_{i=0}^{n-1}(F_{i+B}-F_i)=AB,
 \qquad \sum_i z_i=nw-AB.                                    \tag{D.FI2}\label{maps:eq:D.FI2}
\]
The identity holds for every integer $B$, by division by $n$ and
then telescoping the remaining finite interval. The actual valuation
vectors of the two long periods are $(n,-A)$ and $(B,-w)$: the latter
uses the inverse of the lift $M_{Qz}\to M_z$. Reversing the second
short-circle coordinate changes them to $(n,A),(B,w)$. Their index
is therefore $|nw-AB|$. For an effective nonzero local divisor this
equals $\sum_i z_i$. A standard unimodular triangular Mumford
cellulation has twice this many zero-dimensional torus strata, and
all its positive-dimensional strata have Euler number zero. Hence
\[
                  \chi(U_{I_n})=2\sum_i z_i.                  \tag{D.FI3}\label{maps:eq:D.FI3}
\]
This proves the local equality directly from the extending line
frames, rather than identifying an arbitrary support function with
the line. For example, $n=2,p=B=1,w=0$ gives
$F=(1,0)$, $A=-1$, $(z_0,z_1)=(1,0)$, and index one, precisely
the component-supported zero of \ref{maps:src:D.2.11}.

\item \emph{2. Local rational facts used in the rank obstruction.}
Let $F=T^4$, $A_i$ its monodromy at a deleted point, $L_i$ the
mapping-torus boundary, and $U_i$ a filling. Put
\[
 f_i=\dim H_1(F;\mathbb Q)_{A_i},\quad
 g_i=\dim H_2(F;\mathbb Q)_{A_i},\quad h_i=b_1(U_i).
\]
For every allowed local type, nearby-fibre specialization generates
$H_1(U_i;\mathbb Q)$; a multiple meridian is a fibre loop. Thus
$h_i\le f_i$. Here are the actual local models and the degree-two
fact needed below.

\item A rank-one Mumford cap is topologically a cycle of nodal elliptic
components times $T^2$: the degree-zero ruled line has topologically
trivial Chern class, and its translation gluing is isotopic to the
identity. The surviving circle of the nodal curve, its fibre-area
class, and the two transverse circles show that
$H^2(U_i;\mathbb Q)\to H^2(F;\mathbb Q)^{A_i}$ is onto.
For a rank-two triangular cap, take the universal toric model with
rays $(v,1)$, $v\in\mathbb Z^2$, and the standard periodic triangular
fan. Its period lattice $L\subset\mathbb Z^2$ acts by
\[
 (z_1,z_2,t)\longmapsto
 (u_{\ell,1}(t)t^{\ell_1}z_1,
  u_{\ell,2}(t)t^{\ell_2}z_2,t),\qquad \ell\in L.
\]
Here $u:L\to\mathcal O^*(\Delta)^2$ is a character. It extends to
$\mathbb Z^2$: put the inclusion in Smith form and take holomorphic
roots of its two nowhere-zero unit values on the disc. The extended
period action uses the same fan and shifts its dual complex by
$\mathbb Z^2$. Its quotient is a minimal hexagonal cap. The original
cap covers this quotient with deck group $\mathbb Z^2/L$. This
action is free even on the central fibre, because its action on
$\mathbb R^2/L$ is a nontrivial translation. This is an ordinary
finite covering, not the toric power map discussed in C.1.
Pull back the four independent specialization classes of \ref{maps:src:D.1.2}. On the regular torus
the pullback is an isomorphism over $\mathbb Q$ onto the invariant
space, of dimension four; hence it again gives surjectivity.
Subdividing the toric cellulation preserves these pulled-back classes.
For free finite affine quotients, rational transfer and averaging
give the same assertion. This includes higher logarithmic orders,
the good-reduction caps, and the semistable covers of odd $I_n^*$.

\item There is a separate check for an original cap; it cannot be replaced
without checking its component Chern class. Its central model is
the circle bundle of $M$ over the mapping torus of the component
translation of the Kodaira curve. Since the linearization is a unit
there, its multidegree vector $c$ is invariant under that translation.
If $c\ne0$, the rational circle-bundle transgression kills the
$\delta$ circle, so $h_i=f_i-1$. If $c=0$, let $\rho$ be the number
of component orbits. On the invariant component subspace, the
Kodaira intersection matrix is negative semidefinite with kernel
the whole fibre and rank $\rho-1$. The map
$H_2(U_i)\to H_2(U_i,L_i)$ has this rank: pair its image with the
component four-cycles (the component times the two circle factors).
Thus the cokernel of $H_2(L_i)\to H_2(U_i)$ has dimension at least
$\rho-1$. If $c=0$ and $\rho=1$, the invariant fibre-area class
and the circle classes give the same degree-two surjectivity as
above. These statements also follow directly for a multiplicative
curve from its cycle of spheres times $T^2$; for an additive curve
use its tree of spheres. Finite component permutations are handled
by averaging, which is legitimate over $\mathbb Q$.

\item \emph{3. The global rank budget.}
Delete $N$ discs containing every singular fibre, every zero of the
linearization, and every extra logarithmic site, and write $E$ for
the resulting torus bundle. Its fibre coinvariants in degree one
have dimension $r=1$. To verify this without assuming a special
section pair, observe that $\mathcal O(aG-O)$ differs from
$\mathcal O(G-O)^{\otimes a}$ by a vertical line, trivial on this
punctured base. The elliptic-plus-$\delta$ monodromy block therefore
scales its last row by $a$. For $a\ne0$ this does not change its
rational image rank. The exact twenty-one canonical blocks saved
by the program all have rank three. All full monodromy differences
lie in these first three coordinates, and hence have precisely
this image. Vertical modifications and the $Q$ column cannot
change the conclusion.

\item Also $v_2=\dim H_2(F;\mathbb Q)_{\pi_1}\ge1$. Explicitly,
in the marking $(e_1,e_2,\delta,\gamma)$ the actual monodromy
has the following blocks, where $z$ is the section cocycle of $G$:
\[
 A=\begin{pmatrix}
 B&0&bz\\ -az^tJB&1&c\\0&0&1
 \end{pmatrix},\qquad
 J=\begin{pmatrix}0&1\\-1&0\end{pmatrix},\qquad B\in SL_2(\mathbb Z).
\]
The same vertical-line trivialization gives the row multiplied by
$a$, while the Abel--Jacobi lift of $Q=bG$ gives the column $bz$.
The central entry $c$ is unrestricted. Direct block multiplication
shows that $A^t\Omega A=\Omega$, where
\[
 \Omega=\begin{pmatrix}aJ&0\\
 0&\begin{matrix}0&-b\\ b&0\end{matrix}\end{pmatrix}.
\]
In the only potentially nonzero cross term,
$B^tJz=(-z^tJB)^t$, so the two terms cancel. Thus
$\Omega=a e_1^*e_2^*+b\gamma^*\delta^*$ is a nonzero invariant,
independently of $c$ and of any vertical modification. The program
also verifies this identity symbolically. The case $a=0$ or $b=0$ cannot have an effective
nonempty $D$: norm over a component-preserving translation makes
such a divisor a base divisor of degree $\langle P,Q\rangle=0$.
It therefore cannot give Euler number two.

\item The graph-base homology sequence gives
\[
 b_1(E)=N,\qquad b_2(E)=v_2+4N-7,\qquad b_2(L_i)=g_i+f_i.
\]
If $H_1(X;\mathbb Q)=H_2(X;\mathbb Q)=0$, the first MV map is
onto and its kernel has dimension $\sum_i(f_i-h_i)$. This kernel
is a quotient of $H_2(X;\mathbb Q)$, so every $h_i=f_i$.
In particular every original cap has $c=0$. The degree-two map
from the boundary is also onto each cap; the component intersection
calculation then forces $\rho=1$ for every original cap.
All remaining caps consequently satisfy the degree-two invariant
surjectivity proved in step 2, and so $b_2(U_i)\ge g_i$.
After first mapping onto $\bigoplus H_2(U_i)$, the boundary kernel
has dimension at most $\sum f_i$. It must still surject onto
$H_2(E)$, since $H_2(X)=0$. Therefore
\[
 v_2+4N-7\le\sum_i f_i,
 \qquad\sum_i\operatorname{rank}(A_i-I)\le7-v_2\le6.          \tag{D.FI4}\label{maps:eq:D.FI4}
\]

\item Let $F_a,F_m$ count additive and multiplicative fibres of the
rational elliptic surface. The Shioda--Tate formula, $\rho(S)=10$,
and Mordell--Weil rank one give root rank seven. Comparing this
with $\chi(S)=12$ gives
\[
                            2F_a+F_m=5.
\]
This is also checked independently for every profile in the program;
see \cite[Theorem 6.3 and Corollary 6.13]{manSS10}.
Every additive monodromy difference has rank at least two, and
every multiplicative one rank at least one. A rank-two Mumford
zero at a multiplicative fibre raises its rank to two; a zero at
a regular fibre contributes an additional rank one. All original
and logarithmic caps have Euler number zero, as do rank-one caps.
Thus $\chi(X)=2$ and \eqref{maps:eq:D.FI3} force at least one rank-two zero,
already consuming the sole extra rank allowed by \eqref{maps:eq:D.FI4}.
There are no regular zeros and no second multiplicative zero.
Equation \eqref{maps:eq:D.FI3} now says the effective divisor at that fibre has
total coefficient one. It is one reduced component, as asserted.

\item \emph{4. Norm, height, and finite enumeration.}
Choose an integer $R$ such that $RQ$ meets the identity component
of every reducible fibre. Norm the line identity over the first
$R$ iterates of $t_Q$. The resulting divisor has zero component
multidegrees, and is therefore a sum of whole fibres. Since it is
effective and supported on the orbit of the single component,
$Q$ must act transitively on this $I_n$. Its norm is $(R/n)F$.
The normed line is
$M\otimes t_{RQ}^*M^{-1}$; the vertical correction cancels because
$RQ$ fixes every component. The height formula with a narrow second
section gives the base degree
$\langle P,RQ\rangle=R\langle P,Q\rangle$; all its local correction
terms vanish. Hence \eqref{maps:eq:D.FI1}. The divisor and height comparison is
integral: a vertical divisor of zero multidegree is a whole-fibre
multiple, by the corank-one component intersection matrix, and
$\operatorname{Pic}^0(S)=0$.

\item Here is the entire finite algorithm, with no implicit parameter
bound. For each multiplicative $I_n$ with $n\mid d$, list the
positive ordered divisors $ab=d/n$. Require that the component
of $Q$ generate $\mathbb Z/n$. For each Kodaira fibre let $I$ be
its integral component intersection matrix, $T$ the component
pullback permutation of $Q$, and $c_P=e_{P\text{-component}}-e_0$.
For a candidate divisor vector $z$, solve
\[
              (1-T)I V=Iz-(1-T)c_P.                         \tag{D.FI5}\label{maps:eq:D.FI5}
\]
At the chosen zero fibre, $z$ is one component; elsewhere $z=0$.
For the natural construction require $V=0$ and test every possible
zero component. For the larger construction solve \eqref{maps:eq:D.FI5} by
unimodular Smith operations. Since $I$ is invertible modulo the
invariant whole fibre, its solution kernel is exactly $\ker(1-T)$.
No additional numerical condition is omitted: \eqref{maps:eq:D.FI5} makes
$J\otimes\mathcal O(-D)$ a base pullback, and its norm has degree
zero by \eqref{maps:eq:D.FI1}, so the base pullback is trivial. Thus every retained
row actually has this effective linearization divisor.

\item The program checks the affine component diagrams, their primitive
multiplicity kernels and the translation permutations, then solves
every \eqref{maps:eq:D.FI5} integrally. The following counts retain marked fibre
positions; the first number uses $V=0$, the second allows $V$:
\[
\begin{array}{c|rr@{\qquad}c|rr}
43&6&6&43/II&2&2\\
45&8&17&45/II&4&9\\
46&6&9&46/II&2&3\\
47&14&16&47/II&6&8\\
47/III&8&10&49&10&10\\
49/II&2&2&49/III&4&4\\
50&10&15&50/II&2&3\\
50/IV&4&6&55&14&23\\
55/II&6&11&56&26&28\\
56/II&10&12&56/III&14&16\\
56/IV&14&16&\text{total}&172&226
\end{array}
\]
For $n>1$, the natural component equation in fact forces
$P$ and $Q$ to meet the same adjacent component, labelled $1$ or
$-1$, and the zero to be $C_0$; the code tests this directly,
including the doubled edges for $I_2$. The broader count normalizes
the zero component to $C_0$ by a power of $Q$; translating $M$
changes it only by an allowed vertical divisor. This explains
the different marking convention in the two counts.
\end{enumerate}
\end{proof}

\paragraph{Scope after \ref{maps:src:D.2.22}.}
The primary questions have different breadth. Viaclovsky's
Question 4 \cite[Section 1.1]{manViaclovskyLogs26} asks which starting configurations admit infinitely
many sphere multiplicities. Engel's Remark 8.5 \cite{manEngel26} expressly allows
an arbitrary fibrewise degree-zero line bundle, not only
$\mathcal O(P-O)$. We keep those scopes distinct. \ref{maps:src:D.2.22} removes
any need for an unexplained bound on section pairs or zero
positions when searching for spheres in the stated local sector.
The remaining task is the integral, all-order gluing calculation
for the finitely many initial rows and all allowed filling choices.
The $172$ and $226$ counts do not by themselves complete that task
or classify all nonsphere outputs.


\begin{theorem}[uniform all-order fundamental groups on the natural catalogue]\label{maps:src:D.2.23}
\StatementRecord{maps:src:D.2.23}{proved}{Historical events: Claude Fable 5.2 (Anthropic), effort high, 16 November 2026: carried as background without fault found. GPT-6 Astra (OpenAI), ChatGPT Pro mode, effort Pro, S6\_09\_PAPERS, 17 November 2026: targeted proof check. Also checked in this writing. Pre-continuation writing events: GPT-6 Astra (OpenAI), ChatGPT Pro mode, effort Pro, S6\_10\_PAPERS, 18 November 2026: targeted proof check. Continuation metadata events: GPT-6 Astra (OpenAI), ChatGPT Pro mode, effort Pro, S6\_10\_PAPERS continuation, 19 November 2026: used and rechecked in this run; status-vocabulary reconciliation; metadata reconciliation only; no new mathematical check. \emph{Pre-continuation recorded proof scope.} GPT-6 Astra (OpenAI, ChatGPT Pro mode), effort Pro, 2026-11-17 UTC; Claude Fable 5.2 (Anthropic), effort high, 16 November 2026; also checked in this writing by GPT-6 Astra (OpenAI), ChatGPT Pro mode, effort Pro, 18 November 2026}{the hypotheses, notation and definitions stated here; \cref{maps:src:D.2.22}; \cref{maps:src:D.2.19}}


For every one of the $172$ natural initial rows of \ref{maps:src:D.2.22}, retain any
of the stated original or Mumford fillings and perform any admissible
free affine logarithmic replacements. At an odd $I_n^*$, a logarithmic
replacement here means the free cyclic quotient of its rank-one
semistable Mumford cover. Regular logarithmic sites are allowed.
Let $m_i\ge1$ be their orders, $v_i$ their integral power vectors in
$H_1(T^4;\mathbb Z)$, and put $n_i=\gamma^*(v_i)$.
For a chosen system of peripheral lifts, let $e$ be the integer in
the global relation of their surviving $\gamma$ coordinates.
Then the complete fundamental group is
\[
 \pi_1(X)=\left\langle \gamma,c_1,\ldots,c_s\ \left|\
 [\gamma,c_i]=1,\ c_i^{m_i}=\gamma^{n_i},\
 c_1\cdots c_s=\gamma^e\right.\right\rangle .                 \tag{D.FG1}\label{maps:eq:D.FG1}
\]
The canonical lifts from the construction, obtained by unrolling the
translation quotient before filling, have $e=0$. Changing $c_i$ to
$c_i\gamma^{t_i}$ replaces $(n_i,e)$ by
$(n_i+m_it_i,e+\sum t_i)$ and leaves \eqref{maps:eq:D.FG1} unchanged.
Thus the formula also applies directly to other marked twist conventions.
The first integral homology is the Smith cokernel of the $(s+1)$ square
presentation matrix with rows
\[
 (-n_i,0,\ldots,m_i,\ldots,0)\quad(1\le i\le s),
 \qquad (-e,1,\ldots,1).                                    \tag{D.FG2}\label{maps:eq:D.FG2}
\]
There is no multiplicity bound and no restriction to invariant rational
representatives of the affine translations: their actual power vectors
are used.

After eliminating order-one sites by $c_i=\gamma^{n_i}$, replace $e$
by $e-\sum_{m_i=1}n_i$. If at least three sites of order greater than
one remain, the group is nontrivial. With zero, one, or two remaining
sites, it is abelian and is trivial precisely when, respectively,
\[
 |e|=1,\qquad |m_1e-n_1|=1,\qquad
 |m_2n_1+m_1n_2-em_1m_2|=1.                                \tag{D.FG3}\label{maps:eq:D.FG3}
\]
In the two-site case its invariant factors are the gcd of the entries
of the two-by-two relation matrix and its determinant divided by that
gcd; when the determinant vanishes use one free summand. In particular,
simple connectivity forces coprime remaining multiplicities and
$\gcd(m_i,n_i)=1$ at every remaining site.

\end{theorem}


\begin{proof}
\begin{enumerate}[label=\arabic*.,leftmargin=*]
\item We supply the integral vanishing comparison, since rational monodromy
alone would not prove this statement. The exact program
\texttt{manifolds\_sphere\_gluing.py} constructs all $172$ integral
monodromy tuples. Every zero-free point is locally narrow for at least
one of $P,Q$, so its lift is the one already proved in \ref{maps:src:D.2.19}.
The matrix at the one component-supported zero is independently fixed
by the global product relation. It agrees integrally with the normalized
weight $1/n$ derived in \ref{maps:src:D.2.22}. The program verifies every product,
all local Smith forms, and the preserved alternating form.

\item At a multiplicative Mumford cap the primitive vanishing sublattice is
$\operatorname{sat}\operatorname{im}(A_i-I)$: its individual primitive
circles bound the toric vanishing thimbles. The zero cap has rank two.
Every other multiplicative cap in all $172$ rows has rank one and
Smith width one; there is no nonprimitive necklace hidden in this
assertion. If instead the original zero-free quotient is retained,
these same circles are still nullhomotopic. When the line multidegree
is zero, the local line trivialization identifies the elliptic vanishing
thimble with this primitive circle. When it is nonzero, the natural
line has component degree vector $e_{P\text{-component}}-e_0$;
its entries have gcd one, so the component two-spheres kill the
$\delta$ circle integrally. The elliptic vanishing thimble then kills
its primitive elliptic coordinate as well, and hence kills the required
mixed circle. These comparisons do not replace an integral image by
its saturation without a geometric bounding disc.

\item Close these vanishing lattices under all monodromy matrices. Exact
unimodular operations show that their normal closure is
$\ker\gamma^*$ in $170$ rows. In the two rows
\[
 (46/II,a,b)=(46/II,1,4),\ (46/II,4,1)
\]
its quotient in $\ker\gamma^*$ is $\mathbb Z/3$. The $I_3^*$ cap
kills this last quotient. An original cap kills its elliptic vanishing
loops. For a semistable logarithmic cap, take its cyclic base cover.
The upstairs rank-one Mumford model has an actual disc bounding the
primitive circle in $\operatorname{im}(A_*^2-I)\otimes\mathbb Q$.
A selected regular fibre upstairs maps diffeomorphically to the
regular fibre downstairs: all nonidentity deck transformations move
its base root. Projecting this disc therefore bounds the same primitive
fibre circle downstairs, rather than a multiple. Appending this column
gives three unit Smith factors in each exceptional row. This proves
that all the fibre subgroup except the primitive $\gamma$ is killed,
for every allowed filling choice.

\item Over the punctured base the torus bundle has a section up to homotopy,
since the base retracts onto a graph. Its group is the corresponding
semidirect product. Modulo the vanishing normal closure just proved,
every monodromy acts trivially and its fibre group is exactly
$\mathbb Z\gamma$. A free affine log cap has a cyclic cover whose
fundamental group is its torus lattice, modulo the primitive vanishing
line in the semistable case. The covering extension has presentation
$c_i\lambda c_i^{-1}=A_i\lambda$, $c_i^{m_i}=v_i$; no further relation
is introduced. After the preceding quotient this becomes exactly
$[c_i,\gamma]=1$, $c_i^{m_i}=\gamma^{n_i}$.

\item For an original cap the translation quotient can be unrolled, and the
remaining group is $\mathbb Z\gamma$, with the base meridian trivial.
The same unrolling extends across a Mumford cap by unrolling its long
period in the infinite toric model. These are restrictions of the
original unrolled translation construction, so their global peripheral
relation has zero $\gamma$ coordinate. Equivalently, transition
functions of the line bundle contribute only the $\delta$ coordinate.
This proves $e=0$ in the canonical marking. Other meridian choices
produce exactly the stated change of $(n_i,e)$. Van Kampen now gives
\eqref{maps:eq:D.FG1}, including its completeness, and abelianization gives \eqref{maps:eq:D.FG2}.

\item Quotienting \eqref{maps:eq:D.FG1} by its central fibre yields the sphere-orbifold
presentation with cone orders $m_i>1$. With at least three such orders
it is nontrivial. One direct verification is the polygon-reflection
construction: a polygon with angles $\pi/m_i$ exists in the sphere,
plane or hyperbolic plane according to the angle sum, and products of adjacent side reflections give rotations of exact
orders $m_i$ whose cyclic product is the identity. They therefore
give a nontrivial representation of the orbifold presentation. The familiar finite spherical triangle
cases have the same nontrivial rotations. For at most two sites,
eliminate the product relation; the remaining generators commute.
The zero-, one-, and two-site integral relation matrices give \eqref{maps:eq:D.FG3}.
Reducing the last determinant modulo each multiplicity proves the
final coprimality assertions.
\end{enumerate}
\end{proof}

\begin{theorem}[both $I_3^*$ profiles are excluded from the natural sphere table]\label{maps:src:D.2.24}
\StatementRecord{maps:src:D.2.24}{proved}{Historical events: Claude Fable 5.2 (Anthropic), effort high, 16 November 2026: carried as background without fault found. GPT-6 Astra (OpenAI), ChatGPT Pro mode, effort Pro, S6\_09\_PAPERS, 17 November 2026: targeted proof check. Also checked in this writing. Pre-continuation writing events: GPT-6 Astra (OpenAI), ChatGPT Pro mode, effort Pro, S6\_10\_PAPERS, 18 November 2026: targeted proof check. Continuation metadata events: GPT-6 Astra (OpenAI), ChatGPT Pro mode, effort Pro, S6\_10\_PAPERS continuation, 19 November 2026: used and rechecked in this run; status-vocabulary reconciliation; metadata reconciliation only; no new mathematical check. \emph{Pre-continuation recorded proof scope.} GPT-6 Astra (OpenAI, ChatGPT Pro mode), effort Pro, 2026-11-17 UTC; Claude Fable 5.2 (Anthropic), effort high, 16 November 2026; also checked in this writing by GPT-6 Astra (OpenAI), ChatGPT Pro mode, effort Pro, 18 November 2026}{the hypotheses, notation and definitions stated here; \cref{maps:src:D.2.22}; \cref{maps:eq:D.IS2}; \cref{maps:eq:D.FI4}}


Neither OS profile $46$ nor $46/II$ can give a homotopy six-sphere
in the natural section-line construction and local filling sector
of \ref{maps:src:D.2.22}, for any admissible logarithmic orders or affine twists.
This includes arbitrary choices of original and semistable logarithmic
fillings at $I_3^*$.
\end{theorem}


\begin{proof}
\begin{enumerate}[label=\arabic*.,leftmargin=*]
\item \ref{maps:src:D.2.22} reduces these two profiles to the eight natural marked rows
in its finite catalogue. In every one of them, at $I_3^*$,
\[
 N=A_*^2-I,\qquad N^2=0,\qquad \operatorname{rank}N=1,
 \qquad \operatorname{Smith}(N)=(3).                         \tag{D.IS1}\label{maps:eq:D.IS1}
\]
The exact matrices also give
\[
 A_*N=-N,\qquad
 A_*=-1\text{ on }\Lambda/\ker N,\qquad
 A_*=1\text{ on }(\ker N/\operatorname{sat}\operatorname{im}N)
                         \otimes\mathbb Q.                 \tag{D.IS2}\label{maps:eq:D.IS2}
\]
These are the vanishing, dual-complex, and transverse factors of the
actual rank-one degeneration, respectively.

\item An original cap has more than one component orbit, so was already
excluded by the component intersection defect in \ref{maps:src:D.2.22}. For a log
cap of order $m=2k$, its semistable cover has monodromy
$A_*^{2k}=I+kN$. Thus its minimal rank-one Mumford cycle has exactly
$3k$ components. Its second rational homology has filtration with
component space $\mathbb Q^{3k}$, transverse area $\mathbb Q$, and
the two-dimensional product of the dual-complex circle with
$H_1(T^2;\mathbb Q)$. The finite deck group acts by permutation on
the first space, trivially on the area, and by minus one on the last
space, by \eqref{maps:eq:D.IS2}. Translation parts and unit gluing change neither
the component permutation nature nor these linear actions. Since the
group is finite, the filtration splits equivariantly over $\mathbb Q$;
rational transfer therefore gives
\[
                         b_2(U_*)=\rho+1,
\]
where $\rho$ is the number of component orbits. A group of order
$2k$ cannot act transitively on $3k$ objects: each orbit has size at
most $2k$, so $\rho\ge2$. Toric refinements can only add component
classes and do not remove this inequality. On the other hand
$g_*=\dim H_2(T^4;\mathbb Q)^{A_*}=2$, directly from the displayed
blocks. Hence $b_2(U_*)-g_*\ge1$.

\item The rank budget of \ref{maps:src:D.2.22} is saturated: its baseline is five, the
unique rank-two zero adds one, and $v_2\ge1$. Retaining an excess
$b_2(U_i)-g_i$ in that same MV calculation strengthens \eqref{maps:eq:D.FI4} to
\[
 \sum_i\operatorname{rank}(A_i-I)
 +\sum_i\bigl(b_2(U_i)-g_i\bigr)\le7-v_2.
\]
The left side is at least seven, while the right side is at most
six. This contradiction excludes a sphere. All comparisons took
place on the actual semistable cover and its free finite quotient;
no bound on $k$ or on the translation vector was used.
\end{enumerate}
\end{proof}


\begin{theorem}[all initial rows and all regular logarithms in the semistable sector]\label{maps:src:D.2.25}
\StatementRecord{maps:src:D.2.25}{proved}{Historical events: Claude Fable 5.2 (Anthropic), effort high, 16 November 2026: carried as background without fault found. GPT-6 Astra (OpenAI), ChatGPT Pro mode, effort Pro, S6\_09\_PAPERS, 17 November 2026: targeted proof check. Also checked in this writing. Pre-continuation writing events: GPT-6 Astra (OpenAI), ChatGPT Pro mode, effort Pro, S6\_10\_PAPERS, 18 November 2026: targeted proof check. Continuation metadata events: GPT-6 Astra (OpenAI), ChatGPT Pro mode, effort Pro, S6\_10\_PAPERS continuation, 19 November 2026: used and rechecked in this run; status-vocabulary reconciliation; metadata reconciliation only; no new mathematical check. \emph{Pre-continuation recorded proof scope.} GPT-6 Astra (OpenAI, ChatGPT Pro mode), effort Pro, 2026-11-17 UTC; Claude Fable 5.2 (Anthropic), effort high, 16 November 2026; also checked in this writing by GPT-6 Astra (OpenAI), ChatGPT Pro mode, effort Pro, 18 November 2026}{the hypotheses, notation and definitions stated here; \cref{maps:src:D.2.22}; \cref{maps:src:D.2.23}; \cref{maps:eq:D.FG2}; \cref{maps:eq:D.FG1}; \cref{maps:eq:D.FG3}; \cref{maps:src:D.2.14}; \cref{maps:src:D.2.12}; \cref{maps:src:D.2.13}; \cite[Theorem A]{manGR25}}


Consider any of the $62$ natural initial rows of \ref{maps:src:D.2.22} on the four
all-multiplicative profiles $45,47,55,56$. Write $P=aG,Q=bG$, with
$a,b>0$. Use the primitive Mumford model at each multiplicative fibre.
An original zero-free filling with zero line multidegree and one
component orbit may be used in its place; it has the same marked
primitive rank-one model. Perform any finite number of regular
logarithmic modifications, with data
\[
 m_i>0,\qquad v_i=(x_i,y_i,z_i,n_i)\in\mathbb Z^4,
 \qquad\gcd(m_i,x_i,y_i,z_i,n_i)=1.                           \tag{D.SC1}\label{maps:eq:D.SC1}
\]
The coordinates are $(e_1,e_2,\delta,\gamma)$. Put
$\omega=a e_1^*e_2^*+b\gamma^*\delta^*$. The initial rows have
$\gcd(a,b)=1$. Retain the peripheral integer $e$ of \ref{maps:src:D.2.23};
one may keep every order-one site and use the canonical value $e=0$.

Here is a complete integral homology algorithm without a bound on
orders, number of sites, or translation coordinates. Let $S$ be the
Seifert relation matrix \eqref{maps:eq:D.FG2}, and write
\[
 \operatorname{coker}S=\mathbb Z^r\oplus A_1,\qquad r\in\{0,1\}.
\]
Form an integral presentation $\mathcal C$ with generators
$s$ and $t_i(u)$, linear in $u\in\Lambda=\mathbb Z^4$, and relations
\[
 \begin{split}
 m_i t_i(u)&=\omega(u,v_i)s\quad(u=e_1,e_2,\delta,\gamma),\\
 t_i(v_i)&=0,\\
 \sum_i t_i(\delta)+be\,s&=0.
 \end{split}                                                \tag{D.SC2}\label{maps:eq:D.SC2}
\]
Write $T=\operatorname{Tor}\operatorname{coker}\mathcal C$.
Then $\operatorname{coker}\mathcal C$ has free rank $r$ and
\[
\begin{array}{c|c}
 j&H_j(X;\mathbb Z)\\\hline
 0,6&\mathbb Z\\
 1&\mathbb Z^r\oplus A_1\\
 2,3&\mathbb Z^{2r}\oplus T\\
 4&\mathbb Z^{2r}\oplus A_1\\
 5&\mathbb Z^r.
\end{array}                                                 \tag{D.SC3}\label{maps:eq:D.SC3}
\]
The fundamental group is exactly \eqref{maps:eq:D.FG1}, rather than merely the
abelian group in degree one. Smith reduction of the explicit matrix
\eqref{maps:eq:D.SC2} computes $T$, including all extensions.

If this group is trivial, let
\[
 M=\prod_i m_i,\quad
 X_0=\sum_i(M/m_i)x_i,\quad Y_0=\sum_i(M/m_i)y_i,\quad
 g=\gcd(b,M,X_0,Y_0).
\]
For no displayed sites use $M=1$ and $X_0=Y_0=0$.
Then
\[
 H_2(X)=H_3(X)=
 \mathbb Z/g\oplus\mathbb Z/M\oplus\mathbb Z/(bM/g).          \tag{D.SC4}\label{maps:eq:D.SC4}
\]
Consequently a sphere occurs precisely when the group criterion
\eqref{maps:eq:D.FG3} holds and $b=M=1$. In the canonical marking with all sites
retained, this says that every regular logarithmic order is one and
$|\sum_i n_i|=1$.

There are exactly $17$ permitted initial rows with $b=1$:
\[
\begin{array}{c|c|c|c}
 \text{profile}&\text{zero fibre}&(a,b)&\text{marked positions}\\\hline
 45&I_1&(8,1)&4\\
 47&I_2&(7,1)&1\\
 47&I_1&(14,1)&3\\
 55&I_4&(5,1)&1\\
 55&I_1&(20,1)&3\\
 56&I_5&(6,1)&1\\
 56&I_3&(10,1)&1\\
 56&I_2&(15,1)&1\\
 56&I_1&(30,1)&2.
\end{array}                                                 \tag{D.SC5}\label{maps:eq:D.SC5}
\]
This is the complete natural sphere table for these four profiles,
including original-versus-Mumford choices: any additional original
cap not covered by the first paragraph is excluded by the positive
local defect in \ref{maps:src:D.2.22}.

More generally, every simply connected output in these $17$ rows,
with arbitrary data \eqref{maps:eq:D.SC1}, is oriented diffeomorphic to
\[
                              \widetilde\Sigma_0N_M,         \tag{D.SC6}\label{maps:eq:D.SC6}
\]
with the twisted framing and the Smale five-manifold of \ref{maps:src:D.2.14}.
Thus its full smooth type depends only on $M$, including when $M$
is even. For $b>1$, \eqref{maps:eq:D.SC2}--\eqref{maps:eq:D.SC4} give the complete integral
homology and exclude a sphere; no smooth classification of all
those torsion types is asserted here.

\end{theorem}


\begin{proof}
\begin{enumerate}[label=\arabic*.,leftmargin=*]
\item \emph{1. The actual marked caps and the constant integral calculation.}
At the component zero the valuation lattice has index one by \ref{maps:src:D.2.22},
so its cap is the primitive hexagonal model. Every other
multiplicative monodromy in these $62$ rows has rank-one difference
and Smith width one, as checked in \ref{maps:src:D.2.23}. Its local Mumford cap is
the marked one-node model times $T^2$. The integral specialization
onto its $H_2$, and that of the hexagonal cap, are onto; their
boundary kernels are exactly the normal tori swept by
$\operatorname{Fix}(A_i)$. The toric angular coordinates give these
as pure swept classes, with zero fibre-area term. The same statement
for an original zero-degree, one-orbit cap follows by quotienting
its component cycle by $Q$, using the unit-factor isotopy of \ref{maps:src:D.2.12}.

\item These normal markings agree with the section-line construction.
Rigidify $M$ along $O$ on the punctured base. Put the divisor of this
rigidification at the chosen zero; its integral change there is
multiplication of the line coordinate by $t^h$. On the infinite
toric model this translates the fan by an integral lattice vector,
and hence extends over the cap. At all other points the
rigidification is a unit and the earlier relative unit isotopy
applies. Thus this normalization does not assume that equality of
vertical forms alone determines mixed swept classes.

\item The exact program \texttt{manifolds\_semistable\_catalogue.py}
performs the graph and exterior calculations for all $62$ rows.
The degree-two fibre coinvariants are freely generated by $s$,
with primitive detector $\omega$; all its Smith factors are one.
Adjoin one regular puncture, with monodromy the identity. In the
notation of \ref{maps:src:D.2.13}, the $14$ old normal columns form a primitive
rank-$14$ sublattice of the rank-$17$ graph kernel. Their quotient
is freely generated by the regular sweeps of $e_1,e_2,\gamma$,
while the sweep of $\delta$ is zero. All $62$ left inverses and
quotient matrices are saved, so this is an exact integral
certificate, not a numerical rank inference.

\item Adding another regular puncture adds four free graph coordinates.
The old normal lattice remains primitive by projection onto the
previous graph. Since every monodromy fixes $\delta$, the peripheral
Fox maps satisfy $\sum j_i(\delta)=0$. Induction therefore gives
all regular sweeps freely, with the sole relation
$\sum_i t_i(\delta)=0$. The degree-two exterior summand remains
$\mathbb Zs$. A change of global peripheral lift by $e\gamma$
adds $be\,s$ to this relation: the correction to a swept $\delta$
cycle is its wedge with that translation, and
$\omega(\delta,\gamma)=-b$. Equivalently, under
$v_i\mapsto v_i+m_i h_i\gamma$ the new sweep is
$t_i'(u)=t_i(u)+\omega(u,h_i\gamma)s$;
$e\mapsto e+\sum h_i$ makes \eqref{maps:eq:D.SC2} invariant. This checks its
sign against \eqref{maps:eq:D.FG1}.

\item \emph{2. The complete regular-log kernel.}
At a regular log point its reduced torus has lattice
$\Lambda'=\Lambda+\mathbb Z(v_i/m_i)$. The nearby fibre inclusion
and the base circle send a boundary bivector respectively to
$z$ and $u\wedge(v_i/m_i)$ in $\bigwedge^2\Lambda'$.
This map is onto. The condition in \eqref{maps:eq:D.SC1} means $v_i$ is primitive
modulo $m_i$; the integral normal lattice is exactly
$m_i\Lambda+\mathbb Z v_i$. This can be checked in a basis in which
$v_i$ modulo $m_i$ is the first basis vector, using the Euclidean
algorithm. Its generators give the full normal columns
\[
 (-\omega(u,v_i),\ m_it_i(u)),\qquad (0,t_i(v_i)).
\]
They are precisely the first two lines of \eqref{maps:eq:D.SC2}. The argument
identifies the full integral kernel, including its index, not only
its rational span. Since specialization is onto every cap in degree
two, Mayer--Vietoris reduces to this cokernel.

\item The remaining first MV kernel is free: all boundary $H_1$ groups
are free in these primitive toric models. Its rank is $r$, because
the domain and target of the first boundary map have equal rational
dimension and its cokernel is $H_1(X)$. Hence it contributes a
split $\mathbb Z^r$ to $H_2(X)$. Over $\mathbb Q$, eliminate the
sweeps in \eqref{maps:eq:D.SC2}; its sole residual relation is
\[
                 b\left(e-\sum_i n_i/m_i\right)s=0.
\]
Thus its free rank is the same $r$ as the Seifert matrix.
Only the primitive hexagonal cap contributes Euler characteristic,
so $\chi(X)=2$. Integral Poincare duality and the universal coefficient
theorem give \eqref{maps:eq:D.SC3}, including the identical torsion in $H_2,H_3$.
This also proves the formula when the rational Euler number is zero;
that case is not obtained by deleting a nonzero-order torsion factor.

\item For one site in canonical marking, quotienting the zero $\delta$
sweep displays the concrete four-by-five matrix
\[
 \begin{pmatrix}
 -ay&ax&bn&-bz&0\\
 m&0&0&0&x\\0&m&0&0&y\\0&0&0&m&n
 \end{pmatrix}.                                             \tag{D.SC7}\label{maps:eq:D.SC7}
\]
Its maximal minors are, up to signs,
$bm^2n\,(m,n,z,y,x)$. In particular, whenever $n\ne0$,
the finite degree-two cokernel has order $bm^2|n|$.
The implementation accepts the general multi-site matrix \eqref{maps:eq:D.SC2},
and verifies its comparison with all constant graph presentations.

\item \emph{3. The simply connected Smith form.}
Assume \eqref{maps:eq:D.FG1} is trivial. After removing order-one sites there are
at most two multiplicities greater than one, they are coprime, and
every $n_i$ is a unit modulo its $m_i$. Put
\[
                  D=M\left(\sum_i n_i/m_i-e\right).
\]
The group criterion gives $D=\pm1$. Work over $\mathbb Z_p$.
If $p$ does not divide $M$, all sweeps can be eliminated and the
cokernel is $\mathbb Z_p/(b)$. If $p\mid M$, it divides exactly
one multiplicity, say $m_i$. For every other site $m_j$ is a unit,
so its $\delta$ sweep equals $-bn_j/m_j$ times $s$. The global
relation then eliminates $t_i(\delta)$. Its own normal relation
becomes $bD/(M/m_i)$ times $s$, and hence is equivalent to $bs=0$.
Since $n_i$ is a unit, its relation $t_i(v_i)=0$ eliminates
$t_i(\gamma)$, whose remaining normal relation is redundant by
alternation of $\omega$. The three remaining generators have
presentation
\[
 \begin{pmatrix}
 b&0&0\\ -ay_i&m_i&0\\ ax_i&0&m_i
 \end{pmatrix}.                                             \tag{D.SC8}\label{maps:eq:D.SC8}
\]
Its determinantal divisors are $g_i,m_i g_i,bm_i^2$, where
$g_i=\gcd(b,m_i,x_i,y_i)$ in $\mathbb Z_p$; we used
$\gcd(a,b)=1$. The Chinese remainder theorem packages all these
primary invariant factors as \eqref{maps:eq:D.SC4}, because
$X_0,Y_0$ reduce at $p\mid m_i$ to unit multiples of $x_i,y_i$.
The invariant factors divide in the stated order, since $g$ divides
both $b$ and $M$. This proves the formula for every integer twist,
not by extrapolating from the finite checks.

\item The group in \eqref{maps:eq:D.SC4} vanishes exactly for $b=M=1$. Simple connectivity,
Euler number two, and integral duality then give a homotopy sphere;
$\Theta_6=0$ identifies it with the standard $S^6$. The seventeen
initial rows in \eqref{maps:eq:D.SC5} are the complete $b=1$ filter of \ref{maps:src:D.2.22}.
An excluded original cap has nonzero Chern degree or multiple
component orbits and would already force $b_2>0$, so cannot supply
additional sphere rows.

\item \emph{4. The smooth type when $b=1$.}
We verify the hypotheses of the geometric proof \ref{maps:src:D.2.14} rather than
inferring a smooth type from homology. At every zero-free
multiplicative point the component equation forces $P$ to be narrow
when $Q=G$; the component action of $Q$ is primitive and transitive. The code
checks these component statements and, more directly, that adjoining
$\delta$ to the actual primitive vanishing vector gives two unit
Smith factors at every such site. Thus $\delta$ stays primitive
modulo the vanishing line and its circle is free on the cap; width
one alone is not being used to infer this. At the zero, index one makes the valuation lattice the full original
short lattice, in which $\delta$ is the second coordinate. The code
checks its primitive inclusion in that lattice in all seventeen rows.
The actual hexagonal fan weights are therefore the same $(-1,0,1)$ as in \ref{maps:src:D.2.14}, with one fixed rational curve and
normal bundle $\mathcal O(-1)_{+1}\oplus\mathcal O(-1)_{-1}$.
For every regular log cap, $\gcd(m_i,n_i)=1$ implies that a scalar
angle lying in $\Lambda+\mathbb Z(v_i/m_i)$ is integral. Hence no
new circle stabilizers occur.

\item The fixed curve lies in the cyclic fibre summand generated by $s$.
The $\delta$ normal relations together with the global relation,
when $b=1$, are the unimodular Seifert matrix because $D=\pm1$.
They kill $s$ integrally. The fixed sphere is therefore nullhomologous
and nullhomotopic. Its equivariant surgery gives
$Y\cong X\#(S^3\times S^3)$ with a free, projectably spin circle
action, exactly as in \ref{maps:src:D.2.14}. The same orientation convention in its
normal-coordinate conjugation is retained. Its spin simply connected
quotient has $H_2=\mathbb Z\oplus(\mathbb Z/M)^2$, with primitive
Euler class. Smale--Barden and the spin branch of
\cite[Theorem A]{manGR25} give
$Y\cong(S^3\times S^3)\#\widetilde\Sigma_0N_M$.
Cancellation by the full Zhubr invariants, as justified in \ref{maps:src:D.2.14},
proves \eqref{maps:eq:D.SC6}, including the two-primary refined invariant.
\end{enumerate}
\end{proof}


\begin{theorem}[complete natural sphere table without $I_n^*$]\label{maps:src:D.2.26}
\StatementRecord{maps:src:D.2.26}{proved}{Historical events: Claude Fable 5.2 (Anthropic), effort high, 16 November 2026: carried as background without fault found. GPT-6 Astra (OpenAI), ChatGPT Pro mode, effort Pro, S6\_09\_PAPERS, 17 November 2026: targeted proof check. Also checked in this writing. Pre-continuation writing events: GPT-6 Astra (OpenAI), ChatGPT Pro mode, effort Pro, S6\_10\_PAPERS, 18 November 2026: targeted proof check. Continuation metadata events: GPT-6 Astra (OpenAI), ChatGPT Pro mode, effort Pro, S6\_10\_PAPERS continuation, 19 November 2026: used and rechecked in this run; status-vocabulary reconciliation; metadata reconciliation only; no new mathematical check. \emph{Pre-continuation recorded proof scope.} GPT-6 Astra (OpenAI, ChatGPT Pro mode), effort Pro, 2026-11-17 UTC; Claude Fable 5.2 (Anthropic), effort high, 16 November 2026; also checked in this writing by GPT-6 Astra (OpenAI), ChatGPT Pro mode, effort Pro, 18 November 2026}{the hypotheses, notation and definitions stated here; \cref{maps:src:D.2.22}; \cref{maps:eq:D.SC1}; \cref{maps:eq:D.SC3}; \cref{maps:eq:D.FG3}; \cref{maps:src:D.2.25}; \cref{maps:src:D.2.23}; \cite[Section 2.4, equation (2.2)]{manUU21}; \cref{maps:src:D.2.12}; \cref{maps:eq:D.HY6}; \cref{maps:eq:D.SC8}; \cref{maps:src:D.2.14}}


For the sixteen profiles in \ref{maps:src:D.2.22} containing no fibre $I_n^*$,
the natural section-line construction has exactly $59$ marked
initial/additive-filling patterns that can yield a sphere.
They lie on $50$ distinct initial rows. The complete table,
including the section pair, zero component, marked fibre positions,
and every additive original/log choice, is
\path{records/manifolds_nonstar_sphere_patterns.csv}.
At the zero the index-one Mumford filling is required; at every
other multiplicative fibre in these sphere patterns either the
original filling or the primitive Mumford filling is permitted.
The pattern counts are
\[
\begin{array}{c|r@{\qquad}c|r@{\qquad}c|r@{\qquad}c|r}
43&3&43/II&1&45&4&45/II&4\\
47&4&47/II&4&47/III&4&49&3\\
49/II&1&49/III&2&55&4&55/II&4\\
56&5&56/II&6&56/III&5&56/IV&5.
\end{array}                                                 \tag{D.HY1}\label{maps:eq:D.HY1}
\]
These are parameter families, not multiplicity samples. The exact
unrestricted parameter conditions for every row follow below.

More generally, take a natural initial row without $I_n^*$ and
retain only original caps of zero line multidegree with one
component orbit, using the primitive Mumford caps elsewhere.
A finite good-reduction log cap can occur in a simply connected
output only when its integral fixed lattice contains a vector
$\ell_i$ with $\gamma^*(\ell_i)=1$. Write its power vector
$v_i=n_i\ell_i+r_i\delta$. The following calculation applies
whenever $\gcd(m_i,n_i)=1$, a necessary condition for simple
connectivity. Set
\[
\begin{array}{c|ccccc}
 \text{finite type}&II&III&IV&III^*&IV^*\\\hline
 d_i\text{ (monodromy order)}&6&4&3&4&3\\
 q_i\text{ (vertical index)}&1&2&3&2&3.
\end{array}                                                 \tag{D.HY2}\label{maps:eq:D.HY2}
\]
Allow every order $m_i=d_i k_i$, $k_i>0$, every such power vector,
and every free affine lift with that power. Additional regular logs
have the unrestricted data \eqref{maps:eq:D.SC1}.

Here is their full integral degree-two presentation. For each finite
log retain a generator $d_i$; for each regular log retain its four
sweeps $t_j(\Lambda)$; retain the fibre generator $s$. Impose
\[
\begin{aligned}
 (m_i/q_i)d_i+(bn_i/q_i)s&=0 &&\text{(finite logs)},\\
 m_jt_j(u)-\omega(u,v_j)s&=0,\quad t_j(v_j)=0
                                      &&\text{(regular logs)},\\
 \sum_i d_i+\sum_jt_j(\delta)+be\,s&=0. &
\end{aligned}                                                 \tag{D.HY3}\label{maps:eq:D.HY3}
\]
All entries are integers. With $T$ the torsion of this cokernel and
$r,A_1$ obtained from the complete Seifert matrix, the full homology
is exactly \eqref{maps:eq:D.SC3}. Thus this is also an all-order integral algorithm
for the nonsimply-connected members satisfying the stated finite
coprimality condition. No conclusion for the omitted noncoprime
finite twists is silently inferred.

For a simply connected output put
\[
 Q_f=\prod_{i\text{ finite}}q_i,\quad c=b/Q_f,\quad
 R=\prod_{j\text{ regular}}m_j,
 \quad X_0=\sum_j(R/m_j)x_j,\quad Y_0=\sum_j(R/m_j)y_j.
\]
Empty products equal one. Then $c$ is a positive integer and
\[
 H_2(X)=H_3(X)=
 \mathbb Z/g\oplus\mathbb Z/R\oplus\mathbb Z/(cR/g),
 \qquad g=\gcd(c,R,X_0,Y_0).                                \tag{D.HY4}\label{maps:eq:D.HY4}
\]
Consequently the complete sphere criterion is
\[
 \pi_1(X)=1\text{ by \eqref{maps:eq:D.FG3}},\qquad b=Q_f,
 \qquad m_j=1\text{ at every regular log}.                  \tag{D.HY5}\label{maps:eq:D.HY5}
\]
For one finite log this includes every multiple of its order; for
two finite logs it requires their multiplicities to be coprime and
the displayed Seifert determinant to be $\pm1$. The only two-finite
sphere pattern is the original $IV^*,III,I_1$ pair. The finite
numerators $r_i$ are unrestricted. All twelve profiles in \eqref{maps:eq:D.HY1}
which contain additive fibres therefore admit infinitely many
allowable sphere multiplicities; the four all-multiplicative
profiles have the order-one conclusion of \ref{maps:src:D.2.25}.

Whenever $c=1$ and the output is simply connected, its exact oriented
smooth type is $\widetilde\Sigma_0N_R$. For $c>1$ the groups \eqref{maps:eq:D.HY4}
are exact, but a general smooth identification of these remaining
torsion types is not asserted by this theorem.

\end{theorem}


\begin{proof}
\begin{enumerate}[label=\arabic*.,leftmargin=*]
\item \emph{1. Necessary local choices and the integral graph lattice.}
The sphere reduction \ref{maps:src:D.2.22} already excludes every nonzero-degree
original cap and every original cap with multiple component orbits.
At a finite good-reduction point the image of $\gamma^*$ on the
integral fixed lattice is the component order of $Q$, namely one,
two or three in these profiles. If it is greater than one, that
prime divides both $m_i$ and every power-vector coordinate $n_i$;
\ref{maps:src:D.2.23} then excludes simple connectivity. These are exact fixed-lattice
calculations, included for all $172$ initial rows in
\texttt{manifolds\_sphere\_normal\_lattices.py}.

\item That program also computes the complete boundary fixed-sweep map
\[
 \bigoplus_i\operatorname{Fix}(A_i)
       \longrightarrow H_1(B^\circ;H_1(T^4;\mathbb Z)).
\]
For every row it is onto with unit Smith factors. Its kernel has
rank one and is precisely the primitive diagonal copy of $\delta$,
by the peripheral Fox identity. Thus the graph part is the direct
sum of these fixed lattices with the sole relation
$\sum t_i(\delta)=0$. Adding regular sites has the same integral
inductive comparison as \ref{maps:src:D.2.25}. The fibre coinvariants are
$\mathbb Zs$, detected primitively by $\omega$.

\item There is an integral check for the permitted original additive
caps, beyond their rational invariant cycles. A zero-degree original
cap is topologically $S^1_\delta$ times the mapping torus of the
component translation of its Kodaira tree. Here zero multidegree
trivializes the underlying line on the tree of rational curves.
In such a trivialization the remaining equivariant multiplier is a
unit. Its circle-valued part has a logarithm because the tree has
$H^1=0$; interpolation of that logarithm gives an isotopy to the
identity. The rigidification along the zero section fixes the
constant phase. This isotopy extends over the normal collar, and
its unit multiplier has zero meridian winding. Consequently it
preserves the marked normal tori used below and proves the stated
product description, rather than merely a rational homology
comparison. Its homology is free;
with one orbit, $H_2$ has rank two and the component intersection
map to $H_2(U,\partial U)$ is zero. Poincare--Lefschetz duality
therefore makes $H_2(\partial U)\to H_2(U)$ onto integrally.
The $\delta$ normal torus bounds the section disc times $S^1_\delta$.
For the other fixed loop, take the least power of $Q$ which is
narrow and its local Neron logarithm. The resulting family of loops
extends over the disc, so its normal torus also bounds. Local
logarithms at additive reduction are described in
\cite[Section 2.4, equation (2.2)]{manUU21}.
The linearization is a unit and the line has zero component degrees,
so the rigidified norm lift introduces no extra fibre term.
These two tori project onto the entire integral fixed lattice.
They are therefore the full free part of the kernel; any boundary
torsion maps to zero in both free targets and changes nothing in
MV. This proves that each unchanged minimal cap simply sets its
fixed sweeps to zero, with the asserted markings.

\item \emph{2. The actual finite log kernel.}
Let $\Lambda_0=\ker\gamma^*$. At every finite slot with primitive
$\gamma$, the exact vertical coinvariant calculation is
\[
 (\Lambda_0)_{A_i}=\mathbb Z\beta_i,
 \qquad [\delta]=q_i[\beta_i],
 \qquad \operatorname{Fix}(A_i)=\mathbb Z\delta\oplus\mathbb Z\ell_i.
                                                               \tag{D.HY6}\label{maps:eq:D.HY6}
\]
The table \eqref{maps:eq:D.HY2} follows either from the component group or the
saved primitive Smith bases. The code verifies all $68$ such slots.
For $m=m_i$, its norm satisfies
$N_m\beta_i=(m/q_i)\delta$ and $N_m\ell_i=m\ell_i$.
Invariance of $\omega$ then proves the exact, signed identity
\[
                     -\omega(\beta_i,\ell_i)=b/q_i.
\]
In particular $q_i\mid b$.

\item An arbitrary affine good-reduction lift with the prescribed integral
power is translation-conjugate, over the root disc, to the invariant
representative. To see the integral issue, the usual sequence of
period lattices and holomorphic vector sections identifies the
translation-conjugacy obstruction with
$\Lambda^{A_i}/N_m\Lambda$; it is exactly the power vector modulo
norms. The vector-section cohomology vanishes by finite averaging.
Thus equality of the power class gives an actual holomorphic
translation conjugacy extending over the smooth central torus.
This justifies the invariant representative for every affine lift
under discussion, rather than restricting the theorem to a chosen
constant lift.

\item The reduced finite fibre has first homology generated by
$\beta_i,\ell_i,c_i$ with
$m c_i=n_i\ell_i+r_iq_i\beta_i$. Choose
$\alpha m+\beta n_i=1$. Its first homology is free, with basis
$\beta_i$ and $\alpha\ell_i+\beta c_i$. Its flat normal complex
line has torsion Chern class, hence zero integral Chern class because
$H^2$ of this fibre is free. Consequently the boundary map in degree
two is onto and its kernel is exactly the two normal circle tori
of this basis. The actual affine collar maps give their images in
the exterior as
\[
\begin{split}
 &\bigl(-\omega(\beta_i,v_i),\ t_i(N_m\beta_i)\bigr),\\
 &\bigl(-\alpha\omega(\ell_i,v_i),\
               t_i(\alpha N_m\ell_i+\beta v_i)\bigr).
\end{split}                                                  \tag{D.HY7}\label{maps:eq:D.HY7}
\]
This is the marked rectangle-sweep computation of \ref{maps:src:D.2.12}; the basis
calculation here verifies that these are the full integral kernel,
not a sublattice of finite index. Substituting \eqref{maps:eq:D.HY6} gives
\[
 (bn_i/q_i)s+(m_i/q_i)t_i(\delta),\qquad
 -\alpha b r_i s+t_i(\ell_i)+\beta r_i t_i(\delta).
\]
The second relation eliminates $t_i(\ell_i)$ with coefficient one.
The first is exactly the finite relation in \eqref{maps:eq:D.HY3}. In particular
all $r_i$ disappear from the degree-two group for an integral reason.
The regular kernel was proved in \ref{maps:src:D.2.25}. Together with the one
global $\delta$ relation this proves \eqref{maps:eq:D.HY3}. All cap $H_1$ groups
in this coprime range are free, so the remaining MV kernel is the
split free group of rank $r$. Rational elimination and integral
duality give the whole table \eqref{maps:eq:D.SC3}.

\item \emph{3. All-primary Smith calculation and sphere enumeration.}
Suppose the Seifert group is trivial. Its nontrivial multiplicities
are pairwise coprime, and all numerators are units modulo them.
Thus the $q_i$ are pairwise coprime and $Q_f\mid b$.
Write $M$ for the product of all multiplicities, including the
finite ones, and $D=M(\sum n_i/m_i-e)=\pm1$.
Eliminate the $\delta$ sweeps over $\mathbb Z_p$ exactly as in
\ref{maps:src:D.2.25}. At a prime dividing a finite multiplicity $m_i$, its row
has been divided by $q_i$, so the residual fibre relation is
$(b/q_i)D/(M/m_i)$ times $s$. All other multiplicities are units
at that prime. At a prime dividing a regular multiplicity, the
residual relation is $b$ times a unit; the finite $q_i$ are units
there. At a prime dividing no multiplicity, it is again $b$ times
a unit. These are precisely the primary fibre orders of
$c=b/Q_f$.

\item The only extra generators at a prime dividing a regular multiplicity
are its two elliptic sweeps. Their three-generator matrix is
\eqref{maps:eq:D.SC8}, with the fibre order replaced by $c$. Because
$\gcd(a,c)=1$, its invariant factors are
$g_i,m_i,cm_i/g_i$. Chinese remainders yield \eqref{maps:eq:D.HY4}, with the
stated $X_0,Y_0$. This proves \eqref{maps:eq:D.HY5} for all integers.

\item The finite program now has no arbitrary search bound: list every
natural row, every additive original/log choice satisfying the
local conditions, and the vertical indices \eqref{maps:eq:D.HY2}. There are
$164$ minimal-cap patterns in this coprime finite range. Two have
an unavoidable orbifold obstruction from noncoprime base orders;
$162$ can satisfy simple connectivity. Exactly $59$ have $b=Q_f$.
Their complete data are saved in the CSV and JSON tables. Each
of them is realized: choose a numerator $\pm1$ for one finite
site, Bezout numerators for the coprime $(3k,4l)$ pair, or total
order-one regular numerator $\pm1$ if there is no finite site.
The free affine quotient exists by its primitive $\gamma$ value.
This also gives infinitely many orders for every listed pattern
with at least one finite site. Conversely every sphere was forced
into these choices by \ref{maps:src:D.2.22} and \ref{maps:src:D.2.23}, proving completeness in
all sixteen profiles.

\item \emph{4. The exact smooth type for $c=1$.}
The catalogue directly checks that at every zero-free multiplicative
site of a $b=Q_f$ pattern, $P$ is narrow, $Q$ acts transitively on
components, and adjoining $\delta$ to the actual vanishing vector
has two unit Smith factors. Thus the scalar circle is free there.
It is free on every permitted original additive cap by its trivial
component Chern class, and on each finite log cap because
$\gcd(m_i,n_i)=1$. At the zero its lattice is the full original
short lattice, retaining $\delta$ as the second coordinate, so its
single fixed sphere and normal weights are those of \ref{maps:src:D.2.14}.

\item The square subsystem on $s$ and the $\delta$ sweeps in \eqref{maps:eq:D.HY3}
has determinant $bD/Q_f=cD=\pm1$. Hence $s$ is zero integrally.
The fixed sphere is nullhomotopic, and the equivariant surgery,
projectable spin quotient, Smale--Barden classification, twisted
suspension formula and full Zhubr cancellation from \ref{maps:src:D.2.14} apply
with $H_2=(\mathbb Z/R)^2$. They give precisely
$X\cong\widetilde\Sigma_0N_R$. This supplies the asserted smooth
identification and all its refined spin data.
\end{enumerate}
\end{proof}


\begin{theorem}[the three $I_1^*$ profiles give no natural sphere]\label{maps:src:D.2.27}
\StatementRecord{maps:src:D.2.27}{proved}{Historical events: Claude Fable 5.2 (Anthropic), effort high, 16 November 2026: carried as background without fault found. GPT-6 Astra (OpenAI), ChatGPT Pro mode, effort Pro, S6\_09\_PAPERS, 17 November 2026: targeted proof check. Also checked in this writing. Pre-continuation writing events: GPT-6 Astra (OpenAI), ChatGPT Pro mode, effort Pro, S6\_10\_PAPERS, 18 November 2026: targeted proof check. Continuation metadata events: GPT-6 Astra (OpenAI), ChatGPT Pro mode, effort Pro, S6\_10\_PAPERS continuation, 19 November 2026: used and rechecked in this run; status-vocabulary reconciliation; metadata reconciliation only; no new mathematical check. \emph{Pre-continuation recorded proof scope.} GPT-6 Astra (OpenAI, ChatGPT Pro mode), effort Pro, 2026-11-17 UTC; GPT-6 Astra (OpenAI, ChatGPT Pro mode), S6\_06\_BLIND, effort Pro, 2026-11-17; Claude Fable 5.2 (Anthropic), effort high, 16 November 2026; also checked in this writing by GPT-6 Astra (OpenAI), ChatGPT Pro mode, effort Pro, 18 November 2026}{the hypotheses, notation and definitions stated here; \cref{maps:src:D.2.22}; \cref{maps:src:D.2.24}; \cite[Sections 7.6--7.7]{manSS10}; \cref{maps:src:D.2.23}; \cref{maps:eq:D.ST3}; \cref{maps:eq:D.ST4}; \cref{maps:eq:D.ST8}; \cref{maps:eq:D.HY3}; \cref{maps:eq:D.ST10}; \cref{maps:src:D.2.26}; \cref{maps:eq:D.ST11}; \cref{maps:eq:D.ST2}; \cref{maps:eq:D.ST12}}


None of the natural section-line inputs on
\[
 (I_1^*,I_3,I_1,I_1),\qquad (I_1^*,II,I_3),\qquad
 (I_1^*,IV,I_1)
\]
produces a homotopy six-sphere, for any permitted original,
logarithmic or Mumford filling, multiplicity or affine twist in
the sector of \ref{maps:src:D.2.22}. The obstruction is integral. In the only
remaining minimal star caps, the full normal-kernel indices are
\[
\begin{array}{c|c|c|c}
 m_*&\text{component action}&H_2(U_*;\mathbb Z)
       &\kappa_*\ \text{(normal projection index)}\\\hline
 2&\text{one component}&\mathbb Z^2&1\\
 4&\text{two components interchanged}&
                   \mathbb Z^2\oplus\mathbb Z/2&4.
\end{array}                                                   \tag{D.ST1}\label{maps:eq:D.ST1}
\]
If the complete output is simply connected and its other caps
have the minimal rational ranks, the order of $H_2$ has exactly
one factor of $2$ in the first row, and exactly two factors of $2$
in the second row. In particular $H_2\ne0$.

\end{theorem}


\begin{proof}
\begin{enumerate}[label=\arabic*.,leftmargin=*]
\item \emph{1. Exhausting the possible star caps.}
The sixteen natural initial rows on these three profiles are
listed in the certificate
\path{records/manifolds_i1star_obstruction.json}.
No original star cap has both zero component multidegree and one
component orbit, so \ref{maps:src:D.2.22} excludes every such original cap from
a sphere. A semistable root cover has order $m_*=2k$, and
$A_*^{2k}-1=kN$, where $N$ has rank one, square zero and primitive
image. Thus its minimal rank-one fibre has exactly $k$ components.
This is the integral version of the $I_n^*$ semistable reduction
used in \ref{maps:src:D.2.24}; see also \cite[Sections 7.6--7.7]{manSS10}.

\item On the dual-complex circle the residual action reverses
orientation. Its action on the $k$ components therefore has the
form $i\mapsto r-i$ and has at least $\lceil k/2\rceil$ orbits.
The transverse two-torus has trivial linear action after the
vanishing circle is collapsed. The component classes contribute
one class per orbit, the transverse area contributes one, and
the dual-circle cross terms are anti-invariant. Consequently
\[
 b_2(U_*)=\rho+1,\qquad
 \dim H_2(T^4;\mathbb Q)_{A_*}=2.
\]
The saturated rank budget of \ref{maps:src:D.2.22} forces $\rho=1$.
Hence $k=1$, or $k=2$ with the two components interchanged.
Further toric subdivisions add a component defect and cannot
restore this equality. This proves that \eqref{maps:eq:D.ST1} exhausts the
remaining cap types without an order bound imposed externally.

\item There is a second necessary integral condition. On the quotient
of the fibre lattice by the primitive vanishing circle, an
actual affine power vector is fixed by the induced involution.
In eight of the sixteen rows the image of this fixed lattice
under $\gamma^*$ is $2\mathbb Z$. These are precisely the ordered
section pairs
\[
 (a,b)=(4,1),\ (4,3),\ (12,1).
\]
Their star numerator is even, while their star multiplicity is
even. The Seifert presentation \ref{maps:src:D.2.23} therefore precludes simple
connectivity. This uses the fixed lattice \emph{after} collapsing
the vanishing circle, not the smaller fixed lattice before that
collapse.

\item For the other eight rows the pairs are
\[
 (a,b)=(1,4),\ (1,12),\ (3,4),                             \tag{D.ST2}\label{maps:eq:D.ST2}
\]
and there is an integral basis $(\alpha,\beta,\lambda,\ell)$
with $\gamma^*(\ell)=1$ and
\[
\begin{gathered}
 A\alpha=-\alpha,\quad A\beta=-\beta-\lambda,\quad
 A\lambda=\lambda-\alpha,\quad A\ell=\ell,\\
 \delta=-\alpha+2\lambda,\qquad
 A^2-1=\alpha\otimes\beta^*.
\end{gathered}                                               \tag{D.ST3}\label{maps:eq:D.ST3}
\]
Here $\alpha$ is primitive. The exact integral change of basis
for every marked row is in the certificate. For example, in a
standard marking put
\[
\begin{aligned}
 \alpha&=2e_1+a\delta,&
 \lambda&=e_1+\tfrac{a+1}{2}\delta,\\
 \beta&=e_2+\tfrac{3a-1}{2}\lambda,&
 \ell&=-\tfrac{3b}{4}e_1+\tfrac b2e_2+\gamma.
\end{aligned}
\]
The exceptional marking on $50/II$ is also checked integrally.
In this basis the degree-two fibre coinvariants are freely
generated by
\[
 x=[\beta\wedge\lambda],\qquad y=[\beta\wedge\ell],
\]
with relations
\[
 [\alpha\wedge\beta]=-2x,\quad
 [\lambda\wedge\ell]=-2y,\quad
 [\alpha\wedge\lambda]=[\alpha\wedge\ell]=0.                \tag{D.ST4}\label{maps:eq:D.ST4}
\]
In particular the boundary group $H_2(\partial U_*;\mathbb Z)$
is free of rank four; there is no boundary torsion that can be
discarded. The exterior fibre functional is still
$\omega=a e_1^*e_2^*+b\gamma^*\delta^*$, and
\[
 \omega(x)=-a,\qquad \omega(y)=b/4.                         \tag{D.ST5}\label{maps:eq:D.ST5}
\]

\item \emph{2. The normalization and its integral parity.}
We compute the retraction of $U_*$ to its central fibre by its
normalization and conductor. Since a simply connected output
has $\gcd(m_*,n_*)=1$, its star numerator is odd. All quotient
maps on the transverse elliptic curves below are therefore
unramified, with the displayed degrees.

\item For $m_*=2$ the normalization is an oriented $S^2$ bundle $N$
over a two-torus. Its conductor is a connected bisection $C$,
which maps with degree two to the double curve $D$. The relation
$A\lambda=\lambda-\alpha$ in \eqref{maps:eq:D.ST3} determines the parity of
this bundle. In a transverse circle coordinate $u$ its
pole-interchanging clutch can be represented by
\[
 z\longmapsto e^{-2\pi i u}/z.
\]
The corresponding $\mathrm{PGL}_2$ loop has representative
$\left(\begin{smallmatrix}0&e^{-2\pi i u}\\1&0\end{smallmatrix}\right)$.
Its determinant has odd winding, so its lift to the double cover
of $\mathrm{SO}(3)$ changes sign. Hence $w_2(N)$ evaluates to one
on the base. The remaining affine translations and unit
multipliers change the flat part of the clutch, not this winding;
the odd numerator changes the primitive base marking by an odd
index and preserves the parity. This connects the parity
calculation to every actual affine cap, not only a selected
linear model.

\item Let $f$ be the sphere fibre and choose a section class $\sigma$
with $\sigma^2=1$. Such a choice follows by adding multiples of
$f$ to any odd-square section. The bisection has zero
self-intersection, so
\[
 [C]=2\sigma-f.
\]
The normalization pushout gives
\[
 H_2(U_*)=
 \langle\sigma,f,D\mid 2\sigma-f-2D=0\rangle
          \cong\mathbb Z^2.                               \tag{D.ST6}\label{maps:eq:D.ST6}
\]
There is no additional $H_1$-kernel term: the conductor map to
$D$ is injective on $H_1$. The regular torus
$\alpha\wedge\beta$ maps to $f$, since it sweeps the single
semistable component with degree one. The regular
$\lambda\wedge\ell$ torus lies in the smooth part of the
normalization, covers its base twice and is disjoint from $C$.
Its class is consequently $2\sigma-f$, hence $2D$ in the
pushout. By \eqref{maps:eq:D.ST4},
\[
 x\longmapsto-\sigma+D,\qquad y\longmapsto-D.
\]
These are an integral basis. Thus the fibre-coinvariant map has
cokernel zero in this case.

\item For $m_*=4$ with the two components interchanged, choose one
component upstairs. Its stabilizer is generated by the square
of the deck generator. That square fixes
$\alpha,\lambda,\ell$. The open $\mathbb C^*$ bundle of its
normalization therefore has an abelian period lattice. More
explicitly it has generators $\alpha,\lambda,\ell,c^2$ and one
relation
\[
 2c^2=t\alpha+r\lambda+n_*\ell.
\]
Because $n_*$ is odd, $\alpha$ is primitive in this lattice.
The associated complex line over the quotient two-torus has
zero Chern class. Its two pole sections are thus homologous;
write their common class as $\sigma$. They map with degree two
to two different double curves $D_0,D_1$. Hence
\[
\begin{split}
 H_2(U_*)&=
 \langle f,\sigma,D_0,D_1\mid
                  \sigma-2D_0=\sigma-2D_1=0\rangle\\
 &\cong\mathbb Z f\oplus\mathbb Z D_0
                       \oplus\mathbb Z/2\,(D_1-D_0).
\end{split}                                                 \tag{D.ST7}\label{maps:eq:D.ST7}
\]
Again the conductor map on $H_1$ is injective into the sum of
its double-curve targets, so this is the complete group.

\item The regular $\alpha\wedge\beta$ torus sweeps both semistable
components, each with degree one onto the chosen component
fibre downstairs, and therefore maps to $2f$. The regular
$\lambda\wedge\ell$ torus covers the normalization base twice.
It avoids both pole sections, so its class is $2\sigma=4D_0$.
Thus, modulo possible multiples of the displayed order-two
class,
\[
 x\longmapsto-f,\qquad y\longmapsto-2D_0.
\]
The free part of this map has determinant two. Including the
order-two summand, its full cokernel has order four, independent
of the two possible torsion corrections. This is precisely the
integral effect missed by the rational vertical coinvariant
identity $[\delta]=4[\beta]$.

\item For completeness, the conductor has no third homology, and
its degree-two map into the direct sum of the double-curve
groups is injective: each conductor component maps with degree
two to its corresponding double curve. Thus the normalization
sequence identifies $H_3(U_*)$ with $H_3(N)=\mathbb Z^2$.
It also gives $H_4(U_*)=H_4(N)=\mathbb Z$, generated by the
reduced central component. A positive multiple of this divisor is principal,
being the divisor of the base parameter. Its intersection with
every two-cycle is therefore zero. Under Poincare--Lefschetz
duality this is the entire map
$H_2(U_*)\to H_2(U_*,\partial U_*)=H^4(U_*)=\mathbb Z$.
It follows that
\[
 H_2(\partial U_*)\longrightarrow H_2(U_*)
 \quad\hbox{is onto integrally}.                            \tag{D.ST8}\label{maps:eq:D.ST8}
\]

\item Let $C_*=H_2(T^4)_{A_*}$ and let $K_*$ be the full kernel in
\eqref{maps:eq:D.ST8}. The boundary Wang sequence gives
\[
 0\longrightarrow C_*\longrightarrow H_2(\partial U_*)
        \longrightarrow\operatorname{Fix}(A_*)
        \longrightarrow0.
\]
The fibre map $C_*\to H_2(U_*)$ is injective, and the preceding
calculations therefore give
\[
 [\operatorname{Fix}(A_*):\operatorname{proj}K_*]
 =\bigl|H_2(U_*)/\operatorname{im}C_*\bigr|
 =\begin{cases}1,&m_*=2,\\4,&m_*=4\text{ with swap}.
   \end{cases}                                             \tag{D.ST9}\label{maps:eq:D.ST9}
\]
This proves \eqref{maps:eq:D.ST1}, retaining the normalization torsion and
both full normal directions.

\item \emph{3. The full normal determinant and the global obstruction.}
Here is a useful integral determinant calculation. Suppose all
caps have the minimal rational ranks and their boundary maps
in degree two are onto. Put
\[
 \kappa_i=[\operatorname{Fix}(A_i):
                         \operatorname{proj}K_i],
\]
where $K_i$ is the full integral boundary kernel. The primitive
graph calculation for all $172$ rows gives the exterior group
as $\mathbb Zs$ plus the direct sum of the fixed sweep lattices,
with the sole relation $\sum_i t_i(\delta)+be\,s=0$.
Over $\mathbb Q$, each $K_i$ is the graph of a linear fibre
functional on its fixed lattice. Its value on $\delta$ is fixed
without choosing any of its other coefficients: the
$\delta$ normal torus bounds on the root cover, and transfer
with the actual affine power vector gives the exact relation
\[
 m_i t_i(\delta)+bn_i s=0.                                 \tag{D.ST10}\label{maps:eq:D.ST10}
\]
The sign is the same oriented convention as \eqref{maps:eq:D.HY3}.

\item Choose integral bases of the projected kernels. Their
covolumes are $\kappa_i$. Eliminate all independent sweep
coordinates over $\mathbb Q$ and then the single diagonal
$\delta$ relation. The determinant of the complete square
Mayer--Vietoris matrix is
\[
 \left(\prod_i\kappa_i\right)
                   b\left(e-\sum_i n_i/m_i\right).         \tag{D.ST11}\label{maps:eq:D.ST11}
\]
All longitudinal and mixed fibre coefficients cancel: only
the values \eqref{maps:eq:D.ST10} on the diagonal vector enter the last
Schur complement. Multiplying the rational calculation by the
integral covolumes is an equality of integer determinants,
not a saturation assumption. The certificate also verifies
this identity symbolically with arbitrary two-by-two normal
projection matrices and arbitrary longitudinal coefficients.

\item For an unchanged minimal cap or an index-one multiplicative
cap, $\kappa_i=1$. For a finite good-reduction cap \ref{maps:src:D.2.26} gives
$\kappa_i=m_i/q_i$. For a regular logarithm its projected
kernel is $m_i\Lambda+\mathbb Z v_i$, of index $m_i^3$;
this is the full lattice because the affine action is free.
These statements include every cap that could remain in a
sphere after the \ref{maps:src:D.2.22} rank obstruction.

\item If $X$ is simply connected, the Seifert determinant is $\pm1$,
so
\[
 \left|e-\sum_i n_i/m_i\right|=(\prod_i m_i)^{-1}.
\]
All degree-one groups in this Mayer--Vietoris map are free.
For the star boundary this follows from \eqref{maps:eq:D.ST3}; for the
remaining boundaries the local matrices have only unit
nonzero Smith factors, as checked in every retained row.
The star cap has the abelian presentation obtained from the
collapsed lattice and the relation $m_*c=v_*$; its relation
is primitive because $\gcd(m_*,n_*)=1$. The same primitive
relation proves freeness for the finite and regular log caps.
An unchanged minimal additive cap has the product/mapping-torus
model of \ref{maps:src:D.2.26} and free first homology. The exterior group is
free as well. The degree-one Mayer--Vietoris source and target
have equal rank, while their cokernel is $H_1(X)=0$.
Consequently that map is an isomorphism and has no additional
degree-one bridge kernel.
Thus \eqref{maps:eq:D.ST11} gives the exact order
\[
 |H_2(X)|=\frac{b}{q_*Q_o}\,R^2,\qquad
 q_*:=m_*/\kappa_*=
 \begin{cases}2,&m_*=2,\\1,&m_*=4\text{ with swap},\end{cases} \tag{D.ST12}\label{maps:eq:D.ST12}
\]
where $Q_o$ is the product of the indices $q_i$ of the other
finite caps and $R$ is the product of regular logarithmic
multiplicities.

\item Simple connectivity permits at most two multiple fibres and
forces their multiplicities to be coprime. Since $m_*$ is
even, every other multiplicity is odd. The only possible
additional finite index is $3$ from the $IV$ fibre of $50/IV$;
an order-six $II$ log cannot coexist with the star log in a
simply connected output. In every row of \eqref{maps:eq:D.ST2},
$b/Q_o$ has precisely two factors of $2$, and $R$ is odd.
Equation \eqref{maps:eq:D.ST12} now has two-primary order $2$ for $m_*=2$
and $4$ for $m_*=4$. Neither can be a homology sphere.
All excluded cap choices already have a positive rational
second Betti number or nontrivial fundamental group. This
completes the exhaustion.
\end{enumerate}
\end{proof}

\begin{corollary}[all twenty-one natural sphere profiles]\label{maps:src:D.2.28}
\StatementRecord{maps:src:D.2.28}{proved}{Historical events: Claude Fable 5.2 (Anthropic), effort high, 16 November 2026: carried as background without fault found. GPT-6 Astra (OpenAI), ChatGPT Pro mode, effort Pro, S6\_09\_PAPERS, 17 November 2026: targeted proof check. Also checked in this writing. Pre-continuation writing events: GPT-6 Astra (OpenAI), ChatGPT Pro mode, effort Pro, S6\_10\_PAPERS, 18 November 2026: targeted proof check. Continuation metadata events: GPT-6 Astra (OpenAI), ChatGPT Pro mode, effort Pro, S6\_10\_PAPERS continuation, 19 November 2026: used and rechecked in this run; status-vocabulary reconciliation; metadata reconciliation only; no new mathematical check. \emph{Pre-continuation recorded proof scope.} GPT-6 Astra (OpenAI, ChatGPT Pro mode), effort Pro, 2026-11-17 UTC; Claude Fable 5.2 (Anthropic), effort high, 16 November 2026; also checked in this writing by GPT-6 Astra (OpenAI), ChatGPT Pro mode, effort Pro, 18 November 2026}{the hypotheses, notation and definitions stated here; \cref{maps:src:D.2.22}; \cref{maps:eq:D.HY3}; \cref{maps:eq:D.HY5}; \cref{maps:src:D.2.23}; \cref{maps:src:D.2.24}; \cref{maps:src:D.2.27}; \cref{maps:src:D.2.25}; \cref{maps:src:D.3.2}}


For the natural line $M=\mathcal O(P-O)$, with all section pairs,
component-supported linearization zeros, permitted original and
Mumford fillings, additive and regular logarithms, arbitrary
multiplicities and affine twists in the sector of \ref{maps:src:D.2.22}, the
complete profile classification is as follows. A pattern means
a marked initial section/zero/additive-filling row, not an
individual complex structure or a sample of its multiplicities.
\[
\begin{array}{c|l|r|l}
\text{profile}&\text{singular fibres}&\text{sphere patterns}
                     &\text{sphere multiplicities}\\\hline
43&III^*,3I_1&3&\text{unbounded}\\
43/II&III^*,II,I_1&1&\text{unbounded}\\
45&I_8,4I_1&4&\text{all orders }1\\
45/II&I_8,II,2I_1&4&\text{unbounded}\\
46&I_3^*,3I_1&0&\text{none}\\
46/II&I_3^*,II,I_1&0&\text{none}\\
47&I_7,I_2,3I_1&4&\text{all orders }1\\
47/II&I_7,I_2,II,I_1&4&\text{unbounded}\\
47/III&I_7,III,2I_1&4&\text{unbounded}\\
49&IV^*,I_2,2I_1&3&\text{unbounded}\\
49/II&IV^*,II,I_2&1&\text{unbounded}\\
49/III&IV^*,III,I_1&2&\text{unbounded}\\
50&I_1^*,I_3,2I_1&0&\text{none}\\
50/II&I_1^*,II,I_3&0&\text{none}\\
50/IV&I_1^*,IV,I_1&0&\text{none}\\
55&I_5,I_4,3I_1&4&\text{all orders }1\\
55/II&I_5,I_4,II,I_1&4&\text{unbounded}\\
56&I_5,I_3,I_2,2I_1&5&\text{all orders }1\\
56/II&I_5,I_3,I_2,II&6&\text{unbounded}\\
56/III&I_5,I_3,III,I_1&5&\text{unbounded}\\
56/IV&I_5,I_2,IV,I_1&5&\text{unbounded}.
\end{array}                                                  \tag{D.ST13}\label{maps:eq:D.ST13}
\]
Thus twelve profiles admit unbounded sphere multiplicities,
four admit spheres only with all logarithmic orders one, and
five admit none. There are exactly $59$ marked sphere patterns
on $50$ distinct natural initial rows. Their full individual
conditions, including all affine parameters, are \eqref{maps:eq:D.HY3}--\eqref{maps:eq:D.HY5}
and the complete table
\path{records/manifolds_nonstar_sphere_patterns.csv}.
The profile summary is also saved as
\path{records/manifolds_complete_natural_sphere_profiles.csv}.
Every sphere in this classification is diffeomorphic to the
standard $S^6$.

\end{corollary}


\begin{proof}
\begin{enumerate}[label=\arabic*.,leftmargin=*]
\item \ref{maps:src:D.2.22} is the exhaustive finite initial reduction; \ref{maps:src:D.2.23} gives
the full fundamental group for every order. \ref{maps:src:D.2.24} eliminates
the two $I_3^*$ profiles, \ref{maps:src:D.2.27} eliminates the three $I_1^*$
profiles, and \ref{maps:src:D.2.25}--\ref{maps:src:D.2.26} compute every remaining sphere pattern
and all its unrestricted multiplicity conditions. Their
simply connected homology-sphere outputs are homotopy spheres
by Hurewicz and Whitehead, and $\Theta_6=0$ gives the stated
smooth type, as in \ref{maps:src:D.3.2}. The exact program
\path{code/III/manifolds_i1star_obstruction.py} checks all sixteen
star blocks, all eight integral canonical forms, both
normalization presentations, every possible torsion marking
in the order-four fibre map, the symbolic determinant identity,
and the twenty-one-profile tally.
\end{enumerate}
\end{proof}

The scope here is the natural section-line sphere question.
The optional vertically corrected catalogue of $226$ rows in
\ref{maps:src:D.2.22} has not been substituted for it. Formula \eqref{maps:eq:D.ST13} concerns the sphere outputs. The unrestricted
finite-log homology is supplied in \ref{maps:src:D.2.29}--\ref{maps:src:D.2.33}, and the entire specified
minimal simply connected smooth sector is completed by \ref{maps:src:D.2.35}--\ref{maps:src:D.2.36},
\ref{maps:src:D.2.38}--\ref{maps:src:D.2.39}, \ref{maps:src:D.2.42} and \ref{maps:src:D.2.44}--\ref{maps:src:D.2.45}. The remaining weighted and
additional-regular-log smooth classification is scoped in \ref{maps:src:D.2.4}.


\begin{theorem}[all free finite twists in the primitive-$\gamma$ sector]\label{maps:src:D.2.29}
\StatementRecord{maps:src:D.2.29}{proved}{Historical events: Claude Fable 5.2 (Anthropic), effort high, 16 November 2026: carried as background without fault found. GPT-6 Astra (OpenAI), ChatGPT Pro mode, effort Pro, S6\_09\_PAPERS, 17 November 2026: targeted proof check. Also checked in this writing. Pre-continuation writing events: GPT-6 Astra (OpenAI), ChatGPT Pro mode, effort Pro, S6\_10\_PAPERS, 18 November 2026: targeted proof check. Continuation metadata events: GPT-6 Astra (OpenAI), ChatGPT Pro mode, effort Pro, S6\_10\_PAPERS continuation, 19 November 2026: used and rechecked in this run; status-vocabulary reconciliation; metadata reconciliation only; no new mathematical check. \emph{Pre-continuation recorded proof scope.} GPT-6 Astra (OpenAI, ChatGPT Pro mode), effort Pro, 2026-11-17 UTC; Claude Fable 5.2 (Anthropic), effort high, 16 November 2026; also checked in this writing by GPT-6 Astra (OpenAI), ChatGPT Pro mode, effort Pro, 18 November 2026}{the hypotheses, notation and definitions stated here; \cref{maps:src:D.2.26}; \cref{maps:eq:D.HY2}; \cref{maps:eq:D.SC2}; \cref{maps:src:D.2.23}; \cref{maps:eq:D.UF6}; \cref{maps:eq:D.UF8}; \cref{maps:eq:D.HY3}; \cref{maps:eq:D.UF9}}


The restriction $\gcd(m_i,n_i)=1$ in the homology part of \ref{maps:src:D.2.26}
can be removed. The resulting exact formula retains torsion in
the finite cap itself. We keep the same natural initial rows,
minimal original and Mumford fillings, and the finite sites
whose integral fixed lattice contains $\ell_i$ with
$\gamma^*(\ell_i)=1$. These are the $68$ marked finite positions
certified in \ref{maps:src:D.2.26}; their orders and indices are $(d_i,q_i)$
from \eqref{maps:eq:D.HY2}.

At such a position write
\[
 m_i=d_i k_i,\qquad v_i=n_i\ell_i+r_i\delta,\qquad
 t_i=\gcd(m_i,n_i,q_i r_i),\qquad k_i>0.
\]
Every integer power vector and affine lift in this sector is
allowed exactly when
\[
 \gcd(m_i,n_i,r_i)=1,\qquad d_i\mid m_i/t_i.                 \tag{D.UF1}\label{maps:eq:D.UF1}
\]
This is an exact freeness criterion for all powers of the
cyclic action. Under it $t_i$ is either $1$ or $q_i$; the latter
case, when $q_i>1$, is permitted precisely when $q_i\mid k_i$.
The finite central quotient surface $F_i$ and its complex
normal line $L_i$ satisfy
\[
\begin{gathered}
 H_1(F_i)=\mathbb Z^2\oplus\mathbb Z/t_i,\qquad
 H_2(F_i)=\mathbb Z^2\oplus\mathbb Z/t_i,\\
 H_3(F_i)=\mathbb Z^2,\qquad H_0(F_i)=H_4(F_i)=\mathbb Z,\\
 \operatorname{ord}c_1(L_i)=t_i.
\end{gathered}                                               \tag{D.UF2}\label{maps:eq:D.UF2}
\]
Here order one means the zero class. Thus the normal line need
not be topologically trivial.

The full global degree-two presentation has a fibre generator
$s$, two fixed sweeps $D_i,L_i'$ for every finite log, and the
four sweeps $t_j(\Lambda)$ for every regular log. The finite
relations are
\[
\begin{aligned}
 (bn_i/q_i)s+(m_i/q_i)D_i&=0,\\
 -br_i s+m_iL_i'&=0,\\
 r_iD_i+n_iL_i'&=0.
\end{aligned}                                                \tag{D.UF3}\label{maps:eq:D.UF3}
\]
The regular-log relations remain those of \eqref{maps:eq:D.SC2}, and the
single global relation is
\[
 \sum_iD_i+\sum_jt_j(\delta)+be\,s=0.                       \tag{D.UF4}\label{maps:eq:D.UF4}
\]
There is no coprimality assumption on $m_i,n_i$ beyond the
actual freeness criterion \eqref{maps:eq:D.UF1}. The regular logs still have
arbitrary free affine vectors.

Let $T$ be the torsion of the cokernel of this matrix. Let
$\mathbb Z^r\oplus A_1$ be the Smith cokernel of the full
Seifert matrix of \ref{maps:src:D.2.23}, including these finite numerators
and all regular numerators. Then $r=0$ or $1$, and the complete
integral homology is
\[
\begin{array}{c|c}
 j&H_j(X;\mathbb Z)\\\hline
0,6&\mathbb Z\\
1&\mathbb Z^r\oplus A_1\\
2,3&\mathbb Z^{2r}\oplus T\\
4&\mathbb Z^{2r}\oplus A_1\\
5&\mathbb Z^r.
\end{array}                                                  \tag{D.UF5}\label{maps:eq:D.UF5}
\]
The complete fundamental group is still the Seifert
presentation, rather than its abelianization. In particular
this extends the non-star minimal-filling homology algorithm
to every free affine twist at each of its primitive-$\gamma$
finite sites.

\end{theorem}


\begin{proof}
\begin{enumerate}[label=\arabic*.,leftmargin=*]
\item \emph{1. Exact freeness and the cap groups.}
Let $A$ be the finite monodromy, of order $d$, and omit the
site index. Its nontrivial eigenvalues are the two eigenvalues
of the corresponding finite elliptic monodromy. Consequently
$A^j$ has the same rational fixed subspace as $A$ whenever
$d\nmid j$. Write $P$ for the rational projection to this fixed
subspace. The integral certificates of \ref{maps:src:D.2.26} give
\[
 \operatorname{Fix}(A)=\mathbb Z\delta\oplus\mathbb Z\ell,
 \qquad P\Lambda=\mathbb Z(\delta/q)\oplus\mathbb Z\ell,
 \qquad N_m\beta=(m/q)\delta,                               \tag{D.UF6}\label{maps:eq:D.UF6}
\]
where $N_m=1+A+\cdots+A^{m-1}$ and $\beta$ is the actual
primitive vertical coinvariant used there. The same certificates
give
$\Lambda/(A-1)\Lambda=\mathbb Z\beta\oplus\mathbb Z\ell$,
with $\delta=q\beta$ in that quotient.

\item Every affine action with the prescribed power is conjugate
on the good-reduction root disc to the invariant translation
form $h(x)=Ax+v/m$. Indeed its moving translation part has
zero finite-group average and is a translation coboundary;
the averaging solution is holomorphic on the disc. This is
the actual affine comparison of \ref{maps:src:D.2.26}, and it does not use
$\gcd(m,n)=1$. Constant unit choices extend over the cap and
preserve the marked normal comparison.

\item If $g=\gcd(m,n,r)>1$, then $h^{m/g}$ has integral translation
part and fixes the origin of its linear action. Freeness
therefore requires $g=1$. Assume this now. A power with
$d\mid j$ has identity linear part and nonintegral translation,
so is free. For $d\nmid j$, its moving linear map is surjective
on the moving torus. A fixed point exists precisely when its
fixed quotient translation is zero. By \eqref{maps:eq:D.UF6} this is
\[
 m\mid jn,\qquad m\mid jqr.
\]
The relevant powers are exactly the nonzero multiples of
$m/t$ below $m$, where $t=\gcd(m,n,qr)$. Since $g=1$ and
$q\in\{1,2,3\}$, $t$ divides $q$. Those powers all have
identity linear part exactly when $d\mid m/t$. This proves
\eqref{maps:eq:D.UF1} and the stated alternative for $t$. Freeness on the
central compact torus implies freeness on a sufficiently small
root neighborhood for this finite group, so the criterion is
also sufficient for the smooth holomorphic filling.

\item The quotient $F=T^4/\langle h\rangle$ has the affine fundamental
group presentation with the extra generator $c$ and the power
relation $mc=v$. On abelianizing, \eqref{maps:eq:D.UF6} gives
\[
 H_1(F)=\langle\beta,\ell,c\mid
                   mc-n\ell-qr\beta=0\rangle
        =\mathbb Z^2\oplus\mathbb Z/t.                     \tag{D.UF7}\label{maps:eq:D.UF7}
\]
It is an oriented four-manifold of Euler number zero and
rational first Betti number two. Integral Poincare duality and
the universal coefficient theorem give the other groups in
\eqref{maps:eq:D.UF2}, including the order-$t$ summand of $H_2$.

\item The normal line is the associated line for the root rotation
character. On $H_1(F)$ that character sends $c$ to $1/m$ in
$\mathbb R/\mathbb Z$ and the original fibre generators to
zero. The torsion generator
\[
 (m/t)c-(n/t)\ell-(qr/t)\beta
\]
has character $1/t$. The Bockstein of this character is
$c_1(L)$, so its order is exactly $t$. This proves the last
assertion of \eqref{maps:eq:D.UF2} without treating a torsion Chern class as
zero.

\item \emph{2. The full integral Thom kernel.}
The boundary is the marked mapping torus of $A$ and has
$H_1=\mathbb Z^3$, $H_2=\mathbb Z^4$. These statements follow
from the integral local Smith matrices, whose nonzero factors
are all units, in every one of the $68$ positions. The cap is
the disk bundle of $L\to F$. Since $c_1(L)$ is torsion, the
homology Gysin map $H_2(F)\to H_0(F)$ is zero. The pair sequence
therefore makes $H_2(\partial U)\to H_2(U)=H_2(F)$ onto,
including its torsion summand.

\item The full kernel is the image of the Thom map
\[
 H_3(U,\partial U)=H_1(F)\longrightarrow H_2(\partial U).
\]
The normal circle has the central group representative
$\mu=c^m v^{-1}$. For a lift of each of $\beta,\ell,c$, sweep
that lift by $\mu$. These three normal tori generate the
entire Thom image. On the boundary mapping cylinder, their
normal projections and their fibre components are respectively
\[
\begin{array}{c|c|c}
\text{generator}&\text{normal projection}&
                          \text{fibre component}\\\hline
\beta&N_m\beta&-\beta\wedge v\\
\ell&m\ell&-\ell\wedge v\\
c&v&0.
\end{array}                                                  \tag{D.UF8}\label{maps:eq:D.UF8}
\]
For the first row, subdividing the $m$-fold normal circle into
its $m$ mapping-cylinder intervals gives the sum $N_m\beta$;
the final power-vector edge supplies the displayed fibre
parallelogram. The second row is the same construction for an
invariant loop. The third is the torus of the two commuting
words $c$ and $c^m v^{-1}$, so its only normal sweep is $v$.
These are integral marked cycles, not invariant rational
forms.

\item Evaluation by the primitive exterior functional $\omega$,
using $-\omega(\beta,\ell)=b/q$ and
$\omega(\ell,\delta)=b$, turns \eqref{maps:eq:D.UF8} into the three columns
\[
 \begin{pmatrix}bn/q\\m/q\\0\end{pmatrix},\qquad
 \begin{pmatrix}-br\\0\\m\end{pmatrix},\qquad
 \begin{pmatrix}0\\r\\n\end{pmatrix}                       \tag{D.UF9}\label{maps:eq:D.UF9}
\]
on $(s,D,L')$. The target $H_2(\partial U)$ is free, so the
order-$t$ class of $H_1(F)$ maps to zero. Accordingly the sole
primitive relation among these columns is
\[
 (m/t)K_c-(n/t)K_\ell-(qr/t)K_\beta=0.
\]
No further saturation is allowed or needed: they are images
of generators of the whole Thom source. The projected
lattice has the exact index
\[
 \gcd\left(m^2/q,mn/q,mr\right)=(m/q)t.                    \tag{D.UF10}\label{maps:eq:D.UF10}
\]
In particular a torsion cap multiplies this index by $t$.
When $\gcd(m,n)=1$, the same columns reduce by Bezout to
\eqref{maps:eq:D.HY3}, recovering the previous theorem.

\item \emph{3. Global gluing and duality.}
Every degree-two boundary map is onto for these finite caps,
for the regular logs and for the retained minimal original and
Mumford caps. Eliminate their degree-two groups in
Mayer--Vietoris. The exterior fixed-sweep lattice is the
primitive direct sum with its single diagonal $\delta$
relation, as proved for all $172$ rows in \ref{maps:src:D.2.26}.
The full kernels \eqref{maps:eq:D.UF9}, the regular-log kernels and that
one diagonal relation give exactly \eqref{maps:eq:D.UF3}--\eqref{maps:eq:D.UF4}.
This calculation retains the cap torsion through its full
Thom image; it does not assume that the cap or its normal
line has torsion-free invariants.

\item The rational rank of this cokernel is $r$: after solving the
local sweep relations, the only remaining scalar equation is
$b(e-\sum n_i/m_i)s=0$, with the regular numerators included.
It has the same zero criterion as the Seifert matrix. The
kernel of the degree-one Mayer--Vietoris map is a subgroup
of the free boundary group and is therefore free. Its rank
is also $r$, because the source and target have equal free
rank and their cokernel has free rank $r$. Thus its extension
by the degree-two cokernel splits as an abelian group, giving
$H_2=\mathbb Z^{2r}\oplus T$. All caps still have the same
Euler contribution, so $\chi(X)=2$. Duality and the universal
coefficient theorem now give every remaining entry of \eqref{maps:eq:D.UF5}.
This includes $e-\sum n_i/m_i=0$ and does not suppress its
free groups or its torsion.
\end{enumerate}
\end{proof}

For example, take the $IV^*,II,I_2$ sphere initial row
$(a,b)=(1,3)$ of \ref{maps:src:D.2.11}, keep $II$ unchanged and use
$(q,m,n,r)=(3,9,3,1)$ at $IV^*$. The action is free, although
$\gcd(m,n)=3$. The result has
\[
\begin{gathered}
 \pi_1=\mathbb Z/3,\qquad H_1=H_4=\mathbb Z/3,\\
 H_2=H_3=(\mathbb Z/3)^2,\qquad H_5=0.
\end{gathered}
\]
Extrapolating the former primitive-numerator formula would
miss one order-three summand. For a $III^*$ sphere initial
row with $b=2$, the free data $(q,m,n,r)=(2,8,0,1)$ instead
give $\pi_1=\mathbb Z$, middle groups
$H_2=H_3=\mathbb Z^2\oplus\mathbb Z/2$, and
$H_4=\mathbb Z^2$, $H_5=\mathbb Z$.

The exact program
\path{code/III/manifolds_unrestricted_finite_logs.py} compares
\eqref{maps:eq:D.UF1} with every actual affine-power fixed-point test for
all $68$ marked matrices, $1\le k\le6$ and every pair
$(n,r)\bmod m$: $150059$ exact tests. It also compares $305$
full unreduced graph/normal Smith matrices with the reduced
presentation, retaining $99$ torsion-cap instances and $55$
zero-rational-Euler instances. These bounds describe the
independent checks; the proof and the matrix algorithm above
hold for every positive $k$ and every integer twist.


\begin{theorem}[every finite additive lattice and every free twist]\label{maps:src:D.2.30}
\StatementRecord{maps:src:D.2.30}{proved}{Historical events: Claude Fable 5.2 (Anthropic), effort high, 16 November 2026: carried as background without fault found. GPT-6 Astra (OpenAI), ChatGPT Pro mode, effort Pro, S6\_09\_PAPERS, 17 November 2026: targeted proof check. Also checked in this writing. Pre-continuation writing events: GPT-6 Astra (OpenAI), ChatGPT Pro mode, effort Pro, S6\_10\_PAPERS, 18 November 2026: targeted proof check. Continuation metadata events: GPT-6 Astra (OpenAI), ChatGPT Pro mode, effort Pro, S6\_10\_PAPERS continuation, 19 November 2026: used and rechecked in this run; status-vocabulary reconciliation; metadata reconciliation only; no new mathematical check. \emph{Pre-continuation recorded proof scope.} GPT-6 Astra (OpenAI, ChatGPT Pro mode), effort Pro, 2026-11-17 UTC; Claude Fable 5.2 (Anthropic), effort high, 16 November 2026; also checked in this writing by GPT-6 Astra (OpenAI), ChatGPT Pro mode, effort Pro, 18 November 2026}{the hypotheses, notation and definitions stated here; \cref{maps:src:D.2.29}; \cref{maps:eq:D.HY2}; \cref{maps:eq:D.UF2}; \cref{maps:src:D.2.25}; \cref{maps:eq:D.UF4}; \cref{maps:eq:D.UF5}; \cref{maps:src:D.2.23}; \cref{maps:src:D.2.28}; \cref{maps:src:D.2.26}; \cref{maps:eq:D.AF8}; \cref{maps:eq:D.UF8}}


The fixed-$\gamma$ restriction in \ref{maps:src:D.2.29} can also be removed.
All $102$ marked finite additive positions on the $172$ natural
candidate rows have the following integral description.
Let $d$ be the finite monodromy order, $q$ its index from
\eqref{maps:eq:D.HY2}, and put
\[
 \gamma^*(\operatorname{Fix}A)=g\mathbb Z,\qquad
 g\in\{1,q\},\qquad q_v=q/g.
                                                               \tag{D.AF1}\label{maps:eq:D.AF1}
\]
Here $g$ denotes the positive generator of that image.
There are $68$ positions with $g=1$ and $34$ with $g=q>1$.
Choose integral bases of the fixed and coinvariant lattices
so that
\[
\begin{gathered}
 \operatorname{Fix}A=\mathbb Z\delta\oplus\mathbb Z\ell,
 \qquad \gamma^*(\ell)=g,\\
 \Lambda/(A-1)\Lambda=\mathbb Z\beta\oplus\mathbb Z\eta,
 \qquad [\delta]=q_v\beta,\quad [\ell]=g\eta,\\
 N_m\beta=(m/q_v)\delta,
 \qquad N_m\eta=(m/g)\ell.
\end{gathered}                                                \tag{D.AF2}\label{maps:eq:D.AF2}
\]
The lifts $\beta,\eta$ are part of the retained integral
marking. In particular $\gamma^*(\eta)=1$.

For $m=dk>0$, every power vector is $v=r\delta+n\ell$.
Put
\[
 t=\gcd(m,q_vr,gn).
\]
The exact free-action criterion is
\[
 \gcd(m,n,r)=1,\qquad d\mid m/t.                            \tag{D.AF3}\label{maps:eq:D.AF3}
\]
The finite quotient and its normal line have precisely the
groups and the order-$t$ Chern class in \eqref{maps:eq:D.UF2}.
Their full normal Thom kernel, on the fibre generator $s$
and the sweeps $D,L'$ of $\delta,\ell$, has the three columns
\[
 \begin{pmatrix}bgn/q_v\\m/q_v\\0\end{pmatrix},\qquad
 \begin{pmatrix}-br\\0\\m/g\end{pmatrix},\qquad
 \begin{pmatrix}0\\r\\n\end{pmatrix}.                     \tag{D.AF4}\label{maps:eq:D.AF4}
\]
All entries are integers. The actual Seifert numerator is
$gn$, rather than $n$.

Consequently all $198$ minimal original/log additive-filling
patterns on all $148$ nonstar rows of the natural candidate
catalogue have a complete integral homology algorithm, for
every logarithmic multiplicity and every free affine twist.
Use \eqref{maps:eq:D.AF4} at each finite site, the full regular-log columns
of \ref{maps:src:D.2.25} at arbitrarily many regular sites, and the single
relation \eqref{maps:eq:D.UF4}. Take the torsion $T$ of the resulting Smith
cokernel, and take $r,A_1$ from the Seifert matrix with the
actual numerators $g_i n_i$. The answer is exactly \eqref{maps:eq:D.UF5},
including its zero-rational-Euler case. The complete
fundamental group remains \ref{maps:src:D.2.23}. The data of all $198$
patterns and all $102$ integral finite lattices are saved in
\path{records/manifolds_all_finite_logs.json}.

For one finite log and no regular logs there is also a closed
Smith formula, retaining an arbitrary integer clutch $e$.
Set
\[
 n'=n-(m/g)e,
 \qquad h=\gcd(b,m/g,n').                                  \tag{D.AF5}\label{maps:eq:D.AF5}
\]
If $n'\ne0$, then
\[
\begin{gathered}
 \pi_1(X)=\mathbb Z/(g|n'|),\qquad
 H_1(X)=H_4(X)=\mathbb Z/(g|n'|),\\
 H_2(X)=H_3(X)=
 \mathbb Z/h\oplus\mathbb Z/\bigl(bg|n'|t/(qh)\bigr),
 \qquad H_5(X)=0.
\end{gathered}                                               \tag{D.AF6}\label{maps:eq:D.AF6}
\]
When $n'=0$, instead
\[
\begin{gathered}
 \pi_1(X)=\mathbb Z,\qquad H_1(X)=H_5(X)=\mathbb Z,\\
 H_2(X)=H_3(X)=\mathbb Z^2\oplus
                         \mathbb Z/\gcd(b,m/g),\qquad
 H_4(X)=\mathbb Z^2.
\end{gathered}                                               \tag{D.AF7}\label{maps:eq:D.AF7}
\]
In both displays $H_0=H_6=\mathbb Z$, and a cyclic summand of
order one is omitted. Every case with $g>1$ is necessarily
nonsimply connected, so this extension does not alter \ref{maps:src:D.2.28}.

\end{theorem}


\begin{proof}
\begin{enumerate}[label=\arabic*.,leftmargin=*]
\item The $g=1$ basis and its full marked comparison were proved in
\ref{maps:src:D.2.26} and \ref{maps:src:D.2.29}. For the remaining positions the exact
matrices give
\[
 g=q,\qquad
 \Lambda/(A-1)\Lambda=\mathbb Z[\delta]\oplus\mathbb Z[\gamma],
 \qquad
 \ell=(q/d)N_d\gamma\in\Lambda.                            \tag{D.AF8}\label{maps:eq:D.AF8}
\]
Thus take $\beta=\delta$ and $\eta=\gamma$.
The vectors $\delta,\ell$ are an integral basis of the fixed
lattice, and their coinvariant images are $\delta,q\gamma$.
These assertions are checked by unimodular Smith operations
for all $34$ actual local matrices, not only by rational
conjugacy. They prove \eqref{maps:eq:D.AF2}. In all $102$ positions the
nonzero Smith factors of both $A-1$ and
$\bigwedge^2A-1$ are units. Hence the boundary groups used
in \ref{maps:src:D.2.29} are again free of the stated ranks.

\item The invariant exterior form satisfies
\[
 -\omega(\beta,\ell)=bg/q_v,\qquad
 \omega(\eta,\delta)=b,\qquad
 \omega(\eta,\ell)=0.
\]
For $g=q$ these follow directly from \eqref{maps:eq:D.AF8}, invariance of
$\omega$, and $\omega(\delta,\ell)=-bq$.
For $g=1$ they are the earlier identities. The required
divisibility $q_v\mid b$ is part of the same integral
catalogue, so the first coefficient of \eqref{maps:eq:D.AF4} is integral.

\item The fixed quotient torus has lattice $P\Lambda$, and the
coordinates of $v$ in it are $(q_vr,gn)$ by \eqref{maps:eq:D.AF2}.
Repeating the actual-power fixed-point argument of \ref{maps:src:D.2.29}
therefore gives \eqref{maps:eq:D.AF3}. The integer $t$ divides $q$ under
this freeness condition. The quotient surface has
\[
 H_1(F)=\langle\beta,\eta,c\mid
                   mc-q_vr\beta-gn\eta=0\rangle
       =\mathbb Z^2\oplus\mathbb Z/t.
\]
Its normal root character has value $1/t$ on the torsion
generator, exactly as before. Thus the normal line has Chern
class of order $t$, and duality gives all the cap groups.

\item The three full Thom generators in \eqref{maps:eq:D.UF8}, with $\eta$ in
place of the former $\ell$ coinvariant generator, now have
normal projections $(m/q_v)\delta,(m/g)\ell,v$.
Evaluation of their actual fibre parallelograms by $\omega$
gives \eqref{maps:eq:D.AF4}. Their primitive relation is
\[
 (m/t)K_c-(q_vr/t)K_\beta-(gn/t)K_\eta=0,
\]
and their normal-projection index is still $(m/q)t$.
The boundary-to-cap map in degree two remains onto, including
the torsion summand, because the normal Chern class is
torsion. Thus the full Mayer--Vietoris and duality argument of
\ref{maps:src:D.2.29} applies with the actual numerator $gn$ at every site.
This proves the all-order matrix algorithm.

\item To check its scope, each finite additive site can now be
log-transformed, including those with $g>1$; for example a
power along the primitive fixed quotient direction gives a
free action. An original cap is an additional choice exactly
when its line multidegree is zero and its component action
has one orbit. Those are precisely the original caps whose
full integral comparison is already proved in \ref{maps:src:D.2.26}.
Enumerating these choices gives $198$ patterns on all $148$
nonstar initial rows. It introduces no further multiplicity
restriction. Zero-free multiplicative sites and the unique
rank-two zero retain the marked minimal models already used
in that catalogue.

\item For the one-log formula eliminate $D=-be\,s$ in \eqref{maps:eq:D.AF4},
and subtract $e$ times the second column from the third.
On $(s,L')$ the matrix becomes
\[
 \begin{pmatrix}
  bgn'/q_v&-br&0\\
  0&m/g&n'
 \end{pmatrix}.                                            \tag{D.AF9}\label{maps:eq:D.AF9}
\]
Changing $n$ by $(m/g)e$ is exactly a peripheral lift change
by the coinvariant $\eta$. Thus \eqref{maps:eq:D.AF3} is preserved and
$t=\gcd(m,q_vr,gn')$. If $n'\ne0$, the gcd of the entries is
$h=\gcd(b,m/g,n')$: any prime dividing both $m/g$ and $n'$
cannot divide $r$ by freeness. The gcd of the two-by-two
minors is
\[
 bg|n'|t/q.
\]
These are the two determinantal divisors and prove the
middle groups in \eqref{maps:eq:D.AF6}. The Seifert presentation gives the
cyclic group of order $g|n'|$. If $n'=0$, the sole remaining
relation is $-brs+(m/g)L'=0$, whose torsion divisor is
$\gcd(b,m/g)$ since $r$ is prime to $m$ in this free case.
Its free rank one acquires the one free degree-one bridge
class, giving the rank two in \eqref{maps:eq:D.AF7}. Duality gives the
remaining groups in both cases.

\item Finally, if $g=q>1$, then $q$ divides both $m$ and its actual
numerator $gn$. Such a multiple fibre cannot occur in a
simply connected Seifert output by \ref{maps:src:D.2.23}. This verifies
that the new sector consists entirely of nonsimply-connected
members, without assuming that the group is abelian when
there are several logarithmic sites.
\end{enumerate}
\end{proof}

The exact program \path{code/III/manifolds_all_finite_logs.py}
checks all $34$ additional lattice bases, compares the
freeness criterion with $38675$ actual-matrix power tests
for $1\le k\le6$ and all residues, and compares $175$ full
unreduced normal/graph Smith matrices to \eqref{maps:eq:D.AF4}, including
$35$ torsion-cap instances. It also checks $725$ one-log
Smith forms with five different integer clutch values against
\eqref{maps:eq:D.AF6}--\eqref{maps:eq:D.AF7}. The finite checks certify the explicit
constant matrices and examples; the proof covers all orders
and all free integer twists. Initial weights outside the
finite sphere-candidate catalogue and nonminimal original
caps remain separate parts of the broader D.2 request.


\begin{theorem}[all zero-degree original caps and their free defects]\label{maps:src:D.2.31}
\StatementRecord{maps:src:D.2.31}{proved}{Historical events: Claude Fable 5.2 (Anthropic), effort high, 16 November 2026: carried as background without fault found. GPT-6 Astra (OpenAI), ChatGPT Pro mode, effort Pro, S6\_09\_PAPERS, 17 November 2026: targeted proof check. Also checked in this writing. Pre-continuation writing events: GPT-6 Astra (OpenAI), ChatGPT Pro mode, effort Pro, S6\_10\_PAPERS, 18 November 2026: targeted proof check. Continuation metadata events: GPT-6 Astra (OpenAI), ChatGPT Pro mode, effort Pro, S6\_10\_PAPERS continuation, 19 November 2026: used and rechecked in this run; status-vocabulary reconciliation; metadata reconciliation only; no new mathematical check. \emph{Pre-continuation recorded proof scope.} GPT-6 Astra (OpenAI, ChatGPT Pro mode), effort Pro, 2026-11-17 UTC; Claude Fable 5.2 (Anthropic), effort high, 16 November 2026; also checked in this writing by GPT-6 Astra (OpenAI), ChatGPT Pro mode, effort Pro, 18 November 2026}{the hypotheses, notation and definitions stated here; \cref{maps:src:D.2.30}; \cref{maps:src:D.2.23}; \cref{maps:src:D.2.14}; \cref{maps:src:D.2.26}; \cref{maps:src:D.2.16}; \cite[Section 2.4, equation (2.2)]{manUU21}; \cref{maps:eq:D.OD4}; \cref{maps:eq:D.OD5}; \cite[Theorems 2.2--2.3]{manBarden65}; \cite[Theorem A]{manGR25}}


In the natural candidate catalogue we may retain every original
additive cap at which $P$ is narrow, without requiring a single
component orbit. At $I_n^*$ this theorem uses the original cap;
finite good-reduction caps at the other additive sites have all
the choices and twists of \ref{maps:src:D.2.30}. The zero-free multiplicative
caps and the unique rank-two zero retain their preceding marked
models. These choices give $228$ additive-filling patterns on
$160$ natural initial rows. Exactly $30$ patterns have a positive
original component defect. The complete data are
\path{records/manifolds_original_zero_degree.json}.

For an original additive cap let $\rho_i$ be the number of
orbits of $Q$ on its components, and put
\[
 D=\sum_{i\text{ original additive}}(\rho_i-1).
\]
The only nonzero local defects in this catalogue are
\[
\begin{array}{c|rrrr}
\text{type}&III^*&IV^*&I_3^*&I_1^*\\\hline
\rho_i&5&3&3&2\\
\rho_i-1&4&2&2&1.
\end{array}                                                   \tag{D.OD1}\label{maps:eq:D.OD1}
\]
Every other permitted original cap has $\rho_i=1$.

Compute the Seifert group and its Smith cokernel
$\mathbb Z^r\oplus A_1$ as before, treating an unchanged cap
as an order-one meridian with zero numerator in the canonical
marking. Let $T$ be the torsion of the matrix in \ref{maps:src:D.2.30},
using its finite and regular columns and setting both fixed
sweeps of every original additive cap to zero. Then the full
integral groups are
\[
\begin{array}{c|c}
j&H_j(X;\mathbb Z)\\\hline
0,6&\mathbb Z\\
1&\mathbb Z^r\oplus A_1\\
2&\mathbb Z^{2r+D}\oplus T\\
3&\mathbb Z^{2r+2D}\oplus T\\
4&\mathbb Z^{2r+D}\oplus A_1\\
5&\mathbb Z^r.
\end{array}                                                   \tag{D.OD2}\label{maps:eq:D.OD2}
\]
The complete fundamental group is \ref{maps:src:D.2.23}, including every
order and affine twist. Thus the component defect adds free
homology, with no uncomputed torsion extension.

There is a full smooth identification in the $b=1$ sector.
For every simply connected member, let $M$ be the product of
its regular logarithmic multiplicities. Then
\[
 X\cong
 \mathop{\#}^{D}(S^2\times S^4)
 \ \#\ \mathop{\#}^{D}(S^3\times S^3)
 \ \#\ \widetilde\Sigma_0 N_M,                             \tag{D.OD3}\label{maps:eq:D.OD3}
\]
with the oriented twisted suspension and the spin five-manifold
$N_M$ defined in \ref{maps:src:D.2.14}. The allowed free defects in this sector
are $D=0,1,2,4$. The identification retains the full refined
spin data, including even $M$; it is not inferred from
additive homology alone.

\end{theorem}


\begin{proof}
\begin{enumerate}[label=\arabic*.,leftmargin=*]
\item \emph{1. The actual original cap.}
At an additive fibre the condition that $P$ is narrow means
that $\mathcal O(P-O)$ has degree zero on every component.
Its underlying line on the rational Kodaira tree is trivial.
The unit equivariant multiplier can be removed by its
rigidified logarithm isotopy, exactly as in \ref{maps:src:D.2.26}; the tree
has $H^1=0$, and the isotopy has zero meridian winding on the
normal collar. The cap consequently retracts, with its marked
angular circle, onto
\[
 S^1_\delta\times\operatorname{MT}_Q(K),
\]
where $K$ is the underlying Kodaira tree and $Q$ acts by its
actual component translation. The nontrivial contraction
modulus is retained as the mapping-torus direction; only its
unit circle part is straightened. This is the same relative
comparison proved for the odd-star caps in \ref{maps:src:D.2.16}, and uses
only narrowness of $P$, not the number of component orbits.

\item The tree has $H_1(K)=0$ and one free degree-two class per
component. The permutation invariants and coinvariants are
both free of rank $\rho$. The mapping-torus sequence and
Kunneth therefore give
\[
 H_1(U)=\mathbb Z^2,\quad H_2(U)=\mathbb Z^{\rho+1},\quad
 H_3(U)=\mathbb Z^{2\rho},\quad H_4(U)=\mathbb Z^\rho.
                                                               \tag{D.OD4}\label{maps:eq:D.OD4}
\]
There is no permutation torsion in these groups.

\item The two full normal tori are still primitive. The $\delta$
normal torus bounds the section disc times its circle. For
the other, take the least power $Q^q$ that is narrow and its
local Neron logarithm. Its marked loop extends across the
base disc; the rigidified norm of the unit linearization has
no meridian winding, so it introduces no fibre correction.
The local existence and extension of this logarithm are as
in \cite[Section 2.4, equation (2.2)]{manUU21} and \ref{maps:src:D.2.16}.
For each actual matrix the fixed lattice has basis
$\delta,\ell$, with $\gamma^*(\ell)=q$. Thus these two
normal tori project to the \emph{entire} integral fixed
lattice, not a finite-index sublattice. Their values under
the exterior fibre functional are zero in the stated
rigidified marking.

\item We also retain the component cokernel. Since \eqref{maps:eq:D.OD4} is free,
Poincare--Lefschetz duality gives
$H_2(U,\partial U)=H^4(U)=\mathbb Z^\rho$.
The map from $H_2(U)$ is the folded component-intersection
map. On the component classes its rank is $\rho-1$:
the Kodaira intersection matrix is negative semidefinite
with one-dimensional radical generated by the fibre
multiplicity vector, and averaging to permutation invariants
preserves that radical. The exact folded matrices are also
retained in the certificate. The principal fibre relation
bounds the full rank by $\rho-1$, so the additional
$\delta$--mapping-torus class cannot increase it.

\item It follows from the pair sequence that
\[
 \operatorname{im}\bigl(H_2(\partial U)\to H_2(U)\bigr)
 \text{ is primitive of rank }2,
 \qquad
 H_2(U)/\operatorname{im}H_2(\partial U)
                 \cong\mathbb Z^{\rho-1}.                 \tag{D.OD5}\label{maps:eq:D.OD5}
\]
Primitivity is integral: the quotient is the image of a map
into the free group $H_2(U,\partial U)$.
The boundary $H_2$ is free of rank four in every such marked
additive case. Its kernel has rank two, and the two primitive
fixed-loop normal tori just constructed therefore exhaust it.
Indeed their projection is already the whole fixed lattice,
and the fibre part is injective rationally. This proves that
the only change in the local Mayer--Vietoris data is the
free summand in \eqref{maps:eq:D.OD5}.

\item At a zero-free multiplicative site in this natural catalogue,
$P$ narrow implies that $Q$ is transitive on its components.
The direct component data and the primitive vanishing
matrices verify this at every occurrence. Hence an allowed
zero-degree original multiplicative cap has $\rho=1$ and
is already covered by the earlier primitive comparison.
The additive component permutations give exactly \eqref{maps:eq:D.OD1}.

\item \emph{2. Global integral homology.}
The subgroups in \eqref{maps:eq:D.OD5} are primitive, so choose integral
splittings in all the original cap groups. Eliminating their
boundary images in Mayer--Vietoris leaves their direct sum
$\mathbb Z^D$, together with exactly the torsion presentation
of \ref{maps:src:D.2.30}. The other cap boundary maps are onto as proved
there. The original caps have the same first Betti number as
their finite or minimal comparisons, so the degree-one bridge
kernel is still free of rank $r$. Consequently
$H_2=\mathbb Z^{2r+D}\oplus T$.
Every added original cap has Euler number zero, and
$\chi(X)=2$ remains unchanged. Duality gives all of \eqref{maps:eq:D.OD2},
including $b_3=2r+2D$ and the paired middle torsion.
This proves the formula without an implicit integral
spectral-sequence splitting.

\item There are $396$ exact local zero-degree original certificates.
Enumerating ``original'' precisely when $P$ is narrow,
and ``finite log'' at every finite good-reduction site,
with an original cap required at $I_n^*$ in the present
sector, gives the stated $228$ patterns on $160$ rows.
No logarithmic order is bounded in this enumeration.

\item \emph{3. Smooth classification when $b=1$.}
For a simply connected member all remaining finite logs have
primitive actual gamma numerator. Since $b=1$, their indices
$q_i$ are one; a site of fixed-gamma index greater than one
would contradict simple connectivity. \ref{maps:src:D.2.30} and the
regular-log calculation therefore give
\[
 H_2(X)=\mathbb Z^D\oplus(\mathbb Z/M)^2,
 \qquad H_3(X)=\mathbb Z^{2D}\oplus(\mathbb Z/M)^2.
\]
The scalar circle is free on every original cap, by the
product description in part 1, and free on every finite
or regular logarithmic cap by the primitive gamma numerator.
At each zero-free multiplicative cap, $Q=G$ is transitive
and the pair consisting of its actual primitive vanishing
vector and $\delta$ has unit Smith factors. The program
checks this pair at every such site in the relevant patterns.
The rank-two zero has the same full original short lattice
and the same unique fixed sphere as \ref{maps:src:D.2.14}, with normal
bundle $\mathcal O(-1)_{+1}\oplus\mathcal O(-1)_{-1}$.
Thus this is the actual semifree circle action on $X$.

\item The fixed sphere is nullhomotopic. Its specialization comes
from the common fibre class $s$; it has no contribution to
the extra component cokernels in \eqref{maps:eq:D.OD5}. The submatrix on
$s$ and the delta sweeps is the unimodular Seifert matrix
when $b=1$ and $\pi_1(X)=1$. It therefore kills $s$
integrally. Hurewicz gives the assertion about the sphere.

\item We check spin despite the new free second homology. On every
original cap the holomorphic volume form is the product of
a nonvanishing canonical two-form on the local minimal
elliptic surface with $dv/v$. Translation by $Q$ preserves
the first factor, and the change of $dv/v$ is a base-surface
one-form whose wedge with that two-form is zero. The volume
form descends, so the canonical line of that cap is trivial,
as in \ref{maps:src:D.2.16}. The free part of $H_2(X)$ is generated by the
component cokernels of these caps. Hence $c_1(X)$ evaluates
to zero on it. Since $X$ is simply connected,
$H^2(X;\mathbb Z)$ is free, so $c_1(X)=0$ and $X$ is spin.

\item The equivariant surgery of \ref{maps:src:D.2.14} now gives
$Y=X\#(S^3\times S^3)$ with a free circle action and a
projectable spin structure. Its orbit five-manifold is
simply connected and spin, and the homotopy/Gysin sequences
give
\[
 (B,e)\cong
 \left(\mathop{\#}^{D+1}(S^2\times S^3)\#N_M,
                        (x,0,\ldots,0)\right).
\]
Here $e$ is primitive; the spin Smale--Barden classification
and its realization theorem identify this pair, as in
\cite[Theorems 2.2--2.3]{manBarden65}.
Apply the spin branch of \cite[Theorem A]{manGR25}
successively, retaining the zero Euler class on the other
summands. Together with its suspension formula this gives
\[
 Y\cong\mathop{\#}^{D}(S^2\times S^4)
       \#\mathop{\#}^{D+1}(S^3\times S^3)
       \#\widetilde\Sigma_0N_M.
\]
The full Zhubr classification permits the same cancellation
of one $S^3\times S^3$ summand as in \ref{maps:src:D.2.14} and \ref{maps:src:D.2.16};
it retains the refined class $p$ and the full class
$\mu\in H_6(K(H_2,2);\mathbb Z)$, including all free cubic,
mixed and finite data. The added $S^3\times S^3$ maps
constantly to this Eilenberg--Mac Lane space and spin-bounds
$S^3\times D^4$. It therefore changes neither bordism
invariant and adds only two to $b_3$. The diffeomorphism of
the stabilized manifolds identifies these complete data,
which classify the unstabilized manifolds as well.
This proves \eqref{maps:eq:D.OD3}, including its oriented framing and
all two-primary data.
\end{enumerate}
\end{proof}

The exact program \path{code/III/manifolds_original_zero_degree.py}
checks all $396$ local original-cap lattices and folded
intersection ranks, the $228$ additive patterns, and $144$
actual primitive scalar-circle pairs in the $b=1$ patterns
that can be simply connected. It recovers the original
odd-star free defects of \ref{maps:src:D.2.16}--\ref{maps:src:D.2.17} and adds the same uniform
comparison for every permitted narrow-$P$ section pair.
Original caps with nonzero component Chern vector and higher
initial linearization weights remain outside this theorem.


\begin{theorem}[all original additive choices on the natural catalogue]\label{maps:src:D.2.32}
\StatementRecord{maps:src:D.2.32}{proved}{Historical events: Claude Fable 5.2 (Anthropic), effort high, 16 November 2026: carried as background without fault found. GPT-6 Astra (OpenAI), ChatGPT Pro mode, effort Pro, S6\_09\_PAPERS, 17 November 2026: targeted proof check. Also checked in this writing. Pre-continuation writing events: GPT-6 Astra (OpenAI), ChatGPT Pro mode, effort Pro, S6\_10\_PAPERS, 18 November 2026: targeted proof check. Continuation metadata events: GPT-6 Astra (OpenAI), ChatGPT Pro mode, effort Pro, S6\_10\_PAPERS continuation, 19 November 2026: used and rechecked in this run; status-vocabulary reconciliation; metadata reconciliation only; no new mathematical check. \emph{Pre-continuation recorded proof scope.} GPT-6 Astra (OpenAI, ChatGPT Pro mode), effort Pro, 2026-11-17 UTC; Claude Fable 5.2 (Anthropic), effort high, 16 November 2026; also checked in this writing by GPT-6 Astra (OpenAI), ChatGPT Pro mode, effort Pro, 18 November 2026}{the hypotheses, notation and definitions stated here; \cref{maps:src:D.2.30}; \cref{maps:src:D.2.31}; \cref{maps:src:D.2.23}; \cref{maps:src:D.2.20}; \cref{maps:eq:D.OC6}}


Every original additive cap on the $172$ natural candidate
rows can now be included. At $I_n^*$ we retain the original
cap; at every finite good-reduction additive site both the
original cap and every free logarithmic cap are allowed.
Keep the preceding marked Mumford models at multiplicative
sites, with the zero-degree original alternatives already
proved. These choices give exactly $282$ additive-filling
patterns on all $172$ initial rows. Of them, $54$ involve an
original additive cap with nonzero component Chern class.
The complete table is
\path{records/manifolds_original_chern.json}.

At a nonzero-Chern original cap let $s_i$ be the number of
Kodaira components, and let $q_i^c$ be its component-group
order. Its full boundary kernel adds the fibre relation
\[
 h_i s=0,\qquad h_i=b/q_i^c,                               \tag{D.OC1}\label{maps:eq:D.OC1}
\]
in addition to its two full fixed-loop normal relations.
Here $s$ is the common exterior fibre generator.
The local cokernel is free of rank $s_i-2$.

Let $L$ be the number of such nonzero-Chern original caps,
and put
\[
 \beta=
 \sum_{\text{zero-Chern originals}}(\rho_i-1)
 +\sum_{\text{nonzero-Chern originals}}(s_i-2).
                                                               \tag{D.OC2}\label{maps:eq:D.OC2}
\]
Take $T$ from the exact Smith matrix of \ref{maps:src:D.2.30}, adding
\eqref{maps:eq:D.OC1} for every nonzero-Chern original and setting both
fixed sweeps to zero at every unchanged additive cap.
Take $\mathbb Z^r\oplus A_1$ from the same complete Seifert
matrix, with the actual finite and regular numerators.
When $L=0$ the answer is \ref{maps:src:D.2.31}. When $L>0$ it is
\[
\begin{array}{c|c}
j&H_j(X;\mathbb Z)\\\hline
0,6&\mathbb Z\\
1&\mathbb Z^r\oplus A_1\\
2&\mathbb Z^{\beta+L+r}\oplus T\\
3&\mathbb Z^{2\beta+2L}\oplus T\\
4&\mathbb Z^{\beta+L+r}\oplus A_1\\
5&\mathbb Z^r.
\end{array}                                                   \tag{D.OC3}\label{maps:eq:D.OC3}
\]
All positive logarithmic orders and all free affine twists
are included, with arbitrary regular logarithms. The full
fundamental group is still \ref{maps:src:D.2.23}. In particular the extra
degree-one Mayer--Vietoris kernel has been retained, including
when the rational Seifert Euler number is zero.

\end{theorem}


\begin{proof}
\begin{enumerate}[label=\arabic*.,leftmargin=*]
\item \emph{The marked principal two-torus bundle.}
At an original additive site in the natural catalogue, either
$P$ or $Q$ is narrow. If $P$ is narrow, \ref{maps:src:D.2.31} applies.
Otherwise $Q$ is narrow, and the component Chern vector of
$M=\mathcal O(P-O)$ is
\[
 c=e_{\overline P}-e_0.
\]
It is primitive, and its pairing with the Kodaira
multiplicity vector is zero. There are exactly $46$ marked
occurrences of this case. The local Neron logarithm of $Q$
contracts its translation to the identity. Its lift to the
line bundle is straightened by the rigidified unit isotopy,
while the line bundle with its nonzero Chern vector is
retained. This is the actual relative construction of
D.2.20: the original cap retracts to the principal $T^2$
bundle over the Kodaira tree with Chern pair $(c,0)$, in
circle directions $\delta,\gamma$.

\item For clarity, the isotopy does not require the line bundle
to be trivial on the tree. Its Chern vector is fixed
throughout. A base isotopy lifts to this circle bundle,
and the difference between two lifts is a unit-valued
function. Since the tree has $H^1=0$, the unit has a
logarithm; rigidification fixes its constant phase and its
normal-meridian winding. Thus the transported boundary
marking is the one used in the following calculation.

\item The two nonzero transgressions in the integral bundle
spectral sequence are
\[
 x\longmapsto c(x)\delta,\qquad
 x\otimes y\longmapsto c(x)\delta\wedge y.
\]
Their images are primitive. As in \ref{maps:src:D.2.20}, all surviving
groups are free and
\[
 \pi_1(U)=\mathbb Z\gamma,\quad
 H_1(U)=\mathbb Z,\quad H_2(U)=\mathbb Z^{s_i-1},\quad
 H_3(U)=\mathbb Z^{2s_i-1},\quad H_4(U)=\mathbb Z^{s_i}.
                                                               \tag{D.OC4}\label{maps:eq:D.OC4}
\]
The projection $H_2(U)\to H_2(K)$ identifies it with
$\ker c$.

\item \emph{The entire fibre and normal kernel.}
Let $B_i$ be the upper-left elliptic block of $A_i$, and
let $Q_i$ be its elliptic gamma column. Narrowness of $Q$
gives the integral vector $z_i$ satisfying
\[
 Q_i=(B_i-I)z_i.
\]
The logarithm comparison identifies the elliptic projection
of a regular fibre by
$e_1\mapsto e_1$, $e_2\mapsto e_2$, $\delta\mapsto0$ and
$\gamma\mapsto z_i$. Its area covector is therefore
\[
 w_i=(1,0,(z_i)_2,0,-(z_i)_1,0).                            \tag{D.OC5}\label{maps:eq:D.OC5}
\]
The regular elliptic fibre specializes to the primitive
component-multiplicity vector in $\ker c$. Hence the full
map $H_2(T^4)\to H_2(U)$ is $w_i$ times that primitive
vector. Its kernel is exactly the integral $\ker w_i$.
Primitivity of the multiplicity vector inside $\ker c$
is checked directly, and also follows from its primitivity
in the ambient free component group.

\item The two principal circle actions sweep the origin disc.
Their normal vectors are
\[
 \delta,\qquad \ell_i=(-z_i,0,1).
\]
They are an integral basis of $\operatorname{Fix}(A_i)$.
Their actual fibre-area corrections are zero by the
rigidified logarithm comparison. Every relevant local
$A_i-I$ and $\bigwedge^2 A_i-I$ has only unit nonzero Smith
factors, so the boundary groups have the same free ranks
as in \ref{maps:src:D.2.20}. The two normal tori split the Wang projection
integrally. The remaining boundary image is the primitive
regular-fibre image above, and therefore
\[
 H_2(U)/\operatorname{im}H_2(\partial U)
                        =\mathbb Z^{s_i-2}.                \tag{D.OC6}\label{maps:eq:D.OC6}
\]
The full boundary kernel consists of both normal tori and
$\ker w_i$; its fibre part cannot be omitted.

\item Now evaluate $\ker w_i$ by the primitive global functional
$\omega=a e_1^*e_2^*-b\delta^*\gamma^*$.
An integral basis of the kernel has values
\[
 0,-a(z_i)_2,0,a(z_i)_1,-b.
\]
Thus its exact image in the exterior fibre group is
\[
 h_i\mathbb Zs,\qquad
 h_i=\gcd\bigl(b,a(z_i)_1,a(z_i)_2\bigr).
                                                               \tag{D.OC7}\label{maps:eq:D.OC7}
\]
All natural pairs in this catalogue have $\gcd(a,b)=1$,
so this is $\gcd(b,(z_i)_1,(z_i)_2)$.
The actual $46$ integral $z_i$ vectors give
$h_i=b/q_i^c$ in every case. The component orders are
$2$ for $III,III^*$, $3$ for $IV,IV^*$ and $4$ for
$I_1^*,I_3^*$. These are the finite constant lattice
calculations in the certificate; no assertion about
unbounded twists is obtained by enumerating them.
This proves the exact fibre relation \eqref{maps:eq:D.OC1}.

\item \emph{The global bridge and integral duality.}
All local cokernels in \eqref{maps:eq:D.OC6} and \ref{maps:src:D.2.31} are free,
so split them integrally. Their sum contributes
$\mathbb Z^\beta$ to the degree-two cokernel.
The remaining part is exactly the exterior fixed-sweep
matrix with the full finite/regular columns and the fibre
relations \eqref{maps:eq:D.OC1}. If $L>0$, at least one nonzero integer
multiple of $s$ is killed. Over $\mathbb Q$, all fixed
sweeps then vanish as well. This matrix consequently has
rational cokernel rank zero, even when the Seifert Euler
number is zero. Its integral torsion is precisely $T$.

\item Each nonzero-Chern original cap has first Betti number
one instead of the two of its boundary fixed lattice.
The degree-one Mayer--Vietoris source therefore has free
rank $L$ more than its target. Its cokernel is
$H_1(X)=\mathbb Z^r\oplus A_1$, so its kernel is free of
rank $L+r$. It is a subgroup of the free boundary group;
finite torsion in a logarithmic cap does not introduce a
hidden torsion kernel. The extension of this free bridge
by the degree-two cokernel splits as an abelian group.
It follows that
$H_2=\mathbb Z^{\beta+L+r}\oplus T$.
Every original cap has Euler number zero, so $\chi(X)=2$.
Duality gives the other entries of \eqref{maps:eq:D.OC3}, in particular
$b_3=2b_2-2b_1=2\beta+2L$.
This proves the complete formula, including the separate
role of the bridge when $r=1$.

\item Finally every finite additive site permits its two choices,
original or arbitrary free finite log, whereas $I_n^*$ is
kept original in this sector. Enumerating these choices
on all $172$ rows gives $282$ patterns. The $228$ patterns
with no nonzero-Chern original are exactly the preceding
zero-degree sector; the additional $54$ are those to which
\eqref{maps:eq:D.OC3} applies. This exhausts the stated additive choices.
\end{enumerate}
\end{proof}

For the untwisted original Engel input $P=G,Q=6G$, the
$IV^*$ and $III$ original caps have $s_i=7,2$ and
$h_i=2,3$. Hence $\beta=5$, $L=2$, $T=0$ and $r=1$.
Formula \eqref{maps:eq:D.OC3} gives the integral Betti numbers
$(1,1,8,14,8,1,1)$, independently recovering \cref{push:src:A.4.2} and
\ref{maps:src:D.2.20}. The one-original/one-finite cases similarly recover
the five all-order formulas of \ref{maps:src:D.2.21}. The theorem extends
these checks to every ordered section pair in the natural
candidate catalogue and every allowed additive choice.

The exact program \path{code/III/manifolds_original_chern.py}
retains all $46$ primitive Chern vectors, actual logarithm
vectors, full area kernels and primitive fibre images,
and the full $282$-pattern table. It also compares the
unreduced graph/normal/fibre-kernel Smith matrices with
the reduced presentation for every new pattern: $74$ full
Smith comparisons, retaining $9$ torsion-cap instances and
$54$ zero-rational-Euler instances. Original multiplicative caps with nonzero
component Chern class, logarithmic Mumford star caps and
initial weights outside this finite catalogue remain
separate parts of the broader request.


\begin{theorem}[all original multiplicative choices on the natural catalogue]\label{maps:src:D.2.33}
\StatementRecord{maps:src:D.2.33}{proved}{Historical events: Claude Fable 5.2 (Anthropic), effort high, 16 November 2026: carried as background without fault found. GPT-6 Astra (OpenAI), ChatGPT Pro mode, effort Pro, S6\_09\_PAPERS, 17 November 2026: targeted proof check. Also checked in this writing. Pre-continuation writing events: GPT-6 Astra (OpenAI), ChatGPT Pro mode, effort Pro, S6\_10\_PAPERS, 18 November 2026: targeted proof check. Continuation metadata events: GPT-6 Astra (OpenAI), ChatGPT Pro mode, effort Pro, S6\_10\_PAPERS continuation, 19 November 2026: used and rechecked in this run; status-vocabulary reconciliation; metadata reconciliation only; no new mathematical check. \emph{Pre-continuation recorded proof scope.} GPT-6 Astra (OpenAI, ChatGPT Pro mode), effort Pro, 2026-11-17 UTC; Claude Fable 5.2 (Anthropic), effort high, 16 November 2026; also checked in this writing by GPT-6 Astra (OpenAI), ChatGPT Pro mode, effort Pro, 18 November 2026}{the hypotheses, notation and definitions stated here; \cref{maps:src:D.2.32}; \cref{maps:src:D.2.23}; \cref{maps:eq:D.OC1}; \cref{maps:src:D.2.31}; \cref{maps:eq:D.OM5}; \cref{maps:eq:D.OM8}; \cref{maps:eq:D.OM7}; \cref{maps:eq:D.OM6}; \cref{maps:eq:D.OM10}; \cref{maps:src:D.2.30}; \cref{maps:eq:D.OC3}; \cref{maps:eq:D.OM3}}


At every zero-free multiplicative site in the $172$ natural
initial rows, either the original cap or the marked primitive
Mumford cap may be used. Together with all the original and
finite logarithmic additive choices of \ref{maps:src:D.2.32}, retaining
original caps at $I_n^*$, this gives $495$ distinct integral homology-presentation
patterns, after identifying the boundary-kernel agreements
for zero-degree original/Mumford choices proved below.
Every original/Mumford choice represented by such an
agreement is also listed explicitly in
\path{records/manifolds_original_multiplicative.json}.
Every finite and regular logarithmic order and every free
affine twist remains allowed.

The new original caps occur at $120$ marked zero-free
$I_n$ sites with nonzero component Chern vector. Each adds
\[
 h_i s=0,\qquad h_i=b/n,                                   \tag{D.OM1}\label{maps:eq:D.OM1}
\]
to the exact global matrix, has free local degree-two
cokernel of rank $n-2$, and adds one to the free
degree-one Mayer--Vietoris bridge. Consequently \ref{maps:src:D.2.32}
continues to hold verbatim if $L$ counts \emph{all}
nonzero-Chern original caps, additive or multiplicative,
and $\beta$ includes $n-2$ for each new multiplicative
original. In particular, when $L>0$,
\[
\begin{array}{c|c}
j&H_j(X;\mathbb Z)\\\hline
0,6&\mathbb Z\\
1&\mathbb Z^r\oplus A_1\\
2&\mathbb Z^{\beta+L+r}\oplus T\\
3&\mathbb Z^{2\beta+2L}\oplus T\\
4&\mathbb Z^{\beta+L+r}\oplus A_1\\
5&\mathbb Z^r.
\end{array}                                                \tag{D.OM2}\label{maps:eq:D.OM2}
\]
Here $\mathbb Z^r\oplus A_1$ is the exact Seifert
abelianization of \ref{maps:src:D.2.23}, and $T$ is the torsion of the
full finite/regular normal matrix with all the additive
relations \eqref{maps:eq:D.OC1} and multiplicative relations \eqref{maps:eq:D.OM1}.
The full fundamental group is still the Seifert group,
including when its rational Euler number is zero.

\end{theorem}


\begin{proof}
\begin{enumerate}[label=\arabic*.,leftmargin=*]
\item \emph{The actual bundle over the nodal cycle.}
At a zero-free $I_n$ site in these rows, either $P$ is
narrow or $Q$ is narrow. If $P$ is narrow, $Q$ acts
transitively on the components. Its original cap has
one component orbit and the same full integral normal
kernel and primitive fibre image as the marked Mumford
cap, by the multiplicative calculation in \ref{maps:src:D.2.31}.
This proves equality of the recorded integral homology
presentations for these zero-degree choices. It does not
assert a relative diffeomorphism between their caps.

\item Otherwise $n>1$, $Q$ is narrow and the component of $P$
is a unit modulo $n$. Write $K=I_n$ and
\[
 c=e_{\overline P}-e_0\in H^2(K;\mathbb Z)=\mathbb Z^n .
\]
This vector is primitive and its sum is zero. The local
Neron logarithm of $Q$ contracts its translation to the
identity. The lift preserves the line bundle with Chern
vector $c$; this bundle is not replaced by a trivial one.
There is one additional point compared with the tree
argument of D.2.32: the dual graph of $I_n$ has a circle.
On a cycle, translation in the identity component acts
on the generalized Jacobian of a multidegree-$c$ line
by the power given by $\sum_i c_i$. That sum is zero,
so the translation isotopy admits line-bundle lifts.
One can see this directly on the normalization: multiply
the coordinate on every rational component by the same
unit and use its integral powers in the line transition
functions; the final clutch factor is the unit to
$\sum_i c_i=0$.

\item The ratio between the resulting lift and the prescribed
holomorphic lift on $K$ is a global holomorphic unit on
the compact connected nodal cycle, and hence is constant.
It has no winding on the dual circle. Rigidification
removes this constant phase, and the nonvanishing
linearization on the base disc has a logarithm, fixing
the normal-meridian marking as well. Thus no additional
mixed Chern class is introduced. The actual cap retracts,
with its transported boundary marking, to
\[
 U\simeq Y_c\times S^1_\gamma,\qquad
 S^1_\delta\longrightarrow Y_c\longrightarrow K,           \tag{D.OM3}\label{maps:eq:D.OM3}
\]
where the circle bundle has Chern class $c$.

\item The integral Gysin sequence, or the bundle spectral
sequence over $K\simeq S^1\vee\bigvee_{i=1}^nS^2$, gives
\[
\begin{aligned}
 H_1(Y_c)&=\mathbb Z,&
 H_2(Y_c)&=\mathbb Z^n,&
 H_3(Y_c)&=\mathbb Z^n,\\
 H_1(U)&=\mathbb Z^2,&
 H_2(U)&=\mathbb Z^{n+1},&
 H_3(U)&=\mathbb Z^{2n},&
 H_4(U)&=\mathbb Z^n .
\end{aligned}                                             \tag{D.OM4}\label{maps:eq:D.OM4}
\]
All groups are free: the only transgression is the
primitive vector $c$. More precisely,
\[
 0\longrightarrow\mathbb Z[\delta\times\text{dual circle}]
 \longrightarrow H_2(Y_c)
 \longrightarrow\ker(c:\mathbb Z^n\to\mathbb Z)
 \longrightarrow0                                        \tag{D.OM5}\label{maps:eq:D.OM5}
\]
is an exact sequence of free groups. This also identifies
its first generator primitively.

\item \emph{The full marked kernel, including its normal correction.}
Let $u_1$ be the primitive elliptic vanishing circle and
choose $u_2$ with oriented determinant one. As in \ref{maps:src:D.2.32},
narrowness gives an integral vector $z$ satisfying
$(B-I)z=Q_{\rm col}$, and put $\ell=(-z,0,1)$.
In the integral basis $(u_1,u_2,\delta,\ell)$ the actual
monodromy has the form
\[
 A u_2=u_2+\nu u_1+p\delta,\qquad
 A u_1=u_1,\quad A\delta=\delta,\quad A\ell=\ell,
 \qquad |\nu|=n,\quad\gcd(\nu,p)=1.                        \tag{D.OM6}\label{maps:eq:D.OM6}
\]
These are equalities in the marked lattice, not only
a rational normal form. The certificate records the
unimodular basis for each of the $120$ occurrences.
In particular the full fixed lattice is
$\mathbb Zu_1\oplus\mathbb Z\delta\oplus\mathbb Z\ell$,
and both $A-I$ and $\bigwedge^2 A-I$ have only unit
nonzero Smith factors.

\item The regular fibre has primitive rank-three image in
$H_2(U)$. Its generators can be taken as
\[
 u_1\wedge u_2,\qquad u_2\wedge\delta,\qquad
 u_2\wedge\ell.                                           \tag{D.OM7}\label{maps:eq:D.OM7}
\]
The first projects to the all-one component vector in
$\ker c$, which is primitive; a possible addition of
the first term of \eqref{maps:eq:D.OM5} does not affect this assertion.
The second generates that first term primitively, and
the third is the dual-circle class times $\gamma$.
Hence their span is primitive. The full regular-fibre
kernel is exactly
\[
 K_{\rm fib}=
 \mathbb Z(u_1\wedge\delta)\oplus
 \mathbb Z(u_1\wedge\ell)\oplus
 \mathbb Z(\delta\wedge\ell).                              \tag{D.OM8}\label{maps:eq:D.OM8}
\]
For the first two generators, the vanishing disc and
its lifted products provide the required boundaries.
For the third, primitivity of $c$ kills the circle
$\delta$ integrally in $Y_c$. The image of
$\bigwedge^2 A-I$ is a primitive rank-two sublattice
of \eqref{maps:eq:D.OM8}.

\item The normal tori in the $\delta$ and $\ell$ directions
bound the sweeps of the origin disc under the two
principal circle actions. Their marked fibre corrections
are zero. A normal torus in the $u_1$ direction requires
slightly more care. Its projection to the original
elliptic neighbourhood is the usual vanishing-loop
normal torus, which bounds by the local vanishing-disc
construction. Its image in $H_2(U)$ therefore lies in
the first summand of \eqref{maps:eq:D.OM5}; its $\gamma$ coordinate
is zero. It is an integral multiple of the regular
torus $u_2\wedge\delta$. Subtracting that multiple
produces an actual boundary-kernel cycle with normal
projection $u_1$.

\item This correction cannot be discarded at the local level.
However the global primitive alternating functional
satisfies exactly
\[
 \omega(u_2,\delta)=0,\qquad
 \omega=a e_1^*e_2^*-b\delta^*\gamma^*.                    \tag{D.OM9}\label{maps:eq:D.OM9}
\]
Both $u_2$ and $\delta$ have zero $\gamma$ coordinate,
and $\delta$ has zero elliptic projection. Thus the
correction has zero image in the exterior fibre group.
The actual global image of the three normal kernels is
therefore the full marked fixed-sweep lattice with
zero fibre coefficient. This proves the needed marked
statement without guessing the local correction integer.

\item The boundary image is exactly the primitive rank-three
regular-fibre image \eqref{maps:eq:D.OM7}, and consequently
\[
 H_2(U)/\operatorname{im}H_2(\partial U)
                          =\mathbb Z^{n-2}.               \tag{D.OM10}\label{maps:eq:D.OM10}
\]
Evaluate \eqref{maps:eq:D.OM8} by $\omega$. Monodromy invariance gives
$\nu\,\omega(u_1,\ell)+p\,\omega(\delta,\ell)=0$,
where $\omega(\delta,\ell)=-b$ and
$\omega(u_1,\delta)=0$. Since $\gcd(\nu,p)=1$,
$n\mid b$ and the exact common divisor of these values
is $b/n$. This proves \eqref{maps:eq:D.OM1}, including its integral
coefficient.

\item \emph{Global gluing and exhaustion of the stated choices.}
The cap has first Betti number two, whereas the fixed
lattice in \eqref{maps:eq:D.OM6} has rank three. Thus it adds precisely
one free bridge, just as a nonzero-Chern additive original
does. The local degree-two cokernel \eqref{maps:eq:D.OM10} is free,
so it splits from the remaining global presentation.
The preceding kernel calculation shows that the latter
is exactly the old fixed-sweep matrix with the new
relation \eqref{maps:eq:D.OM1}. All finite torsion caps and their
nontrivial normal Chern classes are still handled by
the full three-column Thom matrices of \ref{maps:src:D.2.30}.

\item The proof of \eqref{maps:eq:D.OC3} now applies with the enlarged $L$
and $\beta$. In particular, a positive $L$ kills a
nonzero multiple of $s$, so the reduced degree-two
matrix has rational cokernel zero, even when $r=1$.
The degree-one bridge has rank $L+r$. Every original
cap in \eqref{maps:eq:D.OM3} has Euler number zero, and the unique
index-one toric zero still contributes two. Integral
duality therefore gives \eqref{maps:eq:D.OM2}.
At the group level the new original cap only adds
vertical relations, including the vanishing direction
and $\delta$, and adds no relation on $\gamma$.
The complete vertical lattice was already killed in
\ref{maps:src:D.2.23}, so its Seifert presentation is unchanged.

\item Finally, choose original or marked Mumford independently
at each of the nonzero-Chern multiplicative sites in
the $282$ additive patterns. This yields $495$ integral homology-presentation patterns
on all $172$ initial rows, with all the zero-degree
original/Mumford alternatives separately recorded. It
exhausts exactly the choices stated in the theorem.
\end{enumerate}
\end{proof}

The exact program \path{code/III/manifolds_original_multiplicative.py}
verifies all $120$ integral normal forms, primitive fibre
and monodromy kernels, actual values of $\omega$, and
primitive Chern and component-multiplicity vectors.
Their types are $I_8$ in $6$ occurrences, $I_7$ in $14$,
$I_2$ in $31$, $I_3$ in $23$, $I_5$ in $38$ and $I_4$
in $8$. It compares the full unreduced graph/normal/fibre
matrices with the stated reduced presentation for every
new pattern: $213$ full Smith comparisons, retaining
$18$ torsion-cap and $172$ zero-rational-Euler instances.
These are independent checks of an all-order proof.
The logarithmic Mumford choices at star fibres and
initial linearization weights outside the finite natural
catalogue remain separate parts of the broader request.


\begin{theorem}[every positive weight on the original Engel profile]\label{maps:src:D.2.34}
\StatementRecord{maps:src:D.2.34}{proved}{Historical events: Claude Fable 5.2 (Anthropic), effort high, 16 November 2026: carried as background without fault found. GPT-6 Astra (OpenAI), ChatGPT Pro mode, effort Pro, S6\_09\_PAPERS, 17 November 2026: targeted proof check. Also checked in this writing. Pre-continuation writing events: GPT-6 Astra (OpenAI), ChatGPT Pro mode, effort Pro, S6\_10\_PAPERS, 18 November 2026: targeted proof check. Continuation metadata events: GPT-6 Astra (OpenAI), ChatGPT Pro mode, effort Pro, S6\_10\_PAPERS continuation, 19 November 2026: used and rechecked in this run; status-vocabulary reconciliation; metadata reconciliation only; no new mathematical check. \emph{Pre-continuation recorded proof scope.} GPT-6 Astra (OpenAI, ChatGPT Pro mode), effort Pro, 2026-11-17 UTC; Claude Fable 5.2 (Anthropic), effort high, 16 November 2026; also checked in this writing by GPT-6 Astra (OpenAI), ChatGPT Pro mode, effort Pro, 18 November 2026}{the hypotheses, notation and definitions stated here; \cite[Lemma 2.1 and Remark 8.5]{manEngel26}; \cref{maps:src:D.2.5}; \cref{maps:src:D.2.29}; \cref{maps:src:D.1.2}; \cref{maps:eq:D.WE11}; \cref{maps:src:D.2.32}; \cref{maps:eq:D.WE9}; \cref{maps:eq:D.WE12}; \cref{maps:eq:D.WE13}}


On the $IV^*,III,I_1$ rational elliptic surface, let $G$
generate the Mordell--Weil group, retain
$M_S=\mathcal O(G-O)$, and use translation $6dG$ with
linearization zero $dI_1$, for any integer $d>0$.
Use the smooth $d$-period hexagonal Mumford cap at $I_1$.
At either additive site one may retain the original cap
or use any free finite logarithmic cap described below.
This gives an unbounded initial-weight family outside
the finite sphere-candidate catalogue.

For logarithmic caps put
\[
 m_0=3k,\quad m_1=4\ell,\qquad
 v_i=r_i\delta+n_i u_i,\qquad k,\ell>0.
\]
The exact freeness conditions are
\[
\begin{aligned}
 t_0&=\gcd(m_0,n_0,3r_0),&
 &\gcd(m_0,n_0,r_0)=1,\quad 3\mid m_0/t_0,\\
 t_1&=\gcd(m_1,n_1,2r_1),&
 &\gcd(m_1,n_1,r_1)=1,\quad 4\mid m_1/t_1 .
\end{aligned}                                             \tag{D.WE1}\label{maps:eq:D.WE1}
\]
There is no coprimality restriction on $d$.
All noninvariant parts of affine translations are
included through the marked translation conjugacy.
At an original site set its effective Seifert data to
$\overline m_i=1,\overline n_i=0$; otherwise set them
to $m_i,n_i$. Define
\[
 P=\overline m_1\overline n_0+
   \overline m_0\overline n_1,\qquad
 \mathbb Z^r\oplus A_1=
 \operatorname{coker}
 \begin{pmatrix}
  -\overline n_0&\overline m_0\\
  -\overline n_1&-\overline m_1
 \end{pmatrix}.                                           \tag{D.WE2}\label{maps:eq:D.WE2}
\]
The full fundamental group is this abelian group, and
$r=1$ precisely when $P=0$.

The exact degree-two matrix has target $(s,t,t_0',t_1')$.
For two logarithmic caps it is $B=(B_0\ B_1)$, where
\[
 B_0=
 \begin{pmatrix}
 -2dn_0&-6dr_0&0\\
 -k&0&r_0\\
 0&m_0&n_0\\
 0&0&0
 \end{pmatrix},
 \qquad
 B_1=
 \begin{pmatrix}
 -3dn_1&-6dr_1&0\\
 2\ell&0&-r_1\\
 0&0&0\\
 0&m_1&n_1
 \end{pmatrix}.                                           \tag{D.WE3}\label{maps:eq:D.WE3}
\]
For an original $IV^*$ or $III$ cap replace its three
columns, respectively, by
\[
 O_0=
 \begin{pmatrix}2d&0&0\\0&1&0\\0&0&1\\0&0&0\end{pmatrix},
 \qquad
 O_1=
 \begin{pmatrix}3d&0&0\\0&-1&0\\0&0&0\\0&0&1\end{pmatrix}.
                                                               \tag{D.WE4}\label{maps:eq:D.WE4}
\]
Let $T$ be the torsion of $\operatorname{coker}B$,
let $L$ count original additive caps, and put
$\beta=5$ if $IV^*$ is original, and $\beta=0$ otherwise.
Set
\[
 (b_2,b_3)=
 \begin{cases}
 (d-1+2r,\;2r),&L=0,\\
 (d-1+\beta+L+r,\;2\beta+2L),&L>0.
 \end{cases}                                               \tag{D.WE5}\label{maps:eq:D.WE5}
\]
Then all the integral homology is
\[
\begin{array}{c|c}
j&H_j(X;\mathbb Z)\\\hline
0,6&\mathbb Z\\
1&\mathbb Z^r\oplus A_1\\
2&\mathbb Z^{b_2}\oplus T\\
3&\mathbb Z^{b_3}\oplus T\\
4&\mathbb Z^{b_2}\oplus A_1\\
5&\mathbb Z^r .
\end{array}                                                \tag{D.WE6}\label{maps:eq:D.WE6}
\]
In particular $\chi(X)=2d$ for every choice.

There are useful closed forms. If both sites are
logarithmic and $\gcd(m_i,n_i)=1$, put
$g=\gcd(d,k,2\ell)$. Then
\[
 T=
 \begin{cases}
 \mathbb Z/g\oplus\mathbb Z/(d|P|/g),&P\ne0,\\
 \mathbb Z/g,&P=0.
 \end{cases}                                               \tag{D.WE7}\label{maps:eq:D.WE7}
\]
The fundamental group is $\mathbb Z/|P|$ for $P\ne0$
and $\mathbb Z$ for $P=0$. With at least one original
cap the torsion groups, for every free remaining twist,
are
\[
\begin{array}{c|c}
\text{original sites}&T\\\hline
IV^*\text{ only}&
 \mathbb Z/(d\gcd(2,n_1))\oplus\mathbb Z/\gcd(m_1,n_1)\\
III\text{ only}&
 \mathbb Z/(d\gcd(3,n_0))\oplus\mathbb Z/\gcd(m_0,n_0)\\
\text{both}&\mathbb Z/d .
\end{array}                                                \tag{D.WE8}\label{maps:eq:D.WE8}
\]
These are direct-sum descriptions, not necessarily
ordered invariant factors.

\end{theorem}


\begin{proof}
\begin{enumerate}[label=\arabic*.,leftmargin=*]
\item \emph{The weighted analytic construction.}
Engel's linearization identity for $M_S$ and $6G$
has base degree one
\cite[Lemma 2.1 and Remark 8.5]{manEngel26}.
Iterate that contracting generator $d$ times.
Translations preserve the base, so its linearization
is the product of $d$ copies and has divisor $dI_1$.
Equivalently, replace the long gamma period by $d$
times itself. On the punctured family this is the
ordinary torus cover defined by that sublattice.
The good-reduction models are preserved by it:
the integral matrices below have orders three and
four for every $d$. After the respective root base
changes one can therefore make the same finite
affine quotients as in \ref{maps:src:D.2.5} and \ref{maps:src:D.2.29}.

\item At the cusp retain the same unimodular triangular
fan, now modulo the sublattice generated by one
short-direction period and $d$ times the other.
Its quotient tube is smooth and proper by the same
finite-period toric construction as \ref{maps:src:D.1.2}. Its
normalization has $d$ copies of $dP_6$ and
\[
 H_*(N;\mathbb Z)
 =(\mathbb Z,\mathbb Z^2,\mathbb Z^{d+3},
                         \mathbb Z^2,\mathbb Z^d).         \tag{D.WE9}\label{maps:eq:D.WE9}
\]
The normalization proof of \ref{maps:src:D.1.2} uses
$\operatorname{diag}(1,t-1,t-1)$ over
$\mathbb Z[t]/(t^d-1)$ and is valid for every positive
$d$, including $2$ and $3$.
This constructs the weighted family directly.
It does not assume that the filled threefold is a
free cover of a weight-one finite quotient when $d$
shares a factor with its multiplicities.

\item \emph{Integral exterior and finite-cap data.}
Use the basis $(\delta,a,b,\gamma)$ of \ref{maps:src:D.2.5}, with
the new gamma period. The actual monodromies are
\[
\begin{aligned}
 A_0&=\begin{pmatrix}
 1&-1&0&2d\\0&-1&-1&6d\\0&1&0&0\\0&0&0&1
 \end{pmatrix},\\
 A_1&=\begin{pmatrix}
 1&0&1&-3d\\0&0&1&0\\0&-1&0&6d\\0&0&0&1
 \end{pmatrix},&
 A_\infty&=\begin{pmatrix}
 1&0&0&d\\0&1&1&0\\0&0&1&0\\0&0&0&1
 \end{pmatrix}.
\end{aligned}                                             \tag{D.WE10}\label{maps:eq:D.WE10}
\]
Their product is the identity. The two fixed lattices
are generated by $\delta$ and
\[
 u_0=\gamma+2d(a+b),\qquad
 u_1=\gamma+3d(a+b).
\]
The finite coinvariants still have
$b=-a,\delta=-3a$ and $b=-a,\delta=-2a$,
respectively, and $u_i$ maps to the primitive gamma
class. Thus their fixed-quotient indices and root
characters are exactly those of \ref{maps:src:D.2.29}, independently
of $d$. Its proof gives \eqref{maps:eq:D.WE1}, including every
nonprimitive free numerator, the torsion orders
$t_i$ of the caps and their normal lines, and the
three actual Thom generators.

\item Let $V$ be the regular bundle over a graph joining
the three peripheral circles. The common fibre
coinvariant in degree two is free of rank one:
its primitive functional is
\[
 \xi_d=a^*\wedge b^*+6d\gamma^*\wedge\delta^* .
\]
The integral relations are
$\delta a=\delta b=a\gamma=b\gamma=0$ and
$\delta\gamma=-6d\,ab$, so $s=[ab]$ is primitive.
For the graph part the five columns
\[
 (-\delta,-\delta),\quad(b,-a),\quad
 (\delta,0),\quad(u_0,0),\quad(0,u_1)                     \tag{D.WE11}\label{maps:eq:D.WE11}
\]
are an integral basis of
$\ker[(A_0-I)\ (A_1-I)]$.
They lie in the kernel and have a constant unit
$5$-row minor, using rows $0,2,3,4,7$ in zero-based
indexing. Their rational rank is five, so that unit
minor proves saturation for every $d$.
The common exterior-square relations likewise contain
a unit minor in the displayed primitive kernel of
$\xi_d$; the certificate gives its column indices
as a polynomial identity in $d$.

\item The first two columns of \eqref{maps:eq:D.WE11} are the actual
cusp normal tori in the angular directions
$\delta,a$. Their fibre coefficients are zero in
the rigidified cusp marking. Quotienting by these
two columns leaves the free target
$(s,t,t_0',t_1')$ used above.
The norms and contractions are
\[
\begin{aligned}
 N_0a&=-k\delta,&N_0\gamma&=m_0u_0,&
 \xi_d(a,u_0)&=2d,\\
 N_1a&=-2\ell\delta,&N_1\gamma&=m_1u_1,&
 \xi_d(a,u_1)&=3d,\\
 &&\xi_d(\gamma,\delta)&=6d .
\end{aligned}
\]
A covering-torus loop $\lambda$ has normal column
with graph part $N_i\lambda$ and fibre part
$-\xi_d(\lambda,v_i)$; the affine meridian has
graph part $v_i$ and fibre part zero. These are the
full Thom generators, including the torsion
character case of \ref{maps:src:D.2.29}. Since the delta sweep
at the second site is the negative of the first
modulo \eqref{maps:eq:D.WE11}, they give exactly \eqref{maps:eq:D.WE3}.

\item For an original additive cap, $Q=6dG$ is narrow and
the component Chern vector of $M_S$ is the same
primitive $e_G-e_O$ as at $d=1$. \ref{maps:src:D.2.32}'s actual
principal two-torus model applies with the marked
logarithm multiplied by $d$. The regular elliptic
area covectors, in the order
$(\delta a,\delta b,\delta\gamma,ab,a\gamma,b\gamma)$,
are
\[
 (0,0,0,1,-2d,2d),\qquad
 (0,0,0,1,-3d,3d).
\]
Their full kernels have $\xi_d$ images $2d\mathbb Z$
and $3d\mathbb Z$, and both complete fixed normal
lattices have zero fibre correction. Thus \eqref{maps:eq:D.WE4}
is the exact integral replacement. The local free
cokernels have ranks $7-2=5$ and $2-2=0$, and each
original loses one first-homology rank, giving
one additional free bridge. These statements retain
the primitive Chern class rather than making a
zero-degree approximation.

\item \emph{The cusp boundary torsion and the component meridians.}
Write $E=X\setminus\operatorname{int}N$ and $M=\partial N$.
For $d>1$, $H_1(M)$ has an actual $\mathbb Z/d$ summand:
the monodromy kills $a$ and $d\delta$, whereas the
cusp kills both primitive angular circles. It would
therefore be incorrect to reuse a free-boundary
Mayer--Vietoris splitting from weight one.
Instead use the exact sequence of the pair $(X,E)$.

\item The exterior-square cusp image is primitive of rank
two for every $d$, so $H_2(M)$ is free of rank six.
Its two angular normal tori bound by their toric
section-disc sweeps. They form a primitive sublattice
because their Wang projections are the primitive
fixed basis $\delta,a$.
Furthermore
$H_3(N,M)=H^3(N)=\mathbb Z^2$ by \eqref{maps:eq:D.WE9}.
Exactness bounds the kernel of
$H_2(M)\to H_2(N)$ by rank two. Hence those two
angular tori are its entire integral kernel.

\item Attaching the two finite or original additive caps
to $V$ first shows
\[
 H_2(E)/j_*\ker(H_2(M)\to H_2(N))
 \simeq\operatorname{coker}B\oplus\mathbb Z^{\beta+L}.
                                                               \tag{D.WE12}\label{maps:eq:D.WE12}
\]
Indeed every finite boundary-to-cap degree-two map
is onto, including when the cap and normal line
have torsion; its kernel is the full Thom image.
At an original cap split its free local cokernel.
The degree-one sequence then has a free kernel of
rank $L$, since the first Betti number of $E$ is
one and each original lowers the target rank by one.
The angular cusp classes come from $H_2(V)$, so
they have zero image in this new bridge. Thus
quotienting by them preserves its free splitting.
This proves \eqref{maps:eq:D.WE12} without dropping boundary
torsion or assuming a torsion extension splits.

\item Now $H_2(N,M)=H^4(N)=\mathbb Z^d$, with its basis
given by small positive transverse discs to the
$d$ components. Their boundary classes in $H_1(E)$
are \emph{identical}, not merely proportional.
In the smooth height-one fan a component ray
$(v,1)$ has its primitive meridian $(v,1)$ in
angular-plus-base coordinates. Two such meridians
differ by integral multiples of $a,\delta$.
Both angular classes already vanish in $H_1(E)$:
the unchanged vertical columns of $A_0-I,A_1-I$
generate the entire vertical lattice primitively.
Thus every disc has boundary the same class
$x_0+x_1$, with the common boundary orientation;
at an original site its $x_i$ is zero.

\item The abelian presentation of $H_1(E)$ has generators
$\gamma,x_0,x_1$ and relations
$\overline m_i x_i=\overline n_i\gamma$.
It has free rank one. Its element $x_0+x_1$ has
infinite order precisely when $P\ne0$.
Therefore the kernel of the actual map
\[
 \mathbb Z^d\longrightarrow H_1(E),\qquad
 (z_1,\ldots,z_d)\longmapsto
                 (z_1+\cdots+z_d)(x_0+x_1)
                                                               \tag{D.WE13}\label{maps:eq:D.WE13}
\]
is free of rank $d-1+r$. This remains true when
$H_1(E)$ or the indicated meridian has torsion.
The pair sequence and \eqref{maps:eq:D.WE12} now give a short
exact sequence with this free quotient, hence a
split sequence of abelian groups.

\item The matrix \eqref{maps:eq:D.WE3} has rational corank $r$.
For example a left-kernel vector has its other
three coordinates fixed by its $s$ coordinate;
the remaining equation is $dP$ times that coordinate.
If an original is present, \eqref{maps:eq:D.WE4} kills that
coordinate and the entire rational cokernel, so
the corank is zero. These facts and \eqref{maps:eq:D.WE13} give
the claimed $b_2$ and torsion $T$.

\item Van Kampen, separately, proves the group assertion.
The cusp kills the primitive $a,\delta$ loops;
finite monodromy then kills $b$. Only the central
$\gamma$ and the two meridians remain, with
$x_0x_1=1$ and the two power relations in \eqref{maps:eq:D.WE2}.
Thus the full group is abelian with that exact
integer presentation. Finally the cusp has Euler
number $2d$ by \eqref{maps:eq:D.WE9}, and every other cap has
Euler number zero. Integral duality and universal
coefficients give \eqref{maps:eq:D.WE5}--\eqref{maps:eq:D.WE6}.

\item \emph{Closed Smith forms.}
When both numerators are primitive, the same
Bezout reductions as \ref{maps:src:D.2.5} eliminate two unit
columns of \eqref{maps:eq:D.WE3} and leave
\[
 \begin{pmatrix}
 -2dn_0&-3dn_1\\-k&2\ell
 \end{pmatrix},\qquad \det=-dP .
\]
Its entries have gcd $\gcd(d,k,2\ell)$.
To check this prime by prime, a prime dividing
$k$ cannot divide $n_0$, one dividing $\ell$
cannot divide $n_1$, $3\nmid n_0$, and $2\nmid n_1$.
These observations give precisely the stated gcd
at $2,3$ and at every other prime. Its Smith
form gives \eqref{maps:eq:D.WE7}, including its rank-one case.
For an original $IV^*$, eliminate $t,t_0'$ using
\eqref{maps:eq:D.WE4}. The remaining relations are
$2ds=3dn_1s=0$ and
$m_1t_1'=n_1t_1'=0$, because $6dr_1s$ is already
zero. This is the first row of \eqref{maps:eq:D.WE8}.
The original $III$ case is identical with
$3d$ and $2dn_0$, and with both originals the
remaining fibre relation is $\gcd(2d,3d)s=ds=0$.
\end{enumerate}
\end{proof}

For two logarithmic caps, a simply connected member
has $P=\pm1$, which forces both local numerators
primitive and $\gcd(3k,4\ell)=1$. Thus $g=1$ and
\[
 H_2=\mathbb Z^{d-1}\oplus\mathbb Z/d,\qquad
 H_3=\mathbb Z/d,\qquad H_4=\mathbb Z^{d-1}.                \tag{D.WE14}\label{maps:eq:D.WE14}
\]
This includes even and three-divisible $d$, where the
free-cover reduction in D.1 is unavailable. At minimal
orders and $d$ prime to six it recovers D.1's groups.
The exact smooth classification beyond the cases
already settled there is not inferred from \eqref{maps:eq:D.WE14}.
With both original caps, \eqref{maps:eq:D.WE6} instead gives
\[
 \pi_1=\mathbb Z,\qquad
 H_2=\mathbb Z^{d+7}\oplus\mathbb Z/d,\qquad
 H_3=\mathbb Z^{14}\oplus\mathbb Z/d,\qquad
 H_4=\mathbb Z^{d+7},\quad H_5=\mathbb Z.
\]
At $d=1$ this is exactly \cref{push:src:A.4.2}; for every $d$ it is the
actual iterated-translation family.

The executable \path{code/III/manifolds_weighted_engel.py}
accepts arbitrary integer input, checks \eqref{maps:eq:D.WE1}, and
returns the full group and homology presentation.
It verifies polynomial identities for the monodromies,
primitive exterior-square quotient and graph basis,
and compares the unreduced graph and full area kernels
with \eqref{maps:eq:D.WE3}--\eqref{maps:eq:D.WE4}. Its retained output includes
nonprimitive free twists, nontrivial cap torsion,
every original/log choice and zero-Euler cases.


\begin{theorem}[the fifteen minimal $c=2$ patterns]\label{maps:src:D.2.35}
\StatementRecord{maps:src:D.2.35}{proved}{Historical events: Claude Fable 5.2 (Anthropic), effort high, 16 November 2026: carried as background without fault found. GPT-6 Astra (OpenAI), ChatGPT Pro mode, effort Pro, S6\_09\_PAPERS, 17 November 2026: targeted proof check. Also checked in this writing. Pre-continuation writing events: GPT-6 Astra (OpenAI), ChatGPT Pro mode, effort Pro, S6\_10\_PAPERS, 18 November 2026: targeted proof check. Continuation metadata events: GPT-6 Astra (OpenAI), ChatGPT Pro mode, effort Pro, S6\_10\_PAPERS continuation, 19 November 2026: used and rechecked in this run; status-vocabulary reconciliation; metadata reconciliation only; no new mathematical check. \emph{Pre-continuation recorded proof scope.} GPT-6 Astra (OpenAI, ChatGPT Pro mode), effort Pro, 2026-11-17 UTC; Claude Fable 5.2 (Anthropic), effort high, 16 November 2026; also checked in this writing by GPT-6 Astra (OpenAI), ChatGPT Pro mode, effort Pro, 18 November 2026}{the hypotheses, notation and definitions stated here; \cref{maps:eq:D.HY3}; \cref{maps:src:D.2.8}; \cref{maps:src:D.2.26}; \cref{maps:eq:D.C24}; \cref{maps:eq:D.C25}; \cref{maps:eq:D.C27}; \cref{maps:eq:D.C26}; \cref{maps:src:D.2.14}; \cref{maps:eq:D.HY7}; \cref{maps:eq:D.C215}; \cref{maps:eq:D.C212}; \cref{maps:eq:D.C213}; \cite[Theorems 3, 2.10 and Section 3.12]{manZhubr75}}


Take one of the fifteen patterns with $c=2$ in
\path{records/root_remaining_smooth_profiles.json}, use the primitive
Mumford filling at every multiplicative site, and make no additional
regular logarithmic transformations. Retain the additive original/log
choice of that pattern and the ordinary clutching $e\gamma$ of
\eqref{maps:eq:D.HY3}. Every simply connected output, for every allowed finite
multiplicity and affine twist, is orientation-preservingly diffeomorphic
to the conic surgery $M_{\mathrm{conic}}$ of \ref{maps:src:D.2.8}. Equivalently, if
$\alpha\ne0$ in $H^2(X;\mathbb F_2)$, then
\[
 H_2(X;\mathbb Z)=H_3(X;\mathbb Z)=\mathbb Z/2,
 \qquad \langle\alpha^3,[X]_2\rangle=1.                 \tag{D.C21}\label{maps:eq:D.C21}
\]
The marked patterns, with zero positions numbered from zero as in the
source tables, are
\[
\begin{array}{c|r|r|l|l}
\text{OS/profile}&a&b&\text{zero positions}&\text{additive choice}\\\hline
47&7&2&2,3,4&\text{none}\\
47/II&7&2&2&II\text{ original or log}\\
49&1&6&2,3&IV^*\text{ log}\\
56&5&2&1&\text{none}\\
56&15&2&3,4&\text{none}\\
56/II&5&2&3&II\text{ original or log}\\
56/IV&1&6&0&IV\text{ log}\\
56/IV&5&6&3&IV\text{ log}\\
56/IV&15&2&3&IV\text{ original}.
\end{array}                                                   \tag{D.C22}\label{maps:eq:D.C22}
\]
Thus these fifteen patterns occupy thirteen distinct marked initial
rows. The statement includes the original additive caps in this table;
it does not replace the stipulated multiplicative Mumford caps by
other caps, or include additional regular logarithmic transformations.

\end{theorem}


\begin{proof}
\begin{enumerate}[label=\arabic*.,leftmargin=*]
\item \emph{1. The common piece and its integral markings.}
Throughout, use the actual lattice basis $(u,w,\delta,\gamma)$ and
write
\[
 \Omega=a\,u^*\wedge w^*+b\,\gamma^*\wedge\delta^*.
                                                               \tag{D.C23}\label{maps:eq:D.C23}
\]
Here $a$ is odd, $b=2$ or $6$, and $\gcd(a,b)=1$.
Remove the additive caps and one ordinary disc carrying the clutching;
fill every multiplicative site by its stated Mumford cap, and denote
the resulting common piece by $W$. Denote its regular torus-bundle
part by $V$. In particular all finite affine twists are outside $W$,
and $V$ has the linear section used in the marked construction.

\item There is an exact integral fact about these thirteen rows. The columns
of their actual multiplicative angular vanishing lattices, without
saturation, generate precisely
\[
 \Lambda_0=\ker\gamma^*=\mathbb Zu\oplus\mathbb Zw
                    \oplus\mathbb Z\delta.                    \tag{D.C24}\label{maps:eq:D.C24}
\]
The same is true of all columns of $A_i-I$. Both assertions are checked
by integer Hermite forms, not ranks or component labels, in the
certificate below. Filling a Mumford cap kills its base meridian and
its angular lattice. The usual torus-bundle presentation therefore
gives
\[
 \pi_1(W)=\mathbb Z\gamma\times F_{r-1},
                                                               \tag{D.C25}\label{maps:eq:D.C25}
\]
where $r\ge1$ is the number of unfilled discs. Every remaining
monodromy is the identity on the quotient $\mathbb Z\gamma$.
The section makes the displayed extension split. In this table
$r\le2$, but only torsion-freeness is needed below.

\item We also need an integral cohomology comparison. Each primitive
Mumford cap has a surjective degree-two boundary map, as proved in
\ref{maps:src:D.2.26}. If there are $N$ original punctures, then $b_1(V)=N$ by the
second assertion in \eqref{maps:eq:D.C24}. A filled boundary has first Betti number
one greater than its cap. In the degree-one Mayer--Vietoris map the
target rank minus the source rank is consequently $r$, exactly
$b_1(W)$ by \eqref{maps:eq:D.C25}. Its kernel has rank zero and is a subgroup of a
free group, hence is zero. Eliminating the surjective cap degree-two
maps now shows that
\[
 H_2(V;\mathbb Z)\longrightarrow H_2(W;\mathbb Z)
                  \text{ is onto}.                            \tag{D.C26}\label{maps:eq:D.C26}
\]
Both $H_1(V)$ and $H_1(W)$ are free. Thus restriction from $W$ to $V$
is injective in degree-two integral cohomology and also with
$\mathbb F_2$ coefficients. These facts will prevent a swept-class
ambiguity in the auxiliary-line calculation.

\item \emph{2. The actual symmetric line and its Chern class.}
On a multiplicative chart let $h_1$ be the elliptic period, let $h_G$
represent the section $G$, and let $R$ be the contraction multiplier.
The section pair is $P=aG,Q=bG$. On an auxiliary line with coordinate
$\zeta$ put
\[
\begin{aligned}
 h(z,v,\zeta)&=(h_1z,h_G^av,(-1)^a z^{-a}\zeta),\\
 F(z,v,\zeta)&=(h_G^{-b}z,Rv,
                 h_G^{-ab/2}R^{b/2}v^b\zeta),\\
 J(z,v,\zeta)&=(z^{-1},v^{-1},z^{-a}\zeta).
\end{aligned}                                                   \tag{D.C27}\label{maps:eq:D.C27}
\]
All exponents are integers. Direct multiplication gives
$hF=Fh$, $J^2=1$, $JhJ=h^{-1}$ and $JFJ=F^{-1}$; the integral
logarithmic constant $a$ in the third identity exponentiates to one.
The displayed $J$ is normalized on the rigidifying origin.
The alternating Chern form is $\Omega$. The normalized half-period
signs are $(-1)^a$ on $u,w$ and $+1$ on $\delta,\gamma$; thus its
mod-two semicharacter is
\[
                    q=a(u+w+uw).                              \tag{D.C28}\label{maps:eq:D.C28}
\]
Every actual monodromy preserves this semicharacter. The certificate
checks this on all sixteen vectors modulo two at every marked site.
The symmetric rigidified Appell--Humbert identifications consequently
patch on overlap markings: an isomorphism preserving the rigidification
is unique, so their compositions satisfy the cocycle condition.
This constructs the holomorphic auxiliary line on $V$ from the actual
factors \eqref{maps:eq:D.C27}.
Since $b$ is even, the half-$\delta$ translation lifts in these factors
as $(z,v,\zeta)\mapsto(z,-v,\zeta)$. It commutes with $h,F$.
The overlap identifications preserve this lift: in the invariant
prequantum marking the lift depends only on $\Omega$ and $\delta/2$,
and each monodromy fixes $\delta$ and preserves $\Omega$ and $q$.
Thus the line has the global involution linearization used below.

\item Compare it with the flat prequantum family having connection potential
$\pi i\Omega(x,dx)$ and the same semicharacter, as in the fully marked
comparison in \ref{maps:src:D.2.8}. The monodromies preserve this family. Both lines
are symmetric and rigidified with alternating form $\Omega$; their
ratio is therefore fibrewise flat of order dividing two. The normalized
square-trivializations on the fibres are unique and patch, making the
ratio globally two-torsion. Since $H^2(V;\mathbb Z)$ is free, its Chern
class is zero. The actual line has precisely the pure fibre class
$\Omega$ and zero value on the marked swept cycles. This argument
controls the mixed values as well as the alternating fibre form.

\item These factors extend to a line $\mathcal L$ on $W$. Here is the
integral local extension, including the two parities that we need.
In a genuine $SL_2(\mathbb Z)$ elliptic marking the actual Tate
valuations are
\[
 (\operatorname{ord}h_1,\operatorname{ord}h_G,
       \operatorname{ord}R)=(n,j,\rho),\qquad
 P_{\rm val}=\begin{pmatrix}n&-bj\\aj&\rho\end{pmatrix}.
                                                               \tag{D.C29}\label{maps:eq:D.C29}
\]
The determinant is one at the zero and zero at every other
multiplicative site. These are raw valuations from the actual
monodromy; no residue of $j$ modulo $n$ is substituted.
At the zero, assign the following integer values to the lattice
vertices and interpolate linearly on the given unimodular fan:
\[
 \varphi(x,y)=
 \frac{-a\rho x^2-2abjxy+bn y^2+ax}{2}.                        \tag{D.C210}\label{maps:eq:D.C210}
\]
Indeed $n\rho+abj^2=1$, so $n,\rho$ are odd, and the expression is
integral at every integer vertex. Its two period differences are
\[
\begin{aligned}
 \varphi(x+n,y+aj)-\varphi(x,y)&=-ax,\\
 \varphi(x-bj,y+\rho)-\varphi(x,y)
          &=by+(b\rho-abj)/2.
\end{aligned}                                                   \tag{D.C211}\label{maps:eq:D.C211}
\]
They are exactly the valuations of the line factors in \eqref{maps:eq:D.C27}.
The usual monomial frames with these integral slopes therefore extend
the line and its commuting lattice identifications over every toric
chart. This uses an integral piecewise affine support function; no
positivity of the auxiliary line is required.

\item At a rank-one point write the two valuation periods as
$H=\alpha\nu$, $F=\beta\nu$, where
$\nu=(\nu_1,\nu_2)$ is primitive and
$\gcd(\alpha,\beta)=1$. Thus
$n=\alpha\nu_1$, $aj=\alpha\nu_2$,
$-bj=\beta\nu_1$ and $\rho=\beta\nu_2$.
The identity $-a\beta\nu_1=b\alpha\nu_2$ shows that
$k=a\nu_1/\alpha$ is an integer. On the valuation line use
\[
             \varphi_1(t)=-k t(t-\alpha)/2.
\]
Its two exact period differences are
$-a\nu_1t$ and
$b\nu_2t+(b\rho-abj)/2$, respectively, as required by
\eqref{maps:eq:D.C27}. At every odd-isotropy site $\nu_1$ is odd and
$\beta\nu_1=-bj$ is even, so $\beta$ is even and $\alpha$
is odd. The values of $\varphi_1$ at integer vertices are therefore
integral. At the unique even-isotropy site the valuation vector is
$(2,aj)$ with $aj$ odd, so $\alpha=1$ and the same integrality
holds. Integral monomial frames on the one-dimensional fan, retaining
the invertible second short coordinate, give the required extension.
The two displayed period identities prove that both actual deck
factors descend this extension; it is not enough to match only one
of them.
Holomorphic unit factors extend in the same frames and affect neither
integral slope nor stabilizer character. The local extensions agree
with the line on their annuli and hence glue to $\mathcal L$ on $W$.
By \eqref{maps:eq:D.C26}, its Chern class has exactly the pure fibre and zero swept
values already computed on $V$.

\item \emph{3. The fixed curves and the two line characters.}
Let $\iota$ be the scalar involution $(z,v)\mapsto(z,-v)$.
The exact primitive angular vectors show that its fixed locus on $W$
is the disjoint union of one elliptic curve $E$ at the unique $I_2$
site and the rational vertical-edge curve $C$ at the rank-two zero.
All other rank-one scalar stabilizers have odd order; this is checked
from the actual vanishing vectors in the certificate. The zero has
valuation lattice all of $\mathbb Z^2$, and $\delta=(0,1)$ in that
short lattice. Thus there is one vertical-edge orbit, with the same
normal model
$\mathcal O(-1)_{+1}\oplus\mathcal O(-1)_{-1}$ as in \ref{maps:src:D.2.14}.

\item At the zero the vertical slope difference is
\[
 \varphi(x,y+1)-\varphi(x,y)=-abjx+bny+bn/2.
\]
Its two nonconstant coefficients are even and its constant is odd,
since $b/2$ and $n$ are odd. The induced line character along $C$
is therefore nontrivial. The degree along its adjacent toric wall is,
up to the immaterial orientation convention for the divisor,
$-a\rho+abj$, which is odd. The other diagonal convention changes
this expression by an even integer, so
\[
 \operatorname{char}_{\iota}(\mathcal L|_C)=1,
          \qquad \deg(\mathcal L|_C)\equiv1\pmod2.            \tag{D.C212}\label{maps:eq:D.C212}
\]
At the $I_2$ site, $j$ is odd and $\rho=-abj^2/2$. The coordinates
\[
 x=z^{(1-aj)/2}v,\qquad y=z^{-aj}v^2
\]
are an integral unimodular change of the short torus coordinates;
they send the primitive vanishing vector $(2,aj)$ to $(1,0)$.
The scalar involution acts as $(-x,y)$, giving sign $-1$ on both
normal directions to $E$. The $h$ multiplier is, up to a unit,
$(-1)^a x^{-2a}y^a$. After replacing $F$ by $Fh^{bj/2}$, its
$x$ exponent is zero and its $y$ exponent is $b/2$. The support
function $-at(t-1)$ has all slopes even. Hence
\[
                 \operatorname{char}_{\iota}(\mathcal L|_E)=0.
                                                               \tag{D.C213}\label{maps:eq:D.C213}
\]

\item There are no further fixed points in the removed caps when they are
reattached. An original additive cap of this table has zero line
multidegree and a free scalar circle, by its actual product description
in \ref{maps:src:D.2.26}. On a finite cap, a scalar fixed-point equation on the covering
torus would imply $jn_i/m_i=0$ in $\mathbb R/\mathbb Z$ on applying
$\gamma^*$. Simple connectivity requires $\gcd(m_i,n_i)=1$, so the
root power is trivial; translation by $\delta/2$ cannot then fix a
torus point. Away from the central torus the base rotation gives the
same conclusion. The ordinary clutch disc also has free scalar action.
Thus $X^{\iota}=C\sqcup E$ for every multiplicity and twist in the
statement, including even finite multiplicities.

\item \emph{4. The global mod-two attachment.}
\ref{maps:src:D.2.26} gives $H_2(X)=H_3(X)=\mathbb Z/2$, generated in degree two by
the image of the primitive fibre class $s$; there are no additional
regular sweeps. We verify which cohomology class this is on $W$.
For every finite log in this table, $b/q_i=2$. Its complete two normal
columns \eqref{maps:eq:D.HY7} have pure fibre coefficients
\[
                 bn_i/q_i=2n_i,
                \qquad -\alpha_i b r_i,                       \tag{D.C214}\label{maps:eq:D.C214}
\]
where $\alpha_i m_i+\beta_i n_i=1$. Both coefficients are even,
for arbitrary allowed $m_i,n_i,r_i$. The original additive cap normal
columns have zero fibre coefficient in the rigidified marking.
The ordinary clutching $e\gamma$ has fibre coefficients divisible by
$b$, and its global relation is $be\,s$. Thus the functional that is
one on $s$ and zero on all swept generators annihilates every normal
column modulo two. It induces the nonzero class $\alpha$ on $X$.
This uses the full integral normal columns, not a transfer of even
degree. Its restriction to $V$ is exactly the reduction of the Chern
class just constructed, and \eqref{maps:eq:D.C26} now implies
\[
                  \alpha|_W=\rho_2c_1(\mathcal L).             \tag{D.C215}\label{maps:eq:D.C215}
\]
The absence of an $\operatorname{Ext}(H_1(W),\mathbb F_2)$ term is
part of \eqref{maps:eq:D.C25}, not an implicit universal-coefficient choice.

\item \emph{5. The full two-primary smooth invariant.}
Since $H^2(X;\mathbb Z)=0$, the complex manifold has $c_1=0$ and is
spin. Put $G=\langle\iota\rangle$ and let
$\tau\in H^1(BG;\mathbb F_2)$ be its generator. As in \ref{maps:src:D.2.8},
$\alpha$ has a Borel equivariant lift $\widetilde\alpha$:
$H^1(X;\mathbb F_2)=0$ kills the first possible obstruction, and the
fixed point splits the map to $BG$, killing the obstruction into
$H^3(BG;\mathbb F_2)$.

\item We check the relative fixed-point normalization on $W$. The scalar
circle acts trivially on its fundamental group, by an isotopy fixing
a point of $C$. That point splits the Borel extension, so
\[
                 \pi_1(W_G)=\pi_1(W)\times G.
\]
The first factor is torsion-free by \eqref{maps:eq:D.C25}. Consequently every fixed
point gives the same order-two element $(1,\iota)$ in this product.
By \eqref{maps:eq:D.C215}, the difference between $\widetilde\alpha|_{W_G}$ and
the equivariant Chern class of $\mathcal L$ vanishes in ordinary
cohomology. Pull it to the universal cover and use the degree-two
Cartan--Leray sequence: the difference comes from the group cohomology
of $\pi_1(W)\times G$. Its possible contributions are a common
multiple of $\tau^2$ and a mixed class $\tau\xi$ with
$\xi\in H^1(W;\mathbb F_2)$; its ordinary group contribution is
zero, since $H^2(B\pi_1(W))\to H^2(W)$ is injective. After one global
adjustment by $\tau^2$, \eqref{maps:eq:D.C212}--\eqref{maps:eq:D.C213} give
\[
 \widetilde\alpha|_C=\alpha|_C+\tau^2,
 \qquad \widetilde\alpha|_E=\alpha|_E+\tau\eta,
       \quad\eta\in H^1(E;\mathbb F_2).                       \tag{D.C216}\label{maps:eq:D.C216}
\]
In particular two unrelated fixed-point normalizations have not been
chosen.

\item Both complex normal directions have sign $-1$. Their normal first
Chern class is even because $c_1(X)=0$ and the tangent Chern class of
each fixed curve is even. Their equivariant Euler class is therefore
$\tau^4$. Equivariant Thom localization yields
\[
 \langle\alpha^3,[X]_2\rangle
 =\int_C\frac{(\alpha|_C+\tau^2)^3}{\tau^4}
  +\int_E\frac{(\alpha|_E+\tau\eta)^3}{\tau^4}
 =\int_C\alpha|_C=1.
\]
The elliptic contribution is zero since degree-one classes on an
oriented elliptic curve square to zero. Finally the exact primary
calculation
$\Omega^{\mathrm{Spin}}_6(K(\mathbb Z/2,2))=\mathbb Z/2$,
detected by this cube, and Zhubr's complete classification identify
$X$ with $M_{\mathrm{conic}}$; the classification and the conic's cube
were established in \ref{maps:src:D.2.8} directly from
\cite[Theorems 3, 2.10 and Section 3.12]{manZhubr75}.
This retains the two-primary refined invariant, not just ordinary
Pontryagin classes.
\end{enumerate}
\end{proof}

The exact extraction program \path{code/III/manifolds_tate_markings.py}
constructs the $SL_2(\mathbb Z)$ markings from the actual monodromies
and verifies \eqref{maps:eq:D.C29}. The independent certificate
\path{code/III/fourfolds_D2_c2_AH.py}
checks all $15$ patterns and $13$ marked rows, including $928$
semicharacter substitutions, both full integral Hermite lattices per
row, the exact conjugated monodromy, every toric character and degree
parity, and every finite normal-column parity. Its retained output is
\path{code/fourfolds_D2_c2_AH.out}; the detailed records are
\path{records/fourfolds_D2_c2_AH.json}. The symbolic identities prove
all displayed parameter formulas; no finite multiplicity or twist
bound is used in the theorem.


\begin{theorem}[the residual five- and seven-primary smooth types]\label{maps:src:D.2.36}
\StatementRecord{maps:src:D.2.36}{proved}{Historical events: Claude Fable 5.2 (Anthropic), effort high, 16 November 2026: carried as background without fault found. GPT-6 Astra (OpenAI), ChatGPT Pro mode, effort Pro, S6\_09\_PAPERS, 17 November 2026: targeted proof check. Also checked in this writing. Pre-continuation writing events: GPT-6 Astra (OpenAI), ChatGPT Pro mode, effort Pro, S6\_10\_PAPERS, 18 November 2026: targeted proof check. Continuation metadata events: GPT-6 Astra (OpenAI), ChatGPT Pro mode, effort Pro, S6\_10\_PAPERS continuation, 19 November 2026: used and rechecked in this run; status-vocabulary reconciliation; metadata reconciliation only; no new mathematical check. \emph{Pre-continuation recorded proof scope.} GPT-6 Astra (OpenAI, ChatGPT Pro mode), effort Pro, 2026-11-17 UTC; Claude Fable 5.2 (Anthropic), effort high, 16 November 2026; also checked in this writing by GPT-6 Astra (OpenAI), ChatGPT Pro mode, effort Pro, 18 November 2026}{the hypotheses, notation and definitions stated here; \cref{maps:src:D.2.26}; \cref{maps:src:D.1.11}; \cref{maps:eq:D.HY7}; \cref{maps:eq:D.OP3}; \cref{maps:eq:D.OP2}; \cref{maps:eq:D.OP4}; \cref{maps:eq:D.OP7}; \cref{maps:eq:D.OP8}; \cref{maps:src:D.1.10}; \cref{maps:eq:D.OP11}; \cref{maps:eq:D.OP10}; \cite{manZhubr75}}


Consider the $36$ marked additive-filling patterns on the $30$ initial
rows of \ref{maps:src:D.2.26} with residual $c=5$ or $7$ and no orbifold obstruction.
Use the primitive Mumford cap at every multiplicative position, the
specified minimal original or finite cap at each additive position,
and no additional regular logarithmic transformation. Allow every
positive multiple of the finite monodromy orders and every free
affine twist for which the resulting manifold is simply connected.
For $p=c$, every such manifold is spin and has complete Zhubr data
\[
 H_2=\mathbb Z/p,\qquad b_3=0,\qquad \mu=0,
 \qquad \mathfrak p\in\{0,1\}\quad\text{up to }\operatorname{Aut}(\mathbb Z/p).
 \tag{D.OP1}\label{maps:eq:D.OP1}
\]
Here $4\mathfrak p=\operatorname{PD}(p_1)$, as in \ref{maps:src:D.1.11}.
The following entirely local rule determines which value occurs.
There is exactly one zero-weight multiplicative position
whose double curve is fixed by the scalar subgroup $C_p$; it is of
type $I_p$. In its genuine Tate marking put
\[
 (\operatorname{ord}h_1,\operatorname{ord}h_G,
                  \operatorname{ord}R)=(p,r,\rho),\qquad
 w=(ar)^{-1}\pmod p.
\]
Then $\mathfrak p=0$ if and only if $w^2=1$; otherwise its canonical
value is $1$. Thus the $c=5$ patterns give exactly two oriented
diffeomorphism types and the $c=7$ patterns give exactly one.
These types are specified by the full published classifying
invariants \eqref{maps:eq:D.OP1}, rather than by homology alone.

\end{theorem}


\begin{proof}
\begin{enumerate}[label=\arabic*.,leftmargin=*]
\item \emph{Marked data and the class which must be extended.}
\ref{maps:src:D.2.26} gives $H_2=H_3=\mathbb Z/p$ and zero middle Betti number
under the stated simple-connectivity condition. Its primitive fibre
class $s$ generates $H_2$. Let $u\in H^2(X;\mathbb F_p)$ evaluate
to one on $s$. In the regular torus lattice the alternating form is
\[
 \Omega=a\,e_1^*e_2^*-b\,\delta^*\gamma^*,\qquad p\mid b,
                                      \quad p\nmid a.
 \tag{D.OP2}\label{maps:eq:D.OP2}
\]
Its values are one on $s$ and zero on the rigidified swept cycles.
The entire finite-cap normal image in \eqref{maps:eq:D.HY7} has fibre
coefficients divisible by $b/q_i$, and hence by $p$ in these
patterns. An allowed unchanged additive cap has zero fibre
coefficient, and the ordinary clutch contributes $be$.
Consequently this functional descends modulo $p$ through every
actual cap relation. This is the integral attachment calculation
which identifies its extension with $u$ for all finite orders
and twists; a finite covering transfer is not used.

\item Remove all additive caps and one ordinary clutch disc, retaining
all multiplicative Mumford caps, and denote the remaining space by
$W$. Its regular part is $V$. The actual multiplicative angular
vanishing columns generate $\ker\gamma^*$ primitively in all
$30$ rows: their Smith factors are $(1,1,1)$. The exact certificate
below uses the full marked columns, including the rank-two cusp.
Van Kampen therefore gives
\[
                  \pi_1(W)=\mathbb Z_\gamma\times F_r
 \tag{D.OP3}\label{maps:eq:D.OP3}
\]
for the free base group with the indicated caps removed. In
particular this group is torsion-free. Each primitive rank-one
Mumford cap and the one-component rank-two cap is generated in
degree two by its boundary. In degree one its sole additional
kernel is the base meridian: the angular lattice is already the
primitive image of the local monodromy minus the identity.
The unfilled ordinary disc makes these peripheral base classes
independent in the graph. Mayer--Vietoris therefore shows that
$H_2(V)\to H_2(W)$ is onto. Together with the free $H_1(W)$ of
\eqref{maps:eq:D.OP3}, this removes any universal-coefficient ambiguity in
comparing degree-two classes on $W$ from their actual values on $V$.

\item \emph{A doubled symmetric line.}
The use of $2\Omega$ makes its Appell--Humbert semicharacter
trivial modulo two, irrespective of the parity of $a,b$.
On an auxiliary holomorphic line use the multiplicative factors
\[
\begin{aligned}
 h(z,v,\zeta)&=(h_1z,h_G^a v,z^{-2a}\zeta),\\
 F(z,v,\zeta)&=(h_G^{-b}z,Rv,
                         h_G^{-ab}R^b v^{2b}\zeta),\\
 J(z,v,\zeta)&=(z^{-1},v^{-1},z^{-2a}\zeta).
\end{aligned}                                                    \tag{D.OP4}\label{maps:eq:D.OP4}
\]
Direct substitution gives $hF=Fh$, $J^2=1$, $JhJ=h^{-1}$
and $JFJ=F^{-1}$. The last identity uses precisely the square
of $h_G^{-ab}R^b$. The alternating Chern form is $2\Omega$.
These symmetric line bundles are rigidified at the origin.
Normalized identifications on overlap markings therefore patch:
an automorphism of a rigidified holomorphic line bundle on a
compact torus is the identity. Comparison with the flat real
prequantum family for $2\Omega$ has globally two-torsion Chern
class. Indeed its fibrewise square is canonically trivial after
rigidification, and the unique normalized trivializations patch.
Reduction modulo the odd prime $p$ kills this difference. Thus
half its Chern class modulo $p$ is exactly \eqref{maps:eq:D.OP2}, including
the zero values on every swept cycle.

\item We give the required extensions and scalar characters explicitly.
At the rank-two position write the two actual valuation periods as
\[
 H=(n,ar),\quad F=(-br,\rho),\qquad n\rho+ab r^2=1.
\]
An integral support function is
\[
 \psi(x,y)=-a\rho x^2-2ab rxy+bn y^2+ax.
 \tag{D.OP5}\label{maps:eq:D.OP5}
\]
Its differences under $H,F$ are respectively
$-2ax$ and $2by+b\rho-ab r$, the valuations of the
two factors in \eqref{maps:eq:D.OP4}. The same unimodular fan consequently
extends the line over the cusp. Its vertical slopes
$-2ab rx+bn(2y+1)$ are divisible by $p$.
The scalar $C_p$ character is therefore zero on its fixed
rational double curve $C$. The wall expression, using the two
opposite vertices $(1,0),(-1,1)$ and edge endpoints
$(0,0),(0,1)$, is
\[
 \psi(1,0)+\psi(-1,1)-\psi(0,0)-\psi(0,1)
                         =-2a(\rho-br).
 \tag{D.OP6}\label{maps:eq:D.OP6}
\]
The divisor convention can reverse the sign of this wall
expression; in either convention its half is nonzero modulo
$p$, since $\rho\equiv n^{-1}\pmod p$ and $p\nmid a$.
Thus the class initially normalized by $u(s)=1$ has
$u[C]\ne0$. For the remainder of the proof, rescale $u$
uniquely so that its evaluation on the complex-oriented
curve $C$ is $1$, and scale the auxiliary Chern class by
the same nonzero field element. This fixes the actual
cohomology-generator convention without a support-function
sign ambiguity. In particular $\alpha_C=1$ below.

\item At a rank-one position write its two periods as
$H=\alpha u,F=\beta u$, with $u=(u_1,u_2)$ primitive and
$\gcd(\alpha,\beta)=1$. Thus $n=\alpha u_1$,
$ar=\alpha u_2$, $-br=\beta u_1$, and $\rho=\beta u_2$.
The identity $a\beta u_1=-b\alpha u_2$ implies
$\alpha\mid au_1$. An explicit integral support is
\[
           \psi_1(j)=-\frac{au_1}{\alpha}\,j(j-\alpha).
 \tag{D.OP5a}\label{maps:eq:D.OP5a}
\]
Its differences under $\alpha,\beta$ are exactly
$-2au_1j$ and $2bu_2j+b\rho-ab r$.
The elliptic combination $\alpha F-\beta H$ has zero
exponent in the degenerating coordinate, by the same identity.
These facts extend both specified factors in \eqref{maps:eq:D.OP4} over the
primitive fan and retain their actual boundary marking.
At the unique $I_p$ position fixed by $C_p$, the short period
is already the primitive vector $(p,ar)$. Choose
integers $A,B$ with $Ap+B ar=1$. In the monomial coordinates
with first row $(A,B)$ and second row $(-ar,p)$, the scalar
action has normal weight $B=w$ and elliptic weight zero
modulo $p$. Its support function can be taken to be
$-ap\,j(j-1)$. All its slopes are divisible by $p$.
The scalar character of the auxiliary line on this elliptic
double curve $E$ is therefore also zero.

\item These are equivariant extensions. On the covering coordinates
the scalar subgroup lifts by $(z,v,\zeta)\mapsto
(z,\lambda v,\zeta)$ for $\lambda^p=1$; it commutes with
\eqref{maps:eq:D.OP4} because $p\mid b$. In the prequantum comparison the
same lift is preserved by every linear marking: each fixes
$\delta$ and $2\Omega$, and the semicharacter is trivial.
The symmetric two-torsion comparison has a unique compatible
odd-order lift. Holomorphic units on the local base discs
have logarithms and contribute no scalar character or wall
degree. Thus the above local lifts patch on $W$.

\item \emph{The fixed set and the global equivariant comparison.}
The primitive rank-one period at a multiplicative cap is the
column $(u_1,u_2)$ obtained by taking the Hermite basis of its
two period columns. Scalar isotropy along its double curve
has order $|u_1|$. The retained exact data show that precisely
one such order is divisible by $p$, and its type is $I_p$.
The rank-two fan has just the vertical rational fixed curve
$C$; its other primitive edge directions are not parallel to
$\delta$ modulo $p$. There are no further fixed strata.
An allowed original additive cap is a product with the free
scalar circle. At a finite cap, primitivity of the actual gamma
numerator in a simply connected member excludes scalar fixed
points: equality of a scalar translate with the $j$th affine
power would force $m\mid j$ from the gamma coordinate, and
then require a nontrivial scalar translation to fix a torus
point. This also covers orders divisible by $p$.
Consequently
\[
                         X^{C_p}=C\sqcup E.
 \tag{D.OP7}\label{maps:eq:D.OP7}
\]

\item The class $u$ admits a Borel equivariant lift. The obstruction
with $H^1(X;\mathbb F_p)$ vanishes by simple connectivity;
the remaining differential into $H^3(BC_p;\mathbb F_p)$
vanishes because any point of $C$ splits the base map.
The whole scalar circle fixes that point and acts trivially
on the based fundamental group of $W$. Hence
\[
 \pi_1(W_{C_p})=(\mathbb Z\times F_r)\times C_p.
\]
Every fixed point determines its same, unique order-$p$
subgroup, by torsion-freeness of the first factor.
The difference between the chosen lift of $u$ and the correspondingly scaled half of
the equivariant Chern class of the auxiliary line has zero
ordinary restriction on $W$. The universal-cover low-degree
sequence places this difference in group cohomology.
Its restrictions to the fixed-point copies of $BC_p$ are
therefore equal. Adjust the lift once by a base class.
Writing $t\in H^2(BC_p;\mathbb F_p)$ and
$\eta\in H^1(BC_p;\mathbb F_p)$, we obtain
\[
 \widetilde u|_C=\alpha_C[C]^*,\qquad
 \widetilde u|_E=\alpha_E[E]^*+\eta\zeta,
                  \quad\zeta\in H^1(E;\mathbb F_p).
 \tag{D.OP8}\label{maps:eq:D.OP8}
\]
In particular both restrictions have zero coefficient of $t$.

\item \emph{All the classifying invariants.}
Since $H^2(X;\mathbb Z)=0$, the complex first Chern class
vanishes and $X$ is spin. The two complex normal weights
at $C$ are $1,-1$, and those at $E$ are $w,-w$.
One can also check the sums from equivariant $c_1$:
it comes from $H^2(BC_p;\mathbb Z)$, and is zero at $C$,
hence zero everywhere. The equivariant Euler classes
have leading terms $-t^2$ and $-w^2t^2$.
Inverting $t$ in the equivariant Thom isomorphism gives
localization on \eqref{maps:eq:D.OP7}. Applying it to $u$, whose integral
has negative degree and is zero, gives
\[
                      \alpha_C+\alpha_E/w^2=0.
 \tag{D.OP9}\label{maps:eq:D.OP9}
\]
No mixed term in \eqref{maps:eq:D.OP8} contributes: it has ordinary
degree one and its product with an ordinary degree-two
Euler correction has degree three on a curve.

\item For each normal bundle with roots $wt+x_1,-wt+x_2$,
the square of their sum is zero on the fixed curve.
The tangent line's square is zero as well. Therefore its
equivariant Pontryagin class divided by the normal Euler
class is exactly $-2$. It follows that
\[
 \langle u\,p_1(X),[X]\rangle
   =-2(\alpha_C+\alpha_E)
   =-2(1-w^2).
 \tag{D.OP10}\label{maps:eq:D.OP10}
\]
Moreover the cube of each restriction in \eqref{maps:eq:D.OP8} is zero:
ordinary degree-two classes square to zero on a curve,
and $\eta^2=0$ at an odd prime. Thus
\[
                        \langle u^3,[X]\rangle=0.
 \tag{D.OP11}\label{maps:eq:D.OP11}
\]
Lemma \ref{maps:src:D.1.10} with $r=0,d=p$ identifies
$H_6(K(\mathbb Z/p,2);\mathbb Z)=\mathbb Z/p$
and proves that \eqref{maps:eq:D.OP11} detects its entire fundamental
class invariant. Consequently $\mu=0$.
Multiplication by four is invertible on $H_2$, so
\eqref{maps:eq:D.OP10} computes the entire refined Pontryagin class.
Because $\alpha_C\ne0$, it is zero exactly when $w^2=1$.
Every nonzero element of $\mathbb Z/p$ is carried to one
by an automorphism; these automorphisms leave $\mu=0$.
Zhubr's Theorem~1, with the definitions on pp.~794--795
\cite{manZhubr75}, now gives precisely \eqref{maps:eq:D.OP1} and the
claimed oriented diffeomorphism classification. Existence
comes from the actual constructed members, so no separate
realization assumption is needed.

\item The local weights in the retained table include both
$w^2=1$ and $w^2\ne1$ for $p=5$, and only $w^2\ne1$
for $p=7$. The code checks all $30$ actual marked rows;
the $36$ additive patterns and all allowed higher orders
use the same local data. This proves the final counts.
\end{enumerate}
\end{proof}

The complete Tate markings are reconstructed by
\path{code/III/manifolds_tate_markings.py}, and
\path{code/III/manifolds_odd_prime_AH.py} checks the doubled
line identities, support functions, primitive angular
lattices, scalar isotropy, unique fixed elliptic curves,
and the exact characteristic numbers for every listed row.
The refined Pontryagin coefficient in the curve-normalized
generator is exactly $(w^2-1)/2$ modulo $p$; its nonzero
orbit is subsequently normalized to one in \eqref{maps:eq:D.OP1}.
Its JSON output records both the signed support expression
and these curve-normalized characteristic coefficients.
No original
multiplicative replacement, extra regular log, or
higher initial weight is included in this theorem.


\begin{theorem}[a uniform weighted algorithm on all natural profiles]\label{maps:src:D.2.37}
\StatementRecord{maps:src:D.2.37}{proved}{Historical events: Claude Fable 5.2 (Anthropic), effort high, 16 November 2026: carried as background without fault found. GPT-6 Astra (OpenAI), ChatGPT Pro mode, effort Pro, S6\_09\_PAPERS, 17 November 2026: targeted proof check. Also checked in this writing. Pre-continuation writing events: GPT-6 Astra (OpenAI), ChatGPT Pro mode, effort Pro, S6\_10\_PAPERS, 18 November 2026: targeted proof check. Continuation metadata events: GPT-6 Astra (OpenAI), ChatGPT Pro mode, effort Pro, S6\_10\_PAPERS continuation, 19 November 2026: used and rechecked in this run; status-vocabulary reconciliation; metadata reconciliation only; no new mathematical check. \emph{Pre-continuation recorded proof scope.} GPT-6 Astra (OpenAI, ChatGPT Pro mode), effort Pro, 2026-11-17 UTC; Claude Fable 5.2 (Anthropic), effort high, 16 November 2026; also checked in this writing by GPT-6 Astra (OpenAI), ChatGPT Pro mode, effort Pro, 18 November 2026}{the hypotheses, notation and definitions stated here; \cref{maps:src:D.2.22}; \cref{maps:src:D.2.33}; \cref{maps:src:D.2.32}; \cref{maps:src:D.2.23}; \cref{maps:eq:D.WP10}; \cref{maps:src:D.2.29}; \cref{maps:src:D.2.34}; \cref{maps:eq:D.WP11}; \cref{maps:eq:D.WP12}}


Choose any of the $172$ natural initial rows of \ref{maps:src:D.2.22},
with $P=aG,Q=bG$, and any integer $d>0$ satisfying
\[
                         \gcd(a,d)=1.                    \tag{D.WP1}\label{maps:eq:D.WP1}
\]
Iterate its actual contracting translation $d$ times.
Thus the new translation is $dbG$ and its linearization
is the norm of the original one under the first $d$
translates. Keep the same unimodular Mumford fans,
quotiented by the new period lattice.

Every original/Mumford choice represented by \ref{maps:src:D.2.33},
every free finite-good additive logarithm, every number
of free regular logarithms and every integer peripheral
clutch is covered by the following exact integer
algorithm. Star caps are retained original. The
$495$ initial presentation patterns provide the finite
input; there is no upper bound on $d$, logarithmic
orders or affine coordinates.

Write $A_i$ for the recorded four-dimensional
monodromies in the basis $(e_1,e_2,\delta,\gamma)$.
Replace them by
\[
 A_i(d)=S_d^{-1}A_iS_d,\qquad
 S_d=\operatorname{diag}(1,1,1,d).
                                                               \tag{D.WP2}\label{maps:eq:D.WP2}
\]
Equivalently multiply their first three gamma-column
entries by $d$; every other entry is unchanged.
The primitive common alternating functional is
\[
 \omega_d=a e_1^*e_2^*-db\,\delta^*\gamma^*.                \tag{D.WP3}\label{maps:eq:D.WP3}
\]
Insert identity matrices for regular logarithms.
An additional ordinary order-one logarithm with power
$-e\gamma$ realizes the clutch integer $e$.

Let $N$ be the number of sites, with
$A_0(d)\cdots A_{N-1}(d)=I$, and compute the integral
kernel
\[
 K=\ker\bigl[(A_0(d)-I)\ \cdots\ (A_{N-2}(d)-I)\bigr].
                                                               \tag{D.WP4}\label{maps:eq:D.WP4}
\]
The target of the degree-two presentation is
$\mathbb Zs\oplus K$. For a fixed vector at site $i$,
let $J_i$ denote its usual oriented peripheral
insertion in $K$: the inclusion in block $i$ for
$i<N-1$, and, at the last site, the blocks
$-A_{j+1}(d)\cdots A_{N-2}(d)$.
All these maps are evaluated in a primitive integral
basis of \eqref{maps:eq:D.WP4}.

At a finite-good site of linear order $h$, compute
a primitive fixed basis $F:\mathbb Z^2\to\Lambda$
and the norm $N_h=I+\cdots+A^{h-1}$. If $H$ is an
integral basis matrix for the image of
$F^{-1}N_h:\Lambda\to\mathbb Z^2$, the projected lattice
is $F(H/h)\mathbb Z^2$. For any invariant power
$v=Fz$ and order $m=hk$, put
\[
 \begin{gathered}
 w=hH^{-1}z\in\mathbb Z^2,\qquad t=\gcd(m,w_1,w_2),\\
 \text{the action is free iff}\quad
 \gcd(m,v_1,v_2,v_3,v_4)=1,\quad h\mid m/t.
 \end{gathered}
                                                               \tag{D.WP5}\label{maps:eq:D.WP5}
\]
For every free such logarithm include the five columns
\[
 \bigl(-\omega_d(\lambda,v),\,J_iN_m\lambda\bigr)
       \quad(\lambda=e_1,e_2,\delta,\gamma),\qquad
                         (0,J_iv),                       \tag{D.WP6}\label{maps:eq:D.WP6}
\]
where $N_m=(m/h)N_h$. At a regular logarithm use the
same formula with $A=I,N_m=mI$ and the sole freeness
condition $\gcd(m,v_1,v_2,v_3,v_4)=1$.

At every unchanged original or Mumford cap include
the full marked fixed-normal columns
\[
                       (0,J_iF_i),\qquad
                   F_i\text{ a basis of }\ker(A_i(d)-I).
                                                               \tag{D.WP7}\label{maps:eq:D.WP7}
\]
For a nonzero-Chern original also include
$(dh_i,0)$, where $h_i$ is its recorded
$b/q_i^c$ or $b/n_i$ from \ref{maps:src:D.2.32}--\ref{maps:src:D.2.33}.
These are the entire global images of the boundary
kernels, including any torsion in a boundary.
Let $B$ be the resulting integer matrix and
$T=\operatorname{Tor}\operatorname{coker}B$.

For the component correction define $\beta$ as follows.
A zero-free rank-one site with $P$ narrow contributes
$\gcd(d,n_i)-1$; if $P$ is non-narrow, its marked
Mumford cap contributes zero. A zero-degree original
additive cap contributes $\rho_i(d)-1$, where
$\rho_i(d)$ is the number of component orbits under
$Q^d$, computed by raising the recorded permutation
to the $d$th power. A nonzero-Chern original contributes
its unchanged $s_i-2$ or $n_i-2$.
Let $L$ count these last originals. The weighted
rank-two zero has $d$ normalized components and is
not included in $\beta$.

The full fundamental group is the Seifert presentation
of \ref{maps:src:D.2.23} with the actual gamma coordinates of the
power vectors, the displayed orders and clutch.
Write its abelianization as $\mathbb Z^r\oplus A_1$.
Then
\[
 b_2=
 \begin{cases}
 d-1+\beta+2r,&L=0,\\
 d-1+\beta+L+r,&L>0,
 \end{cases}
 \qquad
 b_3=
 \begin{cases}
 2\beta+2r,&L=0,\\
 2\beta+2L,&L>0.
 \end{cases}                                               \tag{D.WP8}\label{maps:eq:D.WP8}
\]
The complete integral groups are
\[
\begin{array}{c|c}
j&H_j(X;\mathbb Z)\\\hline
0,6&\mathbb Z\\
1&\mathbb Z^r\oplus A_1\\
2&\mathbb Z^{b_2}\oplus T\\
3&\mathbb Z^{b_3}\oplus T\\
4&\mathbb Z^{b_2}\oplus A_1\\
5&\mathbb Z^r .
\end{array}                                                \tag{D.WP9}\label{maps:eq:D.WP9}
\]
In particular $\chi(X)=2d$. The executable
\path{code/III/manifolds_weighted_profiles.py} implements
all steps, including actual freeness, the full group
presentation and both Smith calculations.

\end{theorem}


\begin{proof}
\begin{enumerate}[label=\arabic*.,leftmargin=*]
\item \emph{The geometric meaning of the new weights.}
Iteration of the original lifted translation replaces
its line multiplier by the product of its first $d$
translates. If its initial zero is one component of
$I_n$, the new zero is its orbit sum, with repetitions.
It need not be a whole fibre. This is the prescribed
norm linearization, so no unsupported line-bundle
modification is being introduced.
Outside the finite replacements the new space is
the connected gamma cover defined by $S_d$ of the
old original or Mumford model. It follows directly
that the good-reduction monodromies are \eqref{maps:eq:D.WP2};
the same construction on their root discs permits
all the new free finite quotients.

\item At the unique rank-two zero the old valuation
lattice was the full short lattice. Multiplying
one long period by $d$ gives index $d$. The same
smooth triangular fan therefore has $d$ copies
of the old normalized $dP_6$ in its quotient.
At a zero-free rank-one site, write the old
monodromy as $I+\alpha\beta^*$ with $\alpha$
primitive and $\beta^*$ primitive. The cover
replaces $\beta^*$ by $\beta^*S_d$.
The exact factorization certificate retains all $436$
zero-free rank-one matrices as $\alpha\beta^*$, with
both vectors primitive. If $P$ is narrow, the content
of the first three entries of $\beta^*$ is exactly
$n_i$, and its gamma entry is prime to $n_i$. Thus
the new content is $\gcd(d,n_i)$. At every non-narrow
$P$ site that vertical content is exactly one, so
its new content is one for every $d$. These actual
vectors and entries are recorded in
\path{records/manifolds_weighted_rankone_certificates.json}.
Thus the respective component counts are precisely
those in the statement. This also follows directly
by lifting the component cycle through the gamma
cover. Zero-degree original multiplicative caps
have the same component count and the same global
kernel image; no relative diffeomorphism is asserted.

\item An original zero-degree additive cap is
$S^1_\delta$ times the mapping torus of the Kodaira
tree under $Q^d$. It has $\rho_i(d)$ fundamental
component cycles. At a nonzero-Chern original, $Q$
was already narrow and remains narrow; its
primitive Chern vector is unchanged. Its component
count and its principal-torus description from
\ref{maps:src:D.2.32} or \ref{maps:src:D.2.33} are therefore unchanged.
All these caps have torsion-free top homology and
Euler number zero. The rank-two cusp has Euler
number $2d$ by its normalization complex.

\item \emph{Why the regular integral lattice is free in the stated sector.}
For each original row, \ref{maps:src:D.2.23} gives
\[
 \Lambda_{\{A_i\}}=\mathbb Z\gamma,\qquad
 (\Lambda^2\Lambda)_{\{A_i\}}=\mathbb Zs,
\]
with primitive functional \eqref{maps:eq:D.WP3} at $d=1$.
There is a finite exact certificate for all other
weights. Set $d=0$ in the polynomial matrices
\eqref{maps:eq:D.WP2}, meaning that their first three gamma
entries are zero. For every one of the $172$ rows
the common degree-one relation matrix still has
rank three, the degree-two matrix has rank five,
and every prime in their nonunit Smith factors
divides $a$. The complete constant matrices and
Smith factors are retained in
\path{records/manifolds_weighted_zero_certificates.json}.

\item Here is why this finite certificate proves the
unbounded claim. At a prime not dividing $d$,
$S_d$ is an integral local isomorphism, so the
relations have the original saturated ranks.
At a prime dividing $d$, condition \eqref{maps:eq:D.WP1} makes
$a$ a unit, and reduction modulo that prime is
the just-certified $d=0$ matrix, still of full
rank three or five. Hence no prime divides a
nonunit Smith factor at the new weight.
The two common coinvariants are therefore free
of ranks one and one for every allowed $d$.
Moreover \eqref{maps:eq:D.WP3} is primitive since
$\gcd(a,db)=1$, so it is their actual integral
degree-two functional.

\item The graph Serre sequence consequently gives
\[
 H_1(V)=\mathbb Z\gamma\oplus\mathbb Z^{N-1},
 \qquad H_2(V)=\mathbb Zs\oplus K.                         \tag{D.WP10}\label{maps:eq:D.WP10}
\]
The splitting used in \eqref{maps:eq:D.WP10} is marked. On the
original regular piece the integral fibre class
has value zero on every fixed-loop peripheral
sweep, by the unrolled comparison and the single
diagonal relation of \ref{maps:src:D.2.23}--\ref{maps:src:D.2.26}. Pull it back
through the actual gamma cover. It restricts to
$\omega_d$ and still has value zero on every
new fixed sweep: that sweep maps to the old
sweep of $S_d\lambda$.
Primitivity of $\omega_d$ makes this a genuine
integral splitting. This is a pullback of an
already marked class, not an assumption that
an invariant differential form descends integrally.

\item \emph{All local boundary-kernel images.}
For a finite-good matrix, all nonidentity linear
powers have the same fixed real two-plane.
Its projected lattice is
$P\Lambda=(N_h/h)\Lambda=F(H/h)\mathbb Z^2$.
The translation $v/m$ on that quotient has exact
order $m/t$. A nonidentity linear power can have
a fixed point exactly when its exponent is
divisible by $m/t$. Requiring all such exponents
to have identity linear part is $h\mid m/t$.
For the identity-linear powers one must additionally
have the full translation of order $m$, which is
the primitive-vector condition in \eqref{maps:eq:D.WP5}.
This proves the stated freeness criterion for
arbitrary coordinates and orders.

\item The full normal Thom image is generated by the
four covering-torus loops and the affine meridian.
The cut of their swept normal tori into the regular
fibre is respectively $N_m\lambda$ and $v$.
Integration of the actual pulled-back marked
class from \eqref{maps:eq:D.WP10} gives respectively
$-\omega_d(\lambda,v)$ and zero, exactly as in
\ref{maps:src:D.2.29}. These are complete integral coordinates
in the free target \eqref{maps:eq:D.WP10}, so a local torsion
ambiguity cannot be lost by this computation.
In particular no claim of integral form descent
on the finite quotient is needed.
This proves \eqref{maps:eq:D.WP6}, including finite cap and
normal-line torsion.

\item For a Mumford or zero-degree original cap, every
new fixed normal torus bounds. Its image under
the covering is an old fixed normal torus; the
old bounding section-disc sweep lifts because
its loop is in the covering lattice. Equivalently
use the same toric or Neron-logarithm section
disc and its actual invariant circle actions.
These are all the normal projections, since
the Wang projection takes values in the full
fixed lattice.
After subtracting them, any further kernel
element comes from a fibre class. Rationally
its image is zero exactly on the monodromy-image
subspace. At the rank-two zero, the two primitive
angular normal tori exhaust the boundary kernel:
$H_2(\partial N)=\mathbb Z^6$ and
$H^3(N)=\mathbb Z^2$ by \ref{maps:src:D.2.34}. The regular-fibre
coinvariant of rank four therefore injects. At a
rank-one cap the nodal cycle times $T^2$ supplies
four independent fibre images (whole fibre, two
mixed surviving-circle tori, and transverse area).
At an original zero-degree additive cap the
Kodaira-tree permutation model supplies the whole
fibre and the transverse area, the full two
dimensional coinvariant. These explicit models
prove the required ranks for every weight.
The functional $\omega_d$ annihilates that
subspace. Since the global fibre coinvariant
in \eqref{maps:eq:D.WP10} is free, every such fibre-kernel
class has zero global image, even when it is
torsion in the boundary. This proves the entire
global image \eqref{maps:eq:D.WP7}.

\item For a nonzero-Chern original, the marked logarithm
vectors scale by $d$. The exact area-kernel
calculations of \ref{maps:src:D.2.32}--\ref{maps:src:D.2.33} therefore multiply
their fibre values by $d$, giving $dh_i\mathbb Zs$.
Their complete fixed normal projections still
bound, and the third multiplicative correction
is still annihilated by $\omega_d$.
This proves the remaining column, integrally.

\item \emph{The relative sequence and the free component contribution.}
First attach only the finite and regular logarithmic
caps to $V$, and call the result $E$. There is at
least one unfilled cap, the rank-two zero.
The relative degree-two group of each attached
disc bundle is $\mathbb Z$, whose boundary has
its nonzero multiplicity in its independent base
meridian coordinate. These relative maps into
$H_1(V)$ are injective. Thus $H_2(E)$ is exactly
the quotient of \eqref{maps:eq:D.WP10} by their full Thom
images, including when the cap has torsion.

\item Now attach all the remaining original and Mumford
caps $N_j$. Their relative degree-two groups are
\[
 H_2(N_j,\partial N_j)=H^4(N_j),
\]
free on their fundamental component cycles.
The degree-three boundary images are exactly
the kernels already computed above. Consequently
the pair sequence has the form
\[
 0\longrightarrow\operatorname{coker}B
 \longrightarrow H_2(X)
 \longrightarrow
 \ker\!\left\{\bigoplus_j H^4(N_j)\to H_1(E)\right\}
 \longrightarrow0.                                      \tag{D.WP11}\label{maps:eq:D.WP11}
\]
Its last group is a subgroup of a free group,
so the sequence splits as a sequence of abelian
groups. This avoids all potentially nonprimitive
degree-one boundary splittings.

\item If there are $s$ unfilled caps, $H_1(E)$ has free
rank $s$: the finite/regular power relations
remove their independent base directions.
The final rank is $r$, so the map in \eqref{maps:eq:D.WP11}
has rank $s-r$. Its domain has rank equal to
the sum of the component counts. Subtracting
one per cap gives
\[
 \operatorname{rank}\ker
   =d-1+\beta+L+r.                                        \tag{D.WP12}\label{maps:eq:D.WP12}
\]
This is an exact rank calculation in a free
relative group; torsion meridians can change
its index but cannot change its freeness or
introduce an extension in \eqref{maps:eq:D.WP11}.

\item Over the rationals \eqref{maps:eq:D.WP2} is conjugate to the
original monodromy. Thus the fixed-sweep graph
has the same single diagonal relation as in
\ref{maps:src:D.2.23}--\ref{maps:src:D.2.26}. With no nonzero-Chern original,
the remaining scalar equation is the actual
Seifert Euler equation, so
$\operatorname{corank}_{\mathbb Q}B=r$.
If $L>0$, a nonzero fibre multiple is killed,
and then every fixed sweep is killed rationally;
the corank is zero. Combining this with
\eqref{maps:eq:D.WP12} proves $b_2$ and its full torsion $T$.

\item The fundamental-group argument is also unchanged
integrally. The primitive angular vanishing
vectors are vertical, so they survive in the
cover without a multiple. Restrictions of
$A_i(d)$ to the vertical lattice are the
original restrictions. Their normal closures
therefore kill the same entire vertical
lattice as in \ref{maps:src:D.2.23}, including its original
star thimble when required. Each unchanged
cap has an actual transverse origin disc, so
adds no gamma twist beyond the recorded
peripheral clutch. The finite/regular power
relations then give exactly the full Seifert
group, not merely its abelianization.
Finally $\chi=2d$ and integral duality give
the remainder of \eqref{maps:eq:D.WP8}--\eqref{maps:eq:D.WP9}.
\end{enumerate}
\end{proof}

The program verifies all $172$ weight-zero
prime-support certificates and computes $1626$
full weighted graph presentations on the $495$
patterns, of which $1131$ have weight greater
than one. The all-parameter proof is the
prime-local argument and marked relative
calculation above, not a finite weight search.
Weights sharing a factor with $a$, additional
star logarithmic Mumford choices and arbitrary
initial pairs not obtained by these translation
iterates remain outside this theorem. Its
homology algorithm makes no new smooth
classification assertion.

The shared-even-factor case excluded from the stated formula is handled by the complete integral algorithm of \ref{maps:src:D.2.41}.

\begin{theorem}[the mixed ten- and fourteen-torsion smooth types]\label{maps:src:D.2.38}
\StatementRecord{maps:src:D.2.38}{proved}{Historical events: Claude Fable 5.2 (Anthropic), effort high, 16 November 2026: carried as background without fault found. GPT-6 Astra (OpenAI), ChatGPT Pro mode, effort Pro, S6\_09\_PAPERS, 17 November 2026: targeted proof check. Also checked in this writing. Pre-continuation writing events: GPT-6 Astra (OpenAI), ChatGPT Pro mode, effort Pro, S6\_10\_PAPERS, 18 November 2026: targeted proof check. Continuation metadata events: GPT-6 Astra (OpenAI), ChatGPT Pro mode, effort Pro, S6\_10\_PAPERS continuation, 19 November 2026: used and rechecked in this run; status-vocabulary reconciliation; metadata reconciliation only; no new mathematical check. \emph{Pre-continuation recorded proof scope.} GPT-6 Astra (OpenAI, ChatGPT Pro mode), effort Pro, 2026-11-17 UTC; Claude Fable 5.2 (Anthropic), effort high, 16 November 2026; also checked in this writing by GPT-6 Astra (OpenAI), ChatGPT Pro mode, effort Pro, 18 November 2026}{the hypotheses, notation and definitions stated here; \cref{maps:src:D.2.36}; \cref{maps:src:D.2.26}; \cref{maps:eq:D.C26}; \cref{maps:eq:D.C27}; \cref{maps:eq:D.MP3}; \cref{maps:eq:D.MP5}; \cref{maps:eq:D.C216}; \cref{maps:src:D.2.35}; \cref{maps:src:D.2.8}; \cref{maps:eq:D.OP4}; \cref{maps:eq:D.MP4}; \cref{maps:eq:D.OP8}; \cref{maps:src:D.1.10}; \cite[p.796]{manZhubr75}; \cref{maps:eq:D.MP9}; \cite[p.795]{manZhubr75}; \cref{maps:eq:D.MP7}}


Take one of the twelve minimal patterns in the following table,
with the primitive Mumford filling at every multiplicative position,
the indicated additive choice, and no additional regular logarithmic
transformation. Retain the ordinary clutch $e\gamma$ and allow every
finite multiplicity and affine twist for which the output $X$ is
simply connected. Zero positions use the zero-based source indexing.
Let $Z_{p,\epsilon}$ denote the oriented smooth type of \ref{maps:src:D.2.36} with
$H_2=\mathbb Z/p$, $b_3=0$, $\mu=0$, and refined Pontryagin
orbit $\mathfrak p=\epsilon$, for $\epsilon=0$ or $1$ when that
type occurs there. Then
\[
                  X\cong M_{\mathrm{conic}}\# Z_{p,\epsilon}
                     \quad\hbox{orientation-preservingly},
                                                        \tag{D.MP1}\label{maps:eq:D.MP1}
\]
where the parameters are
\[
\begin{array}{c|r|r|l|l|r|r}
\text{OS/profile}&a&b&\text{zero positions}&\text{additive choice}&p&\epsilon\\\hline
56&1&10&1&\text{none}&5&0\\
56&3&10&3,4&\text{none}&5&1\\
56/II&1&10&3&II\text{ original or log}&5&0\\
56/IV&1&30&3&IV\text{ log}&5&0\\
56/IV&3&10&3&IV\text{ original}&5&1\\
47&1&14&2,3,4&\text{none}&7&1\\
47/II&1&14&2&II\text{ original or log}&7&1.
\end{array}                                               \tag{D.MP2}\label{maps:eq:D.MP2}
\]
There are seven $c=10$ patterns and five $c=14$ patterns on ten
distinct marked initial rows. Thus they give exactly two oriented
smooth types for $c=10$ and one for $c=14$.

\end{theorem}


\begin{proof}
\begin{enumerate}[label=\arabic*.,leftmargin=*]
\item \emph{The actual classes and local extensions.}
\ref{maps:src:D.2.26} gives $H_2(X;\mathbb Z)=H_3(X;\mathbb Z)=\mathbb Z/(2p)$,
with $b_3=0$. In particular $H^2(X;\mathbb Z)=0$, so the complex
first Chern class is zero and $X$ is spin. Use the lattice basis
$(u,w,\delta,\gamma)$ and the primitive alternating form
\[
             \Omega=a\,u^*w^*+b\,\gamma^*\delta^*.
                                                        \tag{D.MP3}\label{maps:eq:D.MP3}
\]
Here $a$ is odd, $b/2$ is odd, and $p\mid b$. The coefficient
$b/q_i$ in every finite-cap normal relation is divisible by $2p$:
the only possibilities are $b/q_i=10$ or $14$, with the $IV$
logarithmic cap using $b/q_i=30/3=10$. The other coefficient is
a multiple of $b$. The allowed original additive caps have zero
fibre coefficient, and the ordinary clutch has coefficient $be$.
Thus the pure fibre functional, zero on the rigidified sweeps,
descends both modulo two and modulo $p$ through the actual
attachments for every allowed finite order and twist.

\item Let $W$ retain all multiplicative Mumford caps and omit the additive
caps and one ordinary disc, and let $V$ be its regular part. The
exact marked angular lattices in all ten rows have Hermite basis
$(u,w,\delta)$; so do the full images of $A_i-I$. Hence
$\pi_1(W)=\mathbb Z_\gamma\times F_{r-1}$, where $r$ is the number
of unfilled discs. The boundary-to-cap degree-two maps are onto.
As in the rank calculation \eqref{maps:eq:D.C26}, the independent peripheral
base classes, with the ordinary disc still unfilled, give
$H_2(V)\twoheadrightarrow H_2(W)$. Both first homology groups are
free. Therefore restriction from $W$ to $V$ is injective in degree
two with integral, mod-two, or mod-$p$ coefficients. This proves
the needed comparison with the globally descended classes.

\item Use the unsquared symmetric factors \eqref{maps:eq:D.C27} for $\Omega$.
All their exponents are integral since $b$ is even. Their
semicharacter $q=a(u+w+uw)$ is preserved by every actual
monodromy, as the exact certificate checks on all sixteen
vectors at each site. Thus the rigidified factors patch.
Their comparison with the marked prequantum family has globally
two-torsion Chern class; $H^2(V;\mathbb Z)$ is free, so that class
is zero. The resulting line has exactly \eqref{maps:eq:D.MP3} on $V$,
including its zero swept values.

\item For completeness the local extension uses the raw Tate periods
\[
 H=(n,ar),\qquad F=(-br,\rho).
\]
At the rank-two zero they have determinant one, and
\[
 \phi(x,y)=\frac{-a\rho x^2-2ab rxy+bn y^2+ax}{2}
                                                        \tag{D.MP4}\label{maps:eq:D.MP4}
\]
has period differences $-ax$ and $by+(b\rho-ab r)/2$.
Here $n,\rho$ are odd. Therefore $\phi$ is integral on every
lattice point. Its vertical slope has odd constant term
$bn/2$, whereas every coefficient of that slope is divisible
by $p$. The line's characters on the rational fixed double
curve $C$ are consequently one for the scalar involution and
zero for $C_p$. The wall degree is odd and a unit modulo $p$,
since its signed value is $-a(\rho-br)$.

\item At a rank-one site take its actual primitive period $u_0=(u_1,u_2)$
and write $H=\alpha u_0,F=\beta u_0$ with coprime
$\alpha,\beta$. The integral support is
\[
       \phi_1(j)=-\frac{a u_1}{2\alpha}j(j-\alpha).
                                                        \tag{D.MP5}\label{maps:eq:D.MP5}
\]
The extracted period identities give $\alpha\mid a u_1$;
either $\alpha$ is odd or $a u_1/\alpha$ is even. Hence \eqref{maps:eq:D.MP5}
is integral. Its two differences are precisely
$-a u_1j$ and $b u_2j+(b\rho-ab r)/2$. The elliptic
zero-period combination has zero exponent in the degenerating
coordinate. These are extensions of both actual factors, with
their given marking, over the primitive toric fans.

\item For the scalar involution there is exactly one further fixed
curve $E_2$, at an $I_2$ site. Its short period $(2,ar)$ is
primitive and \eqref{maps:eq:D.MP5} becomes $-a j(j-1)$, with even slopes.
For $C_p$ there is exactly one further fixed curve $E_p$,
at an $I_p$ site. Its short period $(p,ar)$ is primitive and
the doubled support becomes $-ap j(j-1)$, with slopes
divisible by $p$. The corresponding normal weight is
\[
                         w=(ar)^{-1}\pmod p.             \tag{D.MP6}\label{maps:eq:D.MP6}
\]
There are no other fixed strata for either subgroup. These
claims use the actual primitive two-dimensional period columns,
not merely the Kodaira types. They are checked for every row
in the retained certificate. The original additive caps have
free scalar action. In each simply connected finite cap the
actual gamma numerator is primitive modulo its full order;
applying gamma to an equality with a scalar translate forces
that affine power to be the identity. Thus these caps add no
fixed points even when their orders are divisible by $2$ or $p$.

\item \emph{The two complete primary invariants.}
The mod-two generator $v$ of $H^2(X;\mathbb F_2)$ has a Borel
lift whose characters on $C,E_2$ are respectively one and zero.
To justify a simultaneous choice, the based scalar action on
$W$ gives $\pi_1(W_{C_2})=\pi_1(W)\times C_2$. Its unique
order-two subgroup is the one seen at every fixed point.
The difference from the just-constructed line, being ordinary
zero on $W$, comes from group cohomology. Its restrictions
differ only by the same base class, and possibly a mixed
degree-one term on $E_2$. One base adjustment therefore gives
exactly the restrictions \eqref{maps:eq:D.C216}. The normal Euler class on
both curves is $\tau^4$, because both normal signs are negative
and their total Chern degree is even. The same Thom-localization
calculation as in \ref{maps:src:D.2.35} yields
\[
                         \langle v^3,[X]_2\rangle=1.     \tag{D.MP7}\label{maps:eq:D.MP7}
\]
Consequently the reduction of the spin-bordism class of $X$
to $K(\mathbb Z/2,2)$ is the conic-surgery class, by the full
calculation $\Omega^{\mathrm{Spin}}_6(K(\mathbb Z/2,2))
=\mathbb Z/2$ used in \ref{maps:src:D.2.8}.

\item For the odd reduction take the doubled line \eqref{maps:eq:D.OP4}; its
equivariant characters on $C,E_p$ are both zero. Normalize
$z\in H^2(X;\mathbb F_p)$ by $z[C]=1$. This is possible
because the wall degree in \eqref{maps:eq:D.MP4} is a unit. The torsion-free
group $\pi_1(W)$ again makes the base adjustment common to
both fixed curves. Thus the restrictions have exactly the
form \eqref{maps:eq:D.OP8}. The weights are $(1,-1)$ and $(w,-w)$.
Localization in degree two gives $z[E_p]=-w^2$, the ratio
of the equivariant $p_1$ to the normal Euler class is $-2$,
and the odd-degree mixed term squares to zero. Therefore
\[
 \langle z^3,[X]\rangle=0,\qquad
 \langle z\,p_1(X),[X]\rangle=-2(1-w^2).                 \tag{D.MP8}\label{maps:eq:D.MP8}
\]
\ref{maps:src:D.1.10} identifies the entire target
$H_6(K(\mathbb Z/p,2);\mathbb Z)=\mathbb Z/p$ with this cube
evaluation. Hence the odd fundamental-class invariant is zero,
and the refined Pontryagin coefficient in the stated generator
is $(w^2-1)/2$. Theorem~5 of \cite[p.796]{manZhubr75} says
that these two invariants inject the entire spin-bordism group;
it applies to an arbitrary map to $K(\mathbb Z/p,2)$, so no
isomorphism $H_2(X)\cong\mathbb Z/p$ is being assumed here.
Up to an automorphism its class is exactly that of
$Z_{p,0}$ if $w^2=1$, and $Z_{p,1}$ otherwise.

\item \emph{Why these primary computations identify the whole manifold.}
We supply the splitting, rather than assuming that ordinary
characteristic classes exclude a mixed refinement. Put
$K_2=K(\mathbb Z/2,2)$ and $K_p=K(\mathbb Z/p,2)$.
All their positive integral homology is respectively two-primary
and $p$-primary. One direct verification uses the path-loop
fibration: after inverting the relevant prime, the positive
homology of the fibre $K(\mathbb Z/q,1)$ vanishes by the
cyclic-group resolution. Its Serre sequence then identifies
the base homology with that of the contractible path space.
Thus every positive integral homology element of the base is
killed by a power of $q$, as claimed.

\item All tensor and Tor products between positive homology of
$K_2$ and $K_p$ vanish. The integral reduced Kunneth sequence
therefore makes $K_2\wedge K_p$ acyclic, and equivalently
the inclusion
\[
                K_2\vee K_p\longrightarrow K_2\times K_p
                                                        \tag{D.MP9}\label{maps:eq:D.MP9}
\]
an integral homology equivalence. Both spaces are simply
connected CW complexes. The relative Hurewicz theorem,
applied to the first possible nonzero relative homotopy group,
and then Whitehead's theorem make \eqref{maps:eq:D.MP9} a homotopy equivalence.
Since $\Omega^{\mathrm{Spin}}_6(\mathrm{pt})=0$, the wedge
axiom for reduced spin bordism consequently gives
\[
 \Omega^{\mathrm{Spin}}_6(K(\mathbb Z/(2p),2))
 \cong\Omega^{\mathrm{Spin}}_6(K_2)
       \oplus\Omega^{\mathrm{Spin}}_6(K_p).               \tag{D.MP10}\label{maps:eq:D.MP10}
\]
Here the comparison to the product is the actual Chinese
remainder isomorphism on $H_2$, and the two components are
its two projections. Its automorphism group is the product
of the corresponding primary automorphism groups.

\item The connected sum on the right of \eqref{maps:eq:D.MP1} maps to this product
by collapsing its separating sphere and using the canonical
classifying maps of its two summands. Its $H_2$ map is the
direct-sum isomorphism. A single oriented spin one-handle
bordism compares the connected sum with the disjoint union,
so its two bordism components are precisely the two just
computed for $X$. Both manifolds are simply connected, spin,
have this same $H_2$, and have $b_3=0$. Zhubr's Theorem~3
\cite[p.795]{manZhubr75} now gives the required oriented
diffeomorphism. In particular the possible two-primary
refinement has been determined by \eqref{maps:eq:D.MP7}, not discarded.

\item Finally the actual Tate weights are $w=4$ modulo five for
the three $\epsilon=0$ rows, $w=3$ modulo five for the three
$\epsilon=1$ rows, and $w=4$ modulo seven for all four
remaining rows. This gives exactly \eqref{maps:eq:D.MP2}. Distinct choices
of $\epsilon$ have distinct complete primary bordism
invariants, and both occur among the constructed members.
\end{enumerate}
\end{proof}

The exact certificate \path{code/III/manifolds_mixed_prime_AH.py}
retains all twelve patterns and ten rows, reconstructs their
Tate markings, and verifies the primitive angular lattices,
$736$ semicharacter substitutions, every period difference,
integrality and scalar character, both unique fixed elliptic
curves, and all odd normal weights. It imposes no finite
multiplicity or twist cutoff. The topological prime splitting
is proved above; it is not inferred from those finite checks.


\begin{theorem}[the full nine-valued invariant in the $c=3$ sector]\label{maps:src:D.2.39}
\StatementRecord{maps:src:D.2.39}{proved}{Historical events: Claude Fable 5.2 (Anthropic), effort high, 16 November 2026: carried as background without fault found. GPT-6 Astra (OpenAI), ChatGPT Pro mode, effort Pro, S6\_09\_PAPERS, 17 November 2026: targeted proof check. Also checked in this writing. Pre-continuation writing events: GPT-6 Astra (OpenAI), ChatGPT Pro mode, effort Pro, S6\_10\_PAPERS, 18 November 2026: targeted proof check. Continuation metadata events: GPT-6 Astra (OpenAI), ChatGPT Pro mode, effort Pro, S6\_10\_PAPERS continuation, 19 November 2026: used and rechecked in this run; status-vocabulary reconciliation; metadata reconciliation only; no new mathematical check. \emph{Pre-continuation recorded proof scope.} GPT-6 Astra (OpenAI, ChatGPT Pro mode), effort Pro, 2026-11-17 UTC; Claude Fable 5.2 (Anthropic), effort high, 16 November 2026; also checked in this writing by GPT-6 Astra (OpenAI), ChatGPT Pro mode, effort Pro, 18 November 2026}{the hypotheses, notation and definitions stated here; \cite[Sections 3.11--3.14 and Theorem 5]{manZhubr75}; \cite[equation (2.8), p.267]{ffThomas56}; \cref{maps:src:D.2.26}; \cref{maps:src:D.2.36}; \cref{maps:eq:D.C26}; \cref{maps:eq:D.OP4}; \cref{maps:eq:D.OP5}; \cref{maps:eq:D.OP5a}; \cref{maps:eq:D.T35}; \cite[equations (2.3), (2.5), (2.7), p.267]{ffThomas56}; \cref{maps:eq:D.T34}; \cref{maps:eq:D.T37}; \cref{maps:eq:D.T39}; \cite[p.795]{manZhubr75}}


Use the primitive Mumford filling at every multiplicative site in one
of the following eight minimal patterns, retain the indicated additive
choice and the ordinary clutch $e\gamma$, and make no additional
regular logarithmic transformations:
\[
\begin{array}{c|r|r|l|l}
\text{OS/profile}&a&b&\text{zero positions}&\text{additive choice}\\\hline
56&5&3&2&\text{none}\\
56&10&3&3,4&\text{none}\\
56/II&5&3&0&II\text{ original or log}\\
56/III&1&6&0&III\text{ log}\\
56/III&5&6&3&III\text{ log}\\
56/III&10&3&3&III\text{ original}.
\end{array}                                                   \tag{D.T31}\label{maps:eq:D.T31}
\]
Zero positions are numbered from zero. For every finite multiplicity
and affine twist producing a simply connected manifold, the exact
oriented diffeomorphism type is
\[
 M_{3,0}=\bigl((S^2\times S^4)\setminus
          \operatorname{int}(S^2\times D^4)\bigr)
               \mathop{\cup}_{S^2\times S^3}(D^3\times S^3),     \tag{D.T32}\label{maps:eq:D.T32}
\]
where the embedded sphere represents three times the generator of
$H_2(S^2\times S^4;\mathbb Z)$ and has its oriented normal framing.
For example it is the graph of the degree-three map
$z\mapsto z^3$ on $\mathbb {CP}^1=S^2$ and an equatorial embedding
$S^2\hookrightarrow S^4$.
The complete Zhubr data are
\[
 H_2=H_3=\mathbb Z/3,\qquad b_3=0,\qquad
 \mathfrak p=0,\qquad
 \mu=0\in H_6(K(\mathbb Z/3,2);\mathbb Z)=\mathbb Z/9.          \tag{D.T33}\label{maps:eq:D.T33}
\]
In particular the assertion about $\mu$ determines its order-three
subgroup as well as its ordinary mod-three cubic reduction.

\end{theorem}


\begin{proof}
\begin{enumerate}[label=\arabic*.,leftmargin=*]
\item \emph{1. The invariant which must be computed.}
We first identify its complete target from the primary classification
and the actual Pontryagin operation. Zhubr's epimorphism in Section~3.11
and his free calculation in Section~3.12 show that
$\Omega^{\mathrm{Spin}}_6(K(\mathbb Z/3,2))$ is generated by the
images of $\epsilon,\zeta$, where
\[
 (P,h)(\epsilon)=(6,0),\qquad
 \zeta=[\mathbb {CP}^3,\text{hyperplane}].
\]
The image of $\epsilon$ has both invariants zero: its first is six
modulo three, and its second is the image of zero. The injectivity
in Zhubr's Theorem~5 therefore kills that image.
Lemma~3.14 gives $N(\mathbb Z/3)=0$, so the fundamental-class map
identifies this cyclic bordism group with
$H_6(K(\mathbb Z/3,2);\mathbb Z)$; the mod-two obstruction target
is zero because all positive homology is three-primary.
These facts use \cite[Sections 3.11--3.14 and Theorem 5]{manZhubr75}.

\item The cyclic generator has order dividing nine. On
$\mathbb {CP}^1\times\mathbb {CP}^2$ take the integral class
$3x+y$, with $x,y$ the two hyperplane classes. Its mod-three
reduction factors through $\mathbb {CP}^2$, and hence sends the
six-dimensional fundamental class to zero. Before reduction its
cube evaluates to nine. Naturality therefore gives nine times the
universal generator equal to zero.
For the opposite bound, Thomas's Pontryagin third power is a natural
operation
\[
 \mathcal P_3:H^2(-;\mathbb Z/3)\longrightarrow
                              H^6(-;\mathbb Z/9).
\]
Its value on the reduced hyperplane class of $\mathbb {CP}^3$
evaluates to a generator of $\mathbb Z/9$
\cite[equation (2.8), p.267]{ffThomas56}. Thus the cyclic group has
order exactly nine. Evaluation of $\mathcal P_3$ detects its entire
generator, and hence the whole fundamental-class invariant. This is
stronger than evaluating the ordinary cube modulo three.

\item \emph{2. The common-piece comparison and its circle geometry.}
\ref{maps:src:D.2.26} gives $H_2(X)=H_3(X)=\mathbb Z/3$ and $b_3=0$ for every
simply connected member in \eqref{maps:eq:D.T31}. In particular $H^2(X;\mathbb Z)=0$,
so $c_1(X)=0$ and $X$ is spin. Let
$u\in H^2(X;\mathbb F_3)$ be the generator normalized below.
All the geometric hypotheses of the doubled-line calculation in
\ref{maps:src:D.2.36} hold here. We record the exact inputs to ensure that no
prime-greater-than-three conclusion is being imported.

\item Remove the additive caps and one ordinary clutch disc to form $W$.
In all seven marked initial rows the actual multiplicative angular
columns generate $\ker\gamma^*$ with three unit Smith factors.
Consequently $\pi_1(W)=\mathbb Z_\gamma\times F_r$ is torsion-free.
The integral degree-one bridge calculation and surjective degree-two
cap boundary maps give $H_2(V)\twoheadrightarrow H_2(W)$, as in
\eqref{maps:eq:D.C26}. The pure fibre functional is
$\Omega=a\,u_1^*u_2^*+b\,\gamma^*\delta^*$, with $3\mid b$
and $3\nmid a$. Its finite normal coefficients are divisible by
$b/q_i=3$, its original additive coefficients are zero, and the
ordinary clutch coefficient is $be$. Thus its mod-three reduction
is the actual global degree-two generator, with zero marked swept
values. The doubled factors \eqref{maps:eq:D.OP4} and the exact integral supports
\eqref{maps:eq:D.OP5}, \eqref{maps:eq:D.OP5a} produce the corresponding line on $W$.

\item The scalar circle has precisely two strata whose stabilizer order is
divisible by three: the rational curve $C$ fixed by the whole circle,
and one elliptic double curve $E$ with stabilizer exactly $C_3$.
The latter is at the unique $I_3$ position. On its complement all
stabilizers are finite of order prime to three. This includes the
finite logarithmic caps: primitivity of their actual gamma numerators
makes the scalar circle free there, for every allowed order.
The original additive caps have free scalar action as well.
The seven exact Tate markings in the certificate verify the remaining
multiplicative assertions, including the primitive short period
$(3,aj)$ at $E$ and the two normal characters
$w,-w$, where $w=(aj)^{-1}=\pm1$ in $\mathbb F_3$.

\item The doubled auxiliary line has zero $C_3$ character at both $C$ and
$E$. Its degree on $C$ is a unit modulo three. Normalize $u$ by
$\langle u,[C]\rangle=1$. The torsion-free group of $W$ gives the
same simultaneous fixed-point comparison as in D.2.36: once an
equivariant lift has zero point character on $C$, it also has zero
point character on $E$. Indeed the two fixed-point copies of $BC_3$
represent the same order-three subgroup in
$\pi_1(W_{C_3})=\pi_1(W)\times C_3$. The group-cohomology difference
from half the auxiliary-line class has a common degree-two base
coefficient; the possible mixed degree-one term does not change
that coefficient. This comparison uses the complete normal columns,
including arbitrary finite twists.

\item \emph{3. The full-circle lift and the elliptic Borel space.}
The generator $u$ admits a lift
\[
                 \widehat u\in H^2_{S^1}(X;\mathbb F_3).
\]
In the Borel spectral sequence the only possible differentials out
of ordinary degree two have targets involving $H^1(X;\mathbb F_3)$
or $H^3(BS^1;\mathbb F_3)$; both are zero. Adding a multiple of the
base generator normalizes its restriction on the fixed curve to
\[
                      \widehat u|_{C_{S^1}}=[C]^* .           \tag{D.T34}\label{maps:eq:D.T34}
\]
Restriction $H^2(BS^1;\mathbb F_3)\to H^2(BC_3;\mathbb F_3)$
is nonzero. The point-character comparison in part~2 consequently
makes the coefficient of the degree-two $BC_3$ class on $E$ zero
as well.

\item The effective action $S^1/C_3$ on the elliptic curve $E$ is free.
Its quotient is a circle. Every principal circle bundle over a
circle is trivial, so equivariantly $E$ is the product of a free
$S^1/C_3$ orbit and a circle. Therefore
\[
 E_{S^1}\simeq BC_3\times S^1,
 \qquad
 \widehat u|_{E_{S^1}}=\eta\times\zeta,                    \tag{D.T35}\label{maps:eq:D.T35}
\]
where $\eta\in H^1(BC_3;\mathbb F_3)$ and
$\zeta\in H^1(S^1;\mathbb F_3)$, allowing a scalar in either
factor. Formula \eqref{maps:eq:D.T35} retains the possible nonzero mixed class;
it does not assert that $u$ descends to the coarse quotient $X/C_3$.
In fact the mixed term is generally essential.

\item For odd primes Thomas proves exact external-product multiplicativity
and vanishing on odd-dimensional classes, and independently proves
vanishing on suspensions
\cite[equations (2.3), (2.5), (2.7), p.267]{ffThomas56}.
These are statements with the full target coefficients, so
\[
 \mathcal P_3(\widehat u)|_{E_{S^1}}
   =\mathcal P_3(\eta)\times\mathcal P_3(\zeta)=0
                       \quad\text{in }H^6(E_{S^1};\mathbb Z/9).
                                                               \tag{D.T36}\label{maps:eq:D.T36}
\]
There is also no contribution on $C_{S^1}$: \eqref{maps:eq:D.T34} is pulled back
from the two-dimensional curve, whose degree-six cohomology is zero.
The familiar nonvanishing phenomena for the Pontryagin square of
suspended classes are irrelevant here; the cited formulas explicitly
assume an odd prime.

\item \emph{4. Why the full order-nine class vanishes globally.}
Put $R=\mathbb Z/9$. Choose disjoint invariant tubular neighborhoods
$N$ of $C\sqcup E$ and let $Y=X\setminus\operatorname{int}N$.
The circle acts locally freely on $Y$, with stabilizer orders prime
to three. The quotient $Y/S^1$ is a compact five-dimensional orbifold
with boundary. The Borel projection to this quotient, for both $Y$
and its boundary, has local cohomology fibres $BH$ with $H$ finite
of order invertible in $R$. Finite averaging gives
$H^q(BH;R)=0$ for $q>0$. The Leray comparison and the long exact
sequences of the pairs therefore give
\[
 H^k_{S^1}(Y,\partial Y;R)
             \cong H^k(Y/S^1,\partial(Y/S^1);R)=0
                            \qquad(k>5).                     \tag{D.T37}\label{maps:eq:D.T37}
\]
The dimension bound can be read from an orbifold triangulation,
or from a finite good orbifold atlas and its relative cellular
cohomology. No division by three occurs in this comparison.

\item Equivariant excision identifies the left side of \eqref{maps:eq:D.T37} with
$H^k_{S^1}(X,N;R)$. At $k=6$ it is zero. Thus restriction
\[
 H^6_{S^1}(X;R)\longrightarrow H^6_{S^1}(C\sqcup E;R)
                         \text{ is injective}.               \tag{D.T38}\label{maps:eq:D.T38}
\]
The vanishing of both restrictions proved in part~3 now implies
$\mathcal P_3(\widehat u)=0$ in the full integral-coefficient group
with coefficients $R$. Restricting to the ordinary fibre $X$ gives
\[
                        \langle\mathcal P_3(u),[X]\rangle=0
                                     \quad\text{in }\mathbb Z/9.
                                                               \tag{D.T39}\label{maps:eq:D.T39}
\]
Part~1 proves that this is the entire invariant $\mu$. Since
$N(\mathbb Z/3)=0$, the spin-bordism class itself is zero, and its
refined Pontryagin class is zero as well. This also agrees with
ordinary localization, which yields $p_1=0$ because $w^2=1$;
that latter congruence alone would not have proved \eqref{maps:eq:D.T39}.

\item \emph{5. The stated concrete smooth model.}
The graph embedding in \eqref{maps:eq:D.T32} represents three times the generator.
Its oriented real rank-four normal bundle is trivial: the ambient
manifold is spin and $w_2(TS^2)=0$, and $w_2$ classifies such bundles
on $S^2$. The framing is unique up to homotopy because
$\pi_2SO(4)=0$. The framed surgery preserves simple connectivity
and spin. Its integral sequence and duality give
$H_2=H_3=\mathbb Z/3$, with $b_3=0$.
The mod-three class from the first $S^2$ extends across the surgery
trace, since it evaluates to zero on the attaching sphere.
The original pair $(S^2\times S^4,u)$ spin-bounds
$S^2\times D^5$ over $K(\mathbb Z/3,2)$.
Thus the resulting model has zero spin-bordism class over this
Eilenberg--Mac Lane space. Zhubr's Theorem~3
\cite[p.795]{manZhubr75} identifies it, orientation-preservingly,
with every manifold in \eqref{maps:eq:D.T31}, since both classifying maps induce
isomorphisms on $H_2$ and both have $b_3=0$.
\end{enumerate}
\end{proof}

The program \path{code/III/manifolds_tate_markings.py} was executed with
residual three; the exact extraction output is
\path{code/fourfolds_D2_c3_tate.out}.
The independent certificate \path{code/III/fourfolds_D2_c3_AH.py}
checks all eight patterns on seven marked rows, all 27 adapted
monodromy identities, 14 full four-row angular and global-image
Hermite forms, all three finite normal rows with both arbitrary-twist
factors, and both original additive choices. It also checks every
doubled-line period identity, the unique three-primary isotropy
curve and its normal weights. Its retained output and full
records are \path{code/fourfolds_D2_c3_AH.out} and
\path{records/fourfolds_D2_c3_AH.json}. The order-nine computation
uses the exact primary operation and the relative Borel argument
above; it is not extrapolated from a finite mod-three test.


\begin{theorem}[shared odd factors in every weighted profile]\label{maps:src:D.2.40}
\StatementRecord{maps:src:D.2.40}{proved}{Historical events: Claude Fable 5.2 (Anthropic), effort high, 16 November 2026: carried as background without fault found. GPT-6 Astra (OpenAI), ChatGPT Pro mode, effort Pro, S6\_09\_PAPERS, 17 November 2026: targeted proof check. Also checked in this writing. Pre-continuation writing events: GPT-6 Astra (OpenAI), ChatGPT Pro mode, effort Pro, S6\_10\_PAPERS, 18 November 2026: targeted proof check. Continuation metadata events: GPT-6 Astra (OpenAI), ChatGPT Pro mode, effort Pro, S6\_10\_PAPERS continuation, 19 November 2026: used and rechecked in this run; status-vocabulary reconciliation; metadata reconciliation only; no new mathematical check. \emph{Pre-continuation recorded proof scope.} GPT-6 Astra (OpenAI, ChatGPT Pro mode), effort Pro, 2026-11-17 UTC; Claude Fable 5.2 (Anthropic), effort high, 16 November 2026; also checked in this writing by GPT-6 Astra (OpenAI), ChatGPT Pro mode, effort Pro, 18 November 2026}{the hypotheses, notation and definitions stated here; \cref{maps:src:D.2.37}; \cref{maps:eq:D.WP8}; \cref{maps:eq:D.WP9}; \cref{maps:eq:D.WP6}; \cref{maps:eq:D.OW5}; \cref{maps:src:D.2.29}; \cref{maps:eq:D.WP5}; \cref{maps:src:D.2.32}; \cref{maps:src:D.2.34}; \cref{maps:src:D.2.33}; \cref{maps:eq:D.WP11}; \cref{maps:eq:D.OW6}}


The uniform weighted algorithm of \ref{maps:src:D.2.37} extends from
$\gcd(a,d)=1$ to
\[
                     \gcd(a,d)\text{ odd}.                \tag{D.OW1}\label{maps:eq:D.OW1}
\]
In particular it covers \emph{every} positive translation
iterate of every row with $a$ odd. All its original,
Mumford, finite-good and regular logarithmic choices,
freeness tests and full Seifert groups remain valid.
Star caps are still original. The free ranks and Euler
number are still exactly \eqref{maps:eq:D.WP8}--\eqref{maps:eq:D.WP9}, with
$\chi=2d$.

The torsion must now be computed with the full fibre
presentation rather than one alternating functional.
Use the target
\[
                       \Lambda^2\Lambda\oplus K
\]
and include all columns
\[
 ((\Lambda^2A_i(d)-I)\eta,0)\qquad
                  (\eta\text{ a standard bivector}).     \tag{D.OW2}\label{maps:eq:D.OW2}
\]
For a finite or regular logarithm replace \eqref{maps:eq:D.WP6} by
\[
 (-\lambda\wedge v,\;J_iN_m\lambda)
       \quad(\lambda=e_1,e_2,\delta,\gamma),\qquad
                         (0,J_iv).                       \tag{D.OW3}\label{maps:eq:D.OW3}
\]
At every other cap include its full fixed-normal
columns $(0,J_iF_i)$ and the complete fibre kernel
\[
 {\cal K}_i(d)=
 \bigl\{\eta\in\Lambda^2\Lambda:
          (\Lambda^2S_d)\eta\in
          {\cal K}_i(1)\otimes\mathbb Q\bigr\}.             \tag{D.OW4}\label{maps:eq:D.OW4}
\]
Here $\mathcal K_i(1)$ is the recorded actual fibre kernel
for a nonzero-Chern original. In every other case its
rational span is
$\operatorname{im}(\Lambda^2A_i-I)$.
Formula \eqref{maps:eq:D.OW4} is evaluated as an ordinary saturated
integer kernel, not by discarding the nonunit Smith
factors of \eqref{maps:eq:D.OW2}.

Let $T$ be the torsion cokernel of this full matrix.
Then \eqref{maps:eq:D.WP9}, with this $T$, is the complete integral
homology for every datum in \eqref{maps:eq:D.OW1}.
The executable \path{code/III/manifolds_odd_weighted_profiles.py}
implements this extension and retains all shared-factor
torsion.

\end{theorem}


\begin{proof}
\begin{enumerate}[label=\arabic*.,leftmargin=*]
\item \emph{The integral involution splitting.}
Put $g=\gcd(a,d)$. The prime-local proof in \ref{maps:src:D.2.37}
shows, without assuming $g=1$, that every torsion
prime of
\[
 C=(\Lambda^2\Lambda)_{\{A_i(d)\}}
\]
divides $g$. Its free rank is one. Under \eqref{maps:eq:D.OW1}
we may write $C\simeq\mathbb Z\oplus C_{\rm tor}$,
with $C_{\rm tor}$ of odd order. The regular graph
sequence is an exact sequence of abelian groups
\[
                 0\longrightarrow C\longrightarrow
           H_2(V)\longrightarrow K\longrightarrow0.      \tag{D.OW5}\label{maps:eq:D.OW5}
\]
Fibre inversion is an actual involution of this
linear torus bundle. It acts by $+1$ on $C$ and
by $-1$ on $K$.

\item Let $A:H_2(V)\to\mathbb Z$ be evaluation by the
pullback of the marked weight-one class used in
\ref{maps:src:D.2.37}. It is invariant under inversion: this is
true on the old fibre and on all its generating
fixed sweeps, and therefore remains true on the
cover. Its restriction to $C$ is
\[
 g\,\omega',\qquad
 \omega'=(a/g)e_1^*e_2^*-(db/g)\delta^*\gamma^*,
\]
where $\omega'$ is the primitive free functional.
For any $x\in H_2(V)$, $x+\iota x$ belongs to $C$.
Hence
$2A(x)=A(x+\iota x)$ is divisible by $g$.
Since $g$ is odd, $A/g$ is an integral homomorphism.
Call it $B$.

\item We spell out the resulting splitting to avoid a
formality or invariant-descent assumption. Given
$k\in K$, choose a lift $x$ with $B(x)=0$, subtracting
a fibre generator if necessary. Then
$x+\iota x\in C_{\rm tor}$. Multiplication by two
is an automorphism of this finite group, so
\[
                  y=x-\tfrac12(x+\iota x)                 \tag{D.OW6}\label{maps:eq:D.OW6}
\]
is a well-defined integral class with
$\iota y=-y$ and image $k$. If the lift $x$ is
changed by $c\in C_{\rm tor}$, the half term
changes by $c$, so $y$ is unchanged.
This constructs an additive section of \eqref{maps:eq:D.OW5}.
It is the unique section with $B=0$ and
anti-invariant image. Thus
$H_2(V)=C\oplus K$ with an actual integral marking,
including the odd torsion of $C$.

\item Every literal fixed-loop peripheral torus is
anti-invariant, since inversion reverses its
fibre circle and fixes the base circle. It
therefore has exactly its graph coordinate and
zero fibre coordinate in this splitting.
This remains true for the norm sweep of any
covering-torus loop: take the pushforward of
$\lambda\times S^1$ from the $m$-fold mapping-torus
cover, where $A^m=I$. It is anti-invariant and
has graph coordinate $N_m\lambda$.
This is the step that controls integral torsion
which a vertical alternating form alone cannot see.

\item \emph{The full finite and regular normal columns.}
On the $m$-fold boundary cover the normal meridian
is $m\theta-v$. The shear
$(x,t)\mapsto(x-tv,t)$ sends the product torus
of $\lambda$ and the covering base circle to
that norm sweep minus the regular torus
$\lambda\wedge v$. Pushing it to the ordinary
mapping torus and then to $V$, the preceding
anti-invariant description gives exactly the
first four columns of \eqref{maps:eq:D.OW3}.
There is no unknown torsion component in the
norm sweep.

\item For the affine meridian, use its commuting lifts
$\theta$ and $\mu=\theta^m v^{-1}$ in the boundary
group. Their torus is the fixed torus of $v$
and $\theta$, with the normal orientation
convention of \ref{maps:src:D.2.29}. Hence its image is
$(0,J_iv)$, the last column in \eqref{maps:eq:D.OW3}.
These loops generate the complete normal Thom
image, including torsion loops in the central
finite quotient. The proof of the general
freeness test \eqref{maps:eq:D.WP5} is unchanged.

\item \emph{Every nonfinite fibre kernel and its marked normal.}
All the unchanged caps in this construction have
torsion-free $H_2$. For the rank-two cusp this is
the explicit normalization calculation. For a
rank-one gamma cover the normalized components
are $S^2$ bundles of degree zero over a covered
two-torus, and consecutive double curves are
identified by translations; straightening those
translations gives the same free homology as a
cycle of spheres times $T^2$. Original
zero-degree additive caps have the permutation
mapping-torus homology of a Kodaira tree times
$S^1$, which is free. The nonzero-Chern originals
have the primitive Gysin presentations of
\ref{maps:src:D.2.32}--\ref{maps:src:D.2.33}, also free.

\item The covering square on regular fibres multiplies
bivectors by $\Lambda^2S_d$. The needed ranks come
from the actual caps, as follows. At the rank-two
zero the short-period map has Smith factors $(1,d)$.
In that basis its exterior-square difference has
primitive image of rank two, so the boundary has
$H_2=\mathbb Z^6$. The normalization calculation of
\ref{maps:src:D.2.34} gives $H^3(N)=\mathbb Z^2$. The two angular
normal tori are a primitive sublattice of boundary
homology and bound by their section-disc sweeps.
The pair sequence bounds its whole kernel by rank
two; hence they are exactly that kernel. In
particular the regular-fibre coinvariant of rank
four injects into the cap. This proves that the
rational fibre kernel is precisely the monodromy
image, without a local invariant-cycle assumption.

\item For the rank-one cap, the explicit cycle of
$\rho$ spheres times $T^2$ has four independent
nearby-fibre classes: the sum of the component
spheres, the surviving nodal-circle class times
each transverse circle, and the transverse
$T^2$ area. This is the full four-dimensional
fibre coinvariant, so again its kernel is exactly
the monodromy image over $\mathbb Q$. For an
original zero-degree additive cap, the permutation
mapping torus of the Kodaira tree gives the two
independent images: the nonzero multiplicity-weighted
whole-fibre class in the component-orbit space,
and the transverse $\delta\gamma$ area. These
are the full two-dimensional fibre coinvariant,
also for the original odd-star fibres. The
zero-degree multiplicative original has the same
four explicit classes as the nodal-cycle model.

\item Thus in every zero-degree case the rational
fibre kernel is exactly the inverse image of
the old one. For a nonzero-Chern original the
actual area and principal-bundle formulas give
the same statement directly. Since the target
$H_2$ is torsion free, the integral kernel is the
saturation of that rational kernel. This is
precisely \eqref{maps:eq:D.OW4}.

\item The full fixed-normal projections lift the old
bounding section-disc sweeps, as in \ref{maps:src:D.2.37}.
There is one previously isolated correction:
at a nonzero-Chern original multiplicative cap,
the vanishing-loop normal can differ from a
kernel cycle by an integral multiple of
$u_2\wedge\delta$. In \ref{maps:src:D.2.33} only its zero
alternating-form value was needed. Here it
vanishes in the full global fibre coinvariant.
Indeed, for all $172$ rows the elliptic matrices
$B_i$ satisfy the exact unit Smith certificate
\[
       \mathbb Z^2/\sum_i(B_i-I)\mathbb Z^2=0.             \tag{D.OW7}\label{maps:eq:D.OW7}
\]
These $172$ matrices and their unit factors
are retained in
\path{records/manifolds_elliptic_coinvariants.json}.
Since every $A_i(d)$ fixes $\delta$,
\[
 (\Lambda^2A_i(d)-I)(\delta\wedge z)
                 =\delta\wedge(B_i-I)z
       \quad(z\in\mathbb Z^2).
\]
Consequently $\delta\wedge e_1$ and
$\delta\wedge e_2$ are zero integrally in $C$,
for every weight. This kills the correction
without any assumption on its integer value.

\item Subtracting these complete marked normals
from a boundary-kernel element leaves a
regular-fibre element. Its full image is
exactly \eqref{maps:eq:D.OW4}. Thus \eqref{maps:eq:D.OW2}--\eqref{maps:eq:D.OW4}
give the entire integral global kernel;
in particular a primitive fibre relation
that was only a torsion class in the boundary
has not been omitted.

\item \emph{Global homology and the remaining scope.}
The relative pair argument \eqref{maps:eq:D.WP11} does not
require the regular coinvariant to be free.
Finite and regular caps still have free
relative degree-two groups, and their
nonzero base-meridian coefficients make the
first relative map injective. Their full
degree-three images are now \eqref{maps:eq:D.OW3}.
The remaining caps still have free
$H^4$, so the final kernel in \eqref{maps:eq:D.WP11} is
free and the abelian-group extension splits.
The top component counts and rational ranks
are unchanged. Therefore \eqref{maps:eq:D.WP8} and the
entire duality table hold, with exactly the
full torsion cokernel computed here.

\item The vertical normal closures are also unchanged,
so the same all-order Seifert presentation
computes the full fundamental group.
The argument has not used a bound on weights
or logarithmic parameters. It uses oddness
only in dividing $g$ in the invariant functional
and in the unique half of \eqref{maps:eq:D.OW6}. When both
$a$ and $d$ are even, a two-primary ambiguity
can remain in that marking; this theorem
does not discard it.
\end{enumerate}
\end{proof}

The exact implementation computes the saturated
preimages \eqref{maps:eq:D.OW4} before forming any Smith quotient.
Its overlap checks compare this full fibre-and-graph
matrix with \ref{maps:src:D.2.37}'s independent scalar-functional
matrix whenever $g=1$. Shared odd-factor cases
retain the nontrivial regular fibre coinvariant
and all its relations. The printed certificates
and explicit sample matrices are retained in
\path{records/manifolds_odd_weighted_profiles.json}.

The shared-even-factor case outside the odd-iterate statement is included in \ref{maps:src:D.2.41}.

\begin{theorem}[every positive weight by integral bar chains]\label{maps:src:D.2.41}
\StatementRecord{maps:src:D.2.41}{proved}{Historical events: Claude Fable 5.2 (Anthropic), effort high, 16 November 2026: carried as background without fault found. GPT-6 Astra (OpenAI), ChatGPT Pro mode, effort Pro, S6\_09\_PAPERS, 17 November 2026: targeted proof check. Also checked in this writing. Pre-continuation writing events: GPT-6 Astra (OpenAI), ChatGPT Pro mode, effort Pro, S6\_10\_PAPERS, 18 November 2026: targeted proof check. Continuation metadata events: GPT-6 Astra (OpenAI), ChatGPT Pro mode, effort Pro, S6\_10\_PAPERS continuation, 19 November 2026: used and rechecked in this run; status-vocabulary reconciliation; metadata reconciliation only; no new mathematical check. \emph{Pre-continuation recorded proof scope.} GPT-6 Astra (OpenAI, ChatGPT Pro mode), effort Pro, 2026-11-17 UTC; Claude Fable 5.2 (Anthropic), effort high, 16 November 2026; also checked in this writing by GPT-6 Astra (OpenAI), ChatGPT Pro mode, effort Pro, 18 November 2026}{the hypotheses, notation and definitions stated here; \cref{maps:src:D.2.37}; \cref{maps:src:D.2.40}; \cref{maps:eq:D.WP8}; \cref{maps:eq:D.WP9}; \cref{maps:eq:D.OW2}; \cref{maps:eq:D.OW4}; \cref{maps:eq:D.WP5}; \cite[Chapter I, Section 5]{manBrown82}; \cref{maps:eq:D.BC9}; \cref{maps:eq:D.OW3}; \cref{maps:eq:D.BC8}; \cref{maps:eq:D.BC10}; \cref{maps:eq:D.BC14}; \cref{maps:eq:D.BC13}; \cref{maps:src:D.2.29}; \cref{maps:eq:D.OW7}; \cref{maps:eq:D.WP11}}


The restriction on $\gcd(a,d)$ in \ref{maps:src:D.2.37} and \ref{maps:src:D.2.40} can
be removed. Every positive iterate $Q\mapsto dQ$ of
every one of the $172$ natural initial rows has an
exact integral homology algorithm, for all $495$
original/Mumford presentation patterns, every free
finite-good affine logarithm, and arbitrary regular
logarithms and clutching integer. Star caps remain
original. The full Seifert fundamental group and
the free ranks \eqref{maps:eq:D.WP8}--\eqref{maps:eq:D.WP9} are unchanged.

Here is the required all-parity correction to the
full matrix \eqref{maps:eq:D.OW2}--\eqref{maps:eq:D.OW4}. Index the boundary
matrices as $A_0,\ldots,A_r$, with
$A_0\cdots A_r=I$, and use the first $r$ base loops
as independent generators. All $A_i$ below mean
the actual weighted matrices $A_i(d)$.
For $u,v\in\Lambda=\mathbb Z^4$ set
\[
 p(u,v)=\sum_{\alpha<\beta}u_\alpha v_\beta\,
                  e_\alpha\wedge e_\beta,\qquad
 D_A(u,v)=p(Au,Av)-(\Lambda^2A)p(u,v).
                                                        \tag{D.BC1}\label{maps:eq:D.BC1}
\]
The bilinear map $D_A$ is symmetric. Define its
integral quadratic primitive by
\[
 f_A(u)=-\sum_\alpha D_A(e_\alpha,e_\alpha)
                       \binom{u_\alpha}{2}
        -\sum_{\alpha<\beta}D_A(e_\alpha,e_\beta)
                       u_\alpha u_\beta .
                                                        \tag{D.BC2}\label{maps:eq:D.BC2}
\]
Negative coordinates use the same integral polynomial
$\binom n2=n(n-1)/2$.
Define the peripheral correction
\[
 \begin{split}
 c_i(u)&=f_{A_i}(u) \quad (i<r),\\
 c_r(u)&=-\sum_{j<r}
       f_{A_j}(A_{j+1}\cdots A_{r-1}A_r u).
 \end{split}                                            \tag{D.BC3}\label{maps:eq:D.BC3}
\]
The transported reversed word in the second line
is part of the formula.

Keep every fibre relation \eqref{maps:eq:D.OW2} and every full
saturated cap kernel \eqref{maps:eq:D.OW4}. Replace the five
finite or regular logarithmic normal columns by
\[
 \begin{gathered}
 \left(\sum_{\nu=0}^{m-1}c_i(A_i^\nu\lambda)
                  -\lambda\wedge v,\quad J_iN_m\lambda\right)
       \quad(\lambda=e_1,e_2,\delta,\gamma),\\
                       (c_i(v),J_iv).
 \end{gathered}                                          \tag{D.BC4}\label{maps:eq:D.BC4}
\]
At every other cap the fixed-normal columns are
\[
                         (c_i(v),J_iv),
                         \qquad v\text{ in a basis of }\operatorname{Fix} A_i.
                                                        \tag{D.BC5}\label{maps:eq:D.BC5}
\]
The finite freeness test is still exactly \eqref{maps:eq:D.WP5}.
If $A_i$ has order $h\mid m$, the sum in \eqref{maps:eq:D.BC4}
is $(m/h)\sum_{\nu=0}^{h-1}c_i(A_i^\nu\lambda)$;
thus this is a finite algorithm for unbounded orders.
Let $T$ be the torsion cokernel of this full matrix.
Then the complete integral homology is \eqref{maps:eq:D.WP9}
with this $T$, for every $d>0$, including all
shared even factors.

\end{theorem}


\begin{proof}
\begin{enumerate}[label=\arabic*.,leftmargin=*]
\item \emph{An integral chain comparison with an explicit marking.}
Use the normalized bar chains of $\Lambda$ with
trivial coefficients. Their low-dimensional
boundaries are
\[
 \partial[u]=0,\qquad
 \partial[u|v]=[v]-[u+v]+[u].
                                                        \tag{D.BC6}\label{maps:eq:D.BC6}
\]
The bar complex is a functorial chain model of
$B\Lambda$; see the standard resolution
\cite[Chapter I, Section 5]{manBrown82}.
We give all the comparison formulas used here.
The maps
\[
             p_1([u])=u,\qquad p_2([u|v])=p(u,v)
                                                        \tag{D.BC7}\label{maps:eq:D.BC7}
\]
annihilate boundaries in the appropriate degree.
In degree two the cycle $[u|v]-[v|u]$ maps to
$u\wedge v$. The coordinate subtori give an
integral basis, so $p_2$ induces exactly
$H_2(T^4;\mathbb Z)\simeq\Lambda^2\Lambda$.

\item Antisymmetrizing \eqref{maps:eq:D.BC1} gives zero, proving
the symmetry of $D_A$. The binomial identity
$\binom{x+y}{2}=\binom{x}{2}+\binom{y}{2}+xy$
then gives the exact integral identity
\[
 f_A(u)+f_A(v)-f_A(u+v)
             =p(Au,Av)-(\Lambda^2A)p(u,v).
                                                        \tag{D.BC8}\label{maps:eq:D.BC8}
\]
Equivalently, extending $f_A([u])=f_A(u)$
linearly on bar one-chains,
$f_A\partial=p_2A-(\Lambda^2A)p_2$ on bar
two-chains.

\item The regular piece retracts to its torus bundle
over a rose with $r$ petals. Its integral bar
mapping-cone complex has
\[
 {\cal C}_n=B_n(\Lambda)\oplus
                   \bigoplus_{i<r}B_{n-1}(\Lambda),
 \qquad
 \partial(c;g_i)=
 \left(\partial c+\sum_{i<r}(A_i-I)g_i;\
                                     -\partial g_i\right).
                                                        \tag{D.BC9}\label{maps:eq:D.BC9}
\]
This is a model for the actual linear torus bundle,
not an assertion of equivariant formality. Its
universal cover is a torus-universal-cover bundle
over a tree, hence contractible. The bar model
has the same semidirect-product fundamental group,
and their maps on every peripheral torus agree
when their marked loops agree.

\item The sign in \eqref{maps:eq:D.BC9} uses the existing inverse
endpoint convention
$(x,1)\sim(A_i^{-1}x,0)$. Label a cylinder by its
terminal fibre chain $g$; its initial label is
$A_i g$. With fibre-first product orientation
its boundary is $(A_i-I)g$. This is the same cut
orientation used in \eqref{maps:eq:D.OW3}.

\item In degree two define
\[
 P_2(c;g_i)=
 \left(p_2c+\sum_{i<r}f_{A_i}(g_i);\
                                  (p_1g_i)_{i<r}\right).
                                                        \tag{D.BC10}\label{maps:eq:D.BC10}
\]
The target is the small complex whose degree-two
homology is
\[
 C\oplus K,\qquad
 C=\Lambda^2\Lambda/
          \sum_{i<r}(\Lambda^2A_i-I)\Lambda^2\Lambda,
 \quad
 K=\ker[A_0-I\ \cdots\ A_{r-1}-I].
                                                        \tag{D.BC11}\label{maps:eq:D.BC11}
\]
No saturation or torsion quotient is taken in $C$.
Formula \eqref{maps:eq:D.BC8} proves that $P_2$ carries a
degree-three boundary to a small-complex boundary:
its fibre coordinate is
\[
 \sum_{i<r}\bigl(p_2(A_i-I)h_i-f_{A_i}\partial h_i\bigr)
       =\sum_{i<r}(\Lambda^2A_i-I)p_2h_i .
                                                        \tag{D.BC12}\label{maps:eq:D.BC12}
\]
The contribution $p_2\partial c_3$ is zero.
Similarly $P_2$ commutes with the differential
to degree one. Thus it induces a homomorphism
on $H_2$.

\item Both cone homologies have the short exact
sequence with left term the degree-two fibre
coinvariant and right term $K$. On the left
$P_2$ is the integral isomorphism induced by
$p_2$, and on the right it is the identity
induced by $p_1$. The short five lemma proves
that \eqref{maps:eq:D.BC10} is an integral isomorphism
$H_2(V)\simeq C\oplus K$. This supplies the
missing two-primary marking directly. Neither
division by two nor descent of an invariant
differential form is used.

\item \emph{Fixed sweeps, norm sweeps, and the last boundary.}
If $A_i v=v$ and $i<r$, the literal fixed-loop
torus is the bar cycle with $g_i=[v]$ and
all other entries zero. Formula \eqref{maps:eq:D.BC10} sends
it to $(f_{A_i}(v),J_iv)$.
For an arbitrary loop $\lambda$ with
$A_i^m\lambda=\lambda$, the pushforward of
its torus in the $m$-fold mapping-torus cover
has edge chain
\[
                        \sum_{\nu=0}^{m-1}[A_i^\nu\lambda].
                                                        \tag{D.BC13}\label{maps:eq:D.BC13}
\]
Its boundary telescopes to zero. Hence its
coordinates are the norm term in \eqref{maps:eq:D.BC4}
before subtracting $\lambda\wedge v$.

\item For the last boundary the base word is
$t_r=t_{r-1}^{-1}\cdots t_0^{-1}$.
The local bar one-chain $[u]$ maps to the
independent edge chains
\[
         -[A_{j+1}\cdots A_{r-1}A_ru],
                            \qquad 0\le j<r .
                                                        \tag{D.BC14}\label{maps:eq:D.BC14}
\]
This identity is at the chain level: applying
$\sum(A_j-I)$ makes the bars telescope to
$[A_ru]-[u]$, with no leftover fibre two-chain.
For a fixed vector, its linear part is exactly
the previous marked $J_rv$. Applying
\eqref{maps:eq:D.BC10} gives the second line of \eqref{maps:eq:D.BC3}.
Applying \eqref{maps:eq:D.BC14} to every summand of
\eqref{maps:eq:D.BC13} gives the full transported last-boundary
norm correction. In particular no reordering
of the last peripheral segments has been
silently discarded.

\item \emph{The actual cap maps.}
For a finite or regular logarithm the boundary
group has its normal meridian
$\mu=t_i^m v^{-1}$, with $A_i^m=I$ and
$A_i v=v$. The loops $\lambda,t_i^m,v$ lie
in one abelian subgroup. Bilinearity of its
torus classes therefore gives the exact
integral relation
\[
              [\lambda,\mu]=[\lambda,t_i^m]-[\lambda,v].
                                                        \tag{D.BC15}\label{maps:eq:D.BC15}
\]
The first term is \eqref{maps:eq:D.BC13}, and the second is
the fibre bivector $\lambda\wedge v$.
This proves the first four columns of \eqref{maps:eq:D.BC4}
even when their images have two-primary
torsion. The affine meridian generator gives
$[t_i,\mu]=[v,t_i]$, the fixed sweep already
computed, proving the fifth column. These
are exactly the full normal Thom generators
from \ref{maps:src:D.2.29}--\ref{maps:src:D.2.30}, including the torsion
normal-line cases.

\item Every fixed-normal torus at an unchanged cap
is still the same literal peripheral sweep,
so its image is \eqref{maps:eq:D.BC5}. The fibre kernels
are exactly the saturated lattices proved
from the explicit local models in \ref{maps:src:D.2.40}.
The possible original multiplicative
correction $\delta\wedge u_2$ is zero in
the full global fibre coinvariant by the
$172$ unit elliptic certificates \eqref{maps:eq:D.OW7}.
Thus the listed columns give the entire
marked kernel image at every allowed cap.

\item Finally the relative pair sequence \eqref{maps:eq:D.WP11}
has free final quotient for every weight;
its proof did not use the parity of $g$.
Its ranks depend only on the actual component
counts already proved in \ref{maps:src:D.2.37}. It follows
that the corrected full matrix computes
exactly the middle torsion, while the Seifert
group, free ranks and duality table stay
unchanged. All formulas are integral
polynomials, binomial polynomials, saturated
integer kernels and finite Smith operations.
They impose no upper bound on $d$, the
logarithmic orders or affine coordinates.
\end{enumerate}
\end{proof}

The executable
\path{code/III/manifolds_all_weighted_profiles.py}
implements \eqref{maps:eq:D.BC1}--\eqref{maps:eq:D.BC5}, records the full
quadratic coefficient vectors, and compares
its results with the independent inversion
method whenever $\gcd(a,d)$ is odd.
This theorem concerns translation iterates
of the finite natural initial catalogue.
It does not assert that all broader initial
linearizations or star logarithmic caps are
such iterates, and it does not infer an
uncomputed smooth classification from the
homology.


\begin{theorem}[the mixed six-, fifteen-, and thirty-torsion types]\label{maps:src:D.2.42}
\StatementRecord{maps:src:D.2.42}{proved}{Historical events: Claude Fable 5.2 (Anthropic), effort high, 16 November 2026: carried as background without fault found. GPT-6 Astra (OpenAI), ChatGPT Pro mode, effort Pro, S6\_09\_PAPERS, 17 November 2026: targeted proof check. Also checked in this writing. Pre-continuation writing events: GPT-6 Astra (OpenAI), ChatGPT Pro mode, effort Pro, S6\_10\_PAPERS, 18 November 2026: targeted proof check. Continuation metadata events: GPT-6 Astra (OpenAI), ChatGPT Pro mode, effort Pro, S6\_10\_PAPERS continuation, 19 November 2026: used and rechecked in this run; status-vocabulary reconciliation; metadata reconciliation only; no new mathematical check. \emph{Pre-continuation recorded proof scope.} GPT-6 Astra (OpenAI, ChatGPT Pro mode), effort Pro, 2026-11-17 UTC; Claude Fable 5.2 (Anthropic), effort high, 16 November 2026; also checked in this writing by GPT-6 Astra (OpenAI), ChatGPT Pro mode, effort Pro, 18 November 2026}{the hypotheses, notation and definitions stated here; \cref{maps:src:D.2.39}; \cref{maps:src:D.2.36}; \cref{maps:src:D.2.26}; \cref{maps:src:D.2.35}; \cref{maps:eq:D.OP4}; \cref{maps:eq:D.MT2}; \cite[equations (2.3), (2.5), p.267]{ffThomas56}; \cref{maps:eq:D.T37}; \cref{maps:eq:D.T38}; \cref{maps:eq:D.C27}; \cref{maps:src:D.2.8}; \cref{maps:src:D.1.10}; \cref{maps:eq:D.MP9}; \cref{maps:eq:D.MP10}; \cite[p.795]{manZhubr75}}


Consider the fourteen minimal patterns in \eqref{maps:eq:D.MT1}, with the primitive
Mumford cap at every multiplicative site, the indicated additive
choice, and no additional regular logarithmic transformation.
Allow every finite multiplicity, affine twist and ordinary clutch
$e\gamma$ producing a simply connected manifold. Let
$M_{3,0}$ be the explicitly framed surgery of \ref{maps:src:D.2.39}, and
$Z_{5,\epsilon}$ the oriented type of \ref{maps:src:D.2.36}. Define
\[
\begin{aligned}
 Y_6&=M_{\mathrm{conic}}\# M_{3,0},\\
 Y_{15,\epsilon}&=M_{3,0}\# Z_{5,\epsilon},\\
 Y_{30}&=M_{\mathrm{conic}}\# M_{3,0}\# Z_{5,0}.
\end{aligned}
\]
The orientation-preserving diffeomorphism type is given by the
following table; zero positions have the source's zero-based indexing.
\[
\begin{array}{c|r|r|l|l|l}
\text{OS/profile}&a&b&\text{zero positions}&\text{additive choice}&\text{type}\\\hline
56&1&6&0&\text{none}&Y_6\\
56&5&6&3,4&\text{none}&Y_6\\
56/II&1&6&2&II\text{ original or log}&Y_6\\
56&1&15&2&\text{none}&Y_{15,0}\\
56&2&15&3,4&\text{none}&Y_{15,1}\\
56/II&1&15&0&II\text{ original or log}&Y_{15,0}\\
56/III&1&30&3&III\text{ log}&Y_{15,0}\\
56/III&2&15&3&III\text{ original}&Y_{15,1}\\
56&1&30&3,4&\text{none}&Y_{30}.
\end{array}                                               \tag{D.MT1}\label{maps:eq:D.MT1}
\]
Thus the five $c=6$ patterns give one oriented type, the seven
$c=15$ patterns give two, and the two $c=30$ patterns give one.

\end{theorem}


\begin{proof}
\begin{enumerate}[label=\arabic*.,leftmargin=*]
\item \emph{The geometric input for every primary calculation.}
\ref{maps:src:D.2.26} gives $H_2=H_3=\mathbb Z/c$ and $b_3=0$ in these
simply connected cases. As $H^2(X;\mathbb Z)=0$, each is spin.
There are twelve distinct marked initial rows. Write
$\Omega=a\,u^*w^*+b\,\gamma^*\delta^*$ in their actual markings.
The exact certificate below verifies its invariance and verifies
that both the multiplicative angular columns and the full images
of $A_i-I$ have integral Hermite lattice $\ker\gamma^*$.
Thus, after omitting the additive caps and one ordinary disc,
the common piece $W$ has torsion-free group
$\mathbb Z_\gamma\times F_r$. The primitive Mumford boundary
surjections and the independent remaining peripheral loops give
$H_2(V)\twoheadrightarrow H_2(W)$, with free first homology.
Consequently the degree-two comparisons on $W$ used in
\ref{maps:src:D.2.35}--\ref{maps:src:D.2.36} apply to the actual marked classes here as well.

\item The three possible finite patterns have normal fibre coefficients
$6$, $15$, and $30/2=15$, respectively. All are divisible by
the stated residual $c$, as is the second normal coefficient,
a multiple of $b$. Every allowed original additive choice has
zero coefficient, and the clutch coefficient is $be$.
Hence the primitive fibre functional with zero swept values
descends modulo every prime dividing $c$, independently of
the allowed finite order and affine twist.

\item We record the local inputs to this comparison. The actual Tate
periods are $H=(n,ar),F=(-br,\rho)$; at the zero
$n\rho+ab r^2=1$. The doubled symmetric line \eqref{maps:eq:D.OP4} has
integral support
\[
          \psi=-a\rho x^2-2ab rxy+bn y^2+ax,
                                                        \tag{D.MT2}\label{maps:eq:D.MT2}
\]
with the two exact period differences $-2ax$ and
$2by+b\rho-ab r$. For an odd prime $p\mid c$, all its
vertical slopes are divisible by $p$, and half its signed
wall degree is the unit $-a(\rho-br)$ modulo $p$.
At a rank-one site write $H=\alpha u_0,F=\beta u_0$, with
$u_0=(u_1,u_2)$ primitive and $\gcd(\alpha,\beta)=1$.
Its integral doubled support is
$-(a u_1/\alpha)j(j-\alpha)$; the extracted identities prove
$\alpha\mid a u_1$ and give both prescribed period differences.
Thus the auxiliary line extends on all the specified caps,
with the actual boundary marking.

\item For each $p\mid c$ there is exactly one additional
$C_p$-fixed elliptic double curve $E_p$, at an $I_p$ site.
Its primitive short period is $(p,ar)$ and its doubled support
is $-ap j(j-1)$. The line therefore has zero $C_p$ character
on both the rational fixed curve $C$ and $E_p$, for odd $p$.
The normal weights at $E_p$ are $w,-w$, where
$w=(ar)^{-1}\pmod p$. All these statements, including the
uniqueness of each fixed elliptic curve, are verified using
the actual primitive period columns in the twelve rows.
The finite caps have free scalar action by primitivity of
their gamma numerators; the original additive caps have
free scalar action. No extra prime-isotropy stratum is omitted.

\item \emph{The full three-primary class.}
For every row take $u_3\in H^2(X;\mathbb F_3)$ normalized
to have degree one on $C$. It lifts to $S^1$-equivariant
cohomology because $H^1(X;\mathbb F_3)=0$ and
$H^3(BS^1;\mathbb F_3)=0$. Normalize its point character
on $C$ to zero. The torsion-free group of $W$ and the
zero characters of \eqref{maps:eq:D.MT2} then make its $C_3$ point
character on $E_3$ zero as well, exactly as in \ref{maps:src:D.2.39}.
The restrictions of this lift are a class pulled back from
$C$ and an external product of degree-one classes on
$BC_3\times S^1\simeq(E_3)_{S^1}$.

\item Thomas's exact odd-prime external-product formula and
odd-degree vanishing make the Pontryagin third power of
both restrictions zero in coefficients $\mathbb Z/9$
\cite[equations (2.3), (2.5), p.267]{ffThomas56}.
The complementary locally free circle action has all
stabilizers of order prime to three. Its Borel pair has
the cohomology of a five-dimensional orbifold pair with
these coefficients, so its relative degree-six group
is zero. Thus restriction to $C\sqcup E_3$ is injective
in equivariant degree six, as proved in \eqref{maps:eq:D.T37}--\eqref{maps:eq:D.T38}.
The full Pontryagin third power on $X$ is consequently zero.

\item The target calculation in part~1 of \ref{maps:src:D.2.39} and Zhubr's
Theorem~5 now show that the entire spin-bordism image of
$X\to K(\mathbb Z/3,2)$ is zero. This use concerns an
arbitrary classifying map: it does not assume that
$H_2(X)$ itself equals $\mathbb Z/3$. It determines the
hidden subgroup of order three in $\mathbb Z/9$, in
addition to the ordinary mod-three cube. The primary
class is therefore exactly the class of $M_{3,0}$.

\item \emph{The two- and five-primary classes.}
When $c$ is even, $a$ is odd and $b/2$ is odd in every
row of \eqref{maps:eq:D.MT1}. The unsquared line \eqref{maps:eq:D.C27} has the
monodromy-invariant semicharacter $a(u+w+uw)$.
Half of \eqref{maps:eq:D.MT2} and half of its rank-one support are
integral. Their scalar involution characters are one
on $C$ and zero on $E_2$, and their degree on $C$ is odd.
The common-piece comparison and the same simultaneous
Borel normalization as in \ref{maps:src:D.2.35} thus give
$\langle u_2^3,[X]_2\rangle=1$. Its entire two-primary
bordism class is the conic-surgery class of \ref{maps:src:D.2.8}.

\item For $c=15$ or $30$, the five-primary fixed curves have
the zero characters supplied by \eqref{maps:eq:D.MT2}. Normalize
$u_5[C]=1$. The actual localization of \ref{maps:src:D.2.36} gives
\[
 \langle u_5^3,[X]\rangle=0,\qquad
 \mathfrak p_5=(w^2-1)/2\pmod5.                           \tag{D.MT3}\label{maps:eq:D.MT3}
\]
The first equality detects the entire five-primary
fundamental class by \ref{maps:src:D.1.10}, and the second is the
complete refined Pontryagin class since four is a
unit modulo five. The actual weight is $w=4$ for
all $a=1$ rows, and $w=2$ for the $a=2$ rows.
Thus their complete bordism orbits are respectively
those of $Z_{5,0}$ and $Z_{5,1}$, by Zhubr's Theorem~5.

\item \emph{The complete oriented comparison.}
Iterate the prime splitting \eqref{maps:eq:D.MP9}--\eqref{maps:eq:D.MP10}.
Positive integral homology of the distinct primary
Eilenberg--Mac Lane factors has zero cross tensor and
Tor products. The wedge-to-product inclusion is therefore
a simply connected homology equivalence and hence a
homotopy equivalence. It identifies the full spin-bordism
class over $K(\mathbb Z/c,2)$ with the primary classes
just computed, through the actual Chinese remainder
projections. The automorphism group splits in the same way.

\item The stated connected sums have exactly these classes:
their natural maps obtained by collapsing the separating
spheres induce the direct-sum isomorphism on $H_2$, and
the spin one-handle bordism identifies connected sum
with addition. They are simply connected and spin, have
$H_2=\mathbb Z/c$ and $b_3=0$. The uniqueness assertion
of Zhubr's Theorem~3 \cite[p.795]{manZhubr75} consequently
gives each diffeomorphism in \eqref{maps:eq:D.MT1}. The two $c=15$
types are distinct because their five-primary refined
classes have different automorphism orbits.
\end{enumerate}
\end{proof}

The extraction program was run at residuals $6,15,30$.
The independent exact local certificate is
\path{code/III/manifolds_three_mixed_AH.py}; its output records
all fourteen patterns, twelve marked rows, $464$
semicharacter substitutions and all three finite normal
coefficients. The full three-primary value is justified
by the primary operation and the Borel comparison above,
not by those finite algebraic checks alone.


\begin{theorem}[a general triangular cusp and its full boundary kernel]\label{maps:src:D.2.43}
\StatementRecord{maps:src:D.2.43}{proved}{Historical events: Claude Fable 5.2 (Anthropic), effort high, 16 November 2026: carried as background without fault found. GPT-6 Astra (OpenAI), ChatGPT Pro mode, effort Pro, S6\_09\_PAPERS, 17 November 2026: targeted proof check. Also checked in this writing. Pre-continuation writing events: GPT-6 Astra (OpenAI), ChatGPT Pro mode, effort Pro, S6\_10\_PAPERS, 18 November 2026: targeted proof check. Continuation metadata events: GPT-6 Astra (OpenAI), ChatGPT Pro mode, effort Pro, S6\_10\_PAPERS continuation, 19 November 2026: used and rechecked in this run; status-vocabulary reconciliation; metadata reconciliation only; no new mathematical check. \emph{Pre-continuation recorded proof scope.} GPT-6 Astra (OpenAI, ChatGPT Pro mode), effort Pro, 2026-11-17 UTC; Claude Fable 5.2 (Anthropic), effort high, 16 November 2026; also checked in this writing by GPT-6 Astra (OpenAI), ChatGPT Pro mode, effort Pro, 18 November 2026}{the hypotheses, notation and definitions stated here; \cref{maps:src:D.2.22}; \cref{maps:src:E.2.19}; \cref{maps:eq:D.GC9}; \cref{maps:eq:D.GC11}; \cref{maps:eq:D.GC10}; \cref{maps:eq:D.GC8}; \cref{maps:eq:D.GC13}}


Let $N_L$ be the smooth same-fan triangular Mumford cusp
whose short-period lattice is a full sublattice
$L\subset\mathbb Z^2$. Write its Smith factors as
$s\mid t$, and put $D=st$. Holomorphic unit factors
in the two periods are allowed. Let $M_L=\partial N_L$,
and let $\Lambda=\mathbb Z^4$ be the regular torus
lattice. Its unipotent monodromy is $A=I+N$,
where $N^2=0$ and the map from the long lattice to
the primitive short lattice has Smith factors $(s,t)$.
Put $F=\ker(A-I)$, the primitive short lattice.

Then $\pi_1(N_L)=L\simeq\mathbb Z^2$, and its ordinary
integral homology is
\[
 \begin{array}{c|ccccc}
 j&0&1&2&3&4\\ \hline
 H_j(N_L)&\mathbb Z&\mathbb Z^2&
       \mathbb Z^{D+3}\oplus\mathbb Z/s&
       \mathbb Z^2&\mathbb Z^D .
 \end{array}                                             \tag{D.GC1}\label{maps:eq:D.GC1}
\]
All higher groups vanish. In particular $\chi(N_L)=2D$.
The full regular-fibre kernel is the generally
non-saturated lattice
\[
 \ker\{H_2(T^4)\to H_2(N_L)\}
       =(\Lambda^2A-I)\Lambda^2\Lambda+\Lambda^2F .
                                                        \tag{D.GC2}\label{maps:eq:D.GC2}
\]
Its image is $\mathbb Z^4\oplus\mathbb Z/s$, with
free cokernel of rank $D-1$ in $H_2(N_L)$.
The two actual fixed angular normal tori bound
and give the entire fixed-normal projection.
Together with \eqref{maps:eq:D.GC2} they specify the full
boundary kernel. More explicitly,
\[
 \begin{split}
 H_1(M_L)&=\mathbb Z^3\oplus\mathbb Z/s\oplus\mathbb Z/t,\\
 H_2(M_L)&=\mathbb Z^6\oplus(\mathbb Z/s)^2,\\
 \ker\{H_2(M_L)\to H_2(N_L)\}
             &\simeq\mathbb Z^2\oplus\mathbb Z/s .
 \end{split}                                             \tag{D.GC3}\label{maps:eq:D.GC3}
\]
The torsion killed by this map is the primitive
angular area modulo the monodromy image. The other
boundary torsion summand maps isomorphically to
the actual cap torsion.

There is also an exact component-meridian marking.
Use the positive $q$-meridian $\theta$ and index
the $D$ normalized components by $v\in\mathbb Z^2/L$.
The positive transverse disc to that component has
boundary
\[
                     \theta+\bar v
           \quad\hbox{in }\mathbb Z\theta\oplus\mathbb Z^2/L
                     \subset H_1(M_L).                  \tag{D.GC4}\label{maps:eq:D.GC4}
\]
These are the complete columns of the relative
map $H_2(N_L,M_L)=H^4(N_L)=\mathbb Z^D\to H_1(M_L)$.
In particular the component differences need not
vanish before global gluing.

\end{theorem}


\begin{proof}
\begin{enumerate}[label=\arabic*.,leftmargin=*]
\item \emph{The actual covering and its two-variable normalization.}
The extension of period units and the free same-fan
cover construction in \ref{maps:src:D.2.22} identify $N_L$ with
the finite cover of the minimal hexagonal cusp
associated to $L$. The units can also be removed
through the relative smooth family of \ref{maps:src:E.2.19}.
The normalized components are $dP_6$ surfaces,
the double curves are spheres, and the triple
strata are points. Their dual complex is
$\mathbb R^2/L$. Van Kampen gives
$\pi_1(N_L)=L$, because each component and every
intersection component is simply connected.

\item View the common universal cover with the full deck
action of the minimal cusp, and put
$R=\mathbb Z[x^{\pm1},y^{\pm1}]$ and
$I=(x-1,y-1)$. In the cyclic toric order
$(E_1,H-E_1-E_2,E_2,H-E_2-E_3,E_3,H-E_1-E_3)$,
the three opposite-pair translations have
weights $x^{-1},y,xy$. Indeed the primitive
fan rays are $(1,0),(0,1),(-1,1)$, and the
opposite pairing translates by their negatives;
choose the second period generator with the
corresponding sign. Simultaneously inverting
the deck variables changes neither the ideal
$I^2$ nor the assertions below.

\item With intersection form $\operatorname{diag}(1,-1,-1,-1)$,
the degree-two normalization matrix is
\[
 B=\begin{pmatrix}
 -x^{-1}&-1&-x^{-1}&-x^{-1}\\
 1&1&1&y\\
 -xy&-xy&-1&-xy
 \end{pmatrix}.                                         \tag{D.GC5}\label{maps:eq:D.GC5}
\]
The homology differential is its transpose with
the corresponding deck involution. One unit
row/column cancellation reduces its transpose to
\[
 F=\begin{pmatrix}
 1-x&0\\0&xy-1\\y-1&xy(y-1)
 \end{pmatrix}.                                         \tag{D.GC6}\label{maps:eq:D.GC6}
\]
For clarity there is an explicit Laurent-unimodular
comparison, not just equality of ranks or minors.
Let
\[
 T=\begin{pmatrix}
 0&0&-x^{-1}\\
 -x^{-1}y^{-2}&x^{-1}y^{-1}&-x^{-2}y^{-1}\\
 -y^{-1}&0&0
 \end{pmatrix},\qquad
 S=\begin{pmatrix}y&xy^2\\x&0\end{pmatrix}.
\]
Their determinants are Laurent units, and
\[
 T^{-1}F=
 \begin{pmatrix}-(y-1)&0\\x-1&-(y-1)\\0&x-1\end{pmatrix}S.
                                                        \tag{D.GC7}\label{maps:eq:D.GC7}
\]
The middle matrix is the exact presentation
of $I^2$ by the generators
$(x-1)^2,(x-1)(y-1),(y-1)^2$.
Its exactness follows directly from the regular
sequence $x-1,y-1$: the two displayed relations
are all the syzygies among those three monomials.
Thus
\[
                H_2(\widetilde N_L;\mathbb Z)\simeq I^2
                         \quad\hbox{as an }R\hbox{-module}.
                                                        \tag{D.GC8}\label{maps:eq:D.GC8}
\]
Here the universal cover is the same for every $L$.
The degree-zero normalization row is the cellular
complex of the plane, and the degree-four row is
$R$ in normalization degree zero. The degree-two
differential is injective over the domain $R$.
No higher differential is possible, because
normalization degree is at most two and the odd
homology of every stratum vanishes.

\item \emph{The integral finite-cover groups.}
Write $G=\mathbb Z^2/L$ and $R_G=\mathbb Z[G]$.
The degree-two normalization cokernel is
$I^2\otimes_R R_G$. We compute its torsion exactly.
The two short exact sequences
\[
 0\to I\to R\to\mathbb Z\to0,\qquad
 0\to I^2\to I\to\mathbb Z^2\to0
                                                        \tag{D.GC9}\label{maps:eq:D.GC9}
\]
use the trivial $R$-action on $I/I^2=\mathbb Z^2$.
The Koszul resolution for $x-1,y-1$ identifies
$\operatorname{Tor}^R_j(\mathbb Z,R_G)$ with
$H_j(L;\mathbb Z)$.
Consequently $I\otimes_R R_G$ is free of rank $D+1$:
it is an extension of the free augmentation ideal
of $\mathbb Z[G]$, of rank $D-1$, by $L\simeq\mathbb Z^2$.
Its map to $\mathbb Z^2$ in \eqref{maps:eq:D.GC9} is onto,
since $x-1$ and $y-1$ map to the two standard
generators. Its kernel is free of rank $D-1$.

\item The relevant Tor segment from the second sequence is
\[
 \mathbb Z\ \xrightarrow{\;\alpha_L\;}\ \mathbb Z^4
 \longrightarrow I^2\otimes_RR_G
 \longrightarrow I\otimes_RR_G
 \longrightarrow\mathbb Z^2\longrightarrow0 .
                                                        \tag{D.GC10}\label{maps:eq:D.GC10}
\]
The source is $H_2(L)$ and the target is
$H_1(L)\otimes\mathbb Z^2$.
If $u,v$ are an oriented basis of $L$, the map is
\[
                  \alpha_L(1)=u\otimes v-v\otimes u,
                                                        \tag{D.GC11}\label{maps:eq:D.GC11}
\]
where the second tensor factor uses the inclusion
$L\hookrightarrow\mathbb Z^2$.
One can check \eqref{maps:eq:D.GC11} directly in the Koszul
complex. In Smith coordinates
$L=\langle(s,0),(0,t)\rangle$, let
$N_x=1+\cdots+x^{s-1}$ and
$N_y=1+\cdots+y^{t-1}$.
The map on resolutions sends the relation
$(-(y-1),x-1)$ of $I$ to
$-e_y\otimes e_x+e_x\otimes e_y$.
The invariant cycle $N_xN_y$ therefore maps,
in $H_1(L)\otimes\mathbb Z^2$, to a vector with
nonzero coordinates $-s,t$.
Its Smith factor is exactly $s=\gcd(s,t)$.
The augmentation ideal and its square are invariant
under changes of basis of $\mathbb Z^2$, so this
also proves \eqref{maps:eq:D.GC11} for an arbitrary $L$.

\item Thus $\alpha_L$ is injective and its cokernel is
$\mathbb Z^3\oplus\mathbb Z/s$.
The free last kernel in \eqref{maps:eq:D.GC10} gives
\[
            I^2\otimes_RR_G
                    \simeq\mathbb Z^{D+2}\oplus\mathbb Z/s.
                                                        \tag{D.GC12}\label{maps:eq:D.GC12}
\]
The preceding Tor segment gives
$\operatorname{Tor}_1^R(I^2,R_G)=\mathbb Z^2$.
The degree-zero normalization row now has the
homology of the dual torus, adding one free
summand to total degree two and two to degree
one. The degree-four row contributes $\mathbb Z^D$.
These observations prove \eqref{maps:eq:D.GC1}; the extension
by the dual torus's degree-two $\mathbb Z$ splits
as an abelian-group extension.

\item \emph{The actual boundary map and the nonsaturated kernel.}
Choose bases with
$A b_1=b_1+s a_1$, $A b_2=b_2+t a_2$ and
$A a_i=a_i$. The exterior-square difference has
image generated by
\[
 s\,a_1\wedge a_2,\qquad
 s\,a_1\wedge b_2-t\,a_2\wedge b_1
                         +st\,a_1\wedge a_2 .
                                                        \tag{D.GC13}\label{maps:eq:D.GC13}
\]
Its two Smith factors are $(s,s)$.
The Wang sequence gives the first two boundary
groups in \eqref{maps:eq:D.GC3}.

\item The minimal cusp's two fixed angular normal
tori have actual section-disc bounding sweeps.
In the unitless toric chart at the dense locus
of the component indexed by zero, take the
section with its two nonvanishing torus
coordinates equal to one and its normal
coordinate equal to $q$. Acting on this
section by either short circle gives a map
$S^1\times D^2\to N$ whose boundary is the
literal fixed angular normal torus. Its
fundamental circle lies in a simply connected
normalized component, so its image in the
cusp fundamental group is zero. The whole
map therefore lifts to every period cover
$N_L$, with its marked boundary. Transport
through the relative unit-removing family
preserves these boundary classes. Thus the
two primitive fixed-normal projections still
bound, and their actual classes split off
the full Wang projection $H_2(M_L)\to F$.

\item The primitive angular area also maps to zero.
In the universal cover the lifted regular
fibre is $T^2\times\mathbb R^2$. Its degree-two
group is a trivial $R$-module. Its image in
the module $I^2$ of \eqref{maps:eq:D.GC8} is zero, since
$I^2$ has no nonzero element fixed by both
$x$ and $y$. This nullhomology projects to
every finite period cover.

\item Poincare--Lefschetz duality gives
\[
 H_3(N_L,M_L)=H^3(N_L)=\mathbb Z^2\oplus\mathbb Z/s,
 \qquad H_2(N_L,M_L)=H^4(N_L)=\mathbb Z^D .
                                                        \tag{D.GC14}\label{maps:eq:D.GC14}
\]
The already bounding angular normal tori exhaust
the rational rank of the boundary kernel.
Hence the four-dimensional rational fibre
coinvariant injects into the cap.
Moreover the cokernel of
$H_2(M_L)\to H_2(N_L)$ injects into the free
group $H^4(N_L)$, so it is free.
Because the fixed-normal projections bound,
the boundary image equals the regular-fibre image.
It therefore contains all of the cap torsion.
There is no hidden free-to-torsion extension
here: if a fibre-coinvariant element maps to
a torsion class, rational injectivity forces
that element itself to be torsion. Hence
the cap torsion is the image of the boundary
fibre torsion subgroup, not of a free class
whose rational image has been overlooked.

\item The two boundary fibre torsion summands in
\eqref{maps:eq:D.GC13} have order $s$. The angular-area
one is killed. The other must map onto
$\operatorname{Tor}H_2(N_L)=\mathbb Z/s$ and
hence maps isomorphically. No free fibre
class can enter the kernel, since the rational
fibre coinvariant injects. This proves exactly
\eqref{maps:eq:D.GC2}, not its saturation, and gives all
the remaining assertions in \eqref{maps:eq:D.GC3}.
The free cokernel has rank $(D+3)-4=D-1$.

\item Finally the smooth fan has primitive rays
$(v,1)$. A positive small normal disc to
the associated component has this primitive
angular-plus-base meridian. Reducing the
angular coordinate modulo the period lattice
gives \eqref{maps:eq:D.GC4}, with the complex orientation
fixing the common sign. Every coset of $L$
occurs once. Their differences generate
$\mathbb Z^2/L$, and the zero coset gives
$\theta$; thus the displayed columns give
the entire kernel of $H_1(M_L)\to H_1(N_L)$.
\end{enumerate}
\end{proof}

The complete Laurent-unimodular operations,
finite HNF-lattice normalization matrices,
boundary exterior-square matrices and all
component-meridian columns are implemented in
\path{code/III/manifolds_general_cusp.py}.
When $s=1$, the fibre kernel is saturated
and the result recovers the local input used
in \ref{maps:src:D.2.34}--\ref{maps:src:D.2.41}. When $s>1$, replacing \eqref{maps:eq:D.GC2}
by its saturation would incorrectly erase
the actual cap torsion $\mathbb Z/s$.


\begin{theorem}[the seven semistable residual-four and residual-eight smooth types]\label{maps:src:D.2.44}
\StatementRecord{maps:src:D.2.44}{proved}{Historical events: Claude Fable 5.2 (Anthropic), effort high, 16 November 2026: carried as background without fault found. GPT-6 Astra (OpenAI), ChatGPT Pro mode, effort Pro, S6\_09\_PAPERS, 17 November 2026: targeted proof check. Also checked in this writing. Pre-continuation writing events: GPT-6 Astra (OpenAI), ChatGPT Pro mode, effort Pro, S6\_10\_PAPERS, 18 November 2026: targeted proof check. Continuation metadata events: GPT-6 Astra (OpenAI), ChatGPT Pro mode, effort Pro, S6\_10\_PAPERS continuation, 19 November 2026: used and rechecked in this run; status-vocabulary reconciliation; metadata reconciliation only; no new mathematical check. \emph{Pre-continuation recorded proof scope.} GPT-6 Astra (OpenAI, ChatGPT Pro mode), effort Pro, 2026-11-17 UTC; Claude Fable 5.2 (Anthropic), effort high, 16 November 2026; also checked in this writing by GPT-6 Astra (OpenAI), ChatGPT Pro mode, effort Pro, 18 November 2026}{the hypotheses, notation and definitions stated here; \cref{maps:src:D.2.8}; \cref{maps:src:D.2.26}; \cref{maps:src:D.2.41}; \cref{maps:src:D.2.35}; \cite[Sections 3.11--3.14]{manZhubr75}; \cite[Theorem 5 and Sections 3.11--3.14]{manZhubr75}; \cite{ffZhubr13}; \cite[pp.369--372]{ffBrown72}; \cite[Section 5]{ffZhubr13}; \cite[Definition 1.7, Remark 1.8, Theorem 3.1]{mapsKaplan79}; \cite[Theorem 5]{manZhubr75}; \cite[Theorem 3]{manZhubr75}}


Consider the default semistable rows OS55 with $(a,b)=(5,4)$ and zero
index $2,3,4$, and OS45 with $(a,b)=(1,8)$ and zero index $1,2,3,4$.
Use the primitive Mumford cap at every multiplicative site, no additional
regular logarithmic transforms, and either simply connected ordinary
clutch $e\gamma$, $e=\pm1$. Write $c=b$.
The resulting oriented six-manifolds have
\[
 H_2(X;\mathbb Z)=H_3(X;\mathbb Z)=\mathbb Z/c,
 \qquad \lambda(X)=0,\qquad \mu(X)=0,
 \qquad
 P(X)=\begin{cases}2,&\mathrm{OS55},\\0,&\mathrm{OS45}.
 \end{cases}                                                   \tag{D44.1}\label{maps:eq:D44.1}
\]
Here $P$ is Zhubr's refined quarter-Pontryagin class, transported to
$\mathbb Z/c$ by any generator; both displayed values are independent
of that choice. In particular all three OS55 manifolds are oriented
diffeomorphic to $Z_{4,1}$, and all four OS45 manifolds are oriented
diffeomorphic to $Z_{8,0}$, with the following concrete surgery models.
The manifold $Z_{c,0}$ is surgery on an embedded two-sphere representing
$c$ times the generator of $H_2(S^2\times S^4)$.
The manifold $Z_{c,1}$ is surgery in
$M_{\rm conic}\#(S^2\times S^4)$ on an embedded two-sphere representing
$(t,-c/2)\in\mathbb Z/2\oplus\mathbb Z$, where $t$ is the generator
of the conic-surgery manifold of \ref{maps:src:D.2.8}. These are spin surgeries with
their unique normal framings. This theorem makes no assertion about
the remaining additive profiles or additional regular logarithms.

\end{theorem}
\begin{proof}
\begin{enumerate}[label=\arabic*.,leftmargin=*]
\item The marked homology calculation \ref{maps:src:D.2.26} and its all-weight verification
\ref{maps:src:D.2.41} give the stated integral groups, with no other reduced homology
except $H_6=\mathbb Z$. The full integral support-function and
semicharacter certificate
\path{code/III/fourfolds_D2_2power_AH.py}
verifies the integral line-descent and actual-angle hypotheses
used in \ref{maps:src:D.2.35} for these rows. Here are the complete remaining
localization steps. The subgroup $G=C_c$ of the scalar circle has
fixed locus the rational curve $C$ and the one exceptional
elliptic double curve $E$ at $I_c$; it acts freely off them.
The actual normal characters are $(1,-1)$ on $C$ and $(w,-w)$
on $E$, with $w^2=1$ in $\mathbb Z/c$. The two normal lines on
$C$ have degree $-1$, and those on $E$ have degree zero.
The symmetric auxiliary line of \ref{maps:src:D.2.35} has character $c/2$ on
$C$, zero on $E$, and degree coprime to $c$ on $C$.
Its identification with the global mod-$c$ class is the same
integral common-piece argument as there: the boundary-to-cap
surjections and primitive degree-one map make $H_2(W)\to H_2(X)$
onto; the line has zero swept values and kills the ordinary
clutch relation modulo $c$. Moreover $\pi_1(W_G)$ has the
unique torsion subgroup $G$. The difference between equivariant
lifts therefore has the same point character on $C$ and $E$.
After one base-character adjustment, the stated auxiliary
characters are the characters of a global equivariant lift.

\item Normalize its ordinary class $\alpha\in H^2(X;\mathbb Z/c)$
to have degree one on $C$. Put $t\in H^2(BG;\mathbb Z/c)$,
$x=[C]^*$ and $y=[E]^*$. Its restrictions have the form
\[
 \alpha|_C=x+(c/2)t,\qquad
 \alpha|_E=A y+\eta\zeta,\qquad
 e_G(N_C)=-t^2,\qquad e_G(N_E)=-w^2t^2,                       \tag{D44.0a}\label{maps:eq:D44.0a}
\]
where $\eta$ has degree one on $BG$ and $\zeta$ lifts from
integral degree-one cohomology of $E$. Localization over
$\mathbb Z/c$ is legitimate: the complement is free, so its
cohomology vanishes after inverting $t$. The degree-one
localization identity gives $A=-w^2$. The elliptic restriction
has cube zero, because $y\zeta=0$ and $\zeta^2=0$; the rational
curve's contribution to the cube is $-3(c/2)^2=0\pmod c$.
Thus $\langle\alpha^3,[X]\rangle=0$ in the full coefficient
ring, not merely modulo two.
The equivariant complex determinant is trivial: its ordinary
class is zero and its total character on $C$ is zero.
Hence $\lambda_G=-c_{2,G}$. On either fixed curve,
$c_{2,G}(TX)=e_G(N)$, since the square of the tangent-line
class vanishes. Therefore
\[
 \langle\alpha\lambda,[X]\rangle=-(1+A)=-(1-w^2)=0.           \tag{D44.0b}\label{maps:eq:D44.0b}
\]
The normalized class detects the cyclic degree-four group, so
$\lambda=0$. Zhubr's Sections 3.11--3.14 identify the complete
group $H_6(K(\mathbb Z/c,2);\mathbb Z)=\mathbb Z/c$ by the cube,
giving $\mu=0$ \cite[Sections 3.11--3.14]{manZhubr75}.
The free-coefficient calculation
and coefficient surjectivity in Zhubr's Sections 3.11--3.14 give
$2P=\operatorname{PD}(\lambda)$
\cite[Theorem 5 and Sections 3.11--3.14]{manZhubr75}.
Thus the two remaining possibilities
are exactly $P=0,c/2$. The calculation below distinguishes them by
the actual tangent framing, without halving a congruence or using
the conjectural change-of-three-frame formula of \cite{ffZhubr13}.

\item \emph{A framing parity that depends only on the three-cycle.}
We first record the precise primary construction used to compare a
regular three-torus with a torsion three-sphere.
Let a closed oriented six-manifold $Y$ have a stable tangent framing.
For an embedded oriented three-manifold $N\subset Y$ with a chosen
normal framing, denote by $[N]_{\rm fr}\in\Omega^{\rm fr}_3
\cong\mathbb Z/24$ its induced stable framed-bordism class. Then
\[
 [N]_{\rm fr}\pmod2
   =q_Y\bigl(\operatorname{PD}_{\mathbb F_2}[N]\bigr)             \tag{D44.2}\label{maps:eq:D44.2}
\]
for a function of the indicated mod-two cohomology class alone.
This statement includes torsion classes in integral homology.
Here is a derivation directly from Brown's construction
\cite[pp.369--372]{ffBrown72}. Set $K_3=K(\mathbb F_2,3)$ and let
$j:S^3\to K_3$ represent its fundamental class. Brown's $F(u,0)$
is literally the stabilization of the map $u:Y\to K_3$.
The framed Pontryagin--Thom collapse $Y\to S^3$, followed by $j$,
therefore represents $F(\operatorname{PD}[N],0)$.
Apply the framed duality map $W_*A_\alpha$ of his Theorem 1.7 and
Corollary 1.11. Its value is the image of $[N]_{\rm fr}$ under
\[
 j_*:\pi_3^{\rm S}\longrightarrow\pi_6^{\rm S}K_3
          \cong\mathbb Z/2.                                  \tag{D44.3}\label{maps:eq:D44.3}
\]
The target is Brown's Lemma 1.2. The remark immediately following
his Corollary 1.13 identifies $j_*\nu$ with its nonzero generator:
the coefficient is the mod-two Hopf invariant, and the quaternionic
Hopf element $\nu$ has Hopf invariant one. Dividing Brown's
$\mathbb Z/4$-valued function, whose values here are $0,2$, by two
gives \eqref{maps:eq:D44.2}. This is a Pontryagin--Thom argument about the specified
normal framing; no compression of an arbitrary framed torus to a
sphere is being assumed.

\item For clarity we also fix the sphere convention. The homomorphism
$\pi_3SO(3)\to\pi_3SO(6)$ has image $2\mathbb Z$.
Consequently changing the normal framing of an embedded $S^3$
does not change the parity of its induced stable framed class.
For the standard sphere bounding a ball, the induced stable framed
class is zero and its tangent three-frame represents the nonzero
class in $\pi_3V_{6,3}=\mathbb Z/2$ relative to a constant ambient
three-frame. Equivalently, the tangent Lie framing on $S^3$ has
class $\nu$, and a change of ambient frame by degree $d$ changes
that class by $d\nu$, while the relative three-frame class changes
by $d\bmod2$. Thus Zhubr's required nontrivial relative
three-frame is equivalent to \emph{even} induced stable framed
class on the sphere. These conventions agree with the quaternionic
calculation in \cite[Section 5]{ffZhubr13}.

\item Since $\lambda(X)=0$, $TX$ is stably trivial: the only obstruction
for a spin bundle on a six-dimensional complex is its class in
$H^4(-;\pi_3\mathrm{Spin})$, namely $\lambda$.
Moreover $TX$ is actually trivial on $X^\circ=X\setminus\{p\}$.
For example its complex bundle has $c_1=c_2=0$ and its last
obstruction lies in $H^6(X^\circ;\mathbb Z)=0$.
Any framing there extends stably across $p$, since
$\pi_5SO=0$. All such full framings have the same relative
three-frame parity on an embedded torsion $S^3$: their difference
lifts to $\mathrm{Spin}(6)$ and its degree-three class is a class
of $H^3(X^\circ;\mathbb Z)$, which evaluates to zero on a torsion
three-cycle.
An embedded sphere in the integral class of the regular torus
$T=\langle e_1,e_2,\gamma\rangle$ used below can also be
constructed explicitly, so its existence is not an additional
representability assumption. Apply Kaplan's construction to
the spin structure induced on the normally framed $T$:
it gives a four-dimensional filling with one zero-handle and
only two-handles, preserving that almost framing
\cite[Definition 1.7, Remark 1.8, Theorem 3.1]{mapsKaplan79}.
Delete the zero-handle and reverse the cobordism. Realize the
resulting two-handle attachments as ambient surgeries in
$X^\circ$. Their attaching circles bound embedded disks
with interiors disjoint from $T$ and from one another, since
$\pi_1(X^\circ)=0$, $2+3<6$ and $2+2<6$.
Prescribe a boundary collar along the first normal vector.
The two tangent framing directions of each attachment and
the other two normal directions give a four-frame around
the boundary of its disk. Its extension obstruction in
$\pi_1SO(4)=\mathbb Z/2$ is exactly the spin obstruction
of that two-handle, and vanishes. Thickening these disks
therefore performs the prescribed surgeries inside $X^\circ$.
The resulting embedded $S^3$ has the same integral homology
class as $T$. Brown's identity already supplies its parity;
this construction additionally verifies its geometric
realization without changing the prescribed spin structure.
Zhubr's Theorem 1 now implies
\[
 P(X)=0\quad\Longleftrightarrow\quad
 \text{a full ambient framing is admissible on the torsion }S^3.
                                                               \tag{D44.4}\label{maps:eq:D44.4}
\]
Indeed the forward implication is not an assumption about arbitrary
three-frames. Take the admissible three-frame supplied by that
theorem. If $P=0$, its spin rank-three complement has vanishing
primary obstruction $c_2$; the only possible further obstruction
on the punctured six-manifold is in $H^5(X^\circ;\mathbb Z/2)=0$.
The complement is therefore trivial, and the three-frame extends
to a full framing. Conversely an admissible full framing has
trivial complement and hence refined class zero. Independence
of the full framing proves \eqref{maps:eq:D44.4} for the particular framing
used below.

\item \emph{The reference frame and the actual critical-cycle defect.}
Write $f:X\to\mathbb P^1$ for the torus fibration and
$\mathcal V=\ker(df)$ for its holomorphic differential kernel.
This is a locally free rank-two sheaf, including at the critical
curves. In the nodal chart $f=xy$ it has basis
$x\partial_x-y\partial_y,\partial_z$; in the triangular cusp
chart $f=xyz$ it has basis
$x\partial_x-y\partial_y,y\partial_y-z\partial_z$.
The scalar $\delta$-action defines a nowhere-zero section of this
\emph{abstract} vector bundle, although its image as a vector
field vanishes on the fixed curve. In a toric chart its two
coefficients are those of the nonzero primitive scalar lattice
vector, so they cannot both vanish. This description is unchanged
by the translation clutch.
As $H^2(X;\mathbb Z)=0$, the quotient line $\mathcal V/\mathcal O\delta$
is topologically trivial; hence $\mathcal V$ has a global smooth
frame whose first vector is $\delta$. The line
$\mathcal L=f^*T\mathbb P^1$ is also topologically trivial.

\item Let $Z$ denote the scheme defined by the coefficients of $df$.
There is an exact sequence
\[
 0\longrightarrow\mathcal V\longrightarrow TX
  \xrightarrow{df}\mathcal L
  \longrightarrow\mathcal L\otimes\mathcal O_Z
  \longrightarrow0.                                         \tag{D44.5}\label{maps:eq:D44.5}
\]
Its codimension-two cycle is the literal sum $C_{\rm crit}$ of
the complex double curves, each with coefficient one. This
includes the three axes at a triple point: the ideal
$(xy,xz,yz)$ has multiplicity one on each axis. Since
$\mathcal V$ and $\mathcal L$ are topologically trivial,
\eqref{maps:eq:D44.5} gives $c_2(TX)=[C_{\rm crit}]$. Thus this cycle is
integrally nullhomologous.

\item Off $Z$, splitting \eqref{maps:eq:D44.5} identifies $TX$ with the framed
bundle $\mathcal V\oplus\mathcal L$. Its comparison with a
global frame on $X^\circ$ has an integral degree-three class
$a$ whose local residue on every smooth double curve is one,
up to the same overall sign. In the model $xy$ this is the
explicit Bott matrix
\[
 \frac1{\sqrt{|x|^2+|y|^2}}
 \begin{pmatrix}x&\overline y\\-y&\overline x\end{pmatrix}
       \in SU(2).                                            \tag{D44.6}\label{maps:eq:D44.6}
\]
It has degree one on the normal three-sphere; adding the
tangent line to the critical curve sends this to the stable
generator, not twice that generator. The latter factor two
belongs to the different map $SO(3)\to SO$, as is also seen
from $\lambda=-c_2$ for this complex rank-two normal calculation.
Removing the finitely many triple points does not affect this
degree-three residue calculation. The integral relative
cohomology sequence, or equivalently the Thom intersection
construction, therefore gives
\[
 \langle a,[T]\rangle
      =\pm\operatorname{lk}(T,C_{\rm crit})                   \tag{D44.7}\label{maps:eq:D44.7}
\]
for any regular three-torus $T$ disjoint from the critical
cycle. The ambiguity of adding a class from $H^3(X^\circ;\mathbb Z)$
evaluates to zero on our torsion three-cycle. Thus the integer
in \eqref{maps:eq:D44.7}, not merely a residue modulo $c$, is well defined.

\item Take specifically $T=\langle e_1,e_2,\gamma\rangle$ in a regular
fibre inside the ordinary clutch disc. The reference frame
$\mathcal V\oplus\mathcal L$ has zero degree-three difference
from the flat translation frame there. Indeed the vertical
frames have the same first vector $\delta$; their change lies
in the stabilizer $\mathbb C\rtimes\mathbb C^*$, which retracts
to a circle. The base line contributes another $\mathbb C^*$,
and changing the horizontal splitting contributes a contractible
space. None contributes a degree-three class. The induced flat
framing of $T^3$ is its product Lie framing, with stable class
$\eta^3=12\nu$, an even element of $\mathbb Z/24$.
A possible determinant-circle transition cannot change this
parity. A map $T^3\to S^1$ is homotopic to a character and,
after an integral change of torus basis, depends on just one
circle. Its real rotation can be placed stably in that circle
and one normal direction. It changes just that circle's framed
class, so the resulting product-framed torus has class either
$0$ or $\eta^3$, both even. After splitting off the determinant,
the remaining $SU(3)$ transition is a degree-$k$ map on a
three-dimensional complex and changes the framed class by
$\pm k\nu$. Thus its exact parity is the integer in \eqref{maps:eq:D44.7}.
Equations \eqref{maps:eq:D44.2} and \eqref{maps:eq:D44.7} consequently give
\[
 q_X(\operatorname{PD}[T])
       =\operatorname{lk}(T,C_{\rm crit})\pmod2.              \tag{D44.8}\label{maps:eq:D44.8}
\]
The torus represents an odd generator of $H_3(X)=\mathbb Z/c$
by the explicit linking calculation below. Its reduction is
therefore the same as that of an embedded torsion generator
$S^3$. Combine \eqref{maps:eq:D44.4} and \eqref{maps:eq:D44.8} to obtain the actual decision
\[
 P(X)=\frac c2
       \bigl(\operatorname{lk}(T,C_{\rm crit})\bmod2\bigr).
                                                               \tag{D44.9}\label{maps:eq:D44.9}
\]

\item \emph{The integral linking calculation.}
All local preimages and bounding coefficients here are those
of the actual marked caps, not an invariant-form calculation
on an unmarked torus. At each site choose the recorded
determinant-one elliptic basis $(u_1,u_2)$ and use the short
and long planes $(u_1,\delta)$ and $(u_2,\gamma)$. The exact
valuation matrix is
\[
 L=\begin{pmatrix}n&-bH\\aH&R\end{pmatrix}.                  \tag{D44.10}\label{maps:eq:D44.10}
\]
At a rank-one site write $L=\alpha r$ with $\alpha$ primitive.
Choose a short complement $\sigma$ with
$\det(\alpha,\sigma)=1$ and long vectors $(\beta,\ell)$ with
$r\beta=\rho>0$, $r\ell=0$, and $\det(\beta,\ell)=1$.
The literal transverse complex elliptic curve has regular-fibre
preimage $\sigma\wedge\ell$. The product nodal chart supplies
the specialization chain with degree one. These signs follow
by orienting the vanishing elliptic factor as $(\alpha,\beta)$
and the transverse factor as $(\sigma,\ell)$.

\item At the unique rank-two cusp $\det L=1$. Put
$(b_1,b_2)=(u_2,\gamma)L^{-1}$ and $(a_1,a_2)=(u_1,\delta)$,
so $Nb_j=a_j$. The three integral support functions
$r(r-1)/2,rs,s(s-1)/2$ have degree matrix
\[
 \begin{pmatrix}0&-1&0\\1&1&-1\\-1&0&0\end{pmatrix}          \tag{D44.11}\label{maps:eq:D44.11}
\]
on the three complex opposite-pair double curves in cyclic
ray order $(1,0),(0,1),(-1,1)$. This unimodular matrix gives
degree vector $(-1,1,-1)$ for their sum. The fibre vector
\[
 a_1\wedge b_1+b_1\wedge a_2+a_2\wedge b_2                  \tag{D44.12}\label{maps:eq:D44.12}
\]
has precisely those degrees and zero long--long component.
The latter is the degree-two class of the dual torus, so these
four coordinates identify its specialization with the actual
sum of the three double curves integrally. In particular no
undetected fourth homology component is being suppressed.

\item Transport the five preimages to the common regular fibre along
the fixed base paths. Their sum is a representative of
$C_{\rm crit}$ in the exact marked presentation of \ref{maps:src:D.2.41}.
All transport and cap relations can be realized by chains
outside the ordinary disc. In that disc the four normal
Thom chains are the products of a lattice loop with the
disc, and the redundant fifth chain has loop vector
$(0,0,0,-e)$. The intersection of a bounding relation with
$T=e_1e_2\gamma$ is consequently its clutch $\delta$
coefficient. Integral Smith solutions give
\[
 \begin{array}{c|c|c|c}
  \mathrm{OS}&c&\text{zero indices}&
       \operatorname{lk}(T,C_{\rm crit})\text{ for }e=1\\\hline
  45&8&1,2,3,4&4,4,4,4\\
  55&4&2,3,4&37,7,27.
 \end{array}                                                  \tag{D44.13}\label{maps:eq:D44.13}
\]
For example, in OS55 at zero index $2$, the five cap vectors
in ordered bivector coordinates
$(e_1e_2,e_1\delta,e_1\gamma,e_2\delta,e_2\gamma,\delta\gamma)$
are
\[
 \begin{gathered}
 (0,0,0,-12,0,5),\quad(3,0,1,-12,0,4),\\
 (-19,-23,-10,0,5,6),\quad(0,0,0,0,0,1),\quad(0,4,0,0,0,1).
 \end{gathered}                                               \tag{D44.14}\label{maps:eq:D44.14}
\]
Their sum has an integral bounding relation with clutch loop
$(0,0,37,0)$. The same Smith calculation for $c$ times a
two-dimensional torsion generator gives a clutch coefficient
coprime to $c$ in every row. Nondegeneracy of the torsion
linking pairing proves that $T$ is an odd generator of
$H_3(X)$. Every integral relation among presentation columns
has zero clutch $\delta$ coefficient, so the displayed
linking does not depend on a selected Smith solution.
Changing the stated short complements or the cusp mixed
preimage by an actual cap-kernel class also changes that
coefficient by zero. For $e=-1$ the exact same calculation
with its actual clutch vector reverses the displayed integer
signs and retains the unit torsion-linking test.

\item The executable
\path{code/III/manifolds_even_critical_cycles.py}
retains all fourteen rows, the full integral coefficient
vectors, local bases, alternative preimages and both clutch
signs in
\path{records/manifolds_even_critical_cycles.json}.
Its integral computation is independent of the differential
and framing arguments above. Formula \eqref{maps:eq:D44.9} and table
\eqref{maps:eq:D44.13} prove \eqref{maps:eq:D44.1}.
The separate executable
\path{code/III/fourfolds_D2_frame_certificate.py}
checks the local kernel and Bott-matrix identities, the
unimodular cusp degree matrix and, for all fourteen marked
rows, the independent degree identity
$\Omega(C_{\rm crit})=-be\,\operatorname{lk}(T,C_{\rm crit})$
for the actual pure fibre functional
$\Omega=a e_1^*e_2^*+b\gamma^*\delta^*$.
It retains its inputs' hash and every local degree in
\path{records/fourfolds_D2_frame_certificate.json}.

\item \emph{Recognition of the surgery models.}
The normal rank-four bundle of any embedded sphere in these
spin six-manifolds is trivial, and its framing is unique
since $\pi_2SO(4)=0$. The stated surgeries have $H_2=H_3=\mathbb Z/c$
and zero free third Betti number. The classifying map to
$K(\mathbb Z/c,2)$ extends across each surgery trace.
The model $Z_{c,0}$ has refined class and fundamental-class
invariant both zero. For $Z_{c,1}$, the conic summand's refined
class maps to $c/2$, and the complete cubic maps to zero for
$c=4,8$, since $(c/2)^3=0\pmod c$. Thus their full Zhubr
invariants are respectively $(P,\mu)=(0,0)$ and $(c/2,0)$,
using naturality in \cite[Theorem 5]{manZhubr75} and the full
conic invariant established in \ref{maps:src:D.2.8}.
The classification theorem \cite[Theorem 3]{manZhubr75}, including
its finite invariants, gives the asserted oriented diffeomorphisms. No conclusion
here rests on ordinary homology alone.\hfill$\square$
\end{enumerate}
\end{proof}


\begin{theorem}[the remaining minimal residual-four, eight and twenty types]\label{maps:src:D.2.45}
\StatementRecord{maps:src:D.2.45}{proved}{Historical events: Claude Fable 5.2 (Anthropic), effort high, 16 November 2026: carried as background without fault found. GPT-6 Astra (OpenAI), ChatGPT Pro mode, effort Pro, S6\_09\_PAPERS, 17 November 2026: targeted proof check. Also checked in this writing. Pre-continuation writing events: GPT-6 Astra (OpenAI), ChatGPT Pro mode, effort Pro, S6\_10\_PAPERS, 18 November 2026: targeted proof check. Continuation metadata events: GPT-6 Astra (OpenAI), ChatGPT Pro mode, effort Pro, S6\_10\_PAPERS continuation, 19 November 2026: used and rechecked in this run; status-vocabulary reconciliation; metadata reconciliation only; no new mathematical check. \emph{Pre-continuation recorded proof scope.} GPT-6 Astra (OpenAI, ChatGPT Pro mode), effort Pro, 2026-11-17 UTC; Claude Fable 5.2 (Anthropic), effort high, 16 November 2026; also checked in this writing by GPT-6 Astra (OpenAI), ChatGPT Pro mode, effort Pro, 18 November 2026}{the hypotheses, notation and definitions stated here; \cref{maps:src:D.2.30}; \cref{maps:src:D.2.44}; \cref{maps:src:D.2.36}; \cref{maps:src:D.2.26}; \cref{maps:src:D.2.41}; \cite[Definition 1.7, Remark 1.8 and Theorem 3.1]{mapsKaplan79}; \cite{ffZhubr13}; \cite[Section 5]{ffZhubr13}; \cite[Sections 3.11--3.14]{manZhubr75}; \cref{maps:src:D.2.38}; \cite{manZhubr75}}


Retain the primitive Mumford filling at every multiplicative site,
no additional regular logarithmic transformations, and the indicated
choice at the type-$II$ site in the following eleven patterns:
\[
\begin{array}{c|c|c|c}
 \mathrm{OS/profile}&(a,b)&\text{zero indices}&\text{type-}II\text{ choice}\\\hline
 45/II&(1,8)&2,3&\text{original or good log}\\
 55/II&(5,4)&3&\text{original or good log}\\
 55/\mathrm{default}&(1,20)&2,3,4&\text{none}\\
 55/II&(1,20)&3&\text{original or good log}.
\end{array}                                                    \tag{D45.1}\label{maps:eq:D45.1}
\]
Indices are zero-based. In an original-only pattern the ordinary
clutch is $e\gamma$ with $e=\pm1$. At a good $II$ logarithm allow
every order $m=6k$, $k>0$, every integer ordinary clutch $e$, and
every power vector
\[
 v=r\delta+n\ell,\qquad r\in\mathbb Z,
 \qquad n-me=\sigma\in\{1,-1\},                             \tag{D45.2}\label{maps:eq:D45.2}
\]
where $\ell$ is the actual primitive fixed lift with
$\gamma^*(\ell)=1$. These are exactly the simply connected data
in this one-log sector; they are free by \ref{maps:src:D.2.30}.
All the resulting manifolds have $H_2=H_3=\mathbb Z/b$ and
zero free third Betti number. Their full Zhubr invariants, in
the primitive fibre coordinate specified below, are
\[
 \begin{array}{c|c|c|c|c}
 b&P&\lambda&\mu&\text{oriented diffeomorphism type}\\\hline
 4&2&0&0&Z_{4,1}\\
 8&0&0&0&Z_{8,0}\\
 20&14&8&0&Z_{4,1}\#Z_{5,1}.
 \end{array}                                                 \tag{D45.3}\label{maps:eq:D45.3}
\]
Here $Z_{4,1},Z_{8,0}$ are the explicit framed two-sphere surgery
models of \ref{maps:src:D.2.44}, and $Z_{5,1}$ is the nonzero refined-class,
zero fundamental-class model of \ref{maps:src:D.2.36}. In particular \ref{maps:src:D.2.44}--\ref{maps:src:D.2.45}
identify all eighteen minimal patterns of residual order
$4,8,20$, including every simply connected finite multiplicity
and twist in \eqref{maps:eq:D45.2}. No assertion about additional regular
logs or higher translation weights is included.

\end{theorem}
\begin{proof}
\begin{enumerate}[label=\arabic*.,leftmargin=*]
\item The homology and fundamental-group statements are \ref{maps:src:D.2.26},
\ref{maps:src:D.2.30} and \ref{maps:src:D.2.41}. In particular $H^2(X;\mathbb Z)=0$, so
$c_1(X)=0$ and $X$ is spin. The differential construction of
\ref{maps:src:D.2.44} extends to the stated caps with two precise changes.
At an original type-$II$ elliptic cusp, the smooth total
surface has base map
\[
 q=y^2-x^3,\qquad
 \ker(df)=\mathcal O(2y,3x^2,0)\oplus\mathcal O(0,0,1).
                                                               \tag{D45.4}\label{maps:eq:D45.4}
\]
To see that this is the map, not just the singular zero set,
solve the Weierstrass equation for the base using its nonzero
linear base coefficient, split off its nondegenerate $y$
direction, and take the cube root of the remaining holomorphic
unit. The extra scalar coordinate is a local product coordinate.
The critical ideal $(x^2,y)$ has transverse length two.
Its literal critical cycle is therefore $2E_{II}$.
The normal comparison is the normalized unitary matrix
\[
 \frac1{\sqrt{4|y|^2+9|x|^4}}
 \begin{pmatrix}2y&-3\overline{x}^{2}\\3x^2&2\overline y\end{pmatrix}.
                                                               \tag{D45.5}\label{maps:eq:D45.5}
\]
Its degree is two: the holomorphic map $(x,y)\mapsto(2y,3x^2)$
has two positive local preimages of a small regular value.
Thus the relative obstruction, as well as the scheme cycle,
has coefficient two. The scalar $\delta$ is nonzero in the
second abstract kernel summand, including on $E_{II}$.

\item At a good logarithmic cap the total cyclic quotient is free,
so its good cover is etale over the smooth quotient total
space. On that cover $q=s^m$. If $D$ is the reduced multiple
fibre, remove the common divisorial factor of the differential
and use the target line
\[
 \mathcal L=f^*T\mathbb P^1\otimes
                   \mathcal O(-(m-1)D).                     \tag{D45.6}\label{maps:eq:D45.6}
\]
The corrected differential pulls back to $m\,ds$ and is
surjective. The division is deck-equivariant, and a unit
change $q=u s^m$ gives $mu\,ds+s\,du$, still surjective.
Consequently a good logarithmic cap has no corrected critical
cycle. Counting the uncorrected divisorial zero would give
the wrong relative bundle. Its vertical kernel is a subbundle
and has the descended nowhere-zero scalar section.

\item Write $C$ for the complete corrected critical cycle: coefficient
one on each nodal elliptic curve and each triangular double
curve, coefficient two on an original $II$ elliptic curve,
and no contribution at a good logarithm. The locally free
kernel $\mathcal V$ has the global nonzero abstract section
$\delta$. Both $\mathcal V/\mathcal O\delta$ and $\mathcal L$
are topologically trivial because $H^2(X;\mathbb Z)=0$.
The framed bundle $\mathcal V\oplus\mathcal L$ therefore
provides a specified full reference frame $\rho$ of $TX$
away from $C$. The local matrices in \ref{maps:src:D.2.44} and \eqref{maps:eq:D45.5} give
the exact relative class
\[
                  \lambda(TX,\rho)=-\operatorname{PD}(C).
                                                               \tag{D45.7}\label{maps:eq:D45.7}
\]
This specifies local coefficients, not merely the absolute
identity $\lambda=-c_2$. In particular we do not assume
that $C$ is nullhomologous; it has order five when $b=20$.

\item \emph{A relative-characteristic formula retaining the rational link.}
We prove the needed general fact. Let $N\subset X\setminus C$
be an oriented normally framed three-manifold generating
$H_3(X)=\mathbb Z/b$, and suppose that the stable tangent
frame induced by $\rho$ on $N$ is an even multiple of
$\nu\in\Omega^{\rm fr}_3=\mathbb Z/24$.
Choose a four-chain $B$ with $\partial B=bN$, transverse to
the curve strata and avoiding their finitely many singular
points. Set
\[
                     \mathcal L(N,C)=\frac{B\cdot C}{b}
                             \in\mathbb Q.                  \tag{D45.8}\label{maps:eq:D45.8}
\]
This is the actual rational number for the marked cycles,
not a residue modulo integers. It is independent of $B$,
since $[C]$ is torsion. Then
\[
       \operatorname{lk}(N,\operatorname{PD}P)
                          =-\frac12\mathcal L(N,C)
                             \quad\text{in }\mathbb Q/\mathbb Z.
                                                               \tag{D45.9}\label{maps:eq:D45.9}
\]

\item Here is a proof that includes the missing parity. Remove a
point away from all chains. Write $b=2^a q$ with $q$ odd.
Join $q$ parallel copies of $N$ by normally framed tubes.
Kaplan's prescribed-spin construction, followed by deletion
of its zero-handle and reversal, gives a cobordism to a
sphere using only two-handles
\cite[Definition 1.7, Remark 1.8 and Theorem 3.1]{mapsKaplan79}.
As in \ref{maps:src:D.2.44}, its attaching disks and normal four-frames can
be realized in the punctured simply connected six-manifold:
general position gives disjoint disks and the spin condition
is exactly their $\pi_1SO(4)$ extension obstruction.
Thus the outgoing sphere is actually embedded and represents
$q[N]$, the two-primary generator.
Now choose Zhubr's admissible three-frame $v$ for this
particular embedded sphere, as permitted by his Theorem 1
\cite{ffZhubr13}, and put $\xi=v^\perp$. Its spin framing
obstruction is $o_4(\xi)=P$, and
$\lambda(\xi)=2o_4(\xi)$ with his normalization.
Choose a spin frame $\tau$ of $\xi|_N$.
It extends over the tube and two-handle trace: the relative
obstruction groups are $\pi_0\mathrm{Spin}(3)$ and
$\pi_1\mathrm{Spin}(3)$, both zero. Hence the framings
induced by $v+\tau$ on the incoming and outgoing manifolds
are stably bordant. On the outgoing sphere, admissibility
means even framed parity, as normalized in \ref{maps:src:D.2.44} and
\cite[Section 5]{ffZhubr13}. Since $q$ is odd,
$[N]_{v+\tau}$ is even as well.

\item The complement of the real codimension-four curve strata
is simply connected, by general position applied to disks.
Both $\rho|_N$ and $v+\tau$ therefore induce the same spin
structure. Their difference lifts to $\mathrm{Spin}(6)$.
Let $a$ be its normalized degree-three integral class and
$k=\langle a,[N]\rangle$. The stable $J$ image of the
$\pi_3\mathrm{Spin}(6)$ generator is $\nu$, so the two
even framed parities imply $k$ is even. The integral
transgression of $\lambda$ gives, with consistent
transition-map direction, the identity of relative classes
\[
 \lambda(TX,\rho|_N)=2o_4(\xi,\tau)+\delta a
                         \quad\text{in }H^4(X\setminus\{p\},N;\mathbb Z).
                                                               \tag{D45.10}\label{maps:eq:D45.10}
\]
If $Q=\langle o_4(\xi,\tau),B\rangle$, evaluation gives
the identity of integers
\[
                       -B\cdot C=2Q+bk.                    \tag{D45.11}\label{maps:eq:D45.11}
\]
By the integral definition of the torsion pairing,
$Q/b\bmod\mathbb Z=\operatorname{lk}(N,\operatorname{PD}P)$.
Dividing \eqref{maps:eq:D45.11} and using the evenness of $k$ proves
\eqref{maps:eq:D45.9}. This is not division of a class in an even cyclic
group. Its integer correction was fixed by actual spin
frames before division. The proof uses Zhubr's proved
Theorem 1, and does not use his conjectural general
change-of-three-frame formula. It also permits a nonzero
odd-primary $\lambda$.

\item \emph{The finite-core generator and its reference parity.}
For a single good $II$ logarithm the regular three-torus
can be an even class in $H_3(X)$. We use the following
actual core three-manifold instead. The central four-manifold
is an elliptic-torus bundle over the rank-two lattice
\[
 \langle\delta,\ell,g\mid mg=r\delta+n\ell\rangle.
                                                               \tag{D45.12}\label{maps:eq:D45.12}
\]
Its elliptic monodromy is $B$ on $g$, trivial on the other
two generators, where $B$ has order six and
$\det(B-I)=1$. With $\sigma=n-me$, define the primitive
integer character
\[
 \chi(\delta)=1,\qquad \chi(\ell)=-\sigma r,
                       \chi(g)=-\sigma re.                 \tag{D45.13}\label{maps:eq:D45.13}
\]
The relation in \eqref{maps:eq:D45.12} is respected. Let $N$ be the
connected level set of the resulting map of the core to
$S^1$, cooriented by $\chi$. Put
$h=g+\sigma re\delta$ and $k_0=\ell+\sigma r\delta$.
These lie in $\ker\chi$ and satisfy $mh=nk_0$.
Since $\gcd(m,n)=1$, their primitive generator has elliptic
monodromy $B^\sigma$: an integer inverse of $n$ modulo
$m$ is congruent to $\sigma$ modulo six. Thus $N$ is an
actual affine elliptic mapping torus, with
\[
                        H_1(N;\mathbb Z)=\mathbb Z.        \tag{D45.14}\label{maps:eq:D45.14}
\]
Translations do not change this Wang sequence.

\item For completeness, any actual tangent parallelization of
a closed oriented three-manifold $M$ has framed parity
\[
        [M,\tau]_{\rm fr}\bmod2
                    =1+\dim H_1(M;\mathbb F_2).             \tag{D45.15}\label{maps:eq:D45.15}
\]
Use Kaplan's spin-compatible filling $W$ with one zero-handle
and $j$ two-handles. Its tangent bundle is trivial, since
its spine is two-dimensional and $w_2=0$. On the boundary
use the outward normal followed by $\tau$.
The relative stable framing obstruction is $\lambda$;
its $J$ image is $[M,\tau]$. Modulo two it is the relative
$w_4$, hence the relative Euler class, whose evaluation is
$\chi(W)=1+j$. The mod-two intersection matrix of the
spin filling is alternating, so has even rank. Its
cokernel is $H_1(M;\mathbb F_2)$; hence
$j\equiv\dim H_1(M;\mathbb F_2)$. This proves \eqref{maps:eq:D45.15}
while retaining the prescribed spin structure.

\item The reference frame on the finite-core $N$ is the
stabilization of an actual tangent frame. Indeed $\delta$
is transverse to $N$ by $\chi(\delta)=1$. For the global
complex vertical frame $(\delta,v_2)$, project
$J\delta,\operatorname{Re}v_2,\operatorname{Im}v_2$ to $TN$
along $\delta$. They form a tangent frame. The normal
frame is $\delta$ together with the two real directions
of the corrected target line $\mathcal L$. Projection
and horizontal splitting are upper triangular homotopies.
Equation \eqref{maps:eq:D45.14}, then \eqref{maps:eq:D45.15}, therefore gives even
reference parity. This conclusion is not inferred from
the even-degree torus cover of $N$.

\item \emph{The actual marked cycles and the unbounded chain identity.}
Use the same integral regular-fibre coordinate as D.2.41:
\[
                \Omega=a e_1^*e_2^*-b\delta^*\gamma^*.
                                                               \tag{D45.16}\label{maps:eq:D45.16}
\]
The full degree-two monodromy cokernel is $\mathbb Z$,
primitively measured by $\Omega$. Choose an integral
fibre two-cycle $s$ with $\Omega(s)=1$; its image generates
$H_2(X)=\mathbb Z/b$. The nodal and triangular preimages
are those proved in \ref{maps:src:D.2.44}. At an original $II$ site put
\[
 z=(B-I)^{-1}q_0\in\mathbb Z^2,\qquad
 \ell=(-z,0,1),\qquad [E_{II}]=\delta\wedge\ell,             \tag{D45.17}\label{maps:eq:D45.17}
\]
where $q_0$ is the full elliptic long-column vector.
The full matrix fixes $\ell$, not just its quotient.
The principal $\delta,\ell$ torus through the elliptic
cusp is the reduced critical elliptic curve; its regular
torus specializes with degree one. The comparison takes
place within that cap. Thus $2\delta\wedge\ell$ is the
literal critical preimage, with the complex torus orientation.
For OS45/II, $\ell=(-8,0,0,1)$; for the OS55/II rows,
$\ell=(0,0,0,1)$.

\item Transport all these preimages along the fixed base paths
and let $K=\Omega(C)$. The exact integers are
\[
\begin{array}{c|c|c|r|c}
 \text{profile}&(a,b)&\text{zero and }II\text{ choice}&K&
                                  \operatorname{ord}[C]\\\hline
 45/II&(1,8)&2,3;\ \mathrm{original}&-32&1\\
 45/II&(1,8)&2,3;\ \mathrm{log}&-16&1\\
 55/II&(5,4)&3;\ \mathrm{original}&-108&1\\
 55/II&(5,4)&3;\ \mathrm{log}&-100&1\\
 55/II&(1,20)&3;\ \mathrm{original}&-68&5\\
 55/II&(1,20)&3;\ \mathrm{log}&-28&5\\
 55/\mathrm{default}&(1,20)&2&-108&5\\
 55/\mathrm{default}&(1,20)&3,4&-68&5.
\end{array}                                                    \tag{D45.18}\label{maps:eq:D45.18}
\]
The difference $C-Ks$ bounds a chain in the regular
complement, after the specified local specialization
chains are included. This uses the full integral
monodromy image, whose five nonzero Smith factors are
units; it does not replace an image by its saturation.
For a logarithmic pattern all these chains avoid its
finite core, since that cap has no critical curve.

\item Here is the exact all-parameter relation that fixes
both the generator and the rational linking.
Every monodromy fixes $\delta$. Its quadratic bar
correction on that single loop is zero, including
the transported last peripheral word, and the
peripheral graph maps satisfy $\sum_iJ_i\delta=0$.
Take the finite normal Thom column for $\delta$,
add $m$ times the ordinary-clutch $\delta$ column,
and add $m$ times the other filled-cap $\delta$ sweeps.
The graph terms cancel. Its actual fibre part is
\[
 -\delta\wedge(r\delta+n\ell)+me\delta\wedge\gamma,
 \qquad \Omega=b(n-me)=b\sigma.                            \tag{D45.19}\label{maps:eq:D45.19}
\]
The remaining difference from $b\sigma s$ is an
integral monodromy relation, supported away from $N$.
All other cap sweeps also bound away from $N$.
The finite normal Thom chain is the disc bundle
over its actual central loop. Its intersection
with the level set $N$ is its character value,
which here is $\chi(\delta)=1$.
Thus \eqref{maps:eq:D45.19} gives the actual pairings
\[
 \operatorname{lk}(N,s)=\frac{\sigma}{b}\pmod{\mathbb Z},
 \qquad \mathcal L(N,C)=\frac{\sigma K}{b}.                 \tag{D45.20}\label{maps:eq:D45.20}
\]
The first is a unit pairing, proving that $N$ is a
full $H_3$ generator, including when $b=20$.
The fifth affine-meridian column, if used in another
Smith bounding chain, pairs by its actual value
$\chi(g)=-\sigma re$, not zero. This retained term
makes the full-chain comparison agree with \eqref{maps:eq:D45.20}.
There is no additional factor of $m$: the covering
order is already present in the boundary norm column.

\item For an original-only pattern use the ordinary
$T^3=\langle e_1,e_2,\gamma\rangle$ and $e=\pm1$.
The same peripheral relation, now using its ordinary
normal Thom chain, has fibre value $-be$ and
intersection one with $T^3$. Therefore
\[
 \operatorname{lk}(T^3,s)=-\frac e b\pmod{\mathbb Z},
 \qquad \mathcal L(T^3,C)=-\frac{eK}{b}.                    \tag{D45.21}\label{maps:eq:D45.21}
\]
It is again a full generator with even reference parity,
as proved in \ref{maps:src:D.2.44}. If $b=20$, these equations can
equivalently be computed by bounding $5C$ and dividing
its literal intersection by five. No nullhomology
of $C$ is asserted.

\item Substitute \eqref{maps:eq:D45.20} or \eqref{maps:eq:D45.21} into the already
proved relative formula \eqref{maps:eq:D45.9}. Since $K$ is an
even integer in every row of \eqref{maps:eq:D45.18}, cancellation
of the displayed unit pairing gives
\[
                       P=-K/2\quad\text{in }\mathbb Z/b.
                                                               \tag{D45.22}\label{maps:eq:D45.22}
\]
Here the division was justified by \eqref{maps:eq:D45.10}--\eqref{maps:eq:D45.11}
using the actual cycles and framings. It is not
division of the congruence $2P=\lambda$.
The table gives $P=2,0,14$ as in \eqref{maps:eq:D45.3}, and the
absolute identity $\lambda=-[C]$ gives $0,0,8$.

\item \emph{The complete fundamental-class invariant and recognition.}
The global auxiliary-line comparison and localization
of \ref{maps:src:D.2.44} apply over $\mathbb Z/4$ and $\mathbb Z/8$,
respectively. All actual original-$II$ coefficients
are zero and all finite-$II$ normal coefficients
are $bn$ and $-br$; hence the same mod-$b$ class
extends for every parameter in \eqref{maps:eq:D45.2}.
For $b=20$ apply that calculation separately over
$\mathbb Z/4$ and $\mathbb Z/5$. The $C_4$ action
has only the rational fixed curve and the $I_4$
elliptic fixed curve, with characters $2,0$.
The $C_5$ action similarly has the rational and
$I_5$ curves, with characters $0,0$; all remaining
stabilizers are prime to the relevant coefficient
order. The actual Tate normal weights have square
one modulo four, and weight two modulo five.
The cube formula is zero in each coefficient ring:
on the rational curve it is $-3$ times the square
of its point character, and on the elliptic curve
the possible mixed degree-one term has square zero.
Thus the full fundamental-class invariant is zero,
since $H_6(K(\mathbb Z/b,2);\mathbb Z)=\mathbb Z/b$
is detected by the cube for $b=4,8,20$
\cite[Sections 3.11--3.14]{manZhubr75}.

\item The finite data therefore agree with the first two
surgery models of \ref{maps:src:D.2.44}. For $b=20$, the two
primary parts of $P=14$ are $2\in\mathbb Z/4$
and a nonzero element of $\mathbb Z/5$.
The inclusion
$K(\mathbb Z/4,2)\vee K(\mathbb Z/5,2)
\to K(\mathbb Z/4,2)\times K(\mathbb Z/5,2)$
is a homology equivalence: all cross tensor and
Tor terms vanish because their positive-degree
homology has coprime primary orders. Both spaces
are simply connected, so it is a homotopy equivalence.
The reduced spin-bordism invariant consequently
splits into these two summands, as in \ref{maps:src:D.2.38}.
The complete invariants agree with
$Z_{4,1}\#Z_{5,1}$, and the free third Betti numbers
also agree. Zhubr's Theorems 3 and 5
\cite{manZhubr75} give all the oriented
diffeomorphisms in \eqref{maps:eq:D45.3}.

\item \emph{Exact computations and their bounds.}
The local differential identities are checked in
\path{code/III/maps_D2_II_frame_checks.py}.
The line/support, full angular-image, actual
adapted-matrix and arbitrary finite-normal
coefficients for all eighteen patterns are checked in
\path{code/III/fourfolds_D2_remaining_AH.py}.
This comprises fourteen initial rows, $1056$ exact
semicharacter checks, $62$ full adapted matrices,
four symbolic arbitrary-twist finite normal rows
and four original-$II$ choices.
The actual curve vectors, complete Smith solutions,
finite root-loop characters and literal intersections
are retained by
\path{code/III/manifolds_extended_critical_cycles.py}.
Its $158$ full Smith replays comprise the two original
clutch signs and, for each of the four logarithmic
patterns, $k=1,2,3$, $r=0,1$, $e=-1,0,1$ and
$\sigma=\pm1$. This explicitly bounded replay is
supplemented by the eleven constant certificates in
\path{records/manifolds_extended_critical_symbolic.json}:
integral preimages of $-\delta\ell-bs$,
$\delta\gamma+bs$, and $C-Ks$, together with the
zero bar corrections and exact peripheral graph
identity. Their linear combinations prove \eqref{maps:eq:D45.19}
for every integer parameter in \eqref{maps:eq:D45.2}, with no
unproved finite search bound.\hfill$\square$
\end{enumerate}
\end{proof}


\begin{theorem}[arbitrary whole-fibre initial pairs and weights]\label{maps:src:D.2.46}
\StatementRecord{maps:src:D.2.46}{proved}{Historical events: Claude Fable 5.2 (Anthropic), effort high, 16 November 2026: carried as background without fault found. GPT-6 Astra (OpenAI), ChatGPT Pro mode, effort Pro, S6\_09\_PAPERS, 17 November 2026: targeted proof check. Also checked in this writing. Pre-continuation writing events: GPT-6 Astra (OpenAI), ChatGPT Pro mode, effort Pro, S6\_10\_PAPERS, 18 November 2026: targeted proof check. Continuation metadata events: GPT-6 Astra (OpenAI), ChatGPT Pro mode, effort Pro, S6\_10\_PAPERS continuation, 19 November 2026: used and rechecked in this run; status-vocabulary reconciliation; metadata reconciliation only; no new mathematical check. \emph{Pre-continuation recorded proof scope.} GPT-6 Astra (OpenAI, ChatGPT Pro mode), effort Pro, 2026-11-17 UTC; Claude Fable 5.2 (Anthropic), effort high, 16 November 2026; also checked in this writing by GPT-6 Astra (OpenAI), ChatGPT Pro mode, effort Pro, 18 November 2026}{the hypotheses, notation and definitions stated here; \cref{maps:src:D.2.19}; \cref{maps:eq:D.BC1}; \cref{maps:eq:D.BC14}; \cref{maps:eq:D.BC4}; \cref{maps:src:D.2.29}; \cref{maps:eq:D.BC5}; \cref{maps:src:D.2.22}; \cref{maps:src:D.2.23}; \cite[Section 11.8]{manSS10}; \cref{maps:src:D.2.37}; \cref{maps:src:D.2.43}; \cref{maps:src:D.2.31}; \cref{maps:src:D.2.32}; \cref{maps:src:D.2.40}; \cref{maps:eq:D.GI9}; \cref{maps:src:D.2.41}}


Fix any of the $21$ marked rational elliptic profiles with
$\operatorname{MW}=\mathbb ZG$ in \ref{maps:src:D.2.19}, and write
$\langle G,G\rangle=1/d_0$. Let
\[
 P=aG,\quad Q=bG,\quad a,b>0,
 \qquad \overline P_i=0\ \hbox{or}\ \overline Q_i=0
                         \quad\hbox{at every fibre}.       \tag{D.GI1}\label{maps:eq:D.GI1}
\]
Choose any nonnegative integral whole-fibre linearization
weights $w_i$, of total $ab/d_0$, supported on multiplicative
or regular fibres. Distinct regular support points are allowed.
At a positive weight use the marked periodic Mumford filling.
At a zero-free multiplicative fibre choose either its original
cap or its marked Mumford cap. At a finite good-reduction
additive fibre choose either its original cap or any free affine
logarithmic cap. Retain the original cap at $I_n^*$.
Arbitrary regular logarithms and an arbitrary ordinary clutching
integer are allowed.

There is a finite exact integral algorithm for the fundamental
group and every integral homology group of every such threefold.
It includes pairs not obtained by iterating any of the $172$
initial rows, several zero fibres, noncyclic toric period
lattices, and all their boundary torsion. Its implementation is
\path{code/III/manifolds_general_initial.py}; no bound on
$a,b,w_i$, logarithmic orders or affine numerators occurs in the
algorithm or proof.

Here are the complete data of the algorithm. Use the lattice
$\Lambda=\mathbb Z\langle e_1,e_2,\delta,\gamma\rangle$.
The recorded elliptic monodromies are $B_i$, and the section
$G$ has integral cocycle $g_i$. Put $p_i=ag_i$, $q_i=bg_i$,
choose any rational $z_i$ with $(B_i-I)z_i=q_i$, and set
\[
 r_i=((p_i)_2,-(p_i)_1)B_i,\qquad
 A_i=\begin{pmatrix}
 B_i&0&q_i\\ r_i&1&r_i z_i+w_i\\0&0&1
 \end{pmatrix}.                                           \tag{D.GI2}\label{maps:eq:D.GI2}
\]
The block sizes are $2,1,1$. At a regular zero use
$A_i=I+w_i\delta\gamma^*$. At an ordinary or regular
logarithm the boundary matrix is $I$. These are integral
matrices, their ordered product is $I$, and they are the
actual marked monodromies.

Apply the integral bar comparison \eqref{maps:eq:D.BC1}--\eqref{maps:eq:D.BC14} to these
matrices. Thus the regular-piece degree-two group is
$C\oplus K$, with its full unsaturated fibre presentation
and its transported last-boundary marking. At each finite or
regular logarithm use all five columns \eqref{maps:eq:D.BC4}, with its
actual order $m$ and fixed power $v$, and the exact freeness
test of \ref{maps:src:D.2.29}--\ref{maps:src:D.2.30}. At every other cap use the full fixed-normal
columns \eqref{maps:eq:D.BC5} and the following fibre kernel $F_i$ in
$\Lambda^2\Lambda$. Here $R_i=\Lambda^2A_i-I$ and
$\operatorname{sat}$ denotes saturation only where it is
explicitly written.
\[
\begin{array}{c|c|c}
\hbox{cap}&F_i& h_i=\operatorname{rank}H^4(U_i)\\ \hline
w_i>0\hbox{ at }I_n, n>0&
 \operatorname{im}R_i+\Lambda^2\operatorname{Fix}A_i&nw_i\\
\hbox{rank-one Mumford}&\operatorname{sat}(\operatorname{im}R_i)&\rho_i\\
\hbox{original, }\overline P_i=0&
 \operatorname{sat}(\operatorname{im}R_i)&\rho_i\\
\hbox{original additive, }\overline P_i\ne0&\ker\omega_i&s_i\\
\hbox{original }I_n,\ \overline P_i\ne0&
 \langle u_1\delta,u_1\ell_i,\delta\ell_i\rangle&n .
\end{array}                                                \tag{D.GI3}\label{maps:eq:D.GI3}
\]
Products in the last row are exterior products.
For a rank-one Mumford cap, $\rho_i$ is the unique nonzero
Smith factor of $A_i-I$. For a zero-Chern original cap it is
the number of orbits of translation by $Q$ on the Kodaira
components. The actual marked permutations were given in
\ref{maps:src:D.2.22}. In the last two rows, $Q$ is narrow, so $z_i$ may
and must be integral, and
\[
 \ell_i=(-z_i,0,1),\qquad
 \omega_i=(1,0,(z_i)_2,0,-(z_i)_1,0).
                                                               \tag{D.GI4}\label{maps:eq:D.GI4}
\]
Here $u_1$ is the primitive elliptic vanishing circle and
$s_i$ is the number of Kodaira components. At a positive
regular zero the second row applies, with $\rho_i=w_i$.

Let $\mathcal B$ be the resulting full integral degree-two
matrix, in $\Lambda^2\Lambda\oplus K$, and write
$\operatorname{coker}\mathcal B=\mathbb Z^g\oplus T$.
Let $E$ be the regular piece with all finite and regular
logarithmic caps already filled. Compute $H_1(E)$ from its
literal graph-bundle presentation and the relations
$\tau_i^{m_i}=v_i$. Compute $H_1(X)=\mathbb Z^r\oplus A_1$
by additionally imposing, at each other cap, $\tau_i=1$
and the following fibre relations:
\[
\begin{array}{c|c}
\hbox{cap}&\hbox{killed subgroup of }\Lambda\\\hline
\hbox{positive rank-two Mumford}&\operatorname{Fix}A_i\\
\hbox{rank-one Mumford or zero-Chern original}&
                  \operatorname{sat}(\operatorname{im}(A_i-I))\\
\hbox{nonzero-Chern additive original}&
                  \mathbb Z\langle e_1,e_2,\delta\rangle\\
\hbox{nonzero-Chern multiplicative original}&
                  \mathbb Z\langle u_1,\delta\rangle .
\end{array}                                                \tag{D.GI5}\label{maps:eq:D.GI5}
\]
Set
\[
 q=\sum_{\text{nonlog caps}}h_i-b_1(E)+r,\qquad
 b_2=g+q,\qquad
 \chi=2\sum_{I_n,\ n>0}nw_i,\qquad
 b_3=2+2b_2-2r-\chi .                                    \tag{D.GI6}\label{maps:eq:D.GI6}
\]
Then $q\ge0$ and the complete answer is
\[
\begin{array}{c|c}
j&H_j(X;\mathbb Z)\\\hline
0,6&\mathbb Z\\
1&\mathbb Z^r\oplus A_1\\
2&\mathbb Z^{b_2}\oplus T\\
3&\mathbb Z^{b_3}\oplus T\\
4&\mathbb Z^{b_2}\oplus A_1\\
5&\mathbb Z^r .
\end{array}                                                \tag{D.GI7}\label{maps:eq:D.GI7}
\]
For the nonabelian fundamental group, retain four commuting
fibre generators, all peripheral generators $\tau_i$, the
relations $\tau_i x\tau_i^{-1}=A_i x$ and
$\prod_i\tau_i=1$, all finite-power relations, and
\eqref{maps:eq:D.GI5}. This is a complete finite presentation, not just
its abelianization.

Except in the case described next, the entire vertical
subgroup $\langle e_1,e_2,\delta\rangle$ is killed and this
presentation reduces to the Seifert group of \ref{maps:src:D.2.23}, with
the actual gamma numerators. The exceptional possibility is
\[
\begin{gathered}
 \mathrm{OS49/II}=(II,IV^*,I_2),\quad a\text{ odd},\quad
 b\text{ even},\quad\text{both additive caps logarithmic},\\
 w_{I_2}=0,\quad I_2\text{ filled by Mumford},\quad
 \text{all positive weights at regular fibres}.
\end{gathered}                                             \tag{D.GI8}\label{maps:eq:D.GI8}
\]
The compatibility condition \eqref{maps:eq:D.GI1} remains required; it
can impose an additional factor $3$ on $a$ or $b$.
Before the logarithmic power relations the remaining
vertical group in \eqref{maps:eq:D.GI8} is $(\mathbb Z/2)^2$. The two
additive monodromies act by inverse order-three
automorphisms. With no additional regular logarithms,
the group is the split semidirect product of this Klein
four group by the two-fibre Seifert group. With arbitrary
regular logarithms use the full presentation above.
Every group in \eqref{maps:eq:D.GI8} surjects onto $\mathbb Z/3$, so none
is simply connected.

\end{theorem}


\begin{proof}
\begin{enumerate}[label=\arabic*.,leftmargin=*]
\item \emph{The initial data are actual section-line data.}
Let $M=\mathcal O(P-O)$ and
$J=M\otimes(T_Q^*M)^{-1}$. On a smooth elliptic fibre,
translation acts trivially on $\operatorname{Pic}^0$, so
$J$ is a vertical line bundle. At a singular fibre its
component multidegree is the difference of that of $M$
and its $Q$-translate. It is zero if $P$ is narrow, and
also if $Q$ is narrow. Thus \eqref{maps:eq:D.GI1} makes every local
multidegree zero. The intersection matrix of each
Kodaira fibre has kernel generated by its full
multiplicity vector; hence $J$ is pulled back from the
base. The height calculation in \ref{maps:src:D.2.19}, equivalently
the height pairing with all local mixed corrections
zero \cite[Section 11.8]{manSS10}, gives
\[
                         J\simeq f^*\mathcal O(ab/d_0).
\]
This proves integrality of $ab/d_0$ in the compatible
sector. A base section with precisely the chosen
nonnegative divisor gives the required linearization.
This does not introduce a component-supported vertical
correction to $M$.

\item For completeness, the monodromy formula is also an
integral geometric formula. The rigidified Poincare
bundle has elliptic section cocycles $p_i,q_i$; its
commutator pairing gives the row
$r_i=((p_i)_2,-(p_i)_1)B_i$. A local logarithm of $Q$
gives the central entry $r_i z_i$, and a zero of the
linearization adds $w_i\delta\gamma^*$. This is the
same marked local calculation used in \ref{maps:src:D.2.19} and
\ref{maps:src:D.2.37}, now before restricting either section coefficient.
Changing $z_i$ by a fixed vector changes no entry:
$r_i$ annihilates $\ker(B_i-I)$.

\item If $Q$ is narrow, $z_i$ is integral. If instead $P$
is narrow, write $p_i=(B_i-I)y_i$ with $y_i$ integral.
For $J_2=\left(\begin{smallmatrix}0&1\\-1&0\end{smallmatrix}\right)$,
symplectic invariance gives
\[
 r_i z_i=-y_i^t(B_i-I)^tJ_2B_i z_i
                         =y_i^tJ_2q_i\in\mathbb Z.
\]
Thus \eqref{maps:eq:D.GI2} is integral in both cases. At zero weights
its ordered product is the central matrix with entry
$-ab/d_0$; the local Poincare commutators are bilinear
in $a,b$. This identity is checked already for the
single generator pair on each of the $21$ recorded
elliptic tuples and follows for all coefficients by
bilinearity. Adding the weights changes that entry by
$\sum w_i=ab/d_0$, proving the product relation.
The implementation reconstructs \eqref{maps:eq:D.GI2} directly from
the saved cocycles and verifies integrality and product,
rather than treating a numerical monodromy output as
a topology theorem.

\item \emph{All local kernels, with their integral markings.}
At a positive $I_n$ weight, use a primitive elliptic
vanishing/long basis, followed by $\delta,\gamma$.
The nilpotent part has image in the primitive short
plane $\langle u_1,\delta\rangle$, squares to zero,
and has determinant $nw_i$ on long-to-short lattices.
Its Smith factors $s\mid t$ satisfy $st=nw_i$.
The actual linearization changes only the corresponding
integral long period; the same triangular periodic fan
is preserved by this period sublattice. The unit factors
are removed by the root-and-unit comparison of \ref{maps:src:D.2.22}.
Consequently this is exactly the cap of \ref{maps:src:D.2.43}, with
\[
 H_2(U_i)=\mathbb Z^{nw_i+3}\oplus\mathbb Z/s,
 \qquad H_3(U_i)=\mathbb Z^2,\qquad
 H_4(U_i)=\mathbb Z^{nw_i}.
\]
Its full fibre kernel is the first row of \eqref{maps:eq:D.GI3},
\emph{without saturation}. The two primitive short
normal tori bound their actual orbit discs. In
particular its angular-area torsion relation, and
the other boundary torsion class which maps onto
$\operatorname{Tor}H_2(U_i)$, remain distinguished.
The component meridians are the literal rays $(v,1)$
of \ref{maps:src:D.2.43}; thus no change from the original period
marking to a saturated lattice has occurred.

\item At rank one write $A-I=\rho\alpha\beta^*$ with
$\alpha$ and $\beta^*$ primitive. The local toric
model is the $\rho$-component nodal-cycle model
times the transverse torus, with the same unit
comparison and primitive vanishing disc as in
\ref{maps:src:D.2.37}. Its homology is free, its fourth rank is
$\rho$, and its full fibre kernel is precisely
the saturation of $\operatorname{im}R_i$.
The full fixed lattice supplies the normal tori.
This argument includes a regular zero, for which
$\alpha=\delta$, $\beta^*=\gamma^*$ and $\rho=w_i$.

\item For an original zero-Chern cap, the line bundle is
trivial on the Kodaira normalization and the prescribed
unit lift has no residual meridian winding. Its
retraction is $S^1_\delta$ times the mapping torus
of the actual component permutation by $Q$.
For an additive fibre the curve configuration is
a tree; for $I_n$ its dual graph is a circle.
The latter has no extra line-bundle winding: the
component degrees are all zero, the translation
lift on the compact connected cycle differs from
the rigidified lift by a constant unit, and the
nonvanishing base linearization has a logarithm.
The component permutation therefore contributes
one free fourth class per orbit. The explicit
normalization or tree-chain complex has free
homology and gives exactly the third row of
\eqref{maps:eq:D.GI3}, with all fixed normal tori. This is the
permutation calculation of \ref{maps:src:D.2.31}; it uses the
actual residues of $bG$, not primitivity of $b$.

\item Suppose next that the original cap has
$\overline P_i\ne0$. Then $Q$ is narrow, and the
Chern vector of $M$ on the components is
\[
                         c=e_{\overline P_i}-e_0.
\]
It is primitive even if the component of $P$
is not a generator of the component group.
Its pairing with the Kodaira multiplicity vector
is zero. The integral logarithm of $Q$ and the
rigidified unit comparison contract the translation
while retaining this Chern class. On a tree the
unit comparison has a logarithm because $H^1=0$.
On a nodal cycle its final clutch exponent is
$\sum c_i=0$, so the same comparison has no
dual-circle winding. Thus the actual retract is
the principal two-torus bundle with Chern pair
$(c,0)$, or $Y_c\times S^1_\gamma$ in the cycle
case, precisely as in \ref{maps:src:D.2.32}--\ref{maps:src:D.2.33}.

\item The primitivity of $c$, not the stronger and
unneeded assumption that $\overline P_i$ generates
the component group, makes all these local homology
groups free. In the tree case the regular elliptic
fibre maps to the primitive multiplicity vector in
$\ker c$. Its exact area functional is $\omega_i$,
giving the fourth row of \eqref{maps:eq:D.GI3}. In the cycle case,
the three regular fibre images
\[
                       u_1u_2,\quad u_2\delta,\quad u_2\ell_i
\]
remain primitive: respectively they project to the
primitive all-one component vector, the fibre-circle
times the dual-circle generator, and the dual circle
times the second principal circle. This proves the
last row of \eqref{maps:eq:D.GI3} even when $A_i-I$ has nonunit
Smith factors. Its fixed lattice is exactly
$\langle u_1,\delta,\ell_i\rangle$.

\item The $\delta$ and $\ell_i$ normal tori bound with
zero fibre correction. The $u_1$ normal torus in
the cycle case can differ from a bounding normal
torus by an integral multiple of $u_2\delta$.
As in \ref{maps:src:D.2.40}, this correction is zero in the
\emph{full} global fibre coinvariant, including its
torsion. Indeed the relations on $\delta\wedge\Lambda$
are the actual elliptic columns $(B_i-I)\mathbb Z^2$,
and their common Smith factors are $(1,1)$ for
each of the $21$ profiles. This assertion is
independent of $a,b,w_i$ and is retained as an
explicit constant certificate. Hence all the normal
columns in the theorem have their actual global
marking; no rational replacement is involved.

\item The same local retractions give the fundamental
group kernels \eqref{maps:eq:D.GI5}. All meridians of these
nonlogarithmic caps bound their normal discs.
At a finite logarithm the principal meridian is
$\tau_i^m v^{-1}$, with its full integral normal
line torsion and the five Thom columns already
proved in \ref{maps:src:D.2.29}--\ref{maps:src:D.2.30}. These local statements prove
the claimed van Kampen presentation.

\item \emph{The normal closure and the exceptional profile.}
There is at least one positive weight. Its primitive
vanishing lattice contains $\delta$, so $\delta=1$
in the global fundamental group. Modulo $\delta$,
a zero-free rank-one $I_n$ Mumford cap kills
\[
                   \frac{n}{\gcd(n,a)}\,u_i,              \tag{D.GI9}\label{maps:eq:D.GI9}
\]
where $u_i$ is its primitive elliptic vanishing vector.
To see the coefficient, in its Tate basis the
vanishing vector before primitivization is
$nu_i+aH_i\delta$, with
$\gcd(H_i,n)=1$. The generator component of $G$
has that unit residue on every recorded $I_n$.
If $P$ is narrow the projection in \eqref{maps:eq:D.GI9} is
already primitive. If $P$ is non-narrow, the
compatibility condition makes $Q$ narrow; changing
its integral logarithm changes the long coordinate
but not this projected vanishing line.

\item The elliptic normal closure of these lines under
the finitely many $B_i$ depends only on
$a\bmod d_0$. The original star caps additionally
kill the full elliptic lattice after $\delta$ is
killed: a nonzero Chern vector kills it directly,
while a narrow $P$ supplies its integral graph
trivialization. Thus there are precisely $284$
constant residue cases to check. Their recorded
Hermite matrices are full $\mathbb Z^2$ in every
case except OS49/II with $a$ odd, where the matrix
is $2I_2$. The certificate also computes every
one of the $912$ actual local projected vanishing
lines and verifies \eqref{maps:eq:D.GI9} directly; this is an
exhaustive residue calculation, not a bounded
search over unbounded section pairs.

\item In that sole residual case, a positive $I_2$
weight adds its primitive short circle, and an
original $I_2$ cap adds the primitive elliptic
vanishing circle because its Chern vector is
nonzero. Either one and its $B_i$ translates
generate $\mathbb Z^2$. An original additive
cap does the same. Hence only \eqref{maps:eq:D.GI8} remains.
The two additive matrices modulo $2$ have order
three, are inverse, and $B_i-I$ is invertible
modulo $2$. Since $b$ is even, their gamma
columns vanish modulo $2$. Their fixed integral
power vectors consequently have even elliptic
coordinates. With no extra regular logarithms,
all power and base-product relations thus lie
in the Seifert factor, proving the asserted
split semidirect product. Additional regular
affine powers can have nonzero Klein-four
coordinates, which is why the full presentation
is retained in that case. In all cases set
the entire fibre group and additional regular
peripherals to $1$, and send the two additive
peripherals to opposite generators of $\mathbb Z/3$.
Their orders are multiples of $6$ and $3$.
This is a well-defined surjection and proves
the non-simply-connected assertion.

\item \emph{Integral relative gluing, including cap torsion.}
The bar comparison of \ref{maps:src:D.2.41} applies to any integral
tuple \eqref{maps:eq:D.GI2}; its proof made no scalar-iterate
assumption. It gives $H_2$ of the regular piece
with all fibre torsion and every fixed, norm and
last-peripheral correction. Filling the finite
caps quotients by precisely their five columns,
so it gives the stated $H_2(E)$ presentation.

\item Every remaining cap retracts to a complex of
dimension at most four and has free $H^4$ of
rank $h_i$. For the general cusp this follows
from \ref{maps:src:D.2.43}, which gives free $H_3$ as well as
free $H_4$; for the other caps it follows from
the explicit free groups above. Excision and
Poincare--Lefschetz duality give
\[
 H_2(X,E)=\bigoplus_iH^4(U_i)=\mathbb Z^{\sum h_i},
                    \qquad H_1(X,E)=0 .                  \tag{D.GI10}\label{maps:eq:D.GI10}
\]
The group $H_3(U_i,\partial U_i)=H^3(U_i)$ need
\emph{not} be free: it contains the torsion dual
to $H_2(U_i)$ at a general cusp. Its boundary
image is nevertheless exactly
$\ker[H_2(\partial U_i)\to H_2(U_i)]$.
\ref{maps:src:D.2.43} identifies that image with the two literal
normal sweeps and the nonsaturated first fibre
kernel in \eqref{maps:eq:D.GI3}. In particular the angular-area
$\mathbb Z/s$ is included, while the other
$\mathbb Z/s$ maps isomorphically to cap torsion.
This accounts for the torsion relative boundary
classes with their actual marking; freeness of
cap $H_2$ is not assumed.

\item It follows that the image of $H_3(X,E)$ in
$H_2(E)$ is exactly the remaining columns of
$\mathcal B$. The pair sequence therefore has
the exact form
\[
 0\longrightarrow\operatorname{coker}\mathcal B
 \longrightarrow H_2(X)
 \longrightarrow\ker[\mathbb Z^{\sum h_i}\to H_1(E)]
 \longrightarrow0.                                      \tag{D.GI11}\label{maps:eq:D.GI11}
\]
The final term is a free subgroup of a free
group; its rank is $q$ because
$H_1(E)\to H_1(X)$ is surjective. The sequence
splits as abelian groups, proving the degree-two
formula with precisely the torsion $T$.
All caps except positive rank-two cusps have
Euler number zero; a cusp of period index $nw_i$
has Euler number $2nw_i$. Hence \eqref{maps:eq:D.GI6} gives
the Euler number and the middle Betti number.
Integral duality and the universal coefficient
theorem give every remaining group in \eqref{maps:eq:D.GI7}.
This includes zero rational Euler number and
arbitrary first Betti number, without a Seifert
rank assumption in the exceptional sector.
\end{enumerate}
\end{proof}

The replay \path{code/III/manifolds_general_initial_audit.py}
stores all $284$ normal-closure cases, $912$ local
rank-one projections, $21$ generator height products
and elliptic unit lattices, $90$ complete overlaps
with \ref{maps:src:D.2.41}, and $148$ full Smith examples in the
larger parameter sector. For example, OS45/default
with $P=2G$, $Q=8G$, weight $2$ at $I_8$, all other
weights zero, all Mumford caps and ordinary clutch
$e=1$ is simply connected and has
\[
 H_2=\mathbb Z^{15}\oplus\mathbb Z/8,\qquad
 H_3=\mathbb Z/8,\qquad H_4=\mathbb Z^{15},\qquad\chi=32.
\]
Its rank-two cusp has period index $16$, so it is
outside the scalar iterations of the natural
index-one initial catalogue. This theorem concerns
the full stated whole-fibre sector. Component-supported
linearizations outside \eqref{maps:eq:D.GI1}, arbitrary extra
vertical modifications, logarithmic star Mumford
caps, and general smooth classification of its
nonsphere outputs are not asserted here.


\begin{theorem}[the general triangular-cusp component intersection form]\label{maps:src:D.2.47}
\StatementRecord{maps:src:D.2.47}{proved}{Historical events: Claude Fable 5.2 (Anthropic), effort high, 16 November 2026: carried as background without fault found. GPT-6 Astra (OpenAI), ChatGPT Pro mode, effort Pro, S6\_09\_PAPERS, 17 November 2026: targeted proof check. Also checked in this writing. Pre-continuation writing events: GPT-6 Astra (OpenAI), ChatGPT Pro mode, effort Pro, S6\_10\_PAPERS, 18 November 2026: targeted proof check. Continuation metadata events: GPT-6 Astra (OpenAI), ChatGPT Pro mode, effort Pro, S6\_10\_PAPERS continuation, 19 November 2026: used and rechecked in this run; status-vocabulary reconciliation; metadata reconciliation only; no new mathematical check. \emph{Pre-continuation recorded proof scope.} GPT-6 Astra (OpenAI, ChatGPT Pro mode), effort Pro, 2026-11-17 UTC; Claude Fable 5.2 (Anthropic), effort high, 16 November 2026; also checked in this writing by GPT-6 Astra (OpenAI), ChatGPT Pro mode, effort Pro, 18 November 2026}{the hypotheses, notation and definitions stated here; \cref{maps:src:D.2.43}}


Let $N_L$ be the smooth triangular Mumford cap of \ref{maps:src:D.2.43},
with period sublattice $L\subset\mathbb Z^2$ of finite index
$D>0$. Write $V=\mathbb Z^2/L$ for its vertices and
$E_i$, $i\in V$, for its reduced irreducible central divisors.
Their normalizations $\nu_i:P_i\to E_i\subset N_L$ are the
degree-six toric del Pezzo surface, even when a divisor has
self-gluings. Use the cyclic ray order
\[
 (r_0,\ldots,r_5)=
       (e_1,e_2,e_2-e_1,-e_1,-e_2,e_1-e_2).                 \tag{D47.1}\label{maps:eq:D47.1}
\]
For a coefficient vector $u=(u_i)_{i\in V}$ put
$E(u)=\sum_i u_iE_i$ and
$a_{i,r}(u)=u_{i+r}-u_i$. The entire integral trilinear
form on these compact divisor classes is
\[
\begin{split}
 I_L(u,v,w)=\sum_{i\in V}u_i\Bigl(
 &-\sum_{r=0}^5 a_{i,r}(v)a_{i,r}(w)\\
 &+\sum_{r=0}^5\bigl(a_{i,r}(v)a_{i,r+1}(w)
                    +a_{i,r+1}(v)a_{i,r}(w)\bigr)\Bigr),
\end{split}                                                   \tag{D47.2}\label{maps:eq:D47.2}
\]
where ray indices are cyclic. It is symmetric in all three
arguments, and the constant vector annihilates it.
Equivalently, the cubic is
\[
\begin{split}
 I_L(u,u,u)={}&6\sum_{i\in V}u_i^3
       -3\sum_{\{i,j\}\in\mathcal E}
                      (u_i^2u_j+u_iu_j^2)\\
       &+6\sum_{(i,j,k)\in\mathcal T}u_iu_ju_k.
\end{split}                                                   \tag{D47.3}\label{maps:eq:D47.3}
\]
Here $\mathcal E$ is the set of the $3D$ quotient edge
\emph{cells}, indexed by $i$ and the first three directions
in \eqref{maps:eq:D47.1}, and $\mathcal T$ consists of the $2D$ triangle
cells. Loops, multiple edges and repeated vertex labels in
a triangle are retained with their actual cell multiplicities.
They are not replaced by a simple graph or a simplicial
complex with distinct labels.
The Chern and Pontryagin evaluations on every component are
\[
 \left\langle c_2(TN_L),\nu_{i*}[P_i]\right\rangle=0,
 \qquad
 \left\langle p_1(TN_L),\nu_{i*}[P_i]\right\rangle=0.         \tag{D47.4}\label{maps:eq:D47.4}
\]
All these formulas remain the same when the cap is included
in a closed smooth complex threefold.

For $L=\mathbb Ze_1\oplus d\mathbb Ze_2$, the entire
component trilinear form is zero for every $d>0$.
It is not zero for arbitrary period lattices: for
$L=2\mathbb Z^2$ and its four vertices,
\[
 E_i^3=6,\qquad E_i^2E_j=-2\ (i\ne j),\qquad
 E_iE_jE_k=2\ (i,j,k\text{ distinct}).                      \tag{D47.5}\label{maps:eq:D47.5}
\]
This is the component part of the cohomology ring and
the indicated characteristic pairings. It is not asserted
to determine torsion products, the full fundamental-class
invariant, or a global diffeomorphism type.

\end{theorem}
\begin{proof}
\begin{enumerate}[label=\arabic*.,leftmargin=*]
\item Denote the six toric boundary curves of $P_i$ by $C_{i,r}$.
The fan \eqref{maps:eq:D47.1} is the fan of $\mathbb P^2$ subdivided
at its three torus-fixed points. Consecutive rays have
determinant one and $r_{j-1}+r_{j+1}=r_j$. Consequently
\[
 C_{i,r}^2=-1,\qquad
 C_{i,r}C_{i,r+1}=1,\qquad
 C_{i,r}C_{i,s}=0\quad(s\not\equiv r,r\pm1).               \tag{D47.6}\label{maps:eq:D47.6}
\]
These also follow directly from the monomial coordinates
on the six smooth affine toric charts.

\item For $j\ne i$, the pullback of the Cartier divisor $E_j$
to $P_i$ is precisely the sum of those boundary curves
whose adjacent label is $j$. The total reduced central
divisor is $\sum_jE_j=\operatorname{div}(q)$, so its
absolute line bundle on the cap is trivial. Restricting
this identity determines the remaining pullback and gives
the equality of line bundles, or their divisor classes,
\[
 \nu_i^*\mathcal O(E(u))
       =\mathcal O_{P_i}\left(
                   \sum_{r=0}^5(u_{i+r}-u_i)C_{i,r}\right).
                                                               \tag{D47.7}\label{maps:eq:D47.7}
\]
When an edge returns to the same component, its coefficient
is zero. In local normal-crossing coordinates, the same
irreducible Cartier divisor then contains both branches;
the second branch contributes exactly the correction to
the normal line of the first branch. Thus \eqref{maps:eq:D47.7} does
not apply embedded-divisor adjunction incorrectly to a
self-glued component.

\item The fundamental cycle of $E_i$ is $\nu_{i*}[P_i]$ with
degree one. Intersection with the compact divisor $E(u)$,
then the projection formula, therefore gives
\[
 E(u)E(v)E(w)=\sum_i u_i
           \int_{P_i}c_1\nu_i^*\mathcal O(E(v))
                       c_1\nu_i^*\mathcal O(E(w)).          \tag{D47.8}\label{maps:eq:D47.8}
\]
Equations \eqref{maps:eq:D47.6}--\eqref{maps:eq:D47.7} give \eqref{maps:eq:D47.2} integrally.
It is symmetric because these are products of actual
Cartier classes on the smooth threefold. A constant
coefficient vector has trivial restriction on every
$P_i$, so annihilates the form.

\item There is a relative-cohomology distinction here. The
divisors have compact support and define classes in
$H^2(N_L,\partial N_L;\mathbb Z)$; their triple product
is evaluated in relative degree six. Although
$\sum_iE_i$ has a trivial absolute line bundle on the
cap, its compact-support class need not be zero: the
boundary value of $q$ winds around the origin.
The assertion is that the constant vector annihilates
the component intersection form, not that the sum of
the component fundamental cycles vanishes in homology.
For an inclusion into a closed threefold, the same
compact cycles and local divisor restrictions compute
the same products.

\item \emph{The quotient cell formula and examples.}
Set $v=w=u$ in \eqref{maps:eq:D47.2} and expand. A vertex contributes
$6u_i^3$. Pairing the two directed incidences of every
edge gives $-3(u_i^2u_j+u_iu_j^2)$. Each triangle has
three corner incidences and contributes $6u_iu_ju_k$.
This proves \eqref{maps:eq:D47.3}, with all incidence multiplicities
retained, even when labels coincide.
One explicit set of triangle cells is
\[
 (i,i+e_1,i+e_2),\qquad
 (i+e_1+e_2,i+e_1,i+e_2)\quad(i\in V).                    \tag{D47.9}\label{maps:eq:D47.9}
\]
The integral trilinear form is already specified by
\eqref{maps:eq:D47.2}; no divisibility assumption based on a rational
polarization is needed.

\item If $e_1\in L$ and $V$ is the cyclic sequence of length
$d$ in the $e_2$ direction, the $e_1$ loop terms cancel
the vertex terms. The other two edge families have the
same adjacent labels and cancel the two triangle
families in \eqref{maps:eq:D47.9}. Hence the cubic, and by \eqref{maps:eq:D47.2}
the whole trilinear form, vanish for every $d$, including
$d=1,2$. For $L=2\mathbb Z^2$, each vertex sees the
other three labels twice in its hexagon. Substitution
in \eqref{maps:eq:D47.2} gives \eqref{maps:eq:D47.5}, so a vanishing conclusion
cannot be extrapolated from the cyclic period case.

\item \emph{The canonical volume and characteristic pairings.}
On the torus cover away from $q=0$ consider
\[
                    \omega=d\log X\wedge d\log Y\wedge dq.
                                                               \tag{D47.10}\label{maps:eq:D47.10}
\]
The period translations multiply $X,Y$ by integral
powers of $q$ and by nonvanishing functions of $q$.
The added logarithmic differential terms are multiples
of $dq$, so $\omega$ is invariant and descends.
In every unimodular triangular chart one can choose
coordinates with $q=xyz$, $X=x$, $Y=y$ up to the
unimodular monomial change. Then
\[
                    d\log x\wedge d\log y\wedge d(xyz)
                                 =dx\wedge dy\wedge dz.
                                                               \tag{D47.11}\label{maps:eq:D47.11}
\]
Thus the descended form extends without a zero or pole
over the central fibre, including self-intersections.
The cap has a trivial canonical bundle near all its
compact components.

\item The normalization $\nu_i$ is a holomorphic immersion:
each preimage of a crossing lies on one smooth local
coordinate branch. Even when its image has self-gluings,
the immersion normal line is therefore
$\mathcal N_{\nu_i}=K_{P_i}$, by the adjunction sequence
and the nonvanishing volume form. This normal line
must not be confused with $\nu_i^*\mathcal O(E_i)$,
which includes the other self-intersection branches.
The exact sequence
\[
       0\longrightarrow TP_i\longrightarrow\nu_i^*TN_L
                        \longrightarrow K_{P_i}\longrightarrow0
                                                               \tag{D47.12}\label{maps:eq:D47.12}
\]
now gives
\[
 \nu_i^*c_2(TN_L)=c_2(TP_i)-c_1(TP_i)^2,\qquad
 \nu_i^*p_1(TN_L)=2c_1(TP_i)^2-2c_2(TP_i).                 \tag{D47.13}\label{maps:eq:D47.13}
\]
The hexagonal surface has Euler number six (one for
each torus-fixed point) and $c_1^2=6$, also obtained
by squaring $\sum_rC_{i,r}$ using \eqref{maps:eq:D47.6}.
Both evaluations in \eqref{maps:eq:D47.13} are therefore zero,
proving \eqref{maps:eq:D47.4}. The calculation uses the normal line
of the immersion and is valid for every quotient
lattice, not only for embedded central components.

\item \emph{Exact verification and bounds.}
The program
\path{code/III/fourfolds_D2_component_ring.py}
enumerates all $127$ column Hermite forms
$\left(\begin{smallmatrix}a&b\\0&d\end{smallmatrix}\right)$
with $ad\le12$ and $0\le b<a$. For every one it
computes the integral triple products directly from
the six boundary restrictions, verifies symmetry
in all three indices and the constant annihilator,
and compares every cubic coefficient with the
quotient cell formula while retaining loops and
repeated labels. It also verifies the toric ray,
intersection and canonical-volume identities and
the examples \eqref{maps:eq:D47.5}. All matrices, cells and
nonzero integral products are retained in
\path{records/fourfolds_D2_component_ring.json}.
The proof above applies to arbitrary finite index;
the finite computation tests the conventions and
does not supply an unstated global index bound.
\hfill$\square$
\end{enumerate}
\end{proof}


\begin{theorem}[complete smooth classification of the simply connected weighted Engel family]\label{maps:src:D.2.48}
\StatementRecord{maps:src:D.2.48}{proved}{Historical events: Claude Fable 5.2 (Anthropic), effort high, 16 November 2026: carried as background without fault found. GPT-6 Astra (OpenAI), ChatGPT Pro mode, effort Pro, S6\_09\_PAPERS, 17 November 2026: targeted proof check. Also checked in this writing. Pre-continuation writing events: GPT-6 Astra (OpenAI), ChatGPT Pro mode, effort Pro, S6\_10\_PAPERS, 18 November 2026: targeted proof check. Continuation metadata events: GPT-6 Astra (OpenAI), ChatGPT Pro mode, effort Pro, S6\_10\_PAPERS continuation, 19 November 2026: used and rechecked in this run; status-vocabulary reconciliation; metadata reconciliation only; no new mathematical check. \emph{Pre-continuation recorded proof scope.} GPT-6 Astra (OpenAI, ChatGPT Pro mode), effort Pro, 2026-11-17 UTC; Claude Fable 5.2 (Anthropic), effort high, 16 November 2026; also checked in this writing by GPT-6 Astra (OpenAI), ChatGPT Pro mode, effort Pro, 18 November 2026}{the hypotheses, notation and definitions stated here; \cref{maps:src:D.2.34}; \cref{maps:src:D.1.11}; \cref{maps:eq:D.WE3}; \cref{maps:eq:D.WI5}; \cref{maps:eq:D.WI4}; \cref{maps:src:D.2.44}; \cref{maps:src:D.2.45}; \cref{maps:eq:D.WI7}; \cite[Definition 1.7 and Theorem 3.1, pp.239,243--246]{mapsKaplan79}; \cref{maps:eq:D.WI9}; \cref{maps:eq:D.WI11}; \cref{maps:eq:D.WI13}; \cref{maps:eq:D.WI10}; \cref{maps:eq:D.WI12}}


Take the weight-$d$ Engel construction of \ref{maps:src:D.2.34}, for any $d>0$,
with both additive sites replaced by good logarithmic caps.  Retain
all the actual marked data
\[
 m_0=3k,\qquad m_1=4\ell,\qquad
 v_i=r_i\delta+n_i u_i,\qquad
 P=m_1n_0+m_0n_1=\pm1 .
 \tag{D.WI1}\label{maps:eq:D.WI1}
\]
The manifold is simply connected.  Its complex first and second
Chern classes, spin characteristic class, and refined class satisfy
\[
 c_1(X)=c_2(X)=\lambda(X)=0,\qquad P_{\mathrm{ref}}(X)=0.
 \tag{D.WI2}\label{maps:eq:D.WI2}
\]
Here $P_{\mathrm{ref}}$ denotes Zhubr's class, to distinguish it from
the Seifert determinant $P$ in \eqref{maps:eq:D.WI1}.

The classes $x_i=\operatorname{PD}(D_i)$ of the $d$ cusp components,
indexed cyclically, generate $H^2(X;\mathbb Z)$ with sole relation
$\sum_i x_i=0$.  The complex ruling $f=H-E_1$ has image
$s=[a\wedge b]$, a generator of $\operatorname{Tor}H_2(X)=\mathbb Z/d$.
Put $\kappa=\operatorname{PD}(f)$.  With repetitions retained when
$d=2$, the complete degree-two product is
\[
 x_i x_j=
 \bigl(\mathbf1_{j=i+1}+\mathbf1_{j=i-1}-2\mathbf1_{j=i}\bigr)\kappa.
 \tag{D.WI3}\label{maps:eq:D.WI3}
\]
For $d=1$ the sole component class is zero.  For $d=2$ the two
off-diagonal incidences both occur, and every product in \eqref{maps:eq:D.WI3}
is zero since $2\kappa=0$.  The integral cubic is zero for every $d$.
Here is the complete fundamental-class invariant, including its
two- and three-primary coordinates. Put
$r=d-1$, $D_2=d\gcd(d,2)$, and $D_3=d\gcd(d,3)$.
For $d>1$ there is a class $u\in H^2(X;\mathbb Z/d)$
which, together with $x_0,\ldots,x_{d-2}$, fixes the splitting
$H_2(X)=\mathbb Z^r\oplus\mathbb Z/d$ described below. In
\[
 H_6(K(\mathbb Z^r\oplus\mathbb Z/d,2);\mathbb Z)
 \simeq
 \mathbb Z^{\binom{r+2}{3}}\oplus
 (\mathbb Z/d)^{\binom{r+1}{2}}\oplus
 (\mathbb Z/D_2)^r\oplus\mathbb Z/D_3
 \tag{D.WI14}\label{maps:eq:D.WI14}
\]
the complete class $\mu$ is determined by
\[
 \begin{aligned}
 \langle x_i x_j x_k,[X]\rangle&=0,\\
 \langle u x_i x_j,[X]\rangle
    &=\mathbf1_{j=i+1}+\mathbf1_{j=i-1}-2\mathbf1_{j=i}
                                                     &&\pmod d,\\
 \langle \mathcal Q_2(u)x_i,[X]\rangle&=-12             &&\pmod {D_2},\\
 \langle \mathcal Q_3(u),[X]\rangle&=0                 &&\pmod {D_3}.
 \end{aligned}
 \tag{D.WI15}\label{maps:eq:D.WI15}
\]
The indices in the second line are cyclic modulo $d$; those
used as independent coordinates run from zero to $d-2$.
Here $\mathcal Q_2$ is the ordinary square if $d$ is odd and
the Pontryagin square with values modulo $2d$ if $d$ is even.
Likewise $\mathcal Q_3$ is the ordinary cube if $3\nmid d$
and Thomas's Pontryagin third power with values modulo $3d$
if $3\mid d$.

Consequently, for every $d>0$ all the simply connected data
\eqref{maps:eq:D.WI1} give one oriented diffeomorphism type $M_d$, namely
the unique type with
\[
 H_2=\mathbb Z^{d-1}\oplus\mathbb Z/d,\qquad b_3=0,\qquad
 P_{\rm ref}=0,\qquad \mu\ \text{given by \eqref{maps:eq:D.WI15}}.
 \tag{D.WI16}\label{maps:eq:D.WI16}
\]
One has $M_1=S^6$, and for $d$ prime to six this is the
type already identified in \ref{maps:src:D.1.11}. The types for different
$d$ are distinct, since their Euler numbers are $2d$.

\end{theorem}


\begin{proof}
\begin{enumerate}[label=\arabic*.,leftmargin=*]
\item \emph{Primitive component classes.}
Write $E$ for the complement of the interior of the cusp tube.
The actual finite-cap presentation before adjoining its cusp meridian is
\[
 H_1(E)=
 \langle\gamma,x_0',x_1'\mid m_0x_0'=n_0\gamma,\ 
                                   m_1x_1'=n_1\gamma\rangle .
\]
The determinant condition implies
$\gcd(m_0,m_1)=\gcd(m_i,n_i)=1$ and $H_1(E)=\mathbb Z$.
The cusp meridian is $x_0'+x_1'$ up to the common orientation sign.
It is primitive because adjoining its relation gives $H_1(X)=0$.
Angular differences of adjacent toric meridians are killed in
$H_1(E)$.  Thus the relative Thom map is
\[
 H^1(E)\longrightarrow H^2(X,E)=\mathbb Z^d,\qquad
                   1\longmapsto\pm(1,\ldots,1).
 \tag{D.WI4}\label{maps:eq:D.WI4}
\]
The cokernel of the resulting injection
$\mathbb Z^d/\mathbb Z(1,\ldots,1)\hookrightarrow H^2(X)$
embeds in $H^2(E)$, which is torsion-free by the integral universal
coefficient theorem.  It has rank zero by \ref{maps:src:D.2.34}, so it vanishes.
In particular the entire fibre class is
$f_X^*[\mathrm{point}]=\sum_i x_i=0$.
For a reduced multiple fibre $D$ one has
$m[D]=f_X^*[\mathrm{point}]=0$; the group $H^2(X)$ is free,
so $[D]=0$ as well.

\item The common ruling in the $d$-period cusp is the specialization
of $a\wedge b$, with its complex orientation.  This is an integral
unit-period annulus calculation: that torus remains in one gamma
component and limits to the two adjacent external curves
$E_2+E_3=H-E_1$.  The three support functions
$r(r-1)/2$, $rs$, and $s(s-1)/2$ evaluate on that ruling as
$(-1,0,0)$, exactly their actual multiplier windings on
$a\wedge b$; its dual-torus coordinate is zero.
There is no division by $d$ or saturation at this step.
The primitive functional of \ref{maps:src:D.2.34} evaluates it as
$\xi_d(a\wedge b)=1$.  The matrix \eqref{maps:eq:D.WE3} shows that $s$ has
exact order $d$ in the simply connected sector.  The actual
external intersections and the principal sum of components
therefore give \eqref{maps:eq:D.WI3}.  The adjacent-pair formula counts both
edges when they have the same pair of labels.  Since every
product is torsion, its further product with $H^2(X)$ vanishes
in $H^6(X)=\mathbb Z$.

\item \emph{The differential and its literal critical cycle.}
The normalization immersion $j_i:dP_6\longrightarrow X$ has
normal line $K_{dP_6}$.  Indeed the principal local central
divisor is the product of the normal coordinates to all its
branches; on a fixed normalized branch its normal transition
is the inverse of the six toric boundary transitions.  This
includes both conductor branches of each self-identification.
Consequently
\[
 j_i^*c_1(TX)=c_1(TdP_6)+c_1(K_{dP_6})=0,\qquad
 \langle j_i^*c_2(TX),[dP_6]\rangle=6-6=0 .
 \tag{D.WI5}\label{maps:eq:D.WI5}
\]
Restriction of $H^2(X)$ to the normalized components is injective.
For a component combination with coefficients $u_i$, its
restriction to component $i$ is
$\sum_r(u_{i+r}-u_i)C_r$.  Pairing with the anticanonical ample
class, which has degree one on each of the six boundary curves,
gives the component-graph Laplacian.  The graph is connected, so
zero restrictions force the $u_i$ to be constant.  That is
precisely the already removed principal relation.  Equation
\eqref{maps:eq:D.WI5} now gives $c_1(TX)=0$ integrally.

\item Remove the divisorial multiple-fibre factors from $df_X$.  Its
corrected target is the line
\[
 L=f_X^*T\mathbb P^1\otimes
                      \mathcal O\bigl(-\sum (m_i-1)D_i^{\rm fin}\bigr).
\]
It is topologically trivial by \eqref{maps:eq:D.WI4}.  The abstract rank-two
kernel $V$ is locally free, including at $xy$ and $xyz$ charts,
and has the nowhere-zero abstract section given by $\delta$.
At a node the kernel columns are $(x,-y,0)$ and $(0,0,1)$;
at a triple point they are $(x,-y,0)$ and $(x,0,-z)$.
The constant scalar-circle weight vector is a nonzero constant
combination in these abstract bases.  At the good finite caps
the corrected differential is a submersion and $\delta$ remains
a nonzero tangent vector.  Since
$c_1(V)=c_1(TX)-c_1(L)=0$, its quotient by this section is a
trivial complex line.  Thus $V$ and $L$ admit full smooth frames.
They give the reference frame $\rho$ on the complement of the
literal critical curve cycle $C$.

\item The local Bott clutch at $xy$ has coefficient one:
its two normal columns are
$\bigl(\begin{smallmatrix}x&\overline y\\-y&\overline x\end{smallmatrix}\bigr)$.
The $xyz$ Hilbert--Burch minors give multiplicity one on each
of its three axes and no further codimension-two contribution.
The good finite logarithmic caps have no residual critical
cycle after correcting the target.  Hence, with the complex
orientation and the spin normalization of \ref{maps:src:D.2.44}--\ref{maps:src:D.2.45},
\[
                 \lambda(TX,\rho)=-\operatorname{PD}(C).
 \tag{D.WI6}\label{maps:eq:D.WI6}
\]

\item Only a local covering is used to compute $[C]$.
In the positive period basis
$N b=a$, $N\gamma_1=\delta$, the weight-one cusp's literal
three-curve sum has regular-fibre preimage
$a\wedge b+b\wedge\delta+\delta\wedge\gamma_1$.
The weight-$d$ cusp is its degree-$d$ etale period cover, with
$\gamma_d=d\gamma_1$.  Transfer therefore gives the actual
preimage of its full $3d$-curve sum as
\[
 d\,a\wedge b+d\,b\wedge\delta+\delta\wedge\gamma_d.
 \tag{D.WI7}\label{maps:eq:D.WI7}
\]
The angular vector $\delta$ here is fixed by
$N\gamma_d=d\delta$; this convention removes the possible
opposite choice for the original scalar-circle parameter.
Since $\xi_d$ has values $1,0,-6d$ on these three bivectors,
\eqref{maps:eq:D.WI7} has image $-5d\,s=0$.  Thus $C$ is integrally
nullhomologous and $c_2(TX)=\lambda(TX)=0$.
No covering of either finite logarithmic cap is assumed.

\item \emph{A finite-cap detector with its actual reference parity.}
If $d$ is odd, $2P_{\rm ref}=\lambda=0$ already gives
$P_{\rm ref}=0$.  Suppose $d$ is even.  In the order-four
finite cap, put $\beta=[a]$ and $\eta=[u_1]$ in its two
central base directions.  Its marked base lattice is
\[
 \langle\beta,\eta,g\mid m_1g=-2r_1\beta+n_1\eta\rangle .
 \tag{D.WI8}\label{maps:eq:D.WI8}
\]
Choose integers $B,C$ with $m_1C=-2r_1+n_1B$ and define the
primitive character by $\chi(\beta)=1$, $\chi(\eta)=B$,
$\chi(g)=C$.  Such integers exist because $\gcd(m_1,n_1)=1$.
Its connected fibre $N$ in the central four-manifold is an
elliptic-torus mapping torus.  In its base kernel
$h=g-C\beta$ and $q=\eta-B\beta$ obey $m_1h=n_1q$.
A primitive generator has the root-loop exponent inverse to
$n_1$ modulo $m_1$, hence odd modulo four.  The elliptic
monodromy of $N$ is therefore the order-four matrix or its
inverse.  Its Wang sequence gives
\[
             H_1(N;\mathbb Z)=\mathbb Z\oplus\mathbb Z/2.
 \tag{D.WI9}\label{maps:eq:D.WI9}
\]
Moreover $\chi(\delta)=-2$, so the actual scalar tangent
vector is transverse to $N$ in the reduced finite fibre.
Projecting $J\delta,\operatorname{Re}v_2,\operatorname{Im}v_2$
to $TN$ along $\delta$ identifies the stable tangent frame
induced by $\rho$ with an actual tangent parallelization.
The normal three-frame is $\delta$ followed by the two real
directions of $L$.

\item For an actual parallelization of a closed oriented
three-manifold, its stable framed coefficient modulo two is
$1+\dim H_1(N;\mathbb F_2)$.  The spin-compatible Kaplan
zero-/two-handle proof was given in the finite-II frame lemma:
the relative $\lambda$ reduces to the relative Euler number,
and the even intersection matrix has even mod-two rank.
It applies to every prescribed spin structure
\cite[Definition 1.7 and Theorem 3.1, pp.239,243--246]{mapsKaplan79}.
Equation \eqref{maps:eq:D.WI9} therefore gives reference parity
\[
                         \epsilon=[N]_\rho=1\pmod2 .
 \tag{D.WI10}\label{maps:eq:D.WI10}
\]

\item Retain the actual Thom chains, not only their matrix classes.
Let $K_{a0},K_{a1}$ be the first columns of the two blocks in
\eqref{maps:eq:D.WE3}.  Directly in its four target coordinates,
\[
                    d s=-2\ell P\,K_{a0}-kP\,K_{a1}.
 \tag{D.WI11}\label{maps:eq:D.WI11}
\]
The chain at cap zero misses $N$, whereas the $a$-Thom
chain in cap one pairs with $N$ by $\chi(a)=1$.
Thus $\operatorname{lk}(N,s)=-kP/d$.
The determinant equation implies that $k$ is odd, so $N$
generates the two-primary summand of $H_3(X)$.
The specialization chains for \eqref{maps:eq:D.WI7} are supported in the
cusp and miss $N$.  Equations \eqref{maps:eq:D.WI7}--\eqref{maps:eq:D.WI11} then give
the actual integer link
\[
                          \ell(N,C)=5kP,
 \tag{D.WI12}\label{maps:eq:D.WI12}
\]
up to the common harmless orientation sign.  It is odd.

\item Finally the relative refined-class computation extends with
reference parity $\epsilon$ in the form
\[
 \operatorname{lk}(N,\operatorname{PD}P_{\rm ref})
                 =-\frac{\ell(N,C)}2-\frac{\epsilon}2
                                      \quad\text{in }\mathbb Q/\mathbb Z .
 \tag{D.WI13}\label{maps:eq:D.WI13}
\]
Indeed, choose Zhubr's admissible three-frame on a genuinely
embedded torsion sphere obtained from an odd multiple of $N$
by spin-compatible ambient two-handle surgery.  Its complement
spin three-frame extends over the trace.  The induced stable
framed class on $N$ is even.  The comparison with $\rho$ has
Spin(6) degree $q_0\equiv\epsilon\pmod2$.
For a bounding four-chain $B_4$ with boundary
$\operatorname{ord}(N)N$, the exact relative transgression is
$-B_4\cdot C=2q+\operatorname{ord}(N)q_0$.
Division of this integer identity gives \eqref{maps:eq:D.WI13}, retaining
the actual rational link before reduction modulo one.
This is the proof of the relative lemma in \ref{maps:src:D.2.45} with its
reference-parity term restored, not an application of the
unproved general change-of-frame conjecture.
Free $H_2$ causes no change: $C$ is nullhomologous, and
$2P_{\rm ref}=\lambda=0$ already kills the free part.

\item The two odd halves in \eqref{maps:eq:D.WI13} cancel by \eqref{maps:eq:D.WI10} and
\eqref{maps:eq:D.WI12}.  The two-primary part of $P_{\rm ref}$ is zero,
and its odd-primary part was already zero from
$2P_{\rm ref}=0$.  This proves \eqref{maps:eq:D.WI2}.
\end{enumerate}
\end{proof}

\begin{proof}[Proof of the complete invariant and classification]
\begin{enumerate}[label=\arabic*.,leftmargin=*]
\item \emph{A global finite-coefficient class, retaining boundary torsion.}
At weight one the regular bundle $V_1$ over the peripheral
graph has free $H_1$ and $H_2$. Its degree-two group has the
primitive fibre quotient generated by $s$, together with the
five marked graph generators in \eqref{maps:eq:D.WE11}. Choose the integral
cohomology class $\alpha_1$ which evaluates as $\xi_1$ on the
fibre part and zero on every graph generator.

\item On the weight-one cusp use the actual toric line with support
\[
                 \psi(r,s)=-\frac{r(r-1)}2+3s(s-1).
 \tag{D.WI17}\label{maps:eq:D.WI17}
\]
Its multiplier windings give $\xi_1$ in the positive period
marking used in \eqref{maps:eq:D.WI7}. Its boundary value agrees integrally
with $\alpha_1$: their fibre restrictions agree, and both
vanish on the two angular Wang sweeps, which bound in the
cusp. These are all the remaining generators of $H_2(M_1)$.
Also $H_1(M_1)$ is free, so the integral universal coefficient
theorem leaves no hidden $\operatorname{Ext}$ difference.

\item Before the finite caps are attached, the core and cusp have
the genuine degree-$d$ period covering defined by
$\gamma_d=d\gamma_1$. Pull back both classes along that
cover. They then agree on the entire boundary $M_d$, not
only on its regular-fibre restriction. In particular this
argument retains the $\mathbb Z/d$ term in $H_1(M_d)$ and
the resulting torsion in $H^2(M_d)$. The core pullback
$\alpha_d$ has fibre functional $\xi_d$ and zero values on
the five graph generators: the covering fixes the first
three such generators and multiplies the two $u_i$ generators
by $d$.

\item At a finite cap let $\pi:Y_i\longrightarrow Q_i$ be its
actual normal circle bundle. The obstruction for
$\alpha_d\bmod d$ to extend across the disk bundle is its
Gysin image
\[
                  \pi_!(\alpha_d\bmod d)\in
                         H^1(Q_i;\mathbb Z/d).
 \tag{D.WI18}\label{maps:eq:D.WI18}
\]
Its values on the complete marked generators
$a,\gamma,g$ of $H_1(Q_i)$ are precisely the first-row
coefficients of \eqref{maps:eq:D.WE3}:
\[
 \begin{array}{c|ccc}
       &a&\gamma&g\\ \hline
 i=0&-2dn_0&-6dr_0&0\\
 i=1&-3dn_1&-6dr_1&0.
 \end{array}
 \tag{D.WI19}\label{maps:eq:D.WI19}
\]
The root-loop column $g$ is retained. This is the Thom
definition of fibre integration: a loop in $Q_i$ sweeps its
normal circle, and pairing with $\alpha_d$ gives the displayed
actual normal-column coefficient. Every coefficient is zero
modulo $d$. Since $H^1(Q_i;\mathbb Z/d)$ is
$\operatorname{Hom}(H_1(Q_i),\mathbb Z/d)$, \eqref{maps:eq:D.WI18} vanishes,
even when the integral normal line has torsion. The circle
Gysin sequence therefore supplies an extension over each
finite cap. Mayer--Vietoris now glues them with the pulled-back
cusp line to a class $u\in H^2(X;\mathbb Z/d)$.
Only the core and cusp were covered; neither finite fill was
assumed to be a cover of a weight-one fill.

\item On every normalized cusp component the integral representative
of its restriction is
\[
              \eta=-5H+6E_1-E_3,\qquad
       \eta(E_1,E_2,E_3)=(-6,0,1),\qquad
       \eta(f)=1,\qquad \eta^2=-12.
 \tag{D.WI20}\label{maps:eq:D.WI20}
\]
The three degrees follow directly by taking the differences
of \eqref{maps:eq:D.WI17} across the internal, diagonal, and scalar-fixed
walls, as in \ref{maps:src:D.1.9}. In particular $u(s)=1$.
Choose free lifts in $H_2(X)$ annihilated by $u$, subtracting
appropriate multiples of $s$. This specifies the splitting
in \eqref{maps:eq:D.WI14}--\eqref{maps:eq:D.WI16}. The integral classes $x_i$ form its
dual free basis, and $u$ is the standard cyclic coordinate.

\item \emph{The full mixed coordinates.}
The first mixed value in \eqref{maps:eq:D.WI15} follows from \eqref{maps:eq:D.WI3} and
$\langle u\kappa,[X]\rangle=1$. The Pontryagin square of the
reduction of an integral degree-two class is its ordinary
square with the full target modulus $2d$; this is the
$p=2$ case of the actual cochain construction
\cite[equations (2.1), (2.8), and (3.4), pp.267--268]{ffThomas56}.
By naturality and the projection formula for the normalized
component fundamental class,
\[
 \langle \mathcal Q_2(u)x_i,[X]\rangle
       =\langle\mathcal Q_2(j_i^*u),[dP_6]\rangle
       =\eta^2=-12\pmod {D_2}.
 \tag{D.WI21}\label{maps:eq:D.WI21}
\]
This computes the coefficient modulo $2d$ when $d$ is even;
reducing the ordinary square merely modulo $d$ would lose
precisely the extra order-two information. The relation
$\sum_i x_i=0$ is consistent with \eqref{maps:eq:D.WI21}, since
$12d$ is divisible by $D_2$.

\item \emph{The pure coordinate.}
The following actual scalar-circle calculation evaluates
the full Pontryagin third power. It uses the class just
constructed and its specified component restrictions.
\subsubsection*{The full Pontryagin cube of the weighted Engel torsion class}
We use the actual class $u\in H^2(X_d;\mathbb Z/d)$ whose
restrictions to the normalized cusp components are specified
in \ref{maps:src:D.2.48}. In particular its degree on every scalar-fixed
rational curve $C_j$, $j\in\mathbb Z/d$, is one, and its
component restriction is represented by the integral toric
support
\[
                  \psi(r,s)=-\frac{r(r-1)}2+3s(s-1).
                                                               \tag{D.TP1}\label{maps:eq:D.TP1}
\]
The calculation below uses these actual restrictions. Merely
choosing an arbitrary degree-two class that is a unit on
torsion would not fix its pure coordinate in the presence
of the free summand.

\item The scalar circle action is semifree. On a finite cap of
order $m_i$ its gamma translation has primitive numerator
$n_i$, with $\gcd(m_i,n_i)=1$. A nontrivial root power cannot
be compensated by the scalar direction, because that would
require $j n_i/m_i$ integral. Thus these caps are scalar-free.
The cusp has exactly the $d$ fixed curves $C_j$; their normal
bundles are $\mathcal O(-1)\oplus\mathcal O(-1)$ with weights
$(1,-1)$, and the other cusp orbits are free. These are the
marked geometric inputs of the weighted family.

\item Because $X_d$ is simply connected, $u$ has an equivariant
lift $\widehat u\in H^2_{S^1}(X_d;\mathbb Z/d)$.
In the Borel spectral sequence the only possible
differentials out of ordinary degree two have targets
involving $H^1(X_d;\mathbb Z/d)$ or $H^3(BS^1;\mathbb Z/d)$,
and both vanish. Write
\[
         \widehat u|_{C_j}=x+\lambda_j t,
         \qquad x=[C_j]^*,\quad t\in H^2(BS^1;\mathbb Z/d).
                                                               \tag{D.TP2}\label{maps:eq:D.TP2}
\]
The difference of any two equivariant lifts of a fixed
ordinary class on the simply connected normalized
$dP_6$ is a multiple of $t$. Hence the differences of
the characters on its two scalar-fixed curves can be
read from any toric equivariant lift of \eqref{maps:eq:D.TP1}.
The fixed edges are the translates of the edge from
$(0,j)$ to $(0,j+1)$; the support slope in their scalar
direction is
\[
               \psi(0,j+1)-\psi(0,j)=6j.
                                                               \tag{D.TP3}\label{maps:eq:D.TP3}
\]
The fibre character of the associated monomial
trivialization is the negative of this slope, up to
one global character convention. Equivalently the
intervening weight-one orbit curve has degree $-6$.
Thus, for one common sign $\epsilon=\pm1$,
\[
                    \lambda_j=\lambda_0+6\epsilon j
                                 \quad\text{in }\mathbb Z/d.
                                                               \tag{D.TP4}\label{maps:eq:D.TP4}
\]
This equation is compatible with the period $j\mapsto j+d$.
It does not assert that the integral support line itself
descends with an integral global circle linearization:
the needed character equality is modulo $d$ and follows
from the existing global equivariant cohomology lift.

\item Suppose now that $3\mid d$. Thomas's operation is defined
with precisely the required full coefficients:
\[
 \mathcal P_3:H^2(-;\mathbb Z/d)
                    \longrightarrow H^6(-;\mathbb Z/(3d)).
                                                               \tag{D.TP5}\label{maps:eq:D.TP5}
\]
This is \cite[equation (2.1), p.267]{ffThomas56} with
$p=3$ and $m=d/3$. On the reduction of an integral
class $z$, it equals the reduction of $z^3$ modulo
$3d$. Indeed the correction term in Thomas's actual
cochain formula (3.4) is zero for an integral cocycle,
and its first term is the ordinary threefold diagonal
product \cite[equations (3.4)--(3.6), p.268]{ffThomas56}.
Each restriction \eqref{maps:eq:D.TP2} has such an integral lift.

\item Apply equivariant localization with coefficients
$\mathbb Z/(3d)$ to $\mathcal P_3(\widehat u)$.
No field assumption is needed here: the complement
is scalar-free, and each normal Euler class becomes
a unit after inverting $t$. On a fixed curve it is
\[
                    (-x+t)(-x-t)=-t^2,
\]
because $x^2=0$. Integrating the term of ordinary
curve degree two gives
\[
 \left\langle\mathcal P_3(u),[X_d]\right\rangle
                      =-3\sum_{j=0}^{d-1}\lambda_j^2
                                  \quad\text{in }\mathbb Z/(3d).
                                                               \tag{D.TP6}\label{maps:eq:D.TP6}
\]
The right side is independent of the integer
representatives of the characters: $3\mid d$ makes
the cube of an integral lift well defined modulo
$3d$. The base cohomology polynomial ring injects
into its localization at $t$, so the degree-zero
identity is an identity in the original coefficient
ring, not just a localized statement.

\item For arbitrary integers $d>0,\lambda_0$ and either
sign in \eqref{maps:eq:D.TP4}, the exact identity is
\[
 \sum_{j=0}^{d-1}(\lambda_0+6\epsilon j)^2
   =d\left(\lambda_0^2+6\epsilon\lambda_0(d-1)
                         +6(d-1)(2d-1)\right).              \tag{D.TP7}\label{maps:eq:D.TP7}
\]
Substitution into \eqref{maps:eq:D.TP6} proves
\[
                 \left\langle\mathcal P_3(u),[X_d]\right\rangle=0
                                  \quad\text{in }\mathbb Z/(3d).
                                                               \tag{D.TP8}\label{maps:eq:D.TP8}
\]
Thus the complete pure cyclic fundamental-class coordinate
is zero, including the order-three subgroup invisible to
the ordinary mod-$d$ cube. For $3\nmid d$ the same
localization calculation with the ordinary cube gives
zero modulo $d$. Mixed coordinates are separate from
this calculation; in particular it does not discard
the order-$2d$ mixed Pontryagin-square terms when $d$
is even.

\item The universal integer identities, support slope,
single-curve Euler/cube calculation and an explicitly
bounded supplementary check are executed in
\path{code/III/fourfolds_D48_thomas.py}.
The supplementary range is $d=3,6,\ldots,120$,
$\lambda_0=0,\ldots,d-1$ and both signs, giving
$4920$ exact checks. The arbitrary-$d$ conclusion
comes from \eqref{maps:eq:D.TP7}, not that finite range.

\item \emph{Why these detectors are complete.}
We include the low-degree integral target calculation to
exclude an undetected finite subgroup. For a cyclic group,
\[
 \begin{aligned}
 H_2(K(\mathbb Z/d,2);\mathbb Z)&=\mathbb Z/d,\\
 H_4(K(\mathbb Z/d,2);\mathbb Z)&=\mathbb Z/D_2,\\
 H_6(K(\mathbb Z/d,2);\mathbb Z)&=\mathbb Z/D_3 .
 \end{aligned}
 \tag{D.WI22}\label{maps:eq:D.WI22}
\]
The degree-four assertion follows from Whitehead's exact
sequence: on $K(A,2)$ its two adjacent homotopy groups
vanish, so $H_4=\Gamma(A)$
\cite[Theorem 1 and Section 2]{mapsWhitehead50}.
For $A=\mathbb Z/d$, the universal quadratic generator
$q(1)$ has relations $d^2q(1)=2dq(1)=0$.
Conversely $q(j)=j^2$ modulo $\gcd(d^2,2d)=D_2$
is well defined and quadratic. Thus these are its complete
relations and the image of $\mathbb {CP}^2$ generates.
The ordinary or Pontryagin square evaluates on that image
as one modulo $D_2$, so it detects the entire group.

\item For degree six, Zhubr's epimorphism in Section 3.11 and
his two free bordism generators in Section 3.12 show that
the spin-realizable part for a cyclic coefficient group
is generated in homology by the image of
$[\mathbb {CP}^3]$: his other generator has zero
fundamental-class image. For odd-primary cyclic groups
this is all of $H_6$, since the only obstruction in his
Lemma 3.7 has mod-two target. For a two-power cyclic
group Section 3.12 explicitly gives
$\ker\alpha=H_6=\mathbb Z/2^a$
\cite[Lemmas 3.7 and 3.11, and Section 3.12,
pp.805--807]{manZhubr75}.
The decomposition of $\mathbb Z/d$ into its relatively
prime primary factors gives the corresponding product
of Eilenberg--MacLane spaces. All cross tensor and
$\operatorname{Tor}$ terms vanish, since their positive
homology has relatively prime orders. Hence $H_6$ is
cyclic and is generated by the image $z$ of
$[\mathbb {CP}^3]$ for every $d$.

\item The order of $z$ divides both $d^3$ and $3d$.
Indeed the class $d h$ on $\mathbb {CP}^3$ becomes zero
modulo $d$, while its integral cube is $d^3$.
On $\mathbb {CP}^1\times\mathbb {CP}^2$ the class
$d x+y$ reduces to a class factoring through
$\mathbb {CP}^2$, whose degree-six image is zero,
while its integral cube evaluates to $3d$.
Thus the order divides
$\gcd(d^3,3d)=d\gcd(d,3)=D_3$.
The opposite bound is detected by the ordinary cube when
$3\nmid d$, or by Thomas's operation when $3\mid d$:
on the reduced hyperplane class of $\mathbb {CP}^3$
the full target value is one
\cite[equation (2.8), p.267]{ffThomas56}.
This proves the last assertion of \eqref{maps:eq:D.WI22} with its
stated perfect detector.

\item Finally
$K(\mathbb Z^r\oplus\mathbb Z/d,2)
 =(\mathbb {CP}^{\infty})^r\times K(\mathbb Z/d,2)$.
The first factor has torsion-free homology in even
degrees and zero homology in odd degrees. The integral
Kunneth formula has no $\operatorname{Tor}$ term and
gives exactly \eqref{maps:eq:D.WI14}. Integral cubic monomials,
$u x_i x_j$, $\mathcal Q_2(u)x_i$, and
$\mathcal Q_3(u)$ are the respective perfect coordinate
detectors. Therefore \eqref{maps:eq:D.WI15} computes all of $\mu$,
including the extra factors two and three.

\item \emph{Application of smooth classification.}
The manifold is simply connected and spin, has
$H_2=\mathbb Z^{d-1}\oplus\mathbb Z/d$ and middle Betti
number zero by \ref{maps:src:D.2.34}, refined class zero by the first
proof, and the complete fundamental-class invariant
\eqref{maps:eq:D.WI15}. These are exactly the four invariants in
Zhubr's Theorem 1, with the definitions on pp.794--795
\cite[Theorem 1]{manZhubr75}. Their values are independent
of every marked parameter in \eqref{maps:eq:D.WI1}.
Existence is supplied by the constructed manifolds,
and the injective classification gives one oriented
diffeomorphism type for each $d$. This proves \eqref{maps:eq:D.WI16}
by an applied complete classification, rather than by
ordinary cohomology or an asserted connected-sum model.
\end{enumerate}
\end{proof}



\begin{theorem}[unbounded canonical star logs with the full integral boundary map]\label{maps:src:D.2.49}
\StatementRecord{maps:src:D.2.49}{proved}{Historical events: Claude Fable 5.2 (Anthropic), effort high, 16 November 2026: carried as background without fault found. GPT-6 Astra (OpenAI), ChatGPT Pro mode, effort Pro, S6\_09\_PAPERS, 17 November 2026: targeted proof check. Also checked in this writing. Pre-continuation writing events: GPT-6 Astra (OpenAI), ChatGPT Pro mode, effort Pro, S6\_10\_PAPERS, 18 November 2026: targeted proof check. Continuation metadata events: GPT-6 Astra (OpenAI), ChatGPT Pro mode, effort Pro, S6\_10\_PAPERS continuation, 19 November 2026: used and rechecked in this run; status-vocabulary reconciliation; metadata reconciliation only; no new mathematical check. \emph{Pre-continuation recorded proof scope.} GPT-6 Astra (OpenAI, ChatGPT Pro mode), effort Pro, 2026-11-17 UTC; GPT-6 Astra (OpenAI, ChatGPT Pro mode), S6\_06\_BLIND, effort Pro, 2026-11-17; Claude Fable 5.2 (Anthropic), effort high, 16 November 2026; also checked in this writing by GPT-6 Astra (OpenAI), ChatGPT Pro mode, effort Pro, 18 November 2026}{the hypotheses, notation and definitions stated here; \cref{maps:src:D.2.27}; \cref{maps:src:D.2.41}; \cref{maps:src:D.2.46}; \cref{maps:eq:D.GI10}; \cref{maps:src:D.2.26}; \cref{maps:src:D.2.24}; \cref{maps:eq:D.SL11}; \cref{maps:eq:D.SL10}; \cref{maps:eq:D.SL12}; \cref{maps:eq:D.SL13}; \cref{maps:eq:D.SL15}; \cref{maps:eq:D.SL14}; \cref{maps:eq:D.SL19}; \cref{maps:eq:D.SL17}; \cref{maps:eq:D.SL18}; \cref{maps:eq:D.SL16}; \cref{maps:eq:D.SL20}; \cref{maps:eq:D.SL22}; \cref{maps:eq:D.SL24}}


Let a marked zero-free star boundary have the integral basis
$(\alpha,\beta,\lambda,\ell)$ and monodromy
\[
\begin{gathered}
 A\alpha=-\alpha,\quad A\beta=-\beta-n\lambda,\quad
 A\lambda=\lambda-\alpha,\quad A\ell=\ell,\\
 \delta=-\alpha+2\lambda,\qquad\gamma^*(\ell)=1,
 \qquad n\in\{1,3\}.
\end{gathered}                                               \tag{D.SL1}\label{maps:eq:D.SL1}
\]
There are twelve actual marked star positions of this type
in the natural catalogue: the eight primitive-gamma $I_1^*$
positions of \ref{maps:src:D.2.27} and the four primitive-gamma $I_3^*$
positions. Their unimodular bases are retained in the
certificate below. The theorem also applies to any further
input for which the displayed integral marked basis is supplied.

Take a root of order $m=2k$, with $k>0$. Its rank-one
semistable Mumford cap has $q=nk$ components. Let its affine
deck action on their cyclic labels be
$i\mapsto j-i$, with a marked integer lift $j\in\mathbb Z$.
The component labels use $j$ modulo $q$; its integer lift is
retained in the peripheral marking. Write its collapsed
power as
\[
 v_{\rm col}=R\lambda+N\ell,\qquad R=2r-j,
                \qquad \gcd(N,2k)=1 .                     \tag{D.SL2}\label{maps:eq:D.SL2}
\]
Every $j,r$ and every such $N$ is realized by an actual free
affine root action. All its holomorphic unit factors and
integer affine twists are retained. Put
\[
 f=\#\{i:2i=j\pmod q\},\qquad \rho=(q+f)/2.
\]
Thus $f=1$ for odd $k$, $f=2$ for even $k$ and even $j$,
and $f=0$ for even $k$ and odd $j$.
For the quotient cap $U$, the complete integral groups are
\[
\begin{array}{c|c}
a&H_a(U;\mathbb Z)\\\hline
0&\mathbb Z\\
1&\mathbb Z^2\\
2&\mathbb Z^{\rho+1}\oplus
                    (\mathbb Z/2\text{ if }f=0)\\
3&\mathbb Z^{2\rho}\\
4&\mathbb Z^\rho,
\end{array}                                                \tag{D.SL3}\label{maps:eq:D.SL3}
\]
and higher groups vanish.

The boundary has $H_2(\partial U)=\mathbb Z^4$. Its
degree-two fibre coinvariant has the primitive basis
\[
 x=[\beta\lambda],\qquad
 w=[\beta\ell]+\frac{n-1}{2}[\lambda\ell],\qquad
 [\alpha\beta]=-2x,\quad [\beta\ell]=nw,\quad
 [\lambda\ell]=-2w.                                     \tag{D.SL4}\label{maps:eq:D.SL4}
\]
Let $D,L$ be its literal fixed-normal tori in the
$\delta,\ell$ directions. In coordinates $(x,w,D,L)$,
the \emph{full} boundary kernel is as follows:
\[
\begin{array}{c|c}
f&\ker[H_2(\partial U)\longrightarrow H_2(U)]\\\hline
1&\{(0,b,c,d):kb=2Nc-Rd\}\\
2&\{(0,b,c,d):(k/2)b=Nc-(R/2)d\}\\
0&\{(0,b,c,d):kb=2Nc-Rd,\ b\equiv0\pmod2\}.
\end{array}                                                \tag{D.SL5}\label{maps:eq:D.SL5}
\]
These are unsaturated marked kernels where appropriate.
Their projected indices in the full fixed lattice are
respectively $k,k/2,2k$. The boundary cokernel is free of
rank $\rho-1$. In particular the last row retains the
actual cap $\mathbb Z/2$ and its nonzero fibre correction;
it is not a conclusion from rational ranks or determinant
orders alone.

These columns can be inserted directly into the integral
bar algorithm of \ref{maps:src:D.2.41} or \ref{maps:src:D.2.46}. A vector $(0,b,c,d)$
contributes the actual regular-fibre representative of
$bw$, together with $c$ times the marked $\delta$ sweep
and $d$ times the marked $\ell$ sweep. The cap contributes
$\rho$ to the free relative rank in \eqref{maps:eq:D.GI10} and zero to
the Euler number. Thus all other already proved cap
choices and full finite Thom columns can be retained.

There is also a closed all-order result when this is the
sole logarithmic cap and every other cap has the minimal
natural marking of \ref{maps:src:D.2.26}, with no extra regular logs.
Write the initial coefficient as $b_0=4B$, and let $e$
be the ordinary clutching integer. Put
\[
                         s=N-2ke .                       \tag{D.SL6}\label{maps:eq:D.SL6}
\]
It is an odd nonzero integer, prime to $2k$. Then
\[
 \pi_1(X)=\mathbb Z/|s|,\qquad
 T=\operatorname{Tor}H_2(X)=
 \begin{cases}
 \mathbb Z/(2B|s|),&f=1,\\
 \mathbb Z/(B|s|),&f=2,\\
 \mathbb Z/2\oplus\mathbb Z/(2B|s|),&f=0.
 \end{cases}                                             \tag{D.SL7}\label{maps:eq:D.SL7}
\]
The other groups are
\[
\begin{gathered}
 H_2=\mathbb Z^{\rho-1}\oplus T,\qquad
 H_3=\mathbb Z^{2\rho-2}\oplus T,\\
 H_1=\mathbb Z/|s|,\qquad
 H_4=\mathbb Z^{\rho-1}\oplus\mathbb Z/|s|,\qquad H_5=0,
 \qquad H_0=H_6=\mathbb Z .
\end{gathered}                                             \tag{D.SL8}\label{maps:eq:D.SL8}
\]
Cyclic summands of order one are omitted. For example,
the six default natural rows with $b_0=4$ admit the
simply connected choice $s=\pm1$, $k=2h$, $j$ even.
It has
\[
          H_2=H_4=\mathbb Z^{nh},\qquad
                         H_3=\mathbb Z^{2nh}.              \tag{D.SL9}\label{maps:eq:D.SL9}
\]
Here $n=1$ on OS50/default and $n=3$ on OS46/default.
These are nonminimal star-log outputs, and no smooth
classification is inferred from their homology.

\end{theorem}


\begin{proof}
\begin{enumerate}[label=\arabic*.,leftmargin=*]
\item \emph{The actual unbounded affine actions.}
From \eqref{maps:eq:D.SL1}, $A^{2k}-I=q\alpha\beta^*$, so the
minimal semistable root model has exactly $q$ components.
This is the same integral Tate model used in \ref{maps:src:D.2.24}
and \ref{maps:src:D.2.27}, now at arbitrary root order. Its canonical
deck action reverses the component circle and fixes
the chosen zero component. Flat sections of the root
torus which are relevant to the affine action are
\[
                t=\frac{j}{q}\beta+\frac r k\lambda
                                      +\frac{N}{2k}\ell.
                                                               \tag{D.SL10}\label{maps:eq:D.SL10}
\]
They are single-valued because
$(A^{2k}-I)t=j\alpha$ is integral. The first term
translates the Mumford fan by $j$ components and
therefore extends over the toric cap. The other
terms are ordinary unit translations in the
transverse periods. Their full norm is
\[
         \sum_{a=0}^{2k-1}A^a t
                     =-r\alpha+(2r-j)\lambda+N\ell.
                                                               \tag{D.SL11}\label{maps:eq:D.SL11}
\]
It is integral, so composing the canonical deck
map with this translation still has order $2k$.
Using the inverse endpoint convention changes
only the coefficient of $\alpha$ in \eqref{maps:eq:D.SL11};
its collapsed vector \eqref{maps:eq:D.SL2} is unchanged, and
the corresponding exterior differences vanish
in the fibre coinvariant \eqref{maps:eq:D.SL4}.

\item The gamma displacement of a nonidentity deck
power is nonintegral because $N$ is prime to
$2k$. Hence the action is free on every stratum,
including fixed components and fixed double
curves. Conversely, in this canonical primitive-
gamma sector the flat translation conditions
give $q t_\beta\in\mathbb Z$, $k t_\lambda\in\mathbb Z$
and $2k t_\ell\in\mathbb Z$; a remaining angular
translation is removed by conjugation since
$A\alpha=-\alpha$. This is the marking \eqref{maps:eq:D.SL10}.
The usual root-and-unit isotopy keeps these
integral characters and all component meridians.
Thus the following topological calculation
applies to each actual action, not merely to
an abstract permutation of components.

\item \emph{The normalization and conductor at every order.}
The cap retracts equivariantly to its central
rank-one toric fibre. The quotient of its
$q$-gon dual graph by the reflection is an
interval. There are $\rho$ vertex orbits and
$\epsilon$ edge orbits, with
$\rho+\epsilon=q+1$; the number of fixed vertices
plus fixed edges is two. The normalization
and double-curve incidence graph is the
barycentric subdivision of this interval,
with exactly $q$ conductor incidences.

\item A paired vertex has stabilizer generated by
the square of the deck map. Its open
$\mathbb C^*$ bundle has an abelian period
lattice. The angular circle is primitive
because $\gcd(N,k)=1$. Thus its normalization
is a ruled surface over a torus with zero
Chern class, and its two pole sections have
the same class $\sigma$.
A fixed vertex instead has a pole-interchanging
clutch. Around the transverse $\lambda$ circle
the relation $A\lambda=\lambda-\alpha$ gives
the projective loop
$z\mapsto e^{-2\pi i u}/z$.
Its determinant winding is odd. Passing to
the primitive quotient gamma loop uses an
odd power of this map, since $N$ is odd, and
does not change the parity. Its normalization
is therefore the odd oriented $S^2$ bundle
over a torus. Choose a section $\sigma$ with
$\sigma^2=1$ and denote its complex sphere
fibre by $F$. Its connected conductor
bisection has class $2\sigma-F$; its
self-intersection is zero because the pole
normal line becomes trivial on the cover.

\item Every conductor flag has the even subgroup
as stabilizer. Its map to a normalization
base has degree two at a fixed vertex and
one at a paired vertex. Its map to the
double curve has degree two at a fixed edge
and one at a paired edge. Consequently
$H_2(U)$ has the following \emph{integral}
presentation. For each vertex orbit use
$F_v,\sigma_v$, and for each edge orbit use
its complex double-curve class $D_e$.
For each conductor incidence impose
\[
 \begin{cases}
 2\sigma_v-F_v-d_eD_e=0,&v\text{ fixed},\\
 \sigma_v-d_eD_e=0,&v\text{ paired},
 \end{cases}
 \qquad d_e=\begin{cases}2,&e\text{ fixed},\\1,&e\text{ paired}.
 \end{cases}                                             \tag{D.SL12}\label{maps:eq:D.SL12}
\]
There is no hidden degree-one bridge in this
normalization sequence: the conductor map
on $H_1$ is injective, by leaf elimination
on the incidence tree, since every local
isogeny is injective on $H_1$. The map on
$H_2$ is likewise injective. Hence
$H_3(U)$ is the direct sum of the
$\mathbb Z^2$ groups of the normalizations,
and $H_4(U)$ is their direct sum of top
classes.

\item All conductor section classes in \eqref{maps:eq:D.SL12}
are a single common class $S$. Internal
paired vertices have $\sigma_v=S$;
fixed vertices satisfy $2\sigma_v-F_v=S$;
a fixed endpoint edge satisfies $2D_e=S$.
Leaf elimination now gives free rank
$\rho+1$. If a vertex is fixed its sphere
coefficient is a unit and no torsion
survives. If no vertex is fixed, the two
fixed endpoint edges give
\[
 S=2D_0=2D_1,\qquad
       \tau=D_1-D_0\text{ of exact order two}.             \tag{D.SL13}\label{maps:eq:D.SL13}
\]
This proves the degree-two part of \eqref{maps:eq:D.SL3}.

\item For clarity, the fundamental group of the
cap follows directly from the free-cover
description after collapsing $\alpha$:
\[
\begin{gathered}
 \Gamma=\langle\beta,\lambda,\ell,g\mid
 [\beta,\lambda]=[\beta,\ell]=[\lambda,\ell]=1,\\
 g\beta g^{-1}=-\beta-n\lambda,\quad
 [g,\lambda]=[g,\ell]=1,\quad
                    2kg=R\lambda+N\ell\rangle .
\end{gathered}                                             \tag{D.SL14}\label{maps:eq:D.SL14}
\]
Additive notation in the last relation is
within the abelian subgroup generated by
$g,\lambda,\ell$. Its abelianization is free
of rank two: the relation $2\beta+n\lambda=0$
is primitive because $n$ is odd, and the
remaining relation is primitive because
$N$ is prime to $2k$. This proves $H_1(U)$
and finishes the local groups. As an
independent calculation, choose a primitive
quotient gamma circle. The cap is a mapping
torus of $I_q\times S^1_\lambda$; on its
component spheres the return is the
reflection, while on the actual specialized
torus $t=[\beta\lambda]$ naturality of the
exterior-square action gives
$t\mapsto-t-\sum_iF_i$. Its integral Wang
sequence gives the same groups, including
the order-two summand in \eqref{maps:eq:D.SL13}.

\item \emph{The regular-fibre map and its missing parity.}
The coinvariant relations \eqref{maps:eq:D.SL4} follow
by direct integral exterior-square reduction
of \eqref{maps:eq:D.SL1}; all four nonzero Smith factors
are units. The regular torus $\alpha\beta$
specializes to the sum of the $q$ complex
component spheres, each with degree one.
The transverse torus $\lambda\ell$ maps
to $kS$. Indeed its chosen component has
stabilizer of order $k$ in the paired case,
so its base cover has degree $k$ and its
class is $k\sigma$. In the fixed case its
base cover has degree $2k$; it is disjoint
from the bisection and has class
$k(2\sigma-F)$. Both give $kS$ in \eqref{maps:eq:D.SL12}.
Thus
\[
        2x=-\sum_v|v|F_v,\qquad 2w=-kS.                  \tag{D.SL15}\label{maps:eq:D.SL15}
\]
When $f>0$ the cap group is free, so these
equations determine the images uniquely.
Put
$T_0=\sum_{v\,\mathrm{paired}}F_v+
       \sum_{v\,\mathrm{fixed}}\sigma_v$.
For $f=1$ let $D_0$ be the sole fixed-edge
class. For $f=2$ retain the common $S$.
Then the images are
\[
\begin{array}{c|c|c}
 f&x&w\\\hline
1&-T_0+D_0&-kD_0\\
2&-T_0+S&-(k/2)S.
\end{array}                                               \tag{D.SL16}\label{maps:eq:D.SL16}
\]
For $f=0$, equations \eqref{maps:eq:D.SL15} alone leave
two genuine order-two ambiguities. We
resolve them using actual conductor loops.

\item Choose integers $a_0,b_0$ with
$a_0N+b_0(2k)=1$ and put
$\kappa=a_0g+b_0\ell$. The classifying
group \eqref{maps:eq:D.SL14} becomes the mapping-torus
group of the two-torus $(\beta,\lambda)$
with return
\[
           B=\begin{pmatrix}-1&0\\-n&1\end{pmatrix}.
\]
In fact
$\ell=2k\kappa-a_0R\lambda$ and
$g=N\kappa+b_0R\lambda$.
Here $a_0$ is odd. Its integral degree-two
group is
\[
 H_2(\Gamma)=\mathbb Zz\oplus\mathbb Z/2\,t,
       \qquad z=[\lambda\kappa],\quad t=[\beta\lambda].
                                                               \tag{D.SL17}\label{maps:eq:D.SL17}
\]
This group is computed by the same integral
bar comparison as \ref{maps:src:D.2.41}. For the two-dimensional
return $B$, its quadratic correction is
$f_B(u)=-n\binom{u_\beta}{2}t$.
Thus a $2k$-fold $\beta$ norm torus has
free part $-knz$ and torsion part $-knt$.
The actual regular tori therefore map as
\[
 [\beta\ell]\longmapsto-knz-(kn+a_0R)t,
 \qquad [\lambda\ell]\longmapsto2kz.                     \tag{D.SL18}\label{maps:eq:D.SL18}
\]

\item First choose $0\le j<q$ and label the two fixed endpoint edges so that
their stabilizer loops are $g$ and $\beta g$.
These are the actual two reflected edge
lifts, separated by half the dual circle.
Their primitive transverse periods are
$\kappa$ and
$\beta\kappa-\frac{n(a_0-1)}2\lambda$, respectively.
Their complex torus classes consequently
map to
\[
                        D_0\mapsto z,\qquad D_1\mapsto z-t.
                                                               \tag{D.SL19}\label{maps:eq:D.SL19}
\]
Both use the positive quotient gamma period,
so their orientations agree. It follows
that $\tau$ maps to the nonzero $t$.
After quotienting $H_2(U)$ by its component
spheres, \eqref{maps:eq:D.SL19} is an isomorphism onto
\eqref{maps:eq:D.SL17}. In particular it detects precisely
the torsion ambiguity left in \eqref{maps:eq:D.SL15}.

\item Since $k$ is even and $j,R,a_0$ are odd
in this case, \eqref{maps:eq:D.SL18} sends the primitive
$w$ of \eqref{maps:eq:D.SL4} to $-kz+t$; it sends $x$
to $t$. Hence the exact, marked images are
\[
       x=-\sum_vF_v+\tau,\qquad
                         w=-kD_0+\tau.                   \tag{D.SL20}\label{maps:eq:D.SL20}
\]
Both order-two terms are forced by the
actual classifying map. In particular the
fibre cokernel has cyclic torsion
$\mathbb Z/(2k)$ in this case; replacing
the second image by $-kD_0$ would give a
different marked map.

\item \emph{The literal fixed-normal tori.}
The root meridian is $g^{2k}v^{-1}$,
with \eqref{maps:eq:D.SL11}. Sweeping its bounding
normal disc by $\delta$ and $\ell$ gives
the exact relations in $H_2(U)$
\[
                 2kD+4Nw=0,\qquad
                             2kL-2Rw=0.                  \tag{D.SL21}\label{maps:eq:D.SL21}
\]
The omitted $\alpha$ terms vanish by
\eqref{maps:eq:D.SL4}. These relations determine the
free parts of both images. When $f>0$
there is no torsion ambiguity. When $f=0$,
the classifying map fixes it as well:
the cap loops are $\delta=2\lambda$,
$g=N\kappa+b_0R\lambda$ and
$\ell=2k\kappa-a_0R\lambda$, so their
actual commuting tori map to
\[
                 D\mapsto2Nz,\qquad L\mapsto-Rz
                            \quad\text{in }H_2(\Gamma).
\]
There is no $t$ term. In the same notation
as \eqref{maps:eq:D.SL16}--\eqref{maps:eq:D.SL20}, the full images are
\[
\begin{array}{c|c|c}
f&D&L\\\hline
1&2ND_0&-RD_0\\
2&NS&-(R/2)S\\
0&2ND_0&-RD_0.
\end{array}                                               \tag{D.SL22}\label{maps:eq:D.SL22}
\]
Together with the fibre images these
prove \eqref{maps:eq:D.SL5} by solving an integral
linear relation. For $f=0$ the coefficient
of $x$ must vanish in the independent
component-sphere direction; the coefficient
of $w$ must then be even by \eqref{maps:eq:D.SL20}.
The free coefficient gives the displayed
equation. This explicitly retains the
relative torsion information rather than
saturating the normal lattice.

\item The cokernel of the full boundary map is
free of rank $\rho-1$, as one sees directly
from \eqref{maps:eq:D.SL12} and \eqref{maps:eq:D.SL16}--\eqref{maps:eq:D.SL22}.
Equivalently, Poincare--Lefschetz duality
puts it into the free group $H^4(U)$.
The projected-kernel indices are $k,k/2,2k$;
the computations for $m=2$ and the swapped
$m=4$ in \ref{maps:src:D.2.27} are recovered exactly.
The full local kernel can therefore be
used in the global bar presentation,
with all its existing endpoint corrections.

\item The integer lift of $j$ must also be retained.
Write $j=j_0+hq$, $0\le j_0<q$. The corresponding
peripheral generator changes from $g_0$ to
$g=\beta^h g_0$, and $R=R_0-hq$.
Its literal normal tori satisfy
\[
                    D=D_0^{\rm norm},\qquad
                    L=L_0-hn w.                           \tag{D.SL24}\label{maps:eq:D.SL24}
\]
Here the first symbol denotes the old normal
torus, not a double-curve class. Indeed
$[\delta\beta]=0$ and $[\beta\ell]=nw$
in \eqref{maps:eq:D.SL4}. Substituting \eqref{maps:eq:D.SL24} into
the kernel formulas replaces $R_0$ by
$R$ exactly. In the $f=0$ case it changes
the displayed image of $L$ in \eqref{maps:eq:D.SL22} to
$-RD_0+h\tau$, but the equation in
\eqref{maps:eq:D.SL5} already forces its coefficient
$d$ to be even; the kernel parity is
therefore still precisely $b$ even.
Thus integer affine lifts are included
without suppressing their actual torsion
correction.

\item \emph{The closed one-star Smith calculation.}
For the natural minimal exterior, the
fibre functional has
$\Omega(x)=-a$ and $\Omega(w)=b_0/4=B$.
The primitive graph calculation has the
single diagonal relation $D=-b_0e\,s_0$,
where $s_0$ is its regular-fibre generator.
Because every vector in \eqref{maps:eq:D.SL5} has zero
$x$ coordinate, its global relation on
$(s_0,L)$ is
\[
                         (B(b-4ec),d).                   \tag{D.SL23}\label{maps:eq:D.SL23}
\]
Set $u=b-4ec$ and $s=N-2ke$.
For $f=1$ its defining condition becomes
$ku+Rd=2sc$. Since $k$ is odd and prime
to $2s$, $d$ is unrestricted and
$u$ is a fixed residue modulo $2s$.
The Smith cokernel is cyclic of order
$2B|s|$.
For $f=2$, the condition is
$(k/2)u+(R/2)d=sc$.
Here $k/2$ is prime to $s$, so the same
argument gives a cyclic cokernel of
order $B|s|$.
For $f=0$, $u$ and $d$ are even; write
$u=2u_0,d=2d_0$. The condition becomes
$ku_0+Rd_0=sc$. Its relation matrix has
Smith factors $2$ and $2B|s|$.
This proves the three torsion groups
in \eqref{maps:eq:D.SL7} for every parameter.

\item The one-log Seifert presentation gives
$\pi_1=\mathbb Z/|s|$. Its rational
Euler number cannot vanish because
$N$ is prime to $2k$. The free local
boundary defect is $\rho-1$ and there
is no additional rational degree-one
bridge. All other caps in the stated
corollary have their minimal ranks.
Thus $b_2=\rho-1$. The unique primitive
rank-two zero still contributes Euler
number two, while this star cap has
Euler number zero. Integral duality
gives \eqref{maps:eq:D.SL8}, and the simply connected
specialization gives \eqref{maps:eq:D.SL9}.
\end{enumerate}
\end{proof}

The exact script \path{code/III/manifolds_star_log_caps.py}
retains the twelve actual canonical bases,
the full normalization matrices, the
literal fibre and normal maps, and their
complete integral kernels. It checks
the conductor calculation independently
against the mapping-torus sequence and
compares the global Smith formula on
unbounded symbolic lattices and a finite
set of exact instances. The proof above,
including the order-two marking, is
independent of the finite parameter
bounds used in those corroborating runs.
Root independently read the full proof and replayed all
$12$ actual marked bases, $95$ normalization comparisons
and $1638$ global Smith calculations successfully.
An independent adversarial review by the fourfolds
agent also checked \eqref{maps:eq:D.SL17}--\eqref{maps:eq:D.SL24}, including the
nonzero order-two correction and the integer phase lift.
Root actions with nonprimitive gamma
power, other integral star blocks and
positive-weight star degenerations are
outside this theorem's stated scope.


\begin{theorem}[all vertical corrections and new simply connected fillings]\label{maps:src:D.2.50}
\StatementRecord{maps:src:D.2.50}{proved}{Historical events: Claude Fable 5.2 (Anthropic), effort high, 16 November 2026: carried as background without fault found. GPT-6 Astra (OpenAI), ChatGPT Pro mode, effort Pro, S6\_09\_PAPERS, 17 November 2026: targeted proof check. Also checked in this writing. Pre-continuation writing events: GPT-6 Astra (OpenAI), ChatGPT Pro mode, effort Pro, S6\_10\_PAPERS, 18 November 2026: targeted proof check. Continuation metadata events: GPT-6 Astra (OpenAI), ChatGPT Pro mode, effort Pro, S6\_10\_PAPERS continuation, 19 November 2026: used and rechecked in this run; status-vocabulary reconciliation; metadata reconciliation only; no new mathematical check. \emph{Pre-continuation recorded proof scope.} GPT-6 Astra (OpenAI, ChatGPT Pro mode), effort Pro, 2026-11-17 UTC; GPT-6 Astra (OpenAI, ChatGPT Pro mode), S6\_06\_BLIND, effort Pro, 2026-11-17; Claude Fable 5.2 (Anthropic), effort high, 16 November 2026; also checked in this writing by GPT-6 Astra (OpenAI), ChatGPT Pro mode, effort Pro, 18 November 2026}{the hypotheses, notation and definitions stated here; \cref{maps:eq:D.GI2}; \cref{maps:eq:D.VC3}; \cref{maps:src:D.2.19}; \cref{maps:src:D.2.46}; \cref{maps:eq:D.VC8}; \cref{maps:eq:D.VC6}; \cref{maps:eq:D.VC2}; \cref{maps:src:D.2.43}; \cref{maps:eq:D.GI3}; \cref{maps:eq:D.GI5}; \cref{maps:src:D.2.41}; \cref{maps:eq:D.GI6}}


Let $f:S\to\mathbb P^1$ be one of the twenty-one recorded
rational elliptic surfaces with Mordell--Weil group $\mathbb ZG$,
zero section $O$, and height $\langle G,G\rangle=1/d_0$.
For integers $a,b$, put $P=aG$, $Q=bG$ and
\[
 M_D=\mathcal O_S(P-O+D),\qquad
 J_D=M_D\otimes T_Q^*M_D^{-1},
\]
where $D$ is an integral divisor supported on fibre components.
Here $T_Q:S\to S$ is the actual section-translation automorphism.
Some such $D$ makes
$J_D$ the pullback of a base line bundle if and only if
$d_0\mid ab$. For every such correction the exact identity is
\[
                  J_D\simeq f^*\mathcal O_{\mathbb P^1}(ab/d_0).
                                                        \tag{D.VC1}\label{maps:eq:D.VC1}
\]
Thus the compatible initial line bundles are not exhausted
by requiring $P$ or $Q$ narrow at every fibre.
All corrections are determined by the integral local equations
below, with precisely the invariant component orbit sums as ambiguity.
For $a,b>0$ and any nonnegative whole-fibre weights of total $ab/d_0$,
the corrected punctured family has exactly the marked monodromies
\eqref{maps:eq:D.GI2}. This does not identify its newly permitted original caps.

There are additional explicitly marked simply connected examples.
On OS45/default, choose $(a,b)=(2,4)$ or $(4,2)$, a simple
linearization zero at any of its four $I_1$ fibres, multiplicative
Mumford fillings at every singular fibre, no extra affine boundary
translations, and ordinary gamma clutch $e=\pm1$. Use the actual
vertical corrections specified in the proof. All sixteen choices have
\[
 \begin{gathered}
 \pi_1(X)=1,\qquad H_0(X)=H_6(X)=\mathbb Z,\\
 H_2(X)=H_3(X)=\mathbb Z/b,\qquad
 H_1(X)=H_4(X)=H_5(X)=0 .
 \end{gathered}                                             \tag{D.VC7}\label{maps:eq:D.VC7}
\]
These are new initial and marked filling data beyond the earlier
narrowness condition. Their smooth diffeomorphism classes are not
inferred from these groups.

\end{theorem}


\begin{proof}
\begin{enumerate}[label=\arabic*.,leftmargin=*]
\item \emph{The integral local equation.}
Index the components of a singular fibre by $C_0,\ldots,C_{v-1}$,
with $C_0$ the component meeting $O$. Let $m$ be its positive
primitive multiplicity vector, $J=(C_iC_j)$ its intersection
matrix and $\sigma$ the permutation matrix induced by $Q$.
If $P$ meets component $C_p$, the actual multidegree of
$M_D$ on this fibre is
\[
                    c_D=e_p-e_0+Jz,
                                                        \tag{D.VC2}\label{maps:eq:D.VC2}
\]
where $z$ is the vector of coefficients of $D$ there.
All section components have multiplicity one, so $m^tc_D=0$.
The vanishing of the component degrees of $J_D$ is exactly
\[
                 (\sigma-I)Jz=-(\sigma-I)(e_p-e_0).
                                                        \tag{D.VC3}\label{maps:eq:D.VC3}
\]
Using the inverse permutation for pullback gives the same
invariance condition and the same solution set.
The equation is integral: it may not be replaced by its
rational solution or by the saturation of its image.

\item For completeness, its entire ambiguity also has an exact
description. If $z,z'$ solve \eqref{maps:eq:D.VC3}, then
$J(\sigma-I)(z-z')=0$, since $J$ commutes with the permutation.
The integral kernel of $J$ is $\mathbb Zm$, so
$(\sigma-I)(z-z')=h m$ for some integer $h$.
Apply $I+\sigma+\cdots+\sigma^{q-1}$, where $\sigma^q=I$.
Because $\sigma m=m$, the result is $0=qh m$, hence $h=0$.
Conversely every invariant vector can be added to a solution.
Thus all corrections are one particular solution plus
\[
                         \ker(\sigma-I)\subset\mathbb Z^v.
                                                        \tag{D.VC4}\label{maps:eq:D.VC4}
\]
In particular the ambiguity consists of integral orbit sums
of components, including the whole-fibre ambiguity.

\item \emph{Why vanishing component degrees is sufficient globally.}
On the generic elliptic fibre, translation acts trivially
on $\operatorname{Pic}^0$, so $J_D$ is generically trivial.
It therefore has a rational section whose divisor is vertical:
one extends a generic trivialization as a rational section
on the smooth surface. If \eqref{maps:eq:D.VC3} holds at every fibre,
this vertical divisor has intersection zero with every
component. On each connected fibre the negative-semidefinite
intersection matrix has kernel exactly its multiplicity
vector. The divisor on that fibre is an integral multiple
of the whole fibre, since the multiplicities have gcd one.
Consequently $J_D=f^*\mathcal O(t)$ for an integer $t$.
This proves sufficiency for the actual line bundle, not
only for its rational first Chern class. Necessity follows
by restricting a pullback to the components.

\item We determine $t$ without suppressing a local correction.
Choose $h>0$ with $hQ$ narrow at every singular fibre.
Then $T_{hQ}$ fixes each component divisor. Telescoping
the $h$ translated copies of $J_D$ gives
\[
 \bigotimes_{j=0}^{h-1} T_Q^{j*}J_D
       =M_D\otimes T_{hQ}^*M_D^{-1}
       =\mathcal O(P-O)\otimes
                    T_{hQ}^*\mathcal O(P-O)^{-1}.          \tag{D.VC5}\label{maps:eq:D.VC5}
\]
The left side is $f^*\mathcal O(ht)$, since $T_Q$ preserves
the base. The last expression is in the already proved
narrow-translation case of \ref{maps:src:D.2.19} and \ref{maps:src:D.2.46}. Its degree
is the integral height pairing
$\langle P,hQ\rangle=hab/d_0$.
Pullback on base Picard groups is injective, by restriction
to the zero section. Hence $ht=hab/d_0$, proving \eqref{maps:eq:D.VC1}
and the necessary divisibility $d_0\mid ab$.

\item \emph{The actual punctured monodromies.}
Let the chosen base linearization have whole-fibre weights $w_i$.
Choose $h$ as above and iterate the actual contracting generator
$h$ times. The canonical nonzero section of $\mathcal O(D)$
off the singular fibres telescopes in this iteration, because
$T_{hQ}D=D$ as an actual divisor. Thus the iterated punctured
family is the uncorrected narrow-translation construction with
sections $P,hQ$ and weights $h w_i$. A constant nonzero scalar
in this identification changes no monodromy.

\item In the actual period lattice its long generator is
$\gamma_h=h\gamma$; the elliptic generators and $\delta$ are
unchanged. If $A_i$ denotes the original monodromy and $A_{i,h}$
the iterated one, put $S=\operatorname{diag}(1,1,1,h)$. Then
\[
                         A_iS=SA_{i,h}.                    \tag{D.VC8}\label{maps:eq:D.VC8}
\]
The narrow case \eqref{maps:eq:D.GI2} has $q_{i,h}=h q_i$,
$z_{i,h}=h z_i$, and central entry $r_i(hz_i)+h w_i$.
Solving \eqref{maps:eq:D.VC8} gives exactly the unscaled matrix \eqref{maps:eq:D.GI2}.
This proves its integrality in every corrected case as well.
The argument concerns the punctured family and its actual long
period. It does not assume a covering of a filled original or
logarithmic cap, nor suppress any subsequently chosen affine
boundary marking.

\item \emph{Solvability for every recorded fibre.}
For an $I_n$ fibre the component lattice of total degree
zero modulo the image of $J$ is $\mathbb Z/n$; its quotient
map is $c\mapsto\sum_i i c_i\pmod n$.
This follows by the integral successive-difference reduction
of the cycle Laplacian, whose nonzero Smith factors are
$1,\ldots,1,n$. For $n=1$ the lattice is zero; for $n=2$
the two distinct edges are both retained.
Write $q$ for the component of $Q$, $g=\gcd(n,q)$ and
$t=n/g$. A $q$-translation-invariant multidegree has one
coefficient $u_j$ on each of its $g$ orbits. Its total degree
is zero exactly when $\sum_j u_j=0$, and its component
class is $t\sum_j j u_j\pmod n$.
Thus the subgroup represented by invariant multidegrees
is exactly $t\mathbb Z/n$. The class of $e_p-e_0$ lies in it
if and only if $t\mid p$, equivalently $n\mid pq$.
This proves integral solvability of \eqref{maps:eq:D.VC3} for every
multiplicative fibre.

\item The only other types in the twenty-one profiles are
$II,III,IV,III^*,IV^*,I_1^*,I_3^*$.
The respective component-group orders are $1,2,3,2,3,4,4$.
For these finite lists, the exact marked intersection
matrices and component permutations are supplied in
\path{code/III/root_vertical_corrections.py}.
Their multiplicity kernels and permutation invariance are
checked there before any equation is solved. The integral
Smith calculation for every $p,q$ in each component group
gives the same criterion
\[
              \operatorname{ord}(\Phi_f)\mid pq.           \tag{D.VC6}\label{maps:eq:D.VC6}
\]
For $III$ and $IV$ this is also the same cycle-matrix
calculation just given. The remaining star cases involve
only $4$, $9$, $16$ and $16$ pairs respectively; their
full solutions and all integer kernel generators are
retained in the output. This is an exhaustive finite
proof for those types, not a sampling of section multiples.

\item For every recorded profile the component of $G$ is a unit
in each nontrivial component group, and the least common
multiple of their orders is $d_0$. These finite marked
data are checked in the same executable against the
twenty-one source records. Hence \eqref{maps:eq:D.VC6} for $P=aG,Q=bG$
at all fibres is equivalent to $d_0\mid ab$.
Choose a Smith solution independently at each fibre and
sum the corresponding component divisors. This constructs
the required global $D$ and proves sufficiency.
All integers $a,b$ are covered because their component
classes depend only on their residues modulo $d_0$.
The supplemental check enumerates all $5696$ such residue
pairs, using $227$ distinct local component pairs, and
finds no discrepancy with \eqref{maps:eq:D.VC1}.

\item \emph{A concrete correction and the remaining topology.}
At an $I_4$ fibre with $P$ and $Q$ both meeting component
two, the original multidegree is $(-1,0,1,0)$.
Taking $D=-C_0$ changes it to $(1,-1,1,-1)$, which is
invariant under translation by two. Neither section is
narrow. The full solution set is obtained by adding
integer multiples of $C_0+C_2$ and $C_1+C_3$, as the
actual Smith kernel confirms.
Globally the profiles with such additional residue pairs
are precisely OS45, OS46, OS50 and OS55, with the listed
additive variants retained. The new initial data cannot
be discarded merely because they violate the earlier OR
condition. Their original caps carry the actual corrected
multidegree \eqref{maps:eq:D.VC2}, which can be nonprimitive and can
be permuted nontrivially by $Q$. Computing those marked
cap maps is a separate remaining part of D.2.

\item \emph{The sixteen new all-Mumford examples.}
The actual OS45/default data are
\[
 \begin{gathered}
 (I_8,I_1,I_1,I_1,I_1),\qquad d_0=8,\qquad
 \overline G_{I_8}=3\pmod8,\\
 g_i=(3,0),(0,0),(2,-1),(1,-2),(0,0),\\
 (B_0,\ldots,B_4)=\begin{pmatrix}1&8\\0&1\end{pmatrix},
 \begin{pmatrix}1&1\\0&1\end{pmatrix},
 \begin{pmatrix}3&4\\-1&-1\end{pmatrix},
 \begin{pmatrix}3&1\\-4&-1\end{pmatrix},
 \begin{pmatrix}0&1\\-1&2\end{pmatrix}.
 \end{gathered}                                             \tag{D.VC9}\label{maps:eq:D.VC9}
\]
At $I_8$ one may take the following literal coefficient vectors:
\[
\begin{array}{c|c|c}
 (a,b)&(D_{C_0},\ldots,D_{C_7})&c_D\\ \hline
 (2,4)&(-1,-1,-1,0,0,0,0,0)&(0,0,1,-1,0,0,1,-1)\\
 (4,2)&(-1,0,0,1,1,1,0,0)&(1,-1,1,-1,1,-1,1,-1).
\end{array}                                                \tag{D.VC10}\label{maps:eq:D.VC10}
\]
They are invariant respectively under translation by four and by six;
no correction is needed at the four irreducible fibres. Both height
pairings equal one. Use the punctured marking \eqref{maps:eq:D.VC8}, one weight one
at a selected $I_1$, weight zero elsewhere, and exactly the affine
boundary data in the statement.

\item Each positive-weight cap is the primitive index-one triangular cap
of \ref{maps:src:D.2.43}. Every other cap is the rank-one Mumford cap of its actual
matrix. Thus all local filling kernels in \eqref{maps:eq:D.GI3} and \eqref{maps:eq:D.GI5} apply
without a narrowness assumption: their proofs concern these toric
monodromies and their marked primitive angular discs. The common
piece and every ordinary boundary map are the same integral bar
models as in \ref{maps:src:D.2.41}. No newly corrected original cap is used.

\item For completeness the entire finite calculation is retained in
\path{code/III/fourfolds_corrected_Mumford_examples.py} and its JSON output.
The input range is exactly the two pairs, four zero positions,
and two clutch signs specified above. In all sixteen cases the
unsaturated normal closure of the toric vanishing vectors is the
whole vertical lattice $\mathbb Z\langle e_1,e_2,\delta\rangle$.
All singular peripheral loops are killed by their Mumford fillings;
the remaining ordinary relation kills $\gamma$ with coefficient
$\pm1$. This proves full simple connectivity, not only $H_1=0$.

\item The degree-two bar presentation has $23$ generators and $70$
relations. Its exact Smith form has $22$ unit factors and last
factor $b$. The additional relative free rank in \eqref{maps:eq:D.GI6} is zero:
there are five component classes and the finite-filled exterior
has first Betti number five. The only Euler contribution is the
index-one positive cap, so $\chi=2$. Integral duality and universal
coefficients now give every group in \eqref{maps:eq:D.VC7}.
These actual matrices, vertical corrections, and boundary markings
are retained for review. Invariant orbit sums added to $D$ are not
silently identified with these marked filled spaces; their induced
boundary rigidifications must also be transported.
\end{enumerate}
\end{proof}

The independent certificate
\path{code/fourfolds_vertical_correction_audit.py}
checks all $227$ local pairs with unimodular Smith transformations,
all $5696$ global residues, and the $82$ actual elliptic
monodromy/section-cocycle orders against the recorded component
maps. It also records all sixteen additional positive height-one
section pairs across the full profile list. Those latter sixteen
pairs are initial data, distinct from the sixteen fully filled
OS45 examples in \eqref{maps:eq:D.VC7}.


\begin{theorem}[invariant-Chern original additive caps and their integral extension]\label{maps:src:D.2.51}
\StatementRecord{maps:src:D.2.51}{proved}{Historical events: Claude Fable 5.2 (Anthropic), effort high, 16 November 2026: carried as background without fault found. GPT-6 Astra (OpenAI), ChatGPT Pro mode, effort Pro, S6\_09\_PAPERS, 17 November 2026: targeted proof check. Also checked in this writing. Pre-continuation writing events: GPT-6 Astra (OpenAI), ChatGPT Pro mode, effort Pro, S6\_10\_PAPERS, 18 November 2026: targeted proof check. Continuation metadata events: GPT-6 Astra (OpenAI), ChatGPT Pro mode, effort Pro, S6\_10\_PAPERS continuation, 19 November 2026: used and rechecked in this run; status-vocabulary reconciliation; metadata reconciliation only; no new mathematical check. \emph{Pre-continuation recorded proof scope.} GPT-6 Astra (OpenAI, ChatGPT Pro mode), effort Pro, 2026-11-17 UTC; Claude Fable 5.2 (Anthropic), effort high, 16 November 2026; also checked in this writing by GPT-6 Astra (OpenAI), ChatGPT Pro mode, effort Pro, 18 November 2026}{the hypotheses, notation and definitions stated here; \cref{maps:src:D.2.50}; \cref{maps:src:D.2.31}; \cref{maps:src:D.2.16}}


Let $K$ be the reduced central fibre of an original additive
Kodaira neighbourhood, of type
$II,III,IV,IV^*,III^*,I_1^*$ or $I_3^*$.
Write its rational components as $C_i$, its primitive
multiplicity vector as $\mu=(\mu_i)$, and let $\sigma$
be the actual component permutation induced by the section $Q$.
Let the original line bundle have invariant component degrees
$c_i$, so that
\[
                  \sigma c=c,\qquad \sum_i\mu_i c_i=0.
                                                        \tag{D.IC1}\label{maps:eq:D.IC1}
\]
These include every corrected line bundle in \ref{maps:src:D.2.50}.
Let the component orbits be $O_1,\ldots,O_\rho$, with
lengths $l_j$ and common degrees $c_j$ on $O_j$.
For $c\ne0$ put
\[
                  g=\gcd_j(c_j),\qquad
                  h=\gcd_j(l_jc_j),
                                                        \tag{D.IC2}\label{maps:eq:D.IC2}
\]
where both gcds are positive. For the resulting original
six-dimensional cap $U$, there is an exact integral cellular
cone, with its actual circle-fibre and regular-area classes,
\[
\begin{gathered}
 C_0=\mathbb Z,\quad C_1=\mathbb Z^2,\quad
 C_2=\mathbb Z^v\oplus\mathbb Zt,\quad
 C_3=\mathbb Z^v\oplus\mathbb Z^v,\quad C_4=\mathbb Z^v,\\
 d_2=\begin{pmatrix}c^t&0\\0&0\end{pmatrix},\qquad
 d_3=\begin{pmatrix}0&\sigma-I\\0&c^t\end{pmatrix},\qquad
 d_4=\begin{pmatrix}I-\sigma\\0\end{pmatrix}.
\end{gathered}                                             \tag{D.IC3}\label{maps:eq:D.IC3}
\]
The degree-one generators are the scalar circle $\delta$
and a positive long circle $\gamma$. In particular
\[
\begin{array}{c|ccccc}
i&0&1&2&3&4\\\hline
H_i(U;\mathbb Z)&\mathbb Z&
\mathbb Z\oplus\mathbb Z/g&
\mathbb Z^{\rho-1}\oplus\mathbb Z/h&
\mathbb Z^{2\rho-1}&\mathbb Z^\rho .
\end{array}                                               \tag{D.IC4}\label{maps:eq:D.IC4}
\]
Moreover $\pi_1(U)=\mathbb Z\oplus\mathbb Z/g$.
When $c=0$, these groups are instead
\[
 (H_0,H_1,H_2,H_3,H_4)=
 (\mathbb Z,\mathbb Z^2,\mathbb Z^{\rho+1},
                         \mathbb Z^{2\rho},\mathbb Z^\rho).
                                                        \tag{D.IC5}\label{maps:eq:D.IC5}
\]
The actual regular-fibre torus $\delta\wedge\gamma$
maps to the class of $t$ in \eqref{maps:eq:D.IC3}, which generates
\emph{all} of the cyclic torsion $\mathbb Z/h$ when
$c\ne0$. The elliptic tori crossed with $\delta$ map
to zero. The complex elliptic area maps to the literal
cycle $(\mu,0)$ in \eqref{maps:eq:D.IC3}.

Let $B$ be the elliptic monodromy, $q$ its marked
$Q$-cocycle, and $z=(B-I)^{-1}q$. Put
\[
        q_Q=\gcd_j(l_j\mu_j)
             =\operatorname{ord}(\overline Q).
                                                        \tag{D.IC6}\label{maps:eq:D.IC6}
\]
The actual logarithm $z$ has exact denominator $q_Q$.
For $c\ne0$, the full regular-fibre image in $H_2(U)$
is $\mathbb Z\oplus\mathbb Z/h$ with free cokernel
of rank $\rho-2$. For $c=0$ it is a primitive
$\mathbb Z^2$ with free cokernel of rank $\rho-1$.
In both cases this is already the full boundary image.
Consequently the full boundary kernel projects onto
all of $\operatorname{Fix}(A)$ in the Wang sequence;
there is no omitted finite normal index.

If $h=1$, these facts determine the complete marked
local kernels. In the regular basis
$(e_1,e_2,\delta,\gamma)$ the actual fibre kernel is
\[
 \ker\bigl(q_Q(1,0,z_2,0,-z_1,0):\mathbb Z^6\to\mathbb Z\bigr),
                                                        \tag{D.IC7}\label{maps:eq:D.IC7}
\]
where the six coordinates have their usual increasing
exterior order. Both literal fixed-normal tori bound.
The cap kills every vertical fundamental-group
coordinate and retains the primitive $\gamma$.
This includes the canonical new corrected star degrees
\[
 \begin{array}{c|c}
 I_1^*&(1,1,0,0,-1,0)\\
 I_3^*&(1,1,0,0,-1,0,0,0)
 \end{array}
                                                        \tag{D.IC8}\label{maps:eq:D.IC8}
\]
for component classes $P=Q=2$ in $\mathbb Z/4$.
Their actual permutation has two paired outer orbits
and fixed inner components, and $g=h=1$.

For $h>1$, \eqref{maps:eq:D.IC3}--\eqref{maps:eq:D.IC6} do not by themselves specify
the torsion corrections of the two mixed
elliptic--gamma tori and the second fixed-normal torus.
The theorem does not replace those marked columns by
their saturation. In particular it does not supply an
unrestricted global gluing algorithm for those caps.

\end{theorem}


\begin{proof}
\begin{enumerate}[label=\arabic*.,leftmargin=*]
\item \emph{The actual original cap and the circle bundle.}
The reduced additive fibre is simply connected and
is homotopy equivalent to a wedge of its $v$ rational
component spheres. The tangency in type $III$ and the
triple point in type $IV$ do not add a graph loop;
the same assertion follows directly by taking the
union of the rational components and their actual
connected intersections. The other listed fibres
have a tree of rational components.
The elliptic neighbourhood retracts onto this fibre.
The unit original linearization makes its compact
cap the mapping torus of the principal scalar-circle
bundle $P_c\to K$, with the lift of $Q$.
This is the same actual original filling used in
\ref{maps:src:D.2.31}--\ref{maps:src:D.2.32}, now retaining the full degree vector.

\item A principal circle bundle over this wedge is determined
by its $v$ Chern evaluations. Its cellular complex has
one zero-cell, the scalar one-cell, $v$ two-cells with
boundary degrees $c_i$, and $v$ three-cells with zero
boundary. The lift of $Q$ acts by the component
permutation on the two- and three-cells. There is no
additional homotopy class of a gauge transformation:
$[K,S^1]=H^1(K;\mathbb Z)=0$. More explicitly, after
choosing a fixed point of the induced component-tree
map and a fibre point there, the component maps are
homotopic to their orientation-preserving permutations.
The invariant vector $c$ lifts this homotopy to the
circle bundle. A comparison unit has a continuous
logarithm on $K$, so its lift is homotopic through
bundle maps to the identity. A constant scalar at
the chosen point is removed by a scalar isotopy.
This fixes the scalar circle and its product with
the mapping-torus direction throughout the comparison.
The signs in \eqref{maps:eq:D.IC3} use the literal product orientation
$t=\delta\wedge\gamma$: the boundary of a two-cell
crossed with the long interval contributes $+c_i t$.
This is the usual mapping cone after changing the signs
of its shifted summands. Thus \eqref{maps:eq:D.IC3} is an actual model, not
merely a rational model with the same permutation.

\item The circle-bundle homotopy sequence gives
$\pi_1(P_c)=\mathbb Z/g$ for $c\ne0$.
The return fixes the scalar generator; the mapping-torus
extension by the infinite cyclic group therefore splits.
This proves the fundamental-group statement and $H_1$.

\item \emph{The integral extension.}
In $C_2/\operatorname{im}d_3$, successive cells in an
orbit have the relations
\[
                  e_{i+1}-e_i+c_jt=0.
                                                        \tag{D.IC9}\label{maps:eq:D.IC9}
\]
Eliminating them leaves one free $e_j$ for each orbit,
and the remaining relations are $l_jc_jt=0$.
Hence this quotient is
$\mathbb Z^\rho\oplus\mathbb Z/h$.
Its remaining boundary sends $e_j$ to $c_j\delta$
and sends $t$ to zero. Its kernel is therefore
$\mathbb Z^{\rho-1}\oplus\mathbb Z/h$.
This proves the full nonsplit extension in $H_2$.
For example, two interchanged components of degree
$2$ contribute a relation $4t=0$, not two unrelated
order-two summands. The fibre product torus
$\delta\wedge\gamma$ is the cone generator $t$ itself.

\item The degree-three cone has the direct sum of the
coinvariants of the component three-cells, of rank
$\rho$, and the invariant component two-cycles with
Chern evaluation zero, of rank $\rho-1$.
Both groups are free. Degree four is the invariant
component three-cell group of rank $\rho$.
The same calculation with $c=0$ gives \eqref{maps:eq:D.IC5}.

\item There is also an intrinsic check of the extension.
Let $Y$ be the mapping torus of $K$ and $Q$.
The cap is the principal $S^1_\delta$ bundle over $Y$.
The Euler class evaluates as $c_j$ on the orbit
representatives in $H_2(Y)$, and as $l_jc_j$ on
invariant orbit sums in $H_3(Y)$ paired with the
positive $\gamma$. The integral Gysin sequence
therefore injects $\mathbb Z/h$ into $H_2(U)$ by
$\gamma\mapsto\delta\wedge\gamma$.
The quotient is the free kernel of the first Chern
functional on $H_2(Y)$. This verifies both the
extension and the actual generator without a
choice of an abstract splitting.

\item \emph{The regular area and the full boundary image.}
In the single circle bundle $P_c\to K$ the map
$H_2(P_c)\to H_2(K)$ is injective, with image
$\ker(c^t)$. The regular elliptic area projects to
$\mu$, the multiplicity cycle of the central fibre,
and hence has its unique lift $(\mu,0)$.
The elliptic loops crossed with the scalar circle
vanish because $H_1(K)=0$.
After passing to $Y$, this area projects to the
orbit vector $(l_j\mu_j)$. Its content is $q_Q$.
The equality with the order of the component of
$Q$ follows directly from the marked component
permutations and multiplicities for the seven listed
types; the retained exact certificate checks every
component class. The actual section-cocycle model
has the same order in $\operatorname{coker}(B-I)$,
so the logarithm $z$ has exactly this denominator.

\item To identify the rational mixed area coefficients,
pass to the $q_Q$-fold long-circle cover.
The section $q_QQ$ is narrow; its marked local
logarithm straightens the translation through the
original elliptic neighbourhood, as in \ref{maps:src:D.2.16} and
\ref{maps:src:D.2.31}. Its corrected long direction is
$q_Q(\gamma-z)$. An elliptic loop crossed with this
direction is zero in the base mapping torus because
the elliptic loop bounds in the Kodaira neighbourhood.
Thus the two mixed free coefficients are $z_2$ and
$-z_1$ times the elliptic-area coefficient.
The integers $q_Q,q_Qz_1,q_Qz_2$ have gcd one.
Consequently these actual fibre tori generate the
primitive free orbit vector
$(l_j\mu_j)/q_Q$.
The torus $\delta\wedge\gamma$ already generates
all cap torsion. This proves the claimed fibre
image and its free cokernel.

\item Finally $H_2(U,\partial U)=H^4(U)=\mathbb Z^\rho$,
since $H_3(U)$ is free. The boundary first Betti
number is three, while that of $U$ is one for
$c\ne0$ and two for $c=0$. The boundary maps onto
$H_1(U)$: it contains the scalar and long generators.
The relative sequence therefore makes the full
boundary degree-two image have rational rank one
or two, respectively. Those are already the ranks
of the primitive fibre images just proved. Their
cokernels are free, so no further finite enlargement
is possible. The images are equal integrally.
The Wang projection of the full kernel is therefore
all of the fixed lattice.

\item For completeness, the scalar fixed-normal torus
bounds the scalar-circle bundle over the zero-section
disc. The other fixed-normal torus has zero projection
to $H_2(Y)$: after the same narrow cover and marked
straightening, it is the boundary of the zero-section
disc crossed with the long circle. Since $H_2(Y)$
is free, it was zero before the cover as well.
Its image in $H_2(U)$ is consequently torsion.
If $h=1$ both normal images vanish, and the full
fibre map is the primitive integral area functional
in \eqref{maps:eq:D.IC7}. This proves the complete marked assertion
in that subsector.
\end{enumerate}
\end{proof}

The exact implementation is
\path{code/III/manifolds_invariant_chern_caps.py}.
It constructs the full cone, retains the actual
$\delta\gamma$ and multiplicity cycles, checks their
orders and cokernels, and verifies the component-order
and logarithm-denominator markings against every
recorded additive site. Its finite checks corroborate
the all-integer elimination in \eqref{maps:eq:D.IC9}; they are not
bounds on the permitted Chern coefficients.


\begin{theorem}[all height-integral initial pairs with corrected original additive caps]\label{maps:src:D.2.52}
\StatementRecord{maps:src:D.2.52}{proved}{Historical events: Claude Fable 5.2 (Anthropic), effort high, 16 November 2026: carried as background without fault found. GPT-6 Astra (OpenAI), ChatGPT Pro mode, effort Pro, S6\_09\_PAPERS, 17 November 2026: targeted proof check. Also checked in this writing. Pre-continuation writing events: GPT-6 Astra (OpenAI), ChatGPT Pro mode, effort Pro, S6\_10\_PAPERS, 18 November 2026: targeted proof check. Continuation metadata events: GPT-6 Astra (OpenAI), ChatGPT Pro mode, effort Pro, S6\_10\_PAPERS continuation, 19 November 2026: used and rechecked in this run; status-vocabulary reconciliation; metadata reconciliation only; no new mathematical check. \emph{Pre-continuation recorded proof scope.} GPT-6 Astra (OpenAI, ChatGPT Pro mode), effort Pro, 2026-11-17 UTC; Claude Fable 5.2 (Anthropic), effort high, 16 November 2026; also checked in this writing by GPT-6 Astra (OpenAI), ChatGPT Pro mode, effort Pro, 18 November 2026}{the hypotheses, notation and definitions stated here; \cref{maps:src:D.2.50}; \cref{maps:eq:D.IC2}; \cref{maps:src:D.2.46}; \cref{maps:eq:D.IC8}; \cref{maps:eq:D.VC3}; \cref{maps:eq:D.IC7}; \cref{maps:eq:D.GI2}; \cref{maps:src:D.2.41}; \cref{maps:src:D.2.43}; \cref{maps:src:D.2.51}; \cref{maps:src:D.2.31}; \cref{maps:src:D.2.33}; \cref{maps:eq:D.GI6}; \cref{maps:eq:D.CI1}}


For any of the twenty-one recorded Mordell--Weil rank-one
profiles, take $P=aG,Q=bG$ with $a,b>0$ and
$d_0\mid ab$. Choose any nonnegative whole-fibre
linearization weights of total $ab/d_0$, supported
at multiplicative or additional regular fibres.
Use the actual corrected line bundle of \ref{maps:src:D.2.50}.
The following filling choices have an exact all-order
fundamental-group presentation and complete integral
homology algorithm:
\begin{enumerate}[label=(\roman*)]
\item every marked same-fan rank-one or triangular
Mumford cap at a multiplicative or regular site;
\item every free finite-good additive logarithmic cap,
with all its affine powers and all five integral Thom
columns retained;
\item every original additive cap whose corrected
multidegree has $h=1$ in \eqref{maps:eq:D.IC2}, and every unchanged
zero-degree original additive cap;
\item every original multiplicative cap with its
uncorrected degree vector, as in \ref{maps:src:D.2.46};
\item arbitrarily many regular logarithmic caps and
any ordinary integral clutch.
\end{enumerate}
The odd-star fibres are filled originally in this
algorithm. In particular, the recorded Smith
corrections from \ref{maps:src:D.2.50} admit every positive pair
with integral height: all newly non-narrow additive
positions are the component pair $(2,2)$ at an odd
star, and their corrected degrees are \eqref{maps:eq:D.IC8}.
Newly non-narrow multiplicative positions may all
use the marked Mumford filling.
Thus this theorem removes the narrowness-OR restriction
on the initial sections in \ref{maps:src:D.2.46}.
It does not identify original multiplicative caps
with their nonzero vertical corrections, nor the
unresolved marked $h>1$ additive caps.

The exact implementation is
\path{code/III/manifolds_corrected_initial.py}.
It accepts the same unbounded data as \ref{maps:src:D.2.46} and
optional literal component-divisor vectors. A supplied
vector is checked against the integral equation
\eqref{maps:eq:D.VC3}; no rational correction is substituted.
For an original additive site it computes the
orbit lengths and $h$, and applies the complete
marked kernel \eqref{maps:eq:D.IC7} only when $h=1$.
A nonzero whole-fibre change at a zero-degree
original cap is also retained as outside this
implementation until its relative meridian marking
is supplied.

\end{theorem}


\begin{proof}
\begin{enumerate}[label=\arabic*.,leftmargin=*]
\item \ref{maps:src:D.2.50} constructs the actual initial line bundle
and proves that its punctured monodromy is exactly
\eqref{maps:eq:D.GI2}, without the narrowness-OR assumption.
All entries are integral and their ordered product
is the identity. The common punctured torus bundle
therefore has exactly the same integral fibre-and-graph
bar model as in \ref{maps:src:D.2.41} and \ref{maps:src:D.2.46}, including its
quadratic comparison and literal last-boundary word.

\item Every selected Mumford cap depends on the actual
marked unipotent period matrix and the selected fan.
The rank-one calculation and \ref{maps:src:D.2.43} therefore apply
to these newly corrected inputs verbatim: the full
nonsaturated angular-area kernel, primitive normal
discs and all component meridians remain as stated.
No extension of the vertical divisor over a different
original cap is being inferred from this fact.
The finite-good and regular logarithmic caps use
the actual affine quotient action, so their freeness
criterion, norm columns and fifth Thom generator
are unchanged as well.

\item At an original additive site with $h=1$, \ref{maps:src:D.2.51}
supplies every required integral marking. Its
fundamental group is generated by the primitive
long circle; all three vertical directions are
killed. The fibre degree-two kernel is the actual
primitive area kernel \eqref{maps:eq:D.IC7}, and both literal
fixed-normal tori bound. Its component contribution
to $H_4$ is the number $\rho$ of component orbits,
not the number of components before the permutation.
The unchanged zero-degree case and the uncorrected
original multiplicative case are already proved
in \ref{maps:src:D.2.31}, \ref{maps:src:D.2.33} and \ref{maps:src:D.2.46}.

\item The implementation therefore adds these exact
local columns to the common bar presentation.
Its fundamental-group presentation retains the
commuting fibre generators, the conjugation matrices,
the actual ordered peripheral product, each finite
power relation, and each nonfinite killed vector.
The normal closure of the vertical vectors is
computed by repeatedly adjoining their conjugates
and taking integral Hermite normal form until the
lattice stabilizes. The ascending-chain condition
on submodules of $\mathbb Z^3$ proves termination.
No saturation or assumed abelian quotient replaces
this presentation.

\item For degree-two homology, form the full integer
relation matrix $M$ and retain its Smith cokernel.
If $E$ denotes the common piece after the finite
caps have been filled, the remaining relative
free rank is exactly
\[
           \sum_i b_4(U_i)-b_1(E)+b_1(X).
                                                        \tag{D.CI1}\label{maps:eq:D.CI1}
\]
The proof is the same integral pair sequence as
\eqref{maps:eq:D.GI6}: each $H_2(U_i,\partial U_i)=H^4(U_i)$
is free, and its actual boundary columns are
already included. The cap torsion permitted by
\ref{maps:src:D.2.43} remains in its nonsaturated fibre kernel;
the new original additive caps have no degree-two
torsion because $h=1$. The quotient in \eqref{maps:eq:D.CI1}
is free, so it splits as an abelian-group extension.
This gives $H_2(X)$ with all torsion.
The Euler number is twice the sum of the indices
of the positive rank-two caps. The exact $H_1$,
this Euler number, integral duality and universal
coefficients give every remaining group, including
the Euler-zero cases.

\item Finally, the canonical correction list in \ref{maps:src:D.2.50}
has $45$ additional residue pairs over the full
profile list. At additive sites the component
orders are prime except for the two odd stars.
The criterion $\operatorname{ord}(\Phi)\mid pq$
therefore forces $p=0$ or $q=0$ except for
$p=q=2$ in $\mathbb Z/4$.
The displayed canonical correction at those
stars has a fixed inner coefficient $-1$,
so $h=1$. This proves that the algorithm applies
to every positive height-integral initial pair
with these retained corrections and the stated
filling choices. The conclusion is unbounded in
section multiples, weights, logarithmic orders
and every permitted affine twist.
\end{enumerate}
\end{proof}

The independent replay script
\path{code/III/manifolds_corrected_initial_audit.py}
compares the full matrices against the earlier
\ref{maps:src:D.2.46} implementation and the sixteen independently
computed \ref{maps:src:D.2.50} examples. It also runs every one
of the $45$ new residue pairs, with all listed
multiplicative zero positions, both available
additive filling patterns, and additional
higher-weight and regular-log inputs.
Those finite executions check the implementation;
the actual local models and integral pair argument
above prove the stated all-parameter scope.


\begin{theorem}[complete primitive cubic and fundamental-class invariant of the weighted OS45 example]\label{maps:src:D.2.53}
\StatementRecord{maps:src:D.2.53}{proved}{Historical events: Claude Fable 5.2 (Anthropic), effort high, 16 November 2026: carried as background without fault found. GPT-6 Astra (OpenAI), ChatGPT Pro mode, effort Pro, S6\_09\_PAPERS, 17 November 2026: targeted proof check. Also checked in this writing. Pre-continuation writing events: GPT-6 Astra (OpenAI), ChatGPT Pro mode, effort Pro, S6\_10\_PAPERS, 18 November 2026: targeted proof check. Continuation metadata events: GPT-6 Astra (OpenAI), ChatGPT Pro mode, effort Pro, S6\_10\_PAPERS continuation, 19 November 2026: used and rechecked in this run; status-vocabulary reconciliation; metadata reconciliation only; no new mathematical check. \emph{Pre-continuation recorded proof scope.} GPT-6 Astra (OpenAI, ChatGPT Pro mode), effort Pro, 2026-11-17 UTC; Claude Fable 5.2 (Anthropic), effort high, 16 November 2026; also checked in this writing by GPT-6 Astra (OpenAI), ChatGPT Pro mode, effort Pro, 18 November 2026}{the hypotheses, notation and definitions stated here; \cref{maps:eq:D.PC11}; \cref{maps:eq:D.PC12}; \cref{maps:eq:D.OF1}; \cref{maps:eq:D.OF14}; \cref{maps:eq:D.OF17}; \cref{maps:eq:D.OF18}}


Take the marked OS45/default rational elliptic surface, the initial
sections \(P=2G,Q=8G\), whole-fibre weights
\((2,0,0,0,0)\) supported at \(I_8\), all multiplicative Mumford
caps, and the ordinary clutch \(e=1\). The resulting simply
connected manifold \(X\) has
\[
 H_2(X;\mathbb Z)=\mathbb Z^{15}\oplus\mathbb Z/8,\qquad
 H_3(X;\mathbb Z)=\mathbb Z/8,\qquad
 c_1(X)=c_2(X)=\lambda(X)=0.                              \tag{D.OS1}\label{maps:eq:D.OS1}
\]
Its sixteen normalized cusp components are indexed by
\((r,s)\in\mathbb Z/8\times\mathbb Z/2\), with \(i=2r+s\).
Their divisor classes generate an index-two sublattice in
\(H^2(X;\mathbb Z)\). The full primitive lattice is generated
by \(x_0,\ldots,x_{13},y\), with
\[
             2y=\sum_{r=0}^7x_{r,1},\qquad
             \sum_{r,s}x_{r,s}=0.                        \tag{D.OS2}\label{maps:eq:D.OS2}
\]
The complete integral cubic is \eqref{maps:eq:D.PC11}--\eqref{maps:eq:D.PC12} below.
There is a primitive cyclic coordinate \(u\in H^2(X;\mathbb Z/8)\)
with the full boundary marking \eqref{maps:eq:D.OF1}--\eqref{maps:eq:D.OF14}. In the splitting
fixed by \(u\) and this primitive basis,
\[
 H_6(K(H_2(X),2);\mathbb Z)
   =\mathbb Z^{680}\oplus(\mathbb Z/8)^{120}
                       \oplus(\mathbb Z/16)^{15}
                       \oplus\mathbb Z/8,                \tag{D.OS3}\label{maps:eq:D.OS3}
\]
and the complete fundamental-class invariant is specified by
the integral cubic and \eqref{maps:eq:D.OF17}--\eqref{maps:eq:D.OF18}.
In particular
\[
 \begin{gathered}
 y^2=0,\qquad
 \langle\mathcal P_2(u)x_i,[X]\rangle=7\pmod{16}
                   \ (0\le i<14),\\
 \langle\mathcal P_2(u)y,[X]\rangle=0,\qquad
 \langle u^3,[X]\rangle=0.
 \end{gathered}                                           \tag{D.OS4}\label{maps:eq:D.OS4}
\]
This theorem does not assert a completed smooth classification:
the refined characteristic class is still one of \(0,4\)
in the cyclic order-eight summand of \(H^4(X;\mathbb Z)\).
The full finite-coefficient invariant here does not choose
between those two possibilities.

\end{theorem}
\begin{proof}
\begin{enumerate}[label=\arabic*.,leftmargin=*]
\item \textbf{The primitive component lattice and cubic in the sixteen-component OS45 example}
We first record a saturation calculation which uses the actual component
meridians. Let $X$ be simply connected, let $N$ be a single rank-two
Mumford cusp with $D$ components, and put $E=X\setminus\operatorname{int}N$.
Assume $b_2(X)=D-1$, $H_3(N;\mathbb Z)$ is free, and
\[
 H_1(E;\mathbb Z)=\mathbb Z\oplus A,\qquad
                       \mu_i=(1,a_i),                      \tag{D.PC1}\label{maps:eq:D.PC1}
\]
where $A$ is finite and $\mu_i$ is the positively oriented meridian of
component $i$. The component classes $x_i$ then generate a sublattice
\[
 L=\mathbb Z^D/\mathbb Z(1,\ldots,1)
                       \subset H^2(X;\mathbb Z)
\]
whose full primitive extension, inside $L\otimes\mathbb Q$, is
\[
 H^2(X;\mathbb Z)
   =L+\left\{\sum_i\widetilde{\chi(a_i)}x_i:
                         \chi\in\operatorname{Hom}(A,\mathbb Q/\mathbb Z)
        \right\}.                                        \tag{D.PC2}\label{maps:eq:D.PC2}
\]
A tilde means any rational lift of a value in $\mathbb Q/\mathbb Z$;
changing it adds an element of $L$. In particular the component
lattice has index $|A|$.

\item Poincare--Lefschetz duality identifies
$H^2(X,E;\mathbb Z)=H_4(N;\mathbb Z)=\mathbb Z^D$,
with its positive component Thom basis, and
$H^3(X,E;\mathbb Z)=H_3(N;\mathbb Z)$ is free. The connecting
map from $H^1(E;\mathbb Z)$ evaluates on the meridians \eqref{maps:eq:D.PC1},
so its image is precisely $\mathbb Z(1,\ldots,1)$.
The relative exact sequence therefore injects $L$ into $H^2(X)$.
Since $X$ is simply connected, the latter group is free; its rank is
$D-1$ by hypothesis. Its quotient by $L$ is finite and embeds in
$H^2(E)$, whose torsion is
$\operatorname{Ext}(H_1(E),\mathbb Z)=\operatorname{Hom}(A,\mathbb Q/\mathbb Z)$.
All this torsion lifts to $H^2(X)$ because the next relative group is
free. Thus the quotient has exactly that order.

\item Here is also its precise embedding, rather than only its index.
For a character $\chi$ of order $m$, extend it by zero on the free
summand of $H_1(E)$. Regard it as a class in
$H^1(E;\mathbb Z/m)$ by multiplying its values by $m$.
Its relative coboundary has component coordinates
$m\widetilde{\chi(a_i)}\pmod m$. Choose the displayed integers
as an integral relative lift. Its image in $H^2(X;\mathbb Z/m)$
is zero by exactness. Hence its integral image is divisible by $m$
in $H^2(X;\mathbb Z)$, giving the class in \eqref{maps:eq:D.PC2}.
The coefficient exact sequences identify its exterior restriction
with the corresponding torsion Bockstein, up to the harmless common
connecting-map sign. Since these characters give the full torsion
of $H^2(E)$, they give the full primitive extension. This proves
\eqref{maps:eq:D.PC2}, including every finite denominator.

\item Now take the exact marked data
\[
 \mathrm{OS45/default},\quad P=2G,\quad Q=8G,\quad
 (w_0,w_1,w_2,w_3,w_4)=(2,0,0,0,0),\quad e=1,              \tag{D.PC3}\label{maps:eq:D.PC3}
\]
with all multiplicative fibres filled by their marked Mumford caps.
\ref{maps:src:D.2.46} gives
\[
 H_2(X)=\mathbb Z^{15}\oplus\mathbb Z/8,\qquad
 H_3(X)=\mathbb Z/8,\qquad \pi_1(X)=1.
                                                               \tag{D.PC4}\label{maps:eq:D.PC4}
\]
The elliptic monodromy at $I_8$ is
$\left(\begin{smallmatrix}1&8\\0&1\end{smallmatrix}\right)$
and its section cocycle is $(3,0)$. Formula \eqref{maps:eq:D.GI2} therefore gives
\[
 A_0=\begin{pmatrix}
 1&8&0&24\\0&1&0&0\\0&-6&1&-16\\0&0&0&1
 \end{pmatrix},\qquad
 L_{\rm per}=\left\langle(8,-6),(24,-16)\right\rangle
                 =8\mathbb Z\oplus2\mathbb Z.             \tag{D.PC5}\label{maps:eq:D.PC5}
\]
The short basis is $(e_1,\delta)$ and the long basis is $(e_2,\gamma)$.
This is an equality of sublattices, not merely a comparison of
indices. The determinant is sixteen and the Smith factors are two
and eight. Label the normalized components by
\[
             (r,s)\in\mathbb Z/8\times\mathbb Z/2,
             \qquad i=2r+s.                              \tag{D.PC6}\label{maps:eq:D.PC6}
\]

\item The actual cusp complement satisfies
\[
 H_1(E)=\mathbb Z\oplus\mathbb Z/2,
                      \qquad \mu_{r,s}=(1,s\bmod2).      \tag{D.PC7}\label{maps:eq:D.PC7}
\]
To verify the marking, retain all peripheral generators and the
ordinary clutch in the exact degree-one presentation. Its Smith
form has seven unit factors, one factor two, and one free generator.
The torsion character in the original nine coordinates is
\[
                  (0,-2,1,2,0,0,2,-2,2)\pmod2,
                                                               \tag{D.PC8}\label{maps:eq:D.PC8}
\]
which is precisely $\delta^*$ modulo two. The free character
evaluates to one on the cusp meridian. A height-one ray at
$(r,s)$ has literal meridian $\theta+r e_1+s\delta$; it
therefore gives exactly \eqref{maps:eq:D.PC7}. The full original presentation,
Smith change of basis and all sixteen meridian images are saved
in the executable below.

\item Applying \eqref{maps:eq:D.PC2}, there is a unique integral class $y$ with
\[
                    2y=\sum_{r=0}^{7}x_{r,1}.             \tag{D.PC9}\label{maps:eq:D.PC9}
\]
The component lattice has exact index two. In particular it would
be incorrect to declare any fifteen of the component classes a
primitive basis. A convenient actual basis is
\[
 x_0,x_1,\ldots,x_{13},y,\qquad
 x_{14}=-2y-\sum_{\substack{i<14\\i\ \mathrm{even}}}x_i,
 \qquad
 x_{15}=2y-\sum_{\substack{i<14\\i\ \mathrm{odd}}}x_i.
                                                               \tag{D.PC10}\label{maps:eq:D.PC10}
\]
The integral matrix of the sixteen component classes in this basis
has Smith factors $1^{14},2$ and kernel generated by the all-one
vector. These are the asserted primitive markings.

\item Here is the complete cubic on that basis. For component classes,
\ref{maps:src:D.2.47} specializes to the following symmetric integral trilinear
form. With distinct vertices as appropriate,
\[
 \begin{aligned}
 T(x_i,x_i,x_i)&=6,\\
 T(x_i,x_i,x_j)&=
 \begin{cases}
 -2,&i=(r,s),\ j=(r,1-s),\\
 -1,&i=(r,s),\ j=(r\pm1,t),\quad t=0,1,\\
 0,&\text{otherwise},
 \end{cases}\\
 T(x_i,x_j,x_k)&=
 \begin{cases}
 1,&\{i,j,k\}\text{ is three of the four vertices in}\\[-2pt]
   &\hspace{12pt}\{(r,0),(r,1),(r+1,0),(r+1,1)\}
                   \text{ for some }r,\\
 0,&\text{otherwise}.
 \end{cases}
 \end{aligned}                                             \tag{D.PC11}\label{maps:eq:D.PC11}
\]
All horizontal indices are cyclic modulo eight. Repeated geometric
edges and triangle cells are retained in this formula.
The remaining primitive entries are
\[
 \begin{aligned}
 T(y,x_{r,s},x_{r',s'})
 &=(-1)^s\left(
       \mathbf1_{s'=s,\ r'=r+1}
      +\mathbf1_{s'=s,\ r'=r-1}
      -2\mathbf1_{(r',s')=(r,s)}\right),\\
 T(y,y,x_i)&=0,\qquad T(y,y,y)=0.
 \end{aligned}                                             \tag{D.PC12}\label{maps:eq:D.PC12}
\]
These formulas determine every entry on the primitive basis
\eqref{maps:eq:D.PC10}, and every resulting entry is integral.

For a geometric verification of the half-class entries, on a
normalized $dP_6$ take the cyclic boundary classes
\[
 E_1,\ H-E_1-E_2,\ E_2,\ H-E_2-E_3,\ E_3,\ H-E_1-E_3
                                                               \tag{D.PC13}\label{maps:eq:D.PC13}
\]
in the ray order of \ref{maps:src:D.2.47}. The four nonhorizontal boundary classes
sum to $2f$, where $f=H-E_1$ and $f^2=0$. Restricting \eqref{maps:eq:D.PC9}
therefore gives
\[
                         y|_{P_{r,s}}=(-1)^s f.            \tag{D.PC14}\label{maps:eq:D.PC14}
\]
This division is valid in the torsion-free group $H^2(dP_6;\mathbb Z)$.
The fibre $f$ has degree one on each horizontal boundary curve
and zero on the other four. The Cartier restriction formula in
\ref{maps:src:D.2.47} now gives the first line of \eqref{maps:eq:D.PC12}. Squaring \eqref{maps:eq:D.PC14}
gives $y^2x_i=0$; multiplying \eqref{maps:eq:D.PC9} by $y^2$ gives $y^3=0$.
Thus no rational cubic has been silently substituted for the
integral form.

Finally $c_1(X)=0$ and both $c_2$ and $p_1$ pair to zero with
every class in the primitive free basis. For $c_1$, the normalized
component tangent calculation is zero and the component restriction
is rationally injective by the connected toric graph Laplacian;
$H^2(X)$ is torsion-free, so this proves integral vanishing.
For $c_2,p_1$, \ref{maps:src:D.2.47} gives zero on every component. The component
classes have finite index, and the evaluations are integers, so they
also vanish on $y$. These statements concern the free characteristic
pairings; they do not determine the torsion class of $c_2$, the
refined class, or the full finite-coefficient fundamental-class
invariant.

\path{code/III/fourfolds_OS45_primitive_cubic.py}
reconstructs the actual complement before using its torsion
character, checks all sixteen meridians, computes the exact
index-two primitive lattice, and verifies all $15^3=3375$ primitive
triples and all $16^3=4096$ recovered component triples. The
normalization restriction \eqref{maps:eq:D.PC14}, the actual period basis, and
all sparse integral entries are retained in its JSON output.
The preliminary component-only calculation is retained separately
in \path{code/III/fourfolds_OS45_component_ring.py}.

\paragraph{Integral vanishing of the critical class.}
The primitive-lattice argument above gives the literal relation
\(\sum x_{r,s}=0\). Therefore the target line of \(df\) is
topologically trivial. All fibres in this example are reduced
semistable fibres. The exact complex of \ref{maps:src:D.2.44} consequently
applies without any additive correction: its abstract kernel
has the nowhere-zero scalar section, and its quotient line is
topologically trivial because \(c_1(X)=0\). The local unit Bott
clutch gives
\[
                     c_2(X)=\operatorname{PD}(C_{\rm crit}).
                                                               \tag{D.OS5}\label{maps:eq:D.OS5}
\]
Here every actual complex double curve occurs once.

We retain the actual integral specialization of this complete cycle.
In the ordered bivector basis
\((e_1e_2,e_1\delta,e_1\gamma,e_2\delta,e_2\gamma,\delta\gamma)\),
the five fibre preimages, in their actual site order, are
\[
 \begin{array}{c|rrrrrr}
 I_8\text{ weighted cusp}&-16&0&6&8&0&2\\
 I_1&0&0&0&0&0&1\\
 I_1&0&8&0&0&0&1\\
 I_1&0&0&0&8&0&1\\
 I_1&0&0&0&0&0&1
 \end{array}.                                             \tag{D.OS6}\label{maps:eq:D.OS6}
\]
For the first row, the actual local period cover has matrix
\(L=\left(\begin{smallmatrix}8&24\\-6&-16\end{smallmatrix}\right)\)
and degree sixteen. Transfer the minimal cusp's literal
three-curve preimage
\(a_1b_1+b_1a_2+a_2b_2\), expressing \(b_1,b_2\) by \(L^{-1}\)
in the actual long basis, and multiply by sixteen. This gives
the first integral row of \eqref{maps:eq:D.OS6}, without division in homology
or an assumed covering of the global filled manifolds.
Each other row is the determinant-one transverse torus of the
corresponding primitive node, as in \ref{maps:src:D.2.44}.
Their sum is
\[
                        z=(-16,8,6,16,0,6).              \tag{D.OS7}\label{maps:eq:D.OS7}
\]
In the actual twenty-three-row relation matrix \(M\) of
\ref{maps:src:D.2.46}, augment this vector by seventeen zeros. With columns
numbered from zero, a literal integer preimage is the column
combination with nonzero coefficients
\[
 \begin{array}{c|rrrrrrrr}
 j&6&8&10&15&16&18&19&20\\\hline
 c_j&-6&2&2&-28&2&-2&-52&-8\\
 \end{array}
 \qquad
 \begin{array}{c|rrrrrrr}
 j&26&38&41&51&56&61&66\\\hline
 c_j&2&10&10&10&10&10&10 .
 \end{array}                                             \tag{D.OS8}\label{maps:eq:D.OS8}
\]
Direct multiplication gives \(Mc=(z,0)\).
All columns of \(M\) are the actual core prism and cap Thom
relations, so this proves that the literal critical cycle is
integrally nullhomologous. Equations \eqref{maps:eq:D.OS5}--\eqref{maps:eq:D.OS8} prove
the absolute Chern and spin-characteristic vanishing in
\eqref{maps:eq:D.OS1}. They do not by themselves compute linking with an
odd middle-dimensional generator: local specialization chains
must be retained for that separate refined calculation.

\textbf{The actual torsion coordinate and complete fundamental-class
invariant of the sixteen-component OS45 manifold}
We use exactly the datum \eqref{maps:eq:D.PC3}, its period marking \eqref{maps:eq:D.PC5}, and the
primitive free basis \eqref{maps:eq:D.PC10}. Put \(k=\mathbb Z/8\).
There is a primitive torsion coordinate
\(u\in H^2(X;k)\) whose complete cusp support is
\[
 \Psi(r,s)=4\left\lfloor r/8\right\rfloor
       +3(r\bmod8)\left\lfloor s/2\right\rfloor
                                      \quad\hbox{in }k.       \tag{D.OF1}\label{maps:eq:D.OF1}
\]
Here and below \(r\bmod8\) means the representative between zero and
seven. The assertion includes the entire boundary cohomology class,
including its coefficient-extension part. We prove this before using
any support-function evaluations.

\paragraph{The specified core cochain.}
Use the actual degree-two core model of \ref{maps:src:D.2.46}. Its fibre cells are
ordered
\[
 e_1e_2,\ e_1\delta,\ e_1\gamma,\ e_2\delta,\
                                      e_2\gamma,\ \delta\gamma.
\]
For the six boundary monodromies, including the ordinary identity cap,
let
\[
 D=(A_0-I,\ldots,A_4-I),\qquad K=\ker_{\mathbb Z}D.
\]
Choose the integral cycle basis and the integral left inverse \(K^+\)
retained by the exact program. The ordinary degree-two cycles are the
six fibre cells followed by this graph-cycle basis. The actual
relation matrix \(M\) has twenty-three rows. If
\(UMV\) is its stored Smith form, the unique nonunit factor is eight.
Let \(\varphi\) be the negative of the corresponding row of \(U\).
Then
\[
 \varphi M=0\pmod8,\qquad
 \varphi|_{\Lambda^2\Lambda}=(1,0,0,0,0,0).                \tag{D.OF2}\label{maps:eq:D.OF2}
\]
Its value on the stored Smith torsion generator is \(-1\).
This is already a primitive functional on the order-eight image of
the core in \(H_2(X)\).

We specify the cochain, rather than merely its homology functional.
The inclusion \(K\subset\mathbb Z^{20}\) is primitive because it is
the kernel of an integral map. Extend the graph row by
\(\varphi_{\rm graph}K^+\), and use the fibre row in \eqref{maps:eq:D.OF2}.
This gives a literal degree-two cochain on the core mapping-cylinder
model. The degree-three boundary columns are the actual exterior
monodromy relations, so \eqref{maps:eq:D.OF2} says that it is closed modulo eight.
The fixed bar-prism comparisons in \ref{maps:src:D.2.46} identify this cochain with
one on the actual torus bundle. In particular its restriction is not
left unspecified by the possible \(2\)-torsion in \(H_1\) of the core.

At the rank-two cusp the actual peripheral inclusion has graph
matrix \(J=K^+J_{\rm raw}\). The retained cochain gives
\[
 \varphi_{\rm graph}J=(0,52,0,148)
       \quad\hbox{on }(e_1,e_2,\delta,\gamma).             \tag{D.OF3}\label{maps:eq:D.OF3}
\]
The quadratic correction on each of these four basis loops is zero
at this cusp. Thus the two long-strip constants are
\[
                      (c_1,c_2)=(4,4)\quad\hbox{in }k.   \tag{D.OF4}\label{maps:eq:D.OF4}
\]
The full Smith row, the chosen left inverse, and the actual
peripheral matrices are reproducible by
\path{code/III/maps_weighted_os45.py}. In particular \eqref{maps:eq:D.OF3} is not
obtained by restricting only the regular-fibre alternating form.

\paragraph{A full boundary comparison for finite-coefficient supports.}
We give the comparison explicitly, since this cusp has torsion in
its boundary \(H_1\). Let \(\widetilde N\) be the common universal
cover of the triangular Mumford cusp and let
\(\widetilde M\) be its boundary. Its normalized components are
indexed by \(v\in\mathbb Z^2\). The boundary cover retracts to the
three-torus generated by the two short angular loops and the positive
\(q\)-meridian \(\theta\). \ref{maps:src:D.2.43} gives
\[
 H_1(\widetilde N;\mathbb Z)=0,\quad
 H_2(\widetilde N;\mathbb Z)=I^2,\quad
 \operatorname{im}\{H_2(\widetilde M;\mathbb Z)\to
                         H_2(\widetilde N;\mathbb Z)\}=0. \tag{D.OF5}\label{maps:eq:D.OF5}
\]
Indeed the three boundary generators are the
angular area and the two angular-normal tori, all with the actual
bounding sweeps described there.

The relative component Thom cochains identify
\[
 H^2(\widetilde N,\widetilde M;k)
                  =\operatorname{Map}(\mathbb Z^2,k).
\]
One can obtain this directly from locally finite
component chains: a degree-four locally finite cycle is an arbitrary
family of fundamental cycles of the normalized components, and
degree-four intersection strata do not occur. Equivalently,
Poincare--Lefschetz duality uses Borel--Moore degree-four homology.
No compact-support convention replaces the displayed product of
coefficient groups by a direct sum.

The positive meridian of component \(v\) is
\(\theta+v_1e_1+v_2\delta\). Therefore the relative connecting map
from \(H^1(\widetilde M;k)\) sends a covector \((c,m_1,m_2)\) to
the affine function \(c+m_1v_1+m_2v_2\). By \eqref{maps:eq:D.OF5}, the relative
sequence is exactly
\[
 0\longrightarrow\operatorname{Aff}(\mathbb Z^2,k)
 \longrightarrow\operatorname{Map}(\mathbb Z^2,k)
 \longrightarrow H^2(\widetilde N;k)\longrightarrow0.     \tag{D.OF6}\label{maps:eq:D.OF6}
\]
Here \(\operatorname{Aff}\) consists of these three-parameter
affine functions. This also identifies the quotient with
\(\operatorname{Hom}(I^2,k)\), consistently with the actual
normalization computation of \ref{maps:src:D.2.43}.

Let the actual period lattice \(L\) act by translation on the
argument \(v\). If a function \(\Psi\) satisfies
\[
                  \Psi(v+l)-\Psi(v)=a_l(v)
                       \in\operatorname{Aff},           \tag{D.OF7}\label{maps:eq:D.OF7}
\]
its absolute class on \(\widetilde N\) is \(L\)-invariant.
It descends to \(N_L\): in total degree two the covering
cochain double complex has \(H^1(\widetilde N;k)=0\), and
the only possible obstruction lies in \(H^3(L;k)=0\).
The choices of descent form a torsor for \(H^2(L;k)=k\).

There is a useful precise version of the boundary part of this
argument. The boundary group is the split extension
\[
       \pi_1(M_L)=\mathbb Z^3\rtimes L,                    \tag{D.OF8}\label{maps:eq:D.OF8}
\]
where the normal \(\mathbb Z^3\) is generated by the two shorts
and \(\theta\), and the two long loops give the specified splitting.
Its dual module \(H^1(\widetilde M;k)\) is exactly
\(\operatorname{Aff}\): evaluating a covector on
\(\theta+v\) produces its affine function, and translating \(v\)
gives the deck action.
For boundary classes with zero restriction to \(\widetilde M\),
the degree-two filtration is the short exact sequence
\[
 0\longrightarrow H^2(L;k)\longrightarrow
 \ker\{H^2(M_L;k)\to H^2(\widetilde M;k)\}
 \longrightarrow H^1(L;\operatorname{Aff})\longrightarrow0.
                                                               \tag{D.OF9}\label{maps:eq:D.OF9}
\]
The splitting of \eqref{maps:eq:D.OF8} kills the degree-two transgression;
there is no degree-three group cohomology of \(L\).
These two observations also prove \eqref{maps:eq:D.OF9} directly from the
two-dimensional free resolution of \(L=\mathbb Z^2\).
Restriction to the actual long two-torus splits its leftmost term.

The image in the rightmost group of the boundary value of a
descended support is represented by the full affine cocycle
\(l\mapsto a_l\) in \eqref{maps:eq:D.OF7}. This follows either by chasing
\eqref{maps:eq:D.OF6} in that resolution, or by choosing the relative Thom
cochain \(\Psi\): its deck difference is the boundary one-cocycle
whose value on \(\theta+v\) is precisely \(a_l(v)\).
Thus the gradients, constants, and one long-torus value determine
the entire boundary class. This statement retains every
\(\operatorname{Ext}(H_1(M_L),k)\) ambiguity.

In our positive fan convention a core cochain with short-long
matrix \(E\), zero angular area and zero short-normal strip values
has affine cocycle
\[
 a_{l_j}(v)=-E(v,b_j)+c_j,\qquad j=1,2,                  \tag{D.OF10}\label{maps:eq:D.OF10}
\]
where \(b_1=e_2,b_2=\gamma\) and \(l_j=(A-I)b_j\).
For a direct verification, translate the marked \(\theta,b_j\)
strip from component zero to component \(v\). Its change is the
oriented short-long torus with value \(-E(v,b_j)\).
The initial strip contributes \(c_j\). There is no additional
angular-area term, since that coefficient is zero. The literal
basis-loop prism correction is the one already included in
\eqref{maps:eq:D.OF3}. This gives \eqref{maps:eq:D.OF10} at the cochain level, including
its constant term.

For \eqref{maps:eq:D.OF2}--\eqref{maps:eq:D.OF4}, \(E=\left(\begin{smallmatrix}1&0\\0&0
\end{smallmatrix}\right)\). The required affine functions are
therefore
\[
             a_{(8,-6)}=-r+4,\qquad a_{(24,-16)}=4.
                                                               \tag{D.OF11}\label{maps:eq:D.OF11}
\]
The support \eqref{maps:eq:D.OF1} has exactly these differences.
This is an identity for all integers, not a bounded numerical
test: write \(r=8a+r_0,s=2b+s_0\) with \(0\le r_0<8\) and
\(0\le s_0<2\), and substitute. The two differences are
\(4-9r_0\) and \(12-24r_0\), respectively, which reduce to
\eqref{maps:eq:D.OF11}.
Choose the remaining descent parameter so its value on the
actual long torus \(\langle e_2,\gamma\rangle\) is zero,
as it is for \eqref{maps:eq:D.OF2}.
Equations \eqref{maps:eq:D.OF9}--\eqref{maps:eq:D.OF11} now prove equality on the full
cusp boundary, including both its finite-coefficient and
integral-torsion contributions.

For an independent check of the short-long marking, retain the
actual annulus map from \ref{maps:src:D.2.43}, not only its augmentation.
If \(x^py^q-1=(x-1)S_x+(y-1)S_y\), the angular vector \((a,b)\)
closes to the \(I^2\)-class with coefficients
\[
             (aS_x,\ bS_x+aS_y,\ bS_y)                   \tag{D.OF12}\label{maps:eq:D.OF12}
\]
on \(((x-1)^2,(x-1)(y-1),(y-1)^2)\).
Its literal Picard columns are
\[
 x^{-1}y^{-2}E_2+y^{-1}E_3,\quad
       -x^{-1}y^{-1}E_2,\quad
       x^{-1}E_1+x^{-2}y^{-1}E_2.                        \tag{D.OF13}\label{maps:eq:D.OF13}
\]
In the present positive component labels \(x\) translates by
\((-1,0)\), and \(y\) by \((0,1)\); the path and angular exponents
are changed simultaneously. Substituting \eqref{maps:eq:D.OF1} into all the
shifted columns gives the integer values
\[
                         \begin{pmatrix}9&24\\0&0\end{pmatrix}
                  \equiv E\pmod8.                       \tag{D.OF14}\label{maps:eq:D.OF14}
\]
The path dependence in \(S_x,S_y\) is exactly an
Hilbert--Burch relation, so this checks the actual annulus
marking rather than a rational replacement.

\paragraph{Extension through the other caps.}
Each of the four other multiplicative fibres has primitive
rank-one monodromy. Its boundary \(H_1\) is free; so is the
ordinary cap boundary \(H_1\).
The regular-fibre specialization is onto the degree-two group
of each such cap: in a primitive nodal marking this is the
specialization from \(T^2\times T^2\) to
\((S^2\vee S^1)\times T^2\); for the ordinary cap it is the
identity on the central \(T^4\).
Their actual degree-two boundary kernels
are precisely the columns used in \(M\), including the
ordinary cap's affine-meridian column.
Equation \eqref{maps:eq:D.OF2} therefore makes the specified core cochain
extend across each such cap. The free boundary \(H_1\)
ensures that agreement of the resulting integral-cycle
functional is agreement of the full finite-coefficient
cohomology class. These choices and the cusp comparison
glue to the asserted \(u\in H^2(X;k)\).
Its torsion value is the unit \(-1\), unchanged from \eqref{maps:eq:D.OF2};
we choose the opposite torsion generator to normalize it to one.
Choose free lifts of the homology basis dual to \eqref{maps:eq:D.PC10}
which are annihilated by \(u\). This fixes the splitting
\[
                         H_2(X)=\mathbb Z^{15}\oplus\mathbb Z/8
                                                               \tag{D.OF15}\label{maps:eq:D.OF15}
\]
used for the fundamental-class coordinates.

\paragraph{The complete coordinate values.}
Let \(R_0,\ldots,R_5\) be the six positive fan rays in cyclic
order and define the actual wall values
\[
 d_j(v)=\Psi(v)+\Psi(v+R_j)
                    -\Psi(v+R_{j-1})-\Psi(v+R_{j+1}).
\]
An integral representative for \(u|_{P_v}\) is
\[
 \eta_v=(d_0+d_1+d_2)H-d_0E_1-d_2E_2-d_4E_3.             \tag{D.OF16}\label{maps:eq:D.OF16}
\]
Opposite wall restrictions agree modulo eight. Changes of integral
representative by eight times a class preserve its square modulo
sixteen. The full exact component table is retained in
\path{records/maps_weighted_os45_mu.json}.

For component labels \(i,j\), all mixed ordinary coordinates are
specified without a long matrix display by
\[
 \begin{aligned}
 \langle ux_i^2,[X]\rangle
                    &=-\sum_{a=0}^{5}d_a(v_i)\pmod8,\\
 \langle ux_ix_j,[X]\rangle
                    &=\sum_{\{a:\ v_i+R_a=v_j\ {\rm mod}\ L\}}
                                       d_a(v_i)\pmod8\quad(i\ne j),\\
 \langle ux_i y,[X]\rangle
                    &=(-1)^{s_i}\eta_i\cdot(H-E_1)\pmod8.
 \end{aligned}                                             \tag{D.OF17}\label{maps:eq:D.OF17}
\]
The last value is four for component labels \(0,1,14,15\)
and zero for the other twelve. These formulas use actual
Cartier restrictions and retain the repeated vertical edges.
On the primitive free basis \(x_0,\ldots,x_{13},y\), the remaining
power coordinates are
\[
 \begin{split}
 \langle\mathcal P_2(u)x_i,[X]\rangle&=7\pmod{16}
                                             \quad(0\le i<14),\\
 \langle\mathcal P_2(u)y,[X]\rangle&=0\pmod{16},\\
 \langle u^3,[X]\rangle&=0\pmod8,\qquad
 \langle uy^2,[X]\rangle=0\pmod8.
 \end{split}                                               \tag{D.OF18}\label{maps:eq:D.OF18}
\]
The first line is \(\eta_i^2\) on the normalized component.
The two omitted component labels have square fifteen modulo
sixteen, consistently with their primitive-basis relations.
The remaining three evaluations, including the full
order-sixteen value which cannot be obtained by division,
are proved in the following localization input.
\textbf{The remaining finite-coefficient evaluations for OS45}
Let $u\in H^2(X;\mathbb Z/8)$ be the actual primitive torsion
coordinate constructed by the complete boundary comparison, with
cusp support
\[
 \Psi(r,s)=4\left\lfloor\frac r8\right\rfloor
       +3(r\bmod8)\left\lfloor\frac s2\right\rfloor
                           \quad\pmod8.                  \tag{D.OL1}\label{maps:eq:D.OL1}
\]
Here $r\bmod8$ is the representative from zero to seven. This
subsection uses that comparison on the whole cusp boundary;
matching only the regular-fibre bilinear form would not suffice.
The actual affine changes are
\[
 \Delta_{(8,-6)}\Psi=-r+4,\qquad
 \Delta_{(24,-16)}\Psi=4\quad\pmod8.                     \tag{D.OL2}\label{maps:eq:D.OL2}
\]
In particular both have zero slope in the scalar $s$ direction.
Let $y$ be the primitive class in \eqref{maps:eq:D.PC9}.
We prove the three full evaluations
\[
 \langle u^3,[X]\rangle=0\pmod8,\qquad
 \langle\mathcal P_2(u)y,[X]\rangle=0\pmod{16},\qquad
 \langle uy^2,[X]\rangle=0\pmod8.                         \tag{D.OL3}\label{maps:eq:D.OL3}
\]
The second formula includes the order-two part which cannot be
recovered by halving the sum of the component evaluations.

The scalar circle action is semifree. Its fixed locus consists of
the sixteen vertical toric curves $C_{r,s}$ joining vertices
$(r,s)$ and $(r,s+1)$, with $(r,s)$ indexed as in \eqref{maps:eq:D.PC6}.
Each has normal bundle $\mathcal O(-1)\oplus\mathcal O(-1)$
with weights $1,-1$. In every other triangular toric chart the
nonzero weights are also primitive. At each of the four zero-free
$I_1$ caps the primitive angular collapsing vector has first
coordinate one in its elliptic/scalar plane; the scalar circle
therefore projects primitively to the transverse torus and is
free. The ordinary torus cap is free as well. Thus there are no
unrecorded finite-stabilizer strata in this localization.

Since $X$ is simply connected, both $u$ and $y$ have circle-equivariant
lifts. Indeed the only outgoing differentials from ordinary degree
two in the Borel spectral sequence have targets involving
$H^1(X)$ or $H^3(BS^1)$, both zero. Let $x$ be the positive
generator of $H^2(C_{r,s};\mathbb Z)$ and let $t$ be the generator
of $H^2(BS^1;\mathbb Z)$. Write
\[
 \widehat u|_{C_{r,s}}=d_{r,s}x+\lambda_{r,s}t,
 \qquad
 \widehat y|_{C_{r,s}}=\mu_{r,s}t.                         \tag{D.OL4}\label{maps:eq:D.OL4}
\]
There is no $x$ term in the second expression: \eqref{maps:eq:D.PC14} and
$f\cdot C_{r,s}=0$ give $y(C_{r,s})=0$.

For clarity the ordinary degrees and the fibre characters are
computed separately. The degree of the support on a vertical wall is
\[
 d_{r,s}=\Psi(r,s)+\Psi(r,s+1)
              -\Psi(r+1,s)-\Psi(r-1,s+1)\pmod8 .         \tag{D.OL5}\label{maps:eq:D.OL5}
\]
Choosing representatives from zero to seven gives
\[
\begin{array}{c|cc}
 r&d_{r,0}&d_{r,1}\\\hline
 0,7&4&7\\
 1,\ldots,6&0&3.
\end{array}                                                \tag{D.OL6}\label{maps:eq:D.OL6}
\]
The scalar fibre character supplied by the support is its vertical
slope, up to a single common sign. Thus one may take
\[
                 \lambda_{r,0}=0,\qquad
                 \lambda_{r,1}=3r .                      \tag{D.OL7}\label{maps:eq:D.OL7}
\]
This is periodic under the actual period lattice because of
\eqref{maps:eq:D.OL2}. To justify the common character, rather than making a
separate choice on each normalized component, compare along the
actual toric orbit curves. Their degrees determine the differences
of the characters at their fixed endpoints. Horizontal and diagonal
orbit curves connect the sixteen fixed curves in a connected graph:
they join $(r,s)$ to $(r+1,s-1)$ and to $(r-1,s)$ respectively.
The support realizes all these differences. Consequently any global
equivariant lift has the characters in \eqref{maps:eq:D.OL7}, with a common sign
convention and one common additive character $L\in\mathbb Z/8$.
No global integral line lifting $u$ is assumed.

The integral component divisors have their canonical equivariant
linearizations. Along $C_{r,s}$ the two adjacent component normal
weights are $-1,+1$ in the lower/upper vertex order. Thus the weight
of the odd-component sum in \eqref{maps:eq:D.PC9} is $+1$ for $s=0$ and $-1$
for $s=1$. The difference
$2\widehat y-\sum_r\widehat x_{r,1}$ is a global multiple $ct$.
The existence of integral fixed-point weights forces $c$ odd;
changing the lift of $y$ changes $c$ by an even integer. Normalize
$c=-1$. Then
\[
                       \mu_{r,0}=0,\qquad
                       \mu_{r,1}=-1.                     \tag{D.OL8}\label{maps:eq:D.OL8}
\]
Every other lift adds one common integer $V$.

Each restriction of $\widehat u$ in \eqref{maps:eq:D.OL4} lifts integrally on
$C_{r,s}\times BS^1$. Thomas's cochain formula therefore gives
$\mathcal P_2(\widehat u)=(d x+\lambda t)^2$ with the full
modulus sixteen \cite[equations (2.1), (3.4), pp.267--268]{ffThomas56}.
The ordinary cube uses modulus eight. The equivariant normal Euler
class is
\[
                           (-x+t)(-x-t)=-t^2.
\]
Localization is valid over these coefficient rings because the
complement of the fixed locus is free and all normal weights are
units. Taking the coefficient of $x$ after division by $-t^2$ gives
\[
 \begin{aligned}
 \langle u^3,[X]\rangle&=-3\sum_{r,s}d_{r,s}\lambda_{r,s}^2
                                                        &&\pmod8,\\
 \langle\mathcal P_2(u)y,[X]\rangle
       &=-2\sum_{r,s}d_{r,s}\lambda_{r,s}\mu_{r,s}
                                                        &&\pmod{16},\\
 \langle uy^2,[X]\rangle&=-\sum_{r,s}d_{r,s}\mu_{r,s}^2
                                                        &&\pmod8.
 \end{aligned}                                             \tag{D.OL9}\label{maps:eq:D.OL9}
\]
The degree-zero base ring injects into its localization, so these
are identities before inverting $t$ as well.

The representatives \eqref{maps:eq:D.OL6}--\eqref{maps:eq:D.OL8} satisfy the exact integer sums
\[
 \sum d_{r,s}=40,\qquad \sum d_{r,s}\lambda_{r,s}=336,
 \qquad \sum d_{r,s}\lambda_{r,s}^2=5544,
 \qquad \sum_r d_{r,1}=32.                               \tag{D.OL10}\label{maps:eq:D.OL10}
\]
The three values in \eqref{maps:eq:D.OL9} are consequently $-16632$, $672$,
and $-32$, proving \eqref{maps:eq:D.OL3} with their full moduli.
For completeness, under the replacements
$\lambda\mapsto L+\epsilon\lambda$ and $\mu\mapsto V+\mu$,
where $\epsilon=\pm1$, the exact values are
\[
 \begin{aligned}
 -120L^2-2016\epsilon L-16632,\\
 -80LV+64L-672\epsilon V+672\epsilon,\\
 -40V^2+64V-32.
 \end{aligned}                                             \tag{D.OL11}\label{maps:eq:D.OL11}
\]
Every coefficient is divisible by eight, sixteen, and eight,
respectively. Thus the results are independent of all global
linearization choices, not only of the displayed normalization.

The same support gives all component mixed evaluations. For a
vertex $v$, let $R_j$ be the six cyclic rays of \ref{maps:src:D.2.47} and put
\[
 d_j(v)=\Psi(v)+\Psi(v+R_j)
                  -\Psi(v+R_{j-1})-\Psi(v+R_{j+1}).
\]
An integral representative of its component restriction is
\[
 \eta_v=(d_0+d_1+d_2)H-d_0E_1-d_2E_2-d_4E_3.
                                                               \tag{D.OL12}\label{maps:eq:D.OL12}
\]
Changing the representative by eight times an integral class
changes neither its square modulo sixteen nor any of its needed
pairings modulo eight. In the primitive basis of \eqref{maps:eq:D.PC10},
\[
 \langle\mathcal P_2(u)x_i,[X]\rangle=7\pmod{16}
                     \quad(0\le i<14),
\]
while the two omitted components $i=14,15$ have value fifteen.
Also $\langle ux_i y,[X]\rangle$ equals four for $i=0,1,14,15$
and zero for all the other component labels. The remaining
$\langle ux_i x_j,[X]\rangle$ are exactly
$\eta_i\cdot\nu_i^*D_j\pmod8$, with the explicit Cartier
restriction in \ref{maps:src:D.2.47}. Together with \eqref{maps:eq:D.OL3} these determine the
finite-coordinate evaluations once the support comparison is
established.

The exact program
\path{code/III/fourfolds_OS45_localization.py}
retains every wall value and every $\eta_v$, verifies the actual
quasiperiods on all sixteen residues, and proves the six
all-character polynomial divisibility identities \eqref{maps:eq:D.OL11}.
It explicitly records the required geometric support comparison
as a separate hypothesis.

In particular \(y^2=0\) integrally. Its free part is zero by
\eqref{maps:eq:D.PC12}, and its torsion part is detected injectively by the
primitive \(u\); the last equality of \eqref{maps:eq:D.OF18} kills that part.

Finally \ref{maps:src:D.2.48}'s primary coefficient-group calculation and
the integral Kunneth decomposition give
\[
 H_6(K(\mathbb Z^{15}\oplus\mathbb Z/8,2);\mathbb Z)
       =\mathbb Z^{680}\oplus(\mathbb Z/8)^{120}
                         \oplus(\mathbb Z/16)^{15}
                         \oplus\mathbb Z/8.              \tag{D.OF19}\label{maps:eq:D.OF19}
\]
The free coordinates are exactly the primitive integral cubic
\eqref{maps:eq:D.PC11}--\eqref{maps:eq:D.PC12}. Equations \eqref{maps:eq:D.OF17}--\eqref{maps:eq:D.OF18} give all
remaining coordinates: the first mixed group is detected by
\(ux_ix_j\), the second by the full Pontryagin square, and the
last by \(u^3\). Hence this determines the complete
fundamental-class invariant \(\mu\), not merely its image
under ordinary triple cup products.
The refined class remains a separate \(0\) versus \(4\)
calculation until the actual odd middle-dimensional generator
and its reference framing have been evaluated.

\path{code/III/maps_weighted_os45_mu.py} checks the exact Smith
cochain, both full affine strip constants, all sixteen wall
restrictions, the full shifted annulus map, every primitive
mixed entry and the complete moduli in \eqref{maps:eq:D.OF18}.
The independent fourfolds programs check all primitive cubic
entries and all-character localization identities.
\end{enumerate}
\end{proof}
\subsection{Moving the zero of the linearization}
\begin{theorem}[the global moving-zero family]\label{maps:src:D.3.1}
\StatementRecord{maps:src:D.3.1}{proved}{Historical events: Claude Fable 5.2 (Anthropic), effort high, 16 November 2026: carried as background without fault found. GPT-6 Astra (OpenAI), ChatGPT Pro mode, effort Pro, S6\_09\_PAPERS, 17 November 2026: targeted proof check. Also checked in this writing. Pre-continuation writing events: GPT-6 Astra (OpenAI), ChatGPT Pro mode, effort Pro, S6\_10\_PAPERS, 18 November 2026: targeted proof check. Continuation metadata events: GPT-6 Astra (OpenAI), ChatGPT Pro mode, effort Pro, S6\_10\_PAPERS continuation, 19 November 2026: used and rechecked in this run; status-vocabulary reconciliation; metadata reconciliation only; no new mathematical check. \emph{Pre-continuation recorded proof scope.} GPT-6 Astra (OpenAI, ChatGPT Pro mode), effort Pro, 2026-11-17 UTC; Claude Fable 5.2 (Anthropic), effort high, 16 November 2026; also checked in this writing by GPT-6 Astra (OpenAI), ChatGPT Pro mode, effort Pro, 18 November 2026}{the hypotheses, notation and definitions stated here; \cite{manEngel26}}


Let $S\to\mathbb P^1$ be Engel's rational elliptic surface, with
fibres $IV^*,III,I_1$ at $0,1,\infty$, respectively. Put
$L=\mathcal O_S(P-O)$, $T=t_{6P}$ and
$J=L\otimes T^*L^{-1}\simeq\pi^*\mathcal O_{\mathbb P^1}(1)$,
with the direction of $T$ fixed in \cref{push:src:A.1.2}. We write $X^{\log}(u)$ for the logarithmically transformed sphere, and $X(u)$ for the untwisted fibre of \cref{push:src:A.1.2} and \cref{push:src:A.2.2}. Choose sections
$\phi_0,\phi_1$ of $\mathcal O(1)$ with divisors $\infty,0$,
respectively, normalized so that near infinity
$\phi_0=q\phi_1$, $q=1/t$. Set
$\phi_s=\pi^*(\phi_0-s\phi_1)$.
For every sufficiently small fixed $u\ne0$ there is a proper flat
holomorphic family $\mathcal Y_u\to\Delta_s$, with a holomorphic
map to $\mathbb P^1$, having the following properties. We perform
the same logarithmic transformations of multiplicities $3,4$ and
the same admissible sphere twist data as for the chosen $X^{\log}(u)$.
The central threefold $Y_{u,0}$ has exactly one ordinary double point
and is the contraction of the diagonal $(-1,-1)$-curve of $X^{\log}(u)$.
Every fibre $Y_{u,s}$ for $s\ne0$ is smooth. Its local change at the
contracted point is
\[
 xy=u(zv-s),\qquad q=zv.
\]
This is the conifold transition described in Engel's Remark 8.5.
It is transverse to the $u\to0$ degeneration of \cref{push:src:A.1.2}.
\end{theorem}


\begin{proof}
\begin{enumerate}[label=\arabic*.,leftmargin=*]
\item The equality defining $J$ and its normalization are checked in
\cref{push:src:A.1.2}, using Engel's Lemma~2.1 and Corollary~2.2 \cite{manEngel26}. We retain
that equality, but allow the section of $J$ to vary. Here are the
charts and their identifications. Choose Hermitian norms and
$0<r<1/4$. On the compact surface $S$, after restricting both
parameters there is a uniform inequality
$|u|\|\phi_s(x)\|<r^2$. Take a core
\[
 C=\{(x,a,s):a\in L_x,\ r<\|a\|<r^{-1}\}
\]
and a node chart
\[
 N=\{(x,a,b,s):a\in L_x,\ b\in L^{-1}_{Tx},\
           a\otimes b=u\phi_s(x),\ \|a\|,\|b\|<2r\}.
\]
On the collar $r<\|a\|<2r$, identify $N$ with $C$ by
$(x,a,b,s)\mapsto(x,a,s)$. On $r<\|b\|<2r$, identify it by
$(x,a,b,s)\mapsto(Tx,b^{-1},s)$. These maps and their inverses are
holomorphic; on the first collar $b=u\phi_s/a$, and on the second
$a=u\phi_s b^{-1}$. The two node collars are disjoint by the
uniform inequality. A change of local line frames multiplies both
sides of the node equation by the same unit, so the construction is
intrinsic.

\item The quotient is Hausdorff. At each end of an attaching collar one
chart contains the limiting circle and the other omits it; there is
no pair of distinct unpaired limits of the gluing graph. The two
node collars are separated, so no new chain of identifications
changes this conclusion. For properness, the closed core
$3r/2\leq\|a\|\leq2/(3r)$ and the closed node subset
$\|a\|,\|b\|\leq7r/4$ cover the quotient. Over any compact set
of parameters their preimages are compact because $S$ is compact.
This proves properness, also for the joint map to
$\mathbb P^1\times\Delta_s$.

\item On the core, flatness over $s$ is immediate. On the node chart, the
local rings are hypersurface rings over the smooth surface $S$;
the equation $ab=u\phi_s$ is not divisible by $s-s_0$, so
multiplication by $s-s_0$ is injective. A coherent module over the
one-dimensional regular local parameter ring is flat precisely
when it is torsion free; hence the family is flat. Alternatively,
for the joint map to $(t,s)$ the fibres all have pure dimension
two, and the Cohen--Macaulay regular-sequence criterion gives
flatness over the smooth two-dimensional base.

\item Away from the zero of $\phi_s$, the collar relation gives exactly
$L^\times/\langle u\phi_s\rangle$, with the contracting generator
covering $T^{-1}$. The uniform norm inequality ensures that every
orbit meets a core or node representative, so this is an actual
biholomorphic comparison. Near $0$ and $1$ the ratio
$\phi_s/\phi_0$ is a holomorphic unit, uniformly for small $s$.
A local logarithm changes the scalar period by a holomorphic
function. Thus the same good-reduction torus cover and the same
finite affine translation data extend holomorphically in $s$.
Their finite cyclic actions stay free, by the unchanged primitive
invariant translation coordinate. We replace the two corresponding
pieces by these smooth logarithmic-transform models and glue on
the punctured overlaps by the just-described period comparison.
This is precisely the original finite gluing when $s=0$.
Properness and flatness persist, since the replacements are proper
families with the same compact annular boundaries and are locally
finite free quotients of smooth families.

\item At the node of $S_\infty$ write $q=zv$. A unit change in a line
coordinate gives $xy=u(zv-s)$. The relative gradient in
$(x,y,z,v)$ is $(y,x,-uv,-uz)$, so the only possible critical point
is the origin, and it lies on the fibre exactly when $s=0$.
At a smooth point of a zero fibre of $\phi_s$, its zero is simple,
so $xy=u(q-s)$ is smooth. These observations exhaust the singular
points: the other nodes are absent because the elliptic surface
has only its stated three exceptional fibres, and the two finite
ones were replaced by the smooth free quotients. At $s=0$ the
quadratic form $xy-uzv$ is nondegenerate and gives one ordinary
double point. Its small resolution is the diagonal subdivision of \cref{push:src:A.2.1} and \cref{push:src:A.2.2}. In one chart write $x=az$, $uv=ay$; in the other write
$z=bx$, $y=buv$, with $b=a^{-1}$. The two normal coordinates
transform by multiplication by $a$, giving
$\mathcal O(-1)\oplus\mathcal O(-1)$ over the exceptional
$\mathbb P^1$. There is one such diagonal modulo the deck periods.
The comparison with \cref{push:src:A.1.2} and \cref{push:src:A.2.2} identifies its small resolution with the
chosen $X^{\log}(u)$.

\item Before the finite logarithmic replacements, both constructions sit in
the global core/node plumbing with the explicit two-parameter equation
$xy=u(zv-s)$. The checkpoint follows the slice $s=0$ and varies
$u$ towards zero; the present construction fixes $u\ne0$ and varies
$s$. At $u=0$ the node equation is $xy=0$ and recovers the same
untwisted pinched projective bundle for every $s$. The finite
logarithmic replacements used for the sphere have only been constructed
here for $u\ne0$; no extension of those replacements across $u=0$
is asserted. The common cusp chart and the untwisted global plumbing
are the precise comparison with \cref{push:src:A.1.2} and \cref{push:src:A.2.2}.
\end{enumerate}
\end{proof}

Engel's general zero/pole construction is stated in \cite[Remark~8.5]{manEngel26}; his Remark~8.1 concerns the separate degeneration as $u$ tends to zero. The full moving-section plumbing and its geometric comparisons are proved above.

The complete credited statement and re-derivation are in \cref{maps:src:D.3.2}.
\begin{theorem}[the complex invariants of the smoothing]\label{maps:src:D.3.3}
\StatementRecord{maps:src:D.3.3}{proved}{Historical events: Claude Fable 5.2 (Anthropic), effort high, 16 November 2026: carried as background without fault found. GPT-6 Astra (OpenAI), ChatGPT Pro mode, effort Pro, S6\_09\_PAPERS, 17 November 2026: targeted proof check. Also checked in this writing. Pre-continuation writing events: GPT-6 Astra (OpenAI), ChatGPT Pro mode, effort Pro, S6\_10\_PAPERS, 18 November 2026: targeted proof check. Continuation metadata events: GPT-6 Astra (OpenAI), ChatGPT Pro mode, effort Pro, S6\_10\_PAPERS continuation, 19 November 2026: used and rechecked in this run; status-vocabulary reconciliation; metadata reconciliation only; no new mathematical check. \emph{Pre-continuation recorded proof scope.} GPT-6 Astra (OpenAI, ChatGPT Pro mode), effort Pro, 2026-11-17 UTC; Claude Fable 5.2 (Anthropic), effort high, 16 November 2026; also checked in this writing by GPT-6 Astra (OpenAI), ChatGPT Pro mode, effort Pro, 18 November 2026}{the hypotheses, notation and definitions stated here; the explicit constructions and calculations in the proof}


For $u\ne0$ sufficiently small and $0<|s|\ll1$, let $g:Y_{u,s}\to B$
be the above map and $S'_2=(g^{-1}(1))_{\mathrm{red}}$. Then
\[
 K_{Y_{u,s}}\simeq g^*\mathcal O_B(-1)\otimes\mathcal O_{Y_{u,s}}(2S'_2),
 \quad a(Y_{u,s})=1,
 \quad (h^{0,0},h^{0,1},h^{0,2},h^{0,3})=(1,1,0,0).
\]
In particular $K^{\otimes12}\simeq g^*\mathcal O_B(-6)$ and the
Kodaira dimension is $-\infty$.
\end{theorem}


\begin{proof}
\begin{enumerate}[label=\arabic*.,leftmargin=*]
\item We first compute the canonical bundle by actual transition forms.
On the core over a surface chart, the canonical form is
$\Omega_S\wedge da/a$. Changing a line frame adds a base one-form
to $da/a$, which wedges to zero with $\Omega_S$. On the second
plumbing collar the transition covers $T$ and changes the line
coordinate by the scalar $u\phi_s$; its logarithmic derivative
again contributes only a base form. Translation on the elliptic
surface preserves its invariant relative differential. Thus the
canonical transition functions are those of $K_S=\pi^*\mathcal O(-1)$.
On $xy=u(q-s)$ the residue expression
$dq\wedge dr\wedge dx/x$ becomes, using the equation,
a nonzero constant multiple of $dy\wedge dr\wedge dx$; it extends
across the node-chart loci where a line coordinate vanishes.
The expression over $q=zv$ has the same residue interpretation.
For $s\ne0$ it is regular wherever the total threefold is smooth.
Consequently moving $\phi$ introduces no new canonical divisor.

\item The two finite modifications give the same precise corrections as
in the original threefold. Indeed, for good reduction at $IV^*$
write $t=\zeta^3$, $x=\zeta^4X$, $y=\zeta^6Y$; then
$dt\wedge dx/y$ has order $2-2=0$ in $\zeta$. At $III$ write
$t-1=\zeta^4$, $x=\zeta^2X$, $y=\zeta^3Y$; its order is
$3-1=2$. Wedging with the logarithmic line coordinate gives the
threefold form, and the free quotient is etale as a map of total
spaces. These orders therefore descend respectively to zero and
$2S'_2$. This proves the asserted canonical formula. Since
$4S'_2=g^*(1)$, its twelfth power is $g^*\mathcal O(-6)$.
Connected fibres imply $g_*\mathcal O_Y=\mathcal O_B$, so none of
its positive powers has a section; if a power of $K$ had a section,
its twelfth power would contradict this fact.

\item For the algebraic dimension, use a marked regular fibre with
period matrix $\left(\begin{smallmatrix}6\mu&\tau&1&0\\
\beta_s&\mu&0&1\end{smallmatrix}\right)$.
Moving the zero changes only the scalar period $\beta_s$. Its
monodromy around the new zero $q=s$ is
$\beta_s\mapsto\beta_s\pm1$, while $\tau,\mu$ are unchanged.
If constants $a,b,c,d,e$ satisfy
$a+b\tau+c\mu+d\beta_s+e(6\mu^2-\tau\beta_s)=0$
identically, continuation around that zero gives $d-e\tau=0$.
The elliptic period $\tau$ is nonconstant, so $d=e=0$.
Continuation around infinity gives $b=0$ because
$\tau\mapsto\tau-1$ and $\mu$ is unchanged. Finally $\mu$ is
nonconstant: a constant satisfying
$\mu\circ g_1=(1-\mu)/\tau$ would force a nonconstant $\tau$ to
be constant, or give $0=1$. Hence $c=a=0$.
The type-$(1,1)$ equation for an integral alternating form now
shows that, outside a countable union of discrete subsets of the
base, the Neron--Severi group is generated by
$q=u\wedge w+6\gamma\wedge\delta$. Its associated Hermitian form
is indefinite: it is nondegenerate, its signature varies
continuously on the connected smooth-fibre locus, and at $s=0$
it has signature $(1,1)$ by the explicit period inequality.
It therefore supplies neither an ample class nor a nonzero
square-zero semipositive class. The algebraic reduction of a
complex torus is an abelian quotient; an abelian surface quotient
would supply the former class, and an elliptic quotient the latter.
Thus a very general fibre has algebraic dimension zero.
Every meromorphic function on $Y$ is constant on these fibres.
Its vertical differential consequently vanishes identically on
the regular locus, by the identity theorem, and it descends to a
meromorphic function on $B$. At an exceptional value this descent
extends meromorphically by restriction to a transverse holomorphic
disc and the bounded-pole criterion. Therefore
$\mathbb C(Y)=g^*\mathbb C(B)$ and $a(Y)=1$.

\item It remains to compute the structure-sheaf cohomology. The small
resolution $X^{\log}(u)\to Y_{u,0}$ has
$\pi_*\mathcal O_X=\mathcal O_{Y_{u,0}}$ and
$R^i\pi_*\mathcal O_X=0$ for $i>0$. For the latter assertion,
formal neighbourhoods of the exceptional curve have successive
quotients $\operatorname{Sym}^n(\mathcal O(1)^{\oplus2})$;
all their positive cohomology groups vanish, and formal functions
and coherence give the stated vanishing. Thus
$H^q(Y_{u,0},\mathcal O)=H^q(X^{\log}(u),\mathcal O)$, whose dimensions
are $(1,1,0,0)$, as computed by the supplied direct-image formulas.
Upper semicontinuity gives $h^{0,2}=h^{0,3}=0$ nearby. In a proper
flat analytic family the holomorphic Euler characteristic is
constant, as follows from the finite complex representing the
proper derived direct image over a small parameter disc. Its
value here is zero; connectedness gives $h^{0,0}=1$, and hence
$h^{0,1}=1$.
\end{enumerate}
\end{proof}

The complete credited statement and re-derivation are in \cref{maps:src:D.3.4}.

\begin{theorem}[Credited re-derivation: the diffeomorphism type of every cover]\label{maps:src:D.1.11}
\StatementRecord{maps:src:D.1.11}{proved}{Historical events: Claude Fable 5.2 (Anthropic), effort high, 16 November 2026: carried as background without fault found. GPT-6 Astra (OpenAI), ChatGPT Pro mode, effort Pro, S6\_09\_PAPERS, 17 November 2026: targeted proof check. Also checked in this writing. Pre-continuation writing events: GPT-6 Astra (OpenAI), ChatGPT Pro mode, effort Pro, S6\_10\_PAPERS, 18 November 2026: targeted proof check. Continuation metadata events: GPT-6 Astra (OpenAI), ChatGPT Pro mode, effort Pro, S6\_10\_PAPERS continuation, 19 November 2026: used and rechecked in this run; status-vocabulary reconciliation; metadata reconciliation only; no new mathematical check. \emph{Pre-continuation recorded proof scope.} GPT-6 Astra (OpenAI, ChatGPT Pro mode), effort Pro, 2026-11-17 UTC; Claude Fable 5.2 (Anthropic), effort high, 16 November 2026; also checked in this writing by GPT-6 Astra (OpenAI), ChatGPT Pro mode, effort Pro, 18 November 2026}{the hypotheses, notation and definitions stated here; \cref{maps:src:D.1.9}; \cite{manZhubr75}; \cref{maps:src:D.1.4}; \cref{maps:src:D.1.8}; \cref{maps:src:D.1.6}; \cref{maps:src:D.1.10}; \cite{manReiser24}}


For each $d\geq1$ prime to six let $M_d$ denote the oriented
simply connected spin six-manifold with the following complete
Zhubr data: $A=\mathbb Z^{d-1}\oplus\mathbb Z/d$, $b_3=0$,
refined Pontryagin class zero, integral cubic zero, and finite
cyclic cubic as in \ref{maps:src:D.1.9} in the basis
$x_0,\ldots,x_{d-2},u$. These data specify a unique oriented
diffeomorphism type. One has
\[
 \widetilde X_p\cong M_{|p|},\qquad M_1=S^6.
\]
For $d>1$ this is a classification by a complete published invariant;
no assertion that $M_d$ is a torsion-free connected sum is made.
\end{theorem}


\begin{proof}
\begin{enumerate}[label=\arabic*.,leftmargin=*]
\item We apply Zhubr's Theorem~1 \cite{manZhubr75}, with the definitions
on pp.~794--795. Its invariant quadruple consists of $H_2$, the
middle Betti number, the refined class $p$ and
$\mu=f_*[M]\in H_6(K(H_2,2);\mathbb Z)$, modulo isomorphism.
\ref{maps:src:D.1.4} gives the first two data. \ref{maps:src:D.1.8} gives $p_1=0$ and the refined
class satisfies $4p=\operatorname{PD}(p_1)$. Multiplication by four
is injective on $A$, so $p=0$. \ref{maps:src:D.1.6} and \ref{maps:src:D.1.9} give precisely all
the evaluations of $\mu$ enumerated in \ref{maps:src:D.1.10}. Thus the entire
quadruple is specified. Existence is supplied by the actual covers
already constructed, so its admissibility conditions require no
separate realization assumption. The injective classification in
Zhubr's theorem gives uniqueness and the displayed diffeomorphism.
For $d=1$ the same data are those of $S^6$. This also explains why
integral cup products alone, without \ref{maps:src:D.1.9}, would not have sufficed;
compare the torsion discussion in \cite{manReiser24} before Remark~4.7.
\end{enumerate}
\end{proof}

\section{The OS49 construction and every free affine cap}\label{sec:III:os49-night}


We retain the explicit marked construction and its complete integral cap calculation. The coprime statement \cref{maps:src:D.2.11} is extended in \cref{night:P3F.34}; its printed coprimality hypothesis remains part of its scope.


\begin{theorem}\label{night:P3F.26}\label{stmt:P3F.26}
\StatementRecord{night:P3F.26}{proved}{Historical events: GPT-6 Astra (OpenAI), ChatGPT Pro mode, effort Pro, S6\_06\_BLIND, 16--18 November 2026: targeted proof check. Also checked in this writing. Pre-continuation writing events: GPT-6 Astra (OpenAI), ChatGPT Pro mode, effort Pro, S6\_10\_PAPERS, 18 November 2026: targeted proof check. Continuation metadata events: GPT-6 Astra (OpenAI), ChatGPT Pro mode, effort Pro, S6\_10\_PAPERS continuation, 19 November 2026: used and rechecked in this run; status-vocabulary reconciliation; metadata reconciliation only; no new mathematical check. \emph{Pre-continuation recorded proof scope.} GPT-6 Astra (OpenAI, ChatGPT Pro mode), S6\_06\_BLIND, effort Pro, 16-18 November 2026 (checkpoint 25); also checked in this writing by GPT-6 Astra (OpenAI), ChatGPT Pro mode, effort Pro, 18 November 2026}{P3F.1}

The explicit $IV^*+II+I_2$ source with $M=\mathcal O(G-O)$ and $Q=3G$ has component-supported linearization divisor $C_0$. Its specified local fillings and order-three cap give a simply connected integral homology six-sphere, hence the standard smooth $S^6$.
\end{theorem}
\begin{proof}
\begin{enumerate}[label=\arabic*.,leftmargin=*]
\item The rational elliptic surface \[
 y^2=x^3-3t^3(t-1)x+2t^4(t-1)(t-\tfrac12)
                                                        \tag{P3F.26.1}
\] has fibres \(IV^*,II,I_2\) at \(0,1,\infty\), and primitive Mordell--Weil generator \(G=(t^2,t^2)\). Put \(M=\mathcal O(G-O)\), and suspend translation by \(-Q=-3G\). There is a linearization whose divisor is the identity component \(C_0\) of the \(I_2\) fibre, with multiplicity one. Its primitive rank-two toric filling at infinity, the original line-complement quotient at \(II\), and the order-three free logarithmic cap at \(IV^*\), with logarithmic vector \((1,-1,0,0)^t\), glue to a smooth compact complex manifold diffeomorphic to \(S^6\). Its integral topology is computed below, including every differential. This is the component-supported target specified in the prompt; we do not identify its linearization with a whole-fibre divisor.

\item \subsubsection{The source, section and component divisor}\label{the-source-section-and-component-divisor}

\item Write \(f=t^3(t-1)\), \(g=t^4(t-1)(t-1/2)\). Direct calculation gives \[
 f^3-g^2=-\tfrac14t^8(t-1)^2,
 \qquad j/1728=4t(1-t).
\] The vanishing orders at zero and one are respectively \((3,4,8)\) and \((1,1,2)\) for \((f,g,\Delta)\). At infinity the minimal coefficients are units and the discriminant has order two. These give exactly the asserted Kodaira fibres. Substitution verifies the displayed section. On the cubic root \(t=s^3\), set \(x=s^4X,y=s^6Y\). The central good-reduction curve is \(Y^2=X^3+1\), and \(G\) specializes to \((0,1)\), a nonzero point of order three. At infinity, with \(q=1/t\), its coordinates \((q^2x,q^3y)=(1,q)\) meet the Weierstrass node \((1,0)\), and the first resolution gives component one of \(I_2\). Its intersections with \(O\) are empty. The height is \[
 2-\frac43-\frac12=\frac16.
\] Shioda--Tate gives rank one. The verified Oguiso--Shioda Main Theorem table used in P3F.1 has exactly one row with root lattice \(E_6+A_1\), namely row49, with trivial torsion and primitive height \(1/6\). Since multiplication of a section multiplies its height by the square of that integer, the displayed \(G\) is a primitive generator. This application has the required rationality: the Weierstrass fundamental line has degree one, the base is \(\mathbb P^1\), and the section excludes multiple fibres.

\item For translation by \(-Q\), the divisor of the line of linearizations is \[
 D=G-O-(-2G)+(-3G).
\] At \(IV^*\), \(Q\) meets the identity component, so all multidegrees cancel. At \(I_2\), its multidegrees on \((C_0,C_1)\) are \((-2,2)\). They equal those of \(C_0\), since \(C_0^2=-2\) and \(C_0C_1=2\). On a rational elliptic surface the difference \(D-C_0\), having zero generic degree and zero component multidegrees, is an integral fibre class. Its intersection with \(O\) is zero: the height formula gives \((-2G)O=(-3G)O=0\), and \(DO=1=C_0O\). Thus \(D\sim C_0\). Choose its canonical section and multiply it by a small nonzero scalar \(\lambda\). This specifies the linearization and its exact zero divisor.

\item \subsubsection{Actual local fillings and their gluing}\label{actual-local-fillings-and-their-gluing}

\item Near \(I_2\), use a Tate coordinate \(Z\) on the source and a nonzero line coordinate \(V\). A logarithmic representative of \(G\) has valuation one. Inverting \(G\) changes every associated marking coherently; choose the sign giving that positive valuation. The two commuting deck periods have valuations \[
 h:(Z,V)\longmapsto(q^2aZ,qbV),\qquad
 g:(Z,V)\longmapsto(q^{-3}cZ,\lambda q^{-1}dV),       \tag{P3F.26.2}
\] where the four coefficients are nonvanishing holomorphic units. To check the second valuation without a guessed toric model, translation by \(-Q\) exchanges \(C_0\) and \(C_1\). The norm of the linearization has divisor \(C_0+C_1=F\). Therefore \(h^3g^2\), whose first valuation is zero, has second valuation one. This gives exactly the exponent \(-1\) in (P3F.26.2).

\item The integral change of deck generators \[
 H=h^{-1}g^{-1},\qquad K=h^3g^2
\] has determinant one, and its valuation columns are respectively \((1,0)^t,(0,1)^t\). Thus the precise local filling is the unit-square Mumford model triangulated by its diagonal, with units retained in the deck maps. For a triangle with lattice vertices \((i,j),(i+1,j),(i,j+1)\), the cone on the three height-one rays has determinant one; its affine chart is \(\mathbb C^3\), with \(q\) the product of the three toric coordinates. The other triangle has the analogous chart. Adjacent charts glue by the integral monomials determined by their common rays. Translations by \(H,K\) preserve the fan and multiply those monomials by holomorphic units. Modulo this full period lattice there are one vertex, three edges and two triangles. The usual annular moduli fundamental domain in \((\log|Z|,\log|V|)\), together with these boundary charts, is compact over a closed sufficiently small disc. It proves properness and the Hausdorff property. The unimodular charts prove smoothness. The central normalization is the standard hexagonal toric surface, with opposite boundary curves identified by the retained unit periods. Its Euler number is two.

\item At the type \(II\) fibre the linearization is invertible. Use the actual smooth total line complement \(M^\times\) over the original elliptic neighbourhood, divided by the contracting lift of \(-Q\). A Hermitian norm and a sufficiently small scalar make that lift strictly decrease the norm, uniformly on this compact neighbourhood. Its iterates therefore act freely and properly discontinuously, with a compact annular fundamental domain. The quotient is smooth, proper over the base disc, and agrees with the regular torus family on the punctured disc. No replacement of this additive source fibre is made.

\item At \(IV^*\), cubic good reduction gives a smooth elliptic family. The point \(Q=3G\) specializes to zero. The degree-zero line and the nonvanishing linearization consequently give a smooth two-torus family on the root chart. The linear homology deck matrix is \(\widehat T_A=(T_A^{-1})^t\), displayed below. The vector \(\ell=(1,-1,0,0)^t\) is invariant. The affine cyclic action is translation by \(\ell/3\) followed by \(\widehat T_A\). The invariant gamma covector evaluates to one on \(\ell\), so its nonidentity powers have gamma displacement \(1/3\) and \(2/3\), and are free. Its quotient is the smooth finite cap. Its punctured gluing is the holomorphic logarithmic translation \[
 (z,s)\longmapsto
 \left(z+\frac{\log s}{2\pi i}\Pi(s)\ell,s\right),
 \qquad t=s^3.
\] Changing the branch of the logarithm adds an integral period; the cubic deck changes it by the displayed affine translation. This verifies descent and gives its inverse by subtraction.

\item Glue these models to the regular family over the three-punctured sphere by these logarithmic and toric chart identifications. On every compact parameter subannulus the discs can be shrunk once to give uniform compact fundamental domains. Their union with a compact regular-base part proves compactness of each threefold and properness of the holomorphic family over that subannulus. Exhaust the punctured parameter disc by such subannuli; the constructions agree on their common ordinary locus and by the same toric charts, so they glue to the family on the whole punctured disc. This does not impose a parameter-independent cusp radius at \(\lambda=0\). Smoothness, holomorphy and the Hausdorff property follow from the explicit local charts and their annular overlap maps.

\item \subsubsection{Marked monodromy and its completeness}\label{marked-monodromy-and-its-completeness}

\item The elliptic cohomology matrices can be marked \[
 A=\begin{pmatrix}-2&-3\\1&1\end{pmatrix},\quad
 B=\begin{pmatrix}1&1\\-1&0\end{pmatrix},\quad
 C=\begin{pmatrix}1&2\\0&1\end{pmatrix},\qquad ABC=I.
\] This is a marking of the displayed surface. Its exact degree-two map \(j/1728=4t(1-t)\) has permutations \(S=(12)\) over one, \(R=1\) over zero, and cusp permutation \((12)\). The associated index-two modular cover therefore has two order-three elliptic points and one cusp of width two. Taking the local lifts of types \(IV^*\) and \(II\) gives the displayed order-three and order-six matrices, with the positive cusp transvection \(C\). There is no additional quadratic twist: an odd valuation changes one of these fixed local types, and a rational function on \(\mathbb P^1\) with all valuations even is a square.

\item There is also no unverified integral choice of the section cocycle. Since \(B-I\) is unimodular, choose the unique Abel--Jacobi lift which makes its \(B\)-cocycle zero. The \(I_2\) component implies the remaining cocycle is \((r,0)\), with \(r\) odd. The product relation gives the \(A\)-cocycle \((-r,0)\). For an elliptic matrix \(E=(\begin{smallmatrix}a&b\\c&d\end{smallmatrix})\), cocycle \((j,h)\), and scalar entry \(z\), the line/suspension weights \((1,3)\) give \[
 T(E,(j,h),z)=
 \begin{pmatrix}
 1&3(cj-ah)&3(dj-bh)&z\\
 0&a&b&j\\0&c&d&h\\0&0&0&1
 \end{pmatrix}.
\] The order-three condition gives \(z_A=-r^2\). The norm calculation above gives \(z_C=(3r^2-1)/2\). The central entry of the global product is \((r^2-1)/2\). Its only integral zeros are \(r=\pm1\), interchanged by \(G\mapsto-G\). Hence, in the positive marking, the actual cohomological tuple is \[
 T_A=\begin{pmatrix}1&-3&-3&-1\\0&-2&-3&-1\\0&1&1&0\\0&0&0&1\end{pmatrix},
 \qquad
 T_B=\begin{pmatrix}1&0&0&0\\0&1&1&0\\0&-1&0&0\\0&0&0&1\end{pmatrix},
\] \[
 T_C=\begin{pmatrix}1&0&3&1\\0&1&2&1\\0&0&1&0\\0&0&0&1\end{pmatrix}.
                                                        \tag{P3F.26.3}
\] Their product and \(T_A^3\) are the identity. These and the normal-function normalization equation are executed exactly in \path{code/part3/os49_component.py}.

\item \subsubsection{Integral topology}\label{integral-topology}

\item The local type \(II\) neighbourhood retracts onto its cuspidal rational curve, which is homeomorphic to \(S^2\): its normalization is a continuous bijection from a compact space to a Hausdorff space. The restriction of \(M\) has degree zero and is topologically trivial. A local logarithm of \(Q\) gives a translation isotopy through its identity-component group; on the central cuspidal curve this is its additive translation, fixing the cusp. Its suspension, including the line, thus retracts onto \(S^2\times T^2\). The collapse from a smooth fibre kills the two elliptic one-cycles and has degree one on its oriented two-cell. It follows integrally that its specialization image is the complete invariant lattice of \(\bigwedge^qT_B\) in every degree. At the rank-two cusp the unimodular fan calculation is exactly the normalization/double-curve complex of E26 Lemma7.4: the checked period basis has index one, so its specialization also has full integral invariant image. Its two vanishing circles are primitive.

\item In van Kampen, these two nonfinite fillings kill the section meridians and the images of \(\widehat T_B-I,\widehat T_C-I\). The eight-column relation matrix has Smith factors \((1,1,1)\); its quotient is generated by the invariant gamma direction. The base relation kills the remaining finite meridian, and its cube relation kills \(\ell\). Adding that column gives Smith factors \((1,1,1,1)\). The boundary maps are surjective onto the local filled fundamental groups, by the indicated collapse and finite quotient presentations. Since the original torus group is abelian, these are all van Kampen generators and relations; hence \(\pi_1(X)=1\).

\item For completeness the section used here is obtained from a meromorphic frame of \(M|_O\), whose degree is one, with divisor at the cusp. On the rank-two chart its valuation is an integral period vector and is subtracted by that deck period. It therefore extends to the cusp. It is invertible at \(II\). At \(IV^*\) take it as the punctured origin before the specified logarithmic translation. Thus no unrecorded affine clutch is present in these relations.

\item Use the parabolic complex, with the full matrices \(P3F.26.3\), \[
 L_q\xrightarrow{d^0}\bigoplus_{i=A,B,C}(E_i-I)L_q
       \xrightarrow{d^1}L_q,
 \qquad E_i=\bigwedge^qT_i,
\] \[
 d^0v=((E_i-I)v)_i,
 \qquad d^1(y_A,y_B,y_C)=y_A+E_Ay_B+E_AE_By_C.
\] The executed program prints the integral Hermite bases and every entry of both maps. In degrees one, two and three the local image ranks are respectively \((2,2,2),(4,4,2),(2,2,2)\), and both differentials have ranks three, five and three. Their ambient stacked and cokernel Smith forms have only unit nonzero entries. Thus the parabolic middle cohomology is zero integrally, and the last cokernels are free of rank one.

\item The global invariant generators are \[
 1,\quad\psi,\quad\omega=uw+3\psi\delta,
 \quad\phi=\psi uw,\quad\nu=\psi uw\delta.
\] The coinvariant generators are \[
 1,\quad[\delta],\quad\eta,\quad[uw\delta],\quad[\nu],
 \qquad[\psi\delta]=\eta,\quad[uw]=3\eta.
\] All other exterior basis classes map to zero.

\item At the finite cap the five-generator affine relation matrix \((\begin{smallmatrix}\widehat T_A-I&-\ell\\0&3\end{smallmatrix})\) has Smith factors \((1,1,1)\). Thus its reduced bielliptic surface has \(H_1=\mathbb Z^2\), and Poincaré duality makes its integral cohomology torsion-free. Its degree-two intersection form is unimodular. The invariant first-cohomology basis upstairs is \(\psi,h=-w+3\delta\), and its descending lattice is \(\mathbb Z(3\psi)\oplus\mathbb Zh\). The full invariant degree-two basis is \[
 u_2=-\psi w+3\psi\delta,\qquad v_2=\psi w+uw,
\] with intersection matrix \((\begin{smallmatrix}0&3\\3&0\end{smallmatrix})\), and \(\omega=u_2+v_2\). The degree-three covering multiplies intersection numbers by three, so its downstairs unimodular degree-two lattice has pullback determinant nine. This is already the determinant of the full invariant lattice; hence every invariant degree-two class descends. Perfect pairing with degree one gives degree-three index one, and degree four has index three. The global specialization indices are exactly \((3,1,1,3)\). The global invariant generators reach both nontrivial local cokernels, so the corresponding skyscraper corrections create no additional first-cohomology groups.

\item The integral Leray page consequently has only columns zero and two. Its zeroth column is generated by \(1,3\psi,\omega,\phi,3\nu\), and its second column by the five coinvariant generators above times the positive base class \(H_B\). To verify that no intrinsic differential is omitted, replace the affine order-three cap by the coarse quotient of the same root torus family under its linear action. This auxiliary compact topological filling has rational fibre cohomology equal to the invariant cohomology of the root torus: take an equivariant barycentric subdivision and average cochains. Inversion about the global origin extends over that linear quotient and the centrally symmetric toric fan; at \(II\) it is the group inversion and dual line map. It acts by \((-1)^q\) on the rational direct images, so its rational \(d_2\) is zero. The actual affine boundary differs by the explicit translation path \(z\mapsto z+t\ell/3\); three lifted prisms make the integral prism for \(\ell\). Rational transfer gives one third of its contraction. The already computed actual source and target groups are free, so this determines the integral differential uniquely. The logarithmic translation prism consequently gives \[
 d_2(a)=\frac{H_B}{3}[\iota_\ell a].
\] On \(\psi,\omega,\phi,\nu\) the contractions in coinvariants are respectively \(1,3[\delta],3\eta,[uw\delta]\). After the indicated source indices, all four integral differentials are multiplication by one. No higher differential is possible. Thus the only cohomology groups are \(H^0=H^6=\mathbb Z\), with no torsion or extension ambiguity. Simple connectivity, Hurewicz and Whitehead give a homotopy sphere, and \(\Theta_6=0\) identifies it with the standard smooth sphere.

\item The unit-circle action is semifree. The \(II\) line-complement quotient is free by strict norm contraction; the finite cap is free by its nonintegral gamma displacement; and the cusp has the standard primitive vertical-wall fan, so its only fixed set is one \(\mathbb P^1\) of normal bundle \(\mathcal O(-1)\oplus\mathcal O(-1)\), with weights \((-1,+1)\). All of these checks concern the actual component divisor filling.
\end{enumerate}
\end{proof}

\begin{theorem}\label{night:P3F.27}\label{stmt:P3F.27}
\StatementRecord{night:P3F.27}{proved}{Historical events: GPT-6 Astra (OpenAI), ChatGPT Pro mode, effort Pro, S6\_06\_BLIND, 16--18 November 2026: targeted proof check. Also checked in this writing. Pre-continuation writing events: GPT-6 Astra (OpenAI), ChatGPT Pro mode, effort Pro, S6\_10\_PAPERS, 18 November 2026: targeted proof check. Continuation metadata events: GPT-6 Astra (OpenAI), ChatGPT Pro mode, effort Pro, S6\_10\_PAPERS continuation, 19 November 2026: used and rechecked in this run; explicitly scoped metadata reconciliation; metadata reconciliation only; no new mathematical check. GPT-6 Astra (OpenAI), ChatGPT Pro mode, effort Pro, S6\_10\_PAPERS continuation, 19 November 2026: used and rechecked in this run; explicitly scoped metadata reconciliation; metadata reconciliation only; no new mathematical check. GPT-6 Astra (OpenAI), ChatGPT Pro mode, effort Pro, S6\_10\_PAPERS continuation, 19 November 2026: used and rechecked in this run; status-vocabulary reconciliation; metadata reconciliation only; no new mathematical check. \emph{Pre-continuation recorded proof scope.} GPT-6 Astra (OpenAI, ChatGPT Pro mode), S6\_06\_BLIND, effort Pro, 16-18 November 2026 (checkpoint 25); also checked in this writing by GPT-6 Astra (OpenAI), ChatGPT Pro mode, effort Pro, 18 November 2026}{P3F.31; \cref{iv:correction:P3K.6}; P3K.65}

The positive OS49 loop over $\lambda=r^3$ has actual cusp angular cycle $K=3u+2\gamma$ and retained line clutch of winding six on the sixth-root cusp cover. The exact change of page marking is given below. Its original direct graph-line comparison remains an auxiliary route; P3F.31 proves positive-fan index zero, and P3K.65 gives negative-fan class \(-3\) for this \(r\)-loop.
\end{theorem}
\begin{proof}
\begin{enumerate}[label=\arabic*.,leftmargin=*]
\item In the above component-supported construction, the primitive circle which extends as the rotation of the second nodal factor is \(K=3u+2\gamma\). It is not the original suspension cycle \(\gamma\). The change of its page boundary marking is \[
 H=-u-\gamma,\qquad K=3u+2\gamma,
 \qquad u=2H+K,\qquad\gamma=-3H-K.               \tag{P3F.27.1}
\] The boundary angle for a \(K\)-rotation evaluates to one on \(u\) and to minus one on \(\gamma\), and vanishes on \(H,w\). These statements follow by inversion of the displayed unimodular matrix, and the program executes them exactly.

\item On writing \(\lambda=r^3\), \(s=\log q/(2\pi i)\) and \(v=\log r/(2\pi i)\), the leading period columns are \[
 u=(2s,s)^t+U(q),\qquad
 \gamma=(-3s,-s+3v)^t+C(q),
\] where \(U,C\) are holomorphic. In particular \[
 H=(s,-3v)^t-U-C,\qquad
 K=(0,s+6v)^t+3U+2C.                             \tag{P3F.27.2}
\] The positive closing deck on the ordinary locus is multiplication of the line coordinate by \(t\). At the finite cap its logarithmic translation is \(\log t\,\delta/(2\pi i)\), exactly the change in that cap's gluing when \(c=\log\lambda/(2\pi i)\) increases by three. At infinity it is multiplication by \(q^{-1}\); this extends through the cusp by translation of the toric fan by its integral second basis vector. At \(II\) it is a holomorphic unit. Thus these maps close the actual holomorphic family over \(\lambda=r^3\). Its cusp difference is \(-s\delta\). The bounded extending origin cancels that difference by the integral period \(K\), because \(-s\delta+K=6v\delta+3U+2C\) is holomorphic at \(q=0\). In full detail, put \(D=3U+2C\). The two origins on \(q=Q_0^6\) are \[
 \sigma(q,v)=vD(q)+3v(v-1)\delta,
 \qquad
 o(q,v)=-\frac{sD(q)}6-\frac{s^2}{12}\delta-sv\delta.
\] The first is bounded at \(q=0\), and its characters extend by units in the cusp charts. The second has parameter difference \(-s\delta\), and subtracting it trivializes that deck. At fixed \(q\ne0\), the suspension multiplier has order three in \(r\), so the ordinary three-component rank-one parameter chain supplies the extending zero origin. Their difference and deck identities are \[
 F=\sigma-o=\frac{D(s+6v)}6+
       \left(\frac{(s+6v)^2}{12}-3v\right)\delta,
\] \[
 \begin{aligned}
 F(s,v+\theta)-F(s,v)&=\theta K+3\theta(\theta-1)\delta,\\
 F(s+6,v)-F(s,v)&=K+3\delta.
 \end{aligned}
\] The latter is an integral period. In particular the two origins descend on the specified cyclic cover. The extra real delta gauge has equal zero endpoint values, while its \(Q_0\)-deck winding is six. This proves that the positive family loop has the stated actual \(K\)-angular clutch, with its line gauge retained.

\item The existing regular-binding page of the original source has boundary cycles \((u,w,\gamma)\). On the actual cusp product \(E'_{I_1}\times S^1_K\), the source torus has cycles \((H,w)\). The old gamma angle winds on \(H\). Consequently the phase avoidance step used in P3K.6 for the original family does not apply to this compression without modification: its projected compression meets every old gamma phase. The graph map from the page to the parameter disc still extends by a radial cutoff, for any boundary angle, so that extension is not the obstruction. The unproved step is the identification of the normal-section Euler line on that graph with the old source line through a homotopy avoiding the rotating compression.

\item This direct graph-line comparison was the unresolved route at the first writing of P3F.27. P3F.31 subsequently constructed the actual positive-fan page through its line-complement subfan. P3K.65 constructs the full native C2 return and the oriented Cox ruling arrow, giving negative-fan class \(-3\) while the positive fan has class zero. That proof retains the outer normal-line marking and does not use the incomplete direct negative-fan graph homotopy.

\item The old degree-three compact elliptic pullback has three \(II\) and three \(I_2\) fibres, \(\chi=1\), signature \(-8\), height \(1/2\), total component correction \(3/2\), and \(G\cdot O=0\). These numbers are checked. Their actual positive-fan page identification is P3F.31. P3K.65 relates the actual two sphere returns by its native relative diffeomorphism and oriented Cox arrow; it does not identify the two integer indices. The earlier direct graph homotopy remains an uncompleted auxiliary route. The boundary marking and the obstruction to that particular homotopy recorded above are retained exactly.
\end{enumerate}
\end{proof}

\begin{theorem}\label{night:P3F.28}\label{stmt:P3F.28}
\StatementRecord{night:P3F.28}{proved}{Historical events: GPT-6 Astra (OpenAI), ChatGPT Pro mode, effort Pro, S6\_06\_BLIND, 16--18 November 2026: targeted proof check. Also checked in this writing. Pre-continuation writing events: GPT-6 Astra (OpenAI), ChatGPT Pro mode, effort Pro, S6\_10\_PAPERS, 18 November 2026: targeted proof check. Continuation metadata events: GPT-6 Astra (OpenAI), ChatGPT Pro mode, effort Pro, S6\_10\_PAPERS continuation, 19 November 2026: used and rechecked in this run; status-vocabulary reconciliation; metadata reconciliation only; no new mathematical check. \emph{Pre-continuation recorded proof scope.} GPT-6 Astra (OpenAI, ChatGPT Pro mode), S6\_06\_BLIND, effort Pro, 16-18 November 2026 (checkpoint 25); also checked in this writing by GPT-6 Astra (OpenAI), ChatGPT Pro mode, effort Pro, 18 November 2026}{P3F.26; P3F.27; P3F.31; P3K.65}

The actual OS49 linearization gives the exact sequence $0\to T^*M\to M\to\mathcal O_{C_0}(-1)\to0$. Its iterates alternate the two components, and all the translated source-line direct images vanish. The direct mixed-circle symbol comparison remains an auxiliary route. The actual positive-fan index is computed in P3F.31--32, and P3K.65 computes the corrected negative-fan sphere class by the native Cox arrow.
\end{theorem}
\begin{proof}
\begin{enumerate}[label=\arabic*.,leftmargin=*]
\item For the actual OS49 source and line of P3F.26, let \(T:p\mapsto p+Q\). Its inverse is the suspended translation \(-Q=-3G\). A lift covering that inverse is, at the target point \(p\), precisely a map \(M_{p+Q}\to M_p\), so its linearization is a morphism \(T^*M\to M\). Then \[
 T^*M=M(-C_0),\qquad T^{2*}M=M(-F),
                                                        \tag{P3F.28.1}
\] and the actual linearization gives the exact sequence \[
 0\longrightarrow T^*M\xrightarrow{\phi}M
       \longrightarrow\mathcal O_{C_0}(-1)
       \longrightarrow0.                              \tag{P3F.28.2}
\] The next translated sequence has quotient \(\mathcal O_{C_1}(-1)\). Every further translated sequence has the same form, alternating these two components. All of their quotients have zero derived global sections. Moreover \[
 R\pi_*M=0,
 \qquad R\pi_*(T^{j*}M)=0\quad(j\in\mathbb Z).          \tag{P3F.28.3}
\] These are exact statements about the sheaf morphisms in the actual construction. The direct spin-c symbol comparison pursued here remains an auxiliary route. P3F.31--32 give positive-fan index zero; P3K.65 gives the negative-fan class \(-3k\) for \(\lambda=r^{3k}\), retaining the actual closing and normal marking.

\item \textbf{Proof.} The divisor calculation in P3F.26 gives \(\operatorname{Hom}(T^*M,M)=\mathcal O(C_0)\), and the chosen linearization has its simple zero along \(C_0\). Multiplication by that section is injective between line bundles on the smooth source surface. Its cokernel is the restriction \(M|_{C_0}\). Since \(G\) meets \(C_1\), while \(O\) meets \(C_0\), this restriction has degree \(-1\) on \(C_0\simeq\mathbb P^1\). This proves (P3F.28.2), including the precise map and support of its cone. Translation exchanges \(C_0,C_1\), and their sum is the full reduced \(I_2\) fibre. Taking its second iterate gives \(P3F.28.1\). The next cokernel is \(T^*M|_{C_1}\), whose degree is \(M\cdot C_1-C_0\cdot C_1=1-2=-1\). Full-fibre shifts do not change these restrictions. This proves the assertion for every iterate, in both directions. The standard two affine charts on \(\mathbb P^1\) compute \(H^0(\mathcal O(-1))=H^1(\mathcal O(-1))=0\), so the derived global sections of every displayed cone vanish.

\item There is also a direct global calculation of the source index. The line \(M\) is nontrivial of degree zero on the generic elliptic fibre. Thus it has no global section. The same argument applies to \(K_S\otimes M^{-1}\), because \(K_S=-F\) restricts trivially to that fibre. Serre duality gives \(H^2(S,M)=0\). Since \(M^2=-2\) and \(M\cdot K_S=0\), surface Riemann--Roch gives \[
 \chi(S,M)=1+\frac{M\cdot(M-K_S)}2=0.
\] It follows that \(H^1(S,M)=0\) as well. The sheaf \(\pi_*M\) is torsion-free and is generically zero, hence is zero. The coherent sheaf \(R^1\pi_*M\) is supported at finitely many base points. Leray therefore identifies its space of global sections with \(H^1(S,M)=0\). A nonzero coherent sheaf with finite support has nonzero global sections, so it is zero. There are no higher direct images because the fibres have dimension one. Translation is an automorphism over the base, which proves \(P3F.28.3\) for every iterate. In particular no assumption about generic fibres alone is used to conclude vanishing at the nonreduced \(IV^*\) fibre.

\item The same calculation applies to the precise degree-three compact regular-binding pullback discussed in P3F.27. Its section line has square \(-2\), is nontrivial on the generic fibre, and its canonical class is a fibre class. Its three identity components \(C_{0,i}\) satisfy \[
 M\cdot C_{0,i}=-1,\qquad C_{0,i}^2=-2,
 \qquad C_{0,i}\cdot C_{0,j}=0\ (i\ne j).
\] Consequently the explicitly defined sheet-translation line \(M(-\sum_iC_{0,i})\) has exactly the same square as \(M\): \[
 \left(M-\sum_iC_{0,i}\right)^2
 =-2-2(-3)-6=-2.
\] Each elementary modification is the pullback of \(P3F.28.2\), with its actual quotient \(\mathcal O_{C_{0,i}}(-1)\). This proves equality of their Dolbeault indices by an explicit exact sequence. Identifying this exact sequence with the direct negative-fan graph-normal homotopy remains uncompleted in this auxiliary approach. P3K.65 supplies the actual sphere-return comparison by the native Cox ruling map. It gives different fan classes and assigns no equality of the two integer indices.
\end{enumerate}
\end{proof}

\begin{theorem}\label{night:P3F.29}\label{stmt:P3F.29}
\StatementRecord{night:P3F.29}{proved}{Historical events: GPT-6 Astra (OpenAI), ChatGPT Pro mode, effort Pro, S6\_06\_BLIND, 16--18 November 2026: targeted proof check. Also checked in this writing. Pre-continuation writing events: GPT-6 Astra (OpenAI), ChatGPT Pro mode, effort Pro, S6\_10\_PAPERS, 18 November 2026: targeted proof check. Continuation metadata events: GPT-6 Astra (OpenAI), ChatGPT Pro mode, effort Pro, S6\_10\_PAPERS continuation, 19 November 2026: used and rechecked in this run; status-vocabulary reconciliation; metadata reconciliation only; no new mathematical check. \emph{Pre-continuation recorded proof scope.} GPT-6 Astra (OpenAI, ChatGPT Pro mode), S6\_06\_BLIND, effort Pro, 16-18 November 2026 (checkpoint 25); also checked in this writing by GPT-6 Astra (OpenAI), ChatGPT Pro mode, effort Pro, 18 November 2026}{P3F.26; P3F.31; P3F.34; P3K.65}

For every OS49 cap order $m=3k$ and invariant vector $\ell=n(1,-1,0,0)+d\delta$ with $\gcd(n,m)=1$, the specialization indices are $(m,k,k,m)$ and all four integral Leray differentials are multiplication by $n$. The fundamental group and homology in degrees one through four are $\Z/|n|$. Every $n=\pm1$ member is the standard $S^6$, for every $k$ and $d$.
\end{theorem}
\begin{proof}
\begin{enumerate}[label=\arabic*.,leftmargin=*]
\item Keep the actual component-supported construction of P3F.26. Let \(m=3k\), with \(k>0\), and choose the logarithmic vector \[
 \ell=n(1,-1,0,0)^t+d(0,0,0,1)^t,
 \qquad \gcd(n,m)=1.                                  \tag{P3F.29.1}
\] The root-chart action of order \(m\), with the same order-three linear part, is free. The resulting smooth compact complex threefold has \[
 \pi_1(X)=\mathbb Z/|n|,
 \qquad H_i(X;\mathbb Z)=\mathbb Z/|n|
       \quad(1\leq i\leq4),
 \qquad H_5(X;\mathbb Z)=0,
 \qquad H_0=H_6=\mathbb Z.                            \tag{P3F.29.2}
\] Consequently every \(n=1\) or \(n=-1\) member is the standard smooth six-sphere, for every \(k\) and every delta twist \(d\). This statement determines topology in its coprime sector. P3F.31--32 construct and compute the actual positive-fan page and return; P3K.65 gives negative-fan class \(-3k\) for the positive \(r\)-loop and \(-k\) for the shorter \(w\)-loop, for both norm signs and every fixed delta character; the positive-fan index is zero on both parameters. P3F.34 supplies the additional torsion factor in the more general noncoprime free-cap sector.

\item \textbf{Proof: the actual finite cap.} Write \(A=\widehat T_A\) for the homology matrix in P3F.26. It fixes \(\ell\), and \(A^3=I\). On the \(m\)-th-root chart use the affine action \[
 \alpha(z)=Az+\ell/m.
\] Its \(m\)-th power is translation by the integral vector \(\ell\). The invariant integral covector \(\psi\) evaluates to \(n\) on \(\ell\). Thus the nonidentity powers have nonintegral gamma displacements \(jn/m\), for \(0<j<m\). They have no fixed points. The punctured gluing is the logarithmic translation already displayed in P3F.26, now with \(t=s^m\). Its inverse and branch independence are the same explicit formulas. This constructs the specified smooth cap and its attachment to the unchanged ordinary, type \(II\), and component cusp pieces.

\item The invariant first-cohomology lattice upstairs has basis \[
 \psi,\qquad h=-w+3\delta.
\] The integer map \(\Lambda\to\mathbb Z^2\) given by these two covectors is onto and has kernel \((A-I)\Lambda\): the matrix \(A-I\) has primitive image of rank two. After adjoining the affine generator, the abelianized cap group therefore has the presentation \[
 \mathbb Z^3/\langle(n,3d,-m)\rangle.
\] It is free of rank two because \(\gcd(n,m)=1\). Poincare duality and the universal coefficient theorem show that the cap surface's integral cohomology is torsion-free, with unimodular degree-two intersection form.

\item Its integral first-cohomology pullback is exactly \[
 L_1=\{a\psi+b h:na+3db\equiv0\pmod m\}.
\] If \(\bar n n\equiv1\pmod m\), a basis is \(m\psi,\ h-3d\bar n\psi\). In particular the global invariant \(\psi\) has specialization index \(m\).

\item \textbf{Proof: all specialization lattices.} Keep the degree-two invariants of P3F.26, \[
 u_2=\psi h=-\psi w+3\psi\delta,
 \qquad v_2=(\psi+u)w=\psi w+uw,
 \qquad\omega=u_2+v_2.
\] The subgroup generated by \(\alpha^3\) is the pure translation subgroup of order \(k\); it enlarges the original torus lattice by \(\ell/k\). An integral alternating form descends to that torus precisely when contraction by \(\ell\) is divisible by \(k\). Here the contractions are exactly \[
 \iota_\ell u_2=nh-3d\psi,
 \qquad \iota_\ell v_2=0.
\] Since \(n\) is invertible modulo \(k\), the invariant degree-two forms on the intermediate torus have basis \(ku_2,v_2\). Their intersection matrix, computed using that torus's fundamental class, is \(3U\): the original pairing \(u_2v_2=3\) is multiplied by \(k\) in the form and divided by \(k\) by the covering degree.

\item The remaining covering has degree three. Its downstairs intersection lattice is unimodular, so its degree-two pullback has determinant nine. The full invariant lattice \(3U\) already has determinant of absolute value nine. The index formula for a sublattice therefore forces this pullback to be the whole invariant lattice. In the original torus the precise degree-two specialization lattice is consequently \[
 L_2=\mathbb Z(ku_2)\oplus\mathbb Z v_2.
\] Thus \(\omega\) has index \(k\), and its image generates the local specialization cokernel \(\mathbb Z/k\).

\item For degree three the full invariant lattice has basis \[
 \phi=\psi uw,
 \chi=(\psi+u)w\delta.
\] The nonzero pairings with \(\psi,h\), in the original positive volume \(\nu=\psi uw\delta\), are \(\psi\chi=\nu\) and \(h\phi=-3\nu\). Perfect pairing downstairs and the covering degree \(m\) therefore identify its pullback lattice exactly as \[
 L_3=\{a\phi+b\chi:a+d\bar n b\equiv0\pmod k\}.
\] Indeed pairing with \(m\psi\) imposes no condition, while pairing with \(h-3d\bar n\psi\) requires \(m\mid3(a+d\bar n b)\). The class \(\phi\) has index \(k\) and generates the cokernel. In degree four the covering degree gives index \(m\). We have proved the four global specialization indices \[
                        (m,k,k,m).                  \tag{P3F.29.3}
\] The global invariant reaches each full local cokernel. The skyscraper exact sequence therefore creates no extra first cohomology of the direct-image sheaves. Their parabolic complexes and their free coinvariant groups are exactly the executed ones in P3F.26; the local system and the nonfinite fillings have not changed.

\item \textbf{Proof: fundamental group and every differential.} The two nonfinite fillings give the same primitive relation lattice as P3F.26, with quotient generated by gamma. The remaining base meridian dies by the base relation. Its \(m\)-th power supplies the relation \(\ell=0\), which becomes \(n\gamma=0\), since the delta generator already dies in the nonfinite quotient. These are all the boundary generators, so van Kampen gives the asserted fundamental group.

\item The inversion and coarse linear-cap comparison of P3F.26 applies without change. The actual translation prism now gives \[
 d_2(a)=\frac{H_B}{m}[\iota_\ell a].
\] In the four coinvariant targets the contractions of \(\psi,\omega,\phi,\nu\) are respectively \[
 n,\quad3n[\delta],\quad3n\eta,\quad n[uw\delta].
\] Every term involving \(d\) maps to zero. Multiplication by the four source indices \((m,k,k,m)\), followed by division by \(m=3k\), gives four differentials equal to multiplication by \(n\). The integral targets and sources are free, so the rational prism calculation determines these integral maps uniquely. There are no other differentials. It follows that \(H^2,H^3,H^4,H^5\) are \(\mathbb Z/|n|\), while \(H^1=0\) and \(H^0=H^6=\mathbb Z\). The universal coefficient theorem and Poincare duality give \(P3F.29.2\), without extension ambiguity. For \(|n|=1\), simple connectivity and these integral groups imply a homotopy sphere, and \(\Theta_6=0\) gives the standard smooth sphere.

\item The symbolic contractions, pairings, invariant bases and presentation projection are executed in \path{code/part3/os49_all_orders.py}; its finite sample checks are additional checks, not a replacement for the arbitrary-integer lattice argument above.
\end{enumerate}
\end{proof}

\begin{theorem}\label{night:P3F.31}\label{stmt:P3F.31}
\StatementRecord{night:P3F.31}{proved}{Historical events: GPT-6 Astra (OpenAI), ChatGPT Pro mode, effort Pro, S6\_06\_BLIND, 16--18 November 2026: targeted proof check. Also checked in this writing. Pre-continuation writing events: GPT-6 Astra (OpenAI), ChatGPT Pro mode, effort Pro, S6\_10\_PAPERS, 18 November 2026: targeted proof check. Continuation metadata events: GPT-6 Astra (OpenAI), ChatGPT Pro mode, effort Pro, S6\_10\_PAPERS continuation, 19 November 2026: used and rechecked in this run; explicitly scoped metadata reconciliation; metadata reconciliation only; no new mathematical check. GPT-6 Astra (OpenAI), ChatGPT Pro mode, effort Pro, S6\_10\_PAPERS continuation, 19 November 2026: used and rechecked in this run; explicitly scoped metadata reconciliation; metadata reconciliation only; no new mathematical check. GPT-6 Astra (OpenAI), ChatGPT Pro mode, effort Pro, S6\_10\_PAPERS continuation, 19 November 2026: used and rechecked in this run; status-vocabulary reconciliation; metadata reconciliation only; no new mathematical check. \emph{Pre-continuation recorded proof scope.} GPT-6 Astra (OpenAI, ChatGPT Pro mode), S6\_06\_BLIND, effort Pro, 16-18 November 2026 (checkpoint 25); also checked in this writing by GPT-6 Astra (OpenAI), ChatGPT Pro mode, effort Pro, 18 November 2026}{P3F.22; P3F.24; P3F.26; P3F.27; P3F.28; P3F.9; P3K.5; \cref{iv:correction:P3K.6}; P3K.65; P3K.7}

The explicit positive-diagonal OS49 filling over $\lambda=r^3$ is a sphere family with positive-loop return index zero, oriented homotopy-sphere class zero, and product mapping torus. Its actual page carries the source line through both $I_2$ components and nodes. P3K.65 gives the actual native C2 flop comparison and the negative-diagonal \(r\)-loop class \(-3\).
\end{theorem}
\begin{proof}
\begin{enumerate}[label=\arabic*.,leftmargin=*]
\item This statement concerns the positive-diagonal filling defined below. It leaves the negative-diagonal family of P3F.26 intact. The two fans have the common torus family and are related by the specified toric flop; they are not identified as the same family.

\item This positive-diagonal filling gives an actual sphere family over \(\lambda=r^3\), with zero positive-loop return index. The actual page, circle-line comparison, boundary support, and cap marking are constructed below. The named historical checks of these maps and of the relative cohomology argument are retained in the statement record.

\item \subsubsection{1. The exact line subfan}\label{the-exact-line-subfan}

\item Keep the periods \(h=(2,1)\), \(g=(-3,-1)\) of P3F.26 and use the periodic triangulation whose two triangles in a unit square are \[
 (0,0),(1,0),(1,1),\qquad (0,0),(0,1),(1,1).          \tag{P3F.31.1}
\] All cones at height one are unimodular, and the full period lattice has index one. The compact fundamental domains of P3F.26 prove smoothness, compactness and properness of this actual new filling. The finite cap and the type-II line-complement quotient are unchanged. Its integral topology can be compared by actual morphisms. Put \(u_j\) for the height-one rays over the four corners of a unit square. The common refinement adds the ray \[
 v=u_1+u_2=u_0+u_4=(1,1,2).
\] It is the blowup of the diagonal \(\mathbb P^1\) in either small resolution. The normal bundle of each centre is \(\mathcal O(-1)^{\oplus2}\), and the exceptional divisor is \(\mathbb P^1\times\mathbb P^1\), with the two morphisms contracting its two rulings. The equivariant period identifications give these global morphisms on the two compact fillings; there is one orbit of diagonal curves.

\item For this blowup, integral cohomology is the direct sum of pullback cohomology and the Gysin image of \(H^{q-2}(\mathbb P^1;\mathbb Z)\). One can verify the splitting by the long exact sequence of the complement of the centre and the Thom isomorphism; integration on the exceptional ruling, whose exceptional-divisor degree is \(-1\), proves that the added generator is primitive. Thus the only additional groups are one free generator in degrees two and four. Van Kampen preserves the fundamental group: both local pieces and their sphere-bundle boundary are simply connected. Apply these facts to the two morphisms from the same refinement. Since the negative-diagonal filling of P3F.26 is a simply connected integral homology sphere, so is this positive-diagonal filling. It is therefore the standard smooth six-sphere. This comparison does not identify the two families' return maps.

\item Here is the additional property needed for the page. On the infinite Tate chain of the source \(I_2\), put \[
 h_j=\lceil j/2\rceil,
 \qquad h_{j+2}=h_j+1.
\] The fan of the actual line complement \(M^\times\) has rays \((j,h_j,1)\) and cones on consecutive pairs of these rays. Projection to \((j,1)\) is the source Tate-chain projection. On each such cone the total space is \(\mathbb C^2\times\mathbb C^*\); the last factor is the delta circle's complexification. The line degree on component \(j\) is \[
 -(h_{j-1}+h_{j+1}-2h_j).
\] The minus sign follows directly from fibre-coordinate transitions: if the two adjacent slopes are \(s_-,s_+\), the fibre coordinates on the two component charts differ by \(Z^{s_--s_+}\), which is the inverse of the line-frame transition. The degrees are therefore \((-1,+1)\), exactly those of \(M=\mathcal O(G-O)\). The edges are alternately \((1,1)\) and \((1,0)\), so this is a literal subfan of (P3F.31.1). Its quotient by \(h\) is the actual \(M^\times\) over the whole source \(I_2\) neighbourhood, including both nodes. This identifies its map into the filling by toric charts, not merely by a boundary lattice.

\item It also checks the direction of the lift. Under \(g\), the graph height defect is \[
 h_j-1-h_{j-3}=0\quad(j\text{ even}),\qquad
 h_j-1-h_{j-3}=1\quad(j\text{ odd}).
\] Thus the lift is invertible over source \(C_0\) and acquires its simple zero over source \(C_1\), which is sent to target \(C_0\). This is precisely the target-point linearization \(\phi:T^*M\to M\) of P3F.28. Holomorphic units are retained in the line frames and in both periods; they do not change any of these chart inclusions.

\item For comparison, the negative-diagonal fan of P3F.26 has horizontal, vertical, and negative-diagonal edges, and does not contain this positive-diagonal line subfan. Flipping the diagonal of each unit square gives the actual toric birational comparison, through the common refinement. No diffeomorphism or equality of return indices is concluded from that comparison here.

\item \subsubsection{2. A page carrying the actual source line}\label{a-page-carrying-the-actual-source-line}

\item Choose a Hermitian norm on the source line which, on the regular annulus bounding the cusp neighbourhood, is unitary flat on every elliptic fibre. Such fibrewise metrics vary smoothly on that annulus; interpolate their logarithms to any smooth metric on a smaller neighbourhood of the central fibre. No flat extension across the nonzero component degrees is imposed. The lift covering \(-Q\), before the zero is removed, is a global holomorphic line map on the compact source. Its operator norm is bounded. After choosing \(|\lambda|\) sufficiently small, it decreases every nonzero norm on its domain by a factor bounded above by a constant \(c<1\). On a source cusp neighbourhood choose three radii \(R_1<R_2<R_3\) with \(cR_3<R_1\). The three norm shells \(\|v\|=R_j\) map to the filling through the actual line subfan. No nonzero iterate of the contracting lift identifies two points of those shells: every forward iterate that stays in the line complement has norm at most \(cR_3\), and the inverse case is the same assertion reversed. A zero cannot leave this open subset and re-enter it through a later toric translate. Indeed the exact height defect of \(g^N\), for \(N>0\), is \[
 h_j-N-h_{j-3N}=
 \begin{cases}\lfloor N/2\rfloor&j\text{ even},\\
 \lceil N/2\rceil&j\text{ odd}.
 \end{cases}
\] It is positive after the first zero; the same cone calculation at a source node uses both endpoint defects. Thus no later positive translate is a point of the line-complement subfan. The period \(h\) has already been divided out in the source line bundle. The resulting maps are therefore injective immersions. On a compact closed cusp neighbourhood the shells are compact, so these maps are embeddings.

\item Their delta quotients are three embedded copies of the actual source \(I_2\) neighbourhood. Each misses the fixed sphere, because the delta action on \(M^\times\) is free and the toric open subset just specified misses the vertical fixed edge. On each copy the actual delta bundle is exactly the chosen unit circle bundle of \(M\), by the displayed map. This gives the line identification through both nodes and both components.

\item Away from the cusp use the marked regular-binding page of the affine circle quotient. The order-three finite cap has gamma displacement \(1/3\); the corresponding circle base is the weighted Hopf three-sphere of weights \((1,3)\). Its regular binding page projects with degree three to the elliptic base, branched at the order-three value and at the regular binding. The type-II neighbourhood uses the original line-complement quotient, so its source surface is retained. On the cubic root of \(IV^*\), the source has the already verified smooth good reduction. These are the same explicit root-chart and affine translations as in P3F.26.

\item On this regular annulus the norm multiplier \(c(q)\) of the lift is constant on each elliptic fibre. Indeed translation pulls its unitary flat metric to another flat metric, and the logarithm of the norm ratio of the holomorphic line isomorphism is harmonic on the compact elliptic fibre, hence constant. Consequently \(\log\|v\|/\log c(q)\pmod1\) is the actual real suspension height, with the marked gamma cycle. The norm shells can be interpolated to levels \(\|v\|=c(q)^{a_j}\), where the three distinct \(a_j\) have the prescribed cyclic order. Choose them in a sufficiently small interval that all shell-separation inequalities remain strict on this compact annulus. An orientation-preserving reparametrization of the suspension circle then carries these levels to the three regular-binding page levels. This reparametrization is over the unchanged source and base coordinates, so it is a graph isotopy of the actual suspension charts. There is an exact lift: put \(t=\log\|v\|/\log c(q)\), lift the circle reparametrization to \(F_q:\mathbb R\to\mathbb R\) with \(F_q(t+1)=F_q(t)+1\), and send \[
 v\longmapsto c(q)^{F_q(t)-t}v.
\] Its multiplier is positive and one-periodic in \(t\). It commutes with the contracting lift, keeps the source and base coordinates fixed, and is delta-equivariant. Thus it carries the actual circle line, including its gluing. Keep the inner norm shells and the outer page fixed during the interpolation. This glues the three source cusp neighbourhoods to the outside page with the stated boundary maps.

\item This constructs an embedded four-dimensional page \(V\subset Y\) which misses the fixed sphere and has the original marked binding \(T^3\). Its projection to the degree-three base page is still the source elliptic fibration: the annular interpolation is over the unchanged base coordinate. Its compact regular-fibre cap has three \(II\) and three \(I_2\) fibres. The actual delta line on \(V\) is the pullback section line \(M\), up to a base line. To see the last assertion without discarding component data, the comparison is the exact norm-shell map at each \(I_2\), the original quotient map at each \(II\), and the good-reduction root map at the finite cap. On the smooth pieces normalize the fibrewise line isomorphism at the zero section; its overlap ratios are base functions. The base page has boundary, so its line is topologically trivial. The affine torsor origins can be chosen over this base because its punctured-cell model is a graph and the obstruction has no degree-two cochains.

\item \subsubsection{3. The full graph-to-page comparison by regular bands}\label{the-full-graph-to-page-comparison-by-regular-bands}

\item Use the relative cusp product with primitive source cycles \((H,w)\) and angular cycle \(K\). Its fixed-edge charts again have the opposite normal weights and bundle \(\mathcal O(-1)^{\oplus2}\). In detail, the two positive-fan charts along a vertical edge have coordinates \[
 (A,B,Z)=(q/V,V/Z,Z),\qquad
 (A',B',Z')=(qZ/V,V,1/Z),
\] with \((A',B')=(ZA,ZB)\), and delta weights \((-1,+1,0)\). For the opposite small resolution the orbit-product calculation of P3F.9 is obtained by exchanging the two moving conifold coordinates. Both local orbit coordinates change sign, so their transition derivative stays positive. The same relative line-coordinate construction therefore supplies the product over the unchanged nodal elliptic factor, with boundary marking \((H,w,K)\). The exact real handle of P3K.6 gives a compression \(W\) from a regular \((H,w)\)-torus at \(q=\epsilon\) to the fixed sphere, whose projection lies over the real segment \([0,\epsilon]\). Choose the regular binding at \(q=\epsilon\), outside the region where the cusp page was modified above. The K-rotation of this compression gives the same actual trace as the angular clutch \(\theta K+3\theta(\theta-1)\delta\) of P3F.27.

\item The incoming map \(E'_\epsilon\times S^1_K\to T^3\) is the actual diffeomorphism with the unimodular boundary marking \(u=2H+K\), \(\gamma=-3H-K\). On the page boundary its inverse angle is \(b-g_\epsilon(p)\), where \(g_\epsilon\) is the real-valued incoming graph of the compression. Multiplying the corresponding circle-valued function by a radial cutoff gives a map from \(V\) to the inner parameter disc. The graph of that map is embedded, and attaches to the rotating compression exactly as in P3K.6. This gives the actual trace four-manifold \(X\).

\item The page need not miss the entire projected compression. The following argument controls every remaining intersection. At \(q=0\) the compression is the fixed sphere, and the constructed page is disjoint from it. Compactness therefore gives \(q_a>0\) such that their intersection has no point over \([0,q_a]\). All remaining base values lie in \([q_a,\epsilon]\) and are regular values of the source elliptic fibration. The three lifts of this interval to the base page are disjoint regular bands ending on its boundary.

\item Here is the boundary check for a normal pushoff, which is needed because the incoming torus meets the whole binding. Near the unchanged binding write the page collar as \[
 q-\epsilon=\rho e^{i\alpha(\gamma)},\qquad \rho\geq0,
\] where \(\alpha\) winds three times. Choose its normal pushoff inside these unchanged base-coordinate annular and binding charts, with base coordinate \[
 q-\epsilon=(\rho+i\eta)e^{i\alpha(\gamma)},
 \qquad\eta>0.
\] Intersection with the real compression requires \(\rho=-\eta\cot\alpha\geq0\) and \(q-\epsilon=-\eta/\sin\alpha\leq0\). Thus it is confined to the three closed boundary arcs where \(\pi/2\leq\alpha\leq\pi\). They are proper arcs of the boundary circle. Enlarge the three lifted regular bands to contain slightly larger arcs than the actual finite-collar intersections; \(\alpha=\pi\) is only a limiting endpoint of the displayed formula, and the three enlarged arcs still have disjoint gaps. Each enlarged band is a half-disc attached to the boundary, and it still contains no critical source value. The pushoff can be taken small enough that the remainder of the page misses the entire projected compression. More explicitly, on the complement of the three bands the original page and the compact projected compression are disjoint; retain a positive separation there and make the pushoff smaller than this separation. Near the binding the displayed formula gives the required separation outside the three proper arcs. The radial homotopy keeps this pushoff's \(Y\)-coordinate fixed, so its only possible outer-annulus intersection is with the unrotated compression at parameter angle zero. All such points lie in the three bands. This gives actual neighbourhoods \[
 U_i\cong T^2\times D_i,
 \qquad U_i\cap\overline{V\setminus U_i}
       \cong T^2\times A_i,
\] where \(D_i\) is a half-disc and \(A_i\) its inner semicircle.

\item For the orientation of the meridian-one Hopf line used in Paper~IV, its comparison with the source line is the explicit fibre map $D_M:M^\times\to(M^{-1})^\times$, $D_M(v)(w)=\langle v,w\rangle/\langle v,v\rangle$. Here the Hermitian pairing is complex linear in its second argument. It satisfies $D_M(e^{i\theta}v)=e^{-i\theta}D_M(v)$, so the meridian line is $M^{-1}$ and $b|_V=-2c_1(M)$.  We keep this dual-line comparison in the graph construction below; its square is $4c_1(M)^2$.

\item Let \(x\) be the unique meridian-one class on the complement of the embedded trace in \(Y\times D^2\), as in P3K.6. Pull it back by the framed graph-normal section. Outside the displayed bands the exact radial homotopy of P3K.6 reaches the original parameter-zero page pushoff without meeting the trace. Hence this pullback equals $-c_1(M)$ there, the Chern class of the dual page line. But \[
 H^2(T^2\times D_i,T^2\times A_i;\mathbb Z)=0,
\] because \(A_i\hookrightarrow D_i\) is a homotopy equivalence. Excision and the relative long exact sequence therefore make restriction to the complement of all these bands injective in degree two. The graph-normal meridian class and the dual page-line class agree on the whole page. Possible differences of relative lifts are boundary gauges; their squares and their products with the page line vanish because its boundary Chern class is zero. This proves the actual graph-normal line comparison, including the boundary restriction.

\item \subsubsection{4. Characteristic calculation and the scope of the result}\label{characteristic-calculation-and-the-scope-of-the-result}

\item The compact elliptic page has \(\chi=1\), signature \(-8\), and \(G\cdot O=0\). Consequently its section line has relative square \(-2\), and the relative Pontryagin pairing is \(-24\). Cap the trace as in P3K.7. The different incoming torus marking is retained: the compression cap kills the primitive vanishing cycle \(w\), while the regular elliptic cap kills the old gamma cycle. The exact change of the remaining two cycles is \[
 \begin{pmatrix}u&\gamma\end{pmatrix}
 =\begin{pmatrix}H&K\end{pmatrix}
   \begin{pmatrix}2&-3\\1&-1\end{pmatrix},
 \qquad\det\begin{pmatrix}2&-3\\1&-1\end{pmatrix}=1.
\] It extends as a diffeomorphism of the two-torus factor of the compression cap. With the original marked boundary the two caps are exactly \[
 H_w=D^2_w\times S^1_u\times S^1_\gamma,\qquad
 H_\gamma=S^1_w\times S^1_u\times D^2_\gamma.
\] Their closed difference is \(S^1_u\times S^3\), by the two solid-torus decomposition in the \((w,\gamma)\) coordinates. In particular both caps have signature zero. The restriction \(H^2(H_w;\mathbb Z)\to H^2(T^3;\mathbb Z)\) is injective, so the zero boundary Chern class forces the cap line class to be zero. A change of its gluing to the page is the image of a boundary gauge; its square and product with the page line vanish by the relative cup-product formula of P3K.5. Novikov additivity and the signature theorem compute the Pontryagin pairing while retaining both cap markings. The exact graph-normal comparison above therefore gives \[
 b^2=4(-2)=-8,\qquad p_1(T\bar X)=-24.
\] The sphere-bundle calculation of P3K.7, with both return-orientation reversals retained, gives \[
 A=\frac{3b^2-p_1(T\bar X)}{24}=0.
\] The positive \(r\)-loop has the actual K-clutch of P3F.27 for this fan as well: its integer fan translations preserve \(P3F.31.1\), and its bounded origin extends in these charts. P3F.22 consequently gives zero class in \(\Theta_7\), and P3F.24 gives the product mapping torus.

\item This calculation constructs the explicit positive-diagonal family. P3K.65 subsequently proves that the negative-diagonal family in P3F.26 has \(r\)-loop class \(-3\), by its actual native C2 return and oriented Cox ruling map. Its outer normal-line marking is retained, and no equality of the two integer indices is claimed. The supplied E26 Mumford proposition permits arbitrary periodic unimodular triangulations and does not select one of these diagonals. Its literal Remark8.5 requires a whole-fibre divisor, so bare OS49 with component divisor \(C_0\) is the additional configuration separately requested by the prompt, rather than an instance of that particular hypothesis.
\end{enumerate}
\end{proof}

\begin{theorem}\label{night:P3F.32}\label{stmt:P3F.32}
\StatementRecord{night:P3F.32}{proved}{Historical events: GPT-6 Astra (OpenAI), ChatGPT Pro mode, effort Pro, S6\_06\_BLIND, 16--18 November 2026: targeted proof check. Also checked in this writing. Pre-continuation writing events: GPT-6 Astra (OpenAI), ChatGPT Pro mode, effort Pro, S6\_10\_PAPERS, 18 November 2026: targeted proof check. Continuation metadata events: GPT-6 Astra (OpenAI), ChatGPT Pro mode, effort Pro, S6\_10\_PAPERS continuation, 19 November 2026: used and rechecked in this run; status-vocabulary reconciliation; metadata reconciliation only; no new mathematical check. \emph{Pre-continuation recorded proof scope.} GPT-6 Astra (OpenAI, ChatGPT Pro mode), S6\_06\_BLIND, effort Pro, 16-18 November 2026 (checkpoint 25); also checked in this writing by GPT-6 Astra (OpenAI), ChatGPT Pro mode, effort Pro, 18 November 2026}{P3F.11; P3F.22; P3F.24; P3F.29; P3F.31; P3K.65}

For the positive-diagonal OS49 filling, every cap order $m=3k$, numerator $n=\pm1$, and integral delta twist gives a sphere family with zero positive-loop return index and product mapping torus. The actual angular clutch is $n\theta K+nm\theta(\theta-1)\delta$, and its root-deck line gauge has winding $2mn$.
\end{theorem}
\begin{proof}
\begin{enumerate}[label=\arabic*.,leftmargin=*]
\item Use the positive-diagonal fan of P3F.31, the cap order \(m=3k\), \(k>0\), and the vector \[
 \ell=n(1,-1,0,0)^t+d\delta,\qquad n\in\{1,-1\},\quad d\in\mathbb Z.
\] The exact construction below extends the positive-fan page and angular origins to every such order and twist. Every resulting sphere family has zero positive-loop return index and a product mapping torus. P3K.65 gives the corrected negative-diagonal sphere class \(-3k\) for this positive \(r\)-loop, or \(-k\) for the shorter \(w\)-loop, for both norm signs and every fixed delta character.

\item \subsubsection{1. The actual closing map and its origins}\label{the-actual-closing-map-and-its-origins}

\item Put \(\lambda=r^m\), \(v=\log r/(2\pi i)\) and \(s=\log q/(2\pi i)\). The periods are \[
 u=(2s,s)^t+U(q),\qquad
 \gamma=(-3s,-s+mv)^t+C(q),\qquad D=3U+2C,
\] so \[
 H=-u-\gamma,\qquad K=3u+2\gamma=D+(s+2mv)\delta.
\] The invariant cap covector evaluates to \(n\) on \(\ell\). When \(v\) increases by one, its affine logarithmic origin changes by \(n\log t\,\delta/(2\pi i)\). The fixed extra displacement \(d\delta/m\) has no parameter difference. Consequently multiplication of the ordinary line coordinate by \(t^n\) closes the family. It is a unit at the type-II value and is \(q^{-n}\) at the cusp. Its toric map is the integral vertical translation by \(n\), which preserves the positive-diagonal fan.

\item On \(q=Q^{2m}\), choose the two extending origins \[
 \begin{aligned}
 \sigma(s,v)&=nvD+nmv(v-1)\delta,\\
 o(s,v)&=-\frac{nsD}{2m}
          -\frac{ns^2}{4m}\delta-nsv\delta.
 \end{aligned}
\] The first is bounded in the cusp charts; the second has parameter difference \(-ns\delta\), exactly the logarithm of the closing map at the cusp. Their difference is \[
 F(s,v)=\frac{nD(s+2mv)}{2m}
       +\left(\frac{n(s+2mv)^2}{4m}-nmv\right)\delta,
\] and direct subtraction gives \[
 \begin{aligned}
 F(s,v+\theta)-F(s,v)
    &=n\theta K+nm\theta(\theta-1)\delta,\\
 F(s+2m,v)-F(s,v)&=nK+nm\delta.
 \end{aligned}
\] The cyclic descent is therefore by an integral period. Its line gauge has winding \(2mn\), while the displayed real angular gauge has zero values at both endpoints. The actual positive \(r\)-loop has \(n\) turns of the primitive \(K\)-circle. For \(n=-1\) the parametrized rotating trace is reversed; the index changes sign by P3F.11.

\item \subsubsection{2. The page, including the finite-cap line}\label{the-page-including-the-finite-cap-line}

\item The finite cap is the actual free affine quotient of P3F.29. Its gamma displacement is \(n/m\). For \(n=\pm1\) its circle base is the weighted Hopf three-sphere with one exceptional orbit of order \(m\), with the gamma orientation reversed in the negative case. The regular-binding page maps to the complex elliptic base with degree \(m\), with full branching at the finite logarithmic value and at the regular binding value. It is the explicit cyclic cover obtained by taking an \(m\)-th root of a rational coordinate whose zero and pole are those two values. The genus of its compact base is zero.

\item Here the cap hypothesis is verified by the actual root charts: its linear part has order three and fixes \(\ell\), and its \(j\)-th power has gamma displacement \(jn/m\). For \(0<j<m\) this displacement is nonintegral because \(\gcd(n,m)=1\). It is therefore free. No nonidentity delta rotation compensates such a displacement. Thus the delta action is free on these cap charts as well, exactly as required in the trace construction.

\item The finite-value root is a multiple of three, so the supplied order-three good reduction extends over that root as a smooth elliptic surface. The line on this root is the pulled-back degree-zero section line. The extra \(d\delta/m\) in the affine deck multiplies its circle coordinate by a constant root of unity. It acts as the identity on the elliptic source coordinate. Thus its change of gluing is a scalar transition on the base: on a cut annulus it is the constant character \(\exp(2\pi i d/m)\). After the page has gone around all \(m\) sheets its total character is \(\exp(2\pi i d)=1\). Choosing the scalar transitions on the individual sheet overlaps gives a line pulled back from the base page. This describes the actual additional gauge; it does not change the Picard class on an elliptic fibre. A different cap trivialization can change the compact extension only by an integral fibre class.

\item At the cusp choose \(m\) norm shells in the actual \(M^\times\) subfan, with radii \(R_1<\cdots<R_m\) satisfying \(cR_m<R_1\). The exact height defects and no-reentry proof in P3F.31 apply to every one of these shells. Its explicit positive-multiplier annular lift glues them to the \(m\) outer page sheets. At the type-II values use the original line-complement map. This is an embedded page \(V_m\subset Y\setminus F\) carrying the actual delta line \(M\), up to the base line just described.

\item Its regular compact cap has \(m\) type-II and \(m\) type-\(I_2\) fibres. The section and zero section remain disjoint under the base cover and its good-reduction extension. More explicitly, the degree-\(3k\) rational coordinate power factors through the degree-three cover of P3F.31; the further degree-\(k\) cover branches only at its smooth good-reduction fibre and at the regular binding fibre. It pulls back two already disjoint sections, so it introduces no section intersection. Hence \[
 \chi(\mathcal O)=k,\qquad \sigma=-8k,\qquad
 M_{\mathrm{rel}}^2=-2k,\qquad p_{\mathrm{rel}}=-24k.
\] These identities are unaffected by an integral fibre-class change: the section line has degree zero on a fibre, and the fibre has square zero.

\item \subsubsection{3. The actual graph line and return}\label{the-actual-graph-line-and-return}

\item For the orientation of the meridian-one Hopf line used in Paper~IV, its comparison with the source line is the explicit fibre map $D_M:M^\times\to(M^{-1})^\times$, $D_M(v)(w)=\langle v,w\rangle/\langle v,v\rangle$. Here the Hermitian pairing is complex linear in its second argument. It satisfies $D_M(e^{i\theta}v)=e^{-i\theta}D_M(v)$, so the meridian line is $M^{-1}$ and $b|_V=-2c_1(M)$.  We keep this dual-line comparison in the graph construction below; its square is $4c_1(M)^2$.

\item Use the \(K\)-compression and its exact inverse boundary angle. The annular lift above leaves the binding base coordinate unchanged. The boundary phase now winds \(m\) times, so the normal pushoff equation of P3F.31 confines its intersections with the unrotated compression to \(m\) proper boundary arcs. The page misses the fixed sphere through its actual line subfan; therefore all remaining intersections are over a compact regular base interval. Enlarge the \(m\) lifts to disjoint half-disc bands. Each pair is \[
 (T^2\times D_i,\ T^2\times A_i),
 \qquad H^2(T^2\times D_i,T^2\times A_i;\mathbb Z)=0.
\] The exact radial homotopy on the complement, followed by this restriction injectivity, identifies the graph-normal meridian line with the dual source line $M^{-1}$ on all of \(V_m\).

\item The cap marking still has \((u,\gamma)=(2H+K,-3H-K)\). Its determinant is one, and its vanishing primitive cycle is \(w\). The two marked caps and their line restrictions are precisely those in P3F.31. Thus the actual sphere-bundle bordism has \[
 b^2=-8k,\qquad p_1(T\bar X)=-24k,
\] and one primitive \(K\)-turn has \[
 A_K=\frac{3(-8k)-(-24k)}{24}=0.
\] The explicitly computed positive parameter loop has \(n\) turns and the endpoint-zero angular delta gauge, so its index is \(nA_K=0\). The topology of the positive-fan filling follows from P3F.29 and the actual common blowup of P3F.31, which adds the same primitive cohomology summands to both fillings. It is the standard smooth six-sphere. The Hopf transfer P3F.22 gives oriented return class zero, and P3F.24 gives the product mapping torus for every \(k,d\) and both signs of \(n\).
\end{enumerate}
\end{proof}

\begin{theorem}\label{night:P3F.33}\label{stmt:P3F.33}
\StatementRecord{night:P3F.33}{proved}{Historical events: Claude Opus 5.5 (Anthropic), effort max, sweep 1, scope s6, 18 November 2026: used and rechecked in this run. GPT-6 Astra (OpenAI), ChatGPT Pro mode, effort Pro, S6\_06\_BLIND, 16--18 November 2026: targeted proof check. Also checked in this writing. Pre-continuation writing events: GPT-6 Astra (OpenAI), ChatGPT Pro mode, effort Pro, S6\_10\_PAPERS, 18 November 2026: targeted proof check. Continuation metadata events: GPT-6 Astra (OpenAI), ChatGPT Pro mode, effort Pro, S6\_10\_PAPERS continuation, 19 November 2026: used and rechecked in this run; status-vocabulary reconciliation; metadata reconciliation only; no new mathematical check. \emph{Pre-continuation recorded proof scope.} GPT-6 Astra (OpenAI, ChatGPT Pro mode), S6\_06\_BLIND, effort Pro, 16-18 November 2026 (checkpoint 25); also checked in this writing by GPT-6 Astra (OpenAI), ChatGPT Pro mode, effort Pro, 18 November 2026}{P3F.26; P3F.29; P3F.31; P3F.32}

The explicit order-$3k$ OS49 affine cap is free exactly when $\gcd(k,n,d)=1$ and either $3\nmid n$ or $3\mid k$. Every free cap gives fundamental group $\Z/|n|$, so its sphere sector is exactly $n=\pm1$.
\end{theorem}
\begin{proof}
\begin{enumerate}[label=\arabic*.,leftmargin=*]
\item Keep the explicit marked OS49 root charts, put \(m=3k\), \(k>0\), and use \[
 \alpha(z)=Az+\ell/m,\qquad
 \ell=n(1,-1,0,0)^t+d\delta,\qquad n,d\in\mathbb Z.
\] Here \(A\) is exactly the integral order-three matrix of P3F.26. This cap action is free exactly for \[
 \gcd(k,n,d)=1,\qquad
 3\nmid n\ \text{or}\ 3\mid k.                 \tag{P3F.33.1}
\] For every cap in this explicit free sector, the completed threefold has fundamental group \(\mathbb Z/|n|\), with the convention \(\mathbb Z/0=\mathbb Z\). Its sphere members are consequently exactly \(n=\pm1\). Their topology is already proved in P3F.29; all their positive-diagonal returns are computed in P3F.32. This does not assert the still-uncomputed higher integral homology for the remaining non-coprime free caps.

\item \subsubsection{1. All stabilizers on the actual root chart}\label{all-stabilizers-on-the-actual-root-chart}

\item A nonidentity base rotation has no fixed point away from the central root fibre, so only that fibre needs to be checked. For \(j=3r\), \(A^j=I\), and \(\alpha^j\) is translation by \(r\ell/k\). It has a fixed point precisely when \[
 k\mid rn,\qquad k\mid rd.
\] No \(1\leq r<k\) satisfies both conditions precisely when \(\gcd(k,n,d)=1\).

\item For \(3\nmid j\), the images of \(A^j-I\) and \(A-I\) are the same. For \(j=2\) the exact identity is \(A^2-I=-A^2(A-I)\). The primitive quotient lattice \(\Lambda/(A-I)\Lambda\) has the two coordinates \[
 \psi,\qquad h=-w+3\delta
\] of P3F.29. Their values on \(\ell\) are \(n\) and \(3d\). Thus the affine fixed-point equation is soluble exactly when \[
 m\mid jn,\qquad m\mid3jd.                     \tag{P3F.33.2}
\] This is an equivalence: when the two coordinates are integral, surjectivity of \(\Lambda\to\mathbb Z^2\) supplies a lattice vector with those coordinates; subtract it and solve in the real image of \(A^j-I\). The resulting point lies in the same underlying complex torus and is an actual fixed point.

\item Put \(g=\gcd(m,n,3d)\). The integers satisfying (P3F.33.2) are the multiples of \(m/g\). There is no such integer \(j<m\) with \(3\nmid j\) precisely when \(3\mid m/g\). Together with \(\gcd(k,n,d)=1\), this is equivalent to \(P3F.33.1\). Indeed every prime other than three dividing \(g\) would divide \(k,n,d\); if \(3\mid n\), freeness requires \(3\mid k\), and the first gcd condition then forces \(3\nmid d\). This checks every group element, without a bounded search.

\item In particular the non-coprime free sector is nonempty: \(m=9,n=3,d=1\) satisfies (P3F.33.1). The gamma numerator alone does not detect its free action; the delta displacement is essential.

\item \subsubsection{2. The full fundamental-group presentation}\label{the-full-fundamental-group-presentation}

\item The cap universal cover has deck generators \(\Lambda\) and the lifted root rotation \(\alpha\), with exact relations \[
 [\Lambda,\Lambda]=1,\qquad
 \alpha\lambda\alpha^{-1}=A\lambda,\qquad
 \alpha^m=\ell.
\] This presentation is valid throughout the free sector just classified. Attach it using the same punctured torus charts as in P3F.26. The unchanged type-II and cusp fillings impose the already computed primitive relation lattice, whose quotient is generated by gamma: \(u,w,\delta\) vanish. Their base meridians vanish as well, and the base relation then kills the remaining meridian \(\alpha\). Its power relation becomes exactly \(n\gamma=0\). All its conjugation relations become trivial because the gamma covector is \(A\)-invariant. Thus the displayed presentation reduces to \[
 \langle\gamma\mid n\gamma=0\rangle.
\] This proves the fundamental group for every free cap, including the non-coprime ones. For \(|n|\ne1\) it excludes a sphere. For \(n=\pm1\) the cap satisfies the coprime hypotheses of P3F.29, which prove simple connectivity and the full integral homology of a six-sphere. The actual common blowup transfers that topology to the positive fan, as proved in P3F.31. There are therefore no additional sphere members among the non-coprime free affine caps.
\end{enumerate}
\end{proof}

\begin{theorem}\label{night:P3F.34}\label{stmt:P3F.34}
\StatementRecord{night:P3F.34}{proved}{Historical events: Claude Opus 5.5 (Anthropic), effort max, sweep 1, scope s6, 18 November 2026: computation or formal certificate. GPT-6 Astra (OpenAI), ChatGPT Pro mode, effort Pro, S6\_06\_BLIND, 16--18 November 2026: targeted proof check. Also checked in this writing. Pre-continuation writing events: GPT-6 Astra (OpenAI), ChatGPT Pro mode, effort Pro, S6\_10\_PAPERS, 18 November 2026: targeted proof check. Continuation metadata events: GPT-6 Astra (OpenAI), ChatGPT Pro mode, effort Pro, S6\_10\_PAPERS continuation, 19 November 2026: used and rechecked in this run; status-vocabulary reconciliation; metadata reconciliation only; no new mathematical check. \emph{Pre-continuation recorded proof scope.} GPT-6 Astra (OpenAI, ChatGPT Pro mode), S6\_06\_BLIND, effort Pro, 16-18 November 2026 (checkpoint 25); also checked in this writing by GPT-6 Astra (OpenAI), ChatGPT Pro mode, effort Pro, 18 November 2026}{P3F.26; P3F.29; P3F.33}

For every free affine OS49 cap with $n\ne0$, let $g=\gcd(3k,n,3d)\in\{1,3\}$. Then $H_1=H_4=\Z/|n|$ and $H_2=H_3=\Z/|n|\oplus\Z/g$. The free cap $(m,n,d)=(9,3,1)$ has $H_2=H_3=(\Z/3)^2$, correcting the general-n CSV value. The complete $n=0$ groups are also computed below.
\end{theorem}
\begin{proof}
\begin{enumerate}[label=\arabic*.,leftmargin=*]
\item \textbf{Proved; correction to the general-n CSV value.} Use the exact free sector of P3F.33, and put \[
 m=3k,\qquad g=\gcd(m,n,3d).
\] The calculation below gives, for \(n\ne0\), \[
 \begin{aligned}
 H_1(X)&=H_4(X)=\mathbb Z/|n|,\\
 H_2(X)&=H_3(X)=\mathbb Z/|n|\oplus\mathbb Z/g,\\
 H_5(X)&=0,\qquad H_0(X)=H_6(X)=\mathbb Z.
 \end{aligned}                                      \tag{P3F.34.1}
\] In the free sector \(g\) is one or three. The free example \((m,n,d)=(9,3,1)\) therefore has \(H_2=H_3=(\mathbb Z/3)^2\). This is not the general-n value stated in CSV row D.2.11. The complete integral Gysin and exterior-restriction calculation was independently audited before this comparison was marked final.

\item For \(n=0\), the free criterion forces \(3\mid k\) and \(\gcd(k,d)=1\). The same calculation gives \[
 H_1=H_5=\mathbb Z,\quad
 H_2=H_3=\mathbb Z^2\oplus\mathbb Z/3,\quad
 H_4=\mathbb Z^2,\quad H_0=H_6=\mathbb Z.             \tag{P3F.34.2}
\]

\item \subsubsection{1. The fixed exterior and the actual boundary circle}\label{the-fixed-exterior-and-the-actual-boundary-circle}

\item Remove a closed finite-cap neighbourhood. Its exterior \(V\) does not depend on \(m,n,d\): fix the same scalar \(\lambda\) and use the same ordinary, type-II and cusp charts. Its boundary \(E\) is the marked mapping torus of the order-three linear monodromy \(A\). The affine part of a torus mapping torus is removed by an explicit translation isotopy on its base interval. For the convention \((x,1)\sim(Ax+a,0)\), the map to the linear mapping torus is \[
                    (x,t)\longmapsto(x+tA^{-1}a,t).
\] At the two ends, \(a=A(A^{-1}a)\) verifies that this map descends. Choose the marked base loop by the lifted path \(x(t)=-tA^{-1}a\); its image is the zero-section base loop, and every fibre period is unchanged. The translation lift is the one supplied by the root gluing, so no integral period is inserted in this comparison. This is the same boundary trivialization used in P3F.26--29. In particular every pair \((V,E)\) below is this one fixed marked pair. Only the finite-cap logarithmic gluing changes with \((k,n,d)\); the displayed isotopy uses its chosen translation lift and adds no integral period. Thus the exterior restriction lattice used in Section 3 is independent of all three cap integers.

\item Write \(\alpha\) for its base-loop generator, and let \(\eta\in H^1(E;\mathbb Z)\) take value one on \(\alpha\) and zero on the fibre lattice. The other integral generators of \(H^1(E)\) are \[
 \psi,\qquad h=-w+3\delta.
\] The root-cap construction makes \(E\to B_{m,n,d}\) a principal normal circle bundle. Its positively marked circle, up to the common orientation convention, is the actual loop \[
             \beta_{m,n,d}=m\alpha-\ell.             \tag{P3F.34.3}
\] It goes once around the normal root disc while its image in the ordinary base winds \(m\) times. The compensating fibre translation is \(-\ell\), exactly as in the root gluing. Its values under \((\psi,h,\eta)\) are \[
                         (-n,-3d,m).
\] The free action criterion ensures that this is a free geometric circle action, even when its abelianized homology class is not primitive.

\item The cap retracts to the compact flat four-manifold \(B_{m,n,d}=T^4/\langle\alpha\rangle\). Its rational cohomology is the \(A\)-invariant exterior algebra, so its Betti numbers are \(1,2,2,2,1\) and its Euler characteristic is zero. Abelianizing its exact deck presentation gives \[
 H_1(B_{m,n,d})=\mathbb Z^3/\langle(-n,-3d,m)\rangle
               =\mathbb Z^2\oplus\mathbb Z/g.
\] In particular \[
 H^1(B_{m,n,d})=
 L_{m,n,d}:=\ker\bigl[(-n,-3d,m):\mathbb Z^3\to\mathbb Z\bigr].
                                                               \tag{P3F.34.4}
\]

\item \subsubsection{2. An integral basis of the boundary two-cohomology}\label{an-integral-basis-of-the-boundary-two-cohomology}

\item Keep the fibre forms \[
 u_2=\psi h,\qquad v_2=(\psi+u)w
\] of P3F.29, and put \[
                 p=\eta\psi,\qquad q=\eta h/3.
\] The four classes \(u_2,v_2,p,q\) form an integral basis of \(H^2(E;\mathbb Z)\). The division by three in the last class must be retained; replacing \(q\) by \(\eta h\) would lose the torsion in \(P3F.34.1\).

\item Here is a proof using the already constructed cap \((m,n,d)=(3,1,0)\). Denote its projection by \(\pi_0\). The normal circle is \(3\alpha-\ell_0\), and the integral class \(-\psi\) evaluates to one on it. Its Euler class therefore vanishes by the integral Gysin sequence, and this circle bundle is topologically trivial. The exact first-cohomology pullback from this cap has basis \[
                      h,\qquad 3\psi+\eta.
\] P3F.29 proves that the cap's integral two-cohomology restricts to the full fibre invariant lattice with basis \(u_2,v_2\). Fibre integration, or direct contraction by \(3\alpha-\ell_0\), identifies its actual boundary pullback as \[
                    u_2+q,\qquad v_2.
\] Indeed these invariant forms are horizontal; the invariant forms on the finite torus cover compute all rational cohomology by averaging. The kernel of restriction is spanned rationally by \(p,q\), and horizontality forces the displayed correction uniquely. The product splitting with the class \(-\psi\) then supplies the other two integral classes \[
 (-\psi)h=-u_2,\qquad
 (-\psi)(3\psi+\eta)=p.
\] Thus \(v_2,u_2+q,u_2,p\) is an integral basis, equivalently \(u_2,v_2,p,q\).

\item \subsubsection{3. Determine the exterior restriction without guessing a splitting}\label{determine-the-exterior-restriction-without-guessing-a-splitting}

\item Fill \(V\) by the \((3,1,0)\) cap. The resulting manifold is the proved sphere of P3F.26. Its integral Mayer--Vietoris sequence shows that \(H^2(V)\cong\mathbb Z^2\) and that restriction identifies it with a direct summand of \(H^2(E)\). Its complement is the cap pullback just computed. The homology sequence similarly gives \[
 H_1(V)=\mathbb Z,\qquad H_2(V)=\mathbb Z^2.
\] The generator of \(H_1(V)\) is gamma, as in the exact van Kampen calculation of P3F.26.

\item Let two basis vectors of the exterior restriction have coordinates in the basis \(u_2,v_2,p,q\), and project them modulo \(v_2\) to vectors \(U_1,U_2\in\mathbb Z^3\) in the ordered basis \(u_2,p,q\). Write \[
                       W=U_1\times U_2.
\] No component of this integral vector is assumed in advance. For the circle \(P3F.34.3\), integration of the invariant forms gives the exact one-classes \[
 \begin{aligned}
 (\pi_\beta)_*u_2&=3d\psi-nh,\\
 (\pi_\beta)_*v_2&=0,\\
 (\pi_\beta)_*p&=m\psi+n\eta,\\
 (\pi_\beta)_*q&=kh+d\eta.
 \end{aligned}                                      \tag{P3F.34.5}
\] The circle has period one, so these are integral fibre-integration maps, not unnormalized contractions.

\item For every \(k>0\) and \(d\in\mathbb Z\), the cap with \(n=1\) gives the already proved sphere of P3F.29. The pair sequence for attaching its disc bundle shows that its map \[
 H_1(B_{3k,1,d})\longrightarrow H_2(V)
\] is an integral isomorphism. Its dual is precisely (P3F.34.5) restricted to \(H^2(V)\). Taking cross products of its two images gives determinant \[
                 \pm\bigl(W_1 k-W_2d+W_3\bigr).
\] It has absolute value one for every positive \(k\) and every integer \(d\). Therefore \[
                   W_1=W_2=0,\qquad W_3=\pm1.
\] This proves integrally that the exterior image projects isomorphically onto \(\mathbb Z u_2\oplus\mathbb Z p\) and has no \(q\)-component. Its possible \(v_2\)-terms vanish under \(P3F.34.5\), so they do not affect any map being computed. This argument uses actual proved sphere fillings to determine the exact restriction lattice; it does not assume a boundary normalization.

\item \subsubsection{4. The two Smith invariants and every homology group}\label{the-two-smith-invariants-and-every-homology-group}

\item After an integral change of basis in \(H^2(V)\), the dual transfer map has the two columns \[
       v_1=(3d,-n,0),\qquad v_2=(m,0,n)
                 \quad\text{in }L_{m,n,d}.            \tag{P3F.34.6}
\] Their cross product is \[
              v_1\times v_2=n(-n,-3d,m).
\] The primitive normal of \(L_{m,n,d}\) is \((-n,-3d,m)/g\), so their index in this rank-two lattice is \(g|n|\). The gcd of all their entries is \(g\). Because \(L_{m,n,d}\) is primitive in \(\mathbb Z^3\), divisibility of a vector in it is the same as divisibility in the ambient lattice. Thus the first Smith invariant is \(g\), and the second is \(|n|\). Freeness implies \(g\mid n\): all primes other than three are excluded by \(\gcd(k,n,d)=1\), and a nontrivial factor three occurs exactly in the already classified \(3\mid n\) case. In fact \(g\in\{1,3\}\).

\item Use the pair sequence of \(X=V\cup N(B)\), with the Thom identification \(H_j(N(B),E)=H_{j-2}(B)\). Its degree-two portion is \[
 H_1(B)\longrightarrow H_2(V)\longrightarrow H_2(X)
 \longrightarrow H_0(B)\xrightarrow{\ n\ }H_1(V).
\] For \(n\ne0\), the last arrow is injective. Torsion in \(H_1(B)\) maps to zero in the free group \(H_2(V)\). The transpose of the free transfer has the Smith invariants just computed, and therefore \[
                   H_2(X)=\mathbb Z/g\oplus\mathbb Z/|n|.
\] P3F.33 already gives \(H_1(X)=\mathbb Z/|n|\). Euler additivity gives \(\chi(X)=2\), because \(V\) and \(E\) are fixed, every flat cap has Euler zero, and the reference filling is \(S^6\). The rational groups in degrees one and two vanish. Poincare duality and Euler characteristic then force every middle rational group to vanish. Integral duality and the universal coefficient theorem give \[
 H_3(X)\cong\operatorname{Ext}(H_2(X),\mathbb Z),\quad
 H_4(X)\cong\operatorname{Ext}(H_1(X),\mathbb Z),\quad
 H_5(X)=0.
\] This proves all groups in \(P3F.34.1\).

\item For \(n=0\), the two vectors in (P3F.34.6) generate \(3\mathbb Z\psi\) in the lattice \(\mathbb Z\psi\oplus\mathbb Z(kh+d\eta)\). The pair sequence therefore has cokernel \(\mathbb Z\oplus\mathbb Z/3\), followed by a free \(\mathbb Z\) kernel of the now-zero last arrow. It gives \(H_2=\mathbb Z^2\oplus\mathbb Z/3\). Duality and Euler characteristic give \(P3F.34.2\).

\item The executable check builds a primitive integral basis of \(P3F.34.4\), expresses \(P3F.34.6\) in that basis, and computes its Smith form for all sampled free triples. The arbitrary-integer proof is the cross-product and divisibility argument above.
\end{enumerate}
\end{proof}
\section{The fourfold and transverse fibre products}

In this section $X=X_c$ has the canonical sphere twists
$v_1=(1,2,-4,0)^t$, $v_2=(-1,-3,3,0)^t$, the untwisted cusp collar,
and an arbitrary admissible parameter $c$.  Orient all complex manifolds by
their complex structures.  Cohomology is integral unless a coefficient
field is displayed.  Put
\[
 L=\mathcal O_X(S_1-S_2).
\]
The divisor comparisons $3S_1\sim4S_2\sim W$ give
$L^4=\mathcal O_X(S_1)$, $L^6=\mathcal O_X(2S_2)$ and
$L^{12}=\mathcal O_X(W)$.  The comparisons use the meromorphic functions
$t$, $t-1$ and their quotient; thus these are comparisons of actual
holomorphic lines, not just of Chern classes.

\begin{proposition}[Counterexample: the proposed normalization]\label{maps:src:E.1.1}
\StatementRecord{maps:src:E.1.1}{proved}{Historical events: Claude Fable 5.2 (Anthropic), effort high, 16 November 2026: carried as background without fault found. GPT-6 Astra (OpenAI), ChatGPT Pro mode, effort Pro, S6\_09\_PAPERS, 17 November 2026: targeted proof check. Also checked in this writing. Pre-continuation writing events: GPT-6 Astra (OpenAI), ChatGPT Pro mode, effort Pro, S6\_10\_PAPERS, 18 November 2026: targeted proof check. Continuation metadata events: GPT-6 Astra (OpenAI), ChatGPT Pro mode, effort Pro, S6\_10\_PAPERS continuation, 19 November 2026: used and rechecked in this run; status-vocabulary reconciliation; metadata reconciliation only; no new mathematical check. \emph{Pre-continuation recorded proof scope.} GPT-6 Astra (OpenAI, ChatGPT Pro mode), effort Pro, 2026-11-17; Claude Fable 5.2 (Anthropic), effort high, 16 November 2026; also checked in this writing by GPT-6 Astra (OpenAI), ChatGPT Pro mode, effort Pro, 18 November 2026}{the hypotheses, notation and definitions stated here; the explicit constructions and calculations in the proof}


The local space $(F\times_{\widetilde\Delta}E)/\mu_m$ in Lead 6 is not
the normalization of the literal fibre product $X\times_B S$ at an
additive fibre.  It is a different proper model of its punctured family.

\end{proposition}


\begin{proof}
\begin{enumerate}[label=\arabic*.,leftmargin=*]
\item The reduced fibres of $S$ of types IV* and III have respectively seven
and two irreducible components.  In the literal fibre product the inverse
image of each such component contains an irreducible component which
dominates it: at a general smooth point of the component, the base-change
equation is $s^m=u^a$, times smooth factors.  Each branch of its
normalization still dominates that component.  A finite normalization
cannot contract any of these components.  Consequently the normalized
central fibre has at least seven, respectively two, components.  In the
proposed replacement the reduced central fibre is
$(F_0\times E_0)/\mu_m$, which is irreducible.  The component counts
disprove the asserted identification.  Good reduction of the elliptic
surface uses contraction as well as normalization; omitting that
contraction is the exact error in the reader's argument.
\end{enumerate}
\end{proof}

\begin{theorem}[a corrected compact construction]\label{maps:src:E.1.2}
\StatementRecord{maps:src:E.1.2}{proved}{Historical events: Claude Fable 5.2 (Anthropic), effort high, 16 November 2026: carried as background without fault found. GPT-6 Astra (OpenAI), ChatGPT Pro mode, effort Pro, S6\_09\_PAPERS, 17 November 2026: targeted proof check. Also checked in this writing. Pre-continuation writing events: GPT-6 Astra (OpenAI), ChatGPT Pro mode, effort Pro, S6\_10\_PAPERS, 18 November 2026: targeted proof check. Continuation metadata events: GPT-6 Astra (OpenAI), ChatGPT Pro mode, effort Pro, S6\_10\_PAPERS continuation, 19 November 2026: used and rechecked in this run; status-vocabulary reconciliation; metadata reconciliation only; no new mathematical check. \emph{Pre-continuation recorded proof scope.} GPT-6 Astra (OpenAI, ChatGPT Pro mode), effort Pro, 2026-11-17; Claude Fable 5.2 (Anthropic), effort high, 16 November 2026; also checked in this writing by GPT-6 Astra (OpenAI), ChatGPT Pro mode, effort Pro, 18 November 2026}{the hypotheses, notation and definitions stated here; \cref{maps:eq:E.1}; \cref{maps:eq:E.4}; \cref{maps:eq:E.3}}


There is a smooth compact complex fourfold $Z$ with proper holomorphic
maps $g:Z\to B$ and $\Pi:Z\to X$, with $g=f\Pi$, satisfying the local
models and the marked punctured-fibre product specified here.  The map $\Pi$
has a holomorphic section.  The construction is a small crepant
resolution of a Weierstrass model over $X$, and is bimeromorphic to
$X\times_B S$.

\end{theorem}


\begin{proof}
\begin{enumerate}[label=\arabic*.,leftmargin=*]
\item Let $a$ be a nonzero scalar multiple of the divisor section of $L^4$
and let $b$ be a nonzero scalar multiple of the square of the divisor
section of $\mathcal O_X(S_2)=L^3$.  Choose the scalars so that
$4a^3+27b^2$ is the pullback of the section of $\mathcal O_B(1)$
vanishing at $p_0$.  This is possible because $a^3$ and $b^2$ have
divisors $3S_1$ and $4S_2$, which correspond to a basis of
$H^0(B,\mathcal O_B(1))$.  The Weierstrass equation
\[
 y^2z=x^3+axz^2+bz^3
 \tag{E.1}\label{maps:eq:E.1}
\]
in the projective bundle associated to
$\mathcal O_X\oplus L^2\oplus L^3$ defines a proper complex space
$\mathscr W\to X$.  Its section at infinity is disjoint from its
singular locus.  Its $j$-function is $1728t$.

\item Over $B^\circ$, equation \eqref{maps:eq:E.1} is the product with Engel's elliptic
curve: both have that marked Weierstrass family.  At $p_1$, pass to the
actual root atlas $t=s^3$.  The coefficient $a$ has a simple zero in $s$
and $b$ is a unit.  At $p_2$, on $t-1=s^4$, $a$ is a unit and $b$ has
order two.  These are the smooth good-reduction families $E$.
The deck action on $F\times E$ is diagonal, with the specified
$A_j\oplus C_j$ linear part and translation $(v_j,0)/m_j$.
It is free because its projection to $F$ is free.  Its quotient is
therefore smooth and is exactly the local model asserted in the theorem.
The section at infinity is invariant.  Notice that the full fibres here
have multiplicities three and four; only their reductions are smooth.

\item At the cusp, use the angular--deck marking of the supplied toroidal
bridge, and write the two-torus period as
$Z_2=sB_0+C(t_c)$, where
\[
 B_0=\begin{pmatrix}0&1\\-1&0\end{pmatrix},\quad
 P=\begin{pmatrix}1&-1&0\\0&1&-1\end{pmatrix},\quad
 Q=B_0^{-1}P,\quad
 S=\begin{pmatrix}1&1\\0&1\\0&0\end{pmatrix}.
\]
Here $PS=I_2$.  Put $k=(1,1,1)^t$, $r=(0,0,1)$, and write the Tate
period of the elliptic surface as $s+c_E(t_c)$, with $c_E$ holomorphic
and bounded.  Define
\[
 Z_3=sI_3+C_3(t_c),\qquad
 C_3=SCQ+kc_Er.
 \tag{E.2}\label{maps:eq:E.2}
\]
For $A=\binom{P}{r}$ and $B=\binom{Q}{r}$ we have
\[
 \det A=\det B=1,\qquad
 AZ_3B^{-1}=\operatorname{diag}(Z_2,s+c_E).
 \tag{E.3}\label{maps:eq:E.3}
\]
This proves an integral, holomorphic product identification on the
punctured overlap, including the bounded correction which the untwisted
Tate curve alone would miss.

\item Take the fan over the periodic Kuhn tetrahedra
\[
 [m,m+e_{\pi_1},m+e_{\pi_1}+e_{\pi_2},m+(1,1,1)],
 \quad m\in\mathbb Z^3,\quad\pi\in S_3.
 \tag{E.4}\label{maps:eq:E.4}
\]
Every cone is unimodular, since subtracting its first vertex reduces
its determinant to that of a permuted triangular integral matrix with
diagonal entries one.  Every maximal chart is therefore $\mathbb C^4$,
with $t_c=z_0z_1z_2z_3$.  The action
\[
 (x,t_c)\longmapsto
 \bigl(e^{2\pi iC_3(t_c)\lambda}t_c^\lambda x,t_c\bigr),
 \qquad\lambda\in\mathbb Z^3,
 \tag{E.5}\label{maps:eq:E.5}
\]
extends to the fan by translation of its cells.  On the boundary a
fixed point would require a bounded cell to equal its nonzero translate.
On a punctured fibre, invertibility of $\operatorname{Im}Z_3$ excludes
a fixed point.  If $y=\operatorname{Im}s$ and
$K=\sup\|\operatorname{Im}C_3\|$, shrinking the disc gives
\[
 \tfrac12\|\lambda\|\leq
 \|(I+y^{-1}\operatorname{Im}C_3)\lambda\|
 \leq\tfrac32\|\lambda\|.
\]
This proves proper discontinuity by tropical coordinates.  Rounding
the three coordinates of
$(I+y^{-1}\operatorname{Im}C_3)^{-1}\operatorname{trop}(x)$ puts every
orbit in finitely many toric charts with bounded coordinates.  Over a
smaller closed disc their closures are compact.  Thus the quotient
$N_3^C$ is smooth and proper over that disc.

\item The map $(P,r)$ takes \eqref{maps:eq:E.4} to a subdivision of the product of an
$A_2$ triangle and an interval.  There are three tetrahedra in each
triangular prism.  No ray is added.  Consequently the induced map is
a small toric resolution of the local equation
$xy=w_0w_1w_2$ in \eqref{maps:eq:E.1}; over a double curve it resolves
$xy=w_0w_1$ times a smooth factor.  All rays have height one, so the
resolution is crepant.  The formula $PZ_3=Z_2Q$ proves equivariance
of its map to the cusp of $X$.  The slice $x_3=1$ gives the zero
section: its lattice map is $S$, and $S$ takes each $A_2$ triangle
to a face of a Kuhn tetrahedron.  Thus the section extends over every
double curve and both triple points, rather than only over the
punctured cusp.

\item We glue $N_3^C$ and the two finite models to the punctured product by
\eqref{maps:eq:E.3} and their root-atlas identifications.  These overlap maps are
holomorphic isomorphisms and satisfy the cocycle condition because
they are identifications with the same marked product.  Properness
over the three closed discs and compactness over the remaining compact
base give compactness of the glued Hausdorff space.  Smoothness has
been checked in all its charts.  The maps to $X$ and the zero section
agree on their overlaps.  Since both this model and the literal
fibre product agree on a dense open set, they are bimeromorphic.
\end{enumerate}
\end{proof}

\begin{theorem}[the integral cusp calculation]\label{maps:src:E.1.3}
\StatementRecord{maps:src:E.1.3}{proved}{Historical events: Claude Fable 5.2 (Anthropic), effort high, 16 November 2026: carried as background without fault found. Claude Opus 5.5 (Anthropic), effort max, sweep 1, scope s6, 18 November 2026: computation or formal certificate. GPT-6 Astra (OpenAI), ChatGPT Pro mode, effort Pro, S6\_09\_PAPERS, 17 November 2026: targeted proof check. Also checked in this writing. Pre-continuation writing events: GPT-6 Astra (OpenAI), ChatGPT Pro mode, effort Pro, S6\_10\_PAPERS, 18 November 2026: targeted proof check. Continuation metadata events: GPT-6 Astra (OpenAI), ChatGPT Pro mode, effort Pro, S6\_10\_PAPERS continuation, 19 November 2026: used and rechecked in this run; status-vocabulary reconciliation; metadata reconciliation only; no new mathematical check. \emph{Pre-continuation recorded proof scope.} GPT-6 Astra (OpenAI, ChatGPT Pro mode), effort Pro, 2026-11-17; GPT-6 Astra (OpenAI, ChatGPT Pro mode), S6\_06\_BLIND, effort Pro, 2026-11-17; Claude Fable 5.2 (Anthropic), effort high, 16 November 2026; also checked in this writing by GPT-6 Astra (OpenAI), ChatGPT Pro mode, effort Pro, 18 November 2026}{the hypotheses, notation and definitions stated here; \cref{maps:eq:E.5}; \cref{maps:eq:E.7}; \cite[Proposition~8.3.5, pp.~42--43]{mapsDeligne1974HodgeIII}; \cite[Theorem B]{ffChen}}


For the central fibre $W_3$ of \eqref{maps:eq:E.5},
\[
 H_q(W_3;\mathbb Z)=\mathbb Z^{b_q},\qquad
 (b_0,\ldots,b_6)=(1,3,9,9,10,3,1).
 \tag{E.6}\label{maps:eq:E.6}
\]
Integral specialization is onto the monodromy invariant lattice in
every degree.  Its only kernel is a copy of $\mathbb Z$ in degree four.

\end{theorem}


\begin{proof}
\begin{enumerate}[label=\arabic*.,leftmargin=*]
\item The normalization is the three-dimensional permutohedral toric variety.
Its paired boundary surfaces consist of four copies of $\mathrm{dP}_6$
and three copies of $\mathbb P^1\times\mathbb P^1$.  There are twelve
normalized triple curves and six quadruple points.  The normalization
complex for the constant sheaf is the alternating branch complex.
Local exactness follows from the augmented simplex on the branches
of $z_0z_1z_2z_3=0$.  Its nonzero rows after taking cohomology are
\[
\begin{array}{c|rrrr|rr}
q& p=0&p=1&p=2&p=3&\operatorname{rank}d^0&
 \operatorname{rank}d^1\ \text{or }d^2\\ \hline
0&1&7&12&6&0&4,5\\
2&11&22&12&0&5&9\\
4&11&7&0&0&4&0\\
6&1&0&0&0&0&0
\end{array}
\tag{E.7}\label{maps:eq:E.7}
\]
Every nonzero Smith invariant in these matrices is one.

\item Here is a reproducible specification of the integral matrices, in
addition to their complete arrays in the accompanying code output.
The fourteen rays are $\pm v$, with
$v\in\{0,1\}^3\setminus\{0\}$.  The twenty-four maximal cones
are the translated Kuhn tetrahedra containing zero.  This enumeration
requires exactly the eight cubes with lower vertex in $\{-1,0\}^3$;
no omitted cube can contain zero.  The Picard group of the normalization
is the group of ray coefficients modulo the three character relations.
We set the coefficients at $e_1,e_2,e_3$ to zero, giving an integral
basis of rank eleven.  On a boundary surface $D_v$, choose adjacent
rays $u_1,u_2$ for which $(v,u_1,u_2)$ is unimodular and set the two
remaining character coefficients to zero.  Restriction is evaluated
by subtracting the unique integral linear function with the prescribed
values at these three rays.  The paired restriction is
$\operatorname{res}_{D_v}-\operatorname{res}_{D_{-v}}$, with the
second adjacent ray corresponding to $u-v$.

\item For a triangle $(0,v,u)$ with $0<v<u$ coordinatewise, the next map
is the alternating sum of the degrees on the three boundary curves
of $D_{u-v},D_u,D_v$.  In a surface fan, if the neighbours of a ray
$w$ are $h_1,h_2$ and $h_1+h_2=bw$ modulo $v$, that degree is
$a_{h_1}+a_{h_2}-ba_w$.  In degree four, the restriction map is dual
to the seven divisor differences $[D_v]-[D_{-v}]$.  In degree zero
the matrices are the cellular incidence matrices of the one-vertex
Kuhn triangulation of $T^3$.  These rules define \eqref{maps:eq:E.7} over
$\mathbb Z$, and the program checks both compositions are zero.
Its exact Smith reductions give the ranks and primitivity just stated.

\item For completeness, the possible higher differential of this
normalization spectral sequence vanishes integrally.  Replace $C_3$
by $hC_3$, $0\leq h\leq1$.  The preceding uniform properness estimate
applies simultaneously.  Local parameter vector fields preserving
$t_c$ and each coordinate stratum patch by a partition of unity;
their flow identifies the central fibres and their specialization
maps.  We may therefore use $C_3=0$.  This toroidal degeneration is
projective after passing to the finite cover with deck lattice
$2\mathbb Z^3$.  An explicit coherent polarization is supplied by
\[
 Q(n)=4\sum n_i^2-(\sum n_i)^2.
\]
For the tetrahedron ordered $1,2,3$, its affine interpolant is
$3n_1+n_2-n_3$, and their difference is
\[
 \sum_i n_i(n_i-1)+
 \sum_{i<j}(n_i-n_j)(n_i-n_j-1)\geq0
 \quad(n\in\mathbb Z^3).
\]
Equality holds exactly at its four vertices.  Translation and
permutation give strict coherence on every face.  The finite cover
has simple normal crossings, and its descent complex is a complex
of pure toric cohomology groups.  Over $\mathbb Q$ the higher
differentials vanish by their different weights \cite[Proposition~8.3.5, pp.~42--43]{mapsDeligne1974HodgeIII}.  Taking finite
deck invariants returns the original fibre.  The integral $E_2$
groups of \eqref{maps:eq:E.7} are free, so a higher differential which is zero
over $\mathbb Q$ is zero integrally.  Summing its diagonals proves
\eqref{maps:eq:E.6}, with no extension torsion.

\item The same projective comparison allows the local invariant cycle
theorem, in the form of \cite[Theorem B]{ffChen}, followed by exact
finite-group invariants over $\mathbb Q$.  The cusp monodromy on
$H_1$ is $\left(\begin{smallmatrix}I_3&I_3\\0&I_3\end{smallmatrix}\right)$.
Its exterior-power invariant dimensions are
$(1,3,9,9,9,3,1)$, as checked by exact exterior matrices in the
attached program.  Comparison with \eqref{maps:eq:E.6} proves the last assertion.
We now strengthen this to an integral statement.  Write $b_i$ for the
three degree-one invariant deck classes and $a_i$ for the angular
classes, so the infinitesimal monodromy sends $a_i$ to $b_i$.
At $C_3=0$, every integral symmetric matrix $S$ gives a line bundle
on the toroidal family.  Here is its descent data, so no integral
invariant-cycle theorem is being assumed.  Assign to a fan vertex $n$
the integer
\[
 q_S(n)=\frac{n^tSn-\sum_iS_{ii}n_i}{2}.
\]
Interpolate on each unimodular Kuhn tetrahedron.  This is an integral
piecewise linear support function.  Translation by $\lambda$ changes
it by the integral affine function
$n^tS\lambda+q_S(\lambda)$, giving the deck linearization
$t_c^{-q_S(\lambda)}x^{-S\lambda}$.
The identity
$q_S(\lambda+\mu)=q_S(\lambda)+q_S(\mu)+\lambda^tS\mu$
is exactly the cocycle condition.  On a smooth fibre its Chern class
is $\sum S_{ij}a_i b_j$.  Restriction of this line bundle to the
central fibre therefore supplies every class
\[
 a_i b_i,\qquad a_i b_j+a_j b_i\quad(i<j)
\]
in the specialization image, as well as the $b_i$ and their products.
These elements generate the full integral invariant exterior algebra
for three pairs.  To check saturation explicitly, the accompanying
code enumerates all monomials of weighted degree at most six in these
nine generators and reduces their coordinates in an integral kernel
basis of $\bigwedge^q U-I$.  In every degree the nonzero Smith factors
are all one.  This finite check has only $2^3$ possible products of
the odd generators and at most three even factors.  It can also be
seen directly: in degree three the mixed kernel is the trace-zero
rank-eight lattice, generated by six off-diagonal terms and two
primitive diagonal differences; the complementary degrees are
generated by the corresponding products of two or three $a_i b_i$.
The parameter isotopy already constructed identifies the integral
specialization maps for general $C_3$.  Thus the images are primitive
and surjective in every degree.  Comparing their ranks with \eqref{maps:eq:E.6}
gives exactly the asserted free degree-four kernel.
\end{enumerate}
\end{proof}

\begin{theorem}[fundamental group and rational homology]\label{maps:src:E.1.4}
\StatementRecord{maps:src:E.1.4}{proved}{Historical events: Claude Fable 5.2 (Anthropic), effort high, 16 November 2026: carried as background without fault found. Claude Opus 5.5 (Anthropic), effort max, sweep 1, scope s6, 18 November 2026: computation or formal certificate. GPT-6 Astra (OpenAI), ChatGPT Pro mode, effort Pro, S6\_09\_PAPERS, 17 November 2026: targeted proof check. Also checked in this writing. Pre-continuation writing events: GPT-6 Astra (OpenAI), ChatGPT Pro mode, effort Pro, S6\_10\_PAPERS, 18 November 2026: targeted proof check. Continuation metadata events: GPT-6 Astra (OpenAI), ChatGPT Pro mode, effort Pro, S6\_10\_PAPERS continuation, 19 November 2026: used and rechecked in this run; status-vocabulary reconciliation; metadata reconciliation only; no new mathematical check. \emph{Pre-continuation recorded proof scope.} GPT-6 Astra (OpenAI, ChatGPT Pro mode; recorded reviewer /root/cdp\_audit), effort Pro, 17 November 2026; Claude Fable 5.2 (Anthropic), effort high, 16 November 2026; also checked in this writing by GPT-6 Astra (OpenAI), ChatGPT Pro mode, effort Pro, 18 November 2026}{the hypotheses, notation and definitions stated here; \cref{maps:src:E.1.2}; \cref{maps:src:E.1.3}; \cref{maps:eq:E.10}; \cref{maps:eq:E.11}}


The fourfold in Theorem \ref{maps:src:E.1.2} has
\[
 \pi_1(Z)=1,\qquad \chi(Z)=6,\qquad
 (b_0,\ldots,b_8)=(1,0,2,1,2,1,2,0,1).
 \tag{E.8}\label{maps:eq:E.8}
\]
The first homology groups of its two finite reduced fibres are,
respectively, $\mathbb Z^2\oplus\mathbb Z/3$ and
$\mathbb Z^2\oplus\mathbb Z/2$.

\end{theorem}


\begin{proof}
\begin{enumerate}[label=\arabic*.,leftmargin=*]
\item The section of $\Pi$ splits the fundamental-group sequence over its
smooth locus $X\setminus W$.  Every loop in $Z$ can be moved off
the discriminant inverse image and is a product of a loop in the
elliptic fibre and a section loop.  A section loop becomes trivial
in $Z$ because it bounds in the embedded simply connected $X$.
Thus $\pi_1(Z)$ is a quotient of the coinvariants of
$\mathbb Z^2$ under $C_1,C_2$.  Loops realizing these two
monodromies are the root-circle paths closed by the affine
identifications in $X\setminus W$.  The Smith form of
$(C_1-I\mid C_2-I)$ is $(1,1)$, so these coinvariants vanish.

\item Euler characteristic localizes on the exceptional fibres: all
regular torus fibres and the finite free torus quotients have Euler
characteristic zero.  The cusp has six zero-dimensional torus
strata, and every positive-dimensional torus stratum has Euler
characteristic zero.  Hence $\chi(Z)=6$.

\item Write $V=R^1g_*\mathbb Q$ on the punctured base, and $j$ for its
inclusion.  Finite quotient specialization is an isomorphism onto
invariants over $\mathbb Q$.  Theorem \ref{maps:src:E.1.3} gives
$R^qg_*\mathbb Q=j_*\bigwedge^qV$ except for one kernel
$\mathbb Q_{p_0}$ in degree four.  For any local system $L'$ on
the punctured sphere,
\[
 H^2(B,j_*L')=(L')_{\pi_1},\qquad
 \chi(B,j_*L')=-\operatorname{rank}L'+
       \sum_{j=0}^2\dim (L')^{M_j}.
 \tag{E.9}\label{maps:eq:E.9}
\]
The first equality is compactly supported duality and the second is
Euler additivity on the open stratum and the three point stalks.
The exact exterior-matrix computation gives the following rows of
the actual rational Leray $E_2$ page:
\[
\begin{array}{c|rrrrrrr}
q&0&1&2&3&4&5&6\\ \hline
\dim E_2^{0,q}&1&1&2&3&3&1&1\\
\dim E_2^{1,q}&0&1&0&1&0&1&0\\
\dim E_2^{2,q}&1&1&2&3&2&1&1.
\end{array}
\tag{E.10}\label{maps:eq:E.10}
\]
This is not an assumed degeneration.  The differential ranks
$r_q=\operatorname{rank}(E_2^{0,q}\to E_2^{2,q-1})$ are
\[
 (r_1,\ldots,r_6)=(1,1,2,2,1,1).
 \tag{E.11}\label{maps:eq:E.11}
\]
Indeed $\Pi^*$ transports the four nonzero transgressions of the
sphere fibration to the corresponding direct summands here.
These four transgressions are forced over $\mathbb Q$ by the
one-dimensional rows and $H^*(X;\mathbb Q)=H^*(S^6;\mathbb Q)$;
this also checks the particular reduction of Engel's Lemma 7.6
which we use.  The divisor class of the zero section gives a
permanent class $u$ restricting to the elliptic orientation class.
Multiplication by $u$ gives another independent image in degree
three, and proves $r_3=2$, $r_5=1$, $r_6=1$, and $r_4\geq2$.
The target summands are independent by the marked decomposition
$\Lambda_X\oplus\Lambda_E$.  The resulting formulas for $b_3$
and $b_5$ are $b_3=1$ and $b_5=3-r_4$.  Poincare duality of
the compact smooth fourfold gives $r_4=2$.  No later differential
exists because the base has real dimension two.  Taking diagonal
sums of \eqref{maps:eq:E.10} with \eqref{maps:eq:E.11} proves \eqref{maps:eq:E.8}.

\item Finally, for a finite quotient its affine fundamental group has
generators $\Lambda\oplus\mathbb Z a$ and abelianized relations
$(A_j\oplus C_j-I)\Lambda$ and $m_ja-(v_j,0)$.  The complete
relation matrix is printed by \texttt{fourfolds\_checks.py}.
Its nonzero Smith factors are $(1,1,1,1,3)$ and
$(1,1,1,1,2)$, respectively.  This proves the integral assertions
about the finite fibres.
\end{enumerate}
\end{proof}

\begin{theorem}[integral homology of the fourfold]\label{maps:src:E.1.5}
\StatementRecord{maps:src:E.1.5}{proved}{Historical events: Claude Fable 5.2 (Anthropic), effort high, 16 November 2026: carried as background without fault found. GPT-6 Astra (OpenAI), ChatGPT Pro mode, effort Pro, S6\_09\_PAPERS, 17 November 2026: targeted proof check. Also checked in this writing. Pre-continuation writing events: GPT-6 Astra (OpenAI), ChatGPT Pro mode, effort Pro, S6\_10\_PAPERS, 18 November 2026: targeted proof check. Continuation metadata events: GPT-6 Astra (OpenAI), ChatGPT Pro mode, effort Pro, S6\_10\_PAPERS continuation, 19 November 2026: used and rechecked in this run; status-vocabulary reconciliation; metadata reconciliation only; no new mathematical check. \emph{Pre-continuation recorded proof scope.} GPT-6 Astra (OpenAI, ChatGPT Pro mode), effort Pro, 2026-11-17; Claude Fable 5.2 (Anthropic), effort high, 16 November 2026; also checked in this writing by GPT-6 Astra (OpenAI), ChatGPT Pro mode, effort Pro, 18 November 2026}{the hypotheses, notation and definitions stated here; \cref{maps:src:E.1.2}; \cref{maps:src:E.1.4}; \cref{maps:eq:E.I1}; \cref{maps:src:E.1.3}; \cref{maps:src:D.1.4}; \cref{maps:eq:E.I2}; \cref{maps:eq:E.I3}; \cref{maps:eq:E.I5}; \cref{maps:eq:E.I4}}


The fourfold of Theorem \ref{maps:src:E.1.2} has torsion-free integral homology:
\[
 H_i(Z;\mathbb Z)=\mathbb Z^{b_i},\qquad
 (b_0,\ldots,b_8)=(1,0,2,1,2,1,2,0,1).
\]
In particular the finite-fibre torsion in \ref{maps:src:E.1.4} does not survive
as torsion in the homology of the total space.

\end{theorem}


\begin{proof}
\begin{enumerate}[label=\arabic*.,leftmargin=*]
\item We compute the integral covering cokernels, rather than substituting
integral invariants in the rational calculation. Write
\(\mathcal V=R^1g_*\mathbb Z|_{B^\circ}\), and put
\(\mathcal J_q=j_*\bigwedge^q\mathcal V\).
The cohomology monodromies are the inverse transposes of
\(A_1\oplus C_1,A_2\oplus C_2,M_0\oplus C_0\).
All lattices below are in this marking.

\item First we describe a finite filling without making any equivariant
formality assumption. Since the first coordinate of each twist
\(v_j\) is \(\pm1\), the integral change of basis with first
column \(v_j/(v_j)_\gamma\) and remaining columns the last five
standard basis vectors changes its linear monodromy to
\(\operatorname{diag}(1,D_j)\) and its translation to
\(\pm e_0/m_j\). The reduced finite fibre is the mapping torus
of the actual linear automorphism \(D_j:T^5\to T^5\).
The covering \(T^6\to B_j\) is the \(m_j\)-fold cover of its
suspension circle. Set \(T_j=D_j^{-t}\). The two vertical matrices
on homology, in coordinates \((u,w,\delta,e_1,e_2)\), are
\[
D_1=\begin{pmatrix}
0&1&0&0&0\\-1&-1&0&0&0\\1&0&1&0&0\\
0&0&0&-1&-1\\0&0&0&1&0
\end{pmatrix},\qquad
D_2=\begin{pmatrix}
0&-1&0&0&0\\1&0&0&0&0\\0&1&1&0&0\\
0&0&0&0&1\\0&0&0&-1&0
\end{pmatrix}.
\]
The Wang sequences of the mapping torus and its cover give the
following exact description of the cokernel of covering pullback:
\[
 C_j^q\simeq
 \frac{\ker(\bigwedge^{q-1}T_j-I)}
 {\operatorname{im}\mathcal N_j^{q-1}},\qquad
 \mathcal N_j^{q-1}=\sum_{a=0}^{m_j-1}
                  (\bigwedge^{q-1}T_j)^a .
 \tag{E.I1}\label{maps:eq:E.I1}
\]
Here the target of pullback is the invariant lattice in
\(H^q(T^6;\mathbb Z)\). To verify \eqref{maps:eq:E.I1}, its restriction to
the vertical torus surjects onto
\(\ker(\bigwedge^qT_j-I)\), and restriction of a mapping-torus
class has exactly the same image, by Wang exactness. Thus every
cokernel class has a representative containing the suspension
covector. On this kernel the covering map is the displayed norm:
cut the covered circle into its \(m_j\) consecutive intervals
and transport the vertical cochain through them. This proves both
surjectivity and the precise kernel of the horizontal description.
It does not choose a splitting of either Wang sequence.

\item For \(q=2\), the two quotients in \eqref{maps:eq:E.I1} are respectively zero
and \(\mathbb Z/2\). For \(q=3\), they are
\[
 C_1^3=\mathbb Z/3,\qquad
 C_2^3=(\mathbb Z/2)^2\oplus\mathbb Z/4.
 \tag{E.I2}\label{maps:eq:E.I2}
\]
The same Wang sequences show that the kernels of finite
specialization in degree two are \(\mathbb Z/3\) and
\(\mathbb Z/2\), and that the degree-three and degree-four
kernels vanish. The cusp specialization is integrally onto,
with the sole free degree-four kernel of \ref{maps:src:E.1.3}.

\item We give the integral global maps needed to use these local groups.
The global invariant two-form
\(q_X=u^*w^*+6\gamma\delta^*\) maps onto \(C_2^2\).
Indeed, if its pullback descended to the six-dimensional finite
fibre, restriction to the elliptic zero section would make it
descend to the original bielliptic surface \(S_2\). Its square
there would be \(\pm12/4=\pm3\), which is impossible in the
even intersection form of \(S_2\). That form is even because
\(H^2(S_2;\mathbb Z)\) is torsion free and its canonical class
is torsion, as checked in \ref{maps:src:D.1.4}. Thus the global-to-local map in
degree two is onto its only nonzero covering cokernel.

\item For degree three all three global invariant classes contain
\(\gamma\), and remain purely horizontal in both adapted
finite markings. They are
\[
 \gamma e_1^*e_2^*,\qquad
 \gamma(-u^*e_1^*+w^*e_2^*),\qquad
 \gamma u^*w^* .
 \tag{E.I3}\label{maps:eq:E.I3}
\]
Juxtaposition denotes exterior product. Consequently their images
in \eqref{maps:eq:E.I2} are exactly the vertical two-forms in \eqref{maps:eq:E.I3} modulo the
norm lattices; there is no unspecified lift in this computation.
In the integral invariant bases printed by the accompanying program,
the two norm matrices and this common image matrix are
\[
\begin{split}
N_3&=\begin{pmatrix}
0&0&0&0&0&0&0&0&0&3\\
0&0&-2&1&0&1&1&1&-1&0\\
0&0&-1&2&0&2&-1&0&-1&0\\
3&-1&0&0&2&0&0&0&0&0
\end{pmatrix},\\
N_4&=\begin{pmatrix}
0&0&0&0&0&0&0&0&0&4\\
0&0&-2&0&0&0&2&1&-1&0\\
0&0&0&2&0&2&0&-1&-1&0\\
4&-2&0&0&2&0&0&0&0&0
\end{pmatrix},\qquad
G=\begin{pmatrix}1&0&0\\0&1&0\\0&0&0\\0&0&1\end{pmatrix}.
\end{split}
\]
Every nonzero Smith factor of
\(\bigl(\operatorname{diag}(N_3,N_4)\mid(G,G)^t\bigr)\)
is one; there are eight. Thus the global classes surject onto
the direct sum of the two degree-three covering cokernels.

\item We also compute the integral parabolic groups, retaining torsion.
The degree-two global coboundary matrix, obtained by stacking
the two monodromy-minus-identity matrices, has a single nontrivial
Smith factor six. Its torsion generator restricts with orders
three, two and one at the three boundary circles. Hence no
nonzero torsion class is parabolic. Its rational parabolic rank
is zero, so \(H^1(B,\mathcal J_2)=0\). In degree three the
global coboundary image is primitive and the rational parabolic
rank is one; hence \(H^1(B,\mathcal J_3)=\mathbb Z\).
These assertions use the cocycle description as the kernel of
restriction from the free-group cohomology to the three
peripheral circle groups. The matrices and the restriction
orders are computed exactly in
\texttt{fourfolds\_integral\_audit.py}.

\item Passing from \(\mathcal J_q\) to \(R^qg_*\mathbb Z\)
adds the point-supported kernels just described and replaces
the finite stalks by the pullback sublattices. A point-supported
sheaf has no positive cohomology. The long exact sequences,
together with the two surjections onto covering cokernels,
therefore give
\[
 H^1(B,R^2g_*\mathbb Z)=0,\qquad
 H^1(B,R^3g_*\mathbb Z)=\mathbb Z.
 \tag{E.I4}\label{maps:eq:E.I4}
\]
They also identify degree-two base cohomology with the global
coinvariants. In the relevant degrees these are torsion-free,
with
\[
 H^2(B,R^1g_*\mathbb Z)=\mathbb Z,
 \qquad H^2(B,R^2g_*\mathbb Z)=\mathbb Z^2.
 \tag{E.I5}\label{maps:eq:E.I5}
\]
The first is the corresponding group from \(X\). The second
is the direct sum of that from \(\bigwedge^2\mathcal V_X\)
and the elliptic orientation line; the mixed coinvariants are
zero. These identifications, including primitivity, are also
verified by Smith forms, so no finite index is suppressed.

\item The integral differential onto the first group in \eqref{maps:eq:E.I5} is
surjective by naturality from the sphere fibration: its own
corresponding Leray quotient must vanish because \(H^3(X)=0\).
The differential onto the second group is surjective as well.
The \(X\) summand is killed by the differential transported
from \(H^4(X)=0\); the elliptic orientation summand is killed
by multiplying the primitive degree-one transgression of \(X\)
by the permanent class of the elliptic zero-section divisor.
That divisor has degree one on each elliptic fibre. Thus these
are two primitive independent target generators, not merely two
independent rational images.

\item Now \(H^0(B,R^3g_*\mathbb Z)\) and
\(H^0(B,R^4g_*\mathbb Z)\) are free: in the first degree
specialization has no kernel, and in the second its only
kernel is the free cusp class. There is no later Leray
differential, since the base has cohomological dimension two.
Equation \eqref{maps:eq:E.I4} and the first surjectivity in \eqref{maps:eq:E.I5} show that
\(H^3(Z;\mathbb Z)\) is a free group, of rank one by \ref{maps:src:E.1.4}.
Equation \eqref{maps:eq:E.I4} and the second surjectivity show that
\(H^4(Z;\mathbb Z)\) is an extension of two free groups of
rank one, and hence is \(\mathbb Z^2\).
Simple connectivity makes \(H^2(Z;\mathbb Z)\) torsion free.
The universal coefficient theorem now proves successively that
\(H_2(Z;\mathbb Z)=\mathbb Z^2\) and
\(H_3(Z;\mathbb Z)=\mathbb Z\). Poincare duality determines
the other degrees and gives exactly the displayed groups.
\end{enumerate}
\end{proof}


\begin{theorem}[complex invariants and middle form]\label{maps:src:E.2.1}
\StatementRecord{maps:src:E.2.1}{proved}{Historical events: Claude Fable 5.2 (Anthropic), effort high, 16 November 2026: carried as background without fault found. GPT-6 Astra (OpenAI), ChatGPT Pro mode, effort Pro, S6\_09\_PAPERS, 17 November 2026: targeted proof check. Also checked in this writing. Pre-continuation writing events: GPT-6 Astra (OpenAI), ChatGPT Pro mode, effort Pro, S6\_10\_PAPERS, 18 November 2026: targeted proof check. Continuation metadata events: GPT-6 Astra (OpenAI), ChatGPT Pro mode, effort Pro, S6\_10\_PAPERS continuation, 19 November 2026: used and rechecked in this run; status-vocabulary reconciliation; metadata reconciliation only; no new mathematical check. \emph{Pre-continuation recorded proof scope.} GPT-6 Astra (OpenAI, ChatGPT Pro mode; recorded reviewer /root/cdp\_audit), effort Pro, 17 November 2026; Claude Fable 5.2 (Anthropic), effort high, 16 November 2026; also checked in this writing by GPT-6 Astra (OpenAI), ChatGPT Pro mode, effort Pro, 18 November 2026}{the hypotheses, notation and definitions stated here; \cref{maps:src:E.1.2}; \cref{maps:eq:E.14}; \cref{maps:eq:E.13}; \cref{maps:eq:E.1}; \cite[Theorem 1.1 and Remark 1.2]{ffJiaMeng}; \cite[(1.2)--(1.4)]{ffEsoleKang}; \cref{maps:eq:E.8}; \cite[Theorem 2.2.1]{ffPetersSterk}}


For the fourfold of Theorem \ref{maps:src:E.1.2},
\[
 K_Z=\Pi^*L^{-5},\qquad a(Z)=2,\qquad
 (h^{0,0},\ldots,h^{0,4})=(1,1,0,0,0).
 \tag{E.12}\label{maps:eq:E.12}
\]
It is not in Fujiki's class $\mathcal C$.  Moreover
$c_1(Z)=0$ integrally, $\int_Zc_2^2=2$, $\sigma(Z)=2$, and
the intersection form on $H^4(Z;\mathbb Z)/\mathrm{torsion}$ is
$\langle1\rangle\oplus\langle1\rangle$.

\end{theorem}


\begin{proof}
\begin{enumerate}[label=\arabic*.,leftmargin=*]
\item We first check the structure-sheaf input from the supplied family
calculation.  The orbifold base has roots of orders three and four.
Write its line generator as $\mathcal O(1)$; it pulls back to $L$.
For the coarse map $\rho$,
\[
 \rho_*\mathcal O(k)=\mathcal O_B(d(k)),\qquad
 d(k)=\frac{k}{12}-\frac{k\bmod3}{3}
                    -\frac{-k\bmod4}{4}.
 \tag{E.13}\label{maps:eq:E.13}
\]
This follows by dividing a local invariant section by its necessary
root-coordinate power.  Finite-group invariants over $\mathbb C$
are exact.  In particular $\mathcal O(-1)$ and $\mathcal O(-2)$
are acyclic, since their coarse degree is $-1$.

\item For $\bar f:X\to\mathfrak B$, the actual direct images are
\[
 R^0\bar f_*\mathcal O_X=\mathcal O,\quad
 R^1\bar f_*\mathcal O_X=\mathcal O\oplus\mathcal O(-1),
 \quad R^2\bar f_*\mathcal O_X=\mathcal O(-1).
 \tag{E.14}\label{maps:eq:E.14}
\]
Here is the step of the supplied proof that verifies \eqref{maps:eq:E.14}.
The cusp branch resolution of $\mathcal O_W$ is the cellular
complex of its dual two-torus, and has cohomology dimensions
$(1,2,1)$ with nonzero exterior cup pairing.  Proper flat base
change therefore gives vector bundles with those ranks.  The
relative logarithmic vertical bundle has a fixed line generated
by $\widehat\delta$, and its determinant has twelfth power
$\mathcal O(-p_0)$, by the modular discriminant.  Thus the
determinant is $\mathcal O(-1)$.  Relative duality gives the last
term of \eqref{maps:eq:E.14}.  The invariant covector $\gamma$ gives a nowhere
zero section of $R^1\bar f_*\mathcal O_X$: on a torus its class
in $\Lambda^*_{\mathbb C}/\operatorname{rowspan}\Pi_c$ is
nonzero by the last two columns of $\Pi_c$; at the cusp it is
the nonzero deck cocycle on the dual two-torus.  Its determinant
and the nonzero cusp cup pairing then give the extension
$0\to\mathcal O\to R^1\bar f_*\mathcal O_X
\to\mathcal O(-1)\to0$.  It splits because
$H^1(\mathfrak B,\mathcal O(1))=0$, also from \eqref{maps:eq:E.13}.
The same logarithmic determinant gives $K_X=L^{-6}$.
These checks use the actual root atlases and reduced cusp; they
do not replace the multiple fibres by an unramified base disc.

\item The Weierstrass adjunction calculation gives
$K_{\mathscr W}=\Pi^*(K_X\otimes L)$, and the small toric
resolution is crepant.  This proves the canonical formula in
\eqref{maps:eq:E.12}.  Every fibre of $\Pi$ is a connected genus-one curve;
over the deeper strata its toric resolution is a cycle of rational
curves.  Thus $\Pi$ is flat, by equidimensionality between smooth
spaces, and relative duality gives
$R^1\Pi_*\mathcal O_Z=L^{-1}$ and
$R^0\Pi_*\mathcal O_Z=\mathcal O_X$.
Restriction to the section splits the unit map in the derived
direct image, so
\[
 R\Pi_*\mathcal O_Z\simeq
 \mathcal O_X\oplus L^{-1}[-1].
\]
Projection formula, \eqref{maps:eq:E.13}, and \eqref{maps:eq:E.14} show that
$H^*(X,L^{-1})=0$ and $h^*(X,\mathcal O_X)=(1,1,0,0)$.
This proves the asserted Hodge numbers.

\item A meromorphic trivialization of $L$ makes the coordinate $x$ in
\eqref{maps:eq:E.1} a meromorphic function over $X$, of degree two on each
elliptic fibre.  It gives a generically finite meromorphic map to
$X\times\mathbb P^1$.  Traces and elementary symmetric functions
over its finite generic fibre show that the meromorphic function
field is an algebraic extension of that of the target.  Since
$a(X)=1$, this proves $a(Z)=2$.  If $Z$ were in class
$\mathcal C$, a sequence of blowups along smooth centres would
give a compact Kahler model, by the form of the model theorem
in \cite[Theorem 1.1 and Remark 1.2]{ffJiaMeng}.
Both $b_1$ and $h^{0,1}$ are unchanged by such blowups: the
exceptional fibres have no first homology and no positive-degree
structure-sheaf cohomology.  The Kahler identity
$b_1=2h^{0,1}$ would give $0=2$, a contradiction.

\item Since $L$ is pulled from a line on $S^6$, its integral first Chern
class vanishes, so $c_1(Z)=0$.  Riemann--Roch and the signature
formula, in the explicit normalization of
\cite[(1.2)--(1.4)]{ffEsoleKang}, yield
\[
 0=\chi(\mathcal O_Z)=\frac{3c_2^2-c_4}{720},\qquad
 \sigma(Z)=\frac{3c_2^2+14c_4}{45}.
\]
Here $c_4[Z]=6$, hence $c_2^2[Z]=2$ and $\sigma(Z)=2$.
The integral middle pairing modulo torsion is unimodular by
Poincare duality, has rank two by \eqref{maps:eq:E.8}, and is positive definite
by its signature.  The rank-at-most-four lattice classification
\cite[Theorem 2.2.1]{ffPetersSterk} therefore gives precisely
$\langle1\rangle\oplus\langle1\rangle$.
\end{enumerate}
\end{proof}

The complete credited statement and re-derivation are in \cref{maps:src:E.2.2}.
\begin{theorem}[a shifted second sphere]\label{maps:src:E.2.3}
\StatementRecord{maps:src:E.2.3}{proved}{Historical events: Claude Fable 5.2 (Anthropic), effort high, 16 November 2026: carried as background without fault found. GPT-6 Astra (OpenAI), ChatGPT Pro mode, effort Pro, S6\_09\_PAPERS, 17 November 2026: targeted proof check. Also checked in this writing. Pre-continuation writing events: GPT-6 Astra (OpenAI), ChatGPT Pro mode, effort Pro, S6\_10\_PAPERS, 18 November 2026: targeted proof check. Continuation metadata events: GPT-6 Astra (OpenAI), ChatGPT Pro mode, effort Pro, S6\_10\_PAPERS continuation, 19 November 2026: used and rechecked in this run; status-vocabulary reconciliation; metadata reconciliation only; no new mathematical check. \emph{Pre-continuation recorded proof scope.} GPT-6 Astra (OpenAI, ChatGPT Pro mode; recorded reviewer /root/cdp\_audit), effort Pro, 17 November 2026; Claude Fable 5.2 (Anthropic), effort high, 16 November 2026; also checked in this writing by GPT-6 Astra (OpenAI), ChatGPT Pro mode, effort Pro, 18 November 2026}{the hypotheses, notation and definitions stated here; \cref{maps:src:E.2.2}; \cref{maps:eq:E.9}}


Let $X'$ be a second canonical member and choose $\varphi\in
\operatorname{PGL}_2(\mathbb C)$ so that its three exceptional
values are disjoint from those of $X$.  The smooth fivefold
$P=X\times_B\varphi^*X'$ has Euler characteristic zero and
an infinite fundamental group with presentation
\[
\begin{split}
G=\langle a_1,a_2,b_1,b_2,h_1,h_2\mid\ &
 h_1,h_2\ \text{central},\quad
 a_1^3=h_1,\ a_2^4=h_1^{-1},\\
 &b_1^3=h_2,\ b_2^4=h_2^{-1},\quad
 a_1a_2b_1b_2=1\rangle .
\end{split}
\tag{E.16}\label{maps:eq:E.16}
\]
Its abelianization is $\mathbb Z$, and its rational Betti numbers are
\[
 (1,1,1,20,63,88,63,20,1,1,1).
 \tag{E.17}\label{maps:eq:E.17}
\]
For fixed $X,X'$, all admissible $\varphi$ give the same oriented
diffeomorphism type.

\end{theorem}


\begin{proof}
\begin{enumerate}[label=\arabic*.,leftmargin=*]
\item Theorem \ref{maps:src:E.2.2} proves smoothness.  At each special value the other
factor is a smooth four-torus.  The Euler characteristic of every
fibre is therefore zero, and Euler additivity gives $\chi(P)=0$.
We give van Kampen with the actual cusp relations, rather than
deducing simple connectivity from that of the factors.
Over the six-punctured sphere the fibre group is
$\Lambda_X\oplus\Lambda_{X'}$ and the base is free on five
peripheral loops.  Each monodromy acts on its own direct summand.
Filling the first cusp kills its angular generators
$\widehat w,\widehat\delta$ and its cusp meridian.  Their normal
closure under $A_1$ also kills $\widehat u$, since
$A_1\widehat w=\widehat u-\widehat w$.  The same calculation
in the other summand leaves only $h_1=\widehat\gamma$ and
$h_2=\widehat\gamma'$.  The finite monodromies now act trivially
on them.  The four finite filling relations are precisely those
in \eqref{maps:eq:E.16}, since $\gamma(v_1)=1$ and $\gamma(v_2)=-1$.
Both cusp collars are untwisted; their meridians are killed.
The product of the remaining four meridians gives the last
relation.  These are all relations introduced by the six
fillings, so van Kampen gives exactly \eqref{maps:eq:E.16}.

\item Quotienting by $h_1,h_2$ gives the orbifold group of the sphere
with cone orders $3,4,3,4$.  Its orbifold Euler characteristic
is $-5/6$, so its universal orbifold cover is the disc and the
group is infinite.  Equivalently one can build the Seifert
$T^2$-bundle with these four primitive translation vectors; over
the disc it is a trivial torus bundle.  Its exact sequence shows
that the central subgroup in \eqref{maps:eq:E.16} is $\mathbb Z^2$.
The abelianized five-by-six relation matrix has Smith factors
$(1,1,1,1,1)$, as printed in the code, proving $G_{\rm ab}=\mathbb Z$.

\item For the rational homology there is no extra cusp kernel: each
cusp is the original rank-two cusp times a smooth torus, and its
specialization is an isomorphism onto invariants over $\mathbb Q$.
The two monodromy groups act on independent summands.  In exterior
degrees $q=0,\ldots,8$, the $H^0$ and $H^2$ rows are both
\[
 (1,2,3,4,5,4,3,2,1),
\]
and the $H^1$ row is
\[
 (0,0,18,60,86,60,18,0,0).
\]
To verify this without a matrix-size guess, the local invariant
polynomials of the first factor are
$A(z)=1+2z+2z^2+2z^3+z^4$ at a finite point and
$C(z)=A(z)+2z^2$ at its cusp.  Hence the Euler polynomial of
the direct-image rows is
\[
 -4(1+z)^8+(4A(z)+2C(z))(1+z)^4.
\]
Formula \eqref{maps:eq:E.9}, with six point strata, gives the displayed rows.

\item Each invariant algebra of a factor has one generator in degrees
zero through four.  Its invariant-to-coinvariant map is nonzero
in even degrees and zero in odd degrees.  In the marking used
above, this follows from the invariant form
$6\gamma\wedge\delta+u\wedge w$ and its exterior products:
the odd invariant classes pair to zero with every invariant
complementary class, and the even ones have nonzero square or
top pairing.  Its four transgressions are all nonzero, as the
factor is a sphere.  The derivation rule therefore sends a tensor
of bidegree $(i,j)$ into bidegrees $(i-1,j)$ and $(i,j-1)$,
with the first term nonzero exactly when $i>0$ and $j$ is even,
and the second exactly when $j>0$ and $i$ is even.  The nonzero
column supports are disjoint except for pairs in a single row,
so their scalar normalizations do not affect ranks.  This gives
\[
 (r_1,\ldots,r_8)=(1,2,2,3,3,2,2,1).
\]
The exact program forms these matrices and takes their ranks.
There are no later differentials.  Their diagonal sums give
\eqref{maps:eq:E.17}, which also satisfies Poincare duality and Euler zero.

\item Finally, the admissible set of $\varphi$ is the complement of
finitely many proper complex hypersurfaces in the connected
complex manifold $\operatorname{PGL}_2(\mathbb C)$, so it is
path connected.  The universal pullback over this set is a
proper smooth submersion: transversality was checked pointwise
above.  Parallel transport along a path gives an oriented
diffeomorphism of its fibres.
\end{enumerate}
\end{proof}

\begin{proposition}[the shifted product's smooth type]\label{maps:src:E.2.4}
\StatementRecord{maps:src:E.2.4}{proved}{Historical events: Claude Fable 5.2 (Anthropic), effort high, 16 November 2026: carried as background without fault found. GPT-6 Astra (OpenAI), ChatGPT Pro mode, effort Pro, S6\_09\_PAPERS, 17 November 2026: targeted proof check. Also checked in this writing. Pre-continuation writing events: GPT-6 Astra (OpenAI), ChatGPT Pro mode, effort Pro, S6\_10\_PAPERS, 18 November 2026: targeted proof check. Continuation metadata events: GPT-6 Astra (OpenAI), ChatGPT Pro mode, effort Pro, S6\_10\_PAPERS continuation, 19 November 2026: used and rechecked in this run; status-vocabulary reconciliation; metadata reconciliation only; no new mathematical check. \emph{Pre-continuation recorded proof scope.} GPT-6 Astra (OpenAI, ChatGPT Pro mode; recorded reviewer /root/cdp\_audit), effort Pro, 17 November 2026; Claude Fable 5.2 (Anthropic), effort high, 16 November 2026; also checked in this writing by GPT-6 Astra (OpenAI), ChatGPT Pro mode, effort Pro, 18 November 2026}{the hypotheses, notation and definitions stated here; \cref{maps:eq:E.16}; \cref{maps:src:E.2.21}; \cref{maps:src:E.2.1}; \cref{maps:src:E.2.5}; \cref{maps:src:E.2.18}}


The shifted product in \eqref{maps:eq:E.16} is the closed oriented ten-manifold
$\mathcal H_{10}$ specified by the $2432$ ordinary framed handles
of \ref{maps:src:E.2.21}. Its regular-fibre complement is, relative to its entire
marked boundary, the $76$-handle six-manifold $\mathcal H_6$ of
that theorem. The full attaching embeddings, their ordered normal
coordinates, and every dual belt embedding are prescribed by
finite coordinate formulas. The identification follows by an
actual relative diffeomorphism of the cusp, core and finite caps,
and then by the handle-by-handle comparison of \ref{maps:src:E.2.21} using
Milnor's Theorems~1.4, 3.13 and~5.8. Top attaching maps are
retained, so there is no uncomputed sphere-capping ambiguity.

Together with \ref{maps:src:E.2.1}--\ref{maps:src:E.2.3}, \ref{maps:src:E.2.5}--\ref{maps:src:E.2.20}, this supplies the complex
invariants, the general smooth compact pullback construction,
the product case $S^6\times N$, and the shifted product's
fundamental group and full integral homology. The smooth type
is given by the explicit framed-handle presentation, not by an
inference from its homology or by an asserted connected-sum
normal form. \ref{maps:src:E.2.18}'s actual quadratic obstruction is consistent
with this presentation. The former open paragraph is retained
in the research records.


\end{proposition}
\begin{proof}
\begin{enumerate}[label=\arabic*.,leftmargin=*]
\item The relative diffeomorphism of \cref{maps:src:E.2.20} identifies the cusp, flat core and both finite caps while fixing the full outer torus marking. Apply the ordered block expansion of \cref{maps:src:E.2.21} to this comparison: each full attaching embedding, including its normal coordinates, is carried to the displayed coordinate embedding. The compactly supported collar isotopies in that proof extend the comparison across each successive handle, with the protected boundary fixed. Taking its product with the two regular torus factors and gluing by the explicit normal-angle reversal gives the ten-dimensional comparison. The stated counts are the two expanded handle lists of \cref{maps:src:E.2.21}; every top handle uses its prescribed attaching map. The remaining assertions are exactly the already established results \cref{maps:src:E.2.1,maps:src:E.2.2,maps:src:E.2.3,maps:src:E.2.5,maps:src:E.2.18}.
\end{enumerate}
\end{proof}



\begin{proposition}[the first integral torsion of the shifted product]\label{maps:src:E.2.5}
\StatementRecord{maps:src:E.2.5}{proved}{Historical events: Claude Fable 5.2 (Anthropic), effort high, 16 November 2026: carried as background without fault found. GPT-6 Astra (OpenAI), ChatGPT Pro mode, effort Pro, S6\_09\_PAPERS, 17 November 2026: targeted proof check. Also checked in this writing. Pre-continuation writing events: GPT-6 Astra (OpenAI), ChatGPT Pro mode, effort Pro, S6\_10\_PAPERS, 18 November 2026: targeted proof check. Continuation metadata events: GPT-6 Astra (OpenAI), ChatGPT Pro mode, effort Pro, S6\_10\_PAPERS continuation, 19 November 2026: used and rechecked in this run; status-vocabulary reconciliation; metadata reconciliation only; no new mathematical check. \emph{Pre-continuation recorded proof scope.} GPT-6 Astra (OpenAI, ChatGPT Pro mode; recorded reviewer /root/cdp\_audit), effort Pro, 17 November 2026; Claude Fable 5.2 (Anthropic), effort high, 16 November 2026; also checked in this writing by GPT-6 Astra (OpenAI), ChatGPT Pro mode, effort Pro, 18 November 2026}{the hypotheses, notation and definitions stated here; \cref{maps:src:E.2.3}; \cref{maps:src:E.1.3}; \cref{maps:src:D.1.2}; \cref{maps:eq:E.I1}; \cref{maps:src:E.1.5}; \cref{maps:eq:E.20}}


Let \(Z_\varphi=X\times_B\varphi^*X\) be the transverse
fivefold in \ref{maps:src:E.2.3}, with the two sets of three critical values
disjoint. Its second integral homology is
\[
 H_2(Z_\varphi;\mathbb Z)
   =\mathbb Z\oplus(\mathbb Z/12)^7.
\]
In particular the rational Betti vector in \ref{maps:src:E.2.3} cannot be
interpreted as integral torsion-free homology.
\end{proposition}


\begin{proof}
\begin{enumerate}[label=\arabic*.,leftmargin=*]
\item Use subscripts one and two for the two rank-four cohomology
local systems. At a critical value of one factor the other
factor is a smooth torus bundle over a disc, hence a product
there. Each local specialization is therefore the tensor
product of the original specialization with the identity on
the exterior algebra of \(\mathbb Z^4\). The original finite
reduced surfaces and the original cusp have free cohomology.
Their specialization maps are injective; at the cusp they
are integral isomorphisms onto invariants. For the latter,
apply the support-function construction of \ref{maps:src:E.1.3} in two
coordinates. It supplies the invariant deck classes
\(b_1,b_2\), the three classes
\(a_1b_1,a_2b_2,a_1b_2+a_2b_1\), and their products.
Together with \(b_1b_2\) these give an integral kernel
basis in degree two; their products give integral bases
in degrees three and four. Both the central-fibre and
invariant lattices have ranks \((1,2,4,2,1)\), as follows
also from the normalization matrix in \ref{maps:src:D.1.2} at \(d=1\).
Thus the specialization is onto between free groups of
equal ranks in every degree, proving the claimed isomorphism.
Thus the only difference between \(R^qg_*\mathbb Z\) and
\(\mathcal J_q=j_*\bigwedge^q(\mathcal V_1\oplus\mathcal V_2)\)
is their finite covering cokernel, with no point-supported kernel.

\item The finite Wang calculation \eqref{maps:eq:E.I1}, now with the original
three-dimensional vertical lattice, gives
\[
 C_3^1=\mathbb Z/3,\quad C_4^1=\mathbb Z/4,
 \qquad C_3^2=0,\quad C_4^2=\mathbb Z/2.
\]
The invariant covector \(\gamma\) generates the first two
cokernels, and the invariant two-form \(q_X\) generates the
last, by the same even-intersection argument as in \ref{maps:src:E.1.5}.
Consequently the degree-two covering cokernel at all four
finite values, combined, is
\[
 C^2=\left((\mathbb Z/3)^4\oplus
                 (\mathbb Z/4)^4\oplus\mathbb Z/2\right)^2.
\]
The two monodromy groups act independently on the two factors.
The three global invariant two-forms have integral basis
\(q_1,\gamma_1\gamma_2,q_2\). The first and third map
onto the respective \(\mathbb Z/2\) summands and vanish
in the other summands. The middle form maps to the primitive
\(\gamma\) coordinate of each of the four order-three or
order-four blocks. After the two order-two summands have
been removed, it therefore generates a direct cyclic summand
of order twelve in \((\mathbb Z/12)^8\). Hence
\[
 \operatorname{coker}\{H^0(B,\mathcal J_2)\to C^2\}
      =(\mathbb Z/12)^7.
 \tag{E.20}\label{maps:eq:E.20}
\]
The program \texttt{fourfolds\_shifted\_integral.py} realizes
this as an eighteen-row integer presentation; its Smith form
has eleven unit factors and seven factors twelve.

\item In degree two the global coboundary image is primitive.
Thus the parabolic group, a subgroup of the free-group
cohomology on the punctured base, is torsion free. Its rank
is eighteen: the six local invariant ranks are
\((16,16,18,16,16,18)\), and its global invariant and
coinvariant ranks are both three. The Euler formula on the
six-punctured sphere gives \(3+3+4\cdot28-100=18\).
The primitivity and all these ranks are checked by the exact
exterior matrices in the same program. The sheaf sequence
\(0\to R^2g_*\mathbb Z\to\mathcal J_2\to C^2\to0\)
and \eqref{maps:eq:E.20} now give
\[
 H^1(B,R^2g_*\mathbb Z)
     =\mathbb Z^{18}\oplus(\mathbb Z/12)^7.
\]
Here the quotient by its finite subgroup is free, so this
abelian-group extension splits.

\item The group \(H^2(B,R^1g_*\mathbb Z)\) is the free
rank-two group of coinvariants, one primitive summand from
each factor. Pulling back the integral degree-two
transgression of the original sphere fibration through each
projection kills these two summands separately and
surjectively. Thus \(E_\infty^{2,1}=0\) in the Leray
sequence of \(g\). Its other middle group
\(E_\infty^{0,3}\) is free, since \(R^3g_*\mathbb Z\)
injects into \(\mathcal J_3\). There are no higher
differentials. Consequently
\(\operatorname{Tor}H^3(Z_\varphi;\mathbb Z)
=(\mathbb Z/12)^7\). The universal coefficient theorem
identifies this with \(\operatorname{Tor}H_2(Z_\varphi;\mathbb Z)\).
Finally \ref{maps:src:E.2.3} gives \(b_2=1\), proving the assertion.
This argument does not determine the extensions in the
higher integral Leray filtrations.
\end{enumerate}
\end{proof}


\begin{theorem}[full integral homology of the shifted product]\label{maps:src:E.2.6}
\StatementRecord{maps:src:E.2.6}{proved}{Historical events: Claude Fable 5.2 (Anthropic), effort high, 16 November 2026: carried as background without fault found. GPT-6 Astra (OpenAI), ChatGPT Pro mode, effort Pro, S6\_09\_PAPERS, 17 November 2026: targeted proof check. Also checked in this writing. Pre-continuation writing events: GPT-6 Astra (OpenAI), ChatGPT Pro mode, effort Pro, S6\_10\_PAPERS, 18 November 2026: targeted proof check. Continuation metadata events: GPT-6 Astra (OpenAI), ChatGPT Pro mode, effort Pro, S6\_10\_PAPERS continuation, 19 November 2026: used and rechecked in this run; status-vocabulary reconciliation; metadata reconciliation only; no new mathematical check. \emph{Pre-continuation recorded proof scope.} GPT-6 Astra (OpenAI, ChatGPT Pro mode; recorded reviewer /root/cdp\_audit), effort Pro, 17 November 2026; Claude Fable 5.2 (Anthropic), effort high, 16 November 2026; also checked in this writing by GPT-6 Astra (OpenAI), ChatGPT Pro mode, effort Pro, 18 November 2026}{the hypotheses, notation and definitions stated here; \cref{maps:src:E.2.3}; \cref{maps:eq:E.MV1}; \cref{maps:src:E.2.5}; \cref{maps:eq:E.MV3}; \cref{maps:eq:E.MV2}; \cref{maps:eq:E.MV5}; \cref{maps:eq:E.MV6}; \cref{maps:eq:E.MV7}; \cref{maps:eq:E.MV8}}


For the transverse shifted product $P$ of \ref{maps:src:E.2.3}, write
\[
 A=(\mathbb Z/12)^7,\qquad
 B=(\mathbb Z/12)^{10}\oplus(\mathbb Z/2)^7,\qquad
 C=(\mathbb Z/12)^6\oplus\mathbb Z/24\oplus(\mathbb Z/2)^{17}.
\]
Its complete integral homology is
\[
 \begin{array}{c|l}
 i&H_i(P;\mathbb Z)\\\hline
 0,1,8,9,10&\mathbb Z\\
 2&\mathbb Z\oplus A\\
 3&\mathbb Z^{20}\oplus B\\
 4&\mathbb Z^{63}\oplus C\\
 5&\mathbb Z^{88}\oplus C\\
 6&\mathbb Z^{63}\oplus B\\
 7&\mathbb Z^{20}\oplus A.
 \end{array}
\]
In particular the higher integral filtration extensions are nontrivial:
the middle torsion includes a cyclic factor of order $24$.
\end{theorem}


\begin{proof}
\begin{enumerate}[label=\arabic*.,leftmargin=*]
\item \emph{A decomposition with free groups on every piece.}
Choose a separating circle in $B=\mathbb P^1$ so that one closed
disc contains the three critical values of $X$ and the other
contains the three critical values of $X'$. The smooth torus
family of the other factor is trivial over each disc. Set
\[
 F=T^4,\qquad E=X\setminus\operatorname{int}(F\times D^2),
 \qquad E'=X'\setminus\operatorname{int}(F'\times D^2).
\]
The product trivializations give the actual decomposition
\[
 P=(E\times F')\ \cup_{F\times F'\times S^1}\ (F\times E').
 \tag{E.MV1}\label{maps:eq:E.MV1}
\]
The boundary identification preserves the two marked torus factors
and reverses the normal circle, because the two base discs have
opposite boundary orientations. There is no additional torus
clutching: the same regular-disc trivializations define both the
boundary identifications of $E,E'$ and \eqref{maps:eq:E.MV1}.

\item Since $X,X'$ are six-spheres, the Mayer--Vietoris sequence of
$S^6=E\cup(F\times D^2)$ gives, integrally,
\[
 H_0(E)=\mathbb Z,\qquad H_k(E)=H_{k-1}(F)\quad(1\leq k\leq4),
 \qquad H_k(E)=0\quad(k\geq5).
 \tag{E.MV2}\label{maps:eq:E.MV2}
\]
The isomorphism is the normal-circle sweep
$\alpha\mapsto[S^1]\times\alpha$. Indeed, in degrees one
through four the boundary map to $H_k(E)\oplus H_k(F)$ is
an isomorphism, and its second component is the identity on
the unswept fibre summand and zero on the swept summand.
In degree five the boundary fundamental class is the image
of the sphere fundamental class, so $H_5(E)=0$.
Thus all homology groups of the two pieces and their common
boundary in \eqref{maps:eq:E.MV1} are free abelian. The same statements hold
for $E'$.

\item \emph{Determining every boundary map integrally.}
Let $V=H^1(F;\mathbb Z)$, in the basis $(\gamma,u,w,\delta)$,
and use the positively ordered volume
$v=\gamma u w\delta$. The primitive global invariant classes
of the original torus fibration are
\[
 I_0=1,\quad I_1=\gamma,\quad
 I_2=q=uw+6\gamma\delta,\quad
 I_3=\gamma uw,\quad I_4=v.
\]
The images of $H^k(E)\to H^k(F)$ are exactly
\[
 n_k\mathbb Z I_k,\qquad
 (n_1,n_2,n_3,n_4)=(12,2,2,12).
 \tag{E.MV3}\label{maps:eq:E.MV3}
\]
Here is an integral verification, rather than descent of invariant
real forms. On the disc the only local restrictions on a global
invariant class are the actual finite-cover pullback lattices;
the cusp specialization is an integral isomorphism onto its
invariant lattice by \ref{maps:src:E.2.5}. In degree one the finite images
contain $3\gamma$ and $4\gamma$, giving the intersection
$12\mathbb Z\gamma$. In degree two the order-three image
contains $q$, while the order-four image requires an even
coefficient of $q$, by the finite integral lattices of B.5.13.
In degree three use the perfect integral pairing downstairs:
an invariant class is a pullback precisely when its pairings
with every pulled-back degree-one class are divisible by the
covering degree. The other degree-one generators are
$2u+w+3\delta$ and $u+w+2\delta$. Their pairings with
$t\gamma uw$ are, up to signs, $3t$ and $2t$.
Thus the first cover imposes no further condition, and the
second requires $t$ even. In degree four the pullback degrees
are three and four, giving $12\mathbb Z v$.

\item These local restrictions are also sufficient for the image
from $E$. For a local system on a disc with punctures, the
restriction from its free-group $H^1$ to the direct sum of
the local coinvariants is onto: each puncture has its own
independent free generator. Consequently
$H^2(D,j_*\mathcal L)=0$. Replacing $j_*\mathcal L$ by the
actual direct-image sheaf changes it only by point-supported
cokernels, and so does not change this vanishing.
The Leray sequence over the disc therefore makes
$H^k(E)\to H^0(D,R^kf_*\mathbb Z)$ onto. This proves
\eqref{maps:eq:E.MV3} without an integral formality assumption.

\item Write $\theta$ for the normal-circle covector. The dual of
\eqref{maps:eq:E.MV2} identifies $H^k(E)$ with $\bigwedge^{k-1}V$, and
its boundary restriction has the graph form
\[
 r_k(\alpha)=\phi_k(\alpha)+\theta\alpha,
 \qquad
 \phi_k:\bigwedge^{k-1}V\longrightarrow\bigwedge^kV.
 \tag{E.MV4}\label{maps:eq:E.MV4}
\]
The image of $\phi_k$ is exactly $n_k\mathbb Z I_k$ by
\eqref{maps:eq:E.MV3}. If $k+l=5$, classes extending to $E$ have zero
product evaluated on $\partial E$. Therefore
\[
 (-1)^k\phi_k(\alpha)\beta+
                    \alpha\phi_l(\beta)=0
 \quad\text{in }\bigwedge^4V.
 \tag{E.MV5}\label{maps:eq:E.MV5}
\]
Each $I_k$ is primitive, and wedge pairing between the two
complementary exterior powers is unimodular. Writing
$\phi_k=f_k I_k$, equation \eqref{maps:eq:E.MV5}, tested against a class
pairing to one with $I_k$, proves that $f_k$ is an integral
multiple of $\alpha\mapsto\langle\alpha I_{5-k},[F]\rangle$.
The exact image \eqref{maps:eq:E.MV3} fixes the absolute multiple. Thus
\[
 \begin{split}
 \phi_k(\alpha)&=c_k
       \langle\alpha I_{5-k},[F]\rangle I_k,\\
 (c_1,c_2,c_3,c_4)&=(12\varepsilon,2\eta,-2\eta,-12\varepsilon),
 \qquad \varepsilon,\eta\in\{1,-1\}.
 \end{split}
 \tag{E.MV6}\label{maps:eq:E.MV6}
\]
This proves every boundary map up to two signs. We allow those
signs independently for $E$ and $E'$ below; the answer is the
same for all sixteen possibilities.

\item \emph{The full integral Mayer--Vietoris matrices.}
Let $i_p:H_p(F)\to H_p(E)$ be the transpose of $\phi_p$ in
the sweep bases, with $i_0=1$; let $j_{p+1}$ denote the sweep
isomorphism for $0\leq p\leq3$, and zero for $p=4$.
The analogous maps of the second copy are $i'_q,j'_{q+1}$.
For a fibre tensor $\alpha\otimes\beta$, with degrees $p,q$,
the matrix $D_k$ of \eqref{maps:eq:E.MV1} is
\[
 \alpha\otimes\beta\longmapsto
  (i_p\alpha\otimes\beta,-\alpha\otimes i'_q\beta),
 \qquad p+q=k.
 \tag{E.MV7}\label{maps:eq:E.MV7}
\]
For a swept tensor, place the normal circle first. Then
\[
 \theta\times\alpha\times\beta\longmapsto
 \bigl(j_{p+1}\alpha\otimes\beta,
                 (-1)^p\alpha\otimes j'_{q+1}\beta\bigr),
 \qquad p+q=k-1.
 \tag{E.MV8}\label{maps:eq:E.MV8}
\]
The second sign contains the Koszul sign from moving the circle
past $\alpha$, the reversal of the two boundary-circle
orientations, and the second minus in Mayer--Vietoris. The last
two cancel. This specifies every integer entry of $D_k$.

\item The executable \texttt{code/manifolds\_shifted\_mv.py} constructs
the exterior bases, the four adjoint maps \eqref{maps:eq:E.MV6}, and the full
matrices \eqref{maps:eq:E.MV7}--\eqref{maps:eq:E.MV8}. Exact Smith reduction gives the table
below. The source ranks are $\binom9k$; a zero binomial is
understood outside $0\leq k\leq9$.
\[
 \begin{array}{c|rrrr|l}
 k&\text{target rank}&\text{source rank}&\operatorname{rank}D_k&
       \text{unit factors}&\text{nonunit factors}\\\hline
 0&2&1&1&1&--\\
 1&10&9&9&9&--\\
 2&28&36&27&20&12^{[7]}\\
 3&64&84&53&36&2^{[7]},12^{[10]}\\
 4&114&126&82&58&2^{[17]},12^{[6]},24\\
 5&138&126&94&70&2^{[17]},12^{[6]},24\\
 6&104&84&73&56&2^{[7]},12^{[10]}\\
 7&44&36&35&28&12^{[7]}\\
 8&8&9&8&8&--\\
 9&0&1&0&0&--\\
 10&0&0&0&0&--
 \end{array}
\]
An exponent $[r]$ records repetition, not an integer power.
All sixteen choices of the two signs in each copy give precisely
this table; the retained output and JSON contain every case.

\item Finally the exact homology sequence gives
\[
 0\longrightarrow\operatorname{coker}D_k
 \longrightarrow H_k(P)
 \longrightarrow\ker D_{k-1}\longrightarrow0.
\]
The last group is free, since the boundary homology is free.
This sequence therefore splits as abelian groups, so the Smith
factors in the table are the actual torsion of $H_k(P)$.
Adding the free kernel ranks gives exactly the rational Betti
vector in \ref{maps:src:E.2.3} and the integral groups asserted here.
The resulting torsion also satisfies the required ten-dimensional
duality $\operatorname{Tor}H_k(P)\cong
\operatorname{Tor}H_{9-k}(P)$.
\end{enumerate}
\end{proof}

\begin{proposition}[a unique transverse family type and its stable tangent bundle]\label{maps:src:E.2.7}
\StatementRecord{maps:src:E.2.7}{proved}{Historical events: Claude Fable 5.2 (Anthropic), effort high, 16 November 2026: carried as background without fault found. GPT-6 Astra (OpenAI), ChatGPT Pro mode, effort Pro, S6\_09\_PAPERS, 17 November 2026: targeted proof check. Also checked in this writing. Pre-continuation writing events: GPT-6 Astra (OpenAI), ChatGPT Pro mode, effort Pro, S6\_10\_PAPERS, 18 November 2026: targeted proof check. Continuation metadata events: GPT-6 Astra (OpenAI), ChatGPT Pro mode, effort Pro, S6\_10\_PAPERS continuation, 19 November 2026: used and rechecked in this run; status-vocabulary reconciliation; metadata reconciliation only; no new mathematical check. \emph{Pre-continuation recorded proof scope.} GPT-6 Astra (OpenAI, ChatGPT Pro mode; recorded reviewer /root/cdp\_audit), effort Pro, 17 November 2026; Claude Fable 5.2 (Anthropic), effort high, 16 November 2026; also checked in this writing by GPT-6 Astra (OpenAI), ChatGPT Pro mode, effort Pro, 18 November 2026}{the hypotheses, notation and definitions stated here; \cref{maps:src:E.2.6}}


For fixed marked sphere fibrations, all M\"obius transformations
whose critical triples are disjoint give orientation-preservingly
diffeomorphic shifted products. The same conclusion holds while
either sphere varies in its proper smooth marked period family.
For every such fivefold $P$ there is an isomorphism of real vector
bundles
\[
 TP\oplus\mathbb R^4\simeq\mathbb R^{14}.
 \tag{E.MV9}\label{maps:eq:E.MV9}
\]
Thus all its positive-degree Stiefel--Whitney and Pontryagin
classes vanish. For its given complex tangent bundle,
$c_1(P)=c_2(P)=0$ integrally.
These assertions and \ref{maps:src:E.2.6} do not identify the complete smooth
structure from its ordinary homology and fundamental group.
\end{proposition}


\begin{proof}
\begin{enumerate}[label=\arabic*.,leftmargin=*]
\item The allowed transformations form the complement in
$\operatorname{PGL}_2(\mathbb C)$ of the nine analytic
hypersurfaces imposing an equality of one critical value from
each triple. This complement is path connected: a piecewise
smooth path between two points can be perturbed relative to
its endpoints to avoid these sets, which have real codimension
two. Equivalently, use local complex coordinate balls and
perturb each segment there. Choose a smooth such path.

\item The corresponding fibre products form a compact smooth family
over its parameter interval. Indeed the vertical derivative of
$f_1(x)-f_2(y)$ is surjective at every point, since its two
factors cannot be critical at the same base value. The same
surjectivity solves the equation for a lift of the interval
coordinate vector field. A partition of unity supplies a
global smooth lift; compactness makes its flow exist for the
whole interval. This flow gives the claimed diffeomorphisms.
The complex orientations vary continuously and hence are
preserved. The identical argument applies to smooth proper
period families of either factor, keeping the critical triples
disjoint throughout.

\item View $P$ as the transverse inverse image of the diagonal in
$B\times B$ under $f_1\times f_2$. Its real oriented normal
two-plane bundle is the underlying bundle of $g^*T_B$,
where $g:P\to B$ is the common base map. It is trivial:
its first Chern class is
\[
 c_1(g^*T_B)=2g^*[B]^\vee
    =2\operatorname{pr}_1^*f_1^*[B]^\vee=0,
\]
because $H^2(X;\mathbb Z)=0$, and complex line bundles are
classified by their first Chern class. The tangent bundles
of the two six-spheres become trivial after adding one real
line each. Restriction of their product tangent bundle gives
\[
 (TP\oplus\mathbb R^2)\oplus\mathbb R^2
       \simeq\mathbb R^{14},
\]
proving \eqref{maps:eq:E.MV9} and the characteristic-class assertions for
the real tangent bundle. Finally the holomorphic tangent
sequence is
\[
 0\longrightarrow TP\longrightarrow
 \operatorname{pr}_1^*TX\oplus\operatorname{pr}_2^*TX'
 \longrightarrow g^*T_B\longrightarrow0.
\]
Both sphere factors have integral $c_1=c_2=0$, since their
cohomology in degrees two and four vanishes. The last line
has zero first Chern class, as just proved. Multiplying total
Chern classes in this exact sequence gives
$c_1(P)=c_2(P)=0$ integrally.
\end{enumerate}
\end{proof}

\begin{proposition}[the shifted product bounds with its natural stable framing]\label{maps:src:E.2.8}
\StatementRecord{maps:src:E.2.8}{proved}{Historical events: Claude Fable 5.2 (Anthropic), effort high, 16 November 2026: carried as background without fault found. GPT-6 Astra (OpenAI), ChatGPT Pro mode, effort Pro, S6\_09\_PAPERS, 17 November 2026: targeted proof check. Also checked in this writing. Pre-continuation writing events: GPT-6 Astra (OpenAI), ChatGPT Pro mode, effort Pro, S6\_10\_PAPERS, 18 November 2026: targeted proof check. Continuation metadata events: GPT-6 Astra (OpenAI), ChatGPT Pro mode, effort Pro, S6\_10\_PAPERS continuation, 19 November 2026: used and rechecked in this run; status-vocabulary reconciliation; metadata reconciliation only; no new mathematical check. \emph{Pre-continuation recorded proof scope.} GPT-6 Astra (OpenAI, ChatGPT Pro mode; recorded reviewer /root/cdp\_audit), effort Pro, 17 November 2026; Claude Fable 5.2 (Anthropic), effort high, 16 November 2026; also checked in this writing by GPT-6 Astra (OpenAI), ChatGPT Pro mode, effort Pro, 18 November 2026}{the hypotheses, notation and definitions stated here; \cref{maps:eq:E.MV9}; \cref{maps:src:B.2.1}; \cite[Chapter 3, Section 2, Theorems 2.1 and 2.5]{rootHirsch76}; \cref{maps:eq:E.B2}; \cref{maps:src:E.2.7}; \cref{maps:src:E.2.9}}


Equip the two oriented six-spheres with their standard bounding
stable tangent framings. Trivialize the oriented two-plane bundle
$f_1^*T\mathbb P^1$ on the first sphere, and equip the transverse
shifted product $P$ with the induced stable framing in \eqref{maps:eq:E.MV9}.
Then there is a compact smooth oriented eleven-manifold $V$ with
\[
 \partial V=P,\qquad TV\oplus\mathbb R^3\simeq\mathbb R^{14},
 \qquad [P]=0\quad\hbox{in }\Omega^{\mathrm{fr}}_{10}.
 \tag{E.B1}\label{maps:eq:E.B1}
\]
The boundary framing is the stated one. In particular its induced
spin structure bounds. This supplies a framed bordism invariant;
it does not identify the smooth type from ordinary homology.
\end{proposition}


\begin{proof}
\begin{enumerate}[label=\arabic*.,leftmargin=*]
\item Write $f_2:S^6\to\mathbb P^1$ for the shifted map of the second
factor. By \ref{maps:src:B.2.1} the first map $f_1:S^6\to\mathbb P^1$ is
null-homotopic. It therefore extends to a smooth map
$F:D^7\to\mathbb P^1$, which can be chosen to equal $f_1$ on
each sphere of a collar of the boundary. To see this with the
boundary fixed, first use the null-homotopy to obtain a continuous
extension constant in the collar direction. Smooth its interior
in coordinate charts, with cutoffs supported away from the collar;
on the overlap the extension is already smooth.

\item Let $C_2$ be the three critical values of $f_2$. The restriction
$f_1$ is transverse to each point of $C_2$, since the two critical
triples are disjoint. Perturb $F$ only in the interior to be
transverse to all these points. The applicable finite simultaneous
transversality statement is
\cite[Chapter 3, Section 2, Theorems 2.1 and 2.5]{rootHirsch76}.
For the relative assertion, the proof of that theorem uses local
perturbations multiplied by cutoff functions. Their supports may
be restricted to the complement of the collar, where
transversality has not already been secured; the open condition
on the collar persists. Compactness permits finitely many such
perturbations. Thus no assumption that $F$ is a submersion
everywhere is needed.

\item Define the smooth fibre product
\[
 V=\{(x,y)\in D^7\times S^6:F(x)=f_2(y)\}.
 \tag{E.B2}\label{maps:eq:E.B2}
\]
At a point whose common value is outside $C_2$, the differential
$df_2$ is onto. At a point over $C_2$, the differential $dF$ is
onto. Consequently $dF-df_2$ is onto everywhere on \eqref{maps:eq:E.B2}.
The same argument holds on its boundary, using $df_1$ there.
The implicit-function theorem, including its boundary version,
makes $V$ a smooth eleven-manifold with boundary exactly $P$.
It is compact as a closed subset of $D^7\times S^6$.
The inverse-image orientation using the complex orientation of
$T\mathbb P^1$ induces the complex orientation of $P$ at the
boundary, because the collar keeps the first map fixed.

\item Let $g:V\to\mathbb P^1$ be the common map. The real normal
two-plane bundle of \eqref{maps:eq:E.B2} is $g^*T\mathbb P^1$, identified by
$dF-df_2$. This is the pullback from $D^7$ of $F^*T\mathbb P^1$,
which is trivial because the disc is contractible. A choice of
its oriented framing restricts on the boundary to the chosen
framing of $f_1^*T\mathbb P^1$ up to homotopy: two oriented
framings of this two-plane bundle differ by a map
$S^6\to SO(2)$, and $H^1(S^6;\mathbb Z)=0$.
The tangent exact sequence and the bounding stable framing of
the second sphere now give
\[
 \begin{split}
 TV\oplus\mathbb R^2&\simeq
       T(D^7\times S^6)|_V,\\
 TV\oplus\mathbb R^3&\simeq
       \mathbb R^7\oplus(TS^6\oplus\mathbb R)|_V
       \simeq\mathbb R^{14}.
 \end{split}
\]
At the boundary, the outward normal of $D^7$ supplies the
additional real line in $TV|_P=TP\oplus\mathbb R$.
The resulting framing is precisely the stable tangent framing
obtained from the two sphere framings and the normal framing
in \ref{maps:src:E.2.7}. This proves \eqref{maps:eq:E.B1} by the definition of framed bordism.
It also proves the assertion about the induced spin structure,
by forgetting the framing.

\item This disc-extension construction does not by itself extend a
classifying map $P\to B\pi_1(P)$ over $V$.
The different filling in \ref{maps:src:E.2.9} below does extend that map.
Both are bordism statements; the higher equivariant surgery
data required for the smooth classification remain separate.
\end{enumerate}
\end{proof}


\begin{theorem}[framed null-bordism over the fundamental group]\label{maps:src:E.2.9}
\StatementRecord{maps:src:E.2.9}{proved}{Historical events: Claude Fable 5.2 (Anthropic), effort high, 16 November 2026: carried as background without fault found. GPT-6 Astra (OpenAI), ChatGPT Pro mode, effort Pro, S6\_09\_PAPERS, 17 November 2026: targeted proof check. Also checked in this writing. Pre-continuation writing events: GPT-6 Astra (OpenAI), ChatGPT Pro mode, effort Pro, S6\_10\_PAPERS, 18 November 2026: targeted proof check. Continuation metadata events: GPT-6 Astra (OpenAI), ChatGPT Pro mode, effort Pro, S6\_10\_PAPERS continuation, 19 November 2026: used and rechecked in this run; status-vocabulary reconciliation; metadata reconciliation only; no new mathematical check. \emph{Pre-continuation recorded proof scope.} GPT-6 Astra (OpenAI, ChatGPT Pro mode; recorded reviewer /root/cdp\_audit), effort Pro, 17 November 2026; Claude Fable 5.2 (Anthropic), effort high, 16 November 2026; also checked in this writing by GPT-6 Astra (OpenAI), ChatGPT Pro mode, effort Pro, 18 November 2026}{the hypotheses, notation and definitions stated here; \cref{maps:eq:E.16}; \cref{maps:src:E.2.8}; \cref{maps:src:B.2.1}; \cref{maps:src:E.2.2}; \cite[Theorems 2.3--2.4 and the Euler-number discussion in
Section 2]{rootScott83}; \cref{maps:src:E.2.3}; \cite[Definition 8.16 and Theorem 8.17]{rootDavisKirk01}; \cref{maps:src:B.1.2}; \cite[Corollary~1.3 and Table~1, pp.~112--113]{IsaksenWangXu2023}; \cref{maps:eq:E.B7}}


Let $G$ be the group in \eqref{maps:eq:E.16}. There is a closed smooth
aspherical four-manifold $N$ with
\[
 \pi_1(N)=G,\qquad \widetilde N\simeq\mathbb H^2\times\mathbb R^2,
 \qquad H_*(N;\mathbb Z)\cong H_*(S^1\times S^3;\mathbb Z),
 \tag{E.B3}\label{maps:eq:E.B3}
\]
and a smooth classifying map $c:P\to N$ inducing the displayed
identification of fundamental groups. With the natural stable
framing of \ref{maps:src:E.2.8}, one has
\[
                 [P,c]=0\quad\hbox{in }\Omega^{\mathrm{fr}}_{10}(N).
 \tag{E.B4}\label{maps:eq:E.B4}
\]
Thus $P$ has a compact stably framed eleven-dimensional filling
over $BG$. This is stronger than the filling of \ref{maps:src:E.2.8}; it is
still a bordism statement, not a diffeomorphism classification.
\end{theorem}


\begin{proof}
\begin{enumerate}[label=\arabic*.,leftmargin=*]
\item \emph{The aspherical target.}
The repaired quotient maps of \cref{mapsI:src:A.2.1} are smooth maps
$e_i:S^6\to S^3$ over the base, and \ref{maps:src:B.2.1} gives the exact
factorizations $f_i=h_i e_i$. Here $h_i:S^3\to\mathbb P^1$
is the weighted Hopf map of orders $3,4$, with the second
map shifted by the chosen M\"obius transformation. Define
\[
 N=\{(u,v)\in S^3\times S^3:h_1(u)=h_2(v)\},\qquad
 c(x,y)=(e_1(x),e_2(y)).
 \tag{E.B5}\label{maps:eq:E.B5}
\]
The critical values of each $h_i$ are only its two exceptional
circle values. They are disjoint, so the same transversality
argument as in \ref{maps:src:E.2.2} makes $N$ a compact smooth four-manifold.
The two circle actions on the weighted Hopf spheres give a
principal $T^2$ bundle in the orbifold sense over
\[
                  \mathcal B=S^2(3,4,3,4).
\]
At a cone of order $m$, its uniformizing chart is
$(D^2\times T^2)/C_m$, with the generator rotating the disc
and translating one torus coordinate by a primitive fraction
of denominator $m$. This local description is smooth and is
the same finite-filling description as \cref{mapsI:src:A.1.3}, with the unused
circle retained.

\item The orbifold has Euler characteristic
$2-2(1-1/3)-2(1-1/4)=-5/6$. It is good and has hyperbolic
orbifold universal cover, by
\cite[Theorems 2.3--2.4 and the Euler-number discussion in
Section 2]{rootScott83}. Pulling the principal torus bundle to
$\mathbb H^2$ gives an ordinary principal $T^2$ bundle over a
contractible base, hence the trivial bundle. Its universal
cover is $\mathbb H^2\times\mathbb R^2$. The local chart above
shows that these orbifold pullbacks are ordinary covering maps
of total spaces. This proves asphericity and injectivity of
the fibre subgroup $\mathbb Z^2$.

\item Using the marked fibre generators $h_1,h_2$ and the four
positive peripheral generators gives exactly
\[
 \begin{split}
 \pi_1(N)=\langle a_1,a_2,b_1,b_2,h_1,h_2\mid {}&h_1,h_2
       \text{ central},\quad a_1^3=h_1,\quad a_2^4=h_1^{-1},\\
 &b_1^3=h_2,\quad b_2^4=h_2^{-1},\quad
 a_1a_2b_1b_2=1\rangle.
 \end{split}
\]
The map $c$ sends the surviving regular-fibre generators
$\widehat\gamma,\widehat\gamma'$ to $h_1,h_2$, respectively,
and each peripheral generator to the one with the same name.
This is because $e_i$ is the literal circle quotient on these
fibres and finite collars. Its cusp replacement extends over
a solid torus and kills precisely the cusp meridian in this
presentation. Comparing with the full van Kampen calculation
in \ref{maps:src:E.2.3} therefore shows that $c_*$ is the identity on the
displayed presentation, and in particular is an isomorphism.
As $N$ is aspherical, $c$ is a classifying map.

\item For completeness the abelian relation matrix has Smith
diagonal $(1,1,1,1,1)$; a primitive integer kernel vector in
the order $(a_1,a_2,b_1,b_2,h_1,h_2)$ is
$(4,-3,-4,3,12,-12)$. Thus $H_1(N)=\mathbb Z$.
Every stratum of its base has a torus fibre, so Euler
additivity gives $\chi(N)=0$. Integral Poincare duality and
the universal coefficient theorem give $H_3(N)=\mathbb Z$
and no torsion in $H_2(N)$. Euler characteristic then forces
$H_2(N)=0$, proving the integral homology assertion in \eqref{maps:eq:E.B3}.

\item \emph{The bordism computation before taking the inverse image.}
For a based sphere, the stable splitting
$(S^3)_+\simeq S^0\vee S^3$ identifies
\[
 \Omega^{\mathrm{fr}}_6(S^3)
       =\pi^{\mathrm S}_6\oplus\pi^{\mathrm S}_3.
 \tag{E.B6}\label{maps:eq:E.B6}
\]
We use Pontryagin--Thom in its form for maps to a space,
\cite[Definition 8.16 and Theorem 8.17]{rootDavisKirk01}.
The first component of $[S^6,e_i]$ is obtained by forgetting
the map, so it is zero for the bounding sphere framing.
The second component is the stable class of $e_i$, equivalently
the framed inverse image of a regular point. It is $\eta^3$
by \ref{maps:src:B.1.2}. One can also check the splitting directly: the
Pontryagin--Thom collapse of the framed sphere into its
disjoint-basepoint suspension has degree one on the sphere
summand; composing with $e_i$ gives its stable homotopy class.
Thus
\[
                         [S^6,e_i]=(0,\eta^3).
\]
Smashing the two sphere splittings gives
\[
 (S^3\times S^3)_+\simeq
       S^0\vee S^3\vee S^3\vee S^6.
\]
Under this splitting the external product
$[S^6\times S^6,e_1\times e_2]$ has four components.
The first three each contain a factor $[S^6]=0$; the last
is $\eta^3\eta^3=\eta^6=0$. Indeed $\eta^4$ belongs to
$\pi^{\mathrm S}_4=0$, as recorded in
\cite[Corollary~1.3 and Table~1, pp.~112--113]{IsaksenWangXu2023}. Therefore
\[
 [S^6\times S^6,e_1\times e_2]=0
      \quad\hbox{in }\Omega^{\mathrm{fr}}_{12}(S^3\times S^3).
 \tag{E.B7}\label{maps:eq:E.B7}
\]
This calculation uses the product of the specific bounding
sphere framings, not just vanishing of ordinary characteristic
classes.

\item \emph{The comparison with the requested classifying map.}
By \eqref{maps:eq:E.B7} there is a compact stably framed thirteen-manifold
$Z$ with boundary $S^6\times S^6$ and a map
$E:Z\to S^3\times S^3$ extending $e_1\times e_2$.
Smooth it relative to the boundary and make it transverse
to $N$ there. Transversality on the boundary already holds:
the derivative of its defining equation is $df_1-df_2$,
which is onto. Relative transversality therefore supplies
the interior perturbation without changing that boundary.

\item The normal bundle of $N\subset S^3\times S^3$ is the
restriction of $h_1^*T\mathbb P^1$ pulled back from the first
sphere. This oriented two-plane bundle is trivial because
$H^2(S^3;\mathbb Z)=0$. Choose a framing on it before
restriction to $N$. Its pullback under $e_1\times e_2$
is a framing of $f_1^*T\mathbb P^1$, agreeing with that
in \ref{maps:src:E.2.8} up to homotopy since $H^1(S^6;\mathbb Z)=0$.
It follows that
\[
                  V=E^{-1}(N),\qquad E|_V:V\longrightarrow N
\]
is a compact stably framed eleven-manifold with boundary
$(P,c)$ and the required boundary framing. This proves \eqref{maps:eq:E.B4}.

\item The exact group-matrix and Euler calculation, together with
the two equivalent stable-product reductions, is recorded in
\texttt{code/shifted\_group\_bordism.py}.
The earlier disc-extension filling need not extend the
classifying map; \eqref{maps:eq:E.B7} supplies a different filling that does.
Neither construction computes the higher equivariant surgery
data needed to distinguish smooth ten-manifolds over $N$.
\end{enumerate}
\end{proof}


\begin{proposition}[the fibre complement and a spinning obstruction]\label{maps:src:E.2.10}
\StatementRecord{maps:src:E.2.10}{proved}{Historical events: Claude Fable 5.2 (Anthropic), effort high, 16 November 2026: carried as background without fault found. GPT-6 Astra (OpenAI), ChatGPT Pro mode, effort Pro, S6\_09\_PAPERS, 17 November 2026: targeted proof check. Also checked in this writing. Pre-continuation writing events: GPT-6 Astra (OpenAI), ChatGPT Pro mode, effort Pro, S6\_10\_PAPERS, 18 November 2026: targeted proof check. Continuation metadata events: GPT-6 Astra (OpenAI), ChatGPT Pro mode, effort Pro, S6\_10\_PAPERS continuation, 19 November 2026: used and rechecked in this run; status-vocabulary reconciliation; metadata reconciliation only; no new mathematical check. \emph{Pre-continuation recorded proof scope.} GPT-6 Astra (OpenAI, ChatGPT Pro mode; recorded reviewer /root/cdp\_audit), effort Pro, 17 November 2026; Claude Fable 5.2 (Anthropic), effort high, 16 November 2026; also checked in this writing by GPT-6 Astra (OpenAI), ChatGPT Pro mode, effort Pro, 18 November 2026}{the hypotheses, notation and definitions stated here; \cref{maps:src:E.2.6}; \cref{maps:src:B.5.2}; \cref{maps:eq:E.MV4}; \cref{maps:eq:E.MV6}; \cite[Section 3, immediately before Theorem 3.4]{mapsWinter08}; \cite[Theorem 4.5]{mapsWinter08}; \cref{maps:eq:E.S5}; \cref{maps:eq:E.S4}}


Let $F=T^4$ be a regular fibre of the original sphere fibration and
$E=S^6\setminus\operatorname{int}(F\times D^2)$, with the marking of
\ref{maps:src:E.2.6}. Its group and peripheral homomorphism are
\[
 \begin{split}
 K=\pi_1(E)&=\langle a,b,h\mid h\text{ central},\ a^3=h,
                                      \ b^4=h^{-1}\rangle,\\
 g&\longmapsto h,\qquad u,w,d\longmapsto1,
       \qquad \mu\longmapsto ab.
 \end{split}
 \tag{E.S1}\label{maps:eq:E.S1}
\]
Here $\mu$ is the normal meridian, with the peripheral orientation
chosen so that $[a]=4[\mu]$, $[b]=-3[\mu]$ and $[h]=12[\mu]$.
The image of the boundary group is the embedded copy
$\langle h,\mu\rangle\simeq\mathbb Z^2$. In the unique homotopy
class of normal framings for which every pushed-off loop of $F$
has linking number zero, the image of $g$ is
\[
             \lambda=h\mu^{-12},
 \qquad \ker\{\pi_1(F)\longrightarrow K\}
                 =\mathbb Z\langle u,w,d\rangle.
 \tag{E.S2}\label{maps:eq:E.S2}
\]
In particular the peripheral structure is that of the $(3,4)$
torus knot, with its regular Seifert fibre rather than its preferred
longitude as the original $g$ coordinate. Nevertheless the embedded
four-torus $F\subset S^6$ is not the ordinary threefold closed spin
of that classical knot, even allowing ambient orientation reversal
and arbitrary changes of its normal framing. More strongly, it is
not a single ordinary closed spin of any embedded $T^3\subset S^5$.

\end{proposition}


\begin{proof}
\begin{enumerate}[label=\arabic*.,leftmargin=*]
\item Apply the van Kampen presentation in \ref{maps:src:B.5.2} with the regular base disc
removed. The cusp kills $w,d$ and its peripheral loop. Conjugating
$w$ by the order-three generator also kills $u$. The two finite caps
impose $a^3=h$ and $b^4=h^{-1}$; there is now no relation $ab=1$,
since the outer boundary has not been filled. The outer meridian
is $ab$ in the stated orientation. This proves \eqref{maps:eq:E.S1}, including
the actual peripheral map, rather than only an abstract group
isomorphism. Eliminating $h$ gives $\langle a,b\mid a^3b^4=1\rangle$.
Its quotient by $h$ is $C_3*C_4$. The word $ab$ has infinite order
in that free product, so a relation $h^r\mu^s=1$ forces $s=0$.
Abelianization then forces $r=0$. This proves peripheral injectivity.
It also shows that $\lambda$ has infinite order: its image in
$C_3*C_4$ is $(ab)^{-12}$.

\item Write $\theta\in H^1(\partial E;\mathbb Z)$ for the normal-circle
covector and retain $(\gamma,u,w,\delta)$ for $H^1(F;\mathbb Z)$.
The generator $\xi\in H^1(E;\mathbb Z)$ restricts to
$\theta+12\gamma$. Homotopy classes of normal framings differ by
$H^1(F;\mathbb Z)$. Consequently there is a unique framing for
which the associated section $s:F\to\partial E$ satisfies
$s^*\xi=0$: it has $s^*\theta=-12\gamma$. Its induced peripheral
map is exactly \eqref{maps:eq:E.S2}. Simultaneously reversing the meridian
conventions changes the displayed signs and preserves all the
following assertions.

\item For this section set
\[
 \mathcal R=\operatorname{im}\{H^2(E;\mathbb Z)
                         \longrightarrow H^2(F;\mathbb Z)\},
 \qquad
 \operatorname{rad}(\mathcal R)=
 \{v\in H_1(F;\mathbb Z):\iota_v\omega=0
                                    \text{ for all }\omega\in\mathcal R\}.
 \tag{E.S3}\label{maps:eq:E.S3}
\]
This uses the canonical identification of the cohomology of a torus
with the exterior algebra of its first cohomology. Equations
\eqref{maps:eq:E.MV4}--\eqref{maps:eq:E.MV6} give classes whose boundary restrictions are
$\theta u$ and $\theta w$: their unswept terms are zero because
$u\gamma uw=w\gamma uw=0$. Restricting them along $s$ proves
\[
                -12\gamma u,\ -12\gamma w\ \in\mathcal R.
 \tag{E.S4}\label{maps:eq:E.S4}
\]
The simultaneous contraction kernel of these two forms is
$\mathbb Z d$. Indeed their contractions force successively
$\gamma(v)=u(v)=w(v)=0$. Thus
\[
       \operatorname{rad}(\mathcal R)\subseteq\ker\gamma
        =\ker\{\pi_1(F)\longrightarrow K\}.
 \tag{E.S5}\label{maps:eq:E.S5}
\]
Both the zero-linking framing and this inclusion are invariants of
the embedded pair, with no choice of the original torus coordinates.

\item We define the comparison construction precisely. For a closed
submanifold $L\subset\operatorname{int}D^n$, its ordinary closed
spin is $L\times S^1$ in
$D^n\times S^1\subset\partial(D^n\times D^2)=S^{n+1}$.
Start with the closed classical knot $K_0\subset\operatorname{int}D^3$
and repeat this operation three times, each time placing the
previous compact pair in the interior of a disc. This is the
ordinary torus spinning construction reviewed in
\cite[Section 3, immediately before Theorem 3.4]{mapsWinter08};
no turning, twisting, or family of deformations is inserted.
If $\Sigma$ is a Seifert surface for $K_0$, the same three spins
give the embedded Seifert hypersurface
$\Sigma\times T^3$ with boundary $K_0\times T^3$.
Pushing the latter off using the normal of this hypersurface
gives its zero-linking framing.

\item Let $v$ be the original knot-circle class and let $v_1,v_2,v_3$
be the spinning circles. Each pushed-off two-cycle $v\wedge v_i$
bounds the corresponding copy of $\Sigma\times S^1$ in a parallel
copy of that hypersurface, disjoint from the knot. Hence every
restriction of a complement two-class vanishes on $v\wedge H_1(F)$:
$v$ belongs to $\operatorname{rad}(\mathcal R)$.
On the other hand its pushed-off loop is the original preferred
longitude and is nontrivial for $K_0=T(3,4)$, whereas the spinning
circles are trivial in the complement group. These peripheral
claims also follow directly at each spin: its complement is
\[
 ((D^n\setminus\operatorname{int}\nu L)\times S^1)
       \ \cup_{S^{n-1}\times S^1}\ (S^{n-1}\times D^2),
\]
so van Kampen kills precisely the new spinning circle and preserves
the old complement group. They agree with the primary peripheral
calculation \cite[Theorem 4.5]{mapsWinter08} in the first spin.
Thus the ordinary spun torus violates \eqref{maps:eq:E.S5}. An identification
of embedded pairs would preserve its two sides, giving a
contradiction.

\item For the stronger assertion, let $L=T^3\subset S^5$ be arbitrary
and suppose that its ordinary closed spin were our four-torus.
Write $E_L$ for its exterior. Deleting a ball disjoint from
$L$ changes cohomology only in degrees at least four. In the
displayed spun-complement decomposition, with $n=5$, the cap
and intersection are $S^4\times D^2$ and $S^4\times S^1$.
Mayer--Vietoris and the surjectivity onto the latter first
cohomology therefore give, in degrees two and three, precisely
the cohomology of $(E_L\setminus D^5)\times S^1$.
Let $s$ be the spinning-circle covector and $v_s$ its primitive
homology generator. Put $\mathcal R_L^j=\operatorname{im}
\{H^j(E_L)\to H^j(L)\}$, using the zero-linking framing.
One has $\mathcal R_L^1=0$. Also $\mathcal R_L^3=0$:
the pushed-off $L$ bounds a parallel Seifert hypersurface in
its exterior, so every top-degree restriction evaluates to
zero. The product restrictions consequently satisfy
\[
       \mathcal R^2=\mathcal R_L^2,\qquad
       \mathcal R^3=\mathcal R_L^2\wedge s.
 \tag{E.S12}\label{maps:eq:E.S12}
\]
In particular $v_s\in\operatorname{rad}(\mathcal R^2)$,
and contraction by $v_s$ is injective on $\mathcal R^3$.
For our actual four-torus, \eqref{maps:eq:E.S4} forces $v_s=\pm d$.
But \eqref{maps:eq:E.MV4}--\eqref{maps:eq:E.MV6}, now applied to $\alpha=\gamma\delta$
in degree three, show that $\mathcal R^3$ contains
$\pm2\gamma uw$. Its swept term vanishes in the zero-linking
framing, and this is a nonzero class annihilated by $d$.
This contradicts that injectivity. No ordinary closed spin
is possible. Turned or deform-spun constructions insert
additional data and are not excluded by this argument.
\end{enumerate}
\end{proof}

\begin{theorem}[exact second-homotopy coinvariants]\label{maps:src:E.2.11}
\StatementRecord{maps:src:E.2.11}{proved}{Historical events: Claude Fable 5.2 (Anthropic), effort high, 16 November 2026: carried as background without fault found. GPT-6 Astra (OpenAI), ChatGPT Pro mode, effort Pro, S6\_09\_PAPERS, 17 November 2026: targeted proof check. Also checked in this writing. Pre-continuation writing events: GPT-6 Astra (OpenAI), ChatGPT Pro mode, effort Pro, S6\_10\_PAPERS, 18 November 2026: targeted proof check. Continuation metadata events: GPT-6 Astra (OpenAI), ChatGPT Pro mode, effort Pro, S6\_10\_PAPERS continuation, 19 November 2026: used and rechecked in this run; status-vocabulary reconciliation; metadata reconciliation only; no new mathematical check. \emph{Pre-continuation recorded proof scope.} GPT-6 Astra (OpenAI, ChatGPT Pro mode; recorded reviewer /root/cdp\_audit), effort Pro, 17 November 2026; Claude Fable 5.2 (Anthropic), effort high, 16 November 2026; also checked in this writing by GPT-6 Astra (OpenAI), ChatGPT Pro mode, effort Pro, 18 November 2026}{the hypotheses, notation and definitions stated here; \cref{maps:src:E.2.9}; \cref{maps:src:E.2.6}; \cref{maps:eq:E.S7}; \cref{maps:eq:E.MV7}; \cref{maps:eq:E.MV8}; \cref{maps:eq:E.MV6}; \cref{maps:eq:E.MV4}; \cref{maps:eq:E.S9}; \cref{maps:eq:E.S10}; \cref{maps:eq:E.S11}}


For the classifying map $c:P\to N=K(G,1)$ of \ref{maps:src:E.2.9}, one has
\[
 c_*:H_3(P;\mathbb Z)\longrightarrow H_3(N;\mathbb Z)=\mathbb Z
       \quad\text{equal to zero},
 \qquad
 \bigl(\pi_2(P)\bigr)_G
       \simeq\mathbb Z^2\oplus(\mathbb Z/12)^6.
 \tag{E.S6}\label{maps:eq:E.S6}
\]
The pullback of a generator of $H^3(N;\mathbb Z)$ has exact order
twelve in $H^3(P;\mathbb Z)$. In particular $\pi_2(P)$ is not a
projective $\mathbb ZG$-module. The fivefold cannot be diffeomorphic,
or homotopy equivalent, to $N\times S^6$ connected-summed with any
finite collection of $S^r\times S^{10-r}$, $2\leq r\leq8$, and
any homotopy ten-sphere. This computes coinvariants and a definite
module obstruction; it does not assert a presentation or a
classification of the full $\mathbb ZG$-module.

\end{theorem}


\begin{proof}
\begin{enumerate}[label=\arabic*.,leftmargin=*]
\item Use the separating base circle in \ref{maps:src:E.2.6} also in the target of
\ref{maps:src:E.2.9}. There $M$ denotes the exterior of a regular orbit of the
weighted Hopf map, namely the torus-knot exterior of \cref{mapsI:src:A.1.3}. The
two decompositions, and the restriction of $c$, are
\[
 \begin{array}{ccc}
 (E\times F')\cup(F\times E')&\longrightarrow&
                  (M\times S^1)\cup(S^1\times M')=N,\\
 F\times F'\times S^1&\longrightarrow&S^1\times S^1\times S^1.
 \end{array}
 \tag{E.S7}\label{maps:eq:E.S7}
\]
On the common boundary the map retains exactly the covectors
$\gamma_1,\gamma_2,\theta$. It is the literal quotient there.
Each knot exterior has the integral homology of a circle, so the
third homology and cohomology of each target piece in \eqref{maps:eq:E.S7} vanish.

\item We first check the homological statement, retaining the full
integer boundary maps. In the degree-two source of \eqref{maps:eq:E.MV7}--\eqref{maps:eq:E.MV8}
write $t_1,t_2,t_{12}$ for the coefficients of
$\theta\times g_1$, $\theta\times g_2$, and $g_1\times g_2$.
The image of $i_2:H_2(F)\to H_2(E)$ lies entirely in the
$d$-sweep coordinate, by \eqref{maps:eq:E.MV6}. Consequently the $g_1$-sweep
coordinate in the first $H_2(E)$ summand of the target is just
$t_1$, and the corresponding second coordinate is just $t_2$.
Every vector in the kernel of the boundary matrix therefore has
$t_1=t_2=0$. Its first $H_1(E)\otimes\mathbb Zg_2$ coordinate
then forces $12t_{12}=0$, up to an immaterial orientation sign.
These three coordinates are the only ones retained in
$H_2(T^9)\to H_2(T^3)$. The target Mayer--Vietoris sequence
injects $H_3(N)$ into $H_2(T^3)$; its source-piece third groups
are zero. Naturality now proves $c_*H_3(P)=0$.

\item The cohomological calculation records the extension which this
vanishing alone would miss. On the target boundary write
$v=\gamma_1\gamma_2$. The generators of the two degree-two
piece cohomology groups restrict, with consistent meridian
orientations, to
\[
       \theta\gamma_2+12v,
       \qquad \theta\gamma_1+12v.
 \tag{E.S8}\label{maps:eq:E.S8}
\]
Their image is a primitive rank-two subgroup of $H^2(T^3)$:
the minor in the last two coordinates is a unit. Its quotient
is freely generated by $v$. Since the degree-three piece groups
are zero, the connecting class $\partial v$ is a generator of
$H^3(N;\mathbb Z)$.

\item In the source decomposition both $\theta\gamma_1$ and
$\theta\gamma_2$ already belong to the image of the degree-two
piece restrictions. This is \eqref{maps:eq:E.MV4} with $\alpha=\gamma$,
since $\phi_2(\gamma)=0$; the second normal circle is reversed,
which only changes its sign. Subtracting one of these classes
from the corresponding product of a degree-one class of $E$
with $\gamma_2$ proves that $12v$ is a source boundary image.
Conversely the coefficient of $v$ in every such boundary image
is divisible by twelve. The pure degree-two terms from a single
factor have no $v$ coefficient, and the mixed terms use the
restriction $\theta+12\varepsilon\gamma$ of the degree-one
generator, with $\varepsilon=\pm1$. It follows that
\[
                  c^*(\partial v)=\partial v
        \quad\text{has exact order }12.
 \tag{E.S9}\label{maps:eq:E.S9}
\]
This uses the actual Mayer--Vietoris image; no splitting of a
torsion filtration has been assumed.

\item Put $A=(\mathbb Z/12)^7$ and $L=(\pi_2(P))_G$.
The low-degree homology sequence of the universal-cover
fibration, together with \ref{maps:src:E.2.6} and \ref{maps:src:E.2.9}, gives
\[
             0\longrightarrow\mathbb Z z\longrightarrow L
                \longrightarrow\mathbb Z\oplus A\longrightarrow0.
 \tag{E.S10}\label{maps:eq:E.S10}
\]
The zero on the left uses $c_*H_3(P)=0$, just proved.
To determine this extension, replace $c$ by its mapping cylinder
and use relative groups for $(N,P)$. The induced map on fundamental
groups is an isomorphism and the target has zero second homotopy
group. On universal covers the relative pair is two-connected;
relative Hurewicz gives
$H_3(\widetilde N,\widetilde P)=\pi_2(P)$ and all lower relative
groups vanish. The relative covering homology sequence therefore
identifies $H_3(N,P)=L$ and $H_2(N,P)=0$.
The universal coefficient theorem and the relative cohomology
sequence identify the map
\[
       H^3(N,P)=\operatorname{Hom}(L,\mathbb Z)
         \longrightarrow H^3(N)=\mathbb Z
\]
with evaluation on $z$. Its image is $12\mathbb Z$ by \eqref{maps:eq:E.S9}.

\item The free summand of the quotient in \eqref{maps:eq:E.S10} splits. Lift its
seven order-twelve generators to $x_1,\ldots,x_7$ in the
remaining extension. There are integers $e_i$, defined modulo
twelve, with the complete abelian presentation
\[
        \langle z,x_1,\ldots,x_7\mid
                      12x_i=e_i z\ (1\leq i\leq7)\rangle.
 \tag{E.S11}\label{maps:eq:E.S11}
\]
Evaluation of homomorphisms to $\mathbb Z$ on $z$ has image
$\bigl(12/\gcd(12,e_1,\ldots,e_7)\bigr)\mathbb Z$.
Thus that gcd is one. The relation matrix in \eqref{maps:eq:E.S11} has
Smith factors $1,12,12,12,12,12,12$, giving
$\mathbb Z\oplus(\mathbb Z/12)^6$. Adding the split free
summand proves \eqref{maps:eq:E.S6}.

\item If $\pi_2(P)$ were projective over $\mathbb ZG$, it would be
a direct summand of a free module. Taking coinvariants preserves
that splitting, so $L$ would be a projective, hence free,
abelian group. The torsion in \eqref{maps:eq:E.S6} contradicts this. Finally
$N\times S^6$ has zero second homotopy group. Connected sum in
dimension ten with a simply connected sphere product adds its
second homotopy as a free induced $\mathbb ZG$-module, as seen
by lifting the connected-sum necks and using their simply
connected boundary spheres. Only $r=2,8$ contribute, with one
free generator each; a homotopy sphere contributes none.
Their second homotopy modules are therefore free, so none is
homotopy equivalent to $P$.

\item The program \texttt{maps\_shifted\_pi2.py} checks the simultaneous
contraction kernel, the three homological boundary coordinates,
and the exact-order-twelve image certificate for all sixteen
independent boundary-adjoint sign choices in \ref{maps:src:E.2.6}.
\end{enumerate}
\end{proof}

\begin{theorem}[central localizations of the homotopy module]\label{maps:src:E.2.12}
\StatementRecord{maps:src:E.2.12}{proved}{Historical events: Claude Fable 5.2 (Anthropic), effort high, 16 November 2026: carried as background without fault found. GPT-6 Astra (OpenAI), ChatGPT Pro mode, effort Pro, S6\_09\_PAPERS, 17 November 2026: targeted proof check. Also checked in this writing. Pre-continuation writing events: GPT-6 Astra (OpenAI), ChatGPT Pro mode, effort Pro, S6\_10\_PAPERS, 18 November 2026: targeted proof check. Continuation metadata events: GPT-6 Astra (OpenAI), ChatGPT Pro mode, effort Pro, S6\_10\_PAPERS continuation, 19 November 2026: used and rechecked in this run; status-vocabulary reconciliation; metadata reconciliation only; no new mathematical check. \emph{Pre-continuation recorded proof scope.} GPT-6 Astra (OpenAI, ChatGPT Pro mode; recorded reviewer /root/cdp\_audit), effort Pro, 17 November 2026; Claude Fable 5.2 (Anthropic), effort high, 16 November 2026; also checked in this writing by GPT-6 Astra (OpenAI), ChatGPT Pro mode, effort Pro, 18 November 2026}{the hypotheses, notation and definitions stated here; \cref{maps:src:E.2.9}; \cref{maps:src:E.2.10}; \cref{maps:eq:E.S17}; \cref{maps:src:E.2.3}; \cref{maps:eq:E.S19}; \cref{maps:eq:E.S21}; \cref{maps:src:E.2.11}}


Retain the central generators $h_1,h_2$ of $G$ from \ref{maps:src:E.2.9} and put
\[
 R=\mathbb ZG,\qquad R_i=R[(h_i-1)^{-1}],\qquad
              R_{12}=R[(h_1-1)^{-1},(h_2-1)^{-1}].
 \tag{E.S13}\label{maps:eq:E.S13}
\]
These are flat central localizations. As modules over the indicated
localized group rings, the actual second homotopy group satisfies
\[
 \begin{split}
 R_1\otimes_R\pi_2(P)&\simeq R_1/(h_2-1),\\
 R_2\otimes_R\pi_2(P)&\simeq R_2/(h_1-1),\\
 R_{12}\otimes_R\pi_2(P)&=0.
 \end{split}
 \tag{E.S14}\label{maps:eq:E.S14}
\]
More generally, $H_j(P;R_1)$ is zero outside $2\leq j\leq7$;
in those six degrees it is respectively a direct sum of
\[
                  1,\ 3,\ 4,\ 4,\ 3,\ 1
 \tag{E.S15}\label{maps:eq:E.S15}
\]
copies of $R_1/(h_2-1)$. The analogous statement holds with the
indices interchanged. All homology with coefficients in $R_{12}$
vanishes. For the fibre complement of \ref{maps:src:E.2.10}, with
$S=\mathbb ZK[(h-1)^{-1}]$, one likewise has
\[
             S\otimes_{\mathbb ZK}\pi_2(E)\simeq S,
       \qquad H_4(E;S)\simeq S,
 \tag{E.S16}\label{maps:eq:E.S16}
\]
and its other localized homology groups vanish. These are
localizations of the actual group-ring modules, not a replacement
of $G$ or $K$ by their abelianizations. Moreover, if
$q=(h_1-1)(h_2-1)$, then
\[
                  q^2\pi_2(P)=0,
        \qquad \dim_{\mathbb Q}\bigl(\pi_2(P)\otimes\mathbb Q\bigr)
                     =\infty.
 \tag{E.S20}\label{maps:eq:E.S20}
\]
Consequently the universal cover has no finite homotopy type.
There is no fibration of $P$ over $N$ with compact manifold fibre
and inducing an isomorphism on fundamental groups.

\end{theorem}


\begin{proof}
\begin{enumerate}[label=\arabic*.,leftmargin=*]
\item \emph{The cusp module with its marked deck transformation.}
In the original cusp, the angular generators $w,d$ vanish and
$\pi_1(N_0)=\mathbb Z\langle g,u\rangle$. In $E$ their images
are respectively $h,1$; in the first factor of $P$ they are
$h_1,1$. Thus the connected cover of this cusp occurring inside
the universal cover is the cover $\widehat W$ of its central
fibre in which the $g$ deck transformation is $z$, whereas $u$
is retained as a loop. Write $A=\mathbb Z[z,z^{-1}]$.
Its ordinary homology modules are
\[
 H_0(\widehat W)=H_1(\widehat W)=A/(z-1),\quad
 H_2(\widehat W)=A\oplus(A/(z-1))^2,
 \quad H_3(\widehat W)=0,\quad H_4(\widehat W)=A.
 \tag{E.S17}\label{maps:eq:E.S17}
\]
Here is an integral normalization calculation of this input.
The dual complex of the cover is an infinite cylinder, whose
ordinary homology is $A/(z-1)$ in degrees zero and one and zero
in degree two. The normalization degree-four row is a free
copy of $A$. In degree two, the homology map from the three
double curves to the four component classes is the transpose,
up to the unimodular intersection marking and possible inversion
of $z$, of
\[
 B(z)=\begin{pmatrix}
 -1&-1&-1&-1\\ 1&1&1&z\\-z&-z&-1&-z
 \end{pmatrix}
 \ \sim\
 \begin{pmatrix}1&0&0&0\\0&z-1&0&0\\0&0&z-1&0\end{pmatrix}.
 \tag{E.S18}\label{maps:eq:E.S18}
\]
For example multiply on the left by
$\left(\begin{smallmatrix}-1&0&0\\1&1&0\\-z&0&1\end{smallmatrix}\right)$,
then use as new columns $e_1,e_4-e_1,e_3-e_1,e_2-e_1$.
Both operations are unimodular over $A$. The transpose is
injective and its cokernel is the displayed degree-two module.
All strata have only even homology, and the normalization
direction has length two, so no higher differential is possible.
There is no degree-two extension from the dual complex, whose
second homology is zero. This proves \eqref{maps:eq:E.S17}. Ordinary chains
in this locally finite normalization have compact support and
use direct sums, so no compactly supported cohomology has been
substituted for ordinary homology. The exact matrix certificate
is \texttt{manifolds\_cusp\_alexander.py}, contributed by the
six-manifolds worker and independently rerun here.

\item \emph{Localization and the vanishing of every other piece.}
The central subgroup $\langle h_1,h_2\rangle$ is an injected
$\mathbb Z^2$ by \ref{maps:src:E.2.9}. The group ring is free over its central
Laurent subring, by choosing coset representatives. In particular
$h_i-1$ and their products are non-zero-divisors. Inverting them
is central localization and is flat; it commutes with homology
of the universal-cover chain complex. The analogous assertions
hold for $h$ in $K$.

\item Use the regular torus bundle over a graph and the six closed
cap neighborhoods in \ref{maps:src:E.2.3}. For $R_1$ coefficients, the first
regular torus has its $g_1$ loop acting by $h_1$. The torus
chain complex is the Koszul complex of its commuting loop
actions. Its $h_1-1$ term is invertible, providing a chain
contraction. Thus every regular fibre has zero homology with
these coefficients. The homology sequence of a bundle over
a graph, and the same calculation over each boundary circle,
show that the whole regular piece and every cap boundary are
acyclic. Each cap of the second factor is a product with a
first-factor regular torus over a disc, so it too is acyclic.

\item At a finite cap of the first factor, the affine generator
$a_j$ has $a_j^{m_j}=h_1^{\pm1}$ in its image in $G$.
The connected cover of the reduced flat surface corresponding
to its map onto $\langle a_j\rangle$ is $T^3\times\mathbb R$.
The $m_j$-th power acts by an integral translation with linear
part $A_j^{m_j}=I$, and therefore acts identically on all of
its homology. Thus $h_1-1$ annihilates every such homology
module before localization. Retaining the other regular torus
does not change that assertion, and induction to $G$ commutes
with the central action. Both first-factor finite caps are
therefore acyclic after localization as well. This argument
keeps their actual affine translation, rather than treating
the quotient as an untwisted torus product.

\item \emph{The one surviving cap and induction.}
The first cusp is $N_{0,1}\times F_2$, since the two critical
triples are disjoint. Its image in $G$ is precisely
$H=\langle h_1,h_2\rangle$. A connected component of its
preimage in the universal cover is
\[
                \widehat W\times\mathbb R\times T^3
\]
up to deformation retraction: $g_1$ acts on $\widehat W$ by
$z=h_1$, and $g_2=h_2$ translates the real factor. It acts
trivially on the homology of the retained three-torus. All
components together give induction from $\mathbb ZH$ to
$R$. By \eqref{maps:eq:E.S17}, inverting $h_1-1$ leaves one free
$\mathbb Z[h_1^{\pm1},(h_1-1)^{-1}]$ summand in cusp degrees
two and four and nothing else. The integral product homology
with $T^3$ has no torsion correction. Inducing the result gives
\[
 H_j(N_{0,1}\times F_2;R_1)
   \simeq
 \left(R_1/(h_2-1)\right)^{
           \binom3{j-2}+\binom3{j-4}},
 \tag{E.S19}\label{maps:eq:E.S19}
\]
where a binomial coefficient is zero outside its usual range.
All other pieces and their intersections are acyclic, as
proved above, so Mayer--Vietoris identifies this with
$H_j(P;R_1)$. This proves \eqref{maps:eq:E.S15} and \eqref{maps:eq:E.S19}, and interchanging
the factors proves the second version. Inverting $h_2-1$
also kills \eqref{maps:eq:E.S19}, proving the double-localization assertion.
Since $\widetilde P$ is simply connected, Hurewicz identifies
its second homology with $\pi_2(P)$; flat localization now
gives \eqref{maps:eq:E.S14}.

\item For $E$ the same decomposition has only one cusp and two
finite caps. Every other piece is acyclic after $h-1$ is
inverted by the identical argument. Its cusp image is
$\langle h\rangle$, so inducing \eqref{maps:eq:E.S17} after localization
gives one free $S$ summand in degrees two and four. This
proves \eqref{maps:eq:E.S16}.

\item \emph{A uniform annihilator without localization.}
Put $I=(h_1-1,h_2-1)\subset R$. Before localization, the homology
of a regular eight-torus with $R$ coefficients is
\[
                 (R/I)\otimes_{\mathbb Z}\bigwedge^j\mathbb Z^6
                      \quad(0\leq j\leq6).
 \tag{E.S21}\label{maps:eq:E.S21}
\]
Indeed the two nontrivial loop images are the independent central
generators $h_1,h_2$. Their differences from one form a regular
sequence in the central Laurent ring, and $R$ is free over that
ring. The two-variable Koszul complex therefore has only its
degree-zero quotient as homology; the other six loop differentials
are zero. This proves \eqref{maps:eq:E.S21}.
The homology sequence of the regular bundle over a graph has
only base degrees zero and one. Both filtration quotients are
annihilated by $I$, by \eqref{maps:eq:E.S21}, so $I^2$ annihilates the homology
of the regular piece. The same two-stage Wang argument applies
to every boundary mapping torus. It does not require splitting
either filtration or treating the monodromy as trivial.

\item At a finite cap its connected lift is a product of six torus
directions and two real directions up to deformation retraction.
Each central deck generator acts trivially on homology: the
one belonging to that cap is the finite affine power already
checked above, and the other translates the regular real
direction. Hence $I$ annihilates its induced homology. At the
first cusp, $h_2$ translates the real factor of the other
regular torus, so $h_2-1$ annihilates all of its homology;
the second cusp is annihilated by $h_1-1$. Thus $q$ annihilates
the homology of every piece and every intersection, since
$q\in I^2$. The decomposition has a regular piece and disjoint
caps with disjoint boundary collars, so its Mayer--Vietoris
sequence gives in every degree
\[
 0\longrightarrow\operatorname{coker}D_j
   \longrightarrow H_j(P;R)
   \longrightarrow\ker D_{j-1}\longrightarrow0.
\]
Both end groups are killed by $q$, and therefore the middle
group is killed by $q^2$. Since $H_2(P;R)=\pi_2(P)$, this
proves the first assertion of \eqref{maps:eq:E.S20} integrally.

\item \emph{Infinite rational rank.}
If $\pi_2(P)\otimes\mathbb Q$ were finite dimensional,
localizing its linear automorphism $h_1$ at $h_1-1$ would
remain finite dimensional. Explicitly, this removes the
generalized eigenspace with eigenvalue one and inverts an
automorphism on the remaining finite-dimensional quotient.
But \eqref{maps:eq:E.S14}, tensored with $\mathbb Q$, contains a copy of
$\mathbb Q[h_1^{\pm1},(h_1-1)^{-1}]$, as follows by taking
the identity coset in the free module over the central
Laurent subring and setting $h_2=1$. That vector space is
infinite dimensional, a contradiction. This proves the
second assertion and rules out a finite homotopy type of
the universal cover. A fibration over the aspherical $N$
inducing a fundamental-group isomorphism would have simply
connected fibre, and pulling it to $\widetilde N$ would
identify that fibre up to homotopy with $\widetilde P$.
A compact manifold fibre has finite homotopy type, so such
a fibration is impossible.

\item The quotient $R_1/(h_2-1)$ is nonzero, since the two central
generators are independent, and it is annihilated by the
non-zero-divisor $h_2-1$. A projective $R_1$-module is a
direct summand of a free module and has no such nonzero
annihilated element. Thus \eqref{maps:eq:E.S14} also gives an independent
proof of the nonprojectivity in \ref{maps:src:E.2.11}. It does not determine
the part of $\pi_2(P)$ supported where both central
elements specialize to one, or its remaining extension data.
\end{enumerate}
\end{proof}

\begin{theorem}[an actual circle fibration and infinite-order return]\label{maps:src:E.2.13}
\StatementRecord{maps:src:E.2.13}{proved}{Historical events: Claude Fable 5.2 (Anthropic), effort high, 16 November 2026: carried as background without fault found. GPT-6 Astra (OpenAI), ChatGPT Pro mode, effort Pro, S6\_09\_PAPERS, 17 November 2026: targeted proof check. Also checked in this writing. Pre-continuation writing events: GPT-6 Astra (OpenAI), ChatGPT Pro mode, effort Pro, S6\_10\_PAPERS, 18 November 2026: targeted proof check. Continuation metadata events: GPT-6 Astra (OpenAI), ChatGPT Pro mode, effort Pro, S6\_10\_PAPERS continuation, 19 November 2026: used and rechecked in this run; status-vocabulary reconciliation; metadata reconciliation only; no new mathematical check. \emph{Pre-continuation recorded proof scope.} GPT-6 Astra (OpenAI, ChatGPT Pro mode; recorded reviewer /root/cdp\_audit), effort Pro, 17 November 2026; Claude Fable 5.2 (Anthropic), effort high, 16 November 2026; also checked in this writing by GPT-6 Astra (OpenAI), ChatGPT Pro mode, effort Pro, 18 November 2026}{the hypotheses, notation and definitions stated here; \cref{maps:src:E.2.9}; \cite[Theorem 3.6]{rootScott83}; \cref{maps:eq:E.S26}; \cref{maps:src:E.2.12}; \cref{maps:src:B.6.24}}


The shifted fivefold admits a smooth fibre bundle
\[
                  Y^9\longrightarrow P\stackrel{\rho}{\longrightarrow}S^1
 \tag{E.S22}\label{maps:eq:E.S22}
\]
with connected closed fibre, inducing the primitive character
\[
 \chi(a_1,a_2,b_1,b_2,h_1,h_2)
                    =(4,-3,-4,3,12,-12)
 \tag{E.S23}\label{maps:eq:E.S23}
\]
on the presentation of \ref{maps:src:E.2.9}. Its fibre group is the fundamental
group of the oriented principal circle bundle of Euler number
$-1$ over a surface of genus six. Here the base has its complex
orientation and the positive circle is $h_1h_2$; reversing that
circle changes the Euler number to $+1$. In particular
\[
 \pi_1(Y)=\ker\chi,\qquad H_1(Y;\mathbb Z)=\mathbb Z^{12}.
 \tag{E.S24}\label{maps:eq:E.S24}
\]
The return diffeomorphism of \eqref{maps:eq:E.S22} has infinite order even
among unbased homotopy classes of self-equivalences of $Y$,
although its induced outer automorphism of $\pi_1(Y)$ has
exact order twelve. No connected finite cover of $P$ is
homotopy equivalent to a product $S^1\times Z$.
These assertions give a geometric mapping-torus description;
they do not identify the smooth nine-dimensional fibre or
the full return diffeomorphism up to classification.

\end{theorem}


\begin{proof}
\begin{enumerate}[label=\arabic*.,leftmargin=*]
\item \emph{A closed one-form with a nonzero derivative at every point.}
Use the actual smooth map $c=(e_1,e_2):P\to N$ in \ref{maps:src:E.2.9},
not an arbitrary representative of its classifying homotopy
class. The manifold $N$ has the locally free $T^2$ action
induced by the two weighted Hopf actions. Take a closed
one-form representing the primitive integral class $\chi$
on $N$ and average it over this compact torus. Its average
$\alpha$ is closed, represents the same integral class,
and is invariant. If $V_i$ is the fundamental vector field
with time-one orbit $h_i$, Cartan's formula gives
$d(\alpha(V_i))=0$. Integrating along an ordinary orbit
therefore gives
\[
                  \alpha(V_1)=12,\qquad \alpha(V_2)=-12.
 \tag{E.S25}\label{maps:eq:E.S25}
\]
The integrality of the periods integrates $\alpha$ to a
smooth map $\rho_N:N\to\mathbb R/\mathbb Z$.

\item Choose the two cusp replacements defining $e_i$ in disjoint
base discs, themselves disjoint from the other factor's
critical values. At every point of $P$, at least one factor
is a regular torus at which its $e_i$ is the literal character
quotient of \cref{mapsI:src:A.1.1}. In a local flat marking of that factor,
the constant vertical vector $g_i$ maps to $V_i$. Holding
the other factor fixed gives a tangent vector to $P$, and
\[
                     (c^*\alpha)(g_i)=12\quad\hbox{or }-12.
\]
Thus $c^*\alpha$ is nowhere zero. This is a local derivative
calculation on the actual map, not an inference from its
fundamental-group character. The map $\rho=\rho_Nc$ is a
proper submersion. Equivalently, choose a vector field $W$
with $(c^*\alpha)(W)=1$; its complete flow supplies local
trivializations and a return diffeomorphism. The fibre is
connected because $\chi$ is primitive: the image in
$\pi_1(S^1)$ of a connected component of a bundle's total
space has index equal to the number of fibre components.
The homotopy sequence now proves $\pi_1(Y)=\ker\chi$.

\item \emph{The exact surface cover and its circle Euler class.}
Write $\Gamma=G/\langle h_1,h_2\rangle$ for the oriented
orbifold group of $S^2(3,4,3,4)$. The character modulo
twelve descends to the surjection
\[
 \overline\chi:\Gamma\longrightarrow\mathbb Z/12,
       \qquad(a_1,a_2,b_1,b_2)\longmapsto(4,-3,-4,3).
 \tag{E.S26}\label{maps:eq:E.S26}
\]
Its restriction to every cone group is injective. Hence its
kernel $\Gamma_0$ is torsion free and defines a degree-twelve
ordinary surface cover of the orbifold. Its Euler characteristic
is $12(-5/6)=-10$, so that surface is $\Sigma_6$. The image
of $\ker\chi$ in $\Gamma$ is exactly $\Gamma_0$: a lift of
an element of $\Gamma_0$ has character a multiple of twelve,
which can be cancelled by a power of $h_1$. Its central
kernel is the primitive subgroup
\[
 \ker\chi\cap\langle h_1,h_2\rangle
                =\mathbb Z\langle h_1h_2\rangle.
 \tag{E.S27}\label{maps:eq:E.S27}
\]

\item The circle Euler number can be fixed without an ambiguous
Seifert sign convention. On either weighted Hopf sphere put
$z_1=\cos r\,e^{i\varphi_1}$ and
$z_2=\sin r\,e^{i\varphi_2}$. For the positive action with
weights $(4,3)$ the smooth invariant connection is
\[
 \vartheta=\tfrac14\cos^2r\,d\varphi_1
                    +\tfrac13\sin^2r\,d\varphi_2.
\]
It evaluates to one on the generator with period $2\pi$.
With the boundary orientation of the unit ball in $\mathbb C^2$,
\[
             \frac1{(2\pi)^2}\int_{S^3}\vartheta\wedge d\vartheta
                         =\frac1{12}.
 \tag{E.S28}\label{maps:eq:E.S28}
\]
For the associated complex line whose unit circle is this
positive circle bundle, the convention agreeing with the
tautological Hopf line is $c_1=-[d\vartheta/(2\pi)]$.
Thus each orbifold circle Euler number is $-1/12$.
Both shifted base identifications preserve the complex
orientation. Pullback to $\Sigma_6$ multiplies the two Euler
numbers by twelve, giving $(-1,-1)$; this is also the
covering formula of \cite[Theorem 3.6]{rootScott83}.
In the integral fibre basis
\[
                  v=(1,1),\qquad w=(0,1)
\]
the determinant is one, the Euler vector is $-v$, and
$\chi(v)=0$, $\chi(w)=-12$. Consequently the circle kernel
has Euler number $-1$ with generator $h_1h_2$.

\item One can see this bundle directly in $N$. Invariance of
$\alpha$ implies
$\rho_N((s_1,s_2)x)=\rho_N(x)+12(s_1-s_2)$, where
$s_i\in\mathbb R/\mathbb Z$. A level contains twelve
diagonal circles over an ordinary base point. The peripheral
permutations \eqref{maps:eq:E.S26} give exactly the connected cover
$\Sigma_6$; they have no stabilizers at the cone points.
The diagonal circle acts freely, since at every cone only
one of the two Hopf actions has a nontrivial stabilizer.
The level is therefore the circle bundle just computed.
This identifies its fundamental group with $\ker\chi$
and gives an actual three-dimensional aspherical model.
The homology sequence of a circle bundle of Euler number
$\pm1$ gives its first homology $\mathbb Z^{12}$, proving
the second assertion in \eqref{maps:eq:E.S24}.

\item \emph{Faithful central action and the return obstruction.}
The centre of $G$ is exactly $\langle h_1,h_2\rangle$:
the quotient is a nonelementary oriented hyperbolic orbifold
group with trivial centre. This centre acts faithfully on
$\pi_2(P)$. Indeed if $h_1^a h_2^b$ acts as the identity,
localize first at $h_1-1$ and apply \ref{maps:src:E.2.12}. On
$R_1/(h_2-1)$ its action is multiplication by $h_1^a$.
This module is free over the central ring
$\mathbb Z[h_1^{\pm1},(h_1-1)^{-1}]$ after choosing coset
representatives, so this is the identity only if $a=0$.
The symmetric localization then gives $b=0$.

\item If $H<G$ has finite index, its image in $\Gamma$ is again
nonelementary. Its centralizer in $\Gamma$ is trivial:
commuting with two hyperbolic elements with distinct axes
fixes at least three boundary points and hence is the
identity. It follows that
$Z(H)=H\cap Z(G)$. A homotopy product $S^1\times Z$ for
the corresponding finite cover would supply an infinite
central element acting trivially on $\pi_2$, namely the
circle factor; this contradicts the faithful action just
proved. If a positive power of the return of \eqref{maps:eq:E.S22} were
homotopic to the identity, its mapping torus would be
homotopy equivalent to $S^1\times Y$. That mapping torus
is the corresponding finite cyclic cover of $P$, so this
too is impossible.

\item Finally take $t=a_1a_2$, for which $\chi(t)=1$, as a
stable letter. The image of $\ker\chi$ in $\Gamma$ is
$\Gamma_0$, so the same centralizer argument gives
$C_G(\ker\chi)=Z(G)$. The automorphism
$\operatorname{Ad}(t^n)|_{\ker\chi}$ is inner precisely
when $t^n$ differs from an element of $\ker\chi$ by a
central element. Since $\chi(Z(G))=12\mathbb Z$, this
holds precisely for $12\mid n$. Thus the outer action has
order twelve, while the actual return has infinite homotopy
order. This return is the circle fibration of the shifted
ten-manifold; it is distinct from the six-sphere parameter
returns computed in \ref{maps:src:B.6.24}.

\item The exact integral character, cone orders, primitive kernel
basis and Euler calculation are recorded in
\texttt{code/maps\_shifted\_circle.py}. No identification
of $Y$ with a product or connected sum is used here.
\end{enumerate}
\end{proof}

\begin{theorem}[the cyclic-cover splice and its separating torus twist]\label{maps:src:E.2.14}
\StatementRecord{maps:src:E.2.14}{proved}{Historical events: Claude Fable 5.2 (Anthropic), effort high, 16 November 2026: carried as background without fault found. GPT-6 Astra (OpenAI), ChatGPT Pro mode, effort Pro, S6\_09\_PAPERS, 17 November 2026: targeted proof check. Also checked in this writing. Pre-continuation writing events: GPT-6 Astra (OpenAI), ChatGPT Pro mode, effort Pro, S6\_10\_PAPERS, 18 November 2026: targeted proof check. Continuation metadata events: GPT-6 Astra (OpenAI), ChatGPT Pro mode, effort Pro, S6\_10\_PAPERS continuation, 19 November 2026: used and rechecked in this run; status-vocabulary reconciliation; metadata reconciliation only; no new mathematical check. \emph{Pre-continuation recorded proof scope.} GPT-6 Astra (OpenAI, ChatGPT Pro mode; recorded reviewer /root/cdp\_audit), effort Pro, 17 November 2026; Claude Fable 5.2 (Anthropic), effort high, 16 November 2026; also checked in this writing by GPT-6 Astra (OpenAI), ChatGPT Pro mode, effort Pro, 18 November 2026}{the hypotheses, notation and definitions stated here; \cref{maps:src:E.2.13}; \cref{maps:src:E.2.6}; \cref{maps:src:E.2.10}; \cref{maps:eq:E.S34}; \cref{maps:eq:E.S35}; \cref{maps:eq:E.S37}; \cref{maps:eq:E.S39}; \cref{maps:eq:E.S38}; \cref{maps:eq:E.S41}; \cref{maps:eq:E.S43}; \cref{maps:eq:E.S36}; \cref{maps:eq:E.S40}; \cref{maps:eq:E.S44}; \cref{maps:eq:E.S45}; \cref{maps:src:E.2.7}; \cref{maps:eq:E.S47}; \cref{maps:eq:E.S46}}


Let $E_{12}\to E$ be the connected cyclic cover defined by
reduction modulo twelve of the positive meridian character.
The circle fibration of \ref{maps:src:E.2.13} may be chosen with fibre
\[
 Y=(E_{12}\times T^3)\ \cup_{T^8}\ (T^3\times E_{12}),
 \tag{E.S29}\label{maps:eq:E.S29}
\]
with the following exact boundary marking. Write
$\theta_i'$ for the lifted normal circle and retain the
two regular-torus markings. Then
\[
       \theta_1'=\gamma_2-\gamma_1,\qquad
       \theta_2'=\gamma_1-\gamma_2.
 \tag{E.S30}\label{maps:eq:E.S30}
\]
There is a return representative $\psi$ whose twelfth power
is the separating torus twist
\[
 \begin{gathered}
 T(x,s)=(x+\eta(s)(g_1+g_2),s)\quad\hbox{on }T^8\times[-1,1],
 \\
 T=\operatorname{id}\text{ elsewhere}.
 \end{gathered}
 \tag{E.S31}\label{maps:eq:E.S31}
\]
Here $\eta$ is zero near $-1$ and one near $1$, and circle
coordinates have period one. In particular \eqref{maps:eq:E.S31} is the
identity near both ends of its collar. Its rational
cohomological variation satisfies
\[
 (T^*-I)^2=0,\qquad
 \operatorname{rank}(T^*-I)|_{H^j(Y;\mathbb Q)}=
 \begin{cases}1&j=4,5,\\0&j\ne4,5.\end{cases}
 \tag{E.S32}\label{maps:eq:E.S32}
\]
The rational Betti numbers of the actual fibre are
\[
       (b_0,\ldots,b_9)(Y)
          =(1,12,121,407,693,693,407,121,12,1).
 \tag{E.S33}\label{maps:eq:E.S33}
\]
Every eigenvalue of $\psi^*$ is a twelfth root of unity.
At eigenvalue one there is exactly one Jordan block of
size two in each of degrees four and five; all its other
blocks, and all blocks at other eigenvalues, have size one.
Thus the infinite order is already visible on ordinary
rational cohomology. The smooth complement $E_{12}$ in
\eqref{maps:eq:E.S29} remains the actual finite cover of $E$, not a
replacement inferred from its group or homology.

\end{theorem}


\begin{proof}
\begin{enumerate}[label=\arabic*.,leftmargin=*]
\item \emph{Adapting the fibration to the actual splice.}
Cut the base along the circle used in \ref{maps:src:E.2.6}, with the two
critical triples on opposite sides. Over the first disc the
second factor is a regular marked $T^4$, and over the second
disc the first factor is regular. Let
$\xi_E:E\to\mathbb R/\mathbb Z$ represent the positive
meridian class. Its restriction to the boundary can be
chosen exactly
\[
                         \xi_E=\theta+12\gamma.
 \tag{E.S34}\label{maps:eq:E.S34}
\]
Indeed this is its integral restriction by \ref{maps:src:E.2.10}; homotopy
extension makes the representative equal to that linear
map on a collar. The two maps
\[
        \rho_A=\xi_{E_1}-12\gamma_2,
             \qquad \rho_B=12\gamma_1-\xi_{E_2}
 \tag{E.S35}\label{maps:eq:E.S35}
\]
agree exactly on the common boundary, where
$\theta_2=-\theta_1$. Smooth them with this common product
collar. The regular-factor circle derivative is $-12$ on
the first piece and $12$ on the second, so they give a
proper submersion throughout, including both cusp pieces.
These are also permissible choices in E.2.13: on the
quotient $N$, an invariant closed one-form differs from
these representatives only by the differential of a
single-valued function on each of the two pieces. Absorb
that function by translation in the unused regular circle.
Thus this construction uses the same actual quotient and
primitive character.

\item A level of $\xi_E-12\gamma_2$ is the pullback of the
degree-twelve circle cover along $\xi_E$, times the three
remaining regular circles. The pullback is $E_{12}$.
Its boundary is connected because the normal meridian has
character one. The covering of that boundary has
$\theta=12\theta'$ and unchanged $\gamma,u,w,\delta$.
Equation \eqref{maps:eq:E.S34} in the level condition gives precisely
\eqref{maps:eq:E.S30}. This proves \eqref{maps:eq:E.S29} with its actual smooth
boundary maps.

\item Choose the vector field integrating \eqref{maps:eq:E.S35} to be
$-g_2/12$ on the first piece off its collar and $g_1/12$
on the second. On the common regular collar interpolate by
\[
                     W_s=\frac{\eta(s)g_1-(1-\eta(s))g_2}{12}.
 \tag{E.S36}\label{maps:eq:E.S36}
\]
It has $d\rho(W_s)=1$, and it preserves the collar
coordinate $s$ and the original base point. Its flow
therefore defines a return $\psi$. At time twelve it is
the identity off the collar, since the two indicated
regular circles have period one. On the collar it
translates by $\eta(s)g_1-(1-\eta(s))g_2$; removing the
integral translation $-g_2$ gives \eqref{maps:eq:E.S31}. This comparison
uses a globally regular factor on each piece, so there
is no continuation through a singular cap to justify.

\item \emph{The variation map and the invariant boundary calculation.}
Put $D=T^8$ and let $r:H^j(Y;\mathbb Q)\to H^j(D;\mathbb Q)$
be restriction. Denote the Mayer--Vietoris connecting map
by $\partial:H^{j-1}(D;\mathbb Q)\to H^j(Y;\mathbb Q)$.
Writing $v=g_1+g_2$, the collar pullback of a translation
invariant form is
\[
       T^*\omega=\omega+d\eta\wedge\iota_v\omega.
\]
Average its restriction to the torus and change the global
representative by an exact form to make this equality
applicable on the collar. The difference represents the
connecting class, and hence, with the connecting sign
chosen by the increasing collar coordinate,
\[
                          T^*-I=\partial\,\iota_v\,r.
 \tag{E.S37}\label{maps:eq:E.S37}
\]
As $r\partial=0$, its square is zero.

\item No unknown part of $H^*(E_{12})$ is needed to compute
\eqref{maps:eq:E.S37}. Its cyclic deck action is a translation on the
boundary torus, hence acts trivially on boundary cohomology.
Every noninvariant deck eigenspace therefore restricts to
zero. Rational transfer identifies the invariant part
with $H^*(E)$, whose marked restriction is \ref{maps:src:E.2.6}. Its
pullback to \eqref{maps:eq:E.S30} is
\[
 R_{k,1}(a)=\phi_k(a)+12(\gamma_2-\gamma_1)a,
 \qquad
 R_{k,2}(b)=\phi_k(b)+12(\gamma_1-\gamma_2)b.
 \tag{E.S38}\label{maps:eq:E.S38}
\]
Retain the other factor's three covectors $u,w,\delta$,
and take exterior products to obtain the two subspaces
$U_j^k\subset\bigwedge^k\mathbb Q^8$.
The image of $r$ is exactly $U_1^k\cap U_2^k$, and the
kernel of $\partial$ is $U_1^{k-1}+U_2^{k-1}$.
Consequently the rank in \eqref{maps:eq:E.S32} is exactly the rank of
\[
 U_1^k\cap U_2^k\ \stackrel{\iota_v}{\longrightarrow}
 \bigwedge^{k-1}\mathbb Q^8/(U_1^{k-1}+U_2^{k-1}).
 \tag{E.S39}\label{maps:eq:E.S39}
\]

\item Here are the full dimensions for this small rational
calculation. Let $d_k$ be twice the degree-$k$ dimension
of $H^*(E)\otimes H^*(T^3)$ and let
$r_k=\dim(U_1^k+U_2^k)$:
\[
 \begin{array}{c|rrrrrrrrr}
 k&0&1&2&3&4&5&6&7&8\\ \hline
 d_k&2&8&20&44&70&68&36&8&0\\
 r_k&1&8&19&34&47&46&27&8&0
 \end{array}
 \tag{E.S40}\label{maps:eq:E.S40}
\]
Contraction in \eqref{maps:eq:E.S39} has rank one only for
$k=4,5$ and is zero in every other degree. All entries
are computed by exterior products of \eqref{maps:eq:E.S38} in
\texttt{maps\_shifted\_return.py}; its output retains an
explicit nonzero boundary witness in each of degrees
four and five. One can normalize the remaining sign in
$\phi_2$ by the ring structure: with
$r_1(1)=\theta+12\gamma$, the equality
$r_1(1)r_2(\delta)=r_3(a)$ forces the coefficient
$c_2=-2$, and then $c_3=2$, in the notation of \ref{maps:src:E.2.6}.
This is also the normalization used by that executable.
No rank calculation treats a rational cokernel as an
integral lattice.

\item \emph{The finite-cover cohomology needed for the full Betti numbers.}
For clarity, the full rational cohomology of $E_{12}$
can also be obtained without an uncomputed extension.
Let $\mathcal L_\lambda$ be the rank-one system on $E$
whose meridian acts by $\lambda$, with $\lambda^{12}=1$.
On a regular torus it is trivial. Its three base transports
in vertical degree $q$ are
\[
 \lambda^4\bigwedge^q A_1^{-t},\qquad
 \lambda^{-3}\bigwedge^q A_2^{-t},\qquad
                         \bigwedge^q M_0^{-t}.
 \tag{E.S41}\label{maps:eq:E.S41}
\]
The exponents retain the actual affine meridian character;
using only the finite linear matrices would give a different
cover.

\item At a finite cap, averaging on its finite root cover
identifies the cap cohomology with the invariant part of
\eqref{maps:eq:E.S41}. At the unchanged cusp the system is trivial,
because $g$ has character twelve and $u$ has character zero.
Its rational specialization is an isomorphism onto the
monodromy invariants in every degree. Here is an explicit
check of that last input. The cusp Betti numbers are
$(1,2,4,2,1)$ by its normalization calculation. In degree
two, the integral surjectivity of the specialization on
homology in (B.5.16) gives injectivity on cohomology;
the monodromy-invariant subspace has the same dimension
four. Its image therefore includes $uw$ and $\gamma\delta$.
In degree one the image is $\langle\gamma,u\rangle$;
their products with those two degree-two classes give the
two independent invariant classes in degree three and
the top class in degree four. Thus all the specialization
maps have the asserted ranks, with no additional
skyscraper summand. The paragraph following (B.5.16)
also proves these specialization assertions integrally
in all degrees by the actual honeycomb quotient map.

\item It follows that $R^q f_*\mathcal L_\lambda=j_*L_q$ over
the disc containing the three critical values. On this
disc $H^p(j_*L_q)=0$ for $p\geq2$. To check this directly,
retract the punctured disc to a free graph on three loops.
Restriction of a graph cocycle to the three local
coinvariant spaces is onto: its values on the three free
generators can be chosen independently. The exact
restriction sequence therefore kills $H^2(j_*L_q)$.
If $d_q=\binom4q$, $g_q$ is the common fixed dimension of
\eqref{maps:eq:E.S41}, and $k_{i,q}$ are their three individual fixed
dimensions, this same graph calculation gives
\[
 \dim H^0(j_*L_q)=g_q,\qquad
 \dim H^1(j_*L_q)=2d_q+g_q-\sum_i k_{i,q}.
 \tag{E.S42}\label{maps:eq:E.S42}
\]
The Leray sequence has only these two base degrees and
degenerates. Exact exterior-matrix ranks give the following
table; a row applies to each character of the indicated
order:
\[
 \begin{array}{c|rrrrrr}
 \operatorname{ord}(\lambda)&b_0&b_1&b_2&b_3&b_4&b_5\\ \hline
 1 &1&1&4&6&4&0\\
 2 &0&0&4&6&4&0\\
 3 &0&0&3&4&3&0\\
 4 &0&0&3&4&3&0\\
 6 &0&1&5&6&5&1\\
 12&0&1&4&4&4&1
 \end{array}
 \tag{E.S43}\label{maps:eq:E.S43}
\]
Here $b_j=\dim H^j(E;\mathcal L_\lambda)$ and $b_6=0$.
For the nontrivial characters all $g_q$ are zero; for the
trivial one they are all one. The program
\texttt{maps\_shifted\_cover.py} computes \eqref{maps:eq:E.S41}--\eqref{maps:eq:E.S43}
over the exact cyclotomic fields, represented by rational
companion matrices. It prints every local fixed dimension.
Summing the rows with multiplicities $1,1,2,2,2,4$ gives
\[
             (b_0,\ldots,b_6)(E_{12})=(1,7,46,56,46,6,0).
 \tag{E.S44}\label{maps:eq:E.S44}
\]
Both sides have Euler characteristic $24$, as required by
the degree-twelve covering of $E$.

\item \emph{The Betti and Jordan calculations on the actual fibre.}
The return on each of the two pieces of \eqref{maps:eq:E.S29}, up to
isotopy allowing its boundary to move, is its cyclic deck
transformation times the identity of $T^3$. This follows
also by moving the interpolation in \eqref{maps:eq:E.S36} entirely into
the other piece's collar. It acts trivially on the
cohomology of $D$. The generalized eigenvalue-one part of
the Mayer--Vietoris sequence therefore uses just the
invariant part $H^*(E)$ in \eqref{maps:eq:E.S40}. Its dimensions are
\[
       (1,0,1,19,45,45,19,1,0,1).
 \tag{E.S45}\label{maps:eq:E.S45}
\]
Every other eigenspace is the direct sum of the two piece
eigenspaces, because the boundary has none. In particular
those eigenspaces are semisimple. Their combined Poincare
polynomial is
\[
     2(1+t)^3(6t+42t^2+50t^3+42t^4+6t^5),
 \tag{E.S46}\label{maps:eq:E.S46}
\]
by \eqref{maps:eq:E.S43}--\eqref{maps:eq:E.S44}. Adding \eqref{maps:eq:E.S45} proves \eqref{maps:eq:E.S33}.

\item Finally the Wang sequence for $P=M_\psi$ and the rational
Betti numbers in \ref{maps:src:E.2.7} give the numbers of Jordan blocks
at eigenvalue one as
\[
          (1,0,1,19,44,44,19,1,0,1).
 \tag{E.S47}\label{maps:eq:E.S47}
\]
The only possible block sizes are one and two: the
twelfth power is \eqref{maps:eq:E.S31}, and \eqref{maps:eq:E.S37} squares to zero.
Subtracting \eqref{maps:eq:E.S47} from \eqref{maps:eq:E.S45} gives exactly one block
of size two in each of degrees four and five. This also
independently checks the direct variation ranks \eqref{maps:eq:E.S39}.
All other eigenvalues and blocks are already described
by \eqref{maps:eq:E.S43} and \eqref{maps:eq:E.S46}. None of these rational calculations
claims an integral torsion classification or replaces the
smooth $E_{12}$ by a manifold with the same invariants.
\end{enumerate}
\end{proof}

\begin{theorem}[integral fibre groups and the coefficient of the torus twist]\label{maps:src:E.2.15}
\StatementRecord{maps:src:E.2.15}{proved}{Historical events: Claude Fable 5.2 (Anthropic), effort high, 16 November 2026: carried as background without fault found. GPT-6 Astra (OpenAI), ChatGPT Pro mode, effort Pro, S6\_09\_PAPERS, 17 November 2026: targeted proof check. Also checked in this writing. Pre-continuation writing events: GPT-6 Astra (OpenAI), ChatGPT Pro mode, effort Pro, S6\_10\_PAPERS, 18 November 2026: targeted proof check. Continuation metadata events: GPT-6 Astra (OpenAI), ChatGPT Pro mode, effort Pro, S6\_10\_PAPERS continuation, 19 November 2026: used and rechecked in this run; status-vocabulary reconciliation; metadata reconciliation only; no new mathematical check. \emph{Pre-continuation recorded proof scope.} GPT-6 Astra (OpenAI, ChatGPT Pro mode; recorded reviewer /root/cdp\_audit), effort Pro, 17 November 2026; Claude Fable 5.2 (Anthropic), effort high, 16 November 2026; also checked in this writing by GPT-6 Astra (OpenAI), ChatGPT Pro mode, effort Pro, 18 November 2026}{the hypotheses, notation and definitions stated here; \cref{maps:eq:E.S44}; \cref{maps:eq:E.S33}; \cref{maps:eq:E.S42}; \cref{maps:eq:E.S53}; \cref{maps:eq:E.S29}; \cref{maps:eq:E.S38}; \cref{maps:eq:E.S40}; \cref{maps:eq:E.S37}; \cref{maps:eq:E.S55}}


The groups $H^*(E_{12};\mathbb Z)$ are torsion free, with
the ranks in \eqref{maps:eq:E.S44}. For the actual fibre $Y$, the only
torsion in cohomology is
\[
       H^3(Y;\mathbb Z)_{\rm tors}
       =H^7(Y;\mathbb Z)_{\rm tors}=(\mathbb Z/6)^2.
 \tag{E.S48}\label{maps:eq:E.S48}
\]
Equivalently, its integral homology has the ranks \eqref{maps:eq:E.S33}
and torsion $(\mathbb Z/6)^2$ exactly in degrees two and six.
The twelfth return $T=\psi^{12}$ acts as the identity on
all integral cohomology except in degrees four and five.
In each of those two free groups it has, in an integral
basis, the form
\[
             \begin{pmatrix}1&72\\0&1\end{pmatrix}
                                  \oplus I_{691}.
 \tag{E.S49}\label{maps:eq:E.S49}
\]

For $n\geq1$, let $P_{12n}$ be the connected cyclic cover
of $P$ of degree $12n$ defined by $\chi$. Its Betti numbers,
independent of $n$, are
\[
       (1,13,133,528,1099,1384,1099,528,133,13,1).
 \tag{E.S50}\label{maps:eq:E.S50}
\]
Its integral homology torsion is exactly
\[
 H_j(P_{12n};\mathbb Z)_{\rm tors}=
 \begin{cases}
       (\mathbb Z/6)^2,&j=2,3,6,7,\\
       \mathbb Z/(72n),&j=4,5,\\
       0,&\text{otherwise}.
 \end{cases}
 \tag{E.S51}\label{maps:eq:E.S51}
\]
In particular these covers are pairwise nonhomotopy-equivalent.
These integral conclusions still do not recognize the smooth
$E_{12}$ or replace its boundary marking by homological data.

\end{theorem}


\begin{proof}
\begin{enumerate}[label=\arabic*.,leftmargin=*]
\item \emph{The integral cohomology of the actual cover.}
Apply the proof of \eqref{maps:eq:E.S42} with the coefficient lattice
$\mathbb Z[C_{12}]$, using its regular permutation matrix
$P_{12}$. In vertical degree $q$ set
\[
 \begin{aligned}
 T_{1,q}=P_{12}^4\otimes\bigwedge^q A_1^{-t},\quad
 &T_{2,q}=P_{12}^{-3}\otimes\bigwedge^q A_2^{-t},\\
 T_{0,q}=I_{12}\otimes\bigwedge^q M_0^{-t},
 &V_q=\mathbb Z^{12\binom4q}.
 \end{aligned}
 \tag{E.S52}\label{maps:eq:E.S52}
\]
This coefficient system computes $H^*(E_{12};\mathbb Z)$.
There is no integral averaging step. At the first finite
cap the cover has four components, each the ordinary
degree-three torus root cover; at the second there are
three ordinary degree-four torus root covers. An invariant
vector is determined by its value on one sheet of such
an orbit, with the other values its integral monodromy
translates. This identifies cap cohomology with the full
primitive invariant lattice. The cusp has twelve unchanged
components and the integral specialization of (B.5.16).
Consequently its direct-image sheaf is exactly the integral
$j_*L_q$, including at every critical point.

\item The three graph-cocycle values again surject onto the
three local coinvariant groups, integrally. Thus $H^2(j_*L_q)=0$,
and its remaining cohomology is the kernel and cokernel
of the exact lattice map
\[
 \begin{gathered}
 V_q\longrightarrow
        \bigoplus_{i=0}^2\operatorname{im}(T_{i,q}-I),\\
           v\longmapsto((T_{i,q}-I)v)_i.
 \end{gathered}
 \tag{E.S53}\label{maps:eq:E.S53}
\]
Use Hermite column bases for the three image lattices;
this retains their integral indices rather than replacing
them by rational spans. The source and target ranks,
and the rank of \eqref{maps:eq:E.S53}, are respectively
\[
 \begin{array}{c|rrrrr}
 q&0&1&2&3&4\\ \hline
 \operatorname{rank}V_q&12&48&72&48&12\\
 \operatorname{rank}\bigoplus\operatorname{im}(T_{i,q}-I)
                      &17&92&126&92&17\\
 \operatorname{rank}(\text{E.S53})&11&47&71&47&11
 \end{array}
 \tag{E.S54}\label{maps:eq:E.S54}
\]
Every nonzero Smith entry is one. Hence $H^0(j_*L_q)$
has rank one and $H^1(j_*L_q)$ is free with ranks
$(6,45,55,45,6)$. The two-column Leray filtration has
free end groups and therefore splits as abelian groups.
This proves the torsion-free assertion and \eqref{maps:eq:E.S44}
integrally. The full integer matrices and the exact
Hermite/Smith check are in
\texttt{maps\_shifted\_cover\_integral.py}.

\item \emph{Why the marked rational images now determine the integral images.}
Each piece $A=E_{12}\times T^3$ in \eqref{maps:eq:E.S29} has free
integral homology. By Poincare--Lefschetz duality,
$H^{k+1}(A,\partial A;\mathbb Z)$ is free. The cokernel
of $H^k(A;\mathbb Z)\to H^k(\partial A;\mathbb Z)$
embeds in that free group, so the boundary image is a
primitive lattice. Its rational span has already been
computed in \eqref{maps:eq:E.S38}. Therefore the actual integral
image is precisely the saturation of that displayed
image in $\bigwedge^k\mathbb Z^8$.

\item Choose saturated column bases $\overline A_k,\overline B_k$
for the two images. The combined matrix
$[\overline A_k\ \overline B_k]$ has the ranks in the
second row of \eqref{maps:eq:E.S40}. All its nonzero Smith entries are
one except that, in each of degrees two and six, exactly
two entries are six. The Mayer--Vietoris exact sequence
has a free kernel term on its right, since all piece
cohomology is free. Thus its abelian extension splits,
and these Smith entries give exactly \eqref{maps:eq:E.S48}. The free
ranks remain \eqref{maps:eq:E.S33}.

\item \emph{The integral variation, including possible torsion components.}
Formula \eqref{maps:eq:E.S37} holds integrally. One can use the singular
prism of the loop of translations instead of differential
forms: the action map $S^1\times T^8\to T^8$ has
cohomology formula
\[
 a\longmapsto 1\otimes a+[S^1]^*\otimes\iota_v a.
\]
The difference between the two collar homotopies is
exactly this loop, so its connecting class is
$\partial\iota_v r(a)$. This proves the integral formula
also on torsion and still gives $(T^*-I)^2=0$.

\item For complete lattice control, take an integral kernel
basis of $[\overline A_k\ -\overline B_k]$. Its first
block, multiplied by $\overline A_k$, is a basis of the
common boundary lattice. Contract it by $v$ and reduce
modulo $[\overline A_{k-1}\ \overline B_{k-1}]$ using
the integral Smith transforms. In degrees four and five
the result has rank one and unique nonzero Smith entry
$72$. In degree three its two possible components
modulo six are both zero; in degree seven the common
boundary image is zero. All the remaining reductions
are zero. These computations are performed, with the
unimodular transforms verified, in
\texttt{maps\_shifted\_fibre\_integral.py}.
In degrees four and five the connecting cokernel is
free and embeds as a primitive summand of $H^k(Y)$,
because the following Mayer--Vietoris kernel is free.
Thus $72$ is the divisibility in the whole cohomology
group, not just in a chosen sublattice.

\item A rank-one integral map of Smith entry $72$ has the
form $72u\ell$, where $u$ and $\ell$ are primitive.
Its square is zero precisely when $\ell(u)=0$.
Extend $u$ to a basis of $\ker\ell$ and choose a vector
on which $\ell$ is one. This gives \eqref{maps:eq:E.S49}. The same
calculation proves identity on every other integral
cohomology group, including its torsion.

\item \emph{The finite-cover groups and the two potentially nontrivial extensions.}
The cover $P_{12n}$ is the mapping torus of $T^n$.
Since $(T^*-I)^2=0$, its variation is $n(T^*-I)$,
with Smith entry $72n$ in degrees four and five.
For every $k$ the integral Wang sequence gives
\[
 0\longrightarrow\operatorname{coker}(T^{n*}-I)_{k-1}
 \longrightarrow H^k(P_{12n};\mathbb Z)
 \longrightarrow\ker(T^{n*}-I)_k\longrightarrow0.
 \tag{E.S55}\label{maps:eq:E.S55}
\]
The ranks immediately give \eqref{maps:eq:E.S50}. Whenever the right
group is free, \eqref{maps:eq:E.S55} splits as an abelian sequence.
In particular degrees four and eight have torsion
$(\mathbb Z/6)^2$, contributed by $H^3(Y)$ and $H^7(Y)$
respectively, while degrees five and six have torsion
$\mathbb Z/(72n)$. Every other such degree has no torsion.

\item The only right groups in \eqref{maps:eq:E.S55} that are not free
occur in degrees three and seven. Their extensions
must not be split by assumption. Instead use the
nondegenerate integral torsion linking pairing of the
closed oriented ten-manifold $P_{12n}$: cohomological
torsion in degree $k$ is Pontryagin dual to that in
degree $11-k$. The already fixed torsion in degree
eight therefore forces the degree-three torsion to
be $(\mathbb Z/6)^2$, and the fixed degree-four torsion
forces the same answer in degree seven. The free
ranks are already known. This resolves precisely
the two possibly absorbing extensions. The universal
coefficient theorem now translates these cohomology
groups to \eqref{maps:eq:E.S51}. The order $72n$ in the middle
groups distinguishes the covers for distinct $n$.
\end{enumerate}
\end{proof}

\begin{theorem}[the full complement module and its peripheral map]\label{maps:src:E.2.16}
\StatementRecord{maps:src:E.2.16}{proved}{Historical events: Claude Fable 5.2 (Anthropic), effort high, 16 November 2026: carried as background without fault found. GPT-6 Astra (OpenAI), ChatGPT Pro mode, effort Pro, S6\_09\_PAPERS, 17 November 2026: targeted proof check. Also checked in this writing. Pre-continuation writing events: GPT-6 Astra (OpenAI), ChatGPT Pro mode, effort Pro, S6\_10\_PAPERS, 18 November 2026: targeted proof check. Continuation metadata events: GPT-6 Astra (OpenAI), ChatGPT Pro mode, effort Pro, S6\_10\_PAPERS continuation, 19 November 2026: used and rechecked in this run; status-vocabulary reconciliation; metadata reconciliation only; no new mathematical check. \emph{Pre-continuation recorded proof scope.} GPT-6 Astra (OpenAI, ChatGPT Pro mode; recorded reviewer /root/cdp\_audit), effort Pro, 17 November 2026; Claude Fable 5.2 (Anthropic), effort high, 16 November 2026; also checked in this writing by GPT-6 Astra (OpenAI), ChatGPT Pro mode, effort Pro, 18 November 2026}{the hypotheses, notation and definitions stated here; \cref{maps:eq:E.S1}; \cref{maps:src:E.2.17}; \cref{maps:src:E.2.21}; \cref{maps:src:E.2.12}; \cref{maps:eq:E.C5}; \cref{maps:eq:E.C8}; \cref{maps:eq:E.C11}; \cref{maps:eq:E.C7}; \cref{maps:eq:E.C14}; \cref{maps:eq:E.C12}; \cref{maps:eq:E.C16}; \cref{maps:eq:E.C17}; \cref{maps:src:E.2.6}}


Let $K=\pi_1(E)$ have the marking \eqref{maps:eq:E.S1}, and put
\[
 R=\mathbb ZK,\qquad Q=K/\langle h\rangle=C_3*C_4,
 \qquad \overline R=\mathbb ZQ.
\]
Use right deck modules; conversion to left modules applies the
involution $g\mapsto g^{-1}$ to every coefficient. Then
\[
 \pi_2(E)\simeq R\oplus\overline R^{\,2}\oplus L,
 \qquad
 L=\langle x,t\mid t(a-1)=t(b-1)=0,
                         \ x(1+b)+3t=0\rangle_{\overline R}.
 \tag{E.C1}\label{maps:eq:E.C1}
\]
Equivalently, replacing $x$ by $x+t$ and eliminating $t$ gives
\[
 L\simeq\overline R\big/
       \big((b^2-1)\overline R+(1+b)(a-1)\overline R\big).
 \tag{E.C2}\label{maps:eq:E.C2}
\]
There is a nonsplit exact sequence
\[
 0\longrightarrow\mathbb Z\langle t\rangle
  \longrightarrow L\longrightarrow
  \mathbb Z_{-}\otimes_{\mathbb ZC_4}\overline R
  \longrightarrow0,
 \tag{E.C3}\label{maps:eq:E.C3}
\]
where $b$ acts as $-1$ on $\mathbb Z_{-}$ and the first term
has trivial $Q$ action. Its extension class is the nonzero
element of an order-two extension group. In the connected
lift of the peripheral $T^5$ to the universal cover of $E$,
the degree-two classes $uw,ud,wd$ map respectively to
\[
                           t,\quad0,\quad0.
 \tag{E.C4}\label{maps:eq:E.C4}
\]
Thus the central torsion summand of the module is annihilated
by $h-1$ itself. Formula \eqref{maps:eq:E.C1} is a full integral module
calculation, including the central value $h=1$.  The complete
splice extension is determined separately in \ref{maps:src:E.2.17}, and the
smooth type by full framed handles in \ref{maps:src:E.2.21}; neither conclusion
is inferred from this module calculation alone.

There is also a specified chain-level peripheral comparison.
For the splitting constructed below, let $e$ generate the
free $R$ summand and let $y,z$ generate the two copies of
$\overline R$. The projection of the peripheral degree-two
chain map to \eqref{maps:eq:E.C1} is
\[
 \begin{array}{c|rrrrrr}
 \text{cell}&gu&gw&gd&uw&ud&wd\\ \hline
 \text{value}&0&0&6t&t&0&0
 \end{array}
 \tag{E.C19a}\label{maps:eq:E.C19a}
\]
and, writing $\Theta$ for the outer base circle with
group image $\mu=ab$, use the cone convention: the
notation $v\Theta$ denotes the suspended cell
$\Theta\times v$, with the interval factor first.
In this convention,
\[
 \begin{array}{c|l}
  g\Theta&eh-2x-6yb+z(5-2b)\\
  u\Theta&x-yb+zb\\
  w\Theta&-y(1+b)+z\\
  d\Theta&x+2t+z(b-1).
 \end{array}
 \tag{E.C19b}\label{maps:eq:E.C19b}
\]
These are values on chains, not assertions that all
the displayed cells are cycles in the cover. In
particular the free coefficient $h$ on $g\Theta$
is nonzero even though the free summand receives
zero from the lifted peripheral homology in \eqref{maps:eq:E.C4}.
The original cone and its full peripheral map, not
only the map on homology, are part of the certificate.

\end{theorem}


\begin{proof}
\begin{enumerate}[label=\arabic*.,leftmargin=*]
\item We give the local chain comparison before the global
calculation, because the ordinary homology maps alone would
not determine \eqref{maps:eq:E.C1}.

\item \emph{The cusp boundary and its derived attachment.}
Write $A=\mathbb Z[h,h^{-1}]$ and $q=h-1$. Let
$\widehat W$ be the gamma cover used in \ref{maps:src:E.2.12}, and let
$\widehat M$ be the corresponding boundary cover. Ordinary
chains with the deck action give finite free $A$ chain
models. The homology of $\widehat W$ is
\[
 H_0=H_1=A/(q),\quad H_2=A\oplus(A/(q))^2,
 \quad H_3=0,\quad H_4=A.
 \tag{E.C5}\label{maps:eq:E.C5}
\]
For the boundary, start with the torus chain complex
$C= A\otimes\bigwedge\langle g,u,w,d\rangle$,
whose differential contracts $g$ and multiplies by $q$.
The cellular monodromy has degree-one matrix
\[
 F(M_0)g=g-hd,\quad F(M_0)u=u+w,
 \quad F(M_0)w=w,\quad F(M_0)d=d.
 \tag{E.C6}\label{maps:eq:E.C6}
\]
Its higher matrices are exterior powers. The mapping-torus
complex is the cone of $I-F(M_0)$, with differential
\[
 d(c,e)=\big(dc+(I-F(M_0))e,-de\big).
 \tag{E.C7}\label{maps:eq:E.C7}
\]
For example $d(g\theta)=hd-q\theta$ and
$d(u\theta)=-w$. Cancelling these two unit pairs leaves
the degree-one classes $u$ and $\theta$, with resolutions
$q$ and a unit times $q^2$. In degree two the four surviving
classes are $uw,ud,w\theta,d\theta$, each with resolution
$q$. In particular
\[
 H_0(\widehat M)=A/(q),\quad
 H_1(\widehat M)=A/(q^2)\oplus A/(q),\quad
 H_2(\widehat M)=(A/(q))^4.
 \tag{E.C8}\label{maps:eq:E.C8}
\]
Geometrically this is the Heisenberg three-manifold times
a circle: the deck transformation shears the meridian
by the $d$ circle. The homology map to \eqref{maps:eq:E.C5} is the
identity in degree zero, projection to $u$ in degree one,
and the primitive projection
$(uw,ud,w\theta,d\theta)\mapsto(U,D,0,0)$ in degree two.
The last assertion uses the actual specialization
(B.5.16) and the literal collapsed meridian tori
(B.5.17), rather than ranks alone.

\item The missing extension is fixed by a relative intersection
calculation. Let $S_i$ denote the normalized hexagon
component with gamma index $i$. In its standard Picard
basis $H,E_1,E_2,E_3$, the six boundary curves are
\[
 E_1,\ H-E_1-E_2,\ E_2,\ H-E_2-E_3,
                       \ E_3,\ H-E_1-E_3.
 \tag{E.C9}\label{maps:eq:E.C9}
\]
The first opposite pair is self-glued in the surviving
$u$ direction; the other two pairs join $S_i$ to its
two gamma neighbours. Put $C_0=-H+E_1$, the primitive
generator of the free kernel of the normalization matrix
$B(h)$ in \ref{maps:src:E.2.12}. The divisor $S_i$ restricts on its
own normalization to $2C_0$: the total central divisor
is principal, and the two self-gluing preimages cancel
their terms in the normal restriction. It restricts
to $-C_0$ on each neighbouring component and to zero
on every other component. Thus the free part of the
equivariant intersection map is
\[
 H_4(\widehat W)=A\longrightarrow
 H_4(\widehat W,\widehat M)\longrightarrow A,
 \qquad 2-h-h^{-1}=-h^{-1}q^2.
 \tag{E.C10}\label{maps:eq:E.C10}
\]
Poincare--Lefschetz duality over $A$, with its usual
involution, gives the same coefficient, up to a unit,
for the free part of $H_2(\widehat W)\to
H_2(\widehat W,\widehat M)$. The latter relative group
is $A$: the universal-coefficient sequence in degree
four has $\operatorname{Hom}_A(H_4,A)=A$, zero $H_3$,
and no higher extension terms. Indeed every module
in \eqref{maps:eq:E.C5} has a projective resolution of length at
most one. The relative sequence therefore includes
\[
 0\longrightarrow A\xrightarrow{\,\text{unit}\ q^2\,}A
  \longrightarrow A/(q^2)\longrightarrow0.
 \tag{E.C11}\label{maps:eq:E.C11}
\]
Consequently the boundary map has a unit component in
$\operatorname{Ext}^1_A(A/(q^2),A)=A/(q^2)$.
Every unit of $A/(q^2)=\mathbb Z[\varepsilon]/(\varepsilon^2)$
lifts to a unit $\pm h^n$ of $A$.

\item Here is why this determines the required low chain map
without discarding a peripheral invariant. All the
modules in \eqref{maps:eq:E.C5} and \eqref{maps:eq:E.C8} have projective dimension
at most one. Their complexes split in the derived
category into their two-term projective resolutions:
the obstructions to splitting consecutive truncations
are extension groups in degree at least two and vanish.
In degrees relevant to $H_2$ of a mapping cone, a map
between these split complexes consists of the displayed
homology maps and extension components from $H_i$ to
$H_{i+1}$. Normalize the component in \eqref{maps:eq:E.C11} by a unit
on the free $H_2$ summand of the target. Its components
to the two torsion $H_2$ summands can then be removed
by target maps $A\to(A/(q))^2$: reduction of the unit
modulo $q$ is $\pm1$. The remaining components from
the $u$ summand are removed by target degree-one
extension maps from $H_1(\widehat W)$ to
$H_2(\widehat W)$. Finally the degree-zero component
is removed using the isomorphism on $H_0$ and a target
extension map from $H_0$ to $H_1$.

\item Each change just described is a unit on the diagonal
of the split complex and is triangular with extension
entries off the diagonal. It is therefore invertible.
Replacing the split resolutions by the original finite
free complexes represents it by a chain homotopy
equivalence. No realization by a diffeomorphism or a
map of the cusp space is claimed or needed. The change
is applied to the cusp target and its entire incoming
boundary map simultaneously. The other vertices and
all maps into the regular core are unchanged. It thus
gives a chain homotopy equivalence of the global
mapping-cone complex that is the identity on the
external peripheral complex. In particular the
calculation of \eqref{maps:eq:E.C4} retains the actual boundary
map. Tensoring with $R$ preserves these comparisons,
since $R$ is free over $A$.

\item One resulting cusp model in degrees zero through three
has ranks $(1,2,4,2)$, generators
\[
 (1);\quad(g,u);\quad(gu,E,U,D);\quad(R_U,R_D),
\]
and differentials $dg=q$, $d(gu)=qu$,
$dR_U=qU$, $dR_D=qD$. The map \eqref{maps:eq:E.C7} to this model
is, through degree two,
\[
 \begin{array}{c|ccccc}
  &g&u&w&d&\theta\\ \hline
  i_1&g&u&0&0&0
 \end{array},\qquad
 \begin{array}{c|rrrrrr}
  &gu&gw&gd&uw&ud&wd\\ \hline
  i_2&gu&-hD&E&U&D&0
 \end{array},
 \qquad i_2(v\theta)=0.
 \tag{E.C12}\label{maps:eq:E.C12}
\]
The signs and the factor $h$ satisfy the chain equation
directly. At $h=1$, choose
$E=-\overline C_1-\overline C_2$,
$U=\overline C_2+\overline C_3$, and
$D=\overline C_2$; this is precisely (B.5.16).

\item \emph{Finite caps and the global complex.}
For any marked monodromy
$A=\left(\begin{smallmatrix}1&0\\t&D\end{smallmatrix}\right)$
use the chain lift
$F(A)=\left(\begin{smallmatrix}1&0\\ht&D\end{smallmatrix}\right)$
and its exterior powers. These lifts multiply exactly
and commute with contraction by $qg^*$.

\item This exterior lift represents the actual equivariant
torus map. The gamma cover is
$\mathbb R_g\times T^3$, and the marked linear map is
$(g,v)\mapsto(g,Dv+gt)$. At every integer value of $g$
the translation is integral, so the vertex torus map
is exactly $D$. In a single unit strip, concentrate
the loop of translations $v\mapsto v+st$ at its
terminal endpoint by an equivariant cellular
approximation relative to the two vertex tori.
The translation prism on the minimal cellular complex
of $T^3$ is exterior multiplication by $t$: its value
on each product torus is the homology class of the
literal translation sweep, and all differentials
in that minimal torus complex are zero. The terminal
vertex has deck coefficient $h$. Consequently
\[
 v\longmapsto\bigwedge D(v),\qquad
 g\wedge v\longmapsto
 g\wedge\bigwedge D(v)+ht\wedge\bigwedge D(v),
 \tag{E.C20}\label{maps:eq:E.C20}
\]
which is exactly the exterior lift above. There is
only this one strip parameter. Composing two such
actual maps adds the translation loops by
$t_A+D_At_B$, the same strict multiplication law
as $F(A)F(B)=F(AB)$. Thus no uncomputed higher
coherence is being substituted for the torus map.
A regular
boundary with deck image $s$ is the cone of $sI-F(A)$.
This convention gives the group relation
$s\lambda s^{-1}=A\lambda$.

\item At finite cap $j$, take its adapted first fibre vector
$n_jv_j$, where
\[
 (m_1,n_1,v_1)=(3,1,(1,2,-4,0)),\qquad
 (m_2,n_2,v_2)=(4,-1,(-1,-3,3,0)).
 \tag{E.C13}\label{maps:eq:E.C13}
\]
It has gamma coordinate one and is fixed by $A_j$.
The other three coordinates form $L_0=\langle u,w,d\rangle$,
with monodromy $D_j=A_j|_{L_0}$. The cap is the
mapping torus of this three-torus, and its complex has
differential $sI-\bigwedge^rD_j$ on the swept cells.
The fibre-to-cap map is the identity on $\bigwedge L_0$;
on the adapted long-circle factor it is the integral
norm matrix $N_{m_jn_j}(s,D_j)$ determined by
\[
 (sI-D_j)N_{m_jn_j}(s,D_j)=(h-1)I,
 \qquad s=a\text{ or }b.
 \tag{E.C14}\label{maps:eq:E.C14}
\]
For positive $n_j$ the norm is
$\sum_{r=0}^{m_j-1}s^rD_j^{m_j-1-r}$, and for
$n_j=-1$ it is minus $s^{-m_j}$ times that sum.
In phase degree $r$, use
$N_{m_jn_j}(s,\bigwedge^rD_j)$, rather than the
exterior power of the norm matrix.
The boundary base-circle generator maps to the cap circle, while
the cell containing both that meridian and the adapted
long circle maps to zero. These are actual winding
maps in the mapping-torus coordinates. The linear
three-torus return fixes its origin; hence there is
no additional translation prism in this comparison.

\item Form the core over a bouquet of three circles with
deck images $a,b,1$ and monodromies $A_1,A_2,M_0$.
Attach the two finite complexes by \eqref{maps:eq:E.C14} and the
cusp by \eqref{maps:eq:E.C12}, using their double mapping cylinders.
Through degree three the resulting free $R$ complex
has ranks $(4,20,49,62)$. Its entries and every unit
cancellation are given exactly in
\texttt{code/maps\_complement\_module.py} and its
JSON output. The program uses the integral normal
form $h^r$ times an alternating word in
$a,a^2$ and $b,b^2,b^3$; no finite group quotient or
floating-point approximation is used.

\item For completeness the final reduction can be checked
without reading a long matrix. Cancelling degree-zero
and degree-one unit pairs leaves ranks $(1,2,34,62)$,
with $d_1=(a-1,b-1)$. The only nonzero columns of
$d_2$ are $F$ and $F(b-1)$, where
\[
 F=\begin{pmatrix}1+a+a^2\\h(1+b+b^2+b^3)\end{pmatrix}.
 \tag{E.C15}\label{maps:eq:E.C15}
\]
The first component of $F$ is injective on $R$:
multiplication by $a-1$ converts it to $h-1$, and
$R$ is free over the Laurent ring $A$. Thus removing
the $F$ column after subtracting its multiple gives
a free basis of $\ker d_2$ of rank33. The corresponding
row of $d_3$ is zero by the chain equation. Unit
reductions of the remaining presentation leave five
generators $(x,t,y,z,e)$. One generator $e$ occurs
in no relation. Three pure $t$ relations are
\[
 tq=0,\qquad t(1-a^{-1})=0,
                  \qquad t(b-a^{-1})=0.
 \tag{E.C16}\label{maps:eq:E.C16}
\]
They generate exactly the augmentation relations on
$t$: multiply the second by $a$ to get $t(a-1)=0$,
and subtract it from the third to get $t(b-1)=0$.
Every other pure $t$ coefficient has augmentation
zero and is consequently redundant. Reducing the
remaining $t$ coefficients by augmentation, three
relations give successively $yq=0,zq=0,xq=0$.
The final three relations are
\[
 x(1+b)+3t=0,\quad x(b+b^2)+3t=0,
                   \quad x(1+b+b^2+b^3)+6t=0.
 \tag{E.C17}\label{maps:eq:E.C17}
\]
The last two are respectively the first multiplied
by $b$, and the first multiplied by $1+b^2$.
These identities prove \eqref{maps:eq:E.C1}. The program verifies
the original chain equations, all Schur complements,
and the coefficient list used in \eqref{maps:eq:E.C16}--\eqref{maps:eq:E.C17}.

\item The invariant subgroup generated by $t$ injects in
$L$. One way to see this is to start with the exact
sequence of right $\mathbb ZC_4$ modules
\[
 0\longrightarrow\mathbb Z\longrightarrow
 \mathbb Z[C_4/C_2]\longrightarrow\mathbb Z_-\longrightarrow0,
\]
where the first map is the sum of the two cosets.
Induce to $Q$ and push out its first term by the
augmentation to the trivial module $\mathbb Z$.
The resulting presentation is \eqref{maps:eq:E.C1}, after the
change $x\mapsto x+t$ which replaces the coefficient
$3$ by $1$. Pushout preserves the displayed injection.
Restriction to $C_4$ shows nonsplitting: an invariant
correction of the lift $x$ changes this odd coefficient
by an even number. Equivalently the standard cyclic
resolution computes
$\operatorname{Ext}^1_{\mathbb ZC_4}(\mathbb Z_-,\mathbb Z)
=\mathbb Z/2$. Induction--restriction gives the same
extension group for \eqref{maps:eq:E.C3}.

\item Finally insert the actual fibre cycles $uw,ud,wd$ in
the core before any cancellation and carry their
coordinate columns through the same basis changes.
The five-generator presentation gives respectively
\[
 t,\qquad t(a^{-1}-a),\qquad t(a^{-1}-1).
 \tag{E.C18}\label{maps:eq:E.C18}
\]
Relations \eqref{maps:eq:E.C16} reduce these to \eqref{maps:eq:E.C4}. The external
boundary cover has homotopy type $T^3$ and these are
its three degree-two generators. This also checks
the ordinary coefficient independently: after taking
$K$ coinvariants, $2x+3t=0$, so $t$ is twice a
primitive class, exactly the coefficient2 in \ref{maps:src:E.2.6}.

\item To retain the remaining peripheral data, use the
boundary word $abc$, where the cusp generator $c$
has deck image1. On swept fibre chains its map to
the three core edges has blocks
\[
                 bI,\qquad F(A_1),\qquad F(A_1)F(A_2).
 \tag{E.C21}\label{maps:eq:E.C21}
\]
The identity
\[
 (aI-F(A_1))b+F(A_1)(bI-F(A_2))
       +F(A_1)F(A_2)(I-F(M_0))=(ab-1)I
\]
checks its chain equation. This is the usual prism
subdivision of the actual peripheral word, with
its given ribbon marking. Tracking all ten
degree-two columns through the same reductions,
and using only $t(a-1)=t(b-1)=0$ and
$xq=yq=zq=0$, gives \eqref{maps:eq:E.C19a}--\eqref{maps:eq:E.C19b}.
The unreduced columns are saved in the JSON
certificate. In particular no assertion that a
homologically invisible peripheral component is
zero enters this calculation.
\end{enumerate}
\end{proof}

\begin{theorem}[the full shifted homotopy module]\label{maps:src:E.2.17}
\StatementRecord{maps:src:E.2.17}{proved}{Historical events: Claude Fable 5.2 (Anthropic), effort high, 16 November 2026: carried as background without fault found. GPT-6 Astra (OpenAI), ChatGPT Pro mode, effort Pro, S6\_09\_PAPERS, 17 November 2026: targeted proof check. Also checked in this writing. Pre-continuation writing events: GPT-6 Astra (OpenAI), ChatGPT Pro mode, effort Pro, S6\_10\_PAPERS, 18 November 2026: targeted proof check. Continuation metadata events: GPT-6 Astra (OpenAI), ChatGPT Pro mode, effort Pro, S6\_10\_PAPERS continuation, 19 November 2026: used and rechecked in this run; status-vocabulary reconciliation; metadata reconciliation only; no new mathematical check. \emph{Pre-continuation recorded proof scope.} GPT-6 Astra (OpenAI, ChatGPT Pro mode; recorded reviewer /root/cdp\_audit), effort Pro, 17 November 2026; Claude Fable 5.2 (Anthropic), effort high, 16 November 2026; also checked in this writing by GPT-6 Astra (OpenAI), ChatGPT Pro mode, effort Pro, 18 November 2026}{the hypotheses, notation and definitions stated here; \cref{maps:src:E.2.16}; \cref{maps:src:E.2.13}; \cref{maps:eq:E.P7}; \cref{maps:eq:E.P8}; \cref{maps:eq:E.P11}; \cref{maps:src:E.2.9}; \cref{maps:src:E.2.10}; \cref{maps:eq:E.P9}; \cref{maps:eq:E.P14}; \cref{maps:eq:E.P16}; \cref{maps:src:E.2.11}}


Use the right-module convention of \ref{maps:src:E.2.16}. Write
\[
 \begin{gathered}
 R=\mathbb ZG,\qquad q_i=h_i-1,\qquad S=\mathbb Z\Gamma,\\
 \Gamma=G/\langle h_1,h_2\rangle
  =\langle a,b,c,d\mid a^3=b^4=c^3=d^4=abcd=1\rangle.
 \end{gathered}
 \tag{E.P1}\label{maps:eq:E.P1}
\]
where $a,b$ belong to the first complement, $c,d$ to
the second, and $\mu=ab=(cd)^{-1}$. The full integral
second-homotopy module is
\[
 \pi_2(P)\simeq R/q_2R\ \oplus\ R/q_1R
                    \ \oplus\ T_L\oplus T_R.
 \tag{E.P2}\label{maps:eq:E.P2}
\]
The last two modules are inflated from the following
explicit six-generator presentations over $S$.
Put $N_a=1+a+a^2$, $N_b=1+b+b^2+b^3$, and define
$N_c,N_d$ similarly.

For $T_L$ use generators $p_i,r_i$ for $0\leq i\leq2$,
put $v_i=p_i-r_i$, and define
\[
 \begin{split}
 \Omega_L={}&v_0(1+d-d^2-d^3)
       +v_1(d^3-d)+v_2(d^2+d^3),\\
 p_iN_a={}&0,\qquad r_iN_b=0\quad(0\leq i\leq2),\\
 \Omega_L(a-1)={}&\Omega_L(b-1)
       =\Omega_L(c-1)=\Omega_L(d-1)=0.
 \end{split}
 \tag{E.P3}\label{maps:eq:E.P3}
\]
For $T_R$ use generators $B_i,C_i$, put $w_i=B_id+C_i$,
and define
\[
 \begin{split}
 \Omega_R={}&w_0a(1-b-b^2+b^3)
              +w_1a(b^2-1)-w_2a(1+b^3),\\
 B_iN_c={}&0,\qquad C_iN_d=0\quad(0\leq i\leq2),\\
 \Omega_R(a-1)={}&\Omega_R(b-1)
       =\Omega_R(c-1)=\Omega_R(d-1)=0.
 \end{split}
 \tag{E.P4}\label{maps:eq:E.P4}
\]
These are complete presentations: six generators and
ten relations in each block. In particular
\[
 q_1q_2\pi_2(P)=0,\qquad
 \ker(q_1:\pi_2(P)\to\pi_2(P))
  \cap\ker(q_2:\pi_2(P)\to\pi_2(P))=T_L\oplus T_R.
 \tag{E.P5}\label{maps:eq:E.P5}
\]
This improves the earlier squared annihilator and
specifies its entire central-intersection summand.
Its coinvariants are $(\mathbb Z/12)^6$, so \eqref{maps:eq:E.P2}
recovers $\pi_2(P)_G=\mathbb Z^2\oplus(\mathbb Z/12)^6$.

Let $\Pi=\pi_1(Y)$ and let $\Sigma\simeq\pi_1(\Sigma_6)$
be its quotient by $z=h_1h_2$, as in \ref{maps:src:E.2.13}. Then
\[
 \pi_2(Y)\simeq(\mathbb Z\Pi)^{24}
          \oplus\operatorname{Res}^{\Gamma}_{\Sigma}(T_L\oplus T_R),
 \qquad
 \pi_2(Y)_\Pi\simeq\mathbb Z^{110}\oplus(\mathbb Z/6)^2.
 \tag{E.P6}\label{maps:eq:E.P6}
\]
The surface-group summand is inflated to $\Pi$, so
$z$ acts trivially on it. Formulas \eqref{maps:eq:E.P2}--\eqref{maps:eq:E.P6}
do not identify the smooth type of $P$ or $Y$.

\end{theorem}


\begin{proof}
\begin{enumerate}[label=\arabic*.,leftmargin=*]
\item \emph{A small chain presentation with the actual peripheral map.}
The calculation in \ref{maps:src:E.2.16} permits the following model
through degree three for the first complement. Its
bases have ranks $(1,2,6,6)$, with degree-one basis
$(a,b)$ and degree-two basis $(r,x,t,y,z,e)$.
The first two differentials are
\[
 d_1=(a-1,b-1),\qquad
 d_2(r)=\begin{pmatrix}N_a\\h_1N_b\end{pmatrix},
 \quad d_2(x)=d_2(t)=d_2(y)=d_2(z)=d_2(e)=0.
 \tag{E.P7}\label{maps:eq:E.P7}
\]
The six degree-three columns are the six relations
\[
 t(a-1),\quad t(b-1),\quad xq_1,\quad yq_1,\quad zq_1,
                         \quad x(1+b)+3t.
 \tag{E.P8}\label{maps:eq:E.P8}
\]
The degree-two norm $N_a$ in \eqref{maps:eq:E.P7} is written with
the same expression as in \eqref{maps:eq:E.P3}, but here $a^3=h_1$,
not $1$. The second complement is the identical
model with $(a,b,h_1)$ replaced by $(c,d,h_2)$.

\item The ordinary one-relator complex $C_K$ is the subcomplex
on degrees zero and one and the cell $r$. It is an
aspherical model: it is the double mapping cylinder
of the degree-three and degree-four maps of circles,
and its universal cover is a tree of contractible
lines and strips. Equivalently it is the standard
Seifert spine of the torus-knot exterior. Projection
of \eqref{maps:eq:E.P7}--\eqref{maps:eq:E.P8} onto $C_K$ retains the actual
peripheral map. In degree one this map is
\[
 g\longmapsto(N_a,0),\qquad
 u,w,d\longmapsto0,\qquad
 \Theta\longmapsto(b,1).
 \tag{E.P9}\label{maps:eq:E.P9}
\]
Its degree-two $r$ coefficient is $b-1$ on the swept
cell $g\Theta$ and zero on the other nine cells.
Its remaining components are exactly (E.C19).
These additional columns were tracked in the
\ref{maps:src:E.2.16} certificate, rather than chosen from their
ordinary homology values.

\item Here and below the small complex is only a
presentation through the degree needed for $H_2$.
Replacing redundant degree-three columns by the six
generators of their image in \eqref{maps:eq:E.P8} changes no
degree-two cycle or boundary. Reducing a peripheral
degree-two column modulo that image changes it by
a chain homotopy with values in degree three.
Products with a torus and the mapping cones below
therefore retain the same degree-two homology.
For the comparison with an aspherical target one
may first keep the original, full complex: its
projection to $C_K$ extends by zero in degrees at
least three, since the Fox column in \eqref{maps:eq:E.P7} is
injective and hence the $r$ row of $d_3$ is zero.
Only after taking the kernel is its image in degree
two shortened as in \eqref{maps:eq:E.P8}. Thus no false
acyclicity of a truncated target is used.

\item \emph{The exact kernel complex for the splice.}
Let $T(s)$ be the four-torus chain complex whose
first circle has deck image $s$ and whose other
three circles have trivial image. Its differential
contracts that first circle and multiplies by $s-1$.
The two vertex complexes for the actual splice are
\[
 V_1=C_{E_1}\otimes T(h_2),\qquad
 V_2=C_{E_2}\otimes T(h_1).
 \tag{E.P10}\label{maps:eq:E.P10}
\]
We have reversed the two product factors in the
second vertex only for notation; the Koszul signs
below account for that reversal. The boundary
complex $B$ is the nine-torus complex with ordered
coordinates
\[
 (g_1,u_1,w_1,d_1,g_2,u_2,w_2,d_2,\Theta),
 \quad\text{deck images }(h_1,1,1,1,h_2,1,1,1,\mu).
 \tag{E.P11}\label{maps:eq:E.P11}
\]
Its differential is the usual exterior contraction
by those deck images minus one.

\item For clarity all source signs can be specified
without an implicit boundary convention. Write
a boundary product cell as $\alpha\wedge\beta$
with an optional final $\Theta$, where $\alpha$
uses the first four coordinates and $\beta$ the
second four. Put $\epsilon=1$ when $\Theta$ occurs
and $0$ otherwise. Its first vertex map applies
the \ref{maps:src:E.2.16} peripheral map to $(\alpha,\Theta)$,
keeps $\beta$, and has sign
$(-1)^{\epsilon(|\alpha|+|\beta|)}$.
The second map applies that peripheral map to
$(\beta,\Theta_2)$, keeps $\alpha$, and has sign
$(-1)^{|\alpha|(|\beta|+\epsilon)+\epsilon|\beta|}$.
If $\epsilon=1$, multiply its coefficient on the
right by $-\mu$. This is the exact cellular map
of $\Theta_1\mapsto\Theta_2^{-1}$, since
$\mu_2=cd=\mu^{-1}$. The suspended cells in (E.C19)
have the interval first; these signs convert from
the ordered exterior cells in \eqref{maps:eq:E.P11}.
Denote the two maps by $j_1,j_2$.

\item The chain complex of the splice is the cone of
$(j_1,-j_2):B\to V_1\oplus V_2$.
Project each $C_E$ to $C_K$ and each $T(s)$ to its
first circle, and project $B$ to its three coordinates
$g_1,g_2,\Theta$. The result is the complex of the
aspherical graph of spaces
\[
 (K(K_1,1)\times S^1)\ \cup_{T^3}\
                         (K(K_2,1)\times S^1).
 \tag{E.P12}\label{maps:eq:E.P12}
\]
It is the model $N$ of \ref{maps:src:E.2.9}. Both peripheral maps
on groups are injective by \ref{maps:src:E.2.10}; their common
group is $\langle h_1,h_2,\mu\rangle=\mathbb Z^3$.
Consequently its universal cover is a Bass--Serre
tree of contractible vertex and edge spaces and
is contractible. Its chain complex has ranks
$(2,7,9,5,1)$, including its degree-four cell.

\item The indicated projection is a degreewise split
surjection of the full chain complexes. The
peripheral Fox coefficient in \eqref{maps:eq:E.P9} makes the
boundary square commute; higher degrees commute
as well by the product construction and the
zero higher groups of $C_K$. Let $D$ be its kernel.
The target has zero positive-degree homology, so
the long exact sequence gives
$H_2(D)\simeq H_2(\widetilde P)=\pi_2(P)$.
Using the small models for the relevant images,
the ranks of $D$ through degree three are
\[
                         (0,6,40,123).
 \tag{E.P13}\label{maps:eq:E.P13}
\]
Its six degree-one generators are the unused
regular-factor circles at the two vertices.
The six shifted boundary circles $u_i,w_i,d_i$
map to them by unit entries, because each circle
maps to zero in its own complement and identically
to the opposite regular torus. Cancel these six
unit pairs. This leaves a presentation of $H_2(D)$
with34 generators and123 relations, with no
uncomputed kernel.

\item Here is also an explicit block definition of that
matrix. In each $V_i$ keep the cells killed by
projection to $C_{K_i}\otimes C(S^1)$; their ranks
in degrees two and three are17 and45. In $B_2$
keep the33 cells containing at least one of
$u_1,w_1,d_1,u_2,w_2,d_2$. In those ordered bases
the matrix is
\[
 \mathcal M=
 \begin{pmatrix}
 (d_{V_1})_3&0&(j_1)_2\\
 0&(d_{V_2})_3&-(j_2)_2
 \end{pmatrix}:
 R^{45}\oplus R^{45}\oplus R^{33}
             \longrightarrow R^{17}\oplus R^{17}.
 \tag{E.P14}\label{maps:eq:E.P14}
\]
Every entry is specified by \eqref{maps:eq:E.P7}--\eqref{maps:eq:E.P11} and
(E.C19); \eqref{maps:eq:E.P14} itself is a complete finite
presentation of the desired module. The formula
includes the actual derived peripheral comparison.
In particular a splitting of the two cusp terms
has not been assumed from their localized homology.

\item \emph{Reduction to the compact presentation.}
The exact certificate
\path{code/III/maps_splice_module.py} checks both
peripheral chain equations and $d_2d_3=0$, cancels
the indicated unit pairs, and reduces \eqref{maps:eq:E.P14} to
16 generators and52 nonzero columns. It computes
in the actual amalgam
\[
 G=(K_1\times\mathbb Z)\mathop{*}_{\mathbb Z^3}
                                      (\mathbb Z\times K_2),
 \tag{E.P15}\label{maps:eq:E.P15}
\]
not in a finite quotient. To specify its normal
form, the torus-knot factor uses $h^r$ times an
alternating word in $a,a^2$ and $b,b^2,b^3$.
Modulo the peripheral subgroup, strip initial
$ab$ and $b^3a^2=(ab)^{-1}$, allowing their last
letter to be a partial power; choose $b^3$ in
the sole length-one tie with $a$. Inverting these
left coset representatives gives right coset
representatives. The usual alternating amalgam
normal form, followed by a final element of
$\mathbb Z^3$, is unique. These elementary word
operations are implemented exactly, including
both central exponents.

\item In the resulting16-generator presentation two
generators occur only in their respective single
relations $q_2e_1=0$ and $q_1e_2=0$. They give the
first two direct summands of \eqref{maps:eq:E.P2}. For the other
14 generators, two generators $t_0,t_1$ have
trivial action of $a,b,c,d$. There are three
generators initially killed by $q_2$ whose pure
relations include $h_1N_b$; the identity
$h_1N_b(b-1)=1-h_1$ kills them by $q_1$ as well.
There are three initially killed by $q_1$ with
pure $N_c$ relations, and three more with pure
$h_2N_d$ relations. They are killed by $q_2$.
The remaining generators already have both
central relations. Thus the whole14-generator
summand is a module over $S$, without imposing
an additional quotient. The pure norm columns
are exact differences of the recorded columns
and are checked before setting either central
element equal to1.

\item For a readable final reduction, the14 generators
separate into two seven-generator blocks. The
first has generators $U_0,U_1,U_2,X,Y,Z,t_1$,
with $t_1$ invariant, $U_iN_b=0$, and
\[
 \begin{split}
 (U_0+X-Yd+Zd)N_a&=0,\\
 (U_1-Y(1+d)+Z)N_a&=0,\\
 (U_2+X+2t_1+Z(d-1))N_a&=0,\\
 X(1+d)+3t_1&=0.
 \end{split}
 \tag{E.P16}\label{maps:eq:E.P16}
\]
The second has generators $B_i,C_i,t_0$, with
$B_iN_c=C_iN_d=0$, $t_0$ invariant, and the single
additional relation $t_0=\Omega_R$. These formulas
follow by collecting the displayed16-generator
columns; the certificate retains every column
and its unit reductions.

\item In the first block put
\[
 v_0=X-Yd+Zd,\quad v_1=-Y(1+d)+Z,
 \quad v_2=X+2t_1+Z(d-1),\qquad
 r_i=U_i,\quad p_i=U_i+v_i.
 \tag{E.P17}\label{maps:eq:E.P17}
\]
The linear part of this change on $(X,Y,Z)$ is
\[
 \begin{pmatrix}1&0&1\\-d&-1-d&0\\d&1&d-1\end{pmatrix},
 \quad\det=1,
\]
over $\mathbb Z[d,d^{-1}]$. It is therefore an
integral group-ring basis change. Its inverse gives
\[
 X=v_0(1-d^2)+v_1(d^2-d)+(v_2-2t_1)d^2.
\]
Substitute in the last relation of \eqref{maps:eq:E.P16}, using
the invariance of $t_1$. The result is
$t_1=\Omega_L$. Eliminate $t_1$ with this unit
coefficient. Its four invariance relations become
the four relations in \eqref{maps:eq:E.P3}. Eliminating $t_0$
in the second block gives \eqref{maps:eq:E.P4}. This proves the
compact presentation in \eqref{maps:eq:E.P2}--\eqref{maps:eq:E.P4}, including
the integral central-intersection information.

\item For the two cyclic axis summands, the other
central element is a non-zero-divisor: their
group rings are free over its Laurent polynomial
ring. Thus neither contributes to the common
kernel of $q_1,q_2$. This proves \eqref{maps:eq:E.P5}. On
coinvariants, the four invariance relations in
each block vanish, while the norm relations
become three copies of multiplication by3 and
three copies of multiplication by4. Each block
therefore has coinvariants $(\mathbb Z/12)^3$.
This is an independent recovery of \ref{maps:src:E.2.11}.

\item \emph{Restriction to the actual circle fibre.}
The universal covers of $P$ and $Y\times\mathbb R$
are $\Pi$-equivariantly homotopy equivalent by
\ref{maps:src:E.2.13}. The map $\Pi\to G/\langle h_i\rangle$
is injective, since $\chi(h_i)=\pm12$, and has
index12. Hence each of the first two group-ring
summands in \eqref{maps:eq:E.P2} is free of rank12 over
$\mathbb Z\Pi$. The last two factor through the
index12 surface subgroup $\Sigma<\Gamma$ and
have trivial $z$ action. This proves the module
formula in \eqref{maps:eq:E.P6}.

\item For its exact coinvariants use the regular
12-sheet permutation representation with
exponents $(4,-3,-4,3)$ on $(a,b,c,d)$.
Evaluating the compact12-by20 presentation gives
an integral144-by240 matrix. Its Smith form has
56 unit entries and two entries6; its free
cokernel rank is86. The certificate computes
the same result from the independently retained
14-by26 presentation. Adding the24 free
coinvariants proves the last formula in \eqref{maps:eq:E.P6}.
All these are integral module and homotopy
statements. No smooth surgery or smooth
classification is inferred from them.
\end{enumerate}
\end{proof}


\begin{proposition}[the cusp quadratic attachment]\label{maps:src:E.2.18}
\StatementRecord{maps:src:E.2.18}{proved}{Historical events: Claude Fable 5.2 (Anthropic), effort high, 16 November 2026: carried as background without fault found. GPT-6 Astra (OpenAI), ChatGPT Pro mode, effort Pro, S6\_09\_PAPERS, 17 November 2026: targeted proof check. Also checked in this writing. Pre-continuation writing events: GPT-6 Astra (OpenAI), ChatGPT Pro mode, effort Pro, S6\_10\_PAPERS, 18 November 2026: targeted proof check. Continuation metadata events: GPT-6 Astra (OpenAI), ChatGPT Pro mode, effort Pro, S6\_10\_PAPERS continuation, 19 November 2026: used and rechecked in this run; status-vocabulary reconciliation; metadata reconciliation only; no new mathematical check. \emph{Pre-continuation recorded proof scope.} GPT-6 Astra (OpenAI, ChatGPT Pro mode; recorded reviewer /root/cdp\_audit), effort Pro, 17 November 2026; Claude Fable 5.2 (Anthropic), effort high, 16 November 2026; also checked in this writing by GPT-6 Astra (OpenAI), ChatGPT Pro mode, effort Pro, 18 November 2026}{the hypotheses, notation and definitions stated here; \cref{maps:src:E.2.16}; \cref{maps:eq:E.P6}; \cite[Theorems 1 and 5]{mapsWhitehead50}; \cref{maps:src:D.1.9}; \cref{maps:src:E.2.12}; \cref{maps:eq:E.C12}; \cref{maps:eq:E.Q4}; \cref{maps:eq:E.Q5}; \cref{maps:eq:E.Q3}; \cref{maps:eq:E.Q7}; \cref{maps:eq:E.Q6}; \cref{maps:eq:E.C3}; \cref{maps:eq:E.Q8}; \cref{maps:src:E.2.17}; \cref{maps:eq:E.P2}; \cref{maps:eq:E.P3}; \cref{maps:eq:E.P4}; \cref{maps:eq:E.Q11}; \cref{maps:eq:E.Q10}; \cref{maps:eq:E.Q12}}


Keep the right-module presentation of \ref{maps:src:E.2.16} and write
\[
 R=\mathbb ZK,\qquad M=\pi_2(E)=eR\oplus T,
 \qquad T=\overline R^{\,2}\oplus L.
\]
For an abelian group $A$, let $\Gamma(A)$ be Whitehead's universal
quadratic group, with quadratic generator $q(a)$ and symmetric
bilinear cross effect
$[a,b]=q(a+b)-q(a)-q(b)$.  In particular $[a,a]=2q(a)$.
There is a lift of the normalized cusp component
$S=\mathbb {CP}^2\#3\overline{\mathbb {CP}}^{\,2}$ to the
universal cover $\widetilde E$ whose fundamental class has image
\[
 \pm\bigl([e(h-1),t]-12q(t)\bigr)
       \quad\hbox{under}\quad
 H_4(\widetilde E;\mathbb Z)\longrightarrow\Gamma(M).
 \tag{E.Q1}\label{maps:eq:E.Q1}
\]
The sign is solely the orientation chosen on the normalized
component.  Its mixed component in
\[
 \Gamma(M)=\Gamma(eR)\oplus(eR\otimes_{\mathbb Z}T)
                                      \oplus\Gamma(T)
 \tag{E.Q2}\label{maps:eq:E.Q2}
\]
is nonzero.  Consequently $E$ is not homotopy equivalent to
$E_0\#(S^2\times S^4)$ for any connected compact six-manifold
$E_0$, even without requiring the equivalence to preserve the
boundary.  Thus the free $R$ summand of its second homotopy group
does not come from a sphere-product connected summand.
Likewise the actual nine-dimensional fibre $Y$ is not
homotopy equivalent to $Y_0\#(S^2\times S^7)$ for any
connected closed nine-manifold $Y_0$, despite the twenty-four
free $\mathbb Z\Pi$ summands in \eqref{maps:eq:E.P6}.

\end{proposition}


\begin{proof}
\begin{enumerate}[label=\arabic*.,leftmargin=*]
\item We use the natural, deck-equivariant part of Whitehead's exact
sequence
\[
 H_4(\widetilde E)\xrightarrow{\omega_E}\Gamma(\pi_2(E))
       \longrightarrow\pi_3(E)\longrightarrow
 H_3(\widetilde E)\longrightarrow0.
 \tag{E.Q3}\label{maps:eq:E.Q3}
\]
The exactness and identification of its quadratic term are
\cite[Theorems 1 and 5]{mapsWhitehead50}; the direct-sum rule
\eqref{maps:eq:E.Q2} is proved there in Section~2.  These statements apply
to the simply connected universal cover, with no finiteness
requirement.  Naturality, rather than a choice of a word for
one attaching cell, will specify the element in \eqref{maps:eq:E.Q1}.

\item For the normalization calculation take the negatives of
the complex-oriented Picard homology generators and denote
them by $(H,E_1,E_2,E_3)$.  Their intersection matrix is
still $\operatorname{diag}(1,-1,-1,-1)$.  This common reversal
matches the phase-cell orientations of (B.5.16): the toric
line with class $-(uw+6\gamma\delta)$ has complex curve
degrees $(-6,0,1)$, whereas its values on the corresponding
phase classes are $(6,0,-1)$; the support function $rs$
similarly fixes the middle sign.  The exact wall-degree
calculation in \ref{maps:src:D.1.9} gives this comparison.  With these
orientations the six conductor classes are
\[
 E_1,\ H-E_1-E_2,\ E_2,\ H-E_2-E_3,
                                      E_3,\ H-E_1-E_3.
\]
The first opposite pair is self-glued in the surviving $u$
direction.  In the gamma cover its relation is therefore
$E_1=H-E_2-E_3$.  The other two pairs connect adjacent gamma
translates.  Their relations, after eliminating $H$, are
$(h-1)E_2=(h-1)E_3=0$.  Thus the image of normalization
homology is, with these actual curve generators,
\[
 A E_1\oplus(A/(h-1))E_2\oplus(A/(h-1))E_3,
 \qquad A=\mathbb Z[h,h^{-1}],\qquad H=E_1+E_2+E_3.
 \tag{E.Q4}\label{maps:eq:E.Q4}
\]
This also follows directly from the transpose of the matrix
$B(h)$ of \ref{maps:src:E.2.12}, with the intersection signs in the Picard
basis retained.  It is the whole $H_2$ of the gamma cover:
the degree-zero normalization row is the cellular complex
of an infinite cylinder, whose second homology is zero.

\item We fix the comparison with the cusp basis of \eqref{maps:eq:E.C12}, rather
than identify its free summand only abstractly.  In the ordinary
central fibre the first conductor class $\overline C_1$ is the
self-$u$ pair; the other two are the images of $E_2,E_3$,
in either order.  The chosen cusp basis is
\[
 E=-\overline C_1-\overline C_2,\qquad
 U=\overline C_2+\overline C_3,\qquad D=\overline C_2.
\]
The two torsion coordinates in \eqref{maps:eq:E.Q4} are determined by
this ordinary comparison, because each is annihilated by
$h-1$.  The normalization map is onto the whole group in
\eqref{maps:eq:E.Q4}, and $E_2,E_3$ generate its entire $A$-torsion
submodule.  Consequently the coefficient of $E_1$ in the
free quotient is a generator of $A$, not merely a nonzero
element: it is a unit $\pm h^n$.  The ordinary comparison
fixes its sign.  Translating the entire chosen normalization
component by $h^{-n}$ makes this coefficient $-1$ and leaves
both torsion coordinates fixed.  The common Picard sign
convention above changes every generator simultaneously;
no arbitrary shear of the free generator is introduced.
The resulting
normalization map in the basis used for the full peripheral
comparison is
\[
 E_1\longmapsto-E-D,\qquad
 E_2\longmapsto D,\qquad E_3\longmapsto U-D,
 \qquad H\longmapsto-E+U-D.
 \tag{E.Q5}\label{maps:eq:E.Q5}
\]
Exchanging $E_2,E_3$ has no effect below.  This argument
retains the full normalization map on $H_2$; it does not
suppose that the derived chain changes in \ref{maps:src:E.2.16} are induced
by diffeomorphisms.

\item The intersection form of $S$ is
$\operatorname{diag}(1,-1,-1,-1)$.  Its usual four-cell
description, or the Hopf invariant of its two-dimensional
classes, therefore gives
\[
 \omega_S[S]=q(H)-q(E_1)-q(E_2)-q(E_3).
\]
Substitute $H=E_1+E_2+E_3$ and then \eqref{maps:eq:E.Q5}.  Bilinearity
of the cross effect gives
\[
 \omega_S[S]\longmapsto
 [E_1,E_2]+[E_1,E_3]+[E_2,E_3]
                      =-[E,U]-2q(D).
 \tag{E.Q6}\label{maps:eq:E.Q6}
\]
The normalization map $S\to E$ lifts to $\widetilde E$
because $S$ is simply connected.  Naturality of \eqref{maps:eq:E.Q3}
shows that the last expression is the image of its actual
four-dimensional fundamental class.  In particular no
unrecorded three-cell can invalidate the assertion that
this element belongs to $\operatorname{im}\omega_E$.

\item The exact global cancellation certificate now retains
the three cusp two-cycles in addition to the peripheral
chains.  It gives, in the generators $(x,t,y,z,e)$ of
\ref{maps:src:E.2.16},
\[
 E\longmapsto(1-h)e+6t,\qquad
 U\longmapsto t,\qquad D\longmapsto0.
 \tag{E.Q7}\label{maps:eq:E.Q7}
\]
For clarity, before the last relation reductions its
$t$ coefficient for $E$ is
$-3h+4ha+4ha^2+h^2$, whose augmentation is six;
its free coefficient is $1-h$.  Its coefficient for $D$
is $2h^{-1}a^2-1-a$, which has augmentation zero and hence
annihilates the invariant generator $t$.  The other
coordinates vanish.  The program
\path{code/III/maps_cusp_quadratic.py} checks these identities
using the identical integral group-ring cancellations as
\ref{maps:src:E.2.16}.  Applying $\Gamma$ to \eqref{maps:eq:E.Q7} in \eqref{maps:eq:E.Q6} gives
\[
 -[(1-h)e+6t,t]
                =[e(h-1),t]-12q(t),
\]
as asserted.  In ordinary coinvariants $t=2v$ and
$E=12v$ for a primitive class $v$.  The quadratic image
is then $-48q(v)$, consistent with the degree-two cup
square evaluated on the normalization; this is a check,
not the proof of \eqref{maps:eq:E.Q1}.

\item We finish with the connected-sum obstruction, allowing
every possible free summand.  Multiplication by $h-1$
is injective on $R$: a finitely supported group-ring
element constant along an infinite $h$ orbit is zero.
It annihilates $T$, so
\[
 \operatorname{Ann}_M(h-1)=T,
                       \qquad\operatorname{Hom}_R(T,R)=0.
 \tag{E.Q8}\label{maps:eq:E.Q8}
\]
Suppose that a copy $fR$ is a direct summand of $M$.
Write $f=eu+t_0$ and let a retraction send $e$ to $v$.
It vanishes on $T$, and hence $vu=1$.

\item The ring $\mathbb ZK$ is directly finite.  One elementary
verification uses the index-twelve subgroup of $K$:
its quotient by the central $\mathbb Z$ is the free group
of the punctured genus-three surface, so it is
$F_6\times\mathbb Z$.  This subgroup, and hence $K$, is
residually finite.  In every finite-group quotient the
regular representation is finite dimensional over
$\mathbb Q$, so $vu=1$ implies $uv=1$.  Residual
finiteness separates the finite support of $uv-1$ and
proves the same equality in $R$.  Thus $u$ is a unit,
the kernel of the retraction is exactly $T$, and every
free splitting is obtained by replacing $e$ by
$eu+t_0$.

\item Under such a change the mixed component of \eqref{maps:eq:E.Q1} is
$f u^{-1}(h-1)\otimes t$.  It is nonzero: $t$ has
infinite order by \eqref{maps:eq:E.C3}, and $u^{-1}(h-1)$ is nonzero.
The central generator may also be reversed by a group
automorphism, which multiplies $h-1$ by a unit and
does not affect the argument.

\item If $E\simeq E_0\#(S^2\times S^4)$, its universal cover
has, through dimension four, the homotopy type of
$\widetilde E_0$ wedged with one copy of $S^2\vee S^4$
for each deck element.  This follows by lifting the
connected-sum balls; deleting or regluing their boundary
five-spheres does not alter the four-skeleton.  The new
two-spheres generate a free $R$ summand.  The four-spheres
have trivial attaching maps, and every remaining
four-dimensional Whitehead image lies in
$\Gamma(\pi_2(E_0))$.  There can therefore be no mixed
component between this free summand and its complement.
This contradicts \eqref{maps:eq:E.Q1}--\eqref{maps:eq:E.Q8} for every choice of the
summand, and proves the stated obstruction.

\item For the assertion about $Y$, first retain the actual maps
from the two complements to the splice.  The certificate
\path{code/III/maps_splice_quadratic.py} tracks the four cycles
$(t_0,e_0,t_1,e_1)$ through every unit cancellation used in
\ref{maps:src:E.2.17}.  In its sixteen-generator presentation they become
the four unit vectors in rows $(4,3,12,15)$, respectively
(rows are numbered from zero).  The subsequent seven-generator
block changes therefore give
\[
 t_0=\Omega_R,\qquad t_1=\Omega_L,
 \qquad e_0\in R/q_2R,\quad e_1\in R/q_1R
 \tag{E.Q9}\label{maps:eq:E.Q9}
\]
with exactly the axes in \eqref{maps:eq:E.P2}.  These inclusion maps have
not been selected from an abstract direct-sum decomposition.
Both $t_i$ are invariant under $G$: the four displayed
relations for each $\Omega$ in \eqref{maps:eq:E.P3}--\eqref{maps:eq:E.P4} annihilate
it by $a-1,b-1,c-1,d-1$.  The certificate also checks these
as the eight pure columns numbered
$(7,8,11,13)$ for $t_0$ and $(1,3,17,18)$ for $t_1$
in the retained fourteen-generator central presentation,
before eliminating the invariant generators.
On $\Gamma$-coinvariants
their respective images are
\[
 -2(B_2+C_2),\qquad 2(p_2-r_2),
 \tag{E.Q10}\label{maps:eq:E.Q10}
\]
and each has order exactly six: the relevant two coordinates
have independent orders three and four in \eqref{maps:eq:E.P3}--\eqref{maps:eq:E.P4}.
In particular an integer annihilating either $t_i$ in the
full module must be divisible by six.

\item Put $A_\Pi=\mathbb Z\Pi$ and $z=h_1h_2$.  Choose the twelve
coset representatives giving the free restriction of each
axis in \eqref{maps:eq:E.P6}, and denote the resulting twenty-four generators
by $f_j$.  The normalized cusp surfaces and all their chosen
deck translates are actual maps to $\widetilde P$.  Transfer
them to $\widetilde Y$ using the $\Pi$-equivariant homotopy
equivalence $\widetilde P\simeq\widetilde Y\times\mathbb R$.
On the first axis $h_2$ is the identity, so multiplication
by $h_1-1$ equals multiplication by $z-1$; the analogous
statement holds on the second axis.  Naturality of \eqref{maps:eq:E.Q3}
therefore supplies, for every $j$, an element in the image
of the Whitehead map of $\widetilde Y$ of the form
\[
 \pm\bigl([f_j(z-1),t_{i(j)}]-12q(t_{i(j)})\bigr),
 \tag{E.Q11}\label{maps:eq:E.Q11}
\]
where $i(j)=0$ on the first twelve axes and $1$ on the
remaining twelve.  Invariance of $t_i$ under $G$ removes
any extra coset action in this formula.

\item Write $\pi_2(Y)=F\oplus T_Y$, with $F=A_\Pi^{24}$ and
$(z-1)T_Y=0$ as in \eqref{maps:eq:E.P6}.  Suppose a free copy $fA_\Pi$
is a direct summand, with a retraction $r:\pi_2(Y)\to A_\Pi$
satisfying $r(f)=1$.  Multiplication by $z-1$ is injective
on $A_\Pi$, so $r(T_Y)=0$.  Write $r(f_j)=v_j$ and
$f=\sum_j f_j u_j+t$; then
\[
                         \sum_j v_j u_j=1.
 \tag{E.Q12}\label{maps:eq:E.Q12}
\]
In the splitting $fA_\Pi\oplus\ker r$, the mixed component
of \eqref{maps:eq:E.Q11} is $f v_j(z-1)\otimes t_{i(j)}$.  If all these
mixed components were zero, every integer coefficient of
$v_j(z-1)$ would be divisible by six, by \eqref{maps:eq:E.Q10}.  The
group-ring basis makes this implication immediate even if
$t_i$ has larger or infinite order.  Multiplication by
$z-1$ is also injective on $(\mathbb Z/6)[\Pi]$: along an
infinite $z$ orbit a finitely supported coefficient function
with zero difference is zero.  Hence every $v_j$ belongs
to $6A_\Pi$, contradicting \eqref{maps:eq:E.Q12}.

\item Finally, a connected sum with $S^2\times S^7$ has, through
dimension four in its universal cover, the original space
wedged with a free deck orbit of two-spheres.  All images
of its four-dimensional Whitehead map consequently have
zero mixed component with that free summand.  The preceding
argument rules this out for every free summand of the actual
module.  This proves the assertion for $Y$.
\end{enumerate}
\end{proof}

This computes a specified part of the actual quadratic
attaching data.  It does not identify every higher handle
attachment or the smooth type of the shifted fivefold.


\begin{proposition}[the relative smooth cusp model]\label{maps:src:E.2.19}
\StatementRecord{maps:src:E.2.19}{proved}{Historical events: Claude Fable 5.2 (Anthropic), effort high, 16 November 2026: carried as background without fault found. GPT-6 Astra (OpenAI), ChatGPT Pro mode, effort Pro, S6\_09\_PAPERS, 17 November 2026: targeted proof check. Also checked in this writing. Pre-continuation writing events: GPT-6 Astra (OpenAI), ChatGPT Pro mode, effort Pro, S6\_10\_PAPERS, 18 November 2026: targeted proof check. Continuation metadata events: GPT-6 Astra (OpenAI), ChatGPT Pro mode, effort Pro, S6\_10\_PAPERS continuation, 19 November 2026: used and rechecked in this run; status-vocabulary reconciliation; metadata reconciliation only; no new mathematical check. \emph{Pre-continuation recorded proof scope.} GPT-6 Astra (OpenAI, ChatGPT Pro mode; recorded reviewer /root/cdp\_audit), effort Pro, 17 November 2026; Claude Fable 5.2 (Anthropic), effort high, 16 November 2026; also checked in this writing by GPT-6 Astra (OpenAI), ChatGPT Pro mode, effort Pro, 18 November 2026}{the hypotheses, notation and definitions stated here; \cref{maps:src:E.2.21}; \cref{maps:src:E.2.16}; \cref{maps:eq:E.T3}; \cref{maps:eq:E.T4}; \cref{maps:eq:E.T8}}


For the original marked cusp, let
\[
 Z(q)=\frac{\log q}{2\pi i}B_0+C(q),\qquad
 B_0=\begin{pmatrix}0&1\\-1&0\end{pmatrix},
\]
where $C(q)$ is holomorphic and bounded on a disc about zero.
For a sufficiently small closed disc the actual six-dimensional
cusp tube is orientation-preservingly diffeomorphic, relative to
its full marked torus-bundle boundary, to the unitless toric
quotient defined by $C=0$.  The latter has the following finite
normal-crossings plumbing description.

Start with the canonical complex line bundle of
$S=\mathbb {CP}^2\#3\overline{\mathbb {CP}}^{\,2}$ and its six
toric curves in cyclic order
\[
 \begin{aligned}
 C_0&=E_1,& C_1&=H-E_1-E_2,& C_2&=E_2,\\
 C_3&=H-E_2-E_3,& C_4&=E_3,& C_5&=H-E_1-E_3.
 \end{aligned}
 \tag{E.T1}\label{maps:eq:E.T1}
\]
Pair $C_i$ with $C_{i+3}$ for $i=0,1,2$.  Both the surface-normal
line and the canonical line restricted to every $C_i$ have degree
$-1$.  In the resulting $\mathcal O(-1)\oplus\mathcal O(-1)$
coordinates $(z,a,b)$, with $z=0$ at the endpoint meeting
$C_{i+1}$ and $b$ the surface-normal coordinate, the pairing is
\[
             (z,a,b)\longmapsto(z^{-1},zb,za).
 \tag{E.T2}\label{maps:eq:E.T2}
\]
The function giving the cusp parameter is $q=zab$ in this
coordinate chart.  Formula \eqref{maps:eq:E.T2}, not merely the induced map
of zero sections, specifies the normal framing of each plumbing.
The two triple points are the endpoint orbits
$\{0,2,4\}$ and $\{1,3,5\}$, where endpoint $i$ means
$C_i\cap C_{i+1}$.  Ordered local complex triples undergo a
cyclic permutation around each orbit, with threefold composition
the identity.  The plumbing is formed with collars and rounded
corners; its distinguished outer boundary is $|q|=\eta$.

This statement supplies the relative smooth cusp input for the
ordinary framed handle construction of \ref{maps:src:E.2.21}.  It does not
promote the homological cancellations of \ref{maps:src:E.2.16} to smooth
cancellations.

\end{proposition}


\begin{proof}
\begin{enumerate}[label=\arabic*.,leftmargin=*]
\item We first construct the comparison by a smooth family, rather
than use the logarithmic homeomorphism in (B.5.7).  Set
$C_s(q)=(1-s)C(q)$ for $0\leq s\leq1$.  The affine toric cover
is fixed.  On its dense torus the deck translation for
$\lambda\in\mathbb Z^2$ becomes
\[
 (x,q,s)\longmapsto
 \bigl(q^{B_0\lambda}
       \exp(2\pi i C_s(q)\lambda)x,\ q,\ s\bigr).
 \tag{E.T3}\label{maps:eq:E.T3}
\]
The exponential is componentwise.  The powers of $q$ permute
the toric charts by translation of the $A_2$ triangulation,
while the exponential factors are nowhere-zero analytic units
on every chart.  Thus \eqref{maps:eq:E.T3} defines a smooth family of
holomorphic deck actions on a fixed smooth toric covering
space.  On a central stratum the lattice translation permutes
its bounded fan cells without fixed points.  Off the central
fibre, freeness and properness are exactly the nondegeneracy
of the real period lattice.

\item These assertions hold uniformly after shrinking $\eta$.
Indeed put $T=-\log|q|/(2\pi)$, and let $K$ bound
$\|\operatorname{Im}C_s(q)\|$ on the compact parameter interval
and closed disc.  For $T>K+1$ the matrix
$TB_0+\operatorname{Im}C_s(q)$ is invertible, since $B_0$ is
orthogonal.  A fundamental parallelogram has logarithmic
position
$B_0a+T^{-1}\operatorname{Im}C_s(q)a$ with
$a\in[0,1]^2$; it meets a uniformly bounded number of toric
charts.  More explicitly choose a triangle containing the
unperturbed point $B_0a$ and let $P$ be its unimodular vertex
matrix.  Its barycentric coordinates $t_j$ are nonnegative,
and the corresponding chart coordinates satisfy
$\log|z_j|=t_j\log|q|+
O(2\pi\|P^{-1}\|K\sqrt2)$.
There are only finitely many such triangles, so with
$c=2\pi\sqrt2\max\|P^{-1}\|$ every chosen chart coordinate
has modulus at most $\exp(cK)$, uniformly in $q,s$ and
all compact phases.  Their closed polydiscs have compact
closure in the affine charts.  These finitely many compact
chart images prove
compactness of the quotient over $[0,1]$ and $|q|\leq\eta$.
The central charts and their incidence relations are fixed.
Projection of this quotient to $s$ is consequently a proper
smooth submersion, including on every normal-crossings stratum.
Its only boundary is $|q|=\eta$ together with the parameter
end faces, and its boundary restriction is also a submersion.

\item On $|q|=\eta$ prescribe the marked real coordinates
$x=\exp(2\pi i(Z_s(q)a+b))$, $(a,b)\in T^4$.
They give a smooth product trivialization in $s$ and precisely
the prescribed mapping-torus marking.  Extend its horizontal
vector field over a collar, and then over the compact family.
Local toric product charts provide extensions tangent to
every central stratum; a partition of unity preserves that
property and the lift of $\partial_s$.  Integrating this
field proves the asserted relative diffeomorphism.  Its
orientation is the complex orientation on each fibre.
This is relative Ehresmann transport in an actual smooth
family; no smoothness is inferred from a logarithmic formula
at $q=0$.

\item It remains to identify the endpoint $C=0$ with the framed
plumbing data.  In the lattice $\mathbb Z^2$ use cyclic rays
\[
 (1,0),\ (0,1),\ (-1,1),\ (-1,0),\ (0,-1),\ (1,-1).
 \tag{E.T4}\label{maps:eq:E.T4}
\]
The star of the vertex zero in the $A_2$ triangulation consists
of the six triangles spanned by zero and consecutive rays.
The corresponding three-dimensional cones have primitive rays
$(v,1)$ and determinant one.  Their common divisor at the
central ray is the toric surface $S$.  A principal central
fibre has normal line
\[
             \mathcal O_S(-C_0-\cdots-C_5)=K_S.
 \tag{E.T5}\label{maps:eq:E.T5}
\]
The cyclic surface fan in \eqref{maps:eq:E.T4} also proves this identification
directly: it is the fan of the three torus-fixed blow-ups of
$\mathbb {CP}^2$, and each neighboring-ray relation gives
$C_i^2=-1$.  In the Picard basis \eqref{maps:eq:E.T1},
$K_S=-3H+E_1+E_2+E_3$, so $K_S\cdot C_i=-1$.

\item Near a double curve the two adjacent unimodular charts give
the total space of $\mathcal O(-1)\oplus\mathcal O(-1)$.
Its two base charts have coordinates
\[
 (z,a,b),\qquad
 (w,A,B)=(z^{-1},za,zb),\qquad q=zab=wAB.
 \tag{E.T6}\label{maps:eq:E.T6}
\]
The period translation taking a neighboring component back
to the chosen component exchanges the two normal directions.
It reverses the chosen endpoint coordinate on the opposite
toric curve.  In the first chart this is exactly \eqref{maps:eq:E.T2};
in the target infinity chart it is simply $(z,b,a)$, so it
extends smoothly over the endpoint.  The map is involutive,
preserves $q$, and its complex Jacobian determinant is $1$.
These equalities specify the normal-disk attachment, including
its sign and angular twisting.

\item The endpoint $C_i\cap C_{i+1}$ is sent to
$C_{i+2}\cap C_{i+3}$.  At the first endpoint write the
ordered ambient coordinates as the surface-normal coordinate,
the $C_i$-normal coordinate and the $C_{i+1}$-normal coordinate.
The next endpoint has their order cyclically permuted.
After three crossings the same coordinates return, not just
the same unframed planes.  Thus the two endpoint orbits give
ordinary triple charts $q=z_0z_1z_2$, with no additional
clutching loop at either triple point.

\item Equivalently, for a triangle with vertex matrix $P$ whose
columns are $(v_i,1)$, translation by $n\in\mathbb Z^2$
replaces it by $L_nP$, where
\[
 L_n=\begin{pmatrix}1&0&n_1\\0&1&n_2\\0&0&1\end{pmatrix}.
 \tag{E.T7}\label{maps:eq:E.T7}
\]
It acts as the identity in the corresponding translated
unimodular coordinates, since
$(L_nP)^{-1}L_nP=I$.  This verifies the preceding plumbing
and triple identifications in the actual toric charts.
More explicitly, order the star matrix as
$P_i=((0,1),(r_i,1),(r_{i+1},1))$ and reverse its
columns when using the coordinate order $(z,a,b)$.
For every $i$ the matrix of translation by $-r_i$ from
this chart to the opposite following-endpoint chart is
\[
 \begin{pmatrix}-1&0&0\\1&0&1\\1&1&0\end{pmatrix}.
 \tag{E.T9}\label{maps:eq:E.T9}
\]
Its monomials are exactly \eqref{maps:eq:E.T2}.  In normal-first order,
$P_{i+2}^{-1}L_{-r_i}P_i$ is the cyclic permutation
matrix.  Thus the framing and triple-cycle assertions
are checked for all six actual fan identifications,
not imposed on an abstract plumbing diagram.
Taking collars of the component core, three double-curve
annuli and two triple charts, and rounding their outer
corners, gives the asserted finite plumbing.  These are
subcharts of the smooth toric quotient, so the construction
is Hausdorff and the collar identifications are the indicated
smooth ones; it is not a quotient formed by identifying
arbitrary closed disk-bundle neighborhoods.

\item Finally the period translations for $g,u$ are $-e_2,e_1$.
On $q=\eta e^{-2\pi i\theta}$ the marked expression is
$x=\exp(2\pi i((-\theta+iT)B_0a+b))$.
At $\theta=1$ it agrees with $\theta=0$ and
$b$ replaced by $b-B_0a$.  This is exactly the inverse
endpoint convention for
\[
 M_0=\begin{pmatrix}
 1&0&0&0\\0&1&0&0\\0&1&1&0\\-1&0&0&1
 \end{pmatrix}.
 \tag{E.T8}\label{maps:eq:E.T8}
\]
The comparison therefore preserves the full outer marking,
not merely the isomorphism class of the boundary manifold.
The six unimodular charts, all normal degrees, the plumbing
Jacobian and involution, both triple cycles and \eqref{maps:eq:E.T8} are
checked exactly in \path{code/III/maps_cusp_smooth_plumbing.py}.
\end{enumerate}
\end{proof}


\begin{proposition}[the explicit smooth plumbing presentation]\label{maps:src:E.2.20}
\StatementRecord{maps:src:E.2.20}{proved}{Historical events: Claude Fable 5.2 (Anthropic), effort high, 16 November 2026: carried as background without fault found. GPT-6 Astra (OpenAI), ChatGPT Pro mode, effort Pro, S6\_09\_PAPERS, 17 November 2026: targeted proof check. Also checked in this writing. Pre-continuation writing events: GPT-6 Astra (OpenAI), ChatGPT Pro mode, effort Pro, S6\_10\_PAPERS, 18 November 2026: targeted proof check. Continuation metadata events: GPT-6 Astra (OpenAI), ChatGPT Pro mode, effort Pro, S6\_10\_PAPERS continuation, 19 November 2026: used and rechecked in this run; status-vocabulary reconciliation; metadata reconciliation only; no new mathematical check. \emph{Pre-continuation recorded proof scope.} GPT-6 Astra (OpenAI, ChatGPT Pro mode; recorded reviewer /root/cdp\_audit), effort Pro, 17 November 2026; Claude Fable 5.2 (Anthropic), effort high, 16 November 2026; also checked in this writing by GPT-6 Astra (OpenAI), ChatGPT Pro mode, effort Pro, 18 November 2026}{the hypotheses, notation and definitions stated here; \cref{maps:src:E.2.21}; \cref{maps:src:E.2.19}; \cref{maps:eq:E.S1}; \cref{maps:eq:E.T2}; \cref{maps:eq:E.U7}; \cref{maps:eq:E.U10}; \cref{maps:eq:E.U8}; \cref{maps:eq:E.U9}; \cref{maps:src:B.5.1}; \cref{maps:src:E.2.6}}


The following finite smooth plumbing presentation uses only standard
pieces and specified attaching maps. Its conversion into ordinary
framed handles and the relative smooth identification are given in \ref{maps:src:E.2.21}.

Let $F=\mathbb R^4/\mathbb Z^4$, with basis $(g,u,w,d)$, and use the
two matrices $A_j$ and fixed vectors $v_j$ of \cref{mapsI:src:A.1.1}.  Set
\[
 (m_1,m_2)=(3,4),\quad(\kappa_1,\kappa_2)=(1,-1),\quad
 w_j=\kappa_jv_j,\qquad L=\ker\gamma,
\]
so that $\Lambda=\mathbb Zw_j\oplus L$ and $A_jw_j=w_j$.
On $L$ the two matrices are
\[
 B_1=\begin{pmatrix}0&1&0\\-1&-1&0\\1&0&1\end{pmatrix},
 \qquad
 B_2=\begin{pmatrix}0&-1&0\\1&0&0\\0&1&1\end{pmatrix}.
 \tag{E.U1}\label{maps:eq:E.U1}
\]
Define the standard smooth four-manifolds
\[
 \mathcal S_j=(T^3\times\mathbb R)/
       ((l,\sigma+1)\sim(B_j^{-1}l,\sigma)),
 \qquad \mathcal N_j=\mathcal S_j\times D^2_r.
 \tag{E.U2}\label{maps:eq:E.U2}
\]
In particular the cap's oriented normal disk is part of the
definition, not an unspecified disk bundle.

The cap boundary is identified with $M_{A_j}$ by the following
explicit map.  Write $x=\phi w_j+l$ and use the interval marking
$(x,1)\sim(A_j^{-1}x,0)$ on $M_{A_j}$.  Then
\[
 [x,\theta]\longmapsto
 \left([l,\sigma=m_j\kappa_j\phi+\theta],
                         \ r=e^{2\pi i\kappa_j\phi}\right).
 \tag{E.U3}\label{maps:eq:E.U3}
\]
All angles on the right are taken in their stated mapping-torus
quotients.  In the ordered angular variables $(\phi,\theta)$
and $(\sigma,\alpha)$, where $r=e^{2\pi i\alpha}$, the map is
\[
 \begin{pmatrix}\sigma\\\alpha\end{pmatrix}
 =
 \begin{pmatrix}m_j\kappa_j&1\\\kappa_j&0\end{pmatrix}
 \begin{pmatrix}\phi\\\theta\end{pmatrix},
 \qquad
 \begin{pmatrix}\phi\\\theta\end{pmatrix}
 =
 \begin{pmatrix}0&\kappa_j\\1&-m_j\end{pmatrix}
 \begin{pmatrix}\sigma\\\alpha\end{pmatrix}.
 \tag{E.U4}\label{maps:eq:E.U4}
\]
These are integral unimodular matrices.

To define a compact six-manifold $\mathcal E_{\rm pl}$, take the
flat $F$-bundle over an oriented disc with three open interior
discs removed.  Its three marked inner loop monodromies are
$A_1,A_2,M_0$, with $A_1A_2M_0=I$.  Equivalently cut the planar
surface along three disjoint paths to its outer boundary, take
the product with $F$, and reglue the two banks by these three
integral linear diffeomorphisms.  This specifies smooth band
attachments, including their product normal coordinates.
Fill the first two boundary components by \eqref{maps:eq:E.U2}--\eqref{maps:eq:E.U4}, and
the third by the finite canonical-bundle plumbing of \ref{maps:src:E.2.19},
using its complete marked $M_0$ boundary map.  The remaining
boundary is marked as $F\times S^1_\mu$, with the outer meridian
orientation of \eqref{maps:eq:E.S1}.  The core uses its zero section, so no
affine translation remains unspecified on that boundary.

For the actual original regular-fibre complement there is an
orientation-preserving diffeomorphism
\[
                  E\ \cong\ \mathcal E_{\rm pl}
                 \quad\hbox{relative to }F\times S^1_\mu .
 \tag{E.U5}\label{maps:eq:E.U5}
\]
Consequently the shifted fivefold has the explicit smooth
presentation
\[
 P\ \cong\
 (\mathcal E_{\rm pl}\times T^4)
 \ \cup_{\ (x_1,x_2,\theta)\mapsto(x_1,x_2,-\theta)\ }\
 (T^4\times\mathcal E_{\rm pl}).
 \tag{E.U6}\label{maps:eq:E.U6}
\]
Every piece in this presentation is a standard torus product,
a linear three-torus mapping torus times an oriented disk, or
the canonical disk bundle of the explicitly specified toric
blow-up with the normal-disk plumbings \eqref{maps:eq:E.T2}.  All attachment
maps have been specified on their normal coordinates.

\end{proposition}


\begin{proof}
\begin{enumerate}[label=\arabic*.,leftmargin=*]
\item We verify first that \eqref{maps:eq:E.U2}--\eqref{maps:eq:E.U4} are the actual finite caps,
including the normal framing.  The original cap is
\[
 (F\times D^2_s)/
 \bigl((x,s)\sim(A_jx+v_j/m_j,e^{-2\pi i/m_j}s)\bigr).
 \tag{E.U7}\label{maps:eq:E.U7}
\]
The good-reduction real flat marking is smooth on the entire
root disc, so these are smooth coordinates even at $s=0$.
Because $\gamma A_j=\gamma$ and $\gamma(v_j)=\kappa_j=\pm1$,
\[
                  r=s\,e^{2\pi i\kappa_j\gamma(x)}
 \tag{E.U8}\label{maps:eq:E.U8}
\]
is invariant under the displayed action.  It is single-valued
on $F$, has the same absolute value as $s$, and is complex
linear on every normal disk.  The induced map
$[x,s]\mapsto([x],r)$ identifies \eqref{maps:eq:E.U7} smoothly with
$S_j\times D^2$.  An inverse is obtained by choosing any
representative $x$ and setting
$s=r e^{-2\pi i\kappa_j\gamma(x)}$; invariance makes this
independent of the representative in the quotient.

\item In the integral basis $(w_j,u,w,d)$ the affine generator is
\[
                  (\phi,l)\longmapsto
                     (\phi+\kappa_j/m_j,B_jl).
 \tag{E.U9}\label{maps:eq:E.U9}
\]
Thus $\sigma=m_j\kappa_j\phi$ gives exactly the mapping torus
in \eqref{maps:eq:E.U2}.  This also supplies a direct inverse to the product
description: choose a representative $(l,\sigma)$, put
$x=(\kappa_j\sigma/m_j)w_j+l$ and
$s=r e^{-2\pi i\sigma/m_j}$.  Replacing that representative
by the defining mapping-torus relation changes $(x,s)$ by
the finite affine generator and hence not its class.

\item The actual boundary parametrization from (B.5.3) is
\[
 [x,\theta]\longmapsto
       [x+\theta v_j/m_j,e^{-2\pi i\theta/m_j}].
 \tag{E.U10}\label{maps:eq:E.U10}
\]
Substitute \eqref{maps:eq:E.U10} into \eqref{maps:eq:E.U8}--\eqref{maps:eq:E.U9}.  The coefficient of
$\theta$ in the normal angle is
$-1/m_j+\kappa_j^2/m_j=0$, while the base angle is
$m_j\kappa_j\phi+\theta$.  This proves \eqref{maps:eq:E.U3} and \eqref{maps:eq:E.U4}.
At $\theta=1$ that base angle advances by one, precisely
replacing $l$ by $B_j^{-1}l$; this checks the mapping-torus
seam.  A lattice change of $\phi$ changes the base angle by
$m_j\kappa_j$, and $B_j^{m_j}=I$, so it introduces no further
ambiguity.

\item The function of the cap that gives its original base parameter
is, in these smooth coordinates,
\[
                         q=r^{m_j}e^{-2\pi i\sigma}.
 \tag{E.U11}\label{maps:eq:E.U11}
\]
On the marked boundary it has angle
$m_j\alpha-\sigma=-\theta$, the original clockwise base
convention.  A positive normal meridian at fixed
$(l,\sigma)$ has $(\Delta\phi,\Delta\theta)=(\kappa_j,-m_j)$.
Since $\kappa_jw_j=v_j$, it is exactly $t_{v_j}\rho_j^{-m_j}$,
recovering the actual filling relation
$\rho_j^{m_j}=t_{v_j}$.  The normal circle is therefore not
silently exchanged for a once-around base loop.  The exact
integer checks of all these seam, framing and meridian
identities are in \path{code/III/maps_finite_cap_framings.py}.

\item Over the punctured planar core, the actual family has the
smooth real flat marking with linear monodromies $A_j$ and
its zero section.  The logarithmic finite gluing has exactly
the drift in \eqref{maps:eq:E.U10}; after subtracting it, the remaining
translation has a single-valued real lift and is removed by
the collar isotopy already proved in \ref{maps:src:B.5.1}.  No nonzero
lattice winding is discarded.  At the cusp use the actual
relative smooth diffeomorphism of \ref{maps:src:E.2.19}.  Each comparison
has precisely the prescribed boundary value and can be
made product on a smaller collar.  They therefore glue to
the relative diffeomorphism \eqref{maps:eq:E.U5}.  In particular the
normal framing of the unfilled ordinary fibre is unchanged.

\item Finally use the two hemispheres in the actual shifted
fibre-product decomposition of \ref{maps:src:E.2.6}.  On each hemisphere
the other factor is a regular marked torus bundle over a
disc, hence a product preserving its zero section and flat
marking.  Their common boundary has unchanged torus
coordinates and opposite normal angles.  Taking the product
of \eqref{maps:eq:E.U5} with those regular factors and gluing gives
\eqref{maps:eq:E.U6}, with its induced complex orientation.  This is a
comparison of all smooth pieces and all boundary maps,
not an identification inferred from group or homology.
\end{enumerate}
\end{proof}

\ref{maps:src:E.2.21} explicitly converts the canonical-bundle self-plumbing and
all the exterior attachments into ordinary framed handles.
In the absolute planar-core description used here, the band part
contains three products
$T^4\times D^1\times D^1$ attached by the specified linear maps;
its product handles include index five.  Those high handles
cannot all disappear: if a compact six-manifold with this boundary
had only handles of index at most four, its dual relative
decomposition would make $\pi_1(\partial E)\to\pi_1(E)$ surjective,
contradicting \eqref{maps:eq:E.S1}.  The displayed plumbing presentation is
therefore not being confused with a four-dimensional thickening
or a chain cancellation.


\begin{theorem}[ordinary framed handles and the smooth type]\label{maps:src:E.2.21}
\StatementRecord{maps:src:E.2.21}{proved}{Historical events: Claude Fable 5.2 (Anthropic), effort high, 16 November 2026: carried as background without fault found. GPT-6 Astra (OpenAI), ChatGPT Pro mode, effort Pro, S6\_09\_PAPERS, 17 November 2026: targeted proof check. Also checked in this writing. Pre-continuation writing events: GPT-6 Astra (OpenAI), ChatGPT Pro mode, effort Pro, S6\_10\_PAPERS, 18 November 2026: targeted proof check. Continuation metadata events: GPT-6 Astra (OpenAI), ChatGPT Pro mode, effort Pro, S6\_10\_PAPERS continuation, 19 November 2026: used and rechecked in this run; status-vocabulary reconciliation; metadata reconciliation only; no new mathematical check. \emph{Pre-continuation recorded proof scope.} GPT-6 Astra (OpenAI, ChatGPT Pro mode; recorded reviewer /root/cdp\_audit), effort Pro, 17 November 2026; Claude Fable 5.2 (Anthropic), effort high, 16 November 2026; also checked in this writing by GPT-6 Astra (OpenAI), ChatGPT Pro mode, effort Pro, 18 November 2026}{the hypotheses, notation and definitions stated here; \cref{maps:eq:E.T4}; \cite[Section~4.1, equation~(16), p.~103]{mapsVogel09}; \cref{maps:src:E.2.20}; \cref{maps:eq:E.U3}; \cref{maps:eq:E.U4}; \cref{maps:src:E.2.19}; \cref{maps:eq:E.T6}; \cref{maps:eq:E.H13}; \cite{mapsMilnor65}}


The following finite ordinary framed-handle prescription defines a
compact oriented six-manifold $\mathcal H_6$ with marked boundary
$T^4\times S^1$, and a closed oriented ten-manifold $\mathcal H_{10}$.
There are orientation-preserving diffeomorphisms
\[
 E\cong\mathcal H_6\quad\hbox{relative to the full marked boundary},
 \qquad P\cong\mathcal H_{10}.
 \tag{E.H0}\label{maps:eq:E.H0}
\]
The respective numbers of ordinary handles, in increasing index, are
\[
 (2,5,14,22,21,10,2),\qquad
 (4,31,119,290,485,574,485,290,119,31,4).
\]
The attaching embeddings and their normal coordinates are specified
below; these numbers alone are not the specification.  No handle
cancellation, smooth recognition from a chain complex, or uncomputed
homotopy-sphere summand is used.
The balls and complex normal charts carry their complex orientations;
orient the product blocks by their displayed coordinate embeddings.
For example the dense block below has orientation
$-d\arg X\wedge d\arg Y\wedge d\mu_1\wedge d\mu_2
\wedge d\operatorname{Re}q\wedge d\operatorname{Im}q$.
This fixes the orientation convention even when the chosen order of
angle and moment coordinates differs from the complex coordinate order.

\emph{1. The two triple-point balls.}
Use two complex coordinate balls, denoted $V_+$ and $V_-$, with
ordered coordinates $Z=(Z_0,Z_1,Z_2)$ and $W=(W_0,W_1,W_2)$.
The boundary faces are rounded coordinate faces of
\[
 |Z_i|\leq\epsilon_v,\quad |Z_0Z_1Z_2|\leq\eta,
 \qquad
 |W_i|\leq\epsilon_v,\quad |W_0W_1W_2|\leq\eta.
 \tag{E.H1}\label{maps:eq:E.H1}
\]
They are smooth six-balls.  Throughout the prescription, rounding is
in the nonnegative radii, preserves the angular coordinates, and is
fixed on the attaching collars.  Positive changes of the radii only
give the corresponding collar isotopy.  The complex products in
\eqref{maps:eq:E.H1} are the same function $q$.
The coordinate formulas may first be read on manifolds with corners;
apply the same radial collar reparametrization on the two sides of
each attaching face before smoothing their union.  This convention
also specifies the maps at a corner meeting $|q|=\eta$.

\emph{2. The three annular blocks, with their endpoint framings.}
For $i=0,1,2$ use a block $S^1\times D^1\times D^4$ with coordinates
\[
 z=\exp(s+2\pi i\theta),\quad -R\leq s\leq R,
 \qquad (a,b)\in D^4,\qquad q=zab.
 \tag{E.H2}\label{maps:eq:E.H2}
\]
The normal four-ball here is a rounded radial domain
$|a|,|b|\leq\epsilon_e$, $|ab|\leq\eta/|z|$;
its positive radial trivialization identifies it with $D^4$.
The ends attach to framed tubes about the following coordinate
circles.  A table entry $(Z_j;Z_k,Z_l)$ means that the curve
coordinate is $Z_j$ and the two ordered complex normal coordinates
are $Z_k,Z_l$; the same convention applies to $W$.
\[
\begin{array}{c|c|c}
 i&z=0\hbox{ end: }(z;a,b)&z=\infty\hbox{ end: }(w;A,B)\\ \hline
0&(Z_2;Z_1,Z_0)&(W_0;W_1,W_2)\\
1&(W_2;W_1,W_0)&(Z_1;Z_2,Z_0)\\
2&(Z_0;Z_2,Z_1)&(W_1;W_2,W_0)
\end{array}
\tag{E.H3}\label{maps:eq:E.H3}
\]
The full-dimensional endpoint map between the two end charts is
\[
                 (w,A,B)=(z^{-1},za,zb).
 \tag{E.H4}\label{maps:eq:E.H4}
\]
In particular the second attaching circle has angle $-\theta$,
and its two complex normal vectors each rotate by $+2\pi\theta$.
The factors $|z|$ are the specified positive radial rescalings.
They do not erase either angular rotation.  Every coordinate circle
of the two balls occurs exactly once in \eqref{maps:eq:E.H3}.  The maps are
embeddings of $S^1\times\partial D^1\times D^4$, not just attaching
loops with an unspecified choice of normal framing.

\emph{3. The dense torus block and its full sector maps.}
Attach $T^2\times D^2\times D^2_q$ along
$T^2\times\partial D^2\times D^2_q$.  Here is a completely explicit
choice of its base disk and the attaching map.  Use the rays
$r_i$ of \eqref{maps:eq:E.T4} and the hexagon
\[
 \mathcal P=\{v\in\mathbb R^2:\langle v,r_i\rangle\geq-1\},
 \quad
 \operatorname{Vert}(\mathcal P)=
 \{(-1,-1),(0,-1),(1,0),(1,1),(0,1),(-1,0)\}.
 \tag{E.H5}\label{maps:eq:E.H5}
\]
On the positive real part of $(\mathbb C^*)^2$, use the coordinates
given by the toric moment map
\[
 \mu(X,Y)=
 \frac{\displaystyle\sum_{v\in\operatorname{Vert}(\mathcal P)\cup\{0\}}
                 |X^{v_1}Y^{v_2}|^2v}
      {\displaystyle\sum_{v\in\operatorname{Vert}(\mathcal P)\cup\{0\}}
                 |X^{v_1}Y^{v_2}|^2}.
 \tag{E.H6}\label{maps:eq:E.H6}
\]
Choose a small inward parallel hexagon, with its corners rounded
symmetrically, $\mathcal P_\epsilon\subset\operatorname{int}\mathcal P$.
The inverse image is $T^2\times\mathcal P_\epsilon$; its two circle
coordinates are $\arg X,\arg Y$.  On the boundary disk circle put
$(X,Y)$ equal to the positive inverse of $\mu$, multiplied by these
two phases.  On the six sectors of that boundary, the attaching
map into the preceding blocks is the following monomial map, with
the whole complex normal coordinate $q$ included:
\[
\begin{array}{c|ccc}
 i& n_i&a_i&a_{i+1}\\ \hline
0&q/(XY)&X&Y\\
1&q/Y&XY&X^{-1}\\
2&qX&Y&(XY)^{-1}\\
3&qXY&X^{-1}&Y^{-1}\\
4&qY&(XY)^{-1}&X\\
5&q/X&Y^{-1}&XY
\end{array}
\qquad q=n_i a_i a_{i+1}.
\tag{E.H7}\label{maps:eq:E.H7}
\]
Use the cyclic identifications
$(n_2,a_2,a_3)=(Z_1,Z_2,Z_0)$,
$(n_4,a_4,a_5)=(Z_2,Z_0,Z_1)$ and the analogous two cyclic
permutations of $W$ for vertices $3,5$.  Vertices $0,1$ are $Z,W$.
The chart at an edge is $(z,a,b)=(a_{i+1},a_i,n_i)$, with \eqref{maps:eq:E.H4}
at its other end.  These rules specify \eqref{maps:eq:E.H7} also where a sector
meets an annular block.  The transition on a common sector is the
monomial map $P_{i+1}^{-1}P_i$, where
$P_i=((0,1),(r_i,1),(r_{i+1},1))$; the sum of its exponent rows
is $(1,1,1)$, so it preserves $n_i a_i a_{i+1}$.  Thus the sector maps specify the
normal framing as well as the core map.  They have no unassigned
phase, clutching map, or choice of normal line.

\emph{4. Expansion into ordinary handles.}
We use the following fixed elementary rule on every parameterized
block in this prescription:
\[
 T^r\times D^k\times D^{n-r-k}
 \quad\hbox{attached on}\quad
 T^r\times\partial D^k\times D^{n-r-k}.
 \tag{E.H8}\label{maps:eq:E.H8}
\]
Cut each ordered parameter circle into its two closed semicircles,
with small endpoint collars.  The first semicircle is its zero
handle and the second is its one handle.  For a bit vector
$\varepsilon\in\{0,1\}^r$, the resulting ordinary handle is of
index $k+|\varepsilon|$.  Attach by increasing $|\varepsilon|$,
then lexicographically.  On the original attaching face use
the displayed full-dimensional attaching embedding of the block;
on a newly exposed parameter face use the identity endpoint map
in that circle.  The normal coordinates are the remaining ordered
semicircle coordinates followed by the original normal coordinates.
Corner rounding and the permutation putting negative coordinates
first complete an ordinary framed attaching embedding
$S^{k+|\varepsilon|-1}\times D^{n-k-|\varepsilon|}$.
This is a finite rule for the full maps, not an instruction to choose
some attaching map in a stated homotopy class.  For one parameter
it is the literal two-piece decomposition in
\cite[Section~4.1, equation~(16), p.~103]{mapsVogel09}.
Its iteration requires no classification theorem.

Apply \eqref{maps:eq:E.H8} first to \eqref{maps:eq:E.H2}, with $(r,k,n)=(1,1,6)$, and then
to the core, with $(r,k,n)=(2,2,6)$.  Together with the two balls
this gives cusp handle counts $(2,3,4,2,1)$.

\emph{5. The torus bands and the finite caps.}
Take the marked flat torus bundle of \ref{maps:src:E.2.20} over the disc with
three holes, regarding the cusp boundary as the incoming boundary.
Cut along the two marked disjoint arcs from the two finite boundary
circles to the outer circle.  The cut surface is an annulus;
its torus bundle is the collar of the already constructed cusp
boundary.  Restoring the two strips attaches two blocks
\[
                 T^4\times D^1\times D^1.
 \tag{E.H9}\label{maps:eq:E.H9}
\]
To fix their maps without a planar drawing, use the oriented ribbon
polygon whose successive based inner boundary words are
$A_1,A_2,M_0$ and whose outer word is their product.  The negative
bank of the $j$th cut has coordinate $(x,t)$, the positive bank
has $(A_jx,t)$, and the strip uses the product normal coordinate.
Reversing a bank orientation replaces that matrix by its inverse.
This is the cut-polygon definition of the flat bundle in \ref{maps:src:E.2.20};
the identity $A_1A_2M_0=I$ specifies the unfilled outer boundary.
There is no affine term.  Thus \eqref{maps:eq:E.H9} has its full attaching
embedding, including the transverse interval, fixed by these
linear maps.  Expand each by \eqref{maps:eq:E.H8}, with $(r,k,n)=(4,1,6)$.
The unused boundary circles are precisely the two finite boundaries
and the marked ordinary boundary.
For clarity, the initial bank coordinates on the already attached
cusp boundary are the explicit coordinates of E.2.19:
\[
 \begin{gathered}
 q=\eta e^{-2\pi i\theta},\qquad T=-\log\eta/(2\pi),\\
 (X,Y)=\exp\bigl(2\pi i((-\theta+iT)B_0a+b)\bigr),
 \qquad a=(g,u),\ b=(w,d).
 \end{gathered}
 \tag{E.H9a}\label{maps:eq:E.H9a}
\]
The exponential is componentwise; at an interval seam use
$b\mapsto b-B_0a$.  Equations \eqref{maps:eq:E.H7} and \eqref{maps:eq:E.H9a} therefore
give the actual bank embeddings into the existing handle
boundary, including their derivatives.  They do not merely
identify that boundary abstractly as a mapping torus.

For the $j$th finite cap use its already explicit
$\mathcal S_j\times D^2_r$ and the entire boundary embedding
\eqref{maps:eq:E.U3}.  The mapping torus $\mathcal S_j$ has the following product
block construction: a block $T^3\times D^1$ over the first
semicircle of $\sigma$, followed by a block $T^3\times D^1$ over
the second semicircle, with endpoint maps $I$ and $B_j^{-1}$.
The first is a three-parameter zero handle, the second a
three-parameter one handle.  As a filling relative to
$\mathcal S_j\times S^1_r$, the normal disk supplies two negative
directions.  Consequently the relative cap consists of
\[
 \begin{array}{c|c|c}
 \hbox{base-circle block}&(r,k,n)&\hbox{full attaching data}\\ \hline
 0&(3,2,6)&\hbox{\eqref{maps:eq:E.U3}, the first semicircle, and the $r$-disk}\\
 1&(3,3,6)&\hbox{\eqref{maps:eq:E.U3}, endpoint maps $I,B_j^{-1}$, and the $r$-disk}.
 \end{array}
 \tag{E.H10}\label{maps:eq:E.H10}
\]
More explicitly an attaching face in the normal direction is sent
to \eqref{maps:eq:E.U3} at $r=e^{2\pi i\alpha}$, and an attaching face in a
base direction is the identical face of the preceding base block,
crossed with $D^2_r$.  At the base seam use $B_j^{-1}$ on $l$ and
the identity on $r$.  The inverse angular formula \eqref{maps:eq:E.U4} supplies
the full map back to the boundary being filled.  In particular
the normal meridian is $v_j-m_j\rho_j$, not the base meridian.
Apply \eqref{maps:eq:E.H8} to the two blocks in \eqref{maps:eq:E.H10}.  Each cap adds the
ordinary relative counts $(1,4,6,4,1)$ in indices $2,\ldots,6$.
This completes the definition of $\mathcal H_6$.

\emph{6. The closed ten-dimensional handlebody.}
For each handle of $\mathcal H_6$, take its product with the four
ordered torus circles and apply \eqref{maps:eq:E.H8}.  This gives
$\mathcal H_6\times T^4$ with its full boundary framing.  On the
other side take the dual relative handle decomposition of
$T^4\times\mathcal H_6$: reverse the handle order and exchange
the negative and positive disks of every handle.  A handle of
index $j$ becomes one of index $10-j$, and its attaching map is
the actual belt embedding of the original handle.  Glue its
initial boundary by the explicit map
\[
                 (x_1,x_2,\theta)\longmapsto
                                    (x_1,x_2,-\theta).
 \tag{E.H11}\label{maps:eq:E.H11}
\]
All subsequent dual attaching maps are obtained by this fixed
exchange of the two coordinate disks, in reverse order, on the
previously specified embeddings.  Use the factor order $E,T^4$
on the first half and $T^4,E$ on the second; their interchange
has sign $(-1)^{24}=1$.  The boundary-angle reversal in \eqref{maps:eq:E.H11}
is the orientation-reversing boundary gluing.  This defines
$\mathcal H_{10}$ entirely by ordinary handles.  If
\[
 c(t)=2+5t+14t^2+22t^3+21t^4+10t^5+2t^6,
 \qquad
       c_{10}(t)=(c(t)+t^6c(t^{-1}))(1+t)^4,
 \tag{E.H12}\label{maps:eq:E.H12}
\]
the coefficients are the counts in the statement.  The expanded
lists of all $76$ and $2432$ handles, including their parameter
bits, index, coindex, side, duality and attaching-map macro, are
\path{code/III/maps_ordinary_handles.out.json}.  Equations
\eqref{maps:eq:E.H1}--\eqref{maps:eq:E.H11}, \eqref{maps:eq:E.U3}--\eqref{maps:eq:E.U4} and the displayed integer matrices
are the definition of those macros.

\end{theorem}


\begin{proof}
\begin{enumerate}[label=\arabic*.,leftmargin=*]
\item The main geometric point is that the first three steps give
\emph{actual smooth blocks of the cusp tube}.  We give the collar
argument explicitly.  Work in the smooth toric quotient of \ref{maps:src:E.2.19}.
Choose disjoint small coordinate neighborhoods of the two triple
points.  The three incident coordinate axes in each neighborhood
are exactly the six ends in \eqref{maps:eq:E.H3}.  Remove smaller endpoint disks
from each normalized double curve.  Its remaining annulus has
the two complex normal lines $\mathcal O(-1)|_{\rm annulus}$,
which are trivialized by $(a,b)$ in \eqref{maps:eq:E.T6}.  The resulting normal
domain in the actual tube is exactly the domain in \eqref{maps:eq:E.H2}.
Every ray in $\mathbb C^2$ meets its rounded boundary once and
transversely.  Indeed the three defining radius functions have
positive radial derivative wherever the corresponding boundary
face is active.  Positive radial rescaling therefore gives a
smooth family of four-balls as $z$ ranges over this compact
annulus.  Choose the radial change to be constant linear scaling
on a common sufficiently small ball about the origin and interpolate
monotonically only away from that ball; it is then smooth at the
origin as well as at the boundary and preserves the stated normal
frames.  The same argument for the coordinate and product
radius functions in \eqref{maps:eq:E.H1} gives six-balls at the triple points.
Rounding can be chosen constant on the end attaching collars.

\item There is no unverified compatibility of these end frames.  At
vertex $i$ the normal-first coordinates are $(n_i,a_i,a_{i+1})$.
At vertex $i+2$ the deck identification is their cyclic
permutation, and the threefold composition is the identity by
\ref{maps:src:E.2.19}.  At the two ends of $C_i$ the coordinates are
$(z,a,b)=(a_{i+1},a_i,n_i)$ and
$(w,A,B)=(a_{i-1},a_i,n_{i-1})$.  The ordinary
$\mathcal O(-1)^2$ transition is \eqref{maps:eq:E.H4}.  Substituting the two
cyclic permutations gives exactly \eqref{maps:eq:E.H3}.  This checks all six
ends, not just one representative.  Triviality of a normal bundle
over the annulus is not being used to replace its relative
endpoint framing by a constant framing.

\item We next choose the stratum collars so that their remaining
interface is the entire attaching face of the dense block.
On the normalized toric surface take a small collar of its
hexagonal toric boundary, choosing the corner neighborhoods
first and the edge strips second.  Choose the strips narrower
than the corner collars and invariant under the opposite-edge
inversion and the two cyclic triple identifications.  The
uncovered moment region is the rounded interior disk
$\mathcal P_\epsilon$.  This can be done directly with six
inward parallel affine inequalities in \eqref{maps:eq:E.H5}, smoothing each
corner inside its already chosen corner neighborhood.  Thus the
common frontier in the smooth part of the normalization is
exactly $T^2\times\partial\mathcal P_\epsilon$, not an additional
collection of annuli or disks.

\item The map \eqref{maps:eq:E.H6} gives the claimed coordinates: in logarithmic
positive coordinates it is the gradient of a log-sum-exponential
function.  Its Hessian is twice the covariance matrix of the
seven affinely spanning exponent vectors, hence positive
definite.  Strict convexity and the growth in each supporting
direction give a diffeomorphism onto
$\operatorname{int}\mathcal P$.  The inverse on the compact
disk and its frontier is smooth.  The canonical line on this
free torus locus is trivial; more specifically the toric
function $q$ is a nonzero toric factor times its normal
coordinate.  Hence \emph{$q$ itself}, including its phase,
trivializes that normal line.  The portion of the truncated
tube over this disk is therefore exactly
\[
 (T^2\times\mathcal P_\epsilon)\times D^2_q,
 \quad\hbox{with interface}\quad
 (T^2\times\partial\mathcal P_\epsilon)\times D^2_q.
 \tag{E.H13}\label{maps:eq:E.H13}
\]
In extending the preceding moment cuts off the zero fibre,
hold their equations constant in this $q$ coordinate on an
open neighborhood of the compact frontier.  On their ends
use the product cuts of the annular and triple charts.  A
collar interpolation between these choices is an interpolation
of positive radial cut functions, fixed on each already chosen
end collar.  Their radial derivatives remain positive, so its
flow is a collar isotopy.  The cuts can consequently be chosen
to agree with \eqref{maps:eq:E.H13} exactly.  In particular no twisting of
the $q$ disk is introduced on the six sector seams.

\item These are blocks in a sufficiently small actual truncated
neighborhood of the compact central fibre.  Their central
parts cover that fibre: two triple points, the three open
double annuli and the remaining normalized surface core.
Compactness from \ref{maps:src:E.2.19} permits one common sufficiently small
$\eta$: the union of their open chart neighborhoods contains
the compact zero fibre, so properness implies that it contains
the entire inverse image of a sufficiently small closed
$q$-disc.  The nested cuts partition this entire inverse
image, rather than just a chosen sub-neighborhood of it.
If a larger originally chosen tube is used, the
difference is an exterior collar: on the compact annular
region without critical points choose a vector field lifting
the radial direction of $q$, and integrate it.  On the
prescribed outer collar choose the lift in the marked real
torus coordinates of \ref{maps:src:E.2.19}.  It is then the identity on all
those marked coordinates during radial collar transport.
This proves that no additional component or handle, and no
unrecorded outer-boundary diffeomorphism, is being omitted.

\item The actual embedding in \eqref{maps:eq:E.H13} is \eqref{maps:eq:E.H7}.  To check this,
the exponent matrix for a unimodular chart is $P_i^{-1}$,
applied to $(X,Y,q)$.  This gives each of the six displayed
rows.  Consecutive changes are $P_{i+1}^{-1}P_i$ and telescope
to the identity.  The translations identifying the two
triple orbits are exactly those already used in \eqref{maps:eq:E.H3}.
All six normal coordinate multipliers in \eqref{maps:eq:E.H7} are thereby
fixed, and the sector seams are equal maps on overlaps,
not just homotopic maps.  Positive rescaling of the chosen
radial domains leaves these monomials unchanged before
the collar reparametrization.  This completes the smooth,
fully framed identification of the cusp with steps 1--3.

\item For completeness the handle expansion rule is itself a
geometric equality.  For one circle, write $S^1=I_0\cup I_1$.
First attach $I_0\times D^k\times D^{n-k-1}$; after rounding it
is a $k$-handle.  The remaining block is
$(I_1\times D^k)\times D^{n-k-1}$, a $(k+1)$-handle.
Its attaching face consists of the original negative face
and its two endpoint faces in the first handle.  The stated
maps on these three faces agree on their corners because
they are restrictions of the same product coordinates.
Iterate this argument on the ordered circle factors.
It proves \eqref{maps:eq:E.H8}, including all newly created attaching
faces and their normal framings.

\item The planar core relative to its incoming circle has two
ordinary surface one-handles.  Its flat bundle is a product
over each strip, so restoring the two cut strips is exactly
\eqref{maps:eq:E.H9}.  The zero section and the full lattice marking fix
the normal coordinates on these attachments.  Each cap in
\eqref{maps:eq:E.H10} is the actual smooth product of \ref{maps:src:E.2.20}.  Its relative
handle claim can also be checked by the function
$-|r|^2+\varepsilon f(l,\sigma)$, with a cutoff near the
incoming boundary: choose $f$ by a two-critical-point
function of $\sigma$, perturbed on its two critical tori
by the product Morse function on $T^3$.  On the remaining
arcs its base derivative dominates the perturbation, so
there are no further critical points.  The indices are
two plus the base indices, as in \eqref{maps:eq:E.H10}.  The product
description above fixes its handles without needing this
alternative Morse verification.  Summing the actual blocks
gives
\[
 (2,3,4,2,1,0,0)+2(0,1,4,6,4,1,0)
                     +2(0,0,1,4,6,4,1)
       =(2,5,14,22,21,10,2).
\]
Every comparison has its prescribed boundary value, so the
relative diffeomorphism of \ref{maps:src:E.2.20} gives the first assertion
of \eqref{maps:eq:E.H0}.  Taking the products, the dual relative handles
and the actual boundary flip of \ref{maps:src:E.2.20} gives the second.
This is an identification by explicit full-dimensional
embeddings of ordinary handles.  It neither asserts an
unproved connected-sum normal form nor leaves a cap called
``the original complement'' inside the prescription.

\item The exact checks of all exponent matrices, the six endpoint
placements, both normal angular rotations, all sector seams,
the total linear monodromy, and the expanded handle lists
are in \path{code/III/maps_ordinary_handles.py}.  Its saved output
has Euler numbers $2$ and $0$, respectively, as a check;
the smooth identification follows from the preceding block
and collar proof, not from these numerical equalities.

\item We finish by stating precisely why these data determine the smooth
type, including its top handles.  The relevant relative uniqueness
results are Milnor's Theorem~3.13, which identifies an elementary
cobordism from its full characteristic embedding, Theorem~1.4,
which makes smooth composition unique with the marked boundary
fixed, and the isotopy-extension Theorem~5.8
\cite{mapsMilnor65}.  In the present notation the characteristic
embedding is the whole map
$S^{j-1}\times D^{n-j}\longrightarrow\partial W$, not just
the attaching sphere or its homology class.  Its normal coordinates
are part of \eqref{maps:eq:E.H3}--\eqref{maps:eq:E.H11}.  The elementary consequence needed
here can also be seen directly.  If $\phi_0,\phi_1$ are joined
by the specified isotopy of full attaching embeddings, extend
the velocity field on their images to a slightly larger boundary
collar and multiply it by a cutoff which is zero on the protected
marked boundary.  Its flow extends the attaching isotopy to an
ambient boundary isotopy.  Extending this flow across a product
collar of $W$, and using the identity on the new coordinate handle,
identifies $W\cup_{\phi_0}(D^j\times D^{n-j})$ with
$W\cup_{\phi_1}(D^j\times D^{n-j})$.  The same cutoff argument
applies to the small radial corner changes above.  This derives
the relative, framed version from the stated isotopy-extension
theorem, without assuming an extension of an arbitrary sphere
diffeomorphism over a disk.

\item Apply this comparison successively in the displayed block and
sector order.  The initial ball comparisons are their coordinate
identifications.  At each next stage the actual attaching map
and the prescribed attaching map agree in the computed normal
coordinates, or differ by the explicitly described positive
collar isotopy.  The comparison of that stage is consequently
the identity on the new coordinate handle and the product-collar
extension on the previous stages.  Theorems~3.13 and~1.4 give an
inductive relative diffeomorphism of the whole six-manifold.
The final ordinary boundary coordinates are unchanged.  The
same argument, with the actual belt embeddings and the specified
boundary flip, applies to every handle on the dual ten-dimensional
side.  In particular a top handle retains its entire attaching
map $S^{n-1}\to\partial W$; none is replaced by an arbitrary
capping disk after forgetting that map.  There is therefore no
remaining disk-gluing or exotic-sphere ambiguity.  The computed
ordered full framed attaching data, under this relative uniqueness
theorem, are the classifying data used for \eqref{maps:eq:E.H0}; the preceding
algebraic invariants are not being used as a smooth classification.
\end{enumerate}
\end{proof}





\section{Local covers, global quotients and cusp carriers}\label{sec:III:carriers}

\subsection{The Kummer cover extends across the toric cusp}
\label{wb:cusp-cover-section}

This section determines exactly what is true in the assertion that the
degree-$N^2$ cover of the regular fibres extends by a toric power map.  The
answer is positive for the completed \emph{local} cusp family, but the source
action has to be changed, and two different degrees must not be conflated.

Retain the cusp data
\[
 B_0=\begin{pmatrix}0&1\\-1&0\end{pmatrix},\qquad
 Z(s)=sB_0+C(t),\qquad t=e^{2\pi i s},
\]
where $C$ is holomorphic on a disc $\Delta_t$ about zero.  Put
\[
 c_{\bar\lambda}(t)
 :=\exp\!\bigl(2\pi iC(t)\bar\lambda\bigr)
 \in(\mathbb C^*)^2,
 \qquad \bar\lambda\in\bar\Lambda\simeq\mathbb Z^2.
\]
Let $\Sigma$ be the fan in
$N'=\mathbb Z^2\oplus\mathbb Z$ consisting of the cones over the cells of
the translation-invariant $A_2$ triangulation at height one, and write
$Y=Y_\Sigma$.  On the dense torus, with coordinates $(x_1,x_2,t)$, the
target action is
\begin{equation}
 \Psi_{\bar\lambda}(x,t)
 =\bigl(c_{\bar\lambda}(t)t^{B_0\bar\lambda}x,t\bigr).
 \label{wb:cusp-cover-target-action}
\end{equation}
Its quotient is the completed cusp family
$f_0\colon N_0=Y/\bar\Lambda\to\Delta_t$.

Fix an integer $N\geq1$.  Let $v$ be a coordinate on a second disc and
make the base change
\begin{equation}
 \rho_N\colon\Delta_v\longrightarrow\Delta_t,\qquad
 \rho_N(v)=v^N,\qquad t=v^N.
 \label{wb:cusp-cover-base-change}
\end{equation}
On a copy $Y^{(N)}$ of the same toric variety, with dense-torus coordinates
$(\xi_1,\xi_2,v)$, define
\begin{equation}
 C_N(v):=\frac1N C(v^N),\qquad
 c^{(N)}_{\bar\lambda}(v)
 :=\exp\!\bigl(2\pi iC_N(v)\bar\lambda\bigr)
 \label{wb:cusp-cover-source-factor}
\end{equation}
and
\begin{equation}
 \Psi^{(N)}_{\bar\lambda}(\xi,v)
 :=\bigl(c^{(N)}_{\bar\lambda}(v)
          v^{B_0\bar\lambda}\xi,v\bigr).
 \label{wb:cusp-cover-source-action}
\end{equation}

\begin{theorem}[Exact Kummer extension through the cusp]
\label{wb:cusp-cover-extension-theorem}
\StatementRecord{wb:cusp-cover-extension-theorem}{proved}{Historical events: GPT-6 Astra (OpenAI), ChatGPT Pro mode, effort Pro, S6\_09\_PAPERS, 17 November 2026: targeted proof check. Also checked in this writing. Pre-continuation writing events: GPT-6 Astra (OpenAI), ChatGPT Pro mode, effort Pro, S6\_10\_PAPERS, 18 November 2026: targeted proof check. Continuation metadata events: GPT-6 Astra (OpenAI), ChatGPT Pro mode, effort Pro, S6\_10\_PAPERS continuation, 19 November 2026: used and rechecked in this run; status-vocabulary reconciliation; metadata reconciliation only; no new mathematical check. \emph{Pre-continuation recorded proof scope.} GPT-6 Astra (OpenAI, ChatGPT Pro mode), S6\_09\_PAPERS, effort Pro, 2026-11-17; also checked in this writing by GPT-6 Astra (OpenAI), ChatGPT Pro mode, effort Pro, 18 November 2026}{the hypotheses, notation and definitions stated here; \cref{wb:cusp-cover-source-action}; \cref{wb:cusp-cover-target-action}; \cref{wb:cusp-cover-source-factor}; \cref{wb:cusp-cover-maximal-chart-map}; \cref{wb:cusp-cover-basechange-local-ring}; \cref{wb:cusp-cover-triple-fibre-length}; \cref{wb:cusp-cover-double-fibre-length}; \cref{wb:cusp-cover-additive-fibre-map}; \cref{wb:cusp-cover-equivariance-computation}}

\mbox{}\ 
After shrinking the $v$-disc, the following statements hold.

\begin{enumerate}
\item The maps in \eqref{wb:cusp-cover-source-action} extend to a free,
properly discontinuous action on $Y^{(N)}$.  The quotient
\[
 f_N\colon N_0^{\langle N\rangle}
 :=Y^{(N)}/\bar\Lambda\longrightarrow\Delta_v
\]
is a proper holomorphic map from a smooth threefold.  Its central fibre
$W_N$ is reduced, irreducible, and nonnormal, and its normalization is a
toric del Pezzo surface of degree six.

\item Multiplication by $N$ on the cocharacter lattice defines a finite
flat toric morphism
\begin{equation}
 P_N\colon Y^{(N)}\longrightarrow Y,\qquad
 P_N(\xi_1,\xi_2,v)=(\xi_1^N,\xi_2^N,v^N),
 \label{wb:cusp-cover-toric-power-map}
\end{equation}
of degree $N^3$.  It is equivariant for
\eqref{wb:cusp-cover-source-action} and
\eqref{wb:cusp-cover-target-action}, hence descends to a finite flat map
\begin{equation}
 \kappa_N\colon N_0^{\langle N\rangle}\longrightarrow N_0,
 \qquad f_0\circ\kappa_N=\rho_N\circ f_N,
 \label{wb:cusp-cover-completed-map}
\end{equation}
of degree $N^3$.

\item The map $\kappa_N$ is etale over $v\ne0$, and
\begin{equation}
 \kappa_N^*W=NW_N,
 \qquad
 \operatorname{Ram}(\kappa_N)=(N-1)W_N.
 \label{wb:cusp-cover-direct-ramification}
\end{equation}
For $N>1$, the branch support of the map to the original cusp family
is the whole central fibre. For $N=1$, the displayed formulas give the identity and the branch support is empty.

\item Let
\begin{equation}
 N_0[N]:=N_0\times_{\Delta_t}\Delta_v,
 \qquad t=v^N.
 \label{wb:cusp-cover-base-changed-target}
\end{equation}
Then $N_0[N]$ is normal and $\kappa_N$ factors uniquely as
\begin{equation}
 N_0^{\langle N\rangle}
 \xrightarrow{\ Q_N\ }N_0[N]
 \longrightarrow N_0.
 \label{wb:cusp-cover-relative-factorization}
\end{equation}
The map $Q_N$ has generic degree $N^2$.  For $N>1$ it is quasi-etale but
not flat: it is etale over the smooth locus of $N_0[N]$, and its non-etale
and nonflat loci are exactly the inverse images of the double and triple
strata of the base-changed central fibre.  For $N>1$, its degree $N^2>1$ proves that $Q_N$ is not a
normalization map. For $N=1$ it is the identity of the normal target.

\item On the reduced central fibres, the induced map
$W_N\to W$ has generic degree $N^2$.  It lifts to a commutative diagram
\begin{equation}
\begin{CD}
 E_N\simeq\mathrm{dP}_6 @>{[N]}>> E\simeq\mathrm{dP}_6\\
 @V{\nu_N}VV @VV{\nu}V\\
 W_N @>>> W,
\end{CD}
 \label{wb:cusp-cover-normalization-diagram}
\end{equation}
where $[N]$ is toric multiplication by $N$, of degree $N^2$, ramified
with index $N$ along each of the six toric boundary curves.  The diagram
respects the paired conductor curves and the two triple-point classes.

\item Over $v\ne0$, $Q_N$ is exactly the cover defined by
\[
 \mathbb Z\langle\widehat\gamma,\widehat u,
 N\widehat w,N\widehat\delta\rangle
 \subset
 \mathbb Z\langle\widehat\gamma,\widehat u,
 \widehat w,\widehat\delta\rangle.
\]
Its deck group is $(\mathbb Z/N\mathbb Z)^2$.  The pulled-back cusp
monodromy is $M_0^N$, and in the displayed cover basis it is $M_0$:
\begin{equation}
 S_N^{-1}M_0^NS_N=M_0,\qquad
 S_N=\operatorname{diag}(1,1,N,N).
 \label{wb:cusp-cover-monodromy-conjugacy}
\end{equation}
\end{enumerate}
\end{theorem}



\begin{proof}
\begin{enumerate}[label=\arabic*.,leftmargin=*]
\item We prove the assertions in the order in which the constructions occur.

\item \medskip\noindent
\emph{The source action.}
For $\bar\lambda\in\bar\Lambda$, the integral linear map
\[
 \varphi_{\bar\lambda}(y,y_3)
 =(y+y_3B_0\bar\lambda,y_3)
\]
translates the height-one slice by the integral vector
$B_0\bar\lambda$.  It preserves the $A_2$ triangulation and hence induces
a toric automorphism $\Phi_{\bar\lambda}$ of $Y^{(N)}$.  On the dense
torus this automorphism is
\[
 \Phi_{\bar\lambda}(\xi,v)
 =(v^{B_0\bar\lambda}\xi,v).
\]
Multiplication by the nowhere-zero holomorphic function
$c^{(N)}_{\bar\lambda}(v)$ extends across every toric chart.  Since
$C_N(v)\bar\lambda$ is additive in $\bar\lambda$, the resulting maps
satisfy
$\Psi^{(N)}_{\bar\lambda}\Psi^{(N)}_{\bar\mu}
=\Psi^{(N)}_{\bar\lambda+\bar\mu}$.

\item For completeness, freeness and proper discontinuity can be checked
directly.  Put
\[
 R_N(v):=-2\pi\operatorname{Im}C_N(v).
\]
It is bounded at $v=0$.  At a dense-torus point fixed by
$\bar\lambda$, logarithmic absolute values give
\[
 \log|v|B_0\bar\lambda+R_N(v)\bar\lambda=0.
\]
Choose the disc so small that
$|\log|v||>2\|R_N\|_\infty$.  Because $B_0$ is orthogonal, the last
equality is impossible for $\bar\lambda\ne0$.  On the central fibre,
$\Psi^{(N)}_{\bar\lambda}$ translates every bounded cell by the nonzero
vector $B_0\bar\lambda$, so it fixes no toric orbit.

\item For a dense-torus point set
\[
 y(\xi,v)=\frac{(\log|\xi_1|,\log|\xi_2|)}{\log|v|},
 \qquad
 B_{N,v}=B_0+\frac{R_N(v)}{\log|v|}.
\]
Then
\[
 y(\Psi^{(N)}_{\bar\lambda}(\xi,v))
 =y(\xi,v)+B_{N,v}\bar\lambda,
 \qquad
 |B_{N,v}\bar\lambda|\ge\tfrac12|\bar\lambda|.
\]
% checked carrier correction: compact-chart-control
The logarithmic positions of an entire affine toric chart need not be
bounded. The required assertion is about its compact subsets, and follows
as follows. In a maximal unimodular chart write its height-one rays as
$(b_0,1),(b_1,1),(b_2,1)$ and its coordinates as $w_0,w_1,w_2$.
On its dense torus,
\[
 v=w_0w_1w_2,\qquad
 y=\frac{\log|\xi|}{\log|v|}
   =\sum_{i=0}^2\lambda_i b_i,\qquad
 \lambda_i=\frac{\log|w_i|}{\log|v|},\qquad
 \sum_i\lambda_i=1.
\]
For a compact subset of this chart all $|w_i|$ have a common upper
bound $M\ge1$. When $|v|<1$, this gives
$\lambda_i\ge-\log M/|\log|v||$. Thus its logarithmic positions remain
in a bounded neighbourhood of the fixed triangle as $v\to0$.
On a closed annulus away from zero, the product identity and upper
bounds also bound each $|w_i|$ below. Hence the same boundedness holds
there. Every point of a compact set has a neighbourhood whose closure
is compact in one maximal chart.  Choose a finite such cover;
the intersections with these compact closures give finitely many
compact chart pieces covering the original set.  Applying the
preceding estimate to these pieces bounds all of its logarithmic
positions on the dense-torus locus.

\item For the modified action,
$y\mapsto y+B_{N,v}\bar\lambda$ and
$|B_{N,v}\bar\lambda|\ge|\bar\lambda|/2$.
Consequently only finitely many translates of one compact set can
meet another on the dense torus. At $v=0$, choose finitely many
maximal charts covering the two compact sets. A translate meeting
both must send a face of one of their triangles to a face of another.
There are finitely many pairs of faces and at most one translation
for each pair; thus the same finiteness holds at the central fibre.
This proves proper discontinuity without asserting boundedness of
whole affine charts.

\item To prove that $f_N$ is proper, reduce
$B_{N,v}^{-1}y$ modulo $\mathbb Z^2$ for $v\ne0$. The resulting $y$
lies in one bounded set, uniformly on the chosen small disc. Choose
an $A_2$ triangle containing $y$. Its barycentric coordinates are
nonnegative, and the corresponding toric coordinates satisfy
$|w_i|\le1$. Only finitely many triangles meet that bounded set.
Fix a closed subdisc strictly inside the chosen action disc.
Every orbit over that subdisc has a representative in finitely
many closed unit polydiscs intersected with its inverse image.
These intersections are closed subsets of compact polydiscs
in the ambient smooth Hausdorff toric space, and all lie in the
domain where the modified action is defined. At $v=0$, vertex translations move a component containing the
point to a fixed component; that component is the compact toric
surface with hexagonal star. Adding this compact component gives
a compact set mapping onto the quotient above a closed subdisc.
The inverse image of any compact subset of the base is a closed
subset of this compact image, and is therefore compact.

\item The central-fibre and
normalization assertions follow exactly from the height-one fan: locally
$v$ is the product of the coordinates belonging to the rays through a
cell, every component occurs with multiplicity one, the lattice action is
transitive on vertices, and the star of a vertex is the complete smooth
hexagonal fan.

\item \medskip\noindent
\emph{The toric map and equivariance.}
Multiplication by $N$ on $N'$ sends every cone $\sigma\in\Sigma$ onto
itself as a real cone.  It therefore defines the toric morphism $P_N$.
For a maximal cone, let $z_0,z_1,z_2$ be the target characters dual to its
three primitive height-one ray generators, and let $w_0,w_1,w_2$ be the
corresponding source characters.  The map on convergent local rings is
\begin{equation}
 \mathbb C\{z_0,z_1,z_2\}\longrightarrow
\mathbb C\{w_0,w_1,w_2\},\qquad z_i\longmapsto w_i^N.
 \label{wb:cusp-cover-maximal-chart-map}
\end{equation}
The target-chart preimage is the corresponding source chart, and the
right-hand ring is free of rank $N^3$ over the left-hand ring, with basis
\[
 \{w_0^{a_0}w_1^{a_1}w_2^{a_2}:0\le a_i<N\}.
\]
Thus $P_N$ is finite flat of degree $N^3$ on every maximal chart, hence
globally.

\item The choice \eqref{wb:cusp-cover-source-factor} satisfies equivariance by the following exact identity.
Indeed, writing powers coordinatewise,
\begin{align}
 P_N\Psi^{(N)}_{\bar\lambda}(\xi,v)
 &=\bigl((c^{(N)}_{\bar\lambda}(v))^N
          v^{NB_0\bar\lambda}\xi^N,v^N\bigr)\\
 &=\bigl(c_{\bar\lambda}(v^N)
          (v^N)^{B_0\bar\lambda}\xi^N,v^N\bigr)\\
 &=\Psi_{\bar\lambda}P_N(\xi,v).
 \label{wb:cusp-cover-equivariance-computation}
\end{align}
Consequently $P_N$ descends to \eqref{wb:cusp-cover-completed-map}.  To see
finiteness and flatness after the infinite-lattice quotient without any
implicit descent theorem, take an evenly covered open set in $N_0$ and
one of its lifts $V\subset Y$.  Equivariance gives a unique representative
of every source orbit above that open set whose image lies in $V$; two
such representatives cannot differ by a nontrivial lattice translate.
The descended map is therefore locally identified with
$P_N^{-1}(V)\to V$.  It is finite flat of degree $N^3$ by
\eqref{wb:cusp-cover-maximal-chart-map}.

\item In the same chart the base characters are
\[
 t=z_0z_1z_2,\qquad v=w_0w_1w_2,\qquad P_N^*t=v^N.
\]
The Jacobian determinant of \eqref{wb:cusp-cover-maximal-chart-map} is
\[
 N^3(w_0w_1w_2)^{N-1}=N^3v^{N-1}.
\]
This proves \eqref{wb:cusp-cover-direct-ramification}; the local quotient
maps are biholomorphisms, so the calculation is unchanged downstairs.

\item \medskip\noindent
\emph{The relative degree-$N^2$ factor.}
The equality
\((f_0\circ\kappa_N)([\xi,v])=v^N=(\rho_N\circ f_N)([\xi,v])\)
gives the factorization
\eqref{wb:cusp-cover-relative-factorization}.  At a triple point the exact
local-ring map is
\begin{equation}
 \begin{gathered}
 A_N:=
 \frac{\mathbb C\{z_0,z_1,z_2,v\}}
      {(z_0z_1z_2-v^N)},
 \qquad
 B:=\mathbb C\{w_0,w_1,w_2\},\\
 A_N\longrightarrow B,\qquad
 z_i\longmapsto w_i^N,\qquad
 v\longmapsto w_0w_1w_2.
 \end{gathered}
 \label{wb:cusp-cover-basechange-local-ring}
\end{equation}
Let
\[
 G_N=\{(\varepsilon_0,\varepsilon_1,\varepsilon_2)\in\mu_N^3:
       \varepsilon_0\varepsilon_1\varepsilon_2=1\}
 \simeq\mu_N^2
\]
act diagonally on $B$.  A monomial
$w_0^{a_0}w_1^{a_1}w_2^{a_2}$ is $G_N$-invariant if and only if
\[
 a_0\equiv a_1\equiv a_2\pmod N.
\]
Writing their common residue as $r$, with $0\le r<N$, expresses that
monomial uniquely in the form
\[
 (w_0w_1w_2)^r
 (w_0^N)^{b_0}(w_1^N)^{b_1}(w_2^N)^{b_2}.
\]
The same coefficientwise argument applies to convergent power series.
Hence
\begin{equation}
 A_N=B^{G_N}.
 \label{wb:cusp-cover-invariant-ring}
\end{equation}
It follows that $A_N$ is normal, that the map has generic degree
$|G_N|=N^2$, and that it is a finite quotient rather than a normalization.
The corresponding face-chart calculations prove normality along the
double and smooth strata as well, so $N_0[N]$ is normal everywhere.

\item The stabilizer in $G_N$ is trivial when none or exactly one of the $w_i$
vanishes, is isomorphic to $\mu_N$ when exactly two vanish, and is all of
$G_N$ when all three vanish.  Thus $Q_N$ is etale in codimension one and
has no ramification divisor, but it is not etale above the double and
triple strata.  It is also nonflat at every such point when $N>1$.  At a
triple point the closed fibre of \eqref{wb:cusp-cover-basechange-local-ring}
has length
\begin{align}
 \dim_{\mathbb C}
 \frac{\mathbb C\{w_0,w_1,w_2\}}
      {(w_0^N,w_1^N,w_2^N,w_0w_1w_2)}
 &=N^3-(N-1)^3\\
 &=3N^2-3N+1,
 \label{wb:cusp-cover-triple-fibre-length}
\end{align}
which exceeds the generic rank $N^2$.  At a point of a double stratum,
localizing at the nonzero third coordinate and choosing its holomorphic
$N$th root reduces the ramified factor to
\[
 \frac{\mathbb C\{z_0,z_1,v\}}{(z_0z_1-v^N)}
\longrightarrow\mathbb C\{w_0,w_1\},\qquad
 z_i\longmapsto w_i^N,\quad v\longmapsto w_0w_1.
\]
Its local generic rank is $N$, whereas its closed fibre has length
\begin{equation}
 \dim_{\mathbb C}
 \frac{\mathbb C\{w_0,w_1\}}
      {(w_0^N,w_1^N,w_0w_1)}=2N-1>N.
 \label{wb:cusp-cover-double-fibre-length}
\end{equation}
The remaining transverse nonzero torus coordinate is etale locally.
Equations \eqref{wb:cusp-cover-triple-fibre-length} and
\eqref{wb:cusp-cover-double-fibre-length}, together with etaleness elsewhere,
identify the nonflat locus exactly.

\item \medskip\noindent
\emph{The regular fibres and monodromy.}
Choose a logarithm $r$ with $v=e^{2\pi ir}$.  The target period matrix at
$t=v^N$ is
\[
 Z(Nr)=NrB_0+C(v^N),
\]
whereas the source residual action is determined by
\[
 Z_N(r)=rB_0+\frac1NC(v^N)=\frac1N Z(Nr).
\]
The source fibre is therefore
\[
 \frac{\mathbb C^2}{\mathbb Z^2+Z_N(r)\mathbb Z^2},
\]
and multiplication by $N$ on $\mathbb C^2$ maps it to the target fibre.
Equivalently, in the unscaled additive coordinate it is the identity map
\begin{equation}
 \frac{\mathbb C^2}
      {N\mathbb Z^2+Z(Nr)\mathbb Z^2}
 \longrightarrow
 \frac{\mathbb C^2}
      {\mathbb Z^2+Z(Nr)\mathbb Z^2}.
 \label{wb:cusp-cover-additive-fibre-map}
\end{equation}
The imaginary part of $Z(Nr)$ is invertible on a regular fibre, so its
period lattice meets $\mathbb Z^2$ trivially in the required real sense;
the index in \eqref{wb:cusp-cover-additive-fibre-map} is $N^2$.  Its deck
group is $(\mathbb Z/N\mathbb Z)^2$.  In the ordered lattice basis
$(\widehat\gamma,\widehat u,\widehat w,\widehat\delta)$, this is exactly
the sublattice obtained by $S_N$.  Since $M_0=I+R$ with $R^2=0$ and
$R(a,b,c,d)=(0,0,b,-a)$, one has
\[
 M_0^N=I+NR,\qquad S_N^{-1}M_0^NS_N=M_0,
\]
proving \eqref{wb:cusp-cover-monodromy-conjugacy}.

\item \medskip\noindent
\emph{The central normalization and conductor.}
Multiplication by $N$ fixes every ray of $\Sigma$.  It therefore maps a
chosen source component of the toric central fibre to the corresponding
target component.  The star of that ray is the hexagonal fan, so the
restriction is the toric power map
$[N]\colon\mathrm{dP}_6\to\mathrm{dP}_6$.  On its dense two-torus it has
degree $N^2$; in each toric boundary chart it is an $N$th-power map in the
normal coordinate, hence has ramification index $N$ along all six boundary
curves.  Equation \eqref{wb:cusp-cover-equivariance-computation} says that
the source and target lattice identifications commute with this map.
Consequently opposite boundary curves, their conductor identifications,
and the two translation-orbits of triangle vertices descend compatibly.
This proves \eqref{wb:cusp-cover-normalization-diagram} and completes the
proof.
\end{enumerate}
\end{proof}


\paragraph{Exact comparison with the unmodified source action}
\label{wb:cusp-cover-naive-action-failure}
\StatementRecord{wb:cusp-cover-naive-action-failure}{proved}{Historical events: GPT-6 Astra (OpenAI), ChatGPT Pro mode, effort Pro, S6\_09\_PAPERS, 17 November 2026: targeted proof check. Also checked in this writing. Pre-continuation writing events: GPT-6 Astra (OpenAI), ChatGPT Pro mode, effort Pro, S6\_10\_PAPERS, 18 November 2026: targeted proof check. Continuation metadata events: GPT-6 Astra (OpenAI), ChatGPT Pro mode, effort Pro, S6\_10\_PAPERS continuation, 19 November 2026: used and rechecked in this run; status-vocabulary reconciliation; metadata reconciliation only; no new mathematical check. \emph{Pre-continuation recorded proof scope.} GPT-6 Astra (OpenAI, ChatGPT Pro mode), S6\_09\_PAPERS, effort Pro, 2026-11-17; also checked in this writing by GPT-6 Astra (OpenAI), ChatGPT Pro mode, effort Pro, 18 November 2026}{the hypotheses, notation and definitions stated here; \cref{wb:cusp-cover-source-factor}}

If one reuses the target action with $v$ in place of $t$ at the source,
then the two sides of equivariance have multiplicative factors
\[
 c_{\bar\lambda}(v)^N v^{NB_0\bar\lambda}
 \quad\hbox{and}\quad
 c_{\bar\lambda}(v^N)v^{NB_0\bar\lambda},
\]
respectively. Equivariance of this unmodified action is exactly the equality $c_{\bar\lambda}(v)^N=c_{\bar\lambda}(v^N)$ for every $\bar\lambda$. The displayed source factor \eqref{wb:cusp-cover-source-factor} satisfies the required equality by taking its specified holomorphic logarithm. For $N=1$ the modified and unmodified actions agree. The explicit
logarithm $2\pi iC(v^N)\bar\lambda/N$ makes this root holomorphic; no
unjustified choice of an $N$th root is involved.






\paragraph{Exact scope}
\label{wb:cusp-cover-global-scope}
\StatementRecord{wb:cusp-cover-global-scope}{proved}{Historical events: GPT-6 Astra (OpenAI), ChatGPT Pro mode, effort Pro, S6\_09\_PAPERS, 17 November 2026: targeted proof check. Also checked in this writing. Pre-continuation writing events: GPT-6 Astra (OpenAI), ChatGPT Pro mode, effort Pro, S6\_10\_PAPERS, 18 November 2026: targeted proof check. Continuation metadata events: GPT-6 Astra (OpenAI), ChatGPT Pro mode, effort Pro, S6\_10\_PAPERS continuation, 19 November 2026: used and rechecked in this run; status-vocabulary reconciliation; metadata reconciliation only; no new mathematical check. \emph{Pre-continuation recorded proof scope.} GPT-6 Astra (OpenAI, ChatGPT Pro mode), S6\_09\_PAPERS, effort Pro, 2026-11-17; also checked in this writing by GPT-6 Astra (OpenAI), ChatGPT Pro mode, effort Pro, 18 November 2026}{the hypotheses, notation and definitions stated here; the argument and qualifications immediately following}

The theorem constructs a finite morphism of completed families over the
local cusp discs.  It is a degree-$N^3$ map to the original family and a
degree-$N^2$ map only after the target is base changed.  For $N>1$, its base map $v\mapsto v^N$ differs from the identity on the disc, and the degree $N^2>1$ of the relative factor excludes a normalization map. At $N=1$ both factors are identities.

Nothing here supplies a cover over the entire base $B$.  In fact the
displayed sublattice is not preserved by the original full monodromy when
$N>1$: for example, the printed matrix $A_1$ sends
$\widehat u$ to a vector whose $\widehat w$- and
$\widehat\delta$-coefficients are $-1$ and $1$.  A further global base
cover and compatible extensions through both finite fillings would have
to be constructed separately.  The targets of the displayed Kummer morphisms are exactly $N_0$ and $N_0[N]$. The additional maps to the sphere carriers are constructed explicitly in the subsequent carrier results.





\subsection{Radial period coordinates and the compact quotient}
\label{mapsIII:carrier-radial-prerequisites}

The carrier constructions use the following real quotient of the period
family of \cref{wb:prop:period-input-audit}.  Fix the lattice basis
\[
 \Lambda=\mathbb Z\langle\widehat\gamma,\widehat u,
                         \widehat w,\widehat\delta\rangle.
\]
The cusp coordinate and its holomorphic residual functions give
\[
 \Pi=[Z\mid I_2],\qquad
 Z=s\begin{pmatrix}0&1\\-1&0\end{pmatrix}
       +\begin{pmatrix}6\mu(t)&h(t)\\b(t)-h(t)&\mu(t)\end{pmatrix},
 \qquad t=e^{2\pi i s}.
\]
Here $\mu,b,h$ are holomorphic on a disc containing zero.  On the positive
radial ray put $t=e^{-2\pi T}$ and $s=iT$.  Choose $T_0$ large enough that
this ray lies in that disc for $T\geq T_0$, and set
\begin{equation}
 \begin{aligned}
 \mu_T&=\alpha_T+im_T,& h_T&=\eta_T+iH_T,& b_T&=\xi_T+iB_T,\\
 c_T&=\xi_T-\eta_T,& L_T&=T+H_T,& q_T&=L_T-B_T,\\
 D_T&=L_Tq_T+6m_T^2.
 \end{aligned}
 \label{wb:cusp-spec-radial-notation}
\end{equation}
All these quantities are real.  Write
$m=\operatorname{Im}\mu(0)$, $H=\operatorname{Im}h(0)$ and
$B=\operatorname{Im}b(0)$.  Taylor expansion at zero yields
\begin{equation}
 m_T=m+O(e^{-2\pi T}),\qquad
 H_T=H+O(e^{-2\pi T}),\qquad
 B_T=B+O(e^{-2\pi T}),
 \label{wb:cusp-spec-holomorphic-errors}
\end{equation}
and the analogous estimates for the real parts.  In particular
$D_T=T^2+O(T)>0$ after increasing $T_0$.

Identify $\mathbb C^2$ with $\mathbb R^4$ in the order
$(x_1,y_1,x_2,y_2)$ and equip it with the Euclidean metric.  The induced
flat metric on a quotient is the metric used below.  This metric is an
explicit choice in these coordinates.

\begin{proposition}[Real periods and their dual coordinates]
\label{mapsIII:carrier-real-periods}
\StatementRecord{mapsIII:carrier-real-periods}{proved}{Historical events: GPT-6 Astra (OpenAI), ChatGPT Pro mode, effort Pro, S6\_09\_PAPERS, 17 November 2026: targeted proof check. Also checked in this writing. Pre-continuation writing events: GPT-6 Astra (OpenAI), ChatGPT Pro mode, effort Pro, S6\_10\_PAPERS, 18 November 2026: targeted proof check. Continuation metadata events: GPT-6 Astra (OpenAI), ChatGPT Pro mode, effort Pro, S6\_10\_PAPERS continuation, 19 November 2026: used and rechecked in this run; status-vocabulary reconciliation; metadata reconciliation only; no new mathematical check. \emph{Pre-continuation recorded proof scope.} GPT-6 Astra (OpenAI, ChatGPT Pro mode), effort Pro, 17 November 2026; also checked in this writing by GPT-6 Astra (OpenAI), ChatGPT Pro mode, effort Pro, 18 November 2026}{the hypotheses, notation and definitions stated here; the explicit constructions and calculations in the proof}

The radial fibre and its real period matrix are
\[
 F_T=\mathbb R^4/P_T\mathbb Z^4,\qquad
 P_T=\begin{pmatrix}
 6\alpha_T&\eta_T&1&0\\
 6m_T&L_T&0&0\\
 c_T&\alpha_T&0&1\\
 -q_T&m_T&0&0
 \end{pmatrix}.
\]
Its determinant is $-D_T$.  Its complete inverse transpose is
\begin{equation}
 P_T^{-\mathsf T}=
 \begin{pmatrix}
 0&0&1&0\\[2pt]
 \dfrac{m_T}{D_T}&\dfrac{q_T}{D_T}&
 \dfrac{-6m_T\alpha_T-q_T\eta_T}{D_T}&
 \dfrac{-m_Tc_T-q_T\alpha_T}{D_T}\\[7pt]
 0&0&0&1\\[2pt]
 -\dfrac{L_T}{D_T}&\dfrac{6m_T}{D_T}&
 \dfrac{6(L_T\alpha_T-m_T\eta_T)}{D_T}&
 \dfrac{L_Tc_T-6m_T\alpha_T}{D_T}
 \end{pmatrix}.
 \label{wb:cusp-spec-full-inverse-transpose}
\end{equation}
Thus $(P_T\mathbb Z^4)^*=P_T^{-\mathsf T}\mathbb Z^4$.
\end{proposition}



\begin{proof}
\begin{enumerate}[label=\arabic*.,leftmargin=*]
\item At $s=iT$ the complex period matrix is
\[
 Z_T=\begin{pmatrix}6\mu_T&h_T+iT\\
                          b_T-h_T-iT&\mu_T\end{pmatrix}.
\]
Taking real and imaginary parts of its columns and of the identity
columns gives the displayed $P_T$.  Expansion along its last two
columns gives $\det P_T=-L_Tq_T-6m_T^2=-D_T$.

\item To solve $P_T^{\mathsf T}x=(r,s,n_1,n_2)^{\mathsf T}$, its last two
equations first give $x_1=n_1$ and $x_3=n_2$.  Put
\[
 R=r-6\alpha_Tn_1-c_Tn_2,\qquad
 S=s-\eta_Tn_1-\alpha_Tn_2.
\]
The remaining two equations and their solution are
\[
 \begin{pmatrix}6m_T&-q_T\\L_T&m_T\end{pmatrix}
 \binom{x_2}{x_4}=\binom R S,
 \qquad
 \binom{x_2}{x_4}=
 \frac1{D_T}\begin{pmatrix}m_T&q_T\\-L_T&6m_T\end{pmatrix}
 \binom R S.
\]
Substituting $R,S$ proves every entry in
\eqref{wb:cusp-spec-full-inverse-transpose}.  Finally a real covector
$\ell$ is integral on $P_T\mathbb Z^4$ exactly when
$P_T^{\mathsf T}\ell\in\mathbb Z^4$, proving the dual-lattice formula.
\end{enumerate}
\end{proof}

Let $e_1,\ldots,e_4$ denote the coordinate unit vectors in this real
order and let $v_{\gamma,T},v_{u,T}$ be the first two columns of $P_T$.
Then
\[
 P_T\mathbb Z^4=
 \mathbb Zv_{\gamma,T}\oplus\mathbb Zv_{u,T}
                 \oplus\mathbb Ze_1\oplus\mathbb Ze_3.
\]

\begin{proposition}[Exact compact quotient of the native cusp fibre]
\label{wb:cusp-spec-native-quotient}
\StatementRecord{wb:cusp-spec-native-quotient}{proved}{Historical events: GPT-6 Astra (OpenAI), ChatGPT Pro mode, effort Pro, S6\_09\_PAPERS, 17 November 2026: targeted proof check. Also checked in this writing. Pre-continuation writing events: GPT-6 Astra (OpenAI), ChatGPT Pro mode, effort Pro, S6\_10\_PAPERS, 18 November 2026: targeted proof check. Continuation metadata events: GPT-6 Astra (OpenAI), ChatGPT Pro mode, effort Pro, S6\_10\_PAPERS continuation, 19 November 2026: used and rechecked in this run; status-vocabulary reconciliation; metadata reconciliation only; no new mathematical check. \emph{Pre-continuation recorded proof scope.} GPT-6 Astra (OpenAI, ChatGPT Pro mode), effort Pro, 17 November 2026; also checked in this writing by GPT-6 Astra (OpenAI), ChatGPT Pro mode, effort Pro, 18 November 2026}{the hypotheses, notation and definitions stated here; the explicit constructions and calculations in the proof}

Define the orthogonal projection and its projected period matrix by
\[
 \rho(x_1,x_2,x_3,x_4)=(x_2,x_4),\qquad
 Y_T=\begin{pmatrix}6m_T&L_T\\-q_T&m_T\end{pmatrix}.
\]
Then $\det Y_T=D_T>0$, so $Q_T=\mathbb R^2/Y_T\mathbb Z^2$ is a
compact real torus, and projection induces an exact sequence of compact
real Lie groups
\[
 0\longrightarrow
 \mathbb T^2_{\mathrm{re}}=
 \frac{\mathbb Re_1\oplus\mathbb Re_3}
      {\mathbb Ze_1\oplus\mathbb Ze_3}
 \longrightarrow F_T\xrightarrow{\ \pi_T\ }Q_T
 \longrightarrow0.
\]
This is a Riemannian quotient for the displayed flat metrics.  There is a
constant $C$, independent of $T\geq T_0$, such that
\[
 \|Y_T(a,b)^{\mathsf T}\|\geq(T-C)\sqrt{a^2+b^2},\qquad
 \operatorname{sys}(Q_T)\geq T-C.
\]
For every $R>0$, the radius-$R$ balls of $Q_T$ and $Q_T^2$ are exactly
Euclidean once $T-C>4R$.  Consequently their pointed limits are
$\mathbb R^2$ and $\mathbb R^4$, respectively, both locally isometrically
and smoothly, and in pointed Gromov--Hausdorff topology.
\end{proposition}



\begin{proof}
\begin{enumerate}[label=\arabic*.,leftmargin=*]
\item The determinant calculation gives
$\det Y_T=6m_T^2+L_Tq_T=D_T$.  Projection kills the last two period
columns and sends the first two to those of $Y_T$.  Therefore
$\rho(P_T\mathbb Z^4)=Y_T\mathbb Z^2$, which proves that $\pi_T$ is
well defined and surjective.

\item If $\pi_T([x])=0$, choose $a,b\in\mathbb Z$ such that
$\rho(x)=Y_T(a,b)^{\mathsf T}$.  The vector
$x-av_{\gamma,T}-bv_{u,T}$ belongs to
$\mathbb Re_1\oplus\mathbb Re_3$ and represents the same point of
$F_T$.  Conversely every point in this real plane maps to zero.
If a lattice vector lies in that plane, its first two lattice
coefficients satisfy $Y_T(a,b)^{\mathsf T}=0$.  Invertibility forces
$a=b=0$, so the intersection lattice is exactly
$\mathbb Ze_1\oplus\mathbb Ze_3$.  This proves the asserted kernel and
its injective inclusion.  Since $\rho$ is the orthogonal projection
along that plane, its Euclidean quotient metric is the displayed one.

\item Write
\[
 Y_T=T\begin{pmatrix}0&1\\-1&0\end{pmatrix}
       +\begin{pmatrix}6m_T&H_T\\B_T-H_T&m_T\end{pmatrix}.
\]
The second matrix has uniformly bounded operator norm by holomorphy.
Choose $C$ above this bound.  The first matrix is $T$ times an orthogonal
matrix, and the reverse triangle inequality gives the stated estimate.
A nonzero integer vector has norm at least one, which proves the
systolic bound; it also proves the same bound for the product lattice.

\item For $x,y$ in the Euclidean radius-$R$ ball and a nonzero lattice vector
$\lambda$, the estimate
\[
 \|x-y+\lambda\|\geq\|\lambda\|-\|x-y\|
 \geq(T-C)-2R>2R\geq\|x-y\|
\]
holds when $T-C>4R$.  Thus no lattice translation shortens a distance
between points of this ball.  Every point of the quotient ball has a
lift in the Euclidean ball, by choosing a shortest representative of
its lattice coset.  The covering map therefore gives an isometry onto
the entire quotient ball.  These are smooth local isometries, proving
all the stated pointed convergences.
\end{enumerate}
\end{proof}

\begin{theorem}[The maximal regular-locus
\texorpdfstring{\(\mathbb R^4\)}{R4} quotient and its compactification
obstruction]
\label{wb:cusp-spec-globalization-obstruction}
\StatementRecord{wb:cusp-spec-globalization-obstruction}{proved}{Historical events: GPT-6 Astra (OpenAI), ChatGPT Pro mode, effort Pro, S6\_09\_PAPERS, 17 November 2026: targeted proof check. Also checked in this writing. Pre-continuation writing events: GPT-6 Astra (OpenAI), ChatGPT Pro mode, effort Pro, S6\_10\_PAPERS, 18 November 2026: targeted proof check. Continuation metadata events: GPT-6 Astra (OpenAI), ChatGPT Pro mode, effort Pro, S6\_10\_PAPERS continuation, 19 November 2026: used and rechecked in this run; status-vocabulary reconciliation; metadata reconciliation only; no new mathematical check. \emph{Pre-continuation recorded proof scope.} GPT-6 Astra (OpenAI, ChatGPT Pro mode), S6\_09\_PAPERS, effort Pro, 2026-11-17; also checked in this writing by GPT-6 Astra (OpenAI), ChatGPT Pro mode, effort Pro, 18 November 2026}{the hypotheses, notation and definitions stated here; \cref{wb:cusp-spec-native-quotient}; \cref{wb:cusp-cone-S12-route}; \cref{wb:cusp-cone-S24-local-route}; \cref{wb:cusp-cone-S24-route}; \cite[\S15, Proposition~15.1 and the following remark,
p.~434]{BorelSerre1953}; \cref{wb:cusp-spec-invariant-character-lattice}; \cref{wb:cusp-spec-full-inverse-transpose}; \cref{wb:cusp-spec-radial-notation}; \cref{wb:cusp-spec-holomorphic-errors}; \cref{wb:cusp-spec-invariant-flag-calculation}; \cref{wb:cusp-spec-finite-overlap-identity}; \cref{wb:cusp-spec-cusp-nonextension-chart}}

\mbox{}\ 
Let
\[
 \Lambda=\mathbb Z\langle\widehat\gamma,\widehat u,
                   \widehat w,\widehat\delta\rangle,\qquad
 K_\Lambda=\mathbb Z\langle\widehat w,\widehat\delta\rangle,
\]
and define
\[
 \varepsilon_\gamma
 (a\widehat\gamma+b\widehat u+c\widehat w+d\widehat\delta)=a,\qquad
 L_\Lambda=\ker\varepsilon_\gamma
 =\mathbb Z\langle\widehat u,\widehat w,\widehat\delta\rangle.
\tag{G.1}
\]
Then the following statements hold.

\smallskip
\noindent
\emph{(i) Regular group quotient.}
Over \(B^\circ\), the period family \(J=X^\circ=f^{-1}(B^\circ)\) is a
holomorphic family of complex Lie groups, not merely an affine torus
bundle.  The formula
\begin{equation}
 \chi_\gamma\bigl([z,\Pi(z)x]\bigr)
 =\bigl([z],\varepsilon_\gamma(x)\bmod\mathbb Z\bigr)
 \label{wb:cusp-spec-global-circle-map}
\end{equation}
defines a global smooth real-analytic fibrewise group epimorphism
\begin{equation}
 \chi_\gamma:J\longrightarrow B^\circ\times(\mathbb R/\mathbb Z).
 \label{wb:cusp-spec-global-circle-quotient}
\end{equation}
Write
\begin{equation}
 \theta_\gamma:=\operatorname{pr}_{\mathbb R/\mathbb Z}\circ\chi_\gamma:
 J\longrightarrow\mathbb R/\mathbb Z
 \label{wb:cusp-spec-global-circle-coordinate}
\end{equation}
for its circle coordinate.  Thus
\(\chi_\gamma=(f,\theta_\gamma)\); the two symbols have different
codomains and will not be interchanged below.
Its kernel is the connected real three-torus family
\begin{equation}
 \mathcal H:=
 \frac{\mathfrak h_z^\circ\times
       ((L_\Lambda\otimes\mathbb R)/L_\Lambda)}{\Delta},
 \label{wb:cusp-spec-global-kernel-family}
\end{equation}
% checked carrier correction: deck-group
where \(\Delta\) is the actual triangle deck group of
\(\mathfrak h_z^\circ\to B^\circ\).  It acts diagonally by its deck action
and by the restricted integral monodromy on \(L_\Lambda\).
The deck representation maps \(\pi_1(B^\circ)\) onto \(\Delta\);
these two groups are not identified.  Equivalently, one can replace
\(\mathfrak h_z^\circ\) by the universal cover of \(B^\circ\) and
take the quotient by \(\pi_1(B^\circ)\).
The invariant-lattice calculation uses the image of the monodromy
representation, while finite base covers correspond to finite-index
subgroups of \(\pi_1(B^\circ)\).  Thus every fibre of \(\mathcal H\) is
\((L_\Lambda\otimes\mathbb R)/L_\Lambda\); the family need not be a
product.
The quotient is real and is not a holomorphic quotient of complex tori.

\smallskip
\noindent
\emph{(ii) Exact radial quotient metric.}
On the positive radial cusp, the restriction of
\cref{wb:cusp-spec-global-circle-map} to
\(F_T=\mathbb R^4/P_T\mathbb Z^4\) is
\begin{equation}
 \chi_{\gamma,T}([y])
 =\left\langle P_T^{-\mathsf T}e_\gamma,y\right\rangle
   \bmod\mathbb Z,\qquad e_\gamma=(1,0,0,0)^{\mathsf T}.
 \label{wb:cusp-spec-radial-circle-character}
\end{equation}
For the Riemannian quotient of the Euclidean metric, the target circle has
circumference
\begin{equation}
 \ell_T
 =\frac1{\lVert P_T^{-\mathsf T}e_\gamma\rVert}
 =\frac{D_T}{\sqrt{L_T^2+m_T^2}}
 =T+H-B+O(T^{-1}).
 \label{wb:cusp-spec-global-circle-circumference}
\end{equation}
In particular, \(\ell_T\to\infty\).

\smallskip
\noindent
\emph{(iii) Four copies and minimality.}
The source
\begin{equation}
 J^{[4]}=J\times_{B^\circ}J\times_{B^\circ}J\times_{B^\circ}J
 \label{wb:cusp-spec-fourfold-fibre-product}
\end{equation}
is the actual fourfold fibre product over all of \(B^\circ\).  The product
map
\begin{equation}
 \chi_\gamma^{[4]}=(\chi_\gamma,\chi_\gamma,\chi_\gamma,\chi_\gamma):
 J^{[4]}\longrightarrow B^\circ\times(\mathbb R/\mathbb Z)^4
 \label{wb:cusp-spec-fourfold-circle-map}
\end{equation}
has kernel
\(\mathcal H^{[4]}=\mathcal H\times_{B^\circ}\cdots
\times_{B^\circ}\mathcal H\), whose fibres are
\(\mathbb T^{12}\).  Along the synchronized positive
radial cusp, its quotient fibre, with product quotient metric, is
\begin{equation}
 G_T=(\mathbb R/\ell_T\mathbb Z)^4,\qquad
 \operatorname{sys}(G_T)=\ell_T.
 \label{wb:cusp-spec-four-circle-fibre}
\end{equation}
Consequently
\begin{equation}
 (G_T,0)\longrightarrow(\mathbb R^4,0)
 \label{wb:cusp-spec-four-circle-limit}
\end{equation}
in pointed local-isometric, pointed smooth, and pointed
Gromov--Hausdorff topology.  If \(\ell_T>4R\), the radius-\(R\) ball
about zero is exactly Euclidean.

Equivalently, the external Cartesian product gives a quotient
\[
 J^4\longrightarrow(B^\circ)^4\times(\mathbb R/\mathbb Z)^4,
\]
and \(J^{[4]}\) is its pullback along the diagonal
\(B^\circ\to(B^\circ)^4\).  At four independent cusp parameters its target
lattice is
\(\operatorname{diag}(\ell_{T_1},\ldots,\ell_{T_4})\mathbb Z^4\), so the
same pointed limit holds whenever
\(\min_i\ell_{T_i}\to\infty\).

Among connected torus quotients obtained from \(r\) copies by a primitive
full-monodromy-invariant kernel and whose quotient systole tends to infinity
on the radial cusp, the quotient rank is at most \(r\).  Thus four copies
are minimal for a four-dimensional all-expanding quotient, and
\cref{wb:cusp-spec-fourfold-circle-map} attains the bound.

The cusp-local two-copy map
\begin{equation}
 \pi_T\times\pi_T:F_T\times F_T\longrightarrow Q_T\times Q_T
 \label{wb:cusp-spec-doubled-local-quotient}
\end{equation}
from \cref{wb:cusp-spec-native-quotient} also has pointed limit
\(\mathbb R^4\).  It is not a quotient over the full \(B^\circ\), however,
and no finite base cover makes it one.

\smallskip
\noindent
\emph{(iv) Finite fillings.}
At \(p_j\), use the notation of
\cref{wb:prop:candidate-finite-fillings} and put
\[
 (m_1,v_1,\kappa_1)
 =\bigl(3,\widehat\gamma+2\widehat u-4\widehat w,1\bigr),
\]
\[
 (m_2,v_2,\kappa_2)
 =\bigl(4,-\widehat\gamma-3\widehat u+3\widehat w,-1\bigr),\qquad
 \kappa_j=\varepsilon_\gamma(v_j).
\tag{G.2}
\]
In the un-conjugated flat coordinates, the finite filling generator is
\begin{equation}
 \Phi_j(s,x)=
 \left(e^{-2\pi i/m_j}s,A_jx+\frac{v_j}{m_j}\right).
 \label{wb:cusp-spec-finite-affine-action}
\end{equation}
The raw character is not invariant:
\begin{equation}
 \varepsilon_\gamma\left(A_jx+\frac{v_j}{m_j}\right)
 =\varepsilon_\gamma(x)+\frac{\kappa_j}{m_j}\pmod{\mathbb Z}.
 \label{wb:cusp-spec-finite-character-shift}
\end{equation}
It instead descends equivariantly to the Seifert circle family
\begin{equation}
 C_j=
 \frac{\Delta_j\times(\mathbb R/\mathbb Z)}
 {(s,\theta)\sim
  (e^{-2\pi i/m_j}s,\theta+\kappa_j/m_j)}.
 \label{wb:cusp-spec-finite-Seifert-target}
\end{equation}
The total space \(C_j\) is a smooth solid torus, the map
\(C_j\to D_j\) has a multiple central circle of order \(m_j\), and the
induced map \(N_j\to C_j\) is a smooth proper \(\mathbb T^3\)-bundle.
The target translation in \cref{wb:cusp-spec-finite-Seifert-target} is
exactly compatible with the logarithmic gluing map \(h_j\).

For four synchronized copies, the smooth simultaneous-root extension is
\begin{equation}
 \widetilde N_j^{[4]}
 =\frac{(J'_j)^{\times_{\Delta_j}4}}
 {\langle(\widetilde g_j^{\log})^{\times4}\rangle},
 \label{wb:cusp-spec-finite-fourfold-normalization}
\end{equation}
with target
\begin{equation}
 C_j^{[4]}=
 \frac{\Delta_j\times(\mathbb R/\mathbb Z)^4}
 {(s,\theta_1,\ldots,\theta_4)\sim
  (e^{-2\pi i/m_j}s,
   \theta_1+\kappa_j/m_j,\ldots,\theta_4+\kappa_j/m_j)}.
 \label{wb:cusp-spec-finite-fourfold-target}
\end{equation}
The map \(\widetilde N_j^{[4]}\to C_j^{[4]}\) is a smooth proper
\(\mathbb T^{12}\)-bundle.  There is a canonical finite map
\begin{equation}
 \nu_j:\widetilde N_j^{[4]}
 \longrightarrow N_j\times_{D_j}N_j\times_{D_j}N_j\times_{D_j}N_j
 \label{wb:cusp-spec-finite-normalization-map}
\end{equation}
which is an isomorphism over \(D_j^*\), is the normalization of the
reduced ordinary filled fibre product, and has degree \(m_j^3\) on the
reduced central fibre.  These two degrees are \(27\) and \(64\).  Gluing
these two local normalizations to \(J^{[4]}\) by
\((h_j)^{\times4}\), and gluing their targets to the product
\(B^\circ\times(\mathbb R/\mathbb Z)^4\) by
\((h_j^{S^1})^{\times4}\), gives the corresponding normalization and
Seifert target over \(B\setminus\{p_0\}\).

\smallskip
\noindent
\emph{(v) Cusp nonextension.}
Neither the literal map \(\chi_\gamma\) nor
\(\chi_\gamma^{[4]}\) extends continuously through the toric cusp fibre
\(W\).  Thus the four-copy quotient is global over \(B^\circ\), and its
finite-filling normalization is global over \(B\setminus\{p_0\}\), but
the displayed map \(\chi_\gamma^{[4]}\) has no continuous extension to
the compact fourfold fibre product or to \(X^4\) with the stated
regular-fibre restriction, and it does not factor through
\(X^{\wedge4}\cong S^{24}\).  The cohomology class of
\(\chi_\gamma\) does extend over
\(N_0\); this weaker homotopy statement does not extend the displayed
quotient map or its metric.

\smallskip
\noindent
\emph{(vi) Products, smashes, and sphere constructions.}
The homotopy-sphere recognition in \cref{maps:src:B.5.2}, together with $\Theta_6=0$ \cite[Theorem~5.1 and table, p.~512]{manKM63}, gives a based diffeomorphism $h:X\to S^6$. Its Cartesian powers give
\begin{equation}
 X^2\cong_{\mathrm{diff}}S^6\times S^6\not\simeq S^{12},\qquad
 H_6(X^2;\mathbb Z)\cong\mathbb Z^2,
 \label{wb:cusp-spec-product-not-sphere}
\end{equation}
and
\begin{equation}
 \begin{aligned}
 X^4&\cong_{\mathrm{diff}}(S^6)^4\not\simeq S^{24},\\
 H_6(X^4;\mathbb Z)&\cong\mathbb Z^4,&
 H_{12}(X^4;\mathbb Z)&\cong\mathbb Z^6,&
 H_{18}(X^4;\mathbb Z)&\cong\mathbb Z^4.
 \end{aligned}
 \label{wb:cusp-spec-four-product-not-sphere}
\end{equation}
The synchronized regular fourfold \(J^{[4]}\) has real dimension \(18\),
not \(24\).  Choosing a base point gives based homeomorphisms
\begin{equation}
X\wedge X\cong S^{12},\qquad X^{\wedge4}\cong S^{24},
\label{wb:cusp-spec-smash-spheres}
\end{equation}
but these smash spheres do not carry either \emph{literal fixed-target}
quotient above.  The later cone constructions
\cref{wb:cusp-cone-S12-route} and \cref{wb:cusp-cone-S24-local-route} and \cref{wb:cusp-cone-S24-route}
change the codomain by retaining the radial parameter and collapsing the
obstructing boundary cross-section; they therefore do not contradict this
nonextension theorem.  Moreover,
neither \(S^{12}\) nor \(S^{24}\), with any differentiable structure,
admits an almost-complex structure
\cite[\S15, Proposition~15.1 and the following remark,
p.~434]{BorelSerre1953}; hence neither can be promoted to a compact
complex manifold.  This is an obstruction to a proposed complex
\(S^{12}\)- or \(S^{24}\)-globalization, not to the real quotient or its
pointed metric limit.  Precisely,
\(\pi_T\times\pi_T\) is a map on \(F_T^2\), and it fails to factor
through \(F_T\wedge F_T\); \(\chi_{\gamma,T}^{[4]}\) is a map on
\(F_T^4\), and it fails to factor through \(F_T^{\wedge4}\).
Neither map is a map on \(X^2\) or \(X^4\), and no extension with the
stated fibre restriction exists.

There is one exact but different sphere construction.  If
\(C_T=\mathbb R/\ell_T\mathbb Z\) is based at zero, postcomposition with
the target fat-wedge collapse gives
\begin{equation}
 F_T^4\xrightarrow{\chi_{\gamma,T}^{[4]}}C_T^4
 \longrightarrow C_T^{\wedge4}\cong S^4.
 \label{wb:cusp-spec-target-smash-map}
\end{equation}
This map does factor through \(F_T^{\wedge4}\).  With the quotient length
metric and base point
\([\ell_T/2,\ell_T/2,\ell_T/2,\ell_T/2]\), its pointed limit is
\(\mathbb R^4\), with every radius-\(R\) ball exactly Euclidean once
\(\ell_T>4R\).  The metric is singular at the collapsed fat wedge, and
the map still does not extend to \(X^4\) or \(X^{\wedge4}\).
\end{theorem}



\begin{proof}
\begin{enumerate}[label=\arabic*.,leftmargin=*]
\item \medskip\noindent\emph{Invariant lattices and the regular group law.}\par
The matrices in \cref{wb:prop:monodromy-certificate} give
\[
 (M_0-I)(a,b,c,d)^{\mathsf T}=(0,0,b,-a)^{\mathsf T},
\]
so
\begin{equation}
 K_\Lambda=\ker(M_0-I)=\operatorname{im}(M_0-I).
 \label{wb:cusp-spec-cusp-kernel}
\end{equation}
They also give
\begin{equation}
 \begin{aligned}
 A_1\widehat u&=-\widehat w+\widehat\delta,&
 A_1\widehat w&=\widehat u-\widehat w,&
 A_1\widehat\delta&=\widehat\delta,\\
 A_2\widehat u&=\widehat w,&
 A_2\widehat w&=-\widehat u+\widehat\delta,&
 A_2\widehat\delta&=\widehat\delta.
 \end{aligned}
 \label{wb:cusp-spec-invariant-flag-calculation}
\end{equation}
Thus \(L_\Lambda\) is invariant.  It is the smallest invariant lattice
containing \(K_\Lambda\), because
\(\widehat u=A_1\widehat w+\widehat w\).  Both lattices are primitive.

\item For a row covector \((a,b,c,d)\), invariance under \(A_1\) gives
\(b=2c\) and \(d=3c\), whereas invariance under \(A_2\) gives
\(b=c\) and \(d=2c\).  Hence \(b=c=d=0\), and
\begin{equation}
 \operatorname{Hom}(\Lambda,\mathbb Z)^\Delta
 =\mathbb Z\varepsilon_\gamma.
 \label{wb:cusp-spec-invariant-character-lattice}
\end{equation}

\item The family over the cover is
\((\mathfrak h_z^\circ\times\mathbb C^2)/\Lambda\), and the action used
to form \(J\) is
\[
 [z,\zeta]\longmapsto[gz,R_g(z)\zeta].
\]
It is complex linear in every fibre and fixes the zero section.
Fibrewise addition, inversion, and zero are therefore equivariant and
descend to \(J\).  In the flat real-analytic marking
\[
 \Xi(z,x+\Lambda)=[z,\Pi(z)x],
\]
equation (2.3d) makes the action
\[
 (z,x+\Lambda)\longmapsto(gz,A_gx+\Lambda).
\]
Equation \cref{wb:cusp-spec-invariant-character-lattice} proves that
\cref{wb:cusp-spec-global-circle-map} is well defined.  Primitivity of
\(\varepsilon_\gamma\) proves surjectivity and connectedness of the
kernel.  The induced action on \(\Lambda/L_\Lambda\simeq\mathbb Z\) is
trivial, so the target is the product circle bundle.  Since the kernel has
odd real dimension, this is not a holomorphic quotient.

\item \medskip\noindent\emph{The quotient metric and the four-copy limit.}\par
The first column of
\cref{wb:cusp-spec-full-inverse-transpose} is
\begin{equation}
 P_T^{-\mathsf T}e_\gamma
 =\left(0,\frac{m_T}{D_T},0,-\frac{L_T}{D_T}\right)^{\mathsf T}.
 \label{wb:cusp-spec-circle-dual-vector}
\end{equation}
The quotient metric on \(\mathbb R/\mathbb Z\) is therefore
\(\ell_T^2d\theta^2\).  Equivalently,
\[
 P_T(L_\Lambda\otimes\mathbb R)
 =\operatorname{span}_{\mathbb R}\{e_1,e_3,v_{u,T}\}
\]
has unit normal
\[
 \nu_T=\frac{(0,m_T,0,-L_T)}{\sqrt{L_T^2+m_T^2}},
\]
and
\[
 \operatorname{dist}
 \bigl(v_{\gamma,T},P_T(L_\Lambda\otimes\mathbb R)\bigr)
 =\left|\langle v_{\gamma,T},\nu_T\rangle\right|
 =\frac{L_Tq_T+6m_T^2}{\sqrt{L_T^2+m_T^2}}
 =\ell_T.
\]
The expansions in
\cref{wb:cusp-spec-radial-notation} and \cref{wb:cusp-spec-holomorphic-errors} give
\[
 D_T=T^2+(2H-B)T+O(1),\qquad
 \sqrt{L_T^2+m_T^2}=T+H+O(T^{-1}),
\]
which proves \cref{wb:cusp-spec-global-circle-circumference}.

\item After rescaling the circle coordinate by \(\ell_T\), the product target is
\(\mathbb R^4/\ell_T\mathbb Z^4\).  If
\(x,y\in B_R(0)\) and \(0\ne\lambda\in\ell_T\mathbb Z^4\), then
\[
 \lVert x-y+\lambda\rVert
 \ge\lVert\lambda\rVert-\lVert x-y\rVert
 \ge\ell_T-2R>2R\ge\lVert x-y\rVert
\]
when \(\ell_T>4R\).  Hence the covering projection is an exact isometry
on \(B_R(0)\), proving all three convergences in
\cref{wb:cusp-spec-four-circle-limit}.  The quotient map is global over
\(B^\circ\); only the metric limit is evaluated along the radial cusp.
The separately chosen Euclidean fibre metric is not claimed to be a
canonical global metric away from that cusp.

\item \medskip\noindent\emph{Minimality and the finite-base-cover obstruction.}\par
Let \(N\subset\Lambda^{\oplus r}\) be a primitive invariant kernel whose
connected quotient has radial systole tending to infinity.  Every element
of \(K_\Lambda^{\oplus r}\) has uniformly bounded physical length, since
\[
 P_T\widehat w=e_1,\qquad P_T\widehat\delta=e_3.
\]
If such an element survived in the quotient lattice, its quotient norm
would be no greater than its source norm and would give a nonzero uniformly
bounded period.  Therefore \(K_\Lambda^{\oplus r}\subset N\).
Invariance and
\cref{wb:cusp-spec-invariant-flag-calculation} then force
\(L_\Lambda^{\oplus r}\subset N\), and
\[
 \operatorname{rank}(\Lambda^{\oplus r}/N)\le4r-3r=r.
\]
This proves the lower bound \(r\ge4\), while
\(N=L_\Lambda^{\oplus4}\) proves sharpness.

\item For the finite-cover assertion, put
\begin{equation}
 H_*=A_2A_1
 =\begin{pmatrix}
 1&0&0&0\\
 6&1&1&0\\
 0&0&1&0\\
 -5&0&-1&1
 \end{pmatrix}=I+N_*.
 \label{wb:cusp-spec-no-finite-cover-matrix}
\end{equation}
Direct multiplication gives \(N_*^2=0\), and hence
\begin{equation}
 H_*^n\widehat w
 =\widehat w+n(\widehat u-\widehat\delta),\qquad
 H_*^n\widehat\delta=\widehat\delta.
 \label{wb:cusp-spec-no-finite-cover-orbit}
\end{equation}
Thus \(H_*^nK_\Lambda\ne K_\Lambda\) for every \(n\ne0\).
Every finite-index subgroup of the base fundamental group contains a
positive power of the loop with monodromy \(H_*\).  Such a subgroup cannot
preserve \(K_\Lambda\), proving that no finite base cover globalizes the
native two-torus quotient.

\item \medskip\noindent\emph{The affine finite fillings and the actual overlap
maps.}\par
Equation \cref{wb:cusp-spec-invariant-character-lattice} gives
\(\varepsilon_\gamma A_j=\varepsilon_\gamma\), so
\cref{wb:cusp-spec-finite-character-shift} follows.  Define
\[
 \overline\Phi_j(s,\theta)
 =\left(e^{-2\pi i/m_j}s,\theta+\frac{\kappa_j}{m_j}\right).
\]
Then \((s,x)\mapsto(s,\varepsilon_\gamma(x))\) is
\((\Phi_j,\overline\Phi_j)\)-equivariant; denote the descended map by
\(\overline\chi_j:N_j\to C_j\).  A nonidentity power which fixes
the base must lie over \(s=0\), where it translates the circle by
\(k\kappa_j/m_j\).  Since \(\gcd(\kappa_j,m_j)=1\), this is nonintegral
for \(0<k<m_j\).  The target action is free; the source action is free by
\cref{wb:prop:candidate-finite-fillings}.

\item For an explicit solid-torus chart, put
\(\lambda_j=e^{2\pi i/m_j}\), \(u=e^{2\pi i\theta}\), and choose
\(r_j\in\mathbb Z\) with \(r_j\kappa_j\equiv1\pmod {m_j}\).  Then
\[
 w=su^{r_j},\qquad U=u^{m_j}
\]
are invariant and identify \(C_j\) real-analytically with a solid torus.
In these coordinates
\[
 t_j=s^{m_j}=w^{m_j}U^{-r_j},
\]
which proves the central multiplicity.  Equivariance and the primitive
kernel give the proper \(\mathbb T^3\)-bundle \(N_j\to C_j\).

\item It remains to check the actual logarithmic gluing.  Let
\(\Pi_{\mathbb R}(z):\Lambda\otimes\mathbb R\to\mathbb C^2\) be the
underlying real isomorphism.  On a logarithm branch \(\ell\), define
\(b_j^\ell(z)\) uniquely by
\begin{equation}
 \Pi_{\mathbb R}(z)b_j^\ell(z)
 =\frac{\ell(z)}{2\pi i}\Pi(z)v_j,\qquad
 \beta_j(z)=\varepsilon_\gamma(b_j^\ell(z))\bmod\mathbb Z.
 \label{wb:cusp-spec-finite-target-conjugacy}
\end{equation}
Changing \(\ell\) by \(2\pi in\) adds \(nv_j\) to \(b_j^\ell\), hence
adds \(n\kappa_j\in\mathbb Z\) to its
\(\varepsilon_\gamma\)-value.  Thus \(\beta_j\) is branch independent.
In flat coordinates the source conjugacy (2.9e) is translation by
\(b_j^\ell\).  Applying \(\varepsilon_\gamma\) to that conjugacy gives
\begin{equation}
 \beta_j(g_jz)=\beta_j(z)-\frac{\kappa_j}{m_j}\pmod{\mathbb Z}.
 \label{wb:cusp-spec-finite-target-cocycle}
\end{equation}
Therefore
\[
 h_j^{S^1}(z,\theta)=(z,\theta+\beta_j(z))
\]
conjugates \(\overline\Phi_j\) to the linear target action
\((z,\theta)\mapsto(g_jz,\theta)\), and the overlap identity is
\begin{equation}
 \chi_\gamma\circ h_j=h_j^{S^1}\circ\overline\chi_j.
 \label{wb:cusp-spec-finite-overlap-identity}
\end{equation}

\item The diagonal action on
\((J'_j)^{\times_{\Delta_j}4}\) is free, and the same coprimality argument
proves that \(C_j^{[4]}\) is smooth.  Every fibre of the resulting map is
a \(\mathbb T^{12}\)-torsor; over \(D_j^*\) its structure torus is the
fourfold fibre product of \(\mathcal H\).  No group zero is asserted at
the affine filling.
In local etale coordinates along the central fibre, the ordinary filled
fibre product has equations
\begin{equation}
 s_i^{m_j}=s_1^{m_j},\qquad i=2,3,4.
 \label{wb:cusp-spec-finite-fibre-product-branches}
\end{equation}
It has \(m_j^3\) smooth branches along the reduced central fibre.  The
simultaneous-root model separates them, is smooth and hence normal, and is
an isomorphism away from the central fibre.  The map \(\nu_j\) is proper:
both spaces are proper over \(D_j\), and the inverse image of a compact
subset of the target is closed in the inverse image of its compact image
in \(D_j\).  It is quasi-finite, since its fibres are singletons away from
the central fibre and have at most \(m_j^3\) points over the central
fibre.  Hence \(\nu_j\) is finite.  The universal property of
normalization now proves that it is the normalization of the reduced
ordinary fibre product.  On the central fibre the relative root choices form
\((\mathbb Z/m_j)^4/\operatorname{diag}(\mathbb Z/m_j)\), of order
\(m_j^3\), which proves the asserted degrees.

\item On \(D_j^*\), the source identification is \((h_j)^{\times4}\), and
\cref{wb:cusp-spec-finite-overlap-identity} identifies the target overlap
with \((h_j^{S^1})^{\times4}\).  These are precisely the fourfolds of
the one-copy overlap maps and therefore satisfy the cocycle identities.
They glue the two local normalizations and Seifert targets to
\(J^{[4]}\) and to
\(B^\circ\times(\mathbb R/\mathbb Z)^4\), respectively.  Equivalently,
the source gluing follows from uniqueness of normalization.

\item The fixed product-target trivialization does not extend over \(X\), even if
base winding is allowed.  Indeed, in the presentation of
\cref{mapsI:top:prop:geometric-pi1}, let a homomorphism to \(\mathbb Z\) restrict to
\(n\varepsilon_\gamma\) on \(\Lambda\) and take values \(a,b\) on
\(\widehat\rho_1,\widehat\rho_2\).  The three attachment relations give
\[
 3a=n,\qquad4b=-n,\qquad a+b=0.
\]
The first two give \(a=n/3\), \(b=-n/4\), while the third gives
\(n/12=0\).  Hence \(n=0\).  This is the exact global obstruction behind
the fact that the two Seifert targets do not combine into a fixed-circle
map on \(X\).

\item \medskip\noindent\emph{The explicit toric cusp obstruction.}\par
On the positive radial ray, write
\[
 Z_T=iTB_0+C(t_c),\qquad
 B_0=\begin{pmatrix}0&1\\-1&0\end{pmatrix},\qquad
 t_c=e^{-2\pi T}.
\]
For \(a=(a_1,a_2)^{\mathsf T}\in\mathbb R^2\), take the point of the
regular fibre represented by \(Z_Ta\), with zero short-period coordinates,
and exponentiate it to the dense torus of the cusp model:
\[
 x_T(a)=\exp(2\pi iZ_Ta)\in(\mathbb C^*)^2.
\]
Its logarithmic position is
\begin{equation}
 y_T(a)=
 \frac{(\log|x_{T,1}(a)|,\log|x_{T,2}(a)|)}{\log|t_c|}
 =B_0a+\frac{\operatorname{Im}C(t_c)a}{T}.
 \label{wb:cusp-spec-cusp-nonextension-position}
\end{equation}
The circle character has the constant value
\begin{equation}
 \chi_{\gamma,T}(x_T(a))=a_1\bmod\mathbb Z.
 \label{wb:cusp-spec-cusp-nonextension-character}
\end{equation}

\item Use the affine toric chart of
\(\sigma=\operatorname{conv}(0,e_1,e_2)\).  Its coordinates are
\begin{equation}
 z_0=\frac{t_c}{x_1x_2},\qquad z_1=x_1,\qquad z_2=x_2.
 \label{wb:cusp-spec-cusp-nonextension-chart}
\end{equation}
If \(B_0a=(y_1,y_2)\) is in the interior of \(\sigma\), then
\(y_1>0\), \(y_2>0\), and \(1-y_1-y_2>0\).  For all sufficiently large
\(T\), the three absolute values in
\cref{wb:cusp-spec-cusp-nonextension-chart} are
\[
 |t_c|^{1-y_{T,1}-y_{T,2}},\qquad
 |t_c|^{y_{T,1}},\qquad |t_c|^{y_{T,2}},
\]
and all tend to zero.  Every such path converges to the same toric fixed
point and hence to the same triple point of \(W\).

\item Take
\[
 a=\left(-\frac14,\frac14\right),\qquad
 a'=\left(-\frac12,\frac14\right).
\]
Then
\[
 B_0a=\left(\frac14,\frac14\right),\qquad
 B_0a'=\left(\frac14,\frac12\right)
\]
are both in the interior of \(\sigma\).  The paths converge to the same
point of \(W\), whereas
\[
 \chi_{\gamma,T}(x_T(a))=\frac34,\qquad
 \chi_{\gamma,T}(x_T(a'))=\frac12.
\]
No continuous extension of the literal map \(\chi_\gamma\) exists.
Choosing identical marked radial paths in the other three factors and
using these two paths in one factor proves the same assertion for
\(\chi_\gamma^{[4]}\).

\item The radial family of translations by \(\widehat u/2\) in \(\mathcal H\)
fails to extend as well.  Take two paths
whose limiting logarithmic positions are
\[
 y=\left(\frac18,\frac18\right),\qquad
 y'=\left(\frac38,\frac38\right).
\]
Both approach the fixed point of
\(\operatorname{conv}(0,e_1,e_2)\).  Translation by
\(\widehat u/2\in\mathcal H_T\) adds the exact displacement
\[
 \frac12B_0e_2+
 \frac1{2T}\operatorname{Im}C(t_c)e_2
 \longrightarrow\left(\frac12,0\right)
\]
to logarithmic position.  The translated limits are
\[
 \left(\frac58,\frac18\right)
 \quad\hbox{and}\quad
 \left(\frac78,\frac38\right).
\]
The first lies in the interior of
\(\operatorname{conv}(0,e_1,e_2)\), while the second lies in the interior
of \(\operatorname{conv}(e_1,e_2,e_1+e_2)\).  These two triangle types are
the two translation orbits in the \(A_2\) triangulation and give the two
distinct triple points of \(W\), as proved in
\cref{wb:thm:candidate-cusp-geometry}.  Thus one regular-fibre translation
sends two paths with the same limit to paths with different limits; it has
no continuous extension to \(N_0\).

\item This pointwise obstruction must not be confused with the cohomology class.
By \cref{mapsI:top:prop:geometric-pi1},
\[
 \pi_1(N_0)=\Lambda/K_\Lambda,
\]
and \(\varepsilon_\gamma\) kills \(K_\Lambda\).  It therefore defines an
integral \(H^1\)-class on \(N_0\), hence a map to the Eilenberg--Mac Lane
space \(S^1\) whose restriction is homotopic to \(\chi_\gamma\).  The
explicit paths prove that no such representative agrees with
\(\chi_\gamma\) on the whole punctured cusp.  The phase torus
\((K_\Lambda\otimes\mathbb R)/K_\Lambda\) extends as the compact real torus
action of the toric model, generally with boundary stabilizers; the full
radial \(\widehat u/2\)-translation and the full quotient do not.

\item \medskip\noindent\emph{Products and smash quotients.}\par
The integral Kunneth theorem gives
\[
 H_k(X^n;\mathbb Z)\cong
 \begin{cases}
  \mathbb Z^{\binom nj},&k=6j,\quad0\le j\le n,\\
  0,&\text{otherwise},
 \end{cases}
 \qquad \chi(X^n)=2^n,
\]
which proves
\cref{wb:cusp-spec-product-not-sphere} and \cref{wb:cusp-spec-four-product-not-sphere}.
The synchronized regular \(n\)-fold fibre product has dimension
\[
 6n-2(n-1)=4n+2,
\]
which gives \(18\) for four copies.

\item A product map factors through a smash quotient if and only if it is
constant on the corresponding fat wedge.  On
\(F_T\times\{0\}\), the local two-copy quotient is
\[
 x\longmapsto(\pi_T(x),0),
\]
which is nonconstant.  Thus the literal conclusion is that
\(\pi_T\times\pi_T\) does not factor through \(F_T\wedge F_T\); it would
be a type error to call it a map on \(X\wedge X\).  The restriction of
\(\chi_{\gamma,T}^{[4]}\) to \(F_T\times\{0\}^3\) is similarly
\[
 x\longmapsto(\chi_{\gamma,T}(x),0,0,0),
\]
so it does not factor through \(F_T^{\wedge4}\).

\item There is no global extension which repairs the domain mismatch:
\(H^1(X^n;\mathbb Z)=0\), whereas both local quotient maps pull nonzero
target \(H^1\)-classes back to nonzero classes on the cusp fibre.  Hence no
map on \(X^2\) or \(X^4\) has the required restriction, and precomposition
with the smash collapses in \cref{wb:cusp-spec-smash-spheres} cannot produce
the local quotient maps.

\item Identify the metric quotient circle with
\(C_T=\mathbb R/\ell_T\mathbb Z\) by
\(\theta\mapsto\ell_T\theta\).  For
\cref{wb:cusp-spec-target-smash-map}, target-smash collapse is constant
whenever one source coordinate is the identity, so the modified map does
factor through \(F_T^{\wedge4}\).  Its chosen base point has distance
\(\ell_T/2\) from the collapsed target fat wedge.  If two points lie
within distance \(R\) of it, a route through the collapsed point has
length at least \(\ell_T-2R\), while their direct Euclidean distance is at
most \(2R\).  When \(\ell_T>4R\), neither the collapse nor circle
wrap-around shortens distances in that ball.

\item Let \(u_i\in H^1(F_T^4;\mathbb Z)\) be the pullback of the generator from
the \(i\)-th quotient circle.  The orientation class of
\(S^4=C_T^{\wedge4}\) pulls back under
\cref{wb:cusp-spec-target-smash-map} to
\[
 \pm u_1\smile u_2\smile u_3\smile u_4\ne0.
\]
But \(H^4(X^4;\mathbb Z)=0\).  Thus even this modified \(S^4\)-target map
does not extend to \(X^4\), and a fortiori does not arise from a map on
\(X^{\wedge4}\cong S^{24}\).

\item For completeness, Borel--Serre prove that \(S^{2n}\) admits no
almost-complex structure for \(n\ge4\).  Their argument applies to every
differentiable structure on the sphere: on \(S^{2n}\), all lower Chern
classes vanish, while their reduced-power calculation forces the top
Chern class to vanish modulo a suitable odd prime \(p<n\); on the other
hand that class is the Euler class, namely \(2h\) for a generator
\(h\in H^{2n}(S^{2n};\mathbb Z)\), a contradiction.  Taking \(n=6\) and
\(n=12\) gives the asserted obstructions for \(S^{12}\) and \(S^{24}\).
This topological obstruction does not itself prohibit the real
four-circle target \(G_T\), the target-smash \(S^4\), or their metric
limits.

\item Finally,
\(X*X\cong S^{13}\) and \(X^{*4}\cong S^{27}\), so joins do not produce
the proposed \(S^{12}\) or \(S^{24}\) dimensions.  The based suspension map $\Sigma^k h:\Sigma^k X\to\Sigma^k S^6\cong S^{6+k}$ gives $S^{12}$ at $k=6$ and $S^{24}$ at $k=18$. The displayed fixed-target fibre quotients are obstructed from extending by the preceding cohomology calculations; the subsequent cone constructions give the specified comparison maps with modified targets.
\end{enumerate}
\end{proof}

\subsection{Cone-compactified sphere carriers for the native cusp quotient}
\label{wb:cusp-cone-section}

This subsection answers a topological question which is strictly different
from the finite-cover construction.  No period is enlarged by passing to a
cover.  Instead, the native expanding quotient is retained on the cusp, the
parts on which that quotient is unavailable are collapsed, and the resulting
based map is passed through a smash product.  The outcome is an exact global
map with source \(S^{12}\), and there is an analogous \(S^{24}\) map from the
full-monodromy-invariant circle quotient.  The targets are conically
stratified spaces.  The Euclidean limits below use the intrinsic metrics
inherited from the flat quotient tori on the regular slices; they are not
ambient metric limits in one fixed metric
on the compact targets.

Fix a closed coordinate disc
\begin{equation}
 D_0=\{t\in\mathbb C:|t|\leq\epsilon\}\subset B,\qquad
 p_0=\{t=0\},\qquad N_0=f^{-1}(D_0),\qquad W=f^{-1}(0),
 \label{wb:cusp-cone-disc}
\end{equation}
small enough that it contains no singular value other than \(p_0\).  Put
\begin{equation}
 K_\Lambda=\mathbb Z\langle\widehat w,\widehat\delta\rangle,\qquad
 \overline\Lambda=\Lambda/K_\Lambda,\qquad
 \overline T^2=
 \frac{\overline\Lambda\otimes\mathbb R}{\overline\Lambda}.
 \label{wb:cusp-cone-quotient-lattice}
\end{equation}
The ordered basis of \(\overline\Lambda\) is the image of
\((\widehat\gamma,\widehat u)\).

\begin{lemma}[The native quotient over the whole punctured cusp]
\label{wb:cusp-cone-native-descent}
\StatementRecord{wb:cusp-cone-native-descent}{proved}{Historical events: GPT-6 Astra (OpenAI), ChatGPT Pro mode, effort Pro, S6\_09\_PAPERS, 17 November 2026: targeted proof check. Also checked in this writing. Pre-continuation writing events: GPT-6 Astra (OpenAI), ChatGPT Pro mode, effort Pro, S6\_10\_PAPERS, 18 November 2026: targeted proof check. Continuation metadata events: GPT-6 Astra (OpenAI), ChatGPT Pro mode, effort Pro, S6\_10\_PAPERS continuation, 19 November 2026: used and rechecked in this run; status-vocabulary reconciliation; metadata reconciliation only; no new mathematical check. \emph{Pre-continuation recorded proof scope.} GPT-6 Astra (OpenAI, ChatGPT Pro mode), S6\_09\_PAPERS, effort Pro, 2026-11-17; also checked in this writing by GPT-6 Astra (OpenAI), ChatGPT Pro mode, effort Pro, 18 November 2026}{the hypotheses, notation and definitions stated here; \cref{wb:cusp-spec-native-quotient}}

There is a well-defined real-analytic fibrewise group epimorphism
\begin{equation}
 \overline\pi:N_0\setminus W
 \longrightarrow D_0^*\times\overline T^2
 \label{wb:cusp-cone-native-family-map}
\end{equation}
whose kernel is the connected two-torus family defined by \(K_\Lambda\).
The target local system is trivial.  On the positive radial fibre \(F_T\),
\(t=e^{-2\pi T}\), the isomorphism
\begin{equation}
 \iota_T:\overline T^2\longrightarrow Q_T,\qquad
 [(a,b)]\longmapsto[Y_T(a,b)^{\mathsf T}]
 \label{wb:cusp-cone-radial-isometry}
\end{equation}
satisfies
\begin{equation}
 \pi_T=\iota_T\circ\overline\pi|_{F_T}.
 \label{wb:cusp-cone-native-restriction}
\end{equation}
If \(\overline T^2\) is given the pullback by \(\iota_T\) of the
Euclidean quotient metric on \(Q_T=\mathbb R^2/Y_T\mathbb Z^2\), then
\(\iota_T\) is an isometry and
\begin{equation}
 \operatorname{sys}(\overline T^2,g_T)
 =\operatorname{sys}(Q_T)\geq T-C.
 \label{wb:cusp-cone-native-systole}
\end{equation}
\end{lemma}



\begin{proof}
\begin{enumerate}[label=\arabic*.,leftmargin=*]
\item On the universal cover of \(D_0^*\), use the real-analytic flat marking
\[
 (z,x+\Lambda)\longmapsto[z,\Pi(z)x].
\]
The deck generator acts by
\[
 (z,x+\Lambda)\longmapsto(g_0z,M_0x+\Lambda).
\]
The displayed matrix \(M_0\) gives
\begin{equation}
 (M_0-I)(a,b,c,d)^{\mathsf T}=(0,0,b,-a)^{\mathsf T},\qquad
 \operatorname{im}(M_0-I)=K_\Lambda.
 \label{wb:cusp-cone-monodromy-image}
\end{equation}
Consequently \(M_0\) preserves \(K_\Lambda\) and induces the identity on
\(\overline\Lambda\).  The formula
\[
 [z,\Pi(z)x]\longmapsto
 \bigl([z],[x\bmod K_\Lambda]\bigr)
\]
is invariant under the deck action and proves
\cref{wb:cusp-cone-native-family-map}.  The lattice \(K_\Lambda\) is primitive,
so the kernel is connected.  On the radial fibre, projection to the
\((e_2,e_4)\)-coordinates sends
\(a\widehat\gamma+b\widehat u\) to \(Y_T(a,b)^{\mathsf T}\) and kills
\(K_\Lambda\).  This proves
\cref{wb:cusp-cone-radial-isometry} and \cref{wb:cusp-cone-native-restriction}.  The metric
and systolic assertion is the quotient calculation in
\cref{wb:cusp-spec-native-quotient}.
\end{enumerate}
\end{proof}

Define
\begin{equation}
 \mathcal A=
 \frac{[0,\epsilon]\times S^1\times\overline T^2}
 {\{0\}\times S^1\times\overline T^2\sim v_0,\quad
  \{\epsilon\}\times S^1\times\overline T^2\sim v_1}.
 \label{wb:cusp-cone-A-definition}
\end{equation}
Thus
\begin{equation}
 \mathcal A\cong
 \operatorname{Susp}(S^1\times\overline T^2)
 \cong\operatorname{Susp}(T^3),
 \label{wb:cusp-cone-A-suspension}
\end{equation}
where \(\operatorname{Susp}\) denotes the unreduced suspension.  The extra
circle is the angular coordinate of the punctured base.  If one deliberately
forgets that angle, the evident map
\(\mathcal A\to\operatorname{Susp}(\overline T^2)\) gives a smaller target
with the same restriction to every fixed radial fibre.  We retain
\(\mathcal A\) because it records the entire punctured-cusp quotient.

\begin{theorem}[Exact localized \(S^{12}\) route]
\label{wb:cusp-cone-S12-route}
\StatementRecord{wb:cusp-cone-S12-route}{proved}{Historical events: GPT-6 Astra (OpenAI), ChatGPT Pro mode, effort Pro, S6\_09\_PAPERS, 17 November 2026: targeted proof check. Also checked in this writing. Pre-continuation writing events: GPT-6 Astra (OpenAI), ChatGPT Pro mode, effort Pro, S6\_10\_PAPERS, 18 November 2026: targeted proof check. Continuation metadata events: GPT-6 Astra (OpenAI), ChatGPT Pro mode, effort Pro, S6\_10\_PAPERS continuation, 19 November 2026: used and rechecked in this run; status-vocabulary reconciliation; metadata reconciliation only; no new mathematical check. \emph{Pre-continuation recorded proof scope.} GPT-6 Astra (OpenAI, ChatGPT Pro mode), S6\_09\_PAPERS, effort Pro, 2026-11-17; also checked in this writing by GPT-6 Astra (OpenAI), ChatGPT Pro mode, effort Pro, 18 November 2026}{the hypotheses, notation and definitions stated here; \cref{wb:cusp-cone-A-definition}; \cref{wb:cusp-cone-A-suspension}; \cref{wb:cusp-cone-native-descent}; \cref{wb:cusp-cone-native-family-map}; \cite[Chapter~0, Smash Product]{Hatcher2002}; \cref{wb:cusp-cone-native-restriction}}

Choose \(x_0\in W\), and base \(\mathcal A\) at \(v_0\).  There is a
continuous surjective based quotient map
\begin{equation}
 q_2:(X,x_0)\longrightarrow(\mathcal A,v_0)
 \label{wb:cusp-cone-q2}
\end{equation}
with the following exact description.  For \(x\in N_0\setminus W\), write
\[
 t(f(x))=re^{i\phi},\qquad 0<r<\epsilon,
\]
and let \(\overline\pi_f(x)\in\overline T^2\) be the fibre component of
\(\overline\pi(x)\).  Then
\begin{equation}
 q_2(x)=[r,e^{i\phi},\overline\pi_f(x)],\qquad
 q_2(W)=v_0,\qquad
 q_2\bigl(X\setminus\operatorname{int}N_0\bigr)=v_1.
 \label{wb:cusp-cone-q2-formula}
\end{equation}
On \(f^{-1}(|t|=\epsilon)\), the value is \(v_1\).

The smash square is a global based map
\begin{equation}
 q_2^{\wedge2}:X\wedge X\longrightarrow
 \mathcal A\wedge\mathcal A.
 \label{wb:cusp-cone-S12-map}
\end{equation}
Using \(X\cong_{\mathrm{diff}}S^6\), a based choice of the diffeomorphism
identifies its source with
\begin{equation}
 X\wedge X\cong S^6\wedge S^6\cong S^{12}.
 \label{wb:cusp-cone-S12-source}
\end{equation}
For \(0<e^{-2\pi T}<\epsilon\), the canonical product-to-smash quotient
maps restrict to embeddings on \(F_T^2\) and on the corresponding interior
target-product cross-section.  Under those embeddings and \(\iota_T^2\),
the induced smash map restricts exactly to
\begin{equation}
 \pi_T\times\pi_T:F_T\times F_T\longrightarrow Q_T\times Q_T.
 \label{wb:cusp-cone-S12-fibre-map}
\end{equation}
With the Euclidean \(\ell^2\)-product flat quotient metric,
\begin{equation}
 \operatorname{sys}(Q_T\times Q_T)\geq T-C,\qquad
 (Q_T\times Q_T,(0,0))\longrightarrow(\mathbb R^4,0)
 \label{wb:cusp-cone-S12-R4-limit}
\end{equation}
in pointed local-isometric, pointed smooth, and pointed
Gromov--Hausdorff topology.  Once \(T-C>4R\), the radius-\(R\) ball is
exactly Euclidean.

The integral homology of the targets is
\begin{equation}
 \widetilde H_k(\mathcal A;\mathbb Z)\cong
 \begin{cases}
  \mathbb Z^3,&k=2,3,\\
  \mathbb Z,&k=4,\\
  0,&\text{otherwise},
 \end{cases}
 \label{wb:cusp-cone-A-homology}
\end{equation}
and
\begin{equation}
 \widetilde H_k(\mathcal A\wedge\mathcal A;\mathbb Z)\cong
 \begin{cases}
  \mathbb Z^9,&k=4,\\
  \mathbb Z^{18},&k=5,\\
  \mathbb Z^{15},&k=6,\\
  \mathbb Z^6,&k=7,\\
  \mathbb Z,&k=8,\\
  0,&\text{otherwise}.
 \end{cases}
 \label{wb:cusp-cone-A-smash-homology}
\end{equation}
In particular, the target is not \(S^4\), \(S^{12}\), or a compact
four-manifold.
\end{theorem}



\begin{proof}
\begin{enumerate}[label=\arabic*.,leftmargin=*]
\item The equivalence relation in \cref{wb:cusp-cone-A-definition} is closed: it is
the union of the diagonal with the squares of the two closed end
cross-sections.  A compact Hausdorff space modulo a closed equivalence
relation is compact Hausdorff.  Rescaling \(r\) to the unit interval gives
exactly the unreduced suspension in \cref{wb:cusp-cone-A-suspension}.

\item Well-definedness of the interior formula follows from
\cref{wb:cusp-cone-native-descent}.  A full turn in \(\phi\) changes the
marking by \(M_0\), but \(M_0\) is the identity modulo \(K_\Lambda\).

\item We check both collapsed strata.  If \(U\) is an open neighborhood of
\(v_0\), its inverse image in the cylinder is open and contains the compact
set \(\{0\}\times S^1\times\overline T^2\).  Hence, for some \(\delta>0\),
it contains
\([0,\delta)\times S^1\times\overline T^2\).  Therefore
\(q_2^{-1}(U)\) contains \(f^{-1}(|t|<\delta)\), a neighborhood of \(W\).
The same compactness argument at the other end shows that every neighborhood
of \(v_1\) contains the image of
\((\epsilon-\delta,\epsilon]\times S^1\times\overline T^2\).  Its preimage
together with \(X\setminus\operatorname{int}N_0\) is a neighborhood of the
outer boundary.  Continuity on the interior is
\cref{wb:cusp-cone-native-family-map}; the closed pasting lemma gives
continuity on \(X\).

\item Every interior point has a preimage because \(\overline\pi\) is fibrewise
surjective.  The nonempty sets \(W\) and
\(X\setminus\operatorname{int}N_0\) map onto \(v_0\) and \(v_1\),
respectively.  Thus \(q_2\) is surjective.  A continuous surjection from
compact \(X\) to Hausdorff \(\mathcal A\) is a quotient map.

\item The map is based, so functoriality gives \cref{wb:cusp-cone-S12-map}.  The
standard CW quotient gives \(S^m\wedge S^n\cong S^{m+n}\)
\cite[Chapter~0, Smash Product]{Hatcher2002}, proving
\cref{wb:cusp-cone-S12-source}.

\item The base point lies in \(W\), so \(F_T^2\) is disjoint from the source fat
wedge.  Its image lies in an interior cross-section of
\(\mathcal A^2\), disjoint from the target fat wedge.  The source and target
fat wedges are closed in the compact Hausdorff products \(X^2\) and
\(\mathcal A^2\), so both smash quotients are Hausdorff.  Their restrictions
to the relevant compact subspaces are continuous and one-to-one, and are
therefore embeddings.
Equations \cref{wb:cusp-cone-native-restriction} and \cref{wb:cusp-cone-q2-formula} give
\cref{wb:cusp-cone-S12-fibre-map}.

\item The period lattice of the product target is
\(\operatorname{diag}(Y_T,Y_T)\mathbb Z^4\).  Its systole is the systole of
\(Y_T\mathbb Z^2\), proving the first assertion in
\cref{wb:cusp-cone-S12-R4-limit}.  If \(x,y\in B_R(0)\subset\mathbb R^4\)
and \(0\ne\lambda\) is a product period, then
\[
 \|x-y+\lambda\|\geq\|\lambda\|-\|x-y\|
 \geq(T-C)-2R>2R\geq\|x-y\|
\]
when \(T-C>4R\).  The quotient projection is an exact isometry on the
Euclidean radius-\(R\) ball.  Conversely, any point of quotient distance
less than \(R\) has a shortest lattice representative of norm less than
\(R\), so this embedded ball is the entire quotient radius-\(R\) ball.
This proves all three pointed convergences.

\item Reduced suspension shifts reduced homology by one.  Since
\[
 H_1(T^3)=\mathbb Z^3,\qquad
 H_2(T^3)=\mathbb Z^3,\qquad H_3(T^3)=\mathbb Z,
\]
\cref{wb:cusp-cone-A-homology} follows.  All groups are free, so the Kunneth
theorem for smash products has no Tor terms.  The ranks in degrees
\(4,5,6,7,8\) are
\[
 3\cdot3,\quad 2(3\cdot3),\quad
 2(3\cdot1)+3\cdot3,\quad 2(3\cdot1),\quad1,
\]
which proves \cref{wb:cusp-cone-A-smash-homology}.
\end{enumerate}
\end{proof}

\paragraph{The localized one-circle construction.}
The full-monodromy-invariant character also gives a smaller cusp-supported
target.  Define
\begin{equation}
 \mathcal B=
 \frac{[0,\epsilon]\times S^1\times S^1}
 {\{0\}\times S^1\times S^1\sim b_0,\quad
  \{\epsilon\}\times S^1\times S^1\sim b_1}
 \cong\operatorname{Susp}(T^2),
 \label{wb:cusp-cone-B-definition}
\end{equation}
where the two circle coordinates in an interior cross-section are the base
angle and the value of \(\theta_\gamma\).

\begin{proposition}[Localized \(S^{24}\) route with target
\(\operatorname{Susp}(T^2)^{\wedge4}\)]
\label{wb:cusp-cone-S24-local-route}
\StatementRecord{wb:cusp-cone-S24-local-route}{proved}{Historical events: GPT-6 Astra (OpenAI), ChatGPT Pro mode, effort Pro, S6\_09\_PAPERS, 17 November 2026: targeted proof check. Also checked in this writing. Pre-continuation writing events: GPT-6 Astra (OpenAI), ChatGPT Pro mode, effort Pro, S6\_10\_PAPERS, 18 November 2026: targeted proof check. Continuation metadata events: GPT-6 Astra (OpenAI), ChatGPT Pro mode, effort Pro, S6\_10\_PAPERS continuation, 19 November 2026: used and rechecked in this run; status-vocabulary reconciliation; metadata reconciliation only; no new mathematical check. \emph{Pre-continuation recorded proof scope.} GPT-6 Astra (OpenAI, ChatGPT Pro mode), S6\_09\_PAPERS, effort Pro, 2026-11-17; also checked in this writing by GPT-6 Astra (OpenAI), ChatGPT Pro mode, effort Pro, 18 November 2026}{the hypotheses, notation and definitions stated here; \cref{wb:cusp-cone-S24-R4-limit}; \cref{wb:cusp-cone-S12-route}}

There is a continuous surjective based quotient map
\begin{equation}
 q_{1,\mathrm{cusp}}:(X,x_0)\longrightarrow(\mathcal B,b_0)
 \label{wb:cusp-cone-q1-local}
\end{equation}
given by
\begin{equation}
 q_{1,\mathrm{cusp}}(x)
 =[r,e^{i\phi},\theta_\gamma(x)]
 \quad(0<|t|=r<\epsilon),\qquad
 q_{1,\mathrm{cusp}}(W)=b_0,\qquad
 q_{1,\mathrm{cusp}}
 \bigl(X\setminus\operatorname{int}N_0\bigr)=b_1.
 \label{wb:cusp-cone-q1-local-formula}
\end{equation}
Its fourth smash power has source \(X^{\wedge4}\cong S^{24}\), and on
four synchronized radial fibres it is exactly
\begin{equation}
 \chi_{\gamma,T}^{\times4}:F_T^4\longrightarrow C_T^4.
 \label{wb:cusp-cone-S24-local-fibre-map}
\end{equation}
Consequently its intrinsic target slices satisfy the Euclidean limit
\cref{wb:cusp-cone-S24-R4-limit}.  The target homology is
\begin{equation}
 \widetilde H_k(\mathcal B;\mathbb Z)\cong
 \begin{cases}
  \mathbb Z^2,&k=2,\\
  \mathbb Z,&k=3,\\
  0,&\text{otherwise},
 \end{cases}
 \label{wb:cusp-cone-B-homology}
\end{equation}
and
\begin{equation}
 \widetilde H_k(\mathcal B^{\wedge4};\mathbb Z)\cong
 \begin{cases}
  \mathbb Z^{16},&k=8,\\
  \mathbb Z^{32},&k=9,\\
  \mathbb Z^{24},&k=10,\\
  \mathbb Z^8,&k=11,\\
  \mathbb Z,&k=12,\\
  0,&\text{otherwise}.
 \end{cases}
 \label{wb:cusp-cone-B-smash-homology}
\end{equation}
\end{proposition}



\begin{proof}
\begin{enumerate}[label=\arabic*.,leftmargin=*]
\item The monodromy invariance
\(\varepsilon_\gamma M_0=\varepsilon_\gamma\) makes the interior formula
single-valued after a turn in the base angle.  The two endpoint-continuity
arguments are identical to those in the proof of
\cref{wb:cusp-cone-S12-route}, with compact cross-section \(T^2\) in place of
\(T^3\).  Fibrewise surjectivity of \(\theta_\gamma\), together with the
nonempty inner fibre and exterior, proves surjectivity.  Compact-to-Hausdorff
makes the map a quotient map.

\item Since \(x_0\in W\), the fourth smash power exists.  The radial regular
fibres and their circle targets avoid the two closed fat wedges.  The
restrictions of the Hausdorff smash quotients to these compact subspaces are
continuous and one-to-one, hence embeddings, and the fibre map is exactly
\cref{wb:cusp-cone-S24-local-fibre-map}.  The intrinsic metric statement is
therefore the four-circle calculation
\cref{wb:cusp-cone-S24-R4-limit}.  Finally, suspension shifts the free homology
of \(T^2\), giving \cref{wb:cusp-cone-B-homology}.  In the fourth smash power,
choosing \(j\) of the four factors in degree \(3\) and the other
\(4-j\) factors in degree \(2\) gives degree \(8+j\) and rank
\(\binom4j2^{4-j}\).  This is \cref{wb:cusp-cone-B-smash-homology}.
\end{enumerate}
\end{proof}

\paragraph{The full-base circle target before coning.}
Choose pairwise disjoint closed coordinate discs \(D_0,D_1,D_2\) about the
three singular values and put
\begin{equation}
 P=B\setminus\bigcup_{j=0}^2\operatorname{int}D_j.
 \label{wb:cusp-cone-pair-of-pants}
\end{equation}
For \(j=1,2\), let \(C_j\) be the solid-torus Seifert target in
\cref{wb:cusp-spec-finite-Seifert-target}, and let \(h_j^{S^1}\) be the
target overlap map in \cref{wb:cusp-spec-finite-overlap-identity}.  Glue
\begin{equation}
 \mathcal M=(P\times S^1)
 \underset{h_1^{S^1}}{\cup}C_1
 \underset{h_2^{S^1}}{\cup}C_2.
 \label{wb:cusp-cone-M-definition}
\end{equation}
The gluing is along boundary tori over \(\partial D_1,\partial D_2\).
The remaining boundary is
\begin{equation}
 \partial\mathcal M=\partial D_0\times S^1\cong T^2.
 \label{wb:cusp-cone-M-boundary}
\end{equation}
Define
\begin{equation}
 \mathcal C=\mathcal M\cup_{\partial\mathcal M}
 \operatorname{Cone}(\partial\mathcal M),
 \label{wb:cusp-cone-C-definition}
\end{equation}
and write \(c_0\) for the cone vertex.
This changes the codomain at the cusp: the whole boundary two-torus,
including the target circle, is collapsed to \(c_0\).  It therefore does
not contradict the proved nonextension of the literal map
\(\chi_\gamma\) to the fixed product target \(B\times S^1\).

\begin{theorem}[Exact cone-compactified \(S^{24}\) route]
\label{wb:cusp-cone-S24-route}
\StatementRecord{wb:cusp-cone-S24-route}{proved}{Historical events: GPT-6 Astra (OpenAI), ChatGPT Pro mode, effort Pro, S6\_09\_PAPERS, 17 November 2026: targeted proof check. Also checked in this writing. Pre-continuation writing events: GPT-6 Astra (OpenAI), ChatGPT Pro mode, effort Pro, S6\_10\_PAPERS, 18 November 2026: targeted proof check. Continuation metadata events: GPT-6 Astra (OpenAI), ChatGPT Pro mode, effort Pro, S6\_10\_PAPERS continuation, 19 November 2026: used and rechecked in this run; status-vocabulary reconciliation; metadata reconciliation only; no new mathematical check. \emph{Pre-continuation recorded proof scope.} GPT-6 Astra (OpenAI, ChatGPT Pro mode), S6\_09\_PAPERS, effort Pro, 2026-11-17; also checked in this writing by GPT-6 Astra (OpenAI), ChatGPT Pro mode, effort Pro, 18 November 2026}{the hypotheses, notation and definitions stated here; \cref{wb:cusp-spec-finite-Seifert-target}; \cref{wb:cusp-spec-global-circle-circumference}; \cref{wb:cusp-cone-S12-route}}

There is a continuous surjective based quotient map
\begin{equation}
 q_1:(X,x_0)\longrightarrow(\mathcal C,c_0),\qquad x_0\in W,
 \label{wb:cusp-cone-q1}
\end{equation}
defined as follows:
\begin{enumerate}
\item over \(P\), it is
\((f,\theta_\gamma):f^{-1}(P)\to P\times S^1\), equivalently the
restriction of \(\chi_\gamma\) to \(f^{-1}(P)\);
\item over \(D_j\), \(j=1,2\), it is the finite-filling map
\(N_j\to C_j\) from \cref{wb:cusp-spec-finite-Seifert-target};
\item on \(N_0\setminus W\), writing \(t=re^{i\phi}\), it is
\begin{equation}
 x\longmapsto[r,(e^{i\phi},\theta_\gamma(x))]
 \in\operatorname{Cone}(\partial D_0\times S^1),
 \label{wb:cusp-cone-q1-cusp-formula}
\end{equation}
and every point of \(W\) maps to \(c_0\).
\end{enumerate}
The fourth smash power is a global based map
\begin{equation}
 q_1^{\wedge4}:X^{\wedge4}\longrightarrow\mathcal C^{\wedge4}.
 \label{wb:cusp-cone-S24-map}
\end{equation}
Using \(X\cong_{\mathrm{diff}}S^6\), its source is
\begin{equation}
 X^{\wedge4}\cong(S^6)^{\wedge4}\cong S^{24}.
 \label{wb:cusp-cone-S24-source}
\end{equation}
On four synchronized radial fibres, the canonical product-to-smash quotient
maps restrict to embeddings on the source fibre product and the interior
target-product cross-section.  Under those embeddings,
\cref{wb:cusp-cone-S24-map} restricts exactly to
\begin{equation}
 \chi_{\gamma,T}^{\times4}:F_T^4\longrightarrow C_T^4,\qquad
 C_T=\mathbb R/\ell_T\mathbb Z.
 \label{wb:cusp-cone-S24-fibre-map}
\end{equation}
With its Euclidean \(\ell^2\)-product flat quotient metric,
\begin{equation}
 \operatorname{sys}(C_T^4)=\ell_T
 =T+H-B+O(T^{-1}),\qquad
 (C_T^4,0)\longrightarrow(\mathbb R^4,0)
 \label{wb:cusp-cone-S24-R4-limit}
\end{equation}
in pointed local-isometric, pointed smooth, and pointed
Gromov--Hausdorff topology.

The target homology is
\begin{equation}
 \widetilde H_k(\mathcal C;\mathbb Z)\cong
 \begin{cases}
  \mathbb Z,&k=2,3,\\
  0,&\text{otherwise},
 \end{cases}
 \label{wb:cusp-cone-C-homology}
\end{equation}
and
\begin{equation}
 \widetilde H_k(\mathcal C^{\wedge4};\mathbb Z)\cong
 \begin{cases}
  \mathbb Z,&k=8,12,\\
  \mathbb Z^4,&k=9,11,\\
  \mathbb Z^6,&k=10,\\
  0,&\text{otherwise}.
 \end{cases}
 \label{wb:cusp-cone-C-smash-homology}
\end{equation}
\end{theorem}



\begin{proof}
\begin{enumerate}[label=\arabic*.,leftmargin=*]
\item The compact spaces \(P\times S^1,C_1,C_2\) are glued along closed boundary
tori by homeomorphisms, producing a compact orientable three-manifold
\(\mathcal M\) with boundary \(\partial D_0\times S^1\).  Attaching the cone
along this closed boundary gives a compact Hausdorff finite CW space
\(\mathcal C\).

\item The overlap identity
\[
 \chi_\gamma\circ h_j=h_j^{S^1}\circ\overline\chi_j
\]
proves that the first two pieces of \(q_1\) agree.  On \(\partial D_0\), the
cone base is glued to the remaining boundary of \(P\times S^1\), so the
first and third pieces agree.  If \(U\) is a neighborhood of \(c_0\), its
inverse image in the cone contains \([0,\delta)\times T^2\) for some
\(\delta>0\), by compactness of the cone base.  Therefore \(q_1^{-1}(U)\)
contains \(f^{-1}(|t|<\delta)\), a neighborhood of \(W\).  This proves
continuity at the new stratum; continuity elsewhere follows from the input
maps and the closed pasting lemma.

\item Every regular circle quotient and both finite Seifert targets are
surjective, and \(W\) maps onto the cone vertex.  Thus \(q_1\) is a
continuous surjection from compact \(X\) to Hausdorff \(\mathcal C\), hence
a quotient map.  It is based, so its fourth smash power exists.  Iterating
\(S^m\wedge S^n\cong S^{m+n}\) gives \cref{wb:cusp-cone-S24-source}.

\item The base point lies in \(W\), while each radial regular fibre is disjoint
from \(W\).  Hence \(F_T^4\) and its product-circle image avoid the closed
source and target fat wedges.  The canonical quotient maps from the source
product and the indicated interior target-product cross-section to their
smash products restrict to embeddings.  Under those embeddings, the induced
smash map restricts to \(\chi_{\gamma,T}^{\times4}\), which itself retains
its displayed nontrivial torus kernel.
The quotient-metric calculation in
\cref{wb:cusp-spec-global-circle-circumference} gives \(\ell_T\).
The product lattice is \(\ell_T\mathbb Z^4\), so the radius-\(R\) argument,
including the shortest-representative proof that the embedded Euclidean ball
is the whole quotient ball, in the proof of \cref{wb:cusp-cone-S12-route} proves
\cref{wb:cusp-cone-S24-R4-limit}.

\item For homology, let \(\rho_0,\rho_1,\rho_2\) be the boundary loops of the pair
of pants, with \(\rho_0+\rho_1+\rho_2=0\), and let \(\eta\) be the regular
target-circle class.  With the boundary-loop and circle orientations used
in the finite-filling presentation, the meridian relations are
\begin{equation}
 3\rho_1-\eta=0,\qquad4\rho_2+\eta=0.
 \label{wb:cusp-cone-Seifert-meridians}
\end{equation}
Reversing the circle orientation reverses both displayed \(\eta\)-signs
and changes none of the following groups.  Thus
\begin{equation}
 H_1(\mathcal M;\mathbb Z)
 =\frac{\mathbb Z\langle\eta,\rho_1,\rho_2\rangle}
 {\langle3\rho_1-\eta,\;4\rho_2+\eta\rangle}
 \cong\mathbb Z\langle g\rangle,
 \label{wb:cusp-cone-M-H1}
\end{equation}
where
\begin{equation}
 \rho_0=-g,\qquad\rho_1=4g,\qquad
 \rho_2=-3g,\qquad\eta=12g.
 \label{wb:cusp-cone-M-boundary-map}
\end{equation}
In particular,
\(H_1(\partial\mathcal M)=\mathbb Z\langle\rho_0,\eta\rangle
\to H_1(\mathcal M)\) is surjective.  The orientable three-manifold
\(\mathcal M\) has one torus boundary, \(H_3(\mathcal M)=0\), and
\(\chi(\mathcal M)=\frac12\chi(\partial\mathcal M)=0\); hence
\(H_2(\mathcal M;\mathbb Q)=0\).  To obtain the integral statement without
suppressing torsion, Poincare--Lefschetz duality and the universal
coefficient theorem give
\[
 H_2(\mathcal M;\mathbb Z)
 \cong H^1(\mathcal M,\partial\mathcal M;\mathbb Z)
 \cong \operatorname{Hom}
 \bigl(H_1(\mathcal M,\partial\mathcal M;\mathbb Z),\mathbb Z\bigr),
\]
because \(H_0(\mathcal M,\partial\mathcal M;\mathbb Z)=0\).  Thus
\(H_2(\mathcal M;\mathbb Z)\) is free; its rank is zero, so it vanishes.

\item The cone attachment is the homotopy cofiber of
\(\partial\mathcal M\hookrightarrow\mathcal M\).  Since the boundary
inclusion is a cofibration,
\[
 \widetilde H_k(\mathcal C;\mathbb Z)
 \cong H_k(\mathcal M,\partial\mathcal M;\mathbb Z).
\]
Poincare--Lefschetz duality and
\(H^0(\mathcal M)=H^1(\mathcal M)=\mathbb Z\),
\(H^2(\mathcal M)=H^3(\mathcal M)=0\), give
\cref{wb:cusp-cone-C-homology}.  The homology is free, so the smash Kunneth
theorem says that four factors of degree \(2\) or \(3\) produce ranks
\(\binom4j\) in degree \(8+j\), proving
\cref{wb:cusp-cone-C-smash-homology}.
\end{enumerate}
\end{proof}

\begin{definition}[The three radial target-slice embeddings]
\label{wb:cusp-cone-target-slice-embeddings}
\StatementRecord{wb:cusp-cone-target-slice-embeddings}{proved}{Historical events: GPT-6 Astra (OpenAI), ChatGPT Pro mode, effort Pro, S6\_09\_PAPERS, 17 November 2026: targeted proof check. Also checked in this writing. Pre-continuation writing events: GPT-6 Astra (OpenAI), ChatGPT Pro mode, effort Pro, S6\_10\_PAPERS, 18 November 2026: targeted proof check. Continuation metadata events: GPT-6 Astra (OpenAI), ChatGPT Pro mode, effort Pro, S6\_10\_PAPERS continuation, 19 November 2026: used and rechecked in this run; status-vocabulary reconciliation; metadata reconciliation only; no new mathematical check. \emph{Pre-continuation recorded proof scope.} GPT-6 Astra (OpenAI, ChatGPT Pro mode), S6\_09\_PAPERS, effort Pro, 2026-11-17; also checked in this writing by GPT-6 Astra (OpenAI), ChatGPT Pro mode, effort Pro, 18 November 2026}{the hypotheses, notation and definitions stated here; the argument and qualifications immediately following}

Put \(r_T=e^{-2\pi T}\) and let
\[
 \vartheta_T:\mathbb R/\mathbb Z\longrightarrow
 C_T=\mathbb R/\ell_T\mathbb Z,\qquad [a]\longmapsto[\ell_Ta].
\]
For \(y_1,y_2\in Q_T\), define
\[
 j_T(y_1,y_2)=
 \bigl[[r_T,1,\iota_T^{-1}(y_1)],
       [r_T,1,\iota_T^{-1}(y_2)]\bigr]_{\wedge}
 \in\mathcal A^{\wedge2}.
\]
For \(c_1,\ldots,c_4\in C_T\), define
\[
 k_T^{\mathcal B}(c_1,\ldots,c_4)=
 \bigl[[r_T,1,\vartheta_T^{-1}(c_1)],\ldots,
       [r_T,1,\vartheta_T^{-1}(c_4)]\bigr]_{\wedge}
 \in\mathcal B^{\wedge4},
\]
and, using the cusp-cone piece of \(\mathcal C\),
\[
 k_T^{\mathcal C}(c_1,\ldots,c_4)=
 \bigwedge_{j=1}^4[r_T,(1,\vartheta_T^{-1}(c_j))]
 \in\mathcal C^{\wedge4}.
\]
These are embeddings because every displayed factor lies away from the smash
basepoint, so equality in the smash quotient is equality in every factor;
compact-to-Hausdorff injectivity then gives a topological embedding.
\end{definition}



\begin{proposition}[Intrinsic expansion versus ambient collapse]
\label{wb:cusp-cone-metric-scope}
\StatementRecord{wb:cusp-cone-metric-scope}{proved}{Historical events: GPT-6 Astra (OpenAI), ChatGPT Pro mode, effort Pro, S6\_09\_PAPERS, 17 November 2026: targeted proof check. Also checked in this writing. Pre-continuation writing events: GPT-6 Astra (OpenAI), ChatGPT Pro mode, effort Pro, S6\_10\_PAPERS, 18 November 2026: targeted proof check. Continuation metadata events: GPT-6 Astra (OpenAI), ChatGPT Pro mode, effort Pro, S6\_10\_PAPERS continuation, 19 November 2026: used and rechecked in this run; status-vocabulary reconciliation; metadata reconciliation only; no new mathematical check. \emph{Pre-continuation recorded proof scope.} GPT-6 Astra (OpenAI, ChatGPT Pro mode), S6\_09\_PAPERS, effort Pro, 2026-11-17; also checked in this writing by GPT-6 Astra (OpenAI), ChatGPT Pro mode, effort Pro, 18 November 2026}{the hypotheses, notation and definitions stated here; \cref{wb:cusp-cone-S12-R4-limit}; \cref{wb:cusp-cone-S24-R4-limit}; \cref{wb:cusp-cone-native-systole}}

Let \(D_{12}\), \(D_{24}^{\mathcal B}\), and
\(D_{24}^{\mathcal C}\) be arbitrary fixed metrics compatible with the
topologies of \(\mathcal A^{\wedge2}\), \(\mathcal B^{\wedge4}\), and
\(\mathcal C^{\wedge4}\), respectively.  In their corresponding ambient
metrics, the radial target
slices satisfy
\begin{equation}
 \operatorname{diam}_{\mathrm{amb}}
 \bigl(j_T(Q_T^2)\bigr)\longrightarrow0,\qquad
 \operatorname{diam}_{\mathrm{amb}}
 \bigl(k_T^{\mathcal B}(C_T^4)\bigr)\longrightarrow0,\qquad
 \operatorname{diam}_{\mathrm{amb}}
 \bigl(k_T^{\mathcal C}(C_T^4)\bigr)\longrightarrow0.
 \label{wb:cusp-cone-ambient-collapse}
\end{equation}
By contrast, in the intrinsic quotient metrics,
\begin{equation}
 \operatorname{diam}(Q_T^2)\geq\tfrac12(T-C)\longrightarrow\infty,\qquad
 \operatorname{diam}(C_T^4)\geq\tfrac12\ell_T\longrightarrow\infty.
 \label{wb:cusp-cone-intrinsic-expansion}
\end{equation}
Thus \cref{wb:cusp-cone-S12-R4-limit} and \cref{wb:cusp-cone-S24-R4-limit} assert
pointed local-isometric limits for the specified flat quotient metrics on the varying slices. The identity from each such slice to its embedded image pulls back the fixed ambient metric to a different metric: its diameter tends to zero by the displayed calculation.
\end{proposition}



\begin{proof}
\begin{enumerate}[label=\arabic*.,leftmargin=*]
\item Every neighborhood of the relevant cone vertex contains every sufficiently
small radial cross-section, uniformly in all compact angular and torus
coordinates.  In a compatible metric, every such cross-section eventually
lies in an arbitrarily small ball about the vertex.  Products followed by
the smash quotient preserve this uniform convergence, proving
\cref{wb:cusp-cone-ambient-collapse}.

\item For a flat torus with systole \(s\), choose a shortest primitive period
\(\lambda\).  The class of \(\lambda/2\) has distance \(s/2\) from zero:
for every period \(\mu\), the nonzero period \(\lambda-2\mu\) has length at
least \(s\).  Hence the diameter is at least \(s/2\).
Equations \cref{wb:cusp-cone-native-systole} and \cref{wb:cusp-cone-S24-R4-limit} prove
\cref{wb:cusp-cone-intrinsic-expansion}.
\end{enumerate}
\end{proof}

The obstruction to an almost-complex structure on \(S^{12}\),
with all three complete classical proofs, is given in
\cref{wb:cusp-cone-S12-almost-complex-obstruction}.


\begin{theorem}[Sharp obstructions to stronger interpretations]
\label{wb:cusp-cone-obstructions}
\StatementRecord{wb:cusp-cone-obstructions}{proved}{Historical events: GPT-6 Astra (OpenAI), ChatGPT Pro mode, effort Pro, S6\_09\_PAPERS, 17 November 2026: targeted proof check. Also checked in this writing. Pre-continuation writing events: GPT-6 Astra (OpenAI), ChatGPT Pro mode, effort Pro, S6\_10\_PAPERS, 18 November 2026: targeted proof check. Continuation metadata events: GPT-6 Astra (OpenAI), ChatGPT Pro mode, effort Pro, S6\_10\_PAPERS continuation, 19 November 2026: used and rechecked in this run; status-vocabulary reconciliation; metadata reconciliation only; no new mathematical check. \emph{Pre-continuation recorded proof scope.} GPT-6 Astra (OpenAI, ChatGPT Pro mode), S6\_09\_PAPERS, effort Pro, 2026-11-17; also checked in this writing by GPT-6 Astra (OpenAI), ChatGPT Pro mode, effort Pro, 18 November 2026}{the hypotheses, notation and definitions stated here; \cref{wb:cusp-cone-S12-route}; \cref{wb:cusp-cone-S12-almost-complex-obstruction}; \cite[\S15, Proposition~15.1 and the following remark,
p.~434]{BorelSerre1953}; \cref{wb:cusp-cone-kernel-contained}; \cref{wb:cusp-cone-finite-index-orbit}}

The sphere routes above have the following exact limitations.
\begin{enumerate}
\item Let \(r\geq1\), and let \(N\subset\Lambda^{\oplus r}\) be a
primitive kernel defining a connected fibrewise torus quotient over the
full regular base, or over any finite cover of that base.  Equip the quotient
with the norm induced by the unrescaled marked Euclidean metric on the
original fibres.  If its quotient systole tends to infinity on the radial
cusp, then
\begin{equation}
 \operatorname{rank}(\Lambda^{\oplus r}/N)\leq r.
 \label{wb:cusp-cone-rank-bound}
\end{equation}
Thus two copies cannot give a full-base four-torus quotient.  The
\(S^{12}\) map in \cref{wb:cusp-cone-S12-route} evades the obstruction only by
localizing to the cusp and collapsing the exterior.

\item No nontrivial finite base cover, including one extending with
branching over the original singular values, makes \(K_\Lambda\) a
sub-local system over the full regular base.  The local cyclic change
\(t=s^M\) used for the auxiliary degree-\(M^2\) cover cannot extend as a
connected cover of \(B\cong\mathbb P^1\) branched only at \(p_0\).

\item There is no map from \(X^2\), \(X\wedge X\), \(X^4\), or
\(X^{\wedge4}\) to a fixed \(T^4\) whose restriction retains all four
target character classes of the displayed product quotient.  There is no
map to a fixed \(S^4\) whose restriction is a degree-one collapse of
\(Q_T^2\), or of the four-circle target, for even one regular fibre.  The
conical targets evade these obstructions because the fibre character and
orientation classes do not extend as global target classes.

\item Cartesian products do not produce the named spheres:
\begin{equation}
 X^2\simeq S^6\times S^6\not\simeq S^{12},\qquad
 X^4\simeq(S^6)^4\not\simeq S^{24}.
 \label{wb:cusp-cone-products-not-spheres}
\end{equation}
The two- and four-fold synchronized fibre products have real dimensions
\(10\) and \(18\), respectively, and generally have singular completed
fibres.  Joins give
\begin{equation}
 X*X\cong S^{13},\qquad X^{*4}\cong S^{27}.
 \label{wb:cusp-cone-join-dimensions}
\end{equation}

\item Since \(X\cong S^6\) is simply connected, every connected unbranched
finite cover of \(X\), \(X^r\), \(X\wedge X\cong S^{12}\), or
\(X^{\wedge4}\cong S^{24}\) is trivial.  The local degree-\(M^2\)
unspooling cover is not the restriction of a nontrivial connected
unbranched cover of any of these compact sources.  The finite flat normal covers of \cref{maps:src:C.1.5} use the explicit base cover and normalization recorded there; the present assertion concerns unbranched covers of the listed compact sources.

\item By \cref{wb:cusp-cone-S12-almost-complex-obstruction}, no smooth manifold
homeomorphic to \(S^{12}\) admits an almost-complex structure.  Nor does
\(S^{24}\), with any differentiable structure, by the Borel--Serre theorem
\cite[\S15, Proposition~15.1 and the following remark,
p.~434]{BorelSerre1953}.  Thus the smash sources cannot be promoted to
compact complex spheres.
\end{enumerate}
\end{theorem}



\begin{proof}
\begin{enumerate}[label=\arabic*.,leftmargin=*]
\item Every element of \(K_\Lambda^{\oplus r}\) has uniformly bounded physical
length on the radial cusp.  If one survived in the quotient lattice, it
would be a nonzero period of bounded quotient norm.  Hence
\begin{equation}
 K_\Lambda^{\oplus r}\subset N.
 \label{wb:cusp-cone-kernel-contained}
\end{equation}
For the full monodromy group,
\(A_1\widehat w=\widehat u-\widehat w\) and invariance force every copy of
\(\widehat u\) into \(N\).  Therefore
\[
 L_\Lambda^{\oplus r}
 =\mathbb Z\langle\widehat u,\widehat w,\widehat\delta\rangle^{\oplus r}
 \subset N,
\]
which proves \cref{wb:cusp-cone-rank-bound}.

\item The same conclusion holds after a finite base cover.  Let \(H\) be its
finite-index subgroup and put \(H_*=A_2A_1=I+N_*\), \(N_*^2=0\).
Every finite-index subgroup contains \(H_*^n\) for some \(n>0\), while
\begin{equation}
 H_*^n\widehat w
 =\widehat w+n(\widehat u-\widehat\delta).
 \label{wb:cusp-cone-finite-index-orbit}
\end{equation}
Equations
\cref{wb:cusp-cone-kernel-contained} and \cref{wb:cusp-cone-finite-index-orbit} give
\(n\widehat u\in N\).  Primitivity gives \(\widehat u\in N\), proving the
same rank bound.  The orbit also gives
\(H_*^nK_\Lambda\ne K_\Lambda\) for every \(n>0\), so no finite cover of
the regular base preserves the native kernel.  A compact cover branched
over the singular values restricts to such a cover.  Finally, a connected
finite cover of \(\mathbb P^1\) branched only at \(p_0\) would restrict to
a connected unbranched finite cover of
\(\mathbb P^1\setminus\{p_0\}\cong\mathbb C\), hence has degree one.

\item For a fixed \(T^4\) target, \(H^1(X^r;\mathbb Z)=0\), as does the first
cohomology of the smash spheres.  A global map therefore pulls every target
character class back to zero.  On a regular fibre the quotient homomorphisms
inject the target character lattice into the character lattice of the
source torus, giving four independent nonzero pullbacks, a contradiction.

\item For the fixed \(S^4\) obstruction, let
\(\omega_T\in H^2(Q_T;\mathbb Z)\) be the orientation class.  Since
\(\pi_T^*:H^1(Q_T)\to H^1(F_T)\) injects two independent classes and torus
cohomology is their exterior algebra,
\begin{equation}
 (\pi_T\times\pi_T)^*(\omega_T\times\omega_T)
 \ne0\quad\text{in }H^4(F_T^2;\mathbb Z).
 \label{wb:cusp-cone-fibre-H4-nonzero}
\end{equation}
A degree-one collapse \(Q_T^2\to S^4\) pulls the orientation generator
back to \(\pm(\omega_T\times\omega_T)\).  But
\[
 H^4(X^2;\mathbb Z)=H^4(S^{12};\mathbb Z)=0,
\]
so no global extension has this restriction.  The four-circle case is the
same: its orientation class pulls back to the exterior product of four
independent circle characters, while
\(H^4(X^4)=H^4(S^{24})=0\).

\item The Kunneth theorem gives \(H_6(X^2)=\mathbb Z^2\), whereas
\(H_6(S^{12})=0\), and gives nonzero groups of ranks \(4,6,4\) in degrees
\(6,12,18\) for \(X^4\), whereas \(S^{24}\) has none.  The real dimension
of an \(r\)-fold synchronized fibre product is
\(6r-2(r-1)=4r+2\), giving \(10\) and \(18\).  The join formula
\(S^m*S^n\cong S^{m+n+1}\) proves
\cref{wb:cusp-cone-join-dimensions}.  The remaining assertions follow from
simple connectivity, \cref{wb:cusp-cone-S12-almost-complex-obstruction}, and
the cited Borel--Serre theorem for \(S^{24}\).
\end{enumerate}
\end{proof}

\paragraph{Canonicity and exact scope.}
The lattice \(K_\Lambda=\ker(M_0-I)=\operatorname{im}(M_0-I)\) and the
punctured-cusp quotient are intrinsic to the marked cusp monodromy.
The compact map \(q_2\) additionally depends on the chosen cusp coordinate,
closed cusp disc, and the decision to collapse its outer boundary and the
exterior.  The identifications
\(X^{\wedge2}\cong S^{12}\) and \(X^{\wedge4}\cong S^{24}\) depend on a
based choice of \(X\cong S^6\).  None of these choices is claimed to be
holomorphic or canonical in the algebraic category.

The map \(q_2^{\wedge2}\) is the direct two-copy construction from the
native \(T^2\times\mathbb R^2\) cusp: it is globally defined on the smash
sphere, and its radial regular-fibre quotients are \(Q_T\times Q_T\), which
converge pointed local-isometrically, smoothly, and Gromov--Hausdorff to
\(\mathbb R^4\).  It is not a torus quotient over the full
regular base.  The map \(q_{1,\mathrm{cusp}}^{\wedge4}\) is the parallel
cusp-supported construction with target
\(\operatorname{Susp}(T^2)^{\wedge4}\).  The stronger map
\(q_1^{\wedge4}\) uses the quotient over every regular and finite-filled
fibre before the cusp cone is attached; four copies are minimal for that
full-base all-expanding group quotient.
Neither construction gives a compact complex \(S^{12}\) or \(S^{24}\), an
ambient Euclidean limit in a fixed compact metric target, or a claim about
the interacting Yang--Mills spectrum.

The argument does not rule out a separately constructed branched or
toroidal compactification of the auxiliary finite covers after additional
branching away from \(p_0\).  Such a construction would have to specify its
global base cover, normalize and resolve every pulled-back singular fibre,
track the lifted metric, and determine the topology of the resulting total
space.  No such smooth compact sphere construction follows from the present
data.


\subsubsection{Optimal sphere exponents and the exact Euclidean boundary}
\label{wb:cusp-cone-optimal-strengthening}

Within the class of native connected fibrewise torus quotients equipped with
the induced, unscaled quotient of the marked Euclidean metric, the two sphere
dimensions in the preceding constructions are forced by two different
integral quotient problems.  Cusp-locally, each copy has an expanding
quotient of rank two, so two copies are necessary and sufficient within that
class.  Over the full regular base, monodromy leaves at most one expanding
quotient direction per copy, so four copies are necessary and sufficient
within the same class.  We now prove this scoped optimality, identify the
compact conical carriers, and construct a noncompact comparison manifold in
which the limiting \(\mathbb R^4\) is an actual boundary fibre.  There is only
a continuous comparison map from this manifold to the compact conical
carrier: it set-theoretically collapses the new boundary, is a homeomorphism
away from that boundary, and is nonproper and nonquotient.

Throughout this subsubsection,
\[
 B_R^Z(z):=\{w\in Z:d_Z(z,w)\leq R\}
\]
denotes the closed radius-\(R\) ball; when the ambient space and center are
clear, superscripts and centers are suppressed.

\begin{theorem}[Optimality within the native connected torus-quotient class]
\label{wb:cusp-cone-optimal-exponents}
\StatementRecord{wb:cusp-cone-optimal-exponents}{proved}{Historical events: GPT-6 Astra (OpenAI), ChatGPT Pro mode, effort Pro, S6\_09\_PAPERS, 17 November 2026: targeted proof check. Also checked in this writing. Pre-continuation writing events: GPT-6 Astra (OpenAI), ChatGPT Pro mode, effort Pro, S6\_10\_PAPERS, 18 November 2026: targeted proof check. Continuation metadata events: GPT-6 Astra (OpenAI), ChatGPT Pro mode, effort Pro, S6\_10\_PAPERS continuation, 19 November 2026: used and rechecked in this run; status-vocabulary reconciliation; metadata reconciliation only; no new mathematical check. \emph{Pre-continuation recorded proof scope.} GPT-6 Astra (OpenAI, ChatGPT Pro mode), S6\_09\_PAPERS, effort Pro, 2026-11-17; also checked in this writing by GPT-6 Astra (OpenAI), ChatGPT Pro mode, effort Pro, 18 November 2026}{the hypotheses, notation and definitions stated here; the explicit constructions and calculations in the proof}

Consider only native connected fibrewise torus quotients of copies of the
original marked fibres, and equip every such quotient with the induced,
unscaled Euclidean quotient metric.

\begin{enumerate}
\item For a native connected torus quotient of \(r\) copies over the
punctured cusp, if the quotient systole tends to infinity on the positive
radial ray, then its lattice rank is at most \(2r\).  A rank-four quotient
therefore requires \(r\geq2\).  Equality is attained for \(r=2\) by
\[
 (\Lambda/K_\Lambda)^{\oplus2},
 \qquad K_\Lambda=
 \mathbb Z\langle\widehat w,\widehat\delta\rangle,
\]
whose physical period lattice is
\[
 L_T^{(12)}=\operatorname{diag}(Y_T,Y_T)\mathbb Z^4.
\]
The corresponding smash source is
\(X^{\wedge2}\cong S^{12}\).

\item For a native connected torus quotient of \(r\) copies over the full
regular base, or over a finite cover of that base, radial divergence of the
quotient systole forces its lattice rank to be at most \(r\).  A rank-four
quotient therefore requires \(r\geq4\).  Equality is attained by the four
copies of the full-monodromy-invariant character \(\chi_\gamma\); on the
radial fibre its period lattice is
\[
 L_T^{(24)}=\ell_T\mathbb Z^4,
 \qquad
 \ell_T=T+H-B+O(T^{-1}),
\]
and the corresponding smash source is
\(X^{\wedge4}\cong S^{24}\).
\end{enumerate}

Thus, among precisely these native connected torus quotients with their
induced unscaled marked Euclidean metrics, \(S^{12}\) is the smallest smash
sphere produced by the cusp-local rank-two quotient, and \(S^{24}\) is the
smallest smash sphere produced by a quotient which survives the full
regular-base monodromy.  No optimality is asserted here among arbitrary
continuous quotient maps or after a \(T\)-dependent rescaling of the metric.
\end{theorem}



\begin{proof}
\begin{enumerate}[label=\arabic*.,leftmargin=*]
\item For \(r\) copies, every element of
\[
 K_\Lambda^{\oplus r}
 =\mathbb Z\langle\widehat w,\widehat\delta\rangle^{\oplus r}
\]
has uniformly bounded physical length on the radial cusp.  Let
\(N\subset\Lambda^{\oplus r}\) be a primitive sublattice defining a
connected torus quotient.  If some \(v\in K_\Lambda^{\oplus r}\) had
nonzero class in \(\Lambda^{\oplus r}/N\), that nonzero quotient period
would have quotient norm at most the original norm of \(v\), uniformly in
\(T\).  The quotient systole could not diverge.  Hence
\[
 K_\Lambda^{\oplus r}\subset N,
\]
and
\[
 \operatorname{rank}(\Lambda^{\oplus r}/N)
 \leq 4r-2r=2r.
\]
For \(r=2\), the primitive kernel \(K_\Lambda^{\oplus2}\) attains this
bound.  The cusp monodromy is the identity modulo \(K_\Lambda\), so the
quotient descends over the punctured cusp.  Its radial lattice is exactly
\(\operatorname{diag}(Y_T,Y_T)\mathbb Z^4\), and
\[
 \operatorname{sys}(L_T^{(12)})
 =\operatorname{sys}(Y_T\mathbb Z^2)\geq T-C\longrightarrow\infty.
\]

\item For a quotient over the full regular base, monodromy invariance gives a
stronger bound.  Since \(\widehat w\in N\) and
\[
 A_1\widehat w=\widehat u-\widehat w,
\]
one has \(\widehat u\in N\).  Together with
\(\widehat w,\widehat\delta\in N\), this places the rank-three lattice
\[
 \mathbb Z\langle\widehat u,\widehat w,
 \widehat\delta\rangle^{\oplus r}
\]
inside \(N\), and hence
\[
 \operatorname{rank}(\Lambda^{\oplus r}/N)\leq r.
\]
The same conclusion holds after a finite base cover.  Indeed, if
\(\Gamma_{\mathrm{mon}}\) is the corresponding finite-index monodromy
subgroup and \(G_*=A_2A_1=I+N_*\), then \(N_*^2=0\), and
\(\Gamma_{\mathrm{mon}}\) contains \(G_*^m\) for
some \(m>0\).  The exact orbit formula
\[
 G_*^m\widehat w
 =\widehat w+m(\widehat u-\widehat\delta)
\]
implies \(m\widehat u\in N\); primitivity of \(N\) implies
\(\widehat u\in N\).

\item The invariant character \(\chi_\gamma\) has precisely the complementary
rank-one quotient in each copy.  The full-base construction, including the
two finite Seifert fillings and the cone at the unipotent cusp, proves that
four copies attain rank four.  Finally, after choosing the based
diffeomorphism \(X\cong_{\mathrm{diff}}S^6\) already reconstructed in this
workbench,
\(X^{\wedge r}\cong(S^6)^{\wedge r}\cong S^{6r}\), which gives the two
claimed sphere dimensions.
\end{enumerate}
\end{proof}

\begin{theorem}[Exact sphere-to-sphere carriers with moving four-spheres]
\label{wb:cusp-cone-angle-forgotten-sphere-carriers}
\StatementRecord{wb:cusp-cone-angle-forgotten-sphere-carriers}{proved}{Historical events: GPT-6 Astra (OpenAI), ChatGPT Pro mode, effort Pro, S6\_09\_PAPERS, 17 November 2026: targeted proof check. Also checked in this writing. Pre-continuation writing events: GPT-6 Astra (OpenAI), ChatGPT Pro mode, effort Pro, S6\_10\_PAPERS, 18 November 2026: targeted proof check. Continuation metadata events: GPT-6 Astra (OpenAI), ChatGPT Pro mode, effort Pro, S6\_10\_PAPERS continuation, 19 November 2026: used and rechecked in this run; status-vocabulary reconciliation; metadata reconciliation only; no new mathematical check. \emph{Pre-continuation recorded proof scope.} GPT-6 Astra (OpenAI, ChatGPT Pro mode), S6\_09\_PAPERS, effort Pro, 2026-11-17; also checked in this writing by GPT-6 Astra (OpenAI), ChatGPT Pro mode, effort Pro, 18 November 2026}{the hypotheses, notation and definitions stated here; the explicit constructions and calculations in the proof}

There are continuous surjective based quotient maps
\[
 \widehat q_3:(X,x_0)\longrightarrow(S^3,*),
 \qquad
 \widehat q_\gamma:(X,x_0)\longrightarrow(S^2,*),
\]
whose smash powers are continuous surjective quotient maps
\begin{equation}
 \widehat q_3^{\wedge2}:S^{12}\twoheadrightarrow S^6,
 \qquad
 \widehat q_\gamma^{\wedge4}:S^{24}\twoheadrightarrow S^8.
 \label{wb:cusp-cone-sphere-to-sphere-maps}
\end{equation}
Use the marked presentation \(\overline T^2=\mathbb R^2/\mathbb Z^2\) and
put
\[
 K^{(1)}
 =((\mathbb R/\mathbb Z)\times\{0\})
  \cup(\{0\}\times(\mathbb R/\mathbb Z)),
 \qquad
 c:\overline T^2\longrightarrow\overline T^2/K^{(1)}\cong S^2.
\]
Here \(c\) is the quotient collapsing the coordinate one-skeleton.  For
every sufficiently large \(T\), put
\begin{align*}
 r_T&=e^{-2\pi T},\\
 K_T^{(1)}&:=\iota_T(K^{(1)})\subset Q_T,\\
 S_T^2&:=\overline T^2/K^{(1)}\cong Q_T/K_T^{(1)},\\
 c_T&:=c\circ\iota_T^{-1}:Q_T\longrightarrow S_T^2,
\end{align*}
where the identification of the two displayed quotients is induced by
\(\iota_T\).  Within this theorem and its proof, \(c_T\) denotes the
displayed collapse map; the earlier scalar \(c_T=\xi_T-\eta_T\) does not
enter.
At reduced-suspension cross-section height \(r_T\), the first target contains
a based topologically embedded four-sphere, equipped with its intrinsic
quotient length metric,
\[
 \mathbb S_{12,T}^4=S_T^2\wedge S_T^2\hookrightarrow S^6,
\]
and the second target contains, at the same cross-section height, a based
topologically embedded four-sphere, equipped with its intrinsic quotient
length metric,
\[
 \mathbb S_{24,T}^4=C_T^{\wedge4}\hookrightarrow S^8.
\]
On synchronized radial fibres, the exact restrictions factor as
\begin{align}
 F_T^2&\xrightarrow{\ \pi_T\times\pi_T\ }
 Q_T^2\xrightarrow{\ c_T\times c_T\ }(S_T^2)^2
 \xrightarrow{\ \rho_{12,T}\ }S_T^2\wedge S_T^2
 =\mathbb S_{12,T}^4,
 \label{wb:cusp-cone-S12-S4-radial-map}\\
 F_T^4&\xrightarrow{\ \chi_{\gamma,T}^{\times4}\ }
 C_T^4\xrightarrow{\ \rho_{24,T}\ }C_T^{\wedge4}
 =\mathbb S_{24,T}^4,
 \label{wb:cusp-cone-S24-S4-radial-map}
\end{align}
Here \(c_T\) collapses \(K_T^{(1)}\), and \(\rho_{12,T}\) and
\(\rho_{24,T}\) are the smash quotient maps.  Thus the
full collapsed locus in \(Q_T^2\) is exactly
\begin{equation}
 W_{12,T}
 =\bigl(K_T^{(1)}\times Q_T\bigr)
  \cup\bigl(Q_T\times K_T^{(1)}\bigr),
 \label{wb:cusp-cone-S12-exact-collapsed-locus}
\end{equation}
and the full collapsed locus in \(C_T^4\) is exactly
\begin{equation}
 W_{24,T}
 =\{(x_1,x_2,x_3,x_4)\in C_T^4:
       x_j=0\pmod{\ell_T}\ \text{for some }j\}.
 \label{wb:cusp-cone-S24-exact-collapsed-locus}
\end{equation}
The maps
\(\rho_{12,T}\circ(c_T\times c_T):Q_T^2\to
S_T^2\wedge S_T^2\) and
\(\rho_{24,T}:C_T^4\to C_T^{\wedge4}\) send exactly \(W_{12,T}\) and
\(W_{24,T}\), respectively, to the smash basepoint and are injective on
their complements.  The full composites from \(F_T^2\) and \(F_T^4\)
still identify every pair of points identified by
\(\pi_T\times\pi_T\) and \(\chi_{\gamma,T}^{\times4}\), respectively; in
particular, their first arrows retain their nontrivial torus kernels.
Their exact smash-basepoint preimages in the full radial sources are
\[
 \bigl(\pi_T^{-1}(K_T^{(1)})\times F_T\bigr)
 \cup\bigl(F_T\times\pi_T^{-1}(K_T^{(1)})\bigr)
\]
and
\[
 \{(x_1,x_2,x_3,x_4)\in F_T^4:
   \chi_{\gamma,T}(x_j)=0\ \text{in }C_T\ \text{for some }j\},
\]
respectively.  Outside these basepoint preimages, two source points have the
same final image if and only if their images under the corresponding first
arrow in \cref{wb:cusp-cone-S12-S4-radial-map} and \cref{wb:cusp-cone-S24-S4-radial-map}
are equal.

Choose the Euclidean Riemannian \(\ell^2\)-product length metrics on
\(Q_T^2\) and \(C_T^4\) before taking the smash quotients, and use the
successive quotient length metrics specified in the proof.  With the
specified midpoint basepoints,
\begin{equation}
 (\mathbb S_{12,T}^4,\zeta_{12,T})
 \longrightarrow(\mathbb R^4,0),
 \qquad
 (\mathbb S_{24,T}^4,\zeta_{24,T})
 \longrightarrow(\mathbb R^4,0)
 \label{wb:cusp-cone-moving-S4-R4-limits}
\end{equation}
in pointed local-isometric and pointed Gromov--Hausdorff topology.  Moreover,
the convergence is smooth on the eventually Euclidean basepoint
neighborhoods.  For every \(R>0\), once \(T-C>4R\), respectively once
\(\ell_T>4R\), the corresponding quotient map restricts to an isometry from
the radius-\(R\) ball in \(Q_T^2\), respectively \(C_T^4\), \emph{onto} the
entire radius-\(R\) ball in the indicated smash sphere.  The domain ball, and
hence the quotient ball, is exactly a Euclidean four-ball.
\end{theorem}



\begin{proof}
\begin{enumerate}[label=\arabic*.,leftmargin=*]
\item Write \(\Sigma_r(Z,z_0)\) for the reduced suspension
\[
 \frac{[0,\epsilon]\times Z}
 {(\{0,\epsilon\}\times Z)\cup([0,\epsilon]\times\{z_0\})},
\]
where the displayed union is collapsed to one basepoint.

\item Use the marked presentation
\(\overline T^2=\mathbb R^2/\mathbb Z^2\), let
\[
 K^{(1)}=
 ((\mathbb R/\mathbb Z)\times\{0\})
 \cup(\{0\}\times(\mathbb R/\mathbb Z)),
\]
and let
\[
 c:\overline T^2\longrightarrow
 \overline T^2/K^{(1)}\cong S^2
\]
be the degree-one quotient obtained by collapsing the one-skeleton.  The
formula
\[
 \widehat p_2:\mathcal A\longrightarrow
 \Sigma_r(\overline T^2/K^{(1)},c(K^{(1)}))\cong S^3,
 \qquad
 [r,e^{i\phi},u]\longmapsto[r,c(u)]
\]
is well defined: the two suspension ends of \(\mathcal A\) and the entire
cylinder over \(K^{(1)}\) map to the reduced-suspension basepoint.  It is a
continuous surjection.  Set
\(\widehat q_3=\widehat p_2\circ q_2\).

\item For every sufficiently large \(T\), define the transported one-skeleton and
its collapse by
\begin{equation}
 \begin{aligned}
 K_T^{(1)}&:=\iota_T(K^{(1)})\subset Q_T,\\
 S_T^2&:=\overline T^2/K^{(1)}\cong Q_T/K_T^{(1)},\\
 c_T&:=c\circ\iota_T^{-1}:Q_T\longrightarrow S_T^2.
 \end{aligned}
 \label{wb:cusp-cone-transported-skeleton-collapse}
\end{equation}
The displayed target identification is the one induced by \(\iota_T\), so
\(c_T^{-1}(*)=K_T^{(1)}\), where
\(*=c(K^{(1)})=c_T(K_T^{(1)})\).

\item Give \(Q_T\) its Euclidean quotient length metric and \(S_T^2\) the quotient
length metric under \(c_T\).  Put
\[
 y_*=[(1/2,1/2)]\in\overline T^2,
 \qquad y_T:=\iota_T(y_*)\in Q_T,
 \qquad z_T:=c_T(y_T)\in S_T^2.
\]
Equivalently, the metric on \(\overline T^2\) used below is
\(g_T:=\iota_T^*g_{Q_T}\).  For every \(v\in\mathbb R^2\), the exact
decomposition
\[
 Y_T=T\begin{pmatrix}0&1\\-1&0\end{pmatrix}+E_T,
 \qquad\|E_T\|_{\mathrm{op}}\leq C,
\]
gives
\[
 \|Y_Tv\|\geq(T-C)\|v\|.
\]
The standard flat distance from \(y_*\) to \(K^{(1)}\) is \(1/2\);
therefore
\begin{equation}
 d_T:=\operatorname{dist}_{Q_T}(y_T,K_T^{(1)})
 =\operatorname{dist}_{g_T}(y_*,K^{(1)})
 \geq\frac{T-C}{2}.
 \label{wb:cusp-cone-distance-to-skeleton}
\end{equation}
For a length-space quotient which collapses a closed subset \(K\) to one
point, the distance between two noncollapsed classes represented by \(x,y\)
is
\[
 \min\{d(x,y),d(x,K)+d(y,K)\}.
\]
The two displayed terms are upper bounds, obtained respectively from a
direct path and from paths from \(x\) and \(y\) to \(K\).  Conversely, a
rectifiable quotient path either avoids the collapsed point, in which case
its compact image lifts to \(X\setminus K\) and has length at least
\(d(x,y)\), or meets that point, in which case its two pieces have total
length at least \(d(x,K)+d(y,K)\).  Taking infima proves the formula.
If \(T-C>4R\), then \(d_T>2R\).  The displayed formula and the systolic
bound show that
\[
 c_T:B_R^{Q_T}(y_T)\longrightarrow B_R^{S_T^2}(z_T)
\]
is an isometry onto the entire target ball, and both balls are exactly
Euclidean two-balls.  Indeed, the second term in the minimum cannot realize
a distance at most \(2R\), and any point of the target ball has a unique
noncollapsed representative; the same formula puts that representative in
the source ball.

\item At reduced-suspension coordinate \(r_T=e^{-2\pi T}\), \(S_T^2\) is a based
embedded subspace of \(S^3\): its point \(*\) is the global
reduced-suspension basepoint.  Hence
\[
 \mathbb S_{12,T}^4=S_T^2\wedge S_T^2
\]
embeds in \(S^3\wedge S^3=S^6\).  Equip \(Q_T^2\) first with the Euclidean
Riemannian \(\ell^2\)-product length metric, equip \((S_T^2)^2\) with the
\(\ell^2\)-product of the two quotient length metrics, and let
\[
 \rho_{12,T}:(S_T^2)^2\longrightarrow S_T^2\wedge S_T^2
\]
carry the quotient length metric obtained by collapsing the fat wedge
\[
 \mathcal W_{12,T}
 =(S_T^2\times\{*\})\cup(\{*\}\times S_T^2).
\]
Since \(c_T^{-1}(*)=K_T^{(1)}\), the inverse image of
\(\mathcal W_{12,T}\) under \(c_T\times c_T\) is exactly \(W_{12,T}\) in
\cref{wb:cusp-cone-S12-exact-collapsed-locus}.  On
\(Q_T^2\setminus W_{12,T}\), the map \(c_T\times c_T\) is injective and its
image avoids \(\mathcal W_{12,T}\); the map \(\rho_{12,T}\) is injective on
the complement of \(\mathcal W_{12,T}\).  Hence their composite identifies
all of \(W_{12,T}\) to the smash basepoint and makes no other
identifications.

\item The point \(\zeta_{12,T}=[z_T,z_T]\) has distance \(d_T\) from this fat
wedge.  The one-factor onto isometry above gives an onto isometry from the
radius-\(R\) ball about \((y_T,y_T)\) in \(Q_T^2\) to the corresponding
ball about \((z_T,z_T)\) in \((S_T^2)^2\).  Applying the same collapsed-set
distance formula to \(\rho_{12,T}\), and using \(d_T>2R\), shows that
\begin{equation}
 \rho_{12,T}\circ(c_T\times c_T):
 B_R^{Q_T^2}((y_T,y_T))
 \xrightarrow{\ \cong_{\mathrm{isom}}\ }
 B_R^{\mathbb S_{12,T}^4}(\zeta_{12,T})
 \label{wb:cusp-cone-S12-onto-exact-ball}
\end{equation}
is onto as well as isometrically embedded.  The systolic bound
\(\operatorname{sys}(Q_T^2)\geq T-C>4R\) makes its domain exactly a
Euclidean four-ball.  On a radial fibre,
\[
 (\widehat p_2\circ q_2)|_{F_T}
 =[r_T,c(\iota_T^{-1}(\pi_T(\,\cdot\,)))]
 =[r_T,c_T(\pi_T(\,\cdot\,))],
\]
and taking two factors followed by \(\rho_{12,T}\) proves the entire
factorization in \cref{wb:cusp-cone-S12-S4-radial-map}.

\item For the four-copy construction, choose \(1\in S^1\) as character basepoint
and define
\[
 \widehat p_1:\mathcal B\longrightarrow\Sigma_r(S^1,1)\cong S^2,
 \qquad
 [r,e^{i\phi},\theta]\longmapsto[r,\theta].
\]
The formula forgets the base angle and collapses the two suspension ends
and the entire cylinder with \(\theta=1\) to the reduced-suspension
basepoint.  It is a continuous surjection.  Put
\(\widehat q_\gamma=\widehat p_1\circ q_{1,\mathrm{cusp}}\).

\item At the interior suspension coordinate \(r_T=e^{-2\pi T}\), identify the
character circle \(C_T=\mathbb R/\ell_T\mathbb Z\) with
\(\mathbb R/\mathbb Z\subset S^2\) by \([x]\mapsto[x/\ell_T]\).  It is a
based embedded circle, so its fourth smash is an embedded
\[
 \mathbb S_{24,T}^4=C_T^{\wedge4}\subset(S^2)^{\wedge4}=S^8.
\]
Equip \(C_T^4\) first with the Euclidean Riemannian \(\ell^2\)-product
length metric, and let
\[
 \rho_{24,T}:C_T^4\longrightarrow C_T^{\wedge4}
\]
carry the quotient length metric obtained by collapsing exactly the fat
wedge \(W_{24,T}\) of
\cref{wb:cusp-cone-S24-exact-collapsed-locus}.  Away from that fat wedge the
smash quotient is injective, so no further points are identified.  The
midpoint
\[
 \zeta_{24,T}
 =[(\ell_T/2,\ell_T/2,\ell_T/2,\ell_T/2)]
\]
has distance \(\ell_T/2\) from that fat wedge.  If \(\ell_T>4R\), neither a
nonzero period nor a path through the collapsed fat wedge can shorten a
distance between points in the radius-\(R\) ball of \(C_T^4\).  Conversely,
the collapsed-set distance formula shows that every point of the
radius-\(R\) quotient ball has its unique noncollapsed representative in
that product ball.  Consequently
\begin{equation}
 \rho_{24,T}:
 B_R^{C_T^4}((\ell_T/2,\ell_T/2,\ell_T/2,\ell_T/2))
 \xrightarrow{\ \cong_{\mathrm{isom}}\ }
 B_R^{\mathbb S_{24,T}^4}(\zeta_{24,T})
 \label{wb:cusp-cone-S24-onto-exact-ball}
\end{equation}
is onto as well as isometrically embedded, and its domain is exactly a
Euclidean four-ball because \(\operatorname{sys}(C_T^4)=\ell_T>4R\).  On
the radial source the formula discards only the base angle and then performs
this smash quotient, giving
\cref{wb:cusp-cone-S24-S4-radial-map}.

\item Both one-copy maps are continuous surjections from compact \(X\) to
Hausdorff spheres, hence quotient maps, and they are based.  Their smash
powers are surjective; compact-to-Hausdorff again makes them quotient maps.
Using the same chosen based diffeomorphism \(X\cong_{\mathrm{diff}}S^6\),
the homeomorphisms
\[
 X^{\wedge2}\cong S^{12},\quad(S^3)^{\wedge2}\cong S^6,
 \qquad
 X^{\wedge4}\cong S^{24},\quad(S^2)^{\wedge4}\cong S^8
\]
prove \cref{wb:cusp-cone-sphere-to-sphere-maps}.  The exact Euclidean-ball
statements prove all three convergences in
\cref{wb:cusp-cone-moving-S4-R4-limits}.
\end{enumerate}
\end{proof}

\begin{corollary}[The two-copy rank-two sphere carrier is null-homotopic]
\label{wb:cusp-cone-S12-sphere-map-null}
\StatementRecord{wb:cusp-cone-S12-sphere-map-null}{proved}{Historical events: GPT-6 Astra (OpenAI), ChatGPT Pro mode, effort Pro, S6\_09\_PAPERS, 17 November 2026: targeted proof check. Also checked in this writing. Pre-continuation writing events: GPT-6 Astra (OpenAI), ChatGPT Pro mode, effort Pro, S6\_10\_PAPERS, 18 November 2026: targeted proof check. Continuation metadata events: GPT-6 Astra (OpenAI), ChatGPT Pro mode, effort Pro, S6\_10\_PAPERS continuation, 19 November 2026: used and rechecked in this run; status-vocabulary reconciliation; metadata reconciliation only; no new mathematical check. \emph{Pre-continuation recorded proof scope.} GPT-6 Astra (OpenAI, ChatGPT Pro mode), S6\_09\_PAPERS, effort Pro, 2026-11-17; also checked in this writing by GPT-6 Astra (OpenAI), ChatGPT Pro mode, effort Pro, 18 November 2026}{the hypotheses, notation and definitions stated here; \cref{wb:cusp-cone-sphere-to-sphere-maps}; \cref{wb:cusp-s12-full-carrier-null}}

After the based identifications in
\cref{wb:cusp-cone-sphere-to-sphere-maps}, the continuous surjective quotient
\[
 \widehat q_3^{\wedge2}:S^{12}\twoheadrightarrow S^6
\]
is based null-homotopic.  In fact, \(g\wedge g:S^{12}\to S^6\) is based
null-homotopic for every based map \(g:S^6\to S^3\).
\end{corollary}



\begin{proof}
\begin{enumerate}[label=\arabic*.,leftmargin=*]
\item Apply the full external-smash calculation in
\cref{wb:cusp-s12-full-carrier-null} to
\(g=\widehat q_3\).  That theorem proves the stronger assertion that
\(x\wedge y:S^{12}\to S^6\) is null-homotopic for every pair of based
maps \(x,y:S^6\to S^3\).
\end{enumerate}
\end{proof}

\begin{corollary}[Every repeated smash of the circle-quotient carrier is null]
\label{wb:cusp-cone-S24-sphere-map-null}
\StatementRecord{wb:cusp-cone-S24-sphere-map-null}{proved}{Historical events: GPT-6 Astra (OpenAI), ChatGPT Pro mode, effort Pro, S6\_09\_PAPERS, 17 November 2026: targeted proof check. Also checked in this writing. Pre-continuation writing events: GPT-6 Astra (OpenAI), ChatGPT Pro mode, effort Pro, S6\_10\_PAPERS, 18 November 2026: targeted proof check. Continuation metadata events: GPT-6 Astra (OpenAI), ChatGPT Pro mode, effort Pro, S6\_10\_PAPERS continuation, 19 November 2026: used and rechecked in this run; status-vocabulary reconciliation; metadata reconciliation only; no new mathematical check. \emph{Pre-continuation recorded proof scope.} GPT-6 Astra (OpenAI, ChatGPT Pro mode), S6\_09\_PAPERS, effort Pro, 2026-11-17; also checked in this writing by GPT-6 Astra (OpenAI), ChatGPT Pro mode, effort Pro, 18 November 2026}{the hypotheses, notation and definitions stated here; \cite[Corollary~4.24, p.~360]{Hatcher2002}; \cite[Table~1, p.~113]{IsaksenWangXu2023}; \cref{wb:cusp-cone-S24-smash-square-factorization}}

Put \(f=\widehat q_\gamma:S^6\to S^2\) after the chosen based
identification \(X\cong S^6\).  For every \(r\geq2\),
\begin{equation}
 \bigl[f^{\wedge r}\bigr]=0
 \quad\text{in}\quad
 \pi_{6r}(S^{2r}).
 \label{wb:cusp-cone-S24-null-maps}
\end{equation}
In particular,
\begin{equation}
 \bigl[\widehat q_\gamma^{\wedge2}\bigr]
 =0\in\pi_{12}(S^4),
 \qquad
 \bigl[\widehat q_\gamma^{\wedge4}\bigr]
 =0\in\pi_{24}(S^8).
 \label{wb:cusp-cone-S24-null-special-cases}
\end{equation}
Thus the displayed surjective quotient
\(\widehat q_\gamma^{\wedge4}:S^{24}\twoheadrightarrow S^8\) is globally
null-homotopic.  This does not alter either its point-set surjectivity or the
exact intrinsic \(\mathbb R^4\) limit on the moving four-spheres
\(\mathbb S_{24,T}^4\): those are properties of the specified representative
and its radial restrictions, whereas null-homotopy classifies that
representative only up to based homotopy.
\end{corollary}



\begin{proof}
\begin{enumerate}[label=\arabic*.,leftmargin=*]
\item Associativity of the smash product gives the
exact factorization
\begin{equation}
 \begin{tikzcd}[column sep=large]
 S^6\wedge S^6
 \arrow[r,"f\wedge\operatorname{id}_{S^6}"]&
 S^2\wedge S^6
 \arrow[r,"\operatorname{id}_{S^2}\wedge f"]&
 S^2\wedge S^2 .
 \end{tikzcd}
 \label{wb:cusp-cone-S24-smash-square-factorization}
\end{equation}
Under the standard based sphere identifications, the first arrow is the
sixfold suspension of \(f\), up to the harmless twist homeomorphisms dictated
by the left-versus-right suspension convention:
\begin{equation}
 f\wedge\operatorname{id}_{S^6}
 =\Sigma^6f:S^{12}\longrightarrow S^8.
 \label{wb:cusp-cone-S24-sixfold-suspension}
\end{equation}
Freudenthal's suspension theorem is already in its isomorphism range here,
because \(12\leq 2\cdot8-2\), and remains in that range after every further
suspension.  It therefore identifies \(\pi_{12}(S^8)\) with the fourth
stable stem.  The low-stem calculation gives
\phantomsection\label{qa:III:layout:11}
\begin{equation}
 \pi_{12}(S^8)\cong\pi_4^{\mathrm S}=0.
 \label{wb:cusp-cone-fourth-stable-stem}
\end{equation}
Here we use \cite[Corollary~4.24, p.~360]{Hatcher2002} and
\cite[Table~1, p.~113]{IsaksenWangXu2023}.
Hence \(\Sigma^6f\) is based null-homotopic, so the composite in
\cref{wb:cusp-cone-S24-smash-square-factorization}, namely
\(f\wedge f\), is based null-homotopic.  For \(r\geq2\), the canonical
associativity homeomorphisms identify
\begin{equation}
 f^{\wedge r}=(f\wedge f)\wedge f^{\wedge(r-2)},
\end{equation}
with the convention \(f^{\wedge0}=\operatorname{id}_{S^0}\).  Smashing a
based null-homotopy with any based map is a based null-homotopy.  This proves
\cref{wb:cusp-cone-S24-null-maps} and \cref{wb:cusp-cone-S24-null-special-cases}.
\end{enumerate}
\end{proof}

\begin{corollary}[Literal \(\mathbb R^4\) open cells in the sphere targets]
\label{wb:cusp-cone-literal-R4-open-cells}
\StatementRecord{wb:cusp-cone-literal-R4-open-cells}{proved}{Historical events: GPT-6 Astra (OpenAI), ChatGPT Pro mode, effort Pro, S6\_09\_PAPERS, 17 November 2026: targeted proof check. Also checked in this writing. Pre-continuation writing events: GPT-6 Astra (OpenAI), ChatGPT Pro mode, effort Pro, S6\_10\_PAPERS, 18 November 2026: targeted proof check. Continuation metadata events: GPT-6 Astra (OpenAI), ChatGPT Pro mode, effort Pro, S6\_10\_PAPERS continuation, 19 November 2026: used and rechecked in this run; status-vocabulary reconciliation; metadata reconciliation only; no new mathematical check. \emph{Pre-continuation recorded proof scope.} GPT-6 Astra (OpenAI, ChatGPT Pro mode), S6\_09\_PAPERS, effort Pro, 2026-11-17; also checked in this writing by GPT-6 Astra (OpenAI), ChatGPT Pro mode, effort Pro, 18 November 2026}{the hypotheses, notation and definitions stated here; \cref{wb:cusp-cone-S12-onto-exact-ball}; \cref{wb:cusp-cone-S24-onto-exact-ball}}

Let \(*\) denote the smash basepoint and set
\[
 U_{12,T}:=\mathbb S_{12,T}^4\setminus\{*\},
 \qquad
 U_{24,T}:=\mathbb S_{24,T}^4\setminus\{*\}.
\]
The embedded moving four-spheres restrict to topological embeddings
\[
 U_{12,T}\cong\mathbb R^4\hookrightarrow S^6,
 \qquad
 U_{24,T}\cong\mathbb R^4\hookrightarrow S^8.
\]
More exactly, the quotient-stage maps restrict to homeomorphisms
\begin{align}
 (Q_T\setminus K_T^{(1)})^2
 &\xrightarrow{\ \cong\ }U_{12,T},
 \label{wb:cusp-cone-U12-homeomorphism}\\
 (C_T\setminus\{0\})^4
 &\xrightarrow{\ \cong\ }U_{24,T}.
 \label{wb:cusp-cone-U24-homeomorphism}
\end{align}
These copies of \(\mathbb R^4\) are nonclosed topological open cells: they
are open in their moving four-spheres, not open in the ambient \(S^6\) or
\(S^8\), and their closures in those ambient spheres are
\(\mathbb S_{12,T}^4\) and \(\mathbb S_{24,T}^4\), respectively.  They are
not globally isometric to Euclidean \(\mathbb R^4\) for either the assigned
compact quotient metrics or the restriction of any fixed compatible metric
on the compact sphere targets.  The exact midpoint balls in
\cref{wb:cusp-cone-S12-onto-exact-ball} and \cref{wb:cusp-cone-S24-onto-exact-ball} give the
precise metric statement: as \(T\to\infty\), every fixed Euclidean
radius-\(R\) ball occurs isometrically and onto the corresponding pointed
ball.
\end{corollary}



\begin{proof}
\begin{enumerate}[label=\arabic*.,leftmargin=*]
\item For a quotient which collapses a closed subset \(K\) to a point, restriction
to the complement gives a homeomorphism
\[
 X\setminus K\xrightarrow{\ \cong\ }(X/K)\setminus\{*\}.
\]
Indeed, the complement of the quotient point is open, its inverse image is
the saturated open set \(X\setminus K\), and restriction of a quotient map
over an open subset of its target is again a quotient map.  The restricted
map is bijective, and a bijective quotient map is a homeomorphism.
Consequently
\[
 Q_T\setminus K_T^{(1)}
 \cong S_T^2\setminus\{*\}.
\]
The marked coordinate one-skeleton \(K^{(1)}\) is the \(1\)-skeleton of the
standard one-vertex CW structure on \(\overline T^2\), with one \(0\)-cell,
two \(1\)-cells, and one \(2\)-cell.  Its complement is therefore the open
\(2\)-cell \((0,1)^2\cong\mathbb R^2\).  Transport by
\(\iota_T\) gives
\[
 Q_T\setminus K_T^{(1)}\cong\mathbb R^2.
\]
Likewise \(C_T\setminus\{0\}\) is an open interval and hence is homeomorphic
to \(\mathbb R\).

\item The complement of the smash basepoint is the product of the complements of
the factor basepoints.  Therefore
\[
 (S_T^2\wedge S_T^2)\setminus\{*\}
 \cong(S_T^2\setminus\{*\})^2\cong\mathbb R^4,
\]
and
\[
 C_T^{\wedge4}\setminus\{*\}
 \cong(C_T\setminus\{0\})^4\cong\mathbb R^4.
\]
The full moving four-spheres are topologically embedded in their compact
sphere targets, so these restrictions are topological embeddings.  In each
case the omitted smash basepoint lies in the closure of the open cell, which
proves the closure assertions. These cells are not ambient open: an ambient open set in $S^6$ or $S^8$ has local integral homology concentrated in degree $6$ or $8$, respectively, whereas the displayed homeomorphisms from $\mathbb R^4$ give local homology concentrated in degree $4$. Finally, every compatible metric on a
compact target is bounded, as is each compact quotient metric assigned to a
moving four-sphere, whereas Euclidean \(\mathbb R^4\) is unbounded.  Thus no
global isometry is possible; the exact metric content is precisely the
onto-ball calculation already proved.
\end{enumerate}
\end{proof}


\paragraph{The moving four-spheres are null-homotopic in the targets}
\label{wb:cusp-cone-moving-S4-null-homotopy}
\StatementRecord{wb:cusp-cone-moving-S4-null-homotopy}{proved}{Historical events: GPT-6 Astra (OpenAI), ChatGPT Pro mode, effort Pro, S6\_09\_PAPERS, 17 November 2026: targeted proof check. Also checked in this writing. Pre-continuation writing events: GPT-6 Astra (OpenAI), ChatGPT Pro mode, effort Pro, S6\_10\_PAPERS, 18 November 2026: targeted proof check. Continuation metadata events: GPT-6 Astra (OpenAI), ChatGPT Pro mode, effort Pro, S6\_10\_PAPERS continuation, 19 November 2026: used and rechecked in this run; status-vocabulary reconciliation; metadata reconciliation only; no new mathematical check. \emph{Pre-continuation recorded proof scope.} GPT-6 Astra (OpenAI, ChatGPT Pro mode), S6\_09\_PAPERS, effort Pro, 2026-11-17; also checked in this writing by GPT-6 Astra (OpenAI), ChatGPT Pro mode, effort Pro, 18 November 2026}{the hypotheses, notation and definitions stated here; the argument and qualifications immediately following}

For a fixed interior reduced-suspension coordinate, the inclusion
\(S_T^2\hookrightarrow S^3\) contracts through based maps by sliding that
coordinate to a suspension endpoint; throughout the homotopy the collapsed
meridian remains at the global basepoint.  Hence
\(S_T^2\wedge S_T^2\hookrightarrow S^6\) is null-homotopic.  The identical
sliding argument contracts \(C_T\hookrightarrow S^2\), and therefore also
contracts \(C_T^{\wedge4}\hookrightarrow S^8\).  Thus the moving metric
four-spheres do not contradict the fixed-target orientation obstruction:
their fundamental classes die in the ambient targets.






\paragraph{Homotopical status of the two sphere-to-sphere maps}
\label{wb:cusp-cone-sphere-map-homotopy-status}
\StatementRecord{wb:cusp-cone-sphere-map-homotopy-status}{proved}{Historical events: GPT-6 Astra (OpenAI), ChatGPT Pro mode, effort Pro, S6\_09\_PAPERS, 17 November 2026: targeted proof check. Also checked in this writing. Pre-continuation writing events: GPT-6 Astra (OpenAI), ChatGPT Pro mode, effort Pro, S6\_10\_PAPERS, 18 November 2026: targeted proof check. Continuation metadata events: GPT-6 Astra (OpenAI), ChatGPT Pro mode, effort Pro, S6\_10\_PAPERS continuation, 19 November 2026: used and rechecked in this run; status-vocabulary reconciliation; metadata reconciliation only; no new mathematical check. \emph{Pre-continuation recorded proof scope.} GPT-6 Astra (OpenAI, ChatGPT Pro mode), S6\_09\_PAPERS, effort Pro, 2026-11-17; also checked in this writing by GPT-6 Astra (OpenAI), ChatGPT Pro mode, effort Pro, 18 November 2026}{the hypotheses, notation and definitions stated here; \cref{wb:cusp-cone-S24-sphere-map-null}; \cref{wb:cusp-cone-S12-sphere-map-null}; \cref{wb:cusp-null:thm}; \cref{mapsIII:angle-forgotten-one-copy-null}}

The maps \(S^{12}\to S^6\) and \(S^{24}\to S^8\) induce zero on all
positive-degree integral homology and cohomology groups and have zero
ordinary Hurewicz image.  They have no classical degree because source and
target dimensions differ.  The full four-copy map
\(\widehat q_\gamma^{\wedge4}:S^{24}\to S^8\) is based null-homotopic by
\cref{wb:cusp-cone-S24-sphere-map-null}.  The distinct two-copy map
\(\widehat q_3^{\wedge2}:S^{12}\to S^6\) is also based null-homotopic, by
\cref{wb:cusp-cone-S12-sphere-map-null}.  In both cases the vanishing is a
theorem about the full specified sphere-valued map, not merely its Hurewicz
image or the inclusion of a moving four-sphere. % checked carrier correction: one-copy-status
The one-copy class
\([\widehat q_3]\in\pi_6(S^3)\) is zero by the complete framed-torus
proof in \cref{wb:cusp-null:thm}.  Both this class and
\([\widehat q_\gamma]\in\pi_6(S^2)\) also vanish by the extension
homotopy in \cref{mapsIII:angle-forgotten-one-copy-null}.  That proof
uses the original formulas after forgetting the base angle and fixes
the chosen point of the central fibre throughout.  It supplies based
nullhomotopies of all their positive smash powers.  The retained-angle
map has a different target and is outside that assertion.





\begin{proposition}[The compact cone-on-\(T^4\) carriers]
\label{wb:cusp-cone-explicit-carriers}
\StatementRecord{wb:cusp-cone-explicit-carriers}{proved}{Historical events: GPT-6 Astra (OpenAI), ChatGPT Pro mode, effort Pro, S6\_09\_PAPERS, 17 November 2026: targeted proof check. Also checked in this writing. Pre-continuation writing events: GPT-6 Astra (OpenAI), ChatGPT Pro mode, effort Pro, S6\_10\_PAPERS, 18 November 2026: targeted proof check. Continuation metadata events: GPT-6 Astra (OpenAI), ChatGPT Pro mode, effort Pro, S6\_10\_PAPERS continuation, 19 November 2026: used and rechecked in this run; status-vocabulary reconciliation; metadata reconciliation only; no new mathematical check. \emph{Pre-continuation recorded proof scope.} GPT-6 Astra (OpenAI, ChatGPT Pro mode), S6\_09\_PAPERS, effort Pro, 2026-11-17; also checked in this writing by GPT-6 Astra (OpenAI), ChatGPT Pro mode, effort Pro, 18 November 2026}{the hypotheses, notation and definitions stated here; the explicit constructions and calculations in the proof}

Fix \(T_0\) sufficiently large and put
\(r_0=e^{-2\pi T_0}<\epsilon\).  Define
\[
 \mathcal D_{12}
 =\frac{[0,r_0]\times(\overline T^2)^2}
 {\{0\}\times(\overline T^2)^2\sim d_{12}},
 \qquad
 \mathcal D_{24}
 =\frac{[0,r_0]\times(S^1)^4}
 {\{0\}\times(S^1)^4\sim d_{24}}.
\]
Both spaces are homeomorphic to \(\operatorname{Cone}(T^4)\).  There are
topological embeddings
\begin{align}
 \delta_{12}:\mathcal D_{12}&\hookrightarrow
 \mathcal A\wedge\mathcal A,
 &
 [r,u_1,u_2]&\longmapsto
 [\,[r,1,u_1],[r,1,u_2],],
 \label{wb:cusp-cone-D12-embedding}\\
 \delta_{24}:\mathcal D_{24}&\hookrightarrow
 \mathcal C^{\wedge4},
 &
 [r,\theta_1,\ldots,\theta_4]&\longmapsto
 \bigwedge_{j=1}^4[r,(1,\theta_j)].
 \label{wb:cusp-cone-D24-embedding}
\end{align}
Here \(1\in S^1\) fixes the positive radial base angle.  For every
\(T\geq T_0\), put \(r_T=e^{-2\pi T}\).  The restrictions of the sphere
quotient maps to the synchronized radial fibres are exactly the
cross-sections at height \(r_T\):
\begin{align}
 F_T^2&\xrightarrow{\ \pi_T\times\pi_T\ }
 Q_T^2\cong\{r_T\}\times T^4
 \subset\mathcal D_{12},
 \label{wb:cusp-cone-D12-slice}\\
 F_T^4&\xrightarrow{\ \chi_{\gamma,T}^{\times4}\ }
 C_T^4\cong\{r_T\}\times T^4
 \subset\mathcal D_{24}.
 \label{wb:cusp-cone-D24-slice}
\end{align}
Every one of these embedded four-torus slices is null-homotopic in its
carrier, and hence in the full smash target.

At the cone vertex \(d_\nu\), \(\nu\in\{12,24\}\), the local integral
homology is
\begin{equation}
 H_k(\mathcal D_\nu,
       \mathcal D_\nu\setminus\{d_\nu\};\mathbb Z)
 \cong
 \begin{cases}
  \mathbb Z^4,&k=2,\\
  \mathbb Z^6,&k=3,\\
  \mathbb Z^4,&k=4,\\
  \mathbb Z,&k=5,\\
  0,&\text{otherwise}.
 \end{cases}
 \label{wb:cusp-cone-carrier-local-homology}
\end{equation}
In particular, \(d_\nu\) is not an integral homology-manifold point, with or
without boundary, of any dimension.  Consequently it is not a topological
manifold point, with or without boundary, of any dimension.
\end{proposition}



\begin{proof}
\begin{enumerate}[label=\arabic*.,leftmargin=*]
\item The two displayed domains are cones on
\((\overline T^2)^2\cong T^4\) and \((S^1)^4\cong T^4\).  At \(r=0\),
every factor on the right of
\cref{wb:cusp-cone-D12-embedding} and \cref{wb:cusp-cone-D24-embedding} is the relevant
smash basepoint, so the formulae descend from the cylinders to the cones.
For \(0<r\leq r_0\), none of the factors is the smash basepoint.  Equality
in the smash quotient is then equality in every factor, and the coordinate
descriptions of the suspension and cusp cone give equality of \(r\) and of
all torus coordinates.  Thus each \(\delta_\nu\) is injective.  Its domain
is compact and its target is Hausdorff, so it is an embedding.

\item The already proved one-factor restriction formulae, with
\(r_T=e^{-2\pi T}\), are
\[
 q_2(x)=[r_T,1,\iota_T^{-1}(\pi_T(x))]
 \quad(x\in F_T)
\]
and, in the cusp-cone piece of \(\mathcal C\),
\[
 q_1(x)
 =[r_T,(1,\vartheta_T^{-1}(\chi_{\gamma,T}(x)))]
 \quad(x\in F_T).
\]
Applying these identities in two and four factors gives
\cref{wb:cusp-cone-D12-slice} and \cref{wb:cusp-cone-D24-slice}.  Sliding the cone coordinate
from \(r_T\) to zero contracts each slice to the cone vertex.  Finally,
excision and the long exact sequence of the pair give, for \(L=T^4\),
\[
 H_k(\operatorname{Cone}(L),
     \operatorname{Cone}(L)\setminus\{v\};\mathbb Z)
 \cong\widetilde H_{k-1}(L;\mathbb Z).
\]
The ranks \(4,6,4,1\) of \(H_1,H_2,H_3,H_4\) of \(T^4\) prove
\cref{wb:cusp-cone-carrier-local-homology}.  At an interior point of an
integral homology \(n\)-manifold the local integral homology is \(\mathbb Z\)
in degree \(n\) and zero in every other degree, whereas at a boundary point
of an integral homology manifold with boundary all these local groups
vanish.  The simultaneous nonzero groups in degrees \(2,3,4,5\), including
ranks greater than one, agree with neither pattern for any \(n\).  Every
topological manifold is locally an integral homology manifold, so the final
assertion follows as well.
\end{enumerate}
\end{proof}

\begin{theorem}[Point-set quotient maps and exact homotopy types]
\label{wb:cusp-cone-target-homotopy-types}
\StatementRecord{wb:cusp-cone-target-homotopy-types}{proved}{Historical events: GPT-6 Astra (OpenAI), ChatGPT Pro mode, effort Pro, S6\_09\_PAPERS, 17 November 2026: targeted proof check. Also checked in this writing. Pre-continuation writing events: GPT-6 Astra (OpenAI), ChatGPT Pro mode, effort Pro, S6\_10\_PAPERS, 18 November 2026: targeted proof check. Continuation metadata events: GPT-6 Astra (OpenAI), ChatGPT Pro mode, effort Pro, S6\_10\_PAPERS continuation, 19 November 2026: used and rechecked in this run; status-vocabulary reconciliation; metadata reconciliation only; no new mathematical check. \emph{Pre-continuation recorded proof scope.} GPT-6 Astra (OpenAI, ChatGPT Pro mode), S6\_09\_PAPERS, effort Pro, 2026-11-17; also checked in this writing by GPT-6 Astra (OpenAI), ChatGPT Pro mode, effort Pro, 18 November 2026}{the hypotheses, notation and definitions stated here; the explicit constructions and calculations in the proof}

The smash maps are continuous surjective quotient maps
\[
 q_2^{\wedge2}:S^{12}\twoheadrightarrow\mathcal A^{\wedge2},
 \qquad
 q_1^{\wedge4}:S^{24}\twoheadrightarrow\mathcal C^{\wedge4}.
\]
Their targets have homotopy types
\begin{align}
 \mathcal A^{\wedge2}
 &\simeq
 \bigvee^9S^4\vee\bigvee^{18}S^5\vee
 \bigvee^{15}S^6\vee\bigvee^6S^7\vee S^8,
 \label{wb:cusp-cone-A-smash-wedge}\\
 \mathcal C^{\wedge4}
 &\simeq
 S^8\vee\bigvee^4S^9\vee\bigvee^6S^{10}
 \vee\bigvee^4S^{11}\vee S^{12}.
 \label{wb:cusp-cone-C-smash-wedge}
\end{align}
Both maps induce zero on every positive-degree integral homology and
cohomology group and have zero ordinary Hurewicz image.  Neither map has a
continuous homotopy section.  These conclusions do not determine whether
the unstable homotopy class of either map is zero.
\end{theorem}



\begin{proof}
\begin{enumerate}[label=\arabic*.,leftmargin=*]
\item Let \(f:(K,k_0)\to(L,\ell_0)\) be a surjective based map between compact
Hausdorff spaces.  Every non-basepoint class of \(L^{\wedge r}\) has a
representative \((y_1,\ldots,y_r)\) with each \(y_i\ne\ell_0\).  Choosing
\(x_i\in f^{-1}(y_i)\) proves surjectivity of \(f^{\wedge r}\); the smash
basepoint is also in its image.  The source smash is compact.  The target
fat wedge is closed, hence the target smash is Hausdorff.  A continuous
surjection from a compact space to a Hausdorff space is a quotient map.
This proves the first assertion for \(q_2\) and \(q_1\).

\item The suspension splitting
\[
 \Sigma(X_1\times\cdots\times X_n)
 \simeq
 \bigvee_{\varnothing\ne I\subseteq\{1,\ldots,n\}}
 \Sigma\!\left(\bigwedge_{i\in I}X_i\right)
\]
gives
\[
 \mathcal A=\Sigma T^3
 \simeq\bigvee^3S^2\vee\bigvee^3S^3\vee S^4.
\]
Finite smash products distribute over finite wedges.  The ordered degree
sums give multiplicities \(9,18,15,6,1\) in dimensions
\(4,5,6,7,8\), proving \cref{wb:cusp-cone-A-smash-wedge}.

\item For \(\mathcal C\), write \(x=\rho_1\), \(y=\rho_2\), and
\(h=\eta\).  The pair-of-pants product and the two exact Seifert fillings
give
\[
 \pi_1(\mathcal M)=
 \langle x,y,h\mid[x,h]=[y,h]=1,
 x^3h^{-1}=1,\ y^4h=1\rangle.
\]
The remaining boundary torus is generated by \(h\) and
\(\rho_0=(xy)^{-1}\).  Coning this torus kills their normal closure, and
therefore
\[
 \pi_1(\mathcal C)
 \cong\langle x,y\mid x^3=1,\ y^4=1,\ xy=1\rangle=1,
\]
because \(y=x^{-1}\) and \(\gcd(3,4)=1\).  The previously computed
relative homology gives
\[
 \widetilde H_2(\mathcal C;\mathbb Z)=\mathbb Z,
 \qquad
 \widetilde H_3(\mathcal C;\mathbb Z)=\mathbb Z,
\]
and no other reduced homology.  % checked carrier correction: relative-hurewicz
Since \(\mathcal C\) is simply connected, Hurewicz gives a based
map \(a:S^2\to\mathcal C\) inducing an isomorphism on \(H_2\)
and on \(\pi_2\).  Replace this map by its mapping-cylinder
inclusion \(S^2\hookrightarrow M_a\), with
\(M_a\simeq\mathcal C\).  Both spaces are simply connected and
the map on \(\pi_2\) is an isomorphism, so the pair is
\(2\)-connected.  Relative Hurewicz gives
\[
 \pi_3(M_a,S^2)\cong H_3(M_a,S^2;\mathbb Z).
\]
The homology exact sequence identifies the right side with
\(H_3(M_a;\mathbb Z)=\mathbb Z\): the preceding group
\(H_3(S^2)\) is zero, and the following map on \(H_2\) is an
isomorphism.  The homotopy exact sequence makes
\(\pi_3(M_a)\to\pi_3(M_a,S^2)\) surjective, since its next
map has image in the zero kernel of
\(\pi_2(S^2)\to\pi_2(M_a)\).  Naturality of Hurewicz therefore
gives a based map \(b:S^3\to\mathcal C\) representing an integral
generator of \(H_3(\mathcal C)\).  The wedge \(a\vee b\) is an
integral homology equivalence in every degree.  The homological
Whitehead theorem, with its complete proof in
\cref{mapsI:pf:theorem:5.1}, makes this map a homotopy equivalence.  Thus
\[
 \mathcal C\simeq S^2\vee S^3.
\]
Taking the fourth smash power yields \cref{wb:cusp-cone-C-smash-wedge} with
the binomial multiplicities \(1,4,6,4,1\).

\item The first target has no reduced homology above degree eight and the second
has none above degree twelve.  The source spheres have no positive homology
below degrees twelve and twenty-four.  Hence all positive-degree homology
and cohomology homomorphisms induced by the two maps vanish.  Their
Hurewicz images lie in \(H_{12}(\mathcal A^{\wedge2})=0\) and
\(H_{24}(\mathcal C^{\wedge4})=0\), respectively.  A section would split a
nonzero lower homology group of the target into the zero lower homology of
the source sphere, which is impossible.  Higher homotopy groups of the
displayed wedges need not vanish, so none of these arguments decides the
unstable null-homotopy question.
\end{enumerate}
\end{proof}

\begin{theorem}[Exact decompactification spaces with boundary
\(\mathbb R^4\)]
\label{wb:cusp-cone-exact-decompactification-space}
\StatementRecord{wb:cusp-cone-exact-decompactification-space}{proved}{Historical events: GPT-6 Astra (OpenAI), ChatGPT Pro mode, effort Pro, S6\_09\_PAPERS, 17 November 2026: targeted proof check. Also checked in this writing. Pre-continuation writing events: GPT-6 Astra (OpenAI), ChatGPT Pro mode, effort Pro, S6\_10\_PAPERS, 18 November 2026: targeted proof check. Continuation metadata events: GPT-6 Astra (OpenAI), ChatGPT Pro mode, effort Pro, S6\_10\_PAPERS continuation, 19 November 2026: used and rechecked in this run; status-vocabulary reconciliation; metadata reconciliation only; no new mathematical check. \emph{Pre-continuation recorded proof scope.} GPT-6 Astra (OpenAI, ChatGPT Pro mode), S6\_09\_PAPERS, effort Pro, 2026-11-17; also checked in this writing by GPT-6 Astra (OpenAI), ChatGPT Pro mode, effort Pro, 18 November 2026}{the hypotheses, notation and definitions stated here; the explicit constructions and calculations in the proof}

For \(\nu\in\{12,24\}\), set
\[
 A_T^{(12)}=\operatorname{diag}(Y_T,Y_T),
 \qquad
 A_T^{(24)}=\ell_T I_4,
 \qquad s_0=T_0^{-1}.
\]
On \([0,s_0]\times\mathbb R^4\), define an equivalence relation
\(\sim_\nu\) fibrewise over \(s\) by
\begin{align*}
 (0,x)\sim_\nu(0,y)&\iff x=y,\\
 (s,x)\sim_\nu(s,y)&\iff
 x-y\in A_{1/s}^{(\nu)}\mathbb Z^4
 \quad(0<s\leq s_0).
\end{align*}
Let
\[
 \mathfrak E_\nu
 =([0,s_0]\times\mathbb R^4)/\!\sim_\nu,
 \qquad p_\nu:\mathfrak E_\nu\longrightarrow[0,s_0].
\]
Then \(\mathfrak E_\nu\) is a Hausdorff, second-countable topological
five-manifold with boundary, and
\begin{align}
 p_{12}^{-1}(s)&\cong Q_{1/s}^2 &(s>0),&
 &p_{12}^{-1}(0)&\cong\mathbb R^4,
 \label{wb:cusp-cone-E12-fibres}\\
 p_{24}^{-1}(s)&\cong C_{1/s}^4 &(s>0),&
 &p_{24}^{-1}(0)&\cong\mathbb R^4.
 \label{wb:cusp-cone-E24-fibres}
\end{align}
The other boundary component is the compact four-torus over \(s=s_0\).

Give every positive fibre its Euclidean quotient metric and the zero fibre
its Euclidean metric.  For every \(R>0\), there is \(s_R>0\) such that,
for \(0\leq s<s_R\), writing
\(q_{\nu,s}(x)=[s,x]\), the quotient map restricts to an exact isometry
\[
 q_{\nu,s}:B_R^{\mathbb R^4}(0)
 \xrightarrow{\ \cong_{\mathrm{isom}}\ }
 B_R^{p_\nu^{-1}(s)}([s,0])
\]
onto the entire quotient radius-\(R\) ball, not merely onto an embedded
subset of that ball.
Consequently the positive fibres converge to the zero fibre in pointed
local-isometric, pointed smooth, and pointed Gromov--Hausdorff topology.

There are continuous surjective comparison maps which set-theoretically
collapse the zero boundary fibres,
\[
 \begin{tikzcd}[column sep=large]
 \mathbb R^4=p_{12}^{-1}(0)
 \arrow[r,hook]&
 \mathfrak E_{12}
 \arrow[r,"\beta_{12}"]&
 \mathcal D_{12}
 \arrow[r,hook,"\delta_{12}"]&
 \mathcal A^{\wedge2},
 \end{tikzcd}
\]
and
\[
 \begin{tikzcd}[column sep=large]
 \mathbb R^4=p_{24}^{-1}(0)
 \arrow[r,hook]&
 \mathfrak E_{24}
 \arrow[r,"\beta_{24}"]&
 \mathcal D_{24}
 \arrow[r,hook,"\delta_{24}"]&
 \mathcal C^{\wedge4}.
 \end{tikzcd}
\]
For \(s>0\), \(\beta_\nu\) is the specified marked homeomorphism from
\(p_\nu^{-1}(s)\) to the radial torus slice at \(r=e^{-2\pi/s}\).  Its
restriction
\[
 \beta_\nu:\mathfrak E_\nu\setminus p_\nu^{-1}(0)
 \xrightarrow{\ \cong\ }\mathcal D_\nu\setminus\{d_\nu\}
\]
is a homeomorphism, while
\[
 \beta_\nu^{-1}(d_\nu)=p_\nu^{-1}(0)\cong\mathbb R^4.
\]
Neither \(\beta_\nu\) is proper or a quotient map.  It factors through a
continuous bijection
\[
 \overline\beta_\nu:
 \mathfrak E_\nu/p_\nu^{-1}(0)\longrightarrow\mathcal D_\nu,
\]
but this bijection is not a homeomorphism; equivalently, the topology of
\(\mathcal D_\nu\) is not the quotient topology induced by \(\beta_\nu\).
Thus the limiting \(\mathbb R^4\) is an actual boundary fibre of
\(\mathfrak E_\nu\), whereas the continuous comparison map sends that entire
fibre to the cone vertex.
\end{theorem}



\begin{proof}
\begin{enumerate}[label=\arabic*.,leftmargin=*]
\item We give a common proof.  Write
\[
 L_s=A_{1/s}^{(\nu)}\mathbb Z^4\quad(s>0),
 \qquad L_0=\{0\}.
\]
The matrices vary continuously for \(s>0\), and
\[
 \operatorname{sys}(L_s)\longrightarrow\infty
 \qquad(s\downarrow0)
\]
by \(\operatorname{sys}(Y_T\mathbb Z^2)\geq T-C\) in the
\(S^{12}\) case and \(\operatorname{sys}(\ell_T\mathbb Z^4)=\ell_T\)
in the \(S^{24}\) case.

\item The equivalence relation is closed.  Indeed, suppose
\((s_j,x_j)\sim_\nu(s_j,y_j)\) and
\((s_j,x_j,y_j)\to(s,x,y)\).  If \(s>0\), write
\(y_j-x_j=A_{1/s_j}^{(\nu)}n_j\).  On a compact neighborhood of \(s\),
the inverse matrices are uniformly bounded, so boundedness of
\(y_j-x_j\) forces \(n_j\) to lie in a finite set.  Passing to a
subsequence makes \(n_j\) constant and proves \(x\sim_\nu y\).  If
\(s=0\), any nonzero \(n_j\) would give
\[
 \|y_j-x_j\|\geq\operatorname{sys}(L_{s_j})\longrightarrow\infty,
\]
contrary to convergence.  Hence \(x=y\).

\item The quotient projection \(q_\nu:[0,s_0]\times\mathbb R^4\to
\mathfrak E_\nu\) is open.  If \(U\) is open, its saturation is the union
of \(U\) with the images, for \(n\in\mathbb Z^4\), of
\[
 U\cap((0,s_0]\times\mathbb R^4)
 \quad\text{under}\quad
 (s,x)\longmapsto(s,x+A_{1/s}^{(\nu)}n),
\]
and every one of these images is open.  A closed equivalence relation with
an open quotient map has Hausdorff quotient: two inequivalent points admit
a product neighborhood disjoint from the relation, and the saturations of
the two factor neighborhoods are disjoint open sets.  Since the source is
second countable and the quotient map is open, images of a countable base
form a countable base of the quotient.

\item Fix \(x_0\in\mathbb R^4\).  Choose \(a>0\) and \(s_1>0\) such that
\(\operatorname{sys}(L_s)>2a\) for \(0<s<s_1\).  The quotient projection
is one-to-one on
\([0,s_1)\times B_a^\circ(x_0)\), where
\(B_a^\circ(x_0)=\{x:\|x-x_0\|<a\}\), and its restriction is an open continuous
bijection onto an open subset of \(\mathfrak E_\nu\); it is therefore a
homeomorphism.  These are half-space charts at \(s=0\).  The usual evenly
covered small balls give product charts for \(s>0\), and the same argument
at \(s=s_0\) gives the second boundary component.  This proves the
manifold assertion and the fibre identifications.

\item For the metric assertion, choose \(s_R\) so that
\(\operatorname{sys}(L_s)>4R\) for \(s<s_R\).  If
\(x,y\in B_R(0)\) and \(0\ne\lambda\in L_s\), then
\[
 \|x-y+\lambda\|
 \geq\|\lambda\|-\|x-y\|
 >4R-2R\geq\|x-y\|.
\]
Thus no translated representative shortens the Euclidean distance on that
ball, so \(q_{\nu,s}\) is an exact isometric embedding there.  It is also
onto the full quotient ball.  Indeed, if \(\xi=[s,x]\) satisfies
\(d([s,0],\xi)\leq R\), discreteness of \(L_s\) makes the infimum in
\[
 d([s,0],[s,x])=\min_{\lambda\in L_s}\|x+\lambda\|
\]
attained.  For a minimizing \(\lambda_0\), the representative
\(x+\lambda_0\) lies in \(B_R^{\mathbb R^4}(0)\) and maps to \(\xi\).  At
\(s=0\), \(q_{\nu,0}\) is the identity.  This proves the asserted onto
isometry, and since it holds for every fixed \(R\), it proves all three
pointed convergences.

\item It remains to define these comparison maps, which set-theoretically collapse
the zero fibres.  For \(s>0\), put
\(T=1/s\) and \(r=e^{-2\pi T}\).  In the \(S^{12}\) case, write
\(x=(x_1,x_2)\in\mathbb R^2\oplus\mathbb R^2\) and set
\[
 \beta_{12}([s,x])
 =\bigl[r,
 \iota_T^{-1}([x_1]_{Y_T\mathbb Z^2}),
 \iota_T^{-1}([x_2]_{Y_T\mathbb Z^2})\bigr].
\]
In the \(S^{24}\) case, write
\(x=(x_1,\ldots,x_4)\) and set
\[
 \beta_{24}([s,x])
 =\bigl[r,
 [x_1/\ell_T],\ldots,[x_4/\ell_T]\bigr].
\]
These formulae are well defined on the positive fibres and are continuous
there.  Define \(\beta_\nu([0,x])=d_\nu\) for every \(x\in\mathbb R^4\).
Every neighborhood of \(d_\nu\) in the cone contains all cross-sections
with sufficiently small \(r\).  Since \(e^{-2\pi/s}\to0\) as
\(s\downarrow0\), continuity at the zero fibre is uniform in \(x\).
Surjectivity follows because each positive fibre maps onto its complete
cross-section and the zero fibre maps onto the vertex.  The restriction away
from \(p_\nu^{-1}(0)\) is a homeomorphism because
\(s\mapsto e^{-2\pi/s}\) is a homeomorphism from \((0,s_0]\) to
\((0,r_0]\), and the displayed marked-torus identifications are fibrewise
homeomorphisms varying continuously with \(s\).  The vertex fibre is exactly
\(p_\nu^{-1}(0)\), so \(\beta_\nu\) is not proper.

\item It is also not a quotient map.  Choose
\(a\in\mathbb R^4\setminus\mathbb Z^4\), for example
\(a=(1/2,1/2,1/2,1/2)\), choose \(s_j\downarrow0\), and put
\[
 z_j=[s_j,A_{1/s_j}^{(\nu)}a]\in\mathfrak E_\nu.
\]
For \(\nu=12\),
\[
 \inf_{n\in\mathbb Z^4}
 \|A_{1/s_j}^{(12)}(a+n)\|
 \geq(1/s_j-C)\operatorname{dist}(a,\mathbb Z^4)
 \longrightarrow\infty,
\]
and for \(\nu=24\) the same infimum equals
\[
 \ell_{1/s_j}\operatorname{dist}(a,\mathbb Z^4)
 \longrightarrow\infty.
\]
More explicitly, for every fixed \(x_0\in\mathbb R^4\),
\[
 \inf_{n\in\mathbb Z^4}
 \|A_{1/s_j}^{(\nu)}(a+n)-x_0\|
 \geq
 \inf_{n\in\mathbb Z^4}\|A_{1/s_j}^{(\nu)}(a+n)\|-\|x_0\|
 \longrightarrow\infty.
\]
Hence \(\{z_j:j\geq1\}\) is closed in \(\mathfrak E_\nu\): it has no
accumulation in a positive fibre and the last estimate shows that it
eventually misses every half-space chart about every point of the zero
fibre.  But
\[
 \beta_\nu(z_j)=[e^{-2\pi/s_j},a]\longrightarrow d_\nu,
\]
so \(\{\beta_\nu(z_j):j\geq1\}\) is not closed in \(\mathcal D_\nu\).
Since \(\beta_\nu\) is injective off the zero fibre and the latter maps to
\(d_\nu\), the inverse image of this nonclosed set is the closed set
\(\{z_j:j\geq1\}\).  Thus \(\beta_\nu\) is not a quotient map.  This proves
the two diagrams and their exact topological scope.  Since \(\beta_\nu\) is
constant exactly on \(p_\nu^{-1}(0)\) and otherwise injective, it induces the
continuous bijection \(\overline\beta_\nu\) displayed in the theorem.  Were
that bijection a homeomorphism, \(\beta_\nu\) would be the composite of the
quotient projection with a homeomorphism and hence a quotient map, contrary
to the counterexample.  Therefore \(\overline\beta_\nu\) is not a
homeomorphism.
\end{enumerate}
\end{proof}

\begin{corollary}[Sharp global meaning of the Euclidean limit]
\label{wb:cusp-cone-sharp-global-R4-scope}
\StatementRecord{wb:cusp-cone-sharp-global-R4-scope}{proved}{Historical events: GPT-6 Astra (OpenAI), ChatGPT Pro mode, effort Pro, S6\_09\_PAPERS, 17 November 2026: targeted proof check. Also checked in this writing. Pre-continuation writing events: GPT-6 Astra (OpenAI), ChatGPT Pro mode, effort Pro, S6\_10\_PAPERS, 18 November 2026: targeted proof check. Continuation metadata events: GPT-6 Astra (OpenAI), ChatGPT Pro mode, effort Pro, S6\_10\_PAPERS continuation, 19 November 2026: used and rechecked in this run; status-vocabulary reconciliation; metadata reconciliation only; no new mathematical check. \emph{Pre-continuation recorded proof scope.} GPT-6 Astra (OpenAI, ChatGPT Pro mode), S6\_09\_PAPERS, effort Pro, 2026-11-17; also checked in this writing by GPT-6 Astra (OpenAI), ChatGPT Pro mode, effort Pro, 18 November 2026}{the hypotheses, notation and definitions stated here; \cref{wb:cusp-cone-literal-R4-open-cells}; \cref{wb:cusp-cone-explicit-carriers}; \cref{wb:cusp-cone-exact-decompactification-space}; \cref{wb:cusp-cone-carrier-local-homology}; \cref{wb:sphere-metric-fixed-realization}}

The native routes constructed above exhibit the underlying space
\(\mathbb R^4\) in the following three precise senses:
\begin{enumerate}
\item as each of the nonclosed topologically embedded open cells
\(U_{12,T}\subset S^6\) and \(U_{24,T}\subset S^8\);
\item as the pointed intrinsic Euclidean limit of the four-torus slices in the
global quotient carriers of \(S^{12}\) and \(S^{24}\); and
\item as the Euclidean decompactification boundary fibre of
\(\mathfrak E_{12}\) and \(\mathfrak E_{24}\).
\end{enumerate}
It does not produce \(\mathbb R^4\) as the whole compact sphere, as the
whole compact conical target, or as a continuous quotient target of either
compact space.  Nor does the unrescaled restriction of any fixed compatible
ambient target metric have \(\mathbb R^4\) as its slice limit, and the open
cells in item \(1\) are not globally Euclidean-isometric in such a metric.
\end{corollary}



\begin{proof}
\begin{enumerate}[label=\arabic*.,leftmargin=*]
\item The three positive assertions are
\cref{wb:cusp-cone-literal-R4-open-cells} and \cref{wb:cusp-cone-explicit-carriers} and \cref{wb:cusp-cone-exact-decompactification-space}.
For the negative assertions, a continuous image of a compact sphere or
compact conical target is compact, so it cannot be all of \(\mathbb R^4\).
The local homology calculation
\cref{wb:cusp-cone-carrier-local-homology} shows that the carrier vertex is not
an integral homology-manifold point, and hence not a manifold point, of any
dimension, with or without boundary.  Finally, in every fixed compatible
metric on a compact cone target, all sufficiently deep radial slices lie
uniformly in an arbitrarily small neighborhood of the vertex, so their ambient
diameters tend to zero.  Their intrinsic diameters tend to infinity because
their intrinsic systoles do.  Hence without a \(T\)-dependent rescaling no
fixed ambient metric realizes the intrinsic Euclidean limit on those slices.
The fixed cusp-adapted blow-ups are treated separately in
\cref{wb:sphere-metric-fixed-realization}.  The comparison diagrams record the
unrescaled distinction exactly: they collapse the new \(\mathbb R^4\)
boundary rather than identify it with the compact target.
\end{enumerate}
\end{proof}

\paragraph{Scope.}
The constructions above use only the original, unscaled fibre metrics and
the quotient maps derived from the native \(S^6\) cusp.  They do not use the
auxiliary finite-cover unspooling.  The conical carriers depend on the
chosen cusp coordinate, positive radial ray, and based identification
\(X\cong S^6\).  The spaces \(\mathfrak E_\nu\) additionally remember the
marked Euclidean quotient metrics.  None of these data defines a compact
complex structure on \(S^{12}\) or \(S^{24}\), a four-dimensional physical
spacetime, or an interacting Yang--Mills Hamiltonian.

\subsection{The full \texorpdfstring{\(S^{12}\to S^6\)}{S12-to-S6}
carrier is null-homotopic}
\label{wb:cusp-s12-carrier-homotopy}

This subsection determines the based homotopy class of the full two-copy
sphere carrier.  It is important that the calculation concerns the full
quotient map, not merely the inclusion of one of its moving four-spheres.

Recall the exact maps from
\cref{wb:cusp-cone-S12-route} and \cref{wb:cusp-cone-angle-forgotten-sphere-carriers}.  The
cusp quotient is
\[
 q_2:(X,x_0)\longrightarrow
 \mathcal A=\operatorname{Susp}(S^1\times\overline T^2),
\]
and, for the coordinate one-skeleton
\[
 K^{(1)}=
 ((\mathbb R/\mathbb Z)\times\{0\})
 \cup(\{0\}\times(\mathbb R/\mathbb Z))\subset\overline T^2,
\]
the degree-one quotient
\(c:\overline T^2\to\overline T^2/K^{(1)}\cong S^2\) defines
\[
 \widehat p_2([r,e^{i\phi},u])=[r,c(u)]
 \in\Sigma_r S^2\cong S^3.
\]
Thus
\[
 g:=\widehat q_3=\widehat p_2\circ q_2:(X,x_0)\longrightarrow(S^3,*)
\]
is the specified continuous surjective based quotient.  After the fixed
based identification \(X\cong_{\mathrm{diff}}S^6\), its smash square is the
specified continuous surjective based quotient
\[
 g^{\wedge2}:S^{12}=S^6\wedge S^6
 \longrightarrow S^3\wedge S^3=S^6.
\]

\begin{theorem}[The full external smash pairing is zero]
\label{wb:cusp-s12-full-carrier-null}
\StatementRecord{wb:cusp-s12-full-carrier-null}{proved}{Historical events: GPT-6 Astra (OpenAI), ChatGPT Pro mode, effort Pro, S6\_09\_PAPERS, 17 November 2026: targeted proof check. Also checked in this writing. Pre-continuation writing events: GPT-6 Astra (OpenAI), ChatGPT Pro mode, effort Pro, S6\_10\_PAPERS, 18 November 2026: targeted proof check. Continuation metadata events: GPT-6 Astra (OpenAI), ChatGPT Pro mode, effort Pro, S6\_10\_PAPERS continuation, 19 November 2026: used and rechecked in this run; status-vocabulary reconciliation; metadata reconciliation only; no new mathematical check. \emph{Pre-continuation recorded proof scope.} GPT-6 Astra (OpenAI, ChatGPT Pro mode), S6\_09\_PAPERS, effort Pro, 2026-11-17; also checked in this writing by GPT-6 Astra (OpenAI), ChatGPT Pro mode, effort Pro, 18 November 2026}{the hypotheses, notation and definitions stated here; \cite[opening table, p.~339]{Hatcher2002}; \cite[Table~1, p.~113]{IsaksenWangXu2023}; \cref{wb:cusp-s12-toda-smash-formula}}

For every pair of based maps \(x,y:S^6\to S^3\),
\[
 [x\wedge y]=0\qquad\text{in}\qquad\pi_{12}(S^6).
\]
In particular,
\(\widehat q_3^{\wedge2}:S^{12}\twoheadrightarrow S^6\) is simultaneously
surjective as a point-set quotient and based null-homotopic.
\end{theorem}



\begin{proof}
\begin{enumerate}[label=\arabic*.,leftmargin=*]
\item Write \(E\) for reduced suspension, using the ordered smash coordinates on
spheres.  The external smash product has the exact factorization
\begin{equation}
 x\wedge y
 =
 (x\wedge\operatorname{id}_{S^3})
 \circ
 (\operatorname{id}_{S^6}\wedge y):
 S^6\wedge S^6\longrightarrow S^6\wedge S^3
 \longrightarrow S^3\wedge S^3.
 \label{wb:cusp-s12-smash-factorization}
\end{equation}
Under the displayed ordered sphere identifications, the two factors, in
order of application, represent \(E^6y\in\pi_{12}(S^9)\) and
\(E^3x\in\pi_9(S^6)\).  Equivalently, Toda's
Barratt--Hilton reduced-join formula gives, with
\(p=q=k=h=3\),
\begin{equation}
 [x\wedge y]=(E^3[x])\circ(E^6[y]).
 \label{wb:cusp-s12-toda-smash-formula}
\end{equation}
Here Toda's sign is \((-1)^{(p+k)h}=(-1)^{18}=+1\).  Reversing the
two displayed factors introduces the separate twist sign
\((-1)^{ph}=-1\); no such reversal is used here.

\item The three exact groups needed below are
\begin{equation}
 \pi_6(S^3)\cong\mathbb Z/12,
 \qquad
 \pi_{12}(S^9)\cong\mathbb Z/24,
 \qquad
 \pi_{12}(S^6)\cong\mathbb Z/2.
 \label{wb:cusp-s12-three-stem-groups}
\end{equation}
The first and third are the integral low unstable groups displayed in
Hatcher's table extracted from Toda; the middle group is in the stable
three-stem
\cite[opening table, p.~339]{Hatcher2002};
\cite[Table~1, p.~113]{IsaksenWangXu2023}.

\item Put \(\beta=E^6[y]\in\pi_{12}(S^9)\).  Since
\(12[y]=0\), functoriality of suspension gives \(12\beta=0\).  In the
cyclic group \(\mathbb Z/24\), the elements annihilated by \(12\) are
exactly the even elements.  Hence there is
\(\gamma\in\pi_{12}(S^9)\) such that \(\beta=2\gamma\).
Postcomposition with \(E^3[x]:S^9\to S^6\) induces a group homomorphism
\[
 (E^3[x])_*:\pi_{12}(S^9)\longrightarrow\pi_{12}(S^6).
\]
The target has exponent two, so \cref{wb:cusp-s12-toda-smash-formula} gives
\[
 [x\wedge y]
 =(E^3[x])_*(\beta)
 =2(E^3[x])_*(\gamma)
 =0.
\]
This proves the vanishing of the entire external smash pairing and hence
the theorem.
\end{enumerate}
\end{proof}

\begin{corollary}[All higher smash powers]
\label{wb:cusp-s12-all-higher-smashes-null}
\StatementRecord{wb:cusp-s12-all-higher-smashes-null}{proved}{Historical events: GPT-6 Astra (OpenAI), ChatGPT Pro mode, effort Pro, S6\_09\_PAPERS, 17 November 2026: targeted proof check. Also checked in this writing. Pre-continuation writing events: GPT-6 Astra (OpenAI), ChatGPT Pro mode, effort Pro, S6\_10\_PAPERS, 18 November 2026: targeted proof check. Continuation metadata events: GPT-6 Astra (OpenAI), ChatGPT Pro mode, effort Pro, S6\_10\_PAPERS continuation, 19 November 2026: used and rechecked in this run; status-vocabulary reconciliation; metadata reconciliation only; no new mathematical check. \emph{Pre-continuation recorded proof scope.} GPT-6 Astra (OpenAI, ChatGPT Pro mode), S6\_09\_PAPERS, effort Pro, 2026-11-17; also checked in this writing by GPT-6 Astra (OpenAI), ChatGPT Pro mode, effort Pro, 18 November 2026}{the hypotheses, notation and definitions stated here; the explicit constructions and calculations in the proof}

For every based \(g:S^6\to S^3\) and every \(r\geq2\),
\[
 g^{\wedge r}:S^{6r}\longrightarrow S^{3r}
\]
is based null-homotopic.
\end{corollary}



\begin{proof}
\begin{enumerate}[label=\arabic*.,leftmargin=*]
\item With the fixed associativity homeomorphisms for smash products,
\[
 g^{\wedge r}=g^{\wedge2}\wedge g^{\wedge(r-2)}.
\]
Smashing a based null-homotopy with a based map gives a based
null-homotopy.
\end{enumerate}
\end{proof}


\paragraph{The one-copy invariant and its computed value}
\label{wb:cusp-s12-one-copy-invariant}
\StatementRecord{wb:cusp-s12-one-copy-invariant}{proved}{Historical events: GPT-6 Astra (OpenAI), ChatGPT Pro mode, effort Pro, S6\_09\_PAPERS, 17 November 2026: targeted proof check. Also checked in this writing. Pre-continuation writing events: GPT-6 Astra (OpenAI), ChatGPT Pro mode, effort Pro, S6\_10\_PAPERS, 18 November 2026: targeted proof check. Continuation metadata events: GPT-6 Astra (OpenAI), ChatGPT Pro mode, effort Pro, S6\_10\_PAPERS continuation, 19 November 2026: used and rechecked in this run; status-vocabulary reconciliation; metadata reconciliation only; no new mathematical check. \emph{Pre-continuation recorded proof scope.} GPT-6 Astra (OpenAI, ChatGPT Pro mode), S6\_09\_PAPERS, effort Pro, 2026-11-17; also checked in this writing by GPT-6 Astra (OpenAI), ChatGPT Pro mode, effort Pro, 18 November 2026}{the hypotheses, notation and definitions stated here; \cite[Theorem~5.41(b), p.~573, and pp.~577--578]{HatcherSpectralSequences}; \cref{wb:cusp-null:thm}}

The proof does not need to determine the one-copy class
\([g]\in\pi_6(S^3)\cong\mathbb Z/12\).  Toda's suspension calculation
identifies \(\pi_6(S^3)\) with the even subgroup of the stable three-stem;
Hatcher gives an integral Serre-spectral-sequence derivation of the same
injectivity and identifies twice the stable Hopf--\(\nu\) generator with a
generator of the suspended \(\mathbb Z/12\)-summand
\cite[Theorem~5.41(b), p.~573, and pp.~577--578]{HatcherSpectralSequences}.
Thus, after choosing a stable generator \(\nu_9\), there is a unique
\(a\in\mathbb Z/12\) such that
\begin{equation}
 [E^6g]=2a\nu_9\in\pi_{12}(S^9)\cong\mathbb Z/24.
 \label{wb:cusp-s12-sixfold-suspension-class}
\end{equation}
Because \(a\mapsto2a\) embeds \(\mathbb Z/12\) as the even subgroup of
\(\mathbb Z/24\), \(E^6g\) is null-homotopic if and only if \(g\) is
null-homotopic.  Thus the one-copy class and its sixfold suspension are the
same integer modulo \(12\). It is now zero by Theorem~\ref{wb:cusp-null:thm}.

The exact geometric invariant which computes \(a\) is also identifiable.
Choose \(0<r_0<\epsilon\) and
\(u_0\in\overline T^2\setminus K^{(1)}\), and let
\(y_0=[r_0,c(u_0)]\) lie in the open three-cell of
\(\Sigma_rS^2\).  From the literal formula for \(q_2\),
\begin{equation}
 g^{-1}(y_0)
 =\{x\in f^{-1}(|t|=r_0):\overline\pi_f(x)=u_0\}.
 \label{wb:cusp-s12-regular-inverse-image}
\end{equation}
The fibre of \(\overline\pi\) is the two-torus defined by
\(K_\Lambda=\mathbb Z\langle\widehat w,\widehat\delta\rangle\), and the
cusp monodromy fixes both \(\widehat w\) and \(\widehat\delta\).  Therefore
\eqref{wb:cusp-s12-regular-inverse-image} is the mapping torus of an affine translation of the kernel torus,
which is diffeomorphic to \(S^1\times T^2\cong T^3\): if the translation is
\(k\mapsto k+v\), choose a real lift \(\widetilde v\); the map
\([\theta,k]\mapsto(e^{2\pi i\theta},k+\theta\widetilde v)\) respects
\((1,k)\sim(0,k+v)\) and has the inverse obtained by subtracting
\(\theta\widetilde v\). After choosing a smooth chart in the open
target cell, the differential of \(g\) supplies its normal three-framing.
Pontryagin--Thom identifies the stabilized class of \(g\) with the framed
bordism class of this specifically framed \(T^3\).  Since stabilization on
\(\pi_6(S^3)\) is injective, computing that framing class computes \(a\)
exactly.  The unframed diffeomorphism type \(T^3\), point-set surjectivity,
and the zero ordinary homology map do not compute \(a\).






\paragraph{Scope}
\label{wb:cusp-s12-nullity-scope}
\StatementRecord{wb:cusp-s12-nullity-scope}{proved}{Historical events: GPT-6 Astra (OpenAI), ChatGPT Pro mode, effort Pro, S6\_09\_PAPERS, 17 November 2026: targeted proof check. Also checked in this writing. Pre-continuation writing events: GPT-6 Astra (OpenAI), ChatGPT Pro mode, effort Pro, S6\_10\_PAPERS, 18 November 2026: targeted proof check. Continuation metadata events: GPT-6 Astra (OpenAI), ChatGPT Pro mode, effort Pro, S6\_10\_PAPERS continuation, 19 November 2026: used and rechecked in this run; status-vocabulary reconciliation; metadata reconciliation only; no new mathematical check. \emph{Pre-continuation recorded proof scope.} GPT-6 Astra (OpenAI, ChatGPT Pro mode), S6\_09\_PAPERS, effort Pro, 2026-11-17; also checked in this writing by GPT-6 Astra (OpenAI), ChatGPT Pro mode, effort Pro, 18 November 2026}{the hypotheses, notation and definitions stated here; the argument and qualifications immediately following}

The theorem changes no point-set or metric statement.  A surjective quotient
map may be null-homotopic, and the embedded moving four-spheres and their
intrinsic pointed \(\mathbb R^4\)-limits remain exactly as constructed.  The
theorem concerns the sphere-valued postcomposition
\(\widehat q_3^{\wedge2}\); it does not imply null-homotopy of the earlier
conical-target map
\(q_2^{\wedge2}:S^{12}\to\mathcal A^{\wedge2}\).





\subsection{The original cusp coefficient: an explicit framed nullbordism}


Let \(X\) be the smooth compact source of the original construction,
\(f:X\to B\) its map, and \(t\) the specified cusp coordinate.
Put \(W=f^{-1}(0)\), choose the original \(x_0\in W\), and retain the
ordered quotient lattice basis induced by
\((\widehat\gamma,\widehat u)\), with kernel generated by
\((\widehat w,\widehat\delta)\). Thus the quotient torus is
\(\overline T^2=\mathbb R^2/\mathbb Z^2\). The original periods are
\begin{equation}
 Z(s)=sB_0+C(t),\qquad e^{2\pi i s}=t,\qquad
 B_0=\begin{pmatrix}0&1\\-1&0\end{pmatrix},\qquad
 C(t)=\begin{pmatrix}6\mu(t)&h(t)\\b(t)-h(t)&\mu(t)\end{pmatrix}.
 \label{wb:cusp-null:eq:periods}
\end{equation}
Every entry of \(C\) is holomorphic at \(t=0\).
The dense torus of the original toric cusp has coordinates
\((x_1,x_2,t)\), with
\(x=\exp(2\pi i\zeta)\). The residual lattice action is exactly
\begin{equation}
 \Psi_\lambda(x,t)=
 \bigl(t^{B_0\lambda}\exp(2\pi iC(t)\lambda)x,t\bigr),
 \qquad\lambda\in\mathbb Z^2.
 \label{wb:cusp-null:eq:action}
\end{equation}
It extends freely and properly discontinuously across the original
unimodular toric filling. In particular the quotient projection is a local
biholomorphism at every toric fixed point. These are the reconstructed
data and proved local properties of the original source
(the reconstructed original manuscript, \S4, source pp.~24--29; its complete cusp data are displayed here).

Let
\[
 K^{(1)}=(\mathbb R/\mathbb Z\times\{0\})\cup(\{0\}\times\mathbb R/\mathbb Z),\qquad
 c:\overline T^2\longrightarrow \overline T^2/K^{(1)}\cong S^2
\]
be the specified degree-one cell quotient. Its restriction to the open
two-cell is taken with the smooth coordinates used in the original
regular-fibre description. Both radial endpoints and the complete
basepoint cylinder are collapsed in
\[
 \Sigma_r(S^2,*)=
 \frac{[0,\epsilon]\times S^2}
 {(\{0,\epsilon\}\times S^2)\cup([0,\epsilon]\times\{*\})}
 \cong S^3.
\]
The original based map is
\begin{equation}
 g(x)=[|t|,c(\overline\pi_f(x))]\quad(0<|t|<\epsilon),\qquad
 g(W)=g(X\setminus\operatorname{int}N_0)=*.
 \label{wb:cusp-null:eq:g}
\end{equation}

\begin{theorem}[The original one-copy cusp map is null-homotopic]
\label{wb:cusp-null:thm}
\StatementRecord{wb:cusp-null:thm}{proved}{Historical events: GPT-6 Astra (OpenAI), ChatGPT Pro mode, effort Pro, S6\_09\_PAPERS, 17 November 2026: targeted proof check. Also checked in this writing. Pre-continuation writing events: GPT-6 Astra (OpenAI), ChatGPT Pro mode, effort Pro, S6\_10\_PAPERS, 18 November 2026: targeted proof check. Continuation metadata events: GPT-6 Astra (OpenAI), ChatGPT Pro mode, effort Pro, S6\_10\_PAPERS continuation, 19 November 2026: used and rechecked in this run; status-vocabulary reconciliation; metadata reconciliation only; no new mathematical check. \emph{Pre-continuation recorded proof scope.} GPT-6 Astra (OpenAI, ChatGPT Pro mode), S6\_09\_PAPERS, effort Pro, 2026-11-17; also checked in this writing by GPT-6 Astra (OpenAI), ChatGPT Pro mode, effort Pro, 18 November 2026}{the hypotheses, notation and definitions stated here; \cref{wb:cusp-null:eq:g}; \cref{wb:cusp-null:eq:moduli}; \cref{wb:cusp-null:eq:chart}; \cref{wb:cusp-null:eq:J}; \cref{wb:cusp-null:eq:spatial-bound}; \cref{wb:cusp-null:eq:bordism}; \cref{wb:cusp-null:eq:P}; \cite[\S7, Theorem B and Lemma 4,
pp.~43, 48--50]{CuspNullMilnor1965}}

The based map \(g:(X,x_0)\to(S^3,*)\) in \eqref{wb:cusp-null:eq:g} is based
null-homotopic. More precisely, a regular inverse image with its actual
ordered normal framing induced by \((r,c_1,c_2)\) is framed nullbordant
inside one original cusp chart times \([0,1]\), disjoint from
\(\{x_0\}\times[0,1]\).
\end{theorem}



\begin{proof}
\begin{enumerate}[label=\arabic*.,leftmargin=*]
\item The proof constructs the complete inverse image, compares its actual
framing with an explicitly extendible framing, and gives the extension.

\item \paragraph{1. The exact complete inverse image.}
Retain the lift
\[
 u^0=\begin{pmatrix}-1/4\\1/4\end{pmatrix},\qquad
 \overline u^0=(3/4,1/4)\pmod{\mathbb Z^2},\qquad
 B_0u^0=\begin{pmatrix}1/4\\1/4\end{pmatrix}.
\]
Thus \(\overline u^0\notin K^{(1)}\), and \(B_0u^0\) lies in the
interior of the original triangle \(\operatorname{conv}(0,e_1,e_2)\).
Use its original unimodular chart
\begin{equation}
 z_0=t/(x_1x_2),\qquad z_1=x_1,\qquad z_2=x_2,
 \qquad t=z_0z_1z_2.
 \label{wb:cusp-null:eq:chart}
\end{equation}
Set
\[
 \ell=\log|t|<0,\qquad R(t)=-2\pi\operatorname{Im}C(t),\qquad
 A(t)=\ell B_0+R(t).
\]
For a representative \(\zeta=Z(s)u+k\), where \(u,k\in\mathbb R^2\),
\(\operatorname{Im}s=-\ell/(2\pi)\) gives the exact identity
\begin{equation}
 \begin{pmatrix}\log|x_1|\\\log|x_2|\end{pmatrix}
 =A(t)u.
 \label{wb:cusp-null:eq:moduli}
\end{equation}
The quotient value is \(u\bmod\mathbb Z^2\).
Since \(B_0\) is orthogonal, for all \(v\in\mathbb R^2\),
\begin{equation}
 \|A(t)v\|\geq (|\ell|-\|R(t)\|)\|v\|.
 \label{wb:cusp-null:eq:invertibility}
\end{equation}
The boundedness of \(R\) therefore makes \(A(t)\) invertible on a
sufficiently small punctured disc.

\item For any point of the fibre quotient over \(\overline u^0\), subtracting
the appropriate original period, equivalently applying
\(\Psi_{-\lambda}\), gives a representative with the exact lift
\(u=u^0\). This representative is unique: another residual translate
has logarithmic moduli differing by \(A(t)\lambda\), which vanishes
only for \(\lambda=0\). Conversely every dense-torus point satisfying
\(\log|x|=A(t)u^0\) has this quotient value, since the remaining real
parts of \(\zeta\) are exactly the kernel coordinates \(k\).
For each fixed \(t\), varying \(k\bmod\mathbb Z^2\) gives all independent
phases of \(x_1,x_2\).

\item The following change of coordinates uses the full original period units:
\begin{align}
 h_j^0(t)&=\exp\bigl(-2\pi i(C(t)u^0)_j\bigr),\qquad j=1,2,\\
 \mathcal F(z_0,z_1,z_2)
 &=(w_0,w_1,w_2)
 =\left(\frac{z_0}{h_1^0(t)h_2^0(t)},\,
 h_1^0(t)z_1,\,h_2^0(t)z_2\right).
 \label{wb:cusp-null:eq:straightening}
\end{align}
The functions \(h_j^0\) are nowhere-zero holomorphic units, including
at \(t=0\). The product \(w_0w_1w_2\) is exactly \(t\).
The inverse of \(\mathcal F\) is
\[
 (w_0,w_1,w_2)\longmapsto
 \left(w_0h_1^0(t)h_2^0(t),\,
 \frac{w_1}{h_1^0(t)},\,\frac{w_2}{h_2^0(t)}\right),
 \qquad t=w_0w_1w_2.
\]
Thus this is a biholomorphism of the original chart restricted to the
cusp disc, not a replacement of \(C(t)\).

\item For \(r_0\in(0,\epsilon)\), the entire inverse image of
\(y_0=[r_0,c(\overline u^0)]\) becomes exactly
\begin{equation}
 T_{r_0}=
 \{|w_0|=\rho_0,\ |w_1|=\rho_1,\ |w_2|=\rho_2\},\qquad
 (\rho_0,\rho_1,\rho_2)=
 (r_0^{1/2},r_0^{1/4},r_0^{1/4}).
 \label{wb:cusp-null:eq:torus}
\end{equation}
Indeed, \(\log|h^0(t)|=2\pi\operatorname{Im}C(t)u^0=-R(t)u^0\),
so \eqref{wb:cusp-null:eq:moduli} gives \(\log|w_j|=\ell(B_0u^0)_j\) for
\(j=1,2\); the product identity gives the exponent \(1/2\) for
\(w_0\). Conversely these moduli recover \(u=u^0\) by invertibility
of \(A(t)\). The angular map is explicitly
\phantomsection\label{qa:III:layout:1}
\[
\begin{gathered}
 (\arg t,\arg w_1,\arg w_2)\longleftrightarrow
 (\arg w_0,\arg w_1,\arg w_2),\\
 \arg w_0=\arg t-\arg w_1-\arg w_2\pmod{2\pi},
\end{gathered}
\]
with inverse \(\arg t=\arg w_0+\arg w_1+\arg w_2\).
Its integral matrix has determinant one. This includes the complete
angular mapping torus and all kernel phases, with their original signs.

\item Let \(o\) be the origin of \eqref{wb:cusp-null:eq:chart}. The original free properly
discontinuous action supplies a neighborhood \(V\) of \(o\) on which
the quotient projection is injective. To see the local injectivity
directly, first choose a relatively compact neighborhood; proper
discontinuity leaves only finitely many translates that can intersect it.
None of those nonidentity translates fixes \(o\). Shrinking finitely
many times separates the neighborhood from all of them. The quotient
map is a local biholomorphism, so its restriction to \(V\) is a chart
of the original source.
Choose a polydisc about \(0\) whose closure lies in \(\mathcal F(V)\).
For sufficiently small \(r_0\), both \(T_{r_0}\) and
\begin{equation}
 \{|w_0|\leq\rho_0,\ |w_1|=\rho_1,\ |w_2|=\rho_2\}
 \label{wb:cusp-null:eq:spatial-bound}
\end{equation}
lie in this polydisc. This proves that the whole inverse image embeds
in one chart and supplies space for the bounding manifold below.

\item \paragraph{2. The actual normal framing, including its angle dependence.}
Put \(L_j=\log|w_j|\), and use the unique real lift of the quotient
coordinate near \(u^0\). From the exact change of coordinates,
\begin{equation}
 \begin{pmatrix}L_1\\L_2\end{pmatrix}
 =\ell B_0u^0+A(t)(u-u^0),\qquad
 L_0=\ell-L_1-L_2.
 \label{wb:cusp-null:eq:nearby-moduli}
\end{equation}
Differentiating on \(u=u^0\) annihilates the term
\((dA(t))(u-u^0)\). Hence no derivative of \(C(t)\) has been omitted:
the exact differential there is
\begin{equation}
 \begin{pmatrix}dL_0\\dL_1\\dL_2\end{pmatrix}
 =J(t)\begin{pmatrix}d\ell\\du_1\\du_2\end{pmatrix},\quad
 J(t)=
 \begin{pmatrix}
 1/2&-A_{11}-A_{21}&-A_{12}-A_{22}\\
 1/4&A_{11}&A_{12}\\
 1/4&A_{21}&A_{22}
 \end{pmatrix}.
 \label{wb:cusp-null:eq:J}
\end{equation}
Adding the second and third rows to the first gives \((1,0,0)\),
and expansion in that row proves \(\det J(t)=\det A(t)\).

\item Let \(D=Dc_{\overline u^0}\) be the actual two-by-two derivative in
the ordered open-cell target coordinates \((c_1,c_2)\). It is a fixed
invertible matrix; all its entries and signs are retained. Let
\[
 H=\begin{pmatrix}r_0&0&0\\0&D_{11}&D_{12}\\0&D_{21}&D_{22}\end{pmatrix}.
\]
Let \(n_j\) be the real radial vector \(w_j\) in the \(j\)-th complex
coordinate, with other components zero. Then \(dL_i(n_j)=\delta_{ij}\).
Writing \(N=(n_0,n_1,n_2)\) as the matrix with these vector columns, the actual
\(dg\)-dual normal framing is, in the normal quotient bundle,
\begin{equation}
 \mathfrak n(t)=N J(t)H^{-1}.
 \label{wb:cusp-null:eq:actual-frame}
\end{equation}
Indeed, \(dg(N)=HJ(t)^{-1}\), so multiplication by \(J(t)H^{-1}\)
gives the identity. In the Euclidean metric of this chart the radial
vectors are perpendicular to the product torus; a different source
metric changes their representatives by the canonical projection to its
normal bundle and preserves this normal-quotient framing.

\item Put \(\ell_0=\log r_0\). Decrease \(r_0\) further, if necessary, so
that \(|\ell_0|>2\sup_{|t|\leq r_0}\|R(t)\|\). For
\(0\leq\sigma\leq1\), define only a homotopy of frames by
\[
 A_\sigma(t)=\ell_0B_0+\sigma R(t),
\]
and let \(J_\sigma(t)\) be \eqref{wb:cusp-null:eq:J} with \(A_\sigma\) in place
of \(A\). For every \(v\in\mathbb R^2\),
\[
 \|A_\sigma(t)v\|\geq
 (|\ell_0|-\|R(t)\|)\|v\|>0\quad(v\ne0).
\]
Thus all matrices are invertible. At \(\sigma=0\),
\begin{equation}
 J_0=
 \begin{pmatrix}
 1/2&\ell_0&-\ell_0\\
 1/4&0&\ell_0\\
 1/4&-\ell_0&0
 \end{pmatrix},\qquad \det J_0=\ell_0^2>0.
 \label{wb:cusp-null:eq:J0}
\end{equation}
Continuity also proves \(\det J_\sigma(t)>0\) for every \(\sigma,t\).
The homotopy \(N J_\sigma(t)H^{-1}\), run from \(\sigma=1\) to
\(\sigma=0\), compares the actual frame with the exact constant-matrix
frame \(N P\), where
\begin{equation}
 P=J_0H^{-1}.
 \label{wb:cusp-null:eq:P}
\end{equation}
The original embedded torus and original functions \(C(t)\) have not
changed. A framing homotopy is a framed cylinder: choose a smooth
parameter function constant near its endpoints and use the displayed
frames on \(T_{r_0}\times[0,1]\).

\item \paragraph{3. A collared framed nullbordism with the exact constant matrix.}
Use a new time coordinate \(v\in[0,1]\). Define
\begin{equation}
 q(v)=
 \begin{cases}
 0,&0\leq v\leq1/4,\\
 \rho_0^2\exp\bigl(4-1/(v-1/4)\bigr),&1/4<v\leq1.
 \end{cases}
 \label{wb:cusp-null:eq:q}
\end{equation}
This function is smooth: all right derivatives of
\(\exp(-1/(v-1/4))\) at \(1/4\) vanish. It is strictly increasing
for \(v>1/4\), satisfies \(q(1/2)=\rho_0^2\), and
\(q'(1/2)=16\rho_0^2>0\).
In the chart times \([0,1]\), put
\begin{equation}
 \mathcal W=
 \{|w_0|^2+q(v)=\rho_0^2,\quad
       |w_1|^2=\rho_1^2,\quad|w_2|^2=\rho_2^2\}.
 \label{wb:cusp-null:eq:bordism}
\end{equation}
No point has \(v>1/2\). The three defining differentials are independent:
the latter two have nonzero gradients in their separate complex factors;
the first has nonzero \(w_0\)-gradient whenever \(w_0\ne0\), and at
\(w_0=0\) it has the nonzero \(v\)-derivative \(q'(1/2)\).
Thus \(\mathcal W\) is a smooth compact four-manifold. Its only
boundary is \(T_{r_0}\times\{0\}\), and on \(0\leq v\leq1/4\)
it is the exact product collar \(T_{r_0}\times[0,1/4]\).
The north locus \(w_0=0,v=1/2\) is an interior cap: the implicit
function theorem expresses \(v\) smoothly in terms of
\((\operatorname{Re}w_0,\operatorname{Im}w_0)\) there. The first
factor is a disk obtained by rotating its radial profile with this smooth
cap and one boundary circle. Therefore
\(\mathcal W\cong D^2\times T^2\).
Its spatial coordinates lie in \eqref{wb:cusp-null:eq:spatial-bound}, so the complete
manifold lies in the injective original chart.

\item In the Euclidean chart product metric define
\[
 V_0=(w_0,0,0,q'(v)/2),\qquad
 V_1=(0,w_1,0,0),\qquad V_2=(0,0,w_2,0).
\]
These are one half of the gradients of the three defining equations
in \eqref{wb:cusp-null:eq:bordism}. They are mutually orthogonal and never zero.
For example, if \(F_0=|w_0|^2+q(v)-\rho_0^2\), then
\[
 dF_0(V_0)=2|w_0|^2+\frac12q'(v)^2>0
\]
at every point of \(\mathcal W\). On the entire collar \(q'=0\), so
\((V_0,V_1,V_2)=(N,0)\) exactly. Consequently
\begin{equation}
 (V_0,V_1,V_2)P
 \label{wb:cusp-null:eq:extended-frame}
\end{equation}
is a framing of the normal three-plane bundle of \(\mathcal W\),
and its restriction to the collar is exactly the framing \(NP\) from
\eqref{wb:cusp-null:eq:P}, with every original target-coordinate sign retained.
Concatenate the framing cylinder of the preceding paragraph with this
nullbordism, rescaling the two interval lengths. Their framings agree
exactly on the common product collar, so the concatenation is a smooth
framed nullbordism of the actual inverse image.

\item \paragraph{4. The based map represented by this framed inverse image.}
Choose the injective chart to exclude \(x_0\) if its center is not
\(x_0\). If its center is \(x_0\), the constructed nullbordism still
excludes \(x_0\), since \(|w_1|=\rho_1>0\) and
\(|w_2|=\rho_2>0\) everywhere on it. Compactness therefore gives
tubular neighborhoods for the framed cylinder and cap disjoint from
\(\{x_0\}\times[0,1]\).

\item The original continuous quotient map is smooth and a submersion near
its entire inverse image of \(y_0\). First make it constant in a
neighborhood of \(x_0\), by a based homotopy that changes nothing near
that inverse image. Explicitly, take a target chart
\(\kappa:O\to\mathbb R^3\) with \(\kappa(*)=0\), whose image is a ball,
and \(y_0\notin O\). By continuity choose a neighborhood \(U\) of
\(x_0\) with \(g(U)\subset O\), and a smooth bump function
\(\beta\) supported in \(U\) and equal to one near \(x_0\).
For \(x\in U\), replace \(g(x)\) during this homotopy by
\(\kappa^{-1}((1-a\beta(x))\kappa(g(x)))\), \(0\leq a\leq1\),
and leave \(g\) fixed off \(U\). Convexity of the target chart image
makes the formula well-defined, the vanishing of \(\beta\) near the
boundary of its support proves continuity of the pasted formula, and
\(\kappa(g(x_0))=0\) makes the homotopy based. It never introduces
an inverse image of \(y_0\). Its endpoint is smooth near \(x_0\) and
agrees with \(g\) near the entire old inverse image.

\item Denote the endpoint of the preceding homotopy by \(g_b\). Choose one
closed tubular neighborhood \(U_1\) of the inverse image, contained in
the smooth submersion region and with the entire inverse image contained
in \(\operatorname{int}U_1\). The compact set
\(X\setminus\operatorname{int}U_1\) has image under \(g_b\) missing
\(y_0\). In a fixed target metric put
\[
 d_0=\min_{x\in X\setminus\operatorname{int}U_1}
       d(g_b(x),y_0)>0.
\]
Relative smooth approximation, fixed on the whole closed set \(U_1\)
and on a neighborhood of \(x_0\), gives a smooth based representative
\(g_{\mathrm{sm}}\) whose error on
\(X\setminus\operatorname{int}U_1\) is less than \(d_0/2\).
One explicit relative approximation is obtained by embedding $S^3$ in $\mathbb R^4$, smoothing the coordinate functions by convolution in finitely many source charts and a partition of unity, and interpolating this approximation with $g_b$ using a smooth cutoff equal to one on the protected neighborhoods and supported where $g_b$ is smooth. Radial projection $v\mapsto v/\|v\|$ returns the result to $S^3$. A sufficiently small uniform error keeps $v$ nonzero, fixes the protected values exactly, and has the required target error; short target geodesics give the relative homotopy. It therefore has the same inverse image and differential on \(U_1\),
and \(d(g_{\mathrm{sm}}(x),y_0)>d_0/2\) on its complement. Thus there
is no uncontrolled intervening region and no additional inverse image.
Make the approximation also closer than a fixed target injectivity
radius; its short geodesics give the based homotopy.

\item The relative Pontryagin construction now applies to the explicit
framed nullbordism. It produces a based homotopy to a map missing
\(y_0\); stereographic contraction of \(S^3\setminus\{y_0\}\cong\mathbb R^3\)
to the image of \(*\) is based and contracts that last map.
For completeness, the dependence on the framed inverse image in this
argument is precisely the local product-neighborhood construction and
the interpolation lemma of \cite[\S7, Theorem B and Lemma 4,
pp.~43, 48--50]{CuspNullMilnor1965}. Their relative versions fix the source
basepoint because all tubular modifications are away from it and
stereographic interpolation between equal values keeps that value
fixed. A framed cylinder realizes a framing homotopy; the empty
Pontryagin manifold realizes a map missing the chosen regular value.
This proves the claimed based null-homotopy of the exact map \(g\).
\end{enumerate}
\end{proof}

\begin{corollary}\label{wb:cusp_s12_carrier_homotopy:L567}
\StatementRecord{wb:cusp_s12_carrier_homotopy:L567}{proved}{Historical events: GPT-6 Astra (OpenAI), ChatGPT Pro mode, effort Pro, S6\_09\_PAPERS, 17 November 2026: targeted proof check. Also checked in this writing. Pre-continuation writing events: GPT-6 Astra (OpenAI), ChatGPT Pro mode, effort Pro, S6\_10\_PAPERS, 18 November 2026: targeted proof check. Continuation metadata events: GPT-6 Astra (OpenAI), ChatGPT Pro mode, effort Pro, S6\_10\_PAPERS continuation, 19 November 2026: used and rechecked in this run; status-vocabulary reconciliation; metadata reconciliation only; no new mathematical check. \emph{Pre-continuation recorded proof scope.} GPT-6 Astra (OpenAI, ChatGPT Pro mode), S6\_09\_PAPERS, effort Pro, 2026-11-17; also checked in this writing by GPT-6 Astra (OpenAI), ChatGPT Pro mode, effort Pro, 18 November 2026}{the hypotheses, notation and definitions stated here; the explicit constructions and calculations in the proof}

Under the original fixed based identification \(X\cong S^6\), the
one-copy coefficient is
\[
 [g]=0\in\pi_6(S^3)\cong\mathbb Z/12,\qquad a=0\pmod{12},\qquad
 [E^6g]=0\in\pi_{12}(S^9)\cong\mathbb Z/24.
\]
The previously defined homomorphism
\(\Theta_g:\mathbb Z/12\to\pi_6(S^3)\), \(\Theta_g([n])=n[g]\), and
its composite with the marked peripheral character are zero.
\end{corollary}


\begin{proof}
\begin{enumerate}[label=\arabic*.,leftmargin=*]
\item The theorem states that the actual based map is null-homotopic.
Precomposition with the fixed based identification preserves a based
null-homotopy. Suspension preserves it as well. Thus the unique
coefficient of \([g]\) in any generator of \(\mathbb Z/12\) is zero.
For every \([n]\), \(\Theta_g([n])=n\cdot0=0\), and every
composite with this homomorphism is zero.
\end{enumerate}
\end{proof}

\paragraph{Exact scope.}
The proof computes this sphere-valued postcomposition. It does not
compute the homotopy class of the earlier conical-target map
\(q_2:X\to\operatorname{Susp}(S^1\times\overline T^2)\), and it
does not assert a null-homotopy for that different target. Point-set
surjectivity of \(g\) and all of its stated metric restrictions remain
properties of the original map: a homotopy-class computation does not
alter the defining formula \eqref{wb:cusp-null:eq:g}.
The retained-angle companion is proved separately in
Theorem~\ref{maps:src:B.3.1}: $p_3q_2:X\to S^4$ represents
$\eta_4\eta_5\ne0$, so $q_2$ is nonnull. Its Hopf smash
represents $12\nu_6\ne0$ in $\pi_9(S^6)$. These positive
computations use the actual four-framed inverse torus of that different
postcomposition; they do not reverse or contradict the preceding
null-homotopy through $\widehat p_2$.




\subsection{Fixed-target metric realization and its sharp dimensional scope}
\label{wb:sphere-metric-limit-note}

This note concerns the codomains of the three based maps
\[
 q_2^{\wedge 2}:X^{\wedge 2}\cong S^{12}
   \longrightarrow Z_{12}:=\mathcal A^{\wedge 2},
\]
\[
 q_{1,\mathrm{cusp}}^{\wedge 4}:X^{\wedge 4}\cong S^{24}
   \longrightarrow Z_{24}^{\mathcal B}:=\mathcal B^{\wedge 4},
 \qquad
 q_1^{\wedge 4}:X^{\wedge 4}\cong S^{24}
   \longrightarrow Z_{24}^{\mathcal C}:=\mathcal C^{\wedge 4}.
\]
In particular, \(S^{12}\) and \(S^{24}\) are the sources, not the spaces
in which the metric limits occur.  We construct one fixed compatible
length metric on each codomain which realizes the native
four-dimensional flat limits on the target slices.  The realization is
sharp: after the same rescaling, the full ambient limits have dimensions
\(8\) and \(12\), and the collapsed smash base point goes to infinity.

Put
\[
 r_T=e^{-2\pi T},\qquad
 J=\begin{pmatrix}0&1\\-1&0\end{pmatrix},\qquad
 Y_T=TJ+E_T,
\]
where \(E_T\) is the bounded matrix in the radial cusp expansion.  Choose
and fix \(T_0\) and \(C>0\) so that
\(\|E_T\|_{\mathrm{op}}\leq C\) for every \(T\geq T_0\), in the fixed
Euclidean coordinates used below.  This notation avoids confusing the
error matrix with the quotient circle \(C_T\) used below.  On the
fixed torus
\[
 \overline T^2=(\overline\Lambda\otimes\mathbb R)/\overline\Lambda
              \cong\mathbb R^2/\mathbb Z^2
\]
let
\[
 h_T:=\iota_T^*g_{Q_T}=dx^{\mathsf T}Y_T^{\mathsf T}Y_Tdx.
\]
Thus
\begin{equation}
 T^{-2}h_T\longrightarrow h_\infty:=dx^{\mathsf T}dx
 \label{wb:sphere-metric-normalized-Q-limit}
\end{equation}
in every \(C^k\) norm in the fixed flat coordinates.  Indeed, the cusp
coefficients are holomorphic in \(t=e^{-2\pi T}\), so their \(T\)-derivatives
are exponentially small, while \(T^{-1}Y_T=J+T^{-1}E_T\).

Use normalized angular coordinates
\[
 u,\theta\in\mathbb R/\mathbb Z;
\]
\(u\) is the argument of the base coordinate and \(\theta\) is the value of
the invariant character.  Increase this same \(T_0\) once more, if
necessary, so that all radial cusp formulas used below hold for
\(T\geq T_0\) and \(e^{-2\pi T_0}<\epsilon\), and put
\[
 r_0=e^{-2\pi T_0}.
\]
For \(0<r\leq r_0\), write
\[
 T(r):=-\frac{1}{2\pi}\log r.
\]

\paragraph{Target-slice notation.}
We use the embeddings
\(j_T,k_T^{\mathcal B},k_T^{\mathcal C}\) defined in
\cref{wb:cusp-cone-target-slice-embeddings}; the normalized angle \(u=0\) used
below is the same positive radial base angle denoted by \(1\in S^1\) there.

\begin{theorem}[Topology-forced collapse for every fixed compatible metric]
\label{wb:sphere-metric-uniform-collapse}
\StatementRecord{wb:sphere-metric-uniform-collapse}{proved}{Historical events: GPT-6 Astra (OpenAI), ChatGPT Pro mode, effort Pro, S6\_09\_PAPERS, 17 November 2026: targeted proof check. Also checked in this writing. Pre-continuation writing events: GPT-6 Astra (OpenAI), ChatGPT Pro mode, effort Pro, S6\_10\_PAPERS, 18 November 2026: targeted proof check. Continuation metadata events: GPT-6 Astra (OpenAI), ChatGPT Pro mode, effort Pro, S6\_10\_PAPERS continuation, 19 November 2026: used and rechecked in this run; status-vocabulary reconciliation; metadata reconciliation only; no new mathematical check. \emph{Pre-continuation recorded proof scope.} /root, effort Pro, 2026-11-17; also checked in this writing by GPT-6 Astra (OpenAI), ChatGPT Pro mode, effort Pro, 18 November 2026}{the hypotheses, notation and definitions stated here; the explicit constructions and calculations in the proof}

Let \(D_{12},D_{24}^{\mathcal B},D_{24}^{\mathcal C}\) be arbitrary fixed
metrics compatible with the respective compact target topologies, and let
\(*\) denote the corresponding smash base point.  Then
\begin{align}
 d_H^{D_{12}}\bigl(j_T(Q_T^2),\{*\}\bigr)&\longrightarrow0,\\
 d_H^{D_{24}^{\mathcal B}}
  \bigl(k_T^{\mathcal B}(C_T^4),\{*\}\bigr)&\longrightarrow0,\\
 d_H^{D_{24}^{\mathcal C}}
  \bigl(k_T^{\mathcal C}(C_T^4),\{*\}\bigr)&\longrightarrow0.
 \label{wb:sphere-metric-Hausdorff-collapse}
\end{align}
Consequently, with the metrics restricted from any one fixed ambient
target, all three target slices converge in Gromov--Hausdorff topology to
a point, not to \(\mathbb R^4\).
\end{theorem}



\begin{proof}
\begin{enumerate}[label=\arabic*.,leftmargin=*]
\item We prove the first assertion; the other two are identical.  Let
\[
 G:[0,r_0]^2\times(T^3)^2\longrightarrow Z_{12}
\]
be the product of the two inner-cone quotient maps followed by the smash
quotient.  It is continuous and
\(G(0,0,z_1,z_2)=*\) for every \((z_1,z_2)\in(T^3)^2\).  Since the domain
is compact, uniform continuity with respect to \(D_{12}\) gives
\[
 \sup_{(z_1,z_2)\in(T^3)^2}
 D_{12}\bigl(G(r,r,z_1,z_2),*\bigr)\longrightarrow0.
\]
The set \(j_T(Q_T^2)\) is contained in such a cross-section, which proves
the first Hausdorff convergence.  Replacing \(T^3\) by \(T^2\) and taking
four factors proves the other two.  Hausdorff convergence to a singleton
implies Gromov--Hausdorff convergence to that singleton.
\end{enumerate}
\end{proof}

The preceding theorem is unrescaled.  It does not prevent a blow-up of one
fixed metric.  The following construction gives the strongest natural
positive statement.

\begin{theorem}[One fixed cusp-adapted metric realizes the vertical
\(\mathbb R^4\) limits]
\label{wb:sphere-metric-fixed-realization}
\StatementRecord{wb:sphere-metric-fixed-realization}{proved}{Historical events: GPT-6 Astra (OpenAI), ChatGPT Pro mode, effort Pro, S6\_09\_PAPERS, 17 November 2026: targeted proof check. Also checked in this writing. Pre-continuation writing events: GPT-6 Astra (OpenAI), ChatGPT Pro mode, effort Pro, S6\_10\_PAPERS, 18 November 2026: targeted proof check. Continuation metadata events: GPT-6 Astra (OpenAI), ChatGPT Pro mode, effort Pro, S6\_10\_PAPERS continuation, 19 November 2026: used and rechecked in this run; status-vocabulary reconciliation; metadata reconciliation only; no new mathematical check. \emph{Pre-continuation recorded proof scope.} /root, effort Pro, 2026-11-17; also checked in this writing by GPT-6 Astra (OpenAI), ChatGPT Pro mode, effort Pro, 18 November 2026}{the hypotheses, notation and definitions stated here; \cref{wb:sphere-metric-normalized-Q-limit}}

There are fixed compatible length metrics
\[
 D_{12}\text{ on }Z_{12},\qquad
 D_{24}^{\mathcal B}\text{ on }Z_{24}^{\mathcal B},\qquad
 D_{24}^{\mathcal C}\text{ on }Z_{24}^{\mathcal C}
\]
with the following properties.

On the inner cone of \(\mathcal A\), the metric is induced by
\begin{equation}
 g_{\mathcal A}
 =dr^2+r^2du^2+\frac{r^2}{T(r)^2}h_{T(r)}.
 \label{wb:sphere-metric-A-fixed-metric}
\end{equation}
On the cusp cone of either \(\mathcal B\) or \(\mathcal C\), it is induced
by
\begin{equation}
 g_{\mathcal D}=dr^2+r^2(du^2+d\theta^2),
 \qquad\mathcal D\in\{\mathcal B,\mathcal C\}.
 \label{wb:sphere-metric-BC-fixed-metric}
\end{equation}
The metrics on the smash powers are the metric quotients of the product
length metrics after the closed fat wedge is collapsed.  Write
\(g_{Z_{12}}\) and \(g_{Z_{24}^{\mathcal D}}\) for the resulting smooth
Riemannian tensors on the open regular strata.

Let
\[
 p_T^{12}:=j_T(0,0),\qquad
 p_T^{24,\mathcal D}:=k_T^{\mathcal D}(0,0,0,0).
\]
Put
\[
 \Sigma_T^{12}:=j_T(Q_T^2),\qquad
 \Sigma_T^{24,\mathcal D}:=k_T^{\mathcal D}(C_T^4).
\]
Let \(d_{T,12}^{\mathrm{slice}}\) and
\(d_{T,24,\mathcal D}^{\mathrm{slice}}\) be the intrinsic Riemannian path
metrics of the first fundamental forms on these embedded slices.  These
must be distinguished from the restricted ambient distance functions
\[
 d_{T,12}^{\mathrm{amb}}
 :=D_{12}|_{\Sigma_T^{12}\times\Sigma_T^{12}},\qquad
 d_{T,24,\mathcal D}^{\mathrm{amb}}
 :=D_{24}^{\mathcal D}|_{\Sigma_T^{24,\mathcal D}
                         \times\Sigma_T^{24,\mathcal D}}.
\]
In general the latter can be smaller, because an ambient path is allowed
to leave the slice.

Then the induced first fundamental forms on the target slices satisfy the
exact identities
\begin{equation}
 \left(\frac{T}{r_T}\right)^2
  j_T^*g_{Z_{12}}
 =g_{Q_T}\oplus g_{Q_T},
 \label{wb:sphere-metric-A-slice-exact}
\end{equation}
and
\begin{equation}
 \left(\frac{\ell_T}{r_T}\right)^2
  (k_T^{\mathcal D})^*g_{Z_{24}^{\mathcal D}}
 =g_{C_T}^{\oplus4}.
 \label{wb:sphere-metric-BC-slice-exact}
\end{equation}
Consequently,
\begin{align}
 \left(\Sigma_T^{12},\frac{T}{r_T}d_{T,12}^{\mathrm{slice}},
       p_T^{12}\right)
 &\longrightarrow(\mathbb R^4,0),\\
 \left(\Sigma_T^{24,\mathcal D},
       \frac{\ell_T}{r_T}d_{T,24,\mathcal D}^{\mathrm{slice}},
       p_T^{24,\mathcal D}\right)
 &\longrightarrow(\mathbb R^4,0)
 \label{wb:sphere-metric-slice-fixed-blowups}
\end{align}
in pointed local-isometric and pointed Gromov--Hausdorff topology, and
smoothly on the exhausting Euclidean slice balls.

The full ambient pointed blow-ups are instead
\begin{align}
 \left(Z_{12},\frac{T}{r_T}D_{12},p_T^{12}\right)
 &\longrightarrow(\mathbb R^8,0),\\
 \left(Z_{24}^{\mathcal D},
       \frac{\ell_T}{r_T}D_{24}^{\mathcal D},
       p_T^{24,\mathcal D}\right)
 &\longrightarrow(\mathbb R^{12},0).
 \label{wb:sphere-metric-ambient-fixed-blowups}
\end{align}
as pointed metric spaces in pointed Gromov--Hausdorff topology.  More
precisely, for every fixed radius the corresponding rescaled target ball is
eventually contained in an open regular-stratum coordinate neighborhood,
and the pulled-back Riemannian tensors converge there in every \(C^k\) norm.
Thus smooth convergence is asserted on an exhaustion by regular-stratum
neighborhoods, not globally on the singular compact targets.  Under these
limits, the slice \(\mathbb R^4\) is respectively the coordinate subspace
obtained by fixing the two radial and two base-angle coordinates, or by
fixing the four radial and four base-angle coordinates.
\end{theorem}



\begin{proof}
\begin{enumerate}[label=\arabic*.,leftmargin=*]
\item Equation \eqref{wb:sphere-metric-normalized-Q-limit} shows that the metrics on
the compact cross-sections of \eqref{wb:sphere-metric-A-fixed-metric} remain
uniformly bounded and converge to a positive definite flat metric.
Therefore the cross-section diameter is \(O(r)\).  Moreover the radial
coordinate is \(1\)-Lipschitz and a radial segment from level \(r\) to the
vertex has length \(r\).  The length completion consequently adds exactly
one point at \(r=0\), and its topology is the prescribed cone topology.
The same assertions are immediate for
\eqref{wb:sphere-metric-BC-fixed-metric}.  Interpolation on a closed annulus
away from the vertex extends these metrics to fixed compatible length
metrics on the compact spaces \(\mathcal A,\mathcal B,\mathcal C\).

\item If a compact metric space \((E,d)\) and a closed subset \(W\subset E\) are
given, the metric quotient which collapses \(W\) to one point is
\[
 d_{E/W}([x],[y])
 =\min\{d(x,y),d(x,W)+d(y,W)\}.
\]
Applying this to the fat wedges in the product length spaces gives the
claimed fixed compatible metrics on the smash targets.

\item On the \(r_T\)-slice of \(\mathcal A\), with \(u=0\), the first fundamental
form is
\[
 \frac{r_T^2}{T^2}h_T.
\]
The definition \(h_T=\iota_T^*g_{Q_T}\) proves
\eqref{wb:sphere-metric-A-slice-exact}.  On the \(r_T\)-slice of
\(\mathcal B\) or \(\mathcal C\), with \(u=0\), the form is
\(r_T^2d\theta^2\).  Under
\(\vartheta_T:[\theta]\mapsto[\ell_T\theta]\), the native quotient-circle
metric is \(\ell_T^2d\theta^2\).  This proves
\eqref{wb:sphere-metric-BC-slice-exact}.  The native systole estimates
\[
 \operatorname{sys}(Q_T^2)\geq T-C,
 \qquad
 \operatorname{sys}(C_T^4)=\ell_T
\]
then prove \eqref{wb:sphere-metric-slice-fixed-blowups}: every prescribed
radius ball is eventually exactly a Euclidean four-ball.

\item It remains to compute the ambient limits.  In one \(\mathcal A\)-factor,
put
\[
 \lambda_T=\frac{T}{r_T},\qquad
 s=\lambda_T(r-r_T),\qquad
 \eta=Tu,\qquad z=Y_Tx
\]
on the universal cover of the torus directions.  On every fixed bounded
set in \((s,\eta,z)\), one has
\[
 \frac r{r_T}=1+\frac sT,
 \qquad
 T(r)=T+O(T^{-1}),
 \qquad
 Y_{T(r)}Y_T^{-1}=I+O(T^{-2}),
\]
with the same estimates after any fixed number of derivatives.  Thus
\(\lambda_T^2g_{\mathcal A}\) converges in these coordinates to
\[
 ds^2+d\eta^2+dz_1^2+dz_2^2.
\]
The inward radial boundary is at distance \(T\), the \(u\)-period is \(T\),
and the \(z\)-period lattice is \(Y_T\mathbb Z^2\), whose systole is at
least \(T-C\).  The outward edge of the fixed cusp chart is also at scaled
distance
\[
 \frac{T}{r_T}(r_0-r_T)
 =T\left(\frac{r_0}{r_T}-1\right)\longrightarrow\infty.
\]
Thus the inward vertex, all periods, and the outward edge recede to
infinity.  A single factor converges on exhausting regular-stratum
coordinate neighborhoods to \(\mathbb R^4\), and two factors converge to
\(\mathbb R^8\).

\item For one \(\mathcal B\)- or \(\mathcal C\)-factor, put
\[
 \lambda_T=\frac{\ell_T}{r_T},\qquad
 s=\lambda_T(r-r_T),\qquad
 \eta=\ell_Tu,\qquad \zeta=\ell_T\theta.
\]
Then
\[
 \lambda_T^2g_{\mathcal D}
 =ds^2+\left(\frac r{r_T}\right)^2
       (d\eta^2+d\zeta^2),
\]
which converges smoothly on bounded coordinate sets to the Euclidean metric
on \(\mathbb R^3\).  The two circle periods and the inward radial boundary
are all at scaled distance \(\ell_T\to\infty\), while the outward edge is
at scaled distance
\[
 \frac{\ell_T}{r_T}(r_0-r_T)
 =\ell_T\left(\frac{r_0}{r_T}-1\right)\longrightarrow\infty.
\]
Four factors therefore converge on exhausting regular-stratum coordinate
neighborhoods to \(\mathbb R^{12}\).

\item Finally, a level-\(r_T\) factor point has distance exactly \(r_T\) from its
cone vertex.  Indeed, the radial segment has that length.  A competing
curve contained in the inner chart has length at least \(r_T\), while one
which leaves has inner-chart portions of total length at least
\((r_0-r_T)+r_0>r_T\).  The Riemannian-\(\ell^2\)-product distance from a
synchronized point to the fat wedge is the minimum of the factor
distances, and hence is also \(r_T\).

\item For the chosen scale put
\[
 b_T=\lambda_Tr_T,\qquad a_T=\lambda_T(r_0-r_T),
\]
and let \(\sigma_T\) be the least scaled period in a factor.  Thus
\(\sigma_T\geq T-C\) in the \(\mathcal A\) calculation and
\(\sigma_T=\ell_T\) in the \(\mathcal B,\mathcal C\) calculation.  The
preceding computations give
\[
 a_T\longrightarrow\infty,\qquad
 b_T\longrightarrow\infty,\qquad
 \sigma_T\longrightarrow\infty.
\]
Fix \(R\).  A curve of rescaled length at most \(R\) from \(p_T\) meets
neither the fat wedge nor \(r=r_0\) for all large \(T\), so the
radius-\(R\) target ball lies in the open regular stratum.  If \(x,y\) lie
in that ball, a competitor through the fat wedge has length at least
\[
 (b_T-R)+(b_T-R)>2R,
\]
and a competitor which exits the inner chart and returns has length at
least
\[
 2(a_T-R)>2R.
\]
Every remaining competitor of length at most \(2R\) stays in
\(r/r_T=1+o_R(1)\).  After lifting it to the angular and torus covers, a
nonzero deck translation has length at least
\((1-o_R(1))(\sigma_T-2R)>2R\) for all large \(T\).  Since the path through
\(p_T\) has length at most \(2R\), none of these three kinds of global
shortcut computes the distance between \(x\) and \(y\).  Thus the metric
on every fixed target ball is computed entirely in one lifted
regular-stratum chart.  The \(C^\infty\) convergence of the pulled-back
tensors on these exhausting charts implies the pointed
Gromov--Hausdorff limits in
\eqref{wb:sphere-metric-ambient-fixed-blowups}.
\end{enumerate}
\end{proof}

\begin{theorem}[Complete intrinsic scalar-rescaling trichotomy]
\label{wb:sphere-metric-scaling-trichotomy}
\StatementRecord{wb:sphere-metric-scaling-trichotomy}{proved}{Historical events: GPT-6 Astra (OpenAI), ChatGPT Pro mode, effort Pro, S6\_09\_PAPERS, 17 November 2026: targeted proof check. Also checked in this writing. Pre-continuation writing events: GPT-6 Astra (OpenAI), ChatGPT Pro mode, effort Pro, S6\_10\_PAPERS, 18 November 2026: targeted proof check. Continuation metadata events: GPT-6 Astra (OpenAI), ChatGPT Pro mode, effort Pro, S6\_10\_PAPERS continuation, 19 November 2026: used and rechecked in this run; status-vocabulary reconciliation; metadata reconciliation only; no new mathematical check. \emph{Pre-continuation recorded proof scope.} /root, effort Pro, 2026-11-17; also checked in this writing by GPT-6 Astra (OpenAI), ChatGPT Pro mode, effort Pro, 18 November 2026}{the hypotheses, notation and definitions stated here; \cref{wb:sphere-metric-fixed-realization}; \cref{wb:sphere-metric-outward-shortcut-exclusion}}

For the fixed metrics of
\cref{wb:sphere-metric-fixed-realization}, let \(\lambda_T>0\) be arbitrary
and put
\[
 \beta_T:=\lambda_Tr_T.
\]
Choose one of the following triples, retaining the indicated ambient
metric:
\[
 \begin{split}
  (\Sigma_T,d_T^{\mathrm{slice}},d_T^{\mathrm{amb}},p_T;D)
  &=(\Sigma_T^{12},d_{T,12}^{\mathrm{slice}},
      d_{T,12}^{\mathrm{amb}},p_T^{12};D_{12}),\\
  (\Sigma_T,d_T^{\mathrm{slice}},d_T^{\mathrm{amb}},p_T;D)
  &=(\Sigma_T^{24,\mathcal D},
      d_{T,24,\mathcal D}^{\mathrm{slice}},
      d_{T,24,\mathcal D}^{\mathrm{amb}},
      p_T^{24,\mathcal D};D_{24}^{\mathcal D}),
      \qquad \mathcal D\in\{\mathcal B,\mathcal C\}.
 \end{split}
\]
After passing to a subsequence so that \(\beta_T\) has a limit in
\([0,\infty]\), the intrinsic pointed slice limit is
\begin{equation}
 (\Sigma_T,\lambda_Td_T^{\mathrm{slice}},p_T)
 \longrightarrow
 \begin{cases}
  \{\mathrm{pt}\},&\beta_T\longrightarrow0,\\
  ((\mathbb R/L\mathbb Z)^4,0),
      &\beta_T\longrightarrow L\in(0,\infty),\\
  (\mathbb R^4,0),&\beta_T\longrightarrow\infty.
 \end{cases}
 \label{wb:sphere-metric-scaling-trichotomy-result}
\end{equation}
The convergence in
\eqref{wb:sphere-metric-scaling-trichotomy-result} is pointed
Gromov--Hausdorff in all three cases.  In the finite positive case the flat
metrics converge smoothly under the fixed marked-torus identifications; in
the unrolling case the convergence is pointed local-isometric and smooth on
the exhausting Euclidean balls.
In the middle line the lattice is orthogonally equivalent to the square
lattice \(L\mathbb Z^4\).  At the same time,
\begin{equation}
 \lambda_TD(p_T,*)=\beta_T
 \label{wb:sphere-metric-vertex-distance}
\end{equation}
for every one of the displayed choices.

Only in the unrolling regime \(\beta_T\to\infty\) do we compare the
intrinsic metric with the restricted ambient distance.  In that regime, for
every \(R\in(0,\infty)\),
\begin{equation}
 \sup_{x,y\in
  \overline B_{\lambda_Td_T^{\mathrm{slice}}}(p_T,R)}
 \left|
  \lambda_Td_T^{\mathrm{amb}}(x,y)
  -\lambda_Td_T^{\mathrm{slice}}(x,y)
 \right|\longrightarrow0.
 \label{wb:sphere-metric-intrinsic-ambient-agreement}
\end{equation}
Moreover, there are \(\varepsilon_{T,R}\to0\) such that
\begin{equation}
 \overline B_{\lambda_Td_T^{\mathrm{slice}}}(p_T,R)
 \subset
 \overline B_{\lambda_Td_T^{\mathrm{amb}}}(p_T,R)
 \subset
 \overline B_{\lambda_Td_T^{\mathrm{slice}}}
             (p_T,(1+\varepsilon_{T,R})R).
 \label{wb:sphere-metric-intrinsic-ambient-ball-comparison}
\end{equation}
Thus the restricted ambient metric has the same pointed
\(\mathbb R^4\) limit on these fixed slice balls.  No agreement is asserted
in the regimes \(\beta_T\to0\) or \(\beta_T\to L\in(0,\infty)\), where
radial or fat-wedge shortcuts need not disappear.  Every rescaling which
unrolls the intrinsic slice sends the cone/fat-wedge base point to infinity,
so no such limit is a tangent cone based at that point.

In the same regime \(\beta_T\to\infty\), the corresponding full ambient
bubbles are
\begin{align}
 (Z_{12},\lambda_TD_{12},p_T^{12})
 &\longrightarrow(\mathbb R^8,0),\\
 (Z_{24}^{\mathcal D},\lambda_TD_{24}^{\mathcal D},
  p_T^{24,\mathcal D})
 &\longrightarrow(\mathbb R^{12},0),
 \qquad\mathcal D\in\{\mathcal B,\mathcal C\}.
 \label{wb:sphere-metric-arbitrary-unrolling-ambient-bubbles}
\end{align}
These are pointed Gromov--Hausdorff limits of the singular metric targets.
The convergence of tensors is \(C^\infty\) only under coordinate
identifications on exhausting open regular-stratum neighborhoods; no
global smooth convergence of the singular targets is claimed.
\end{theorem}



\begin{proof}
\begin{enumerate}[label=\arabic*.,leftmargin=*]
\item For the \(Q_T^2\) route, the rescaled period matrix in each two-dimensional
factor is
\[
 \frac{\lambda_Tr_T}{T}Y_T
 =\beta_T\left(J+\frac{E_T}{T}\right).
\]
If \(\beta_T\to L\in(0,\infty)\), this lattice converges to
\(LJ\mathbb Z^2\); two factors give the square four-torus in
\eqref{wb:sphere-metric-scaling-trichotomy-result}.  If
\(\beta_T\to0\), the lengths of a lattice basis tend to zero, so the
diameter tends to zero.  If \(\beta_T\to\infty\), the reverse triangle
inequality gives systole
\[
 \frac{\beta_T}{T}(T-C)
 =\beta_T\left(1-\frac CT\right)\longrightarrow\infty,
\]
so every fixed-radius ball is eventually Euclidean.

\item For a \(C_T^4\) slice, the four periods after rescaling have the common
length
\[
 \lambda_Tr_T=\beta_T,
\]
which gives the same three intrinsic cases directly.

\item Consider first one factor.  A curve from a level-\(r_T\) point to the cone
vertex which remains in the inner chart has length at least \(r_T\), because
the \(dr^2\)-term makes \(r\) \(1\)-Lipschitz there.  If it leaves the
chart, it has a first crossing of \(r=r_0\) and, before ending at the
vertex, a last re-entry through \(r=r_0\); its two inner-chart portions
have total length at least
\((r_0-r_T)+r_0>r_T\).  The radial segment has length \(r_T\), so the
factor distance to its based vertex is exactly \(r_T\).  The
Riemannian-\(\ell^2\)-product distance from every synchronized slice point
to the fat wedge is the minimum of these factor distances and hence is
also \(r_T\).  The metric-quotient formula in the preceding proof therefore
gives \eqref{wb:sphere-metric-vertex-distance}.

\item It remains to prove the restricted-ambient comparison, so assume
\(\beta_T\to\infty\) and fix \(R\).  Two points in
\(\overline B_{\lambda_Td_T^{\mathrm{slice}}}(p_T,R)\) have intrinsic
distance at most \(2R/\lambda_T\).  The slice path is an ambient competitor,
so
\[
 d_T^{\mathrm{amb}}(x,y)\leq d_T^{\mathrm{slice}}(x,y).
\]
The quotient shortcut through the fat wedge has rescaled length
\(2\lambda_Tr_T=2\beta_T\), and hence is longer than \(2R\) for all large
\(T\).  A path which leaves the fixed cusp chart must reach \(r=r_0\).
Its outward cost is controlled by
\begin{equation}
 \lambda_T(r_0-r_T)
 =\beta_T\left(\frac{r_0}{r_T}-1\right)\longrightarrow\infty,
 \label{wb:sphere-metric-outward-shortcut-exclusion}
\end{equation}
so no path of rescaled length at most \(2R\) can use the interpolating
annulus or the global part of a factor.

\item Consider now an ambient competitor of length at most
\(2R/\lambda_T+o(\lambda_T^{-1})\) which stays in the cusp chart.  In every
factor its radial coordinate satisfies
\[
 |r-r_T|\leq \frac{2R+o(1)}{\lambda_T},
 \qquad
 \frac r{r_T}=1+O_R(\beta_T^{-1}).
\]
On this tube, the horizontal tensors in the directions tangent to the slice
are uniformly \((1+\varepsilon_{T,R})\)-bilipschitz to their values at
\(r=r_T\), where \(\varepsilon_{T,R}\to0\).  For
\(\mathcal B\) and \(\mathcal C\) this follows directly from
\((r/r_T)^2=1+O_R(\beta_T^{-1})\).  For \(\mathcal A\), the same conclusion
follows, without changing the fixed coordinates or the matrix \(Y_T\), from
\[
 \begin{aligned}
  T(r)-T&=O_R(\beta_T^{-1}),\\
  \frac{T(r)}T&=1+O_R((\beta_TT)^{-1}),\\
  Y_{T(r)}Y_T^{-1}&=I+O_R((\beta_TT)^{-1}).
 \end{aligned}
\]
Projecting the competitor to the torus or character directions, and
discarding its radial and base-angle components, produces a path between
the same slice endpoints.  The bilipschitz estimate therefore gives
\[
 d_T^{\mathrm{slice}}(x,y)
 \leq(1+\varepsilon_{T,R})d_T^{\mathrm{amb}}(x,y).
\]
Together with the opposite inequality and
\(\lambda_Td_T^{\mathrm{slice}}(x,y)\leq2R\), this proves
\eqref{wb:sphere-metric-intrinsic-ambient-agreement}.  The first inclusion in
\eqref{wb:sphere-metric-intrinsic-ambient-ball-comparison} follows from
\(d_T^{\mathrm{amb}}\leq d_T^{\mathrm{slice}}\).  For a point \(x\) in
the middle ball, apply the same fat-wedge and outward exclusions to a
near-minimizing ambient path from \(p_T\) to \(x\), whose rescaled length is
at most \(R+o(1)\), and project that path to the slice directions.  This
gives
\[
 \lambda_Td_T^{\mathrm{slice}}(p_T,x)
 \leq(1+\varepsilon_{T,R})R
\]
after enlarging \(\varepsilon_{T,R}\) by a quantity tending to zero, and
proves the second inclusion.

\item For completeness, the full ambient bubbles at an arbitrary unrolling scale
are obtained without changing any coordinate or orientation.  In one
\(\mathcal A\)-factor use
\[
 s=\lambda_T(r-r_T),\qquad
 \eta=\beta_Tu,\qquad
 z=\frac{\beta_T}{T}Y_Tx.
\]
Writing \(A_T(r)=Y_{T(r)}Y_T^{-1}\), one obtains the exact expression
\[
 \lambda_T^2g_{\mathcal A}
 =ds^2+\left(\frac r{r_T}\right)^2d\eta^2
  +\left(\frac r{r_T}\right)^2
   \left(\frac{T}{T(r)}\right)^2
   dz^{\mathsf T}A_T(r)^{\mathsf T}A_T(r)dz.
\]
On every fixed bounded coordinate set,
\[
 \frac r{r_T}=1+O_R(\beta_T^{-1}),\qquad
 T(r)-T=O_R(\beta_T^{-1}),\qquad
 A_T(r)=I+O_R((\beta_TT)^{-1}),
\]
with the same estimates after every fixed number of derivatives.  Hence
the tensor converges smoothly to the Euclidean metric on \(\mathbb R^4\).
The \(\eta\)-period is \(\beta_T\), and the \(z\)-period lattice is
\[
 \frac{\beta_T}{T}Y_T\mathbb Z^2,
 \qquad
 \operatorname{sys}\left(
   \frac{\beta_T}{T}Y_T\mathbb Z^2\right)
 \geq\beta_T\left(1-\frac CT\right)\longrightarrow\infty.
\]
In one \(\mathcal B\)- or \(\mathcal C\)-factor use
\[
 s=\lambda_T(r-r_T),\qquad
 \eta=\beta_Tu,\qquad \zeta=\beta_T\theta.
\]
Then
\[
 \lambda_T^2g_{\mathcal D}
 =ds^2+\left(\frac r{r_T}\right)^2
       (d\eta^2+d\zeta^2),
\]
and both angular periods are \(\beta_T\to\infty\).

\item In all three cases the inward boundary and the fat wedge are at scaled
distance \(\beta_T\), while a lift starting at radius \(r_T\) incurs the
one-way radial cost in
\eqref{wb:sphere-metric-outward-shortcut-exclusion} before its first outward
crossing of \(r=r_0\).  Fix \(R\).  For large \(T\), all these quantities
and all nonzero period lengths exceed \(4R\).
If \(x,y\) lie in the ambient rescaled radius-\(R\) ball, their distances
from the fat wedge are at least \(\beta_T-R\); hence a quotient path through
the smash base point has length at least \(2(\beta_T-R)>2R\).  Every
remaining competitor lifts to the complement of the fat wedge in the
product.  If such a lift exits the inner product chart, then some factor
has a first crossing of \(r=r_0\); the radial variation before that crossing
is at least \(\lambda_T(r_0-r_T)-R\), and returning to the lift of \(y\)
incurs the same lower bound.  Its length is therefore greater than \(2R\).
Finally, a path of length at most \(2R\) remains in the tube
\(r/r_T=1+O_R(\beta_T^{-1})\); there every nonzero deck translation has
length tending to infinity, so such a path cannot wrap around an angular
or torus period.  The displayed coordinate neighborhoods therefore exhaust
the pointed balls without quotient or global shortcuts.
Their smooth tensor convergence proves the pointed Gromov--Hausdorff
limits in
\eqref{wb:sphere-metric-arbitrary-unrolling-ambient-bubbles}, with smooth
convergence confined to this regular-stratum exhaustion.
\end{enumerate}
\end{proof}

\begin{corollary}[Unit-volume normalization does not give
\(\mathbb R^4\)]
\label{wb:sphere-metric-unit-volume}
\StatementRecord{wb:sphere-metric-unit-volume}{proved}{Historical events: GPT-6 Astra (OpenAI), ChatGPT Pro mode, effort Pro, S6\_09\_PAPERS, 17 November 2026: targeted proof check. Also checked in this writing. Pre-continuation writing events: GPT-6 Astra (OpenAI), ChatGPT Pro mode, effort Pro, S6\_10\_PAPERS, 18 November 2026: targeted proof check. Continuation metadata events: GPT-6 Astra (OpenAI), ChatGPT Pro mode, effort Pro, S6\_10\_PAPERS continuation, 19 November 2026: used and rechecked in this run; status-vocabulary reconciliation; metadata reconciliation only; no new mathematical check. \emph{Pre-continuation recorded proof scope.} /root, effort Pro, 2026-11-17; also checked in this writing by GPT-6 Astra (OpenAI), ChatGPT Pro mode, effort Pro, 18 November 2026}{the hypotheses, notation and definitions stated here; the argument and qualifications immediately following}

The four-volume of \(Q_T^2\) is \((\det Y_T)^2\sim T^4\).  Its unit-volume
length factor is therefore \((\det Y_T)^{-1/2}\sim T^{-1}\), and its
unit-volume limit is a square flat four-torus, not \(\mathbb R^4\).
Likewise, the four-volume of \(C_T^4\) is \(\ell_T^4\), and its unit-volume
normalization is exactly the square unit four-torus.
\end{corollary}



\begin{proposition}[What vertex tangents and compact asymptotic cones can
mean]
\label{wb:sphere-metric-tangent-scope}
\StatementRecord{wb:sphere-metric-tangent-scope}{proved}{Historical events: GPT-6 Astra (OpenAI), ChatGPT Pro mode, effort Pro, S6\_09\_PAPERS, 17 November 2026: targeted proof check. Also checked in this writing. Pre-continuation writing events: GPT-6 Astra (OpenAI), ChatGPT Pro mode, effort Pro, S6\_10\_PAPERS, 18 November 2026: targeted proof check. Continuation metadata events: GPT-6 Astra (OpenAI), ChatGPT Pro mode, effort Pro, S6\_10\_PAPERS continuation, 19 November 2026: used and rechecked in this run; status-vocabulary reconciliation; metadata reconciliation only; no new mathematical check. \emph{Pre-continuation recorded proof scope.} /root, effort Pro, 2026-11-17; also checked in this writing by GPT-6 Astra (OpenAI), ChatGPT Pro mode, effort Pro, 18 November 2026}{the hypotheses, notation and definitions stated here; \cref{wb:sphere-metric-fixed-realization}; \cref{wb:sphere-metric-C-boundary-quotient-homeomorphism}; \cref{wb:cusp-cone-S24-route}; \cref{wb:sphere-metric-boundary-link-local-shift}}

The compact target topology alone determines no metric tangent cone.
For an exact conic length metric built from fixed piecewise-Riemannian
metrics on the links, however, the vertex links have topological and
Hausdorff dimensions \(7\) for \(Z_{12}\) and \(11\) for both
\(Z_{24}\) targets.  The resulting vertex cones have topological and
Hausdorff dimensions \(8\) and \(12\), and are not \(\mathbb R^4\).

More precisely, writing \(L_3=T^3\) and \(L_2=T^2\), one has the based
topological models
\[
 Z_{12}\cong\operatorname{Susp}(L_3*L_3),
 \qquad
 Z_{24}^{\mathcal B}
 \cong\operatorname{Susp}(L_2*L_2*L_2*L_2).
\]
At the suspension base point their conic links are respectively
\(L_3*L_3\) and \(L_2^{*4}\).  Also
\[
 (\mathcal C,c_0)
 \cong
 (\mathcal M/\partial\mathcal M,[\partial\mathcal M]),
 \qquad
 (Z_{24}^{\mathcal C},*)
 \cong
 (\mathcal M^4/\partial(\mathcal M^4),
  [\partial(\mathcal M^4)]),
\]
so an exact conic metric at the collapsed boundary has link
\(\partial(\mathcal M^4)\).  At the collapsed point of
\(Z_{24}^{\mathcal C}\), the complete integral local-homology calculation is
\begin{equation}
 H_i\bigl(Z_{24}^{\mathcal C},
          Z_{24}^{\mathcal C}\setminus\{*\};\mathbb Z\bigr)
 \cong
 \begin{cases}
  \mathbb Z^{\binom4{i-1}},&2\leq i\leq5,\\
  0,&i=6,7,\\
  \mathbb Z^{\binom4{i-8}},&8\leq i\leq12,\\
  0,&\text{otherwise}.
 \end{cases}
 \label{wb:sphere-metric-C-smash-local-homology}
\end{equation}
In particular, its local homology is not concentrated in top degree.  For
the first two targets the top local homology occurs in degree \(8\) and
\(12\), respectively; for the third target
\eqref{wb:sphere-metric-C-smash-local-homology} includes a top
\(\mathbb Z\) in degree \(12\).  None has the local homology of
\(\mathbb R^4\).

Finally, the standard asymptotic cone obtained by blowing down any fixed
compact metric target is a point.  Indeed, if \(R_n\to\infty\), then
\[
 \operatorname{diam}(Z,R_n^{-1}D)
 =R_n^{-1}\operatorname{diam}(Z,D)\longrightarrow0.
\]
Thus the limits in \cref{wb:sphere-metric-fixed-realization} are
moving-base-point blow-ups, often called bubbles.  They are neither
vertex-based tangent cones nor asymptotic cones of the compact targets.
\end{proposition}



\begin{proof}
\begin{enumerate}[label=\arabic*.,leftmargin=*]
\item For compact spaces \(V_i\) with closed subspaces \(A_i\), there is a
canonical homeomorphism
\[
 (V_1/A_1)\wedge\cdots\wedge(V_m/A_m)
 \cong
 \frac{V_1\times\cdots\times V_m}
      {\bigcup_iV_1\times\cdots\times A_i\times\cdots\times V_m}.
\]
First reverse the suspension coordinate, $r\mapsto\epsilon-r$, so that the chosen inner basepoint is carried to the outer suspension vertex. Take $V_i=\operatorname{Cone}(L_i)$ and $A_i=L_i$, its outer boundary.
The actual map from the product cones to the cone on the iterated join uses radial coordinate $\rho=\max_i r_i$ and join weights $t_i=r_i/\sum_jr_j$ away from the common vertex. Its inverse is $r_i=\rho t_i/\max_jt_j$; zero weights discard exactly the irrelevant cone-link coordinates. These maps extend continuously at the vertex, and the union of outer faces goes to $\rho=1$.
Collapsing that link gives its suspension.  This proves the first two
models and the link dimensions
\[
 \dim(T^3*T^3)=3+3+1=7,
 \qquad
 \dim((T^2)^{*4})=4\cdot2+3=11.
\]

\item We next prove the based homeomorphism for \(\mathcal C\), rather than
identifying the two presentations without a map.  Put
\(L=\partial\mathcal M\), write the attached cone as
\[
 \operatorname{Cone}(L)
 =L\times[0,1]/(L\times\{0\}),
\]
with base \(L\times\{1\}\), and choose a collar
\(\kappa:L\times[0,1]\hookrightarrow\mathcal M\) satisfying
\(\kappa(x,0)=x\).  There is a homeomorphism
\[
 \Phi:\mathcal C\setminus\{c_0\}
       \longrightarrow\operatorname{int}\mathcal M
\]
given on the punctured cone, the inner collar, and the remaining core by
\[
 \Phi([x,t])=\kappa(x,t/e)\quad(0<t\leq1),\qquad
 \Phi(\kappa(x,s))=\kappa(x,e^{s-1})\quad(0\leq s\leq1),
\]
\[
 \Phi(q)=q
 \quad\text{for}\quad
 q\in\mathcal M\setminus\kappa(L\times[0,1)).
\]
The first two formulas agree at the glued points
\([x,1]=\kappa(x,0)\), and the second and third agree at
\(\kappa(x,1)\).  Their collar-coordinate images are respectively
\(L\times(0,e^{-1}]\) and \(L\times[e^{-1},1]\); hence the displayed
piecewise map and its piecewise inverse are continuous.

\item The compact Hausdorff spaces \(\mathcal C\) and
\(\mathcal M/L\) are both one-point compactifications of
\(\operatorname{int}\mathcal M\): for \(\mathcal C\) use \(\Phi\), and for
\(\mathcal M/L\) use the quotient map on
\(\mathcal M\setminus L=\operatorname{int}\mathcal M\).  Uniqueness of the
one-point compactification therefore extends \(\Phi\) to a based
homeomorphism
\begin{equation}
 (\mathcal C,c_0)
 \xrightarrow{\ \cong\ }
 (\mathcal M/\partial\mathcal M,[\partial\mathcal M]).
 \label{wb:sphere-metric-C-boundary-quotient-homeomorphism}
\end{equation}
Only after establishing
\eqref{wb:sphere-metric-C-boundary-quotient-homeomorphism} do we apply the
quotient identity at the start of the proof.  It gives
\[
 \mathcal C^{\wedge4}
 \cong
 \frac{\mathcal M^4}
 {\displaystyle\bigcup_{a=1}^4
  \mathcal M^{a-1}\times\partial\mathcal M
  \times\mathcal M^{4-a}}
 =\mathcal M^4/\partial(\mathcal M^4),
\]
and carries the smash base point to the collapsed boundary.

\item Put \(N=\mathcal M^4\).  Rounding its product corners preserves the
homeomorphism type of the pair \((N,\partial N)\) and makes
\(\partial N\) an eleven-dimensional closed manifold.  A boundary collar
descends in \(N/\partial N\) to a cone neighborhood of the collapsed
point, so the local link is exactly \(\partial N\) and
\begin{equation}
 H_i(N/\partial N,(N/\partial N)\setminus\{*\};\mathbb Z)
 \cong\widetilde H_{i-1}(\partial N;\mathbb Z).
 \label{wb:sphere-metric-boundary-link-local-shift}
\end{equation}

\item The integral homology calculation in the proof of
\cref{wb:cusp-cone-S24-route} for the specific compact orientable
three-manifold \(\mathcal M\) gives
\[
 H_j(\mathcal M;\mathbb Z)\cong
 \begin{cases}
  \mathbb Z,&j=0,1,\\
  0,&\text{otherwise}.
 \end{cases}
\]
All groups are free, so the integral Kunneth theorem gives
\[
 H_j(N;\mathbb Z)\cong
 \begin{cases}
  \mathbb Z^{\binom4j},&0\leq j\leq4,\\
  0,&\text{otherwise}.
 \end{cases}
\]
Poincare--Lefschetz duality, with the product orientation retained, gives
\[
 H_j(N,\partial N;\mathbb Z)
 \cong H^{12-j}(N;\mathbb Z)
 \cong
 \begin{cases}
  \mathbb Z^{\binom4{j-8}},&8\leq j\leq12,\\
  0,&\text{otherwise}.
 \end{cases}
\]
The long exact sequence of \((N,\partial N)\) now has disjoint absolute
and relative degree ranges.  It gives
\[
 \widetilde H_j(\partial N;\mathbb Z)\cong
 \begin{cases}
  \mathbb Z^{\binom4j},&1\leq j\leq4,\\
  0,&j=5,6,\\
  \mathbb Z^{\binom4{j-7}},&7\leq j\leq11,\\
  0,&\text{otherwise}.
 \end{cases}
\]
Substitution in
\eqref{wb:sphere-metric-boundary-link-local-shift} proves
\eqref{wb:sphere-metric-C-smash-local-homology}.

\item To see explicitly why topology alone does not fix the tangent metric, compare the cone tensors $dr^2+r^2h$ and $dr^2+4r^2h$ on the same link of dimension $d$. The identity is a homeomorphism, while their unit-ball Hausdorff volumes differ by the factor $2^d$, so the marked metric cones are not isometric. An exact metric cone is invariant under every radial blow-up at its vertex,
and its Hausdorff dimension is one plus that of its fixed
piecewise-Riemannian link.  For the first target,
\[
 H_8\bigl(\operatorname{Cone}(T^3*T^3),T^3*T^3;\mathbb Z\bigr)
 \cong\mathbb Z;
\]
the analogous top local group for the \(\mathcal B\) target is
\(\mathbb Z\) in degree \(12\), while the full \(\mathcal C\) calculation
is \eqref{wb:sphere-metric-C-smash-local-homology}.  Euclidean
\(\mathbb R^4\) has local homology \(\mathbb Z\) only in degree \(4\), so
none of these conic vertex models is \(\mathbb R^4\).  The diameter
calculation proves the final assertion.
\end{enumerate}
\end{proof}

\paragraph{Consequences for the three sphere routes.}
The exact positive statement is that one fixed metric on each compact
codomain contains the native \(\mathbb R^4\) decompactification as a
vertical four-plane in a moving-point Euclidean bubble.  The whole bubbles
are \(\mathbb R^8\) for the \(S^{12}\) route and \(\mathbb R^{12}\) for the
\(S^{24}\) routes.  The factor \(\beta_T=\lambda_Tr_T\) proves that this
dimension discrepancy cannot be removed by a different scalar rescaling
of the same fixed cusp-adapted metric: the slice unrolls only when the
smash base point has already gone to infinite scaled distance.

No metric on a source sphere is supplied by these quotient constructions.
If an independent smooth Riemannian metric is placed on \(S^{12}\) or
\(S^{24}\), its ordinary tangent spaces are \(\mathbb R^{12}\) or
\(\mathbb R^{24}\).  Pulling a codomain metric back through
\(q_2^{\wedge2}\) or \(q_1^{\wedge4}\), on the regular strata where the
maps are differentiable, gives a degenerate quadratic form along the
positive-dimensional quotient kernels, not a Riemannian source metric.

The moving four-sphere carriers
\[
 \mathbb S_{12,T}^4=S_T^2\wedge S_T^2\subset S^6,\qquad
 \mathbb S_{24,T}^4=C_T^{\wedge4}\subset S^8
\]
of \cref{wb:cusp-cone-angle-forgotten-sphere-carriers} carry separately
assigned metrics.  Namely, \(S_T^2\) first receives the quotient length
metric obtained from
\((\overline T^2,h_T)\cong(Q_T,g_{Q_T})\) by collapsing
\(K^{(1)}\subset\overline T^2\), equivalently by collapsing
\(\iota_T(K^{(1)})\subset Q_T\).  Then
\(\mathbb S_{12,T}^4\) receives the quotient length metric of the
Riemannian \(\ell^2\)-product \((S_T^2)^2\) after its fat wedge is
collapsed.  Likewise, \(\mathbb S_{24,T}^4\) receives the quotient length
metric of the Riemannian \(\ell^2\)-product \(C_T^4\) after its fat wedge is
collapsed.  These metrics are assigned directly to the moving carriers;
they are unrescaled intrinsic quotient metrics, not restrictions of
\(D_{12}\) or \(D_{24}^{\mathcal B}\), whose domains are the conical
targets, and no fixed ambient metric on \(S^6\) or \(S^8\) is used to
define them.  Their Euclidean balls are
based at the midpoint classes
\[
 \begin{aligned}
  y_*&=[(1/2,1/2)]\in\overline T^2,\qquad z_T=c(y_*),\\
  \zeta_{12,T}&=[z_T,z_T],\\
  \zeta_{24,T}
  &=[(\ell_T/2,\ell_T/2,\ell_T/2,\ell_T/2)].
 \end{aligned}
\]
not at the images of \(p_T^{12}\) or \(p_T^{24,\mathcal B}\), which belong
to the subsequently collapsed loci.  In particular, the trichotomy in
\cref{wb:sphere-metric-scaling-trichotomy} is the trichotomy of the
uncollapsed intrinsic flat-torus slices; it is not a finite-\(\beta_T\)
description of the singular quotient-metric moving four-spheres.  Indeed,
the restriction of any one compatible metric on the compact sphere
\(S^6\) or \(S^8\) has diameter bounded by the diameter of that sphere.  By
contrast, in the separately assigned moving metrics the midpoint distances
to the collapsed loci are at least \((T-C)/2\) and exactly
\(\ell_T/2\), respectively, and hence tend to infinity.  Thus the assigned
unrescaled metrics cannot themselves be restrictions of one fixed ambient
metric.  This does not preclude a separately proved realization after a
scalar blow-up of another fixed quotient metric.

Finally, the rank obstruction for a full-base torus quotient must retain
its actual metric hypothesis: \emph{quotient systole tends to infinity}
means the quotient norm induced by the unrescaled Euclidean fibre metric.
Without that qualification the assertion is false, since an arbitrary
\(T\)-dependent scalar multiplication can make the systole of any fixed
nonzero quotient tend to infinity.

\subsection{Expanded proof of the retained-angle cusp class}
\label{wb:h24:ret:section}

\paragraph{The exact object.}
We use the fixed based smooth identification of the family member $X$ with
$S^6$, its cusp filling, and the complete original toric coordinates of
\cref{mapsI:src:A.1.2}. The native quotient and compact extension are
\cref{wb:cusp-cone-native-descent} and \cref{wb:cusp-cone-S12-route}.
The local parameter $t$ in the following formulas is the cusp parameter
$q$ of the family convention; $B_0=J$ in the radial metric notation.
All original period terms and native lattices are retained. The
calculation is independent of the angle-forgetting map to $S^3$.

\paragraph{Original coordinates and the specified sphere-valued map}

The ordered quotient lattice is
\[
 \overline\Lambda=\mathbb Z\langle\overline{\widehat\gamma},
              \overline{\widehat u}\rangle,
 \qquad
 K_\Lambda=\mathbb Z\langle\widehat w,\widehat\delta\rangle,
 \qquad
 \overline T^2=\mathbb R^2/\mathbb Z^2.
\]
In these ordered bases the original period data are
\begin{equation}
 B_0=\begin{pmatrix}0&1\\-1&0\end{pmatrix},\qquad
 Z(s)=sB_0+C(t),\qquad t=e^{2\pi i s},\qquad
 C(t)=\begin{pmatrix}6\mu(t)&h_{\rm per}(t)\\
 b(t)-h_{\rm per}(t)&\mu(t)\end{pmatrix}.
 \label{wb:h24:ret:eq:original-periods}
\end{equation}
Here \(b=\beta+\tau\), \(s=\tau-h_{\rm per}(t)\), and every entry of
\(C\) is holomorphic at \(t=0\), exactly as in the original source.
The subscript in \(h_{\rm per}\) distinguishes that period function
from the sphere-valued map below.  Write
\begin{equation}
 R(t)=-2\pi\operatorname{Im}C(t),\qquad
 \ell=\log|t|,\qquad A(t)=\ell B_0+R(t).
 \label{wb:h24:ret:eq:A}
\end{equation}
Choose a closed cusp disc \(|t|\leq\epsilon\) inside the original
coordinate disc and small enough for the quotient construction.  On a
slightly larger closed disc put \(M=\sup\|R(t)\|<\infty\).  We will
choose a regular value with \(0<r_0<\epsilon\) and
\(|\log r_0|>M\).  This chooses a point of the already specified map;
it does not change that map or the cutoff radius \(\epsilon\).

The original covering model on its dense torus is
\begin{equation}
 x_j=e^{2\pi i\zeta_j},\qquad
 \Psi_n(x,t)=
 \bigl(e^{2\pi i C(t)n}t^{B_0n}x,t\bigr),\qquad n\in\mathbb Z^2.
 \label{wb:h24:ret:eq:deck}
\end{equation}
For a branch of \(s\), writing
\(\zeta=Z(s)u+k\) with \(u,k\in\mathbb R^2\) gives
\begin{equation}
 \log|x|=A(t)u.
 \label{wb:h24:ret:eq:logx}
\end{equation}
Indeed, \(\operatorname{Im}s=-\ell/(2\pi)\), so
\(-2\pi\operatorname{Im}(Z(s)u+k)=\ell B_0u+R(t)u\).
Since \(B_0\) is orthogonal,
\begin{equation}
 \bigl|\bigl(\ell B_0+vR(t)\bigr)\xi\bigr|
 \geq (|\ell|-vM)|\xi|>0
 \quad(0\leq v\leq1,\ \xi\ne0)
 \label{wb:h24:ret:eq:invertibility}
\end{equation}
whenever \(|\ell|>M\).  Thus \(A(t)\) is invertible there.  Equation
\eqref{wb:h24:ret:eq:deck} changes \(A(t)^{-1}\log|x|\) by \(n\), so the native
quotient is exactly
\begin{equation}
 \overline\pi_f([x,t])=[A(t)^{-1}\log|x|]\in\mathbb R^2/\mathbb Z^2.
 \label{wb:h24:ret:eq:native}
\end{equation}
This also follows from the flat marking in \cref{wb:cusp-cone-native-descent}.

Let
\[
 T^3=S^1_\phi\times\overline T^2,
 \qquad S^1_\phi=\mathbb R/(2\pi\mathbb Z),
\]
with coordinate basepoints \(\phi=0\), \(u_1=0\), \(u_2=0\).
Its coordinate two-skeleton is the following specified closed subset:
\begin{equation}
 K^{(2)}=
 \bigl(\{0\}\times\overline T^2\bigr)
 \cup\bigl(S^1_\phi\times\{u_1=0\}\bigr)
 \cup\bigl(S^1_\phi\times\{u_2=0\}\bigr).
 \label{wb:h24:ret:eq:two-skeleton}
\end{equation}
The product CW structure has one open three-cell, with coordinates
\((\phi,u_1,u_2)\in(0,2\pi)\times(0,1)^2\).  Consequently
\[
 c_3:T^3\longrightarrow T^3/K^{(2)}\cong S^3
\]
is the degree-one quotient, where the orientation of \(S^3\) is fixed
by the displayed coordinate order.  One explicit identification of its
open cell with \(\mathbb R^3\) is
\[
 (\phi,u_1,u_2)\longmapsto
 \bigl(-\cot(\phi/2),-\cot(\pi u_1),-\cot(\pi u_2)\bigr).
\]
Each derivative in its own variable is strictly positive, and approach
to any boundary face leaves every compact subset of \(\mathbb R^3\).
Extending by its point at infinity gives the stated based sphere
identification and verifies its continuity on the collapsed skeleton.

The exact conical target and map are
\begin{align}
 \mathcal A&=
 \frac{[0,\epsilon]\times T^3}
 {\{0\}\times T^3\sim v_0,\quad
  \{\epsilon\}\times T^3\sim v_1},
 \label{wb:h24:ret:eq:A-target}\\
 q_2(x)&=[r,e^{i\phi},\overline\pi_f(x)]
       &&(t(f(x))=re^{i\phi},\ 0<r<\epsilon),\notag\\
 q_2(W)&=v_0,\qquad
 q_2\bigl(X\setminus\operatorname{int}N_0\bigr)=v_1.
 \label{wb:h24:ret:eq:q2}
\end{align}
Fix the same source basepoint \(x_{\rm base}\in W\).  Define
\begin{equation}
 p_3:\mathcal A\longrightarrow\Sigma_r(S^3,*),\qquad
 p_3([r,z])=[r,c_3(z)],\qquad
 H=p_3q_2:(X,x_{\rm base})\longrightarrow(S^4,*).
 \label{wb:h24:ret:eq:H}
\end{equation}
The reduced suspension here collapses both endpoint copies of \(S^3\)
and the entire interval over \(*\) to one point.  Thus \(p_3\) sends
both \(v_0\) and \(v_1\) to \(*\) and is well-defined and based.
It is continuous by the defining quotient topology.  An explicit
identification of its open four-cell with \(\mathbb R^4\) is
\begin{equation}
 (r,\phi,u_1,u_2)\longmapsto
 \bigl(-\cot(\pi r/\epsilon),-\cot(\phi/2),
       -\cot(\pi u_1),-\cot(\pi u_2)\bigr).
 \label{wb:h24:ret:eq:four-cell}
\end{equation}
This fixes the target orientation.  Replacing \(r\) locally by
\(\ell=\log r\) preserves this orientation, since \(dr/d\ell=r>0\).
The ordered coordinates in \eqref{wb:h24:ret:eq:four-cell} specify the sphere
identification used throughout this calculation.

\paragraph{The statement proved.} We give the complete expanded proof of \cref{maps:src:B.3.1}, including the full differential, the relative smoothing and the spin bordism calculation.


\paragraph{The entire regular inverse image lies in one original chart}

Choose the actual lift
\begin{equation}
 u_0=(-1/4,1/4)^{\mathsf T},\qquad
 p_0=B_0u_0=(1/4,1/4)^{\mathsf T},\qquad \phi_0=\pi.
 \label{wb:h24:ret:eq:chosen-value}
\end{equation}
Modulo \(\mathbb Z^2\), \(u_0\) is \((3/4,1/4)\) in the open cell.  Put
\(t_0=r_0e^{i\phi_0}\) and
\(y_0=[r_0,c_3(e^{i\phi_0},[u_0])]\).  The numbers \(r_0\) and
\(t_0\) will remain nonzero throughout the framing calculation.

The original triangle \(\{0,e_1,e_2\}\) has toric chart
\begin{equation}
 \mathcal U_\sigma\cong\mathbb C^3,\qquad
 (z_0,z_1,z_2)=(t/(x_1x_2),x_1,x_2),\qquad
 t=z_0z_1z_2.
 \label{wb:h24:ret:eq:z-chart}
\end{equation}
These characters are dual to \((0,0,1),(1,0,1),(0,1,1)\), respectively,
as direct pairing verifies.  This is the chart of
\cref{mapsI:src:A.1.2}, with its original coordinates.

For \(j=1,2\), introduce the nowhere vanishing holomorphic functions
\begin{equation}
 h_j(t)=\exp\bigl(-2\pi i(C(t)u_0)_j\bigr).
 \label{wb:h24:ret:eq:units}
\end{equation}
The following is an exact coordinate map, with the exact inverse
displayed alongside it:
\begin{align}
 F(z)&=(w_0,w_1,w_2)
 =\bigl(z_0/(h_1(t)h_2(t)),\ h_1(t)z_1,\ h_2(t)z_2\bigr),
 \label{wb:h24:ret:eq:F}\\
 F^{-1}(w)&=
 \bigl(w_0h_1(t)h_2(t),\ w_1/h_1(t),\ w_2/h_2(t)\bigr),
 \quad t=w_0w_1w_2.
 \label{wb:h24:ret:eq:Finv}
\end{align}
Both maps are holomorphic wherever \(|t|\) is in the period disc,
both preserve \(t\), and substitution proves that both composites are
the identity.  In particular the maps extend across the origin and
across every coordinate divisor.  No period matrix is replaced by a
different matrix in this coordinate change.

\begin{lemma}\label{wb:h24:ret:chart-lemma}
\StatementRecord{wb:h24:ret:chart-lemma}{proved}{Historical events: GPT-6 Astra (OpenAI), ChatGPT Pro mode, effort Pro, S6\_09\_PAPERS, 17 November 2026: targeted proof check. Also checked in this writing. Pre-continuation writing events: GPT-6 Astra (OpenAI), ChatGPT Pro mode, effort Pro, S6\_10\_PAPERS, 18 November 2026: targeted proof check. Continuation metadata events: GPT-6 Astra (OpenAI), ChatGPT Pro mode, effort Pro, S6\_10\_PAPERS continuation, 19 November 2026: used and rechecked in this run; status-vocabulary reconciliation; metadata reconciliation only; no new mathematical check. \emph{Pre-continuation recorded proof scope.} /root, effort Pro, 2026-11-17; also checked in this writing by GPT-6 Astra (OpenAI), ChatGPT Pro mode, effort Pro, 18 November 2026}{the hypotheses, notation and definitions stated here; \cref{wb:h24:ret:eq:H}; \cref{wb:h24:ret:eq:q2}; \cref{wb:h24:ret:eq:deck}; \cref{wb:h24:ret:eq:z-chart}}

For sufficiently small \(r_0>0\), the whole inverse image
\(M=H^{-1}(y_0)\) is represented injectively inside one contractible
coordinate neighborhood in \(X\).  In its \(w\)-coordinates it is
exactly the embedded two-torus
\begin{equation}
 w(a,b)=
 \bigl(\rho_0e^{i(\phi_0-a-b)},\rho_1e^{ia},\rho_2e^{ib}\bigr),
 \qquad
 \rho_0=r_0^{1/2},\quad \rho_1=\rho_2=r_0^{1/4},
 \quad (a,b)\in\bigl(\mathbb R/(2\pi\mathbb Z)\bigr)^2.
 \label{wb:h24:ret:eq:M}
\end{equation}
\end{lemma}


\begin{proof}
\begin{enumerate}[label=\arabic*.,leftmargin=*]
\item The point \(y_0\) lies off the collapsed subset of the target, so
\eqref{wb:h24:ret:eq:H} and \eqref{wb:h24:ret:eq:q2} imply
\[
 M=\{x:t(f(x))=t_0,\ \overline\pi_f(x)=[u_0]\}.
\]
Any lift of such a point to \((\mathbb C^*)^2\times\{t_0\}\) has
\(A(t_0)^{-1}\log|x|=u_0+n\) for one \(n\in\mathbb Z^2\).
Applying the actual transformation \(\Psi_{-n}\) from
\eqref{wb:h24:ret:eq:deck} produces a lift with
\(\log|x|=A(t_0)u_0\).  If two such lifts represent the same point,
\(\Psi_n\) relates them and changes their lifted quotient coordinate
by \(n\); hence \(n=0\), and the lifts are identical.  This proves
both completeness and injectivity of this chosen representative.

\item Since
\(\log|h_j(t)|=2\pi\operatorname{Im}(C(t)u_0)_j
=-(R(t)u_0)_j\), its \(w\)-coordinates satisfy
\[
 (\log|w_1|,\log|w_2|)^{\mathsf T}
 =A(t_0)u_0-R(t_0)u_0=\ell_0p_0,
 \qquad \ell_0=\log r_0.
\]
Thus \(|w_1|=|w_2|=r_0^{1/4}\), and the unchanged product
\(w_0w_1w_2=t_0\) forces \(|w_0|=r_0^{1/2}\) and the displayed
phase.  Conversely every pair \((a,b)\) in \eqref{wb:h24:ret:eq:M} gives one
such lift.  The original kernel coordinates \(k_j\) had period one
in \(x_j=e^{2\pi i\zeta_j}\); the phase coordinates \(a,b\) have
period \(2\pi\), with no cover or change of lattice.

\item The covering \(Y_\epsilon\to N_0\) is injective on a neighborhood of
the origin in \eqref{wb:h24:ret:eq:z-chart}, because the original action is free
and properly discontinuous.  Choose a ball \(V\) centered at that
origin with closure inside such a neighborhood and the period disc.
Choose a ball \(V_w\) about the origin with closure contained in
\(F(V)\).  Every coordinate in \eqref{wb:h24:ret:eq:M} tends uniformly to zero
as \(r_0\downarrow0\).  Therefore for sufficiently small \(r_0\)
the entire torus lies in \(V_w\).  The inverse image of \(V_w\)
under \(F\), followed by the injective covering projection, is the
required contractible chart in \(X\).
\end{enumerate}
\end{proof}

\paragraph{The full differential and its ordered normal framing}

We calculate in a neighborhood of the compact torus \eqref{wb:h24:ret:eq:M}.
The quotient coordinates admit a unique local lift near \(u_0\).
Set
\[
 v=(\log|w_1|,\log|w_2|)^{\mathsf T},\quad
 \ell=\log|w_0w_1w_2|,\quad
 \phi=\arg(w_0w_1w_2)
\]
using the branch of \(\phi\) near \(\phi_0=\pi\).
Equations \eqref{wb:h24:ret:eq:logx} and \eqref{wb:h24:ret:eq:F} give the exact identity
\begin{equation}
 v=\ell B_0u+R(t)(u-u_0),\qquad
 u=u_0+A(t)^{-1}(v-\ell p_0).
 \label{wb:h24:ret:eq:exact-u}
\end{equation}
Differentiate this identity before restricting to the inverse image:
\begin{equation}
 du=-A(t)^{-1}(dA(t))A(t)^{-1}(v-\ell p_0)
        +A(t)^{-1}(dv-p_0\,d\ell).
 \label{wb:h24:ret:eq:full-derivative}
\end{equation}
On \(M\) the vector \(v-\ell p_0\) is exactly zero.  Therefore
\begin{equation}
 du|_M=A_0^{-1}(dv-p_0\,d\ell),\qquad A_0=A(t_0).
 \label{wb:h24:ret:eq:derivative-on-M}
\end{equation}
This explicitly accounts for all derivatives of the original bounded
period terms; none are omitted or estimated away.

Identify \(T_w\mathbb C^3\) with \(\mathbb C^3\) in the given coordinate chart,
and define real tangent vectors
\begin{equation}
 L_j=(0,\ldots,w_j,\ldots,0),\qquad
 Q_j=(0,\ldots,iw_j,\ldots,0)\qquad(0\leq j\leq2).
 \label{wb:h24:ret:eq:LQ}
\end{equation}
Thus \(d\log|w_i|(L_j)=\delta_{ij}\),
\(d\arg w_i(Q_j)=\delta_{ij}\), and the crossed evaluations vanish.
The ordered tangent vectors of the parametrization \eqref{wb:h24:ret:eq:M} are
\begin{equation}
 T_a=-Q_0+Q_1,\qquad T_b=-Q_0+Q_2.
 \label{wb:h24:ret:eq:tangents}
\end{equation}
They are independent because their \(w_1\) and \(w_2\) components,
respectively, are nonzero.  For the target order
\((\ell,\phi,u_1,u_2)\), define the following four lifts of its
coordinate vectors:
\begin{align}
 N_\ell&=\tfrac12L_0+\tfrac14L_1+\tfrac14L_2,
       &N_\phi&=Q_0,\label{wb:h24:ret:eq:normal-first}\\
 N_i&=-\bigl((A_0)_{1i}+(A_0)_{2i}\bigr)L_0
           +(A_0)_{1i}L_1+(A_0)_{2i}L_2
       &&(i=1,2).\label{wb:h24:ret:eq:normal-u}
\end{align}
Let \(P=(\ell,\phi,u_1,u_2)\).  Evaluation using
\eqref{wb:h24:ret:eq:derivative-on-M} gives the entire differential table
\begin{equation}
\begin{array}{c|rrrrrr}
 &N_\ell&N_\phi&N_1&N_2&T_a&T_b\\\hline
 d\ell&1&0&0&0&0&0\\
 d\phi&0&1&0&0&0&0\\
 du_1&0&0&1&0&0&0\\
 du_2&0&0&0&1&0&0
\end{array}.
 \label{wb:h24:ret:eq:differential-table}
\end{equation}
In particular \(P\), and therefore \(H\) in its open target cell,
is a submersion along \(M\).  The classes of
\((N_\ell,N_\phi,N_1,N_2)\) in \(TX|_M/TM\) give the specified
normal four-framing.  Orthogonal representatives are not required for
this quotient-bundle definition; projecting the displayed lifts along
\(TM\) to any chosen orthogonal complement preserves the quotient
frame.

If the normal frame is specified using the literal four real coordinates
in \eqref{wb:h24:ret:eq:four-cell}, their differential at the chosen value is
\[
 dH_{\rm cell}=K\,dP,\qquad
 K=\operatorname{diag}\!\left(
 \frac{\pi r_0}{\epsilon}\csc^2\frac{\pi r_0}{\epsilon},
 \frac12,2\pi,2\pi\right).
\]
The first factor includes \(dr=r_0\,d\ell\), the second uses
\(\phi_0=\pi\), and the final two use the actual representatives
\((3/4,1/4)\) of \(u_0\). Thus, writing
\(\mathcal N=(N_\ell,N_\phi,N_1,N_2)\), the literal dual frame is
\(\mathcal N K^{-1}\). All four entries of \(K\) are positive
constants on \(M\). The explicitly positive diagonal matrices
\((1-v)K^{-1}+vI_4\), \(0\leq v\leq1\), give a homotopy
from that literal frame to \(\mathcal N\). This retains the original
target coordinates and establishes that either coordinate description
gives the same framed bordism class.

For completeness, the original period matrix admits the following
explicit homotopy of this normal framing.  Put
\begin{equation}
 A_{0,v}=\ell_0B_0+vR(t_0),\qquad 0\leq v\leq1,
 \label{wb:h24:ret:eq:A-homotopy}
\end{equation}
and replace \(A_0\) by \(A_{0,v}\) only in
\eqref{wb:h24:ret:eq:normal-u}.  Equation \eqref{wb:h24:ret:eq:invertibility} proves
invertibility for every \(v\).  Their images under the actual
\((du_1,du_2)\) are the columns of \(A_0^{-1}A_{0,v}\), so these
are normal frames throughout the homotopy.  This is a homotopy of
the computed frame on the same embedded torus, not a replacement of
the original degeneration by a different one.

\begin{lemma}[The exact orientation sign]\label{wb:h24:ret:orientation-lemma}
\StatementRecord{wb:h24:ret:orientation-lemma}{proved}{Historical events: GPT-6 Astra (OpenAI), ChatGPT Pro mode, effort Pro, S6\_09\_PAPERS, 17 November 2026: targeted proof check. Also checked in this writing. Pre-continuation writing events: GPT-6 Astra (OpenAI), ChatGPT Pro mode, effort Pro, S6\_10\_PAPERS, 18 November 2026: targeted proof check. Continuation metadata events: GPT-6 Astra (OpenAI), ChatGPT Pro mode, effort Pro, S6\_10\_PAPERS continuation, 19 November 2026: used and rechecked in this run; status-vocabulary reconciliation; metadata reconciliation only; no new mathematical check. \emph{Pre-continuation recorded proof scope.} /root, effort Pro, 2026-11-17; also checked in this writing by GPT-6 Astra (OpenAI), ChatGPT Pro mode, effort Pro, 18 November 2026}{the hypotheses, notation and definitions stated here; \cref{wb:h24:ret:eq:normal-u}; \cref{wb:h24:ret:eq:A-homotopy}}

Use the complex real-coordinate orientation, whose ordered coordinates are
\[
 (\operatorname{Re}w_0,\operatorname{Im}w_0,
  \operatorname{Re}w_1,\operatorname{Im}w_1,
  \operatorname{Re}w_2,\operatorname{Im}w_2).
\]
Then one has
\begin{equation}
 \det(N_\ell,N_\phi,N_1,N_2,T_a,T_b)
       =-\det(A_0)|w_0w_1w_2|^2<0.
 \label{wb:h24:ret:eq:sign}
\end{equation}
Hence the tangent orientation induced by the ordered normal frame is
the orientation of \((-T_a,T_b)\).
\end{lemma}


\begin{proof}
\begin{enumerate}[label=\arabic*.,leftmargin=*]
\item The ordered frame \((L_0,Q_0,L_1,Q_1,L_2,Q_2)\) has determinant
\(|w_0|^2|w_1|^2|w_2|^2>0\) in real coordinates.  Regroup its rows
as \((L_0,L_1,L_2,Q_0,Q_1,Q_2)\); this permutation has three
inversions and sign \(-1\).  Regroup the columns of the frame in
\eqref{wb:h24:ret:eq:sign} as
\((N_\ell,N_1,N_2,N_\phi,T_a,T_b)\); this has two inversions and
sign \(+1\).  The angular block is
\(\left(\begin{smallmatrix}1&-1&-1\\0&1&0\\0&0&1\end{smallmatrix}\right)\),
of determinant one.  The radial block has first column
\((1/2,1/4,1/4)^{\mathsf T}\) and the other columns given by
\eqref{wb:h24:ret:eq:normal-u}.  Adding its last two rows to its first row gives
first row \((1,0,0)\), with remaining lower-right block \(A_0\).
Its determinant is \(\det A_0\).  This proves the equality.
Equation \eqref{wb:h24:ret:eq:A-homotopy} joins \(A_0\) to \(\ell_0B_0\)
through invertible matrices, so
\(\operatorname{sign}\det A_0=\operatorname{sign}\det(\ell_0B_0)=+1\).
The inequality and the tangent orientation follow.
\end{enumerate}
\end{proof}

\paragraph{An explicit contraction of the full frame}

Define a complex-linear matrix, also regarded as a real-linear matrix,
\begin{equation}
 D(a,b)=\operatorname{diag}(e^{-i(a+b)},e^{ia},e^{ib})\in\operatorname{SU}(3).
 \label{wb:h24:ret:eq:D}
\end{equation}
By \eqref{wb:h24:ret:eq:M}, every \(L_j\) and \(Q_j\) at \((a,b)\) is
\(D(a,b)\) times its value at \((0,0)\).  The coefficients in
\eqref{wb:h24:ret:eq:normal-first}, \eqref{wb:h24:ret:eq:normal-u}, and
\eqref{wb:h24:ret:eq:tangents} are constant on \(M\), since \(t=t_0\).
It follows that the positively oriented full frame is exactly
\begin{equation}
 (N_\ell,N_\phi,N_1,N_2,-T_a,T_b)_{(a,b)}
       =D(a,b)G_0,
 \label{wb:h24:ret:eq:full-frame}
\end{equation}
where \(G_0\in\operatorname{GL}^+(6,\mathbb R)\) is its single value at \((0,0)\).
Thus neither an unspecified phase nor a position-dependent
real-linear matrix remains.

Here is a periodic contraction of \(D\), with all parameters
specified.  For \(v\in[0,1]\) put
\[
 s_v=\sin(\pi v/2),\qquad c_v=\cos(\pi v/2),\qquad
 Q_v(\theta)=
 \begin{pmatrix}s_ve^{-i\theta}&-c_v\\c_v&s_ve^{i\theta}\end{pmatrix},
 \qquad K_v(\theta)=Q_v(\theta)Q_v(0)^*.
\]
The two columns of \(Q_v\) have norm one and zero Hermitian inner
product, and \(\det Q_v=s_v^2+c_v^2=1\).  Thus
\(Q_v,K_v\in\operatorname{SU}(2)\).  Direct substitution gives
\[
 K_0(\theta)=I,\quad
 K_1(\theta)=\operatorname{diag}(e^{-i\theta},e^{i\theta}),\quad
 K_v(0)=I.
\]
Let \(\iota_{01},\iota_{02}:\operatorname{SU}(2)\to\operatorname{SU}(3)\) act in the
ordered coordinate planes \((0,1)\) and \((0,2)\), respectively,
and as the identity in the unused coordinate.  Then
\begin{equation}
 D_v(a,b)=\iota_{01}(K_v(a))\,\iota_{02}(K_v(b))
 \label{wb:h24:ret:eq:D-homotopy}
\end{equation}
is a well-defined continuous map on
\([0,1]\times(\mathbb R/(2\pi\mathbb Z))^2\), with
\(D_0=I\) and \(D_1=D\).  This proves the needed null-homotopy
directly, including both phase seams.  In particular
\(D_v(a,b)G_0\) contracts the comparison matrix in
\eqref{wb:h24:ret:eq:full-frame} to one constant positive matrix.

\begin{proposition}\label{wb:h24:ret:framing-proposition}
\StatementRecord{wb:h24:ret:framing-proposition}{proved}{Historical events: GPT-6 Astra (OpenAI), ChatGPT Pro mode, effort Pro, S6\_09\_PAPERS, 17 November 2026: targeted proof check. Also checked in this writing. Pre-continuation writing events: GPT-6 Astra (OpenAI), ChatGPT Pro mode, effort Pro, S6\_10\_PAPERS, 18 November 2026: targeted proof check. Continuation metadata events: GPT-6 Astra (OpenAI), ChatGPT Pro mode, effort Pro, S6\_10\_PAPERS continuation, 19 November 2026: used and rechecked in this run; status-vocabulary reconciliation; metadata reconciliation only; no new mathematical check. \emph{Pre-continuation recorded proof scope.} /root, effort Pro, 2026-11-17; also checked in this writing by GPT-6 Astra (OpenAI), ChatGPT Pro mode, effort Pro, 18 November 2026}{the hypotheses, notation and definitions stated here; \cref{wb:h24:ret:eq:sign}; \cref{wb:h24:ret:eq:full-frame}; \cref{wb:h24:ret:eq:D-homotopy}}

The stable tangential framing induced by the actual normal four-frame
of \(H^{-1}(y_0)\) agrees, up to homotopy and the orientation sign
\eqref{wb:h24:ret:eq:sign}, with the Lie framing
\((\partial_a,\partial_b)\) on the native two-torus.
\end{proposition}


\begin{proof}
\begin{enumerate}[label=\arabic*.,leftmargin=*]
\item In the coordinate chart, the stable tangent trivialization is the
isomorphism
\[
 \mathbb R^4\oplus TM\longrightarrow T\mathbb C^3|_M,
 \qquad (v,V)\longmapsto
 v_1N_\ell+v_2N_\phi+v_3N_1+v_4N_2+V.
\]
Using the intrinsic frame \((-\partial_a,\partial_b)\) on \(TM\),
its matrix is precisely \eqref{wb:h24:ret:eq:full-frame}.  The homotopy
\eqref{wb:h24:ret:eq:D-homotopy} therefore makes this comparison constant.
A constant positive matrix changes no homotopy class of stable
framing: Gram--Schmidt deforms it to an element of \(SO(6)\), and
an orthogonal matrix of determinant one can be reduced to the identity
by finitely many oriented plane rotations, each joined to the
identity by varying its angle.  This proves the asserted comparison
within the chart.

\item It remains to compare the chart with the ambient sphere framing.
Pull back the usual stable tangent trivialization
\(TS^6\oplus\mathbb R\cong\mathbb R^7\) under the fixed based diffeomorphism
\(X\cong S^6\).  On the contractible \(w\)-ball containing \(M\),
it differs from the coordinate trivialization plus one real line by
a continuous map into \(\operatorname{GL}(7,\mathbb R)\).  Contracting the ball to its
center contracts this comparison map to a constant.  Thus its
restriction to \(M\) contributes no varying stable framing class.
If the chosen sphere orientation disagrees with the complex
orientation, the remaining constant has negative determinant and
reverses the resulting framed bordism class.  Both that sign and the
explicit sign \eqref{wb:h24:ret:eq:sign} disappear in the order-two class
computed below.  This argument compares actual ambient
trivializations on a contractible chart; it does not suppose that
\(\mathbb C^3\) globally trivializes \(TX\).
\end{enumerate}
\end{proof}

\paragraph{Nonzero framed bordism and the exact unstable class}

We give the nonvanishing test for the resulting torus explicitly.
A stable framing gives a spin structure by lifting its oriented frame
to the trivial \(\operatorname{Spin}\)-bundle.  On the Lie-framed
circle this is the disconnected double cover, whereas the spin
structure induced from the boundary of a framed disc is the
connected double cover: as the boundary tangent turns once, its
\(2\pi\)-rotation lifts from \(1\) to \(-1\) in
\(\operatorname{Spin}(2)\).  We call the former nonbounding and the
latter bounding.

For a closed oriented spin surface \(S\), the spin quadratic
refinement
\[
 q_S:H_1(S;\mathbb F_2)\longrightarrow\mathbb F_2
\]
takes a class represented by an embedded circle to zero or one
according to whether the induced spin structure, using its normal
line in the surface, is bounding or nonbounding.  It satisfies
\begin{equation}
 q_S(x+y)=q_S(x)+q_S(y)+(x\mathbin{\cdot}y).
 \label{wb:h24:ret:eq:quadratic}
\end{equation}
This is the usual spin quadratic form: resolving transverse
intersections adds one tangent turn per intersection, and the
\(2\pi\)-spin-lift calculation above gives the last term.  The
well-defined correspondence between spin structures and such
quadratic refinements is given in
\cite[Theorem 3.2 and Remark 3.3]{DebrayGunningham2018Root}; its original
surface proof is \cite{Johnson1980Root}.

For a symplectic basis \((e_i,f_i)_{i=1}^g\), put
\begin{equation}
 \operatorname{Arf}(S)=\sum_{i=1}^gq_S(e_i)q_S(f_i)\in\mathbb F_2.
 \label{wb:h24:ret:eq:Arf}
\end{equation}
The basis independence can be seen without selecting a preferred
basis: factor the finite sum over the symplectic planes to obtain
\[
 \sum_{x\in H_1(S;\mathbb F_2)}(-1)^{q_S(x)}
   =\prod_{i=1}^g
     \sum_{a,b\in\mathbb F_2}(-1)^{a q_S(e_i)+b q_S(f_i)+ab}
   =2^g(-1)^{\sum_i q_S(e_i)q_S(f_i)}.
\]
The left side has no choice of basis, proving the assertion.

Here is also the bordism-invariance argument needed for the present
obstruction.  A spin bordism has a handle decomposition.  An index
zero or three handle creates or removes a spin sphere, whose first
homology and Arf invariant are zero.  An index one handle either
connects two components, in which case the first-homology quadratic
forms add orthogonally, or increases the genus of one component.
In the latter case the new symplectic pair contains the meridian
of the handle.  That meridian bounds the cocore disc with its
induced spin structure, so its quadratic value is zero; the pair
contributes zero to \eqref{wb:h24:ret:eq:Arf}, regardless of the value on its
dual.  Existing cycles can be represented outside the attaching
discs, so their quadratic values are preserved.  An index two
handle is the reverse of an index one handle.  Therefore each
handle preserves the total Arf invariant of the boundary, and a
spin boundary has Arf invariant zero.  This also proves that a
stably framed boundary has Arf invariant zero, because forgetting
the framing gives a spin bordism.  The full spin bordism statement
is recorded in \cite[Theorem 3.6]{DebrayGunningham2018Root}.

The intrinsic Lie-torus calculation is the known computation in \cite[Example~3.8, p.~6]{DebrayGunningham2018Root}; we retain its direct derivation. For the intrinsic Lie framing on
\(M=(\mathbb R/(2\pi\mathbb Z))^2\), take
\[
 e=[a\mapsto(a,0)],\qquad f=[b\mapsto(0,b)].
\]
Both the tangent and normal vectors in the surface are constant in
the Lie frame along either circle.  Thus both circles have the
nonbounding spin structure, and
\begin{equation}
 e\mathbin{\cdot}f=1,\qquad q_M(e)=q_M(f)=1,
 \qquad q_M(xe+yf)=x+xy+y,
 \qquad \operatorname{Arf}(M)=1.
 \label{wb:h24:ret:eq:actual-Arf}
\end{equation}
Reversing the first Lie tangent vector changes the orientation but
does not change these mod-two values.  In particular the
specifically framed inverse image computed above is not a framed
boundary.

To identify the element, let \(\eta\in\pi_1^{\rm S}\) denote the
stable class of the complex Hopf map.  Its Pontryagin--Thom inverse
image is the circle with its Lie framing.  This can be verified
directly on the Hopf map \(S^3\to\mathbb{CP}^1\): near the fibre
\((z_1,z_2)=(e^{ia},0)\), the target coordinate is \(z_2/z_1\), so
the two horizontal normal vectors are \((0,e^{ia})\) and
\((0,ie^{ia})\).  Add the tangent vector \((ie^{ia},0)\) and the
outward radial vector \((e^{ia},0)\).  In \(\mathbb C^2\) the resulting
full frame makes one complete rotation in each of its two real
coordinate planes.  This loop is contractible explicitly: for the
fixed real-linear involution \(J(z_1,z_2)=(z_1,\overline z_2)\),
\[
 J\operatorname{diag}(e^{ia},e^{ia})J=\operatorname{diag}(e^{ia},e^{-ia})=K_1(-a),
\]
and \(JK_v(-a)J\) contracts the original loop to the identity using
the periodic homotopy already constructed.  Thus the induced stable
tangent frame on the circle is the Lie frame.  This is the basic
framed-circle identification in \cite[Theorems~20--21, pp.~96--101]{Pontryagin1955Root}.

The product of two normally framed inverse images represents the
product in stable homotopy: its tubular collapse is the smash of
the two tubular collapse maps, as follows immediately from the
product tubular neighborhood \(\nu_1\times\nu_2\) and its ordered
normal coordinates.  Consequently the product of the two Lie
circles represents \(\eta^2\).  Equation \eqref{wb:h24:ret:eq:actual-Arf}
proves \(\eta^2\ne0\) by an explicit bordism invariant.  The
stable group \(\pi_2^{\rm S}=\mathbb Z/2\) is Pontryagin's computation
\cite[pp.~957--959]{Pontryagin1950Root} and \cite[Theorem~24, pp.~109--112]{Pontryagin1955Root}; the stable table in
\cite[opening table, p.~339, and Corollary~4.24, pp.~360--361]{Hatcher2002} gives the same group.  Thus the sign from
\eqref{wb:h24:ret:eq:sign} satisfies \(-\eta^2=\eta^2\).

\begin{proof}[Expanded proof of \cref{maps:src:B.3.1}]
\begin{enumerate}[label=\arabic*.,leftmargin=*]
\item The globally continuous quotient \(H\) is smooth on a neighborhood
of its compact inverse image \(M\), by
\eqref{wb:h24:ret:eq:exact-u} and \eqref{wb:h24:ret:eq:four-cell}. First make it constant
near the source basepoint, before choosing any approximation bound.
Choose a target chart \(\kappa:O\to\mathbb R^4\) about \(*\) with
\(\kappa(*)=0\), convex ball image, and \(y_0\notin O\).
Choose a source neighborhood \(U\) of \(x_{\rm base}\), disjoint
from \(M\), with \(H(U)\subset O\). Let \(\beta\) be a smooth
function supported in \(U\), equal to one near \(x_{\rm base}\),
with \(0\leq\beta\leq1\). On \(U\) use the based homotopy
\[
 H_s(x)=\kappa^{-1}\bigl((1-s\beta(x))\kappa(H(x))\bigr),
 \qquad 0\leq s\leq1,
\]
and leave \(H\) unchanged off \(U\). Its endpoint \(H_b\) is
smooth and constant near the basepoint; it agrees with \(H\)
near \(M\) and has no new inverse image of \(y_0\).

\item Choose one closed tubular neighborhood \(U_1\) of \(M\), with
\(M\subset\operatorname{int}U_1\), contained in the unchanged
smooth region and disjoint from the basepoint neighborhood.
Compactness gives the positive bound
\[
 d_0=\min_{x\in X\setminus\operatorname{int}U_1}
             d\bigl(H_b(x),y_0\bigr)>0.
\]
Relative smooth approximation, fixing the same closed set \(U_1\)
and a neighborhood of the basepoint, produces a smooth based map
\(H_{\rm sm}\) whose error from \(H_b\) is less than \(d_0/2\)
on its complement. Choose the approximation closer also than a
fixed target injectivity radius; the short target geodesics give a
based homotopy. Thus the exact inverse image and its differential
are unchanged on \(U_1\), and the complement acquires no new
inverse image. The relative Pontryagin construction applies as in
\cite[\S7, Theorem B and Lemma 4, pp.~43, 48--50]{CuspNullMilnor1965}.

\item Pontryagin--Thom therefore identifies the stabilization of
\([H]\) with the bordism class of \(M\) carrying the four normal
vectors \eqref{wb:h24:ret:eq:normal-first}--\eqref{wb:h24:ret:eq:normal-u}.  The comparison
\eqref{wb:h24:ret:eq:full-frame} and its explicit contraction identify that
class with the Lie torus, with the harmless order-two orientation
sign already computed.  Hence
\[
 E^\infty[H]=\eta^2\ne0.
\]
Freudenthal suspension is an isomorphism
\[
 E:\pi_{n+2}(S^n)\longrightarrow\pi_{n+3}(S^{n+1})
 \qquad(n\geq4),
\]
because \(n+2\leq2n-2\).  Therefore stabilization from
\(\pi_6(S^4)\) is an isomorphism onto \(\pi_2^{\rm S}\).
The composite \(\eta_4\circ\eta_5:S^6\to S^4\) stabilizes to
\(\eta^2\), so injectivity gives the exact equality
\cref{maps:src:B.3.1}.  Finally, a based null-homotopy of \(q_2\)
would postcompose with \(p_3\) to a based null-homotopy of \(H\),
contradicting \cref{maps:src:B.3.1}.
\end{enumerate}
\end{proof}

\paragraph{Exact scope of the comparison with the angle-forgetting map.}
The map here retains the angular coordinate before collapsing the
coordinate two-skeleton.  Its regular inverse image is the specific
normally four-framed \(T^2\) above.  The angle-forgetting
postcomposition has a normally three-framed \(T^3\) inverse image.
The calculation above is independent of its class and supplies an
actual obstruction to null-homotopy of the original \(q_2\), even
when that other postcomposition is null-homotopic.  The exact
morphism carrying this obstruction is
\[
 (p_3)_*:[X,\mathcal A]_*\longrightarrow[X,S^4]_*,\qquad
 [q_2]\longmapsto[H]=\eta_4\eta_5.
\]
No conclusion about an arbitrary earlier quotient follows merely
from the diffeomorphism type of one inverse image; the specified
four-framing is what proves this value.

\begin{corollary}[The exact order-two quotient detected by the original cusp map]
\label{wb:h24:ret:q2-detected-quotient}
\StatementRecord{wb:h24:ret:q2-detected-quotient}{proved}{Historical events: GPT-6 Astra (OpenAI), ChatGPT Pro mode, effort Pro, S6\_09\_PAPERS, 17 November 2026: targeted proof check. Also checked in this writing. Pre-continuation writing events: GPT-6 Astra (OpenAI), ChatGPT Pro mode, effort Pro, S6\_10\_PAPERS, 18 November 2026: targeted proof check. Continuation metadata events: GPT-6 Astra (OpenAI), ChatGPT Pro mode, effort Pro, S6\_10\_PAPERS continuation, 19 November 2026: used and rechecked in this run; status-vocabulary reconciliation; metadata reconciliation only; no new mathematical check. \emph{Pre-continuation recorded proof scope.} /root, effort Pro, 2026-11-17; also checked in this writing by GPT-6 Astra (OpenAI), ChatGPT Pro mode, effort Pro, 18 November 2026}{the hypotheses, notation and definitions stated here; \cref{maps:src:B.3.1}}

Let \(d:(S^6,*)\to(X,x_{\mathrm{base}})\) be the fixed based
diffeomorphism, and put
\[
\alpha=[q_2\circ d]\in\pi_6(\mathcal A,v_0).
\]
The original top-cell postcomposition induces a surjective homomorphism
\[
\varepsilon=(p_3)_*:\pi_6(\mathcal A,v_0)
  \longrightarrow\pi_6(S^4,*)\cong\mathbb Z/2,
\qquad \varepsilon(\alpha)=\eta_4\circ\eta_5.
\]
It gives the exact quotient isomorphism
\[
\pi_6(\mathcal A,v_0)/\ker\varepsilon
 \xrightarrow{\ [\beta]\mapsto\varepsilon(\beta)\ }
\mathbb Z/2.
\]
Every odd multiple of \(\alpha\) is nonzero. Its order is either
infinite or an even positive integer; this calculation does not determine
that order or assert a splitting of the displayed quotient.
\end{corollary}


\begin{proof}
\begin{enumerate}[label=\arabic*.,leftmargin=*]
\item Postcomposition distributes over the pinch sum on the source sphere,
so \((p_3)_*\) is a group homomorphism. Theorem~\ref{maps:src:B.3.1}
gives \((p_3)_*(\alpha)=[p_3q_2d]=\eta_4\eta_5\ne0\).
This is the nonzero element of the cyclic group of order two, proving
surjectivity. The first isomorphism theorem gives the stated quotient
with exactly the displayed map. For every integer \(n\),
\[
\varepsilon(n\alpha)=(n\bmod2)\eta_4\eta_5.
\]
An odd multiple therefore has nonzero image and cannot vanish.
If \(m\alpha=0\), applying \(\varepsilon\) forces \(m\)
to be even. These arguments prove all the assertions without inferring
the exact order of \(\alpha\).
\end{enumerate}
\end{proof}





\begin{lemma}[The three-dimensional coned Seifert target]
\label{mapsIII:carrier-three-complex-wedge}
\StatementRecord{mapsIII:carrier-three-complex-wedge}{proved}{No earlier separately identified check is recorded. Checked only in this writing. Pre-continuation writing events: GPT-6 Astra (OpenAI), ChatGPT Pro mode, effort Pro, S6\_10\_PAPERS, 18 November 2026: targeted proof check. Continuation metadata events: GPT-6 Astra (OpenAI), ChatGPT Pro mode, effort Pro, S6\_10\_PAPERS continuation, 19 November 2026: used and rechecked in this run; status-vocabulary reconciliation; metadata reconciliation only; no new mathematical check. \emph{Pre-continuation recorded proof scope.} GPT-6 Astra (OpenAI, ChatGPT Pro mode), S6\_09\_PAPERS, effort Pro, 16-17 November 2026; also checked in this writing by GPT-6 Astra (OpenAI), ChatGPT Pro mode, effort Pro, 18 November 2026}{the hypotheses, notation and definitions stated here; the explicit constructions and calculations in the proof}

Let $C$ be a simply connected CW complex with
$\widetilde H_2(C;\mathbb Z)=\widetilde H_3(C;\mathbb Z)=\mathbb Z$
and no other reduced integral homology. Then
$C\simeq S^2\vee S^3$.
\end{lemma}



\begin{proof}
\begin{enumerate}[label=\arabic*.,leftmargin=*]
\item The Hurewicz isomorphism $\pi_2(C)\to H_2(C;\mathbb Z)$ supplies a
based map $a:S^2\to C$ inducing an isomorphism on $H_2$ and on
$\pi_2$. Replace it by its mapping-cylinder inclusion
$S^2\hookrightarrow M_a$, where $M_a\simeq C$.
The pair $(M_a,S^2)$ is $2$-connected: both spaces are simply
connected, and the map on $\pi_2$ is an isomorphism.
Relative Hurewicz therefore gives
\[
 \pi_3(M_a,S^2)\xrightarrow{\ \cong\ }
 H_3(M_a,S^2;\mathbb Z).
\]
The homology exact sequence and $a_*:H_2(S^2)\cong H_2(M_a)$
identify this relative group with $H_3(M_a;\mathbb Z)$.
The homotopy exact sequence makes
$\pi_3(M_a)\to\pi_3(M_a,S^2)$ surjective, because its next map
lands in the kernel of the isomorphism
$\pi_2(S^2)\to\pi_2(M_a)$.
Naturality of the Hurewicz homomorphism then supplies a map
$b:S^3\to C$ representing an integral generator of $H_3(C)$.
The wedge $a\vee b:S^2\vee S^3\to C$ is an integral homology
equivalence between simply connected CW complexes. The homological
Whitehead theorem, proved in the classical topology appendix of Paper I,
makes it a homotopy equivalence.
\end{enumerate}
\end{proof}

For the coned Seifert target $\mathcal C$, the previously displayed
presentation kills the remaining boundary generators and gives
$\pi_1(\mathcal C)=1$, while Poincaré--Lefschetz duality gives exactly
the homology in this lemma. This provides the needed wedge equivalence
without an additional appeal to Whitehead's exact sequence.



\begin{proposition}[The actual degree-three finite-band map]
\label{mapsaudit:finite-degree-three}
\StatementRecord{mapsaudit:finite-degree-three}{proved}{No earlier separately identified check is recorded. Checked only in this writing. Pre-continuation writing events: GPT-6 Astra (OpenAI), ChatGPT Pro mode, effort Pro, S6\_10\_PAPERS, 18 November 2026: targeted proof check. Continuation metadata events: GPT-6 Astra (OpenAI), ChatGPT Pro mode, effort Pro, S6\_10\_PAPERS continuation, 19 November 2026: used and rechecked in this run; status-vocabulary reconciliation; metadata reconciliation only; no new mathematical check. \emph{Pre-continuation recorded proof scope.} GPT-6 Astra (OpenAI, ChatGPT Pro mode), S6\_09\_PAPERS, effort Pro, 16-17 November 2026; also checked in this writing by GPT-6 Astra (OpenAI), ChatGPT Pro mode, effort Pro, 18 November 2026}{the hypotheses, notation and definitions stated here; the explicit constructions and calculations in the proof}

Let $V=\mathbb Z\langle\gamma,u,w,\delta\rangle=H^1(F;\mathbb Z)$,
with $\operatorname{vol}=\gamma u w\delta$, and write
\[
 b=\gamma uw,\qquad
 c_1=uw\delta-4\gamma u\delta-2\gamma w\delta,\qquad
 c_2=uw\delta-3\gamma u\delta-3\gamma w\delta.
\]
For the degree-three and degree-four finite coverings $\pi_j:F\to S_j$,
\[
 \pi_1^*H^3(S_1;\mathbb Z)=\mathbb Zb\oplus\mathbb Zc_1,
 \qquad
 \pi_2^*H^3(S_2;\mathbb Z)=\mathbb Z(2b)\oplus\mathbb Zc_2.
\]
Choose the unique cohomology bases having those ordered pullbacks and the
Kronecker-dual bases of $H_3(S_j;\mathbb Z)$. Relative to the ordered
homology basis dual to $(\gamma uw,\gamma u\delta,\gamma w\delta,uw\delta)$,
the finite-band Mayer--Vietoris map is
\[
 \alpha_3=(\pi_{1*},-\pi_{2*})=
 \begin{pmatrix}
 1&0&0&0\\
 0&-4&-2&1\\
 -2&0&0&0\\
 0&3&3&-1
 \end{pmatrix}.
\]
Its Smith form has nonzero entries $(1,1,1)$ and its kernel is generated
by $(0,1,1,6)^t$. The order-three pushforward on $H_3$ is surjective;
the order-four pushforward has index two. In particular neither is
identified with its degree-one pushforward by integral changes of bases.
\end{proposition}


\begin{proof}
\begin{enumerate}[label=\arabic*.,leftmargin=*]
\item The integral fixed subgroups of $\bigwedge^3 V$ are
$\langle b,c_1\rangle$ and $\langle b,c_2\rangle$, respectively, by
solving the two integral invariance equations. The degree-one pullback
lattices are $\langle3\gamma,\eta_1\rangle$ and
$\langle4\gamma,\eta_2\rangle$, where
$\eta_1=2u+w+3\delta$ and $\eta_2=u+w+2\delta$.
The pairings of these bases with $(b,c_j)$ have matrices
\[
 \begin{pmatrix}0&3\\-3&0\end{pmatrix},\qquad
 \begin{pmatrix}0&4\\-2&0\end{pmatrix}.
\]
The downstairs $H^1$--$H^3$ pairing is unimodular, because the finite
surfaces have torsion-free integral cohomology and are oriented closed
four-manifolds. Pullback multiplies every pairing by the covering degree.
Thus its degree-three image has index one in the first invariant lattice
and index two in the second. If $rb+sc_2$ descends, pairing with the
descended class $\eta_2$ is divisible by four, so $-2r$ is divisible by
four. Hence its coefficient $r$ is even. This parity sublattice already
has index two, so it is the whole image. This proves the two displayed
pullback formulas, including sufficiency.

\item For $z\in H_3(F)$ and $a\in H^3(S_j)$, naturality of the actual
Kronecker pairing gives $\langle a,\pi_{j*}z\rangle
=\langle\pi_j^*a,z\rangle$. Evaluating the four listed classes therefore
gives the stated matrix, with the minus sign on the second summand.
Its $3\times3$ minors include $1$, and its rank is three. Solving its
kernel equations gives $z_1=0$, $z_3=z_2$, $z_4=6z_2$, proving the
remaining assertions. The first two rows have image all of $\mathbb Z^2$;
the last two have image $2\mathbb Z\oplus\mathbb Z$. In degree one the
corresponding image indices are three and four. These indices are
unchanged by unimodular changes of basis and rule out the asserted
degree-one/degree-three identification in the working proof.
\end{enumerate}
\end{proof}

\subsection{An extension homotopy for the angle-forgotten carriers}
\label{mapsIII:carrier-extension-section}

The literal flat quotient on the punctured cusp need not extend
continuously across the central fibre. Its homotopy class as a map to a
torus does extend. This distinction gives a second proof of the
angle-forgotten rank-two nullhomotopy and settles the one-copy circle
question left open in the working-edition homotopy-status remark.

\begin{lemma}[Suspending a torus-valued map that extends up to homotopy]
\label{mapsIII:carrier-extension-lemma}
\StatementRecord{mapsIII:carrier-extension-lemma}{proved}{Historical events: GPT-6 Astra (OpenAI), ChatGPT Pro mode, effort Pro, S6\_09\_PAPERS, 17 November 2026: targeted proof check. Also checked in this writing. Pre-continuation writing events: GPT-6 Astra (OpenAI), ChatGPT Pro mode, effort Pro, S6\_10\_PAPERS, 18 November 2026: targeted proof check. Continuation metadata events: GPT-6 Astra (OpenAI), ChatGPT Pro mode, effort Pro, S6\_10\_PAPERS continuation, 19 November 2026: used and rechecked in this run; status-vocabulary reconciliation; metadata reconciliation only; no new mathematical check. \emph{Pre-continuation recorded proof scope.} GPT-6 Astra (OpenAI, ChatGPT Pro mode), effort Pro, 2026-11-17; also checked in this writing by GPT-6 Astra (OpenAI), ChatGPT Pro mode, effort Pro, 18 November 2026}{the hypotheses, notation and definitions stated here; the explicit constructions and calculations in the proof}

Let $X$ be a compact Hausdorff space, and let $N\subset X$ be a compact
manifold with boundary which is a closed neighbourhood and satisfies
$\operatorname{int}_X(N)=N\setminus\partial N$.
Thus its frontier in $X$ is exactly its manifold boundary.
Let $r:N\to[0,\epsilon]$ be continuous, with
$W=r^{-1}(0)$, $\partial N=r^{-1}(\epsilon)$, and put
$U=N\setminus W$. Suppose $U$ is connected and locally path connected.
Let $\alpha:U\to T^d=\mathbb R^d/\mathbb Z^d$ be continuous.
Suppose there is a continuous $h:N\to T^d$ for which
$\alpha$ and $h|_U$ induce the same homomorphism to
$\pi_1(T^d)=\mathbb Z^d$. Fix $x_0\in W$, and choose $h(x_0)=0$.
For a based continuous map $c:(T^d,0)\to(Z,z_0)$ to a compact
Hausdorff based space, define
\[
 G(x)=
 \begin{cases}
 [r(x),c(\alpha(x))],&x\in U,\\
 *,&x\in W\cup(X\setminus\operatorname{int}N),
 \end{cases}
 \qquad
 \Sigma_r Z=
 \frac{[0,\epsilon]\times Z}
 {(\{0,\epsilon\}\times Z)\cup([0,\epsilon]\times\{z_0\})}.
\]
Then $G:X\to\Sigma_rZ$ is based nullhomotopic, by a homotopy
stationary on $(X\setminus\operatorname{int}N)\cup\{x_0\}$.
\end{lemma}



\begin{proof}
\begin{enumerate}[label=\arabic*.,leftmargin=*]
\item The difference $\alpha-h|_U:U\to T^d$ induces the zero homomorphism
on fundamental groups. The lifting criterion for the covering
$\mathbb R^d\to T^d$ therefore supplies a continuous function
$L:U\to\mathbb R^d$ with
\[
 \alpha(x)=h(x)+[L(x)].
\]
For $s\in[0,1]$, put
\[
 A_s(x)=h(x)+[(1-s)L(x)],\qquad
 G_s(x)=[r(x),c(A_s(x))]\quad(x\in U),
\]
and set $G_s=*$ on $W$ and outside $\operatorname{int}N$.
These formulas agree on the outer boundary. They are jointly continuous
at $W$ and at $\partial N$, even if $L$ is unbounded toward $W$:
every neighbourhood of the suspension basepoint contains the images of
$[0,\delta)\times Z$ and $(\epsilon-\delta,\epsilon]\times Z$
for some $\delta>0$, by compactness of $Z$. This estimate is uniform
in $s$ and in $A_s(x)$. Thus $G_s$ is a based homotopy from $G$
to the map $x\mapsto[r(x),c(h(x))]$ on $N$.

\item Now use that $h$ is defined on all of $N$. Set
\[
 H_s(x)=[(1-s)r(x)+s\epsilon,c(h(x))]
 \quad(x\in N),
 \qquad H_s(x)=*\quad(x\notin\operatorname{int}N).
\]
This is continuous by the closed pasting lemma. The height on
$\partial N$ is always $\epsilon$, so the formulas agree there.
At $x_0$ the second coordinate is $c(0)=z_0$, so $H_s(x_0)=*$
for every $s$. Finally $H_0=G_1$ and $H_1=*$.
Concatenating the two homotopies proves the assertion.
\end{enumerate}
\end{proof}

\begin{proposition}[Both angle-forgotten one-copy carriers are nullhomotopic]
\label{mapsIII:angle-forgotten-one-copy-null}
\StatementRecord{mapsIII:angle-forgotten-one-copy-null}{proved}{Historical events: GPT-6 Astra (OpenAI), ChatGPT Pro mode, effort Pro, S6\_09\_PAPERS, 17 November 2026: targeted proof check. Also checked in this writing. Pre-continuation writing events: GPT-6 Astra (OpenAI), ChatGPT Pro mode, effort Pro, S6\_10\_PAPERS, 18 November 2026: targeted proof check. Continuation metadata events: GPT-6 Astra (OpenAI), ChatGPT Pro mode, effort Pro, S6\_10\_PAPERS continuation, 19 November 2026: used and rechecked in this run; status-vocabulary reconciliation; metadata reconciliation only; no new mathematical check. \emph{Pre-continuation recorded proof scope.} GPT-6 Astra (OpenAI, ChatGPT Pro mode), effort Pro, 2026-11-17; also checked in this writing by GPT-6 Astra (OpenAI), ChatGPT Pro mode, effort Pro, 18 November 2026}{the hypotheses, notation and definitions stated here; \cref{mapsIII:carrier-extension-lemma}}

For the actual cusp neighbourhood $N_0=f^{-1}(|t|\leq\epsilon)$ and
the original flat marking, the maps
\[
 \widehat q_3:X\longrightarrow S^3,
 \qquad
 \widehat q_\gamma:X\longrightarrow S^2
\]
defined by forgetting the base angle and taking the reduced suspension
are based nullhomotopic. The assertion concerns their original formulas
and basepoint $x_0\in W$. It also gives based nullhomotopies of every
positive smash power of either map.
\end{proposition}



\begin{proof}
\begin{enumerate}[label=\arabic*.,leftmargin=*]
\item Write
\[
 \Lambda=\mathbb Z\langle\widehat\gamma,\widehat u,
                    \widehat w,\widehat\delta\rangle,
 \quad
 K=\mathbb Z\langle\widehat w,\widehat\delta\rangle,
 \quad
 \bar\Lambda=\Lambda/K.
\]
The geometric fundamental-group calculation in
\cref{mapsI:top:prop:geometric-pi1} proves the following
local information before taking any global quotient:
\[
 \pi_1(N_0)=\bar\Lambda,\qquad
 \pi_1(N_0\setminus W)=\Lambda\rtimes_{M_0}\mathbb Z,
 \qquad
 (\lambda,n)\longmapsto\bar\lambda.
\]
Take the radius $\epsilon$ strictly inside the punctured-overlap chart.
Because the fibration is a submersion over a small annulus about
$|t|=\epsilon$, a radial product collar identifies the fundamental
groups of this closed neighbourhood with those of the slightly larger
open neighbourhood used in that calculation, and identifies the
punctured inclusions as well.  This same collar makes $N_0$ a closed
codimension-zero neighbourhood in $X$, with
$\operatorname{int}_X(N_0)=N_0\setminus\partial N_0$ and frontier
$\partial N_0=f^{-1}(|t|=\epsilon)$, as required by the lemma.
Here the generator of the second factor is the zero-section meridian.
It bounds the actual holomorphic disc
$t\mapsto(t,1,1)$ in the toric chart, so there is no additional
translation in the displayed inclusion. Also
$(M_0-I)\Lambda=K$, as follows directly from
$(M_0-I)(a,b,c,d)=(0,0,b,-a)$.

\item The native map
$\alpha:N_0\setminus W\to\bar T^2=
(\bar\Lambda\otimes\mathbb R)/\bar\Lambda$
therefore induces precisely $(\lambda,n)\mapsto\bar\lambda$.
There is a map $h:N_0\to\bar T^2$ inducing the identity on
$\pi_1(N_0)=\bar\Lambda$. To see existence without any extra
classification assumption, take a CW model for the compact smooth
manifold $N_0$, map its one-cells with the specified integral winding
numbers, and extend over each two-cell because its attaching word
has trivial image in $\bar\Lambda$. All higher extensions exist
because the universal cover $\mathbb R^2$ of $\bar T^2$ is
contractible, so $\pi_k(\bar T^2)=0$ for $k\geq2$.
Compose with the CW comparison map to obtain $h$ on $N_0$.
Translation in the target arranges $h(x_0)=0$ and changes none of
the induced fundamental-group homomorphisms.

\item For $\widehat q_3$, apply Lemma~\ref{mapsIII:carrier-extension-lemma}
with $r=|t|$ and with
\[
 c:\bar T^2\longrightarrow
 \bar T^2/\bigl((S^1\times\{0\})\cup(\{0\}\times S^1)\bigr)
 \cong S^2.
\]
Its reduced suspension is exactly the original
$\widehat q_3(x)=[|t|,c(\alpha(x))]$ on the punctured cusp.
For $\widehat q_\gamma$, compose $\alpha$ and $h$ with the first
coordinate character $\bar T^2\to S^1$ and take $c$ to be the
identity of the based circle. The same lemma gives the nullhomotopy.
Both homotopies fix $x_0$, so their smash products with any based maps
are well defined and nullhomotopic.
\end{enumerate}
\end{proof}

\paragraph{Scope and preservation of the other proof.}
This argument uses the proved local fundamental-group inclusion and does
not produce a pointwise extension of the original flat character across
$W$. In particular, it does not alter the nonextension result for that
literal character. The explicit three-framed torus and its embedded
nullbordism provide a distinct proof of $\widehat q_3\simeq*$ and
are retained. The retained-angle map $q_2$ and its nonzero
$S^4$-valued postcomposition have different targets; the present
homotopy does not apply to them because their base-angle coordinate has
nonzero winding on the meridian that is killed in $N_0$.





\section{Geometric realization of two marked collision profiles}
\label{mapsD:repair:modular-collisions}

We use positive base meridians and column vectors for integral elliptic
homology.  A positive nodal monodromy is
\[
 T=\begin{pmatrix}1&1\\0&1\end{pmatrix},\qquad
 S=\begin{pmatrix}0&-1\\1&0\end{pmatrix},\qquad
 R=ST=\begin{pmatrix}0&-1\\1&1\end{pmatrix}.
\]
Thus $s=[S]$ and $r=[R]$ generate
$G=\operatorname{PSL}_2(\mathbb Z)=\langle s,r\mid s^2=r^3=1\rangle$,
and $[T]=sr$.  Products of matrices are written in the order used for
the archived cocycles: $A_1\cdots A_n=I$.  The section cocycle
$(q_i)$ satisfies
\begin{equation}\label{mapsD:repair:cocycle-relation}
 q_i\in\operatorname{sat}\operatorname{im}(A_i-I),\qquad
 \sum_{i=1}^n P_iq_i=0,\qquad P_i=A_1\cdots A_{i-1}.
\end{equation}
Changing one global logarithm by $v\in\mathbb Z^2$ changes $q_i$ by
$(A_i-I)v$.  Saturation here means intersection of the rational image
with $\mathbb Z^2$.

The fibre configurations and the rank-one Mordell--Weil lattices in
the next theorem are contained in Miranda's existence list and
Fukae's Table~3; they are not asserted to be new configurations
\cite[\S4, pp.~206--210]{manMiranda90}
\cite[Table~3, rows~50 and~56]{manFukae02}.
The additional verification is the realization of the particular
based monodromy and section cocycles used in the calculations.

\begin{theorem}[Realization of the archived marked profiles]
\label{mapsD:repair:collision-realization}
\StatementRecord{mapsD:repair:collision-realization}{proved}{Historical events: GPT-6 Astra (OpenAI), ChatGPT Pro mode, effort Pro, S6\_09\_PAPERS: carried as background without fault found. GPT-6 Astra (OpenAI), ChatGPT Pro mode, effort Pro, S6\_09\_PAPERS, 17 November 2026: targeted proof check. Also checked in this writing. Pre-continuation writing events: GPT-6 Astra (OpenAI), ChatGPT Pro mode, effort Pro, S6\_10\_PAPERS, 18 November 2026: targeted proof check. Continuation metadata events: GPT-6 Astra (OpenAI), ChatGPT Pro mode, effort Pro, S6\_10\_PAPERS continuation, 19 November 2026: used and rechecked in this run; status-vocabulary reconciliation; metadata reconciliation only; no new mathematical check. \emph{Pre-continuation recorded proof scope.} GPT-6 Astra (OpenAI, ChatGPT Pro mode; recorded reviewer /root/cdp\_audit), effort Pro, 17 November 2026; GPT-6 Astra (OpenAI, ChatGPT Pro mode), effort Pro, 2026-11-17 UTC; also checked in this writing by GPT-6 Astra (OpenAI), ChatGPT Pro mode, effort Pro, 18 November 2026}{the hypotheses, notation and definitions stated here; \cite[pp.~35--38]{mapsMilne2017}; \cref{mapsD:repair:50-weierstrass}; \cite[\S8.2--8.4, pp.~38--39]{manSS10}; \cref{mapsD:repair:cocycle-relation}; \cite[Lemma~4.1, Theorem~4.2 and Corollaries~4.2--4.3,
pp.~21--22]{manFukae02}; \cite[Lemmas~3.1--3.2, (11), pp.~11--13]{manFukae02}}

There are rational elliptic surfaces with zero section over
$\mathbb P^1$ having, in the indicated order, the following singular
fibres and positive based monodromy matrices:
\[
\begin{array}{c|c|c}
\text{profile}&\text{fibre types}&(A_1,\ldots,A_n)\\ \hline
50/\mathrm{II}&(\mathrm{II},I_3,I_1^*)&
\left(
\begin{pmatrix}3&7\\-1&-2\end{pmatrix},
\begin{pmatrix}7&12\\-3&-5\end{pmatrix},
\begin{pmatrix}-2&-1\\1&0\end{pmatrix}\right)\\[5pt]
56/\mathrm{II}&(I_2,\mathrm{II},I_5,I_3)&
\left(
\begin{pmatrix}5&8\\-2&-3\end{pmatrix},
\begin{pmatrix}27&37\\-19&-26\end{pmatrix},
\begin{pmatrix}61&80\\-45&-59\end{pmatrix},
\begin{pmatrix}4&3\\-3&-2\end{pmatrix}\right).
\end{array}
\]
Their Mordell--Weil groups are infinite cyclic.  A generator can be
chosen with the respective integral logarithm cocycle
\[
 q^{50}=\bigl((-5,2),(-8,4),(0,1)\bigr),\qquad
 q^{56}=\bigl((-2,1),(0,0),(4,-3),(1,-1)\bigr).
\]
Its height is $1/12$ in profile $50/\mathrm{II}$ and $1/30$ in
profile $56/\mathrm{II}$.  These are the exact marked data in the
archived collision-profile file, not merely data with the same fibre
types or the same height.
\end{theorem}


\noindent\textsf{Uses:} the finite-index modular-curve construction
\cite[Propositions~2.16, 2.21; Theorem~2.22; Proposition~4.3 and
Remark~4.4]{mapsMilne2017}; the minimal Weierstrass criterion
\cite[\S4.5, Table~1; \S8.2--8.4]{manSS10}; and the junction
identification \cite[\S3.1, (4), (8), (9); Lemmas~3.1--3.2, 4.1;
Corollaries~4.2--4.3]{manFukae02}.

\begin{proof}
\begin{enumerate}[label=\arabic*.,leftmargin=*]
\item \par\medskip\noindent\textbf{1. Two specified modular covers.}\quad 
We first construct the covers, retaining their markings.  A row of
either of the following tables gives the right action of $s,r,r^{-1}$
on a coset labelled by its first entry:
\[
\begin{array}{c|rrr}
50/\mathrm{II}&s&r&r^{-1}\\\hline
0&1&2&3\\1&0&1&1\\2&3&3&0\\3&2&0&2
\end{array}
\qquad
\begin{array}{c|rrr}
56/\mathrm{II}&s&r&r^{-1}\\\hline
0&1&2&3\\1&0&1&1\\2&4&3&0\\3&5&0&2\\
4&2&6&7\\5&3&8&9\\6&9&7&4\\7&8&4&6\\
8&7&9&5\\9&6&5&8
\end{array}
\]
Each $s$-permutation is an involution, each $r$-permutation has cube
one, and the action is transitive.  These statements can be checked
directly in the displayed tables.  Consequently the stabilizer of
the coset $0$ is a subgroup of index $4$, respectively $10$, in $G$.
Its inverse image in $\operatorname{SL}_2(\mathbb Z)$ contains $-I$.
We use the compact modular curve of this subgroup.  The construction
of its complex structure uses the coordinates $z^e$ at elliptic
points and $\exp(2\pi i\tau/w)$ at a cusp of width $w$;
it applies to every finite-index subgroup and does not require a
congruence assumption \cite[pp.~35--38]{mapsMilne2017}.

\item The cycles of $s$, $r$, and $sr$ are as follows.  Parentheses denote
cycles on the coset labels, including fixed points:
\[
\begin{array}{c|l|l}
&50/\mathrm{II}&56/\mathrm{II}\\\hline
s&(01)(23)&(01)(24)(35)(69)(78)\\
r&(023)(1)&(023)(1)(467)(589)\\
sr&(012)(3)&(01265)(384)(79).
\end{array}
\]
Thus both curves have exactly one elliptic point of order $3$ and
none of order $2$.  Their cusp widths are $(3,1)$ and $(5,3,2)$.
The map $j$ to the $j$-line has respectively degree $4$ and $10$.
Its ramification contributions over $0,1728,\infty$ are
\[
 (2,2,2)\quad\hbox{and}\quad(6,5,7).
\]
Their sums are $2\cdot4-2$ and $2\cdot10-2$.  Riemann--Hurwitz
therefore gives genus zero in both cases.  In particular each base
is $\mathbb P^1$.

\item \par\medskip\noindent\textbf{2. The actual based peripheral loops.}\quad 
We describe a finite graph on each modular curve; this also specifies
the paths used for the matrices.  Put a black vertex at each
$r$-orbit and a white vertex at each $s$-orbit.  For every coset $q$,
join its black and white vertices by an edge labelled $q$, and
subdivide that edge at a midpoint also denoted by $q$.  The clockwise
cyclic order at a black vertex is the order of its $r$-cycle.  White
vertices have valence two.  This is the inverse image of an arc
joining the elliptic values on the modular sphere.  Filling its
complementary cusp discs recovers the compact curve.  Remove a
small disc about the fixed black vertex labelled $(1)$ as well.
The resulting surface is the base with its singular-fibre points
removed.  This description follows directly by lifting the modular
triangles: a three-cycle fills a disc above the order-three value,
an $s$-pair fills a disc above the order-two value, and an $sr$-cycle
is the boundary of a cusp disc.

\item The positive cusp boundary reads $sr=T$.  At the fixed black vertex,
the positive elliptic meridian reads $r^{-1}$.  These signs also
follow from the analytic coordinates: $\tau\mapsto\tau+1$ gives
a positive loop of $\exp(2\pi i\tau)$, while $R$ has derivative
$\exp(-2\pi i/3)$ at its fixed point in the upper half-plane.
We use the oriented elliptic marking in which this positive nodal
meridian acts by $T$; it fixes the choice between the two possible
outer identifications of the modular group with elliptic homology.

\item The base point is the midpoint $0$.  Stems to the midpoints $1,3,7$
are, respectively,
\[
 s,\qquad r^{-1},\qquad rsr^{-1};
\]
the last is needed only in the index-ten cover.  They are paths in
the graph just specified.  A cusp based at midpoint $q$, of width
$w$ and with stem $g$, therefore has based peripheral word
$gT^wg^{-1}$.  The fixed black vertex has word $sr^{-1}s^{-1}$.
Set
\begin{equation}\label{mapsD:repair:canonical-matrices}
 H=SR^{-1}S^{-1}=\begin{pmatrix}0&1\\-1&1\end{pmatrix},\quad
 V=R^{-1}TR=\begin{pmatrix}2&1\\-1&0\end{pmatrix},\quad
 B=V^3=\begin{pmatrix}4&3\\-3&-2\end{pmatrix},
\end{equation}
and
\[
 C=(RSR^{-1})T^2(RSR^{-1})^{-1}
   =\begin{pmatrix}5&2\\-8&-3\end{pmatrix}.
\]
The cusp and elliptic words, in positive boundary order, are
\[
 (T^3,V,H),\qquad (T^5,B,C,H).
\]

\item For completeness, we verify the boundary relation in the fundamental
group of the punctured surface, before imposing the order-three
orbifold relation.  Denote the black and white vertices by $b_i,w_i$
in the order in which their cycles occur in the displayed table.
Orient every half-edge from its midpoint $q$ towards its black or
white endpoint.  Take as a spanning tree every half-edge except
\[
\begin{array}{c|l}
50/\mathrm{II}&3\longrightarrow w_1,\\
56/\mathrm{II}&8\longrightarrow w_4,\quad9\longrightarrow w_3.
\end{array}
\]
The indicated subgraphs are connected and have one fewer edge than
vertices, hence are trees.  In the index-four case let $a$ denote
the remaining oriented half-edge, after the tree is contracted.
In the index-ten case call the two remaining half-edges $a,b$, in
the order displayed.  In both cases let $e$ denote the clockwise
loop at the punctured fixed black vertex.  These are free generators
of the actual punctured-surface group.  Reading the paths above and
discarding tree edges gives
\begin{equation}\label{mapsD:repair:free-peripherals}
\begin{array}{c|ccc}
50/\mathrm{II}&T^3&e a^{-1}&\\[-3pt]
&V&a&\\[-3pt]
&H&e^{-1}&
\end{array}
\qquad
\begin{array}{c|cc}
56/\mathrm{II}&T^5&e b^{-1}\\
&B&a\\
&C&a^{-1}b\\
&H&e^{-1}.
\end{array}
\end{equation}
For example the index-ten width-five cycle reads successively the
midpoints $0,1,2,6,5$; it crosses the punctured vertex once and the
half-edge $9\to w_3$ with negative orientation, giving $eb^{-1}$.
The width-three cycle $3,8,4$ gives $a$, and the width-two cycle
$7,9$ gives $a^{-1}b$.  In the index-four case the cycles $0,1,2$
and $3$ give $ea^{-1}$ and $a$.  A positive loop at the elliptic
puncture is $e^{-1}$ in both cases.  Thus the displayed boundary
products freely reduce to $1$ and their factors generate the whole
free group: in the second case $a=B$, $b=BC$, and $e=H^{-1}$.
The stems can be separated in a thin regular neighbourhood of the
tree, so they are genuine based peripheral loops, with the listed
cyclic order.  In particular this check does not replace the
punctured-surface group by its quotient in $G$.
\par\medskip\noindent\textbf{3. Weierstrass surfaces with these covers.}\quad 
For the index-four cover an explicit model is
\begin{equation}\label{mapsD:repair:50-weierstrass}
\begin{split}
 c_4&=t^2(t+27)(t+3),\\
 c_6&=t^3(t+27)(t^2+18t-27),\\
 \Delta&=t^7(t+27)^2,\\
 y^2&=x^3-\frac{c_4}{48}x-\frac{c_6}{864}.
\end{split}
\end{equation}
Multiplication gives $c_4^3-c_6^2=1728\Delta$.  The coefficients
are sections of $\mathcal O(4)$ and $\mathcal O(6)$ on
$\mathbb P^1$.  Their vanishing orders $(c_4,c_6,\Delta)$ at
$0,-27,\infty$ are, respectively,
\[
 (2,3,7),\qquad(1,1,2),\qquad(0,0,3).
\]
The minimal resolved fibres are therefore $I_1^*,\mathrm{II},I_3$.
No other fibre is singular.  Its $j$-map is
\[
 j(t)=\frac{(t+27)(t+3)^3}{t},\qquad
 j(t)-1728=\frac{(t^2+18t-27)^2}{t}.
\]
It has precisely the index-four passport in Step~1.  That passport
specifies its marked cover up to a relabelling of sheets: after
fixing $r=(023)(1)$, the three possible fixed-point-free involutions
on four letters are conjugate by the centralizer $\langle r\rangle$.
Thus the explicit model can be marked by the first coset table.

For the index-ten cover, the modular curve and its normalized
$j$-map specify the construction without the solution of any
additional polynomial equation.  Write $a$ for the single simple
point over $0$, $b_1,b_2,b_3$ for the triple points over $0$,
$d_1,\ldots,d_5$ for the double points over $1728$, and
$p_5,p_3,p_2$ for the three poles, whose subscripts are their orders.
All these points are distinct.  Choose a nonzero section
$\Delta$ of $\mathcal O(12)$ with divisor
\[
 5p_5+3p_3+2p_2+2a.
\]
The sections $j\Delta$ and $(j-1728)\Delta$ then have divisors
\[
 3(a+b_1+b_2+b_3),\qquad
 2(a+d_1+\cdots+d_5),
\]
respectively.  On $\mathbb P^1$ a section with either displayed
divisor has a cube root in $H^0(\mathcal O(4))$, respectively a
square root in $H^0(\mathcal O(6))$, after multiplying a section
with the divided divisor by a constant.  The required constant
roots exist over $\mathbb C$.  Choose these roots and call them
$c_4,c_6$.  They satisfy identically
\[
 c_4^3=j\Delta,\qquad c_6^2=(j-1728)\Delta,
 \qquad c_4^3-c_6^2=1728\Delta.
\]
The same Weierstrass equation as in
\eqref{mapsD:repair:50-weierstrass} therefore defines an elliptic
surface with a zero section.  At $a$ the orders are $(1,1,2)$,
giving type~II.  At $p_m$ they are $(0,0,m)$, giving $I_m$.
Every other point has nonvanishing discriminant.  The model is
minimal at each point since the orders of its two coefficients
are never simultaneously at least $4$ and $6$.  Its minimal
resolution is projective and has fundamental line bundle
$\mathcal O(1)$.  Both constructions are rational elliptic
surfaces by the minimal degree-$(4,6)$ criterion
\cite[\S8.2--8.4, pp.~38--39]{manSS10}.

\par\medskip\noindent\textbf{4. Lifting the projective marking.}\quad 
On a simply connected open set avoiding the singular values and
the elliptic values of $j$, uniformization identifies the fibre
with $\mathbb C/(\mathbb Z+\tau\mathbb Z)$, where $\tau$ is a
local inverse of the modular $j$-function.  Continuation through
the marked cover gives precisely its projective elliptic
monodromy.  The sign of an integral lift at each singular fibre
is fixed by its Kodaira type: $I_m$ has the positive unipotent
lift, $I_m^*$ its negative, and the type-II lift has trace~$1$.
With the peripheral paths of Step~2 the two ordered integral
tuples are consequently
\[
 (H,T^3,-V),\qquad (C,H,T^5,B).
\]
Their products are $I$, as direct multiplication verifies.  The
first tuple uses the negative cusp lift only at the width-one
cusp.  There is no residual quadratic character to choose: the
meridians generate the punctured-sphere group and their signs
have all been specified.  The cyclic rotations from Step~2
are geometric changes of the starting boundary arc.

Conjugating the first tuple simultaneously by $V$, and the second
by $B$, gives
\[
\begin{aligned}
 &(VHV^{-1},VT^3V^{-1},-V)\\[-1pt]
 &\quad=\left(
 \begin{pmatrix}3&7\\-1&-2\end{pmatrix},
 \begin{pmatrix}7&12\\-3&-5\end{pmatrix},
 \begin{pmatrix}-2&-1\\1&0\end{pmatrix}\right),\\
 &(BCB^{-1},BHB^{-1},BT^5B^{-1},B)\\[-1pt]
 &\quad=\left(
 \begin{pmatrix}5&8\\-2&-3\end{pmatrix},
 \begin{pmatrix}27&37\\-19&-26\end{pmatrix},
 \begin{pmatrix}61&80\\-45&-59\end{pmatrix},
 \begin{pmatrix}4&3\\-3&-2\end{pmatrix}\right).
\end{aligned}
\]
These are the stated matrices, in the stated order, in integral
oriented homology bases.

\par\medskip\noindent\textbf{5. The integral cocycle group and actual sections.}\quad 
We spell out the relation between
\eqref{mapsD:repair:cocycle-relation} and the geometric section
group, so that a rational height computation alone is not used
as an existence argument.  For any of the rational elliptic
surfaces just constructed, let $L$ be the orthogonal complement
of the unimodular sublattice generated by the fibre $F$ and zero section $O$ in integral $H_2$, with matrix $\left(\begin{smallmatrix}0&1\\1&-1\end{smallmatrix}\right)$ in the ordered basis $(F,O)$, and let $V_{\mathrm{fib}}$ be the subgroup generated by
the fibre components not meeting the zero section.  Fukae's
rational junction group $W_2$ consists of integral charges in
the rational local images, with total charge zero.  Its quotient
by integral null junctions $U_1$ is $L/V_{\mathrm{fib}}$;
its quotient by the saturation $U_2$ is the free Mordell--Weil
lattice, with the opposite pairing.  The quotient $U_2/U_1$
is the torsion subgroup
\cite[Lemma~4.1, Theorem~4.2 and Corollaries~4.2--4.3,
pp.~21--22]{manFukae02}.

Here is the exact convention conversion.  Put
\[
 y_i=P_iq_i,\qquad K_i=P_iA_i^{-1}P_i^{-1},
\]
and list the $K_i$ and $y_i$ in reverse order.  Then
$K_1\cdots K_n=I$, so the reversed list satisfies Fukae's
reverse-product convention.  His circuit matrices are inverses
of the ordinary monodromy matrices, as stated in his footnote
on p.~3.  The ordinary peripheral basis underlying the reversed
list is obtained by successive inverse Hurwitz moves:
its elements are $P_iA_iP_i^{-1}$, in reverse order.  This is a
geometric change of disjoint peripheral stems.  Moreover
\[
 y_i\in\operatorname{sat}\operatorname{im}(I-K_i),\qquad
 \sum_i y_i=0.
\]
Thus $(q_i)\mapsto(y_i)$ is an integral isomorphism from the
cocycle lattice to $W_2$.  Its inverse is $q_i=P_i^{-1}y_i$.

For an integral coboundary $q_i=(A_i-I)v$ we obtain
\[
 y_i=(P_{i+1}-P_i)v
     =(I-K_i)P_{i+1}v.
\]
In the reversed list, the product of the preceding Fukae
circuits is $P_{i+1}$.  The displayed formula is therefore
exactly his integral null-junction formula, and allowing
rational $v$ with integral charges gives its saturation
\cite[Lemmas~3.1--3.2, (11), pp.~11--13]{manFukae02}.
Consequently the quotient in
\eqref{mapsD:repair:cocycle-relation} is the full integral
section group, including its possible torsion.

The identification is the one for logarithm cocycles, not just
an abstract lattice isometry.  To see this, cut the punctured
base along its peripheral stems and trivialize the period
lattice on the cut disc.  A section has a continuous elliptic
logarithm there.  Along the two sides of the $i$th cut the
logarithms differ by $q_i$.  Transport to the common reference
fibre changes this charge to $P_iq_i=y_i$.  The boundary of
the swept relative chain is the sum of these charges, so the
closing condition is exactly $\sum_i y_i=0$.  Local rational
prongs fill precisely the rational image of $I-K_i$; retaining
integral charges gives $W_2$.  A change of the common logarithm
is the null sweep just calculated.  Conversely the junction
construction supplies the integral class in $L/V_{\mathrm{fib}}$
with these jumps.  On a rational surface the exponential
sequence identifies integral $H^2$ with the Picard group,
since $H^1(\mathcal O)=H^2(\mathcal O)=0$.  Restricting the
corresponding divisor class to the generic elliptic fibre
gives a degree-zero line bundle, hence its unique point under
the zero-section identification of the elliptic curve with
its degree-zero Picard variety.  That point extends to a
section of the minimal proper elliptic surface.  Vertical
components and the zero-section/fibre adjustment are exactly
the classes already quotiented out.  This reverses the
cut-and-jump construction and proves that every quotient
cocycle is realized by an actual section.

For profile $50/\mathrm{II}$ the first local difference is
\[
 A_1-I=\begin{pmatrix}2&7\\-1&-3\end{pmatrix},\qquad
 \det(A_1-I)=1.
\]
Every integral cocycle has a unique representative with
$q_1=0$, by subtracting one integral coboundary.  For the
displayed cocycle the required vector is $v=(1,-1)$, giving
\[
 q^0=\bigl((0,0),(-2,1),(2,-1)\bigr).
\]
The saturated image at $I_3$ is
$\mathbb Z(-2,1)$.  Once its coefficient $k\in\mathbb Z$
is chosen, the closing relation uniquely determines the
last charge, since $P_3$ is unimodular and the local image
at $I_1^*$ has full rank.  It is $k(2,-1)$.  Hence the
whole quotient is $\mathbb Zq^0$, with no torsion, and the
archived cocycle is a primitive section generator.

For profile $56/\mathrm{II}$ the type-II fibre is the second,
and again $\det(A_2-I)=1$.  Set $q_2=0$ uniquely by an
integral coboundary.  The other three saturated local images
have primitive generators
\[
 (-2,1),\qquad(4,-3),\qquad(1,-1).
\]
After prefix transport their columns are
\[
 (-2,1),\qquad(1,0),\qquad(1,-1).
\]
The first two have determinant $-1$.  The integral kernel
of the resulting $2$-by-$3$ matrix is exactly
$\mathbb Z(1,1,1)$, as the second row gives equality of the
first and third coefficients and the first row then gives
equality of the second.  Thus here too the quotient is
infinite cyclic, without torsion, and its generator is the
archived cocycle.

\par\medskip\noindent\textbf{6. Heights with the precise normalization.}\quad 
Let $z_i\in\mathbb Q^2$ satisfy $(A_i-I)z_i=q_i$ and write
$u\wedge v=u_1v_2-u_2v_1$.  In Fukae's convention a
preimage of $y_i$ under $I-K_i$ is
$P_i(z_i+q_i)$.  The diagonal contribution to the junction
square is therefore $z_i\wedge q_i$.  The reversal of the
ordering changes the cross term to
$-\sum_{i<j}(P_iq_i)\wedge(P_jq_j)$.  This gives
\begin{equation}\label{mapsD:repair:height-formula}
 \langle q,q\rangle_{\mathrm{MW}}
 =-\left\{\sum_i z_i\wedge q_i
       -\sum_{i<j}(P_iq_i)\wedge(P_jq_j)\right\}.
\end{equation}
For a multiplicative fibre this also reproduces Fukae's
$-t_i'^2/n_i$ term in (9): if
$A_i=I+n_i\,v(-v_2,v_1)$, substitution in $(A_i-I)z_i=q_i$
gives the same wedge expression.  Equivalently one may
check it on $A_i=T^{n_i}$, where $q_i=(t_i',0)$ and
$z_i=(0,t_i'/n_i)$, obtaining $-t_i'^2/n_i$.
The formula is independent of the choice of $z_i$, since
the kernel and image of a rank-one unipotent difference
are the same isotropic line.

For $50/\mathrm{II}$ take
\[
 (z_1,z_2,z_3)=
 \bigl((1,-1),(-4/3,0),(1/4,-3/4)\bigr).
\]
The prefix charges and the two sums in braces are
\phantomsection\label{qa:III:layout:2}
\[
\begin{gathered}
 (P_iq_i)=\bigl((-5,2),(4,0),(1,-2)\bigr),\\
 \sum_i z_i\wedge q_i=-3-\frac{16}{3}+\frac14
 =-\frac{97}{12},\qquad
 \sum_{i<j}(P_iq_i)\wedge(P_jq_j)=-8.
\end{gathered}
\]
Thus the junction square is $-1/12$ and the height is $1/12$.
For $56/\mathrm{II}$ take
\[
 (z_1,z_2,z_3,z_4)=
 \bigl((-1/2,0),(0,0),(-1/5,1/5),(1/3,0)\bigr).
\]
Here
\phantomsection\label{qa:III:layout:3}
\[
\begin{gathered}
 (P_iq_i)=\bigl((-2,1),(0,0),(1,0),(1,-1)\bigr),\\
 \sum_i z_i\wedge q_i=-\frac12-\frac15-\frac13
 =-\frac{31}{30},\qquad
 \sum_{i<j}(P_iq_i)\wedge(P_jq_j)=-1.
\end{gathered}
\]
The junction square is $-1/30$ and the height is $1/30$.
All entries in the theorem have now been realized with
the required based and integral markings.
\end{enumerate}
\end{proof}

\phantomsection\label{qa:III:layout:4}
\par\medskip\noindent\textbf{Exact finite verification.}\quad
The current verifier \path{code/III/receiving_final_certificates.py}
reconstructs both coset actions and the free peripheral words directly
from the punctured half-edge graphs. It compares both integral matrix
tuples and both cocycles with the archived marked profiles, checks the
Weierstrass identity, the primitive section quotients, the explicit
Fukae convention morphisms, and both rational heights. Its complete
finite output is \path{analysis/III_final_exact_certificates.json}.
The geometric existence and the actual logarithm-cocycle morphisms
are proved above; the program records the finite exact calculations.

\section{Short receiving applications of the cited classification theorems}


\begin{theorem}[Credited re-derivation: the smooth side is $S^3\times S^3$]\label{maps:src:D.3.2}
\StatementRecord{maps:src:D.3.2}{proved}{Historical events: Claude Fable 5.2 (Anthropic), effort high, 16 November 2026: carried as background without fault found. GPT-6 Astra (OpenAI), ChatGPT Pro mode, effort Pro, S6\_09\_PAPERS, 17 November 2026: targeted proof check. Also checked in this writing. Pre-continuation writing events: GPT-6 Astra (OpenAI), ChatGPT Pro mode, effort Pro, S6\_10\_PAPERS, 18 November 2026: targeted proof check. Continuation metadata events: GPT-6 Astra (OpenAI), ChatGPT Pro mode, effort Pro, S6\_10\_PAPERS continuation, 19 November 2026: used and rechecked in this run; status-vocabulary reconciliation; metadata reconciliation only; no new mathematical check. \emph{Pre-continuation recorded proof scope.} GPT-6 Astra (OpenAI, ChatGPT Pro mode), effort Pro, 2026-11-17 UTC; Claude Fable 5.2 (Anthropic), effort high, 16 November 2026; also checked in this writing by GPT-6 Astra (OpenAI), ChatGPT Pro mode, effort Pro, 18 November 2026}{the hypotheses, notation and definitions stated here; \cite{BHS19}; \cite[p.~512]{manKM63}}


For the sphere twist data and $0<|s|\ll1$, the complex-oriented
smooth manifold $Y_{u,s}$ is orientation-preservingly diffeomorphic
to $S^3\times S^3$. Its integral homology is
$(\mathbb Z,0,0,\mathbb Z^2,0,0,\mathbb Z)$,
its middle intersection form is the rank-two unimodular
skew-symmetric form, and $w_2=p_1=0$.
\end{theorem}


\begin{proof}
\begin{enumerate}[label=\arabic*.,leftmargin=*]
\item An ordinary-double-point smoothing replaces a disc neighbourhood
of its exceptional curve by a disc neighbourhood of the vanishing
three-sphere. In the displayed coordinates this can be read
without a deformation-theoretic inference: the resolution is the
disc bundle of $\mathcal O(-1)^{\oplus2}$ over $S^2$, and the
Milnor fibre of the nondegenerate quadratic equation is the disc
cotangent bundle of $S^3$. The underlying oriented rank-four
bundle over $S^2$ is trivial, since its $w_2$ is the reduction of
$-2$ and oriented rank-four bundles over $S^2$ are classified by
$\pi_1SO(4)=\mathbb Z/2$. The latter disc cotangent bundle is
$D^3\times S^3$ because $S^3$ is parallelizable. Both boundaries
are $S^2\times S^3$.

\item Let $A$ be the complement of the interior of that resolution tube
in $X^{\log}(u)\simeq S^6$. General position makes $A$ simply connected:
a loop and a filling disc in $S^6$ can be kept disjoint from the
embedded two-sphere since $2+2<6$. Alexander duality gives
$\widetilde H_i(A;\mathbb Z)=\mathbb Z$ for $i=3$ and zero for
all other $i$. The normal three-sphere in the boundary generates
that group. In the Mayer--Vietoris sequence for
$Y=A\cup(D^3\times S^3)$, the degree-three map from the boundary
is therefore $\mathbb Z\to\mathbb Z^2$, with both entries units,
and the degree-two boundary group is $\mathbb Z$ with zero image.
It follows that $H_2(Y)=H_4(Y)=0$ and that $H_3(Y)$ is an extension
of $\mathbb Z$ by $\mathbb Z$, hence $\mathbb Z^2$.
Van Kampen gives $\pi_1(Y)=1$. The remaining homology follows by
duality, and the assertions about $w_2,p_1$ follow from the
vanishing of the corresponding cohomology groups.

\item We apply Wall's classification in the precise form of
Burklund--Hahn--Senger, Theorem 2.6, with their Definition 2.5 and
Recollection 2.2 \cite{BHS19}. This theorem says that two
$(n-1)$-connected closed $2n$-manifolds with the same middle
intersection form and normal invariant
$\alpha:H_n\to\pi_{n-1}SO(n)$ differ by a homotopy sphere.
Here $n=3$, the Hurewicz theorem gives $\pi_2(Y)=0$, and
$\pi_2SO(3)=0$, so $\alpha$ is zero. Every unimodular
skew-symmetric rank-two form has an oriented symplectic basis;
thus these invariants equal those of $S^3\times S^3$.
Finally $\Theta_6/bP_7=0$ by the dimension-six entry of
\cite[p.~512]{manKM63}, and $bP_7=0$ by its Theorem~5.1
(equivalently the framed-surgery Theorem~6.6). Thus $\Theta_6=0$
removes the homotopy-sphere ambiguity.
This derives the diffeomorphism assertion from the classification
rather than from Euler characteristic alone.
\end{enumerate}
\end{proof}

\begin{proposition}[Credited re-derivation: comparison with Calabi--Eckmann deformations]\label{maps:src:D.3.4}
\StatementRecord{maps:src:D.3.4}{proved}{Historical events: Claude Fable 5.2 (Anthropic), effort high, 16 November 2026: carried as background without fault found. GPT-6 Astra (OpenAI), ChatGPT Pro mode, effort Pro, S6\_09\_PAPERS, 17 November 2026: targeted proof check. Also checked in this writing. Pre-continuation writing events: GPT-6 Astra (OpenAI), ChatGPT Pro mode, effort Pro, S6\_10\_PAPERS, 18 November 2026: targeted proof check. Continuation metadata events: GPT-6 Astra (OpenAI), ChatGPT Pro mode, effort Pro, S6\_10\_PAPERS continuation, 19 November 2026: used and rechecked in this run; status-vocabulary reconciliation; metadata reconciliation only; no new mathematical check. \emph{Pre-continuation recorded proof scope.} GPT-6 Astra (OpenAI, ChatGPT Pro mode), effort Pro, 2026-11-17 UTC; Claude Fable 5.2 (Anthropic), effort high, 16 November 2026; also checked in this writing by GPT-6 Astra (OpenAI), ChatGPT Pro mode, effort Pro, 18 November 2026}{the hypotheses, notation and definitions stated here; \cite[Example~2.3]{manCDP98}; \cref{maps:src:D.3.3}}


Complex structures with $a=1$ on $S^3\times S^3$ are constructed in
\cite[Example~2.3]{manCDP98}. The present smoothing is not
biholomorphic to the displayed
$a=1$ subfamily of that example.
\end{proposition}


\begin{proof}
\begin{enumerate}[label=\arabic*.,leftmargin=*]
\item That example takes the quotient of
$(\mathbb C^2\setminus\{0\})^2$ by
\[
 t\cdot(x,y,z,w)=(e^t x,e^{at}y,e^{ibt}z,e^{ibt}w),
 \qquad a,b>0,\quad a\notin\mathbb Q.
\]
The invariant map $[z:w]$ is its algebraic reduction. Fixing this
ratio and setting its nonzero scalar coordinate equal to one leaves
the period $t=2\pi k/b$. The fibre is therefore the quotient of
$\mathbb C^2\setminus\{0\}$ by the diagonal expansion
$(x,y)\mapsto(e^{2\pi/b}x,e^{2\pi a/b}y)$, hence a Hopf surface
with first Betti number one. The very general algebraic-reduction
fibre of \ref{maps:src:D.3.3} is a two-torus with first Betti number four.
A biholomorphism induces an isomorphism of meromorphic function
fields and hence identifies very general algebraic-reduction fibres.
The Betti numbers exclude such an identification.
\end{enumerate}
\end{proof}

\begin{theorem}[Credited re-derivation: transverse pullback and the trivial case]\label{maps:src:E.2.2}
\StatementRecord{maps:src:E.2.2}{proved}{Historical events: Claude Fable 5.2 (Anthropic), effort high, 16 November 2026: carried as background without fault found. GPT-6 Astra (OpenAI), ChatGPT Pro mode, effort Pro, S6\_09\_PAPERS, 17 November 2026: targeted proof check. Also checked in this writing. Pre-continuation writing events: GPT-6 Astra (OpenAI), ChatGPT Pro mode, effort Pro, S6\_10\_PAPERS, 18 November 2026: targeted proof check. Continuation metadata events: GPT-6 Astra (OpenAI), ChatGPT Pro mode, effort Pro, S6\_10\_PAPERS continuation, 19 November 2026: used and rechecked in this run; status-vocabulary reconciliation; metadata reconciliation only; no new mathematical check. \emph{Pre-continuation recorded proof scope.} GPT-6 Astra (OpenAI, ChatGPT Pro mode; recorded reviewer /root/cdp\_audit), effort Pro, 17 November 2026; Claude Fable 5.2 (Anthropic), effort high, 16 November 2026; also checked in this writing by GPT-6 Astra (OpenAI), ChatGPT Pro mode, effort Pro, 18 November 2026}{the hypotheses, notation and definitions stated here; the explicit constructions and calculations in the proof}


Let $Y$ be a compact complex manifold and $h:Y\to B$ a holomorphic
map which is a submersion at every point over $p_0,p_1,p_2$.
Then $P=X\times_B Y$ is a compact complex manifold, of dimension
$\dim_{\mathbb C}Y+2$.  If $Y=B\times N$ with its projection,
then $P\simeq X\times N$ holomorphically, is diffeomorphic to
$S^6\times N$, and
\[
 \pi_1(P)=\pi_1(N),\quad
 H_k(P;\mathbb Z)=H_k(N;\mathbb Z)\oplus H_{k-6}(N;\mathbb Z),
 \quad \chi(P)=2\chi(N).
 \tag{E.15}\label{maps:eq:E.15}
\]

\end{theorem}


\begin{proof}
\begin{enumerate}[label=\arabic*.,leftmargin=*]
\item The pullback is the inverse image of the diagonal under
$(f,h):X\times Y\to B\times B$.  At any point, either $df$ is
onto or its base value is exceptional and $dh$ is onto.  Thus
the difference of these two differentials is onto, proving
transversality and smoothness.  It is a closed subset of a
compact product, hence compact.  The maps
$(x,(t,n))\mapsto(x,n)$ and $(x,n)\mapsto(x,(f(x),n))$ give
the asserted holomorphic identification.  The homology statement
is the integral Kunneth theorem with $H_*(S^6)$, which is free;
there are no Tor terms.  Fundamental groups and Euler
characteristics of products give the remaining formulas.
\end{enumerate}
\end{proof}
\section{Limits and open points}
\begin{problem}[Open: the full filling and smooth classification]\label{maps:src:D.2.4}
\StatementRecord{maps:src:D.2.4}{open}{Historical events: Claude Fable 5.2 (Anthropic), effort high, 16 November 2026: carried as background without fault found. Claude Fable 5.2 (Anthropic), effort high, 17 November 2026: carried as background without fault found. GPT-6 Astra (OpenAI), ChatGPT Pro mode, effort Pro, S6\_09\_PAPERS, 17 November 2026: targeted proof check. Also checked in this writing. Pre-continuation writing events: GPT-6 Astra (OpenAI), ChatGPT Pro mode, effort Pro, S6\_10\_PAPERS, 18 November 2026: used and rechecked in this run. Continuation metadata events: GPT-6 Astra (OpenAI), ChatGPT Pro mode, effort Pro, S6\_10\_PAPERS continuation, 19 November 2026: used and rechecked in this run; status-vocabulary reconciliation; metadata reconciliation only; no new mathematical check. \emph{Pre-continuation recorded proof scope.} GPT-6 Astra (OpenAI, ChatGPT Pro mode), effort Pro, 2026-11-17 UTC; Claude Fable 5.2, effort high, 2026-11-17; Claude Fable 5.2, effort high, 2026-11-16; Claude Fable 5.2 (Anthropic), effort high, 16 November 2026; targeted check in this writing by GPT-6 Astra (OpenAI), ChatGPT Pro mode, effort Pro, 18 November 2026}{the hypotheses, notation and definitions stated here; \cref{maps:src:D.2.24}; \cref{maps:src:D.2.29}; \cref{maps:src:D.2.37}; \cref{maps:src:D.2.40}; \cref{maps:src:D.2.41}; \cref{maps:src:D.2.43}; \cref{maps:src:D.2.46}; \cref{maps:src:D.2.44}; \cref{maps:src:D.2.48}; \cref{maps:src:D.2.49}; \cref{maps:src:D.2.50}; \cref{maps:src:D.2.51}; \cref{maps:src:D.2.52}; \cref{maps:src:D.2.53}}


The natural section-line sphere question is complete on all $21$
profiles, with $59$ marked sphere patterns on $50$ initial rows
(\ref{maps:src:D.2.24}--\ref{maps:src:D.2.28}). \ref{maps:src:D.2.29}--\ref{maps:src:D.2.33} compute every allowed finite affine
logarithm and every original/log additive and original/Mumford
multiplicative choice on the natural initial catalogue: $495$
integral presentations on $172$ initial rows. \ref{maps:src:D.2.37}, \ref{maps:src:D.2.40} and
\ref{maps:src:D.2.41} include every positive translation iterate, retaining shared
even factors, arbitrary free finite and regular logarithms and the
ordinary clutch. \ref{maps:src:D.2.43} supplies the full nonsaturated general-cusp
boundary kernel. \ref{maps:src:D.2.46} now removes the scalar-iterate restriction:
all positive section pairs satisfying the fibrewise narrowness
condition and every effective whole-fibre weight distribution on
multiplicative or regular fibres have a complete integral algorithm,
with the original star caps. The exceptional OS49/II Klein-four
vertical subgroup is retained in the actual fundamental group.

The entire specified minimal simply connected smooth sector is now
classified: all $162$ patterns, all their allowed finite logarithmic
orders and affine twists, and every simply connected ordinary clutch,
within the stated primitive-Mumford scope and without additional
regular logarithms. \ref{maps:src:D.2.44}--\ref{maps:src:D.2.45} resolve the last $18$ patterns.
Residual orders $4,8,20$ give respectively
$Z_{4,1}$, $Z_{8,0}$ and $Z_{4,1}\#Z_{5,1}$.
Their full refined classes are $2,0,14$ in the fixed cyclic
coordinates, and their fundamental-class invariants are zero.
The actual relative spin calculation uses the rational linking of
the literal critical cycle. The preferred finite-$II$ core
three-manifold is a full torsion generator even when the ordinary
regular three-torus is not. No division of an even cyclic residue
or unproved general change-of-three-frame formula is used.

\ref{maps:src:D.2.48} gives the complete oriented smooth classification of the
simply connected weighted Engel family for every positive weight,
including its full finite-coefficient fundamental-class invariant.
\ref{maps:src:D.2.49} computes all canonical primitive-long-circle odd-star
logarithmic caps and their actual integral boundary kernels.
\ref{maps:src:D.2.50} proves that component corrections to the initial line exist
exactly when the height product is integral, and constructs sixteen
new simply connected fillings. \ref{maps:src:D.2.51} computes every invariant-Chern
original additive cap, retaining its integral orbit-torsion extension;
its full marked columns are proved when that torsion vanishes.
\ref{maps:src:D.2.52} consequently extends the global integral algorithm to every
positive initial pair with integral height and the stated corrected
filling choices. It does not replace unresolved torsion-valued
boundary columns by rational or saturated maps.

The full D.2 request remains open for nonprimitive-long-circle
and other logarithmic star data, corrected original multiplicative
caps, and the mixed elliptic--$\gamma$ and second fixed-normal
torsion columns of the original additive caps when $h>1$ in \ref{maps:src:D.2.51}.
The complete smooth invariants of the remaining weighted and
additional-regular-log outputs are also undetermined. For example, \ref{maps:src:D.2.46} constructs the
simply connected OS45/default member with $P=2G$, $Q=8G$, weight
$2$ at $I_8$, all other weights zero, all Mumford caps and ordinary
clutch $e=1$. It has
\[
 H_2=\mathbb Z^{15}\oplus\mathbb Z/8,\qquad
 H_3=\mathbb Z/8,\qquad \chi=32.
\]
For this exact example \ref{maps:src:D.2.53} computes the complete primitive
integral cubic and every coordinate of the fundamental-class
invariant. The component lattice has index two; the primitive
extension and the full boundary coefficient-extension data are
retained. Its absolute $c_1,c_2,\lambda$ vanish. The sole remaining
smooth invariant for this example is the refined class $0$ versus
$4$ in $\mathbb Z/8$. An actual order-two angular boundary cycle
bounds on both sides of the cusp and gives an odd class in $H_3(X)$;
its embedded representative, reference framing and critical-specialization
linking are not computed here. Even coefficients in an absolute critical
nullchain do not by themselves settle that local linking.
The unresolved original-cap torsion markings and the complete smooth
invariants of the remaining star-log families are separate open tasks. Earlier failed half-weight, ordinary-characteristic
and nongenerating-three-torus shortcuts are retained with their
corrections in the dated research records.


\end{problem}
\end{document}
